\documentclass[11pt]{article}
\usepackage[T1]{fontenc}
\usepackage{lmodern,amsmath,amssymb,amsthm,mathtools}
\usepackage[margin=1in]{geometry}
\usepackage{booktabs,longtable,array,microtype}
\usepackage[hidelinks]{hyperref}
\usepackage{xurl}
\newtheorem{theorem}{Theorem}[section]
\newtheorem{proposition}[theorem]{Proposition}
\newtheorem{lemma}[theorem]{Lemma}
\newtheorem{corollary}[theorem]{Corollary}
\theoremstyle{remark}\newtheorem{remark}[theorem]{Remark}
\DeclareMathOperator{\Tr}{Tr}
\title{Renewal Yang--Mills Mass-Gap Research Notes\\
Restart after archive f17\\
Version V91: Source Tilts and Native Action-Density Concentration}
\author{Research continuation and mathematical audit}
\date{20 September 2026}
\begin{document}
\maketitle
\begin{abstract}
This shorter working note starts a new cycle after the
hundred-version consolidated archive f17. It records the
major usable results, their limits, and a new calculation.
The interacting four-dimensional continuum Yang--Mills
construction and its positive physical mass gap remain unproved.

The latest advance is not a continuum gap: a one-link high
block and its full return have volume-independent bounds.
A native spin-network calculation then shows why this cannot
be promoted to elimination of an entire retained electric band.
Equal-energy Wilson couplings cross every one-link cutoff.
We keep those resonances and derive the connected second-order
operator on the retained total-electric-energy blocks.
\end{abstract}
\tableofcontents

\section{Context, conventions and the preserved archive}
\label{c18v1:context}

The goal is unchanged: construct a nontrivial interacting
four-dimensional continuum Yang--Mills theory satisfying
the required reconstruction axioms and prove a positive
gap in its physical Hamiltonian. The concrete finite-regulator
results use $\mathrm{SU}(2)$; they do not by themselves cover
every compact simple group.

The previous complete working source is preserved as
\begin{center}\small
\path{Renewal_Yang_Mills_Mass_Gap_Research_Notes_V100_f17.tex}.
\end{center}
Here \emph{f17 V$n$} means that version's appendix in the
archive, not this new V$n$. The f15 and f16.1/f16.2 archives
retain earlier branches and their qualifications. This is
a summary plus a new calculation, not a replacement for
their full proofs. No earlier result is silently upgraded.
The historical corpus has not been independently recertified
line by line.

Work on the homogeneous cubic spatial torus, $N\ge3$,
using product normalized Haar measure and its gauge-invariant
subspace. The Hamiltonian is
\[
 \begin{gathered}
 H=T+b|\mathcal P_N|+V,\qquad T=\kappa\sum_e L_e,\\
 V=-b\sum_{p\in\mathcal P_N}w_p,\qquad
 w_p=\tfrac12\Tr U_p,\qquad \kappa>0,\quad b\ge0 .
 \end{gathered}
 \tag{1.1}\label{c18v1:H}
\]
The positive sphere Casimir has eigenvalues
$\ell(\ell+2)$, multiplicity $(\ell+1)^2$.
The unique positive finite-volume ground is $\Omega$,
with energy $E_0$; set $K=H-E_0$.
The electric constant state $\Omega_{\rm el}=1$ is not
asserted to equal $\Omega$ when $b>0$.

Physical time, auxiliary stochastic time and spatial
smoothing time remain distinct. Fixed-spacing infinite
volume and vanishing-spacing continuum are different
limits. No kinetic or magnetic normalization is changed
to make a desired gap appear.

\section{Major results retained, with their limits}
\label{c18v1:ledger}

\begin{longtable}{@{}p{.19\textwidth}p{.75\textwidth}@{}}
\toprule
Location & Usable result and restriction\\
\midrule
\endhead
f17 V1--V14 &
Exact temporal-link/root integration, a complete fixed-graph
bounded-endpoint expansion, native scores and positive
finite-regulator response tests. These do not supply
uniform long-history or growing-volume expansions.\\[1ex]
f17 V15--V30 &
Temporal purity and sewing identities, positive cubature,
and exact finite decimation/conditional response.
Normalized exterior and shared-message dependence
cannot be dropped in a continuum limit.\\[1ex]
f17 V34--V79 &
Joint-source, Ward/current and compact-corridor calculations,
culminating in a complete first covariance correction
with orientation-dependent contacts and controlled spatial
symbol variation. These are coefficient or stated-regime
results, not a common nonperturbative remainder theorem.\\[1ex]
f17 V81--V84 &
Native energy-source concentration on a stated small
integrated-magnetic-budget strip, then sewn-vacuum,
entropy and positive-time source bounds. Their strip
hypotheses and source normalization remain essential.\\[1ex]
f17 V85--V87 &
A conserved local-energy first-moment sum rule, unweighted
block fluctuation bounds and a complete partial-charge
source bank. These are upper bounds. The neutral return
is nonconstant and still needs lower spectral control.\\[1ex]
f17 V88--V90 &
An isolated magnetic rotor gap, exact conditional geometry,
and a compatible Gibbs trial with a nonnegative neutral
penalty. The trial is proved not to be the native vacuum
in general; its gap or energy-density scale is not the
centered physical gap.\\[1ex]
f17 V91--V98 &
Native conditional scores, variational budgets, spectral
counts and a smooth vacuum-compatible harmonic compression.
The physical shape channel and conditional inverse
losses are retained, not removed by a frozen fibre.\\[1ex]
f17 V99 &
A fixed harmonic cutoff with actual high-block coercivity,
bounded exact coupling, controlled Schur return and
factorial tails for a whole low-energy spectral window,
uniformly in volume.\\[1ex]
f17 V100 &
Physical resonant crossings at every one-link cutoff,
strictly local periodic electric averaging, and an exact
unitary normal form with summable spatial remainder
for $13\pi b/\kappa<1$. This is not the physical
weak-coupling continuum regime.\\
\bottomrule
\end{longtable}

The bottleneck is a uniform estimate on the actual
interacting neutral dynamics in a normalization compatible
with a nontrivial continuum construction. A matrix test,
formal coefficient, conditional sampler gap or auxiliary
comparison gap does not supply that estimate.

\subsection{The high-energy result now available}

Using the exterior ground energy before taking a norm,
f17 V99 proves for one link
\[
 \begin{gathered}
 K=D_e+\kappa L_e+U_e,\quad D_e\ge0,\quad
 U_e=-f(y)\cdot q_e,\quad |f(y)|\le4b,\\
 C_L=Q_LKQ_L\ge d_L:=\kappa(L+1)(L+3)-4b,\\
 B_L=Q_LKF_L=P_{L+1}U_eP_L,\qquad \|B_L\|\le2b .
 \end{gathered}
 \tag{2.1}\label{c18v1:high}
\]
Here $F_L$ includes degrees $0,\ldots,L$, while $P_L$
is just degree $L$. The exterior is not frozen.
For real $z<d_L$,
\[
 0\le\Sigma_L(z)=B_L^*(C_L-z)^{-1}B_L,\qquad
 \|\Sigma_L(z)\|\le\frac{4b^2}{d_L-z}.
 \tag{2.2}\label{c18v1:Sigma}
\]
The head $h_0=F_L\Omega$ lifts with high component
$-C_L^{-1}B_Lh_0$. Its norm uses
$G_L=I+B_L^*C_L^{-2}B_L$.
Vacuum removal in this graph is
$\langle h_0,G_Lh\rangle=0$, not subtraction of a constant.

For $\mathcal P_*=1_{[0,E_*]}(K)$ and
$T_n=\sum_{\ell\ge n}P_\ell$,
\[
 \|T_n\mathcal P_*\|
 \le\frac{2b}{\kappa n(n+2)-4b-E_*}
                    \|T_{n-1}\mathcal P_*\|
 \tag{2.3}\label{c18v1:tail}
\]
when the denominator is positive.
Iteration gives factorial decay. This concerns a total
low-energy window, not separation of the entire $F_L$
and $Q_L$ spectra.

\subsection{Retained energy exchange cannot be dropped}

The f17 V100 theta graph has two plaquette paths of
length two and an exterior path of length four.
Its physical states with path spins
$(\ell/2,(\ell+1)/2,1/2)$ and
$((\ell+1)/2,\ell/2,1/2)$ satisfy
\[
 \begin{gathered}
 E_\ell=\kappa(4\ell^2+12\ell+18),\\
 \langle\Psi'_\ell,V\Psi_\ell\rangle
           =-\frac{b}{2(\ell+1)(\ell+2)}\ne0\quad(b>0).
 \end{gathered}
 \tag{2.4}\label{c18v1:resonance}
\]
They lie on opposite sides of a selected link's
degree-$\ell$ cutoff. The all-$\ell$ coefficient is proved
by normalized invariant-tensor Haar contraction.
It precludes a whole-head inverse of the commutator
with $T$. These high total energies do not contradict
\eqref{c18v1:high}.

A separate exact two-site test in f17 V100 shows that
a fixed-energy Schur return need not be spatially local
even with an independent distant spectator.
Energy-dependent Schur identities and local effective
Hamiltonians are distinct objects.

\section{New continuation: the connected second-order resonant operator}
\label{c18v1:new}

We continue beyond the archive by computing an explicit
operator on every total-electric-energy block, with
connected support retained. This does not yet compute
the interacting vacuum gap.

For $\tau=2\pi/\kappa$ define the strong integrals
\[
 \begin{aligned}
 {\cal R}(A)&=\tau^{-1}\int_0^\tau e^{itT}Ae^{-itT}\,dt,\\
 {\cal J}(A)&=i\tau^{-1}\int_0^\tau
                   (t-\tau/2)e^{itT}Ae^{-itT}\,dt .
 \end{aligned}
 \tag{3.1}\label{c18v1:maps}
\]
Put $R_p={\cal R}(V_p)$, $S_p={\cal J}(V_p)$,
$R=\sum_pR_p$ and $S=\sum_pS_p$.
The archive proves on the electric domain
\[
 [T,S]=V-R,\qquad
 \|R_p\|\le b,\qquad
 \|S_p\|\le s=\frac{\pi b}{2\kappa}.
 \tag{3.2}\label{c18v1:S}
\]
Both maps preserve support and gauge invariance.
The first commutator contribution and its resonant part are
\[
 W_1=\tfrac12[S,V+R],\qquad Z_2={\cal R}(W_1).
 \tag{3.3}\label{c18v1:Zdef}
\]
The subscript denotes order in $b$, not an energy cutoff.

\begin{theorem}[Connected second-order resonant return]
\label{c18v1:Ztheorem}
At every finite native regulator,
\[
 \boxed{\quad
 Z_2=\tfrac12{\cal R}([S,V])
  =\tfrac12\sum_{\substack{p,q\\p\cap q\ne\varnothing}}
                         {\cal R}([S_p,V_q]).
 \quad}
 \tag{3.4}\label{c18v1:Zlocal}
\]
Intersection means a shared link, including $p=q$.
The operator is self-adjoint, gauge invariant and
commutes with $T$. Each term is supported on the
connected union of at most two plaquettes.
For a fixed first plaquette its sum of term norms is at most
\[
                   13sb=\frac{13\pi b^2}{2\kappa}.
 \tag{3.5}\label{c18v1:Zroot}
\]
For every electric eigenvalue $E$ and its projection $P_E$,
\[
 P_EZ_2P_E
 =P_EV(1-P_E)(E-T)^{-1}(1-P_E)VP_E .
 \tag{3.6}\label{c18v1:Zblock}
\]
The inverse is only on $1-P_E$, with norm at most
$1/\kappa$. Equal-energy interactions remain in
$P_ERP_E$ and are not divided by zero.
\end{theorem}

\begin{proof}
In electric matrix elements, $R$ is diagonal in total
energy and $S$ has no equal-energy matrix elements.
Hence \({\cal R}([S,R])=0\).
Disjoint supports commute and the on-site averaging
preserves support. Since $S_p$ is anti-self-adjoint,
$[S_p,V_q]$ is self-adjoint and has norm at most $2sb$.
Averaging is contractive, and at most thirteen
plaquettes meet a fixed plaquette, including itself.
This proves \eqref{c18v1:Zlocal}--\eqref{c18v1:Zroot}.

For $E_m\ne E_n$,
$S_{mn}=V_{mn}/(E_m-E_n)$; at equality it is zero.
For indices $m,n$ of energy $E$,
\[
 \tfrac12[S,V]_{mn}
       =\sum_{k:E_k\ne E}\frac{V_{mk}V_{kn}}{E-E_k}.
\]
This proves \eqref{c18v1:Zblock}. At finite volume
$P_E$ is finite dimensional and $V$ is bounded.
Distinct electric energies differ by at least $\kappa$,
so the reduced inverse and the sum are well-defined.
The identities may first be checked on the finite
Peter--Weyl core and then extended as bounded operators.
No positivity of this signed inverse is asserted at
excited energy $E$.
\end{proof}

The energy-block formula is standard degenerate
perturbation algebra, not a new general principle.
The present calculation retains the complete native
resonant block and its connected plaquette budget.
It obtains a local interaction at this coefficient
order without freezing an exterior resolvent.

\subsection{A fixed-volume meaning with its error stated}

Let $M=|\mathcal P_N|$ and set
\[
 a=Ms,\qquad v=Mb,\qquad w=2av,\qquad
 h=\frac{\pi w}{2\kappa}.
 \tag{3.7}\label{c18v1:finite-constants}
\]
The anti-self-adjoint operator $S_2={\cal J}(W_1)$
has norm at most $h$, connected two-plaquette support,
and $[T,S_2]=W_1-Z_2$. More precisely, each local term
has that support; the sum ranges over the whole lattice.
Then
\[
 \begin{split}
 e^{S_2}e^S(H-bM)e^{-S}e^{-S_2}
       &=T+R+Z_2+\mathcal E_3,\\
 \|\mathcal E_3\|
       &\le v(e^{2a}-1-2a)+2hv+4hw .
 \end{split}
 \tag{3.8}\label{c18v1:finite-error}
\]
This is $O_N(b^3/\kappa^2)$ as $b/\kappa\to0$
at the stated fixed regulator. No uniformity in $M$
is claimed.

\begin{proof}
The exact first conjugation has remainder
$W_1+W_{\ge2}$. The archived commutator series and
$\|S\|\le a$, $\|V\|,\|R\|\le v$ give
$\|W_1\|\le w$ and
$\|W_{\ge2}\|\le v(e^{2a}-1-2a)$.
Also $\|Z_2\|\le w$. The linear commutator of the
second conjugation cancels $W_1-Z_2$.

For bounded $A$, unitary Duhamel integration gives
$\|e^{S_2}Ae^{-S_2}-A\|\le2h\|A\|$.
Apply it to $R$ and $W_1$.
The remainder beyond the first commutator of $T$
is at most $h\|[S_2,T]\|\le2hw$ by the integrated
second-commutator formula. The transformed
$W_{\ge2}$ retains its norm. Summing proves
\eqref{c18v1:finite-error}.
Products, commutators and the periodic maps preserve
the finite Peter--Weyl core. The bounded commutator
$[T,S_2]$ then gives electric-domain preservation
by graph-norm approximation, as in the archive.
\end{proof}

\subsection{Vacuum normalization as a regression}

On $\Omega_{\rm el}$, $w_p$ creates energy $12\kappa$,
has squared norm $1/4$, and different plaquettes
create orthogonal states. Therefore
\[
 \langle\Omega_{\rm el},Z_2\Omega_{\rm el}\rangle
                   =-\frac{Mb^2}{48\kappa}.
 \tag{3.9}\label{c18v1:vacuum-coefficient}
\]
Each of four links contributes fundamental Casimir $3$;
Haar character orthogonality gives the norm.
An edge in only one plaquette makes the cross terms
zero. This reuses f17 V87's free plaquette denominator
as a normalization regression, not as a new mass scale.

At excited $E$, the denominators in
\eqref{c18v1:Zblock} can have either sign.
One cannot apply the negative vacuum shift separately
to every state, omit $R$, or infer a positive
centered gap from \eqref{c18v1:vacuum-coefficient}.
The retained interactions include
\eqref{c18v1:resonance}.

\section{Remaining obligations and direction}
\label{c18v1:next}

The full mass gap remains unproved. The next obligation
is control of the actual interacting neutral energy
exchange and connected higher returns.
The resonance excludes a whole-band nonresonant inverse;
the averaging gives a way to keep those states explicitly.

Two distinct tools are now available:
\begin{enumerate}
\item The all-coupling high Schur block
\eqref{c18v1:high}, with the unknown interacting
retained pencil and graph vacuum.
\item The local electric-resonant normal form.
Its first-step spatial remainder needs
$13\pi b/\kappa<1$; its next coefficient is
\eqref{c18v1:Zlocal}.
\end{enumerate}
Neither supplies a renormalisation trajectory carrying
its hypotheses to physical weak coupling. Repeating
a strong-electric expansion cannot replace that missing
step. Equation \eqref{c18v1:finite-error} is not a
cofinal error estimate.

Local dynamics also does not imply spectral coercivity.
A final proof must identify physical time, control
the actual vacuum on an appropriate dense observable
core, construct a nontrivial continuum limit with
the required positivity and covariance, and keep
a positive physical excitation gap.
Upper source bounds and finite-regulator positivity
alone do not provide these facts.

The V100 literature audit recommends retaining exact
resonances and connected spatial structure when using
effective-Hamiltonian methods. Primary context is
\href{https://arxiv.org/abs/1105.0675}{Bravyi--DiVincenzo--Loss}
and
\href{https://arxiv.org/abs/1810.02428}{Nachtergaele--Sims--Young}.
The companions supply proof-organization lessons,
not missing nonperturbative YM coefficients.

\section{Verification and continuation record}
\label{c18v1:record}

The f17 V100 archive preserves its preceding body
apart from the title. All older archives and companions
remain unchanged. This V1 is the single active source.
Later versions append to it and replace only the
active predecessor. The next companion checkpoint
is V5, the next broad literature checkpoint V10;
at V100 this cycle will again be archived and summarized.

Diagnostics verify the native theta contractions,
an independent fundamental Haar moment, exact spectator
identities, cubic incidence, periodic homological
identities and the connected second-order formula.
Finite qutrit tests check both conjugation error bounds;
they are not surrogate Yang--Mills models or continuum
proofs. The preceding V99 diagnostics are replayed.
Source preservation and rendered layout are checked
separately. Only \texttt{.tex} files are kept in
\texttt{latex\_notes}; build files and diagnostics
remain outside it.

% END CYCLE 18 VERSION 1 BODY
% Append the next version before the final document terminator.

% BEGIN CYCLE 18 VERSION 2 APPEND
\clearpage
\section{V2: A connected cubic error and the centered resonant form}
\label{c18v2:main}

V1 obtained the exact second-order resonant operator but
bounded the next error only at a fixed volume.
This append removes that volume loss in an interaction norm.
It also proves a relative form bound for the complete
vacuum-centered second-order operator, without deleting its
energy-exchange matrix elements. These are finite-spacing
results with constants uniform in the cubic volume; they
do not construct the four-dimensional continuum theory.

The source inventory and V1 body were rechecked before
continuing. The last companion and broad literature
checkpoints are f17 V100; the next regular checks remain
this cycle's V5 and V10. Searches of the existing YM sources
did not locate the particular two-step bound or centered
second-order estimate below. The general commutator and
relative-form methods are standard; no general novelty
claim is made.

\subsection{Parameters, local terms and the relevant norm}

Keep \eqref{c18v1:H} and \eqref{c18v1:maps}. Throughout,
\[
 \beta=\frac b\kappa,\qquad s=\frac{\pi b}{2\kappa},
 \qquad D=13,\qquad q=2Ds=13\pi\beta,\qquad M=|\mathcal P_N|.
 \tag{V2.1}\label{c18v2:parameters}
\]
The symbol $\beta$ here is only the displayed ratio; it is
not an imported Euclidean lattice coupling.
Write $p\sim r$ for two plaquettes sharing a link,
including $p=r$, and put
\[
 \begin{aligned}
 w_{pr}&=\tfrac12[S_p,V_r+R_r],&
 W_1&=\sum_{p\sim r}w_{pr},\\
 t_{pr}&={\cal J}(w_{pr}),&
 S_2&=\sum_{p\sim r}t_{pr},&
 Z_2&={\cal R}(W_1).
 \end{aligned}
 \tag{V2.2}\label{c18v2:pairs}
\]
Then
\[
 \|w_{pr}\|\le2sb,\qquad
 \|t_{pr}\|\le t:=\frac{\pi sb}{\kappa}=2s^2,\qquad
 \sum_{r\sim p}\|w_{pr}\|\le bq.
 \tag{V2.3}\label{c18v2:pair-bounds}
\]
The same rooted bound holds for \({\cal R}(w_{pr})\).
Each term is gauge invariant and supported on the indicated
connected one- or two-plaquette set. All these statements
refer to operators on the original, untruncated Haar Hilbert
space, before restriction to the physical subspace.

For a finite connected plaquette set $Z$, let $E(Z)$ be
its links and let $d_P(Z)$ be its diameter in the overlap
graph restricted to $Z$. In particular $d_P(Z)\le |Z|-1$.
For a specified decomposition $A=\sum_Z A_Z$ define
\[
 \|A\|_{\mu,\mathrm{int}}
 :=\sup_e\sum_{Z:e\in E(Z)}e^{\mu d_P(Z)}\|A_Z\|,
 \qquad \mu\ge0.
 \tag{V2.4}\label{c18v2:interaction-norm}
\]
This is a norm budget for that interaction decomposition,
not the extensive operator norm of $A$.

\subsection{Two connected counting lemmas}

\begin{lemma}[Adding a two-plaquette generator]
\label{c18v2:growth}
If an initial local operator is supported on at most $m$
plaquettes, then the order-$k$ nested commutator with $S_2$,
divided by $k!$, has total term norm at most
\[
 \|A\|\frac{(m/2)_k}{k!}\,x^k,\qquad
 x=4Jt=4q^2,\qquad J=2D^2=338.
 \tag{V2.5}\label{c18v2:pochhammer}
\]
Here $(a)_k=a(a+1)\cdots(a+k-1)$, with $(a)_0=1$.
Each nonzero word has a connected union if the initial
support is connected, and adds at most $2k$ plaquettes.
\end{lemma}

\begin{proof}
A nonzero new commutator must have at least one of its
two plaquettes overlapping the current support.
For a specified plaquette, there are at most $D$ choices
for that member of the ordered pair and at most $D$
for its partner, in either order. Thus there are at most
$Jm$ possible ordered pairs at the first step.
After $j$ steps the support contains at most $m+2j$
plaquettes. Each commutator costs a factor $2t$.
The resulting bound is
\[
 \frac{\|A\|}{k!}(2Jt)^k
            \prod_{j=0}^{k-1}(m+2j)
       =\|A\|\frac{(m/2)_k}{k!}(4Jt)^k.
\]
The relations $t=2s^2$ and $q=2Ds$ give $4Jt=4q^2$.
A pair is connected internally and intersects the previous
support, proving the support statement.
The count is deliberately an upper bound, not a claim
that all listed commutators are nonzero.
\end{proof}

\begin{lemma}[Translation anchoring without a volume factor]
\label{c18v2:anchoring}
Consider a translation-covariant family of rooted connected
words of order $r$, with at most $r+1$ plaquettes per word
and total norm at most $a_r$ for each fixed root.
On the homogeneous cubic torus its contribution to
\eqref{c18v2:interaction-norm} is at most
\[
                   12(r+1)e^{\mu r}a_r .
 \tag{V2.6}\label{c18v2:anchor}
\]
\end{lemma}

\begin{proof}
Choose one root at the origin for each of the three
plaquette orientations and list its relative words.
For a fixed word and a fixed link $e$ of a given
orientation, the number of translations of the word
whose support contains $e$ equals the number of distinct
links of that orientation in the word. It is at most
$4(r+1)$. Translation preserves the term norm.
Sum over words and the three root orientations.
Also $d_P(Z)\le r$. This proves the bound.
On a torus the same counting uses translations modulo $N$;
repeated visits do not increase the number of distinct links.
Grouping words with the same $Z$ cannot increase the estimate.

This argument uses the homogeneous translation-covariant
interaction. It does not silently assert the same constant
for an arbitrary spatially varying coupling.
\end{proof}

\subsection{The complete two-step remainder}

For $0\le u<1/\sqrt5$ define
\[
 \begin{gathered}
 a(u)=(1-4u^2)^{-1/2},\\
 g(u)=a(u)-1+u\bigl(a(u)^2-1\bigr)
           +\frac{a(u)\,[u\,a(u)]^2}{1-u\,a(u)}.
 \end{gathered}
 \tag{V2.7}\label{c18v2:g}
\]
Its Taylor coefficients are nonnegative, and
\[
             g(u)=3u^2+5u^3+13u^4+25u^5+O(u^6).
 \tag{V2.8}\label{c18v2:g-series}
\]

\begin{theorem}[Uniform connected cubic error]
\label{c18v2:uniform-error}
Suppose $q<1/\sqrt5$. At every $N\ge3$ the exact V1
conjugation has a connected, translation-covariant
decomposition
\[
 e^{S_2}e^S(H-bM)e^{-S}e^{-S_2}
                =T+R+Z_2+\mathcal E_3,\qquad
 \mathcal E_3=\sum_Z\mathcal E_{3,Z}.
 \tag{V2.9}\label{c18v2:normal}
\]
For every $\mu\ge0$ such that $q e^\mu<1/\sqrt5$,
\[
 \boxed{\quad
 \|\mathcal E_3\|_{\mu,\mathrm{int}}
       \le 12b\,[g(u)+u g'(u)]\big|_{u=q e^\mu}.
 \quad}
 \tag{V2.10}\label{c18v2:uniform-bound}
\]
Also $\|\mathcal E_3\|\le Mb\,g(q)$.
For fixed $\mu$, the interaction bound is
\[
       108bq^2e^{2\mu}+O(bq^3)
                         =O(b^3/\kappa^2),
 \tag{V2.11}\label{c18v2:cubic}
\]
with constants independent of $N$ as $b/\kappa\to0$.
The interaction tail with $d_P(Z)\ge L$ has unweighted
per-link norm at most $e^{-\mu L}$ times the right side of
\eqref{c18v2:uniform-bound}.
\end{theorem}

\begin{proof}
Let $W_{\ge2}$ be the tail after $W_1$ in f17 V100's
first conjugation. Its order-$n$ terms, $n\ge2$, have
rooted norm at most $bq^n$, connected support on at most
$n+1$ plaquettes, and preserve translation covariance.
Set $\operatorname{ad}_2 A=[S_2,A]$.
The identity $[S_2,T]=-(W_1-Z_2)$ yields exactly
\[
 \begin{split}
 \mathcal E_3={}&e^{\operatorname{ad}_2}W_{\ge2}
      +(e^{\operatorname{ad}_2}-1)R\\
 &+\sum_{k\ge1}
       \frac{k\,\operatorname{ad}_2^k W_1
                       +\operatorname{ad}_2^k Z_2}{(k+1)!}.
 \end{split}
 \tag{V2.12}\label{c18v2:three-families}
\]
At finite volume all generators are bounded.
The first commutators with $T$ are bounded on its domain,
as established in V1. Its graph-norm domain argument and
the integrated conjugation formula justify this identity
for the unbounded $T$. No bound on $[T,V]$ is assumed.

There are three families of rooted words.
For the conjugation of $R$, Lemma~\ref{c18v2:growth}
has $m=1$, $k\ge1$, and rooted bound
$b x^k(1/2)_k/k!$.
For the last line of \eqref{c18v2:three-families},
$m=2$ and the two coefficients sum to $1/k!$.
Its bound is $bq x^k$, $k\ge1$.
For a tail word of first order $n\ge2$, use $m=n+1$
to obtain
\[
                 bq^n x^k\frac{((n+1)/2)_k}{k!},
                 \qquad k\ge0.
 \tag{V2.13}\label{c18v2:tail-family}
\]
The corresponding total word orders are $2k$, $1+2k$,
and $n+2k$, respectively. Each support has at most one
more plaquette than its word order.

Introduce a positive bookkeeping variable $z$ for that
order, and put $u=qz$. Summing the first two families
by the binomial series gives $a(u)-1$ and
$u(a(u)^2-1)$. The third gives
\[
 \sum_{n\ge2}u^n a(u)^{n+1}
                  =\frac{a(u)[u a(u)]^2}{1-u a(u)}.
\]
Here $u a(u)<1$ is equivalent, for nonnegative $u$,
to $5u^2<1$; it also implies $4u^2<1$.
Thus $b g(qz)$ is a convergent rooted majorant.
If $g(u)=\sum_r c_r u^r$, the rooted order-$r$ bound
is $bc_rq^r$. Summing over $M$ roots gives the extensive
operator bound. Lemma~\ref{c18v2:anchoring} gives
\[
 12b\sum_r(r+1)c_r(qe^\mu)^r
             =12b\,[g(u)+u g'(u)]_{u=qe^\mu}.
\]
Absolute convergence permits the grouping by connected
supports. The Taylor expansion and the exponential
weight prove the remaining assertions.
All terms commute with gauge transformations, so the
identities and bounds restrict to the physical subspace.
\end{proof}

This replaces the volume-dependent bound
\eqref{c18v1:finite-error} in a stronger, spatial sense,
not by falsely making the global operator error
volume independent. The sufficient convergence condition
is now $13\pi b/\kappa<1/\sqrt5$.
The constants are not asserted to be sharp.

\subsection{A vacuum-centered bound on every resonant energy block}

Use the more economical local representatives
\[
 z_{pr}=\tfrac12{\cal R}([S_p,V_r]),\qquad
 Z_2=\sum_{p\sim r}z_{pr},\qquad \|z_{pr}\|\le sb.
 \tag{V2.14}\label{c18v2:z-local}
\]
Each $z_{pr}$ commutes with the electric energy on its
one- or two-plaquette support, not just with the full $T$.
This follows because the electric evolution factorizes
over links and the operator is the identity outside
that support.

\begin{proposition}[The full centered second-order form]
\label{c18v2:relative}
For every $b\ge0$ at the stated finite regulator, set
\[
 \widehat Z_2=Z_2+\frac{Mb^2}{48\kappa}I,\qquad
 \rho_2(\beta)=
   \left(\frac{44\pi}{3}+\frac1{36}\right)\beta^2.
\]
Then $\widehat Z_2\Omega_{\rm el}=0$, and on
$\operatorname{Dom}T^{1/2}$,
\[
 |\langle v,\widehat Z_2v\rangle|
                         \le\rho_2(\beta)\langle v,Tv\rangle.
 \tag{V2.15}\label{c18v2:relative-bound}
\]
No small-coupling assumption is needed for this
coefficient identity and form bound.
\end{proposition}

\begin{proof}
Let $P_{0,Z}$ project onto the local constant vector on
the links of $Z$. The local zero-electric-energy space
has rank one, so
\[
 z_{pr}P_{0,Z}=c_{pr}P_{0,Z},\qquad
 c_{pr}=\langle\Omega_{\rm el},z_{pr}\Omega_{\rm el}\rangle,
 \qquad Z=\{p,r\}.
\]
The V1 electric resolvent calculation, or its polarized
matrix element, gives
\[
 c_{pr}=-\operatorname{Re}
       \langle V_p\Omega_{\rm el},T^{-1}V_r\Omega_{\rm el}\rangle
            =-\delta_{pr}\frac{b^2}{48\kappa}.
 \tag{V2.16}\label{c18v2:c}
\]
The inverse acts on the complement of the electric vacuum.
Indeed $V_p\Omega_{\rm el}$ has energy $12\kappa$,
squared norm $b^2/4$, and states from distinct plaquettes
are orthogonal by Haar integration over an unpaired link.

Consequently $\widehat z_{pr}=z_{pr}-c_{pr}I_Z$
annihilates the local constant vector on both sides.
The local electric gap is at least $3\kappa$, so as forms
\[
 \begin{gathered}
 |\langle v,\widehat z_{pr}v\rangle|
 \le\left(sb+\delta_{pr}\frac{b^2}{48\kappa}\right)
                         \langle v,(I-P_{0,Z})v\rangle,\\
 I-P_{0,Z}\le\frac1{3\kappa}
                     \sum_{e\in E(Z)}\kappa L_e .
 \end{gathered}
 \tag{V2.17}\label{c18v2:local-form}
\]
For a fixed link $e$, there are four incident plaquettes.
There are at most $4D$ ordered pairs with their first
plaquette incident to $e$, and at most $4D$ with their
second incident. Their intersection consists of all
$4^2$ ordered pairs from the incident star. Thus at most
\[
                       8D-16=88
\]
ordered overlapping pairs contain $e$. Exactly four
of these are diagonal pairs $p=r$.
Sum \eqref{c18v2:local-form}; the coefficient of $T$ is
\[
 \frac{88sb+4b^2/(48\kappa)}{3\kappa}
          =\left(\frac{44\pi}{3}+\frac1{36}\right)\beta^2.
\]
The scalar sums to $Mb^2/(48\kappa)$.
The annihilation and form assertions follow first on
the core and then on the form domain.
\end{proof}

\begin{corollary}[A genuine bound for the comparison operator]
\label{c18v2:comparison-gap}
Let
\[
 G_2=T+R+\widehat Z_2,\qquad
 \rho(\beta)=\frac43\beta+\rho_2(\beta).
 \tag{V2.18}\label{c18v2:G}
\]
If $\rho(\beta)<1$, $G_2$ has the unique constant ground
and a full-Hilbert-space gap at least
$3\kappa[1-\rho(\beta)]$, uniformly in $N$.
The same lower bound holds on the gauge-invariant subspace.
\end{corollary}

\begin{proof}
The f17 V100 form bound for $R$ and
\eqref{c18v2:relative-bound} give
$G_2\ge(1-\rho)T$ and $G_2\Omega_{\rm el}=0$.
The full electric Hamiltonian has a simple zero eigenvalue
and satisfies $T\ge3\kappa$ on its orthogonal complement.
The lower bound proves the assertions by the spectral theorem.
All operators preserve the gauge-invariant subspace and
the constant vector belongs to it.
\end{proof}

This controls the entire resonant comparison operator,
including the native equal-energy exchanges from f17 V100.
It is not a truncation to the electric vacuum and is not
the gap of the actual $H-E_0$.

\subsection{The remaining vacuum and continuum obligations}

The exact transformed centered Hamiltonian is
\[
 e^{S_2}e^S(H-E_0)e^{-S}e^{-S_2}
 =G_2+\mathcal E_3+
       \left(bM-E_0-\frac{Mb^2}{48\kappa}\right)I .
 \tag{V2.19}\label{c18v2:actual-centered}
\]
Neither its scalar nor its vacuum is replaced by the
second-order comparison value. A small interaction norm
alone is not a direct relative form bound with respect to
$T$. The necessary vacuum condition is elementary:

\begin{lemma}[A pure relative bound must annihilate the vacuum]
\label{c18v2:vacuum-test}
Let $T\ge0$, $T\Omega_{\rm el}=0$, and let $A$ be bounded
and self-adjoint. If
$|\langle v,Av\rangle|\le c\langle v,Tv\rangle$
on $\operatorname{Dom}T^{1/2}$, then $A\Omega_{\rm el}=0$.
\end{lemma}
\begin{proof}
First test $v=\Omega_{\rm el}$. Then test
$v=\Omega_{\rm el}+\varepsilon w$ for real
$\varepsilon$ of both signs. The right side is quadratic
in $\varepsilon$, so the linear coefficient on the left
must vanish. Replacing $w$ by $iw$ gives
$\langle w,A\Omega_{\rm el}\rangle=0$ on a dense domain.
\end{proof}
Subtracting an expectation value removes a scalar, not
off-diagonal vacuum creation. We have not proved that
the remaining interaction $\mathcal E_3$ satisfies this
annihilation condition term by term after a suitable
further dressing.

As a targeted scope check,
\href{https://arxiv.org/abs/math-ph/0412040}{Yarotsky's
relative-perturbation theorem} allows infinite-dimensional
on-site spaces but assumes a classical product reference,
a fixed finite interaction range, and quantitative
perturbative smallness. The original finite-range
Hamiltonian already lies in a familiar strong-electric
perturbative setting. Applying that theorem to the new
infinite-range, exponentially decaying remainder would
require an additional justification; none is presumed here.
Theorem statements and assumptions were reviewed, not
every proof in the paper. No external stability theorem
is used to prove this append's bounds.

There is now a controlled local second-order comparison
and an explicit connected cubic error. The next spectral
step is actual vacuum dressing or a fully verified stability
argument with this error. A further, separate obstacle is
transport to a physical weak-coupling refinement trajectory.
The condition $b/\kappa<1/(13\pi\sqrt5)$ is not such a
trajectory. Improving a strong-electric coefficient
cannot by itself construct a nontrivial four-dimensional
continuum limit or prove a positive physical mass gap.

\subsection{Verification and preservation}

The diagnostic checks the three-family generating series
through order ten, the Pochhammer growth identity,
translation anchoring and the $88$-pair count on three
cubic tori. Three non-Wilson qutrit tests independently
check the exact remainder decomposition and generic
local vacuum-centered form inequality. They are finite
regression tests, not proofs of the uniform statements.
V1's mathematical diagnostics are replayed.

The complete V1 body is preserved before this appendix,
apart from the title. V2 replaces only that active source;
all forty-seven other note files, including f17, are
unchanged. Only \texttt{.tex} files are retained in
\texttt{latex\_notes}. Build products and verification
records are outside that folder. The full interacting
four-dimensional YM construction and mass gap remain
unproved; the goal is unchanged.

% END CYCLE 18 VERSION 2 APPEND

% BEGIN CYCLE 18 VERSION 3 APPEND
\clearpage
\section{V3: Local vacuum dressing and a certified ground-energy error}
\label{c18v3:main}

V2 bounded the whole cubic error in a local interaction
norm, but did not control what that error does to the
vacuum. We now separate its vacuum-creating part and
invert only that source. This gives an exact local
dressing step with a quartic remainder. A separate
variational argument bounds the \emph{actual} ground
energy relative to the twice-dressed trial, and controls
local marginals of the actual twice-transformed vacuum.
No spectral-gap conclusion is inferred from those
ground-state estimates.

The current V2 source and the archive inventory were
checked before this continuation. Earlier notes already
contain strong-coupling lattice-gap results, including
f14 V52; that result is not redone here. The new estimates
concern V2's particular connected remainder and its
vacuum source. The local off-diagonal inverse is a standard
Schrieffer--Wolff construction: compare
\href{https://arxiv.org/abs/1105.0675}
{Bravyi--DiVincenzo--Loss, Section 4.4, (4.19)--(4.22)}.
Our proofs below specify the untruncated electric domain,
the connected interaction bounds and the native constants.
No new general perturbation principle is claimed.

\subsection{A stronger norm available from the existing word expansion}

Retain V2's notation $q=13\pi b/\kappa$, $s=\pi b/(2\kappa)$
and its function $g$. For the same connected plaquette-set
decomposition, define a size-weighted norm
\[
 \|A\|_\alpha^{\mathrm{sz}}
    =\sup_e\sum_{Z:e\in E(Z)}e^{\alpha|Z|}\|A_Z\|,
                   \qquad \alpha\ge0.
 \tag{V3.1}\label{c18v3:size}
\]
This weights the number of plaquettes, not just their
diameter, and is a bound for the specified decomposition.

The proof of Theorem~\ref{c18v2:uniform-error} gives
rooted order-$r$ norm at most $bc_rq^r$, where
$g(u)=\sum_r c_ru^r$, with $|Z|\le r+1$.
Translation anchoring therefore proves
\[
 \|\mathcal E_3\|_\alpha^{\mathrm{sz}}\le\epsilon_\alpha
 :=12b e^\alpha[g(u)+ug'(u)]_{u=qe^\alpha},
 \qquad qe^\alpha<1/\sqrt5 .
 \tag{V3.2}\label{c18v3:epsilon}
\]
In particular $\epsilon_0=O(b^3/\kappa^2)$.
This stronger bound follows from the actual word
expansion; it is not inferred from exponential diameter
decay alone. We choose its self-adjoint,
translation-covariant representatives
$A_Z=\mathcal E_{3,Z}$ throughout.

\subsection{Split the error before attempting an inverse}

Let $\omega_Z$ be the normalized constant vector on the
links $E(Z)$, $P_Z=|\omega_Z\rangle\langle\omega_Z|$,
$Q_Z=I-P_Z$, and
\[
 T_Z=\kappa\sum_{e\in E(Z)}L_e.
\]
These operators are extended by the identity outside
$E(Z)$ when appropriate. The local electric gap is
$T_Z\ge3\kappa Q_Z$.

\begin{lemma}[Scalar, vacuum-preserving and vacuum-creating parts]
\label{c18v3:split}
For a bounded self-adjoint local term $A_Z$, put
\[
 \begin{aligned}
 c_Z&=\langle\omega_Z,A_Z\omega_Z\rangle,\\
 D_Z&=Q_Z(A_Z-c_ZI)Q_Z,\\
 F_Z&=Q_ZA_ZP_Z+P_ZA_ZQ_Z .
 \end{aligned}
 \tag{V3.3}\label{c18v3:split-def}
\]
Then $A_Z=c_ZI+D_Z+F_Z$, $D_ZP_Z=P_ZD_Z=0$, and
\[
 |c_Z|\le\|A_Z\|,\qquad
 \|D_Z\|\le2\|A_Z\|,\qquad
 \|F_Z\|\le\|A_Z\|.
 \tag{V3.4}\label{c18v3:split-norms}
\]
If $A_Z$ is gauge invariant, all these local operators
are gauge invariant.
\end{lemma}

\begin{proof}
Decompose the local Hilbert space as
$\mathbb C\omega_Z\oplus Q_Z\mathcal H_Z$.
The diagonal vacuum block of $A_Z$ is $c_Z$.
Subtracting $c_ZI$ leaves the indicated $Q_Z$ block.
The off-diagonal block $F_Z$ has norm
$\|Q_ZA_Z\omega_Z\|\le\|A_Z\|$.
The other assertions follow by block multiplication.
Constants and electric projections commute with all
left and right link gauge actions, so the same is true
of the displayed combinations.
\end{proof}

For the complete error define
\[
 C=\sum_Zc_Z,\qquad D_{\mathrm v}=\sum_ZD_Z,\qquad
 F=\sum_ZF_Z,\qquad
 \mathcal E_3=CI+D_{\mathrm v}+F.
 \tag{V3.5}\label{c18v3:global-split}
\]
Thus $D_{\mathrm v}\Omega_{\rm el}=0$ and
\[
 \|D_{\mathrm v}\|_\alpha^{\mathrm{sz}}\le2\epsilon_\alpha,
 \qquad \|F\|_\alpha^{\mathrm{sz}}\le\epsilon_\alpha,
 \qquad
 |\langle v,D_{\mathrm v}v\rangle|
       \le\frac{2\epsilon_0}{3\kappa}\langle v,Tv\rangle.
 \tag{V3.6}\label{c18v3:D-form}
\]
For the form estimate, use
$|D_Z|\le2\|A_Z\|Q_Z$ and
$Q_Z\le(3\kappa)^{-1}T_Z$, and then sum over links.
It holds on $\operatorname{Dom}T^{1/2}$, including on
the physical subspace.

\begin{lemma}[A strictly local inverse for the vacuum source]
\label{c18v3:inverse}
Let $T_Z^{-1}Q_Z$ denote the reduced inverse, and define
\[
 S_Z=T_Z^{-1}Q_ZA_ZP_Z-P_ZA_ZQ_ZT_Z^{-1}.
 \tag{V3.7}\label{c18v3:generator}
\]
Then $S_Z^*=-S_Z$, it is local and gauge invariant, and
\[
 [T,S_Z]=F_Z,\qquad
          \|S_Z\|\le\frac{\|A_Z\|}{3\kappa}.
 \tag{V3.8}\label{c18v3:source-inverse}
\]
The commutator identity holds on $\operatorname{Dom}T$.
Consequently $S_{\mathrm v}=\sum_ZS_Z$ satisfies
\[
 [T,S_{\mathrm v}]=F,\qquad
 \|S_{\mathrm v}\|_\alpha^{\mathrm{sz}}
                         \le\frac{\epsilon_\alpha}{3\kappa}.
 \tag{V3.9}\label{c18v3:S-bound}
\]
\end{lemma}

\begin{proof}
Set $\chi_Z=T_Z^{-1}Q_ZA_Z\omega_Z$.
It belongs to $\operatorname{Dom}T_Z$, is orthogonal
to $\omega_Z$, and has norm at most $\|A_Z\|/(3\kappa)$.
Then
\[
 S_Z=|\chi_Z\rangle\langle\omega_Z|
            -|\omega_Z\rangle\langle\chi_Z|.
\]
This anti-self-adjoint rank-at-most-two operator has norm
$\|\chi_Z\|$. Multiplication by $T_Z$, using
$T_Z\chi_Z=Q_ZA_Z\omega_Z$ and $T_Z\omega_Z=0$, gives
the commutator. It maps $\operatorname{Dom}T_Z$ to itself;
tensoring with spectators extends the identity to
$\operatorname{Dom}T$. The finite-volume sum preserves
that domain by its bounded commutator with $T$.
All ingredients commute with gauge transformations.
The interaction bound follows term by term.
\end{proof}

This inverts only the source from local electric energy
zero to positive local energy. It does not divide by a
difference of two positive energies. The native
equal-energy crossings of f17 V100 are therefore not
removed or contradicted. The pure relative estimate
failed in V2 because $F$ was still present; it now holds
for the explicitly separated $D_{\mathrm v}$.

\subsection{A commutator bound with its loss of size weight}

\begin{lemma}[Connected conjugation estimate]
\label{c18v3:commutator}
For $0\le\alpha'<\alpha$, let $\delta=\alpha-\alpha'$.
Bounded connected interactions satisfy
\[
 \|[A,B]\|_{\alpha'}^{\mathrm{sz}}
 \le\frac{16}{\mathrm e\,\delta}
                 \|A\|_\alpha^{\mathrm{sz}}
                 \|B\|_\alpha^{\mathrm{sz}} .
 \tag{V3.10}\label{c18v3:comm-bound}
\]
If $S$ is anti-self-adjoint and
$\eta=16\|S\|_\alpha^{\mathrm{sz}}/\delta<1$, then
\[
 \|(e^{\operatorname{ad}_S}-1)A\|_{\alpha'}^{\mathrm{sz}}
            \le\frac{\eta}{1-\eta}\|A\|_\alpha^{\mathrm{sz}},
 \qquad \operatorname{ad}_S A=[S,A].
 \tag{V3.11}\label{c18v3:conj-bound}
\]
The assertions concern the connected union decomposition.
\end{lemma}

\begin{proof}
Terms with disjoint link supports commute. For a fixed
link in a contributing union $X\cup Y$, split into the
cases where it belongs to $E(X)$ or to $E(Y)$.
For each fixed $X$,
\[
 \sum_{Y:E(Y)\cap E(X)\ne\varnothing}
              e^{\alpha'|Y|}\|B_Y\|
       \le |E(X)|\,\|B\|_{\alpha'}^{\mathrm{sz}}
       \le4|X|\,\|B\|_\alpha^{\mathrm{sz}}.
\]
Use $|X\cup Y|\le|X|+|Y|$, the commutator norm factor
two, and
$|X|e^{-\delta|X|}\le1/(\mathrm e\delta)$.
The two cases give \eqref{c18v3:comm-bound}.

For $k$ nested commutators, distribute the total loss
$\delta$ equally across the $k$ steps.
At each intermediate weight the norm of $S$ is no
greater than its norm at $\alpha$. Thus
\[
 \frac1{k!}\|\operatorname{ad}_S^k A\|_{\alpha'}^{\mathrm{sz}}
 \le
 \frac{k^k}{\mathrm e^k k!}
 \left(\frac{16\|S\|_\alpha^{\mathrm{sz}}}{\delta}\right)^k
                                  \|A\|_\alpha^{\mathrm{sz}}
 \le\eta^k\|A\|_\alpha^{\mathrm{sz}}.
\]
For example $k!\ge(k/\mathrm e)^k$ follows by integrating
$\log x$ below the sum $\sum_{j=1}^k\log j$.
Sum the geometric majorant. Overlapping connected
supports have a connected union; absolute convergence
justifies the grouping. The bounds apply to finite-volume
operators, and their constants do not depend on volume.
\end{proof}

\subsection{An exact vacuum dressing with a quartic local error}

Put
\[
 B=R+\widehat Z_2+D_{\mathrm v},\qquad
 \rho_*=\rho(\beta)+\frac{2\epsilon_0}{3\kappa}.
 \tag{V3.12}\label{c18v3:B}
\]
Then $B\Omega_{\rm el}=0$ and
$|\langle v,Bv\rangle|\le\rho_*\langle v,Tv\rangle$.
For $\alpha$ in \eqref{c18v3:epsilon}, a useful
explicit interaction bound is
\[
 \|B\|_\alpha^{\mathrm{sz}}\le K_\alpha
 :=4be^\alpha+88sb\,e^{2\alpha}
          +\frac{b^2}{12\kappa}e^\alpha+2\epsilon_\alpha .
 \tag{V3.13}\label{c18v3:K}
\]
Here the four- and eighty-eight-incidence counts are
those of V2; diagonal second-order scalar terms have
one-plaquette support.

\begin{theorem}[Vacuum-source reduction, with an explicit residual]
\label{c18v3:dressing}
Choose $0\le\alpha'<\alpha$ with $qe^\alpha<1/\sqrt5$.
If
\[
            \eta=\frac{16\epsilon_\alpha}
                            {3\kappa(\alpha-\alpha')}<1,
 \tag{V3.14}\label{c18v3:eta}
\]
then the exact finite-volume conjugation satisfies
\[
 e^{S_{\mathrm v}}(T+B+F)e^{-S_{\mathrm v}}
                         =T+B+\mathcal E_4,
 \qquad
 \|\mathcal E_4\|_{\alpha'}^{\mathrm{sz}}
      \le\frac{\eta}{1-\eta}(K_\alpha+\epsilon_\alpha).
 \tag{V3.15}\label{c18v3:quartic}
\]
The error is connected and gauge invariant.
For fixed $\alpha,\alpha'$ and $b/\kappa\to0$, its
bound is $O(b^4/\kappa^3)$, uniformly in $N$.
\end{theorem}

\begin{proof}
Using $[S_{\mathrm v},T]=-F$, the exact formula is
\[
 \mathcal E_4=(e^{\operatorname{ad}_{S_{\mathrm v}}}-1)B
     +\sum_{k\ge1}\frac{k}{(k+1)!}
                         \operatorname{ad}_{S_{\mathrm v}}^k F .
 \tag{V3.16}\label{c18v3:E4}
\]
The first commutator with $T$ is bounded; its domain
is preserved by $e^{S_{\mathrm v}}$.
The integrated conjugation identity therefore justifies
the formula without assuming $[T,A_Z]$ bounded.
Apply Lemma~\ref{c18v3:commutator}; for the second sum
$k/(k+1)!\le1/k!$. Equations
\eqref{c18v3:S-bound} and \eqref{c18v3:K} prove the
bound. Since $\epsilon_\alpha=O(b^3/\kappa^2)$
and $K_\alpha=O(b)$ at fixed weights, the displayed
order follows. All constructions preserve connected
supports and gauge invariance.
\end{proof}

The gap of $T+B$ is at least $3\kappa(1-\rho_*)$
when $\rho_*<1$, by its pure relative bound.
The actual transformed Hamiltonian still contains
$\mathcal E_4$, so this comparison gap is not asserted
to be the exact interacting gap.

\subsection{A bound for the actual ground energy, without closing a gap}

The following argument controls the actual ground state
already after the first two conjugations; it does not
require the condition $\eta<1$.
Let $U=e^{S_2}e^S$ and define the exact normalized trial
\[
 \Phi=U^*\Omega_{\rm el},\qquad
 E_{\mathrm{trial}}=\langle\Phi,H\Phi\rangle
             =bM-\frac{Mb^2}{48\kappa}+C .
 \tag{V3.17}\label{c18v3:trial}
\]
The value $C=\langle\Omega_{\rm el},\mathcal E_3
\Omega_{\rm el}\rangle$ is the complete residual
expectation, not a discarded or formally truncated
coefficient.

\begin{theorem}[Certified energy density and dressed electric density]
\label{c18v3:actual-ground}
Assume $q<1/\sqrt5$ and $\rho_*<1$. Let $\Omega$
be the actual normalized finite-volume ground state
and put $\Psi=U\Omega$. Then
\[
 0\le\frac{E_{\mathrm{trial}}-E_0}{M}
       \le\frac{\epsilon_0^2}{12\kappa(1-\rho_*)},
 \tag{V3.18}\label{c18v3:energy}
\]
and
\[
 \frac1M\langle\Psi,T\Psi\rangle
       \le\frac{\epsilon_0^2}
                         {3\kappa(1-\rho_*)^2}.
 \tag{V3.19}\label{c18v3:electric-density}
\]
These bounds are uniform in $N$. Their right sides
are $O(b^6/\kappa^5)$ as $b/\kappa\to0$.
\end{theorem}

\begin{proof}
The exact operator after subtracting the trial energy is
\[
             U(H-E_{\mathrm{trial}})U^*=T+B+F.
 \tag{V3.20}\label{c18v3:shifted}
\]
Set $f_Z=\|F_Z\|$ and $f_\Sigma=\sum_Z f_Z$.
For every $h>0$, Cauchy--Schwarz and
$2ab\le h a^2+h^{-1}b^2$ give
\[
 |\langle v,F_Zv\rangle|
 \le f_Z\bigl(h\|Q_Zv\|^2+h^{-1}\|P_Zv\|^2\bigr).
\]
The per-link sum of $f_Z$ is at most $\epsilon_0$.
Hence, for normalized $v$,
\[
 |\langle v,Fv\rangle|
 \le \frac{h\epsilon_0}{3\kappa}\langle v,Tv\rangle
                         +h^{-1}f_\Sigma .
 \tag{V3.21}\label{c18v3:Young}
\]
Each nonempty plaquette support has at least four links.
The cubic torus has $M$ links as well as $M$ plaquettes,
so summing incidences gives
\[
          4f_\Sigma\le
          \sum_e\sum_{Z:e\in E(Z)}f_Z\le M\epsilon_0 .
 \tag{V3.22}\label{c18v3:fsum}
\]
If $\epsilon_0=0$ the conclusion is immediate.
Otherwise choose $h=3\kappa(1-\rho_*)/\epsilon_0$.
Since $T+B\ge(1-\rho_*)T$,
\[
 T+B+F\ge-\frac{\epsilon_0 f_\Sigma}
                           {3\kappa(1-\rho_*)}I.
\]
Its expectation in $\Omega_{\rm el}$ is zero, so the
variational principle and \eqref{c18v3:fsum} prove
\eqref{c18v3:energy} for the actual $E_0$.

For the electric density choose instead
$h=3\kappa(1-\rho_*)/(2\epsilon_0)$.
The ground energy of $T+B+F$ is nonpositive, whence
\[
 \frac{1-\rho_*}{2}\langle\Psi,T\Psi\rangle
       \le\frac{2\epsilon_0 f_\Sigma}
                          {3\kappa(1-\rho_*)}.
\]
Use \eqref{c18v3:fsum} again.
All forms are evaluated on their electric form domain;
$\Psi$ belongs to it by the domain-preserving unitary
construction. Finally $\epsilon_0=O(b^3/\kappa^2)$
gives the stated asymptotic order.
\end{proof}

In particular this is not merely a variational upper
bound. It bounds the distance of the trial energy from
the true ground energy, per plaquette. It still does not
bound the first excited energy above that ground.

\subsection{Local marginals of the actual twice-transformed vacuum}

\begin{corollary}[Local fidelity, not global overlap]
\label{c18v3:fidelity}
Under the assumptions of Theorem~\ref{c18v3:actual-ground},
let $X$ be any set of $m$ links and let
$\varrho_X$ be the reduced density matrix of
$\Psi=U\Omega$ in the full Haar tensor product.
With $P_{0,X}$ the product constant projection,
\[
 1-\operatorname{Tr}(\varrho_XP_{0,X})
       \le \frac{m\epsilon_0^2}
                        {3\kappa^2(1-\rho_*)^2},
 \tag{V3.23}\label{c18v3:fidelity-bound}
\]
and
\[
 \|\varrho_X-P_{0,X}\|_1
       \le\frac{2\epsilon_0\sqrt m}
                          {\sqrt3\kappa(1-\rho_*)}.
 \tag{V3.24}\label{c18v3:trace}
\]
The right sides may respectively be capped by one and two.
For fixed $m$ they are $O((b/\kappa)^6)$ and
$O((b/\kappa)^3)$.
\end{corollary}

\begin{proof}
Let $Q_e$ denote the projection off constants on one link.
Then $\sum_e Q_e\le T/(3\kappa)$.
The actual ground is translation invariant: it is the
unique normalized positive ground of the homogeneous
finite-volume Hamiltonian. Indeed, the product electric
heat kernel is strictly positive, and bounded real
magnetic multiplication preserves positivity improvement
by the Feynman--Kac formula. Compact resolvent then gives
a simple positive ground. Gauge and translation symmetries
fix this normalized positive vector. The unitary $U$ commutes with
translations, so $\Psi$ is translation invariant too.

There are $N^3$ links of each fixed orientation.
Without assuming equality between different orientations,
\eqref{c18v3:electric-density} therefore implies
\[
 \langle\Psi,Q_e\Psi\rangle
 \le\frac{\langle\Psi,T\Psi\rangle}{3\kappa N^3}
 \le\frac{\epsilon_0^2}{3\kappa^2(1-\rho_*)^2},
\]
using $M=3N^3$.
The commuting one-link projections give
$I-P_{0,X}\le\sum_{e\in X}Q_e$, proving
\eqref{c18v3:fidelity-bound}.

For completeness let $\zeta=
1-\langle\Psi,P_{0,X}\Psi\rangle$.
If $\zeta<1$, the normalized vector
$P_{0,X}\Psi/\sqrt{1-\zeta}$ is a product of the
local constant vector with an exterior vector.
Its overlap with $\Psi$ is $\sqrt{1-\zeta}$.
The trace distance between the two pure states is
$2\sqrt\zeta$, by diagonalizing their rank-at-most-two
difference. Taking a partial trace cannot increase
trace norm. This proves \eqref{c18v3:trace}.
For $\zeta=1$ use the bound two.
\end{proof}

The physical Hilbert space need not itself factorize:
these reduced matrices are defined in the original
Haar tensor product containing the physical state.
Moreover the estimate is for $U\Omega$, not for the
undressed $\Omega$. For a bounded observable $O$ on $X$,
it controls the original dressed observable $U^*OU$:
\[
 \left|\langle\Omega,U^*OU\Omega\rangle
             -\langle\Omega_{\rm el},O\Omega_{\rm el}\rangle\right|
 \le \|O\|\,\|\varrho_X-P_{0,X}\|_1.
\]
A support growing with volume may make the bound trivial;
no volume-uniform overlap with one global product vector
has been asserted.

\subsection{Why a finite dressing is not an infinite construction}

For the actual centered Hamiltonian, the full scalar
identity after the additional dressing is
\[
 \begin{split}
 e^{S_{\mathrm v}}U(H-E_0)U^*e^{-S_{\mathrm v}}
       &=T+B+\mathcal E_4+(E_{\mathrm{trial}}-E_0)I.
 \end{split}
 \tag{V3.25}\label{c18v3:actual}
\]
Thus all vacuum scalars are accounted for.
The positive comparison gap of $T+B$ does not by itself
absorb $\mathcal E_4$ in an extensive operator norm.

The loss $\delta=\alpha-\alpha'$ in
Lemma~\ref{c18v3:commutator} is also essential to the
present proof. Repeating this estimate with a summable
sequence of losses does not automatically prove
convergence. If one keeps using the inverse of $T$,
the next error contains the commutator of the new
generator with the existing $B$; its estimated size is
linear in the new source and proportional to
$\|B\|/\delta$. The bound may cease to contract as the
losses shrink. A convergent construction needs a
stronger inverse or a separate cluster-stability argument,
not an assertion that each finite-order improvement
can be repeated forever.

The targeted literature check also re-examined the
hypotheses of
\href{https://arxiv.org/abs/math-ph/0412040}
{Yarotsky's Theorems 1 and 2}.
They require a product/classical reference with the
specified perturbative control and fixed interaction
range. They are not invoked here for the infinite-range
remainder without an additional verification.
The already established perturbative lattice regime
must remain distinct from the desired weak-coupling
continuum refinement. No estimate in this append
transports $q<1/\sqrt5$ to that refinement trajectory.

The upstream task is now sharper: control the actual
excited spectrum with the retained interaction included
in the vacuum dressing, and establish a physically
normalized scale evolution that reaches a proved regime.
Neither local vacuum fidelity nor sixth-order
energy-density accuracy supplies those statements.
The full interacting four-dimensional construction
and its positive physical mass gap remain unproved.

\subsection{Checks and source preservation}

Diagnostics check local rank-two inverses, the
size-weight and factorial inequalities, overlapping
finite qutrit decompositions, the exact dressing series,
the energy lower bound and the local trace-distance
inequality. The qutrit fixtures are not Wilson lattices
and do not test the translation step of the native
fidelity proof. Native constant substitutions are
reported separately; no numerical test is promoted
to a volume-uniform theorem. The V2 mathematical
diagnostics are replayed.

V3 appends to the complete V2 body, changing only its
version title. All forty-seven other note files are
unchanged. The active V2 source is replaced by V3;
archives remain intact. Only \texttt{.tex} files are
kept in \texttt{latex\_notes}; rendering and diagnostic
files are kept outside that folder. The next companion
checkpoint remains V5 and the next broad literature
sweep remains V10.

% END CYCLE 18 VERSION 3 APPEND

% BEGIN CYCLE 18 VERSION 4 APPEND
\clearpage
\section{V4: Retained-interaction inversion and a sextic vacuum error}
\label{c18v4:main}

V3 left a linear feedback from the retained interaction $B$
whenever the source inverse was taken against $T$ alone.
We now sum that feedback in a different norm, resolving
vectors by their exact sets of excited links. The electric
denominator then cancels the number of places where a
locally vacuum-annihilating term can act. This proves an
all-orders \emph{linear} inverse with a connected tail.
A local regrouping lemma turns its residual into
vacuum-annihilating interactions. One exact additional
conjugation consequently leaves a sextic local error.

The current V3 source and the archive inventory were
checked first. The excitation overlap rule in f7 V40 is
only a second-order statement and explicitly retains the
spectator-denominator obstruction for a nonconstant
retained band. Here the reference vacuum is rank one,
and the inverse is used on vacuum sources, not on that
retained band. The weighted covariance inverse in f7 V45
is also distinct from the operator proved below.

The organization is related to the excitation-set and
creation-operator method in
\href{https://arxiv.org/abs/math-ph/0411042}
{Yarotsky, Section 2, (7)--(15)}, and to the
\href{https://arxiv.org/abs/0707.1894}
{Bravyi--DiVincenzo--Loss Kirkwood--Thomas analysis}.
Those works explain why an excitation norm can succeed
when an extensive operator norm fails. Their ground-state
or spectral theorems are not invoked for our infinite-range
remainder. We prove the needed untruncated inverse and
the quantitative local regrouping directly. No general
priority claim is made.

\subsection{Excitation sectors and the price of local branching}

Write $\Delta=3\kappa$, $\omega=\Omega_{\rm el}$ and
$G=T+B$, where $B$ is exactly \eqref{c18v3:B}.
Use a connected decomposition $B=\sum_X B_X$ in which
\[
 B_XP_X=P_XB_X=0,\qquad
 m_X=|E(X)|\le4|X|.
 \tag{V4.1}\label{c18v4:B-local}
\]
The representatives of $R$, $\widehat Z_2$ and
$D_{\mathrm v}$ have this property by V1--V3.
Their bounds give
$\|B\|_\gamma^{\mathrm{sz}}\le K_\gamma$.
No claim that arbitrary globally centered terms have
\eqref{c18v4:B-local} is being used.

For a set $I$ of links let
\[
 \Pi_I=\bigotimes_{e\in I}Q_e
             \bigotimes_{e\notin I}P_e,\qquad
 P_e=|1\rangle\langle1|,\quad Q_e=\mathbf1-P_e .
\]
Here $\mathbf1$ denotes the identity. The $\Pi_I$ are mutually orthogonal,
sum to the identity, and commute with $T$. On a
nonempty excitation sector,
\[
 T\Pi_I\ge\Delta|I|\Pi_I,\qquad
 \|T^{-1}\Pi_I\|\le(\Delta|I|)^{-1}.
 \tag{V4.2}\label{c18v4:sector-gap}
\]
The inverse is zero on $\mathbb C\omega$.

Define the branching-weighted interaction bound
\[
 J_\alpha(B)=\sup_e\sum_{X:e\in E(X)}
       e^{\alpha|X|}2^{m_X/2}\|B_X\|.
 \tag{V4.3}\label{c18v4:J}
\]
It is finite whenever V3's stronger norm is available:
\[
 J_\alpha(B)\le K_{\alpha+2\log2}.
 \tag{V4.4}\label{c18v4:J-K}
\]
This extra size weight is explicit. Infinite dimension
of a single-link Hilbert space creates no extra factor:
the branching counts vacuum/nonvacuum \emph{sectors},
not a basis inside those sectors.

\begin{lemma}[Excitation-number cancellation]
\label{c18v4:branch}
Let $v_I\in\Pi_I\mathcal H$, $I\ne\varnothing$.
If $E(X)\cap I=\varnothing$, then $B_Xv_I=0$.
Otherwise $B_Xv_I$ has unchanged excitations outside
$E(X)$ and a nonempty excitation set inside $E(X)$.
Its sector components obey
\[
 \sum_J\|\Pi_J B_XT^{-1}v_I\|
 \le
 \frac{2^{m_X/2}\|B_X\|}{\Delta|I|}\|v_I\|.
 \tag{V4.5}\label{c18v4:branch-bound}
\]
\end{lemma}

\begin{proof}
The first assertion follows from $B_XP_X=0$.
The relation $P_XB_X=0$ prevents the output from being
the local constant vector. There are at most $2^{m_X}-1$
possible nonempty output excitation subsets of $E(X)$;
the status of each exterior link is fixed.
Their projections are orthogonal. Cauchy--Schwarz over
this finite list gives the factor
$\sqrt{2^{m_X}-1}\le2^{m_X/2}$, irrespective of the
dimensions of the sectors. Apply
\eqref{c18v4:sector-gap} to the input.
\end{proof}

\subsection{A connected inverse with the full retained interaction}

Fix a nonempty connected root plaquette set $Y$.
Keep vector families $v_{Z,I}$ where $Z$ is connected,
$Y\subseteq Z$, $\varnothing\ne I\subseteq E(Z)$, and
$v_{Z,I}\in\Pi_I\mathcal H$. The cover $Z$ is a
bookkeeping label; $I$ is the actual excitation set.
Allow several labels for the same physical vector and
use the absolute family norm
\[
 \|v\|_{\alpha;Y}
   =\sum_{Z,I}e^{\alpha(|Z|-|Y|)}\|v_{Z,I}\|.
 \tag{V4.6}\label{c18v4:root-norm}
\]
The lift of a family is $\mathcal L v=\sum_{Z,I}v_{Z,I}$.
It is a bounded map from this norm into the Hilbert space.

Define a family map $\mathfrak M$ by acting with
$B_XT^{-1}$ on each component, resolving the result into
its exact excitation sectors, and assigning cover
$Z\cup X$. Terms with $E(X)\cap I=\varnothing$ vanish.
Every remaining union is connected.

\begin{theorem}[Retained-interaction source inverse]
\label{c18v4:inverse}
If
\[
 r_\alpha:=J_\alpha(B)/\Delta<1,
 \tag{V4.7}\label{c18v4:r-exact}
\]
then $\|\mathfrak M\|_{\alpha;Y}\le r_\alpha$.
For a vacuum-orthogonal source $f_Y$ supported on
$E(Y)$ with constant exterior, decompose it into
$\Pi_I f_Y$ with cover $Y$. The family
\[
 \chi^{[Y]}
   =T^{-1}\sum_{n=0}^\infty(-\mathfrak M)^n f_Y
 \tag{V4.8}\label{c18v4:inverse-series}
\]
converges in the rooted norm and in the lifted
electric graph norm. Its lift solves
\[
 G\mathcal L\chi^{[Y]}=f_Y,\qquad
             \mathcal L\chi^{[Y]}\perp\omega.
 \tag{V4.9}\label{c18v4:source-equation}
\]
For any upper bound $r_\alpha\le r<1$,
\[
 \begin{split}
 \|\chi^{[Y]}\|_{\alpha;Y}
   &\le\frac{2^{m_Y/2}\|f_Y\|}{\Delta(1-r)},\\
 \|T\chi^{[Y]}\|_{\alpha;Y}
   &\le\frac{2^{m_Y/2}\|f_Y\|}{1-r}.
 \end{split}
 \tag{V4.10}\label{c18v4:inverse-bounds}
\]
The $T\chi$ family with $n\ge1$ has the same second
bound multiplied by $r$. For $\alpha>0$, the sum of
the norms of terms with $|Z\setminus Y|\ge L$ is at most
\[
 e^{-\alpha L}
       \frac{2^{m_Y/2}\|f_Y\|}{\Delta(1-r)}.
 \tag{V4.11}\label{c18v4:inverse-tail}
\]
All constants are independent of the torus size.
\end{theorem}

\begin{proof}
The output weight relative to the input weight is at
most $e^{\alpha|X|}$. For each fixed $I$,
\[
 \sum_{X:E(X)\cap I\ne\varnothing}
       e^{\alpha|X|}2^{m_X/2}\|B_X\|
       \le |I|J_\alpha(B).
\]
This cancels $|I|^{-1}$ in
\eqref{c18v4:branch-bound}. Sum over all input
components and use the triangle inequality. Thus
$\mathfrak M$ is a contraction without spending a new
weight at each step.

Orthogonality of the initial excitation sectors gives
$\sum_I\|\Pi_I f_Y\|\le2^{m_Y/2}\|f_Y\|$.
The Neumann sum for
$y=\sum_{n\ge0}(-\mathfrak M)^nf_Y$ converges absolutely;
the sum starting at $n=1$ is bounded by $r/(1-r)$
times the initial norm. Applying $T^{-1}$ costs at most
$\Delta^{-1}$. Each component then belongs to
$\operatorname{Dom}T$, and the sum of the norms of
their $T$ images is finite. Closedness of $T$ shows
that the lifted sum lies in its domain.

The lift intertwines $\mathfrak M$ with $BT^{-1}$.
At fixed finite volume $B$ is a bounded operator.
Consequently $(1+BT^{-1})\mathcal L y=f_Y$, which
proves \eqref{c18v4:source-equation}.
For uniqueness, local vacuum annihilation gives
\[
 |\langle v,Bv\rangle|
 \le\frac{\|B\|_0^{\mathrm{sz}}}{\Delta}
                           \langle v,Tv\rangle .
\]
Indeed $|B_X|\le\|B_X\|Q_X$, $Q_X\le T_X/\Delta$,
and one sums the link incidences. The relative constant
is at most $r<1$, so $G$ has kernel $\mathbb C\omega$
and is invertible on its orthogonal complement.
Finally remove the weight from the terms with
$|Z|-|Y|\ge L$ to obtain \eqref{c18v4:inverse-tail}.
\end{proof}

This is a convergent inverse for a specified source,
not a bounded local inverse of the full commutator
$[G,\cdot]$ on arbitrary operators. In particular,
no difference between positive electric energies is
being divided out. The native band resonances remain.

\subsection{From rooted vectors to a homogeneous local generator}

Apply the theorem with
$f_Y=F_Y\omega$, where the $F_Y$ are exactly V3's
self-adjoint local vacuum-source terms.
Write each resulting component as a vector on $E(Z)$
with a constant exterior, and define
\[
 S_{Y,Z,I}
   =|\chi^{[Y]}_{Z,I}\rangle\langle\omega_Z|
       -|\omega_Z\rangle\langle\chi^{[Y]}_{Z,I}|,
 \qquad
 S_B=\sum_{Y,Z,I}S_{Y,Z,I}.
 \tag{V4.12}\label{c18v4:S}
\]
The displayed local vectors include their constant
factors on $E(Z)\setminus I$. Thus
$\|S_{Y,Z,I}\|=\|\chi^{[Y]}_{Z,I}\|$.
Let $S^{(0)}$ denote the $n=0$ contribution.

Choose parameters
\[
 \begin{gathered}
 a\ge0,\quad d_0,d_1>0,\quad
 \tau=a+2\log2,\quad \sigma=\tau+d_1,\\
 \alpha=\sigma+d_0,\qquad \lambda=\alpha+2\log2,
 \qquad qe^\lambda<1/\sqrt5,\\
 r:=K_\lambda/\Delta<1,\qquad
 \mathfrak s:=\frac{3\epsilon_\lambda}
                  {\mathrm e\,d_0\Delta(1-r)} .
 \end{gathered}
 \tag{V4.13}\label{c18v4:parameters}
\]

\begin{lemma}[Translation anchoring and the electric tail]
\label{c18v4:anchoring}
On the homogeneous cubic torus the construction is
translation covariant and satisfies
\[
 \|S_B\|_\sigma^{\mathrm{sz}}\le\mathfrak s,\qquad
 \|[T,S_B-S^{(0)}]\|_\sigma^{\mathrm{sz}}
                              \le\Delta r\mathfrak s,
 \qquad [T,S^{(0)}]=F .
 \tag{V4.14}\label{c18v4:anchor-bounds}
\]
It is gauge invariant, anti-self-adjoint and preserves
$\operatorname{Dom}T$. Moreover
\[
 S_B\omega=G_{\omega^\perp}^{-1}F\omega,\qquad
              [G,S_B]\omega=F\omega.
 \tag{V4.15}\label{c18v4:vacuum-identity}
\]
\end{lemma}

\begin{proof}
Translations preserve the initial decomposition, all
excitation projections, $T^{-1}$ and the family map.
Fix a link of one orientation. Its anchored sum is
constant along an orbit of $N^3$ links and is therefore
at most $N^{-3}$ times the sum over all links.
For each cover, $|E(Z)|\le4|Z|$ and
\[
 |Z|e^{\sigma|Z|}
       \le(\mathrm e\,d_0)^{-1}e^{\alpha|Z|}.
\]
By \eqref{c18v4:inverse-bounds}, this bounds the
anchored norm of $S_B$ by
\[
 \frac{4}{N^3\mathrm e\,d_0\Delta(1-r)}
       \sum_Y e^{\alpha|Y|}2^{m_Y/2}\|F_Y\|.
\]
Every root has at least four links. Since
$M=3N^3$ is the number of links,
\[
 4\sum_Y e^{\lambda|Y|}\|F_Y\|
 \le M\|F\|_\lambda^{\mathrm{sz}}
 \le M\epsilon_\lambda .
\]
Together with $2^{m_Y/2}\le e^{2(\log2)|Y|}$,
this proves the first bound. For a local rank-two
generator the norm of its commutator with $T$ equals
the norm of $T\chi$. The $n\ge1$ part of
\eqref{c18v4:inverse-bounds} proves the second.

At $n=0$, summing the exact excitation components on
the original cover $Y$ gives precisely the local
rank-two generator of V3. Hence $[T,S^{(0)}]=F$
as an operator, not just on $\omega$.
All graph-norm sums are absolutely convergent in finite
volume. Their bounded commutators with $T$ give domain
preservation for $S_B$ and $e^{S_B}$.

Each excitation projector commutes with link gauge
actions. Starting from gauge-invariant sources, the
inverse and all the local rank-two lifts are therefore
gauge invariant. Finally $S_B\omega$ is the sum of
the lifted inverse sources, so
\eqref{c18v4:source-equation} and $G\omega=0$ prove
\eqref{c18v4:vacuum-identity}.
\end{proof}

Only the single anchoring loss $d_0$ was spent after
summing the linear inverse. It was not spent once per
Neumann factor. The distinction between an exact
vacuum identity and a full operator identity in
\eqref{c18v4:vacuum-identity} will now matter.

\subsection{Regroup a globally vanishing vacuum source locally}

\begin{lemma}[Canonical local vacuum regrouping]
\label{c18v4:regroup}
Let $A=\sum_Z A_Z$ be a bounded self-adjoint connected
interaction in finite volume and suppose $A\omega=0$.
There is a connected decomposition
$A=\sum_Z A_Z^\circ$ such that
$A_Z^\circ P_Z=P_ZA_Z^\circ=0$ and
\[
 \|A^\circ\|_a^{\mathrm{sz}}
       \le4\|A\|_{a+2\log2}^{\mathrm{sz}} .
 \tag{V4.16}\label{c18v4:regroup-bound}
\]
Gauge and translation invariance of the original
representatives are preserved. In particular,
\[
 |\langle v,Av\rangle|
 \le\frac{4\|A\|_{a+2\log2}^{\mathrm{sz}}}{\Delta}
                                      \langle v,Tv\rangle .
 \tag{V4.17}\label{c18v4:regroup-form}
\]
The extra size weight is a hypothesis, not a
consequence of $A\omega=0$ alone.
\end{lemma}

\begin{proof}
Set $c_Z=\langle\omega_Z,A_Z\omega_Z\rangle$ and
$f_Z=(A_Z-c_Z)\omega_Z$. Resolve $f_Z$ into its exact
nonempty local excitation subsets $I\subset E(Z)$.
After removing constant spectator factors, denote
these vectors by $f_{Z,I}$ in
$\bigotimes_{e\in I}Q_e\mathcal H_e$.
For such a vector put
\[
 C_I(f_{Z,I})=|f_{Z,I}\rangle\langle\omega_I|,
\]
extended by the identity on other links. Unlike a
rank-two source projected onto $\omega_Z$, this
creation operator puts no vacuum constraint on
$E(Z)\setminus I$. Define the operator on $E(Z)$
\[
 A_Z^\circ=A_Z-c_ZI-
       \sum_{\varnothing\ne I\subset E(Z)}
          \bigl(C_I(f_{Z,I})+C_I(f_{Z,I})^*\bigr).
 \tag{V4.18}\label{c18v4:regroup-def}
\]
It is self-adjoint and annihilates $\omega_Z$.

The identity $A\omega=0$ implies both
$\sum_Zc_Z=0$ and, for each fixed nonempty global
excitation set $I$, $\sum_Zf_{Z,I}=0$.
This follows by projecting onto the mutually
orthogonal global sectors. Since $C_I$ is linear
in its vector and uses the same identity extension
for every $Z$, its total sum is zero. Thus
$\sum_Z A_Z^\circ=A$ as operators. This is why merely
discarding locally off-diagonal terms would have
been incorrect.

Orthogonality gives
$\sum_I\|f_{Z,I}\|\le2^{m_Z/2}\|A_Z\|$.
Consequently
\[
 \|A_Z^\circ\|
   \le(2+2\,2^{m_Z/2})\|A_Z\|
   \le4\,2^{m_Z/2}\|A_Z\|.
\]
Keep cover $Z$ for each term, even if its minimal
support is smaller. This preserves connected labels
and proves \eqref{c18v4:regroup-bound}.
The local vacuum-annihilation form estimate used in
Theorem~\ref{c18v4:inverse} proves
\eqref{c18v4:regroup-form}.
All sums converge absolutely under the stated
weighted bound. Gauge and translation covariance
follow from the equivariance of the projections and
of the creation maps.
\end{proof}

\subsection{One exact Newton-type vacuum step}

Define the \emph{operator}
\[
 D_B=F-[G,S_B].
 \tag{V4.19}\label{c18v4:DB}
\]
It is self-adjoint, and
\eqref{c18v4:vacuum-identity} says $D_B\omega=0$.
It need not have zero local vacuum sources term by
term before regrouping. Put
\[
 \begin{gathered}
 \mathfrak a
   :=\mathfrak s\left(\Delta r+
                       \frac{16K_\sigma}{\mathrm e\,d_1}\right),
 \qquad
 \eta_B:=\frac{16\mathfrak s}{\tau-a},\\
 \nu:=\frac{\eta_B}{1-\eta_B}
                         (\epsilon_\tau+\mathfrak a).
 \end{gathered}
 \tag{V4.20}\label{c18v4:budgets}
\]

\begin{theorem}[Sextic residual with the retained feedback summed]
\label{c18v4:newton}
Assume \eqref{c18v4:parameters} and $\eta_B<1$.
There is an exact domain-preserving identity
\[
 e^{S_B}(T+B+F)e^{-S_B}
            =T+B+D_B+\mathcal E_6.
 \tag{V4.21}\label{c18v4:normal}
\]
The term $D_B$ has a connected, locally
vacuum-annihilating decomposition with
\[
 \|D_B^\circ\|_a^{\mathrm{sz}}\le4\mathfrak a,\qquad
 |\langle v,D_Bv\rangle|
               \le(4\mathfrak a/\Delta)\langle v,Tv\rangle,
 \tag{V4.22}\label{c18v4:DB-bound}
\]
and the connected error satisfies
\[
          \|\mathcal E_6\|_a^{\mathrm{sz}}\le\nu.
 \tag{V4.23}\label{c18v4:E6-bound}
\]
At fixed weights and $b/\kappa\to0$,
\[
 \mathfrak s=O(b^3/\kappa^3),\quad
 \mathfrak a=O(b^4/\kappa^3),\quad
                  \nu=O(b^6/\kappa^5),
 \tag{V4.24}\label{c18v4:orders}
\]
uniformly in $N$. Gauge invariance is preserved.
\end{theorem}

\begin{proof}
Cancel the $n=0$ inverse term using
$[T,S^{(0)}]=F$. Then
\[
 D_B=-[T,S_B-S^{(0)}]-[B,S_B].
\]
The first term has $\sigma$ norm at most
$\Delta r\mathfrak s$ by
\eqref{c18v4:anchor-bounds}. Apply V3's connected
commutator estimate from $\sigma$ to $\tau$ to the
second. Thus there is a self-adjoint connected
decomposition with
\[
                 \|D_B\|_\tau^{\mathrm{sz}}\le\mathfrak a.
 \tag{V4.25}\label{c18v4:DB-before}
\]
Its commutators are bounded; no boundedness of
$[T,B_X]$ is assumed. Since $D_B\omega=0$,
Lemma~\ref{c18v4:regroup}, with $\tau=a+2\log2$,
gives \eqref{c18v4:DB-bound}.

For the conjugation, $[S_B,G]=-F+D_B$.
Its first commutator is bounded and $S_B$ preserves
the electric domain. The integrated conjugation
formula therefore yields the exact series
\[
 \mathcal E_6
   =\sum_{j\ge1}\frac{j}{(j+1)!}
                         \operatorname{ad}_{S_B}^{\,j}F
    +\sum_{j\ge1}\frac{1}{(j+1)!}
                         \operatorname{ad}_{S_B}^{\,j}D_B .
 \tag{V4.26}\label{c18v4:E6-series}
\]
For this estimate use the \emph{unregrouped}
$\tau$-decomposition in \eqref{c18v4:DB-before}.
Both coefficient sequences are bounded by $1/j!$.
V3's conjugation lemma from $\tau$ to $a$, using
$\|S_B\|_\tau^{\mathrm{sz}}\le\mathfrak s$,
proves \eqref{c18v4:E6-bound}. All operator products
are assigned to connected unions.

Finally $K_\gamma=O(b)$ and
$\epsilon_\gamma=O(b^3/\kappa^2)$ at every fixed
admissible weight. Hence $r=O(b/\kappa)$,
$\mathfrak s=O(b^3/\kappa^3)$ and
$\mathfrak a=O(b^4/\kappa^3)$. Substitution into
\eqref{c18v4:budgets} proves the sextic order.
\end{proof}

The linear retained feedback has not been ignored:
it is summed in $S_B$ and reappears in the exactly
vacuum-preserving $D_B$. The remaining vacuum-changing
error is quadratic in the original source size at
fixed small retained interaction.

If $\rho_*+4\mathfrak a/\Delta<1$, then
$T+B+D_B$ has comparison gap at least
$\Delta(1-\rho_*-4\mathfrak a/\Delta)$.
This remains a comparison gap, not the gap after
adding $\mathcal E_6$.

\subsection{Actual ground-energy accuracy after the new conjugation}

Let
\[
 C_6=\langle\omega,\mathcal E_6\omega\rangle,\qquad
 \Phi_4=U^*e^{-S_B}\omega,\qquad
 E_{\mathrm{trial},4}
       =E_{\mathrm{trial}}+C_6,\qquad
 \rho_4=\rho_*+\frac{4\mathfrak a+2\nu}{\Delta}.
 \tag{V4.27}\label{c18v4:trial4}
\]
The scalar $C_6$ is the exact expectation of the
complete remainder. It is neither omitted nor
replaced by a formal coefficient.

\begin{corollary}[Twelfth-order certified trial-energy error]
\label{c18v4:energy}
Under Theorem~\ref{c18v4:newton}'s assumptions and
$\rho_4<1$, the actual native ground energy obeys
\[
 0\le\frac{E_{\mathrm{trial},4}-E_0}{M}
                   \le\frac{\nu^2}{4\Delta(1-\rho_4)}.
 \tag{V4.28}\label{c18v4:energy-bound}
\]
For the actual transformed ground
$\Psi_4=e^{S_B}U\Omega$,
\[
 \frac1M\langle\Psi_4,T\Psi_4\rangle
                  \le\frac{\nu^2}{\Delta(1-\rho_4)^2}.
 \tag{V4.29}\label{c18v4:density-bound}
\]
Both right sides are $O(b^{12}/\kappa^{11})$ at
fixed weights as $b/\kappa\to0$, uniformly in $N$.
\end{corollary}

\begin{proof}
Split each local term of $\mathcal E_6$ as in
Lemma~\ref{c18v3:split}, giving
$\mathcal E_6=C_6I+D_6+F_6$.
The zero-weight norms satisfy
$\|D_6\|_0^{\mathrm{sz}}\le2\nu$ and
$\|F_6\|_0^{\mathrm{sz}}\le\nu$.
Then $B+D_B+D_6$ annihilates $\omega$ and has
pure relative electric form bound at most $\rho_4$.
The exact centered identity is
\[
 e^{S_B}U(H-E_{\mathrm{trial},4})U^*e^{-S_B}
                         =T+B+D_B+D_6+F_6 .
\]
In particular $\Phi_4$ has the claimed trial energy.
For $f_\Sigma=\sum_Z\|F_{6,Z}\|$, connected
nonempty plaquette covers still give
$4f_\Sigma\le M\nu$. V3's Young inequality, now with
$\Delta=3\kappa$, gives for every normalized $v$
\[
 |\langle v,F_6v\rangle|
       \le(h\nu/\Delta)\langle v,Tv\rangle
                                      +h^{-1}f_\Sigma .
\]
Use $h=\Delta(1-\rho_4)/\nu$ for the energy and
$h=\Delta(1-\rho_4)/(2\nu)$ for the electric density.
The case $\nu=0$ is immediate. These are exactly
the variational and form arguments of V3 applied
to the new, explicitly proved residual.
Equation~\eqref{c18v4:orders} gives the stated order.
\end{proof}

For the actual $E_0$, all scalar bookkeeping is thus
visible:
\[
 e^{S_B}U(H-E_0)U^*e^{-S_B}
 =T+B+D_B+D_6+F_6+
                   (E_{\mathrm{trial},4}-E_0)I .
\]
No excited-state estimate follows merely from the
smallness of the last scalar or the electric density.

\subsection{Quantitative scope and the next upstream issue}

For example fix $a=0.1$ and $d_0=d_1=0.5$.
Then $\lambda=0.1+4\log2+1$.
The explicit budgets above are nonempty for small
$b/\kappa$; substitutions at $\kappa=1$ give:
\[
\begin{array}{c|c|c|c}
 b/\kappa&r\text{ (upper bound)}&\eta_B&
                              \nu\text{ (upper bound)}\\ \hline
 10^{-6}&6.43\,10^{-5}&1.71\,10^{-7}&4.02\,10^{-18}\\
 10^{-5}&6.66\,10^{-4}&1.78\,10^{-4}&1.78\,10^{-11}\\
 3\,10^{-5}&2.44\,10^{-3}&5.34\,10^{-3}&4.63\,10^{-8}
\end{array}
\]
The entries are rounded upward illustrative substitutions,
not rigorous floating-point certificates. The exact
inequalities, rather than these numbers, are the
hypotheses and conclusions. Constants are conservative:
a higher asymptotic order does not imply that this
bound is smaller than V3's at every coupling.

The all-orders result in this append is the \emph{linear}
inverse of $T+B$ on a vacuum source. Only one nonlinear
conjugation has been completed. Repeating the construction
as written requires the input weight
$\lambda=a+4\log2+d_0+d_1$ at every step.
The two branching/regrouping costs cannot be paid
infinitely often from a fixed weight supplied by V3.
It would be circular to call the quadratic residual
alone a convergent infinite dressing theorem.

The next useful step is therefore a nonlinear
excitation/creation norm that retains the cancellation
without spending that fixed branching cost at every
Newton step, or an independently verified stability
argument that applies to the resulting interaction.
This is more specific than V3's unsummed linear feedback:
the latter now has an explicit convergent inverse and
an operator-level local regrouping.

Nor has the regime changed to physical weak-coupling
refinement. The strengthened condition
$qe^\lambda<1/\sqrt5$ is still an electric-dominated
fixed-regulator condition. No physical normalization,
scale transport, continuum existence or positive
continuum mass has been proved by this append.
The full interacting four-dimensional Yang--Mills
construction and mass gap remain unproved.

\subsection{Verification and preservation}

Finite diagnostics test the rooted Neumann contraction,
its excitation support and connected tails, the local
rank-two lift, cancellation of a nonzero operator's
global vacuum source by canonical regrouping, and the
exact sextic-error series. Three four-qutrit fixtures
check quadratic scaling in the vacuum-source amplitude;
they are not Wilson lattices and do not verify the
homogeneous translation-anchoring theorem. Native
bound substitutions are reported separately. The
unbounded-domain and volume-uniform arguments are
the proofs above, not consequences of finite tests.
V3's mathematical diagnostics are also replayed.

The full V3 body is preserved, with only its version
title changed. V4 replaces the active V3 source;
all forty-seven other note files and the archives
remain unchanged. Only \texttt{.tex} files are stored
in \texttt{latex\_notes}. Build and verification
products stay outside that folder.
The next companion review is V5 and the next broad
literature sweep is V10.

% END CYCLE 18 VERSION 4 APPEND

% BEGIN CYCLE 18 VERSION 5 APPEND
\clearpage
\section{V5: A nonlinear creation contraction and an actual lattice gap}
\label{c18v5:main}

V4 summed the retained linear feedback but still paid
a fixed size-weight cost at each proposed nonlinear
step. This append goes upstream to the original
four-link Wilson interaction. Pure creation operators
keep every exterior link in a definite vacuum or
nonvacuum sector. Consequently the only sector
branching occurs on the four links of the primitive
plaquette, even when a creation coefficient has a
large connected cover.

An energy-weighted creation norm now yields a
nonlinear contraction with \emph{no loss of spatial
weight per iteration}. Its derivative also excludes
nearby excited eigenvalues of the actual Hamiltonian.
This closes the present small-coupling lattice
construction, not the four-dimensional continuum
problem. Strong-coupling lattice gaps are already
known, including elsewhere in this corpus; no novelty
of that qualitative conclusion is claimed.

This is a direct untruncated application of the
Kirkwood--Thomas/creation-operator strategy. Primary
context is
\href{https://arxiv.org/abs/0707.1894}
{Bravyi--DiVincenzo--Loss, Sections 2.1--2.4}
and
\href{https://arxiv.org/abs/math-ph/0411042}
{Yarotsky, Section 2}.
The finite-range structure, energy norm, connected
cover bookkeeping and numerical constants below are
proved explicitly. We do not import a qubit theorem
as a theorem about infinite-dimensional Haar links.

\subsection{The required companion checkpoint}

A fresh full-file hash comparison found all eight
sources unchanged since their last relevant review.
Their terminal scope and conclusion passages were
reread: NS stability V26, SM--Einstein V72, YM shell
mixing V16, parallel YM V5, perturbative ESM V64,
SM source selection V52, regenerative stationarity
V54, and the exact-action Einstein bridge V45.
The four latter sources are the current supplied
Downloads files. The comparison and paths are
recorded outside the notes folder in
\path{artifacts/ym_c18_v5_companion_review.json}.
This was not an independent recertification of all
inherited proofs.

The NS rank-one kernel constants and SM many-copy,
fixed-time result do not supply a native YM spectral
estimate. Shell mixing V16 still requires the general
noncommuting mixed-curvature calculation and a common
nonperturbative remainder. Parallel YM V5 rejects
volume-uniform tightness of an extensive total
topological charge as a general target.
ESM V64 keeps the populated common critical source
array open. Source selection V52 keeps the actual
preparation/readout bank open. Stationarity V54
distinguishes its linear prepared descriptor response
from a nonlinear original history. Bridge V45 retains
the moving-source dynamics and shrinking selected
interval rather than asserting a continuum bridge.

The useful common lesson is scope, not imported
coercivity: control the actual equation in a fixed
norm and identify the actual spectral object.
No companion coefficient or auxiliary gap is used
in the theorem below. The next companion checkpoint
and broad literature checkpoint are both V10.

\subsection{The original Hamiltonian and the creation space}

Retain the native Hamiltonian \eqref{c18v1:H} and
write
\[
 H=T+bM+V,\qquad
 V=\sum_p V_p,\quad V_p=-bw_p,\quad
 \|V_p\|\le b,\qquad \Delta=3\kappa .
 \tag{V5.1}\label{c18v5:H}
\]
There are four links per plaquette and four plaquettes
incident to each link on the cubic torus $N\ge3$:
\[
       \sup_e\sum_{p:e\in E(p)}\|V_p\|\le4b.
 \tag{V5.2}\label{c18v5:incidence}
\]
The identity shift $bM$ will be restored exactly.

For $\varnothing\ne I$ a set of links, put
\[
 \mathcal H_I^\circ=\bigotimes_{e\in I}Q_e\mathcal H_e,
 \qquad T_I=\kappa\sum_{e\in I}L_e,\qquad
 T_I\ge\Delta|I|\quad\hbox{on }\mathcal H_I^\circ.
\]
For $u\in\mathcal H_I^\circ$ define the pure creation
operator
\[
       \widehat u_I=|u\rangle\langle\omega_I|,
 \tag{V5.3}\label{c18v5:creation}
\]
extended by the identity outside $I$.
The vector $u$ may be entangled between its links.
A connected plaquette cover $Z$ with $I\subseteq E(Z)$
is retained as a label, not as an additional vacuum
projection.

A coefficient family is
$c=(c_{Z,I})$, with
$c_{Z,I}\in\operatorname{Dom}T_I\cap\mathcal H_I^\circ$.
Define
\[
 \begin{split}
 C(c)&=\sum_{Z,I}\widehat{c_{Z,I}}_I,\\
 \|c\|_{a,E}
   &=\sup_e\sum_{Z,I:e\in E(Z)}
                 e^{a|Z|}\|T_Ic_{Z,I}\|,\qquad a\ge0 .
 \end{split}
 \tag{V5.4}\label{c18v5:norm}
\]
All covers are nonempty and connected. At each finite
volume this is a Banach space: there are finitely many
labels, and $T_I^{-1}$ is bounded on the nonvacuum
sector. In particular,
\[
 \|c_{Z,I}\|
       \le\frac{\|T_Ic_{Z,I}\|}{\Delta|I|}.
 \tag{V5.5}\label{c18v5:energy-pay}
\]
The finite-volume operators $C(c)$ and $[T,C(c)]$
are bounded, although their operator norms may grow
with volume. Their interaction norms and the
contraction constants below need not grow.

\begin{lemma}[Creation algebra and finite commutator degree]
\label{c18v5:algebra}
Pure creation operators commute. Their product is zero
if their excitation sets intersect, and otherwise is
the creation operator of their tensor-product vector.
Moreover, on $\operatorname{Dom}T$,
\[
 [T,\widehat u_I]=\widehat{T_Iu}_I,\qquad
       [C(c),[T,C(d)]]=0.
 \tag{V5.6}\label{c18v5:T-algebra}
\]
For arbitrary nine creation operators,
\[
 \operatorname{ad}_{C_1}\cdots
           \operatorname{ad}_{C_9}V_p=0 .
 \tag{V5.7}\label{c18v5:degree}
\]
Thus, with $C=C(c)$,
\[
 \begin{split}
 e^{-C}Te^C&=T+[T,C],\\
 W(c):=e^{-C}Ve^C
    &=\sum_{k=0}^{8}\frac{(-1)^k}{k!}
                              \operatorname{ad}_C^k V .
 \end{split}
 \tag{V5.8}\label{c18v5:similarity}
\]
These are exact domain-preserving identities.
\end{lemma}

\begin{proof}
If $I$ and $J$ intersect, the vacuum bra of one
creation operator meets a nonvacuum factor created
by the other, so either product is zero. For disjoint
sets the operators commute and their product has the
stated tensor vector. This argument uses the
vacuum/nonvacuum projections and does not assume
that $u$ factorizes within $I$.
Since $T_I\omega_I=0$, the first commutator follows.
Its right side is another pure creation operator,
which proves the double-commutator assertion.

Consider a nested commutator with one primitive $V_p$.
The creation derivations commute. A component whose
excitation set misses $E(p)$ can therefore be moved
next to $V_p$, where its commutator is zero.
In any contributing component all nine excitation
sets meet $E(p)$. Expand the commutator into products
with creation factors on the left and right of $V_p$.
Nonzero factors on either side must have pairwise
disjoint excitation sets. Each such set occupies
at least one of the four plaquette links.
There can therefore be at most four creation factors
on either side, and not nine in total. This proves
\eqref{c18v5:degree}. The same argument applies to
arbitrary vector-valued creation coefficients.

Every creation factor maps its electric domain to
itself and has bounded commutator with $T$.
The finite-volume sums preserve the domain.
Also $C^{M+1}=0$, because a nonzero product creates
at least one new excited link per factor.
The exponentials are finite polynomials in $C$,
are inverse bounded operators, and preserve the
electric domain. Formula \eqref{c18v5:T-algebra}
and the bounded-operator conjugation formula for $V$
give \eqref{c18v5:similarity}.
\end{proof}

The operator $e^C$ is not unitary. No unitary or
volume-uniform operator-norm bound for this
similarity will be assumed.

\subsection{Only the primitive plaquette pays for sector branching}

Let $\omega=\Omega_{\rm el}$ and
$Q=\mathbf1-|\omega\rangle\langle\omega|$.
For each term in \eqref{c18v5:similarity}, attach
the connected union
\[
            Z=\{p\}\cup Z_1\cup\cdots\cup Z_k .
\]
Resolve its action on $\omega$ into exact nonempty
excitation sectors, retaining the cover $Z$.
Apply $-T_I^{-1}$ to each resulting coefficient.
This defines a polynomial family map $\mathcal F(c)$
with the lift identity
\[
 \sum_{Z,I}\iota_I T_I\mathcal F(c)_{Z,I}
                     =-QW(c)\omega .
 \tag{V5.9}\label{c18v5:map-lift}
\]
Here $\iota_I$ inserts the constant exterior factors.
The decomposition is fixed by the indicated ordered
commutator expansion; terms with the same $(Z,I)$
may be grouped.

\begin{lemma}[Uniform nonlinear majorant in the same weight]
\label{c18v5:majorant}
Set
\[
 P_8(t)=\sum_{k=0}^{8}
                    \left(1+\frac{k}{4}\right)\frac{t^k}{k!}.
\]
For $\|c\|_{a,E}\le x$,
\[
 \|\mathcal F(c)\|_{a,E}
       \le16be^a P_8(8x/\Delta),\qquad
 \|D\mathcal F(c)\|
       \le\frac{128be^a}{\Delta}P_8'(8x/\Delta).
 \tag{V5.10}\label{c18v5:majorant-bounds}
\]
Both estimates use weight $a$ on input and output.
They do not assume finite-dimensional link spaces.
\end{lemma}

\begin{proof}
Take an ordered $k$-fold commutator of primitive $V_p$
with components $\widehat u_{I_1},\ldots,\widehat u_{I_k}$.
It vanishes unless every $I_j$ meets $E(p)$.
Expand the commutator into at most $2^k$ products.

In a nonzero product, outside $E(p)$ the final
vacuum/nonvacuum status of every link is fixed.
Indeed a right creation produces a nonvacuum factor
there, $V_p$ leaves it unchanged, and a left creation
at the same link would kill the product. Other left
creations produce definite nonvacuum factors.
All remaining sector choices lie on $E(p)$.
There are at most $2^4$ orthogonal output sectors;
Cauchy--Schwarz therefore gives
\[
 \sum_I\|\Pi_I
    (\widehat u_{I_{j_1}}\cdots V_p\cdots
                         \widehat u_{I_{j_k}})\omega\|
             \le4\|V_p\|\prod_{j=1}^k\|u_{I_j}\|.
 \tag{V5.11}\label{c18v5:branch}
\]
The factor four is independent of the size of the
covers or excitation sets. Applying $T_I^{-1}$
followed by the output energy norm exactly cancels
that inverse; no further output sector cost occurs.

For clarity, let the input families have norms
$x_1,\ldots,x_k$. The union weight is at most
$e^a\prod_j e^{a|Z_j|}$. Fix an anchor link $e$
in the union. First place the anchor in $E(p)$.
For every creation slot,
\[
 \sum_{Z,I:I\cap E(p)\ne\varnothing}
            e^{a|Z|}\|c_{Z,I}\|
            \le4x_j/\Delta .
 \tag{V5.12}\label{c18v5:root-slot}
\]
Use \eqref{c18v5:incidence} for the primitive root.
This anchor case costs
$4be^a\,4^k\prod_j(x_j/\Delta)$ before the
branching and commutator factors.

Alternatively place the anchor in one input cover,
say $E(Z_j)$. For a fixed excitation set $I_j$,
\[
 \sum_{p:E(p)\cap I_j\ne\varnothing}\|V_p\|
                          \le4b|I_j|.
\]
The factor $|I_j|$ is cancelled by
\eqref{c18v5:energy-pay} in this particular slot.
Each of the other slots uses
\eqref{c18v5:root-slot}. Thus this case costs
$4be^a\,4^{k-1}\prod_j(x_j/\Delta)$.
There are $k$ choices of this anchor slot.
Overcounting unions with more than one anchor
location is harmless.

Multiplying by $4\,2^k$ from
\eqref{c18v5:branch} and the commutator expansion,
the ordered multilinear coefficient before $1/k!$
is bounded by
\[
 16be^a\left(1+\frac{k}{4}\right)
                          8^k\prod_j(x_j/\Delta).
 \tag{V5.13}\label{c18v5:multilinear}
\]
For $k=0$ the same statement follows by the
factor-four sector bound and the root
incidence. All contributing covers are connected,
because each creation's excitation set meets the
primitive plaquette. Summing $k=0,\ldots,8$ proves
the first bound. Differentiate the bounded
multilinear polynomials and sum over their slots
to obtain the second bound.
\end{proof}

The cost $2^{|E(Z)|/2}$ in V4 is absent here.
It is not being hidden in a renamed norm: only
the four-link primitive is allowed to change
an otherwise definite excitation status.
The energy norm pays the anchor incidence when
the root is attached to a large creation.

\subsection{An actual nonlinear fixed point}

\begin{theorem}[Convergent native creation construction]
\label{c18v5:fixed-point}
If, for some fixed $a\ge0$,
\[
 \theta:=\frac{be^a}{\Delta}\le\frac1{1024},
 \qquad x_*=32be^a,
 \tag{V5.14}\label{c18v5:smallness}
\]
then $\mathcal F$ maps the closed ball
$\|c\|_{a,E}\le x_*$ into itself and is a contraction
there with constant at most
\[
                         \ell_*=\frac7{32}.
 \tag{V5.15}\label{c18v5:ell}
\]
The iterates $c^{(0)}=0$,
$c^{(n+1)}=\mathcal F(c^{(n)})$ converge to a unique
fixed point in that ball, uniformly in the
contraction estimate over all torus sizes.
For $C_*=C(c_*)$,
\[
 (T+V)e^{C_*}\omega
       =\mathcal E_*e^{C_*}\omega,\qquad
 \mathcal E_*=\langle\omega,W(c_*)\omega\rangle .
 \tag{V5.16}\label{c18v5:eigenpair}
\]
The vector and coefficients are gauge invariant.
\end{theorem}

\begin{proof}
Put $t=8x_*/\Delta=256\theta\le1/4$.
The positive coefficients give
\[
 \begin{split}
 P_8(t)&\le e^t(1+t/4)\le17/12,\\
 P_8'(t)&\le e^t(5/4+t/4)\le7/4.
 \end{split}
\]
For the last inequalities use
$e^{1/4}\le(1-1/4)^{-1}=4/3$.
The majorant therefore maps the ball into the
smaller ball of radius $(17/24)x_*$, and its
derivative norm there is at most
$128\theta(7/4)\le7/32$.
The mean-value bound along a segment in the
complex Banach ball proves contraction.
The case $b=0$ gives the zero fixed point directly.
For $b>0$, completeness and the contraction theorem
give the limit and, for example,
\[
 \|c^{(n)}-c_*\|_{a,E}
       \le\frac{\ell_*^n}{1-\ell_*}
                                      \|\mathcal F(0)\|_{a,E}.
 \tag{V5.17}\label{c18v5:iteration}
\]

The fixed-point equation and
\eqref{c18v5:map-lift} imply
$Q([T,C_*]+W(c_*))\omega=0$.
The vacuum expectation of $[T,C_*]$ is zero.
Use \eqref{c18v5:similarity} to obtain
\eqref{c18v5:eigenpair}.
Moreover
$\langle\omega,e^{C_*}\omega\rangle=1$, so the
eigenvector is nonzero. It is in the electric
domain by Lemma~\ref{c18v5:algebra}.

Gauge transformations commute with $T$, the
one-link constant projections, and $V$.
The creation and projection maps are equivariant.
The iterates from zero are therefore invariant,
as are their limit and its vector. No finite
Peter--Weyl cutoff was used.
\end{proof}

For $a>0$, the coefficient tail has the explicit
local bound
\[
 \sup_e\sum_{\substack{Z,I:e\in E(Z)\\|Z|\ge L}}
                   \|c_{Z,I}^*\|
                 \le\Delta^{-1}e^{-aL}x_* .
 \tag{V5.18}\label{c18v5:tail}
\]
This follows directly from the norm and
\eqref{c18v5:energy-pay}. It controls a creation
interaction, not a volume-uniform global overlap.
The normalized vector is
$\Psi_*=e^{C_*}\omega/\|e^{C_*}\omega\|$.
Normalization has not been omitted, and $e^{C_*}$
has not been called a unitary.

\subsection{Why the derivative controls actual excitations}

An eigenvector alone would not select the ground
state or establish a gap. The following argument
uses the linearization of the same nonlinear map.

For fixed finite volume, let
\[
 \mathcal Ld=\sum_{Z,I}\iota_I d_{Z,I}.
\]
Every $v\in Q\operatorname{Dom}T$ has a representation
$\mathcal Ld=v$: use the whole connected plaquette
set as a cover for each exact excitation component.
Define the quotient norm
\[
 |v|_{a,E}
      =\inf_{\mathcal Ld=v}\|d\|_{a,E}.
 \tag{V5.19}\label{c18v5:quotient}
\]
It is a genuine norm, equivalent at each finite
volume to $\|Tv\|$. Indeed
\[
 \|Tv\|\le M\|d\|_{a,E},
\]
whereas the whole-cover representation has norm
at most
$e^{aM}\sqrt{2^M-1}\|Tv\|$.
Thus the kernel of the lift is closed and the
quotient is complete. These equivalence constants
are \emph{not} claimed uniform in volume; the
operator bounds below are uniform without them.
Componentwise inversion gives
\[
                  |T^{-1}v|_{a,E}\le\Delta^{-1}|v|_{a,E}.
 \tag{V5.20}\label{c18v5:quotient-inverse}
\]

Let $\widehat v$ be the canonical creation operator
whose exact-sector coefficients are $\Pi_Iv$ with
their constant exterior removed. It satisfies
$\widehat v\omega=v$ and commutes with every pure
creation operator. The derivative of the lift of
$\mathcal F$ is
\[
 \mathcal L D\mathcal F(c_*)d
       =-T^{-1}Q[W(c_*),\widehat v]\omega,
                  \qquad v=\mathcal Ld .
 \tag{V5.21}\label{c18v5:derivative}
\]
In particular it depends only on $v$, not on the
chosen covers. Denote this descended operator by
$\mathcal J_*$. The contraction derivative bound gives
\[
                         |\mathcal J_*v|_{a,E}
                                      \le\ell_*|v|_{a,E}.
 \tag{V5.22}\label{c18v5:J}
\]

\begin{lemma}[Excited-eigenvalue exclusion]
\label{c18v5:exclusion}
For the eigenpair \eqref{c18v5:eigenpair}, every
linearly independent eigenvector of $T+V$ has
eigenvalue separated from $\mathcal E_*$ by at least
\[
                    (1-\ell_*)\Delta=\frac{25}{32}\Delta.
 \tag{V5.23}\label{c18v5:exclusion-bound}
\]
In particular $\mathcal E_*$ is simple.
\end{lemma}

\begin{proof}
First justify \eqref{c18v5:derivative}. Since creation
operators commute, differentiating
$e^{-C_*-\varepsilon\widehat v}Ve^{C_*+\varepsilon\widehat v}$
at zero gives $[W(c_*),\widehat v]$.
Equation \eqref{c18v5:map-lift} then gives
\eqref{c18v5:derivative}. If $\mathcal Ld=0$, the
sum of its creation maps is zero in each exact
excitation sector, so the derivative lifts to zero.
It therefore descends to the quotient, where taking
the infimum over representations proves
\eqref{c18v5:J}.

Set
\[
 \widetilde H=e^{-C_*}(T+V-\mathcal E_*)e^{C_*}
       =T+[T,C_*]+W(c_*)-\mathcal E_*.
\]
It annihilates $\omega$. For $v\in Q\operatorname{Dom}T$,
the commutation of creation operators implies
\[
 Q\widetilde Hv
    =Tv+Q[W(c_*),\widehat v]\omega .
 \tag{V5.24}\label{c18v5:excitation-equation}
\]
To see this, write $v=\widehat v\omega$ and
commute $\widetilde H$ through $\widehat v$.
The term $[T,C_*]$ is a pure creation operator and
commutes with $\widehat v$, while
$[T,\widehat v]\omega=Tv$. All vectors and operators
are on the electric domain as established above.

Suppose $\phi$ is an eigenvector of $T+V$ with
eigenvalue $\mathcal E_*+\delta$, linearly
independent of $e^{C_*}\omega$.
Then $v=Qe^{-C_*}\phi$ is nonzero.
Its vacuum component contributes nothing to
$Q\widetilde H$; hence \eqref{c18v5:excitation-equation}
gives
\[
                        v=\mathcal J_*v+\delta T^{-1}v.
\]
Equations \eqref{c18v5:quotient-inverse} and
\eqref{c18v5:J} imply
\[
 |v|_{a,E}
      \le\left(\ell_*+\frac{|\delta|}{\Delta}\right)|v|_{a,E}.
\]
This is impossible if $|\delta|<(1-\ell_*)\Delta$.
The quotient norm is nonzero on $v\ne0$, irrespective
of its volume-dependent equivalence constants.
\end{proof}

\begin{theorem}[Actual native lattice ground and gap]
\label{c18v5:actual-gap}
Under \eqref{c18v5:smallness}, the vector $\Psi_*$
is the actual normalized ground state of $H$ and
\[
 E_0=bM+\mathcal E_*,\qquad
     \operatorname{gap}(H)\ge\frac{25}{32}\Delta
                                      =\frac{75}{32}\kappa .
 \tag{V5.25}\label{c18v5:gap}
\]
The ground is simple. The bound is uniform in the
spatial torus size and also holds on the
gauge-invariant subspace.
\end{theorem}

\begin{proof}
Replace $V$ by $sV$, $0\le s\le1$.
The same ball and contraction constant work for
every $s$. The fixed point, creation operator and
eigenvalue $\mathcal E_*(s)$ depend continuously
on $s$, by the uniform contraction and the finite
polynomial formulas.
The self-adjoint operator $T+sV$ has compact
resolvent: $T$ has compact resolvent on the finite
product of compact groups, and $sV$ is bounded.
Its lowest eigenvalue $e_0(s)$ is continuous, since
the variational principle gives
$|e_0(s)-e_0(t)|\le|s-t|\|V\|$ at fixed volume.

By Lemma~\ref{c18v5:exclusion}, the continuous,
nonnegative difference
$\mathcal E_*(s)-e_0(s)$ is either zero or at least
$25\Delta/32$. It is zero at $s=0$.
Continuity on the interval forces it to stay zero.
The exclusion lemma also proves simplicity and
the asserted gap at $s=1$.
Restoring $bM$ changes no gap and gives the
displayed exact energy.

The invariant ground constructed in
Theorem~\ref{c18v5:fixed-point} belongs to the physical
subspace. Gauge restriction is a reducing restriction
of $H$ and cannot introduce an eigenvalue in the
excluded interval above its ground. This proves
the same lower bound for physical excitations.
\end{proof}

This is the actual finite-lattice spectral gap in
the stated electric-dominated regime, not the
comparison gap of $T+B$ in V3 or V4.
No spectral conclusion was inferred from a small
trial-energy error.

\subsection{What is closed and why the programme must go upstream again}

The fixed-weight nonlinear construction now converges.
The original Wilson interaction, rather than a
sequence of larger-support effective terms, is kept
inside a finite-degree creation equation.
The scalar shift, normalization, electric domains
and excited-eigenvalue exclusion have all been
accounted for. Thus V4's fixed-weight-loss
obstruction is bypassed for this explicit small
coupling regime, without declaring its repeated
unitary Newton steps convergent.

For clarity, \eqref{c18v5:smallness} reads
\[
                 \frac b\kappa\le\frac{3e^{-a}}{1024}.
 \tag{V5.26}\label{c18v5:ratio}
\]
The parameter $a$ here is a dimensionless size weight,
not a lattice spacing. Strong-coupling perturbative
lattice gaps were already available in the literature
and in the archives; f14 V52, for example, uses
a different Euclidean transfer convention.
The new point within this current branch is the
complete nonlinear construction and actual spectral
argument with the present native Hamiltonian and
its untruncated domains. It is not a new solution
of the continuum problem.

In particular, the estimate has not been proved
along a physical weak-coupling refinement trajectory.
There is no construction here of the nontrivial
four-dimensional continuum measure or its
reconstructed physical Hamiltonian. Nor does the
bounded similarity at each finite volume by itself
construct an infinite-volume similarity or a
normalized continuum state.

The productive next direction is now the physical
scale bridge: identify an exact block/reference
description on a genuine refinement trajectory
whose dimensionless interaction is in a proved
stability class, with the true neutral, boundary
and source terms retained. More coefficients or
a better strong-electric constant alone do not
supply that bridge. Any proposed flow must derive
its physical time normalization and continuum
noncollapse, not install them by rescaling a
comparison Hamiltonian.

The full interacting four-dimensional Yang--Mills
construction and positive physical mass gap remain
unproved. The small-coupling lattice construction is
a completed benchmark inside that larger goal, not
a replacement for it.

\subsection{Verification and source preservation}

The checker verifies the native four-link/four-incidence
geometry on three cubic tori and the rational
majorant constants. A five-qubit algebra fixture
has nonzero eighth and zero ninth primitive
commutator, testing the actual four-link degree
rather than importing the two-link degree.
Three three-qutrit models check creation commutation,
nonlinear iteration, random-ball Lipschitz estimates,
the exact similarity identity, the descended
excitation-derivative identity, and agreement with
their actual ground eigenvectors and gaps.
They are not native Wilson spectral computations.
Their finite errors and vanishing floating residuals
do not certify the unbounded or volume-uniform
theorems; those rest on the written proofs.
V4's mathematical diagnostics are also replayed.

V5 preserves the complete V4 body apart from its
version title and replaces only the active V4 source.
All forty-seven other note files are unchanged.
Only \texttt{.tex} files remain in
\texttt{latex\_notes}; companion records, checks and
rendering products stay outside it.
The next companion checkpoint and broad literature
sweep are V10.

% END CYCLE 18 VERSION 5 APPEND

% BEGIN CYCLE 18 VERSION 6 APPEND
\clearpage
\section{V6: The boundary-charge scale in the actual weak-electric vacuum}
\label{c18v6:main}

V5 closed an electric-dominated lattice argument.
The next task is not another coefficient in that expansion,
but a reference which can retain the actual vacuum as
$b/\kappa$ grows. Bare time calibration was already treated
in f16.2 V55. Closed and open two-plaquette reference spectra
were treated in f15 V92--V95. Their boundary modes cannot be
replaced by separately imposed neutral constraints.

Here a sharp compact-group estimate determines how many
partial boundary-charge labels the \emph{fully coupled native
vacuum} must retain. It is independent of the volume and of
the unknown interior wavefunction. Combining it with the
existing native local-energy bound also gives a sufficient
cutoff for any fixed bank of boundary vertices. This changes
the proposed block construction: its retained carrier must
grow on a fourth-root scale. It does not prove a gap for the
operator on that carrier.

The partial-gauge actions and their zero and second Wilson
moments already occur in f17 V87; they are not new.
The local-energy estimate used below is f17 V89,
(V89.19)--(V89.20). The new steps are the sharp inverse-deficit
compression, its application to the \emph{restored} vacuum,
and the two-sided boundary-band scale. This differs from
f15 V95's growing charged ground band of an isolated regional
transfer and from f17 V99's one-link high-energy resolvent.

\subsection{The partial charge and the inherited native energy input}

Keep \eqref{c18v1:H}, with $b>0$, and let
$\Omega$ be its actual normalized positive ground.
Set
\[
             \beta_*=\frac b\kappa,\qquad
             A_*=3\sqrt2 .
 \tag{V6.1}\label{c18v6:parameters}
\]
The symbol $\beta_*$ is a ratio of Hamiltonian coefficients,
not the Euclidean Wilson inverse coupling.

Partition the links into $D,D^c$. Choose a crossing
plaquette $p$ and one of its transition vertices $x$:
exactly one of the two links of $p$ incident at $x$ is in
$R$, where $R$ is either $D$ or $D^c$.
Let $\mathsf U_{x,R}(h)$ apply a gauge transformation
$h\in\mathrm{SU}(2)$ only to the links of $R$ at $x$.
Write $d=d_R(x)$ for their number, so $1\le d\le5$.
Its Casimir and spectral cutoff are
\[
 C_x^R=-\sum_{a=1}^3(G_{x,a}^R)^2,\qquad
 P_J^{x,R}=\mathbf1_{[0,J(J+2)]}(C_x^R),
                 \qquad J=0,1,\ldots .
 \tag{V6.2}\label{c18v6:cutoff}
\]
Thus $J$ is twice the largest retained spin.
The Casimir normalization is the same as the one-link
round Laplacian. $P_0^{x,R}$ imposes one additional partial
Gauss constraint, not the original full Gauss constraint.

These projections commute with the original gauge group:
a full gauge transformation at $x$ conjugates $h$,
preserving its Casimir; transformations at other endpoints
commute through the opposite left/right action.
Consequently the cutoffs preserve the physical subspace.
They need not commute with the \emph{restored} Hamiltonian.

We use the following already proved bounds:
\[
 \langle\Omega,(1-w_p)\Omega\rangle
       \le A_*\beta_*^{-1/2},\qquad
 \langle\Omega,L_e\Omega\rangle
       \le A_*\beta_*^{1/2}.
 \tag{V6.3}\label{c18v6:inherited}
\]
For provenance, f17 V89 obtains
$E_0/M\le\min\{b,3\sqrt{2\kappa b}\}$ by a full-Haar
product trial. Positivity and cubic symmetry then bound
each kinetic and magnetic term. The trial is not identified
with $\Omega$. We use that result, not a new variational
assertion or a new physical calibration.

\subsection{A sharp inverse-deficit compression on the round group}

Identify $\mathrm{SU}(2)$ with the unit three-sphere
$q=(q_0,\boldsymbol q)$ with normalized Haar measure.
Let $\Pi_J$ retain all spherical harmonics of degrees
$0,\ldots,J$, equivalently all Peter--Weyl labels in that
range, including their full multiplicities. Put
\[
 d_0(q)=1-q_0,\qquad
 B_J=\frac1{1-\cos\!\bigl(\pi/(2J+3)\bigr)}.
 \tag{V6.4}\label{c18v6:sharp-constant}
\]
The subscript in $d_0$ does not denote a lattice link.

\begin{theorem}[Exact group-band constant]
\label{c18v6:group}
The inverse multiplier is interpreted as a quadratic form.
Its compression to the finite band is bounded and has norm
\[
       \bigl\|\Pi_Jd_0^{-1}\Pi_J\bigr\|=B_J.
 \tag{V6.5}\label{c18v6:compression}
\]
In particular, for every $f\in L^2(\mathrm{SU}(2))$,
\[
          \|\Pi_Jf\|^2\le B_J\int d_0(q)|f(q)|^2\,dq.
 \tag{V6.6}\label{c18v6:group-inequality}
\]
The constant is optimal. The same inequality holds
for Hilbert-space-valued functions.
\end{theorem}

\begin{proof}
Write $q_0=\cos\theta$ and
$\boldsymbol q=\sin\theta\,\boldsymbol n$.
Normalized Haar measure is
$(2/\pi)\sin^2\theta\,d\theta\,d\boldsymbol n$,
where the measure on the unit two-sphere has total mass one.
The angular momentum $l=0,1,\ldots$ of $\boldsymbol n$
decomposes the space into orthogonal radial chains.
A normalized chain basis has the form
\[
 \phi_{n,l,m}(\theta,\boldsymbol n)
  =c_{n,l}(\sin\theta)^l
        C_{n-l}^{\,l+1}(\cos\theta)Y_{l,m}(\boldsymbol n),
                  \quad n=l,l+1,\ldots .
 \tag{V6.7}\label{c18v6:radial-basis}
\]
Here $-l\le m\le l$. These are degree-$n$ spherical
harmonics: substitution in the polar Laplacian gives
eigenvalue $n(n+2)$, and
$\sum_{l=0}^n(2l+1)=(n+1)^2$ accounts for its full
multiplicity. Orthogonality and the Gegenbauer recurrence
give the multiplication matrix of $q_0$:
\[
 q_0\phi_{n,l,m}
    =a_{n,l}\phi_{n+1,l,m}
                  +a_{n-1,l}\phi_{n-1,l,m},\qquad
 a_{n,l}=\frac12
 \sqrt{\frac{(n-l+1)(n+l+2)}{(n+1)(n+2)}}\le\frac12 .
 \tag{V6.8}\label{c18v6:Jacobi}
\]
The lower term is zero at $n=l$.
For example the coefficient follows by inserting
\[
 2(k+\lambda)x C_k^\lambda
   =(k+1)C_{k+1}^\lambda+(k+2\lambda-1)C_{k-1}^\lambda
\]
and taking the ratios of their orthogonality norms, with
$k=n-l,\lambda=l+1$.
These standard polynomial identities can be checked in
\href{https://dlmf.nist.gov/18.9}{DLMF Table 18.9.1}
and \href{https://dlmf.nist.gov/18.3}{Table 18.3.1};
the operator comparison that follows is proved here.

Call the half-line matrix in \eqref{c18v6:Jacobi} $X_l$,
indexing its first row by zero. All its entries are
nonnegative and bounded by those of $X_0$, whose two
off-diagonals equal $1/2$.
For $0<r<1$,
\[
 (I-rX_l)^{-1}=\sum_{k\ge0}r^kX_l^k .
\]
Every path in a matrix entry has nonnegative weight.
It follows entrywise that
$(I-rX_l)^{-1}\le (I-rX_0)^{-1}$.
This is an entrywise comparison, not an unproved
Loewner comparison of the Jacobi matrices.
After compressing the first $s$ rows, it implies an
operator-norm comparison: replace a test vector by its
componentwise absolute values in the quadratic form.

The limit $r\uparrow1$ exists in every such finite
compression. In the integral representation its only
singularity is $d_0^{-1}$ at $\theta=0$.
A finite collection of the smooth harmonics
\eqref{c18v6:radial-basis} is bounded there, and
$d_0^{-1}\sin^2\theta$ is integrable.
For $r\ge1/2$ the regularized multiplier is dominated
by a constant times $d_0^{-1}$.
Dominated convergence therefore justifies this limit.

In the $l=0$ chain $\phi_{n,0,0}$ is the character
$\chi_n=\sin((n+1)\theta)/\sin\theta$.
For indices $i,j=1,\ldots,s$ the limiting matrix is
\[
 \begin{split}
 K^{(s)}_{ij}
 &=\frac2\pi\int_0^\pi
       \frac{\sin(i\theta)\sin(j\theta)}
                         {1-\cos\theta}\,d\theta\\
 &=2\min(i,j).
 \end{split}
 \tag{V6.9}\label{c18v6:min-matrix}
\]
Indeed use the product-to-sum identity and
$\int_0^\pi(1-\cos(k\theta))/(1-\cos\theta)\,d\theta
=\pi k$ for integer $k\ge0$.
The latter follows by expanding the squared finite
geometric sum, whose nonconstant cosine terms integrate
to zero.

The inverse of $K^{(s)}$ is one half the tridiagonal
matrix with off-diagonal entries $-1$, diagonal entries
$2$ except for the last entry $1$.
This is verified by multiplication using successive
differences of $\min(i,j)$.
Its eigenvectors are $\sin(i t_k)$, with
\[
 t_k=\frac{(2k-1)\pi}{2s+1},\quad 1\le k\le s.
\]
They solve the interior second-difference equation and
the boundary conditions $v_0=0$, $v_{s+1}=v_s$.
Thus
\[
       \|K^{(s)}\|=\frac1{1-\cos(\pi/(2s+1))}.
\]
In angular sector $l$ the retained chain has
$s=J-l+1$ rows. Its norm is at most that of $K^{(s)}$,
hence at most $B_J$. The $l=0$ chain attains $B_J$.
This proves \eqref{c18v6:compression} on the full band,
not just on class functions.

For any band vector $u$, weighted Cauchy--Schwarz gives
\[
 |\langle u,f\rangle|^2
 \le\left(\int\frac{|u|^2}{d_0}\right)
                         \left(\int d_0|f|^2\right)
 \le B_J\|u\|^2\int d_0|f|^2 .
\]
Take the supremum over unit band vectors.
For sharpness take a top eigenvector $u$ of the
compression and $f_\varepsilon=(d_0+\varepsilon)^{-1}u$.
Each is in $L^2$; as $\varepsilon\downarrow0$,
$\Pi_Jf_\varepsilon\to B_Ju$ and
$\int d_0|f_\varepsilon|^2\to B_J\|u\|^2$.
The ratio in \eqref{c18v6:group-inequality} tends to $B_J$.
No $L^2$ extremizer for that last inequality is asserted.
Componentwise application, or orthogonal expansion in
the value Hilbert space, proves the vector-valued version.
\end{proof}

In particular
\[
 B_0=2,\qquad
 B_J\sim\frac{2(2J+3)^2}{\pi^2},\qquad
 B_J\le\frac{(2J+3)^2}{2}.
 \tag{V6.10}\label{c18v6:B-growth}
\]
The upper bound uses $1-\cos t\ge2t^2/\pi^2$
on $[0,\pi]$. This calculation is an inverse-deficit
estimate, not a Poincare inequality for an unknown
interacting measure.

\subsection{Transfer to an actual cut plaquette without changing its law}

\begin{theorem}[Boundary-charge/curvature inequality]
\label{c18v6:native}
For the transition vertex and side in
\eqref{c18v6:cutoff}, every vector $\psi$ in the full
native Haar Hilbert space satisfies
\[
       \|P_J^{x,R}\psi\|^2
             \le B_J\langle\psi,(1-w_p)\psi\rangle .
 \tag{V6.11}\label{c18v6:native-inequality}
\]
The same assertion holds on its physical subspace.
It uses no cut ground state, conditional Gibbs law,
or decoupling approximation.
\end{theorem}

\begin{proof}
Choose the plaquette link in $R$ incident at $x$.
It supplies a free coordinate for the partial gauge orbit:
use the partial transformation to set that link to the
identity, retaining all other transformed links as
coordinates $\eta$. Reversing an incoming orientation
replaces the orbit coordinate by its inverse.
Successive left/right changes of the remaining Haar
variables preserve their measure, so the product measure
becomes $dh\,d\eta$.

Only one of the two plaquette links at $x$ transforms.
For fixed $\eta$, the plaquette holonomy is therefore
a left/right translate of $h$ or $h^{-1}$.
Changing from $h$ to that holonomy preserves Haar measure
and every Peter--Weyl degree. In these coordinates
$1-w_p$ is exactly $1-q_0$, and the partial Casimir cutoff
is $\Pi_J$ on the orbit variable.
Apply \eqref{c18v6:group-inequality} on each fibre
and integrate in $\eta$.
This is merely a unitary change of coordinates in the
original full Haar space; $\psi$ remains the original
vector, with all its correlations. The physical
restriction follows from the gauge covariance already
established for $P_J^{x,R}$.
\end{proof}

For $J=0$, this strengthens the inherited neutral
moment identities in a different direction:
\[
           \|P_0^{x,R}\psi\|^2
                  \le2\langle\psi,(1-w_p)\psi\rangle.
 \tag{V6.12}\label{c18v6:singlet}
\]
It bounds how much of a low-curvature vector can lie in
the artificially neutral sector. It does not assert
that the physical vacuum is charged under the
\emph{full} gauge group.

\subsection{Necessary and sufficient boundary bandwidth in the native vacuum}

\begin{theorem}[Two-sided actual-vacuum cutoff bounds]
\label{c18v6:vacuum}
Uniformly in the cubic torus size,
\[
 \begin{split}
 \|P_J^{x,R}\Omega\|^2
       &\le \min\{1,A_*B_J\beta_*^{-1/2}\},\\
 \|(I-P_J^{x,R})\Omega\|^2
       &\le \min\left\{1,
          \frac{A_*d^2\beta_*^{1/2}}{(J+1)(J+3)}\right\}.
 \end{split}
 \tag{V6.13}\label{c18v6:two-sided}
\]
In particular, every fixed partial-charge band has
vanishing vacuum weight as $\beta_*\to\infty$.
For any fixed target retained probability $0<\delta<1$,
a cutoff with $\|P_J^{x,R}\Omega\|^2\ge\delta$ must obey
\[
      2J+3\ge
        \sqrt{\frac{2\delta}{A_*}}\,\beta_*^{1/4}.
 \tag{V6.14}\label{c18v6:necessary}
\]
Conversely, retained probability at least $1-\varepsilon$
is guaranteed whenever
\[
       (J+1)(J+3)\ge
                  \frac{A_*d^2}{\varepsilon}\beta_*^{1/2},
                  \qquad 0<\varepsilon<1.
 \tag{V6.15}\label{c18v6:sufficient}
\]
\end{theorem}

\begin{proof}
The first bound is
\eqref{c18v6:native-inequality} applied to $\Omega$,
using \eqref{c18v6:inherited}.
For the second, Cauchy--Schwarz on the $d$ partial
star derivatives gives the quadratic-form inequality
\[
             C_x^R\preceq d\sum_{\substack{e\in R\\e\ni x}}L_e.
\]
All right and left Casimirs have the same round
normalization, and their signs do not affect this bound.
Thus \eqref{c18v6:inherited} implies
$\langle C_x^R\rangle_\Omega\le A_*d^2\beta_*^{1/2}$.
The first omitted eigenvalue is $(J+1)(J+3)$.
The spectral theorem proves the tail estimate.
Use \eqref{c18v6:B-growth} for the necessary condition;
the sufficient one is a direct substitution in the tail.
\end{proof}

Consequently the smallest cutoff achieving any fixed
retention $\delta\in(0,1)$ is of order
$\beta_*^{1/4}$, with constants depending on $\delta$
and the bounded star degree, but not on $N$.
This is a bandwidth assertion, not an asymptotic
formula for the entire charge distribution.

There is also a useful finite-bank approximation.
For any finite collection $\mathcal A$ of partial actions
from this one link partition, the Casimirs commute:
different vertices act by commuting left/right actions,
and the two sides at one vertex act on disjoint links.
Define
\[
       P_{\mathcal A,J}=\prod_{(x,R)\in\mathcal A}P_J^{x,R}.
\]
Their commuting spectral projections satisfy
\[
 \begin{split}
 \|(I-P_{\mathcal A,J})\Omega\|^2
 &\le \frac{A_*\beta_*^{1/2}}{(J+1)(J+3)}
                       \sum_{(x,R)\in\mathcal A}d_R(x)^2 .
 \end{split}
 \tag{V6.16}\label{c18v6:bank}
\]
Indeed $I-\prod P_i\preceq\sum(I-P_i)$.
If the right side is $\varepsilon<1$, the normalized
projection is defined and obeys
\[
 \left\|\Omega-
   \frac{P_{\mathcal A,J}\Omega}{\|P_{\mathcal A,J}\Omega\|}\right\|
 \le\sqrt{2\varepsilon}.
 \tag{V6.17}\label{c18v6:state-error}
\]
Choose its phase by the positive overlap and use
$2(1-\sqrt{1-\varepsilon})\le2\varepsilon$.
No independence of charges is needed.
For a growing boundary bank the displayed sum is
a real cost and cannot be suppressed.
The cutoff leaves all interior multiplicities intact;
it is not a finite-dimensional truncation of the
entire many-link Hilbert space.

\subsection{Consequence for the next block reference}

Let $P_{\rm sep}$ impose separate Gauss invariance
on both sides at every cut vertex. For a cut with a
crossing plaquette it is a subprojection of
$P_0^{x,R}$ for a transition vertex. Hence
\[
 \|P_{\rm sep}\Omega\|^2
        \le\min\{1,2A_*\beta_*^{-1/2}\},\qquad
 |\langle\Phi,\Omega\rangle|^2
        \le2A_*\beta_*^{-1/2}
 \tag{V6.18}\label{c18v6:neutral-reference}
\]
for every normalized $\Phi\in P_{\rm sep}\mathcal H$.
This includes products of separately neutral block
ground states. It also includes entangled interior
states remaining in that separate-neutral carrier.
The result concerns at least one actual cut plaquette,
so increasing the volume of a closed neutral block
does not remove this particular obstruction.

A mere global overlap obstruction would not exclude
a useful local dressing: V5 already warns against that
inference. Here there is also a bounded diagnostic at
one boundary star. Any separately invariant state has
$\langle P_0^{x,R}\rangle=1$, whereas the actual vacuum
has expectation at most $2A_*\beta_*^{-1/2}$.
Thus no approximation locally close on that star
can keep this extra partial Gauss constraint.
This does not forbid a non-small dressing that restores
the boundary carriers.

In the already recorded Hamiltonian convention
$f16.2$ V55,
\[
 \kappa=\frac{g_0^2}{8a_s},\qquad
 b=\frac4{g_0^2a_s},\qquad
          \beta_*=\frac{32}{g_0^4}.
 \tag{V6.19}\label{c18v6:KS-ratio}
\]
This is a substitution into an inherited bare
convention, not a new derivation of renormalized time.
It puts the necessary and sufficient fixed-bank
charge bandwidth at order $g_0^{-1}$.
The old f15 V95 regional-transfer scale
$n=O(\sqrt\gamma)$ with $\gamma=4/g_0^2$ is consistent
with this, but did not prove the restored vacuum's
weight bounds \eqref{c18v6:two-sided}.
Neither observation derives a physical mass.

The next block construction should therefore retain
these matched boundary charges and their complete
interior multiplicities before attempting a
small-remainder argument. A concrete next target is the
actual Hamiltonian compressed to the growing charge
bank, together with the coupling to its complement.
The vacuum tail \eqref{c18v6:bank} is not a bound on
that coupling, on all excited states, or on its Schur
return. The partial Casimirs are not conserved by
restoring the crossing plaquettes, and the neutral
returns identified in f17 V87 still have to be paid.

The full interacting four-dimensional continuum
construction and positive physical mass gap remain
unproved. This append supplies a native necessary
carrier and a quantitative vacuum-retention estimate
in the weak-electric regime; it does not certify
stability of its interacting effective dynamics.

\subsection{Checks and preservation}

The diagnostic checks the Gegenbauer recurrence and
normalization ratios, the exact radial inverse-deficit
matrix and its tridiagonal inverse, its finite-band
eigenvalues, and the domination of the other angular
sectors by direct quadrature. Smooth compact-group
wavefunctions check the dual inequality and the
small-deficit loss of a fixed band. They are
one-group algebra diagnostics, not computed many-link
Wilson ground states. The latter conclusions follow
from the written theorem and the inherited native
energy bound, not from those fixtures.
V5's mathematical diagnostics are replayed.

V6 appends to the complete V5 body, changing only its
version title and replacing the active filename.
All forty-seven other note sources are preserved.
Only \texttt{.tex} files are stored in
\texttt{latex\_notes}; diagnostic and rendering products
are outside that folder. The last companion checkpoint
was V5. The next companion review and broad literature
sweep are V10.

% END CYCLE 18 VERSION 6 APPEND

% BEGIN CYCLE 18 VERSION 7 APPEND
\clearpage
\section{V7: A cubic-scale complementary block with its boundary source retained}
\label{c18v7:main}

\subsection{What is new, and what is inherited}

V6 determined the charge bandwidth needed to retain the native
vacuum. The present question is different: can the complementary
charge sector be made coercive as an operator, and what coupling
must still be retained? We obtain a volume-independent affirmative
answer to the first question. The exact crossing source remains;
our bounds do not make its return perturbatively small at the
improved cutoff.

Two ingredients are inherited, not reproved as new discoveries.
The tempered one-link trial and its differentiated normalization
are from f17 V89. The one-link electric deletion bound, the
one-shell selection rule and the Feshbach identities are from
f17 V99--V100 and the earlier return analysis. The change here
is to compare with the \emph{classical minimum} exterior operator
before subtracting the vacuum energy. This converts the existing
tempered trial into an all-vector local electric inequality,
then applies it to V6's partial boundary charges. The resulting
cubic-scale complementary bound is not the earlier
square-root-scale deletion estimate.

All statements below concern the actual finite cubic lattice,
$N\ge3$, with the conventions \eqref{c18v1:H}. A scalar trial
comparison does not replace its Wilson interaction or its Haar
measure. The full four-dimensional continuum mass-gap problem
remains open in these notes.

\subsection{Classical-minimum comparison with an explicit local debit}

Separate a link quaternion $q\in S^3$ from all exterior variables
$y$. For nonnegative, possibly inhomogeneous incident couplings
$b_p$, write
\[
 \begin{gathered}
 H=H_{{\rm out},e}+\kappa L_e+B_e-f_e(y)\cdot q,\\
 B_e=\sum_{p\ni e}b_p,\qquad S_{2,e}=\sum_{p\ni e}b_p^2,
 \qquad h_e(y)=|f_e(y)|\le B_e .
 \end{gathered}
 \tag{V7.1}\label{c18v7:split}
\]
Here $H_{{\rm out},e}$ contains the other electric terms and
all nonincident magnetic terms. Each staple contributes a
unit quaternion, with any orientation conjugations included
in $f_e$. Define the exterior operator and its ground energy by
\[
 H_{{\rm cl},e}=H_{{\rm out},e}+B_e-h_e,\qquad
 E_{{\rm cl},e}=\inf\sigma(H_{{\rm cl},e}).
 \tag{V7.2}\label{c18v7:classical}
\]
The subscript \emph{cl} refers only to minimization of the
one-link potential. This is not a classical replacement of
the many-link theory. The potential $h_e$ is bounded and
Lipschitz, so the exterior quadratic form is closed on the
usual $H^1$ space. We never differentiate $h_e$ at a zero of
$f_e$.

\begin{theorem}[All-vector local electric coercivity]
\label{c18v7:comparison}
For every $\tau>0$, set
\[
 D_e(\tau)=\frac{3\kappa B_e}{4\tau}
            +\frac{3\tau}{2}
            +\frac{9\kappa S_{2,e}}{4\tau^2}.
 \tag{V7.3}\label{c18v7:cost}
\]
The actual ground energy satisfies
\[
 0\le E_0-E_{{\rm cl},e}\le
       \min\{B_e,D_e(\tau)\}.
 \tag{V7.4}\label{c18v7:ground-comparison}
\]
Consequently, as closed quadratic forms on the full native
Hilbert space, and hence on its physical subspace,
\[
 K=H-E_0\ \ge\ \kappa L_e-\min\{B_e,D_e(\tau)\}.
 \tag{V7.5}\label{c18v7:local-coercivity}
\]
\end{theorem}

\begin{proof}
The exact decomposition
\[
 H=I\otimes H_{{\rm cl},e}+\kappa L_e+h_e-f_e\cdot q
       \ge I\otimes H_{{\rm cl},e}+\kappa L_e
\]
proves the first inequality in \eqref{c18v7:ground-comparison}.
A constant link function tensored with a normalized exterior
ground state has extra energy $\langle h_e\rangle\le B_e$,
which gives its first upper bound.

For the other upper bound use the normalized smooth trial
\[
 \chi_{\tau,f}(q)=
 \frac{\exp(f\cdot q/(2\tau))}
      {\bigl(\int_{S^3}\exp(f\cdot v/\tau)\,dv\bigr)^{1/2}},
 \qquad (J_\tau\alpha)(q,y)=\chi_{\tau,f_e(y)}(q)\alpha(y).
 \tag{V7.6}\label{c18v7:trial}
\]
Normalized Haar measure is used in both integrals. This
function is smooth in $f$, including $f=0$, and real, so its
normalization gives
$\langle\chi_{\tau,f},\partial\chi_{\tau,f}\rangle=0$.

For completeness the scalar estimates needed from f17 V89
follow directly by integration by parts. Put $a=h/\tau$ and
let $\mu(a)$ be the mean of $t=q_0$ under density proportional
to $e^{at}$. Symmetry gives $0\le\mu(a)\le1$. Integration of
the spherical Laplacian of $t$, for which
$-\Delta t=3t$ and $|\nabla t|^2=1-t^2$, gives
$a\,\mathbb E_a(1-t^2)=3\mu(a)$. The corresponding tilted
one-dimensional measure on $r=1-t\in[0,2]$ has density
proportional to $e^{-ar}\sqrt{r(2-r)}$. Integrating the
derivative of $r e^{-ar}\sqrt{r(2-r)}$ gives
\[
 a\,\mathbb E_a r
 =\frac32-\frac12\,\mathbb E_a\frac{r}{2-r}\le\frac32
 \quad(a>0).
\]
Boundary terms vanish; the last integrand is integrable.
Therefore
\[
 \langle\chi_{\tau,f},\kappa L\chi_{\tau,f}\rangle
       =\frac{3\kappa a\mu(a)}4\le\frac{3\kappa h}{4\tau},
 \qquad
 \langle h-f\cdot q\rangle_{\chi_{\tau,f}^2}\le\frac{3\tau}{2}.
\]
At $a=0$ the assertions follow directly by continuity.

For a direction $v$ in field space,
\[
 \partial_v\chi_{\tau,f}
 =\frac{v\cdot(q-\mathbb E q)}{2\tau}\chi_{\tau,f},
 \qquad
 \|\partial_v\chi_{\tau,f}\|^2\le\frac{|v|^2}{4\tau^2}.
\]
There are twelve distinct exterior staple links for $N\ge3$.
An exterior link occurs in only one of the four staples.
Each of its three orthonormal invariant derivatives has
quaternion norm $b_p$. Thus, pointwise in $y$,
\[
 \sum_{j\ne e,a}|\partial_{j,a} f_e|^2=9S_{2,e},
 \qquad
 W_\tau(y):=\kappa\sum_{j\ne e,a}
       \|\partial_{j,a}\chi_{\tau,f_e(y)}\|^2
       \le\frac{9\kappa S_{2,e}}{4\tau^2}.
 \tag{V7.7}\label{c18v7:derivative-cost}
\]
The zero normalization derivative removes the cross term
with $\partial\alpha$. Hence, on the exterior form domain,
\[
 J_\tau^*HJ_\tau\le H_{{\rm cl},e}+D_e(\tau)
\]
in quadratic-form sense. Testing on the exterior ground
proves the second upper bound in
\eqref{c18v7:ground-comparison}. No differentiability of
that ground beyond $H^1$ is required.
Finally subtract $E_0$ from the exact lower decomposition
and use $H_{{\rm cl},e}\ge E_{{\rm cl},e}$.
\end{proof}

For homogeneous couplings, $B_e=4b$ and $S_{2,e}=4b^2$.
When $b>0$, choose $\tau_0=(12\kappa b^2)^{1/3}$ and define
\[
 \Lambda(\kappa,b)=
 \min\left\{4b,\,
       \frac94(12\kappa b^2)^{1/3}
          +\frac{3\kappa b}{(12\kappa b^2)^{1/3}}\right\}.
 \tag{V7.8}\label{c18v7:Lambda}
\]
Set $\Lambda(\kappa,0)=0$ separately.
The choice $\tau_0$ minimizes the last two terms of
$D_e(\tau)$, not necessarily its whole expression.
The theorem gives, for every link and every vector in
the common form domain,
\[
                K\ge\kappa L_e-\Lambda(\kappa,b).
 \tag{V7.9}\label{c18v7:every-link}
\]
There is no lattice-size factor. With $\beta_*=b/\kappa$,
\[
 \frac{\Lambda}{\kappa}=F(\beta_*):=
 \min\left\{4\beta_*,
       \frac94\,12^{1/3}\beta_*^{2/3}
            +\frac3{12^{1/3}}\beta_*^{1/3}\right\}.
 \tag{V7.10}\label{c18v7:F}
\]
At weak electric coupling this improves the old debit
$4b$ to order $\kappa^{1/3}b^{2/3}$. It is an all-vector
form estimate. The sharper vacuum-only electric moment
from V6 is not superseded.

\subsection{Partial boundary charge and the whole low-energy window}

Fix a nonempty proper set $R$ of the six half-edges
incident to a vertex $x$, as induced by a link partition.
Write $d=|R|$, $C=C_x^R$ for its partial gauge Casimir,
$P=P_J^{x,R}$, $Q=I-P$, and
$\mu_J=(J+1)(J+3)$. The nonnegative integer $J$ uses
V6's twice-spin normalization. These projections commute
with the full Gauss action and with every individual
electric link Laplacian.

\begin{theorem}[Complementary coercivity and window retention]
\label{c18v7:partial}
In the above notation,
\[
 K\ge\frac{\kappa}{d^2}C-\Lambda.
 \tag{V7.11}\label{c18v7:partial-form}
\]
In particular the actual complementary block satisfies
\[
 \mathcal C_J:=QKQ\ge\delta_J Q,\qquad
 \delta_J=\frac{\kappa\mu_J}{d^2}-\Lambda.
 \tag{V7.12}\label{c18v7:high}
\]
For every $E\ge0$, the entire spectral subspace of $K$
in $[0,E]$ obeys
\[
 \|Q\,1_{[0,E]}(K)\|^2
 \le\min\left\{1,\frac{d^2(E+\Lambda)}{\kappa\mu_J}\right\}.
 \tag{V7.13}\label{c18v7:window}
\]
On the physical Hilbert space all three assertions hold
with $d$ replaced by $d_{\rm eff}=\min(d,6-d)$.
\end{theorem}

\begin{proof}
Sum \eqref{c18v7:every-link} over the $d$ selected links
and divide by $d$. Cauchy--Schwarz for the three components
of their infinitesimal gauge generators gives
$C\le d\sum_{e\in R}L_e$ as forms. Combining these two
inequalities gives \eqref{c18v7:partial-form}. This is a
sum of form inequalities, not a maximum or a square of
noncommuting operator inequalities. On $\operatorname{Ran}Q$,
the next partial charge has Casimir eigenvalue $\mu_J$,
which proves \eqref{c18v7:high}. Since $P,Q$ commute with
the electric operator and preserve its domain, adding
the bounded magnetic potential defines the compressed
self-adjoint operators on their reduced electric domains.

For a unit vector $\psi$ in the indicated spectral window,
$\langle\psi,K\psi\rangle\le E$; the form inequality gives
$\langle C\rangle_\psi\le d^2(E+\Lambda)/\kappa$.
Since $C\ge\mu_J Q$, taking the supremum over all such
vectors proves \eqref{c18v7:window}, including arbitrary
superpositions in that window.

On physical vectors the full Gauss generator vanishes:
$G_R+G_{R^c}=0$. The two partial actions commute.
Their Casimirs agree on the physical subspace, so their
spectral projections there agree as well. Apply the
same proof to the smaller of $R$ and $R^c$.
\end{proof}

For example a prescribed lower bound $\delta_J\ge\gamma>0$
is ensured by
\[
 (J+1)(J+3)\ge
 d^2\bigl(F(\beta_*)+\gamma/\kappa\bigr).
 \tag{V7.14}\label{c18v7:choice}
\]
At fixed $\gamma/\kappa$ this is a sufficient cutoff of
order $\beta_*^{1/3}$, instead of $\beta_*^{1/2}$ in the
old deletion bound. It is not a necessary lower bound
on the true complementary spectrum. At fixed $E/\kappa$,
the same order, with a suitable prefactor, ensures any
fixed error in \eqref{c18v7:window}. For the vacuum alone,
V6 already gives the better order $\beta_*^{1/4}$.

\subsection{The exact crossing source remains a top-shell source}

Let $\mathcal X_{x,R}$ be the plaquettes whose two half-edges
at $x$ have opposite membership in $R$. Put
$m=|\mathcal X_{x,R}|$ and
$U_{x,R}=-b\sum_{p\in\mathcal X_{x,R}}w_p$.
Write $\Pi_n$ for the projection onto partial charge $n$.
Plaquettes outside this set commute with the partial
action: either neither of their two links is transformed,
or the two transformations cancel in the trace.

\begin{lemma}[Source shell and its sharp individual coefficient]
\label{c18v7:shell}
The off-diagonal block of the restored Hamiltonian is
\[
 \mathcal B_J:=QKP
       =\Pi_{J+1}U_{x,R}\Pi_J,\qquad
 \|\mathcal B_J\|\le\frac{mb}{2},\qquad m\le8.
 \tag{V7.15}\label{c18v7:source}
\]
For each individual crossing plaquette,
\[
             \|\Pi_{J+1}w_p\Pi_J\|=\frac12,
 \tag{V7.16}\label{c18v7:sharp-source}
\]
on both the full and the physical Hilbert spaces.
No equality is asserted for the norm of the sum.
\end{lemma}

\begin{proof}
The electric terms and the scalar centering commute
with $P$. A crossing trace transforms in the fundamental
partial representation, so it changes the charge label
only by $+1$ or $-1$. Only the transition $J\to J+1$
survives $Q(\cdot)P$.

The free partial-orbit coordinate used in V6 identifies
a crossing $w_p$ with $q_0$ on $S^3$, with the remaining
variables as multiplicities. The raising coefficient
on the angular sector $0\le l\le n$ is
\[
 \frac12\sqrt{\frac{(n-l+1)(n+l+2)}{(n+1)(n+2)}}\le\frac12.
\]
Its maximum is $1/2$ at $l=0$. This proves the upper
bound, including after physical restriction.
For saturation use the gauge-invariant function
$\chi_J(U_p)$, where $\chi_J$ is the $\mathrm{SU}(2)$
character of highest weight $J$. It has norm one
under product Haar measure and is a pure partial-charge
$J$ function at this transition vertex. The identity
\[
 w_p\chi_J(U_p)
       =\tfrac12\chi_{J+1}(U_p)+\tfrac12\chi_{J-1}(U_p),
 \qquad \chi_{-1}=0,
\]
proves the lower bound in \eqref{c18v7:sharp-source}.

There are three opposite pairs among the six half-edges.
If $o$ of these pairs are split by $R$, the crossing
plaquette count is $m=d(6-d)-o$: each nonopposite pair
gives exactly one plaquette for $N\ge3$. For $d=1,5$,
$m=4$; for $d=2,3,4$, $m$ is $6$ or $8$.
Thus the triangle inequality gives \eqref{c18v7:source}.
\end{proof}

In particular, increasing $J$ does not decrease an
individual Wilson raising coefficient. A small vacuum
tail cannot substitute for this off-diagonal block.

\subsection{Exact complementary return, including its norm metric}

Assume $\delta_J>0$, and set $\mathcal A_J=PKP$.
For real $z<\delta_J$, define
\[
 \begin{aligned}
 \Sigma_J(z)&=\mathcal B_J^*(\mathcal C_J-z)^{-1}\mathcal B_J,\\
 0\le\Sigma_J(z)&\le
       \frac{m^2b^2}{4(\delta_J-z)}\,\Pi_J,\\
 0\le\Sigma_J'(z)&\le
       \frac{m^2b^2}{4(\delta_J-z)^2}\,\Pi_J .
 \end{aligned}
 \tag{V7.17}\label{c18v7:return}
\]
These are bounded operators on the retained space.
The support on $\Pi_J$ is part of the statement.
It is stronger than an undifferentiated scalar bound
on all of $P$. Functional calculus and
\eqref{c18v7:source} prove the inequalities.

The exact Schur operator is
\[
 S_J(z)=\mathcal A_J-zP-\Sigma_J(z).
 \tag{V7.18}\label{c18v7:Schur}
\]
For $v$ in its operator domain,
$S_J(z)v=0$ if and only if
$\psi=v-(\mathcal C_J-z)^{-1}\mathcal B_Jv$ is an
eigenvector of $K$ at $z$. This follows by solving
the $Q$ block equation and substituting into the $P$
equation. The diagonal compressions have the domains
already described, and the off-diagonal block is bounded,
so there is no domain extension hidden in this identity.

At $z=0$, square completion on the form domain gives
\[
 \langle v+w,K(v+w)\rangle
 =\langle v,S_J(0)v\rangle+
   \|\mathcal C_J^{1/2}
           (w+\mathcal C_J^{-1}\mathcal B_Jv)\|^2 .
\]
Thus $S_J(0)\ge0$. Its kernel is precisely the span
of $v_0=P\Omega$, which is nonzero because
$\mathcal C_J>0$ and $K\Omega=0$; uniqueness follows
from that of the finite-volume ground. Reconstruction
also retains the exact norm:
\[
 \|\psi\|^2
   =\langle v,(P+\Sigma_J'(z))v\rangle .
 \tag{V7.19}\label{c18v7:vacuum-metric}
\]
This is standard Feshbach algebra applied to the new
native complementary bound, not a new general Schur
theorem. A targeted source check is
Dusson, Sigal and Stamm, \emph{The Feshbach--Schur map
and perturbation theory} (2021),
\url{https://arxiv.org/abs/2105.02058}; no quantitative
Yang--Mills estimate is imported from that abstract.

Neither a gap of $S_J(0)$ above $v_0$ nor locality
of $\Sigma_J(z)$ has been proved. In particular the
resolvent in \eqref{c18v7:return} is a many-link operator,
although the source itself is supported at a single
star. A positive discarded block and a one-dimensional
zero eigenspace are not a volume-uniform retained gap.

\subsection{Three distinct cutoff obligations and the next bottleneck}

Put $c_1=(9/4)12^{1/3}$. If
$J=L\beta_*^{1/3}+O(1)$ and $L^2/d^2>c_1$, then
\[
 \begin{gathered}
 \frac{\delta_J}{\kappa}
     \sim (L^2/d^2-c_1)\beta_*^{2/3},\\
 \frac{mb}{2\delta_J}
     \sim\frac{m}{2(L^2/d^2-c_1)}\beta_*^{1/3}.
 \end{gathered}
 \tag{V7.20}\label{c18v7:non-small}
\]
The second expression is the available \emph{upper
majorant} for
$\|\mathcal C_J^{-1}\mathcal B_J\|$.
Its divergence is not a lower bound on that norm and
is not a theorem that the true return is large.
It states exactly why these estimates alone do not
close a small-return argument.

For comparison $J=L\beta_*^{1/2}+O(1)$ makes the
same majorant tend to $md^2/(2L^2)$, which can be made
small by its prefactor. The relevant obligations are:
\begin{center}
\begin{tabular}{@{}p{.19\textwidth}p{.73\textwidth}@{}}
\toprule
Cutoff order & What has actually been established\\
\midrule
$\beta_*^{1/4}$ &
Necessary and sufficient order for fixed vacuum
retention at a transition vertex (V6).\\
$\beta_*^{1/3}$ &
Sufficient for a coercive complementary charge block
and fixed-error bounded-$E/\kappa$ window retention (V7).\\
$\beta_*^{1/2}$ &
Sufficient for a small crude inverse-times-source
majorant; not a necessary scale for the actual return.\\
\bottomrule
\end{tabular}
\end{center}

The useful next question is therefore a source-sensitive
bound on $\mathcal C_J^{-1}\mathcal B_J$ and its neutral
return, retaining the growing charge bank and the exact
norm metric. It must exploit more than the global
complementary lower bound and the primitive source norm.
Blindly enlarging the raw charge cutoff would not prove
a renormalized local effective dynamics, a retained
uniform gap, or the continuum mass gap.

\subsection{Checks, nonduplication and preservation}

The companion programs were reviewed at V5; the next
five-version checkpoint and the next broad literature
sweep are V10. This append uses the native tempered
trial and charge analysis, with a targeted check of
the already used ultraspherical recurrence at
\url{https://dlmf.nist.gov/18.9}. It does not count this
targeted lookup as a broad sweep.

The diagnostic checks the scalar optimization, all
sixty-two nontrivial partial stars, the raising
coefficient, and cutoff substitutions. A coupled
two-group scalar Galerkin fixture tests the block
identities and norm metric. That fixture is not
a computed Wilson vacuum or a continuum model, and
floating-point checks are not the proof of the
operator inequalities. The proofs are above.
V6's mathematical diagnostics are replayed.

The complete V6 body is preserved, apart from its
version title. The active source becomes V7 and all
forty-seven other note sources are unchanged.
Only \texttt{.tex} files remain in \texttt{latex\_notes};
all build and diagnostic products stay outside it.
The full interacting four-dimensional continuum
construction and physical Yang--Mills mass gap
remain unproved.

% END CYCLE 18 VERSION 7 APPEND

% BEGIN CYCLE 18 VERSION 8 APPEND
\clearpage
\section{V8: Buffered charge decay for an entire low-energy window}
\label{c18v8:main}

\subsection{The source-sensitive improvement and its scope}

V7's unweighted inverse-times-source majorant grows at its
improved cubic cutoff. That does not imply that a low-energy
state has a large source on a more distant charge shell.
We now use the exact nearest-neighbour charge ladder and
self-adjointness to distinguish these questions. Exponential
conjugation costs $\cosh\eta-1$, rather than $e^\eta-1$,
in the real part of the quadratic form. A buffer within
the same cubic order then gives exponentially small tails
for the \emph{whole} native low-energy spectral subspace.

The whole-window assertion is proved by a Sylvester
equation, not by summing separate eigenfunction estimates.
There is no square root of the spectral-window dimension.
We obtain small return columns on that window and an
explicit Ritz comparison with the retained compression.
The latter isolates a concrete remaining task: a
volume-uniform gap of the interacting retained operator.
We do not prove that gap or the continuum construction.

This is a native application of standard weighted-resolvent
and variational methods, not a new general Combes--Thomas
principle. The archive's f1 V16 and f7 V46 use a full
conjugation-norm debit; f15 V94 buffers an isolated magnetic
multiplier using a scalar random-walk kernel. Those statements
do not give the present whole-window bound for the restored
many-link Hamiltonian. V7 supplies the missing native
complementary coercivity, and its sharp individual raising
coefficient supplies the volume-independent hopping constant.

\subsection{The exact charge ladder and its weighted real part}

Keep one partial star and its physical or full-Haar
realization from V7. In the physical realization put
$d=\min(|R|,6-|R|)$; in the full-Haar realization put
$d=|R|$. Let $\Pi_n$ be its charge-$n$ projections,
\[
 P_j=\sum_{n=0}^j\Pi_n,\qquad Q_j=I-P_j,\qquad
 t=\frac{mb}{2},\qquad 4\le m\le8.
 \tag{V8.1}\label{c18v8:notation}
\]
Assume $b>0$ in the nontrivial estimates below.
V7 proves that
\[
 \begin{gathered}
 \Pi_nK\Pi_l=0\quad(|n-l|>1),\qquad
 A_n:=\Pi_{n+1}K\Pi_n,\quad \|A_n\|\le t,\\
 \mathcal C_j=Q_jKQ_j\ge\delta_j Q_j,\qquad
 \delta_j=\frac{\kappa(j+1)(j+3)}{d^2}-\Lambda(\kappa,b),\\
 \mathcal B_j=Q_jKP_j=A_j .
 \end{gathered}
 \tag{V8.2}\label{c18v8:ladder}
\]
The charge-diagonal part can be unbounded and can contain
arbitrarily complicated exterior dynamics. It is not frozen.
The raising sum is bounded, since its domains and ranges
are mutually orthogonal; its norm is at most $t$.
In particular the off-diagonal norm does not count the
interior multiplicities.

Fix a base cutoff $j_0$, and on $\operatorname{Ran}Q_{j_0}$
choose bounded weights
\[
 W_M=\sum_{n\ge j_0+1}
       \exp\!\left(\eta\min\{n-j_0-1,M\}\right)\Pi_n,
 \quad \eta>0,\quad M=1,2,\ldots .
 \tag{V8.3}\label{c18v8:weights}
\]
They commute with every electric Laplacian, preserve the
electric domain, and equal one on the source shell $j_0+1$.

\begin{lemma}[Quadratic conjugation debit]
\label{c18v8:realpart}
Set $\mathcal C=\mathcal C_{j_0}$,
$\delta=\delta_{j_0}$, and
$\mathcal C_M=W_M\mathcal C W_M^{-1}$.
On the operator domain, and hence on the associated form domain,
\[
 \operatorname{Re}\langle u,\mathcal C_Mu\rangle
 \ge \bigl(\delta-q(\eta)\bigr)\|u\|^2,\qquad
 q(\eta)=2t(\cosh\eta-1).
 \tag{V8.4}\label{c18v8:debit}
\]
The bound is independent of $M$, lattice volume and
charge-sector multiplicities.
\end{lemma}

\begin{proof}
Let $\varphi_n=\eta\min\{n-j_0-1,M\}$. An adjacent block
$A_n$ acquires the factor $e^{\varphi_{n+1}-\varphi_n}$,
and its adjoint acquires the reciprocal. The Hermitian
part of their change is
\[
 (\cosh(\varphi_{n+1}-\varphi_n)-1)(A_n+A_n^*).
\]
The raising sum of these blocks has norm at most
$t(\cosh\eta-1)$: each input charge and each output
charge occurs only once. Adding its adjoint bounds
the absolute Hermitian change by $q(\eta)$.
All charge-diagonal terms are unchanged. This proves
\eqref{c18v8:debit} using $\mathcal C\ge\delta$.
The complete, possibly non-self-adjoint difference
$\mathcal C_M-\mathcal C$ is also bounded, uniformly
in $M$ for fixed $\eta$; a bound
$2t(e^\eta-1)$ suffices for domain and semigroup purposes.
That larger bound is not used as the coercivity debit.
\end{proof}

For $\delta-q(\eta)-z>0$, $z\in\mathbb R$, the bounded
similarity and \eqref{c18v8:debit} give
\[
 \|W_M(\mathcal C-z)^{-1}W_M^{-1}\|
       \le\frac1{\delta-z-q(\eta)}.
 \tag{V8.5}\label{c18v8:weighted-resolvent}
\]
Indeed the inverse exists by similarity to $\mathcal C-z$;
Cauchy--Schwarz applied to the real-part lower bound
controls its norm. Also
\[
 \|e^{-s\mathcal C_M}\|
       \le e^{-s(\delta-q(\eta))},\qquad s\ge0 .
 \tag{V8.6}\label{c18v8:semigroup}
\]
One can differentiate the squared norm on the operator
domain and use \eqref{c18v8:debit}, then extend by density.
The semigroup exists because this is a bounded
perturbation of the self-adjoint generator, or directly
because $e^{-s\mathcal C_M}=W_Me^{-s\mathcal C}W_M^{-1}$.
No positivity of the charge-basis matrix entries is assumed.

\subsection{Whole-window localization without a rank factor}

Choose an energy ceiling $E\ge0$ and a base cutoff with
$a:=\delta_{j_0}-E>0$. Define
\[
 \eta=\operatorname{arcosh}\!\left(1+\frac{a}{4t}\right),
 \qquad
 \varepsilon_k=\frac{2t}{a}
          e^{-\eta(k-j_0)},\quad k\ge j_0 .
 \tag{V8.7}\label{c18v8:parameters}
\]
Then $q(\eta)=a/2$. Let $F_E=1_{[0,E]}(K)$.

\begin{theorem}[Buffered native spectral-window tail]
\label{c18v8:window}
For every integer $k\ge j_0$,
\[
          \|Q_kF_E\|\le\min\{1,\varepsilon_k\}.
 \tag{V8.8}\label{c18v8:tail}
\]
For every real $z\le E$, the actual base-cut source
resolvent satisfies
\[
 \|Q_k(\mathcal C_{j_0}-z)^{-1}\mathcal B_{j_0}\|
       \le \varepsilon_k.
 \tag{V8.9}\label{c18v8:column}
\]
Both estimates retain all exterior dynamics and
hold in the full or physical Hilbert space as specified.
\end{theorem}

\begin{proof}
Regard $F_E$ as the isometric inclusion of
$\mathcal H_E=\operatorname{Ran}F_E$, and write
$K_E=K|_{\mathcal H_E}$, so $0\le K_E\le E$.
Put $X=P_{j_0}F_E$, $Y=Q_{j_0}F_E$ and
$\mathcal B=\mathcal B_{j_0}$. The exact $Q$ equation is
\[
                  \mathcal C Y-YK_E=-\mathcal B X .
 \tag{V8.10}\label{c18v8:Sylvester}
\]
The spectral range belongs to the operator domain;
the charge projections preserve it. There is therefore
no undefined product in this equation.

Multiply by $W_M$. Since $W_M\mathcal B=\mathcal B$,
$Y_M=W_MY$ solves
$\mathcal C_MY_M-Y_MK_E=-\mathcal BX$. Its solution is
\[
 Y_M=-\int_0^\infty
       e^{-s\mathcal C_M}\mathcal BX e^{sK_E}\,ds,
 \qquad \|Y_M\|\le\frac{2t}{a}.
 \tag{V8.11}\label{c18v8:integral}
\]
The integral converges in operator norm by
\eqref{c18v8:semigroup}, since
$\delta-q(\eta)-E=a/2$. For completeness, integrating
the derivative of
$e^{-s\mathcal C_M}Y_Me^{sK_E}$ proves the identity;
its terminal term vanishes by the same bound.
Here $\|X\|\le1$. This is an operator-norm calculation
on the entire $\mathcal H_E$, not a separate estimate
for each eigenvector.

When $M\ge k-j_0$, $W_M$ is at least
$e^{\eta(k-j_0)}$ on $Q_k$. Thus
\eqref{c18v8:integral} proves \eqref{c18v8:tail}.
For \eqref{c18v8:column} apply
\eqref{c18v8:weighted-resolvent} to the source shell,
where $W_M=1$, and use $\|\mathcal B\|\le t$.
Its denominator is at least $a/2$.
Again choose $M\ge k-j_0$.
All estimates use bounded weights; no unbounded
similarity transform is asserted.
\end{proof}

If $b=0$, the charge sectors are invariant and the
tails vanish whenever $\delta_{j_0}>E$. This elementary
case is separate from \eqref{c18v8:parameters}.

\subsection{The outer source, exact return and norm metric}

Now choose an outer cutoff $J\ge j_0+1$. Let
$\mathcal A_J=P_JKP_J$ and use the exact
$\mathcal B_J,\mathcal C_J,\Sigma_J(z)$ of V7.
Since $\Pi_J\le Q_{J-1}$, the theorem gives
\[
 \|\mathcal B_JP_JF_E\|\le t\varepsilon_{J-1}
                  =t e^\eta\varepsilon_J .
 \tag{V8.12}\label{c18v8:outer-source}
\]
This is smallness on the actual spectral window,
not a bound making $\|\mathcal B_J\|$ itself small.
For real $z\le E$, functional calculus yields
\[
 \begin{aligned}
 \|(\mathcal C_J-z)^{-1}\mathcal B_JP_JF_E\|
   &\le\frac{t\varepsilon_{J-1}}{\delta_J-E},\\
 \|\Sigma_J(z)P_JF_E\|
   &\le\frac{t^2\varepsilon_{J-1}}{\delta_J-E},\\
 \|F_EP_J\Sigma_J(z)P_JF_E\|
   &\le\frac{t^2\varepsilon_{J-1}^2}{\delta_J-E},\\
 \|F_EP_J\Sigma_J'(z)P_JF_E\|
   &\le\frac{t^2\varepsilon_{J-1}^2}{(\delta_J-E)^2}.
 \end{aligned}
 \tag{V8.13}\label{c18v8:outer-return}
\]
Here $F_E$ on the left can equivalently mean the
orthogonal projection on the original Hilbert space.
Each inequality follows by factoring the exact
return through its source, using
\eqref{c18v8:outer-source} and
$\mathcal C_J-z\ge\delta_J-E>0$.
The last line controls the correction to V7's
reconstruction norm metric on these columns.
It does not identify the exact metric with the identity.

\subsection{A retained Ritz comparison that excludes hidden low states}

For each finite torus, $K$ has compact resolvent.
Its compression $\mathcal A_J$, with its reduced electric
domain, also has compact resolvent: its form-domain
unit ball embeds compactly as a subspace of the original
electric form domain. Denote their eigenvalues,
including multiplicities, by
\[
 0=\lambda_0\le\lambda_1\le\cdots,\qquad
 \nu_0^{(J)}\le\nu_1^{(J)}\le\cdots .
\]
All indices refer consistently to either the physical
or the full Hilbert space.

\begin{theorem}[Quantitative low-energy Ritz enclosure]
\label{c18v8:Ritz}
Suppose $\varepsilon_J<1$, and put
\[
 r_J(E)=
 \frac{(2E+t e^\eta)\varepsilon_J^2}
      {1-\varepsilon_J^2}.
 \tag{V8.14}\label{c18v8:Ritz-error}
\]
For every index $r$ with $\lambda_r\le E$,
\[
        \lambda_r\le\nu_r^{(J)}
                    \le\lambda_r+r_J(E).
 \tag{V8.15}\label{c18v8:Ritz-bound}
\]
In particular $0\le\nu_0^{(J)}\le r_J(E)$. If the
retained compression has gap
$g_J=\nu_1^{(J)}-\nu_0^{(J)}$, then
\[
             \lambda_1\ge\min\{E,g_J-r_J(E)\}.
 \tag{V8.16}\label{c18v8:gap-transfer}
\]
Only when $g_J-r_J(E)>0$ and $E>0$ is this a
positive gap statement. No such uniform retained
gap is supplied by this theorem.
\end{theorem}

\begin{proof}
The lower bound in \eqref{c18v8:Ritz-bound} is the
min--max principle for a restricted form domain.
Let $\mathcal F_r$ span the first $r+1$ eigenvectors
of $K$. Its restriction $L=K|_{\mathcal F_r}$ obeys
$0\le L\le\lambda_r\le E$. On this space set
$X=P_J|_{\mathcal F_r}$ and $Y=Q_J|_{\mathcal F_r}$.
Then
\[
 X^*X=I-Y^*Y,\quad \|Y\|\le\varepsilon_J,\quad
 \|\mathcal B_JX\|\le t\varepsilon_{J-1}.
\]
In particular $X$ is injective and its image has
dimension $r+1$. The exact $P$ equation gives
$\mathcal A_JX+\mathcal B_J^*Y=XL$, hence
\[
 X^*\mathcal A_JX
       =L-Y^*YL-X^*\mathcal B_J^*Y .
\]
For a unit vector $u\in\mathcal F_r$, take real parts
and use the displayed norm bounds. This yields
\[
 \langle Xu,\mathcal A_JXu\rangle
 \le\lambda_r+E\varepsilon_J^2
                 +t\varepsilon_{J-1}\varepsilon_J,\qquad
 \|Xu\|^2\ge1-\varepsilon_J^2.
\]
No commutation of $L$ and $Y^*Y$ was used.
After division, subtracting $\lambda_r$ and using
$\lambda_r\le E$ proves the upper bound
\eqref{c18v8:Ritz-bound} by min--max on
$X\mathcal F_r$.

For \eqref{c18v8:gap-transfer}, if $\lambda_1>E$
there is nothing to prove. Otherwise
$\nu_1^{(J)}\le\lambda_1+r_J(E)$, so
$\lambda_1\ge g_J+\nu_0^{(J)}-r_J(E)
\ge g_J-r_J(E)$.
The same reasoning covers $\lambda_1=E$.
\end{proof}

The retained gap $g_J$ can be computed or estimated
from $P_JHP_J$ without knowing $E_0$, since scalar
centering does not change a gap. The comparison above
does not assume a pre-existing gap of $K$.
It bounds every low eigenvalue in the stated window,
so it does not merely construct a few retained
quasimodes while allowing hidden lower states.

\subsection{Cubic-order buffering and the remaining physical task}

Write $\beta_*=b/\kappa$ and $c_1=(9/4)12^{1/3}$.
Consider bounded $E/\kappa$ as $\beta_*\to\infty$,
and choose fixed $L_0,s>0$ with
$D:=L_0^2/d^2-c_1>0$, setting
\[
 j_0=L_0\beta_*^{1/3}+O(1),\qquad
 J-j_0=s\beta_*^{1/3}+O(1).
 \tag{V8.17}\label{c18v8:scales}
\]
For integer choices with these asymptotics,
V7's explicit $\Lambda$ and \eqref{c18v8:parameters}
give
\[
 \begin{gathered}
 \frac a\kappa\sim D\beta_*^{2/3},\qquad
 \eta\sim\sqrt{\frac Dm}\,\beta_*^{-1/6},\\
 \varepsilon_J=
 \frac mD\,\beta_*^{1/3}
   \exp\!\left(-s\sqrt{\frac Dm}\,
                    \beta_*^{1/6}+o(1)\right)(1+o(1)).
 \end{gathered}
 \tag{V8.18}\label{c18v8:asymptotic}
\]
The $o(1)$ in the exponent is justified by the
next orders in \eqref{c18v7:F}, by the expansion
$\operatorname{arcosh}(1+x)=\sqrt{2x}(1+O(x))$,
and by the $O(1)$ rounding errors. In particular,
for every $c<s\sqrt{D/m}$ and all sufficiently
large $\beta_*$,
\[
 \varepsilon_J\le C\beta_*^{1/3}e^{-c\beta_*^{1/6}},
 \qquad
 \frac{r_J(E)}{\kappa}
       \le C'\beta_*^{5/3}e^{-2c\beta_*^{1/6}} .
 \tag{V8.19}\label{c18v8:asymptotic-bounds}
\]
The constants are independent of the torus size.
All the window-restricted returns in
\eqref{c18v8:outer-return} likewise have a polynomial
prefactor times exponential decay. These estimates
improve the algebraic whole-window tail from V7.
They do not improve V6's \emph{necessary} vacuum
bandwidth or claim optimality of the cubic scale.

There is still no small operator-norm bound on the
full outer return for arbitrary retained vectors.
Its small columns are restricted to actual low-energy
states. The Ritz theorem makes that restriction useful
without turning it into an unrestricted assertion:
a proved retained gap, if available, would transfer
with the displayed error.

This is localization in a \emph{charge label}, not
spatial clustering. A single-star compression leaves
all exterior and interior multiplicities untruncated.
For a finite bank of commuting cutoffs, the tail
bounds imply only
\[
 \left\|\left(I-\prod_iP_{J_i}^{(i)}\right)F_E\right\|^2
       \le\sum_i\varepsilon_{J_i}^{\,2}.
 \tag{V8.20}\label{c18v8:bank}
\]
Indeed $I-\prod_iP_i\le\sum_i(I-P_i)$ as commuting
projections. The number of boundary stars remains
in this estimate; no extensive-bank spectral comparison
is silently deduced from the one-star result.

The next task is therefore the interacting retained
dynamics and its genuinely collective charge/neutral
structure, not another proof that a remote charge
shell has a small low-energy weight. A finite-regulator
retained gap would still need a physically normalized
refinement trajectory and a nontrivial continuum
construction. The full four-dimensional Yang--Mills
mass gap remains unproved.

\subsection{Source audit, diagnostics and preservation}

The targeted external check located Germinet--Klein's
\emph{Operator kernel estimates for functions of
generalized Schr\"odinger operators} (2001),
\url{https://web.ma.utexas.edu/mp_arc-bin/mpa?yn=01-289},
and Klein--L\'eonard--Rosenberger's
\emph{Agmon-type estimates for a class of jump processes}
(2013), \url{https://arxiv.org/abs/1309.3918}.
Their abstracts provide context for weighted-resolvent
and exponential-decay methods; no theorem from those
papers is used without proof here. The complete
block-operator proof needed for V8 is above.
This targeted check is not the broad literature
checkpoint, which remains due at V10 together with
the next companion review.

Diagnostics test the Hermitian conjugation identity
with complex noncommuting matrix blocks, the
whole-window Sylvester identity, source and return
bounds, and the Ritz enclosure. A coupled two-group
scalar Galerkin fixture checks the same statements
after a finite harmonic truncation.
These are finite numerical diagnostics, not native
many-link spectra or rigorous truncation certificates.
The asymptotic constants are also checked symbolically
and by explicit substitutions, and V7's mathematical
diagnostics are replayed.

The complete V7 body is preserved apart from its
version title. V8 replaces the active V7 filename;
all forty-seven other note sources are unchanged.
Only \texttt{.tex} files are kept in
\texttt{latex\_notes}; all diagnostics and rendering
products remain outside it.

% END CYCLE 18 VERSION 8 APPEND

% BEGIN CYCLE 18 VERSION 9 APPEND
\clearpage
\section{V9: The collective electric cutoff and its complete low-spectrum error}
\label{c18v9:main}

\subsection{From one boundary star to a specified retained Hamiltonian}

V8 controlled one partial-charge cutoff. Its retained
space was still infinite dimensional, and its union
bound for a bank did not by itself give a joint
spectral comparison. We now close that particular
assembly step for the complete bank of individual
link-electric cutoffs. The resulting Hamiltonian is
finite dimensional, gauge invariant and local, with
the original electric and magnetic coefficients.

The main new estimate is not another one-link tail.
It controls the \emph{joint} source return in the
Rayleigh quotient. Ordering the first rejected link
assigns each term to an actual plaquette--link
incidence. This keeps the spectral error proportional
to the number of links, not its square.
An explicit cutoff then makes the entire low-energy
spectral error smaller than any prescribed tolerance.

A finite-dimensional retained theory is not a
mass-gap proof. We obtain a complete spectral
comparison, not a uniform lower bound for the
retained gap. The large weak-electric coupling ratio
is unchanged, and the full interacting continuum
construction remains unproved.

\subsection{One-link inputs, joint carrier and exact gauge restriction}

Work on the cubic torus $N\ge3$ and write
$M=|\mathcal E_N|=|\mathcal P_N|=3N^3$.
Let $\Pi_{n,e}$ project onto Peter--Weyl label $n$
of link $e$, so $L_e$ has eigenvalue $n(n+2)$ there.
For integers $J_e\ge0$ define
\[
 P_e=\sum_{n=0}^{J_e}\Pi_{n,e},\quad Q_e=I-P_e,\qquad
 P=\prod_{e\in\mathcal E_N}P_e,\quad Q=I-P .
 \tag{V9.1}\label{c18v9:bank}
\]
All these projections commute with the electric
operator and the original Gauss action.
A single-link partial star has $d=1$, $m=4$ in
V8, and its Casimir is precisely $L_e$.

Fix $b>0$, an energy ceiling $E\ge0$, and a base
cutoff $j_0$ with
\[
 \begin{gathered}
 a=\kappa(j_0+1)(j_0+3)-\Lambda(\kappa,b)-E>0,\\
 \eta=\operatorname{arcosh}\!\left(1+\frac a{8b}\right),
 \qquad
 \varepsilon_k=\frac{4b}{a}e^{-\eta(k-j_0)}
                  \quad(k\ge j_0).
 \end{gathered}
 \tag{V9.2}\label{c18v9:parameters}
\]
V8, with hopping bound $t=2b$, proves for every link
and the \emph{same} whole spectral window
$F_E=1_{[0,E]}(K)$ that
\[
 \|Q_{k,e}F_E\|\le\varepsilon_k,\qquad
 Q_{k,e}=\sum_{n>k}\Pi_{n,e}.
 \tag{V9.3}\label{c18v9:one-link}
\]
No independence of different links is assumed.
Take $J_e\ge j_0+1$ and put
\[
 \varepsilon_e=\varepsilon_{J_e},\qquad
 \widehat\varepsilon_e=\varepsilon_{J_e-1}
                         =e^\eta\varepsilon_e,\qquad
 \Theta=\sum_e\varepsilon_e^2 .
 \tag{V9.4}\label{c18v9:Theta}
\]

The retained full-Haar carrier is the tensor product
of finite link spaces with dimensions
\[
 D_J=\sum_{n=0}^J(n+1)^2
       =\frac{(J+1)(J+2)(2J+3)}6,\qquad
             \dim\operatorname{Ran}P=\prod_eD_{J_e}.
 \tag{V9.5}\label{c18v9:dimension}
\]
This keeps all matrix coefficients in each retained
representation, not only its character.
For a plaquette $p$, set $P_p=\prod_{e\in p}P_e$,
acting on its four-link space. The exact compression
on the retained tensor product is
\[
 \begin{gathered}
 H_{\mathbf J}=PHP|_{\operatorname{Ran}P}
       =\kappa\sum_e L_e^{(J_e)}
               +b\sum_p(I-\overline w_{p,\mathbf J}),\\
 \overline w_{p,\mathbf J}=P_p w_pP_p
       \quad\hbox{on the four retained link factors}.
 \end{gathered}
 \tag{V9.6}\label{c18v9:H}
\]
Identities on other retained factors are implicit.
Each magnetic term is positive, since
$-I\le\overline w_{p,\mathbf J}\le I$.
Its support is exactly within the original plaquette.
No nonlocal interaction is introduced by writing
this compression, although eliminating degrees of
freedom by a Schur return would be different.

The physical retained carrier is
$\operatorname{Ran}P\cap\mathcal H_{\rm phys}$.
Equivalently, average the finite-dimensional gauge
representation over one normalized Haar variable
at each vertex and restrict to the range of that
orthogonal projector. Commutation with $P$ and $H$
makes this the exact restriction of
\eqref{c18v9:H}. The physical spectrum is not
replaced by the full tensor-product spectrum.
The physical dimension is at most
\eqref{c18v9:dimension}; no efficient matrix
diagonalization is claimed.

\subsection{First rejected link: a joint source estimate with local incidence}

We state the useful estimate first with arbitrary
tail bounds. In this subsection suppose
\[
 \|Q_eF_E\|\le u_e,\qquad
 \|\Pi_{J_e,e}F_E\|\le v_e,
 \qquad T_{\rm tail}=\sum_eu_e^2 .
 \tag{V9.7}\label{c18v9:abstract-tails}
\]
For the native choices above one can take
$u_e=\varepsilon_e$ and
$v_e=\widehat\varepsilon_e$.

\begin{lemma}[Joint tail and two-sided source debit]
\label{c18v9:joint-source}
The complete bank satisfies
\[
 \|QF_E\|^2\le T_{\rm tail},\qquad
 \|F_EPKQF_E\|
       \le S_{\rm cross}:=
          \frac b2\sum_p\sum_{e\in p}u_ev_e .
 \tag{V9.8}\label{c18v9:joint-estimate}
\]
In particular the cubic incidence count gives
\[
 \|QF_E\|^2\le\Theta,\qquad
            \|F_EPKQF_E\|\le2b e^\eta\Theta .
 \tag{V9.9}\label{c18v9:native-cross}
\]
The source estimate is on the actual low-energy
window on both sides. It does not bound the
unrestricted norm of $PKQ$ by this small number.
\end{lemma}

\begin{proof}
Commuting projections give
$Q=I-\prod_eP_e\le\sum_eQ_e$. Taking quadratic
forms on the whole spectral window proves the tail
bound.

Order all links once and define
$R_{<e}=\prod_{f<e}P_f$. The exact first-rejection
identity is
\[
                  Q=\sum_e Q_eR_{<e}.
 \tag{V9.10}\label{c18v9:first-rejection}
\]
For example it follows by successively subtracting
adjacent partial products. The summands are mutually
orthogonal projections. No probabilistic independence
is involved.

Write $V_p=-bw_p$. The electric terms and scalar
centering have no off-diagonal block, so
\[
 PKQ=\sum_p\sum_{e\in p}PV_pQ_eR_{<e}.
 \tag{V9.11}\label{c18v9:source-sum}
\]
Indeed terms with $e\notin p$ vanish, because $V_p$
commutes with $P_e$ and $P$ contains $P_e$.
For an incident link the exact single-link selection
rule and V7's sharp raising coefficient give
\[
 P_eV_pQ_e
     =\Pi_{J_e,e}V_p\Pi_{J_e+1,e},\qquad
 \|\Pi_{J_e,e}V_p\Pi_{J_e+1,e}\|\le b/2 .
 \tag{V9.12}\label{c18v9:edge-block}
\]
The remaining projections commute with both
one-link shell projections. Therefore
\[
 \begin{aligned}
 \|F_EPV_pQ_eR_{<e}F_E\|
 &\le\frac b2\,
       \|\Pi_{J_e,e}PF_E\|\,
       \|\Pi_{J_e+1,e}R_{<e}F_E\|\\
 &\le\frac b2\,v_eu_e .
 \end{aligned}
\]
In the last step we moved the commuting contractions
to the outside; we did not assume that they preserve
$\operatorname{Ran}F_E$. Summing proves
\eqref{c18v9:joint-estimate}.
Each cubic link belongs to exactly four plaquettes,
so $\sum_p\sum_{e\in p}\varepsilon_e^2=4\Theta$.
This proves \eqref{c18v9:native-cross}.
\end{proof}

Using only $\|V\|\le Mb$ after the union bound
would lose the local incidence information.
The two window factors in
\eqref{c18v9:joint-estimate} are both essential.
They keep a quadratic tail cost despite a joint
projection over the entire lattice.

\subsection{All low eigenvalues of the finite retained matrix}

Set $A_{\mathbf J}=PKP|_{\operatorname{Ran}P}$.
Use either the full or the physical realization
consistently, and denote the eigenvalues by
\[
 0=\lambda_0\le\lambda_1\le\cdots
          \quad\hbox{for }K,\qquad
 \nu_0\le\nu_1\le\cdots
          \quad\hbox{for }A_{\mathbf J}.
\]
The finite carrier contains the low subspaces
needed below whenever its tail bound is less
than one.

\begin{theorem}[Collective spectral enclosure]
\label{c18v9:spectrum}
Under \eqref{c18v9:abstract-tails}, if
$T_{\rm tail}<1$, define
\[
 R_{\mathbf J}(E)=
       \frac{2E T_{\rm tail}+S_{\rm cross}}
            {1-T_{\rm tail}} .
 \tag{V9.13}\label{c18v9:general-error}
\]
For every index $r$ with $\lambda_r\le E$,
\[
                \lambda_r\le\nu_r
                   \le\lambda_r+R_{\mathbf J}(E).
 \tag{V9.14}\label{c18v9:enclosure}
\]
For the native bounds \eqref{c18v9:Theta}, it is
enough that $\Theta<1$, and one may take
\[
 R_{\mathbf J}(E)=
       \frac{(2E+2b e^\eta)\Theta}{1-\Theta}.
 \tag{V9.15}\label{c18v9:native-error}
\]
There is no assumption of a gap between the
target low eigenvalues.
\end{theorem}

\begin{proof}
Min--max gives $\lambda_r\le\nu_r$ whenever that
retained index exists. To prove existence and the
upper bound, let $\mathcal F_r$ span the first
$r+1$ eigenvectors of $K$, and let
$L=K|_{\mathcal F_r}$, so $0\le L\le\lambda_r\le E$.
Put $X=P|_{\mathcal F_r}$ and $Y=Q|_{\mathcal F_r}$.
Then
\[
 X^*X=I-Y^*Y\ge(1-T_{\rm tail})I .
\]
Thus $X$ is injective and has an $(r+1)$-dimensional
image in the retained carrier. The exact block
equation yields
\[
       X^*A_{\mathbf J}X=L-Y^*YL-X^*PKQY .
\]
Here $QY=Y$, and the last term is the restriction
of the source form in \eqref{c18v9:joint-estimate}.
For a unit vector in $\mathcal F_r$, its real
quadratic form is at most
$\lambda_r+E T_{\rm tail}+S_{\rm cross}$.
Divide by $\|X\psi\|^2\ge1-T_{\rm tail}$ and use
$\lambda_r\le E$. Min--max on $X\mathcal F_r$
proves \eqref{c18v9:enclosure}.
Substituting \eqref{c18v9:native-cross} proves
\eqref{c18v9:native-error}.
No commutation of $L$ and $Y^*Y$ was used.
\end{proof}

This argument uses no coercivity estimate or inverse
for the full joint block $QKQ$. The one-link lower
bounds are not combined by taking a maximum of
operator inequalities. Only the proved window tails,
the exact source decomposition and min--max are used.

In particular, if $E_{0,\mathbf J}$ is the
ground energy of $H_{\mathbf J}$ in the chosen
realization, then
\[
              0\le E_{0,\mathbf J}-E_0
                       \le R_{\mathbf J}(E).
 \tag{V9.16}\label{c18v9:ground}
\]
Let $g_{\mathbf J}=\nu_1-\nu_0$ be the retained
gap, and let $\Delta=\lambda_1$ be the native gap.
When the retained space has at least two states,
\[
 \min\{E,g_{\mathbf J}-R_{\mathbf J}(E)\}
       \le\Delta\le g_{\mathbf J}+R_{\mathbf J}(E).
 \tag{V9.17}\label{c18v9:gap}
\]
If $\Delta\le E$, the stronger symmetric statement is
$|\Delta-g_{\mathbf J}|\le R_{\mathbf J}(E)$.
For the upper bound in \eqref{c18v9:gap}, use
$\lambda_1\le\nu_1=g_{\mathbf J}+\nu_0$ and
\eqref{c18v9:ground}; the lower bound follows from
\eqref{c18v9:enclosure} when $\lambda_1\le E$,
and is automatic otherwise. In the first case
both $\nu_1-\lambda_1$ and $\nu_0$ lie in $[0,R]$,
which proves the symmetric statement.

These comparisons include every state below the
declared ceiling. They do not rely on a finite
observable bank exhausting the physical spectrum.
The distinction matters in view of f16.2 V84's
finite-source-record ambiguity and one-sided Ritz
warning, which remain valid and are not contradicted.

\subsection{An explicit joint cutoff for prescribed spectral accuracy}

For a uniform cutoff $J_e=J$, one has
\[
 \Theta=M\left(\frac{4b}{a}\right)^2
                   e^{-2\eta(J-j_0)} .
 \tag{V9.18}\label{c18v9:uniform-tail}
\]
Given any $\zeta>0$, set $C_E=2E+2b e^\eta$ and choose
\[
 J=j_0+\max\left\{1,\
   \left\lceil\frac1{2\eta}
       \log\!\left[
          M\left(\frac{4b}{a}\right)^2
              \frac{C_E+\zeta}{\zeta}\right]
   \right\rceil\right\}.
 \tag{V9.19}\label{c18v9:cutoff}
\]
Then $\Theta\le\zeta/(C_E+\zeta)<1$, and
\[
                         R_{\mathbf J}(E)\le\zeta .
 \tag{V9.20}\label{c18v9:tolerance}
\]
These implications follow by exponentiating
\eqref{c18v9:cutoff} and substituting into
\eqref{c18v9:native-error}. They also hold when
the logarithm is negative: the extra positive
buffer only improves the inequality.

For the weak-electric asymptotics put
$\beta_*=b/\kappa$, $c_1=(9/4)12^{1/3}$, and
choose $j_0=L_0\beta_*^{1/3}+O(1)$ with
$D=L_0^2-c_1>0$. Keep $E/\kappa$ bounded and
$0<\zeta\le\kappa$. Then
$a/\kappa\sim D\beta_*^{2/3}$ and
$\eta\sim(\sqrt D/2)\beta_*^{-1/6}$.
The explicit prescription gives the sufficient order
\[
 J=O\!\left(
   \beta_*^{1/3}+
   \beta_*^{1/6}
      \left[1+\log M+\log\beta_*+
                         \log(\kappa/\zeta)\right]\right).
 \tag{V9.21}\label{c18v9:scaling}
\]
Indeed $(4b/a)^2=O(\beta_*^{2/3})$ and
$(C_E+\zeta)/\zeta=O(\beta_*\kappa/\zeta)$.
This is an upper prescription, not an optimality
or necessary-cutoff theorem.

At fixed coupling and accuracy, this construction
pays a logarithm of the volume in the charge cutoff.
It does not give a volume-independent cutoff for
the complete global projection. The matrix dimension
still grows as $D_J^M$ before Gauss restriction.
Neither the logarithmic cutoff nor exact finite
dimensionality makes this an efficient classical
calculation of the thermodynamic spectrum.

\subsection{What remains inside the retained problem}

The retained Hamiltonian \eqref{c18v9:H} has not
changed the ratio $b/\kappa$ or supplied a
renormalization step. For example, whenever
$n+1\le\min_{e\in p}J_e$, the gauge-invariant
characters $\chi_n(U_p)$ and $\chi_{n+1}(U_p)$
both belong to the physical retained carrier, and
\[
       \langle\chi_{n+1}(U_p),
             (-b\overline w_{p,\mathbf J})
                       \chi_n(U_p)\rangle=-b/2 .
 \tag{V9.22}\label{c18v9:retained-coupling}
\]
This is the existing Haar character identity,
now simply observed inside the joint carrier.
No decay of the primitive coupling is generated
by the cutoff. Retained electric resonances and
neutral collective modes have not been removed.

Although every uncentered term in
\eqref{c18v9:H} is positive, they have no common
zero vector. The electric terms would force the
constant link state; a plaquette magnetic term
does not annihilate that state, including at $J=0$.
Thus a frustration-free gap criterion cannot be
applied merely because these terms are positive.
Subtracting the extensive ground energy does
not by itself create positive, locally
vacuum-annihilating terms.

A usable next spectral input would be a lower
bound for the \emph{actual} physical gaps
$g_{N,\mathbf J(N)}$ along a cutoff family such as
\eqref{c18v9:cutoff}, uniform in $N$ and in the
required weak-coupling regime. A finite numerical
gap, or a bound at only one cutoff or one volume,
does not supply this input. A certified retained
gap greater than $\zeta$ would give the finite-volume
lower bound in \eqref{c18v9:gap}, but uniformity
must be proved separately.

The next literature and companion checkpoint at
V10 should therefore focus on that interacting
spectral problem and its physically normalized
scale evolution. Repeating truncation estimates
without a new retained spectral mechanism would
not advance the remaining mass-gap obligation.
Even a fixed-spacing thermodynamic gap would
still require the nontrivial four-dimensional
continuum construction and physical normalization.
The full Yang--Mills mass gap remains unproved.

\subsection{Nonduplication, checks and preservation}

The targeted literature check is Tong, Albert,
McClean, Preskill and Su,
\emph{Provably accurate simulation of gauge theories
and bosonic systems}, Quantum 6, 816 (2022),
\url{https://arxiv.org/abs/2110.06942}.
Its abstract includes local-quantum-number truncation
and an eigenstate result for spectrally isolated
energies. No theorem from that paper is imported
here. The present whole-window input is V8's
proved estimate, without an assumed native gap;
the new assembly is \eqref{c18v9:joint-estimate}.
Finite-dimensional truncation and min--max themselves
are not claimed as new methods.
The broad sweep remains due at V10, not completed
by this targeted lookup.

The diagnostic checks first-rejection identities,
the exact cubic plaquette incidence count, the
Peter--Weyl dimension formula, and the explicit
cutoff budget. Coupled finite trit fixtures test
the two-sided source debit and the joint Ritz
and gap comparisons. They are not native Wilson
spectra or a computed physical mass.
V8's mathematical diagnostics are replayed.
The proofs, rather than finite floating-point
tests, establish the native statements.

V9 appends to the complete V8 body, changing only
its version title and active filename.
All forty-seven other note sources are preserved.
Only \texttt{.tex} files are kept in
\texttt{latex\_notes}; verification and rendering
products remain outside that folder.

% END CYCLE 18 VERSION 9 APPEND

% BEGIN CYCLE 18 VERSION 10 APPEND
\clearpage
\section{V10: Physical exchange multiplicities and a frustrated spectral certificate}
\label{c18v10:main}

\subsection{The checkpoint changes the spectral target, not the theory}

V9 supplied a complete finite retained Hamiltonian and a bound
on every low eigenvalue. The next question is how to certify
its \emph{interacting physical} gap without pretending that its
positive primitive terms are frustration free. This append
examines a two-copy energy certificate and computes the
gauge carrier on which a local exchange can actually be odd.

The exterior-square principle is already in f15 V81 and is
not claimed as new. That archive also proved a maximal Gauss
defect and exact returns on isolated cycle characters.
The new calculation below gives the full partial-exchange
operator on arbitrary boundary multiplicities, not just
cycle characters. It identifies a minimal three-link
multiplicity channel, while proving that every one- or
two-link exchange has positive physical compression.
A separate native-vacuum bound quantifies the local
entanglement cost uniformly in volume. None of these
kinematic statements is a positive energy-gap estimate.

The required five-version companion audit found all eight
sources unchanged in full-file hashes since V5. We reread
their concluding scope statements. NS stability V26 and
SM--Einstein V72 retain respectively a rank-one kernel
calculation and a fixed-mode, fixed-time many-copy limit.
Shell mixing V16 leaves the noncommuting measured
renormalization and nonperturbative transport open.
Parallel YM V5 excludes uniform tightness of an extensive
total charge. The supplied perturbative ESM V64, source
selection V52, regenerative stationarity V54, and
exact-action bridge V45 still require their actual critical
array, preparation/readout bank, nonlinear native
remainder, and moving-source dynamics, respectively.
No auxiliary gap or coefficient from those systems is
substituted for the present physical spectrum.
The audit record is
\path{artifacts/ym_c18_v10_companion_review.json}.

The broad literature sweep screened finite-size and
semidefinite gap certificates, zero-free partition functions,
spectral stability/prethermalization, stochastic lattice
gauge inequalities, and Hamiltonian scale flows. Its
primary-source decisions are:
\begin{itemize}
\item S.~Rao,
\emph{A convergent hierarchy of spectral gap certificates
for qubit Hamiltonians},
\url{https://arxiv.org/html/2510.08427v1},
Sections 1--2, motivates an additive two-energy certificate
which need not assume frustration-freeness. We use the
elementary energy-sum identity, proved below; we do not
import its qubit hierarchy, complexity or convergence into
the gauge-constrained link system.
\item Rai and collaborators,
\url{https://arxiv.org/abs/2411.03680}, and
Lemm--Lucia,
\url{https://arxiv.org/abs/2409.09685}, give
frustration-free methods. V9's primitive terms do not
satisfy that hypothesis.
\item Rayudu--Takahashi,
\url{https://arxiv.org/abs/2607.10765}, and
Bravyi--Gosset--Liu--Wong,
\url{https://arxiv.org/abs/2607.14401},
obtain field-dependent Lee--Yang spectral bounds.
A qualifying external-field structure or zero-free region
has not been proved for this pure gauge Hamiltonian.
Its electric Casimir cannot simply be renamed that field.
\item Yin--Lucas,
\url{https://arxiv.org/abs/2209.11242}, control long-lived
low-energy dynamics even when a perturbation closes the
gap; that is not a gap-stability theorem for our retained
operator. Shen--Zhu--Zhu,
\url{https://arxiv.org/abs/2204.12737}, concern a specified
strong-coupling lattice regime, not the missing
weak-coupling continuum step. Leder and collaborators,
\url{https://arxiv.org/abs/1006.5710}, describe Hamiltonian
flows with diagrammatic and ghost-dominance assumptions;
these do not supply a proved remainder here.
\end{itemize}
Except for the specified sections of Rao, this screening
used abstracts and metadata, not a complete verification
of the papers. Search-result claims of a full mass-gap
solution were not adopted. The complete review record is
\path{artifacts/ym_c18_v10_literature_review.json}.
The next companion and broad literature checkpoints are
V15 and V20.

\subsection{An additive certificate with both physical copies retained}

Let $\mathcal V_J$ be V9's \emph{physical} retained carrier,
of dimension at least two, and write $H_J$ for the
restriction of \eqref{c18v9:H} to it. Its ordered eigenvalues
are $e_0\le e_1\le\cdots$, including multiplicities.
On $\mathcal V_J\otimes\mathcal V_J$ let $S$ be the
global exchange and define
\[
 \mathcal C_J=(H_J\otimes I+I\otimes H_J)
           \big|_{\ker(S+I)},\qquad
 \alpha_J=\lambda_{\min}(\mathcal C_J).
 \tag{V10.1}\label{c18v10:two-copy}
\]
For an orthonormal eigenbasis, the vectors
$(u_i\otimes u_j-u_j\otimes u_i)/\sqrt2$, $i<j$,
form an eigenbasis of this restriction. Therefore
$\alpha_J=e_0+e_1$. This includes degeneracy; if
$e_0=e_1$ the resulting gap is zero.

Consequently a proved lower bound $\underline\alpha_J$
for $\mathcal C_J$ and any variational upper bound $U_J$
for $e_0$ give
\[
 g_J=e_1-e_0\ge \underline\alpha_J-2U_J,\qquad
 \Delta_N\ge
 \min\{E,\underline\alpha_J-2U_J-R_J(E)\}.
 \tag{V10.2}\label{c18v10:certificate}
\]
The second assertion is exactly V9's gap comparison,
with its hypotheses and its \emph{native} error retained.
No value of $\underline\alpha_J-2U_J$ is asserted positive.

Here is a precise finite algebraic way to prove the first
lower bound without constructing a basis of $\mathcal V_J$.
Use the full retained Haar tensor product for each replica.
Let $C_x^{(r)}$ be the original full gauge Casimir at vertex
$x$ in replica $r=1,2$. Its kernel enforces that replica's
Gauss law. On the full doubled carrier put
$D_J=H_J^{(1)}+H_J^{(2)}$, where here $H_J$ denotes the
full-Haar compression. An exact matrix identity of the form
\[
\begin{split}
 D_J-\underline\alpha_J I
  ={}&\sum_k A_k^*A_k+
       \sum_{x,r}\bigl(C_x^{(r)}B_{x,r}
                         +B_{x,r}^*C_x^{(r)}\bigr)\\
     &+(I+S)B_*+B_*^*(I+S)
\end{split}
 \tag{V10.3}\label{c18v10:sos}
\]
is sufficient. Evaluate its quadratic form on a vector
in all the gauge kernels and in $\ker(I+S)$: every
constraint term vanishes and the remaining sum is
nonnegative. The matrices in the identity need not
preserve those kernels. The identity itself must be
certified exactly or with a rigorous error debit;
a floating-point positive eigenvalue is not such a
certificate. No fixed support size, feasible solution,
or volume-uniform degree bound is claimed.

Two conventions are essential. The replicas obey
\emph{independent} Gauss laws, not merely the diagonal
gauge action. Also the lift
$\mathfrak d(A)=A\otimes I+I\otimes A$ is a Lie map,
not an algebra homomorphism:
\[
 [\mathfrak d(A),\mathfrak d(B)]
       =\mathfrak d([A,B]),\qquad
 \mathfrak d(A)\mathfrak d(B)-\mathfrak d(AB)
       =A\otimes B+B\otimes A.
 \tag{V10.4}\label{c18v10:lift}
\]
These follow by multiplying the four terms.
Thus the complete Wilson plaquette operator is lifted
as a whole; multiplying individually lifted link factors
would introduce mixed-replica interactions.
The scalar $bM I$ also lifts to $2bM I$.

There is a strict accuracy obligation. Define nonnegative
errors $\delta_0=U_J-e_0$ and
$\delta_\Sigma=e_0+e_1-\underline\alpha_J$. Then
\[
 \underline\alpha_J-2U_J
       =g_J-\delta_\Sigma-2\delta_0 .
 \tag{V10.5}\label{c18v10:accuracy}
\]
A fixed error per link in a ground-energy upper bound
does not control this expression at order one as
$M\to\infty$. Common \emph{exact} scalar shifts cancel,
but these two nonnegative certification errors do not.
This is not a proposed solution by subtracting two
uncontrolled extensive energies.

\subsection{The full physical partial-exchange operator}

We now compute a native carrier needed for local
versions of \eqref{c18v10:sos}. All link cutoffs are finite
and identical between replicas. Split the original links
into $D$ and $D^c$, so the full retained Haar carrier is
$\mathcal H_D\otimes\mathcal H_{D^c}$.
The product of all vertex gauge groups acts on these two
factors, including the induced partial actions at the cut.
Its isotypic decompositions can be written
\[
 \mathcal H_D=\bigoplus_\lambda V_\lambda\otimes M_\lambda,
 \qquad
 \mathcal H_{D^c}=\bigoplus_\lambda
                    V_\lambda^*\otimes N_\lambda,
 \qquad d_\lambda=\dim V_\lambda .
 \tag{V10.6}\label{c18v10:decomposition}
\]
Absent representations have zero multiplicity.
Only matched summands contribute to the invariant
subspace. Choose orthonormal bases in the multiplicities.
Schur's lemma identifies the physical carrier with
\[
 \mathcal V_J\simeq
       \bigoplus_\lambda M_\lambda\otimes N_\lambda,\qquad
 |\lambda,\mu,\nu\rangle_{\rm phys}
   =\frac1{\sqrt{d_\lambda}}\sum_{a=1}^{d_\lambda}
           |\lambda,a,\mu\rangle_D
                  \otimes|\lambda,a,\nu\rangle_{D^c}.
 \tag{V10.7}\label{c18v10:embedding}
\]
The second occurrence of the index denotes the
contragredient basis. This is an invariant maximally
entangled vector in the representation indices,
not an assumption about the interacting vacuum.

Let $\mathbb P=P_G\otimes P_G$ be the two independent
physical projections and let $S_D$ exchange just the
$D$ factors. Write
$\widehat S_D=\mathbb P S_D\mathbb P$ on the physical
doubled carrier.

\begin{theorem}[Exchange acts on multiplicities after gauge projection]
\label{c18v10:exchange}
If the two physical sector labels are unequal,
$\widehat S_D$ is zero on their tensor product.
For a common label $\lambda$,
\[
\begin{split}
 \widehat S_D\bigl(
  |\lambda,\mu,\nu\rangle\otimes
  |\lambda,\mu',\nu'\rangle\bigr)
   =\frac1{d_\lambda}
  |\lambda,\mu',\nu\rangle\otimes
  |\lambda,\mu,\nu'\rangle .
\end{split}
 \tag{V10.8}\label{c18v10:swap-block}
\]
Thus on that block it is $d_\lambda^{-1}$ times
exchange of the two $M_\lambda$ factors, with both
$N_\lambda$ factors fixed. It is positive semidefinite
if and only if every matched $M_\lambda$ has dimension
at most one. If a matched multiplicity has dimension
at least two, it has eigenvalue $-1/d_\lambda$,
including on the globally odd physical sector.
\end{theorem}

\begin{proof}
Apply $S_D$ to the tensor of the two sums in
\eqref{c18v10:embedding}. If the labels differ, each
replica now couples inequivalent representation and
contragredient indices, so projection to the singlet
is zero. For equal labels, the first replica pairs
the second copy's index $b$ with the first copy's
exterior index $a$. The second replica pairs $a$
with $b$. Contraction with the normalized singlets
requires $a=b$; the $d_\lambda$ surviving terms have
total coefficient $d_\lambda/d_\lambda^2$.
The multiplicity indices are exactly those in
\eqref{c18v10:swap-block}. This computation proves
the full matrix identity, including arbitrary
exterior multiplicities.

Exchange on $M_\lambda\otimes M_\lambda$ has
eigenvalues $+1$ and, precisely when
$\dim M_\lambda\ge2$, $-1$. Other label blocks are
zero. If $\nu$ is one fixed exterior multiplicity
vector and $\mu\ne\mu'$, the normalized difference
\[
 |\lambda,\mu,\nu\rangle\wedge
 |\lambda,\mu',\nu\rangle
\]
is globally odd and has the claimed negative
eigenvalue. This proves all statements.
\end{proof}

In particular, gauge compression is not multiplicative:
although $S=\prod_e S_{\{e\}}$ before projection,
one cannot replace its physical restriction by
$\prod_e\widehat S_{\{e\}}$.
The representation factors and their contraction
through the boundary must be kept.

\subsection{One and two links cannot carry the negative channel}

For one link $e=(x,y)$ the local gauge representation
is the multiplicity-free Peter--Weyl decomposition
\[
 \mathcal H_e=\bigoplus_{n=0}^{J_e}
             V_n^{(x)}\otimes (V_n^{(y)})^* .
\]
Here $\dim V_n=n+1$, so the representation of the
\emph{two} endpoint groups has dimension $(n+1)^2$.
The theorem gives the full operator identity
\[
 \widehat S_{\{e\}}
       =\sum_{n=0}^{J_e}\frac1{(n+1)^2}
                     \Pi_{n,e}\otimes\Pi_{n,e}\succeq0
       \quad\hbox{on }\mathcal V_J^{\otimes2}.
 \tag{V10.9}\label{c18v10:single-swap}
\]
The original gauge-invariant projections
$\Pi_{n,e}$ are restricted to $\mathcal V_J$.
No assumption that exterior links are trivial was used.

\begin{corollary}[Two-link positivity]
\label{c18v10:two-link}
On the cubic lattice with $N\ge3$, every set $D$
of at most two distinct links satisfies
$\widehat S_D\succeq0$, at every finite electric cutoff.
\end{corollary}

\begin{proof}
For disjoint links, the endpoint groups are independent
and the tensor of their multiplicity-free decompositions
is again multiplicity free. For adjacent links, the
labels $n,m$ are fixed by the two distinct outer
endpoints. At the common vertex,
\[
 V_n^*\otimes V_m
       \simeq\bigoplus_{\substack{|n-m|\le l\le n+m\\
                                 l\equiv n+m\ ({\rm mod}\ 2)}} V_l
 \tag{V10.10}\label{c18v10:fusion}
\]
has multiplicity one. Hence the full vertex-group
label $(n,l,m)$ occurs at most once.
For completeness, the multiplicity assertion follows
by multiplying the characters
$\chi_n(z)=z^n+z^{n-2}+\cdots+z^{-n}$:
their product is the sum of the characters on the
right, as direct counting of each power verifies.
Character independence proves the decomposition.
Distinct links on this lattice cannot have both
endpoints in common. The preceding theorem now applies.
\end{proof}

Thus the failure of one-link parity projectors is
stronger than just a nonzero commutator: their
physical compressions, and those of two-link
exchanges, have no negative spectrum at all.
This is compatible with the archive's isolated
cycle formulas, which are special cases.

\begin{proposition}[A three-link physical negative-exchange channel]
\label{c18v10:three-link}
For $N\ge4$ and link cutoffs at least $2$, there is
a three-link path $D$ and a normalized globally odd
physical vector $\Xi$ with
\[
             \widehat S_D\Xi=-\frac1{16}\Xi .
 \tag{V10.11}\label{c18v10:negative}
\]
Thus three links are the smallest possible number
for a negative physical partial-exchange eigenvalue
under these cutoff assumptions.
\end{proposition}

\begin{proof}
Use the straight path with vertices
$v_i=(i,0,0)$, $i=0,1,2,3$, interpreted modulo $N$.
Take the partial vertex-group label $1$ (spin $1/2$)
at each of these four vertices and label $0$ elsewhere.
The first and third path links must have label $1$.
The middle link can have label $0$ or $2$:
both $V_1\otimes V_0$ and $V_1\otimes V_2$
contain $V_1$ once, by \eqref{c18v10:fusion}.
No other middle label works at these two internal
vertices. Therefore the local multiplicity has
dimension two and $d_\lambda=2^4=16$.

There is a matching exterior representation. In $D^c$
take fundamental labels along the two disjoint paths
\[
\begin{split}
 &(0,0,0),(0,1,0),(1,1,0),(2,1,0),(3,1,0),(3,0,0),\\
 &(1,0,0),(1,0,1),(2,0,1),(2,0,0).
\end{split}
\]
All their internal vertices carry the unique
degree-two singlet; their four endpoints supply
the required four fundamental labels.
All other exterior links are trivial.
These paths have no shared vertices and avoid $D$.
They give a nonzero $N_\lambda$.
Choose one normalized exterior multiplicity vector
$\nu$, and let $\mu_0,\mu_2$ denote the two middle-link
choices. The vector
$|\lambda,\mu_0,\nu\rangle\wedge
 |\lambda,\mu_2,\nu\rangle$
has the claimed properties by
\eqref{c18v10:swap-block}. Minimality is the
two-link corollary.
\end{proof}

This is a concrete original-gauge spin-network channel.
It is not an eigenvector of the interacting $H_J$,
and $1/16$ is an exchange eigenvalue, not an energy.
Discarding the nontrivial boundary representation would
erase this channel. Conversely, finding this channel
does not provide a positive certificate in
\eqref{c18v10:sos}; its complete Wilson couplings and
energy still have to be controlled.

\subsection{A uniform native-vacuum entanglement debit}

This subsection concerns the actual normalized ground
$\Omega$ of the \emph{untruncated} native Hamiltonian,
not an assumed ground of $H_J$. Set
$\beta=b/\kappa>0$, $A_*=3\sqrt2$, and let $\rho_e$
be its one-link reduced density matrix in the full Haar
tensor product. Gauge invariance at the two distinct
endpoints and partial-trace invariance imply
\[
 \rho_e=\bigoplus_{n\ge0}\frac{p_n}{(n+1)^2}I_{(n+1)^2},
 \qquad p_n=\|\Pi_{n,e}\Omega\|^2,\quad
                       p_n\ge0,\quad\sum_n p_n=1.
 \tag{V10.12}\label{c18v10:density}
\]
Indeed the reduced operator commutes with both endpoint
regular actions; Schur's lemma on their inequivalent
irreducible product representations gives the formula.

\begin{theorem}[Basis-independent local rank cost in the native vacuum]
\label{c18v10:purity}
Uniformly in $N\ge3$ and in the link,
\[
 \|\rho_e\|\le \sigma_\beta,\qquad
 \Tr\rho_e^2\le \sigma_\beta,\qquad
 \sigma_\beta=\min\{1,2A_*\beta^{-1/2}\}.
 \tag{V10.13}\label{c18v10:purity-bound}
\]
Consequently every rank-$r$ link projection $Q_r$ obeys
\[
 \Tr(\rho_e Q_r)\le\min\{1,r\sigma_\beta\}.
 \tag{V10.14}\label{c18v10:rank}
\]
In particular retaining probability at least $1-\delta$
requires $r\ge(1-\delta)/\sigma_\beta$.
The min-entropy $-\log\|\rho_e\|$ and second Renyi entropy
$-\log\Tr\rho_e^2$ are at least
\[
        \max\{0,\tfrac12\log\beta-\log(2A_*)\}.
 \tag{V10.15}\label{c18v10:entropy}
\]
These are reduced-state entropies in the stated Haar
factorization, not an assertion of a physical tensor
factorization across the Gauss constraint.
\end{theorem}

\begin{proof}
V6's \eqref{c18v6:two-sided}, for a single-link partial
action, gives
\[
 p_n\le\sum_{j=0}^n p_j
          \le A_*B_n\beta^{-1/2},\qquad
 B_n=\frac1{1-\cos(\pi/(2n+3))}.
\]
There is the useful sharper ratio bound
$B_n\le2(n+1)^2$. To see it, put $t=n+1\ge1$
and $x=\pi/(4t+2)\le\pi/6$.
Concavity of sine gives $\sin x\ge3x/\pi$ on
$[0,\pi/6]$, and hence
$t\sin x\ge3t/(4t+2)\ge1/2$.
Thus $1-\cos(2x)=2\sin^2x\ge1/(2t^2)$,
which proves the ratio bound.
Each eigenvalue in \eqref{c18v10:density} is therefore
at most $2A_*\beta^{-1/2}$. A density matrix also
has norm at most one.
Then $\Tr\rho_e^2\le\|\rho_e\|\Tr\rho_e=\|\rho_e\|$,
and the rank and entropy inequalities follow by
the spectral theorem.
\end{proof}

The exact swap identity for a reduced density matrix
now gives a quantitative consequence for the native
two-copy vacuum:
\[
 \left\langle\Omega\otimes\Omega,
       \sum_e\frac{I-S_{\{e\}}}{2}
                       \Omega\otimes\Omega\right\rangle
       =\frac12\sum_e(1-\Tr\rho_e^2)
       \ge\frac M2(1-\sigma_\beta).
 \tag{V10.16}\label{c18v10:extensive-floor}
\]
The equality is obtained by contracting the two copies
in a product basis, and the lower bound is
\eqref{c18v10:purity-bound}.
The global exchange still fixes $\Omega\otimes\Omega$.
Thus a globally even vacuum has a local exchange-defect
density approaching $1/2$ as $\beta\to\infty$.
These terms are not vacuum-annihilating excitation
counters. The archive f17 V15 proved a strict
fixed-regulator spatial-swap floor for its finite-time
Wilson transfer vacuum; it did not prove this
volume-uniform weak-electric-coupling bound for
the present differential Hamiltonian.

No uniform native gap has been used to control this
density. Nor can it be transferred silently to the
\emph{actual} ground of $H_J$ using V9's small
ground-energy error alone: a small variational energy
does not control the vector without further spectral
information. This keeps the retained and native
vacua distinct.

\subsection{The remaining spectral task and verification}

The transferable mechanism is now specific.
The additive certificate \eqref{c18v10:certificate}
can be sought on the exact physical odd carrier,
using the multiplicity exchanges in
\eqref{c18v10:swap-block}. The smallest available
local negative channel has been exhibited, including
its nontrivial boundary representation. Replacing it
by one-link or two-link parity projections is excluded.
Using the bare local swap sum as a frustration-free
reference is also excluded quantitatively by
\eqref{c18v10:extensive-floor}. Neither exclusion is
a no-go theorem for all correlated block certificates.

The next useful calculation must control the original
Hamiltonian on these coupled multiplicity sectors,
or prove a different uniform physical spectral
inequality. It must keep the boundary fusion maps,
Wilson couplings, and common ground-energy
normalization. Another unpopulated two-copy
reformulation, isolated negative swap eigenvalue,
or finite numeric gap would not close that task.
The physical scale evolution and nontrivial
four-dimensional continuum reconstruction remain
separate unresolved requirements. The full
Yang--Mills mass gap remains unproved.

Exact diagnostics reconstruct the singlet embeddings
and physical exchange matrices, including nontrivial
local and exterior multiplicities. Independent
character multiplication checks the one- and
two-link multiplicities and the two middle-link
choices; lattice incidence checks verify the
exterior paths. Rational doubled-Hamiltonian fixtures
check additive eigenvalue sums, scalar centering,
and the failure of multiplicative lifting.
These are algebra checks, not a numerical
Wilson gap calculation. Scalar substitutions test
the stated purity majorants; the proofs of the
uniform assertions are above. V9's mathematical
diagnostics were replayed before the append.

V10 preserves the complete V9 body except its version
title and replaces only the active source filename.
All forty-seven other note sources are unchanged.
Only \texttt{.tex} sources remain in
\texttt{latex\_notes}; diagnostic and rendering
products are outside it.
% END CYCLE 18 VERSION 10 APPEND

% BEGIN CYCLE 18 VERSION 11 APPEND
\clearpage
\section{V11: Vacuum-cancelled plaquette dynamics and fusion returns}
\label{c18v11:main}

\subsection{Why this calculation, and what is inherited}

V10 identified the physical exchange multiplicities, but its two-energy
certificate had no interacting energy input. This append computes an
entire native excitation band through second order, including every
intermediate representation and the common vacuum subtraction. It is a
calibration of the required cancellation, not a weak-coupling closure.
In particular, the adjoint shared-link return changes the sign of the
resulting hopping. The smallest electric cutoff misses that effect.

The overlap audit found that f17 V88--V89 already prove the frozen
one-link rotor gap and its moving-fibre derivative obstruction. Neither
is repeated. V1--V2 of this cycle already give the general resonant
second-order operator and its connected interaction bounds; V5 already
proves a volume-uniform strong-electric gap in its stated small-coupling
regime. Here the additional calculation is the complete lowest physical
electric band, its three-orientation symbol, and the exact effect of
discarding its adjoint returns. This is not a claim that strong-coupling
glueball expansions are new mathematics.

As a targeted literature check, Dahmen,
\emph{Strong coupling expansion for scattering phases in hamiltonian
lattice field theories II}, Sections 1.1--1.4, especially
\href{https://arxiv.org/pdf/hep-lat/9412080}{equations (1.26)--(1.29)},
provides a two-spatial-dimensional normalization check below.
The four-dimensional target is not replaced by that model.
The publisher abstract of
\href{https://doi.org/10.1016/0550-3213(86)90567-5}
{\emph{Cluster expansion approach to non-abelian lattice gauge theory
in (3+1)D (I). SU(2)}} also records earlier linked strong-coupling mass
series. No continuum extrapolation from that abstract is imported.
The scheduled companion and broad reviews were completed at V10;
their next checkpoints remain V15 and V20.

\subsection{The complete first electric eigenspace}

Use the original three-dimensional cubic spatial torus with $N\ge5$,
$M=3N^3$, and \eqref{c18v1:H}. The restriction on $N$ excludes
short winding loops from the four-link energy band. Put
\[
 |0\rangle=1,\qquad |p\rangle=\chi_1(U_p)=2w_p,\qquad
 P=\sum_p|p\rangle\langle p|,\qquad Q=I-P .
 \tag{V11.1}\label{c18v11:band}
\]
Here $\chi_n$ is the character of highest label $n$, with Casimir
$n(n+2)$; thus $\chi_1$ is the fundamental trace and $\chi_2$ is the
adjoint character. All projections in this subsection act on the
original Gauss-invariant Hilbert space.

\begin{lemma}[No missing low electric states]
\label{c18v11:first}
The vectors $|p\rangle$ are orthonormal, $T P=12\kappa P$, and
\[
 \operatorname{spec}(T)\cap[0,15\kappa)
       =\{0,12\kappa\},\qquad
 \ker T=\mathbb C|0\rangle,\qquad PVP=0 .
 \tag{V11.2}\label{c18v11:electric}
\]
The multiplicity at $12\kappa$ is exactly $M$.
\end{lemma}
\begin{proof}
In a gauge-invariant Peter--Weyl block, a vertex with exactly one
nontrivially labelled incident link cannot carry an invariant tensor.
Thus the graph of occupied links has no degree-one vertex. Any
nonempty such finite graph contains a cycle. Every occupied link
costs at least $3\kappa$. Below $15\kappa$ there are at most four of
them. The cubic torus at $N\ge5$ has no cycles of length below four,
and each four-cycle is the boundary of an elementary plaquette.
Its bivalent invariant tensors force the same label on all four links
and are unique up to scalar. The only allowed nonzero energy below
$15\kappa$ is consequently $4\cdot3\kappa=12\kappa$, with vector
$\chi_1(U_p)$ for each plaquette. Character orthogonality normalizes
it; a link belonging to just one of two distinct plaquettes makes
their inner product zero. The all-trivial block gives $|0\rangle$.

For the last assertion use the commuting link-centre involutions
$U_e\mapsto-U_e$. Each plaquette trace has odd parity exactly on its
four boundary links, and $T$ preserves every such parity. A nonempty
set of at most four distinct plaquettes cannot have even incidence
on every link: choose one of its plaquettes; each of its four links
would require a different other plaquette, because distinct elementary
squares share at most one link. There are at most three others.
The parity of any product of three plaquette traces is therefore
nontrivial, after cancelling repeated pairs. Its Haar integral
vanishes. This proves $PVP=0$.
\end{proof}

The same parity observation will classify all two-step matrix
elements; it does not require a truncation or a choice of exterior
intertwiner. Degeneracies in higher electric eigenspaces are kept.

\subsection{Both shared-link fusion channels}

\begin{lemma}[Exact intermediate weights]
\label{c18v11:weights}
For two distinct plaquettes $p,q$ sharing a link $e$, let
$\Phi_{pq}=\chi_1(U_p)\chi_1(U_q)$. It has the orthogonal decomposition
\[
 \Phi_{pq}=\Phi_{pq}^{(0)}+\Phi_{pq}^{(2)},\qquad
 \begin{array}{c|cc}
   &\Phi_{pq}^{(0)}&\Phi_{pq}^{(2)}\\ \hline
 \text{label on }e&0&2\\
 \text{electric energy}&18\kappa&26\kappa\\
 \text{squared norm}&1/4&3/4
 \end{array}
 \tag{V11.3}\label{c18v11:channels}
\]
The other six links have label one in both terms. If $p,q$ have
no common link, $\|\Phi_{pq}\|=1$ and $T\Phi_{pq}=24\kappa\Phi_{pq}$.
For $p=q$,
\[
 \chi_1(U_p)^2=1+\chi_2(U_p),\qquad
 \|\chi_2(U_p)\|=1,\qquad T\chi_2(U_p)=32\kappa\chi_2(U_p).
 \tag{V11.4}\label{c18v11:same}
\]
\end{lemma}
\begin{proof}
For the shared link, write the two traces as $\Tr(UA)$ and $\Tr(UB)$,
absorbing orientations and inversions into the three-link path
variables $A,B$. Inversion of a complete trace is harmless for
$\mathrm{SU}(2)$. The two remaining paths use disjoint links, so their
holonomies are independent Haar variables for the norm calculation.
The elementary fundamental matrix integral gives
\[
 \int \Tr(UA)\Tr(UB)\,dU=\tfrac12\Tr(AB^{-1}).
 \tag{V11.5}\label{c18v11:Haar}
\]
For example, write $U,A,B$ as unit quaternions. Each trace is twice
a Euclidean inner product after the same orthogonal change of signs,
and $\int U_\alpha U_\beta\,dU=\delta_{\alpha\beta}/4$.
The right side of \eqref{c18v11:Haar} is half the fundamental trace
of the six-link boundary of the union, hence has squared norm $1/4$.
Integration in $A,B$ first shows $\|\Phi_{pq}\|^2=1$.
The tensor product $1\otimes1=0\oplus2$ on the shared link has no
other summand. Removing its constant projection therefore leaves
squared norm $3/4$ in label two. The six other link labels remain one
because the trace product is linear in each of their matrices.
Their energy is $18\kappa$; label two on $e$ adds $8\kappa$.
Orthogonality is the orthogonality of the link Casimir projections.
This argument does not select only one of the possible intermediate
energies. For disjoint link supports the eight labels are all one
and Haar norms factor. Finally, the character identity
$\chi_1^2=\chi_0+\chi_2$ proves \eqref{c18v11:same}.
\end{proof}

Write $p\sim q$ when distinct plaquettes share a link, and let $A_N$
be this graph's adjacency matrix on the basis \eqref{c18v11:band}.
Every vertex has degree $12$: each of four links lies in three other
plaquettes and none is counted twice. This graph is connected.
Translations within a plaquette plane are adjacency steps; a change
of orientation, followed by another such change, also translates a
plaquette in its normal direction. These operations reach every
plaquette and all three orientations. Set
\[
                      \mathcal L_N^{\rm pl}=12I-A_N .
 \tag{V11.6}\label{c18v11:laplacian}
\]
This is only a finite plaquette graph Laplacian, not the continuum
Yang--Mills Hamiltonian.

\subsection{The full connected coefficient after vacuum subtraction}

\begin{theorem}[Complete native second-order band]
\label{c18v11:coefficient}
The exact second-order coefficient in the physical band is
\[
 \begin{split}
 B_N&:=PVQ(12\kappa-T)^{-1}QVP\\
 &=-\frac{Mb^2}{48\kappa}I
       -\frac{b^2}{\kappa}
                \left(\frac{11}{1680}I+\frac1{336}A_N\right).
 \end{split}
 \tag{V11.7}\label{c18v11:B}
\]
Consequently subtracting the vacuum coefficient
$-Mb^2/(48\kappa)$ gives
\[
 \boxed{\quad
 B_N^{\rm c}=\frac{b^2}{\kappa}
       \left(-\frac{71}{1680}I+\frac1{336}\mathcal L_N^{\rm pl}\right).
 \quad}
 \tag{V11.8}\label{c18v11:centered}
\]
This identity has constants independent of $N$. It is a coefficient
identity, not a remainder estimate at fixed nonzero $b/\kappa$.
\end{theorem}
\begin{proof}
The inverse exists on $Q$ by Lemma~\ref{c18v11:first}. It is signed:
the vacuum denominator is positive, whereas every intermediate
energy in \eqref{c18v11:channels} and the adjoint term of
\eqref{c18v11:same} is above $12\kappa$.

For a diagonal matrix element, parity forces the two perturbing
plaquettes to be identical. Since $V_q=-b\chi_1(U_q)/2$, the
contributions, in units $b^2/\kappa$, are
\[
 \begin{array}{c|c|c}
  q&\text{contribution}&\text{count}\\ \hline
  q=p&\frac14(\frac1{12}-\frac1{20})=\frac1{120}&1\\[2pt]
  q\sim p&-\frac14(\frac{1/4}{6}+\frac{3/4}{14})=-\frac1{42}&12\\[2pt]
  q\ne p,\ q\not\sim p&-\frac1{48}&M-13 .
 \end{array}
 \tag{V11.9}\label{c18v11:diagonal}
\]
The sum is $-M/48-11/1680$.

For $p\ne r$, a term
$\langle r|V_s Q(12\kappa-T)^{-1}QV_q|p\rangle$
can survive only when the mod-two boundaries of $p,q,r,s$ sum to
zero. The four-plaquette observation in the proof of
Lemma~\ref{c18v11:first} forces every plaquette in the multiset to
occur an even number of times. The only choices are
$(q,s)=(p,r)$ and $(q,s)=(r,p)$.

The first is the vacuum transport, of value $b^2/(48\kappa)$:
the adjoint characters on two different plaquettes are orthogonal.
The second passes through the identical vector $\Phi_{pr}$ on
both sides of the reduced resolvent. If the plaquettes have no
common link, it is $-b^2/(48\kappa)$, cancelling the vacuum
transport exactly. If they share a link, it is $-b^2/(42\kappa)$.
Their sum is then
\[
               \frac{b^2}{\kappa}
                    \left(\frac1{48}-\frac1{42}\right)
                    =-\frac{b^2}{336\kappa}.
 \tag{V11.10}\label{c18v11:hop}
\]
There are no other matrix elements. The vacuum second-order
coefficient is \eqref{c18v1:vacuum-coefficient}, so subtraction
gives \eqref{c18v11:centered}. This also directly checks the
connected-support identity \eqref{c18v1:Zblock}.
\end{proof}

The cancellation in \eqref{c18v11:hop} is between physical intermediate
states, not between separately chosen per-volume approximations.
In particular, omitting the vacuum channel would leave nonzero
hopping between arbitrarily distant plaquettes at this order.

\subsection{Three orientations, momentum, and a published normalization check}

Label each plaquette by its centre and its normal direction
$a\in\{1,2,3\}$. In the centre-based Fourier basis, for momenta
$k_i=2\pi m_i/N$, the symbol of $A_N$ is the real symmetric matrix
\[
 \begin{split}
 A(k)_{aa}&=2\sum_{i\ne a}\cos k_i,\\
 A(k)_{ab}&=4\cos(k_a/2)\cos(k_b/2),\qquad a\ne b .
 \end{split}
 \tag{V11.11}\label{c18v11:symbol}
\]
The half-link phases refer to the stated basis; changing a momentum
representative by $2\pi$ merely changes orientation basis phases.

\begin{proposition}[Coherent band bottom and small-momentum coefficient]
\label{c18v11:dispersion}
The minimum of $B_N^{\rm c}$ is $-71b^2/(1680\kappa)$ and is simple
when $b>0$, with constant plaquette amplitudes. The three coefficient
branches are
\[
 12\kappa-\frac{b^2}{\kappa}
       \left(\frac{11}{1680}+\frac{\lambda_j(A(k))}{336}\right).
 \tag{V11.12}\label{c18v11:branches}
\]
For the branch issuing from $\lambda_{\max}(A(0))=12$ the continuous
trigonometric symbol has the expansion
\[
 12\kappa-\frac{71b^2}{1680\kappa}
                  +\frac{b^2}{252\kappa}|k|^2
                  +O\!\left(\frac{b^2}{\kappa}|k|^4\right).
 \tag{V11.13}\label{c18v11:small-k}
\]
These are eigenvalues of the displayed coefficient matrix. They are
not exact dispersions at arbitrary coupling.
\end{proposition}
\begin{proof}
A plaquette has four same-orientation neighbours, displaced within
its plane. For either other orientation the four neighbouring centres
differ by $\pm e_a/2\pm e_b/2$. Summing their Fourier phases proves
\eqref{c18v11:symbol}. On the full graph,
\[
 \langle c,\mathcal L_N^{\rm pl}c\rangle
                     =\sum_{\{p,q\}:p\sim q}|c_p-c_q|^2 .
\]
Connectivity makes its kernel exactly the constant vectors, proving
the minimum and simplicity. At $k=0$, $A(0)$ has every entry equal
to four, hence eigenvalues $12,0,0$. Its top normalized eigenvector
is $(1,1,1)/\sqrt3$. The linear Taylor term vanishes. The diagonal
quadratic changes are $-\sum_{i\ne a}k_i^2$, and the off-diagonal
ones are $-(k_a^2+k_b^2)/2$. Their sum divided by three is
$-4|k|^2/3$. The isolated top eigenvalue therefore equals
$12-4|k|^2/3+O(|k|^4)$: its eigenvector correction affects the
eigenvalue only at fourth order. Substitution proves
\eqref{c18v11:small-k}.
\end{proof}

For a check only, repeat the same local calculation on a
two-dimensional square spatial torus. Its plaquette adjacency
degree is four rather than twelve. The centered diagonal is then
$29b^2/(1680\kappa)$, the hopping is still $-b^2/(336\kappa)$,
and the band-bottom correction is $3b^2/(560\kappa)$.
In Dahmen's normalization set $t=b/(8\kappa)$ and divide energies
by $4\kappa$, with the common magnetic scalar suppressed. The
diagonal, hopping and bottom corrections become respectively
$29t^2/105$, $-t^2/21$, and $3t^2/35$, agreeing with the cited
equations. This imports a normalization test, not a
three-dimensional remainder or a continuum claim.

\subsection{A native cutoff changes the answer, not just the error bar}

Let $P_J^{\rm all}=\prod_e\mathbf1_{[0,J]}(n_e)$ and use exactly
the native compression $H_J=P_J^{\rm all}HP_J^{\rm all}$ on its
physical carrier, as in V9. This retains the original coefficients
and scalar $bM$.

\begin{theorem}[Second-order cutoff exactness and the $J=1$ sign change]
\label{c18v11:cutoff}
For every integer $J\ge2$, both the vacuum coefficient and the entire
$12\kappa$ band coefficient of $H_J$ equal the untruncated ones.
For $J=1$ the vacuum coefficient is still unchanged, but the centered
band coefficient is instead
\[
            B_{N,J=1}^{\rm c}
                 =\frac{b^2}{\kappa}
                         \left(\frac16I+\frac1{96}A_N\right).
 \tag{V11.14}\label{c18v11:J1}
\]
For even $N\ge6$ its minimum is $b^2/(8\kappa)$, in contrast to
the negative full coefficient $-71b^2/(1680\kappa)$.
This comparison is between two specified physical Hamiltonians;
it is not an abstract matrix countermodel.
\end{theorem}
\begin{proof}
Applying one fundamental plaquette multiplication to $|0\rangle$
or $|p\rangle$ creates only link labels $0,1,2$. Thus $J\ge2$
contains every intermediate vector contributing at this order;
its electric inverse has the same eigenvalues on them.

For $J=1$, the adjoint term in \eqref{c18v11:same} disappears,
and the $26\kappa$ term in \eqref{c18v11:channels} disappears.
The same-plaquette return becomes $1/48$, the shared-link return
$-1/96$, and the remote return remains $-1/48$, in units
$b^2/\kappa$. Adding the vacuum subtraction to the diagonal gives
\[
        \frac1{48}-\frac{12}{96}
                          -\frac{M-13}{48}+\frac M{48}=\frac16 .
\]
Off the diagonal the shared-link entry is
$1/48-1/96=1/96$, while the remote cancellation remains exact.
This proves \eqref{c18v11:J1}; the vacuum's first intermediate
states have label one and are unaffected.

Let $C$ be the unsigned link--plaquette incidence matrix,
$C_{ep}=1$ when $e\in p$ and zero otherwise. Since a plaquette
has four links and distinct plaquettes share at most one,
$C^*C=4I+A_N$, so $A_N\ge-4I$. For even $N$, the allowed
momentum $(\pi,\pi,\pi)$ has $A(k)=-4I$ by
\eqref{c18v11:symbol}, attaining this lower bound. Consequently
the minimum in \eqref{c18v11:J1} is
$(1/6-4/96)b^2/\kappa=b^2/(8\kappa)$.
\end{proof}

This does not say that every finite cutoff fails. It shows exactly
why a representation retained in a virtual return cannot be omitted
without recomputing the dynamics. V9's cofinal cutoff estimates and
their coupling-dependent $J$ remain the relevant all-coupling
controls; $J=2$ is exact only for the coefficients just computed.

\subsection{A genuine finite-regulator spectral statement with its loss exposed}

A coefficient alone is not an eigenvalue bound. The following
elementary estimate supplies a finite-volume meaning while recording
the volume loss explicitly. It does not improve V5's
volume-uniform strong-electric gap.

\begin{lemma}[Isolated-cluster Schur error]
\label{c18v11:Schur}
Let a self-adjoint, bounded-below operator $T_*$ have compact
resolvent and an isolated eigenvalue $E$ of finite multiplicity,
with spectral projection $P_*$. Suppose the distance to its
complementary spectrum is at least $d>0$. Let $V_*$ be bounded
self-adjoint, $\|V_*\|\le v\le d/4$, and $P_*V_*P_*=0$.
The eigenvalues of $T_*+V_*$ in the continued $E$ cluster,
in increasing order, differ from those of
\[
 E P_*+P_*V_*Q_*(E-Q_*T_*Q_*)^{-1}Q_*V_*P_*,
       \qquad Q_*=I-P_*,
 \tag{V11.15}\label{c18v11:Schur-coefficient}
\]
by at most
\[
                        r(d,v)=\frac{2v^3}{d(d-2v)} .
 \tag{V11.16}\label{c18v11:Schur-error}
\]
For $v=0$ the error is zero.
\end{lemma}
\begin{proof}
Bounded perturbation and min--max place the cluster in
$[E-v,E+v]$, separated from all other eigenvalues, with unchanged
multiplicity. Throughout this interval set
$D=Q_*(T_*+V_*)Q_*$. Its distance from $\lambda$ is at least
$d-2v$, by the same bounded-perturbation estimate on $Q_*$.
The eigenvalue equation reduces, by solving its $Q_*$ component,
to the finite-dimensional Schur matrix
\[
 F(\lambda)=E P_*+
 P_*V_*Q_*(\lambda-D)^{-1}Q_*V_*P_*-\lambda P_* .
\]
The resolvent identity gives
\[
 \left\|P_*V_*Q_*
  \left[(\lambda-D)^{-1}-(E-Q_*T_*Q_*)^{-1}\right]
                     Q_*V_*P_*\right\|
 \le\frac{v^2(|\lambda-E|+v)}{d(d-2v)}
 \le r(d,v).
\]
All inverses here are actual bounded inverses; none removes an
equal-energy state. Moreover
$F'(\lambda)=-P_*-P_*V_*Q_*(\lambda-D)^{-2}Q_*V_*P_*
\preceq-P_*$. Thus its ordered eigenvalue curves are strictly
decreasing. Comparing them with the ordered eigenvalues of
\eqref{c18v11:Schur-coefficient} minus $\lambda$ gives the claimed
ordered root bounds. To see that the comparison brackets lie in
the chosen interval, the second-order matrix has norm at most
$v^2/d$ and $v^2/d+r(d,v)\le v/2$ when $v\le d/4$.
Equivalently the usual block Gaussian elimination gives the same
eigenvalue count; its complementary block has no spectrum in this
interval. Bounded off-diagonal blocks preserve the stated operator
domains, so no finite-dimensional truncation of $T_*$ is needed.
\end{proof}

\begin{corollary}[Actual gap expansion at fixed volume]
\label{c18v11:actual}
For $N\ge5$ and $v=Mb\le3\kappa/4$, the first $M$ physical
excitation eigenvalues lie in the continued plaquette band.
The actual physical gap $\Delta_N$ satisfies
\[
 \left|\Delta_N-
       \left(12\kappa-\frac{71b^2}{1680\kappa}\right)\right|
       \le r(3\kappa,Mb)+r(12\kappa,Mb)
       \le \frac{17(Mb)^3}{36\kappa^2}.
 \tag{V11.17}\label{c18v11:actual-gap}
\]
The ground energy obeys
\[
       \left|E_0-\left(bM-\frac{Mb^2}{48\kappa}\right)\right|
                              \le r(12\kappa,Mb).
 \tag{V11.18}\label{c18v11:actual-vacuum}
\]
The same estimates hold for the native $H_J$, $J\ge2$.
For $J=1$ and even $N\ge6$, replace the central gap approximation
in \eqref{c18v11:actual-gap} by $12\kappa+b^2/(8\kappa)$.
\end{corollary}
\begin{proof}
Apply Lemma~\ref{c18v11:Schur} to $H-bM=T+V$. The perturbation
has norm at most $Mb$. Lemma~\ref{c18v11:first} supplies distance
$3\kappa$ for the plaquette cluster and $12\kappa$ for the vacuum.
Their second-order operators were computed above. The spectral
separation ensures that no other physical eigenvalue enters below
the plaquette cluster. Subtract the two estimates; the magnetic
scalar cancels exactly. Since $Mb\le3\kappa/4$,
$r(3\kappa,Mb)\le4(Mb)^3/(9\kappa^2)$ and
$r(12\kappa,Mb)\le(Mb)^3/(36\kappa^2)$, which sum to the
displayed bound. Compression has perturbation norm at most $Mb$
and the same unperturbed low eigenspaces and separations. The
cutoff assertions therefore follow from
Theorem~\ref{c18v11:cutoff}.
\end{proof}

The last theorem does not justify taking $N\to\infty$ at fixed
$b>0$: both its smallness hypothesis and its remainder explicitly
lose volume. Nor may \eqref{c18v11:actual-gap} be extrapolated to
$\beta=b/\kappa\to\infty$. The order-one coefficient cancellation
in \eqref{c18v11:centered} is established, but cancellation of all
higher, vacuum-dressed errors has not been established here.
The physical weak-coupling scale trajectory and the nontrivial
four-dimensional continuum construction remain missing.

\subsection{Verification, direction, and scope}

Exact diagnostics integrate the quaternion polynomials underlying
both fusion weights, check the single-cycle character identity,
and recompute every rational energy denominator. Native finite
lattice checks enumerate all distinct two-plaquette parity masks,
verify the degree and connectivity, and compare the centre-based
Fourier action with \eqref{c18v11:symbol}. The quadratic momentum
tensor and two-dimensional normalization conversion are checked
symbolically. Finite Hermitian fixtures test the Schur error bound;
they are explicitly not native YM spectra and their floating-point
checks are not interval certificates. V10's mathematical diagnostics
were replayed before the append. These checks supplement the proofs,
not replace them with numerical evidence.

The useful change of direction is now precise: the next energy
estimate should preserve the cancellation between vacuum transport
and disconnected two-plaquette returns while controlling the
\emph{connected}, vacuum-dressed remainder. It must also retain
the shared-link adjoint channel. Increasing the formal series order,
or repeating a frozen-rotor bound, would not by itself pay that
obligation. A volume-uniform dressed-band estimate would strengthen
the spectral input; it would still need a justified passage to the
physical weak-coupling regime and continuum, not a formal
extrapolation of the present strong-electric coefficients.

V10's entire body is preserved apart from the version title.
This is the single active V11 source; all archives and companions
are unchanged. Only \texttt{.tex} files are retained in
\texttt{latex\_notes}. Build products and diagnostic records remain
outside it. The full interacting four-dimensional Yang--Mills
mass gap remains unproved.
% END CYCLE 18 VERSION 11 APPEND
% BEGIN CYCLE 18 VERSION 12 APPEND
\clearpage
\section{V12: A volume-uniform dressed plaquette band}
\label{c18v12:context}

V11 evaluated the complete lowest physical electric band, but
its spectral remainder used $\|V\|\le Mb$. That estimate could
not survive increasing volume at fixed coupling. This append
removes that particular loss by combining V5's convergent
native creation equation with a contour around the entire
plaquette band. It proves a uniform analytic band and an
exponentially localized cubic remainder, including the actual
vacuum subtraction. It does not construct the continuum theory
or extend the estimate to $\beta=b/\kappa\to\infty$.

The construction is a native application of a known perturbative
architecture, not a claim that weak-perturbation quasiparticles
are new. Yarotsky's
\emph{Quasi-particles in weak perturbations of non-interacting
quantum lattice systems},
\href{https://arxiv.org/abs/math-ph/0411042}{math-ph/0411042},
Sections 1--3, uses creation coordinates, a dressed excitation
operator and a resolvent expansion. Its isolated on-site
one-particle hypothesis does not directly describe our composite,
Gauss-invariant, three-orientation plaquette band. The argument
below supplies the physical carrier, constants and coefficient
identification explicitly. We do not import its scattering
construction or assert a uniform physical Gram matrix.

\subsection{The exact physical dressed operator}

Throughout this append $N\ge5$, $M=3N^3$, and
\[
 H(b)=T+bM+bV_0,\qquad V_0=-\sum_p w_p,\qquad
 E=12\kappa,\qquad Q_0=\mathbf1-|\omega\rangle\langle\omega|.
 \tag{V12.1}\label{c18v12:H}
\]
Here $\omega=1$ and $Q_0$ removes only the electric vacuum.
Let $P$ be the complete physical $E$-eigenspace of $T$.
By Lemma~\ref{c18v11:first},
\[
 P\mathcal H_{\rm phys}
   =\operatorname{span}\{|p\rangle=\chi_1(U_p):p\in\mathcal P_N\},
 \quad \dim P=M,\quad
 \operatorname{spec}(T|_{Q_0\mathcal H_{\rm phys}})
                      \subset\{E\}\cup[15\kappa,\infty).
 \tag{V12.2}\label{c18v12:electric}
\]
The plaquette vectors in this display are orthonormal.

Use the family space of \eqref{c18v5:norm}, restricted to the
closed subspace $\mathcal B_a^G$ in which each
$\iota_I c_{Z,I}$ is invariant under the original full gauge
group. There is no new partial-Gauss constraint. Exact excitation
projections, $T_I$, and the primitive Wilson multiplication
commute with the gauge action. Thus the family map
$\mathcal F_b$ and its derivative preserve this subspace.
Every nonzero component has physical electric energy at least
$E$; its spectral values belong to the set in
\eqref{c18v12:electric}.

Write $\mathcal Ld=\sum_{Z,I}\iota_I d_{Z,I}$ and let
$T^\#$ denote componentwise $T_I$. The lift is onto
$X=Q_0\operatorname{Dom}T\cap\mathcal H_{\rm phys}$:
the whole connected plaquette set supplies a cover for each
exact excitation component. Its kernel is closed, and
\[
 |v|_{a,E}=\inf_{\mathcal Ld=v}\|d\|_{a,E}
 \tag{V12.3}\label{c18v12:quotient}
\]
is a Banach quotient norm. As in V5, it is equivalent at each
finite $N$ to $\|Tv\|$, with constants which may depend on $N$.
We never insert those constants in a uniform estimate.

\begin{lemma}[A fixed complex disk and the native linearization]
\label{c18v12:linearization}
Fix $a\ge0$ and put
\[
 \rho_a=\frac{\kappa e^{-a}}{2688},\qquad
 \Delta=3\kappa,\qquad x_R=32\rho_a e^a.
 \tag{V12.4}\label{c18v12:radius}
\]
For $|b|\le\rho_a$, the V5 map has a unique fixed point $c_b$
in the fixed ball of radius $x_R$. It is holomorphic in the open
disk and continuous on its closure. For
$D_b=D\mathcal F_b(c_b)$,
\[
 \|D_b\|_{\mathcal B_a^G\to\mathcal B_a^G}
 \le\ell_b:=\frac{224|b|e^a}{3\kappa}
 \le\frac1{36}.
 \tag{V12.5}\label{c18v12:derivative}
\]
Define
\[
 S_b=e^{C(c_b)},\qquad
 W_b=S_b^{-1}bV_0S_b,\qquad
 \mathcal E_b=\langle\omega,W_b\omega\rangle .
\]
Then $(T+bV_0)S_b\omega=\mathcal E_bS_b\omega$. The derivative
descends to $J_b$ on $X$, and the actual reduced excitation
operator is
\[
 \begin{split}
 \mathcal L D_b d
    &=J_b v=-T^{-1}Q_0[W_b,\widehat v]\omega,
                      \qquad v=\mathcal Ld,\\
 A_b&=Q_0S_b^{-1}(T+bV_0-\mathcal E_b)S_bQ_0
                  =T(\mathbf1-J_b),\qquad
 \operatorname{Dom}A_b=X .
 \end{split}
 \tag{V12.6}\label{c18v12:A}
\]
For real $0\le b\le\rho_a$, $E_0=bM+\mathcal E_b$ is the
actual ground energy.
\end{lemma}
\begin{proof}
The multilinear estimates of Lemma~\ref{c18v5:majorant}
hold for complex $b$ with $b$ replaced by $|b|$. In the fixed
ball,
\[
 \frac{\rho_a e^a}{\Delta}=\frac1{8064},\qquad
 \frac{8x_R}{\Delta}=\frac2{63}<\frac14.
\]
The bounds $P_8(t)\le17/12$ and $P_8'(t)\le7/4$ from V5
give invariance of that ball and
$\|D\mathcal F_b(c)\|\le224|b|e^a/(3\kappa)\le1/36$.
Iterating from zero gives uniformly convergent holomorphic
Banach-valued polynomials, hence the asserted holomorphic limit.
The exact finite creation algebra proves the eigenpair identity.
Differentiating $W_b$ in a creation direction gives
$[W_b,\widehat v]$, and the lift identity
\eqref{c18v5:derivative} proves descent to the quotient.

For completeness, the conjugated centered operator annihilates
$\omega$. Commute it through $\widehat v$ and use
$[T,\widehat v]\omega=Tv$ and the commutation of pure creations.
Its $Q_0$ component is
$Tv+Q_0[W_b,\widehat v]\omega=T(\mathbf1-J_b)v$.
The domain assertions follow from the bounded finite-volume
creation commutators in Lemma~\ref{c18v5:algebra}.
For real nonnegative $b$, Theorem~\ref{c18v5:actual-gap}
identifies the eigenpair with the actual simple ground.
\end{proof}

Relative to $\mathbb C\omega\oplus Q_0\mathcal H_{\rm phys}$,
the conjugated centered Hamiltonian has block form
$\left(\begin{smallmatrix}0&*\\0&A_b\end{smallmatrix}\right)$.
Consequently $A_b$ describes the genuine nonvacuum spectrum,
not a trial-vacuum compression. Neither $S_b$ nor the induced
coordinate change is claimed unitary or bounded uniformly in
the global Hilbert norm.

\subsection{A genuine contour inverse with no volume factor}

On $\mathcal B_a^G$, use the closed operator
\[
 A_b^\#=T^\#(\mathbf1-D_b),\qquad
 \operatorname{Dom}A_b^\#
       =(\mathbf1-D_b)^{-1}\operatorname{Dom}T^\# .
 \tag{V12.7}\label{c18v12:family-A}
\]
Here $\operatorname{Dom}T^\#$ means that the componentwise
image again has finite family energy norm. At finite $N$ this
is a dense domain: $(\mathbf1-D_b)^{-1}$ is a bounded
isomorphism and $T^\#$ is a closed diagonal operator.
Its inverse at zero is compact, since $(T^\#)^{-1}$ is compact
on the finitely many graph-norm components.

\begin{lemma}[Uniform Riesz projection]
\label{c18v12:contour}
Set
\[
 r=\frac{4\kappa}{3},\qquad
 \Gamma=\{z:|z-E|=r\},\qquad q_b=9\ell_b\le\frac14,
 \qquad G_z=(\mathbf1-z(T^\#)^{-1})^{-1}.
 \tag{V12.8}\label{c18v12:contour-data}
\]
For $z\in\Gamma$,
\[
 \|G_z\|\le9,\qquad
 \|(z-T^\#)^{-1}\|\le r^{-1}.
\]
The following is the actual resolvent, with its indicated order
of factors:
\[
 R_b(z):=(z-A_b^\#)^{-1}
       =-(\mathbf1-G_zD_b)^{-1}G_z(T^\#)^{-1},
 \qquad
 \|R_b(z)-R_0(z)\|
       \le\frac{q_b}{(1-q_b)r}.
 \tag{V12.9}\label{c18v12:resolvent}
\]
For the positively oriented contour define
\[
 P_b^\#=\frac1{2\pi i}\int_\Gamma R_b(z)\,dz,\qquad
 P_0^\#=\mathbf1_{\{E\}}(T^\#).
\]
Then
\[
 \|P_b^\#-P_0^\#\|\le p_b:=\frac{q_b}{1-q_b}\le\frac13.
 \tag{V12.10}\label{c18v12:projection}
\]
All these operators descend through $\mathcal L$ to the
corresponding resolvents and projections of the physical $A_b$.
\end{lemma}
\begin{proof}
For a component spectral value $t=E$, one has
$t/|t-z|=E/r=9$. For $t\ge15\kappa$,
\[
 \frac{t}{|t-z|}
 \le\frac{t}{t-E-r}
 \le\frac{15\kappa}{3\kappa-r}=9.
\]
The distance of every component spectrum from $\Gamma$ is at
least $r$; the spectral theorem gives the two multiplier bounds
in the family norm. Factor the pencil exactly:
\[
 \mathbf1-D_b-z(T^\#)^{-1}
   =(\mathbf1-z(T^\#)^{-1})(\mathbf1-G_zD_b).
\]
The second factor is invertible by a Neumann series with ratio
at most $q_b$. Solving $(z-A_b^\#)v=f$ gives
\eqref{c18v12:resolvent}. In particular the solution has the
required domain, since
$(\mathbf1-D_b)v=z(T^\#)^{-1}v-(T^\#)^{-1}f$.
Subtract the free resolvent and use
$\|G_z(T^\#)^{-1}\|=\|(z-T^\#)^{-1}\|\le1/r$.
Contour integration proves \eqref{c18v12:projection}.

The operators $D_b$ and the spectral multipliers preserve
$\ker\mathcal L$ and intertwine with $J_b$ and the corresponding
multipliers of $T$. Their inverse series therefore descend.
To check that this is a Hilbert-space resolvent, not only a
formal quotient inverse, let $f\in Q_0\mathcal H_{\rm phys}$.
Then $T^{-1}f\in X$. The descended pencil inverse supplies the
unique $v\in X$ solving $(z-A_b)v=f$. Fixed-volume norm
equivalence bounds this solution as a Hilbert-space operator;
no uniform equivalence constant is needed or asserted.
For inputs in $X$, this inverse is exactly the descended
$R_b(z)$, so the contour projections also intertwine.
\end{proof}

The inverse never divides by an energy difference within the
plaquette band. All $M$ resonant states are retained together.
Compactness and the contour give finite-rank projections.
The descended physical projection has rank $M$, by norm-continuous
continuation from $P$ along the disk. Extra copies caused by
different cover labels upstairs are not physical states.

\subsection{Analytic band coordinates and a rooted support bound}

Let
\[
 \begin{split}
 U_b^\#&=P_b^\#P_0^\#
       +(\mathbf1-P_b^\#)(\mathbf1-P_0^\#)\\
       &=\mathbf1+(P_b^\#-P_0^\#)(2P_0^\#-\mathbf1),\\
 F_b^\#&=P_0^\#(U_b^\#)^{-1}A_b^\#P_b^\#P_0^\# .
 \end{split}
 \tag{V12.11}\label{c18v12:F}
\]
The reflection $2P_0^\#-\mathbf1$ is an isometry in the family
norm. Thus $\|U_b^\#-\mathbf1\|\le p_b<1$,
$\|(U_b^\#)^{-1}\|\le(1-p_b)^{-1}$, and
$U_b^\#P_0^\#=P_b^\#U_b^\#=P_b^\#P_0^\#$.
Consequently \eqref{c18v12:F} is exactly the band restriction
of $(U_b^\#)^{-1}A_b^\#U_b^\#$, not a truncated series.
Let $F_N(b)$ be its descended matrix on the orthonormal
physical plaquette basis.

\begin{theorem}[Uniform analytic and spatially localized band]
\label{c18v12:uniform-band}
For $|b|<\rho_a$, $F_N(b)$ is holomorphic,
$F_N(0)=E\mathbf1$, and
\[
 \|F_b^\#-EP_0^\#\|
       \le\frac{r q_b}{1-2q_b}\le\frac{2\kappa}{3}.
 \tag{V12.12}\label{c18v12:family-bound}
\]
Its spectrum is the complete continued physical plaquette
excitation band. For real $0\le b<\rho_a$, this band consists
of the first $M$ nonzero eigenvalues of the original $H-E_0$.

Let $d(p,q)$ be graph distance between plaquettes, where sharing
a link defines adjacency. Uniformly in $N$ and $|b|\le\rho_a$,
\[
 \sup_p\sum_q e^{a d(p,q)}
            |(F_N(b)-E\mathbf1)_{qp}|
       \le\frac{2\kappa e^a}{3}.
 \tag{V12.13}\label{c18v12:columns}
\]
If $F_N(b)=E\mathbf1+\sum_{n\ge1}b^n F_{N,n}$, then
\[
 \sup_p\sum_q e^{a d(p,q)}|F_{N,n}(q,p)|
       \le\frac{2\kappa e^a}{3}\rho_a^{-n}
       \qquad(n\ge1).
 \tag{V12.14}\label{c18v12:coeff-bound}
\]
No assertion of Hermiticity in these band coordinates is needed.
\end{theorem}
\begin{proof}
Functional calculus and the free identity
$(T^\#-E)P_0^\#=0$ give
\[
 (A_b^\#-E)P_b^\#
    =\frac1{2\pi i}\int_\Gamma
                     (z-E)(R_b(z)-R_0(z))\,dz .
\]
Its norm is at most $r q_b/(1-q_b)$.
Multiplication by $P_0^\#(U_b^\#)^{-1}$ and $P_0^\#$ proves
\eqref{c18v12:family-bound}; use
$1-p_b=(1-2q_b)/(1-q_b)$.
All inverse series and contour integrals are holomorphic on the
disk. They intertwine through the lift, and $U_b$ identifies the
physical $P$ with the physical continued band. At real $b$ the
centered spectrum is real. Along $[0,b]$ no eigenvalue crosses
either real endpoint $E-r$ or $E+r$, since these belong to the
resolvent contour throughout. At zero there are exactly $M$
eigenvalues in that interval, no nonvacuum eigenvalues below it,
and all remaining eigenvalues above it. Compact resolvent,
continuous ordered eigenvalues at fixed $N$, and the simple
continued ground prove the assertion about the first $M$.

Here is the spatial assertion, including the cover bookkeeping.
Every output cover of $D_b$ contains its input cover:
in the differentiated commutator expansion it is the connected
union of that cover, the primitive plaquette, and the other
creation covers. Spectral multipliers keep cover labels.
Neumann series, contour integrals, and norm limits therefore
preserve the closed subspace of families all of whose covers
contain a given root plaquette $p$. The same is true of
$F_b^\#-EP_0^\#$.

Represent $|p\rangle$ by the one family component
$(Z,I)=(\{p\},E(p))$. Its norm is $Ee^a$. The output
$y=(F_b^\#-EP_0^\#)d$ lies in $\operatorname{Ran}P_0^\#$.
By \eqref{c18v12:electric}, its nonzero components have the
form $\alpha_{Zq}|q\rangle$, with $I=E(q)$ and $E(q)\subseteq
E(Z)$. All those covers contain $p$. Although $q$ need not
itself belong to $Z$, choose an edge of $q$ and a member of $Z$
containing it. That member is either $q$ or adjacent to $q$.
Connectedness of $Z$ thus gives
\[
                 d(p,q)\le |Z|.
 \tag{V12.15}\label{c18v12:cover-distance}
\]
Fix an edge of $p$. It belongs to every output cover, so the
single anchored sum in the family norm counts every component:
\[
 E\sum_{Z,q}e^{a|Z|}|\alpha_{Zq}|
                  \le\|y\|_{a,E}
                  \le\frac{2\kappa}{3}Ee^a.
\]
Lift, take absolute values, and use
\eqref{c18v12:cover-distance} to obtain
\eqref{c18v12:columns}. Applying the Banach-valued Cauchy
estimate to the same rooted family on circles approaching
$\rho_a$ gives \eqref{c18v12:coeff-bound}.
\end{proof}

\subsection{Identifying the exact second-order coefficient}

\begin{proposition}[The V11 coefficient belongs to this exact band]
\label{c18v12:coefficient}
For the above analytic matrix,
\[
 F_{N,1}=0,\qquad
 F_{N,2}=\frac1\kappa
       \left(-\frac{71}{1680}\mathbf1+
                        \frac{\mathcal L_N^{\rm pl}}{336}\right),
 \qquad \mathcal L_N^{\rm pl}=12\mathbf1-A_N .
 \tag{V12.16}\label{c18v12:coefficient-formula}
\]
Here $A_N$ is precisely the twelve-regular plaquette adjacency
matrix of V11.
\end{proposition}
\begin{proof}
We justify coefficient matching even though $S_b$ and $U_b$
are not unitary. At fixed $N$, let $\Pi_b^{\rm orig}$ be the
original Hilbert-space spectral projection onto the continued
band of $T+bV_0-\mathcal E_b$. For sufficiently small $b$,
\[
       B_b=\Pi_b^{\rm orig} S_b U_b|_{P\mathcal H_{\rm phys}}
\]
is an isomorphism onto that band, with $B_0=P$.
The possible vacuum row of the conjugated operator is killed
by $\Pi_b^{\rm orig}S_b\omega=0$. Hence
\[
        (T+bV_0-\mathcal E_b)B_b=B_bF_N(b).
\]
Expand $B_b=P+bB_1+b^2B_2+\cdots$ and
$\mathcal E_b=e_1b+e_2b^2+\cdots$. The vacuum equation gives
$e_1=0$, $e_2=-M/(48\kappa)$, as already computed in V11.
Taking the $P$ component at first order gives
$F_{N,1}=PV_0P=0$. Taking the complementary component gives
\[
 (\mathbf1-P)B_1
   =(E-T)^{-1}(\mathbf1-P)V_0P .
\]
This is an inverse on the full complement of $P$, which
includes the vacuum. At second order the $P$ component is
therefore
\[
 F_{N,2}
    =PV_0(\mathbf1-P)(E-T)^{-1}
                       (\mathbf1-P)V_0P-e_2P .
 \tag{V12.17}\label{c18v12:signed-return}
\]
Terms involving $PB_1$ vanish because $PV_0P=F_{N,1}=e_1=0$.
Equation~\eqref{c18v11:centered} evaluates exactly this operator
and gives \eqref{c18v12:coefficient-formula}.
This argument needs only a fixed-volume neighborhood to identify
Taylor coefficients; the uniform disk and error were established
separately, without $\|\Pi_b^{\rm orig}S_bU_b\|$ in their constants.
\end{proof}

Thus the disconnected return cancels the vacuum contribution in
the actual analytic band, and the shared-link label-$2$ channel
remains present. Neither return has been discarded to obtain a
favorable coefficient.

\subsection{A uniform cubic error and the actual physical gap}

\begin{theorem}[A volume-independent spectral remainder]
\label{c18v12:remainder}
For $|b|<\rho_a$, put
\[
 \tau=\frac{|b|}{\rho_a},\qquad
 \varepsilon_a(b)=\frac{2\kappa e^a}{3}\frac{\tau^3}{1-\tau}.
 \tag{V12.18}\label{c18v12:epsilon}
\]
Then
\[
 \begin{split}
 F_N(b)&=12\kappa\mathbf1+
       \frac{b^2}{\kappa}
          \left(-\frac{71}{1680}\mathbf1+
                         \frac{\mathcal L_N^{\rm pl}}{336}\right)
                  +\mathcal R_N(b),\\
 \sup_p\sum_q e^{a d(p,q)}|\mathcal R_N(b)_{qp}|
    &\le\varepsilon_a(b),\qquad
 \|\mathcal R_N(b)\|_{\ell^2\to\ell^2}\le\varepsilon_a(b).
 \end{split}
 \tag{V12.19}\label{c18v12:remainder-formula}
\]
The analogous weighted row bound also holds. In particular the
column tail at distances at least $L$ is bounded by
$\varepsilon_a(b)e^{-aL}$. For real $0\le b<\rho_a$, the actual
physical gap obeys
\[
 \left|\Delta_N-
       \left(12\kappa-\frac{71b^2}{1680\kappa}\right)\right|
                    \le\varepsilon_a(b).
 \tag{V12.20}\label{c18v12:actual-gap}
\]
Every constant here is independent of $N$.
\end{theorem}
\begin{proof}
Subtract the first two coefficients and sum the geometric series
in \eqref{c18v12:coeff-bound}, starting at $n=3$. This gives
the weighted column bound and the stated distance tail.
Translations and cubic symmetries preserve the original
Hamiltonian, the family map, and all operations constructing
$F_N(b)$. The fixed point is equivariant by uniqueness.
These symmetries act transitively on the plaquettes and
preserve $d(p,q)$. Thus all weighted column sums are equal,
as are all weighted row sums. The total of the row sums equals
the total of the column sums, so their common values agree.
The unweighted Schur test now gives the $\ell^2$ bound.
No symmetry under matrix transposition was assumed.

Let $B$ be the Hermitian explicit quadratic matrix in
\eqref{c18v12:remainder-formula} and
$m_2=12\kappa-71b^2/(1680\kappa)$. V11 proves
$B\succeq m_2\mathbf1$. If $v$ is a normalized eigenvector of
$F_N(b)$ with real eigenvalue $\lambda$, then
\[
 \lambda=\langle v,Bv\rangle+
                     \langle v,\mathcal R_N(b)v\rangle
                       \ge m_2-\varepsilon_a(b).
\]
Reality follows from the actual physical spectral identification.
This bounds the lowest band eigenvalue from below.

For the upper bound, the subspace invariant under all translations
and cubic rotations in the plaquette permutation representation
is one-dimensional, spanned by the constant vector $u$.
Equivariance implies $F_N(b)u=\lambda_{\rm s}u$.
Also $\mathcal L_N^{\rm pl}u=0$, so
$|\lambda_{\rm s}-m_2|\le\varepsilon_a(b)$.
This is a genuine physical band eigenvalue. The first nonzero
eigenvalue is at most $\lambda_{\rm s}$ and, by
Theorem~\ref{c18v12:uniform-band}, lies in this band.
This proves \eqref{c18v12:actual-gap} without applying a
Hermitian min--max theorem to the non-Hermitian matrix $F_N(b)$.
\end{proof}

For example, at $a=0$ the radius is $\rho_0=\kappa/2688$.
For $b\le\kappa/5376$ the explicit error is at most $\kappa/6$.
More importantly, as $b/\kappa\to0$ it is
$O(b^3/\kappa^2)$ with a volume-independent, albeit very
conservative, constant. This supersedes the $M$-dependent
remainder in \eqref{c18v11:actual-gap}; it is not a claim of
a useful weak-electric convergence radius.

\subsection{Native cutoff uniformity and the missing adjoint return}

\begin{corollary}[Uniform estimate for the original electric compressions]
\label{c18v12:cutoff}
For every native link-label cutoff $J\ge2$, the compressed
physical Hamiltonian $H_J$ satisfies
\eqref{c18v12:remainder-formula} and
\eqref{c18v12:actual-gap}, with the same $\rho_a$ and
$\varepsilon_a$, independently of both $N$ and $J$.
For $J=1$, the explicit quadratic band matrix is instead
\[
             12\kappa\mathbf1+
                    \frac{b^2}{\kappa}
                        \left(\frac{\mathbf1}{6}+\frac{A_N}{96}\right).
 \tag{V12.21}\label{c18v12:J1-band}
\]
It has the same remainder bound. For even $N\ge6$,
\[
 \left|\Delta_{N,J=1}
              -\left(12\kappa+\frac{b^2}{8\kappa}\right)\right|
                            \le\varepsilon_a(b).
 \tag{V12.22}\label{c18v12:J1-gap}
\]
\end{corollary}
\begin{proof}
Use the original on-link truncated spaces with their constant
vectors and electric operators. Compressed primitive Wilson
terms still have four-link support and norm at most $|b|$.
The creation identities, sector branching bound, physical
electric spectral separation, and cover estimates above
therefore apply with unchanged constants. The simple ground
is identified by the V5 continuous-eigenvalue exclusion
argument; pointwise positivity of a compressed kernel is not
needed. The V11 cutoff calculation gives the same quadratic
coefficient for $J\ge2$ and \eqref{c18v12:J1-band} for $J=1$.

For the latter, $A_N\succeq-4\mathbf1$ gives the lower bound
on each real band eigenvalue as in the previous theorem.
For the upper bound on an even torus, the translation momentum
$k=(\pi,\pi,\pi)$ is allowed. Its three-orientation fiber is
invariant under $F_{N,J=1}(b)$, and the V11 symbol gives
$A_N(k)=-4\mathbf1$. Thus the explicit quadratic matrix on that
fiber is the scalar $12\kappa+b^2/(8\kappa)$.
The remainder restricted to it has norm at most
$\varepsilon_a(b)$, so any of its eigenvalues lies within that
distance of the scalar. These are genuine physical eigenvalues
and supply the required upper bound.
\end{proof}

This makes the earlier cutoff sign difference a uniform
small-$b$ asymptotic statement. It does not say that $J=1$
approximates the weak-electric theory. The omitted adjoint
return remains a concrete warning against that extrapolation.

\subsection{Verification and the next upstream obligation}

The exact diagnostic checks the contour constants and rational
tail bounds. A generic six-level Hermitian fixture verifies the
nonunitary ground dressing, the first two Riesz projection
identities, and the signed second-order return including the
vacuum; it is not a native YM spectrum. Native $N=5$ geometry
checks 2319 rooted covers through size four and 9099 covered
plaquette inclusions. Eight of those inclusions have $q\notin Z$,
which is why \eqref{c18v12:cover-distance} pays the last
adjacency step instead of assuming $q\in Z$. A circulant
diagnostic checks the row/column Schur argument and the constant
mode; it is explicitly not used as an example of the physical
real spectrum. The written proofs, not these finite fixtures,
establish the uniform assertions. No proof-assistant
certification is claimed.

The nonduplicative gain is the exact, volume-uniform,
vacuum-dressed band with a controlled localized error.
The finite-volume estimates now hold at fixed small $b/\kappa$
for arbitrarily large tori. This alone is not a construction
of an infinite-volume Hilbert space, a scattering theory, or
a continuum limit.

The next iteration should go upstream rather than merely
compute a higher strong-electric coefficient. The unpaid step
is a justified physical scale transport from the interacting
weak-electric retained carrier to a regime with a coercive
spectral estimate. V6--V9 already show why fixed link labels,
a frozen rotor, or a small trial-energy error cannot supply
that transport. The new band can be a controlled reference
only if the original normalization, all coupled physical
channels, and accumulated scale errors are bounded on an
actual trajectory. No such trajectory is proved here.

The companion checkpoint remains V15 and the next broad
literature checkpoint V20. This append used a targeted
primary-source check, not a new broad sweep. V11's body is
preserved byte for byte apart from the version title.
This is the single active V12 source; archives and companions
are unchanged. Only \texttt{.tex} files are retained in
\texttt{latex\_notes}, with all build products elsewhere.
The full interacting four-dimensional Yang--Mills construction
and positive physical mass gap remain unproved.
% END CYCLE 18 VERSION 12 APPEND
% BEGIN CYCLE 18 VERSION 13 APPEND
\clearpage
\section{V13: Native coarse kinetics, memory, and a weak-electric scale test}
\label{c18v13:context}

V12's uniform strong-electric band is now a controlled reference,
but it does not supply a scale transformation from weak electric
coupling. Here we examine one fully specified transformation:
multiply the fine links along straight coarse edges, and retain
the \emph{actual} vacuum marginal. We derive its exact compressed
kinetic form, the entire discarded-variable source, and a
quantitative restriction on reaching a strong-electric vacuum.
These are fixed-regulator statements valid at arbitrary coupling.
A separate small-coupling calculation certifies that the source
is genuinely nonzero in the native physical model.

The archive already contains exact blocking words (f8 V13),
a cut-comparison coarse kinetic operator (f15 V98), and native
one-link conditional score and Schur identities (f17 V91--V98).
Those are not reproved as new discoveries. The present coarse
law is the pushforward of the \emph{restored interacting ground},
not that cut comparison or a conditional Gibbs substitute.
Shell-mixing V16 was checked again: its complete flat-root
one-loop coefficient still does not provide the noncommuting
Born calculation or a common nonperturbative integral remainder.
No running coupling is inferred from that coefficient here.

The general distinction between a correct equilibrium marginal
and a correct reduced dynamics is established methodology.
For context, the primary abstract of Legoll--Leli\`evre,
\emph{Effective dynamics using conditional expectations},
\href{https://arxiv.org/abs/0906.4865}{arXiv:0906.4865},
describes conditional-expectation dynamics and additional
conditions for dynamical approximation. We use no error theorem
from that paper. The compact-group identities and native
estimates below are proved explicitly.

\subsection{A compact blocking map without a gauge chart}

Let the fine spatial side be $N=\ell n$, with integers
$\ell\ge2$ and $n\ge5$. A coarse oriented edge
$\gamma=(x,i)$, $x\in(\mathbb Z/n\mathbb Z)^3$, has the
straight fine chain
\[
 e_{\gamma,j}=(\ell x+j\widehat i,i),\quad
 0\le j<\ell,\qquad
 g_\gamma=U_{e_{\gamma,0}}\cdots U_{e_{\gamma,\ell-1}}.
 \tag{V13.1}\label{c18v13:block}
\]
Distinct chains have disjoint fine links, although they may
share endpoints. Write $\mathfrak b_\ell(U)=g$.
Every fine gauge transformation cancels at the intermediate
chain vertices, leaving the usual two endpoint transformations.
Thus a coarse gauge-invariant function pulls back to a function
satisfying the original fine Gauss law.

Let $\Omega>0$ be the actual normalized fine ground,
$\nu=\Omega^2dU$, $K=H-E_0$, and
\[
 \mu_\ell=(\mathfrak b_\ell)_*\nu
        =\varrho_\ell(g)\,dg,\qquad
 J_\ell f=\Omega(f\circ\mathfrak b_\ell),\qquad
 P_\ell=J_\ell J_\ell^*,\quad Q_\ell=\mathbf1-P_\ell .
 \tag{V13.2}\label{c18v13:measure}
\]
Here $dg$ is product normalized coarse Haar measure.
The map $J_\ell:L^2(\mu_\ell)\to L^2(dU)$ is an isometry.
For each chain replace its last fine variable by $g_\gamma$,
retaining the preceding variables and every unused fine link.
This is a global smooth product-Haar coordinate change of
Jacobian one. Integrating $\Omega^2$ over the retained details
proves that $\varrho_\ell$ is smooth and strictly positive.
It is coarse gauge invariant. No logarithmic branch, small-field
event, extra Gauss projection, or independent-block vacuum is
being inserted.

Fix the round $\mathrm{SU}(2)$ orthonormal left-derivative frame
$X_a$, normalized by $-\sum_aX_a^2=L$ with eigenvalues
$j(j+2)$. Write $X_{e,a}$ and $X_{\gamma,a}$ for its fine and
coarse versions. If $h_{\gamma,j}$ is the prefix before
$e_{\gamma,j}$ and $R_{\gamma,j}=\operatorname{Ad}h_{\gamma,j}
\in\mathrm{SO}(3)$, the chain rule gives
\[
 X_{e_{\gamma,j},a}(f\circ\mathfrak b_\ell)
   =\sum_c(R_{\gamma,j})_{ca}
                (X_{\gamma,c}f)\circ\mathfrak b_\ell .
 \tag{V13.3}\label{c18v13:derivative}
\]
The prefix does not depend on that same fine link.

\begin{theorem}[The exact native coarse Dirichlet form]
\label{c18v13:kinetic}
For smooth coarse $f,h$,
\[
 \begin{split}
 \langle J_\ell f,KJ_\ell h\rangle
    &=\kappa\ell\int\sum_{\gamma,a}
        \overline{X_{\gamma,a}f}\,X_{\gamma,a}h\,d\mu_\ell
       =:\mathcal E_\ell(f,h),\\
 \mathcal A_\ell
    &=-\kappa\ell\sum_{\gamma,a}
        \left(X_{\gamma,a}^2+
           (X_{\gamma,a}\log\varrho_\ell)X_{\gamma,a}\right).
 \end{split}
 \tag{V13.4}\label{c18v13:form}
\]
The closed form defines a nonnegative self-adjoint coarse
operator, and $J_\ell^*KJ_\ell=\mathcal A_\ell$ on smooth
functions. These statements restrict to the original physical
spaces. The coefficient is exactly $\kappa\ell$ in the original
Hamiltonian time variable.
\end{theorem}
\begin{proof}
The ground-state identity is
$\langle\Omega F,K\Omega G\rangle
=\kappa\int\sum_{e,a}\overline{X_{e,a}F}X_{e,a}G\,d\nu$.
Use \eqref{c18v13:derivative}, orthogonality of each $R$, and
disjointness of the chains. Each coarse derivative pair is
counted exactly $\ell$ times; all unused fine derivatives
vanish. This proves the first line of \eqref{c18v13:form}.
Integration by parts in coarse Haar measure gives the second.
A smooth positive density on a finite compact product makes
this a closed weighted $H^1$ form with smooth core and a
self-adjoint elliptic generator. The coordinate change above
and smoothness of $\Omega$ justify the compression identity
on that core. Gauge equivariance gives the physical restriction.
\end{proof}

This is an operator compression, not an equality between
$K$ and $\mathcal A_\ell$, and not yet a renormalized
Wilson Hamiltonian with a newly assigned magnetic coefficient.
Indeed $J_\ell1=\Omega$, so the physical Rayleigh principle
gives $\operatorname{gap}(K)\le\operatorname{gap}(\mathcal A_\ell)$.
A lower bound for the compressed gap alone therefore has
the wrong direction to close the fine physical gap.
The density $\varrho_\ell$ is unknown in general. Rescaling
the time variable without accounting for its physical meaning
would not eliminate that density or the source derived next.

\subsection{The full discarded source is a conditional covariance}

Put $s_{e,a}=X_{e,a}\log\Omega$ and define the transported sum
and its conditional mean by
\[
 \sigma_{\gamma,c}(U)
     =\sum_{j,a}(R_{\gamma,j})_{ca}s_{e_{\gamma,j},a}(U),
 \qquad
 \bar\sigma_{\gamma,c}(g)
       =\mathbb E_\nu[\sigma_{\gamma,c}\mid\mathfrak b_\ell=g].
 \tag{V13.5}\label{c18v13:sigma}
\]
Indices $\alpha=(\gamma,c)$ will be used for the complete
coarse derivative bank. Define
\[
 \mathsf C_{\alpha\delta}(g)
   =\mathbb E_\nu[
      (\sigma_\alpha-\bar\sigma_\alpha)
      (\sigma_\delta-\bar\sigma_\delta)
                     \mid\mathfrak b_\ell=g].
 \tag{V13.6}\label{c18v13:covariance}
\]
All cross-chain covariances are retained.

\begin{theorem}[Native mean score and exact leakage]
\label{c18v13:leakage}
One has
\[
  \bar\sigma_{\gamma,c}
                =\frac{\ell}{2}X_{\gamma,c}\log\varrho_\ell .
 \tag{V13.7}\label{c18v13:mean}
\]
For the actual off-diagonal source
$\mathcal T_\ell=Q_\ell KJ_\ell$,
\[
 \begin{split}
 \mathcal T_\ell f
    &=-2\kappa\Omega\sum_\alpha
       (\sigma_\alpha-\bar\sigma_\alpha\circ\mathfrak b_\ell)
                     (X_\alpha f)\circ\mathfrak b_\ell,\\
 \|\mathcal T_\ell f\|^2
    &=4\kappa^2\int
             (\nabla f)^*\mathsf C(g)\nabla f\,d\mu_\ell .
 \end{split}
 \tag{V13.8}\label{c18v13:T}
\]
In particular a correct coarse vacuum density alone does not
assert that this source vanishes.
\end{theorem}
\begin{proof}
Consider the fine vector field
$Y_{\gamma,c}=\sum_{j,a}(R_{\gamma,j})_{ca}X_{e_{\gamma,j},a}$.
It has zero Haar divergence: the derivative of its coefficient
in its own fine variable is zero. Also
$Y_{\gamma,c}(h\circ\mathfrak b_\ell)
=\ell(X_{\gamma,c}h)\circ\mathfrak b_\ell$ and
$Y_{\gamma,c}\log\Omega=\sigma_{\gamma,c}$.
Integration by parts against $\Omega^2$ gives
\[
 \ell\int X_{\gamma,c}h\,\varrho_\ell\,dg
       =-2\int h\,\bar\sigma_{\gamma,c}\varrho_\ell\,dg.
\]
A second integration by parts in $dg$ proves
\eqref{c18v13:mean}.

For the second identity in \eqref{c18v13:derivative}, the same
prefix independence and orthogonality also imply
$\sum_{e,a}X_{e,a}^2(f\circ\mathfrak b_\ell)
=\ell\sum_{\gamma,c}(X_{\gamma,c}^2f)\circ\mathfrak b_\ell$.
The exact ground-transformed generator is
$-\kappa\sum_{e,a}(X_{e,a}^2+2s_{e,a}X_{e,a})$.
Apply it to $f\circ\mathfrak b_\ell$, then subtract
$J_\ell\mathcal A_\ell f$ using \eqref{c18v13:mean}.
This gives the first line of \eqref{c18v13:T}.
Conditional expectation proves both its orthogonality to
$\operatorname{Ran}J_\ell$ and the squared-norm formula.
\end{proof}

There are useful finite-scale bounds, but their coupling
dependence must be displayed. Write $\beta=b/\kappa>0$ and
$A_*=3\sqrt2$. The inherited native estimates are
\[
 |\mathbf s_e|\le2\beta,\qquad
 \int|\mathbf s_e|^2\,d\nu
       =\langle\Omega,L_e\Omega\rangle\le A_*\sqrt\beta .
 \tag{V13.9}\label{c18v13:input}
\]
The first is f17 V93 (V93.3), whose compact maximum-principle
proof was rechecked; the second is V6's actual-vacuum energy
bound, not a statement about a trial density.

\begin{corollary}[Explicit local-observable source budgets]
\label{c18v13:source-budget}
If $f$ depends on a set $\mathcal S$ of $m$ coarse edges, then
\[
 \begin{split}
 \|\mathcal T_\ell f\|^2
       &\le \frac{16m\ell b^2}{\kappa}\,\mathcal E_\ell(f,f),\\
 \|\mathcal T_\ell f\|^2
       &\le4A_*\kappa^2m\ell^2\sqrt\beta
              \sum_{\gamma\in\mathcal S}
                            \|\nabla_\gamma f\|_\infty^2 .
 \end{split}
 \tag{V13.10}\label{c18v13:budgets}
\]
The constants do not depend on the fine torus size.
\end{corollary}
\begin{proof}
The orthogonal conditional projection is a contraction, so
one may bound the uncentered
$\sum_{\gamma\in\mathcal S}\sigma_\gamma\cdot\nabla_\gamma f$
in \eqref{c18v13:T}. Cauchy--Schwarz over the $m$ terms and
$|\sigma_\gamma|\le2\ell\beta$ give the first bound using
\eqref{c18v13:form}. Alternatively,
\[
 \int|\sigma_\gamma|^2d\nu
       \le\ell\sum_j\int|\mathbf s_{e_{\gamma,j}}|^2d\nu
       \le A_*\ell^2\sqrt\beta.
\]
Bound each derivative of $f$ in supremum norm to obtain the
second line. At $b=0$ the score and the source are zero.
\end{proof}

Neither estimate becomes a small relative error as
$\beta\to\infty$. The factor $m$ is a local observable support
size, not a uniform bound for an arbitrary extensive function.
A spectral proof would need more than these budgets:
it would need control of the full covariance on the relevant
physical gradients and of the discarded dynamics.

\subsection{A native certificate that instantaneous compression loses dynamics}

For smooth $f$, let
\[
 C_f(t)=\langle J_\ell f,e^{-tK}J_\ell f\rangle,\qquad
 D_f(t)=\langle f,e^{-t\mathcal A_\ell}f\rangle_{\mu_\ell}.
\]
Both use the same fixed-regulator Hamiltonian time.

\begin{proposition}[The first dynamical error is the leakage square]
\label{c18v13:memory}
At each finite regulator and coupling,
\[
 C_f(t)-D_f(t)
           =\frac{t^2}{2}\|\mathcal T_\ell f\|^2+o(t^2)
                         \quad(t\downarrow0).
 \tag{V13.11}\label{c18v13:short-time}
\]
Moreover, for $f(g)=\chi_1(U_C(g))$ on one elementary coarse
plaquette and fixed $N,\ell$,
\[
 \|\mathcal T_{\ell,b}f\|^2
       =\frac{2\ell+1}{6}\,b^2
                         +O_{N,\ell}(b^3/\kappa)
                          \quad(b/\kappa\to0).
 \tag{V13.12}\label{c18v13:native-memory}
\]
In particular this is a nonzero physical source for sufficiently
small positive $b$. The expansion is a calibration of the
blocking operation, not a weak-electric estimate.
\end{proposition}
\begin{proof}
The isometry and form identity give $C_f(0)=D_f(0)$ and
$C_f'(0)=D_f'(0)$. Since
$KJ_\ell f=J_\ell\mathcal A_\ell f+\mathcal T_\ell f$
is an orthogonal decomposition,
\[
 C_f''(0)-D_f''(0)
        =\|KJ_\ell f\|^2-\|\mathcal A_\ell f\|^2
        =\|\mathcal T_\ell f\|^2 .
\]
Smooth functions and the smooth ground have the required
operator domains. Spectral calculus justifies the right
Taylor expansion.

We compute the nonzero coefficient in the original Haar
variables. Write $S(U)=\sum_pw_p(U)$ and
$F=f\circ\mathfrak b_\ell=\chi_1(U_{\partial R})$, where
$R$ is the fine $\ell$-by-$\ell$ square. Ordinary
fixed-volume ground perturbation gives
\[
 \Omega_b=1+\frac{b}{12\kappa}S+O_N(b^2/\kappa^2),
 \qquad T F=12\kappa\ell F .
 \tag{V13.13}\label{c18v13:first-ground}
\]
The expansion holds in every fixed smooth norm: at finite
volume the simple elliptic eigenpair is analytic locally,
and elliptic regularity applied to its differentiated
equation upgrades the expansion successively.

Let $\mathcal W$ be the union of all straight chain links.
Every fine plaquette has at least one link outside
$\mathcal W$ when $\ell\ge2$: its two parallel edges differ
by one unit in a transverse coordinate, so they cannot both
have all transverse coordinates divisible by $\ell$.
Haar integration in such a link
gives
$\mathbb E_{dU}[S\mid\mathfrak b_\ell]=0$.
Consequently $\varrho_{\ell,b}=1+O_N(b^2/\kappa^2)$ and
$\mathcal A_{\ell,b}$ has no term linear in $b$.
Expanding the exact identity
$K_bJ_{\ell,b}f-J_{\ell,b}\mathcal A_{\ell,b}f$ yields
\[
 \mathcal T_{\ell,b}f
       =-\frac b6\sum_{p:p\cap\partial R\ne\varnothing}G_p
                                 +O_{N,\ell}(b^2/\kappa),
 \qquad
 G_p=\sum_{e,a}(X_{e,a}w_p)(X_{e,a}F).
 \tag{V13.14}\label{c18v13:first-leak}
\]
For example this also follows directly from the product rule
for $T(SF)$ in \eqref{c18v13:first-ground}.
Each $G_p$ has an odd link outside $\mathcal W$, so its
conditional Haar projection is zero. Distinct $p$ give
distinct link-centre parity masks
$\partial p\mathbin{\triangle}\partial R$, hence orthogonal
$G_p$'s in Haar $L^2$.

A touching fine plaquette shares either one perimeter link,
or two consecutive links at a corner. Put this number equal
to $r$. On that common path, the product $w_pF$ decomposes
into labels $0$ and $2$, with squared norms $1/16$ and $3/16$.
To verify the normalization directly, absorb the disjoint
remaining paths into two independent unit quaternions $a,c$
and the common path into $u$. The product is
$2(u\cdot a)(u\cdot c)$, whose constant-$u$ part is
$(a\cdot c)/2$. The round Haar second and fourth moments give
the two squared norms above. At an interior vertex of a
two-link common path, the two labels must match by Gauss
invariance, so the same two channels exhaust the product.

Their electric energies divided by $\kappa$ are respectively
$12\ell+12-6r$ and $12\ell+12+2r$. The Laplacian product
rule therefore multiplies these components in $G_p$ by
$3r$ and $-r$. Thus
\[
                 \|G_p\|^2=\frac{3r^2}{4}.
 \tag{V13.15}\label{c18v13:fusion}
\]
Both fusion channels have been retained.

There are four corner plaquettes with $r=2$.
The $4\ell$ perimeter links each belong to four fine
plaquettes; subtracting those eight corner incidences leaves
$16\ell-8$ plaquettes with $r=1$. The side restrictions
exclude wraparound identifications. Orthogonality in
\eqref{c18v13:first-leak} gives
\[
 \frac1{36}\sum_p\|G_p\|^2
   =\frac1{48}\bigl((16\ell-8)+4\cdot4\bigr)
   =\frac{2\ell+1}{6}.
\]
Squaring the remainder proves \eqref{c18v13:native-memory}.
\end{proof}

Thus even an exactly correct instantaneous coarse form does
not make $J_\ell^*e^{-tK}J_\ell=e^{-t\mathcal A_\ell}$.
This failure has been checked on an original physical Wilson
observable, rather than inferred from a generic matrix example.
It is not a lower bound on the source at large $\beta$, and
it is not a mass-gap statement.

\subsection{An actual-vacuum restriction on the amount of blocking}

For a coarse plaquette $C$, let $w_C$ be its normalized trace
and $\mathcal S_C$ its $\ell^2$ elementary fine faces.

\begin{lemma}[A compact surface-defect estimate]
\label{c18v13:surface}
At every fine configuration,
\[
 1-w_C(\mathfrak b_\ell(U))
          \le\ell^2\sum_{p\in\mathcal S_C}(1-w_p(U)).
 \tag{V13.16}\label{c18v13:surface-bound}
\]
Consequently the actual coarse marginal obeys
\[
       \int w_C\,d\mu_\ell
                    \ge1-A_*\ell^4\beta^{-1/2}.
 \tag{V13.17}\label{c18v13:coarse-defect}
\]
\end{lemma}
\begin{proof}
The ordered nonabelian surface word, already used in f8 V13
for $\ell=2$, expresses the large perimeter as a product
of $\ell^2$ rooted conjugates of its elementary plaquettes.
For arbitrary $\ell$ it follows by taking the tree of all
horizontal links and the left vertical boundary: in tree
gauge each row is the product of its plaquettes from right
to left, and rows are multiplied from bottom to top.
Undo the root-fixed gauge change.

For unitary matrices, telescoping the product and using
unitary invariance of the Frobenius norm gives
\[
 \left\|I-\prod_{k=1}^{\ell^2}V_k\right\|_{\rm F}^2
             \le\ell^2\sum_k\|I-V_k\|_{\rm F}^2.
\]
On $\mathrm{SU}(2)$, $\|I-V\|_{\rm F}^2=4(1-\tfrac12\Tr V)$.
Conjugation leaves each norm unchanged, proving
\eqref{c18v13:surface-bound}. Integrate against the actual
fine $\nu$ and use the uniform plaquette-defect bound
$\langle1-w_p\rangle_\nu\le A_*\beta^{-1/2}$ from V6.
\end{proof}

The $\ell^4$ in the expectation bound cannot be improved by
this deterministic surface inequality alone. If all rooted
plaquettes equal $\exp(\theta\tau_3)$, the ratio of the
perimeter defect to the common fine defect tends to
$(\ell^2)^2$ as $\theta\to0$. Such a coherent configuration is
realized in the same tree gauge. Improving the probabilistic
bound would therefore require information on the actual
curvature correlations, not a different order of this
triangle inequality.

We compare with a precisely specified possible endpoint:
let $\nu_{\rm s}$ be the actual vacuum law of a native
homogeneous Wilson Hamiltonian on the coarse $n$-torus,
with coefficients $\kappa_{\rm s}>0$, $b_{\rm s}\ge0$ and
$\theta_{\rm s}=b_{\rm s}/\kappa_{\rm s}$. No assertion that
$\mu_\ell=\nu_{\rm s}$ is made.

\begin{lemma}[A simple strong-electric vacuum moment]
\label{c18v13:strong-moment}
For $0\le\theta_{\rm s}\le1$, every coarse plaquette satisfies
\[
       \int w_C\,d\nu_{\rm s}
                   \le m(\theta_{\rm s})
                :=\sqrt{\theta_{\rm s}/3}+\theta_{\rm s}/3 .
 \tag{V13.18}\label{c18v13:m}
\]
\end{lemma}
\begin{proof}
The constant trial has energy $b_{\rm s}M_{\rm c}$.
The nonnegative magnetic potential and cubic symmetry of
the unique ground give
$\langle L_e\rangle_{\nu_{\rm s}}\le\theta_{\rm s}$.
For the one-link constant projection $P_e$, put
$q=\|(\mathbf1-P_e)\Omega_{\rm s}\|^2\le\theta_{\rm s}/3$.
These projections are used only for an inequality on the
full Haar space, not as extra physical Gauss constraints.

On a link of $C$, $P_ew_CP_e=0$ and $\|w_CP_e\|=1/2$.
The latter follows from the conditional Haar integral
$\int w_C^2\,dU_e=1/4$, with the other links fixed.
Splitting $\Omega_{\rm s}=P_e\Omega_{\rm s}
+(\mathbf1-P_e)\Omega_{\rm s}$ gives
\[
 \langle w_C\rangle_{\nu_{\rm s}}
          \le\sqrt{q(1-q)}+q
          \le\sqrt{\theta_{\rm s}/3}+\theta_{\rm s}/3 .
\]
\end{proof}

\begin{theorem}[A necessary scale for reaching a strong-electric law]
\label{c18v13:scale}
With $d_{\rm TV}(\mu,\nu)=\sup_A|\mu(A)-\nu(A)|$,
\[
 d_{\rm TV}(\mu_\ell,\nu_{\rm s})
       \ge\frac12
          \left[1-A_*\ell^4\beta^{-1/2}
                        -m(\theta_{\rm s})\right]_+ .
 \tag{V13.19}\label{c18v13:TV}
\]
Thus if $d_{\rm TV}(\mu_\ell,\nu_{\rm s})\le\delta$ and
$1-m(\theta_{\rm s})-2\delta>0$, necessarily
\[
 \ell\ge
    \left(\frac{1-m(\theta_{\rm s})-2\delta}{A_*}\right)^{1/4}
                                  \beta^{1/8}.
 \tag{V13.20}\label{c18v13:scale-lower}
\]
In particular this applies to an endpoint in V12's stated
strong-electric disk, $\theta_{\rm s}\le1/2688$.
\end{theorem}
\begin{proof}
A function bounded in $[-1,1]$ has expectation difference
at most $2d_{\rm TV}$. Apply this to the same literal coarse
Wilson plaquette using \eqref{c18v13:coarse-defect} and
\eqref{c18v13:m}, then solve the resulting inequality for
$\ell$. The argument uses one physical observable, so it also
bounds the total variation of the corresponding gauge-orbit
laws or of that observable's one-dimensional marginals.
\end{proof}

This is a necessary condition for this \emph{straight-product,
unchanged-observable} comparison, not a no-go theorem for
other blocking maps, field redefinitions, or extended effective
actions. It neither proves that a scale of order $\beta^{1/8}$
is sufficient nor constructs a running coupling. The same
vacuum-moment mismatch cannot be repaired merely by multiplying
a Hamiltonian by a time-conversion constant, which leaves its
ground state unchanged.

\subsection{What this step supplies, and the next unpaid estimate}

The blocking operation is now specified on the original compact
variables and actual vacuum, with exact kinetic coefficient
$\kappa\ell$. Its entire departure from closed instantaneous
dynamics is the conditional covariance
\eqref{c18v13:covariance}, not a discarded scalar. The native
loop calculation certifies that this departure is real.
The weak-electric moment estimate gives a volume-independent
necessary blocking length for a particular strong-electric
endpoint. None of these results supplies a uniform physical
spectral gap in the weak-electric regime.

The next useful target is quantitative control of that
transported-score covariance on physical coarse gradients,
together with the actual discarded return, at a growing
blocking scale. Alternatively an improved observable map must
come with its native kinetic metric and the analogous source
estimate. Replacing the covariance by its conditional mean,
treating a marginal density as a complete reduced semigroup,
or extrapolating \eqref{c18v13:native-memory} to large $\beta$
would bypass rather than solve this obligation. The existing
one-link and cutoff results remain available; repeating their
conditional form alone is not a new estimate here.

Exact diagnostics verify chain metric factors and gauge
covariance for five lengths, literal noncommuting surface
words for $\ell=2,3$, and the coherent-curvature limiting ratio.
Fine tori of sides $10,15,20$ verify every touching-plaquette
count, distinct parity mask, and unobserved odd link for
$\ell=2,3,4$. Independent quaternion Haar integration verifies
both fusion norms and the gradient-pair norm. A generic
positive matrix checks the short-time square identity only;
it is not a native spectrum. These supplement, but do not
replace, the analytic proofs. V12's full diagnostics and
preservation check were replayed before this append.

V12's full body is preserved except for its version title,
and the single active source is now V13. All 47 other note
files are unchanged; only \texttt{.tex} files are stored in
\texttt{latex\_notes}. The next companion checkpoint remains
V15 and the next broad literature sweep V20. The nontrivial
interacting four-dimensional continuum construction and the
full positive physical Yang--Mills mass gap remain unproved.
% END CYCLE 18 VERSION 13 APPEND
% BEGIN CYCLE 18 VERSION 14 APPEND
\clearpage
\section{V14: Native discarded curvature and a low-energy coverage test}
\label{c18v14:context}

V13 specifies an actual spatial blocking and its exact covariance
source. Before trying to make that source small, we must check
whether the discarded space is a high-energy space. Here a fine
Wilson plaquette, centered in one \emph{actual} conditional vacuum,
supplies a physical vector in that discarded space. Its Rayleigh
quotient has an explicit bound independent of both the fine volume
and the blocking length. There is a polynomial all-coupling bound,
and a bound below $13\kappa$ throughout V12's reference disk
$0\le b/\kappa\le1/2688$. This is an \emph{upper} bound on a
discarded energy, not a lower mass-gap bound.

The one-link conditional identities and sensitivity budgets in
f17 V92--V96 are reused, not renamed as discoveries. In particular
the conditional is not the classical Wilson Gibbs conditional of
f15 V38, and it is not the frozen rotor of f17 V88. The new point
is the original-Gauss-invariant curvature test in the
\emph{spatially} discarded space, with its norm bounded from below
and its entire exterior derivative cost paid. No perturbation
remainder proportional to $Mb$ is used.

For context, the primary abstract of
N\"uske--Koltai--Boninsegna--Clementi,
\emph{Spectral Properties of Effective Dynamics from Conditional
Expectations},
\href{https://arxiv.org/abs/1901.01557}{arXiv:1901.01557},
describes a relative eigenvalue-error estimate in terms of
eigenfunction approximation. We import no error theorem from that
paper; the estimates below are direct native variational and
spectral-measure arguments.

\subsection{A physical fluctuation invisible to the retained variables}

Keep V13's $N=\ell n$, $\ell\ge2$, $n\ge5$, actual positive
ground $\Omega$, $K=H-E_0$, and $P_\ell=J_\ell J_\ell^*$.
Let $\mathcal W$ be the set of fine links used by the straight
chains. Every fine plaquette $p$ has an
$e\in\partial p\setminus\mathcal W$, as proved in V13.
Write $q=U_e$, $y=U_{e^c}$, and
\[
 \rho_e(y)^2=\int\Omega(q,y)^2dq,\qquad
 \zeta_e(q,y)=\Omega(q,y)/\rho_e(y),\qquad
 d\bar\nu_e=\rho_e^2dy .
 \tag{V14.1}\label{c18v14:conditional}
\]
Write $\mathbb E_e$ and $\operatorname{Var}_e$ for conditional
expectation and variance in $\zeta_e(q,y)^2dq$, with $y$ fixed.
For the real physical observable $F_p=\chi_1(U_p)=2w_p$, set
\[
 u_{e,p}(y)=\mathbb E_e F_p,\qquad
 v_{e,p}=\Omega(F_p-u_{e,p}),\qquad
 \mathcal V_{e,p}=\|v_{e,p}\|^2
    =\int\operatorname{Var}_e(F_p)\,d\bar\nu_e .
 \tag{V14.2}\label{c18v14:innovation}
\]
This uses the fully restored Hamiltonian's conditional ground law.

\begin{lemma}[Exact discarded-space membership]\label{c18v14:membership}
The vector $v_{e,p}$ is smooth, nonzero, invariant under the
original full fine gauge group, and
\[
 v_{e,p}\in Q_\ell\mathcal H_{\rm phys},\qquad
 \langle\Omega,v_{e,p}\rangle=0 .
 \tag{V14.3}\label{c18v14:Q}
\]
It is orthogonal to every vector $\Omega h(y)$ with
$h\in L^2(\bar\nu_e)$.
\end{lemma}
\begin{proof}
Compact integration gives smooth positive $\rho_e,\zeta_e$ and
smooth real $u_{e,p}$.
For fixed $y$, the plaquette character has the form
$F_p(q,y)=2a(y)\cdot q$ for a unit quaternion $a(y)$, after a
possible orthogonal change of $q$. It is nonconstant, so its
variance in a strictly positive conditional density is positive.
Conditional centering proves
$\langle\Omega h(y),v_{e,p}\rangle=0$, first for bounded $h$ and
then by density. Every coarse function pulled back by
$\mathfrak b_\ell$ is independent of the unused link $e$.
This proves \eqref{c18v14:Q}, including $h=1$.
A fine gauge transformation acts on $q$ by left and right
multiplication and transforms $y$ compatibly. Haar change of
variables in both integrals defining $\mathbb E_eF_p$ shows that
$u_{e,p}$ is invariant as a function of the transformed exterior.
Hence $v_{e,p}$ is physical. No one-link Gauss constraint is imposed.
\end{proof}

Only the existence of an unused link was needed for orthogonality.
The same vector tests any retained observable family independent
of $e$. This does not assign an independent physical tensor factor
to that link.

\subsection{An exact energy identity with the entire exterior cost}

Let $i=(f,a)$ run over \emph{all} exterior derivative components,
$f\ne e$, and put $s_i=X_i\log\Omega$. Define
\[
 j_e(y)=\sum_i\operatorname{Var}_e(s_i),\qquad
 J_e=\int j_e\,d\bar\nu_e,\qquad
 I_e(y)=\int|\nabla_e\zeta_e|^2dq .
 \tag{V14.4}\label{c18v14:costs}
\]
Here $J_e$ is a scalar sensitivity mean, not V13's coarse isometry.
Differentiating the normalized conditional integral gives the
inherited formula
\[
 X_i u_{e,p}=\mathbb E_e X_iF_p+
                  2\operatorname{Cov}_e(F_p,s_i).
 \tag{V14.5}\label{c18v14:mean-derivative}
\]
Neither $u_{e,p}$ nor $j_e$ is declared local in the exterior.

\begin{proposition}[Complete native innovation energy]\label{c18v14:energy}
One has
\[
 \begin{split}
 \kappa^{-1}\langle v_{e,p},Kv_{e,p}\rangle
  =\int\bigg\{&
     \mathbb E_e|\nabla_eF_p|^2+
     \sum_i\operatorname{Var}_e(X_iF_p)\\
     &+4\sum_i|\operatorname{Cov}_e(F_p,s_i)|^2
                                      \bigg\}\,d\bar\nu_e .
 \end{split}
 \tag{V14.6}\label{c18v14:energy-identity}
\]
Consequently
\[
 \langle v_{e,p},Kv_{e,p}\rangle
       \le\kappa\left(\int|\nabla F_p|^2d\nu+16J_e\right).
 \tag{V14.7}\label{c18v14:energy-bound}
\]
\end{proposition}
\begin{proof}
The ground transform gives
$\langle v,Kv\rangle=\kappa\int|\nabla(F_p-u_{e,p})|^2d\nu$.
The $e$ derivative of $u_{e,p}$ is zero. For exterior $i$,
\[
 X_i(F_p-u_{e,p})=
 (X_iF_p-\mathbb E_eX_iF_p)-2\operatorname{Cov}_e(F_p,s_i).
\]
The first term has zero conditional mean and the second is
exterior measurable. Their cross term integrates to zero exactly.
Summing squares proves \eqref{c18v14:energy-identity}.
Conditional Cauchy--Schwarz gives
$\sum_i|\operatorname{Cov}_e(F_p,s_i)|^2
 \le\operatorname{Var}_e(F_p)j_e\le4j_e$, since $|F_p|\le2$.
Dropping the subtracted mean squares from the derivative variances
gives \eqref{c18v14:energy-bound}. All functions are smooth on
the finite compact regulator, justifying the differentiations
and operator domains.
\end{proof}

Record explicitly the previously paid inputs. Set
\[
 \beta=b/\kappa,\quad A_*=3\sqrt2,\quad
 \alpha=2\pi\beta,\quad R=e^{2\alpha},\quad d=R-1 .
 \tag{V14.8}\label{c18v14:parameters}
\]
The single-link score bound in \eqref{c18v13:input} and the
diameter $\pi$ of round $S^3$ imply
\[
 R^{-1}\le\zeta_e(q,y)^2\le R,\qquad
                      \eta=3\kappa/R
 \tag{V14.9}\label{c18v14:density}
\]
as a valid conditional Poincare floor. Indeed
$\operatorname{osc}_q\log\Omega\le\alpha$, and the ratio of
maximum to minimum conditional density is at most $R$.
Variance minimization and Haar Poincare therefore give the
floor $3\kappa/R$, not $3\kappa/R^2$.

The native mean estimates in f17 V92 and V96 give
\[
 \int I_e\,d\bar\nu_e
   =\langle\Omega,L_e\Omega\rangle\le A_*\sqrt\beta,\qquad
 J_e\le\mathfrak j(\beta)
     :=\min\{4\beta,\ 90\sqrt\beta,\ \beta R(R-1)\}.
 \tag{V14.10}\label{c18v14:budgets}
\]
The first is also \eqref{c18v6:inherited}. The $4\beta$ bound
is the Haar-competitor consequence of f17 (V96.9):
$\kappa\int(I_e+j_e)d\bar\nu_e\le4b$.
The $90\sqrt\beta$ bound is f17 (V96.14), from the tempered
competitor at $\tau=\sqrt{\kappa b}$ and the \emph{actual}
exterior marginal. Its mean proof in f17 V89 was rechecked.
These bounds do not identify a trial conditional with $\zeta_e$.

For completeness, the last term follows from conditional
Poincare and f17's exact mixed-score row (V92.7):
\[
 J_e\le\frac{\kappa}{\eta}
    \sum_{f\ne e,a,c}\mathbb E_\nu
                   |X_{f,c}X_{e,a}\log\Omega|^2
       \le\frac{\kappa}{\eta}\,6\beta\,\mathbb E_\nu w_p .
 \tag{V14.11}\label{c18v14:small-J}
\]
Homogeneity makes the incident plaquette means equal; the same
row identity makes that mean nonnegative. Since
$\int w_p\,dq=0$ and $\int w_p^2dq=1/4$, the density comparison
gives $|\mathbb E_e w_p|\le(R-1)/2$. Substituting
$\eta=3\kappa/R$ proves $J_e\le\beta R(R-1)$.
Thus $J_e=O(\beta^2)$ uniformly in volume at small $\beta$.
At $\beta=0$ all these bounds have their stated zero limits.

\subsection{A polynomial denominator bound from spherical uncertainty}

An upper energy numerator is insufficient if the centered test
vector is too small. The following elementary spectral estimate
pays its norm without an exponential density comparison.

\begin{lemma}[Coordinate variance from kinetic energy]\label{c18v14:variance}
Let $\zeta\in H^1(S^3)$, $\|\zeta\|_2=1$ in normalized Haar
measure, and $I=\int|\nabla\zeta|^2dq$. For unit
$a\in\mathbb R^4$ and $F(q)=2a\cdot q$,
\[
 \operatorname{Var}_{|\zeta|^2dq}(F)
                  \ge\frac1{40500(1+I)^3}.
 \tag{V14.12}\label{c18v14:uncertainty}
\]
Neither the constant nor the power is claimed optimal.
\end{lemma}
\begin{proof}
Rotate $a$ to the first coordinate. Let $\mathsf P$ be the
projection of $L=-\Delta$ onto eigenvalues at most $r=4(I+1)$.
Then $\|(\mathbf1-\mathsf P)\zeta\|_2\le\sqrt{I/r}<1/2$.
The degree-$j$ eigenvalue and multiplicity are $j(j+2)$ and
$(j+1)^2$. Consequently
\[
 D:=\operatorname{rank}\mathsf P
   =\sum_{k=1}^{\lfloor\sqrt{r+1}\rfloor}k^2
   \le(r+1)^{3/2}\le[5(I+1)]^{3/2}.
 \tag{V14.13}\label{c18v14:rank}
\]
The projector diagonal is constant by rotational invariance,
and its Haar integral is $D$. Cauchy--Schwarz in an orthonormal
basis gives $\|\mathsf P\zeta\|_\infty\le\sqrt D$.

Write $m=\mathbb EF$ and $v=\operatorname{Var}(F)>0$.
Positivity follows because a level set of this linear coordinate
has zero Haar measure, while $|\zeta|^2dq$ is an absolutely
continuous probability. For $A=\{|F-m|\le2\sqrt v\}$, Chebyshev
gives probability at least $3/4$. Hence
\[
 \|\mathbf1_A\mathsf P\zeta\|_2
       \ge\sqrt3/2-1/2>1/3,\qquad
 D\,\operatorname{Haar}(A)\ge1/9 .
\]
The Haar density of $F$ on $[-2,2]$ is
$\pi^{-1}\sqrt{1-F^2/4}\le1/\pi$.
Thus $\operatorname{Haar}(A)\le4\sqrt v/\pi\le2\sqrt v$,
and $v\ge1/(324D^2)$.
Use \eqref{c18v14:rank} and $324\cdot125=40500$.
\end{proof}

Apply the lemma in each actual conditional, with its unit
$a(y)$. Jensen's inequality for the convex decreasing function
$(1+x)^{-3}$, followed by \eqref{c18v14:budgets}, gives
\[
 \mathcal V_{e,p}\ge
 \max\left\{R^{-1},\
          \frac1{40500(1+A_*\sqrt\beta)^3}\right\}.
 \tag{V14.14}\label{c18v14:norm-lower}
\]
The first bound follows separately from
$\operatorname{Var}_eF_p=\inf_c\int(F_p-c)^2\zeta_e^2dq$
and the Haar variance one. No independence of conditionals is used.

\subsection{An all-coupling discarded-energy bound without a blocking factor}

Direct differentiation of $F_p=2w_p$ in its four links gives
\[
 |\nabla F_p|^2=16-4F_p^2,\qquad
       \int|\nabla F_p|^2d\nu\le\min\{16,12R\}.
 \tag{V14.15}\label{c18v14:plaquette-gradient}
\]
The second estimate integrates in the selected link: the
conditional Haar integral of $16-4F_p^2$ is $12$. Define
\[
 \begin{split}
 B_{\rm s}(\beta)&=12R^2+16R\,\mathfrak j(\beta),\\
 B_{\rm p}(\beta)&=648000(1+A_*\sqrt\beta)^3
                           (1+\min\{4\beta,90\sqrt\beta\}),\\
 B(\beta)&=\min\{B_{\rm s}(\beta),B_{\rm p}(\beta)\}.
 \end{split}
 \tag{V14.16}\label{c18v14:B}
\]
Thus $B(0)=12$, $B(\beta)=12+O(\beta)$ as $\beta\downarrow0$,
and $B(\beta)=O(\beta^2)$ as $\beta\to\infty$.
The last statement is only an upper-growth estimate, not a
physical eigenvalue asymptotic.

Let $D_\ell$ be the self-adjoint operator of the restricted
form of $K$ on $Q_\ell\mathcal H_{\rm phys}$. At fixed regulator
the conditional projection preserves smooth functions and $H^1$,
with finite constants not claimed uniform in $N$.
Projecting smooth approximants gives a form core in this
subspace; compact embedding gives compact resolvent.
The subspace is perpendicular to $\Omega$, so $D_\ell>0$
at each finite regulator.

\begin{theorem}[Physical discarded curvature at bounded native energy]
\label{c18v14:discarded}
For every allowed $N,\ell,n$, every $b\ge0$, and every selected
$p,e$,
\[
 \frac{\langle v_{e,p},Kv_{e,p}\rangle}{\|v_{e,p}\|^2}
       \le\kappa B(\beta),\qquad
 0<\inf\operatorname{spec}D_\ell\le\kappa B(\beta),\qquad
 \|D_\ell^{-1}\|\ge\frac1{\kappa B(\beta)}.
 \tag{V14.17}\label{c18v14:upper}
\]
All constants are independent of volume and blocking length.
\end{theorem}
\begin{proof}
Combine \eqref{c18v14:energy-bound},
\eqref{c18v14:budgets}, \eqref{c18v14:norm-lower}, and
\eqref{c18v14:plaquette-gradient}.
Using numerator $12R+16J_e$ and denominator $R^{-1}$ gives
$B_{\rm s}$. Using $16+16J_e$ and the polynomial denominator
gives $B_{\rm p}$. Both estimates hold for the same nonzero
physical vector, so their minimum does too.
Its membership in $Q_\ell$ and the Rayleigh principle prove
the spectral upper bound. Positivity and spectral calculus
give the inverse-norm lower bound.
\end{proof}

At fixed $\beta$, a bound $D_\ell\succeq c\kappa\ell I$ with
fixed $c>0$ cannot hold for arbitrarily large $\ell$:
it contradicts \eqref{c18v14:upper} when $c\ell>B(\beta)$.
The corresponding inverse bound $c^{-1}(\kappa\ell)^{-1}$
fails for the same reason. This concerns the original
Hamiltonian clock. It does not exclude a positive floor of
order $\kappa$. Nor does the polynomial bound rule out
particular weak-electric trajectories $\ell=\ell(\beta)$.

\subsection{Quantitative low-energy coverage failure in the reference regime}

\begin{corollary}[A bound below thirteen]\label{c18v14:thirteen}
Throughout $0\le\beta\le1/2688$,
\[
 \inf\operatorname{spec}D_\ell<13\kappa,\qquad
       \|P_\ell(\Omega F_p)\|^2\le(R-1)^2\le1/199^2 .
 \tag{V14.18}\label{c18v14:strong}
\]
\end{corollary}
\begin{proof}
Since $\pi<22/7$, one has $\alpha\le1/400$ in this interval.
For $0\le x<1$, $e^x\le(1-x)^{-1}$, whence
\[
 R\le200/199,\quad R-1\le1/199,\quad R^2\le100/99,\quad
 \mathfrak j(\beta)\le\frac{200}{2688\cdot199^2}.
\]
It follows that
\[
 B_{\rm s}(\beta)\le
       \frac{1200}{99}+\frac{16\cdot200^2}{2688\cdot199^3}<13.
 \tag{V14.19}\label{c18v14:rational}
\]
For the projection estimate, the conditional Haar mean of $F_p$
is zero and its squared Haar norm is one, so $|u_{e,p}|\le R-1$.
The coarse sigma algebra is contained in the exterior one.
The tower property therefore gives
$P_\ell(\Omega F_p)=\Omega\mathbb E_\nu[u_{e,p}\mid
\mathfrak b_\ell]$, with norm at most $R-1$.
This is an absolute projection bound; also
$\|\Omega F_p\|^2\ge R^{-1}$.
\end{proof}

For $\Lambda>0$, put
$\Pi_\Lambda=\mathbf1_{(0,\Lambda]}(K|_{\mathcal H_{\rm phys}})$,
excluding the actual vacuum. The same test also controls the
\emph{original} low-energy spectral window.

\begin{corollary}[Native lower bound on the low-window projection error]
\label{c18v14:coverage}
If $\Lambda>\kappa B(\beta)$, then
\[
 \|Q_\ell\Pi_\Lambda\|=\|\Pi_\Lambda Q_\ell\|
       \ge\sqrt{1-\kappa B(\beta)/\Lambda}.
 \tag{V14.20}\label{c18v14:coverage-bound}
\]
In particular, for $0\le\beta\le1/2688$,
\[
                    \|Q_\ell\Pi_{26\kappa}\|>1/\sqrt2 .
 \tag{V14.21}\label{c18v14:half}
\]
\end{corollary}
\begin{proof}
Normalize $v_{e,p}$ to $v\in Q_\ell$ with $v\perp\Omega$.
The spectral theorem gives
$\|(\mathbf1-\Pi_\Lambda)v\|^2
 \le\langle v,Kv\rangle/\Lambda$; its zero-energy component
vanishes. Hence $\|\Pi_\Lambda Q_\ell\|\ge\|\Pi_\Lambda v\|$
has the asserted lower bound. Taking adjoints gives equal norms.
The strict bound \eqref{c18v14:rational} proves
\eqref{c18v14:half}.
\end{proof}

At $b=0$ the discrepancy is stronger and exact. The entire
$M=3N^3$ dimensional physical $12\kappa$ plaquette eigenspace
of V11 lies in $Q_\ell$, by Haar integration in an unused
link for each plaquette. Yet
$\mathcal A_{\ell,0}=\kappa\ell\sum_\gamma L_\gamma$
has first physical energy $12\kappa\ell$ because $n\ge5$,
and $\mathcal T_{\ell,0}=0$.
Thus zero source and exact coarse equilibrium coexist with
complete loss of the lowest physical band.
The preceding estimates give a volume-uniform continuation
of low-energy omission to positive $b$, not just a free-model
observation.

\subsection{The change of direction and its precise limits}

It is not justified to make the covariance small and then
declare the discarded inverse small from the large coarse
kinetic coefficient. Actual fine-curvature fluctuations
remain in that space; increasing straight-chain length
does not by itself raise their energy.
A scale step must either retain or reconstruct the relevant
curvature channels, or prove an adequate native lower bound
for their discarded dynamics and control their return.
A charged-sector gap supplies neither alternative.

Low-window norm approximation is useful for a spectrally
faithful effective dynamics, but is not logically necessary
for every mass-gap proof. Exact energy-dependent elimination,
or a direct positive bound on the discarded block, could
still work. Likewise this is not a no-go theorem for all
blocking maps: a smeared map depending on every fine link
need not meet the unused-link hypothesis. A curvature-enriched
map would still need its actual kinetic metric, neutral-channel
control, and a proved error estimate; adding a formal observable
list does not prove these properties.

Next, derive a curvature-sensitive retained map using every
fine link, or a source-resolved bound on these discarded channels.
The conditional Poisson framework already exists. Weak-electric
neutral coercivity, physical scale normalization, cumulative
error control, and the nontrivial interacting four-dimensional
continuum construction remain unpaid.

Exact diagnostics check native character gradients, Haar moments,
unused-link geometry, spherical ranks, and rational constants.
A smooth nonproduct, non-Wilson fixture checks the conditional
energy cancellation. These checks supplement the analytic proofs;
no native numerical spectrum is claimed. V13's full diagnostics
and preservation passed before the append.

V13's body is retained byte for byte except its version title.
The single active source is now V14; all 47 other note files
are unchanged. Only \texttt{.tex} files are in
\texttt{latex\_notes}; build products are outside it.
V15 is the next companion checkpoint and V20 the next broad
literature sweep. The full physical Yang--Mills mass gap
remains unproved.
% END CYCLE 18 VERSION 14 APPEND
% BEGIN CYCLE 18 VERSION 15 APPEND
\clearpage
\section{V15: Curvature retention, native loop metric, and kinetic nonclosure}
\label{c18v15:start}

V14 shows that the straight-link map omits physical fine-curvature
fluctuations. We now make a concrete enlargement: retain \emph{every}
elementary plaquette trace, using the actual ground measure and the
original Hamiltonian clock. This retains the entire free lowest physical
band. Its exact kinetic metric, however, contains joined six-link loops.
An explicit local rational-rank certificate proves that one such metric
entry is not determined, even almost everywhere, by all elementary
traces. Thus retaining the right linear observables does not close their
nonlinear dynamics. We identify the missing source, rather than postulate
an autonomous scalar-curvature diffusion.

\subsection{Companion checkpoint and scope of the transfer}

Fresh hashes and concluding passages were checked for NS V26, SM--Einstein
V72, Shell Mixing V16, YM Parallel V5, and the four supplied Downloads
files: perturbative ESM V64, SM Source Selection V52, Native Stationarity
V54, and Exact-Action Einstein Bridge V45. All eight are unchanged since
the V10 checkpoint. Their finite-source, prepared-descriptor, and
normalization results do not supply native weak-electric YM coercivity.
Shell Mixing's flat one-loop calculation still does not bound a general
noncommuting nonperturbative return.

A newer perturbative ESM V69 was also found. We inspected the V65--V68
status passages and V69's local-source closure proof, polar transport,
and scope. Its old kinetic/potential source family misses a first-order
current direction under local operations. The transferable method is
to compute the operation on the \emph{specified} variables before closing
an effective equation. No scalar determinant, frame source, preparation
assumption, or finite-order bound is installed as a YM theorem here.

Loop joining itself is established background, not a new discovery.
For context, Shen--Smith--Zhu's primary abstract and available introductory
HTML describe a Langevin derivation of finite-rank lattice master loop
equations; their measure is a classical Wilson Gibbs law
(\url{https://arxiv.org/abs/2202.00880}). Lemoine's primary abstract concerns
finite-rank Wilson-loop identities for central actions and
$\mathrm U(N)$ (\url{https://arxiv.org/abs/2604.16252}).
Neither supplies the ground-law source estimate needed here.
The formulas below are proved directly in the present $\mathrm{SU}(2)$
normalization. This is a targeted check; the next broad sweep is V20.

\subsection{An all-link physical retained map, without a fictitious density}

Keep $N\ge5$, $\nu=\Omega^2\,dU$, $\beta=b/\kappa$, and $K=H-E_0$.
Let $M=3N^3$, and put
\[
 C(U)=(F_p(U))_{p\in\mathcal P_N},\qquad
 F_p=\Tr U_p=2w_p,\qquad
 \mu=C_\#\nu,\qquad S=C(\mathrm{SU}(2)^{E_N})\subset[-2,2]^M .
 \tag{V15.1}\label{c18v15:map}
\]
Every fine link occurs in this map. Its scalar coordinates are
gauge invariant, but need not separate all physical configurations.
In particular $S$ is not declared a Cartesian product of independent
plaquettes. Write $\mathbb E_C$ for conditioning in $\nu$ on $C$, and
define
\[
 Jf=\Omega f(C),\qquad P=JJ^*,\qquad Q=1-P .
 \tag{V15.2}\label{c18v15:projection}
\]
Here $J:L^2(\mu)\to\mathcal H_{\rm phys}$ is an isometry.
These symbols refer to the new map, not V13's straight-link projection.

For smooth functions on a neighbourhood of $[-2,2]^M$, let
\[
 \Gamma_{pq}(U)=\sum_{e,a}(X_{e,a}F_p)(X_{e,a}F_q),\qquad
 \overline\Gamma=\mathbb E_C\Gamma .
 \tag{V15.3}\label{c18v15:metric}
\]
The closed retained form is the closure, in its form norm, of
\[
 \mathfrak a(f,g)
 =\kappa\int_S \partial f(c)^*
               \overline\Gamma(c)\partial g(c)\,d\mu(c).
 \tag{V15.4}\label{c18v15:form}
\]
No smooth Lebesgue density for $\mu$, global chart, or inverse of
$\overline\Gamma$ is assumed.

\begin{proposition}[Actual retained form and captured free band]
\label{c18v15:retained}
The form \eqref{c18v15:form} is densely defined and closable.
Its closure equals the pullback of the fine ground-state form on the
form-norm closure of these pullbacks. Let $A$ be its nonnegative
self-adjoint operator. Every polynomial in $c$ is in $D(A)$, and
$Af=J^*KJf$ there.
At every $b$, the vectors $\Omega F_p$ and their vacuum-centered
versions belong to $\operatorname{Ran}P$.
At $b=0$, the entire $M$-dimensional physical $12\kappa$ band of V11
is retained, and its vectors have zero source $QKJf$.
\end{proposition}
\begin{proof}
The chain rule and the actual ground transform give
\[
 \langle Jf,KJg\rangle
 =\kappa\int \nabla(f(C))^*\nabla(g(C))\,d\nu
 =\mathfrak a(f,g).
\]
Polynomials restricted to the compact set $S$ separate its points and
contain the constants; their uniform density, and then $L^2(\mu)$
density, follow from Stone--Weierstrass. The fine form is closed.
If a pullback sequence tends to zero in $L^2(\nu)$ and is Cauchy in
energy, fine closability makes its energy tend to zero. This proves
closability and the asserted completion. A polynomial pullback times
$\Omega$ is smooth on the compact fine configuration manifold, hence
is in $D(K)$. Pairing with the form core and using continuity shows
$Af=J^*KJf$. Linear coordinate functions give the membership claims.
At $b=0$, $\Omega=1$ and $KF_p=12\kappa F_p$; V11 identifies these
orthonormal vectors as the full lowest physical band.
\end{proof}

This repairs V14's complete loss of that free band. It proves neither
coverage of a weak-electric spectral window nor autonomous evolution
of all functions of $C$.

\subsection{The native metric contains relative-colour loop information}

Two distinct elementary plaquettes share at most one link when $N\ge5$.
If $e$ is shared, cyclically orient each real trace to use $U_e$ positively:
\[
 F_p=\Tr(U_e A_p),\qquad F_q=\Tr(U_e A_q),\qquad
 G_{pq}^{e}=\Tr(A_pA_q^{-1}).
 \tag{V15.5}\label{c18v15:join}
\]
Reversing a trace does not change it in $\mathrm{SU}(2)$.
Both $A_p,A_q$ are the resulting three-link paths; $G_{pq}^{e}$
is a gauge-invariant six-link loop, with the shared edge canceled.

\begin{lemma}[Exact scalar-curvature kinetic coefficients]
\label{c18v15:metriclemma}
One has
\[
 \Gamma_{pp}=16-4F_p^2,\qquad
 \Gamma_{pq}=
 \begin{cases}
  2G_{pq}^{e}-F_pF_q,&\partial p\cap\partial q=\{e\},\\
  0,&\partial p\cap\partial q=\varnothing .
 \end{cases}
 \tag{V15.6}\label{c18v15:joinmetric}
\]
Pointwise, as quadratic forms on $\mathbb R^M$,
\[
 0\le\Gamma\le64I,\qquad 0\le\overline\Gamma\le64I .
 \tag{V15.7}\label{c18v15:metricupper}
\]
For $\beta>0$ the actual mean diagonal obeys
\[
 \int\overline\Gamma_{pp}\,d\mu
 \le \min\{16,96\sqrt2\,\beta^{-1/2}\}.
 \tag{V15.8}\label{c18v15:metricmean}
\]
\end{lemma}
\begin{proof}
Write $U_eA_p=(h_0,\mathbf h)$ and $U_eA_q=(k_0,\mathbf k)$ as
unit quaternions. The three left-generator derivatives of their traces
are $-2\mathbf h$ and $-2\mathbf k$. Their scalar product is
$4\mathbf h\cdot\mathbf k
 =2\Tr[(U_eA_p)(U_eA_q)^{-1}]-F_pF_q$, giving the off-diagonal
identity after cyclic cancellation. One-link differentiation gives
$4-F_p^2$ on the diagonal; there are four links.
Each off-diagonal entry has magnitude at most $4$ by Cauchy--Schwarz.
Each plaquette has twelve distinct other plaquettes sharing a link,
so the absolute row sum is at most $16+12\cdot4=64$.
The Gram definition gives positivity, and conditioning preserves both
quadratic-form inequalities. Finally
$\Gamma_{pp}=16(1-w_p^2)\le32(1-w_p)$.
Use V6's native bound
$\mathbb E_\nu(1-w_p)\le3\sqrt2\,\beta^{-1/2}$ and the pointwise
bound $16$.
\end{proof}

At a flat configuration $\Gamma=0$. This is not a proof of a zero
retained spectral gap: a degenerate coordinate metric at the image
boundary can still define a gapped operator. Conversely the exact
metric is not a positive constant times a collection of independent
scalar Laplacians. In particular its mean scaling is not obtained by
assigning V13's straight-link coefficient to these new coordinates.

\subsection{The complete metric-and-score source}

Let $s_{e,a}=X_{e,a}\log\Omega$ and set
\[
 h_p=2\sum_{e,a}s_{e,a}X_{e,a}F_p,\qquad
 \overline h_p=\mathbb E_C h_p,\qquad
 Z_f=\sum_{p,q}\Gamma_{pq}\partial_{pq}f(C)
            +\sum_p h_p\partial_p f(C).
 \tag{V15.9}\label{c18v15:sourceingredients}
\]
The score is the actual restored-vacuum score. In particular no
classical Wilson Gibbs drift is substituted.

\begin{theorem}[Native compressed generator and full discarded source]
\label{c18v15:source}
For smooth $f$ as above, and in particular for every polynomial,
\[
 \begin{split}
 Af(c)&=\kappa\left[
 12\sum_p c_p\partial_p f
 -\sum_{p,q}\overline\Gamma_{pq}\partial_{pq}f
 -\sum_p\overline h_p\partial_p f\right](c),\\
 Tf:=QKJf&=-\kappa\Omega(Z_f-\mathbb E_CZ_f),\\
 \|Tf\|^2&=\kappa^2\int_S
                   \operatorname{Var}_\nu(Z_f\mid C=c)\,d\mu(c).
 \end{split}
 \tag{V15.10}\label{c18v15:fullsource}
\]
For complex $f$, variance means the conditional squared modulus.
It retains metric--metric, score--score, and mixed covariances.
If $f$ depends on $m$ plaquette coordinates, then
\[
 \|Tf\|\le\kappa\left(
 16m\|D^2f\|_\infty
 +32\beta\sqrt m\,\|\partial f\|_\infty\right),
 \tag{V15.11}\label{c18v15:localsource}
\]
where the Hessian norm is the operator norm and the gradient norm is
Euclidean on those $m$ coordinates. The displayed norms are suprema
on the cube. No total-volume factor is hidden, but the support and
coupling dependence have not been removed.
\end{theorem}
\begin{proof}
The ground-transformed positive generator is
$\Omega^{-1}K\Omega=-\kappa\sum_{e,a}(X_{e,a}^2+2s_{e,a}X_{e,a})$.
Each $F_p$ has electric Laplacian $-12F_p$, hence the chain rule gives
\[
 KJf=\kappa\Omega(12C\cdot\partial f-Z_f).
\]
The first term is retained. Projection proves the first two identities;
orthogonal conditional centering proves the third.
For the local bound, the principal submatrix $\Gamma_I$ is positive
with trace at most $16m$, so
$|\Gamma_I:D^2f|\le16m\|D^2f\|_\infty$.
The inherited native bound $|\nabla_e\log\Omega|\le2\beta$ on each
of the four links of $p$, together with $|\nabla F_p|\le4$, gives
$|h_p|\le32\beta$. Use Euclidean Cauchy--Schwarz over the $m$
coordinates, the triangle inequality, and the $L^2$ contraction of
conditional centering.
\end{proof}

Even at $b=0$ the metric part of \eqref{c18v15:fullsource} can be
nonzero on nonlinear functions. The next certificate proves that
it actually is, on the original lattice, not on an auxiliary
two-matrix replacement.

\subsection{A finite native certificate of missing information}

Fix the two adjacent $xy$ plaquettes based at $(1,1,1)$ and $(2,1,1)$;
call them $p,q$. Their common edge is $e=((2,1,1),1)$, with directions
numbered $0,1,2$. Use the nonwrapping cube $\{0,1,2,3,4\}^3$, which
embeds in every $N\ge5$ torus, and vary only
\[
 B=\{(x,d):x\in\{1,2,3\}^3,\ d\in\{0,1,2\},\ x_d<3\}.
 \tag{V15.12}\label{c18v15:B}
\]
There are $54$ variable links, hence $162$ left tangent coordinates.
Exactly $108$ elementary plaquettes meet $B$; their union uses $180$
links. No other elementary trace varies when these $54$ links vary.

\begin{lemma}[Rational local rank certificate]
\label{c18v15:rank}
There is an explicitly specified rational unit-quaternion configuration
on these $180$ links such that, with exterior links fixed,
\[
 \operatorname{rank}D_B(F_r:r\cap B\ne\varnothing)=108,\qquad
 \operatorname{rank}D_B(F_r:r\cap B\ne\varnothing,\Gamma_{pq})=109.
 \tag{V15.13}\label{c18v15:ranks}
\]
These ranks are certified by nonzero rational minors, not by a
floating-point singular-value threshold.
\end{lemma}
\begin{proof}[Proof by an exact finite certificate]
Here is the complete data recipe and the two checked minors.
For a link $(x,d)$, form the ASCII string
$x_0,x_1,x_2,d$ with decimal digits, commas, and no spaces.
Let $a_0,a_1,a_2$ be the first three bytes of its SHA--256 digest.
This merely selects integers; no cryptographic property is used.
Put
\[
 v_j=(a_j\bmod7)-3,\qquad
 U_{x,d}=\frac{(1-|v|^2,\,2v_0,\,2v_1,\,2v_2)}{1+|v|^2}.
 \tag{V15.14}\label{c18v15:rationalpoint}
\]
The denominator is at most $28$, and the quaternion norm is exactly
one. Sort $B$ lexicographically by $(x_0,x_1,x_2,d)$ and then list
the generators $\mathbf i,\mathbf j,\mathbf k$ on each link.
Columns are numbered $0,\ldots,161$.
Rows are the $108$ incident plaquettes ordered lexicographically
by $(x_0,x_1,x_2,i,j)$, $i<j$, with word
\[
 U_{x,i}U_{x+\hat i,j}U_{x+\hat j,i}^{-1}U_{x,j}^{-1}.
\]
Their differential matrix is $D$. Append $D_B\Gamma_{pq}$ as its
last row to form $\widetilde D$.
A first derivative replaces a positive occurrence of $U$ by $E_aU$,
or a negative occurrence by $-U^{-1}E_a$, and takes the trace.
The appended row can be computed from \eqref{c18v15:joinmetric},
or independently by differentiating its gradient scalar product.

Let $J_0=\{0,\ldots,161\}\setminus R_0$, where
\[
\begin{split}
R_0=\{&
14,17,23,31,32,35,38,40,41,44,47,53,56,59,62,\\
&70,71,76,77,79,80,85,86,92,93,94,95,98,100,101,\\
&102,103,104,107,109,110,115,116,119,121,122,125,\\
&131,134,137,140,145,146,149,151,152,155,158,161\}.
\end{split}
\]
Columns in a minor are in ascending order. Exact elimination gives
\[
 \begin{array}{c|cc}
 \text{prime}&\det D_{J_0}&
            \det\widetilde D_{J_0\cup\{14\}}\\ \hline
 1009&425&526\\
 1013&468&232
 \end{array}
 \tag{V15.15}\label{c18v15:minors}
\]
All entries are rational with denominators invertible at either
prime. A nonzero modular determinant therefore proves that its
rational determinant is nonzero. Row counts give the matching
upper rank bounds, proving \eqref{c18v15:ranks}.

The reproducible driver
\path{scripts/verify_ym_c18_v15_curvature_map.py}
generates the actual incidence words and every matrix entry from
\eqref{c18v15:rationalpoint}. Both selected determinants are also
checked by a separate scalar elimination. An independent $2$-by-$2$
matrix implementation checks all $648$ supported first-derivative
entries and all $162$ appended derivative entries, the latter directly
from second derivatives of the two plaquette traces.
Over either field it uses
\[
 (w,x,y,z)\longmapsto
 \begin{pmatrix}w+\iota x&y+\iota z\\-y+\iota z&w-\iota x\end{pmatrix},
 \qquad \iota^2=-1,
\]
with $\iota=469$ modulo $1009$ and $\iota=45$ modulo $1013$.
These are exact algebraic checks, not a spectral simulation.
\end{proof}

\begin{theorem}[Elementary traces do not determine the native metric]
\label{c18v15:nonmeasurable}
For every $N\ge5$ and every finite $b\ge0$,
\[
 \delta_{\beta,N}^2:=
 \int\operatorname{Var}_\nu(\Gamma_{pq}\mid C)\,d\nu>0 .
 \tag{V15.16}\label{c18v15:positivevariance}
\]
No volume-uniform or weak-electric lower bound on this positive
number is asserted.
\end{theorem}
\begin{proof}
The certificate lies in a fixed nonwrapping patch. Assign arbitrary
unit quaternions to links outside its $180$-link collar.
The two nonzero minors remain nonzero in an open neighbourhood.
With exterior variables as parameters, the inverse function theorem
gives local coordinates consisting of the $108$ affected traces,
$\Gamma_{pq}$, and $53$ remaining coordinates on the $54$ variable
links. The $109$-minor supplies the required independent directions.
Thus $\Gamma_{pq}$ ranges over an interval while the affected traces
and exterior links are fixed. All unaffected traces depend only
on those exterior links.

The actual $\nu$ has a smooth strictly positive density relative to
product Haar at each finite regulator. In the displayed local
coordinates its density is likewise positive. By Fubini, no function
of the affected traces and exterior links can equal the independent
$\Gamma_{pq}$ coordinate almost everywhere on this neighbourhood.
In particular no measurable function of all of $C$ can do so.
Vanishing conditional variance would assert precisely that equality
almost everywhere, a contradiction. The argument is insensitive to
the positive density's magnitude and so applies to every finite $b$,
but gives no claimed uniform numerical lower bound.
\end{proof}

Two configurations with the same traces at a singular fibre would
not suffice for this theorem. The full-row-rank certificate and the
positive-density neighbourhood are essential. All observables and
metric entries in this argument are original-Gauss invariant;
a conditional independent physical tensor factor is not introduced.

\subsection{A nonzero source at zero magnetic coupling, and the next estimate}

Use the notation
$T_p=QK(\Omega F_p)$, $T_q=QK(\Omega F_q)$, and
$T_{pq}=QK(\Omega F_pF_q)$.
Multiplication by $F_p(C)$ commutes with conditional projection.

\begin{corollary}[The missing six-link channel has an actual source]
\label{c18v15:nonclosure}
For every finite $b\ge0$,
\[
 T_{pq}-F_pT_q-F_qT_p
 =-2\kappa\Omega(\Gamma_{pq}-\mathbb E_C\Gamma_{pq}),
 \tag{V15.17}\label{c18v15:productsource}
\]
and consequently
\[
 2\kappa\delta_{\beta,N}
 \le \|T_{pq}\|+2\|T_p\|+2\|T_q\|.
 \tag{V15.18}\label{c18v15:sourcepositive}
\]
In particular these three sources cannot all vanish.
At $b=0$ one has $T_p=T_q=0$ but
\[
 \|T_{pq}\|=2\kappa\delta_{0,N}>0 .
 \tag{V15.19}\label{c18v15:electricmemory}
\]
The full retained algebra is therefore not invariant under the
actual fine generator, even though it contains every vector in the
free lowest physical band.
\end{corollary}
\begin{proof}
The ground-transformed product rule gives
$\mathcal K(F_pF_q)
 =F_p\mathcal KF_q+F_q\mathcal KF_p-2\kappa\Gamma_{pq}$,
where $\mathcal K=\Omega^{-1}K\Omega$. Apply $Q$ after multiplying
by $\Omega$ to obtain \eqref{c18v15:productsource}.
Use $|F_p|,|F_q|\le2$ and \eqref{c18v15:positivevariance}.
At $b=0$ the individual coordinates are retained electric
eigenvectors, proving the exact norm statement.
\end{proof}

\subsection{A first kinetic closure repair, with its interacting cost}

Let $Y$ consist of all $F_p$ together with all the $G_{pq}^{e}$ for
adjacent plaquettes. Define $J_+f=\Omega f(Y)$ and its projections
$P_+,Q_+$ as in \eqref{c18v15:projection}. This is an observable
enlargement, not an independent-loop measure. On fine configuration
functions put $\mathcal L=\sum_eL_e$, and define
\[
 \mathcal W_2=\operatorname{span}\{1,F_p,F_pF_q,G_{pq}^{e}\}.
 \tag{V15.20}\label{c18v15:W2}
\]
The products include $p=q$; the joined loops are included only for
distinct adjacent plaquettes. This is a finite linear space, not
the full algebra of functions of $Y$.

\begin{proposition}[Quadratic plaquette closure and a score-only return]
\label{c18v15:repair}
One has $\mathcal L\mathcal W_2\subset\mathcal W_2$, with
\[
 \begin{array}{c|c}
 f&\mathcal Lf\\ \hline
 F_p&12F_p\\
 G_{pq}^{e}&18G_{pq}^{e}\\
 F_p^2&32(F_p^2-1)\\
 F_pF_q,\ \partial p\cap\partial q=\varnothing&24F_pF_q\\
 F_pF_q,\ \partial p\cap\partial q=\{e\}
     &26F_pF_q-4G_{pq}^{e}
 \end{array}
 \tag{V15.21}\label{c18v15:closuretable}
\]
For every $f\in\mathcal W_2$ at actual coupling $b$,
\[
 Q_+K(\Omega f)
 =-2\kappa\Omega(1-\mathbb E_Y)(s\cdot\nabla f).
 \tag{V15.22}\label{c18v15:repairedsource}
\]
In particular the enlarged projection has zero source on
$\mathcal W_2$ at $b=0$.
If $f$ depends on $m$ fine links and $A_*=3\sqrt2$, then
\[
 \begin{split}
 \|Q_+K(\Omega f)\|
 &\le4b\sqrt m\,\|\nabla f\|_{L^2(\nu)},\\
 \|Q_+K(\Omega f)\|
 &\le2\kappa\sqrt{mA_*}\,\beta^{1/4}\|\nabla f\|_\infty .
 \end{split}
 \tag{V15.23}\label{c18v15:repairedcost}
\]
The second bound uses V6's mean kinetic estimate.
Neither bound is a weak-electric contraction theorem.
\end{proposition}
\begin{proof}
The six edges of a joined elementary pair are distinct. Its trace
is fundamental in each, so $\mathcal LG=18G$.
The positive Laplacian product rule is
$\mathcal L(fg)=f\mathcal Lg+g\mathcal Lf-2\nabla f\cdot\nabla g$.
Apply \eqref{c18v15:joinmetric} to obtain every remaining row.
For example, the adjacent row is
$24F_pF_q-2(2G_{pq}^{e}-F_pF_q)$.
Thus $\mathcal Lf$ remains a function of $Y$ for $f\in\mathcal W_2$.
The ground transform is
$K(\Omega f)=\kappa\Omega(\mathcal Lf-2s\cdot\nabla f)$;
projecting gives \eqref{c18v15:repairedsource}.

Conditional centering is an $L^2$ contraction.
The pointwise native score bound gives
$|s\cdot\nabla f|\le2\beta\sum_{e\in\operatorname{supp}f}
|\nabla_e f|$, proving the first estimate.
For the second, use
$|s\cdot\nabla f|^2\le
(\sum_{e\in\operatorname{supp}f}|s_e|^2)|\nabla f|^2$
and the inherited actual mean
$\int|s_e|^2\,d\nu=\langle\Omega,L_e\Omega\rangle
 \le A_*\sqrt\beta$.
\end{proof}

The $18\kappa$ and $26\kappa$ fusion energies in this table were
already computed in V11; they are not new spectral coefficients.
Indeed $G_{pq}^{e}$ has energy $18\kappa$ and
$F_pF_q-\tfrac12G_{pq}^{e}$ has energy $26\kappa$.
The new use is a specified global retained map for which the full
quadratic-plaquette space has no free kinetic leakage. At $b>0$
its complete remaining source is \eqref{c18v15:repairedsource},
not an omitted magnetic or metric term.

This gives both a repair and its precise cost. Keeping $G_{pq}^{e}$
removes the proved elementary-trace metric uncertainty and closes
the displayed finite linear space under the electric generator.
It does not prove closure of all functions of $Y$, or justify
replacing their actual conditional measure by a product. Merely
listing successively longer loops does not give a convergent
effective theory.

The next estimate should control the \emph{native resolvent-weighted
return} of these joined-loop and score sources, or a normed retained
loop family whose source tails can actually be bounded. It must keep
the mixed covariance in \eqref{c18v15:fullsource}, the singular image
measure, and the physical clock. The archived conditional
Poisson/Schur identities are already available; their formal statement
is not another missing theorem. The required weak-electric neutral
coercivity, scale trajectory, cumulative error bound, and interacting
four-dimensional continuum construction remain unpaid.

V14's full diagnostics passed before this append. Its body is retained
byte for byte apart from the version title; the other $47$ note files
are unchanged. The single active source is V15. Exact checks include
the native incidence, two modular certificates, independent derivative
and determinant implementations, and sixteen rational loop-joining
tests. These support the computer-assisted finite rank proof; they
do not certify a mass gap or constitute proof-assistant formalization.
Only \texttt{.tex} files are kept in \texttt{latex\_notes}; build products
are outside it. The next companion checkpoint and broad sweep are V20.
The full physical Yang--Mills mass gap remains unproved.

% BEGIN CYCLE 18 VERSION 16 APPEND
\clearpage
\section{V16: Retained-map rank, centre fibres, and weak-field conditioning}
\label{c18v16:start}

The intended next step was to estimate V15's score return. Before
inserting a conditional inverse, we must know what the enlarged
observable map actually discards. A native rank calculation changes
that interpretation. With exterior links fixed, the map retains
every infinitesimal physical variation in a fixed patch, almost
everywhere. On the side-five torus it has the full $750$-dimensional
physical differential rank almost everywhere. This is not yet a
spatial coarse graining: locally it can be a re-expression of the
fine physical variables. It also loses global centre information.
A separate native estimate below proves deterioration of its
unscaled inverse coordinates at large $\beta$, uniformly in volume
for the fixed patch.

The distinction between continuous loop coordinates and discrete
data is established background. Watson's primary abstracts,
\emph{Solution of the SU(2) Mandelstam Constraints} and
\emph{Gauge-Invariant Variables and Mandelstam Constraints in SU(2)
Gauge Theory}, explicitly retain discrete information alongside
continuous invariants
(\url{https://arxiv.org/abs/hep-th/9311126},
\url{https://arxiv.org/abs/hep-th/9408174}).
Only those abstracts were used; no quantitative theorem is imported.
More importantly, f16.2 V86 already proves that specified rooted
single, pair and triple traces separate every compact gauge orbit.
That theorem, its centre-sector warning, and f17 V96--V97's
conditional source/resolvent architecture are inherited.
The new question here concerns the particular short-loop map of
V15, not a new claim that Wilson loops can coordinatize gauge orbits.

\subsection{The precise map and the two finite certificates}

Keep the original finite Hamiltonian, $\nu=\Omega^2\,dU$,
$\beta=b/\kappa$, and $F_p=\Tr U_p$.
Let $Y$ be V15's complete list of elementary $F_p$ and of joined
$G_{pq}^e=\Tr(A_pA_q^{-1})$ for every unordered pair of elementary
plaquettes sharing a link. Rows are labelled even if a functional
dependence exists. For $N\ge5$ there are
\[
 |\mathcal P_N|=3N^3=:M,\qquad
 \#\{G_{pq}^e\}=6M,\qquad \#Y=7M=21N^3 .
 \tag{V16.1}\label{c18v16:counts}
\]
Each link occurs in $4$ elementary and $36$ joined rows, hence in
$40$ rows in total. Each joined word uses six distinct links.

For the local test take the $54$ links $B$ of
\eqref{c18v15:B}. Let $Y_B$ contain all elementary traces meeting
$B$, and all joined words that contain a link of $B$. It has
$108+816=924$ rows and uses $402$ links.
The elementary parents of these rows form a set $\mathcal P_B^*$
of size $288$. All their word vertices lie in
$\{-1,0,\ldots,5\}^3$, so this whole test embeds without vertex
identification in every $N\ge7$ torus. Exterior links are held
fixed when taking $D_BY_B$.

The central vertex $v=(2,2,2)$ has all six incident links in $B$.
Its gauge group acts freely on these variables. Fixing the outgoing
edge $e_*=((2,2,2),0)$ supplies a $159$-dimensional differential
slice: remove its three tangent columns from the $162$ columns
of $B$. The group action, not a numerical rank tolerance, accounts
for these three removed directions.

For the global test set $N=5$. There are $375$ links.
Choose root $(0,0,0)$ and the following $124$-edge tree.
For each nonroot vertex $x$, choose the smallest $d$ with $x_d>0$
and include the edge from $x-\hat d$ to $x$.
The sum of the coordinates decreases along parents, so this is
indeed a spanning tree. It leaves $251$ chords.
The generic physical dimension is
\[
 3\cdot375-3\cdot125=750 .
 \tag{V16.2}\label{c18v16:dimension}
\]

\begin{lemma}[Two native nonzero-minor certificates]
\label{c18v16:certificates}
There are specified rational unit-quaternion configurations such that
the local differential slice has rank $159$, and the global
side-five differential on the physical directions has rank $750$.
The corresponding square minors are
\[
 \begin{array}{c|c|cc}
 \text{test}&\text{minor size}&\bmod1009&\bmod1013\\ \hline
 \text{local fixed-exterior slice}&159&695&978\\
 \text{side-five tree slice}&750&115&689
 \end{array}
 \tag{V16.3}\label{c18v16:determinants}
\]
They certify rank over the rationals, not merely approximate rank
at floating-point sample points.
\end{lemma}
\begin{proof}[Proof by explicit exact certificates]
For the local configuration use V15's full rational recipe
\eqref{c18v15:rationalpoint} on the $402$ links. Negative coordinates
in its ASCII string include the minus sign and no spaces.
Order links by $(x_0,x_1,x_2,d)$ and the three left generators by
$\mathbf i,\mathbf j,\mathbf k$. Delete the columns of $e_*$.
List the $108$ elementary rows in lexicographic
$(x_0,x_1,x_2,i,j)$ order, $i<j$. Next list the $816$ joined rows
by the ordered parent pair and common edge, with the smaller
parent first. The parents can be enumerated on bases
$\{-1,\ldots,4\}^3$; discard joined words that contain no link
of $B$. These conventions determine the $924$-row matrix exactly.

For the global point set tree links to identity. On all chords
except the first two use the same rational recipe, now with
coordinates $0,\ldots,4$. In sorted order the first two chords are
\[
 e_0=((0,0,4),2),\qquad e_1=((0,1,0),2).
\]
Assign
\[
 U_{e_0}=(3/5,4/5,0,0),\qquad
 U_{e_1}=(3/5,0,4/5,0).
 \tag{V16.4}\label{c18v16:globalpoint}
\]
Their common centralizer is $\{\pm I\}$.
Take all three left tangent columns of each chord except
$(e_0,\mathbf j),(e_0,\mathbf k),(e_1,\mathbf k)$.
This gives $750$ columns. List first the $375$ elementary rows,
then the $2250$ joined rows, using the same lexicographic
parent-pair convention, with torus coordinates reduced modulo five.

Entries are obtained by differentiating the literal loop words:
replace a positive occurrence of $U_e$ by $E_aU_e$, or a negative
occurrence by $-U_e^{-1}E_a$, and take the trace.
For full reproducibility, the selected rows are specified below
as hexadecimal integers. Bit $j$, counted from the least significant
bit starting at zero, is one exactly when row $j$ is selected.
Concatenate consecutive displayed lines in their printed order.
The local mask is
\begin{quote}\footnotesize\ttfamily
2000000104004100104004000004000100104000000a00240048010282404280\\
2080208080802100420020404ffffffffffffffffffffffffffffff
\end{quote}
and the side-five mask is
\begin{quote}\footnotesize\ttfamily
10000400010000400002404400412010480412008090041100302200644c00c8\\
9800c43ee8810408204101010110010400410010400202004400410010400410\\
00808011201048041201048020240104400c0880191300322600310fb8108008\\
1101022020450209080940c02140c02140c02140d0104be88948c82158ca21d8\\
ca21d8d9107fb88948c82158ca21d8ca21d8d9107fbcdb15476b76eefb5bbf73\\
daddbbbfdb7fffffffffffffffffffffffffffffffffffffffffffffffffffff\\
fffffffffffffffffffffffffffffffffffffffffffff
\end{quote}
Selected rows and columns are in ascending order.

The driver
\path{scripts/verify_ym_c18_v16_retained_rank.py}
generates these matrices and verifies
\eqref{c18v16:determinants}. A separate $2$-by-$2$ matrix
implementation reconstructs the selected native derivatives:
$1032$ supported derivative terms locally and $6203$ globally
at each prime. Independent column elimination on those matrices
rechecks the determinants obtained by row elimination.
The matrix representation and square roots of $-1$ are those
displayed in V15. Every operation is integer modular arithmetic;
all rational coordinate denominators are at most $28$ and
invertible at both primes. A nonzero reduced minor therefore
proves its rational minor is nonzero.
The row masks select respectively $159$ and $750$ rows, so
the displayed nonzero determinants prove the asserted ranks.
No spectral quantity is computed by this certificate.
\end{proof}

\subsection{What the ranks prove in the actual ground measure}

\begin{theorem}[Local physical information is retained almost everywhere]
\label{c18v16:localrank}
For every $N\ge7$ and every finite $b\ge0$, almost everywhere in
the actual $\nu$,
\[
 \operatorname{rank}D_BY_B=159,\qquad
 \ker D_BY_B=\mathcal G_B,
 \tag{V16.5}\label{c18v16:localrankeq}
\]
where $\mathcal G_B$ is the three-dimensional infinitesimal gauge
action at the central vertex, with exterior links fixed.
The statement holds simultaneously for all translates of this
patch on a fixed finite torus. Locally on its regular set, $159$
selected entries of $Y_B$ are coordinates for the fixed-exterior
physical slice.
\end{theorem}
\begin{proof}
For $\xi\in\mathfrak{su}(2)$ the central gauge action differentiates
outgoing links as $\xi U_e$ and incoming links as $-U_e\xi$.
It is supported in $B$, annihilates every loop trace, and is
injective as a map from $\xi$, already on $e_*$. Thus the kernel
has dimension at least three at every configuration. Fixing $e_*$
gives a complementary tangent slice, since a nonzero such gauge
variation cannot vanish on that edge.

The $159$-minor of Lemma~\ref{c18v16:certificates} is a
nonzero real-analytic function of the collar links. Its zero set
has Haar measure zero. One elementary way to see this here is
to use stereographic coordinates on each sphere, excluding a
single Haar-null point. After multiplication by positive powers
of the coordinate denominators the minor is a nonzero polynomial.
A nonzero polynomial has a null zero set by induction on the
number of variables and Fubini, using the finite zero set of a
nonzero one-variable polynomial. The rational certificate proves
that this polynomial is not identically zero.

The actual $\nu$ is equivalent to product Haar at every fixed
regulator because $\Omega$ is smooth and strictly positive.
Hence the minor is nonzero $\nu$-almost everywhere. The lower
rank $159$ and the gauge upper rank $162-3$ agree. The inverse
function theorem on the fixed-edge gauge slice proves the
local coordinate assertion. Only finitely many translated patches
are involved at any fixed $N$, so intersecting their full-measure
sets proves the simultaneous assertion.
\end{proof}

\begin{theorem}[Full physical differential rank on the side-five torus]
\label{c18v16:globalrank}
At $N=5$, for every finite $b\ge0$, the full map $Y$ has physical
differential rank $750$ almost everywhere in $\nu$.
On the principal gauge stratum its differential kernel is exactly
the tangent to the original gauge orbit. Thus its regular local
physical fibres are zero dimensional.
\end{theorem}
\begin{proof}
The two noncommuting chord values in
\eqref{c18v16:globalpoint} make the configuration irreducible.
On the full-measure irreducible stratum the original gauge orbit
has dimension $3\cdot125=375$: tree transport leaves only the
common centralizer of the chord holonomies, whose Lie algebra is
zero. Reducibility is Haar-null, for instance by testing whether
two transported chord imaginary parts commute. The corresponding
nonzero polynomial test fails only on a null set.

The global certificate is a nonzero minor of the differential of
the full native trace map evaluated on $750$ actual link tangent
columns. Although tree links were fixed to exhibit the point,
this minor is a real-analytic function on the full configuration
space. The same stereographic polynomial argument makes it
nonzero almost everywhere, also in $\nu$.
Gauge invariance gives the upper rank $1125-375=750$ on the
principal stratum. Equality forces the claimed kernel.
A local gauge slice and the inverse function theorem give the
last statement.
\end{proof}

This is a theorem at $N=5$, not a proof of full global rank for
arbitrary $N$. The local patch theorem has the stated uniform
embedding range $N\ge7$. Neither result gives a uniform inverse
bound, a global inverse, a finite number of all fibre branches,
or a lower bound on a physical Hamiltonian.

\subsection{The missing centre data are not removed by local rank}

For each direction $d$, let $Z_d$ change the sign of every
direction-$d$ link crossing the periodic cut $x_d=N-1$.
These three commuting involutions preserve every elementary
and joined contractible loop. They preserve product Haar and
the electric metric, commute with $H$, and fix $\Omega$ by
uniqueness of the normalized positive ground. These symmetry
facts, and their distinction from gauge equivalence, were already
recorded in f16.2 V86 and YM Parallel V5.

\begin{proposition}[The short-loop map still identifies distinct sectors]
\label{c18v16:center}
For every $N\ge5$, $Y(Z_dU)=Y(U)$.
At a generic configuration, its eight centre translates give
eight distinct gauge orbits with exactly the same $Y$.
Consequently full local physical rank cannot make $Y$ globally
injective. The range of $J_+$ is centre neutral, and every
nontrivial centre-character subspace is contained in
$\operatorname{Ran}Q_+$.
\end{proposition}
\begin{proof}
Every elementary loop crosses a cut an even number of times;
the same holds after joining two of them and canceling their
common edge. Thus all components of $Y$ are unchanged.
For each direction choose a fundamental winding loop $W_d$.
Its trace changes sign under $Z_d$ and not under the other
two generators. Generically these three traces are all nonzero:
each is a nonzero analytic function and its zero set is Haar-null.
Any two different centre translates are then distinguished by
at least one of these gauge-invariant traces. They cannot be
gauge equivalent. The equal-\(Y\) statement is exact.

Finally $J_+f=\Omega f(Y)$ is centre invariant because both factors
are. Orthogonality of the characters of $\mathbb Z_2^3$ proves
the assertion about $Q_+$.
\end{proof}

No estimate of the energies of those winding sectors follows.
There may be additional global ambiguities. In particular we do
not infer that the $Q_+$ dynamics is merely an eight-state
problem, or that a neutral-sector restriction already proves
the required infinite-volume theory.

\subsection{A native weak-field conditioning estimate}

For a loop row $\gamma$ write
$d_\gamma=1-\tfrac12\Tr U_\gamma$.
For elementary parents abbreviate
\[
 \delta=\max_{p\in\mathcal P_N}d_p,\qquad
 \delta_B=\max_{p\in\mathcal P_B^*}d_p .
\]
The input tangent norm is the original product of unit-round
link metrics. The output norm is Euclidean on the labelled
trace list; no $\beta$-dependent coordinate rescaling is made.

\begin{lemma}[Derivative collapse for the specified short-loop map]
\label{c18v16:collapse}
For $N\ge5$ one has
\[
 \|DY(U)\|_{\rm op}^2\le7680\,\delta(U),\qquad
 \|D_BY_B(U)\|_{\rm op}^2\le7680\,\delta_B(U).
 \tag{V16.6}\label{c18v16:collapsebounds}
\]
The second statement uses the local bank when $N\ge7$.
Moreover
\[
 \begin{gathered}
 \|DY(U)\|_{\rm HS}^2\le1184\sum_{p}d_p(U),\\
 \int\|DY\|_{\rm HS}^2\,d\nu
 \le M\min\{160,1184A_*\beta^{-1/2}\},
 \qquad A_*=3\sqrt2,\quad\beta>0.
 \end{gathered}
 \tag{V16.7}\label{c18v16:meanmetric}
\]
\end{lemma}
\begin{proof}
For a simple $r$-link fundamental trace $F_\gamma$, each link
gradient has squared norm
$4-F_\gamma^2=4(1-w_\gamma^2)\le8d_\gamma$.
For a tangent $z$,
\[
 |D F_\gamma z|^2
 \le8r\,d_\gamma\sum_{e\in\gamma}|z_e|^2.
\]
A joined loop is the product of its two parent holonomies after
a common transport and a possible inverse. The chordal triangle
inequality in unit quaternions gives
$d_\gamma\le2(d_p+d_q)\le4\delta$; the local version uses
$\delta_B$. Since $r\le6$ and every link is used by $40$ rows,
summing proves $8\cdot6\cdot4\cdot40=7680$ in
\eqref{c18v16:collapsebounds}. Restriction to the local rows and
input columns only decreases the incidence bound.

For the Hilbert--Schmidt norm sum the individual row energies.
The elementary rows contribute at most $32\sum_pd_p$.
Each joined row contributes at most $48\cdot2(d_p+d_q)$.
Every parent is in twelve adjacent unordered pairs, so these
give at most $1152\sum_pd_p$. This proves the first inequality.
The independent pointwise bound $4r$ per row gives
$4(4M+6\cdot6M)=160M$.
Finally use the inherited actual mean
$\mathbb E_\nu d_p\le A_*\beta^{-1/2}$ from V6.
There is no independence assumption.
\end{proof}

On the full-measure regular set of
Theorem~\ref{c18v16:localrank}, restrict $D_BY_B$ to
$\mathcal G_B^\perp$. This is injective, and its inverse on
its range will be denoted $\mathcal I_B$. The orthogonal
complement and its norm use the original fine metric.

\begin{theorem}[Volume-uniform deterioration of the local inverse]
\label{c18v16:inverse}
For every $N\ge7$, $\beta>0$, and $t>1$,
\[
 \nu\left\{
 \|\mathcal I_B\|\ge
 \frac{\beta^{1/4}}{\sqrt{7680\,t\,288\,A_*}}
 \right\}\ge1-\frac1t .
 \tag{V16.8}\label{c18v16:inverseprob}
\]
This is a lower bound on an inverse-coordinate norm, not an
upper bound on a source, a spectral inverse, or a mass.
\end{theorem}
\begin{proof}
Markov's inequality and the $288$ actual parent plaquettes give
\[
 \nu\{\delta_B>\varepsilon\}
 \le\frac{\sum_{p\in\mathcal P_B^*}\mathbb E_\nu d_p}{\varepsilon}
 \le\frac{288A_*}{\varepsilon\sqrt\beta}.
\]
Set $\varepsilon=t\,288A_*/\sqrt\beta$.
On the complementary event the derivative norm is at most
$\sqrt{7680\varepsilon}$ by
\eqref{c18v16:collapsebounds}. For an injective linear map,
its inverse on the range has norm at least the reciprocal of
the map's operator norm. The regular set has probability one,
so this gives \eqref{c18v16:inverseprob}. The set of parent
plaquettes and every constant are independent of $N$.
\end{proof}

At a flat configuration every component of $DY$ vanishes.
The theorem quantifies the conditioning problem on a high-probability
local event in the actual weak-electric ground law, rather than
deducing it only from that exceptional flat point. Its large
constant is not claimed optimal. It does not prevent a gap for
the correctly weighted retained form: the coordinate distribution
and metric can both degenerate together. Changing coordinates
also changes that form and its image measure; it is not a
license to change physical time.

\subsection{The upstream decision after the rank test}

V15's map has useful exact quadratic closure, but the two
rank certificates show why that success must not be read as
a controlled spatial elimination. On the stated regular
fixed-exterior patch it leaves no infinitesimal physical
direction to eliminate, and at $N=5$ it retains all continuous
physical dimensions. Meanwhile its local inverse becomes
ill-conditioned and its global centre ambiguity persists.
Thus replacing the return estimate by a positive gap of a
postulated smooth, uniformly elliptic continuous fibre would
be unjustified.

The next noncircular step must specify a genuine scale-dependent
reduction, or work directly with the original physical form,
and pay the low-energy coverage and native return of the
directions it actually removes. The full-rank short-loop map
can be used as a local coordinate diagnostic, not as an
already proved renormalization step. The original weak-electric
neutral coercivity and physically normalized scale evolution
are still missing. No uniform conditional spectral inverse,
interacting continuum construction, or full YM mass gap has
been obtained in this append.

The previous goal turn was progress: V15 produced its new native
source and rank certificate. V15's complete mathematical body
is preserved here, apart from the version title. Its full
diagnostic passed before the append. The new rank tests use
two primes, an independent matrix-derivative implementation,
and independent determinant eliminations; the incidence and
centre-cut checks also run on sides five, six, and seven.
All $47$ other note files are unchanged. Only the active main
version is replaced, and only \texttt{.tex} files are kept in
\texttt{latex\_notes}. Build products are outside that folder.
The next companion checkpoint and broad literature sweep remain V20.

% BEGIN CYCLE 18 VERSION 17 APPEND
\clearpage
\section{V17: Flow before sampling and a Gaussian discarded-space floor}
\label{c18v17:section}

V16 found that the enlarged short-loop map can retain all continuous
physical information locally, so it is not yet a spatial reduction.
We now specify a reduction: apply the full simultaneous spatial Wilson
flow, and only then take V13's sparse straight-link products. The
compact map is a genuine submersion with positive-dimensional physical
fibres. The substantive new estimate is an all-block-size sampling
inequality for its flat transverse differential. It gives a
volume-uniform conditional spectral floor, and a low-energy coverage
bound on the \emph{full Gaussian oscillator function space}. It does
not give those bounds for the actual interacting ground law.

The order and the distinction matter. Smoothing without sampling is
invertible, as already proved in f16.2 V86. Unflowed sparse blocking
misses the lowest strong-electric physical band by V14. Neither fact
decides what flow followed by sampling does in the weak-field regime.
The calculation below tests precisely that candidate, with the
Hamiltonian coefficients and time normalization explicit.

\subsection{Inherited architecture and the compact candidate}

The f11 V4 and V6 constrained-alias calculation already supplies the
low-alias transverse inverse and a dyadic high-alias extension in four
Euclidean dimensions. We reuse that method, not claim a new general
Schur identity. Here the geometry is the three-dimensional spatial
slice, the block length is arbitrary, the multiplier includes flow, and
the conditional quantum precision is proportional to
$\sqrt{d_1^*d_1}$, not to $d_1^*d_1$. The f17 V7 depth-matched flowed
score is also different: it captures one reference covariance without
sampling away variables. Its discrete-time Gaussian precision is not
substituted for the continuous-Hamiltonian precision used here.

For external context, Sonoda--Suzuki,
\emph{Gradient flow exact renormalization group},
\href{https://doi.org/10.1093/ptep/ptab006}{PTEP 2021, 023B05},
Section 4, combines lattice link blocking and a gauge-covariant
diffusion construction. That construction contains additional
functional differential operators. We do not identify our deterministic
ground-law pushforward with their Wilson action, and import no spectral
or continuum theorem. The flow and blocking architecture is established;
the bound below is proved directly for the specified map.

Keep $N=\ell n$, $\ell\ge2$, $n\ge5$, and set
\[
 E_{\rm sp}(U)=\sum_p(1-w_p(U)),\qquad
 \partial_s\Phi_s(U)=-\nabla E_{\rm sp}(\Phi_s(U)),\qquad
 \mathfrak c_{\ell,s}=\mathfrak b_\ell\circ\Phi_s .
 \tag{V17.1}\label{c18v17:map}
\]
The metric is the original unit-round product metric. The flow time
$s$ is dimensionless spatial smoothing time; it is not physical
Hamiltonian time. In particular the choice below is $s=\ell^2$,
or physical spatial flow time $(a\ell)^2$ at lattice spacing $a$.
No factor of $b/\kappa$ is inserted into this definition.

\begin{proposition}[An actual reduction, with its native metric retained]
\label{c18v17:compact}
At every finite regulator and finite $s\ge0$, $\mathfrak c_{\ell,s}$ is a
smooth surjective submersion onto the coarse link product. It is
equivariant under restriction of the fine gauge transformation to
coarse vertices. Its ordinary fibres have dimension $9(N^3-n^3)$.
At points where both fine and coarse gauge orbits are principal, the
induced physical fibre dimension is $6(N^3-n^3)$.

For the actual law $\nu=\Omega^2dU$, its pushforward has a smooth
strictly positive coarse Haar density. Write
\[
 \mathsf G_{\ell,s}(U)
   =D\mathfrak c_{\ell,s}(U)D\mathfrak c_{\ell,s}(U)^* .
\]
Then
\[
 \ell e^{-32s}I\preceq\mathsf G_{\ell,s}(U)
                    \preceq\ell e^{32s}I.
 \tag{V17.2}\label{c18v17:metric}
\]
For $J_{\ell,s}f=\Omega(f\circ\mathfrak c_{\ell,s})$,
the exact retained form is
\[
 \langle J_{\ell,s}f,KJ_{\ell,s}f\rangle
  =\kappa\int
     (\nabla f)^*
       \mathbb E_\nu[\mathsf G_{\ell,s}\mid\mathfrak c_{\ell,s}=g]
     \nabla f\,d(\mathfrak c_{\ell,s})_*\nu(g).
 \tag{V17.3}\label{c18v17:native-form}
\]
This is not an assertion that the metric stays $\ell I$ after flow.
\end{proposition}
\begin{proof}
The compact smooth vector field is complete in both time directions,
so $\Phi_s$ is a diffeomorphism with inverse $\Phi_{-s}$. Its gauge
equivariance follows from invariance of $E_{\rm sp}$. V13's product
coordinate change proves that $\mathfrak b_\ell$ is a surjective
submersion, and that
$D\mathfrak b_\ell D\mathfrak b_\ell^*=\ell I$.
Composition proves the first assertions and the fibre dimension:
there are $3N^3$ fine and $3n^3$ coarse links, each of dimension three.
On principal strata the gauge orbit dimensions are $3N^3$ and
$3n^3$. The vertical gauge orbit dimension is $3(N^3-n^3)$;
subtract it from the ordinary fibre dimension.

The inherited plaquette Hessian in f16.2 V69 has diagonal link
blocks of norm at most four after summing incident plaquettes.
The off-diagonal block norms have row sum at most twelve. Thus the
full product-metric Hessian has operator norm at most sixteen.
Its variational equation in parallel orthonormal frames gives
\[
 e^{-16s}\|v\|\le\|D\Phi_s v\|\le e^{16s}\|v\|,
 \qquad s\ge0.
\]
Apply this also to the adjoint and use the exact block-product metric
to obtain \eqref{c18v17:metric}. Transporting $\nu$ by the
diffeomorphism gives a smooth positive density in the flowed links;
V13's product coordinate change and compact integration then give
the positive coarse density. Finally the actual ground-state form
identity and the chain rule prove \eqref{c18v17:native-form}.
\end{proof}

These elementary consequences of the inherited flow are included to
identify the map and measure, not as a new gap mechanism.
The estimate \eqref{c18v17:metric} is uniform in volume but becomes
exponentially poor at $s=\ell^2$. No weak-field probability estimate
has been used to improve it.

\subsection{The transverse sampling map and its alias formula}

For the reference calculation only, let $\mathcal E_N$ be the real
space of three-color link fields which are orthogonal to discrete
gradients and have zero spatial mean. Removing the harmonic fields
makes
\[
 \mathsf L=d_1^*d_1\big|_{\mathcal E_N}>0,\qquad
 \lambda(k)=4\sum_{d=1}^3\sin^2(k_d/2)
 \tag{V17.4}\label{c18v17:curl}
\]
strictly positive. This removal is a definition of the Gaussian
reference, not a treatment of the compact theory's torons.

At the identity, the spatial Hessian is $d_1^*d_1$. The differential
of the candidate, followed by the orthogonal coarse transverse
projection $\mathsf P_c$, is therefore
\[
 C_\ell=\mathsf P_c B_\ell e^{-\ell^2\mathsf L},
 \qquad
 (B_\ell A)_{x,d}=\sum_{j=0}^{\ell-1}A_{\ell x+j\widehat d,d}.
 \tag{V17.5}\label{c18v17:C}
\]
Here $\mathsf P_c$ removes only coarse gradients; at zero coarse
momentum it is the identity on three link directions. The fine
harmonic input is absent. The transverse projection expresses the
linearized gauge quotient; no global nonlinear transverse gauge
chart is assumed in \eqref{c18v17:map}.

Use unitary fine and coarse Fourier transforms. For each coarse
momentum $q\in[-\pi,\pi)^3$, the aliases have the unique representatives
\[
 k_m=(q+2\pi m)/\ell\in[-\pi,\pi)^3,\quad m\in\mathcal M_\ell(q),
 \qquad \#\mathcal M_\ell(q)=\ell^3 .
\]
The primary alias is $m=0$. Put
\[
 d(k)=(e^{ik_d}-1)_{d=1}^3,\quad
 D_m=\operatorname{diag}\left(\sum_{j=0}^{\ell-1}e^{ijk_{m,d}}\right),
 \quad P_q=P_{d(q)^\perp},\quad P_0=I.
\]
Summing the literal straight chains and grouping the fine Fourier
momenta gives
\[
 \widehat{C_\ell A}(q)
  =\ell^{-3/2}P_q
      \sum_{m\in\mathcal M_\ell(q)}
               D_m e^{-\ell^2\lambda(k_m)}\widehat A(k_m).
 \tag{V17.6}\label{c18v17:alias}
\]
The factor $\ell^{-3/2}$ follows from the two unitary volume
normalizations and cancels in a kernel constraint. Each nonzero
fine mode satisfies $\widehat A(k)\perp d(k)$. The same formula
holds color by color.

The gauge-compatible identity and bounds are
\[
 D_m d(k_m)=d(q),\qquad
 \|D_m\|\le\ell,\qquad
 \|P_qD_0u\|\ge\frac{2\ell}{\pi}\|u\|
      \quad(u\perp d(k_0),\ q\ne0).
 \tag{V17.7}\label{c18v17:primary}
\]
Indeed the first identity telescopes. For the last,
each primary diagonal modulus satisfies
\[
 |\sin(q_d/2)/\sin(q_d/(2\ell))|\ge2\ell/\pi,
\]
with value $\ell$ when $q_d=0$. Using $D_0d(k_0)=d(q)$,
\[
 \|P_qD_0u\|
 =\inf_z\|D_0(u-zd(k_0))\|
 \ge(2\ell/\pi)\inf_z\|u-zd(k_0)\|
 =(2\ell/\pi)\|u\|.
\]
This is precisely the inherited constraint-preserving inverse
argument, now with the arbitrary-length multiplier displayed.

\subsection{An all-scale lower frequency on the discarded tangent space}

The square root in the next theorem is important. A compressed
lower bound for $\mathsf L$ cannot simply be square-rooted to
obtain the same bound for the compression of $\sqrt{\mathsf L}$.

\begin{theorem}[Flow suppresses low--high alias cancellation]
\label{c18v17:frequency}
Let $\mathcal V_\ell=\ker C_\ell\subset\mathcal E_N$.
For every $N=\ell n$ as above,
\[
 \boxed{\quad
 \langle v,\sqrt{\mathsf L}\,v\rangle
                 \ge\ell^{-1}\|v\|^2,\qquad v\in\mathcal V_\ell .
 \quad}
 \tag{V17.8}\label{c18v17:floor}
\]
The constant is independent of $\ell$ and $n$. In every nonzero
coarse momentum class with $\lambda(k_0)<\ell^{-2}$,
the kernel constraint in fact implies
\[
 \|\widehat v(k_0)\|^2
      \le\rho_*^2\sum_{m\ne0}\|\widehat v(k_m)\|^2,
 \qquad
 \rho_*^2
       =20\big[(1201/1200)^3-1\big]<1/16 .
 \tag{V17.9}\label{c18v17:leakage}
\]
This is a statement about field variations, not yet the native
conditional Hamiltonian.
\end{theorem}
\begin{proof}
In a coarse momentum class the primary $k_0$ minimizes every
$|k_{m,d}|$, hence minimizes $\lambda(k_m)$. Every nonprimary
alias has some $|k_{m,d}|\ge\pi/\ell$, so
\[
 \sqrt{\lambda(k_m)}
       \ge2\sin(\pi/(2\ell))\ge2/\ell\qquad(m\ne0).
 \tag{V17.10}\label{c18v17:high}
\]
If $\lambda(k_0)\ge\ell^{-2}$, all modes in that class already
have frequency at least $\ell^{-1}$, without a kernel constraint.

Suppose $q\ne0$ and $\lambda(k_0)<\ell^{-2}$.
Move the nonprimary terms of \eqref{c18v17:alias} to the other
side and apply \eqref{c18v17:primary} and Cauchy--Schwarz:
\[
 \|\widehat v(k_0)\|^2
 \le\frac{\pi^2}{4}e^{2\ell^2\lambda(k_0)}
      \left(\sum_{m\ne0}e^{-2\ell^2\lambda(k_m)}\right)
      \sum_{m\ne0}\|\widehat v(k_m)\|^2 .
 \tag{V17.11}\label{c18v17:rho}
\]
There is no factor $\ell^3$ in this estimate. To bound the heat
sum uniformly, use $\sin(|k_d|/2)\ge|k_d|/\pi$ and
$|q_d+2\pi m_d|\ge\pi|m_d|$ for every nonzero integer $m_d$.
Thus
\[
 \ell^2\lambda(k_m)\ge4|m|^2,\qquad
 \sum_{m\ne0}e^{-2\ell^2\lambda(k_m)}
 \le\left(1+2\sum_{j=1}^\infty e^{-8j^2}\right)^3-1.
\]
Since $7<e^2<8$, one has $e^8>2401$ and
$\sum_{j\ge1}e^{-8j^2}\le(e^8-1)^{-1}<1/2400$.
Also $(\pi^2/4)e^2<20$, for example from $\pi<22/7$.
Substitution gives \eqref{c18v17:leakage}. Its final strict
comparison is the rational inequality
\[
 20[(1201/1200)^3-1]=4323601/86400000<1/16.
\]
For this class, \eqref{c18v17:high} and \eqref{c18v17:leakage}
give
\[
 \sum_m\sqrt{\lambda(k_m)}\,\|\widehat v(k_m)\|^2
 \ge\frac{2}{\ell}\sum_{m\ne0}\|\widehat v(k_m)\|^2
 \ge\frac{32}{17\ell}\sum_m\|\widehat v(k_m)\|^2.
\]
At $q=0$ the primary fine harmonic field has been removed;
all remaining modes obey \eqref{c18v17:high}, whether or not
the coarse zero-mode constraint is imposed. Summing the three
cases over coarse momenta and colors proves \eqref{c18v17:floor}.
Complex Fourier notation is compatible with the real field condition,
so the resulting inequality is on the stated real space.
\end{proof}

The proof uses smoothing before sampling, not an assertion that
smoothing alone removes modes. It also controls the limit of very
small nonzero $q$: the low-alias inverse in
\eqref{c18v17:primary} stays bounded, so no inverse power of
the smallest torus momentum was inserted.

\subsection{A full Gaussian conditional spectral and coverage theorem}

Here specify the continuous-Hamiltonian quadratic reference:
\[
 H_{\rm G}=-\kappa\Delta_{\mathcal E_N}
                  +\frac b2\langle A,\mathsf L A\rangle,\qquad
 R=\sqrt{2b/\kappa}\sqrt{\mathsf L},\qquad \kappa,b>0 .
 \tag{V17.12}\label{c18v17:HG}
\]
Its normalized positive ground and ground law are
$\Omega_{\rm G}(A)=c\exp(-\langle A,RA\rangle/4)$ and
$d\nu_{\rm G}=\Omega_{\rm G}^2dA$; the covariance is $R^{-1}$.
Direct differentiation gives
\[
 E_{\rm G}=\frac\kappa2\operatorname{Tr}R,\quad
 K_{\rm G}=H_{\rm G}-E_{\rm G},\quad
 \Omega_{\rm G}^{-1}K_{\rm G}\Omega_{\rm G}
            =-\kappa\Delta+\kappa(RA)\cdot\nabla .
 \tag{V17.13}\label{c18v17:generator}
\]
These are exact identities for \eqref{c18v17:HG}; no error estimate
comparing it with the compact Hamiltonian is claimed.

Let $P_{\rm G}$ be conditional expectation onto $C_\ell A$ in
$L^2(\nu_{\rm G})$, transported by multiplication with
$\Omega_{\rm G}$ to the Schr\"odinger space. Put $Q_{\rm G}=1-P_{\rm G}$.
The restricted form on $Q_{\rm G}$ defines $D_{\rm G}$.
Let $\Pi^{\rm G}_E=\mathbf1_{(0,E]}(K_{\rm G})$ exclude the vacuum.

\begin{theorem}[Gaussian discarded floor on all functions]
\label{c18v17:Gaussian}
For every $N=\ell n$ in the stated range,
\[
 \boxed{\quad
 D_{\rm G}\succeq g_\ell I,\qquad
 g_\ell=\frac{\sqrt{2\kappa b}}{\ell},\qquad
 \|D_{\rm G}^{-1}\|\le\frac{\ell}{\sqrt{2\kappa b}} .
 \quad}
 \tag{V17.14}\label{c18v17:DG}
\]
More generally, for every $f$ in the Gaussian form domain,
\[
 \int |f-\mathbb E_{\nu_{\rm G}}[f\mid C_\ell A]|^2d\nu_{\rm G}
 \le\frac{\ell}{\sqrt{2\kappa b}}\,
                       \kappa\int|\nabla f|^2d\nu_{\rm G}.
 \tag{V17.15}\label{c18v17:poincare}
\]
Consequently, on the entire spectral window, not only first chaos,
\[
 \|Q_{\rm G}\Pi^{\rm G}_E\|
       \le\min\{1,\sqrt{E/g_\ell}\}.
 \tag{V17.16}\label{c18v17:coverage}
\]
All statements remain valid after restriction to global color singlets.
\end{theorem}
\begin{proof}
Disintegrate the nondegenerate finite-dimensional Gaussian according
to the linear map $C_\ell$, using its actual range if necessary.
Every conditional is a translate of a Gaussian on
$\mathcal V_\ell=\ker C_\ell$. Its precision on that subspace is
\[
 R_{\rm vert}=P_{\mathcal V_\ell}R|_{\mathcal V_\ell}
             \succeq \frac{\sqrt{2b/\kappa}}{\ell}I
\]
by \eqref{c18v17:floor}. This is a compression of the precision,
not the inverse of a compression of $R^{-1}$.
Diagonalize this positive precision. In one variable of precision
$r$, the normalized Hermite expansion gives
$\operatorname{Var}f\le r^{-1}\int|f'|^2$; the product Hermite
expansion gives the same inequality with the smallest precision
for all smooth functions of all the vertical variables.
Translation of the conditional mean has no effect on this constant.
Integrate over the coarse marginal and bound the squared vertical
gradient by the full squared gradient. This proves
\eqref{c18v17:poincare}, first on smooth tests, then on the closed
Gaussian form domain.

For $Q_{\rm G}\Omega_{\rm G}f=\Omega_{\rm G}f$,
\eqref{c18v17:poincare} proves the form floor
\eqref{c18v17:DG}. Gaussian conditional expectation along a fixed
linear map preserves the form domain with finite-dimensional
constants: its conditional mean is linear and its centered
conditional covariance is constant, so differentiating its Gaussian
integral gives a fixed linear combination of conditional derivatives.
Projected smooth approximations give density in the restricted form.
Thus the restricted closed form and its inverse are well-defined.
For $\psi\in\operatorname{Ran}\Pi^{\rm G}_E$, the same inequality
and $\langle\psi,K_{\rm G}\psi\rangle\le E\|\psi\|^2$ give
\eqref{c18v17:coverage}. Global color rotations commute with
$R,C_\ell$, the Gaussian law, and these projections. Restricting
the inequalities therefore preserves them.
\end{proof}

In particular $E\le g_\ell/4$ gives low-window projection error at
most $1/2$. This is a proved reference coverage criterion with the
native coefficients $\kappa,b$ left in place, not a numerical
single-mode fit. It differs from the strong-electric
$12\kappa\ell$ compressed energy in V14: that was the compact
$b=0$ model, not \eqref{c18v17:HG}.

\subsection{The reference source survives and the infrared remains unpaid}

A conditional floor does not make the compression exact.
Writing $Y=C_\ell A$, $m(Y)=\mathbb E[A\mid Y]$, the actual source
within the Gaussian reference for a linear retained function is
\[
 Q_{\rm G}K_{\rm G}(\Omega_{\rm G}Y)
       =\kappa\Omega_{\rm G}C_\ell R(A-m(Y)).
 \tag{V17.17}\label{c18v17:linear-source}
\]
This formula is vector-valued, component by component. The fine
reference operator, not an auxiliary conditional sampler, is used.

\begin{proposition}[Flowed sampling does not commute with the reference dynamics]
\label{c18v17:nonzero}
For every $\ell\ge2$, $n\ge5$, $\kappa,b>0$, the right side of
\eqref{c18v17:linear-source} is not identically zero. There also
exists a global-color-invariant quadratic retained function with
nonzero discarded source.
\end{proposition}
\begin{proof}
Choose coarse momentum $q=(2\pi/n,0,0)$ and fine aliases
$k_0=q/\ell$, $k_1=(q-2\pi\widehat1)/\ell$.
Both are legitimate representatives and their frequencies satisfy
$0<\lambda_0<\lambda_1$. In polarization $\widehat2$ they are
fine and coarse transverse, and both chain multipliers equal $\ell$.
Ignoring their common positive Fourier factor, choose amplitudes
\[
 v_0=e^{-\ell^2\lambda_1},\qquad
 v_1=-e^{-\ell^2\lambda_0}.
\]
Then $C_\ell v=0$, whereas
\[
 C_\ell Rv
 \ \text{has a nonzero component proportional to}\
 e^{-\ell^2(\lambda_0+\lambda_1)}
          (\sqrt{\lambda_0}-\sqrt{\lambda_1}).
\]
Include conjugate momenta to obtain a real vector.
The conditional Gaussian has strictly positive covariance on the
whole kernel. Hence $C_\ell R(A-m(Y))$ has positive variance in
at least one real coarse component $y$. Do the same construction
in each of the three colors and take $f(Y)=\sum_{a=1}^3 y_a^2$.
Its diffusion term is constant; its discarded drift is
$2\kappa\Omega_{\rm G}\sum_a y_a(cR(A-m(Y)))_a$.
The residual Gaussian is independent of $Y$, and both the variance
of $y_a$ and that of the indicated residual are positive.
The source therefore has positive squared norm. The color
contraction makes $f$ invariant under global color rotations.
\end{proof}

Nor does \eqref{c18v17:DG} generate an infrared mass in the
reference. For a real lowest nonzero transverse Fourier coordinate
$x_a$ in each color, the centered singlet
$\sum_a(x_a^2-\mathbb E x_a^2)$ is an exact eigenfunction of the
ground transform with energy
\[
 2\sqrt{2\kappa b}\sqrt{\lambda_{\min}}
       =4\sqrt{2\kappa b}\sin(\pi/N)\longrightarrow0
       \quad(N\to\infty).
 \tag{V17.18}\label{c18v17:massless}
\]
This follows directly by applying \eqref{c18v17:generator} to
the quadratic polynomial. The ultraviolet conditional floor and
the absence of a volume-uniform full reference gap are compatible.

The native target is now specific: control the actual conditional
variance on the fibres of $\mathfrak c_{\ell,\ell^2}$, with the
original physical energy, and control the nonzero source return.
The Gaussian calculation is a reason to investigate that target,
not a replacement for it. The native fibres are compact nonlinear
objects with gauge and winding structure; the Gaussian reference
has removed harmonic fields. The actual $\Omega$ and the metric
in \eqref{c18v17:native-form} are not Gaussian data.
The all-configuration exponential bound
\eqref{c18v17:metric} does not supply the comparison on growing
blocks. The next upstream work must populate native control of
these quantities, including bad-field and harmonic contributions,
or demonstrate a failure of this particular candidate. Another
unqualified flat calculation would not pay that obligation.

The previous goal turn was progress: V16 produced new exact native
rank and conditioning evidence. Its full diagnostic passed before
this append. The present analytic proof is checked by exact
rational constants and a separate exact conditional-covariance
fixture; floating-point Fourier and alias tests are explicitly
regressions, not all-volume proofs or native spectral data.
V16's body is preserved exactly apart from its title, and all
47 other note files are unchanged. The single active source is
V17; only \texttt{.tex} files are kept in \texttt{latex\_notes}.
The next companion review and broad literature sweep remain V20.
The interacting four-dimensional continuum and its positive
physical Yang--Mills mass gap remain unproved.
% END CYCLE 18 VERSION 17 APPEND
% BEGIN CYCLE 18 VERSION 18 APPEND
\clearpage
\section{V18: Native buffered flow-metric probabilities and moments}
\label{c18v18:section}

V17 supplies a Gaussian conditional floor for flow followed by sampling,
but leaves its nonlinear vacuum comparison unpaid. This continuation
proves an actual-ground-law estimate for one part of that comparison:
the local kinetic metric of the same compact map. No Gaussian law,
harmonic-mode deletion, independent plaquette model, or conditional
spectral gap is assumed. The result has two distinct ranges. A local
metric upper bound holds with high probability on a power-law
mesoscopic range; its bad-event contribution to moments is controlled
on a smaller logarithmic range. Conditioning then gives a genuine
$L^p$ bound for the actual compressed kinetic metric.

These are upper bounds, not the lower spectral inequality needed for
the full mass gap. Keeping that distinction is essential: probability
control alone does not bound an exponentially large bad-event
contribution, and an upper kinetic bound does not prove coercivity.

\subsection{Inputs, source overlap, and the actual flowed law}

Use exactly V17's $N=\ell n$, $\ell\ge2$, $n\ge5$, native
$H$, $\Omega>0$, $K=H-E_0$, $\nu=\Omega^2dU$, and
$\mathfrak c_{\ell,s}=\mathfrak b_\ell\circ\Phi_s$.
Write
\[
 \delta_p(U)=1-w_p(U),\qquad
 \beta=b/\kappa>0,\qquad
 \varepsilon_\beta=\min\{1,A_*\beta^{-1/2}\},\quad A_*=3\sqrt2,
 \qquad W_t=\Phi_t(U).
 \tag{V18.1}\label{c18v18:parameters}
\]
The expectations below are over $U\sim\nu$. In particular the law of
$W_t$ is the pushforward of the actual ground law, not a new Wilson
Gibbs law, and is not declared stationary under spatial flow.

The initial mean-defect estimate is inherited from current V6,
ultimately f17 V89. The exact compact Hessian and its signed
curvature cost are inherited from f16.2 V69. That archive's V70
already has gauge completion, a local occupation kernel,
killed-domain continuation, and a warning about rare amplifying
defects. We do not reintroduce its unpaid uniform-restart
hypothesis as an assumption. V74's transported density and
all-configuration inverse-flow bounds are also already known.
The new step here is to populate a local integrated-curvature
budget in the actual Hamiltonian ground law, pay exterior
influence with a deterministic buffer, and quantify when
bad-event metric moments are integrable along a growing scale.

Finite-range influence estimates are established methodology.
For context only, the primary abstract of Raz--Sims,
\emph{Lieb--Robinson bounds for classical anharmonic lattice systems},
\href{https://arxiv.org/abs/0902.0025}{arXiv:0902.0025},
describes volume-independent classical locality bounds.
We use no theorem about their oscillator systems for the present
gradient flow. The finite-range estimate needed here is proved below.

\begin{lemma}[Native defect budget through the full flow]
\label{c18v18:defect}
For every plaquette and $t\ge0$,
\[
 \mathbb E_\nu\delta_p(W_t)\le\varepsilon_\beta.
 \tag{V18.2}\label{c18v18:mean}
\]
One also has, for every $s\ge0$,
\[
 \mathbb E_\nu\sup_{0\le t\le s}\delta_p(W_t)
                       \le(1+24s)\varepsilon_\beta .
 \tag{V18.3}\label{c18v18:sup}
\]
Neither assertion requires independence or invariance of $\nu$ under flow.
\end{lemma}
\begin{proof}
Along the full gradient flow,
\[
 \frac{d}{dt}E_{\rm sp}(W_t)
                 =-|\nabla E_{\rm sp}(W_t)|^2\le0 .
\]
The initial law is invariant under translations and cubic symmetries
by uniqueness of its positive ground, and the flow is equivariant
under these same symmetries. Hence all plaquette defects have equal
expectation at every time. Divide the expected total-energy inequality
by $3N^3$ and use the inherited initial bound, including the
Haar-trial bound by one, to obtain \eqref{c18v18:mean}.

For the second assertion, put $g_{p,e}=\nabla_e\delta_p$.
If $e$ belongs to $p$, then
$|g_{p,e}|^2=1-w_p^2\le2\delta_p$.
A plaquette has four distinct links, each with three other incident
plaquettes. On $N\ge5$ these twelve neighboring plaquettes are distinct
and share just one link with $p$. Differentiating along the full flow,
discard its nonpositive self term. Cauchy--Schwarz gives
\[
 \left(\frac{d}{dt}\delta_p(W_t)\right)_+
 \le12\delta_p(W_t)+\sum_{q:\,q\sim p}\delta_q(W_t),
\]
because $|g_{p,e}\cdot g_{q,e}|
\le2\sqrt{\delta_p\delta_q}\le\delta_p+\delta_q$.
A nonnegative differentiable function satisfies
$\sup_{t\le s}\delta_p(W_t)\le\delta_p(U)+
\int_0^s(\delta_p'(W_t))_+dt$. Take expectations and use
\eqref{c18v18:mean} on all thirteen terms.
\end{proof}

The main metric estimate will use the time integral directly,
rather than \eqref{c18v18:sup}; this avoids its extra factor $1+s$.

\subsection{A local integrated-curvature budget with centre plaquettes included}

The link graph joins two distinct stored links when they lie in a
common plaquette. For a nonempty link domain $\mathcal D$, let
$\mathcal P(\mathcal D)$ be the distinct plaquettes incident to at
least one of its links, and $M_{\mathcal D}=|\mathcal P(\mathcal D)|$.
Retain the exact cost from f16.2 V69:
\[
 c_p(W)=3\sqrt{1-w_p(W)^2}+4(-w_p(W))_+,\qquad
 v_e(W)=\sum_{p\ni e}c_p(W),\qquad
 v_{\mathcal D}(W)=\max_{e\in\mathcal D}v_e(W).
\]
Define the random integrated local cost
\[
 \mathfrak a_{\mathcal D}(s;U)
                       =\int_0^s v_{\mathcal D}(W_t)\,dt .
 \tag{V18.4}\label{c18v18:cost}
\]
The maximum is over a declared finite buffer, not over the whole
thermodynamic lattice.

\begin{proposition}[A paid native second moment]
\label{c18v18:cost-budget}
For every domain and time,
\[
 \boxed{\quad
 \mathbb E_\nu\mathfrak a_{\mathcal D}(s)^2
                   \le100s^2M_{\mathcal D}\varepsilon_\beta .
 \quad}
 \tag{V18.5}\label{c18v18:cost-second}
\]
In particular
$\nu\{\mathfrak a_{\mathcal D}(s)>\alpha\}
\le\min\{1,100s^2M_{\mathcal D}\varepsilon_\beta/\alpha^2\}$
for $\alpha>0$.
\end{proposition}
\begin{proof}
For $0\le\delta=1-w\le1$,
$c_p^2=9\delta(2-\delta)\le18\delta\le25\delta$.
For $1\le\delta\le2$, Cauchy--Schwarz gives
\[
 c_p^2
 \le(3^2+4^2)\{\delta(2-\delta)+(\delta-1)^2\}
 =25\le25\delta.
\]
This includes $w=-1$, where the cost is four, not zero.
Each link has four incident plaquettes. Consequently
\[
 v_{\mathcal D}(W)^2
 \le100\max_{e\in\mathcal D}\sum_{p\ni e}\delta_p(W)
 \le100\sum_{p\in\mathcal P(\mathcal D)}\delta_p(W).
\]
Cauchy--Schwarz in time, followed by \eqref{c18v18:mean},
proves
\[
 \mathbb E_\nu\left(\int_0^s v_{\mathcal D}(W_t)\,dt\right)^2
 \le s\int_0^s\mathbb E_\nu v_{\mathcal D}(W_t)^2dt
 \le100s^2M_{\mathcal D}\varepsilon_\beta.
\]
Markov's inequality gives the tail statement.
\end{proof}

This is not an exponential occupation moment and says nothing about
a uniform restart law for an auxiliary walk. It is an actual-vacuum
second moment of a specific finite-buffer trajectory cost.

\subsection{Paying the exterior influence rather than deleting it}

Along the full $W_t$, use separate parallel orthonormal tangent
frames on each link, and write $J(t)=D\Phi_t(U)$.
The inherited variational equation and incidence bounds are
\[
 J'=-\mathsf H(t)J,\qquad
 \|\mathsf H(t)\|\le16,\qquad
 \mathsf H(t)\succeq-\operatorname{diag}(v_e(W_t)I_3).
 \tag{V18.6}\label{c18v18:H}
\]
The off-diagonal blocks vanish off the link graph.
For a link domain $\mathcal D$, define $J_{\mathcal D}(t)$
by restricting this \emph{linear variational equation} to its
principal matrix $\mathsf H_{\mathcal D\mathcal D}(t)$.
The background remains the original full trajectory $W_t$,
including every exterior dependence of these entries.
This is not the nonlinear flow with frozen exterior links.

Let $\mathcal B$ be the desired output link set and
\[
 \mathcal D_R=\{e:\operatorname{dist}_{\rm link}(e,\mathcal B)<R\},
 \qquad R\in\mathbb N,\ R\ge1 .
\]
If the buffer fills the torus the ensuing bounds still hold.

\begin{lemma}[Buffered row control with an explicit exterior remainder]
\label{c18v18:buffer}
Using the natural inclusion and projection for
$J_{\mathcal D_R}$, one has
\[
 \begin{split}
 \|\mathbf1_{\mathcal B}
       (J(s)-\iota_{\mathcal D_R}J_{\mathcal D_R}(s)
                                  \pi_{\mathcal D_R})\|
  &\le2\sum_{j\ge R}\frac{(16s)^j}{j!}
    \le2e^{16es-R},\\
 \|\mathbf1_{\mathcal B}J(s)\|
  &\le e^{\mathfrak a_{\mathcal D_R}(s)}+2e^{16es-R}.
 \end{split}
 \tag{V18.7}\label{c18v18:row}
\]
Here $e$ in the exponential denotes Euler's number.
The norm takes \emph{all} initial link tangents as input.
\end{lemma}
\begin{proof}
Expand both propagators in their convergent time-ordered
integral series. A product of fewer than $R$ Hessian factors,
read backward from an output in $\mathcal B$, cannot reach
$\mathcal D_R^c$: each off-diagonal factor moves by at most
one graph edge, and a diagonal factor does not move.
Thus the restricted rows of those terms agree exactly, including
order zero. The norm of either order-$j$ term is at most
$(16s)^j/j!$. Subtraction and summation prove the first bound.
For the exponential bound use
$\mathbf1_{j\ge R}\le e^{j-R}$ and sum the exponential series.
This argument permits arbitrary noncommuting time dependence.

By \eqref{c18v18:H}, solutions of the restricted equation obey
\[
 \frac{d}{dt}|z(t)|^2
       \le2v_{\mathcal D_R}(W_t)|z(t)|^2.
\]
Hence $\|J_{\mathcal D_R}(s)\|
\le\exp\mathfrak a_{\mathcal D_R}(s)$.
The triangle inequality proves the second line.
No exterior return has been replaced by zero.
\end{proof}

For $0<\eta\le1$, choose
\[
 R(s,\eta)=\left\lceil48s+\log(2/\eta)\right\rceil.
 \tag{V18.8}\label{c18v18:R}
\]
Since $e<3$, the exterior remainder in \eqref{c18v18:row}
is at most $\eta$. The buffer grows linearly in dimensionless
flow time, not diffusively. This cost will be retained in
the scale range, rather than called cutoff-independent at
fixed positive physical smoothing time.

\subsection{Native local metric tails on growing blocks}

Let $\Gamma$ be any nonempty subset of the twelve coarse edges
of one elementary coarse cube. Let $\mathcal B_\Gamma$ be
their fine straight-chain links, and define
\[
 G_\Gamma(U)=D(\mathfrak c_{\ell,s})_\Gamma(U)
               D(\mathfrak c_{\ell,s})_\Gamma(U)^*,\qquad
 X_\Gamma(U)=\sqrt{\|G_\Gamma(U)\|/\ell}.
 \tag{V18.9}\label{c18v18:G}
\]
These are in the original orthonormal coarse link frame.
The norm is invariant under its orthogonal gauge rotations.
All input fine directions, including gauge, harmonic and winding
variations, are retained when taking this operator norm.

The chain differential has norm $\sqrt\ell$, since the
chains use disjoint links. Therefore \eqref{c18v18:row} gives
\[
 X_\Gamma\le e^{\mathfrak a_{\mathcal D_R}(s)}+\eta,
 \qquad X_\Gamma\le e^{16s},
 \quad R=R(s,\eta).
 \tag{V18.10}\label{c18v18:X}
\]
The second bound is V17's deterministic global estimate.

\begin{theorem}[Native probability and logarithmic-moment estimates]
\label{c18v18:metric-tail}
For every $\alpha>0$ and $0<\eta\le1$,
\[
 \begin{split}
 \nu\{\|G_\Gamma\|>\ell(e^\alpha+\eta)^2\}
  &\le\min\left\{1,
         \frac{100s^2M_{\mathcal D_R}\varepsilon_\beta}{\alpha^2}
           \right\},\\
 \mathbb E_\nu
    \left[\log_+\frac{X_\Gamma}{1+\eta}\right]^2
  &\le100s^2M_{\mathcal D_R}\varepsilon_\beta .
 \end{split}
 \tag{V18.11}\label{c18v18:metric-probability}
\]
The logarithm is set to zero at $X_\Gamma=0$.
For V17's choice $s=\ell^2$,
\[
 \boxed{\quad
 \nu\{\|G_\Gamma\|>16\ell\}
       \le\min\{1,\,282357600\,A_*\,\ell^{10}\beta^{-1/2}\}.
 \quad}
 \tag{V18.12}\label{c18v18:power-window}
\]
Thus $\ell=o(\beta^{1/20})$ is a sufficient growing-block
high-probability range, uniformly in the surrounding torus.
It is not asserted optimal.
\end{theorem}
\begin{proof}
The first line of \eqref{c18v18:metric-probability} follows
from \eqref{c18v18:X} and Proposition~\ref{c18v18:cost-budget}.
Since $\mathfrak a_{\mathcal D_R}\ge0$,
$e^{\mathfrak a_{\mathcal D_R}}+\eta
\le(1+\eta)e^{\mathfrak a_{\mathcal D_R}}$.
This proves the logarithmic bound with the same second moment.

For the spatial count, translate the coarse cube so its fine
vertices are in $[0,\ell]^3$. Every link-graph step changes
the integer tail coordinates by at most one in each component.
The tails in $\mathcal D_R$ therefore have representatives
in $[-R+1,\ell+R-1]^3$; a plaquette incident to those links
has its base in $[-R,\ell+R]^3$. With three orientations,
\[
 M_{\mathcal D_R}
          \le\min\{3N^3,\,3(\ell+2R+1)^3\}.
 \tag{V18.13}\label{c18v18:count}
\]
Wrapping can identify representatives but cannot increase
this count. The estimate holds for every subset $\Gamma$.

For $\eta=1$, $s=\ell^2$, equation \eqref{c18v18:R} gives
$R=48\ell^2+1$. For $\ell\ge2$,
$\ell+2R+1=96\ell^2+\ell+3\le98\ell^2$.
Set $\alpha=1$, use $(e+1)^2<16$,
$\varepsilon_\beta\le A_*\beta^{-1/2}$, and
$100\cdot3\cdot98^3=282357600$.
This proves \eqref{c18v18:power-window}, with no volume factor.
\end{proof}

The power-law range only gives probability and logarithmic-moment
control. It does not yet control $\mathbb E_\nu\|G_\Gamma\|$:
on the bad event the deterministic bound is still $\ell e^{32s}$.

\subsection{Paying bad-event moments and the actual compressed metric}

Let $\mu=(\mathfrak c_{\ell,s})_*\nu$ be the full coarse
pushforward and define the principal block of its actual
compressed kinetic metric by
\[
 \overline G_\Gamma(g)
        =\mathbb E_\nu[G_\Gamma\mid\mathfrak c_{\ell,s}=g].
 \tag{V18.14}\label{c18v18:conditional}
\]
The conditioning is on \emph{all} retained coarse links,
not only those in $\Gamma$.
For fixed $\beta,N,\ell,s$ it exists and is smooth, as in V17.

\begin{theorem}[Actual metric moments on a logarithmic scale range]
\label{c18v18:moments}
Fix $p\ge1$ and $0<c<1/(64p)$. As $\beta\to\infty$, allow
any integers $\ell\ge2$, $n\ge5$ with
\[
             s=\ell^2\le c\log\beta,\qquad N=\ell n .
 \tag{V18.15}\label{c18v18:log-range}
\]
Uniformly over these choices and the location of the coarse cube,
\[
 \begin{split}
 \mathbb E_\nu\left(\frac{\|G_\Gamma\|}{\ell}-1\right)_+^p
                  &\longrightarrow0,\\
 \int\left(\frac{\|\overline G_\Gamma(g)\|}{\ell}-1\right)_+^p
                      d\mu(g)&\longrightarrow0 .
 \end{split}
 \tag{V18.16}\label{c18v18:overshoot}
\]
In particular
\[
 \limsup\frac{\|\overline G_\Gamma\|_{L^p(\mu)}}{\ell}\le1.
 \tag{V18.17}\label{c18v18:Lp}
\]
These are upper-overshoot limits, not convergence of the metric to
$\ell I$ and not a lower bound on its eigenvalues.
\end{theorem}
\begin{proof}
Write $T=\log\beta>1$, and in \eqref{c18v18:X} take
$\alpha=\eta=T^{-1}$. The chosen radius satisfies
$R\le48cT+\log(2T)+1$.
Since $\ell\le\sqrt{cT}$, \eqref{c18v18:count} gives
$M_{\mathcal D_R}\le C_cT^3$ for all sufficiently large $T$,
with no dependence on $N$.
By \eqref{c18v18:cost-second}, the event
$\mathcal A=\{\mathfrak a_{\mathcal D_R}(s)>T^{-1}\}$ obeys
\[
 \nu(\mathcal A)\le C_c A_*T^7\beta^{-1/2}.
 \tag{V18.18}\label{c18v18:bad}
\]
On its complement,
$X_\Gamma^2\le h_T:=(e^{1/T}+1/T)^2=1+O(T^{-1})$.
On $\mathcal A$ the global bound in \eqref{c18v18:X} gives
$X_\Gamma^{2p}\le e^{32ps}\le\beta^{32pc}$.
Therefore
\[
 \mathbb E_\nu(X_\Gamma^2-1)_+^p
   \le(h_T-1)^p+C_c A_*T^7\beta^{32pc-1/2}
                      \longrightarrow0 .
 \tag{V18.19}\label{c18v18:moment-proof}
\]
The strict condition $32pc<1/2$ is exactly what pays the
bad-event exponential, and explains the smaller range.

The operator norm is convex, and $x\mapsto(x-1)_+^p$ is convex
and nondecreasing on $x\ge0$. Conditional Jensen, in the same
coarse orthonormal frame, gives
\[
 \left(\frac{\|\overline G_\Gamma\|}{\ell}-1\right)_+^p
 \le\mathbb E_\nu\left[
       \left(\frac{\|G_\Gamma\|}{\ell}-1\right)_+^p
                          \,\middle|\,\mathfrak c_{\ell,s}\right].
\]
Integrating proves the second line of \eqref{c18v18:overshoot}.
Finally
$\|\overline G_\Gamma\|/\ell
\le1+(\|\overline G_\Gamma\|/\ell-1)_+$ and the $L^p$
triangle inequality prove \eqref{c18v18:Lp}.
\end{proof}

This theorem is in the actual interacting compact ground law,
with all link directions present. It is not obtained by
conditioning a Gaussian approximation, and it does not assert
that the coarse law is Wilson at a new coupling.

For example, fix $p>1$ and $c<1/(64p)$. If a smooth physical
coarse function $f$ depends only on $\Gamma$, V17's exact
ground-state form and H\"older's inequality give
\[
 \begin{split}
 \langle J_{\ell,s}f,KJ_{\ell,s}f\rangle
 &\le\kappa\|\overline G_\Gamma\|_{L^p(\mu)}
             \|\nabla_\Gamma f\|_{L^{2p/(p-1)}(\mu)}^2\\
 &\le\kappa\ell(1+o(1))
             \|\nabla_\Gamma f\|_{L^{2p/(p-1)}(\mu)}^2 .
 \end{split}
 \tag{V18.20}\label{c18v18:form}
\]
The error is uniform in the surrounding volume and in $f$
relative to the displayed norm. For $p=1$, the corresponding
gradient norm is $L^\infty$.
The higher gradient norm cannot be replaced by $L^2(\mu)$
without another estimate; no density-comparison constant is
silently inserted.

\subsection{What this pays, and what remains before a mass-gap argument}

We have replaced V17's purely exponential local metric envelope
by an actual-vacuum probability bound on one growing range and
an $L^p$ bound for the true compressed metric on another.
The buffer retains the cost of exterior influence, and the
moment theorem retains the cost of bad configurations.
It neither asserts a small global maximum at thermodynamic
volume nor contradicts the archived rare-defect obstruction.

The result concerns each declared coarse cube, uniformly in
its location and in the ambient volume. Simultaneous probability
control over a growing number of cubes would need its union
or correlation cost. No uniformity over arbitrary conditional
exteriors or replica interpolations has been proved.
The logarithmic range is a sufficient mesoscopic estimate;
it does not reach a fixed positive physical blocking length
on an unproved continuum trajectory.

Most importantly, an upper metric estimate is not the missing
native conditional Poincare floor. The latter still needs
control of the actual distribution along the nonlinear fibres
of $\mathfrak c_{\ell,\ell^2}$, together with the original
physical energy. No uniform inverse for that discarded
Hamiltonian, no lower retained metric on physical directions,
and no resolvent-weighted source return has been obtained here.
Harmonic and winding input directions were not removed from
our upper estimates, but their physical energies have not
thereby been bounded below.

The next step should use the now-populated moment bounds in
that native fibre/source problem, or identify a precise further
obstruction. Repeating the Gaussian floor, or treating
\eqref{c18v18:form} as a coercive inequality, would not advance it.
The full interacting four-dimensional continuum construction
and its positive physical Yang--Mills mass gap remain unproved.

The previous goal turn was progress: V17 specified and tested a
real reduction and proved its reference coverage theorem.
Its complete diagnostic passed before this append.
The present checks use exact rational quaternion derivatives,
compact cost constants, noncommuting finite-range word tests,
and native lattice buffer counts. They are neither a simulation
of the unknown ground law nor a spectral proof by testing.
V17's body is preserved exactly except its title; all 47 other
note files are unchanged. Only the active source becomes V18,
and only \texttt{.tex} files are stored in \texttt{latex\_notes}.
The next companion checkpoint and broad literature sweep remain V20.
% END CYCLE 18 VERSION 18 APPEND
% BEGIN CYCLE 18 VERSION 19 APPEND
\clearpage
\section{V19: Native physical loop-metric normalization}
\label{c18v19:section}

V18 controls the full local coarse-link metric from above, but its
scale $\ell$ is not a calibrated scale for a smoothed physical loop.
Before using that estimate in a spectral comparison, we resolve this
normalization issue. For the same compact flow--sampling map, we prove
a two-sided native estimate of order $\ell^{-1}$ in the differential
direction of one coarse Wilson plaquette, on an event with a paid
actual-vacuum probability. We then pass the lower estimate to most
coarse configurations and prove moment upper bounds for the true
conditional metric in this direction.

This is neither a lower bound on all physical covectors nor a
conditional spectral floor along the discarded fibres. In particular,
a lower coefficient multiplying a Wilson-loop gradient does not control
the variance of that loop, let alone of every physical observable.
The result is a normalization input to that remaining problem, with
its scale and exceptional-set restrictions stated explicitly.

\subsection{Inherited inputs and the flat physical normalization}

Use V17--V18's native $N=\ell n$, $\ell\ge2$, $n\ge5$,
$\nu=\Omega^2dU$, $\beta=b/\kappa>0$, and
$\varepsilon_\beta=\min\{1,3\sqrt2\,\beta^{-1/2}\}$.
The smoothing time remains $s=\ell^2$, and
$\mathfrak c=\mathfrak b_\ell\circ\Phi_s$. The Hamiltonian time
and its coefficients are unchanged.

We reuse V17's exact chain/alias identity, f16.2 V69's compact
Hessian, and V18's native flowed-defect and buffered-variation
estimates. Rooted axial gauge and heat localization are established
tools, not new general constructions. Earlier native one-link
Harnack and conditional-score estimates (f17 V91--V98) do not by
themselves give a many-link nonlinear-fibre floor.
The new calculation here combines a physical loop covector with a
local compact-to-flat comparison whose exceptional probability is
controlled in the actual Hamiltonian vacuum.

For external context only, the primary abstract of
M.~L\"uscher, \emph{Properties and uses of the Wilson flow in lattice QCD},
\href{https://arxiv.org/abs/1006.4518}{arXiv:1006.4518},
describes gauge-invariant flowed observables at a smoothing length
of order $\sqrt t$. We import no continuum, vacuum-comparison or
mass-gap theorem from it. Our three-dimensional spatial flow and
Hamiltonian ground law remain the ones defined above.

Let $\mathsf L_0=d_1^*d_1$ act on all real three-color fine link
fields at the identity. In this subsection we do \emph{not} remove
the harmonic fields. Write $d_{0,c}$ for the coarse gradient.
A flat physical covector means $a\in\ker d_{0,c}^*$; the zero
coarse momentum is included. Set
\[
 m=\frac2\pi e^{-3\pi^2},\qquad 0<m<1 .
 \tag{V19.1}\label{c18v19:m}
\]

\begin{lemma}[The flat physical metric has scale $\ell^{-1}$]
\label{c18v19:flat}
For every flat physical coarse covector,
\[
 \frac{m^2}{\ell}\|a\|^2
 \le\|e^{-\ell^2\mathsf L_0}B_\ell^*a\|^2
 \le\frac4\ell\|a\|^2 .
 \tag{V19.2}\label{c18v19:flat-metric}
\]
The constants are independent of $\ell,n$, and all three colors.
\end{lemma}
\begin{proof}
Use the unitary Fourier conventions and the alias matrices $D_m$
of \eqref{c18v17:alias}. To avoid confusion with the constant in
\eqref{c18v19:m}, denote an alias index by $j$ in this proof.
For $a(q)\perp d(q)$, the telescoping identity gives
\[
 d(k_j)^*D_j^*a(q)=d(q)^*a(q)=0.
\]
Thus every fine alias of $B_\ell^*a$ is transverse; the fine
zero mode is allowed and has heat multiplier one.
The contribution of the coarse class $q$ is exactly
\[
 \ell^{-3}\sum_j
       e^{-2\ell^2\lambda(k_j)}|D_j^*a(q)|^2 .
 \tag{V19.3}\label{c18v19:flat-symbol}
\]
For the primary alias, each diagonal modulus of $D_0$ is at
least $2\ell/\pi$, and
$\ell^2\lambda(q/\ell)\le|q|^2\le3\pi^2$.
Its contribution alone gives the lower bound.
For the upper bound, $\|D_j\|\le\ell$ and V17's scalar heat sum
gives, including the primary term,
\[
 \sum_j e^{-2\ell^2\lambda(k_j)}
 \le\left(1+2\sum_{r\ge1}e^{-8r^2}\right)^3
 <(1201/1200)^3<4 .
\]
Sum over $q$ and colors. No inverse smallest momentum or
discarded harmonic field occurs in this proof.
\end{proof}

The factor in \eqref{c18v19:flat-metric} is a statement about the
original norm of a physical covector after flow and sparse sampling.
It is not a change of clock. V18's $\ell$ upper envelope is valid,
but does not locate this smaller physical scale.

\subsection{A local gauge comparison on an embedded buffer}

Fix a coarse plaquette $\square$, with its four coarse edges
denoted by $\Gamma$. Translate and rotate coordinates so that its
fine boundary $\mathcal B$ lies in $[0,\ell]^3$. The link graph
and the strict-radius buffer $\mathcal D_R$ are as in V18. Choose
\[
 \begin{split}
 R&=\left\lceil64\ell^2+\log(16\ell/m)\right\rceil,\\
 L&=2R+\ell+4,\qquad
 h=\frac{m}{3600\,\ell^3L},\qquad s=\ell^2 .
 \end{split}
 \tag{V19.4}\label{c18v19:scales}
\]
Assume, in addition to $N=\ell n$, that
\[
                         N\ge L+1 .
 \tag{V19.5}\label{c18v19:embedding}
\]
This makes the vertex box
$\mathcal Q=[-R-2,\ell+R+2]^3$ an embedded box of side $L$
in the torus. The restriction is important: we have not replaced
a winding neighborhood by a simply connected box.
All ambient-volume uniformity below is under
\eqref{c18v19:embedding}.

Let $\mathcal P(\mathcal Q)$ contain the plaquettes wholly in
this box, and define the gauge-invariant trajectory event
\[
 \mathcal A=
 \left\{U:\ \sup_{0\le t\le s}
          |\Phi_t(U)_p-I|\le h
                   \text{ for every }p\in\mathcal P(\mathcal Q)\right\}.
 \tag{V19.6}\label{c18v19:event}
\]
Here $|\cdot|$ is the Euclidean unit-quaternion norm; in particular
$|U_p-I|^2=2(1-w_p)$, so center plaquettes are not treated as flat.

\begin{lemma}[Fixed local gauge through the actual trajectory]
\label{c18v19:gauge}
For each $U\in\mathcal A$ there is a single, time-independent
fine gauge transformation after which, throughout $0\le t\le s$,
every link needed by a plaquette incident to $\mathcal D_R$ obeys
\[
                  |W_e(t)-I|\le r:=3Lh .
 \tag{V19.7}\label{c18v19:r}
\]
It is permissible to extend this gauge arbitrarily outside
$\mathcal Q$. No derivative of a configuration-dependent gauge
choice is used in the conclusion.
\end{lemma}
\begin{proof}
Relabel the box as $\{0,\ldots,L\}^3$ for this argument.
At time zero take the rooted tree consisting of all direction-$1$
links, direction-$2$ links on $x_1=0$, and direction-$3$ links on
$x_1=x_2=0$. It has $(L+1)^3-1$ links and connects all vertices.
Recursive gauge transport makes precisely these links identity.

Every $2$-link at $x_1=0$ is then identity. A $12$-plaquette
compares the $2$-links at successive $x_1$, because both of its
$1$-links are identity. Multiplication by unit quaternions is
isometric, so the difference of the two $2$-links has norm
at most $h$. Telescoping gives $|W_2(0)-I|\le Lh$.
On $x_1=0$, a $23$-plaquette similarly propagates the bound on
$3$-links from $x_2=0$, giving $Lh$. The $13$-plaquettes then
propagate it in $x_1$, giving $|W_3(0)-I|\le2Lh$.
Consequently every box link is within $2Lh$ of identity initially.
This is a compact group telescoping argument, not an Abelian
replacement of plaquette products.

A link force is a sum of four transported imaginary plaquette
parts. On the event, each such part has norm at most $h$.
For the links in the assertion all their incident plaquettes are
still inside $\mathcal Q$: a link in $\mathcal D_R$ has a tail
in $[-R+1,\ell+R-1]^3$, and the two additional incidence steps
fit in the displayed box. Thus along the full, unfrozen flow
$|\dot W_e(t)|\le4h$ for each needed link. The fixed initial gauge
therefore gives
\[
 |W_e(t)-I|\le(2L+4s)h\le3Lh ,
\]
where $4s\le L$. The gauge can be extended on the remaining
vertices of the torus. Gauge equivariance transports the entire
trajectory and its tangent vectors by orthogonal maps.

For each fixed initial configuration this is a fixed isometry
used to compare derivatives at that point. We do not differentiate
the rule selecting the tree gauge, nor replace the nonlinear
map by a gauge-fixed derivative with an omitted term.
\end{proof}

\subsection{The native variational comparison, including its moving frame}

To compare to a fixed flat matrix, use the same invariant
orthonormal frame on each group factor at every time, rather
than V18's parallel frames. If $J(t)$ is the native flow
differential in these frames, its generator is
\[
 J'(t)=\mathsf A(t)J(t),\qquad
 \mathsf A(t)=-\mathsf H(t)+\mathsf S(t),\qquad
 \mathsf S(t)^*=-\mathsf S(t).
 \tag{V19.8}\label{c18v19:frame}
\]
The frame correction $\mathsf S$ is link-diagonal.
For the unit round group metric its block norm is at most the
link speed: the Levi-Civita correction has one half of the
Lie bracket, and $|[u,v]|\le2|u||v|$.
In particular
\[
 \|\mathsf A(t)\|\le20 ,
 \tag{V19.9}\label{c18v19:A-global}
\]
using the inherited Hessian norm $16$ and the global link-speed
bound $4$. Its off-diagonal support is still the link graph.

\begin{lemma}[Local compact-to-flat generator bound]
\label{c18v19:generator}
In the gauge of Lemma~\ref{c18v19:gauge}, the principal matrix
on $\mathcal D=\mathcal D_R$ satisfies
\[
 \|\mathsf A_{\mathcal D\mathcal D}(t)
                  +(\mathsf L_0)_{\mathcal D\mathcal D}\|
                         \le\gamma:=80r .
 \tag{V19.10}\label{c18v19:gamma}
\]
Only plaquettes incident to this domain are used.
\end{lemma}
\begin{proof}
For one plaquette the diagonal Riemannian Hessian block is
$w_p I_3$, so its difference from the identity block has norm
at most $4r$ when all four links are within $r$ of identity.
For two distinct links, the mixed directional derivative is
the real part of the plaquette word with two imaginary
quaternions inserted, with the appropriate signs for inverses.
Each of its four link factors is unitary and differs from
its identity value by at most $r$. Replacing them one at a
time bounds the difference of the bilinear forms by
$4r|u||v|$. This is an operator-norm bound $4r$, not an
entrywise bound subsequently multiplied by dimension.

Each stored link has four incident plaquettes and twelve
off-diagonal plaquette incidences. The symmetric Hessian
difference consequently has block row and column norm sums
at most $(4+12)4r=64r$ on the principal domain.
The block Schur estimate bounds its operator norm by $64r$.
Also every incident holonomy is within $4r$ of identity, so
the link speed is at most $16r$. The link-diagonal frame
correction is bounded by $16r$. Addition proves
\eqref{c18v19:gamma}. No sign or moving-frame term was dropped.
\end{proof}

Put $J_{\mathcal D}'=\mathsf A_{\mathcal D\mathcal D}J_{\mathcal D}$
along the \emph{same full} trajectory, with initial identity.
Let $T(t)=e^{-t\mathsf L_0}$ and
$T_{\mathcal D}(t)=e^{-t(\mathsf L_0)_{\mathcal D\mathcal D}}$.
The latter is a contraction, since the principal flat matrix
is positive semidefinite. The local comparison and Duhamel give
\[
 \|J_{\mathcal D}(s)\|\le e^{s\gamma},\qquad
 \|J_{\mathcal D}(s)-T_{\mathcal D}(s)\|
                                 \le s\gamma e^{s\gamma}.
 \tag{V19.11}\label{c18v19:duhamel}
\]
The time-ordered word argument of V18, now using the norm bound
$20$, gives
\[
 \begin{split}
 \|\mathbf1_{\mathcal B}(J(s)-T(s))\|
       &\le s\gamma e^{s\gamma}+4e^{60s-R},\\
 \|\mathbf1_{\mathcal B}J(s)\|
       &\le e^{s\gamma}+2e^{60s-R}.
 \end{split}
 \tag{V19.12}\label{c18v19:row}
\]
Indeed each of the two full-versus-principal row errors is at
most $2\sum_{j\ge R}(20s)^j/j!\le2e^{20es-R}\le2e^{60s-R}$.
The order-$j<R$ rows agree by graph support, even for noncommuting
generators. The exterior was used throughout the background
trajectory; it is not frozen or assigned a favorable configuration.

\subsection{A physical loop covector and a native two-sided estimate}

Write
\[
 w_\square(g)=\tfrac12\operatorname{Tr}g_\square,\qquad
 F_\square(U)=w_\square(\mathfrak c(U)),\qquad
 a(g)=\nabla_g w_\square(g).
 \tag{V19.13}\label{c18v19:loop}
\]
The vector $a$ is supported on the four coarse edges.
If $\rho=|\operatorname{Im}g_\square|$, each of its four
link components has norm $\rho$, so $|a|=2\rho$.
The fine differential $z=\nabla_W(w_\square\circ\mathfrak b_\ell)$
at $W=W_s$ has $4\ell$ components, each of norm $\rho$.
It is supported on $\mathcal B$ and has norm $\sqrt\ell\,|a|$.

In the fixed local gauge, let $z_0$ be the flat oriented
boundary current with the common color vector
$-\operatorname{Im}g_\square$ at its basepoint.
It has the same norm as $z$ and can be written
$z_0=B_\ell^*a_0$, where $|a_0|=|a|$ and
$d_{0,c}^*a_0=0$. The exact transported loop derivative gives
\[
 \|z-z_0\|\le10\ell r\sqrt\ell\,|a|.
 \tag{V19.14}\label{c18v19:covector}
\]
For completeness, each component is the common color vector,
with its orientation sign, conjugated by a partial loop product.
The product uses at most $4\ell$ links or their inverses.
Its distance from identity is at most $4\ell r$.
Since $\|\operatorname{Ad}_q-I\|\le2|q-I|$, with any endpoint
factor absorbed in the harmless coefficient $10\ell$, the
component error is at most $10\ell r\rho$.
Summing its square over $4\ell$ links proves
\eqref{c18v19:covector}.
At $\rho=0$ both differentials and the chosen comparison
current vanish, so the unnormalized estimates still apply.
No division by a small curvature component has occurred.

\begin{theorem}[Native physical loop metric on a paid trajectory event]
\label{c18v19:native}
With \eqref{c18v19:scales}--\eqref{c18v19:event}, every
$U\in\mathcal A$ satisfies
\[
 \boxed{\qquad
 \frac{m^2}{4\ell}|a(\mathfrak c(U))|^2
 \le|\nabla_UF_\square(U)|^2
 \le\frac9\ell|a(\mathfrak c(U))|^2 .
 \qquad}
 \tag{V19.15}\label{c18v19:native-metric}
\]
This compares actual physical loop differentials in the original
fine and coarse group metrics.
\end{theorem}
\begin{proof}
The chain rule gives $\nabla_UF_\square=J(s)^*z$.
By \eqref{c18v19:row} and \eqref{c18v19:covector}, with
$\epsilon=e^{60s-R}$,
\[
 \begin{split}
 \|J(s)^*z-T(s)z_0\|
 \le\sqrt\ell\,|a|\big[
 s\gamma e^{s\gamma}+4\epsilon
              +(e^{s\gamma}+2\epsilon)10\ell r\big].
 \end{split}
 \tag{V19.16}\label{c18v19:error}
\]
Here the row estimates are adjointed against vectors supported
on $\mathcal B$; all initial fine tangent directions are still
included.

The chosen constants give $s\gamma=80s(3Lh)=m/(15\ell)<1$
and $\epsilon\le m/(16\ell)<1$. Thus $e^{s\gamma}<3$,
$\ell\le s$, and the square bracket in
\eqref{c18v19:error} is at most
\[
 300sr+4\epsilon
 \le\frac{m}{4\ell}+\frac{m}{4\ell}
 =\frac{m}{2\ell}.
\]
By Lemma~\ref{c18v19:flat},
$m|a|/\sqrt\ell\le\|T(s)z_0\|\le2|a|/\sqrt\ell$.
The direct and reverse triangle inequalities now give
\[
 \frac{m}{2\sqrt\ell}|a|
 \le\|J(s)^*z\|
 \le\frac{2+m/2}{\sqrt\ell}|a|
 \le\frac3{\sqrt\ell}|a|.
\]
Squaring proves the assertion. The case $a=0$ was already
included in the exact loop differentiation.
\end{proof}

\subsection{Actual-vacuum probability and the conditional metric}

V18's proved supremum budget and
$|\mathcal P(\mathcal Q)|\le3(L+1)^3$ yield
\[
 \begin{split}
 \nu(\mathcal A^c)&\le\wp_{\ell,\beta},\\
 \wp_{\ell,\beta}
  &:=\min\left\{1,\
      \frac{6(L+1)^3(1+24\ell^2)\varepsilon_\beta}{h^2}
             \right\}
   \le\min\{1,C_*\,\ell^{18}\varepsilon_\beta\},\\
 C_*&=\frac{150\cdot3600^2\cdot200^5}{m^2}.
 \end{split}
 \tag{V19.17}\label{c18v19:probability}
\]
This is Markov's inequality for the supremum of each
$|W_p-I|^2=2\delta_p(W)$ followed by a union bound.
For the last estimate, $\pi^2<10$ and
$\log(8\pi\ell)\le26\ell$ give $R\le96\ell^2$ for $\ell\ge2$;
hence $L+1\le200\ell^2$ and $1+24\ell^2\le25\ell^2$.
Substitute $h^{-2}=3600^2\ell^6L^2/m^2$.
In particular
\[
 \ell=o(\beta^{1/36})
 \quad\Longrightarrow\quad
 \nu(\mathcal A^c)\longrightarrow0
 \tag{V19.18}\label{c18v19:power}
\]
uniformly over the ambient tori satisfying
\eqref{c18v19:embedding}. The constants are very conservative.
This is a sufficient mesoscopic range, not a physical
continuum scaling theorem.

Let $\mu=\mathfrak c_*\nu$ be the full coarse marginal and
$\overline G(g)=\mathbb E_\nu[D\mathfrak c\,D\mathfrak c^*
\mid\mathfrak c=g]$ its true kinetic metric. Off the set $a(g)=0$,
define the scalar metric in this physical loop direction by
\[
 q_\square(g)=
  \frac{a(g)^*\overline G(g)a(g)}{|a(g)|^2}
  =\mathbb E_\nu\left[
       \frac{|\nabla_UF_\square|^2}{|a(\mathfrak c)|^2}
                         \,\middle|\,\mathfrak c=g\right].
 \tag{V19.19}\label{c18v19:conditional}
\]
All retained links are conditioned upon. The omitted set has
$\mu$-measure zero: the coarse law has a smooth positive Haar
density and $a=0$ exactly when $g_\square=\pm I$.
Ratios may be set to zero there without changing any assertion.

\begin{corollary}[A native lower coefficient on most coarse configurations]
\label{c18v19:coarse-lower}
There is a measurable set $\mathcal C_{\ell,\beta}$ of coarse
configurations such that
\[
 \begin{split}
 \mu(\mathcal C_{\ell,\beta}^c)&\le\min\{1,2\wp_{\ell,\beta}\},\\
 q_\square(g)&\ge\frac{m^2}{8\ell}
                \quad\text{for $\mu$-a.e. }g\in\mathcal C_{\ell,\beta}.
 \end{split}
 \tag{V19.20}\label{c18v19:coarse-floor}
\]
\end{corollary}
\begin{proof}
Take $\mathcal C_{\ell,\beta}$ to be the set on which
$\nu(\mathcal A^c\mid\mathfrak c=g)\le1/2$, omitting the null
zero-gradient set. Conditional Markov bounds its complement
by $2\nu(\mathcal A^c)$. The integrand in
\eqref{c18v19:conditional} is nonnegative everywhere, and is at
least $m^2/(4\ell)$ on $\mathcal A$. This proves the lower bound.
An exceptional configuration is not granted a conditional
probability bound merely because its unconditional probability
is small; the conditional Markov step is essential.
\end{proof}

\begin{theorem}[True loop-direction metric moments at its physical scale]
\label{c18v19:moments}
Fix $p\ge1$ and $0<c<1/(64p)$. Along $\beta\to\infty$ with
$\ell^2\le c\log\beta$ and \eqref{c18v19:embedding},
\[
 \int(\ell q_\square(g)-9)_+^p\,d\mu(g)\longrightarrow0,
 \qquad
 \limsup\ell\|q_\square\|_{L^p(\mu)}\le9.
 \tag{V19.21}\label{c18v19:moment}
\]
The convergence is uniform in the location of the plaquette
and in the surrounding volume under the stated restriction.
\end{theorem}
\begin{proof}
Put $r_\square(U)=|\nabla_UF_\square|^2/
|a(\mathfrak c)|^2$ away from the null set.
Theorem~\ref{c18v19:native} gives $\ell r_\square\le9$
on $\mathcal A$. Globally V17's full metric envelope gives
$r_\square\le\ell e^{32s}$. Therefore
\[
 \mathbb E_\nu(\ell r_\square-9)_+^p
 \le\ell^{2p}e^{32ps}\wp_{\ell,\beta}
 \le C_*\ell^{18+2p}e^{32p\ell^2}\varepsilon_\beta .
 \tag{V19.22}\label{c18v19:moment-debit}
\]
For $T=\log\beta$ the last expression is at most
$C_{p,c}T^{9+p}\beta^{32pc-1/2}$, which tends to zero.
This explicitly pays the bad-event contribution; the power
range \eqref{c18v19:power} alone does not pay it.
Conditional Jensen in \eqref{c18v19:conditional} transfers
the positive-part moment to $q_\square$.
The triangle inequality applied to
$\ell q_\square\le9+(\ell q_\square-9)_+$
proves the second assertion.
\end{proof}

For a smooth real $f(g)=\chi(w_\square(g))$, the exact retained
form is
\[
 \langle J_{\ell,s}f,KJ_{\ell,s}f\rangle
       =\kappa\int q_\square(g)|\nabla_gf|^2\,d\mu(g).
 \tag{V19.23}\label{c18v19:form}
\]
Thus \eqref{c18v19:coarse-floor} supplies a lower bound by
$\kappa m^2/(8\ell)$ times the gradient integral over
$\mathcal C_{\ell,\beta}$. For $p>1$ in the logarithmic range,
H\"older and \eqref{c18v19:moment} give
\[
 \langle J_{\ell,s}f,KJ_{\ell,s}f\rangle
 \le\frac{\kappa}{\ell}(9+o(1))
             \|\nabla_gf\|_{L^{2p/(p-1)}(\mu)}^2 .
 \tag{V19.24}\label{c18v19:form-upper}
\]
The $p=1$ version uses the $L^\infty$ gradient norm.
Neither statement removes the actual marginal $\mu$ or
replaces the original Hamiltonian time.

\subsection{What changed, and the remaining conditional gap}

The new native lower bound is in a specified physical
Wilson-loop direction, not in a gauge direction where a lower
metric would be irrelevant to physical observables.
The relevant kinetic scale in this direction is $\kappa/\ell$,
not the unsmoothed full-link coefficient $\kappa\ell$.
Using the latter as a calibrated flowed kinetic coefficient
would introduce an unproved factor $\ell^2$ into a scale comparison.

Several limitations cannot be suppressed. The good event controls
one embedded neighborhood; it does not fix the global gauge,
discard torons or control all cubes simultaneously. The lower
metric holds on most coarse configurations, not everywhere.
Its exceptional set still requires a capacity or weighted-energy
estimate in a coercivity argument. The loop-gradient itself
vanishes at center holonomy; its variance in the actual marginal
has not been bounded here. Other physical directions and the
distribution along the full nonlinear discarded fibres remain
uncontrolled. No conditional Poincare floor, native discarded
resolvent inverse, or strict source-weighted return estimate
follows from \eqref{c18v19:form}.

The next iteration should attack that remaining distributional
or capacity estimate with the now-correct loop kinetic scale.
It should not treat \eqref{c18v19:coarse-floor} as an already
proved full coarse gap, or promote V17's Gaussian conditional
floor to the interacting ground law. The full nontrivial
four-dimensional continuum construction and positive physical
Yang--Mills mass gap remain unproved.

The preceding goal turn was progress: V18 populated actual-vacuum
metric moments, and its complete diagnostic passed again before
this append. The present diagnostic checks exact rational
compact axial gauges and loop gradients, the mixed Hessian
formula, flat physical alias symbols including their harmonic
primary, and the displayed scale constants. Floating operator
norms and finite configurations are regressions, not proofs of
an unknown vacuum spectrum. The analytic estimates above are
not proof-assistant certified.

The V18 body is preserved exactly apart from its title; all
47 other note files remain unchanged. V19 is the single active
source, and only \texttt{.tex} files are kept in
\texttt{latex\_notes}. V20 is due for both the companion review
and the broad literature sweep before choosing its next attack.
% END CYCLE 18 VERSION 19 APPEND
% BEGIN CYCLE 18 VERSION 20 APPEND
\clearpage
\section{V20: Native bad-set capacity and physical-time persistence}
\label{c18v20:section}

\subsection{Checkpoint, transfer decisions and the precise new target}

V19 obtained a native two-sided metric in one physical Wilson-loop
direction on a quantitatively probable event. Its remaining spectral
obligation is not discharged by probability alone. This appendix proves
a different, usable intermediate statement: an explicit small
\emph{upper} capacity bound for that event's complement, a stationary
physical-time exit bound, and errors for bounded-source correlations
and positive-shift resolvents of the original Hamiltonian. No global
spectral lower bound is inferred from an upper capacity bound.

The five-version checkpoint reread the terminal scope passages and
compared full-file hashes for NS V26, SM--Einstein V72, shell-mixing V16,
parallel YM V5, ESM V64 and its newer V69, SM source-selection V52,
native stationarity V54 and exact-action Einstein V45. All nine sources
are unchanged since V15. This is a targeted transfer review, not an
independent verification of every companion proof. Their useful common
discipline is to retain the actual law, source and clock: a prepared
linear history, finite source bank or fixed-time limit cannot be
promoted to a nonlinear, source-complete, uniform-time theorem.
The parallel YM extensive-charge obstruction continues to favor a
local window, not simultaneous control of the whole spatial torus.

The ten-version primary-literature sweep checked the following distinct
routes. The selections and exclusions matter more than a count of papers.

\begin{enumerate}
\item Lyons--Zheng's
\href{https://archive.numdam.org/item/AST_1988__157-158__249_0/}
{crossing-estimate paper}, especially its forward/backward
decomposition, supplies the method used below. We give its elementary
smooth finite-volume specialization and all native constants.
The probabilistic method is established literature, not a new YM theorem.
\item Leli\`evre--Zhang's
\href{https://arxiv.org/abs/1805.01928}
{nonlinear vectorial effective dynamics} requires, in Assumption 4,
a uniform conditional Poincar\'e inequality. That is still an unpaid
input here. Its pathwise theorem cannot furnish its own hypothesis.
\item Otto--Reznikoff's
\href{https://doi.org/10.1016/j.jfa.2006.10.002}
{two-scale logarithmic Sobolev criterion} uses single-site conditional
constants and mixed Hessian control. It remains a possible assembly
tool after these are proved for the actual vacuum density; the
classical Wilson action is not that density.
\item Ben-Artzi--Einav's
\href{https://arxiv.org/abs/1805.08557}
{weak Poincar\'e estimates without spectral gaps} reinforce the
distinction between useful decay control and a strictly positive
spectral bottom on the vacuum complement.
\item L\"ocherbach--Loukianova--Loukianov's
\href{https://arxiv.org/abs/1103.4622}
{spectral conditions, hitting times and Nash inequalities} connect
exit moments and spectral measures. Their polynomial-decay conclusions
do not turn our finite-time survival estimate into exponential mixing.
\item Le Peutrec's
\href{https://arxiv.org/abs/1607.08714}
{Witten Laplacians on manifolds with boundary} and
Di Ges\`u--Le Peutrec's
\href{https://arxiv.org/abs/1506.04434}
{large-system small-noise spectral asymptotics} suggest boundary-form
and metastability tools. We have not verified their quantitative
hypotheses for the unknown potential $-2\log\Omega$ or for our nonsmooth
good-region boundary. No spectral theorem is imported from them.
\item Shen--Zhu--Zhu's
\href{https://arxiv.org/abs/2204.12737}
{stochastic lattice Yang--Mills at strong coupling} was rechecked.
Its stated coupling regime and invariant Wilson measure do not
supply our weak-coupling quantum ground law. A search including newer
weak-coupling and bootstrap work supplied no verified substitute for
this missing native input.
\end{enumerate}

This changes the immediate attack from another conditional-gap
criterion to a native calculation that needs neither a bound on
$\nabla\log\Omega$ nor conditional mixing. It is not a repetition of
archive f12 V73: that appendix assumed a uniform entrance probability
from every bad starting point to deduce a killed spectral lower bound.
Here we prove a stationary probability of \emph{avoiding} the bad set,
and make no worst-starting-point assertion.

\subsection{The native event and its Lipschitz cutoff}

Use exactly the finite $\mathrm{SU}(2)$ torus, Hamiltonian, metric
normalization and actual ground law $\nu=\Omega^2dU$ of V18--V19.
Write $\beta=b/\kappa>0$, $s=\ell^2$, and
$\varepsilon_\beta=\min\{1,3\sqrt2\,\beta^{-1/2}\}$.
Retain the constants and embedding restriction
\[
 \begin{gathered}
 m=\frac2\pi e^{-3\pi^2},\qquad
 R=\left\lceil64\ell^2+\log\frac{16\ell}{m}\right\rceil,\qquad
 L=2R+\ell+4,\qquad h=\frac{m}{3600\ell^3L},\\
 \ell\ge2,\quad N=\ell n,\quad n\ge5,\quad N\ge L+1.
 \end{gathered}
 \tag{V20.1}\label{c18v20:scales}
\]
Let $\mathcal Q$ be V19's embedded box of side $L$ around the selected
coarse plaquette, and let $\Phi_t$ be the \emph{full} spatial Wilson
gradient flow. It is not frozen outside the box. Define
\[
 \begin{split}
 Z(U)&=\max_{p\in\mathcal P(\mathcal Q)}
          \sup_{0\le t\le s}|(\Phi_t(U))_p-I|,\\
 D&=\{Z<h\},\qquad B=\{Z\ge h\},\\
 \mathfrak b&=\frac{6(L+1)^3(1+24s)\varepsilon_\beta}{h^2},\qquad
 C_h=\frac{3600^2\,200^2}{m^2}.
 \end{split}
 \tag{V20.2}\label{c18v20:event}
\]
Thus $D$ is open, $B$ is closed, and $D$ is contained in V19's
good event $\mathcal A=\{Z\le h\}$. The untruncated quantity
$\mathfrak b$ is an upper budget, not a probability.

\begin{lemma}[A dimension-independent gradient for the trajectory maximum]
\label{c18v20:cutoff}
The function $Z$ is globally $2e^{16s}$-Lipschitz in the full product
Riemannian metric. There is a real gauge-invariant Lipschitz function
$0\le\chi\le1$, zero on $\{Z\le h/2\}$ and one on $B$, such that
\[
 \begin{split}
 \nu\{Z\ge h/2\}&\le\min\{1,4\mathfrak b\},\\
 \|\chi\|_{L^2(\nu)}^2&\le4\mathfrak b,\qquad
 |\nabla\chi|\le\frac{4e^{16s}}h\,
                       {\bf1}_{\{Z\ge h/2\}}\quad\text{a.e.},\\
 \mathcal E(\chi,\chi)
   :=\kappa\int|\nabla\chi|^2\,d\nu
       &\le\frac{64\kappa e^{32s}}{h^2}\mathfrak b .
 \end{split}
 \tag{V20.3}\label{c18v20:cutoff-cost}
\]
These estimates are independent of the surrounding volume subject
to \eqref{c18v20:scales}.
\end{lemma}
\begin{proof}
For $u_p=|U_p-I|=\sqrt{2(1-w_p)}$, the four distinct links of a
plaquette and the sphere Casimir normalization give, where $u_p>0$,
\[
 |\nabla u_p|^2
 =\frac{|\nabla w_p|^2}{2(1-w_p)}
 =\frac{4(1-w_p^2)}{2(1-w_p)}
 =2(1+w_p)\le4.
 \tag{V20.4}\label{c18v20:distance}
\]
There is no singularity in the Lipschitz assertion at $u_p=0$.
Holonomy multiplication is $2$-Lipschitz from its four product
factors into the bi-invariant group metric, and quaternion chord
distance to $I$ is $1$-Lipschitz relative to geodesic distance.
Equivalently one may regularize the square root and integrate its
gradient bound along geodesics.

V17's full derivative bound $\|D\Phi_t\|\le e^{16t}$ makes each
$u_p\circ\Phi_t$ $2e^{16s}$-Lipschitz. A supremum of functions with
the same Lipschitz constant has that same constant, including the
continuous time supremum. Compactness of $[0,s]$ also makes $Z$
continuous. No factor $|\mathcal P(\mathcal Q)|$ is incurred here.

V18's actual-vacuum supremum estimate
\eqref{c18v18:sup}, $u_p^2=2\delta_p$ and a union bound imply
\[
 \nu\{Z\ge a\}
 \le \frac{6(L+1)^3(1+24s)\varepsilon_\beta}{a^2}
 \quad(a>0).
\]
Take $a=h/2$. The required cutoff is
$\chi(U)=\min\{1,(2Z(U)/h-1)_+\}$. The Sobolev chain rule gives
its gradient bound; its nonzero set is contained in $\{Z\ge h/2\}$.
Squaring and integrating proves \eqref{c18v20:cutoff-cost}.
Gauge invariance follows from that of plaquette chord distance
and gauge equivariance of the full flow.
\end{proof}

For $\lambda>0$ with units of inverse physical time, let
$\operatorname{Cap}_{\lambda,\nu}(B)$ denote the capacity of the
closed form $\mathcal E+\lambda\|\cdot\|_2^2$: its admissible
functions are at least one quasi-everywhere on $B$. The continuous
cutoff above is admissible, so
\[
 \operatorname{Cap}_{\lambda,\nu}(B)
 \le 4\lambda\mathfrak b+
           \frac{64\kappa e^{32s}}{h^2}\mathfrak b .
 \tag{V20.5}\label{c18v20:capacity}
\]
The usual equivalent definition using neighborhoods gives the same
upper bound by using $(1+\eta)\chi$ and then $\eta\downarrow0$.
This is an absolute small-set capacity \emph{upper bound}. It is
not a lower bound on capacity divided by mass for arbitrary sets,
which is the direction needed for a Poincar\'e inequality.

\subsection{The process has the original physical Hamiltonian clock}

Multiplication $\mathcal U f=\Omega f$ is unitary from $L^2(\nu)$
to $L^2(dU)$. At every fixed finite regulator, elliptic positivity
gives a smooth strictly positive $\Omega$ and
\[
 \begin{split}
 \mathcal U^{-1}K\mathcal U
   &=-\kappa\big(\Delta+2\nabla\log\Omega\cdot\nabla\big)
     =-\mathcal L,\\
 P_\tau&=\mathcal U^{-1}e^{-\tau K}\mathcal U,\qquad
 \langle f,-\mathcal L f\rangle_\nu=\mathcal E(f,f).
 \end{split}
 \tag{V20.6}\label{c18v20:physical-process}
\]
Here $\Delta$ is the nonpositive Laplace--Beltrami operator.
Expanding $H(\Omega f)$ and using $H\Omega=E_0\Omega$ proves the
identity, first for smooth functions and then by form closure.
The smooth uniformly elliptic diffusion with generator $\mathcal L$
exists globally on the compact group product. Integration by parts
makes it reversible with invariant probability $\nu$. It preserves
the gauge-invariant sector.

Write this diffusion as $X_\tau$, starting with law $\nu$.
Its Euclidean time $\tau$ is the time in $e^{-\tau K}$,
with $\kappa$ retained. It is distinct from the spatial smoothing
parameter $t\le s$ inside $Z(X_\tau)$. In particular the event below
controls a family of full spatial-flow trajectories, one from each
physical-time configuration. It is neither a Wilson-Gibbs
Langevin process nor a proposed Markov process on the coarse links.

\begin{lemma}[Stationary maximal estimate]
\label{c18v20:maximal}
For a real Lipschitz $f$ on this finite product and $T\ge0$,
\[
 \mathbb E_\nu\sup_{0\le\tau\le T}|f(X_\tau)|^2
       \le2\|f\|_2^2+36T\mathcal E(f,f).
 \tag{V20.7}\label{c18v20:maximal-bound}
\]
\end{lemma}
\begin{proof}
For smooth $f$ and $T>0$, set
\[
 \begin{split}
 M_\tau&=f(X_\tau)-f(X_0)-\int_0^\tau\mathcal Lf(X_r)\,dr,\\
 \widehat M_u&=f(X_{T-u})-f(X_T)
                   -\int_{T-u}^T\mathcal Lf(X_r)\,dr .
 \end{split}
\]
Reversibility makes both square-integrable martingales in their
respective forward filtrations, and
$\mathbb E|M_T|^2=\mathbb E|\widehat M_T|^2=2T\mathcal E(f,f)$.
Cancellation of the two drift integrals gives
\[
 f(X_\tau)-f(X_0)
       =\tfrac12\{M_\tau-(\widehat M_T-\widehat M_{T-\tau})\}.
\]
Doob's inequality bounds the $L^2$ norms of the two martingale
suprema by $2\sqrt{2T\mathcal E(f,f)}$ each. Bounding a reversed
increment by twice its supremum therefore gives
\[
 \left\|\sup_{\tau\le T}|f(X_\tau)|\right\|_2
       \le\|f\|_2+3\sqrt{2T\mathcal E(f,f)} .
\]
Square using $(a+b)^2\le2a^2+2b^2$. Smooth approximation on the compact
manifold converges uniformly and in $H^1(\nu)$ for Lipschitz $f$;
uniform convergence and continuous paths pass the bound to $f$.
The case $T=0$ is immediate.
\end{proof}

The drift cancels before any estimate. No uniform control of
$\nabla\log\Omega$, no mixing rate and no spectral lower bound was
assumed to obtain \eqref{c18v20:maximal-bound}.

\begin{theorem}[A paid native physical-time exit bound]
\label{c18v20:exit}
Set $\tau_D=\inf\{\tau\ge0:X_\tau\notin D\}$, including $\tau_D=0$
when $X_0\in B$. Then
\[
 \begin{split}
 \mathbb P_\nu(\tau_D\le T)&\le\eta_{\ell,\beta}(T),\\
 \eta_{\ell,\beta}(T)
   &:=\min\left\{1,\,
       8\mathfrak b+
          \frac{2304\kappa T e^{32\ell^2}}{h^2}\mathfrak b
          \right\},\\
 \mathfrak b&\le C_*\ell^{18}\varepsilon_\beta,\qquad
 h^{-2}\le C_h\ell^{10}.
 \end{split}
 \tag{V20.8}\label{c18v20:exit-bound}
\]
Here $C_*$ is exactly the constant in \eqref{c18v19:probability}.
Consequently, along any family satisfying
\[
 \beta\longrightarrow\infty,\qquad
 \ell^2\le c\log\beta,\qquad
 \kappa T\le\beta^d,\qquad
 c>0,\quad d\ge0,\quad d+32c<\tfrac12,
 \tag{V20.9}\label{c18v20:time-window}
\]
and \eqref{c18v20:scales}, the exit probability tends to zero
uniformly in the ambient volume and the plaquette location.
\end{theorem}
\begin{proof}
Continuity makes $\{\tau_D\le T\}$ the event of visiting $B$ by time
$T$. On it $\sup_{\tau\le T}\chi(X_\tau)=1$. Apply
\eqref{c18v20:maximal-bound} and \eqref{c18v20:cutoff-cost}.
The polynomial bounds follow from $L+1\le200\ell^2$ and the V19
calculation, without any new probability assumption. Explicitly,
\[
 \eta_{\ell,\beta}(T)
 \le\min\left\{1,\,
 C_*\varepsilon_\beta
       \big(8\ell^{18}+2304C_h\kappa T\ell^{28}e^{32\ell^2}\big)
            \right\}.
 \tag{V20.10}\label{c18v20:polynomial}
\]
In the stated range its untruncated right side is
\[
 O_c\!\left((\log\beta)^9\beta^{-1/2}
       +(\log\beta)^{14}\beta^{d+32c-1/2}\right),
 \tag{V20.11}\label{c18v20:rate}
\]
which tends to zero. No large-time limit at fixed $\beta$ is taken.
\end{proof}

For instance $c=1/256$ and $d=1/8$ satisfy the strict inequality.
This is a sufficient asymptotic range with extremely conservative
constants, not a practically sized coupling threshold. If a continuum
trajectory makes $\kappa=\kappa(a)$ grow, the allowed physical horizon
is $T\le\beta^d/\kappa(a)$. It may shrink. Nothing here identifies
$\kappa(a)$ or proves a fixed physical continuum-time window.

On the no-exit event V19's physical loop metric holds at
\emph{every} $X_\tau$, $0\le\tau\le T$:
\[
 \frac{m^2}{4\ell}|a(\mathfrak c(X_\tau))|^2
 \le|\nabla F_\square(X_\tau)|^2
 \le\frac9\ell|a(\mathfrak c(X_\tau))|^2 .
 \tag{V20.12}\label{c18v20:metric-path}
\]
Thus the scale $\kappa/\ell$ persists along most such stationary
physical paths. This remains one physical direction, not a full
coercive metric.

\subsection{Correlation and resolvent errors, without changing the law}

Let $P_T^D$ be the semigroup of this same diffusion killed at $\tau_D$,
extended by zero outside $D$:
\[
 P_T^D G(x)
    =\mathbb E_x\!\left[G(X_T){\bf1}_{\{\tau_D>T\}}\right]
       \quad(T>0).
 \tag{V20.13}\label{c18v20:killed}
\]
It is the part semigroup of the native closed form on $D$, with
Dirichlet boundary condition. The form domain is the closure of
$C_c^\infty(D)$ in $H^1(\nu)$; equivalently functions vanish
quasi-everywhere off $D$. Neither a smooth boundary nor a chosen
conditional distribution is needed. This identification follows by
stopping the continuous diffusion on relatively compact subdomains
and closing the Dirichlet form. Gauge invariance of $D$ makes the
restriction compatible with the physical sector.

\begin{proposition}[Paid errors for bounded physical sources]
\label{c18v20:correlations}
For bounded measurable, possibly complex, $F,G$ and $T>0$,
\[
 \begin{split}
 |\langle F,(P_T-P_T^D)G\rangle_\nu|
    &\le\|F\|_\infty\|G\|_\infty\,\eta_{\ell,\beta}(T),\\
 \|(P_T-P_T^D)G\|_2
    &\le\|G\|_\infty\sqrt{\eta_{\ell,\beta}(T)} .
 \end{split}
 \tag{V20.14}\label{c18v20:correlation-error}
\]
For $\lambda>0$, put
$\mathcal R_\lambda=\int_0^\infty e^{-\lambda T}P_T\,dT$
and $\mathcal R_\lambda^D=\int_0^\infty e^{-\lambda T}P_T^D\,dT$.
Then
\[
 \begin{split}
 |\langle F,(\mathcal R_\lambda-\mathcal R_\lambda^D)G\rangle_\nu|
 &\le\|F\|_\infty\|G\|_\infty
 \min\left\{\frac1\lambda,\,
     \frac{8\mathfrak b}{\lambda}
     +\frac{2304\kappa e^{32\ell^2}\mathfrak b}{h^2\lambda^2}
        \right\}.
 \end{split}
 \tag{V20.15}\label{c18v20:resolvent-error}
\]
Under the scale limit in \eqref{c18v20:time-window}, with
$\kappa T\le\beta^d$ replaced by $\kappa/\lambda\le\beta^d$,
the coefficient of $\|F\|_\infty\|G\|_\infty$ in
\eqref{c18v20:resolvent-error}, multiplied by $\lambda$, tends to zero.
\end{proposition}
\begin{proof}
The difference kernel records paths that have exited by $T$.
With $p_T(x)=\mathbb P_x(\tau_D\le T)$ this gives
$|(P_T-P_T^D)G(x)|\le\|G\|_\infty p_T(x)$.
Integrate against $|F|\nu$ for the first bound; use
$0\le p_T^2\le p_T$ for the second. Laplace integration of the
untruncated affine upper bound in \eqref{c18v20:exit-bound}
gives the second entry of the minimum in
\eqref{c18v20:resolvent-error}; $p_T\le1$ gives the first.
The same polynomial calculation as \eqref{c18v20:rate} proves
the final assertion. All integrals converge for $\lambda>0$.
\end{proof}

Under $\mathcal U$, these are statements about
$\langle\Omega F,e^{-TK}\Omega G\rangle_{dU}$
and $(K+\lambda)^{-1}$, compared with their native Dirichlet
restrictions. They apply, in particular, to bounded gauge-invariant
functions of the flowed coarse Wilson plaquette. The killed
operator is not the coarse compression of $K$, and its resolvent
is not a conditional fibre inverse. The displayed resolvent
estimate is for a \emph{positive shift}, not for $K^{-1}$ on the
vacuum complement. The powers of $\lambda^{-1}$ cannot be dropped.

\subsection{Why this does not prove an excited spectral floor}

Small stationary exit probability is not a small $L^2$ operator
error. Indeed $P_T\to I$ strongly as $T\downarrow0$, whereas the
zero-extended $P_T^D\to{\bf1}_D$ strongly. The set $\{Z>h\}$ contains
a nonempty open set: choose one plaquette holonomy $-I$ at spatial
time zero. The positive density $\nu$ therefore gives $B$ positive
measure. A normalized gauge-invariant indicator supported on $B$
has $P_T^D f=0$ and $\|P_Tf\|_2\to1$. Hence
\[
 \liminf_{T\downarrow0}\|P_T-P_T^D\|_{2\to2}\ge1 .
 \tag{V20.16}\label{c18v20:no-operator}
\]
There is no contradiction with \eqref{c18v20:correlation-error}:
the normalized indicator has $L^\infty$ norm $\nu(B)^{-1/2}$,
which the bounded-source error estimate must retain.

There is also a useful native variational warning.
Let $\mathsf A_D$ be the positive Dirichlet generator of $P_T^D$.
When $4\mathfrak b<1$,
\[
 \inf\sigma(\mathsf A_D)
 \le\frac{\mathcal E(\chi,\chi)}{(1-\|\chi\|_2)^2}
 \le\frac{64\kappa e^{32\ell^2}\mathfrak b}
                 {h^2(1-2\sqrt{\mathfrak b})^2}.
 \tag{V20.17}\label{c18v20:killed-bottom}
\]
To prove this, use $u=1-\chi$. It belongs to the part form domain:
$(u-\delta)_+$ has compact support in $D$ for $\delta>0$, can be
smoothly approximated there, and converges to $u$ in $H^1(\nu)$.
Its energy equals that of $\chi$, and
$\|u\|_2\ge1-\|\chi\|_2$. Rayleigh's principle gives the claim,
also within the gauge-invariant sector.
In the logarithmic range with $c<1/64$, its right side divided by
$\kappa$ tends to zero. This describes a near-vacuum Dirichlet mode;
it neither supplies nor refutes a positive physical excited gap.
Removing the true vacuum before a spectral comparison is indispensable.

The newly paid obligation is therefore
\emph{native small upper capacity and stationary physical path
persistence for one embedded V19 region}. What remains unpaid is
coercivity for arbitrary vacuum-orthogonal $L^2$ sources, including
sources concentrated near exceptional regions, and conditional
variance on the actual nonlinear fibres. A rich physical direction
bank, its true marginal variance, source-weighted return, scale
transport and interacting continuum reconstruction are still
separate obligations. The full Yang--Mills mass gap remains unproved.

The next useful upstream question is whether the true
vacuum-orthogonal low-energy spectral window can concentrate on
the transition layer of $\chi$. Bounded-observable estimates alone
do not answer it. A native energy-dependent localization estimate,
or a capacity-to-mass lower estimate for competing components,
would connect this appendix to an actual spectral lower argument.
Do not repeat the abstract escape criterion without populating
its worst-starting-point or energy-weighted input.

\subsection{Verification and version handoff}

The complete V19 diagnostic passed before this append. The new
diagnostic checks exact compact plaquette distance identities,
cutoff and scale constants, forward/backward cancellation, and
a reversible killed-chain regression for the kernel inequalities.
That chain is a test of algebra and norm conventions, not a
simulation of $\Omega^2dU$. The rigorous claims rely on the proofs
and stated inherited inputs above; there is no proof-assistant
certification or numerical verification of an unknown YM spectrum.

The V19 source is recoverable byte for byte after restoring its
title; all 47 other note files are preserved. V20 is the sole active
source in this cycle. Build and review intermediates are outside
\texttt{latex\_notes}, which contains only \texttt{.tex} files.
Both scheduled checkpoints were performed here. The next companion
checkpoint is V25, the next broad literature sweep V30, and the
next archive/restart remains V100.
% END CYCLE 18 VERSION 20 APPEND
% BEGIN CYCLE 18 VERSION 21 APPEND
\clearpage
\section{V21: Native low-energy localization beyond bounded sources}
\label{c18v21:section}

\subsection{Direction and inherited inputs}

V20 proves stationary exit and bounded-source estimates, but a
normalized low-energy wavefunction need not have a controlled
$L^\infty$ norm. This appendix pays a native, energy-weighted part
of that missing step. A deterministic spatial average of the full-flow
bad events is controlled by the positive magnetic term of the original
Hamiltonian. Translation characters turn this average into a bound
for each specified local region.

We control both squared wavefunction mass and gradient energy in the
bad region, and the error of restricting a low-energy vector to the
good region. The individual-region bounds hold on a fixed-momentum
spectral window. For arbitrary superpositions of momenta, the bounds
hold after spatial averaging. No uniform spectral lower bound is assumed.

The preceding turn was progress, and V20's complete diagnostic passed
before this append. Its companion and broad literature checkpoints
remain current; the next are V25 and V30.
The scalar localization identity below is standard; see Solovej,
\emph{The ionization conjecture in Hartree--Fock theory}, Lemma 2.4,
equation (19),
\href{https://annals.math.princeton.edu/wp-content/uploads/annals-v158-n2-p04.pdf}
{Annals of Mathematics 158 (2003), 509--576}.
No atomic spectral estimate is imported.
Archive f1 V12 already records weighted one-form IMS and its unpaid
costs; f15 V74 already warns about rare compact packets in a different
thermal completed form. The new result is not another abstract
localization identity: it pays a scalar-source cost for this native
Hamiltonian and this nonlinear full-flow event.

Keep V20's $N=\ell n$, $\ell\ge2$, $n\ge5$, $s=\ell^2$, embedded
box $\mathcal Q$ of side $L$, and threshold $h$, with
$N\ge L+1$ and \eqref{c18v20:scales}. Write
\[
 \begin{gathered}
 \mathsf M=3N^3,\qquad V_{\rm sp}=\sum_{p\in\mathcal P_N}\delta_p,
 \qquad H=\kappa\sum_eL_e+bV_{\rm sp},\qquad K=H-E_0,\\
 A=\mathcal U^{-1}K\mathcal U,\qquad \mathcal U f=\Omega f,\qquad
 \mathcal E(f,f)=\kappa\int|\nabla f|^2d\nu,\quad \nu=\Omega^2dU,\\
 E_0/(b\mathsf M)\le\varepsilon_\beta
       =\min\{1,3\sqrt2\,\beta^{-1/2}\},\qquad\beta=b/\kappa>0 .
 \end{gathered}
 \tag{V21.1}\label{c18v21:operators}
\]
The last line is the total-energy full-Haar trial bound recalled in V6,
not merely the vacuum's magnetic expectation. Functions may be complex.

\subsection{A deterministic spatial average}

For all $N^3$ fine translations $x$, define
\[
 \begin{split}
 \mathcal Q_x&=\mathcal Q+x,\qquad
 Z_x(U)=\max_{p\in\mathcal P(\mathcal Q_x)}
                  \sup_{0\le t\le s}|(\Phi_t(U))_p-I|,\\
 B_x&=\{Z_x\ge h\},\quad T_x=\{Z_x\ge h/2\},\quad D_x=\{Z_x<h\},\\
 \chi_x&=\min\{1,(2Z_x/h-1)_+\},\quad u_x=1-\chi_x,\qquad
 \mathfrak C_\ell=\frac{6L^2(L+1)(1+24s)}{h^2}.
 \end{split}
 \tag{V21.2}\label{c18v21:cutoffs}
\]
Here $T_x$ is a set. A translated box remains embedded even if it
crosses the chosen coordinate origin. In using V19's metric, it
corresponds to a translated origin of the coarse sampling map,
not to a new claim about a fixed retained grid.

\begin{lemma}[Deterministic full-flow rare-region average]
\label{c18v21:deterministic}
For every configuration,
\[
 \begin{split}
 \sum_{p\in\mathcal P_N}\sup_{t\le s}\delta_p(\Phi_t(U))
       &\le(1+24s)V_{\rm sp}(U),\\
 \frac1{N^3}\sum_x{\bf1}_{B_x}(U)
       &\le\frac{\mathfrak C_\ell}{\mathsf M}V_{\rm sp}(U),\\
 \frac1{N^3}\sum_x{\bf1}_{T_x}(U)
       &\le\frac{4\mathfrak C_\ell}{\mathsf M}V_{\rm sp}(U).
 \end{split}
 \tag{V21.3}\label{c18v21:average}
\]
\end{lemma}
\begin{proof}
V18's deterministic inequality is
$(\delta_p')_+\le12\delta_p+\sum_{q\sim p}\delta_q$.
After integration and summation, each neighbor is counted twelve
times. Full spatial flow decreases $V_{\rm sp}$ pointwise. The sum
of suprema is thus at most
$V_{\rm sp}(U)+24\int_0^sV_{\rm sp}(\Phi_t(U))dt$.

A box of side $L$ has $L^2(L+1)$ plaquettes of each orientation.
Each fixed torus plaquette occurs in exactly that many translated
boxes. For $a>0$, the indicator union bound gives
\[
 \begin{split}
 \frac1{N^3}\sum_x{\bf1}_{\{Z_x\ge a\}}
 &\le\frac{2}{a^2N^3}\sum_x
       \sum_{p\in\mathcal P(\mathcal Q_x)}
                       \sup_{t\le s}\delta_p(\Phi_t(U))\\
 &=\frac{2L^2(L+1)}{a^2N^3}
                       \sum_p\sup_{t\le s}\delta_p(\Phi_t(U)).
 \end{split}
\]
Take $a=h,h/2$ and use $\mathsf M=3N^3$.
This uses no probability law, independence, or cubic symmetry of
an excited state.
\end{proof}

The constants satisfy
\[
 \mathfrak C_\ell\le\frac{6(L+1)^3(1+24s)}{h^2}
          \le C_*\ell^{18},\qquad h^{-2}\le C_h\ell^{10},
 \tag{V21.4}\label{c18v21:constants}
\]
with V19's $C_*$ and V20's $C_h$. In particular
$\mathfrak C_\ell\varepsilon_\beta\le\mathfrak b$ of V20.

\subsection{Native energy-weighted compression}

For every form-domain $f$, positivity in the original $H$ gives
\[
 b\int V_{\rm sp}|f|^2d\nu
 \le\langle\Omega f,H\Omega f\rangle_{dU}
 =E_0\|f\|_2^2+\mathcal E(f,f).
 \tag{V21.5}\label{c18v21:bare}
\]
The term $E_0$ cannot be dropped.
Let $\mathsf T_y$ denote unitary spatial translations on $L^2(\nu)$.
Uniqueness makes $\Omega$ translation invariant, so $A$ commutes
with every $\mathsf T_y$. Put
\[
 \begin{gathered}
 \mathcal H_k=\{f:\mathsf T_yf=e^{ik\cdot y}f\ \forall y\},\qquad
 k\in(2\pi/N)(\mathbb Z/N\mathbb Z)^3,\\
 P_k=\text{orthogonal projection onto }\mathcal H_k .
 \end{gathered}
 \tag{V21.6}\label{c18v21:momentum}
\]
For $f\in\mathcal H_k$, $|f|^2d\nu$ is translation invariant,
but need not be invariant under rotations or spatial flow.

\begin{theorem}[Actual low-energy wavefunction weights]
\label{c18v21:weighted}
For any form-domain $f\in\mathcal H_k$ and every $x$, put
\[
 Q_\ell[f]=\mathfrak C_\ell
       \left(\varepsilon_\beta\|f\|_2^2+
                            \frac{\mathcal E(f,f)}{b\mathsf M}\right).
 \tag{V21.7}\label{c18v21:Q}
\]
Then
\[
 \begin{split}
 \int_{B_x}|f|^2d\nu&\le Q_\ell[f],\\
 \int_{T_x}|f|^2d\nu&\le4Q_\ell[f],\qquad
                          \|\chi_xf\|_2^2\le4Q_\ell[f],\\
 \kappa\int|\nabla\chi_x|^2|f|^2d\nu
       &\le64\kappa e^{32s}h^{-2}Q_\ell[f].
 \end{split}
 \tag{V21.8}\label{c18v21:weighted-bounds}
\]
For arbitrary form-domain $f$, the same assertions hold with each
left side averaged over $x$.
\end{theorem}
\begin{proof}
Integrate \eqref{c18v21:average} against $|f|^2d\nu$ and apply
\eqref{c18v21:bare} and the trial bound. Use
$\chi_x^2\le{\bf1}_{T_x}$ and V20's
$|\nabla\chi_x|^2\le16e^{32s}h^{-2}{\bf1}_{T_x}$.
For $f\in\mathcal H_k$, translation covariance makes the summands
equal, including the gradient summands since translations are
isometries of the full product metric.
\end{proof}

The sharper bound retaining actual $E_0$ is the quadratic-form
inequality on $\mathcal H_k$
\[
 P_k{\bf1}_{B_x}P_k
 \le\frac{\mathfrak C_\ell}{b\mathsf M}P_k(A+E_0)P_k .
 \tag{V21.9}\label{c18v21:compression}
\]
This does not say that the local indicator commutes with momentum.
For an energy ceiling $\Lambda\ge0$, define
\[
 \begin{gathered}
 \Pi_{\Lambda,k}={\bf1}_{[0,\Lambda]}(A)P_k,\qquad
 q_\Lambda=\mathfrak C_\ell
              \left(\varepsilon_\beta+\frac{\Lambda}{b\mathsf M}\right),\\
 d_\Lambda=64\kappa e^{32s}h^{-2}q_\Lambda .
 \end{gathered}
 \tag{V21.10}\label{c18v21:debits}
\]
These are untruncated budgets; $d_\Lambda$ has units of energy.
For example
\[
 \|{\bf1}_{B_x}\Pi_{\Lambda,k}\|_{2\to2}^2\le\min\{1,q_\Lambda\},
 \qquad
 \|\chi_x\Pi_{\Lambda,k}\|_{2\to2}^2\le\min\{1,4q_\Lambda\}.
 \tag{V21.11}\label{c18v21:restricted}
\]
These are restricted spectral-subspace norms, not the unrestricted
norm ruled out in V20, and have no $\|f\|_\infty$ factor.

\subsection{Gradient-energy leakage and form localization}

For real Lipschitz $v$ and $f\in\operatorname{Dom}A$, expansion of
the two gradients and taking real parts gives the scalar identity
\[
 \mathcal E(vf,vf)=\operatorname{Re}\langle Af,v^2f\rangle_\nu
                    +\kappa\int|\nabla v|^2|f|^2d\nu .
 \tag{V21.12}\label{c18v21:IMS}
\]
Indeed the mixed terms in $\nabla(vf)$ and $\nabla(v^2f)$ agree.
Lipschitz multiplication preserves $H^1$, and the form
representation of $A$ proves the assertion by approximation,
also for complex $f$. No derivative of $\Omega$ is estimated.

\begin{theorem}[Paid norm, form and bad-gradient errors]
\label{c18v21:localization}
For $f\in\operatorname{Ran}\Pi_{\Lambda,k}$ and any $x$,
$u_xf$ belongs to the native Dirichlet form domain on $D_x$ and
\[
 \begin{split}
 \|f-u_xf\|_2^2&\le4q_\Lambda\|f\|_2^2,\\
 (1-4q_\Lambda)\|f\|_2^2
     &\le\|u_xf\|_2^2\le\|f\|_2^2,\\
 |\mathcal E(u_xf,u_xf)-\mathcal E(f,f)|
     &\le(2\Lambda\sqrt{q_\Lambda}+d_\Lambda)\|f\|_2^2,\\
 \mathcal E(u_xf,u_xf)&\le(\Lambda+d_\Lambda)\|f\|_2^2,\\
 \kappa\int_{B_x}|\nabla f|^2d\nu
     &\le(4\Lambda\sqrt{q_\Lambda}+4d_\Lambda)\|f\|_2^2 .
 \end{split}
 \tag{V21.13}\label{c18v21:local-errors}
\]
For an eigenvector $Af=Ef$, $0\le E\le\Lambda$, one has the
sharper statements
\[
 \begin{split}
 \kappa\int_{B_x}|\nabla f|^2d\nu
       &\le(8E q_E+4d_E)\|f\|_2^2,\\
 \mathcal E(u_xf,u_xf)
       &=E\|u_xf\|_2^2+
                   \kappa\int|\nabla\chi_x|^2|f|^2d\nu .
 \end{split}
 \tag{V21.14}\label{c18v21:eigen}
\]
\end{theorem}
\begin{proof}
At fixed finite regulator, bounded spectral subspaces consist
of smooth functions by compact elliptic spectral theory.
As in V20, $(u_x-\delta)_+f$ has compact support in $D_x$ and
converges to $u_xf$ in $H^1(\nu)$, proving the domain assertion.

The norm difference is $\|\chi_xf\|_2$. Moreover
$0\le1-u_x^2\le{\bf1}_{T_x}$ and
$(1-u_x^2)^2\le{\bf1}_{T_x}$. Hence
\[
 \|f\|_2^2-\|u_xf\|_2^2\le4q_\Lambda\|f\|_2^2,\qquad
 \|(1-u_x^2)f\|_2\le2\sqrt{q_\Lambda}\|f\|_2 .
\]
Use $\|Af\|_2\le\Lambda\|f\|_2$ in \eqref{c18v21:IMS} with
$v=u_x$ for the form difference; use
$\|u_x^2f\|_2\le\|f\|_2$ directly for the upper form bound.

Applying the same identity with $v=\chi_x$ and
$\|\chi_x^2f\|_2\le\|\chi_xf\|_2$ gives
\[
 \mathcal E(\chi_xf,\chi_xf)
       \le(2\Lambda\sqrt{q_\Lambda}+d_\Lambda)\|f\|_2^2.
\]
Since $\chi_x=1$ on $B_x$ and, almost everywhere,
\[
 \chi_x^2|\nabla f|^2
   \le2|\nabla(\chi_xf)|^2+2|\nabla\chi_x|^2|f|^2,
\]
integration proves the bad-gradient bound. If $Af=Ef$, instead use
$\operatorname{Re}\langle Af,\chi_x^2f\rangle_\nu
=E\|\chi_xf\|_2^2\le4E q_E\|f\|_2^2$.
The second eigenvector statement is the exact product identity.
\end{proof}

For $4q_\Lambda<1$, multiplication by $u_x$ is injective on the
whole spectral subspace and preserves its dimension. It need not
preserve momentum or vacuum orthogonality. If $f\perp1$, then
\[
 |\langle1,u_xf\rangle_\nu|
   =|\langle\chi_x,f\rangle_\nu|
   \le2\sqrt{q_0}\|f\|_2 .
 \tag{V21.15}\label{c18v21:vacuum-overlap}
\]
Small overlap is not exact removal of the killed operator's ground.

For arbitrary low-energy superpositions put
$\Pi_\Lambda={\bf1}_{[0,\Lambda]}(A)$ and define
\[
 \mathcal Jf=N^{-3/2}(u_xf)_x
         \in\bigoplus_x H_0^1(D_x,\nu).
\]
For $f\in\operatorname{Ran}\Pi_\Lambda$,
\[
 \begin{split}
 (1-4q_\Lambda)\|f\|_2^2
   &\le\|\mathcal Jf\|^2\le\|f\|_2^2,\\
 \left|\frac1{N^3}\sum_x\mathcal E(u_xf,u_xf)-\mathcal E(f,f)\right|
   &\le(2\Lambda\sqrt{q_\Lambda}+d_\Lambda)\|f\|_2^2,\\
 \frac{\kappa}{N^3}\sum_x\int_{B_x}|\nabla f|^2d\nu
   &\le(4\Lambda\sqrt{q_\Lambda}+4d_\Lambda)\|f\|_2^2 .
 \end{split}
 \tag{V21.16}\label{c18v21:averaged-form}
\]
To verify this, replace the individual weights in the preceding proof
by Theorem~\ref{c18v21:weighted}'s averaged weights and use
\[
 \frac1{N^3}\sum_x\|(1-u_x^2)f\|_2
 \le\left(\frac1{N^3}\sum_x\|(1-u_x^2)f\|_2^2\right)^{1/2}
 \le2\sqrt{q_\Lambda}\|f\|_2 .
\]
The same argument with $\chi_x^2f$ proves the gradient assertion.
This is a map into a direct sum of local restrictions, not a single
global good-field event. That direct sum has additional near-vacuum
directions and cannot be assigned the original vacuum multiplicity.

\subsection{A shrinking energy window with relative errors}

\begin{corollary}[Localization at an explicit target energy]
\label{c18v21:scale}
Fix $a>0$, $d\ge0$, $c>0$ with $d+32c<1/2$. Along
\[
 \beta\to\infty,\qquad\ell^2\le c\log\beta,\qquad
 \gamma_\beta=\kappa\beta^{-d},\qquad\Lambda=a\gamma_\beta,
 \tag{V21.17}\label{c18v21:window}
\]
subject to V20's embedding restriction,
$q_\Lambda\to0$ and $d_\Lambda/\gamma_\beta\to0$.
The state-norm errors vanish, and the form and bad-gradient errors
in \eqref{c18v21:local-errors}--\eqref{c18v21:eigen} and
\eqref{c18v21:averaged-form} are
$o(\gamma_\beta)\|f\|_2^2$, uniformly in ambient volume,
momentum sector and box origin.
\end{corollary}
\begin{proof}
Since $b=\kappa\beta$,
$\Lambda/(b\mathsf M)=a\beta^{-1-d}/\mathsf M$.
The explicit bounds are
\[
 \begin{split}
 q_\Lambda&\le C_*\ell^{18}
    \left(3\sqrt2\,\beta^{-1/2}+a\beta^{-1-d}/\mathsf M\right),\\
 d_\Lambda/\gamma_\beta&\le64C_hC_*\ell^{28}e^{32\ell^2}
    \left(3\sqrt2\,\beta^{d-1/2}+a\beta^{-1}/\mathsf M\right).
 \end{split}
 \tag{V21.18}\label{c18v21:rates}
\]
The latter is at most a constant times
$(\log\beta)^{14}
(\beta^{32c+d-1/2}+\beta^{32c-1})$, which tends to zero.
The first also tends to zero. Substitute $\Lambda=a\gamma_\beta$
in the previous bounds.
\end{proof}

The inverse target energy $\beta^d/\kappa$ is also the physical-time
horizon allowed in V20. No clock was changed to produce this agreement.
For example $c=1/256,d=1/8$ yields the power $\beta^{-1/4}$,
times the stated logarithmic factor, in the leading relative cost.

This is a sufficient range, not a bound at an arbitrarily small
target energy. The cost proportional to
$\kappa\ell^{28}e^{32\ell^2}\beta^{-1/2}$ cannot be declared negligible
at a smaller physical scale without a new estimate. No continuum
trajectory for $\kappa,\beta$ or lattice spacing is verified here.

\subsection{The remaining coercive input}

At finite regulator the physical $K$ has compact resolvent, a unique
ground and translation symmetry. Each positive-energy eigenspace
can be decomposed under the finite abelian translation group.
It therefore contains a nonzero eigenvector of some momentum.
An excited $k=0$ vector is orthogonal to the ground; for $k\ne0$
orthogonality is automatic. A candidate small gap can consequently
be tested in a momentum sector without assuming a uniform gap.

The new result excludes concentration of such a normalized
low-energy candidate's squared mass or absolute gradient energy
at scale $\gamma_\beta$ in a specified $B_x$, within the stated
regime. For an eigenvalue much smaller than $\gamma_\beta$,
the error is not asserted to be small relative to that eigenvalue.

The cutoff remains local, not a global gauge chart. The localized
vector $u_xf$ need not depend only on retained coarse links.
V19 controls one Wilson-loop direction, not the variance of every
localized source. A rich physical direction bank and native
conditional variance along the nonlinear fibres are still unpaid.
Nor have we bounded derivatives of conditional expectations or
transferred these fine estimates to V19's coarse exceptional set
$\mathcal C_{\ell,\beta}^c$. The scalar gradient estimate does not
by itself pay the one-form Witten localization cost, which also
weights the cutoff gradient against $|df|^2$.

The next attack should seek actual good-region coercivity for this
newly controlled low-energy source class, retaining its complement
with the proved $o(\gamma_\beta)$ cost. It should not repeat a
small-probability argument, assume conditional Poincar\'e to prove
itself, or identify an artificial Dirichlet near-vacuum with the
original excited floor. Source-weighted return, physical scale
transport, interacting four-dimensional reconstruction and the
full positive physical Yang--Mills mass gap remain unproved.

\subsection{Verification and preservation}

The diagnostic checks exact translated-box multiplicities, the
positive-variation budget, a finite cyclic translation compression
with internal momentum multiplicity, the complex scalar localization
identity and the scale exponents. These finite algebraic regressions
do not simulate the unknown interacting ground, certify the full
native spectrum, or replace the analytic proofs.

V21 preserves the V20 body exactly apart from the title and replaces
only the active V20 source. All 47 other note files are unchanged.
Only \texttt{.tex} files remain in \texttt{latex\_notes}; build products
and diagnostics are outside it. The next companion review is V25,
the next broad literature sweep V30, and the next restart V100.
% END CYCLE 18 VERSION 21 APPEND
% BEGIN CYCLE 18 VERSION 22 APPEND
\clearpage
\section{V22: A native lower metric for every cube-link direction}
\label{c18v22:section}

V21 pays localization errors for actual low-energy states, but a
localized fine state need not be a coarse cylinder. V19, in turn,
normalizes only one Wilson-loop gradient. We now remove the latter
directional restriction locally: the actual flow--sampling metric
has a simultaneous lower bound on all twelve links of one coarse
cube, hence on its entire 36-real-dimensional cotangent space.
The bound also survives conditioning on the \emph{full} retained
coarse configuration, outside a paid coarse exceptional set.

This is an extension of V19's proved local comparison, not a new
assumption about the vacuum and not a many-cube coercivity theorem.
The flat lower estimate works without a transverse restriction;
V19 needed that restriction for its sharper upper estimate.
Keeping these two facts separate avoids promoting a physical
one-loop normalization into an unproved global spectral floor.
No Gaussian replacement of the native ground law is made.

\subsection{Setup, reused inputs and the unrestricted flat lower bound}

Retain V19's $N=\ell n$, $\ell\ge2$, $n\ge5$,
$s=\ell^2$, $m=(2/\pi)e^{-3\pi^2}$, and $R,L,h$ from
\eqref{c18v19:scales}, with $N\ge L+1$.
Retain the full compact spatial flow $\Phi_s$, straight-chain
sampling $\mathfrak b_\ell$, and $\mathfrak c=\mathfrak b_\ell\Phi_s$.
The original Hamiltonian and its ground law are still
\[
 H=\kappa\sum_e\mathcal L_e+bV_{\rm sp},\qquad
 K=H-E_0,\qquad \nu=\Omega^2dU,\qquad \mu=\mathfrak c_*\nu.
\]
The product group metric is the original unit-round $\mathrm{SU}(2)$
metric. All tangent and cotangent norms below use that metric.
Neither a clock change nor a field rescaling is implicit.

Let $\Gamma_\Box$ be the twelve positively oriented edges of the
coarse cube with vertices $\{0,1\}^3$. Let $\mathcal B_\Box$ be
their twelve fine straight chains. These chains share vertices,
but not stored fine links; thus $|\mathcal B_\Box|=12\ell$ and
$\mathcal B_\Box\subset[0,\ell]^3$.
Put $\mathfrak c_\Box=\pi_{\Gamma_\Box}\mathfrak c$.
Use the same box $\mathcal Q=[-R-2,\ell+R+2]^3$ and event
$\mathcal A$ as \eqref{c18v19:event}. In particular, every box
plaquette is controlled through the \emph{whole} time interval
$[0,s]$, not only at its endpoint.

The inherited inputs are V17's chain Fourier identity and native
conditional metric formula, V18's finite-range variation estimates,
V19's local gauge and generator estimates, and V21's source-weighted
localization. They remain finite-volume statements in the actual
Hamiltonian setting. The primary Wilson-flow reference cited in V19,
\href{https://arxiv.org/abs/1006.4518}{L\"uscher, arXiv:1006.4518},
was rechecked for context only; no continuum or spectral theorem
from that reference is used here. The V20 broad literature and
companion review remains the current checkpoint.

Write $\mathsf L_0=d_1^*d_1$ on \emph{all} fine three-color links
and $T=e^{-s\mathsf L_0}$. In particular, gradient and harmonic
directions have not been projected away.

\begin{lemma}[Unrestricted flat lower estimate]
\label{c18v22:flat}
For every coarse link covector $a$, with no divergence constraint,
\[
 \frac{m^2}{\ell}\|a\|^2
 \le \|TB_\ell^*a\|^2
 \le \ell\|a\|^2 .
 \tag{V22.1}\label{c18v22:flat-bound}
\]
The constants are independent of the ambient torus.
\end{lemma}
\begin{proof}
Use V17's unitary Fourier and alias conventions, now without either
transverse projection. The fine symbol is
\[
 \mathsf L_0(k)=\lambda(k)I-d(k)d(k)^*,\qquad
 \lambda(k)=|d(k)|^2.
\]
It has eigenvalue zero on $d(k)$ and eigenvalue $\lambda(k)$
on its orthogonal complement. At $k=0$ it is zero on all three
link directions. Consequently, for every complex vector $v$,
\[
 |e^{-s\mathsf L_0(k)}v|\ge e^{-s\lambda(k)}|v|.
\]
There are three identical color copies. In each coarse momentum
class the squared norm is exactly
\[
 \ell^{-3}\sum_j
   |e^{-s\mathsf L_0(k_j)}D_j^*a(q)|^2 .
 \tag{V22.2}\label{c18v22:flat-symbol}
\]
Every summand is nonnegative. For the primary alias $k_0=q/\ell$,
the chain modulus bound gives
$\sigma_{\min}(D_0)\ge2\ell/\pi$, while
$s\lambda(k_0)\le|q|^2\le3\pi^2$.
Its single summand is therefore at least
$m^2|a(q)|^2/\ell$. Summing over momenta and colors proves the
lower inequality, including zero momentum.

The straight fine chains of all distinct stored coarse edges are
link-disjoint, hence $B_\ell B_\ell^*=\ell I$.
Since $T$ is a contraction, the upper inequality follows.
Complex Fourier notation preserves the norm of real fields.
In contrast to V19's transverse proof, the scalar heat multiplier
has only been used as a \emph{lower} bound for the full matrix heat
multiplier. The transverse upper bound $4/\ell$ has not been
asserted for arbitrary $a$.
\end{proof}

\subsection{Exact compact chain adjoints and the cube buffer}

We describe the compact differential before making a comparison.
In invariant orthonormal coordinates use
$v_i=\delta W_i W_i^{-1}$ on a positive chain
$P=W_1\cdots W_\ell$. Direct product differentiation gives
\[
 \delta P\,P^{-1}
   =\sum_{i=1}^{\ell}\operatorname{Ad}_{P_{i-1}}v_i,
 \qquad P_{i-1}=W_1\cdots W_{i-1},\quad P_0=I.
 \tag{V22.3}\label{c18v22:chain}
\]
Thus the adjoint sends an output covector $a_e$ to
$\operatorname{Ad}_{P_{i-1}}^*a_e$ on the $i$th fine link.

\begin{lemma}[A cube covector comparison without a curvature denominator]
\label{c18v22:adjoint}
Fix $U\in\mathcal A$ and use V19's single fixed local gauge.
Let $W=\Phi_s(U)$ and let $a$ be any covector at
$\mathfrak c_\Box(U)$. In this gauge's invariant coordinates,
extend $a$ by zero off $\Gamma_\Box$, and set
\[
 z=(D\mathfrak b_{\ell,\Box}(W))^*a,\qquad
 z_0=B_\ell^*a.
\]
Then both vectors are supported on $\mathcal B_\Box$, and
\[
 \|z\|=\|z_0\|=\sqrt\ell\,|a|,\qquad
 \|z-z_0\|\le2\ell r\sqrt\ell\,|a|,\quad r=3Lh.
 \tag{V22.4}\label{c18v22:adjoint-bound}
\]
The V19 row comparisons also hold with
$\mathcal B=\mathcal B_\Box$ and its strict-radius buffer:
\[
 \begin{split}
 \|\mathbf1_{\mathcal B_\Box}(J(s)-T)\|
      &\le s\gamma e^{s\gamma}+4\epsilon,\\
 \|\mathbf1_{\mathcal B_\Box}J(s)\|
      &\le e^{s\gamma}+2\epsilon,\qquad
 \gamma=80r,\quad \epsilon=e^{60s-R}.
 \end{split}
 \tag{V22.5}\label{c18v22:rows}
\]
Here $J(s)=D\Phi_s(U)$ in the same invariant frames.
\end{lemma}
\begin{proof}
For the enlarged support, define
$\mathcal D_R=\{e:\operatorname{dist}(e,\mathcal B_\Box)<R\}$
in the link--plaquette adjacency graph.
Every starting link has tail coordinates in $[0,\ell]^3$.
One graph step changes any tail coordinate by at most one.
A link in $\mathcal D_R$ therefore has tail in
$[-R+1,\ell+R-1]^3$. The two extra incidence steps needed to
control the relevant links and their forces fit in $\mathcal Q$,
exactly as in Lemma~\ref{c18v19:gauge}.
That proof initially bounds \emph{every} box link by $2Lh$,
and then bounds its full-flow speed by $4h$ where needed.
It consequently gives the same $r=3Lh$ for this support.

Apply \eqref{c18v22:chain}. Orthogonality of each adjoint action
and link-disjointness of the twelve chains give the two exact
norm identities. The elementary product telescoping bound yields
$|P_{i-1}-I|\le(i-1)r\le\ell r$. For unit quaternions,
$\|\operatorname{Ad}_P-I\|\le2|P-I|$:
write $PvP^{-1}-v=(P-I)vP^{-1}+v(P^{-1}-I)$.
Each component error is at most $2\ell r|a_e|$.
Sum the squared errors over the twelve disjoint chains to prove
\eqref{c18v22:adjoint-bound}. No Wilson-loop gradient, common
curvature color, or division by a possibly vanishing gradient
was required.

V19's principal-generator bound uses at most four incident
plaquettes per stored link, not the cardinality of the output
support. It therefore still gives
$\|\mathsf A_{\mathcal D_R\mathcal D_R}
       +(\mathsf L_0)_{\mathcal D_R\mathcal D_R}\|\le80r$.
Its time-ordered row argument likewise uses the distance to the
complement, an operator norm, and the global bound $20$.
Words of length less than $R$ cannot leave the buffer and
contribute to these rows. The same tail
$2\sum_{j\ge R}(20s)^j/j!\le2e^{60s-R}$ and Duhamel estimate
give \eqref{c18v22:rows}, with no factor of twelve or of volume.
The background trajectory is still the full unfrozen one.

All these comparisons are at a fixed $U$ after a fixed gauge
isometry. Gauge equivariance acts orthogonally on fine and coarse
tangent spaces. The rule selecting the initial gauge is never
differentiated.
\end{proof}

\subsection{The simultaneous native matrix bound and a local lift}

Define the principal cube metric
\[
 G_\Box(U)=D\mathfrak c_\Box(U)D\mathfrak c_\Box(U)^*.
 \tag{V22.6}\label{c18v22:gram}
\]
It is a $36$-dimensional positive endomorphism at
$\mathfrak c_\Box(U)$. The adjoint in this definition includes
\emph{all} initial fine tangent directions, not just the buffer.

\begin{theorem}[Native cube-link metric on the paid event]
\label{c18v22:native}
For every $U\in\mathcal A$, simultaneously for every cube covector,
\[
 \boxed{\quad
 \frac{m^2}{4\ell}I_{36}\ \preceq\ G_\Box(U)\
                   \preceq\ 4\ell I_{36}.
 \quad}
 \tag{V22.7}\label{c18v22:native-bound}
\]
This includes all differentials of gauge-invariant cube
observables, but does not restrict to them.
\end{theorem}
\begin{proof}
The exact adjoint chain rule is
$(D\mathfrak c_\Box)^*a=J(s)^*z$.
Adjoin \eqref{c18v22:rows} against the supported vectors and use
\eqref{c18v22:adjoint-bound} to obtain
\[
 \|J(s)^*z-Tz_0\|
 \le\sqrt\ell\,|a|\big[
       s\gamma e^{s\gamma}+4\epsilon
       +(e^{s\gamma}+2\epsilon)2\ell r\big].
 \tag{V22.8}\label{c18v22:error}
\]
The same explicit parameters as V19 give
$s\gamma=m/(15\ell)<1$ and $\epsilon\le m/(16\ell)<1$.
Using $e^{s\gamma}<3$ and $\ell\le s$, the square bracket is
at most $250sr+4\epsilon$, and hence at most
\[
 300sr+4\epsilon
 \le \frac{m}{4\ell}+\frac{m}{4\ell}
 =\frac{m}{2\ell}.
\]
Lemma~\ref{c18v22:flat}, with $a$ extended by zero, then gives
\[
 \frac{m}{2\sqrt\ell}|a|
 \le \|(D\mathfrak c_\Box)^*a\|
 \le\left(\sqrt\ell+\frac{m}{2\sqrt\ell}\right)|a|
 \le2\sqrt\ell\,|a|.
\]
Squaring proves \eqref{c18v22:native-bound}. The constants are
uniform in $a$, including $a=0$, so this is a matrix inequality,
not a finite collection of scalar tests. Undoing the fixed gauge
preserves it.
\end{proof}

The lower order $\ell^{-1}$ is now available for the entire
local bank. The upper order $\ell$ is a deliberately coarse
envelope; this theorem does not extend V19's physical upper
normalization to all cube directions.

\begin{corollary}[Canonical right inverse for the cube output only]
\label{c18v22:lift}
On $\mathcal A$ the minimal-norm fine lift of a cube tangent
vector $v$ is
\[
 R_\Box v=(D\mathfrak c_\Box)^*G_\Box^{-1}v,\qquad
 D\mathfrak c_\Box R_\Box=I,\qquad
 \|R_\Box v\|^2\le\frac{4\ell}{m^2}|v|^2 .
 \tag{V22.9}\label{c18v22:lift-bound}
\]
This formula is gauge equivariant and needs no gauge selection.
\end{corollary}
\begin{proof}
The matrix bound makes $G_\Box$ invertible. Direct multiplication
gives the right-inverse identity and
$\|R_\Box v\|^2=\langle v,G_\Box^{-1}v\rangle$.
Every other lift differs by an element of
$\ker D\mathfrak c_\Box$, which is orthogonal to the displayed
adjoint range; Pythagoras proves minimality. The norm bound
follows from \eqref{c18v22:native-bound}. The formula uses only
the original metric and equivariant differential, so it is
equivariant.
\end{proof}

Crucially, this lift need not leave the \emph{other} retained coarse
links fixed and need not have fine support inside the buffer.
It is not a right inverse for a prescribed variation of the
entire coarse configuration. Its norm alone also gives no bound
on derivatives of this lift or on conditional scores.

\subsection{Conditioning on the true full coarse configuration}

Set
\[
 \begin{split}
 \overline G_\Box(g)
   &=\mathbb E_\nu[G_\Box(U)\mid\mathfrak c(U)=g],\\
 p_\Box(g)&=\mathbb P_\nu(\mathcal A^c\mid\mathfrak c(U)=g),\\
 \mathcal C_\Box&=\{g:p_\Box(g)\le1/2\}.
 \end{split}
 \tag{V22.10}\label{c18v22:conditional-def}
\]
Invariant frames identify the relevant coarse tangent spaces;
all configurations in a fixed fibre have the same output $g$.
Thus this is an ordinary conditional matrix integral. It is the
principal block of V17's \emph{full} conditional metric, not a
metric computed by replacing $\mu$ with a twelve-link product law.

\begin{theorem}[A simultaneous conditional lower metric]
\label{c18v22:conditional}
With $\wp_{\ell,\beta}$ from \eqref{c18v19:probability},
\[
 \begin{split}
 \mu(\mathcal C_\Box^c)&\le\min\{1,2\wp_{\ell,\beta}\},\\
 \overline G_\Box(g)&\succeq\frac{m^2}{4\ell}
                         (1-p_\Box(g))I_{36},\\
 \overline G_\Box(g)&\succeq\frac{m^2}{8\ell}I_{36}
                  \quad\text{for }\mu\text{-a.e. }g\in\mathcal C_\Box .
 \end{split}
 \tag{V22.11}\label{c18v22:conditional-bound}
\]
In particular, one exceptional set works for every cube direction.
There is no union bound over a direction net.
\end{theorem}
\begin{proof}
The event is unchanged from V19, so
$\nu(\mathcal A^c)\le\wp_{\ell,\beta}$ is already proved in the
actual Hamiltonian vacuum. Integrating $p_\Box$ and applying
Markov gives the first assertion. Pointwise on all fine
configurations,
\[
 G_\Box(U)\succeq\frac{m^2}{4\ell}
                            \mathbf1_{\mathcal A}(U)I_{36},
\]
because $G_\Box$ is positive semidefinite off the event as well.
Conditional integration preserves positive semidefinite order:
test on a countable dense set of vectors, then use continuity of
the finite matrices. This gives the second assertion on one
common full-measure set, and hence the third.
No independence or conditional exceptional-probability assumption
has been added.
\end{proof}

For any smooth complex cube cylinder
$f(g)=F(g_{\Gamma_\Box})$, the original ground-state form gives
the exact identity and lower estimate
\[
 \begin{split}
 \langle\Omega(f\circ\mathfrak c),K\Omega(f\circ\mathfrak c)\rangle
   &=\kappa\int
       \langle\nabla_\Box F,\overline G_\Box\nabla_\Box F\rangle\,d\mu\\
   &\ge\frac{\kappa m^2}{8\ell}
                  \int_{\mathcal C_\Box}|\nabla_\Box F|^2\,d\mu .
 \end{split}
 \tag{V22.12}\label{c18v22:cylinder-form}
\]
For complex $F$, the matrix acts separately on real and imaginary
parts. The formula follows from the chain rule and conditional
integration; constants have zero form. If $F$ is gauge invariant
under the cube vertices, its extension is a physical cylinder
and the same formula holds. It is still a gradient estimate,
not a variance inequality under the native marginal.

The exceptional probability is at most
$2C_*\ell^{18}\varepsilon_\beta$ before truncation at one.
It tends to zero when $\beta\to\infty$ and
$\ell^2\le c\log\beta$ for any fixed $c>0$, subject to the embedded
buffer. When used together with V21's target-energy localization,
retain its stronger condition $d+32c<1/2$.
This statement neither selects a physical continuum trajectory
nor makes the explicit small constant $m$ numerically large.

\subsection{What has advanced, and the remaining assembly problem}

The advance over V17's global exponentially small metric envelope
and V19's single-loop estimate is a native, algebraic-in-$\ell$
lower bound for the \emph{whole local cube bank}, with its actual
vacuum probability, conditional version, and canonical local
right inverse. No new Poincar\'e hypothesis appears in its proof.

It is important not to infer a many-cube lower bound merely from
the principal blocks. Even a positive definite matrix can have
uniformly good diagonal blocks and a collapsing full eigenvalue:
for $0<\eta<1$,
\[
 M_\eta=
 \begin{pmatrix}1+\eta&-1\\-1&1+\eta\end{pmatrix},\qquad
 (M_\eta)_{11},(M_\eta)_{22}\ge1,\qquad
 M_\eta\binom11=\eta\binom11 .
 \tag{V22.13}\label{c18v22:assembly-counterexample}
\]
The other eigenvalue is $2+\eta$, so this example is genuinely
positive definite. It disproves the bare matrix implication,
not a property of the unknown Yang--Mills vacuum. In our problem,
fine images of covectors on different cubes may have cross terms.
A proof controlling those terms or a compatible joint lift would
be additional mathematics, not a consequence of
\eqref{c18v22:native-bound}.

There are also two distinct measure obligations.
First, even \eqref{c18v22:cylinder-form} does not bound a cube
observable's variance without information about the actual
coarse marginal and the exceptional set weighted by its gradient.
Second, an arbitrary V21 localized low-energy fine source is
not a cube cylinder, nor necessarily any retained cylinder.
Its conditional fluctuation along the nonlinear
$\mathfrak c$-fibres requires separate coercivity.
The local right inverse in \eqref{c18v22:lift-bound} does not
fix the full coarse output, and so cannot be silently used as
the horizontal derivative for conditioning on all of $g$.

The next useful upstream targets are therefore explicit
cross-cube metric control or a joint constrained lift, and
native conditional variance for the actual low-energy source
class. The existing scalar localization cost must be retained;
one-form weighted localization, source-weighted return, physical
scale transport and interacting four-dimensional reconstruction
remain unpaid as well. The full positive physical Yang--Mills
mass gap is not proved.

\clearpage
\subsection{Verification and preservation}

The diagnostic tests the unrestricted matrix heat symbol, exact
rational quaternion chain adjoints on all twelve cube paths,
their disjoint-link norm identity, buffer geometry, the constants
in the native comparison, and the positive-matrix assembly
counterexample. These are regression checks, not a computation
of the unknown interacting vacuum or a proof assistant certificate.

The checks include 1296 exact rational differential columns across
eight cube configurations, with chain lengths $2,3,5,8$, both
near identity and unrestricted; 36 Fourier alias matrices including
harmonic and longitudinal directions; nine finite buffer geometries;
199 parameter choices; and 100 exact positive-matrix counterexamples.
The rectangular minimal-norm lift identity is also checked exactly.
At three corner momenta the full Gram eigenvalues were additionally
evaluated with 70 decimal digits. Double-precision Gram sums can
lose eigenvalues below their rounding scale, so their rounded signs
are not used as lower-bound certificates. The proof of
\eqref{c18v22:flat-bound} is the analytic primary-alias inequality,
not a numerical extrapolation from these tests.

V22 preserves the entire V21 body apart from the title, replacing
only the active source. The other 47 note files are unchanged.
Only \texttt{.tex} files are kept in \texttt{latex\_notes};
compilation and inspection products are elsewhere.
The next companion review is V25, the next broad primary-literature
sweep is V30, and the next scheduled restart is V100.
% END CYCLE 18 VERSION 22 APPEND
% BEGIN CYCLE 18 VERSION 23 APPEND
\clearpage
\section{V23: Cross-cube assembly with a controlled diagonal defect}
\label{c18v23:section}

V22 proves a lower metric for every direction on one cube, but
correctly leaves cross-cube cancellation unresolved. Here we
control that cancellation for the \emph{full} retained metric
by a spatial partition and an explicit kernel second moment.
The result is a global matrix lower inequality with a diagonal
bad-patch weight. It does not require every patch to be good.
Conditioning on the full retained configuration preserves the
inequality and pays the mean defect in the actual vacuum.

The distinction from a global spectral floor is essential.
A small mean diagonal defect is not an operator-norm bound on it,
and a large subspace of well-controlled metric directions does
not exclude a small exceptional direction. The new estimate
applies to arbitrary full coarse covectors; the still-missing
step is to control its defect with the relevant source weights
and to prove variance inequalities for the native law.

\subsection{Conventions and the native kernel moment}

Use the same finite spatial torus, $\mathrm{SU}(2)$ metric,
Hamiltonian $H$, ground law $\nu=\Omega^2dU$, and
$K=H-E_0$ as V22. Put $N=\ell n$, $\ell\ge2$, $n\ge5$, $s=\ell^2$,
$\mathfrak c=\mathfrak b_\ell\Phi_s$, and $\mu=\mathfrak c_*\nu$.
Neither the retained map nor physical time is changed.
Write
\[
 Q(U)=D\mathfrak b_\ell(\Phi_s(U)),\qquad
 J(s)=D\Phi_s(U),\qquad
 G(U)=D\mathfrak c(U)D\mathfrak c(U)^*
                   =QJ(s)J(s)^*Q^* .
 \tag{V23.1}\label{c18v23:G}
\]
All matrices use invariant orthonormal link frames. Their entries
below are $3$-by-$3$ real color blocks.

For coarse links $e=(y,i)$ and $f=(z,j)$, let $D_c(e,f)$ be the
periodic $\ell^\infty$ distance between their coarse tails $y,z$;
different orientations at the same tail have distance zero.
For fine links use V18's link--plaquette graph distance $d_f$.
Each step of that graph changes every tail coordinate by at
most one, in the periodic metric.

\begin{lemma}[A uniform native second-moment Schur bound]
\label{c18v23:moment}
For every initial configuration, with no good-event assumption,
\[
 \sup_e\sum_f D_c(e,f)^2\|G_{ef}(U)\|
 \le S_\ell:=
 \frac{e^{40s}}{\ell}\big[(40s+2\ell)^2+40s\big].
 \tag{V23.2}\label{c18v23:moment-bound}
\]
The same bound holds for column sums and is independent of $n$.
\end{lemma}
\begin{proof}
In V19's invariant frames, $J'=\mathsf A(t)J$ with
$\mathsf A=-\mathsf H+\mathsf S$. We need a block Schur bound,
not merely the operator norm bound stated there. A mixed
Hessian block of one plaquette has norm at most one: its
bilinear form is the real part of a product of unit quaternions
with two tangent insertions. Its diagonal block is $w_pI_3$,
also of norm at most one. Four incident plaquettes and twelve
off-diagonal incidences thus bound each Hessian block row sum
by $16$. Symmetry gives the same column bound.
The invariant-frame correction is link diagonal and has norm
at most the link speed, at most $4$. Therefore
\[
 \sup_a\sum_b\|\mathsf A_{ab}(t)\|\le20,\qquad
 \sup_b\sum_a\|\mathsf A_{ab}(t)\|\le20 .
 \tag{V23.3}\label{c18v23:generator-schur}
\]
All off-diagonal blocks remain on the fine link graph.

Expand $J(s)J(s)^*$ in its two absolutely convergent
time-ordered series. A product with total order $r$ uses $r$
graph steps, possibly including diagonal stays. Its nonzero
blocks have $d_f(a,b)\le r$. The sum of its block row norms,
after time integration and summing the two orders, is at most
\[
 \sum_{p+q=r}\frac{(20s)^p}{p!}\frac{(20s)^q}{q!}
                         =\frac{(40s)^r}{r!}.
\]
This estimate also follows by replacing every block by its
norm and multiplying the resulting nonnegative matrices.
It permits noncommuting time dependence and uses the column
bound in \eqref{c18v23:generator-schur} for the adjoint series.
Consequently, for $M=J(s)J(s)^*$ and each fine link $a$,
\[
 \sum_b(d_f(a,b)+2\ell)^2\|M_{ab}\|
 \le\sum_{r\ge0}(r+2\ell)^2\frac{(40s)^r}{r!}
 =e^{40s}\big[(40s+2\ell)^2+40s\big].
 \tag{V23.4}\label{c18v23:fine-moment}
\]
The last equality is the elementary first and second factorial
moment identity for an exponential series; no random path law
has been substituted for the native flow.

Each nonzero block of $Q$ is an orthogonal adjoint transport.
There are $\ell$ such blocks in each coarse row, and every
retained fine chain link occurs in only one coarse chain.
If $a$ is in the chain of $e$ and $b$ in that of $f$, then
\[
 \ell D_c(e,f)\le d_f(a,b)+2\ell.
\]
Indeed the two within-chain tail displacements are each at
most $\ell$, and a fine graph path bounds the periodic
tail distance. Apply the triangle inequality to the blocks
of $QMQ^*$, sum first over the disjoint chains of $f$, and
then over the $\ell$ links of $e$. Equation
\eqref{c18v23:fine-moment}, divided by $\ell^2$, proves
\eqref{c18v23:moment-bound}. Since $G_{fe}=G_{ef}^*$ and
$D_c$ is symmetric, the column statement follows.
\end{proof}

\subsection{A spatial partition with a quantified matrix localization error}

Let $\varrho\ge1$ be an integer with $n\ge2\varrho+1$.
On the cyclic coordinate lattice define
$\phi_\varrho(j)=(1-|j|/\varrho)_+$, using shortest periodic
distance. Put
\[
 z_\varrho=\sum_j\phi_\varrho(j)^2
             =\frac{2\varrho^2+1}{3\varrho},\qquad
 \psi_\varrho(y)=\prod_{i=1}^3\phi_\varrho(y_i),\qquad
 \chi_x(e)=\frac{\psi_\varrho(y-x)}{z_\varrho^{3/2}},
 \quad e=(y,i).
 \tag{V23.5}\label{c18v23:partition}
\]
The same scalar acts on every orientation and color at $y$.
The centers $x$ run over all $n^3$ coarse sites.

\begin{lemma}[Partition identities and full-metric localization]
\label{c18v23:IMS}
For all coarse links,
\[
 \sum_x\chi_x(e)^2=1,\qquad
 \sum_x|\chi_x(e)-\chi_x(f)|^2
       \le\frac{27}{\varrho^2}D_c(e,f)^2 .
 \tag{V23.6}\label{c18v23:partition-bound}
\]
For every $U$ and every full coarse covector $a$,
\[
 \left|\sum_x\langle\chi_xa,G(U)\chi_xa\rangle
                    -\langle a,G(U)a\rangle\right|
 \le E_{\ell,\varrho}\|a\|^2,\qquad
 E_{\ell,\varrho}=\frac{27S_\ell}{2\varrho^2}.
 \tag{V23.7}\label{c18v23:IMS-bound}
\]
\end{lemma}
\begin{proof}
The first equality follows by translating the product sum.
The one-coordinate exact discrete derivative satisfies
$\sum_j|\phi_\varrho(j+1)-\phi_\varrho(j)|^2=2/\varrho$.
After normalizing the three-dimensional product, a unit
coordinate translation therefore has squared difference
$2/(\varrho z_\varrho)\le3/\varrho^2$ in the center-index
$\ell^2$ norm. Telescope along a shortest coarse coordinate
path. Its length is at most $3D_c(e,f)$. The triangle
inequality in that $\ell^2$ space proves the second bound.

The $(e,f)$ block of the localized-metric difference is exactly
\[
 \left(\sum_x\chi_x(e)\chi_x(f)-1\right)G_{ef}
       =-\frac12\sum_x|\chi_x(e)-\chi_x(f)|^2G_{ef}.
\]
The block Schur test, \eqref{c18v23:partition-bound}, and
Lemma~\ref{c18v23:moment} bound its operator norm by
$27S_\ell/(2\varrho^2)$. This proves the quadratic-form
estimate, also for complex covectors after complexification.
No sign has been assigned to the individual cross terms.
\end{proof}

This is a spatial matrix localization identity. The $\chi_x$
depend on lattice position, not on the gauge configuration;
it is different from V21's configuration-space cutoff.
The general localization principle is established machinery:
compare the scalar IMS formulas in J.~P.~Solovej,
\emph{The ionization conjecture in Hartree--Fock theory},
\href{https://annals.math.princeton.edu/wp-content/uploads/annals-v158-n2-p04.pdf}
{Lemma 2.4, Annals of Mathematics 158 (2003)}.
The block identity and its native kernel cost were proved above.
No atomic spectral or variance estimate is imported.

\subsection{Native lower bounds on the larger spatial supports}

Keep V19's $m=(2/\pi)e^{-3\pi^2}$ and choose
\[
 \begin{split}
 R&=\left\lceil64s+\log(16\ell/m)\right\rceil,\\
 L_\varrho&=2\ell\varrho+2R+4,\qquad
 h_\varrho=\frac{m}{3600\ell^3L_\varrho},\qquad
 N\ge L_\varrho+1 .
 \end{split}
 \tag{V23.8}\label{c18v23:patch-scale}
\]
This implies $n\ge2\varrho+1$. For each center $x$, the box
\[
 \mathcal Q_x=\ell x+
       [-\ell\varrho-R-2,\ell\varrho+R+2]^3
\]
is understood as an embedded box in the torus. Let $\mathcal A_x$
be the event that every plaquette wholly in this box satisfies
$|\Phi_t(U)_p-I|\le h_\varrho$ for every $0\le t\le s$.
Let $\Gamma_x$ contain all coarse links whose tails satisfy
$|y_i-x_i|<\varrho$ in every coordinate; it contains the
support of $\chi_x$.

\begin{lemma}[The V22 comparison is independent of output cardinality]
\label{c18v23:patch-lower}
For every $U\in\mathcal A_x$ and every covector supported on
$\Gamma_x$,
\[
 \langle a,G(U)a\rangle\ge\alpha_\ell\|a\|^2,\qquad
                         \alpha_\ell=\frac{m^2}{4\ell}.
 \tag{V23.9}\label{c18v23:local-lower}
\]
Moreover, uniformly in the center and ambient volume under
the embedding assumption,
\[
 \nu(\mathcal A_x^c)\le p_{\ell,\varrho,\beta}:=
 \min\left\{1,\frac{6(L_\varrho+1)^3(1+24s)
                            \varepsilon_\beta}{h_\varrho^2}\right\}.
 \tag{V23.10}\label{c18v23:patch-probability}
\]
\end{lemma}
\begin{proof}
All fine chains of $\Gamma_x$ lie in
$\ell x+[-\ell\varrho,\ell\varrho]^3$ and remain link-disjoint.
Their strict-radius-$R$ graph neighborhood, and the two extra
incidence steps for the needed forces, lie inside $\mathcal Q_x$.
V19's rooted axial tree argument applies to this box of side
$L_\varrho$: it makes every initial box link lie within
$2L_\varrho h_\varrho$ of identity. Since $4s\le L_\varrho$,
the same fixed gauge gives
\[
 |W_e(t)-I|\le r:=3L_\varrho h_\varrho
                    =\frac{m}{1200\ell^3}
\]
on all links needed by the comparison throughout the full flow.
No common gauge between different centers is required.

For $z=Q^*a$ and $z_0=B_\ell^*a$ in that gauge, V22's exact
chain differentiation gives
$\|z\|=\|z_0\|=\sqrt\ell\,\|a\|$ and
$\|z-z_0\|\le2\ell r\sqrt\ell\,\|a\|$.
Its local generator and time-ordered row estimates use block
norms, incidence and distance, not the number of output links.
They thus give the same comparison error
$m\|a\|/(2\sqrt\ell)$ with
$e^{-s\mathsf L_0}z_0$.
V22's unrestricted flat lower bound is valid for any support,
so the reverse triangle inequality proves
\eqref{c18v23:local-lower}. Undo the fixed gauge; no derivative
of the gauge-selection rule has been used.

For the probability, the V18 native mean-supremum estimate is
$\mathbb E_\nu\sup_{t\le s}\delta_p(\Phi_t U)
 \le(1+24s)\varepsilon_\beta$.
There are at most $3(L_\varrho+1)^3$ box plaquettes, and
$|U_p-I|^2=2\delta_p(U)$. Markov and the union bound prove
\eqref{c18v23:patch-probability}. The threshold and larger
box volume are both paid explicitly; V19's old $\ell^{18}$
budget has not been reused with an enlarged unpriced box.
\end{proof}

\subsection{Global assembly with a retained diagonal defect}

For a fixed configuration define the link-scalar weights
\[
 d_e(U)=\sum_x\chi_x(e)^2\mathbf1_{\mathcal A_x^c}(U),
 \qquad
 \mathsf D(U)=\operatorname{diag}_e(d_e(U)I_3).
 \tag{V23.11}\label{c18v23:defect}
\]
Thus $0\preceq\mathsf D\preceq I$. The boxes may overlap
arbitrarily; independence is neither stated nor used.

\begin{theorem}[Full native coarse-metric lower inequality]
\label{c18v23:global}
Under \eqref{c18v23:patch-scale}, for every fine configuration,
\[
 G(U)\succeq
   \alpha_\ell(I-\mathsf D(U))-E_{\ell,\varrho}I .
 \tag{V23.12}\label{c18v23:global-general}
\]
In particular, choose the analysis radius
\[
 \varrho=\varrho_\ell:=
 \left\lceil\frac{12}{m}e^{20s}(40s+2\ell+1)\right\rceil.
 \tag{V23.13}\label{c18v23:radius}
\]
Then $E_{\ell,\varrho}\le\alpha_\ell/2$ and
\[
 \boxed{\quad
 G(U)\succeq\frac{m^2}{8\ell}I
                    -\frac{m^2}{4\ell}\mathsf D(U).
 \quad}
 \tag{V23.14}\label{c18v23:global-defect}
\]
The right side is allowed to be negative on bad directions;
it is not an unconditional positive floor.
\end{theorem}
\begin{proof}
For each center with $U\in\mathcal A_x$, apply
\eqref{c18v23:local-lower} to $\chi_xa$.
For the remaining centers use positivity of $G$.
Sum, then apply \eqref{c18v23:IMS-bound}:
\[
 \begin{split}
 \langle a,Ga\rangle
 &\ge\sum_x\langle\chi_xa,G\chi_xa\rangle
                               -E_{\ell,\varrho}\|a\|^2\\
 &\ge\alpha_\ell\sum_x\mathbf1_{\mathcal A_x}
                                \|\chi_xa\|^2
                               -E_{\ell,\varrho}\|a\|^2\\
 &=\alpha_\ell\langle a,(I-\mathsf D)a\rangle
                               -E_{\ell,\varrho}\|a\|^2.
 \end{split}
\]
This holds for every full covector, proving the matrix inequality.
For the chosen radius, write
$A_s=(40s+2\ell)^2+40s\le(40s+2\ell+1)^2$.
Equations \eqref{c18v23:moment-bound} and
\eqref{c18v23:radius} give
\[
 E_{\ell,\varrho}
 \le\frac{27m^2}{288\ell}
 =\frac{3m^2}{32\ell}
 \le\frac{m^2}{8\ell}=\frac{\alpha_\ell}{2}.
\]
The radius affects only the analysis and its good events, not
the native map or its coefficients.
\end{proof}

\subsection{The actual conditional metric and a large joint direction space}

Use the \emph{full} coarse configuration for every conditional
expectation:
\[
 \overline G(g)=\mathbb E_\nu[G(U)\mid\mathfrak c(U)=g],
 \qquad
 \overline{\mathsf D}(g)=
       \mathbb E_\nu[\mathsf D(U)\mid\mathfrak c(U)=g]
       =\operatorname{diag}_e(\bar d_e(g)I_3).
 \tag{V23.15}\label{c18v23:conditional}
\]
These are the original conditional metric and its explicitly
defined defect, not independent-link replacements.

\begin{corollary}[Conditional assembly and mean defect]
\label{c18v23:conditional-result}
For $\varrho=\varrho_\ell$, almost everywhere under $\mu$,
\[
 \overline G(g)\succeq
      \frac{\alpha_\ell}{2}I
                       -\alpha_\ell\overline{\mathsf D}(g),
 \qquad
 0\le\bar d_e(g)\le1,\qquad
 \int\bar d_e\,d\mu\le p_{\ell,\varrho,\beta}.
 \tag{V23.16}\label{c18v23:conditional-bound}
\]
Let
\[
 \mathcal I(g)=\{e:\bar d_e(g)\le1/4\},\qquad
 \mathcal V(g)=\{a:a_e=0\text{ for }e\notin\mathcal I(g)\}.
\]
Then, simultaneously for all $a\in\mathcal V(g)$,
\[
 \langle a,\overline G(g)a\rangle
                   \ge\frac{m^2}{16\ell}\|a\|^2 .
 \tag{V23.17}\label{c18v23:joint-bank}
\]
Writing $E_c=3n^3$ for the number of coarse links, one has
\[
 \frac1{E_c}\int|\mathcal I(g)^c|\,d\mu
       \le\min\{1,4p_{\ell,\varrho,\beta}\}.
 \tag{V23.18}\label{c18v23:bad-fraction}
\]
\end{corollary}
\begin{proof}
Conditional expectation preserves positive semidefinite order;
a countable dense set of test vectors gives one common
full-measure set. Integrating each diagonal weight and using
\eqref{c18v23:patch-probability} yields
\[
 \int\bar d_e\,d\mu
 =\sum_x\chi_x(e)^2\nu(\mathcal A_x^c)
 \le p_{\ell,\varrho,\beta}\sum_x\chi_x(e)^2
 =p_{\ell,\varrho,\beta}.
\]
For a covector supported on $\mathcal I(g)$, its defect form
is at most $\|a\|^2/4$. Substitute in
\eqref{c18v23:conditional-bound} to get
$\alpha_\ell\|a\|^2/4$, proving \eqref{c18v23:joint-bank}.
Finally,
$\mathbf1_{\{\bar d_e>1/4\}}\le4\bar d_e$;
sum and integrate, also using the trivial fraction bound by one.
\end{proof}

The joint estimate retains cross terms between arbitrarily many
nearby good links; no separation or coloring of them is needed.
It controls a principal compression, however, not its Schur
complement after all other coarse outputs are fixed.
At a fixed $g$, the min--max principle also bounds the number
of eigenvalues of $\overline G(g)$ strictly below $m^2/(16\ell)$
by $3|\mathcal I(g)^c|$, since $\mathcal V(g)$ has that
codimension. This is a metric eigenvalue count, not a Hamiltonian
excitation count.

For any smooth full coarse observable $f$, the original
ground-state form identity and \eqref{c18v23:conditional-bound}
give
\[
 \begin{split}
 \langle\Omega(f\circ\mathfrak c),K\Omega(f\circ\mathfrak c)\rangle
 &\ge\frac{\kappa m^2}{8\ell}
                     \int|\nabla_g f|^2\,d\mu\\
 &\quad-\frac{\kappa m^2}{4\ell}
                  \int\sum_e\bar d_e(g)|\nabla_e f(g)|^2\,d\mu .
 \end{split}
 \tag{V23.19}\label{c18v23:form}
\]
This now has no cube-support restriction. For complex observables
apply the real form separately to real and imaginary parts.
For example, the displayed defect integral is at most
$p_{\ell,\varrho,\beta}\sum_e\|\nabla_e f\|_\infty^2$.
No replacement by $p_{\ell,\varrho,\beta}\int|\nabla f|^2d\mu$
is justified without an additional source-weighted estimate.

\subsection{Paying the radius and locating the remaining obstruction}

The radius in \eqref{c18v23:radius} is intentionally expensive.
It is the price of the unconditional native kernel estimate.
For $\ell\ge2$, the explicit elementary bounds
\[
 \varrho_\ell\le\frac{505}{m}\ell^2e^{20\ell^2},
 \quad R\le96\ell^2,\quad
 L_{\varrho_\ell}+1\le\frac{1200}{m}\ell^3e^{20\ell^2}
 \tag{V23.20}\label{c18v23:radius-growth}
\]
imply
\[
 p_{\ell,\varrho_\ell,\beta}
 \le\min\{1,C_{23}\ell^{23}e^{100\ell^2}\varepsilon_\beta\},
 \qquad
 C_{23}=\frac{150\cdot3600^2\cdot1200^5}{m^7}.
 \tag{V23.21}\label{c18v23:prob-growth}
\]
Indeed, $40\ell^2+2\ell+1\le42\ell^2$ gives the first bound
after accounting for the ceiling. The bound on $R$ is V19's
elementary estimate, and
$2R+5\le100\ell^3$ gives the third with room to spare.
Inserting $h_\varrho^{-2}=3600^2\ell^6L_\varrho^2/m^2$
and $1+24s\le25\ell^2$ in \eqref{c18v23:patch-probability}
proves \eqref{c18v23:prob-growth}.

Thus, for any fixed $0<c<1/200$, if $\beta\to\infty$ and
$\ell^2\le c\log\beta$, the mean defect tends to zero:
the bound is a constant times
$(\log\beta)^{23/2}\beta^{100c-1/2}$.
Ambient-volume uniformity still requires
$N\ge L_{\varrho_\ell}+1$. This is a significantly larger
embedded analysis box than in V19--V22. We have not selected
a physical continuum trajectory or proved that these constants
are useful at moderate coupling.

What has advanced is the cross-cube \emph{matrix} step:
the full native metric and its actual conditional average
have a quantified global lower inequality, with an explicit
diagonal defect rather than omitted off-diagonal terms.
A large joint direction space is controlled, with a vanishing
expected excluded fraction in the stated regime.

What remains cannot be bypassed by the fraction statement.
The matrices $\operatorname{diag}(\eta,1,\ldots,1)$ have one
bad direction, arbitrarily small excluded fraction as dimension
grows, and smallest eigenvalue $\eta$. Nor does a good principal
compression provide a lift holding all complementary retained
outputs fixed. For \eqref{c18v23:form}, the unresolved native
quantity is the defect \emph{weighted by the actual coarse
gradient}. V21 bounds fine-source gradient leakage for its
declared boxes and weights; it does not identify that quantity
with the new, much larger-patch conditional defect. Such a
transfer needs a proof with the new scales and the full map.

Even removal of the metric defect would not supply a Poincar\'e
inequality for the true coarse marginal or conditional variance
on the nonlinear full $\mathfrak c$-fibres. Those are the next
upstream obligations for actual low-energy sources.
One-form localization, source-weighted return, accumulated
physical-scale errors, interacting four-dimensional reconstruction
and the full positive physical Yang--Mills mass gap remain unproved.

\subsection{Verification and preservation}

The diagnostic checks discrete hat normalization and translation
costs, the exact block localization identity, noncommuting
finite-range word moments, the compact plaquette block bound,
conditional diagonal-weight algebra and the radius/probability
constants. Algebra fixtures are not simulations of the unknown
native ground law or substitutes for the proofs above.

V23 preserves V22's body exactly except its title and replaces
only that active source; all other 47 note files remain unchanged.
Only \texttt{.tex} files are kept in \texttt{latex\_notes}.
Build products and diagnostics stay elsewhere.
The next companion checkpoint is V25, the next broad
primary-literature sweep V30, and the scheduled restart V100.
% END CYCLE 18 VERSION 23 APPEND
% BEGIN CYCLE 18 VERSION 24 APPEND
\clearpage
\section{V24: Joint coarse lifts with the complement fixed}
\label{c18v24:section}

V23 pays cross-cube terms in the full coarse metric, but its good-link
principal compression is not a Schur complement. A lift of those
links could still move the remaining retained coordinates. We now
resolve that geometric issue on a smaller, protected set of links:
a matrix harmonic-extension estimate controls the \emph{full inverse}
there, and hence gives an actual fine lift with every complementary
coarse output fixed. The additional spatial neighbourhood is priced
in the native vacuum, rather than silently discarded.

This does not yet pay V23's coarse-gradient-weighted defect.
Nor does it turn V21's error, small at a target energy scale, into
an error relative to every arbitrarily smaller eigenvalue.
The useful advance is a genuine constrained lift, with explicit
norm, spatial and probability costs. Its derivatives, conditional
scores and the required native variance inequalities remain separate
obligations.

\subsection{A harmonic-extension estimate for positive block matrices}

The following elementary lemma isolates the extra inverse step.
It uses the same matrix localization identity as V23. No sign
condition on off-diagonal entries, band truncation, or Markov
generator interpretation is needed.

Let $\mathcal E$ be a finite set of blocks with a symmetric distance
$d$ satisfying the triangle inequality; distinct blocks may have
distance zero. Let $P_T$ denote the block projection for a subset $T$.
All norms use the given Hilbert direct sum. Suppose
\[
 g_-I\preceq M\preceq g_+I,\qquad g_->0,\qquad
 \sup_e\sum_f d(e,f)^2\|M_{ef}\|\le S .
 \tag{V24.1}\label{c18v24:abstract-input}
\]
Let $B$ be the bad blocks, $A=\mathcal E\setminus B$, and assume
$M_{AA}\succeq aI_A$, where $a>0$. For $r>0$ define
\[
 T=\{e:d(e,B)\ge r\}\subset A,
 \qquad d(e,\varnothing)=+\infty .
 \tag{V24.2}\label{c18v24:core}
\]

\begin{lemma}[Full inverse protected by a good region]
\label{c18v24:inverse-lemma}
With these assumptions,
\[
 P_TM^{-1}P_T\preceq
 \left[\frac1a+\frac{S(1+g_+/a)}{2ar^2g_-}\right]P_T .
 \tag{V24.3}\label{c18v24:inverse-abstract}
\]
In particular, $r^2\ge S(1+g_+/a)/(2g_-)$ implies
$P_TM^{-1}P_T\preceq(2/a)P_T$.
\end{lemma}
\begin{proof}
If $T$ is empty the assertion is vacuous. If $B$ is empty,
$M\succeq aI$ and the stronger $M^{-1}\preceq a^{-1}I$ holds.
Otherwise write
\[
 M=\begin{pmatrix}A_0&C\\ C^*&B_0\end{pmatrix},
 \quad A_0=M_{AA},\quad
 X=A_0^{-1}C,\quad
 \Sigma=B_0-C^*A_0^{-1}C .
 \tag{V24.4}\label{c18v24:blocks}
\]
The two uses of positivity that will be needed are
\[
 \Sigma\succeq g_-I_B,\qquad
 C^*A_0^{-1}C\preceq B_0\preceq g_+I_B .
 \tag{V24.5}\label{c18v24:schur-input}
\]
For the first, minimizing the $M$-energy of $(v,b)$ over $v$
gives $\langle b,\Sigma b\rangle\ge g_-\|b\|^2$.
The second follows from positivity of $\Sigma$ and the upper
compression bound on $B_0$.

Given $b$ on $B$, let $u=(-Xb,b)$, its harmonic extension across
$A$: $(Mu)_A=0$. Moreover
\[
 \|u\|^2\le(1+g_+/a)\|b\|^2 ,
 \tag{V24.6}\label{c18v24:harmonic-norm}
\]
because
$\|Xb\|^2\le a^{-1}\langle Xb,A_0Xb\rangle
 =a^{-1}\langle b,C^*A_0^{-1}Cb\rangle$.

Set $\chi_e=\min\{1,d(e,B)/r\}$.
Then $\chi=0$ on $B$, $\chi=1$ on $T$, and
$|\chi_e-\chi_f|\le d(e,f)/r$.
Since $(Mu)_A=0$ and $\chi^2u$ vanishes on $B$,
$\operatorname{Re}\langle Mu,\chi^2u\rangle=0$.
The exact localization identity therefore reads
\[
 \begin{split}
 \langle\chi u,M\chi u\rangle
 &=-\frac12\sum_{e,f}(\chi_e-\chi_f)^2
               \operatorname{Re}\langle u_e,M_{ef}u_f\rangle\\
 &\le\frac{S}{2r^2}\|u\|^2 .
 \end{split}
 \tag{V24.7}\label{c18v24:caccioppoli}
\]
The last step is the block Schur bound. Symmetry of $M$ and
$d$ gives the same weighted column bound as the row bound
in \eqref{c18v24:abstract-input}. This argument does not
assign a favorable sign to any individual summand.

The vector $\chi u$ is supported on $A$, so its energy is at
least $a\|\chi u\|^2$. Equations
\eqref{c18v24:harmonic-norm}--\eqref{c18v24:caccioppoli} imply
\[
 \|P_TX\|^2\le\frac{S(1+g_+/a)}{2ar^2}.
 \tag{V24.8}\label{c18v24:cross-inverse}
\]
Finally, the exact block inverse gives
\[
 (M^{-1})_{AA}=A_0^{-1}+X\Sigma^{-1}X^*.
\]
Its compression to $T$, using
$A_0^{-1}\preceq a^{-1}I$ and
$\Sigma^{-1}\preceq g_-^{-1}I$, proves
\eqref{c18v24:inverse-abstract}.
\end{proof}

The global lower bound $g_-$ was not dropped: it appears in
the collar size needed to suppress the remote Schur return.
The lemma is an inverse-metric statement, not a theorem
about a physical Hamiltonian gap.

\subsection{Populating the lemma with the actual Yang--Mills map}

Keep all of V23's parameters and definitions, including
$s=\ell^2$, the analysis radius $\varrho_\ell$, embedded
box side $L_{\varrho_\ell}$, threshold $h_{\varrho_\ell}$,
and $N\ge L_{\varrho_\ell}+1$.
In particular the retained map $\mathfrak c$, the original
Hamiltonian and its physical clock are unchanged. Write
\[
 \begin{gathered}
 G(U)=D\mathfrak c(U)D\mathfrak c(U)^*,\qquad
 \overline G(g)=\mathbb E_\nu[G(U)\mid\mathfrak c(U)=g],\\
 g_-=\ell e^{-32s},\qquad
 g_+=\ell e^{32s},\qquad
 a_\ell=\frac{m^2}{16\ell}.
 \end{gathered}
 \tag{V24.9}\label{c18v24:native-input}
\]
The deterministic metric bounds $g_-I\preceq G\preceq g_+I$
are V17.2, proved from the actual flow diffeomorphism and the
exact straight-chain differential. They are not assumed
spectral information about $H-E_0$.

For the native fine configuration and the conditional coarse
configuration, respectively, define
\[
 \begin{aligned}
 B_\nu(U)&=\{e:d_e(U)>1/4\},&
 T_\nu(U)&=\{e:D_c(e,B_\nu(U))\ge\mathfrak r_\ell\},\\
 B_\mu(g)&=\{e:\bar d_e(g)>1/4\},&
 T_\mu(g)&=\{e:D_c(e,B_\mu(g))\ge\mathfrak r_\ell\},
 \end{aligned}
 \tag{V24.10}\label{c18v24:native-core}
\]
where $d_e,\bar d_e$ are \emph{precisely} V23's diagonal
defect weights. Choose the integer collar width
\[
 \mathfrak r_\ell=
 \left\lceil\frac4m e^{52s}(40s+2\ell+1)\right\rceil .
 \tag{V24.11}\label{c18v24:collar}
\]
Distances are periodic coarse-tail distances. If a bad set
is empty, its protected set is the full coarse link set.
No additional simply connected box of radius
$\mathfrak r_\ell$ is required; these are distance cutoffs
on the existing torus.

\begin{theorem}[Native and conditional full inverse bounds]
\label{c18v24:native-inverse}
For every fine configuration, and almost everywhere in the
conditional statement,
\[
 \begin{split}
 P_{T_\nu}G(U)^{-1}P_{T_\nu}
                 &\preceq\frac{32\ell}{m^2}P_{T_\nu},\\
 P_{T_\mu}\overline G(g)^{-1}P_{T_\mu}
                 &\preceq\frac{32\ell}{m^2}P_{T_\mu}.
 \end{split}
 \tag{V24.12}\label{c18v24:full-inverse}
\]
Each inverse is the inverse of the \emph{full} metric before
compression. This is stronger than inverting a good principal
compression.
\end{theorem}
\begin{proof}
V23.14 implies $G_{AA}\succeq a_\ell I_A$ on
$A=B_\nu(U)^c$, since $d_e\le1/4$ there.
V23.16 gives the corresponding statement for $\overline G$
on $B_\mu(g)^c$. Both full metrics obey the deterministic
bounds in \eqref{c18v24:native-input}; conditional integration
preserves these matrix inequalities.

Both also obey V23's kernel moment bound with
\[
 S=S_\ell=\ell^{-1}e^{40s}
                       \big[(40s+2\ell)^2+40s\big].
\]
For $\overline G$, this follows from
$\|\mathbb E[G_{ef}\mid g]\|
 \le\mathbb E[\|G_{ef}\|\mid g]$ and summation.
There are finitely many blocks, so one common full-measure
set suffices.

As $g_+/a_\ell=16\ell^2e^{32s}/m^2\ge1$,
\[
 \frac{S_\ell(1+g_+/a_\ell)}{2g_-}
 \le \frac{16}{m^2}e^{104s}
                         \big[(40s+2\ell)^2+40s\big]
 \le \mathfrak r_\ell^2 .
\]
Apply Lemma~\ref{c18v24:inverse-lemma} with $a=a_\ell$.
The resulting constant $2/a_\ell=32\ell/m^2$ proves
both assertions. This is a separate application to each
matrix, not an interchange of expectation and inversion.
\end{proof}

Equivalently, for either protected set $T$, the Schur
complement of its full metric $M$ onto $T$ obeys
\[
 M_{TT}-M_{T,T^c}M_{T^c,T^c}^{-1}M_{T^c,T}
                  \succeq\frac{m^2}{32\ell}I_T .
 \tag{V24.13}\label{c18v24:protected-schur}
\]
Indeed its inverse is the compression $(M^{-1})_{TT}$.
The cases $T=\varnothing$ or $T^c=\varnothing$ have their
usual empty-block interpretations.

\subsection{An actual joint lift with all other outputs fixed}

\begin{corollary}[Complement-preserving native fine lift]
\label{c18v24:lift}
At any fine configuration $U$, let $v$ be a coarse tangent
vector supported on $T_\nu(U)$. Then
\[
 \begin{split}
 \mathcal R_Uv&=D\mathfrak c(U)^*G(U)^{-1}v,\\
 D\mathfrak c(U)\mathcal R_Uv&=v
             \quad\text{on the full retained link product},\\
 \|\mathcal R_Uv\|^2&\le\frac{32\ell}{m^2}\|v\|^2 .
 \end{split}
 \tag{V24.14}\label{c18v24:joint-lift}
\]
In particular every retained output outside $T_\nu(U)$
has zero variation. This is the minimal-norm fine lift,
and its formula is gauge equivariant.
\end{corollary}
\begin{proof}
V17's submersion result makes $G$ invertible everywhere.
Direct multiplication gives the full output identity and
\[
 \|\mathcal R_Uv\|^2=\langle v,G(U)^{-1}v\rangle .
\]
Apply \eqref{c18v24:full-inverse}.
Every other fine lift differs by an element of
$\ker D\mathfrak c$, orthogonal to
$\operatorname{ran}D\mathfrak c^*$, proving minimality.
Gauge equivariance of the map and the isometry of the
product group action prove equivariance of the formula.
\end{proof}

Unlike V22's cube-only right inverse, this lift actually
fixes all complementary retained outputs. It need not have
fine support in the collar. The global formula
$D\mathfrak c^*G^{-1}$ is smooth, but the selected protected
set is only measurable; no derivative of its indicator
has been taken.

There is also an important distinction between the two
lines of \eqref{c18v24:full-inverse}. The set $T_\nu(U)$
may vary inside a fixed coarse fibre. The conditional
set $T_\mu(g)$ controls $\overline G(g)^{-1}$, not
$G(U)^{-1}$ at every point of that fibre. It has not been
used to claim a uniform fine horizontal field on an entire
conditional fibre, or to replace
$\mathbb E[G^{-1}\mid g]$ by $\overline G(g)^{-1}$.
The pointwise native joint lift is exactly the statement
of Corollary~\ref{c18v24:lift}.

\subsection{Paying the protected-set loss without independence}

Let $E_c=3n^3$ and
$p=p_{\ell,\varrho_\ell,\beta}$ from V23.10.
For either native or conditional defect,
$\int d_e\,d\nu\le p$ or $\int\bar d_e\,d\mu\le p$.
The former follows directly from V23's partition definition;
the latter is its conditional expectation.

\begin{proposition}[Native density of links lost to the collar]
\label{c18v24:density}
Both protected sets satisfy
\[
 \begin{split}
 \frac1{E_c}\mathbb E_\nu|T_\nu(U)^c|
       &\le \min\{1,12(2\mathfrak r_\ell+1)^3p\},\\
 \frac1{E_c}\mathbb E_\mu|T_\mu(g)^c|
       &\le \min\{1,12(2\mathfrak r_\ell+1)^3p\}.
 \end{split}
 \tag{V24.15}\label{c18v24:density-bound}
\]
No independence or percolation estimate is assumed.
\end{proposition}
\begin{proof}
For each particular link the probability of being bad is
at most $4p$, by Markov at the threshold $1/4$.
If a link lies outside its protected set, some bad link
has periodic tail distance less than $\mathfrak r_\ell$.
There are at most $3(2\mathfrak r_\ell+1)^3$ such links,
including their three orientations. A union bound gives
the displayed bound for each fixed link. Sum over links
and divide by $E_c$. The same proof works under either law.
If the collar winds around the torus this count only
overestimates the number of available links, so the
argument needs no new embedding assumption.
\end{proof}

For $\ell\ge2$, the ceiling in \eqref{c18v24:collar} gives
$\mathfrak r_\ell\le169\ell^2e^{52\ell^2}/m$.
Since $2\mathfrak r_\ell+1\le3\mathfrak r_\ell$, V23.21
therefore yields the explicit common upper bound
\[
 \min\{1,C_{24}\ell^{29}e^{256\ell^2}\varepsilon_\beta\},
 \qquad
 C_{24}=\frac{324\cdot169^3}{m^3}C_{23}.
 \tag{V24.16}\label{c18v24:density-scale}
\]
The powers are paid: the collar contributes
$\ell^6e^{156\ell^2}$ to V23's
$\ell^{23}e^{100\ell^2}$ probability estimate.
Thus the expected lost fraction tends to zero when
$\beta\to\infty$ and
\[
 \ell^2\le c\log\beta,\qquad 0<c<1/512,
 \qquad N\ge L_{\varrho_\ell}+1 .
 \tag{V24.17}\label{c18v24:log-scale}
\]
This is an asymptotic sufficient regime with very large
constants, not a practical-scale estimate or a chosen
physical continuum trajectory.

\subsection{What is now available, and what is not}

We have progressed from V22's local output-only lift and
V23's good-link principal compression to a full-inverse
bound and an actual simultaneous lift holding the entire
complement fixed. All native coefficients in the inverse
argument were supplied: the crude full metric bounds,
the good-region floor, the kernel moment and the
bad-collar probability. No conditional Poincar\'e
inequality or Hamiltonian mass gap was inserted as an input.

The desired gradient-weighted defect bound still does
not follow from \eqref{c18v24:density-bound}. A low-energy
source can place disproportionate derivative weight in
a small exceptional set; its size alone does not control
that weight. Nor does the norm bound on $\mathcal R_U$
control its divergence against $\nu$, derivatives of
conditional expectations, or a score along the full
coarse fibre. Those quantities are the next useful
inputs for a native conditional-variance argument.

One must also retain the relative-energy qualification.
V21's small localization error at a declared scale
$\gamma_\beta$ is not a bound small relative to an
eigenvalue $E\ll\gamma_\beta$. The present metric inverse
result does not alter that fact. No spectral gap follows
by combining those two statements without the missing
source-weighted and variance estimates.

The native coarse marginal Poincar\'e problem and nonlinear
fibre conditional variance remain open. So do weighted
one-form localization, source-weighted return and physical
scale transport. Interacting four-dimensional reconstruction
and the full positive physical Yang--Mills mass gap are
not proved.

% Keep the last V24 contents entry with the preceding contents lines.
\addtocontents{toc}{\protect\enlargethispage{\baselineskip}}
\subsection{Verification and preservation}

The diagnostic checks the harmonic-extension identity
with exact block Schur algebra, the cutoff energy
identity and moment bound, the full complement-preserving
lift, periodic collar counts, conditional inverse
noncommutation and the native scale constants.
These finite algebra fixtures do not simulate the
unknown interacting ground law or certify its spectrum.

V24 preserves V23's entire body except the title and
replaces only that active source. All other 47 note
files are unchanged. Only \texttt{.tex} files remain
in \texttt{latex\_notes}; build and inspection products
are outside it. The next companion checkpoint is V25,
the next broad primary-literature sweep V30, and the
scheduled restart V100.
% END CYCLE 18 VERSION 24 APPEND
% BEGIN CYCLE 18 VERSION 25 APPEND
\clearpage
\section{V25: Full-fibre scores and a paid geometric contribution}
\label{c18v25:context}

V24 controls the norm of a joint lift on protected directions,
but not its weighted divergence. We now differentiate the
\emph{full} coarse conditioning map using another explicit
smooth lift. Its Jacobian contribution has a volume-independent
mean-square bound, including a bound weighted by actual
fixed-momentum low-energy sources. The remaining vacuum-score
covariance is displayed, not replaced by a fibre gap.

The scheduled companion review found no changes in the nine
sources checked at V20: NS V26, SM--Einstein V72, shell-mixing
V16, parallel YM V5, source-selection V52, stationarity V54,
exact-action V45, and perturbative ESM V64 and V69. Their
terminal scope statements were reread. Fixed-time prepared
response, finite source-frame identities and a flat one-loop
coefficient do not supply the native nonlinear variance
estimate. The original measure, map and physical clock stay fixed.

This is not a rediscovery of a Jacobian or conditional-mean
identity. Archive f16.2 V74 already proves the spatial-flow
Jacobian and a pointwise transported score for a Wilson slice
law; current V13 treats unflowed chain scores; f17 V93 treats
native one-link scores. We use their methods for the present
Hamiltonian ground law and its \emph{flowed, sampled} fibres.
The new quantitative input is the source-weighted geometric
score estimate below. For external comparison, Leli\`evre--Zhang,
\emph{Pathwise estimates for effective dynamics: the case of
nonlinear vectorial reaction coordinates},
\href{https://arxiv.org/html/1805.01928}{arXiv:1805.01928},
Section 2.2, Assumption 4, assumes a uniform conditional
Poincar\'e inequality. That missing hypothesis is not imported
here; their pathwise theorem is not applied to YM.

\subsection{A smooth full-complement lift without a protected mask}

Keep $N=\ell n$, $\ell\ge2$, $n\ge5$, $\beta=b/\kappa>0$,
$\nu=\Omega^2dU$ and
\[
 c=\mathfrak b_\ell\circ\Phi_s,\qquad
 \mu=c_*\nu=\rho(g)\,dg,\qquad
 \Pi F(g)=\mathbb E_\nu[F\mid c=g],\qquad s\ge0 .
 \tag{V25.1}\label{c18v25:setup}
\]
The principal application remains $s=\ell^2$. No large
embedded good box is required for this appendix. Write
$E_{\rm sp}=\sum_p(1-w_p)$ and $W_t=\Phi_t(U)$.

For a coarse link $\gamma$, let $h_{\gamma,j}(W)$ be the
prefix before its $j$th fine chain link, as in V13. For the
unit-round coarse frame $X_{\gamma,a}$ define
\[
 \begin{split}
 Y_{\gamma,a}(W)
 &=\frac1\ell\sum_{j=0}^{\ell-1}\sum_d
       (\operatorname{Ad}h_{\gamma,j})_{ad}
                 X_{e_{\gamma,j},d},\\
 B_{\gamma,a}(U)
 &=D\Phi_{-s}(W_s)Y_{\gamma,a}(W_s).
 \end{split}
 \tag{V25.2}\label{c18v25:lift}
\]
Write $B_\gamma v=\sum_a v_aB_{\gamma,a}$ and $Bv$ for
the complete bank. The notation $B_\gamma F$ denotes the
three-vector $(B_{\gamma,a}F)_a$.

\begin{lemma}[Exact lift and deterministic spatial envelopes]
\label{c18v25:envelope}
The fields are smooth and
\[
 Dc\,Bv=v,\qquad \|Bv\|^2\le\ell^{-1}e^{32s}|v|^2 .
 \tag{V25.3}\label{c18v25:full-lift}
\]
The first identity is on the entire coarse output product.

For $0\le t\le s$ define the tangent map
$B_\gamma^t=D\Phi_t(U)B_\gamma(U)$.
There are deterministic nonnegative weights $k_{\gamma,e}(t)$
such that, for all configurations,
\[
 \|(B_\gamma^t)_e\|\le k_{\gamma,e}(t),\qquad
 \sum_e k_{\gamma,e}(t)=e^{16(s-t)}.
 \tag{V25.4}\label{c18v25:weights}
\]
\end{lemma}
\begin{proof}
Prefix independence gives $\operatorname{div}_{dW}Y_{\gamma,a}=0$.
The chain derivative formula and orthogonality of each adjoint
matrix give $D\mathfrak b_\ell Y_{\gamma,a}=X_{\gamma,a}$.
Different chains have disjoint fine links, so $Y^*Y=\ell^{-1}I$.
Composition with the inverse flow proves the first identity in
\eqref{c18v25:full-lift}; V17's inverse differential norm proves
the second.

For clarity, the sharper block envelope from f16.2 V74 is
recalled with its proof. Let $\mathsf N_{ef}$, $e\ne f$, count
shared plaquettes, and $\mathsf N_{ee}=0$. Every row and
column sums to $12$. In linkwise parallel orthonormal frames,
the variational equation has diagonal block norm at most $4$
and off-diagonal block norm at most $\mathsf N_{ef}$. Iterating
it in either time direction bounds each differential block by
the corresponding entry of
\[
 \mathsf K_v=\exp\big(v(4I+\mathsf N)\big),\qquad v\ge0 .
\]
All its row and column sums are $e^{16v}$. Changing endpoint
frames conjugates blocks by orthogonal matrices and does not
change these norm bounds. No invariant-frame derivative is
being estimated by omitting a frame term.

Since $B_\gamma^t=D\Phi_{t-s}(W_s)Y_\gamma(W_s)$, take
$k_{\gamma,e}(t)=\ell^{-1}
 \sum_{j=0}^{\ell-1}(\mathsf K_{s-t})_{e,e_{\gamma,j}}$.
The adjoint blocks in $Y_\gamma$ have norm $1/\ell$.
The column sum proves \eqref{c18v25:weights}.
\end{proof}

The bank is gauge equivariant in the usual transformed coarse
frame, but is not claimed to be fine-local. It need not minimize
the fine norm. In particular V24's protected polynomial norm
bound belongs to $Dc^*(DcDc^*)^{-1}$, not automatically to $B$.
Here we knowingly retain the global exponential norm cost in
exchange for a tractable smooth divergence on every fibre.

\subsection{The exact score and its nonextensive geometric part}

Let $\mathcal J_s(U)$ be the positive Haar-volume Jacobian of
$\Phi_s$. Set
\[
 Z_{\gamma,a}=B_{\gamma,a}\log\mathcal J_s,\qquad
 S_{\gamma,a}=\operatorname{div}_\nu B_{\gamma,a}.
\]
Both $Z_\gamma$ and $S_\gamma$ are three-vectors. Then
\[
 \begin{split}
 \log\mathcal J_s
   &=-12\int_0^s\sum_p w_p(W_t)\,dt,\\
 Z_{\gamma,a}
   &=12\int_0^s
       dE_{\rm sp}(W_t)[B_{\gamma,a}^t]\,dt,\\
 S_\gamma&=2B_\gamma\log\Omega-Z_\gamma .
 \end{split}
 \tag{V25.5}\label{c18v25:score}
\]
These identities include the volume Jacobian; no coarea
factor is silently discarded.

Indeed $\Delta w_p=-12w_p$, since each trace is linear
in each of its four unit-quaternion variables. For
$\dot W=-\nabla E_{\rm sp}$ this gives
$\operatorname{div}\dot W=-12\sum_pw_p$.
The volume variational equation proves the first line.
Differentiation gives the second. Finally a pullback field
satisfies
\[
 \operatorname{div}_{dU}B_{\gamma,a}
  =(\operatorname{div}_{dW}Y_{\gamma,a})\circ\Phi_s
                  -B_{\gamma,a}\log\mathcal J_s
  =-Z_{\gamma,a}.
\]
This follows, for example, by applying the Lie derivative
to $\Phi_s^*dW=\mathcal J_s\,dU$. Adding
$B_{\gamma,a}\log\Omega^2$ proves the third line.

Put $A_*=3\sqrt2$, $\varepsilon_\beta=\min(1,A_*\beta^{-1/2})$
and
\[
 L_{\beta,s}=4\beta e^{16s}+3(e^{16s}-1),\qquad
 I_{\beta,s}^{1/2}
 =2\sqrt{A_*}\,\beta^{1/4}e^{16s}
       +3\sqrt{2\varepsilon_\beta}(e^{16s}-1).
 \tag{V25.6}\label{c18v25:constants}
\]

\begin{theorem}[Native score budgets for each coarse link]
\label{c18v25:budgets}
Uniformly in spatial volume and the selected coarse link,
\[
 \begin{aligned}
 |Z_\gamma|&\le3(e^{16s}-1),&
 \|Z_\gamma\|_{L^2(\nu)}
       &\le3\sqrt{2\varepsilon_\beta}(e^{16s}-1),\\
 |S_\gamma|&\le L_{\beta,s},&
 \|S_\gamma\|_{L^2(\nu)}^2&\le I_{\beta,s}.
 \end{aligned}
 \tag{V25.7}\label{c18v25:score-bounds}
\]
The three colors are included in each vector norm.
\end{theorem}
\begin{proof}
The inherited actual-vacuum bounds are
\[
 |\nabla_e\log\Omega|\le2\beta,\qquad
 \int|\nabla_e\log\Omega|^2d\nu\le A_*\sqrt\beta .
\]
Their maximum-principle and energy-density proofs in f17 V93
and current V6 were rechecked; they do not identify $\Omega$
with a Gibbs trial. Applying the deterministic weights at $t=0$
gives
\[
 |2B_\gamma\log\Omega|\le4\beta e^{16s},\qquad
 \|2B_\gamma\log\Omega\|_2
          \le2\sqrt{A_*}\,\beta^{1/4}e^{16s}.
\]
The triangle inequality here is for the complete three-vector,
so an additional color factor is unnecessary.

For the geometric term, each link belongs to four plaquettes:
\[
 |\nabla_eE_{\rm sp}|\le4,\qquad
 |\nabla_eE_{\rm sp}|^2
  \le4\sum_{p\ni e}|\nabla_e(1-w_p)|^2
  \le8\sum_{p\ni e}(1-w_p).
 \tag{V25.8}\label{c18v25:force}
\]
V18 gives $\mathbb E_\nu[1-w_p(W_t)]\le\varepsilon_\beta$,
hence $\|\nabla_eE_{\rm sp}(W_t)\|_2\le\sqrt{32\varepsilon_\beta}$.
The matrix-block weights are deterministic; correlation
between the lift and this force requires no independence
assumption. Minkowski's inequality and \eqref{c18v25:weights}
give the bounds
\[
 \begin{split}
 |Z_\gamma|&\le48\int_0^s e^{16(s-t)}dt,\\
 \|Z_\gamma\|_2&\le12\sqrt{32\varepsilon_\beta}
                            \int_0^s e^{16(s-t)}dt .
 \end{split}
\]
Integration and the triangle inequality prove the theorem.
\end{proof}

\subsection{Conditioning on all retained outputs}

The compact product coordinates of V13, after the flow
diffeomorphism, justify smooth fibre integration. All the
following equalities therefore hold for smooth functions
at each finite regulator, and not just for a chosen local
chart or a single-output conditional.

\begin{proposition}[Full conditional differentiation and covariance]
\label{c18v25:conditional}
For any smooth complex $F$ on the fine product,
\[
 \begin{split}
 \nabla_\gamma\log\rho&=\Pi S_\gamma,\\
 \nabla_\gamma\Pi F
    &=\Pi(B_\gamma F)+\Pi[(F-\Pi F\circ c)S_\gamma].
 \end{split}
 \tag{V25.9}\label{c18v25:conditional-derivative}
\]
Let $V_\gamma(g)=\Pi[|S_\gamma-\Pi S_\gamma\circ c|^2](g)$.
Then
\[
 \begin{gathered}
 |\nabla_\gamma\log\rho|\le L_{\beta,s},\qquad
 \int|\nabla_\gamma\log\rho|^2d\mu\le I_{\beta,s},\\
 0\le V_\gamma\le L_{\beta,s}^2,\qquad
 \int V_\gamma\,d\mu\le I_{\beta,s},\\
 |\nabla_\gamma\Pi F-\Pi(B_\gamma F)|^2
   \le V_\gamma\,\Pi[|F-\Pi F\circ c|^2].
 \end{gathered}
 \tag{V25.10}\label{c18v25:conditional-bounds}
\]
No conditional variance inequality for $F$ has been assumed.
\end{proposition}
\begin{proof}
For a smooth coarse test $h$, integration by parts in the
fine measure gives
\[
 \int(\Pi F)X_{\gamma,a}h\,d\mu
   =-\int h\{\Pi(B_{\gamma,a}F)+\Pi(FS_{\gamma,a})\}\,d\mu .
\]
Coarse Haar integration by parts gives the same left hand
side as $-\int h\{X_{\gamma,a}\Pi F+
(\Pi F)X_{\gamma,a}\log\rho\}\,d\mu$.
First set $F=1$, then subtract that identity times $\Pi F$.
This proves \eqref{c18v25:conditional-derivative}.
Jensen and conditional Cauchy--Schwarz prove
\eqref{c18v25:conditional-bounds}. In particular
$V_\gamma=\Pi|S_\gamma|^2-|\Pi S_\gamma|^2$,
which explains both bounds without a spurious factor four.
Complex functions follow by linearity and the complex
Cauchy--Schwarz inequality.
\end{proof}

For a set $\mathcal S$ of $m$ coarse links, put
$QF=F-\Pi F\circ c$. The full-bank norm
\eqref{c18v25:full-lift} and Jensen also give
\[
 \int\sum_{\gamma\in\mathcal S}|\nabla_\gamma\Pi F|^2d\mu
 \le\frac{2e^{32s}}{\ell}\int|\nabla F|^2d\nu
                  +2mL_{\beta,s}^2\|QF\|_2^2 .
 \tag{V25.11}\label{c18v25:Sobolev}
\]
The first term has no $m$ factor: the bank norm is applied
before summing. The second term retains that factor because
no cross-link covariance decay was proved. Extending this
bound to the whole coarse lattice with a volume-independent
error requires additional structure. Bounding $\|QF\|^2$
by a desired fibre gap here would assume an unresolved input.

These identities are independent of the choice of smooth lift:
if $B'$ is another such bank, $Dc(B'-B)=0$ and
$\Pi\operatorname{div}_\nu(B'-B)=0$. Fine integration by
parts also shows that its change in $\Pi(BF)$ cancels its
change in the score covariance. Thus changing lift does not
remove the source. Individual frame components need not be
gauge-invariant observables; the vector identities are
gauge-covariant, and the summed estimates apply to physical
functions without changing the original Gauss law.

\subsection{Paying the geometric score on actual low-energy sources}

The vacuum $L^2$ bound is not a multiplier estimate against
arbitrary sources. The next theorem supplies a precise
weighted upgrade for the geometric term.

Let $A=\Omega^{-1}(H-E_0)\Omega$ on $L^2(\nu)$ and
$\mathsf M=3N^3$. Fix one fine-lattice translation character,
and let $F$ belong to that character subspace and to the
spectral window $\mathbf1_{[0,\Lambda]}(A)$.
Set
\[
 q_\Lambda=\varepsilon_\beta+\frac{\Lambda}{b\mathsf M}.
 \tag{V25.12}\label{c18v25:q}
\]

\begin{theorem}[Energy-weighted geometric score]
\label{c18v25:weighted}
For every coarse link $\gamma$,
\[
 \int |F|^2|Z_\gamma|^2d\nu
 \le54(e^{16s}-1)^2q_\Lambda\|F\|_2^2 .
 \tag{V25.13}\label{c18v25:weighted-bound}
\]
At $s=\ell^2$, $\Lambda\le C\kappa$ and
$\ell^2\le c\log\beta$ with fixed $c<1/64$, the displayed
multiplier bound tends to zero as $\beta\to\infty$,
uniformly over the allowed volumes.
\end{theorem}
\begin{proof}
For $F\ne0$, use the finite measure $|F|^2d\nu$.
Its translation invariance follows from the single-character
hypothesis and invariance of $\Omega$. Spatial flow is
translation equivariant and decreases $E_{\rm sp}$ pointwise.
The original positive kinetic term and the spectral window give
\[
 \begin{split}
 \int |F|^2 E_{\rm sp}(W_t)d\nu
 &\le\int |F|^2 E_{\rm sp}(U)d\nu\\
 &\le b^{-1}\langle\Omega F,H\Omega F\rangle
 \le(E_0+\Lambda)b^{-1}\|F\|_2^2
 \le\mathsf M q_\Lambda\|F\|_2^2 .
 \end{split}
 \tag{V25.14}\label{c18v25:weighted-energy}
\]
No cubic symmetry of $F$ is asserted. Translation invariance
alone makes the weighted expectation of a fixed-orientation
plaquette constant in its position. There are $N^3$ such
plaquettes. Positivity thus bounds each individual weighted
defect by $3q_\Lambda\|F\|_2^2$. Equation
\eqref{c18v25:force} then yields, for every fine link and time,
\[
 \big\||F|\,\nabla_e E_{\rm sp}(W_t)\big\|_2
       \le\sqrt{96q_\Lambda}\,\|F\|_2 .
\]
Apply the same deterministic block envelope and Minkowski
argument as before, now in $|F|^2d\nu$. The square of
$12\sqrt{96}/16$ is $54$, proving
\eqref{c18v25:weighted-bound}. The zero function is immediate.

Finally $(e^{16s}-1)^2\le e^{32s}\le\beta^{32c}$ and
$q_\Lambda\le A_*\beta^{-1/2}+C/(\beta\mathsf M)$.
Both powers decay for $c<1/64$.
\end{proof}

This estimate uses no $L^\infty$ bound for the low-energy
wavefunction and no bad-all-box event. It is already compatible
with V24's stricter $c<1/512$ regime when that result is also
needed. A general superposition of fine momenta need not
satisfy the individual-link version of
\eqref{c18v25:weighted-energy}; the stated character restriction
must not be dropped. Also $QF$ need not stay in the same
fine-momentum spectral window, so
\eqref{c18v25:weighted-bound} cannot simply be inserted with
$QF$ into the conditional covariance bound.

\subsection{Remaining bottleneck and verification}

The new paid quantity is the geometric part $Z_\gamma$ of the
actual smooth full-fibre score. The full source is still
$S_\gamma=2B_\gamma\log\Omega-Z_\gamma$. Its established
$L^2$ budget has a vacuum-score term of order
$\sqrt\beta e^{32s}$ after squaring, and its pointwise bound
is worse. Neither is a small conditional covariance or a
native fibre Poincar\'e theorem. The next useful work is a
source-weighted estimate for that centered vacuum term, or
native nonlinear conditional variance, retaining the
source and lift actually used.

V23's coarse-gradient-weighted diagonal defect is not removed.
V24's protected inverse and the present smooth-score estimates
belong to distinct lifts and cannot be merged without proof.
V21's error small at a declared target energy is still not
small relative to every eigenvalue far below that target.
No physical time rescaling, interacting continuum
reconstruction or full YM mass gap has been established.

The reproducible diagnostic checks literal quaternion-chain
derivatives, native plaquette incidence and Laplacian signs,
non-minimal versus minimal full lifts, a nonconstant-Jacobian
pullback identity, exact normalized conditional derivatives,
and the displayed constants. These are algebraic regression
fixtures, not simulations of the unknown native ground law or
proof-assistant certification. The analytical proofs carry
the uniform assertions.

V24's complete body is preserved except its title; the source
is replaced only after the append. All 47 other note sources
are protected by hashes. Only \texttt{.tex} files are kept in
\texttt{latex\_notes}. The companion checkpoint is now V30,
as is the next broad primary-literature sweep; the archive
restart remains scheduled at V100.
% END CYCLE 18 VERSION 25 APPEND
% BEGIN CYCLE 18 VERSION 26 APPEND
\clearpage
\section{V26: Conditional mean inverses and the score-transport cost}
\label{c18v26:context}

V25 makes the full conditional score explicit, but its raw
vacuum-score bound grows at weak electric coupling. The next
target should not automatically be a small score variance.
A change of horizontal lift can remove that covariance from
the conditional derivative, while transferring its cost to
a vertical correction. We identify that cost for the actual
full flow-sampling map and pay its geometric, inverse-metric
part on a large joint coarse direction space.

The distinction is substantive: V24 controls the inverse of
the conditional mean metric on one selected space, whereas
the mean squared norm of a fine lift involves the
\emph{conditional mean of the inverse}. Those matrices cannot
be interchanged. The first theorem below bounds the latter,
including rare unprotected fine configurations in each fibre.

The underlying conditional Poisson method is already present
in f17 V94, especially V94.8--V94.13, for a one-link split.
The Gaussian harmonic lift is also already present in f8 V18.
We do not count those frameworks as new results. Here the
conditioning is on the whole nonlinear map, its coarea density
is retained, and the geometric part has a newly paid native
bound. The remaining score-Poisson matrix is explicitly
uncontrolled at the needed scales.

V25 was a progress turn and its preservation check passed
again. Its companion checkpoint remains current. The
primary comparison in V25 was rechecked: the conditional
Poincar\'e inequality in Leli\`evre--Zhang,
\href{https://arxiv.org/html/1805.01928}{arXiv:1805.01928},
Assumption 4, is an input to their result, not a consequence
of a right inverse. We import no such inequality for YM.

\subsection{Paying the conditional mean of the full inverse}

Keep all parameters and embedding assumptions of V23--V24,
in particular
\[
 \begin{gathered}
 c=\mathfrak b_\ell\circ\Phi_{\ell^2},\quad
 G=Dc\,Dc^*,\quad \overline G=\mathbb E_\nu[G\mid c=g],
 \quad \nu=\Omega^2dU,\quad \mu=c_*\nu,\\
 m=(2/\pi)e^{-3\pi^2},\qquad
 a=\frac{32\ell}{m^2},\qquad g_-=\ell e^{-32\ell^2}.
 \end{gathered}
 \tag{V26.1}\label{c18v26:parameters}
\]
Here $a$ is an inverse-metric bound, not lattice spacing.
Let $T(U)=T_\nu(U)$ be precisely V24's protected set,
and let $P_{T^c}$ act identically on the three colors of
each excluded coarse link. Define
\[
 \begin{split}
 \mathsf H(g)&=\mathbb E_\nu[G^{-1}\mid c=g],\\
 \pi_e(g)&=\mathbb P_\nu(e\notin T(U)\mid c=g),\\
 \theta_\ell&=ag_-=\frac{32\ell^2}{m^2}e^{-32\ell^2},\qquad
 J(g)=\{e:\pi_e(g)\le\theta_\ell\}.
 \end{split}
 \tag{V26.2}\label{c18v26:mean-inverse}
\]
These definitions condition on all coarse outputs. The new
set $J(g)$ is not V24's set $T_\mu(g)$.

\begin{theorem}[Joint conditional inverse with the rare cost included]
\label{c18v26:inverse}
One has $0<\theta_\ell<1$ for every $\ell\ge2$, and
\[
 \begin{split}
 \mathsf H(g)
   &\preceq 2aI+2g_-^{-1}\operatorname{diag}_e(\pi_e(g)I_3),\\
 P_J\mathsf H(g)P_J&\preceq4aP_J
                  =\frac{128\ell}{m^2}P_J .
 \end{split}
 \tag{V26.3}\label{c18v26:conditional-bound}
\]
If $E_c=3n^3$, then
\[
 \frac1{E_c}\mathbb E_\mu|J(g)^c|
 \le\min\{1,C_{26}\ell^{27}e^{288\ell^2}\varepsilon_\beta\},
 \quad
 C_{26}=\frac{m^2}{32}C_{24}
       =\frac{81\cdot169^3}{8m}C_{23}.
 \tag{V26.4}\label{c18v26:cost}
\]
The expected lost fraction therefore tends to zero when
\[
 \ell^2\le c_0\log\beta,\qquad 0<c_0<1/576,\qquad
 N\ge L_{\varrho_\ell}+1,\qquad \beta\longrightarrow\infty .
 \tag{V26.5}\label{c18v26:scale}
\]
\end{theorem}
\begin{proof}
For any coarse vector $v$, write $v=P_Tv+P_{T^c}v$.
The quadratic Cauchy--Schwarz inequality for the positive
matrix $G^{-1}$, V24's protected compression, and the
global bound $G^{-1}\preceq g_-^{-1}I$ give
\[
 \begin{split}
 v^*G^{-1}v
 &\le2(P_Tv)^*G^{-1}P_Tv
                +2(P_{T^c}v)^*G^{-1}P_{T^c}v\\
 &\le2a|v|^2+2g_-^{-1}|P_{T^c}v|^2 .
 \end{split}
 \tag{V26.6}\label{c18v26:split}
\]
This keeps all cross terms rather than deleting them.
Conditional expectation and
$\mathbb E[P_{T^c}\mid g]=\operatorname{diag}_e(\pi_eI_3)$
prove the first line of \eqref{c18v26:conditional-bound}.
For a vector supported on $J(g)$ the second term is at
most $2g_-^{-1}\theta_\ell|v|^2=2a|v|^2$, proving the
joint bound, with no direction-count factor.

For the threshold, $\pi^2<10$ gives $m^{-2}<3e^{60}$.
The function $x^2e^{-32x^2}$ decreases for $x\ge2$,
so $\theta_\ell<384e^{-68}<1$.
V24's per-link collar bound gives
$\int\pi_e\,d\mu\le C_{24}\ell^{29}e^{256\ell^2}
\varepsilon_\beta$. Markov at $\theta_\ell$ and a sum over
links prove \eqref{c18v26:cost}. The factor
$g_-^{-1}/a$ is paid: it contributes
$m^2e^{32\ell^2}/(32\ell^2)$.
Finally $\varepsilon_\beta\le3\sqrt2\,\beta^{-1/2}$
and $288c_0<1/2$ prove \eqref{c18v26:scale}.
\end{proof}

\begin{corollary}[A conditional whole-fibre squared-norm bound]
\label{c18v26:fibre-lift}
Let $B^0(U)=Dc(U)^*G(U)^{-1}$. For almost every $g$ and
every $v$ supported on $J(g)$,
\[
 Dc(U)B^0(U)v=v\quad(U\in c^{-1}(g)),\qquad
 \mathbb E_\nu[|B^0v|^2\mid g]
                   \le\frac{128\ell}{m^2}|v|^2 .
 \tag{V26.7}\label{c18v26:lift}
\]
\end{corollary}
\begin{proof}
The output identity holds everywhere by the submersion
property, including all complementary outputs.
Also $(B^0)^*B^0=G^{-1}$ pointwise. Apply
\eqref{c18v26:conditional-bound}.
\end{proof}

The norm in \eqref{c18v26:lift} is averaged over the
\emph{entire} conditional fibre, without a fine good-set
restriction. A pointwise polynomial bound throughout that
fibre is not asserted. The full field $B^0$ is smooth;
$J(g)$ selects directions measurably and is not differentiated.
The constants are very large and the sufficient scale is
asymptotic. Neither $\mathsf H$ nor its bound is a
Hamiltonian spectral gap.

\subsection{The actual conditional shape correction}

Here the notation is fixed before using any comparison.
Write $\mu=\rho(g)\,dg$ and disintegrate
$\nu=\int\nu_g\,d\mu(g)$ along $\Sigma_g=c^{-1}(g)$.
Let $\nabla_V$ be the tangent gradient on $\Sigma_g$ for
the induced fine metric. The conditional law includes the
coarea weight: its density relative to induced fibre volume
is proportional to $\Omega^2/\sqrt{\det G}$.
It is not a conditional Wilson trial law.

For every coarse invariant-frame index $\alpha=(e,d)$ set
\[
 \begin{gathered}
 S^0_\alpha=\operatorname{div}_\nu B^0_\alpha,\qquad
 \xi_\alpha=S^0_\alpha-\mathbb E_\nu[S^0_\alpha\mid c],\\
 \mathcal L_g=-\operatorname{div}_{\nu_g}\nabla_V,\qquad
 \phi_\alpha=\mathcal L_g^{-1}\xi_\alpha,\qquad
 \mathsf M_{\alpha\delta}(g)
    =\langle\xi_\alpha,\mathcal L_g^{-1}\xi_\delta
                                      \rangle_{L^2(\nu_g)} .
 \end{gathered}
 \tag{V26.8}\label{c18v26:shape}
\]
The inverse is on mean-zero fibre functions. There is
\emph{no factor of $\kappa$} in $\mathcal L_g$; using
$\kappa\mathcal L_g$ would put a compensating factor
$\kappa$ in the formula for $\mathsf M$.
These are the scores of $B^0$, not V25's different inverse-flow
lift. V25's bounds on that lift's divergence are not assigned
to $S^0$.

We need only finite-regulator existence here. V13's product
coordinate change, composed with the flow diffeomorphism,
identifies every fibre smoothly with a fixed compact
connected product of $\mathrm{SU}(2)$ groups.
The vertical metric and density are smooth and strictly
positive. Comparing on that compact product with its fixed
product metric and Haar density gives a positive Poincar\'e
constant at each finite regulator; compactness of the base
allows finite comparison constants common to $g$.
No uniformity in volume, $\beta$ or $\ell$ follows.
Lax--Milgram on mean-zero functions therefore defines
\eqref{c18v26:shape}. A fixed product harmonic Galerkin basis
gives measurable solutions as limits; the usual smooth
finite-regulator elliptic solutions give the same weak
identities. None of this supplies the quantitative inverse
estimate sought in the programme.

\begin{proposition}[Minimum score-compatible lift cost]
\label{c18v26:minimum}
Put $\widetilde B_\alpha=B^0_\alpha+\nabla_V\phi_\alpha$.
Then, in the weak sense sufficient for fibre integration,
\[
 \begin{gathered}
 Dc\,\widetilde B_\alpha=X_\alpha,\qquad
 \operatorname{div}_\nu\widetilde B_\alpha
                    =(X_\alpha\log\rho)\circ c,\\
 \mathsf W(g):=\mathbb E_\nu[\widetilde B^*\widetilde B\mid g]
                         =\mathsf H(g)+\mathsf M(g).
 \end{gathered}
 \tag{V26.9}\label{c18v26:transport}
\]
Among lifts with these output and divergence conditions,
$\mathsf W$ is minimal in quadratic-form order.
For any constant vector $v$ in the coarse tangent at $g$,
\[
 \begin{split}
 v^*\mathsf M(g)v
 &=\int_{\Sigma_g}|\nabla_V\phi_v|^2d\nu_g\\
 &=\sup_{\substack{h\in H^1(\nu_g),\ \int h\,d\nu_g=0\\
                            \int|\nabla_Vh|^2d\nu_g>0}}
       \frac{|\langle\xi_v,h\rangle_{\nu_g}|^2}
            {\int|\nabla_Vh|^2d\nu_g},
 \end{split}
 \tag{V26.10}\label{c18v26:variational}
\]
where $\xi_v=\sum_\alpha v_\alpha\xi_\alpha$ and
$\phi_v=\sum_\alpha v_\alpha\phi_\alpha$.
\end{proposition}
\begin{proof}
For a vertical field $Z$, ambient weighted integration
by parts disintegrates to fibre integration by parts.
Thus $\operatorname{div}_\nu Z$ equals
$\operatorname{div}_{\nu_g}Z$ on the fibre, with its
actual coarea density. Since
$\operatorname{div}_{\nu_g}\nabla_V\phi_\alpha=-\xi_\alpha$,
the corrected divergence is the conditional mean of
$S^0_\alpha$. V25's test-function proof, applied to the
smooth lift $B^0$, identifies that mean with
$X_\alpha\log\rho$. The correction is vertical, so every
coarse output remains fixed as prescribed.

Pointwise $B^0$ is orthogonal to the vertical tangent
space. Its cross terms with all the corrections vanish.
The weak Poisson equation gives
$\int\langle\nabla_V\phi_\alpha,\nabla_V\phi_\delta\rangle
d\nu_g=\mathsf M_{\alpha\delta}$, proving
\eqref{c18v26:transport}.

Any other correction in direction $v$ differs from
$\nabla_V\phi_v$ by a vertical field $Z$ with zero
$\nu_g$-divergence. Integration by parts gives
$\int\langle\nabla_V\phi_v,Z\rangle d\nu_g=0$.
Its squared norm is therefore the displayed minimum
plus $\|Z\|_2^2$. Finally Cauchy--Schwarz in the
Dirichlet inner product proves
\eqref{c18v26:variational}, with equality at $h=\phi_v$
unless $\xi_v=0$, in which case both sides vanish.
Real fields suffice for the construction; complexification
uses the Hermitian forms.
\end{proof}

This is the full-map counterpart of the inherited
source-specific Poisson metric, not a new assertion
that the Poisson cost is small. The genuine new bound is
\[
 P_J\mathsf W(g)P_J
       \preceq\frac{128\ell}{m^2}P_J+P_J\mathsf M(g)P_J .
 \tag{V26.11}\label{c18v26:remaining-matrix}
\]
It isolates what is still to be estimated. Substituting
$\|\mathcal L_g^{-1}\|$ times a raw score variance for
$\mathsf M$ is allowed only with a separately justified
conditional gap and its actual constants.

\subsection{How the cost enters energy stability}

For smooth fine $F$, the corrected divergence identity gives
\[
 \nabla\Pi F=\mathbb E_\nu[\widetilde B^*\nabla F\mid c=g],
 \qquad \Pi F=\mathbb E_\nu[F\mid c=g].
 \tag{V26.12}\label{c18v26:corrected-derivative}
\]
Indeed V25's conditional covariance term vanishes because
the corrected score is coarse measurable. Equivalently,
integrate $F\xi_\alpha$ by parts using
$\mathcal L_g\phi_\alpha=\xi_\alpha$.
The covariance has become vertical derivative work; it
has not vanished at zero cost.

At fixed $g$, the map
$u\mapsto\overline G(g)^{1/2}
            \mathbb E_{\nu_g}[\widetilde B^*u]$
from fibre $L^2$ vector fields to the coarse tangent has
its product with its adjoint equal to
$\overline G^{1/2}\mathsf W\overline G^{1/2}$.
Consequently the \emph{additional} hypothesis
$\mathsf W(g)\preceq C\,\overline G(g)^{-1}$ almost everywhere
would imply
\[
 \mathcal E((\Pi F)\circ c,(\Pi F)\circ c)
 \le C\,\mathcal E(F,F),\qquad
 \mathcal E(F,F)=\kappa\int|\nabla F|^2d\nu .
 \tag{V26.13}\label{c18v26:stability}
\]
This follows by integrating that operator bound and using
V17's exact coarse form. It is not asserted as a populated
native bound: \eqref{c18v26:remaining-matrix} controls only
one summand, on a selected direction space. Exceptional
directions and the correlation with the actual coarse
gradient still matter.

Even \eqref{c18v26:stability} would be projection energy
stability, not a fine spectral gap. Conditional variance,
the coarse marginal and the physical generator's return
remain separate questions. In particular making a
conditional derivative covariance zero by changing lift
does not make $QKJ$ zero; V17 already exhibits that distinction
for the nontrivial Gaussian flow-sampling reference.

\subsection{A calibration of the correct target, not a new YM model}

The following elementary Gaussian calculation tests the
difference between score size and its Poisson cost.
On $\mathbb R^2$ let
\[
 d\nu(x,y)\ \propto\
       e^{-x^2/2-r(y-a_0x)^2/2}\,dx\,dy,\qquad
 c(x,y)=x,\qquad r>0.
\]
Here $G=\mathsf H=1$, $B^0=\partial_x$, and with
$z=y-a_0x$,
\[
 \xi=ra_0z,\qquad
 \mathcal L_x=-\partial_z^2+rz\partial_z,\qquad
 \phi=a_0z,\qquad
 \mathbb E[\xi^2\mid x]=ra_0^2,\qquad
 \mathsf M=a_0^2,\quad \mathsf W=1+a_0^2 .
 \tag{V26.14}\label{c18v26:calibration}
\]
To verify this, the coarse density is proportional to
$e^{-x^2/2}$, while $\partial_x\log\nu=-x+ra_0z$.
The conditional variance of $z$ is $1/r$, and
$\mathcal L_x(a_0z)=ra_0z$. The correction is
$a_0\partial_y$, so its squared norm is $a_0^2$ and
the corrected score is exactly $-x$.
For fixed nonzero $a_0$, raw conditional score variance
diverges as $r\to\infty$ although the transport cost stays
fixed. Conversely $\mathsf H=1$ does not bound
$\mathsf M$ as $|a_0|\to\infty$.
This is an algebraic calibration of
\eqref{c18v26:transport}, not a compact YM vacuum or a
comparison theorem for it.

\subsection{Next upstream task and verification}

The next native target is the matrix
$\langle\xi_\alpha,\mathcal L_g^{-1}\xi_\delta\rangle$
for the minimal full-map lift, or an explicit vertical
divergence correction giving an upper bound for it.
A useful construction must control the complete direction
bank and its source weights. Another mean Fisher bound,
a gauge-charged gap, or a Gaussian calibration alone would
not pay this physical shape cost.

The unweighted rare-inverse cost is now included in a
conditional whole-fibre lift norm. The gradient-weighted
exceptional direction cost, native conditional variance,
coarse marginal variance, relative-energy localization and
physical scale transport remain open. The interacting
four-dimensional continuum construction and positive
physical YM mass gap remain unproved.

Exact regressions check noncommuting positive block
splittings, conditional averaging with rare states,
whole-output lift identities, the minimum vertical
correction, Gaussian score/Poisson normalization and
the displayed scale constants. These fixtures are
not the unknown native ground law or proof-assistant
certification. V25's entire body is preserved except
its title, and all other 47 note sources are unchanged.
Only \texttt{.tex} files remain in \texttt{latex\_notes};
compilation and inspection products are elsewhere.
The next companion and broad literature checkpoints
remain V30, and the archive restart remains V100.
% END CYCLE 18 VERSION 26 APPEND
% BEGIN CYCLE 18 VERSION 27 APPEND
\clearpage
\section{V27: A uniform full-map Gaussian transport bound}
\label{c18v27:context}

V26 isolates the actual source-specific Poisson matrix, but
does not bound it at the required native scales. Before
trying to transfer another small-score estimate, we test
its proposed matrix target on the \emph{entire specified}
V17 flow-sampling reference. The result is stronger than a
finite-dimensional calibration: the corrected transport
cost and conditional-projection energy are uniformly
bounded in both volume and block length. Coarse zero
momentum is included, despite the reference's removal of
fine harmonic fields.

This is reference progress, not a bound for the interacting
ground law. The normalized Poisson interpretation was
already used in f17 V94--V95 and V26; f8 V18 already gives
the general Gaussian harmonic lift and covariance splitting.
The new calculation is the relative-heat alias estimate for
V17's three-dimensional, arbitrary-length flowed sampler
with Hamiltonian precision proportional to
$\sqrt{d_1^*d_1}$. It pays the full bank, not just one
retained mode or one Wick degree. We then state the exact
native residual certificate that a proposed nonlinear
replacement must satisfy.

The primary definition in Li--Zhao,
\emph{Wasserstein information matrix},
\href{https://arxiv.org/html/1910.11248v5}{arXiv:1910.11248v5},
Section 2.2, again confirms the weighted-Poisson
interpretation. It supplies no YM uniformity estimate.
The archive audit also found the coupling-independent
one-link gauge transport theorem in f17 V95. Its horizontal
physical complement was explicitly unpaid there; it is
not recounted as a solution of the present task.

\subsection{The Gaussian lift, with its actual derivative cost}

Keep precisely V17's real fine transverse, zero-mean
space $\mathcal E_N$, $N=\ell n$, $\ell\ge2$, $n\ge5$,
$\mathsf L>0$, and
\[
 R=\sqrt{2b/\kappa}\sqrt{\mathsf L},\qquad
 C=C_\ell=\mathsf P_cB_\ell e^{-\ell^2\mathsf L},
 \qquad d\nu_{\rm G}\propto e^{-\langle A,RA\rangle/2}\,dA .
 \tag{V27.1}\label{c18v27:reference}
\]
Let $\mathcal Y$ be the coarse transverse space, including
its three link-direction components per color at $q=0$.
The proof below also verifies that $C$ is onto this space.
All operators and adjoints here use the original Euclidean
field metrics and unitary Fourier normalization. Define
\[
 \begin{gathered}
 G=CC^*,\qquad
 \mathsf A=CR^{-1}C^*,\qquad
 T=R^{-1}C^*\mathsf A^{-1},\qquad B^0=C^*G^{-1},\\
 \mathsf W_{\rm G}=T^*T,\qquad
 \mathsf M_{\rm G}=(T-B^0)^*(T-B^0).
 \end{gathered}
 \tag{V27.2}\label{c18v27:lifts}
\]
The matrix $\mathsf A$ is a coarse covariance, not a
fine field. Subscript $\mathrm G$ denotes the reference;
these are not V26's native random matrices.

\begin{lemma}[Identifying the score correction in the inherited Gaussian lift]
\label{c18v27:identify}
The conditional mean is $\mathbb E[A\mid CA=y]=Ty$.
The corrected lift of V26 for this linear Gaussian model
is exactly $T$, and
\[
 CT=I,\qquad
 \mathsf W_{\rm G}=G^{-1}+\mathsf M_{\rm G}
 =\mathsf A^{-1}CR^{-2}C^*\mathsf A^{-1}.
 \tag{V27.3}\label{c18v27:identity}
\]
In particular $T$, $\mathsf W_{\rm G}$ and
$\mathsf M_{\rm G}$ are unchanged when $R$ is multiplied
by any positive scalar.
\end{lemma}
\begin{proof}
Completing the square gives the conditional mean and
the fixed centered conditional covariance
$R^{-1}-T\mathsf A T^*$. The identity $CT=I$ follows by
multiplication. Put $\mathcal V=\ker C$ and $Z=T-B^0$.
Then $Z$ is vertical and $B^0$ is orthogonal to
$\mathcal V$, so $T^*T=G^{-1}+Z^*Z$.

For completeness identify the conditional score, rather
than merely its mean. Write $z=A-Ty\in\mathcal V$.
For a constant coarse vector $v$ the centered weighted
divergence of $B^0v$ is
$\xi_v=-\langle RB^0v,z\rangle$.
The conditional generator without $\kappa$ is
$-\Delta_{\mathcal V}+
(P_{\mathcal V}Rz)\cdot\nabla_{\mathcal V}$.
Since $RT=C^*\mathsf A^{-1}$,
$P_{\mathcal V}RZv=-P_{\mathcal V}RB^0v$.
Thus $\phi_v(z)=\langle Zv,z\rangle$ has mean zero and
solves $\mathcal L_y\phi_v=\xi_v$.
Its vertical gradient is $Zv$. This proves the
score-correction and cost identifications, including
their signs. Multiplying the formula for $T$ out
proves the last equality in \eqref{c18v27:identity};
a scalar multiplying $R$ cancels between its two factors.
\end{proof}

This scalar cancellation is specific and useful:
the raw conditional score covariance generally scales
with $\sqrt{2b/\kappa}$, whereas its Poisson transport
cost does not. No small-score hypothesis is imposed.

\subsection{A relative slow-alias certificate}

We first give a short matrix observation. Suppose in one
momentum block $R=\bigoplus_j r_jI$, $r_j\ge r_*>0$,
and $C=(C_j)_j$. Let
\[
 G=\sum_j C_jC_j^*,\qquad
 D=\sum_{j\in I_*}C_jC_j^*>0,\qquad
 r_j=r_*\quad(j\in I_*).
\]
If $G\preceq\Gamma D$, then
\[
 \mathsf W_{\rm G}\preceq D^{-1}
                         \preceq\Gamma G^{-1}.
 \tag{V27.4}\label{c18v27:slow-certificate}
\]
Indeed
$CR^{-2}C^*\preceq r_*^{-1}\mathsf A$ and
$\mathsf A\succeq r_*^{-1}D$. Congruence by
$\mathsf A^{-1}$ and order reversal under inversion
give the two assertions. The matrices $C_jC_j^*$ need
not commute. Crucially, no global lower bound for $r_*$
is used: it cancels within its momentum block.

Use V17's aliases $k_m=(q+2\pi m)/\ell$ in the fine
Brillouin zone, with
$\lambda_m=4\sum_d\sin^2(k_{m,d}/2)$ and primary $m=0$.
If $q\ne0$, V17's constraint-preserving primary inverse gives
\[
 D=C_0C_0^*
 \succeq\frac4{\pi^2\ell}e^{-2\ell^2\lambda_0}I,
 \qquad
 G\preceq\frac1\ell e^{-2\ell^2\lambda_0}
       \sum_m e^{-2\ell^2(\lambda_m-\lambda_0)}I .
 \tag{V27.5}\label{c18v27:relative-grams}
\]
The primary frequency is minimal, so it is admissible
as $r_*=\sqrt{2b/\kappa}\sqrt{\lambda_0}$ in
\eqref{c18v27:slow-certificate}. The absolute primary
heat factor must be kept on both sides, not bounded
separately by a very small constant.

\begin{lemma}[Uniform relative heat sum]
\label{c18v27:heat}
For every $\ell\ge2$ and $q\in[-\pi,\pi)^3$,
\[
 \sum_{m\in\mathcal M_\ell(q)}
             e^{-2\ell^2(\lambda_m-\lambda_0)}
      <(3+1/1200)^3.
 \tag{V27.6}\label{c18v27:heat-sum}
\]
\end{lemma}
\begin{proof}
In one coordinate put $x=|q_d+2\pi m_d|$,
$y=|q_d|$. These obey $0\le y\le\pi$, $y\le x\le\pi\ell$.
Let $a=x/(2\ell)$ and $b=y/(2\ell)$.
Then $0\le a-b\le\pi/2$ and $0\le a+b\le3\pi/4$.
Concavity of sine, with the endpoints of these intervals,
gives
\[
 \sin(a-b)\ge(2/\pi)(a-b),\qquad
 \sin(a+b)\ge(2/(3\pi))(a+b).
\]
Using $\sin^2a-\sin^2b=\sin(a+b)\sin(a-b)$ and
$x\ge2\pi|m_d|-y$ yields
\[
 \begin{split}
 4\ell^2(\sin^2a-\sin^2b)
 &\ge\frac4{3\pi^2}(x^2-y^2)\\
 &\ge\frac{16}{3}|m_d|(|m_d|-1).
 \end{split}
 \tag{V27.7}\label{c18v27:relative-difference}
\]
For $m_d=0$ the final inequality is equality at zero;
for $|m_d|\ge1$ it uses $y\le\pi$.
Extend the finite alias index set to $\mathbb Z^3$ after
this nonnegative majorization. The one-coordinate sum is
at most
\[
 1+2\sum_{j\ge1}e^{-(32/3)j(j-1)}
 \le3+\frac2{e^{64/3}-1}<3+\frac1{1200}.
\]
Here $j(j-1)\ge2(j-1)$ for $j\ge2$, and V17's
elementary bound $e^8>2401$ suffices for the last strict
inequality. Taking the product proves
\eqref{c18v27:heat-sum}.
\end{proof}

Consequently the nonzero coarse-momentum blocks satisfy
\[
 \mathsf W_{\rm G}(q)\preceq\Gamma_{\ne0}G(q)^{-1},
 \qquad
 \Gamma_{\ne0}:=\frac{\pi^2}{4}(3+1/1200)^3<68.
 \tag{V27.8}\label{c18v27:nonzero}
\]
The last comparison follows from $\pi^2<10$ and
$\frac52(3+1/1200)^3<68$. This includes arbitrarily
small nonzero $q$ and Nyquist degeneracies.

\subsection{Paying coarse zero momentum rather than deleting it}

At $q=0$ the primary \emph{fine} harmonic input was removed
in V17. It cannot be used as the witness in
\eqref{c18v27:relative-grams}. Nor may one silently delete
the coarse zero output. Choose one allowed unit axial
alias on each of the three axes, using $m=-\widehat d$
when $\ell=2$ and $m=\widehat d$ otherwise. Their common
minimal nonzero frequency and scaled energy are
\[
 \lambda_* = 4\sin^2(\pi/\ell),\qquad
 \tau_\ell=\ell^2\lambda_*<4\pi^2<40.
\]
On the axial alias in direction $d$, the transverse plane
is $\widehat d^\perp$, the transverse chain factors are
$\ell$, and its Gram is
$\ell^{-1}e^{-2\tau_\ell}(I-\widehat d\widehat d^*)$.
Thus the three witnesses give exactly
\[
 D=\frac2\ell e^{-2\tau_\ell}I .
 \tag{V27.9}\label{c18v27:zero-witness}
\]
They prove surjectivity of the zero-momentum output;
the primary witnesses already prove surjectivity at
nonzero momentum.

We bound the remaining relative heat sum explicitly.
For a represented integer $j$ with $|j|\le\ell/2$,
\[
 4\ell^2\sin^2(\pi j/\ell)\ge16j^2.
\]
If $|j|\ge2$, subtraction of $\tau_\ell<40$ and
$16j^2-40\ge8(j^2-1)$ give
\[
 4\ell^2\sin^2(\pi j/\ell)-\tau_\ell
                   \ge8(j^2-1).
 \tag{V27.10}\label{c18v27:zero-difference}
\]
For $|j|=1$ the difference and the right side are both
zero. In a nonzero alias choose any coordinate with
$m_d\ne0$. Pay \eqref{c18v27:zero-difference} there,
and the absolute $16m_j^2$ bound in the other two
coordinates. Summing over all three possible choices
only overcounts positive terms. Hence
\[
 \begin{split}
 \sum_{m\ne0}e^{-2\ell^2(\lambda_m-\lambda_*)}
 &\le3\left(2\sum_{j\ge1}e^{-16(j^2-1)}\right)
          \left(1+2\sum_{j\ge1}e^{-32j^2}\right)^2\\
 &<3(2+1/1200)(1201/1200)^2 .
 \end{split}
 \tag{V27.11}\label{c18v27:zero-sum}
\]
In the first factor use $j^2-1\ge3(j-1)$ for $j\ge2$;
its tail is at most $(e^{48}-1)^{-1}<1/2400$.
The other factors use
$\sum_{j\ge1}e^{-32j^2}\le(e^{32}-1)^{-1}<1/2400$.
These estimates also cover $\ell=2$, when some integer
terms in the majorants are not actual aliases.

As in \eqref{c18v27:relative-grams},
$G(0)\preceq\ell^{-1}e^{-2\tau_\ell}$ times this heat
sum. Apply \eqref{c18v27:slow-certificate} with
\eqref{c18v27:zero-witness}. It gives
\[
 \mathsf W_{\rm G}(0)\preceq\Gamma_0G(0)^{-1},
 \qquad
 \Gamma_0:=\frac32(2+1/1200)(1201/1200)^2<4.
 \tag{V27.12}\label{c18v27:zero}
\]
The reference's zero coarse output is supplied by
nonzero fine aliases; it does not restore its missing
fine harmonic or compact toron variables.

\subsection{A uniform whole-function theorem for this reference}

\begin{theorem}[Full-bank transport and projection energy]
\label{c18v27:uniform}
On all of $\mathcal Y$, for every stated volume, block
length and positive $\kappa,b$,
\[
 G^{-1}\preceq\mathsf W_{\rm G}\preceq68G^{-1},
 \qquad 0\preceq\mathsf M_{\rm G}\preceq67G^{-1}.
 \tag{V27.13}\label{c18v27:matrix-bound}
\]
Let $\Pi_{\rm G}F=\mathbb E_{\nu_{\rm G}}[F\mid CA]$,
viewed as a function on $\mathcal Y$.
Every fine Gaussian form-domain function satisfies
\[
 \kappa\int
   |\nabla((\Pi_{\rm G}F)\circ C)|^2\,d\nu_{\rm G}
       \le68\,\kappa\int|\nabla F|^2\,d\nu_{\rm G}.
 \tag{V27.14}\label{c18v27:energy}
\]
The bounds hold for all colors simultaneously and after
restriction to global-color-invariant functions.
\end{theorem}
\begin{proof}
Fourier decomposition, \eqref{c18v27:nonzero} and
\eqref{c18v27:zero} prove the upper matrix bound.
The orthogonal splitting \eqref{c18v27:identity} proves
its lower bound and the assertion about $\mathsf M_{\rm G}$.
All multipliers respect the real Fourier condition;
color merely repeats the same block. There is no
link-count, color-count or inverse lowest-frequency loss.

For smooth $F$ integrate over the translated fixed
conditional Gaussian to obtain
\[
 \nabla_y\Pi_{\rm G}F=\mathbb E[T^*\nabla F\mid CA=y].
\]
Jensen's inequality in this vector formula and the
operator bound
$\|TG^{1/2}\|^2
=\|G^{1/2}\mathsf W_{\rm G}G^{1/2}\|\le68$
give the conditional energy inequality. Integrating it,
and using
$|\nabla(f\circ C)|^2=(\nabla f)^*G\nabla f$,
proves \eqref{c18v27:energy}.
Smooth approximation in the finite-dimensional Gaussian
Sobolev space, the $L^2$ contraction of conditional
expectation and this uniform form bound extend the
result to the full form domain. Global color rotations
commute with every map and measure in this proof,
so the invariant restriction follows.
\end{proof}

Unlike a covariance-norm isometry, this controls the
\emph{original Euclidean-gradient energy}. Unlike V17's
conditional variance floor, it also bounds the energy
of the retained conditional mean. Those are distinct
questions. Neither result proves an infrared gap:
V17's singlet eigenvalue
$4\sqrt{2\kappa b}\sin(\pi/N)$ still tends to zero.
The nonzero Gaussian return $Q_{\rm G}K_{\rm G}J_{\rm G}$
also remains; changing the differentiation lift does
not change that generator.

\subsection{The native residual certificate and the next target}

Return now to V26's actual compact map, ground law,
full fibres and coarea-weighted $\mathcal L_g$.
The following consequence of its minimum-lift theorem
makes precise what a proposed nonlinear replacement
for $T$ would have to pay. It is an exact reformulation
of the residual obligation, not its solution.

Let $A_\alpha(U)$ be any smooth full-output lift:
$Dc\,A_\alpha=X_\alpha$. Define
\[
 \begin{split}
 r_\alpha&=\operatorname{div}_\nu A_\alpha
                -\mathbb E_\nu[\operatorname{div}_\nu A_\alpha\mid c],\\
 \mathsf R^A_{\alpha\delta}(g)
       &=\langle r_\alpha,\mathcal L_g^{-1}r_\delta\rangle_{\nu_g},
 \qquad
 \mathsf N^A(g)=\mathbb E_\nu[A^*A\mid g].
 \end{split}
 \tag{V27.15}\label{c18v27:residual}
\]
Use precisely this lift's centered score. At each finite
regulator these quantities exist as in V26.

Let $\mathcal P_\nabla$ be the orthogonal projection
in vertical $L^2(\nu_g)$ onto fibre gradients. With
$Z=A-B^0$ and $Y=(I-\mathcal P_\nabla)Z$, one has
\[
 \begin{gathered}
 A+\nabla_V\mathcal L_g^{-1}r
                    =\widetilde B+Y,\\
 \mathbb E_\nu[
 (A+\nabla_V\mathcal L_g^{-1}r)^*
 (A+\nabla_V\mathcal L_g^{-1}r)\mid g]
                    =\mathsf W(g)+\mathbb E_\nu[Y^*Y\mid g].
 \end{gathered}
 \tag{V27.16}\label{c18v27:change-lift}
\]
To prove this, $Z$ is vertical and
$r=\xi+\operatorname{div}_{\nu_g}Z$. The weak Poisson
equation gives
$\nabla_V\mathcal L_g^{-1}\operatorname{div}_{\nu_g}Z
=-\mathcal P_\nabla Z$.
The leftover $Y$ is divergence free, so is orthogonal
to the gradient correction in $\widetilde B$; it is
also pointwise orthogonal to $B^0$. Expanding the Gram
therefore proves both equalities. This explicitly
accounts for a change of lift rather than transferring
score bounds between different lifts without a cost.

In particular, if for two \emph{proved} nonnegative
constants $a_*,d_*$, uniform in the required regulator
parameters, one had
\[
 \mathsf N^A(g)\preceq a_*^2\overline G(g)^{-1},
 \qquad
 \mathsf R^A(g)\preceq d_*^2\overline G(g)^{-1}
 \quad\hbox{for almost every }g,
 \tag{V27.17}\label{c18v27:native-hypotheses}
\]
then Minkowski's inequality for each vector in the
full coarse tangent and \eqref{c18v27:change-lift} yield
\[
 \mathsf W(g)\preceq(a_*+d_*)^2\overline G(g)^{-1}.
 \tag{V27.18}\label{c18v27:native-certificate}
\]
V26 would then give the corresponding native projection
energy bound. For the Gaussian lift $A=T$, the centered
residual is zero and the preceding theorem supplies
$a_*=\sqrt{68}$, $d_*=0$. No regulator-uniform native
choice is proved here. The V25 inverse-flow lift is a legitimate
$A$, but its raw score bounds do not bound
$\mathsf R^A$ without an additional source-specific
inverse estimate; its exponential norm bound is also
not assigned to the Gaussian $T$.

The next upstream task is therefore to construct and
test a full-output native trial lift with a paid
residual in the \emph{actual conditional negative
Sobolev norm}, or to bound V26's $\mathsf M$ directly.
A small plaquette-defect probability does not compare
the conditional logarithmic ground derivatives with
their Gaussian values. A locally defined trial field
would also need full-fibre extension and all cutoff
derivative costs. These are concrete missing
requirements, not assumed perturbative remainders.

\subsection{Verification and scope}

V26's exact preservation check passed at the start:
the previous goal turn was progress. The new exact
regressions check a correlated Gaussian precision,
the vertical Poisson sign, all-output identities,
the fourth-degree conditional derivative, coupling
cancellation, the noncommuting slow-bank order
argument, and the change-of-lift residual identity.
Rational arithmetic checks the constants
$68$ and $4$. High-precision scalar and floating
Fourier tests include very small momenta, Nyquist
aliases and every momentum on selected finite tori;
they are regressions, not the all-scale proof above.

The entire V26 body is preserved apart from its title,
and all 47 other note sources are unchanged.
Only \texttt{.tex} files remain in \texttt{latex\_notes};
all diagnostics and render products are outside it.
Both the next companion checkpoint and broad primary
literature sweep remain V30; the archive restart is
V100. Native projection energy stability, conditional
and coarse marginal variance, relative-energy
localization, physical scale transport, the interacting
four-dimensional continuum construction and the
positive physical YM mass gap remain unproved.
% END CYCLE 18 VERSION 27 APPEND
% BEGIN CYCLE 18 VERSION 28 APPEND
\clearpage
\section{V28: A native weighted lift and packed derivative bounds}
\label{c18v28:context}

V27's factor $68$ applies to its transverse, harmonic-free
Gaussian reference. We now return to the actual compact map
and construct a global candidate for its residual problem.
The candidate is smooth, gauge equivariant, and lifts
\emph{every} prescribed coarse-link output. Its whole-fibre
conditional squared norm is polynomially bounded on V26's
joint protected directions. We also pay a packed first
derivative bound for its trial weight, summing over all
fine-link and color derivatives without a volume factor.

This does not pay the full conditional score residual.
The weight is defined from native holonomies, not from the
unknown logarithmic ground Hessian. Its regularization
selects a lift and does not add a mass term to the Hamiltonian.
The original $H,\Omega,\nu$, coarse map and physical time
are unchanged.

The weighted-minimum construction is established linear
algebra, already used in the Gaussian archives. The matrix
loss below is of Kantorovich type; the primary record for
Marshall--Olkin, \emph{Matrix Versions of the Cauchy and
Kantorovich Inequalities} (1990), is
\href{https://statistics.stanford.edu/technical-reports/matrix-versions-cauchy-and-kantorovich-inequalities-0}{Stanford report OLK NSF 275}.
We give the short needed proof, without importing a
spectral theorem from that reference. The packed-derivative
method also has a precedent in f17 V9--V10 for a different
covariant difference and normalized heat filter. Here
the operator is the full plaquette-holonomy derivative and
the function is a regularized inverse square root.
The new assertions are its native constants and their
application to the full nonlinear coarse lift.

\subsection{A global positive trial weight, not the magnetic Hessian}

Use V26's $N=\ell n$, $\ell\ge2$, $n\ge5$,
$c=\mathfrak b_\ell\circ\Phi_{\ell^2}$ and actual law
$\nu=\Omega^2dU$. All link tangents are represented in
left-multiplication orthonormal frames:
$\dot U_e=v_eU_e$, $v_e\in\operatorname{Im}\mathbb H$.
Use one fixed basepoint and orientation for each spatial
plaquette $p$, with holonomy $P_p(U)$. Define
\[
 (\mathcal D_Uv)_p=(DP_p(U)[v])P_p(U)^{-1},\qquad
 \mathsf L_U=\mathcal D_U^*\mathcal D_U,\qquad
 Q_\ell(U)=(\mathsf L_U+\ell^{-2}I)^{1/2},\qquad
 S_\ell=Q_\ell^{-1}.
 \tag{V28.1}\label{c18v28:weight}
\]
This definition is smooth on the whole compact product.
For each plaquette the four-block formula is
\[
 (\mathcal D_Uv)_p=\sum_{e\in p}O_{pe}(U)v_e .
\]
The blocks are signed orthogonal $3\times3$ matrices. A forward
occurrence has block $\operatorname{Ad}_{V}$, where $V$
is its preceding product; a reversed occurrence has
block $-\operatorname{Ad}_{VU_e^{-1}}$.
The orientation signs are not dropped.

\begin{lemma}[Native spectrum and covariance of the trial weight]
\label{c18v28:spectrum}
On every configuration,
\[
 0\preceq\mathsf L_U\preceq16I,\qquad
 \ell^{-1}I\preceq Q_\ell(U)
                  \preceq\sqrt{16+\ell^{-2}}\,I .
 \tag{V28.2}\label{c18v28:spectral-bounds}
\]
All these maps are gauge covariant in their stated
fine and plaquette frames.
At a flat connection $\mathsf L_U$ is the transported
curl Gram; at the identity it is $d_1^*d_1$ on the
\emph{full} tangent space.
\end{lemma}
\begin{proof}
There are four distinct links in each plaquette and four
plaquettes incident to each spatial link. Thus
\[
 \|\mathcal D_Uv\|^2
 \le4\sum_p\sum_{e\in p}|v_e|^2=16\|v\|^2.
\]
The spectral bounds follow by positive functional calculus.
Under a fixed gauge transformation, $v_e$ transforms by
the adjoint rotation at the stored link's tail, whereas
$(DP_p[v])P_p^{-1}$ transforms by the adjoint rotation
at the plaquette basepoint. These are orthogonal maps.
Consequently the Gram transforms by orthogonal conjugacy,
as do its regularized square root and inverse.
At a flat connection the displayed holonomy differential
is precisely the covariant curl.
\end{proof}

Off the flat locus, $\mathsf L_U$ is not the magnetic
Hessian. The exact formula in f16.2 V69 contains the
scalar plaquette weight and signed commutator term,
which are absent from this positive Gram.
Even a gauge-orbit variation can change a noncentral
plaquette by conjugation, so $\mathcal D_U$ need not
annihilate gauge directions there. Gauge covariance
is not gauge-direction annihilation. No such
identification is used below.

\subsection{The complete native weighted lift}

Write $C(U)=Dc(U)$ in the same fine frames and coarse
left-multiplication frames, and set
\[
 \begin{gathered}
 G=CC^*,\qquad
 D_\ell=C S_\ell C^*,\qquad
 A_\ell=S_\ell C^*D_\ell^{-1},\qquad
 B^0=C^*G^{-1},\\
 k_\ell=\sqrt{1+16\ell^2},\qquad
 \mathfrak K_\ell=\frac{(k_\ell+1)^2}{4k_\ell}
                                      \le\ell+1 .
 \end{gathered}
 \tag{V28.3}\label{c18v28:lift-def}
\]
Here $D_\ell$ is a coarse matrix; it is distinct from
the plaquette derivative $\mathcal D_U$.
The inverse exists everywhere because $c$ is a
submersion and $S_\ell>0$.

\begin{theorem}[Full-output identity and its paid native cost]
\label{c18v28:lift}
The field $A_\ell$ is smooth and gauge equivariant, and
\[
 CA_\ell=I,\qquad
 A_\ell^*Q_\ell A_\ell=D_\ell^{-1}.
 \tag{V28.4}\label{c18v28:full-output}
\]
For each coarse vector it is the unique lift minimizing
the $Q_\ell$-weighted squared fine norm. In the original
unweighted metric,
\[
 \begin{gathered}
 G^{-1}\preceq A_\ell^*A_\ell
                 \preceq\mathfrak K_\ell G^{-1},\\
 (A_\ell-B^0)^*(A_\ell-B^0)
                 \preceq(\mathfrak K_\ell-1)G^{-1}.
 \end{gathered}
 \tag{V28.5}\label{c18v28:Kantorovich}
\]
No coarse output or fine harmonic mode has been removed.
\end{theorem}
\begin{proof}
Direct multiplication proves \eqref{c18v28:full-output}.
Also $Q_\ell A_\ell=C^*D_\ell^{-1}$, which is orthogonal
to $\ker C$. Weighted Pythagoras therefore proves the
minimization statement. Smooth positive functional
calculus and matrix inversion prove smoothness.
The covariance from the preceding lemma, together
with equivariance of $c$, proves equivariance.

We prove the norm loss for arbitrary
$0<q_-I\preceq Q\preceq q_+I$ and any surjective $C$.
Put $S=Q^{-1}$, $a=q_+^{-1}$, $b=q_-^{-1}$,
$D=CSC^*$, $A=SC^*D^{-1}$. Functional calculus gives
$S^2\preceq(a+b)S-abI$. Hence
\[
 A^*A\preceq(a+b)D^{-1}-abD^{-1}GD^{-1}.
\]
Congruence by $G^{1/2}$ and completion of the matrix
square, with $X=G^{1/2}D^{-1}G^{1/2}$, give
\[
 (a+b)X-abX^2
 =\frac{(a+b)^2}{4ab}I
       -ab\left(X-\frac{a+b}{2ab}I\right)^2
 \preceq\frac{(a+b)^2}{4ab}I .
 \tag{V28.6}\label{c18v28:matrix-square}
\]
Here no matrices were interchanged as though they
commuted. This proves the upper bound with factor
$(q_++q_-)^2/(4q_+q_-)$.
For \eqref{c18v28:spectral-bounds} that factor is
$\mathfrak K_\ell$.
Since $k_\ell\le4\ell+1$ and $k_\ell^{-1}\le1$,
$\mathfrak K_\ell=(k_\ell+k_\ell^{-1}+2)/4\le\ell+1$.

Finally $A_\ell-B^0$ is vertical and $B^0$ is
pointwise orthogonal to every vertical vector. Thus
$A_\ell^*A_\ell=G^{-1}+(A_\ell-B^0)^*(A_\ell-B^0)$,
which proves the lower and difference bounds.
\end{proof}

The loss in the general matrix argument is real:
$Q=\operatorname{diag}(9,1)$ and $C=(3,1)$ give
$(A^*A)(CC^*)=25/9=(9+1)^2/(4\cdot9)$.
This exact equality fixture is not a claim of
sharpness for the native holonomy family.

\subsection{Whole-fibre conditioning, including rare configurations}

Keep \emph{exactly} V26's
$\mathsf H(g)=\mathbb E[G^{-1}\mid g]$ and protected
direction set $J(g)$, whose threshold uses
$\pi_e(g)=\mathbb P(e\notin T_\nu(U)\mid g)$.
Nothing in \eqref{c18v28:Kantorovich} restricts the
fine configurations included in this expectation.
Consequently
\[
 \begin{split}
 \mathsf N^{A_\ell}(g)
    &:=\mathbb E[A_\ell^*A_\ell\mid g]
                         \preceq\mathfrak K_\ell\mathsf H(g),\\
 P_J\mathsf N^{A_\ell}(g)P_J
    &\preceq\frac{128\ell\mathfrak K_\ell}{m^2}P_J
     \preceq\frac{128\ell(\ell+1)}{m^2}P_J,\\
 P_J\mathbb E[(A_\ell-B^0)^*(A_\ell-B^0)\mid g]P_J
    &\preceq\frac{128\ell(\mathfrak K_\ell-1)}{m^2}P_J .
 \end{split}
 \tag{V28.7}\label{c18v28:conditional}
\]
These are direct conditional expectations of the
pointwise theorem followed by V26. The mean inverse
has not been exchanged with the inverse mean.
The expected lost fraction remains precisely the
V26 bound
\[
 E_c^{-1}\mathbb E_\mu|J^c|
 \le\min\{1,C_{26}\ell^{27}e^{288\ell^2}
                        \varepsilon_\beta\}.
 \tag{V28.8}\label{c18v28:exception}
\]
In particular it vanishes under V26's embedding and
$\ell^2\le c_0\log\beta$, $c_0<1/576$.
The new lift is defined on the entire fibre and fixes
all complementary outputs as well. The set $J$ is
not differentiated. These bounds are neither uniform
bounds on all directions nor a comparison with
$\overline G^{-1}$ on the entire coarse space.

\subsection{Packed native derivatives of the weight}

Let $X_\alpha$, $\alpha=(f,d)$, denote all fine-link
unit left-multiplication vector fields. Derivatives in
this subsection are derivatives of matrix coefficients
in the fixed invariant frames, not an implicit choice
of eigenvectors or a differentiated gauge slice.
Set
\[
 \mathcal D_\alpha=X_\alpha\mathcal D,\qquad
 L_\alpha=X_\alpha\mathsf L,\qquad
 Q_\alpha=X_\alpha Q_\ell,\qquad
 S_\alpha=X_\alpha S_\ell.
\]
The derivative index runs over the whole fine regulator.

\begin{theorem}[Dimension-free packed square-root derivatives]
\label{c18v28:packed}
Pointwise on every configuration,
\[
 \begin{gathered}
 \sum_\alpha\mathcal D_\alpha^*\mathcal D_\alpha
                        \preceq768I,\qquad
 \sum_\alpha\mathcal D_\alpha\mathcal D_\alpha^*
                        \preceq768I,\\
 \sum_\alpha L_\alpha^2\preceq49152I,\\
 \sum_\alpha Q_\alpha^2\preceq12288\ell^2I,\qquad
 \sum_\alpha S_\alpha^2\preceq12288\ell^6I .
 \end{gathered}
 \tag{V28.9}\label{c18v28:packed-bounds}
\]
The constants do not depend on the number of links.
\end{theorem}
\begin{proof}
Each block $O_{pe}$ depends only on the four distinct
links of $p$, each appearing at most once in its
prefix word. When $\alpha=(f,d)$ with $f\in p$,
differentiation of its adjoint rotation inserts a
unit generator conjugated by an isometry. The
quaternion bracket has norm
$|[u,v]|\le2|u||v|$, so
$\|X_\alpha O_{pe}\|\le2$; if $f\notin p$ the
derivative is zero. There are at most twelve such
fine-link/color derivatives for a given plaquette.

For a fine vector $v$, four-term Cauchy--Schwarz gives
\[
 \begin{split}
 \sum_\alpha\|\mathcal D_\alpha v\|^2
 &\le4\sum_p\sum_{e\in p}\sum_{\alpha:\,f\in p}
                |(X_\alpha O_{pe})v_e|^2\\
 &\le4\cdot12\cdot4
              \sum_p\sum_{e\in p}|v_e|^2
  =768\|v\|^2 .
 \end{split}
\]
For a plaquette vector $w$,
$(\mathcal D_\alpha^*w)_e
=\sum_{p\ni e}(X_\alpha O_{pe})^*w_p$.
There are four incident terms; summing the same
block estimate over their twelve derivatives and
over the four edges of each plaquette gives
$\sum_\alpha\|\mathcal D_\alpha^*w\|^2\le768\|w\|^2$.
These prove the first line.

Since
$L_\alpha=\mathcal D_\alpha^*\mathcal D+
\mathcal D^*\mathcal D_\alpha$, the elementary
two-term norm inequality and $\|\mathcal D\|^2\le16$
give
\[
 \begin{split}
 \sum_\alpha\|L_\alpha v\|^2
 &\le2\cdot768\|\mathcal Dv\|^2
        +2\|\mathcal D\|^2\,768\|v\|^2\\
 &\le49152\|v\|^2 .
 \end{split}
 \tag{V28.10}\label{c18v28:packed-L}
\]
All $L_\alpha$ are self-adjoint, so this is the
claimed matrix-square bound.

Differentiate $Q_\ell^2=\mathsf L+\ell^{-2}I$.
The positive Sylvester equation and its convergent
solution are
\[
 Q_\ell Q_\alpha+Q_\alpha Q_\ell=L_\alpha,\qquad
 Q_\alpha=\int_0^\infty
                 e^{-tQ_\ell}L_\alpha e^{-tQ_\ell}\,dt .
 \tag{V28.11}\label{c18v28:Sylvester}
\]
To retain the derivative packing, apply the integral
triangle inequality in the direct sum over $\alpha$:
\[
 \begin{split}
 \left(\sum_\alpha\|Q_\alpha v\|^2\right)^{1/2}
 &\le\int_0^\infty e^{-t/\ell}\sqrt{49152}
                              \|e^{-tQ_\ell}v\|\,dt\\
 &\le\frac{\ell\sqrt{49152}}2\|v\|.
 \end{split}
\]
This proves the $Q_\alpha$ bound without any
derivative-count factor, including at repeated
eigenvalues. Finally $S_\alpha=-S_\ell Q_\alpha S_\ell$
and $\|S_\ell\|\le\ell$, so
$\sum_\alpha\|S_\alpha v\|^2
\le12288\ell^6\|v\|^2$.
\end{proof}

The summation is genuinely over all fine variables.
It is not a per-entry estimate multiplied afterwards
by the dimension. It still does not bound the trace
contraction defining a weighted divergence, nor its
conditional Poisson inverse.

\subsection{Differentiating the trial lift without losing terms}

The matrix derivative is useful for the next source
calculation. For any one coefficient derivative
$\partial$ in these frames, differentiate
\eqref{c18v28:lift-def} to obtain exactly
\[
 \begin{split}
 \partial A_\ell={}&
 -(I-A_\ell C)S_\ell(\partial Q_\ell)A_\ell\\
 &+(I-A_\ell C)S_\ell(\partial C)^*D_\ell^{-1}
                   -A_\ell(\partial C)A_\ell .
 \end{split}
 \tag{V28.12}\label{c18v28:derivative}
\]
This follows by the product rule and
$\partial D_\ell=(\partial C)S_\ell C^*
+C(\partial S_\ell)C^*+CS_\ell(\partial C)^*$,
with $\partial S_\ell=-S_\ell(\partial Q_\ell)S_\ell$.
In particular both differentiated-$C$ terms remain;
$C$ is the actual full flow-sampling differential,
not its flat Fourier symbol.

One part of \eqref{c18v28:derivative} now has an
explicit native packed bound. Put
$F_\alpha=-(I-A_\ell C)S_\ell Q_\alpha A_\ell$.
Then
\[
 \sum_\alpha\|F_\alpha v\|^2
 \le12288\,\mathfrak K_\ell^2\ell^4\,
                                      v^*G^{-1}v .
 \tag{V28.13}\label{c18v28:precision-part}
\]
Indeed $A_\ell C$ is a projection with kernel
$\ker C$. Relative to
$\operatorname{Ran}C^*\oplus\ker C$ it has block
form $\left(\begin{smallmatrix}I&0\\ Z&0\end{smallmatrix}\right)$.
The two complementary projection norms are therefore
both $\sqrt{1+\|Z\|^2}$.
From \eqref{c18v28:Kantorovich},
$\|A_\ell C\|^2\le\mathfrak K_\ell$.
Combine this with $\|S_\ell\|\le\ell$, the packed
$Q_\alpha$ bound and again
\eqref{c18v28:Kantorovich} to prove
\eqref{c18v28:precision-part}.
For a fixed vector supported on $J(g)$, conditional
integration gives the bound
$12288\,\mathfrak K_\ell^2\ell^4
(128\ell/m^2)|v|^2$.
This is the weight-variation part only, not the
whole derivative of $A_\ell$.

\subsection{The actual remaining score and the transfer distinction}

In invariant fine frames every basis vector has
zero Haar divergence. For a coarse frame index
$\gamma$ the actual weighted score of this lift is
\[
 S^{A_\ell}_\gamma
   =\sum_\alpha X_\alpha(A_{\ell,\alpha\gamma})
          +2\sum_\alpha A_{\ell,\alpha\gamma}
                                      X_\alpha\log\Omega .
 \tag{V28.14}\label{c18v28:score}
\]
The first sum requires the \emph{full}
\eqref{c18v28:derivative}; the second uses the actual
positive ground. With
$r_\gamma=S^{A_\ell}_\gamma
-\mathbb E[S^{A_\ell}_\gamma\mid c]$ the remaining
matrix is exactly
\[
 \mathsf R^{A_\ell}_{\gamma\delta}(g)
      =\langle r_\gamma,\mathcal L_g^{-1}
                                      r_\delta\rangle_{\nu_g}.
 \tag{V28.15}\label{c18v28:unpaid}
\]
V27's residual certificate applies to this concrete
candidate. Neither the packed weight derivative
nor \eqref{c18v28:conditional} bounds
\eqref{c18v28:unpaid}. The two lifts in V25 and V28
are different; their score estimates are not
interchanged. The actual coarea density and full
conditioning remain in \eqref{c18v28:unpaid}.

There is also a precise restriction on comparison
with V27. At the identity the weight in
\eqref{c18v28:weight} is
$(d_1^*d_1+\ell^{-2}I)^{1/2}$ on the full fine
tangent space. It retains a positive trial weight
on gauge and harmonic directions. V27 instead
works on the fine transverse harmonic-free space,
uses a square-root precision without this
regularization, and projects the coarse output.
The native $C$ in \eqref{c18v28:lift-def} makes
no such projection. Thus the two lift matrices
cannot simply be identified, even before an
interacting comparison is attempted.
The earlier full-link versus quotient distinction
in f11 V48 was rechecked for this reason.

The next upstream task is to estimate
\eqref{c18v28:unpaid} for this globally defined
candidate, retaining the differentiated full map
and ground-state terms, or to produce a different
candidate with a demonstrably smaller residual.
The $\ell^{-2}I$ in a trial weight is not a
Hamiltonian gap, and no physical mass has been
inserted by this construction.

\subsection{Verification and continuation record}

The previous turn was progress: V27's preserved
source and all 47 companion/archive sources were
revalidated before editing. Exact quaternion
regressions check forward and reversed holonomy
derivatives, gauge covariance and the local packed
inequalities. Independent finite-torus incidence
checks verify the four-by-four counting.
Exact rational matrix tests check the weighted
lift, its endpoint equality example, the
noncommuting Sylvester derivative and all terms
of \eqref{c18v28:derivative}. Floating square-root
covariance and high-precision scale tests are
regressions, not simulations of the native
vacuum or proof-assistant certification.

V27's body is preserved except for its title.
Only the active V28 source replaces it; all
other note sources are unchanged, and
\texttt{latex\_notes} contains only \texttt{.tex}
files. Build and QA products remain outside it.
The next companion and broad literature
checkpoints remain V30, with archive restart at
V100. Native score-residual control, full
projection energy stability, conditional and
coarse marginal variance, the physical return,
relative-energy localization, the interacting
four-dimensional continuum construction and the
positive physical YM mass gap remain unproved.
% END CYCLE 18 VERSION 28 APPEND
% BEGIN CYCLE 18 VERSION 29 APPEND
\clearpage
\section{V29: Packed derivatives of the complete nonlinear lift}
\label{c18v29:context}

V28 constructs a native full-output weighted lift but bounds only
the weight-variation part of its derivative. We now pay both
differentiated-map terms. The calculation uses the actual nonlinear
Wilson flow and the actual disjoint-chain sampling, on every
configuration. All fine derivative indices are summed, with no
factor proportional to the fine volume. Moving invariant frames
are included explicitly.

The result is a geometric derivative estimate, not a uniform
conditional score estimate. Its worst-case constants grow
exponentially with flow time. We also identify precisely why this
packed estimate alone cannot be contracted into a
volume-independent divergence bound. Thus the next step moves
upstream to the centered score in its conditional negative
Sobolev norm, rather than treating one more geometric bound
as a spectral gap.

The first flow differential and its block incidence bounds are
inherited from f16.2 V69/V74 and current V17/V25; they are not
new claims. Archive f14 V66 also bounds third derivatives of a
chart action as trilinear forms. That is not the two-sided packed
coefficient estimate needed here. The new step is to retain
derivatives in both input and output slots through the nonlinear
composition, and then through the full weighted inverse.
For background, the primary record for M.~L\"uscher,
\emph{Properties and uses of the Wilson flow in lattice QCD},
\href{https://arxiv.org/abs/1006.4518}{arXiv:1006.4518},
describes the Wilson-flow framework. No continuum construction,
conditional coercivity theorem, or constants below are imported
from that record; the estimates are proved directly in our
unit-round spatial normalization.

\subsection{The invariant-frame equation and its local derivative tensor}

Retain $N=\ell n$, $\ell\ge2$, $n\ge5$ and the conventions of V28.
A fine index $i=(e,a)$ labels a stored link and one of the three
unit imaginary quaternions $\mathbf e_a$. The invariant field
$X_i$ has tangent $\mathbf e_aU_e$. Write
\[
 E_{\rm sp}=\sum_p(1-w_p),\qquad
 F_i=-X_iE_{\rm sp},\qquad
 \dot W_e=F_e(W)W_e,\qquad W_t=\Phi_t(U).
 \tag{V29.1}\label{c18v29:flow}
\]
Here $F_e\in\operatorname{Im}\mathbb H$ collects the three
coefficients on link $e$. Put
\[
 (\mathsf M(U)v)_e
   =\sum_jv_jX_jF_e(U)+[F_e(U),v_e],\qquad
 T_\alpha(U)=X_\alpha\mathsf M(U).
 \tag{V29.2}\label{c18v29:generator}
\]
These are ordinary coefficient derivatives in the stated frames.

\begin{lemma}[Moving-frame generator and packed local source]
\label{c18v29:local}
The differential $\mathsf J_t(U)=D\Phi_t(U)$, expressed in
invariant frames at both endpoints, satisfies
\[
 \dot{\mathsf J}_t=\mathsf M(W_t)\mathsf J_t,\qquad
 \mathsf J_0=I.
 \tag{V29.3}\label{c18v29:Jacobi}
\]
The symmetric part of $\mathsf M$ is $-\operatorname{Hess}E_{\rm sp}$
in these frames. Consequently the propagator between times
$r\le t$ and its inverse have norm at most $e^{16(t-r)}$.
On every configuration,
\[
 \sum_\alpha T_\alpha^*T_\alpha\preceq b_*^2I,\qquad
 \sum_\alpha T_\alpha T_\alpha^*\preceq b_*^2I,\qquad
 b_*=288\sqrt3,\quad b_*^2=248832.
 \tag{V29.4}\label{c18v29:packed-generator}
\]
\end{lemma}
\begin{proof}
For a variation of a trajectory set $v_e=\delta W_eW_e^{-1}$.
Differentiating this expression gives
\[
 \dot v_e=dF_e(W)[v]+F_ev_e-v_eF_e .
\]
This proves \eqref{c18v29:Jacobi}, including the commutator.
Alternatively the Riemannian variational equation uses
$-\operatorname{Hess}E_{\rm sp}$ in parallel frames; the change
to orthonormal invariant frames adds a skew operator only.
Directly, for every constant coefficient vector $v$,
\[
 \langle v,\mathsf Mv\rangle
   =-\sum_{i,j}v_iv_jX_jX_iE_{\rm sp}
   =-\operatorname{Hess}E_{\rm sp}[v,v].
\]
The last equality holds because the integral curve of this
constant invariant vector on each link is a product-metric
geodesic. The commutator contributes zero to this quadratic
form. The exact plaquette Hessian in f16.2 V69 has each of
its four-by-four link blocks bounded by one. Four plaquettes
per link therefore give
$\|\operatorname{Hess}E_{\rm sp}\|\le16$.
Gronwall in both time directions proves the propagator bounds.

For completeness, a scalar coefficient of any ordered first,
second or third invariant derivative of $w_p$ has absolute
value at most one. The plaquette word contains four distinct
stored links. An invariant derivative replaces its unique
forward occurrence $U_e$ by $\mathbf e_aU_e$, or its unique
reverse occurrence $U_e^{-1}$ by $-U_e^{-1}\mathbf e_a$.
Repeated derivatives insert further unit quaternions in the
correct order; they do not split into a number of terms
growing with the order. The resulting word has norm one,
so its scalar part has absolute value at most one.

Decompose $\mathsf M=\sum_p\mathsf M^p$ using the corresponding
single-plaquette forces. All three indices of
$T^p_{\alpha,ij}=X_\alpha\mathsf M^p_{ij}$ must lie in
the twelve link/color indices $I_p$ of $p$. In the derivative
of the force-gradient term, every entry has magnitude at
most one. For a same-link commutator entry, the structure
constants of $\operatorname{Im}\mathbb H$ are
$2\epsilon_{abc}$; at most one color coefficient contributes
to a fixed matrix entry. Its derivative has magnitude at
most two. Therefore
\[
 |T^p_{\alpha,ij}|\le3\,
      \mathbf1_{\{\alpha,i,j\in I_p\}}.
 \tag{V29.5}\label{c18v29:entry}
\]
No symmetry between these ordered derivative indices is assumed.

For a vector $v$, each output index belongs to four plaquettes,
and each $I_p$ has twelve entries. Two finite
Cauchy--Schwarz inequalities give
\[
 \begin{split}
 \sum_{\alpha,i}\left|\sum_{p:\,i\in I_p}
                      \sum_{j\in I_p}T^p_{\alpha,ij}v_j\right|^2
 &\le4\cdot12\cdot9
       \sum_p\sum_{\alpha,i\in I_p}\sum_{j\in I_p}|v_j|^2\\
 &=4\cdot12\cdot9\cdot12^2\cdot4\,\|v\|^2
   =248832\,\|v\|^2.
 \end{split}
\]
Interchanging the input and output slots in this counting
proves the second inequality of \eqref{c18v29:packed-generator}.
The derivative index is already summed here; no subsequent
dimension factor is required.
\end{proof}

\subsection{Packed second flow variation, with both slots retained}

If a matrix family $T_\alpha$ obeys the two packed bounds
with constant $b^2$, then for any matrix $R$ on its derivative
index space, the family
$\widehat T_\alpha=\sum_\beta R_{\beta\alpha}T_\beta$
obeys the same two bounds with constant $b^2\|R\|^2$.
For each fixed test vector this is the Hilbert--Schmidt
inequality for composing the map
$z\mapsto\sum_\beta z_\beta T_\beta v$ with $R$.
Applying the same argument to $T_\beta^*w$ proves the other
bound. This elementary index-contraction fact will be used
without assuming that $R$ commutes with any $T_\beta$.

Define
\[
 \mathsf K_{\alpha,t}=X_\alpha\mathsf J_t,\qquad
 h(t)=\frac{b_*}{16}e^{16t}(e^{16t}-1)
      =18\sqrt3\,e^{16t}(e^{16t}-1).
 \tag{V29.6}\label{c18v29:h}
\]
The input basis in $\mathsf J_t$ is the invariant basis at
the varying initial point; hence its coefficient matrix at
$t=0$ is identically $I$, and $\mathsf K_{\alpha,0}=0$.

\begin{theorem}[Full nonlinear packed flow variation]
\label{c18v29:second-flow}
For all $t\ge0$ and every initial configuration,
\[
 \sum_\alpha\mathsf K_{\alpha,t}^*\mathsf K_{\alpha,t}
       \preceq h(t)^2I,\qquad
 \sum_\alpha\mathsf K_{\alpha,t}\mathsf K_{\alpha,t}^*
       \preceq h(t)^2I .
 \tag{V29.7}\label{c18v29:second-flow-bounds}
\]
\end{theorem}
\begin{proof}
Differentiating \eqref{c18v29:Jacobi} as a coefficient
identity gives
\[
 \dot{\mathsf K}_{\alpha,t}
 =\mathsf M(W_t)\mathsf K_{\alpha,t}
   +d\mathsf M(W_t)[\mathsf J_t e_\alpha]\mathsf J_t .
 \tag{V29.8}\label{c18v29:second-equation}
\]
The commutator and its derivative are already present in
$\mathsf M$ and $T_\beta$. No further frame derivative is
being set to zero.

Let $V(t,r)$ be the propagator for \eqref{c18v29:Jacobi}.
By variation of constants,
\[
 \mathsf K_{\alpha,t}
  =\int_0^t V(t,r)
       \left(\sum_\beta(\mathsf J_r)_{\beta\alpha}
                                      T_\beta(W_r)\right)
                       \mathsf J_r\,dr .
 \tag{V29.9}\label{c18v29:Duhamel}
\]
For fixed $v$, take the norm in the direct sum over $\alpha$.
The index-contraction fact, \eqref{c18v29:packed-generator}
and $\|\mathsf J_r\|\le e^{16r}$ give
\[
 \left(\sum_\alpha\|\mathsf K_{\alpha,t}v\|^2\right)^{1/2}
 \le b_*\int_0^t e^{16(t-r)}e^{32r}\,dr\,\|v\|
 =h(t)\|v\|.
\]
For the transposed integral, its leftmost factor is
$\mathsf J_r^*$ and its rightmost factor is $V(t,r)^*$.
Use the second packed bound and the same index contraction;
the scalar integral is identical. This proves both statements.
\end{proof}

\subsection{Actual disjoint-chain sampling and the full map}

Let $\mathsf B(W)=D\mathfrak b_\ell(W)$ in invariant frames.
The forward chain for each coarse link has $\ell$ distinct
fine links, and different coarse chains use disjoint fine
links. Thus $\mathsf B\mathsf B^*=\ell I$.
Write $\mathsf B_\beta=X_\beta\mathsf B$.

\begin{lemma}[Packed sampling coefficients]
\label{c18v29:sampling}
On every configuration,
\[
 \sum_\beta\mathsf B_\beta^*\mathsf B_\beta
                       \preceq12\ell^2 I,\qquad
 \sum_\beta\mathsf B_\beta\mathsf B_\beta^*
                       \preceq12\ell^2 I .
 \tag{V29.10}\label{c18v29:sampling-bounds}
\]
\end{lemma}
\begin{proof}
The block on the $j$th chain link is the adjoint rotation
of its prefix. A derivative is zero unless its fine link
lies earlier in the same chain. When nonzero its block
norm is at most two, by the quaternion bracket bound.
For a given coarse row, Cauchy--Schwarz over at most
$\ell$ input blocks, followed by the at most $3\ell$
derivative indices, bounds its packed squared norm by
$12\ell^2$ times the squared norm on that chain.
Disjointness proves the input-side statement after summing
the rows. For the output-side statement, each fine input
link occurs in at most one coarse row, so there are no
cross-row terms. For a fixed coarse vector the sum over
its at most $\ell$ blocks and $3\ell$ derivative indices
is again at most $12\ell^2$ times its squared norm.
\end{proof}

For now allow any $s\ge0$ and put
\[
 \begin{gathered}
 c_s=\mathfrak b_\ell\circ\Phi_s,\quad
 C_s=\mathsf B(W_s)\mathsf J_s,\quad C_{\alpha,s}=X_\alpha C_s,\\
 p_\ell(s)=2\sqrt3\,\ell e^{32s}+\sqrt\ell\,h(s),\qquad
 g_\ell(s)=\ell e^{-32s}.
 \end{gathered}
 \tag{V29.11}\label{c18v29:map-constants}
\]

\begin{theorem}[Two-sided packed differential of the complete map]
\label{c18v29:full-map}
On every configuration,
\[
 \begin{gathered}
 \sum_\alpha C_{\alpha,s}^*C_{\alpha,s}\preceq p_\ell(s)^2I,
 \qquad
 \sum_\alpha C_{\alpha,s}C_{\alpha,s}^*\preceq p_\ell(s)^2I,\\
 C_sC_s^*\succeq g_\ell(s)I.
 \end{gathered}
 \tag{V29.12}\label{c18v29:map-bounds}
\]
These estimates include every coarse output and every fine
derivative, including gauge and harmonic directions.
\end{theorem}
\begin{proof}
The exact coefficient chain rule is
\[
 C_{\alpha,s}
 =\left(\sum_\beta(\mathsf J_s)_{\beta\alpha}
                           \mathsf B_\beta(W_s)\right)\mathsf J_s
       +\mathsf B(W_s)\mathsf K_{\alpha,s}.
 \tag{V29.13}\label{c18v29:map-chain-rule}
\]
For either packed norm, the first summand costs at most
$2\sqrt3\,\ell\|\mathsf J_s\|^2$ and the second at most
$\|\mathsf B\|h(s)=\sqrt\ell\,h(s)$.
The direct-sum triangle inequality proves the first line
of \eqref{c18v29:map-bounds}.
Finally $\mathsf J_s\mathsf J_s^*\succeq e^{-32s}I$,
so multiplication by $\mathsf B$ and its adjoint gives
the lower Gram bound. This last estimate is the inherited
global floor, not a new protected polynomial floor.
\end{proof}

\subsection{Both missing terms in the weighted lift derivative}

Return to $s=\ell^2$, $C=C_s$ and all the definitions
$Q_\ell,S_\ell,D_\ell,A_\ell,G,k_\ell,\mathfrak K_\ell$
of \eqref{c18v28:lift-def}. Abbreviate $p=p_\ell(\ell^2)$,
$g_-=g_\ell(\ell^2)$, $k=k_\ell$ and $K=\mathfrak K_\ell$.
Set
\[
 \mathfrak T_\ell
  =64\sqrt3\,K\ell^2
      +(\sqrt K\,k+K)\,\frac{p}{\sqrt{g_-}}.
 \tag{V29.14}\label{c18v29:lift-constant}
\]

\begin{theorem}[Complete packed native lift derivative]
\label{c18v29:complete-lift}
For every coarse coefficient vector $v$ and every configuration,
\[
 \sum_\alpha\|(X_\alpha A_\ell)v\|^2
              \le\mathfrak T_\ell^2\,v^*G^{-1}v .
 \tag{V29.15}\label{c18v29:complete-bound}
\]
Consequently, for the whole-fibre law and exactly V26's
protected set $J(g)$,
\[
 P_J\,\mathbb E\left[
       \sum_\alpha(X_\alpha A_\ell)^*(X_\alpha A_\ell)
                    \,\middle|\,c=g\right]P_J
       \preceq\frac{128\ell\mathfrak T_\ell^2}{m^2}P_J.
 \tag{V29.16}\label{c18v29:conditional-bound}
\]
The conditional expectation includes rare configurations.
Neither the conditioning map nor the set $J$ is changed
or differentiated.
\end{theorem}
\begin{proof}
Write $P=I-A_\ell C$ and $Q_\alpha=X_\alpha Q_\ell$.
The exact three-term formula \eqref{c18v28:derivative} is
\[
 X_\alpha A_\ell
  =-P S_\ell Q_\alpha A_\ell
     +P S_\ell C_\alpha^*D_\ell^{-1}
     -A_\ell C_\alpha A_\ell .
 \tag{V29.17}\label{c18v29:three-terms}
\]
V28 proves $\|P\|^2\le K$, $\|S_\ell\|\le\ell$,
$A_\ell^*A_\ell\preceq KG^{-1}$, and a packed norm
at most $64\sqrt3\,K\ell^2(v^*G^{-1}v)^{1/2}$
for the first summand.

Put $q_+=\sqrt{16+\ell^{-2}}=k/\ell$.
Since $D_\ell\succeq q_+^{-1}G\succeq q_+^{-1}g_-I$,
\[
 D_\ell^{-2}
   \preceq\|D_\ell^{-1}\|D_\ell^{-1}
   \preceq\frac{q_+^2}{g_-}G^{-1}.
 \tag{V29.18}\label{c18v29:inverse-square}
\]
This does not square a noncommuting Loewner inequality.
Using the output-side packed bound for $C_\alpha$ gives
\[
 \sum_\alpha\|P S_\ell C_\alpha^*D_\ell^{-1}v\|^2
 \le K\ell^2p^2\|D_\ell^{-1}v\|^2
 \le \frac{Kk^2p^2}{g_-}\,v^*G^{-1}v .
 \tag{V29.19}\label{c18v29:second-term}
\]
Using the input-side bound and
$\|A_\ell\|^2\le K/g_-$ gives
\[
 \sum_\alpha\|A_\ell C_\alpha A_\ell v\|^2
 \le\frac K{g_-}p^2\|A_\ell v\|^2
 \le\frac{K^2p^2}{g_-}\,v^*G^{-1}v .
 \tag{V29.20}\label{c18v29:third-term}
\]
The triangle inequality in the direct sum over all $\alpha$
now proves \eqref{c18v29:complete-bound}.
Condition the pointwise matrix inequality over the full fibre
and use $P_J\mathbb E[G^{-1}\mid g]P_J
\preceq128\ell m^{-2}P_J$ from V26.
\end{proof}

All constants are independent of the spatial volume, but
$\mathfrak T_\ell$ is not polynomial in $\ell$. For example,
$p_\ell(s)$ contains $e^{32s}$ and $g_\ell(s)^{-1/2}$
contains $e^{16s}$. At $s=\ell^2$ their product contains
$e^{48\ell^2}$. Nothing in the theorem turns this
deterministic cost into a small score residual. It may be
useful for fixed-regulator differentiability, truncation
estimates or a later probability-weighted refinement; it is
not an all-scale coercivity estimate.

\subsection{Why the score problem has not been solved}

For a fixed coarse index $\gamma$ the required scalar score
is still exactly \eqref{c18v28:score}:
\[
 S^{A_\ell}_\gamma
  =\sum_\alpha (X_\alpha A_\ell)_{\alpha\gamma}
        +2\sum_\alpha A_{\ell,\alpha\gamma}X_\alpha\log\Omega.
 \tag{V29.21}\label{c18v29:true-score}
\]
The first sum is a diagonal trace, not the direct-sum norm
in \eqref{c18v29:complete-bound}. With $d$ fine frame
indices, generic Cauchy--Schwarz loses a factor $d$
in its squared bound. That loss cannot be removed merely
by calling a derivative estimate packed.

There is a compact-product check of this distinction.
On $L$ independent unit-quaternion links let
$f_e(U_e)$ be the first imaginary coordinate and let $V$
have coefficient $f_e(U_e)/\sqrt L$ in the first-color
frame on link $e$, with its other coefficients zero.
At the identity configuration,
\[
 V=0,\qquad
 \sum_\alpha\|X_\alpha V\|^2=1,\qquad
 \operatorname{div}_{dU}V=\sqrt L .
 \tag{V29.22}\label{c18v29:trace-check}
\]
Indeed $X_{e,a}f_e(I)=\delta_{a1}$, and every invariant
frame has zero Haar divergence. This is a counterexample
to a generic dimension-free trace contraction of the
packed bound, not a counterexample for the special
$A_\ell$, and not an obstruction theorem for YM.

Moreover the second term of \eqref{c18v29:true-score}
contains the actual logarithmic ground derivative. None
of the flow geometry replaces that derivative by
$Q_\ell$, by a Gaussian precision, or by a magnetic
Hessian. The remaining target stays
\[
 r_\gamma=S^{A_\ell}_\gamma
              -\mathbb E[S^{A_\ell}_\gamma\mid c],\qquad
 \mathsf R^{A_\ell}_{\gamma\delta}(g)
       =\langle r_\gamma,\mathcal L_g^{-1}r_\delta
                                      \rangle_{\nu_g}.
 \tag{V29.23}\label{c18v29:remaining-target}
\]
A useful next argument must exploit its actual centered
conditional structure, or directly represent it as a
controlled fibre divergence. A generic estimate of the
uncontracted derivative alone is now known to be
insufficient. V27's change-of-lift certificate remains
available, but its residual hypothesis is not proved
by this appendix. No artificial mass, modified law,
coarse projection or discarded fine mode is introduced.

\subsection{Verification and next checkpoint}

The predecessor V28 and the other 47 note sources were
revalidated before editing. The appended estimates are
supported by exact quaternion regressions of ordered
derivatives and the invariant-frame commutator sign,
two-sided local tensor bounds, exact chain-prefix
derivatives, and noncommuting inverse-square and lift
algebra. These are regression tests of the written
proof, not a simulation of the native ground state
or a proof-assistant certificate.

V28's body is preserved except for the title; only
the active V29 file replaces it. All archives and
companions are unchanged. The notes folder contains
only \texttt{.tex} sources; build and inspection files
remain outside it. The next iteration is V30, at which
both the companion review and the broad primary-literature
sweep are due. The last such checkpoints were V25 and
V20 respectively; the targeted background check above
does not replace the scheduled broad sweep.

The present turn makes progress on the full nonlinear
geometric derivative, while leaving the native
score-Poisson residual, full projection-energy stability,
conditional and coarse marginal variance, physical
return and scale transport unpaid. The interacting
four-dimensional continuum construction and positive
physical Yang--Mills mass gap remain unproved.
% END CYCLE 18 VERSION 29 APPEND
% BEGIN CYCLE 18 VERSION 30 APPEND
\clearpage
\section{V30: Native log-ground curvature and the complete score budget}
\label{c18v30:context}

V29 bounded the packed derivative of the complete weighted lift, but
a pointwise divergence estimate lost the number of fine links. We now
move upstream to the \emph{actual} positive ground state. A tensor
minimum argument gives a one-sided bound on its logarithmic Hessian.
Weighted integration by parts then controls the complete lift score
in mean square, without a volume factor. This is an unconditional
score estimate, not the conditional inverse-Poisson estimate needed
to close V27. No trial Gaussian replaces the native law.

\subsection{The V30 companion and primary-literature checkpoints}

The V30 five-version companion review rechecked the current source
inventory, hashes and terminal result/status sections. The NS V26,
SM--Einstein V72, shell-mixing V16 and parallel YM V5 sources are
unchanged since V25. The supplied ESM V64, source-selection V52,
native-stationarity V54 and exact-action V45 are also unchanged;
the newer ESM V69 was included again. Their useful discipline is
to retain the actual source bank, clock, law and remainder. Their
finite-source, fixed-time or finite-shell conclusions do not provide
the missing YM conditional estimate. This checkpoint is not a
recertification of every proof in those manuscripts.

The scheduled ten-version literature sweep covered conditional
transport, Stein kernels, matrix Harnack methods, nonconvex Gibbs
fields, exact factorization, entropic transport and lattice YM.
The following is a targeted primary-source review, not a claim
of exhaustive coverage.
\begin{enumerate}
\item Dolera--Mainini,
\emph{Lipschitz continuity of probability kernels in the optimal
transport framework},
\url{https://arxiv.org/abs/2010.08380}:
Theorem 2.2(iii) and the Poisson discussion relate conditional
transport to a Fisher quantity multiplied by a conditional
Poincare constant. That constant is an input; importing the
conclusion here would not construct the missing inverse.
\item Courtade--Fathi--Pananjady,
\emph{Existence of Stein kernels under a spectral gap, and
discrepancy bound},
\url{https://arxiv.org/abs/1703.07707}:
the primary abstract explicitly assumes a spectral gap.
We do not use this as an independent gap proof.
\item Li--Zhao, \emph{Wasserstein information matrix},
\url{https://arxiv.org/abs/1910.11248}:
the conditional Poisson geometry is compatible with our earlier
f17/V95 and V27 identities. Rechecking it does not constitute a
new estimate or remove the inverse-generator cost.
\item Li--Zhang, \emph{Matrix Li--Yau--Hamilton estimates under
Ricci flow and parabolic frequency},
\url{https://doi.org/10.1007/s00526-024-02668-x}:
the Hessian commutations in Section 2 and the tensor evolution
in Section 3 suggest the useful upstream move. We derive the
fixed-metric, stationary, potential-dependent estimate below;
we do not apply their Ricci-flow theorem to a different equation.
The tensor minimum method is classical, not claimed as a new
general principle.
\item Buchholz--Cotar--Schweiger,
\emph{Gradient Gibbs measures with non-convex potentials and
universality class of Gaussian Free Field},
\url{https://arxiv.org/html/2608.14526v1}:
the abstract, Definition 1.2 and Theorem 1.3 concern scalar
gradient potentials with explicit growth and lower-Hessian
hypotheses, using mixtures of strongly convex potentials.
No corresponding representation of our compact nonabelian
ground-state law has been established.
\item Wiggins, \emph{What solvable models reveal about
Born--Oppenheimer, Born--Huang, and exact factorization},
\url{https://arxiv.org/abs/2608.08668v2}:
the revised abstract distinguishes conditional amplitudes from
frozen eigenstates and emphasizes omitted-channel separation.
This is a useful audit warning, not an imported YM theorem.
\item Conforti, \emph{Weak semiconvexity estimates for
Schr\"odinger potentials and logarithmic Sobolev inequality
for Schr\"odinger bridges},
\url{https://arxiv.org/abs/2301.00083}:
the primary abstract concerns entropic-transport potentials
and prescribed marginals. The shared word
\emph{Schr\"odinger} does not identify those potentials with
the logarithm of our Hamiltonian ground state.
\item Shen--Zhu--Zhu, \emph{A stochastic analysis approach
to lattice Yang--Mills at strong coupling},
\url{https://arxiv.org/abs/2204.12737}:
the stated small-coupling spectral inequalities concern a
Wilson Gibbs law and its stochastic generator. They are
not a weak-coupling quantum-vacuum or physical-clock result.
\end{enumerate}

The nonduplication check is also important. Archive f17/V9 already
contains the weighted nongradient divergence identity used below,
for its finite-history action law. Archive f14/V64 contains
boundary-retaining semiconvex marginal estimates in charts.
We reuse the identity with attribution to that internal result.
The new step here is a global, volume-independent estimate for
the logarithmic Hessian of the \emph{native Hamiltonian ground
state}, followed by its application to V28--V29's complete lift.

\subsection{A global logarithmic-ground Hessian estimate}

Fix the finite spatial regulator of V28--V29, with $N\geq3$,
and the unit-round product metric on
$M=\mathrm{SU}(2)^{E_N}=(S^3)^{E_N}$. Work on this full smooth
manifold, not a gauge quotient. Let $\kappa>0$, $b\geq0$ and
$\beta=b/\kappa$. The positive smooth normalized ground state
$\Omega$ satisfies
\begin{equation*}
 (-\Delta+\mathcal V)\Omega=\epsilon_0\Omega,\qquad
 \mathcal V=\beta E_{\rm sp},\qquad
 u=\log\Omega,\qquad d\nu=\Omega^2\,dU .
 \tag{V30.1}\label{c18v30:ground}
\end{equation*}
Here $dU$ is normalized Haar volume; $\epsilon_0$ is the
ground energy after division by $\kappa$. The exact magnetic
Hessian estimate inherited from V29 gives
\begin{equation*}
 \operatorname{Hess}\mathcal V\leq16\beta\,g .
 \tag{V30.2}\label{c18v30:potential-upper}
\end{equation*}
Mixed sectional curvatures of the product are zero, all
sectional curvatures are nonnegative, and
$\operatorname{Ric}=2g$ is parallel.

\begin{lemma}[Stationary tensor minimum estimate]
Let a compact Riemannian manifold without boundary have nonnegative sectional
curvature and parallel Ricci tensor. Suppose a positive smooth
function $\Omega$ satisfies
$(-\Delta+\mathcal V)\Omega=\epsilon\Omega$ and
$\operatorname{Hess}\mathcal V\leq K_0g$, with $K_0\geq0$.
Then
$\operatorname{Hess}\log\Omega\geq-\sqrt{K_0/2}\,g$.
\end{lemma}

\begin{proof}
Put $s=\nabla u$ and $H=\operatorname{Hess}u$. Taking the
logarithm of the equation gives
\begin{equation*}
 \Delta u+|s|^2=\mathcal V-\epsilon .
 \tag{V30.3}\label{c18v30:log-equation}
\end{equation*}
Use the curvature convention in which, on a unit sphere,
$R_{ikjl}=\delta_{ij}\delta_{kl}-\delta_{il}\delta_{kj}$,
and write $\mathcal R(B)_{ij}=R_{ikjl}B_{kl}$.
For the rough Laplacian on covariant tensors the commutations
are
\begin{align*}
 \operatorname{Hess}\Delta u
   &=\Delta H+2\mathcal R(H)-\operatorname{Ric}H
                         -H\operatorname{Ric},\notag\\
 \operatorname{Hess}|s|^2
   &=2H^2+2\nabla_sH+2\mathcal R(s\otimes s).
 \tag{V30.4}\label{c18v30:commutations}
\end{align*}
They follow by commuting two covariant derivatives of
$du$ and then differentiating $|du|^2$; the derivatives
of Ricci vanish by hypothesis. In particular the
$2\mathcal R(s\otimes s)$ sign is positive in this convention.
Thus
\begin{align*}
 0={}&(\Delta+2\nabla_s)H+2H^2-\operatorname{Hess}\mathcal V
       +2\mathcal R(H)-\operatorname{Ric}H-H\operatorname{Ric}
       +2\mathcal R(s\otimes s).
 \tag{V30.5}\label{c18v30:stationary-tensor}
\end{align*}

Let $\lambda$ be the least eigenvalue of $H$ over the unit
tangent bundle, attained at a point and unit vector $v$.
If $\lambda\geq0$ there is nothing to prove. Otherwise extend
$v$ locally by parallel transport along radial geodesics
from that point. Its first covariant derivatives and rough
Laplacian vanish there. Since $H(v,v)$ has a local minimum,
$((\Delta+2\nabla_s)H)(v,v)\geq0$. Also
$H^2(v,v)=\lambda^2$. The Jacobi operator
$R_v:w\mapsto R(v,w,v,\cdot)$ is positive semidefinite and
$H-\lambda I\geq0$ at the point, so
\[
 \bigl(2\mathcal R(H)-\operatorname{Ric}H
                  -H\operatorname{Ric}\bigr)(v,v)
       =2\Tr\bigl(R_v(H-\lambda I)\bigr)\geq0 .
\]
Finally $\mathcal R(s\otimes s)(v,v)=R(v,s,v,s)\geq0$.
Evaluation of \eqref{c18v30:stationary-tensor} gives
$0\geq2\lambda^2-K_0$. Hence
$\lambda\geq-\sqrt{K_0/2}$, as asserted.
\end{proof}

\begin{corollary}[Native logarithmic semiconvexity]
With $a_\beta=\sqrt{8\beta}$, the actual ground state obeys
\begin{equation*}
 \operatorname{Hess}u\geq-a_\beta g,\qquad
 \operatorname{Hess}(-2u)\leq2a_\beta g
                           =4\sqrt{2\beta}\,g .
 \tag{V30.6}\label{c18v30:Hessian}
\end{equation*}
These constants do not depend on the number of fine links.
\end{corollary}

This is an \emph{upper} bound on the Hessian of the density
potential $-2u$, not the lower bound used in a
Bakry--Emery spectral-gap criterion. In particular we have
\begin{equation*}
 \operatorname{Ric}-2\operatorname{Hess}u
                    \leq(2+2a_\beta)g,
 \tag{V30.7}\label{c18v30:BE-upper}
\end{equation*}
not a proved positive lower bound for this tensor. At
$\beta=0$ the ground state is constant, as it should be.

\subsection{An unconditional full-bank Fisher bound}

Let $\{X_i\}$ be the orthonormal invariant fine-link frame
and $\sigma_i=X_i u$. For any constant real coefficient bank
$z$, the vector field $V_z=\sum_i z_iX_i$ has Haar divergence
zero, norm $|z|$, and $\nabla_{V_z}V_z=0$. Integration by
parts against $\Omega^2dU$ and \eqref{c18v30:Hessian} give
\[
 \int(V_zu)^2\,d\nu
   =-\frac12\int V_z^2u\,d\nu
   =-\frac12\int H(V_z,V_z)\,d\nu
   \leq\frac{a_\beta}{2}|z|^2 .
\]
The means of all $\sigma_i$ are zero. In matrix order,
\begin{equation*}
 \mathbb E_\nu[\sigma\sigma^*]\leq\sqrt{2\beta}\,I,
 \qquad
 \mathbb E_\nu[(2\sigma)(2\sigma)^*]\leq
                                      4\sqrt{2\beta}\,I .
 \tag{V30.8}\label{c18v30:Fisher}
\end{equation*}
This is the operator bound for the full bank, not a trace
bound multiplied by its dimension. A configuration-dependent
bank cannot be substituted here without paying its derivatives.

\subsection{The weighted divergence identity and the connection cost}

For a smooth real vector field $V$, define its complete
native score by
\begin{equation*}
 S_V=\operatorname{div}_\nu V
       =\operatorname{div}_{dU}V+2du(V).
 \tag{V30.9}\label{c18v30:field-score}
\end{equation*}
The identity already used in f17/V9, now for the native law,
is
\begin{equation*}
 \int S_V^2\,d\nu
  =\int\left\{\sum_{i,j}(\nabla_iV^j)(\nabla_jV^i)
        +(\operatorname{Ric}-2H)(V,V)\right\}d\nu .
 \tag{V30.10}\label{c18v30:divergence}
\end{equation*}
For completeness, integrate
$\int S_V^2=-\int V^i\nabla_i(\nabla_jV^j+2u_jV^j)$.
Commuting the two derivatives of $V$ produces
$\operatorname{Ric}(V,V)$. Integrating the remaining
$\nabla_j\nabla_iV^j$ once more produces the displayed
derivative contraction and a term
$2u_jV^i\nabla_iV^j$; this cancels the same term from
differentiating $2u_jV^j$. The remaining term is
$-2H(V,V)$, proving the identity. There is no boundary term.

The derivative contraction need not be nonnegative, but
it is bounded above by $\|\nabla V\|_{\rm HS}^2$. Therefore
\begin{equation*}
 \mathbb E_\nu S_V^2
 \leq\mathbb E_\nu\|\nabla V\|_{\rm HS}^2
              +(2+2a_\beta)\mathbb E_\nu|V|^2 .
 \tag{V30.11}\label{c18v30:score-budget}
\end{equation*}
The Haar divergence and the ground score must be kept
together in this identity; bounding each pointwise would
discard its cancellation.

In our frame $\nabla_{X_a}X_b=-[e_a,e_b]/2$ on one link.
For coefficients $V^i$ this yields
\begin{equation*}
 \|\nabla V\|_{\rm HS}
 \leq\left(\sum_i|X_i(V^\bullet)|^2\right)^{1/2}
                         +\sqrt2\,|V|.
 \tag{V30.12}\label{c18v30:connection}
\end{equation*}
Indeed $\sum_a|[e_a,w]/2|^2=2|w|^2$ on each
$\mathfrak{su}(2)$ block, and the product blocks are
orthogonal. Thus the connection cost also has no volume
factor. The law-dependent input in
\eqref{c18v30:score-budget} is now $a_\beta$, rather
than the finite-history action Hessian from f17/V9.

\subsection{The complete native lift score and conditional centering}

Keep exactly the full map $c$, its matrix $C$, the positive
matrix $G=CC^*$ and the weighted right inverse $A$ from
V28--V29. Keep their constants $K_\ell$, $T_\ell$ and
$g_-=\ell e^{-32\ell^2}$; in particular $CA=I$ and
$G\geq g_-I$. No coarse output, gauge direction or fine
harmonic mode is removed. Write $\mu=c_\#\nu=\rho\,dg$,
$\nu_g$ for the full conditional law and
$\overline H(g)=\mathbb E_{\nu_g}G^{-1}$.
Define
\begin{equation*}
 B_{\ell,\beta}
  =\bigl(T_\ell+\sqrt{2K_\ell}\bigr)^2
                           +(2+2a_\beta)K_\ell .
 \tag{V30.13}\label{c18v30:constants}
\end{equation*}
For each coarse frame index $\gamma$, put
$S_\gamma=\operatorname{div}_\nu A_\gamma$,
$\overline S(g)=\mathbb E_{\nu_g}S$, and
$r=S-\overline S(c)$.

\begin{theorem}[Mean-square score budget for the complete lift]
As inequalities of constant-bank quadratic forms,
\begin{align*}
 \mathbb E_\nu[SS^*]
   &\leq B_{\ell,\beta}\mathbb E_\mu\overline H
                      \leq\frac{B_{\ell,\beta}}{g_-}I,\notag\\
 \mathbb E_\mu\mathbb E_{\nu_g}[rr^*]
   &\leq B_{\ell,\beta}\mathbb E_\mu\overline H .
 \tag{V30.14}\label{c18v30:full-score}
\end{align*}
Moreover $\overline S_\gamma=X_\gamma\log\rho$.
\end{theorem}

\begin{proof}
For a constant coarse coefficient bank $v$, take $V=Av$.
V28 gives $|Av|^2\leq K_\ell v^*G^{-1}v$.
V29 and \eqref{c18v30:connection} give
$\|\nabla(Av)\|_{\rm HS}
 \leq(T_\ell+\sqrt{2K_\ell})\sqrt{v^*G^{-1}v}$.
Apply \eqref{c18v30:score-budget} and note
$S_V=v\cdot S$. This proves the first bound and its
elliptic corollary.

For a smooth coarse test function $f$, the exact identity
$CA=I$ gives $A_\gamma(f\circ c)=(X_\gamma f)\circ c$.
Hence
\[
 \int f(g)\overline S_\gamma(g)\,d\mu(g)
  =-\int X_\gamma f\,d\mu
  =\int f\,X_\gamma\log\rho\,d\mu .
\]
The full compact submersion and positive smooth ground
density make $\rho$ smooth positive, so this identifies
$\overline S$ pointwise. Conditional orthogonality yields
\begin{equation*}
 \mathbb E_\nu[SS^*]
   =\mathbb E_\mu\mathbb E_{\nu_g}[rr^*]
                        +\mathbb E_\mu[\overline S\overline S^*].
 \tag{V30.15}\label{c18v30:decomposition}
\end{equation*}
The second term is positive semidefinite, proving the
remaining claim.
\end{proof}

The same formula bounds the coarse marginal Fisher matrix.
However \eqref{c18v30:full-score} is averaged in $g$.
It does not prove a pointwise conditional bound, and
one cannot insert V26's $g$-dependent protected mask into
a constant-bank inequality.

\subsection{Variable coarse banks: the necessary derivative payment}

There is a useful version for a smooth coarse bank $v(g)$,
provided its derivatives are paid. Put
$R_{\gamma\delta}=X_\delta v^\gamma$ and
\begin{equation*}
 B'_{\ell,\beta}
   =2\bigl(T_\ell+\sqrt{2K_\ell}\bigr)^2
                            +(2+2a_\beta)K_\ell .
 \tag{V30.16}\label{c18v30:variable-constant}
\end{equation*}
Then
\begin{equation*}
 \mathbb E_\nu|v(c)\cdot r|^2
 \leq B'_{\ell,\beta}\int v^*\overline H v\,d\mu
       +2K_\ell e^{64\ell^2}\int\|D_Iv\|_{\rm HS}^2\,d\mu ,
 \tag{V30.17}\label{c18v30:variable-bound}
\end{equation*}
where $D_Iv=(X_\delta v^\gamma)_{\gamma\delta}$ is the
ordered invariant coefficient derivative.

To prove this, use $V=A(v\circ c)$. Differentiating its
coefficients and then adding the connection gives
\[
 \|\nabla V\|_{\rm HS}
 \leq (T_\ell+\sqrt{2K_\ell})
              \sqrt{v(c)^*G^{-1}v(c)}
                         +\|ARC\|_{\rm HS}.
\]
The full-map bounds
$\|A\|^2\leq K_\ell/g_-$ and
$\|C\|^2\leq\ell e^{32\ell^2}$ imply
$\|ARC\|_{\rm HS}^2
 \leq K_\ell e^{64\ell^2}\|R\|_{\rm HS}^2$.
Apply $(x+y)^2\leq2x^2+2y^2$ in
\eqref{c18v30:score-budget}. Finally the product rule and
$CA=I$ give
\[
 S_V=v(c)\cdot S+
          (\operatorname{div}_{dg}v)(c),\qquad
 S_V-\mathbb E[S_V\mid c]=v(c)\cdot r .
\]
Conditional centering is an $L^2$ contraction, proving
\eqref{c18v30:variable-bound}. At each fixed regulator it
extends by smooth approximation to coefficient $H^1$
banks, since all densities and coefficients are smooth
on compact spaces.

In particular $v=\nabla_I f$ incurs ordered second
derivatives of $f$. They are not paid by a first-order
coarse Dirichlet energy. Discontinuous protected masks
need not have the required $H^1$ regularity. This
explicit derivative cost prevents an illicit upgrade
of the constant-bank result.

\subsection{What this gives, and does not give, for the conditional inverse}

Let $L_g=-\operatorname{div}_{\nu_g}\nabla_{\rm vert}$
be the actual conditional generator on the full connected
fibre, with no factor of $\kappa$. For an auxiliary
$\lambda>0$ define
\begin{equation*}
 R_\lambda(g)_{\gamma\delta}
   =\langle r_\gamma,(L_g+\lambda)^{-1}
                                      r_\delta\rangle_{\nu_g}.
 \tag{V30.18}\label{c18v30:shift}
\end{equation*}
The spectral theorem and the score budget imply
\begin{equation*}
 \mathbb E_\mu R_\lambda
  \leq\frac{B_{\ell,\beta}}{\lambda}\mathbb E_\mu\overline H
  \leq\frac{B_{\ell,\beta}}{\lambda g_-}I .
 \tag{V30.19}\label{c18v30:shift-bound}
\end{equation*}
For variable banks the right-hand side of
\eqref{c18v30:variable-bound}, divided by $\lambda$,
similarly bounds
\[
 \int v(g)^*R_\lambda(g)v(g)\,d\mu(g).
\]

This auxiliary shift is a diagnostic resolvent parameter,
not a mass inserted into the Hamiltonian or a change of
the conditional law. At any fixed regulator,
$R_\lambda(g)\uparrow R_0(g)$ as $\lambda\downarrow0$
on centered sources. Our bound diverges in that limit.
A bounded total spectral mass of $r$ does not bound its
inverse-frequency moment. The remaining upstream task
is low-frequency control for this \emph{specific} residual,
or a controlled vertical-divergence representation that
pays its inverse without assuming the desired conditional
Poincare inequality. The full native construction, not
an unrelated measure with a spectral inequality, must
supply that control.

\subsection{Verification, scope and continuation}

Exact quaternion regressions check the curvature
commutations and the stationary tensor identity on a
nonseparable two-link test. Weighted one-link integration
checks the nongradient identity, including a Killing
field and a nonconstant multiple. A radial Wilson-rotor
finite-section calculation provides an additional
numerical diagnostic of the semiconvexity bound.
The engineered two-link density is not a Wilson vacuum;
the rotor is not the full spatial interacting vacuum.
These regressions test signs and normalization in the
written proof. They are not proof-assistant certificates
or numerical evidence for a continuum mass gap.

V29 is recoverable exactly from this source, apart from
the version title change. All other 47 note sources are
preserved. Only the active V30 replaces V29; build and
inspection products remain outside \texttt{latex\_notes}.
The companion and broad-literature checkpoints have both
been completed at V30; their next scheduled versions
are V35 and V40. The cycle still restarts at V100.

The advance is a native, volume-independent-at-fixed-scale
log-ground semiconvexity bound and a complete lift-score
budget, including its conditional centering. The constants
still grow rapidly with $\ell$ through $T_\ell$, and
$a_\beta$ grows with $\sqrt\beta$. No cofinal estimate
follows from volume independence alone. The unshifted
native score-Poisson residual, projection-energy stability,
conditional and coarse variance bounds, physical return
and scale transport remain open. The interacting
four-dimensional continuum Yang--Mills theory and its
positive physical mass gap remain unproved.
% END CYCLE 18 VERSION 30 APPEND
% BEGIN CYCLE 18 VERSION 31 APPEND
\clearpage
\section{V31: Native fibre resampling and a finite-time score correction}
\label{c18v31:context}

V30 controls the native score in mean square, but an unshifted
conditional inverse can amplify its low-frequency part. Here we
construct actual local conditional solvers on the full flowed,
sampled fibre and use them to correct the lift for a finite
auxiliary time. Every coarse output stays fixed under the
correction. The resulting residual is an explicitly defined
heat-bath semigroup applied to the original residual. Its
correction cost is paid without assuming a global fibre gap.
The low-frequency inverse moment is still a missing estimate.

The previous turn was progress: V30's source, preservation
certificate and new proof were revalidated. The nonduplication
audit rechecked current V13's global chain coordinates, V25's
inverse-flow envelopes and Jacobian, and archive f17 V91--V93's
native one-link comparison and logarithmic-gradient estimate.
Those one-link results do not already treat the conditional law
on the present \emph{full flowed fibre}. We retain their methods
and constants where useful, rather than recounting them as new.

For a targeted literature check, Caputo--Menz--Tetali,
\emph{Approximate tensorization of entropy at high temperature},
\url{https://www.numdam.org/item/AFST_2015_6_24_4_691_0.pdf},
Section 1.1, uses the rate-one conditional-expectation generator.
Its Theorem 2.1 requires a quantitative weak-dependence condition.
We use the standard generator, but neither assert that condition
nor import its tensorization conclusion for our vacuum.
This targeted check does not replace the completed V30 broad
review. The next companion and broad checkpoints remain V35
and V40.

\subsection{A pointwise native input and the transported density}

Keep $N=\ell n$, $\ell\geq2$, $n\geq5$, $\beta=b/\kappa\geq0$,
$u=\log\Omega$ and $\nu=\Omega^2dU$. The spatial flow is
$\Phi_s$ of $-\nabla E_{\rm sp}$, and
$c=\mathfrak b_\ell\circ\Phi_s$. First allow any $s\geq0$;
the application to V30 uses $s=\ell^2$.

\begin{lemma}[A square-root-coupling pointwise link score]
At every fine configuration and every link $e$,
\[
 |\nabla_e u|\leq
 \mathfrak s_\beta:=
       \min\{2\beta,\ \pi\sqrt{8\beta}\}.
 \tag{V31.1}\label{c18v31:link-score}
\]
The second bound is a consequence of V30. The first is the
already proved f17/V93 native maximum-principle estimate.
\end{lemma}
\begin{proof}
For a unit tangent $v$ at link $e$, its invariant geodesic
$\gamma(t)=\exp(tv)U_e$ has period $2\pi$ on the unit
three-sphere. Fix every other link and put $f(t)=u(\gamma(t),y)$.
V30 gives $f''(t)\geq-a_\beta$, where $a_\beta=\sqrt{8\beta}$.
Thus $f'(t)\geq f'(0)-a_\beta t$ on $[0,2\pi]$.
Periodicity implies
$0=\int_0^{2\pi}f'(t)dt\geq2\pi f'(0)-2\pi^2a_\beta$.
Reverse $v$ to get $|du(v)|\leq\pi a_\beta$.
Taking the supremum over the unit vectors in that link and
combining with the inherited $2\beta$ bound proves the claim.
This is not a volume-independent bound on the norm of the
entire fine gradient.
\end{proof}

Use flowed variables $W=\Phi_s(U)$ and let
$\mathcal J_{-s}(W)$ be the positive Haar-volume Jacobian of
$\Phi_{-s}$. The exact density in these variables is
\[
 h_s(W)=\Omega(\Phi_{-s}W)^2\mathcal J_{-s}(W).
 \tag{V31.2}\label{c18v31:flowed-density}
\]
It is not a declared Wilson Gibbs density. Define
\[
 \mathfrak d_s
   =2\mathfrak s_\beta e^{16s}+3(e^{16s}-1).
 \tag{V31.3}\label{c18v31:density-constant}
\]

\begin{proposition}[Uniform one-link derivative of the exact flowed density]
For every flowed link $f$,
\[
                     |\nabla_{W_f}\log h_s|\leq\mathfrak d_s .
 \tag{V31.4}\label{c18v31:flowed-score}
\]
The constant is independent of volume and of every exterior.
\end{proposition}
\begin{proof}
Recall V25's block majorant
$\mathsf K_t=\exp(t(4I+\mathsf N))$, where
$\mathsf N_{ef}$ counts common plaquettes for distinct links,
and has zero diagonal. All its row and column sums are
$e^{16t}$. It bounds differential blocks of the flow in
either time direction, in endpoint orthonormal frames.
Consequently a unit perturbation of one flowed link has
$\ell^1$ sum of link tangent norms at most $e^{16s}$ after
$\Phi_{-s}$. Equation \eqref{c18v31:link-score} bounds its
contribution to $2\,d(u\circ\Phi_{-s})$ by
$2\mathfrak s_\beta e^{16s}$.

V25's exact Liouville formula, with time reversed, is
\[
 \log\mathcal J_{-s}(W)
       =12\int_0^s\sum_p w_p(\Phi_{-t}W)\,dt.
\]
Since each link belongs to four plaquettes and each link
gradient of $w_p$ has norm at most one, the link gradient
of $12\sum_pw_p$ has norm at most $48$.
The same block majorant bounds the differentiated integral
by $48\int_0^s e^{16t}dt=3(e^{16s}-1)$.
Both terms in \eqref{c18v31:flowed-density} have been kept.
\end{proof}

\subsection{Exact fibre coordinates and their original metric}

Let $z$ retain the first $\ell-1$ flowed links of every
sampling chain, and all flowed links not in any chain.
There are $m=|E_N|-|E_c|$ group coordinates in $z$.
For each chain put $P_\gamma=z_{\gamma,0}\cdots
z_{\gamma,\ell-2}$ and reconstruct its last link as
$W_{\gamma,\ell-1}=P_\gamma^{-1}g_\gamma$.
Denote this global reconstruction by $W=\Theta(g,z)$ and set
\[
 \chi_g(z)=\Phi_{-s}\Theta(g,z),\qquad
 q_g(z)=\frac{h_s(\Theta(g,z))}{\rho(g)},\qquad
 \rho(g)=\int h_s(\Theta(g,z))\,dz .
 \tag{V31.5}\label{c18v31:fibre-law}
\]
Here $dz$ and $dg$ are product normalized Haar measures.
V13's triangular coordinate change has Haar Jacobian one.
It follows directly by change of variables that
$\mu(dg)=\rho(g)dg$ and
$\nu_g=(\chi_g)_\#(q_gdz)$ are the actual full conditional laws.

These coordinates do not fix a gauge and do not remove
any fibre direction. The Jacobian in $h_s$ already supplies
the required coarea weight. Multiplying $q_g$ by another
induced-volume determinant would change the law.

Write $\mathcal K=D_z\chi_g$ and
$\mathcal M=\mathcal K^*\mathcal K$, with the original
fine metric and the product metric in $z$. Then
\[
 e^{-32s}I\preceq\mathcal M\preceq
                          \ell e^{32s}I .
 \tag{V31.6}\label{c18v31:metric}
\]
Indeed on a chain the unflowed vertical differential has
the form $\binom{I}{-\mathcal R}$, where $\mathcal R$
concatenates $\ell-1$ orthogonal adjoint matrices.
Thus $\mathcal R\mathcal R^*=(\ell-1)I$ and the squared
singular values of this differential lie in $[1,\ell]$.
Unused links give identity blocks. Compose with the
flow differential and its inverse, both bounded by
$e^{16s}$, to obtain \eqref{c18v31:metric}.

For $\widetilde f=f\circ\chi_g$ the original vertical
Dirichlet form, without $\kappa$, is consequently
\[
 \mathcal E_g(f)
   =\int\langle\nabla_z\widetilde f,
                  \mathcal M^{-1}\nabla_z\widetilde f\rangle q_gdz
   \geq\frac{e^{-32s}}{\ell}
                         \int|\nabla_z\widetilde f|^2q_gdz .
 \tag{V31.7}\label{c18v31:energy}
\]
This uses the inverse of the induced metric, not a
replacement of that metric by a product metric.

\subsection{Local conditional inverses on the actual fibre}

For a coordinate $z_j$, let $P_j^g$ be expectation in
that coordinate conditioned on $g$ and all the other
$z$ coordinates, under $q_gdz$. It is an orthogonal
projection in $L^2(q_gdz)$. Put $D_j^g=I-P_j^g$.
Changing $z_j$ changes one flowed link if it is unused
by sampling, and two flowed links if it is in a chain.
Write $\mathfrak m_j=1$ or $2$ accordingly. The differential
on each affected link is an isometry. Hence
\[
 |\nabla_{z_j}\log h_s(\Theta(g,z))|
                  \leq\mathfrak m_j\mathfrak d_s,\qquad
 \operatorname{osc}_{z_j}\log h_s(\Theta(g,z))
                  \leq\pi\mathfrak m_j\mathfrak d_s .
 \tag{V31.8}\label{c18v31:coordinate-osc}
\]
These are uniform conditional bounds, not averages over $g$.

The gap $3$ of unit-round Haar $S^3$ and elementary density
comparison now prove, for every exterior and every smooth $f$,
\[
 \eta_j\operatorname{Var}_j^g(f)
      \leq P_j^g|\nabla_{z_j}f|^2,\qquad
 \eta_j=3e^{-\pi\mathfrak m_j\mathfrak d_s}
      \geq\eta_*:=3e^{-2\pi\mathfrak d_s}>0 .
 \tag{V31.9}\label{c18v31:local-floor}
\]
For clarity, if the conditional density relative to Haar
lies between $p_-$ and $p_+$, then
$\operatorname{Var}_p f\leq
(p_+/3p_-)\int|\nabla f|^2p$, and
$p_+/p_-\leq e^{\pi\mathfrak m_j\mathfrak d_s}$.
The conditional normalizer cancels from this ratio.

Let $\mathcal L_{j,g}$ be the weighted positive
three-sphere Laplacian for this one-coordinate conditional.
On its centered subspace its inverse exists with norm at
most $\eta_j^{-1}$, uniformly in all remaining coordinates.
This is a \emph{proved local inverse}; it is not an inverse
bound for the whole fibre.

\subsection{Comparison with a specified rate-one resampling operator}

Define on the actual fibre coordinates
\[
 \mathsf T_g=\sum_{j=1}^{m}D_j^g,\qquad
 \alpha_s=\frac{e^{-32s}}{\ell}\eta_*
          =\frac3\ell e^{-32s-2\pi\mathfrak d_s}.
 \tag{V31.10}\label{c18v31:heat-bath}
\]
The semigroup $e^{-t\mathsf T_g}$ updates each coordinate
at rate one using its full native conditional distribution.
There is no division by $m$: a total-rate-one normalization
would introduce an additional volume factor.

Let $L_g=-\operatorname{div}_{\nu_g}\nabla_{\rm vert}$
as in V30. Under the isometry induced by $\chi_g$, one has
the closed-form comparison
\[
 L_g\ \succeq\
 \frac{e^{-32s}}{\ell}\sum_j\eta_jD_j^g
 \ \succeq\ \alpha_s\mathsf T_g .
 \tag{V31.11}\label{c18v31:comparison}
\]
Integrate \eqref{c18v31:local-floor}, sum in $j$, and
apply \eqref{c18v31:energy}. Smooth density and compactness
extend this inequality to the form domain. The heat-bath
operators are bounded, and
$\langle f,\mathsf T_gf\rangle=\sum_j\|D_j^gf\|^2$.

At a fixed regulator $\mathsf T_g$ has kernel the constants
and a positive gap. This qualitative assertion does not
use a uniform-in-volume estimate: $q_g$ has a finite positive
minimum and maximum on the compact product. Comparing
the infimum over functions independent of $z_j$ with
product Haar gives
\[
 \sum_j\|D_j^gf\|_{q_g}^2
 \geq(\inf q_g)\sum_j\operatorname{Var}_{{\rm Haar},j}(f)
 \geq\frac{\inf q_g}{\sup q_g}\operatorname{Var}_{q_g}(f).
\]
In the middle expression each conditional variance is
integrated over the remaining Haar coordinates.
The middle inequality is product variance tensorization.
The ratio is not assigned a useful uniform value.

In particular, for a centered source $r$, the variational
formula for inverse energy gives the sufficient comparison
\[
 \langle r,L_g^{-1}r\rangle
       \leq\alpha_s^{-1}
                     \langle r,\mathsf T_g^{-1}r\rangle .
 \tag{V31.12}\label{c18v31:inverse-comparison}
\]
The right-hand inverse moment is not bounded here.
This is a concrete sufficient route, not an equivalence
that would exclude a better direct estimate for $L_g$.

\subsection{Geometric range is not conditional mixing}

The update geometry also has an explicit tail. Let $S_j$
be the one or two flowed links affected by coordinate $j$,
and use the graph joining fine links that share a plaquette.
For a unit coordinate tangent, write
$Z_{j,e}=(D_z\chi_g)_{e,j}$. V25's same majorant gives
\[
 \sum_{e:\,d(e,S_j)\geq R}\|Z_{j,e}\|
 \leq\mathfrak m_j e^{4s}
                \sum_{k\geq R}\frac{(12s)^k}{k!}
 =\mathfrak m_j e^{16s}
       \mathbb P\{\operatorname{Poisson}(12s)\geq R\},
 \quad R\geq1 .
 \tag{V31.13}\label{c18v31:geometric-tail}
\]
In fact $\mathsf N^k$ cannot connect graph distance
greater than $k$, and every column sum of $\mathsf N^k$
is $12^k$. Summing its exponential series proves the
bound, also on a periodic torus. Integrating a shortest
coordinate geodesic shows that the corresponding sum
of fine-link displacements under a finite resampling
is bounded by $\pi$ times the right-hand side.

This tail controls how a chosen coordinate move changes
the fine configuration. It does \emph{not} control how
its conditional probability changes when distant
coordinates are altered. That probability contains the
actual $\Omega$ in \eqref{c18v31:fibre-law}. Thus
$\mathsf T_g$ has geometrically concentrated moves but
is not a proved finite-range or rapidly mixing operator.
No Dobrushin margin, conditional angle estimate or
physical time-decay rate follows from
\eqref{c18v31:geometric-tail}.

\subsection{A paid finite-time correction of the actual lift}

All objects in this subsection first act at fixed $g$;
we suppress that index. For a smooth scalar $f$ define
the product-coordinate vector field
\[
 (\mathcal Q f)_j
      =\nabla_{z_j}\mathcal L_j^{-1}D_jf .
 \tag{V31.14}\label{c18v31:local-flux}
\]
The inverse uses its proved one-coordinate conditional
centering. Conditional integration by parts gives
\[
 \operatorname{div}_{q_g}\mathcal Q f=-\mathsf T_gf,
 \qquad
 \|\mathcal Q f\|_{L^2(q_g;\,\mathrm{product})}^2
       \leq\eta_*^{-1}\langle f,\mathsf T_gf\rangle .
 \tag{V31.15}\label{c18v31:flux-budget}
\]
Indeed each squared gradient norm is
$\langle D_jf,\mathcal L_j^{-1}D_jf\rangle
\leq\eta_*^{-1}\|D_jf\|^2$ after integrating the
exterior. The product vector norm sums these costs,
rather than squaring a sum over $m$ coordinates.
Divergence of the sum uses the exact density $q_g$.

Now let $A$ be V28--V30's complete smooth right inverse,
$S=\operatorname{div}_\nu A$ its column score and
$r=S-\mathbb E[S\mid c]$ its centered residual. For $t\geq0$ put
\[
 \begin{gathered}
 F_t(x)=
  \begin{cases}(1-e^{-tx})/x,&x>0,\\t,&x=0,\end{cases}
 \qquad F_t(\mathsf T)=\int_0^t e^{-\tau\mathsf T}\,d\tau,\\
 Z_t=\mathcal K\,\mathcal Q F_t(\mathsf T_g)r,\qquad
 A_t=A+Z_t .
 \end{gathered}
 \tag{V31.16}\label{c18v31:corrected-lift}
\]
These are columnwise definitions. The spectral quotient
in $F_t$ has a removable singularity and norm at most $t$.
Thus this finite-time construction does not secretly
invert the whole $\mathsf T_g$.

\begin{theorem}[Native finite-time score correction with paid norm]
The correction is vertical, the full output identity
$c_*A_t=I$ holds, and the actual centered score is
\[
 r_t=\operatorname{div}_\nu A_t
       -\mathbb E[\operatorname{div}_\nu A_t\mid c]
                  =e^{-t\mathsf T_g}r .
 \tag{V31.17}\label{c18v31:residual}
\]
For every constant coarse coefficient vector $v$ and every $g$,
\[
 \begin{split}
 \mathbb E_{\nu_g}|Z_tv|^2
 &\leq\alpha_s^{-1}
       \langle v\cdot r,F_t(\mathsf T_g)^2
                           \mathsf T_g(v\cdot r)\rangle_{\nu_g}\\
 &\leq\frac{t}{2\alpha_s}
                  \mathbb E_{\nu_g}|v\cdot r|^2 .
 \end{split}
 \tag{V31.18}\label{c18v31:correction-cost}
\]
The notation $c_*A_t$ means $Dc\,A_t$.
\end{theorem}
\begin{proof}
The range of $\mathcal K$ is the entire vertical tangent
space, hence $Dc\,Z_t=0$. The divergence of a vertical
field in coordinates $(g,z)$ is its conditional weighted
$z$-divergence, because the full density is $\rho(g)q_g(z)$.
Pushforward therefore preserves the divergence in
\eqref{c18v31:flux-budget}. It gives
\[
 \operatorname{div}_\nu Z_t
   =-\mathsf T_gF_t(\mathsf T_g)r
   =-(I-e^{-t\mathsf T_g})r.
\]
Its conditional mean is zero. This proves both the
full-output identity and \eqref{c18v31:residual}.

By \eqref{c18v31:metric} and
\eqref{c18v31:flux-budget}, the squared native norm is
at most $\ell e^{32s}/\eta_*=\alpha_s^{-1}$ times
the displayed spectral quadratic form. For $x\geq0$,
\[
 (1-e^{-x})^2\leq x/2.
\]
One direct proof is to differentiate
$x-2(1-e^{-x})^2$: its derivative is
$(1-2e^{-x})^2\geq0$, and its value at zero is zero.
Apply this with $x=t\lambda$ in spectral calculus
to obtain the second bound.

For domain details, each conditional coefficient is
smooth positive on a compact product. Parameter-dependent
elliptic inversion on the centered one-coordinate space
is smooth at fixed regulator. The conditional projections
are bounded on each fixed $C^k$ space; their finite sum
has an exponential defined there by its convergent
operator series. Thus the construction sends smooth
sources to smooth fields, including their $g$ dependence.
Alternatively all displayed identities hold first weakly,
and the uniform local inverse bound extends
\eqref{c18v31:flux-budget} to $L^2$ sources. No global
conditional gap is needed for finite $t$.
\end{proof}

The construction also preserves gauge equivariance.
A fixed fine gauge transformation acts by fixed left and
right multiplications on each $z_j$ and by endpoint
multiplications on $g$. These are Haar-preserving product
isometries, and $\Theta$ and $\Phi_{-s}$ intertwine them.
The native density, each conditional projection, and each
one-coordinate weighted Laplacian therefore transform
covariantly. Uniqueness of the centered local inverse
makes $\mathcal Q$ covariant as well. Hence an equivariant
initial $A$ gives an equivariant $A_t$, with the usual
coarse tangent transformation on its columns. No gauge
or harmonic variable is projected away.

At $s=\ell^2$, V30 now gives genuinely native
constant-bank matrix bounds
\[
 \begin{split}
 \mathbb E_\nu[Z_t^*Z_t]
       &\preceq\frac{tB_{\ell,\beta}}{2\alpha_{\ell^2}}
                                      \mathbb E_\mu\overline H,\\
 \mathbb E_\nu[A_t^*A_t]
       &\preceq
       \left(\sqrt{K_\ell}
             +\sqrt{\frac{tB_{\ell,\beta}}{2\alpha_{\ell^2}}}
       \right)^2\mathbb E_\mu\overline H .
 \end{split}
 \tag{V31.19}\label{c18v31:averaged-cost}
\]
The second line is Minkowski's inequality for each
constant $v$, using V28's original norm and the first line.
For a variable bank $v(g)$ the first line instead has
$t/(2\alpha_s)$ times the full right side of V30.17,
at $s=\ell^2$. Its derivative payment is not removed.
These averaged bounds are not pointwise-in-$g$ bounds
relative to $\overline G(g)^{-1}$.

\subsection{The high-frequency inverse cost and the remaining low part}

Let $\Sigma(g)=\mathbb E_{\nu_g}[rr^*]$ and, for
$\varepsilon>0$, define the matrix
\[
 M_{<\varepsilon}(g)_{\gamma\delta}
   =\left\langle r_\gamma,
       {\bf1}_{(0,\varepsilon)}(\mathsf T_g)
          \mathsf T_g^{-1}r_\delta\right\rangle_{\nu_g}.
 \tag{V31.20}\label{c18v31:low-moment}
\]
There is no constant-mode contribution because $r$ is
centered. The residual cost for the new full lift obeys
\[
 \mathsf R^{A_t}(g)
 :=\bigl(\langle r_{t,\gamma},L_g^{-1}r_{t,\delta}
                                      \rangle_{\nu_g}\bigr)_{\gamma\delta}
 \preceq\alpha_s^{-1}M_{<\varepsilon}(g)
       +\frac{e^{-2t\varepsilon}}{\alpha_s\varepsilon}\Sigma(g).
 \tag{V31.21}\label{c18v31:split}
\]
To prove it, apply \eqref{c18v31:inverse-comparison} to
$r_t=e^{-t\mathsf T_g}r$. Its heat-bath spectral multiplier
is $e^{-2t\lambda}/\lambda$. On $(0,\varepsilon)$ bound
the exponential by one, and on $[\varepsilon,\infty)$
bound the whole multiplier by
$e^{-2t\varepsilon}/\varepsilon$. This argument does not
assume that $L_g$ and $\mathsf T_g$ commute.

At $s=\ell^2$, averaging and V30 yield
\[
 \mathbb E_\mu\mathsf R^{A_t}
 \preceq\alpha_{\ell^2}^{-1}\mathbb E_\mu M_{<\varepsilon}
       +\frac{B_{\ell,\beta}e^{-2t\varepsilon}}
                    {\alpha_{\ell^2}\varepsilon}
                                      \mathbb E_\mu\overline H .
 \tag{V31.22}\label{c18v31:averaged-split}
\]
The second term has now been paid without an inverse
gap assumption. The first has not. Neither letting
$t\to\infty$ nor choosing a small $\varepsilon$ supplies
a bound for it. At a fixed regulator the residual tends
to zero, but the correction norm need not stay uniformly
bounded in the regulator. The finite-time construction
therefore does not complete V27's certificate.

The parameter $t$ is an auxiliary resampling time.
The original Hamiltonian, its physical clock $\kappa$,
the conditional law, and the coarse map have not changed.
No resolvent shift or artificial Hamiltonian mass is
introduced by \eqref{c18v31:corrected-lift}.

\subsection{Verification and the next upstream decision}

The regressions check noncommuting quaternion chain
coordinates, the complete vertical metric and coarea
determinant identity, and an explicit spherical inverse
flow with its Jacobian and transported density.
Separate finite reversible conditional-expectation
matrices test the flux sign, the finite-time correction,
the $t/2$ cost and the high/low spectral split without
commuting the conditional projections. These finite
fixtures are not the interacting spatial vacuum and
do not certify an infinite-volume spectrum. The
all-regulator assertions above rest on the written proofs.

The single-coordinate inverse is paid for every exterior
on the actual flowed fibre. Its explicit constant is
unfortunately very small:
$-\log\alpha_s=\log(\ell/3)+32s+2\pi\mathfrak d_s$,
and $\mathfrak d_s$ grows like
$e^{16s}\sqrt\beta$ at large coupling. Thus these constants
are not a cofinal physical estimate. A useful next
advance must exploit source-specific low-frequency
cancellation, improve the native conditional comparison,
or control suitable interacting blocks. Repeating
compactness or geometric range arguments cannot pay the
missing conditional dependence. This concrete correction
is an available route, not a requirement to abandon a
sharper direct treatment of $L_g$.

V30's body is retained exactly apart from the title.
Only the active V31 replaces V30; the other 47 sources
are unchanged and the notes folder contains only
\texttt{.tex}. Diagnostics and rendering products stay
outside it. The next companion review is V35, the
next broad literature sweep V40, and the archive restart
V100. The native uniform low-frequency inverse estimate,
projection-energy stability, conditional and coarse
variance, physical return and scale transport remain
unproved. So do the interacting four-dimensional
continuum construction and the positive physical
Yang--Mills mass gap.
% END CYCLE 18 VERSION 31 APPEND
% BEGIN CYCLE 18 VERSION 32 APPEND
\clearpage
\section[V32: compact finite changes and tunnelling risk]{V32: compact finite changes and tunnelling risk}
\label{c18v32:context}

V31 constructs local inverses for the actual conditional fibre law and
a finite-time correction of the full score, but its comparison constant
is too small for cofinal use. This continuation improves the upstream
conditional oscillation estimate. A finite replacement cannot move a
single $S^3$ link by more than $\pi$. Combining this cap with the spatial
propagation kernel gives a polynomial, instead of exponential,
flow-time envelope. This is not yet a conditional-mixing theorem.

Current V25 and archived f16.2/V74 already prove the inverse-flow
Jacobian and differentiated transported-density bounds; those are inputs,
not new results. Archived f6/V83 and V87 already discuss barriers and
the danger of substituting a bare Wilson law for the ground-weighted
law. Our additional $S^3$ example tests the specific curvature inputs
of V30--V31 with a verified Schr\"odinger ground state. It is not
identified with any native Yang--Mills conditional measure.

For targeted context we checked the primary abstracts of H.~Raz and
R.~Sims, \emph{Lieb--Robinson bounds for classical anharmonic lattice
systems}, \href{https://arxiv.org/abs/0902.0025}{arXiv:0902.0025},
and M.~Brooks and G.~Di~Ges\`u, \emph{Sharp Tunneling Estimates for a
Double-Well Model in Infinite Dimension},
\href{https://arxiv.org/abs/1911.03187}{arXiv:1911.03187}.
These concern oscillator dynamics and infinite-dimensional tunnelling,
respectively, not our compact Wilson flow or fibre law.
No theorem is imported from these abstracts: all new estimates are
proved below. This targeted check does not replace the V40 broad sweep.

\subsection{Capped finite-change propagation}

Use V31's spatial regulator $N=\ell n$, $\ell\ge2$, $n\ge5$, with
unit-round $S^3$ fine-link metrics. Let $\mathsf N_{ef}$ count shared
spatial plaquettes for $e\ne f$, with zero diagonal. Its row and column
sums are twelve. Let $d(e,f)$ denote graph distance for its positive
entries. V25's variational bound, in either flow direction, reads
\[
 \|(D\Phi_{\pm t})_{ef}\|\le(\mathsf K_t)_{ef},\qquad
 \mathsf K_t=\exp\{t(4I+\mathsf N)\},\quad t\ge0.
 \tag{V32.1}\label{c18v32:kernel}
\]
Put
\[
 P(t)=3(80t+1)^3+1,\qquad F(t)=\min\{e^{16t},P(t)\}.
 \tag{V32.2}\label{c18v32:F}
\]
\begin{lemma}[A volume-independent finite-change envelope]
\label{c18v32:capped}
If $W,W'$ agree outside a set $S$ of $k$ flowed links, then
\[
 \sum_e d_{S^3}\bigl((\Phi_{\pm t}W)_e,(\Phi_{\pm t}W')_e\bigr)
 \le\pi k F(t).
 \tag{V32.3}\label{c18v32:distance}
\]
This is a finite-distance estimate, not a polynomial differential bound.
\end{lemma}
\begin{proof}
Connect the changed coordinates by geodesics of lengths $\delta_f\le\pi$.
Their image under the flow and the diameter cap give
\[
 d_{S^3}\bigl((\Phi_{\pm t}W)_e,(\Phi_{\pm t}W')_e\bigr)
 \le\min\left\{\pi,\sum_{f\in S}(\mathsf K_t)_{ef}\delta_f\right\}
 \le\pi\min\left\{1,\sum_{f\in S}(\mathsf K_t)_{ef}\right\}.
 \tag{V32.4}\label{c18v32:saturation}
\]
For a later conditional replacement only the endpoints must share the
coarse value. The connecting path used for this global flow bound need
not remain on the fibre.

The column sums of $\mathsf K_t$ are $e^{16t}$, giving the exponential
envelope. Since $(\mathsf N^j)_{ef}=0$ when $d(e,f)>j$,
\[
 \sum_e e^{d(e,f)}(\mathsf K_t)_{ef}
 \le e^{4t}\sum_{j\ge0}\frac{t^j e^j12^j}{j!}
 =e^{(4+12e)t}\le e^{40t}.
 \tag{V32.5}\label{c18v32:weighted-kernel}
\]
Every graph step changes each stored-link tail coordinate by at most
one in periodic distance. There are three orientations per tail.
For integer $R\ge1$ this implies
\[
 \#\{e:d(e,S)<R\}\le3k(2R-1)^3.
 \tag{V32.6}\label{c18v32:ball}
\]
Periodic identifications only reduce the count. Outside this ball
$d(e,f)\ge R$ for every $f\in S$; the weighted kernel bounds the
distance sum there by $\pi k e^{40t-R}$. Inside use the cap $\pi$.
For $t>0$, choose $R=\lceil40t\rceil$. Then $e^{40t-R}\le1$ and
$2R-1\le80t+1$, proving the polynomial envelope.
At $t=0$ the direct bound is $\pi k$ and $F(0)=1$.
\end{proof}

This controls replacement geometry, not remote dependence of a
conditional sampling kernel. It is not a Dobrushin estimate.

\subsection{Oscillation of the exact native density}

Retain the actual ground state $\Omega$, $u=\log\Omega$, and V31's bound
\[
 |\nabla_eu|\le\mathfrak s_\beta,\qquad
 \mathfrak s_\beta=\min\{2\beta,\pi\sqrt{8\beta}\}.
 \tag{V32.7}\label{c18v32:score-input}
\]
The exact flowed density relative to product Haar is
\[
 \begin{aligned}
 h_s(W)&=\Omega(\Phi_{-s}W)^2J_{-s}(W),\\
 \log J_{-s}(W)&=12\int_0^s\sum_p w_p(\Phi_{-t}W)\,dt .
 \end{aligned}
 \tag{V32.8}\label{c18v32:density}
\]
Define
\[
 \widehat d_s=2\mathfrak s_\beta F(s)+48\int_0^sF(t)\,dt .
 \tag{V32.9}\label{c18v32:dhat}
\]
\begin{theorem}[Native finite-replacement oscillation including the Jacobian]
\label{c18v32:oscillation}
If $W,W'$ differ on at most $k$ flowed links, then
\[
 |\log h_s(W)-\log h_s(W')|\le\pi k\widehat d_s,
 \tag{V32.10}\label{c18v32:osc}
\]
where, uniformly in spatial volume and exterior,
\[
 \begin{aligned}
 \widehat d_s
 &\le2\mathfrak s_\beta e^{16s}+3(e^{16s}-1)=d_s^{\rm old},\\
 \widehat d_s
 &\le2\mathfrak s_\beta[3(80s+1)^3+1]
        +\frac9{20}[(80s+1)^4-1]+48s .
 \end{aligned}
 \tag{V32.11}\label{c18v32:explicit}
\]
\end{theorem}
\begin{proof}
Coordinate geodesic telescoping gives
$|u(U)-u(U')|\le\mathfrak s_\beta\sum_e d_{S^3}(U_e,U'_e)$.
Apply \eqref{c18v32:distance} at $t=s$ to the ground factor.

Each incident-link gradient of a normalized plaquette trace has norm
at most one. Telescoping its entries bounds its change by the sum of
its four link distances. Every spatial link lies in four plaquettes,
so the trace-sum change is at most four times total fine-link distance.
Apply \eqref{c18v32:distance} at every integration time in the exact
Jacobian. Its contribution is at most $48\pi k\int_0^sF(t)\,dt$.
This proves the oscillation bound. The first explicit bound uses
$F(t)\le e^{16t}$; the second uses $F(t)\le P(t)$ and
\[
 \int_0^sP(t)\,dt=\frac3{320}[(80s+1)^4-1]+s.
\]
No normalizer was dropped from the full density.
\end{proof}

This improves finite oscillation, not the pointwise logarithmic
derivative. Dividing the finite bound by a vanishing displacement
does not give a derivative estimate. V31's derivative bound is unchanged.

\subsection{Improved native fibre comparison and correction costs}

Use V31's coordinates $W=\Theta(g,z)$,
$\chi_g(z)=\Phi_{-s}\Theta(g,z)$, and exact conditional density
$q_g(z)=h_s(\Theta(g,z))/\rho(g)$ relative to product Haar in $z$.
One coordinate replacement changes $m_j=1$ or $2$ flowed links:
an unused link, or a chain coordinate and its dependent last link.
The endpoints retain the same $g$. Consequently
\[
 \operatorname{osc}_{z_j}\log q_g\le\pi m_j\widehat d_s,\qquad
 \widehat\eta_j=3e^{-\pi m_j\widehat d_s},\qquad
 \widehat\eta_*=3e^{-2\pi\widehat d_s}.
 \tag{V32.12}\label{c18v32:local-floor}
\]
Both $\rho(g)$ and the further coordinate-conditional normalizer cancel
in this oscillation, for every fixed choice of the other coordinates.

Indeed, for a probability density $p$ on unit $S^3$ with log oscillation
at most $D$, put $m=\inf p$, $M=\sup p$. Testing the Haar mean in the
variance infimum and using Haar's gap three proves
\[
 \operatorname{Var}_{p\,dH}(f)
 \le M\operatorname{Var}_{dH}(f)
 \le\frac{M}{3m}\int|\nabla f|^2p\,dH
 \le\frac{e^D}{3}\int|\nabla f|^2p\,dH.
\]
This proves the local floors, with no extra factor two in the exponent.

Let $P_j^g$ be exact coordinate conditional averaging,
$T_g=\sum_j(I-P_j^g)$ with each coordinate updated at rate one, and
$L_g$ the original-metric conditional generator without $\kappa$.
V31's metric bound $(D_z\chi_g)^*D_z\chi_g\le\ell e^{32s}I$
and the summed coordinate inequalities give
\[
 L_g\ge\widehat\alpha_sT_g,\qquad
 \widehat\alpha_s=\frac3\ell e^{-32s-2\pi\widehat d_s}
 \ge\alpha_s^{\rm V31}.
 \tag{V32.13}\label{c18v32:comparison}
\]
The law, metric and operators have not changed.

Retain V31's actual full score residual $r$, local Poisson flux $Q$, and
\[
 Z_t=D_z\chi_g\,Q\int_0^t e^{-\tau T_g}r\,d\tau,\qquad A_t=A+Z_t .
\]
The auxiliary clock $t$ is distinct from the spatial-flow duration $s$
and physical time. V31's proof with the improved local inverse bounds
gives
\[
 \begin{aligned}
 Dc\,A_t&=I,\qquad r_t=e^{-tT_g}r,\\
 \mathbb E_{\nu_g}|Z_t v|^2
 &\le\frac{t}{2\widehat\alpha_s}\mathbb E_{\nu_g}|v\cdot r|^2.
 \end{aligned}
 \tag{V32.14}\label{c18v32:correction}
\]
The local inverses and flux are the same; no new whole-fibre inverse
is assumed in constructing the correction.

Put $\Sigma_g=\mathbb E_{\nu_g}[rr^*]$ and define
\[
 v^*M_{<\varepsilon}(g)v
 =\langle v\cdot r,\,
 {\bf1}_{(0,\varepsilon)}(T_g)T_g^{-1}(v\cdot r)\rangle_{\nu_g}.
\]
For $R^{A_t}_{ij}(g)=\langle(r_t)_i,L_g^{-1}(r_t)_j\rangle_{\nu_g}$,
inverse comparison and spectral splitting give
\[
 R^{A_t}(g)\preceq
 \widehat\alpha_s^{-1}M_{<\varepsilon}(g)
 +\frac{e^{-2t\varepsilon}}{\widehat\alpha_s\varepsilon}\Sigma_g.
 \tag{V32.15}\label{c18v32:low}
\]
Only $T_g$ is split; it need not commute with $L_g$.
The residual is centred, and inverses exclude constants.
At a fixed regulator the smooth positive compact density makes these
inverses legitimate. The uniform low-frequency moment remains unpaid.

At $s=\ell^2$ the explicit bound yields
\[
 -\log\widehat\alpha_{\ell^2}
 \le\log(\ell/3)+32\ell^2
 +O\bigl(\sqrt\beta(1+\ell^2)^3+(1+\ell^2)^4\bigr)
 \tag{V32.16}\label{c18v32:scale}
\]
with absolute implied constant. The floor thus no longer has
double-exponential loss in $\ell^2$ at fixed $\beta$: its negative
logarithm has a polynomial upper envelope. This is not a polynomial
coupling bound or a cofinal closure. Diagnostics compare negative
logarithms rather than underflowed exponentials. V30's complete-score
and variable-bank derivative payments remain unchanged.

\subsection{A matching-input compact two-well test}

The following tests curvature and one-link score inputs alone.
Its scalar potential is engineered, not a Wilson potential or a
gauge-invariant many-link Yang--Mills model.

\begin{theorem}[An actual ground state with an exponentially slow mode]
\label{c18v32:toy}
On unit $S^3$, put $z(q)=q_0$, let $a\ge4$, and use normalized Haar $dH$.
Set
\[
 \begin{aligned}
 V_a&=4a(1-z^2)(2+az^2),& H_a&=-\Delta+V_a,\\
 \Omega_a&=Z_a^{-1/2}e^{az^2},&
 Z_a&=\int e^{2az^2}\,dH,\qquad \beta_a=\tfrac12a^2+a .
 \end{aligned}
 \tag{V32.17}\label{c18v32:toy-def}
\]
Then $V_a\ge0$, $\Omega_a$ is the normalized unique positive ground
state, its energy is $6a$, and
\[
 \begin{aligned}
 \operatorname{Hess}V_a&\le16\beta_a g,&
 \operatorname{Hess}\log\Omega_a&\ge-2ag
                                  \ge-\sqrt{8\beta_a}\,g,\\
 |\nabla\log\Omega_a|&\le a
                      \le\min\{2\beta_a,\pi\sqrt{8\beta_a}\}.
 \end{aligned}
 \tag{V32.18}\label{c18v32:toy-inputs}
\]
Nevertheless the gap $\lambda_a$ of
$-\operatorname{div}_{\Omega_a^2dH}\nabla$ satisfies
\[
 0<\lambda_a\le2\sqrt2\,e^6a^2e^{-2a}.
 \tag{V32.19}\label{c18v32:toy-gap}
\]
It is also the gap of $H_a$ above its ground energy, without an extra
time prefactor.
\end{theorem}
\begin{proof}
The sphere identities $|\nabla z|^2=1-z^2$,
$\operatorname{Hess}z=-zg$ and $\Delta z=-3z$ imply, for $u_a=az^2$,
\[
 \Delta u_a=2a-8az^2,\quad
 |\nabla u_a|^2=4a^2z^2(1-z^2),\quad
 V_a=\Delta u_a+|\nabla u_a|^2+6a.
\]
Hence $H_a\Omega_a=6a\Omega_a$. Integration by parts gives, for every
smooth $f$,
\[
 \int\bigl(|\nabla(\Omega_a f)|^2+V_a|\Omega_a f|^2\bigr)\,dH
 -6a\int|\Omega_a f|^2\,dH
 =\int\Omega_a^2|\nabla f|^2\,dH.
 \tag{V32.20}\label{c18v32:ground-transform}
\]
Every smooth test is $\Omega_a f$, since $\Omega_a$ is smooth and
strictly positive on a compact connected manifold. Thus $6a$ is the
lowest eigenvalue and equality occurs only for constant $f$.
Unitary multiplication by $\Omega_a$ identifies the gaps.
Their strict positivity at fixed $a$ follows also by comparison with
Haar using the positive minimum and finite maximum of $\Omega_a^2$.

Directly $\operatorname{Hess}u_a=2a\,dz\otimes dz-2az^2g\ge-2ag$
and $|\nabla u_a|\le a$. For $x=z^2$, the potential Hessian has two
tangential eigenvalues and one meridional eigenvalue
\[
 \begin{aligned}
 h_{\rm tan}(x)&=(16a-8a^2)x+16a^2x^2,\\
 h_{\rm mer}(x)&=8a^2-16a+(32a-64a^2)x+64a^2x^2 .
 \end{aligned}
\]
With $K_a=8a^2+16a=16\beta_a$,
\[
 \begin{aligned}
 K_a-h_{\rm tan}(x)&=(1-x)(16a^2x+8a^2+16a),\\
 K_a-h_{\rm mer}(x)&=(1-x)(32a+64a^2x).
 \end{aligned}
\]
Both are nonnegative on $[0,1]$; pole values agree by continuity.
This proves the curvature assertions.

For the gap put $\delta=a^{-1/2}$ and take the odd $H^1$ test
\[
 f_a(z)=
 \begin{cases}
 -1,&z\le-\delta,\\
 z/\delta,&|z|<\delta,\\
 1,&z\ge\delta.
 \end{cases}
\]
Its mean is zero for the even density. Smooth odd approximations
justify the Rayleigh test. The $z$-pushforward of Haar is
$(2/\pi)\sqrt{1-z^2}\,dz$, so the unnormalized energy obeys
\[
 \mathcal N_a=\frac2{\pi\delta^2}
   \int_{-\delta}^{\delta}(1-z^2)^{3/2}e^{2az^2}\,dz
 \le\frac{4e^2}{\pi}\sqrt a .
\]
For the squared norm restrict to
$|z|\in[1-1/a,1-1/(2a)]$. For $a\ge4$ these two intervals are outside
$(-\delta,\delta)$, hence $f_a^2=1$. Their combined length is $1/a$,
$\sqrt{1-z^2}\ge1/\sqrt{2a}$, and $2az^2\ge2a-4$. Thus
\[
 \mathcal D_a=\int f_a^2e^{2az^2}\,dH
 \ge\frac{\sqrt2}{\pi a^{3/2}}e^{2a-4}.
\]
The normalizer cancels in $\mathcal N_a/\mathcal D_a$, proving the bound.
\end{proof}

Since $\beta_a\sim a^2/2$, this family excludes a universal
polynomial-in-$\beta$ positive gap lower bound from compactness,
positivity of the actual ground state, the upper potential Hessian
bound and one-link score bounds alone. It satisfies these inputs while
its gap has an exponentially small upper bound on scale
$\sqrt{\beta_a}$. Potential and verified ground state were constructed
together; no arbitrary weight replaced the vacuum of an already fixed
Wilson Hamiltonian.

This is \emph{not} a counterexample to the Yang--Mills gap, evidence that
a native fibre has two wells, or a bound on the actual $T_g$.
It excludes a generic-curvature shortcut. Additional Wilson-specific
structure or source-specific suppression of low spectral modes is
essential for that proposed shortcut.

\subsection{Progress and outstanding obligations}

V32 improves native finite-change and comparison bounds without
changing the law, metric or gauge content. The unresolved
low-frequency source moment in \eqref{c18v32:low} requires native
interaction or source structure, not another compactness argument.

The regression script
\texttt{scripts/verify\_ym\_c18\_v32\_compact\_flow\_oscillation.py}
checks cubic-link kernels, polynomial integration, exact sphere
identities and two-well trial integrals. These and the rerun V31
diagnostics are not proof-assistant certification or a many-link
vacuum computation. V31 is retained exactly except its title;
47 other sources are unchanged. Only V32 replaces V31.
Only \texttt{.tex} files are in \texttt{latex\_notes}; builds stay outside.

Checkpoints remain V35 (companions), V40 (literature), V100 (archive).
Cofinal low-frequency control, projection stability, conditional/coarse
variance, physical return and scale transport are unpaid.
The interacting four-dimensional continuum and positive physical
Yang--Mills mass gap remain \emph{unproved}.
% END CYCLE 18 VERSION 32 APPEND
% BEGIN CYCLE 18 VERSION 33 APPEND
\clearpage
\section[V33: exact internal-gauge reduction of the native score]{V33: exact internal-gauge reduction of the native score}
\label{c18v33:context}

V32 improves a finite-change estimate but leaves the full conditional
low-frequency inverse moment unpaid. This note uses an exact symmetry
of that \emph{particular source}. Gauge transformations equal to the
identity at every coarse vertex preserve each sampled fibre. The
full score residual is invariant under this group, so its native
inverse cost lies entirely in the invariant subspace. A rooted forest
then gives an exact Haar disintegration and a sharper comparison on
the remaining coordinates. The original fine metric, coarse output
and conditional law are retained.

The underlying tree-Haar change of variables is not new: archived
f6/V96 and f15/V10,V45 already establish it in their respective
conditional settings. Archived f17/V95 also proves a native one-link
Ward/Poisson estimate and warns that gauge transport does not control
the physical shape channel. We do not re-present either as a new
general principle. Here the new application is to the complete
flowed sampled fibre, its full score, its induced metric, and an
implementable local-inverse correction on its exact reduced law.

For context, the primary abstract of N.~E.~Ligterink, N.~R.~Walet and
R.~F.~Bishop, \emph{Towards a Many-Body Treatment of Hamiltonian Lattice
SU(N) Gauge Theory},
\href{https://arxiv.org/abs/hep-lat/0001028}{arXiv:hep-lat/0001028},
describes maximal-tree variables for the colourless sector.
Only this abstract was checked in this targeted review; none of its
spectral conclusions is imported. The needed identities are proved
below. The broad literature checkpoint remains V40.

\subsection{A forest compatible with every retained coarse chain}

Retain $N=\ell n$, $\ell\ge2$, $n\ge5$ and
$c=\mathfrak b_\ell\circ\Phi_s$, with actual law $\nu=\Omega^2dU$.
Write $V_{\rm f},E_{\rm f}$ for fine vertices and stored links, and
$V_{\rm c},E_{\rm c}$ for coarse ones. Identify $V_{\rm c}$ with the
fine vertices whose three coordinates are divisible by $\ell$.
For each coarse edge, its sampled chain has $\ell$ links; these
chains have disjoint interiors and disjoint link sets.

Let $\mathcal F_0$ contain the first $\ell-1$ links of every chain
and let $\mathcal B$ contain their last links. The components of
$\mathcal F_0$ are trees, each rooted at one coarse vertex.
Extend $\mathcal F_0$ to a forest $\mathcal F$ covering every fine
vertex with exactly one coarse root in each component, without
using any edge of $\mathcal B$. This can be done by repeatedly
joining an unassigned vertex to an assigned one. As long as an
unassigned vertex remains, connectedness supplies an edge across
that partition. Every edge of $\mathcal B$ already has both
endpoints assigned, so such an extension never needs it.
No step joins two assigned components.

Put $\mathcal C=E_{\rm f}\setminus(\mathcal F\cup\mathcal B)$.
These are the free chord coordinates. The exact counts are
\[
 |\mathcal F|=|V_{\rm f}|-|V_{\rm c}|,\quad
 |\mathcal B|=|E_{\rm c}|,\quad
 |\mathcal C|=2(|V_{\rm f}|-|V_{\rm c}|).
 \tag{V33.1}\label{c18v33:counts}
\]
The last equality uses $|E_{\rm f}|=3|V_{\rm f}|$ and
$|E_{\rm c}|=3|V_{\rm c}|$. The forest is fixed independently of
the configuration and coupling; no changing-rank gauge slice is used.

Let
\[
 \mathcal H_0=\{h\in SU(2)^{V_{\rm f}}:h_x=I\ (x\in V_{\rm c})\}.
\]
For any flowed configuration $W$, recursively choose $k_x$, with
$k_x=I$ at the roots, so that
\[
 W^0_{xy}=k_xW_{xy}k_y^{-1}=I\quad(xy\in\mathcal F).
 \tag{V33.2}\label{c18v33:gauge}
\]
If a tree is traversed against a stored orientation, use the inverse
link in the recursion. On a sampled chain every prefix edge is in
$\mathcal F$, so its last transformed link is precisely its
coarse product $g_e$. Write $y_f=W^0_f$ for $f\in\mathcal C$.

\begin{proposition}[Global native factorization, with all coarse links retained]
\label{c18v33:factor}
The map $W\leftrightarrow(k,g,y)$ is a smooth global bijection, with
\[
 dW=dk\,dg\,dy,\qquad
 W_{xy}=k_x^{-1}W^0_{xy}(g,y)k_y ,
 \tag{V33.3}\label{c18v33:haar}
\]
using normalized product Haar on all three factors. On the fibre
$c(U)=g$, define
\[
 \begin{aligned}
 \chi_g(k,y)&=\Phi_{-s}\bigl(k^{-1}\cdot W^0(g,y)\bigr),\\
 p_g(y)&=\frac{h_s(W^0(g,y))}{\rho(g)},&
 \rho(g)&=\int h_s(W^0(g,y))\,dy .
 \end{aligned}
 \tag{V33.4}\label{c18v33:law}
\]
Here $h_s$ is V32's exact pushed native density, including the
inverse-flow Jacobian. Then
\[
 \nu_g=(\chi_g)_\#\bigl(dk\,p_g(y)\,dy\bigr).
 \tag{V33.5}\label{c18v33:factor-law}
\]
The $\mathcal H_0$ action is free; its coordinates are exactly Haar
and independent of $y$.
\end{proposition}
\begin{proof}
Successive tree-edge Haar translations replace the forest links by
the nonroot vertex variables $k_x$. Each remaining link undergoes
a left and right Haar translation with these variables fixed.
The last-chain links become the independent $g$ variables; all
other nonforest links are the independent $y$ variables.
The explicit inverse in \eqref{c18v33:haar} proves bijectivity
and smoothness, and Fubini proves the measure identity.

The ground state is gauge invariant by positivity, uniqueness and
gauge invariance of its Hamiltonian. The spatial flow intertwines
gauge transformations, which are product-metric isometries
preserving Haar. Its volume Jacobian, and hence $h_s$, is gauge
invariant. Thus $h_s(W)=h_s(W^0(g,y))$ is independent of $k$.
Disintegrate \eqref{c18v33:haar} over $g$ to obtain
\eqref{c18v33:factor-law}. Every density is smooth positive,
so the displayed conditional version is defined for every $g$.
Finally, a stabilizing element of $\mathcal H_0$ is the identity
successively along the forest from each root, proving freeness.
\end{proof}

Only the internal, root-fixed group was factored. Transformations
at the coarse roots and all their stabilizers remain in $g,y$.
There is no quotient by a possibly singular full coarse orbit.
Nor is there an extra orbit-volume or coarea determinant to insert
into \eqref{c18v33:factor-law}; the metric still enters the energy.

\subsection{The score lies in a reducing invariant subspace}

Let $A$ be V28's full equivariant lift, or another smooth equivariant
lift with $Dc\,A=I$. In the fixed coarse left frames put
\[
 S_i=\operatorname{div}_\nu A_i,\qquad
 \overline S(g)=\mathbb E_{\nu_g}S,\qquad r=S-\overline S(c).
 \tag{V33.6}\label{c18v33:score}
\]
Let $\Pi_0$ average over $\mathcal H_0$ on a fibre.
Let $L_g=-\operatorname{div}_{\nu_g}\nabla_{\rm fib}$ use the original
fine induced metric and no factor $\kappa$.

\begin{lemma}[Exact invariant restriction of the source inverse cost]
\label{c18v33:invariant}
Every component of $r$ is $\mathcal H_0$ invariant.
The orthogonal projection $\Pi_0$ commutes with $L_g$ and its spectral
resolvents. Writing $r(\chi_g(k,y))=\bar r_g(y)$, the unitary
\[
 J_g:L^2(p_gdy)\longrightarrow\operatorname{Ran}\Pi_0,\qquad
 (J_gf)(\chi_g(k,y))=f(y)
\]
intertwines a positive self-adjoint reduced generator
$\bar L_g=J_g^*L_gJ_g$ with the invariant restriction of $L_g$.
For every coarse bank $v$ fixed on the fibre,
\[
 \langle v\cdot r,L_g^{-1}(v\cdot r)\rangle_{\nu_g}
 =\langle v\cdot\bar r_g,\bar L_g^{-1}(v\cdot\bar r_g)\rangle_{p_g}.
 \tag{V33.7}\label{c18v33:inverse-equality}
\]
Both inverses are centred.
\end{lemma}
\begin{proof}
For $h\in\mathcal H_0$, the coarse transformation is the identity:
$c(hU)=c(U)$. Equivariance therefore sends each fixed coarse column
$A_i(U)$ to the pushforward $A_i(hU)$. Since $h$ preserves the fine
metric and $\nu$, its weighted divergence obeys $S_i(hU)=S_i(U)$.
The conditional centering is also invariant. This proves the
source assertion without imposing an extra Gauss condition.

The same transformations preserve the fibre, its induced metric and
its conditional law. Their pullbacks preserve its closed Dirichlet
form, hence commute with the associated self-adjoint generator and
resolvents. Averaging gives the reducing projection $\Pi_0$.
The factorization \eqref{c18v33:factor-law} proves unitarity of $J_g$.
At fixed regulator the fibre is a connected compact product in the
coordinates above, with positive smooth density and metric; its
centred Poincar\'e constant is finite. Functional calculus on the
reducing subspace proves \eqref{c18v33:inverse-equality}.
\end{proof}

Thus noninvariant internal-gauge modes have \emph{exactly zero}
source spectral weight. This does not remove any invariant physical
mode, nor prove a gap for the remaining ones.

\subsection{The original metric and a one-link conditional comparison}

Use product-Haar orthonormal frames in $(k,y)$. The metric
$M_g=(D\chi_g)^*D\chi_g$ has blocks $M_{kk}$, $M_{ky}$ and $M_{yy}$.
Write $B_g(k,y)=D_y\chi_g$. At fixed $k,g$, changing the free chords
changes distinct fine links by isometries before inverse flow.
Therefore
\[
 B_g^*B_g=M_{yy}\preceq e^{32s}I .
 \tag{V33.8}\label{c18v33:chord-metric}
\]
In particular the factor $\ell$ from V31's chain-coordinate metric
is absent.

For functions depending only on $y$, the exact reduced energy is
\[
 \begin{aligned}
 \overline{\mathcal E}_g(f,f)
 &=\int(\nabla_yf)^*a_g(y)\nabla_yf\,p_g(y)\,dy,\\
 a_g(y)&=\int [M_g(k,y)^{-1}]_{yy}\,dk .
 \end{aligned}
 \tag{V33.9}\label{c18v33:reduced-energy}
\]
One may retain this average even though gauge isometries make the
relevant block independent of the orbit coordinate in compatible
frames. The Schur complement gives
\[
 [M_g^{-1}]_{yy}
 =(M_{yy}-M_{yk}M_{kk}^{-1}M_{ky})^{-1}
 \succeq M_{yy}^{-1}\succeq e^{-32s}I .
 \tag{V33.10}\label{c18v33:schur}
\]
Thus the reduced kinetic coefficient is not declared to be the
identity merely because the reference measure factors.

Use exactly V32's $\widehat d_s$. A replacement of one $y_f$ at
fixed $g$ and other chords changes just one link of $W^0$.
Its coordinate conditional density consequently has log oscillation
at most $\pi\widehat d_s$. Let $\bar P_f^g$ be its exact averaging
projection, and put
\[
 \bar T_g=\sum_{f\in\mathcal C}(I-\bar P_f^g),\qquad
 \eta_s^{\rm red}=3e^{-\pi\widehat d_s},\qquad
 \alpha_s^{\rm red}=3e^{-32s-\pi\widehat d_s}.
 \tag{V33.11}\label{c18v33:red-alpha}
\]
Every chord updates at rate one.

\begin{theorem}[Metric-faithful comparison for the actual invariant source]
\label{c18v33:comparison}
On $L^2(p_gdy)$,
\[
 \bar L_g\succeq\alpha_s^{\rm red}\bar T_g.
 \tag{V33.12}\label{c18v33:form}
\]
Consequently the native cost in \eqref{c18v33:inverse-equality} is at most
$(\alpha_s^{\rm red})^{-1}
\langle v\cdot\bar r_g,\bar T_g^{-1}(v\cdot\bar r_g)\rangle_{p_g}$.
Compared with V32's scalar prefactor,
\[
 \alpha_s^{\rm red}
   =\ell e^{\pi\widehat d_s}\widehat\alpha_s .
 \tag{V33.13}\label{c18v33:ratio}
\]
This compares prefactors, not the spectra of the two different
auxiliary resampling operators.
\end{theorem}
\begin{proof}
V32's density-oscillation comparison with Haar's gap three gives
the coordinate floor $\eta_s^{\rm red}$. Sum the corresponding
conditional variance inequalities and combine with
\eqref{c18v33:reduced-energy}--\eqref{c18v33:schur}.
This proves the form inequality. The centred inverse inequality
follows from the variational formula for the inverse of a positive
operator, at the fixed regulator where both centred inverses exist.
The ratio is immediate from V32's
$\widehat\alpha_s=(3/\ell)e^{-32s-2\pi\widehat d_s}$.
\end{proof}

In general $\bar T_g$ is not the conjugate of V31's $T_g$ restricted
to invariants: the forest changes which conditional moves are used.
No ordering between their low spectral moments has been proved.
The density $p_g$ is neither a product over chords nor a bare Wilson
density, and its conditionals can depend on remote chords.

\subsection{A reduced local-flux correction with full coarse output}

Let $\bar L_{f,g}$ be the weighted one-chord Laplacian of $p_g$,
with all other chords fixed. Its centred inverse has norm at most
$(\eta_s^{\rm red})^{-1}$. Define
\[
 (\bar Q f)_e=\nabla_e\bar L_{e,g}^{-1}(I-\bar P_e^g)f.
\]
Coordinate integration by parts, with the actual conditional
normalizer retained, gives
\[
 \operatorname{div}_{p_g}\bar Qf=-\bar T_gf,\qquad
 \|\bar Qf\|_{L^2(p_g)}^2
 \le(\eta_s^{\rm red})^{-1}\langle f,\bar T_gf\rangle_{p_g}.
 \tag{V33.14}\label{c18v33:flux}
\]
For $F_t(\lambda)=\int_0^t e^{-\tau\lambda}\,d\tau$, set on the
\emph{original} fibre
\[
 Z_t^{\rm red}(\chi_g(k,y))
    =B_g(k,y)\bar QF_t(\bar T_g)\bar r_g(y),\qquad
 A_t^{\rm red}=A+Z_t^{\rm red}.
 \tag{V33.15}\label{c18v33:correction}
\]

\begin{theorem}[Native correction without a whole-fibre inverse]
\label{c18v33:corrected}
For every finite $t\ge0$,
\[
 \begin{aligned}
 Dc\,A_t^{\rm red}&=I,\\
 r_t^{\rm red}(\chi_g(k,y))
       &=e^{-t\bar T_g}\bar r_g(y),\\
 \mathbb E_{\nu_g}|Z_t^{\rm red}v|^2
       &\le\frac{t}{2\alpha_s^{\rm red}}
                               \mathbb E_{\nu_g}|v\cdot r|^2 .
 \end{aligned}
 \tag{V33.16}\label{c18v33:cost}
\]
Here $r_t^{\rm red}$ is the conditionally centred score of the
new lift, and $t$ is auxiliary time, not spatial-flow or physical time.
\end{theorem}
\begin{proof}
The columns of $B_g$ vary only fibre coordinates, so $Dc\,B_g=0$.
In $(g,k,y)$ coordinates $Z_t^{\rm red}$ has no $g$ or $k$ component.
The full law has density $\rho(g)p_g(y)$ relative to $dg\,dk\,dy$.
Thus its full weighted divergence is exactly the conditional
$y$-divergence, not a product-metric divergence substituted for the
original metric. By \eqref{c18v33:flux} it equals
$-(I-e^{-t\bar T_g})\bar r_g$. Its conditional mean is zero,
which proves the first two identities.

For the norm use \eqref{c18v33:chord-metric} and
\eqref{c18v33:flux}, then the scalar inequality
$(1-e^{-x})^2\le x/2$ for $x\ge0$ proved in V31. The spectral
multiplier is $\lambda F_t(\lambda)^2\le t/2$.
This proves the cost, including the factor
$e^{32s}/\eta_s^{\rm red}=(\alpha_s^{\rm red})^{-1}$.

All one-chord conditional coefficients are smooth positive on
compact products. Parameter-dependent centred elliptic inverses
and the finite-time exponential of the bounded conditional
projections preserve smooth sources at fixed regulator, just as
in V31. No uniform whole-fibre inverse enters the construction.
\end{proof}

The construction remains gauge equivariant. Under a fixed fine
gauge transformation, the reduced chord from $x$ to $x'$ transforms
by fixed left and right multiplications associated with the roots
of their forest components. The coarse links have their usual
endpoint transformation. These are Haar-preserving product
isometries. The conditional projections and local inverse
Laplacians intertwine them, and $D_y\chi_g$ transforms covariantly.
The reconstructed correction thus retains every coarse column.
V30's averaged full-score bounds and variable-bank derivative costs
remain available unchanged.

For $\bar\Sigma_g=\mathbb E_{p_g}[\bar r_g\bar r_g^*]$ and the
matrix $\bar M_{<\varepsilon}$ defined by
\[
 v^*\bar M_{<\varepsilon}(g)v
 =\langle v\cdot\bar r_g,\,
 {\bf1}_{(0,\varepsilon)}(\bar T_g)\bar T_g^{-1}
                  (v\cdot\bar r_g)\rangle_{p_g},
\]
the native corrected inverse cost obeys
\[
 R^{A_t^{\rm red}}(g)\preceq
 (\alpha_s^{\rm red})^{-1}\bar M_{<\varepsilon}(g)
 +\frac{e^{-2t\varepsilon}}{\alpha_s^{\rm red}\varepsilon}
                              \bar\Sigma_g.
 \tag{V33.17}\label{c18v33:remaining}
\]
Use the exact invariant restriction and inverse comparison first,
then spectral calculus only for $\bar T_g$; no commutation with
$\bar L_g$ is assumed. The remaining low-frequency source moment
is now entirely in reduced chord variables, but is still unpaid.

\subsection{Gauge invariance alone does not exclude the two-well test}

There is an exact scope check on a single oriented cycle of $m\ge3$
links. Let $P(U)=U_1\cdots U_m$, $z=\tfrac12\operatorname{Tr}P$,
and use V32's $V_a(z)$ and $\Omega_a(z)$ with $a\ge4$. Define
\[
 H_{m,a}=-\sum_{j=1}^m\Delta_j+mV_a(z),\qquad
 \Omega_{m,a}(U)=Z_a^{-1/2}e^{az^2}.
 \tag{V33.18}\label{c18v33:cycle}
\]
These are gauge invariant on the cycle. The potential is
\emph{not} the standard linear Wilson plaquette potential.

For any smooth function $f(P)$, each link Laplacian acts as the
round group Laplacian on $f$, since fixed left/right multiplications
are isometries. Therefore $H_{m,a}\Omega_{m,a}=6ma\,\Omega_{m,a}$.
The positive ground transform on the connected full link product
proves this is its unique ground state. Haar pushes $P$ to Haar,
so the normalization is exactly V32's $Z_a$.
Its clipped odd test $f_a(z)$ is invariant under conjugation of $P$,
hence satisfies all vertex Gauss constraints. Each link contributes
the same energy as the one-sphere test. The Rayleigh principle on
the physical gauge-invariant subspace gives
\[
 0<\lambda^{\rm phys}_{m,a}
 \le2\sqrt2\,e^6\,m a^2e^{-2a}.
 \tag{V33.19}\label{c18v33:cycle-gap}
\]
This is a proved gauge-invariant compact example, not an assertion
about a Wilson vacuum, a native fibre, or the spectral weight of
our actual source. It shows why the exact internal-gauge reduction
does not by itself dismiss invariant tunnelling.

\subsection{Verification and next obligation}

The new result is the invariant restriction of the actual full score,
its exact sampled-fibre forest law, the paid Schur-complement kinetic
comparison, and a full-output correction using one-link reduced
conditionals. General forest-Haar factorization and one-link Ward
transport are credited to the archives, not counted twice.
The next decisive task is a native bound on
$\bar M_{<\varepsilon}$ in \eqref{c18v33:remaining}, or a direct bound
on the same source under $\bar L_g$ using Wilson-specific structure.

The script
\texttt{scripts/verify\_ym\_c18\_v33\_rooted\_fibre.py}
checks forest counts and coordinate reconstruction, exact nonabelian
gauge identities, nonorthogonal metric Schur complements, invariant
resolvent reduction and local-flux corrections. Fixtures are not a
many-link Wilson vacuum computation or proof-assistant certification.
V32's regressions are rerun. Its body is preserved except the title,
47 other sources are unchanged, and only active V33 replaces V32.
The notes folder contains only \texttt{.tex}; builds remain outside.

Checkpoints remain V35 (companions), V40 (broad literature) and V100
(archive restart). Uniform native low-frequency control, projection
stability, conditional/coarse variance, physical scale transport,
the interacting four-dimensional continuum and its positive
Yang--Mills mass gap remain unproved.
% END CYCLE 18 VERSION 33 APPEND
% BEGIN CYCLE 18 VERSION 34 APPEND
\clearpage
\section[V34: score-preserving gauge projection and exact kinetic gain]{V34: score-preserving gauge projection and exact kinetic gain}
\label{c18v34:context}

V33 retains the internal-gauge/chord cross terms but estimates their
effect only by a scalar floor. We now evaluate the improvement they
actually permit. There are two distinct inverses. The first is a
finite matrix on gauge parameters pinned at all coarse vertices;
a rooted-path estimate pays for it uniformly in the volume. The
second is the conditional generator inverse on physical functions;
that is the unresolved mass-gap input. They must not be identified.

The positive result is an exact, globally smooth projection of an
equivariant lift away from the internal-gauge orbits. It preserves
every coarse output and the entire weighted score, including variable
coarse banks, while decreasing its original-metric norm. An explicit
formula then identifies both the kinetic gain and nonconstant physical
plaquette directions on which that gain is zero.

This is not a new claim of tree-Haar factorization or flat curl
coercivity. The former is credited in V33; archived f15/V11 concerns
a flat quadratic curl estimate with an averaging constraint, rather
than the gauge-parameter matrix below. Archived f16.2/V71 already
warns that physical horizontality does not imply curvature control.
The new statements concern the current nonlinear native lift, its
exact divergence, and the full induced kinetic form.

A targeted primary-source check of T.~Li, Y.-Y.~Li, X.~Wang and
H.~Xing, \emph{Efficient Quantum Simulations of Yang--Mills theory
with Maximal-tree Gauge},
\href{https://arxiv.org/abs/2608.27267}{arXiv:2608.27267v1}
(submitted 27 August 2026), inspected the abstract and the
Hamiltonian/gauge-fixing discussion, especially equations (3)--(8).
Its charge--inverse--charge kinetic term is a useful structural
comparison. Its open-boundary gauge-field formulation, residual
global constraint and quantum simulation estimates are not imported
as theorems about our compact, periodically sampled native law.
All identities used here are proved below. The broad checkpoint
remains V40.

\subsection{A uniformly paid inverse on pinned gauge parameters}

Retain V33's $N=\ell n$, $\ell\ge2$, $n\ge5$, sampled map
$c=\mathfrak b_\ell\circ\Phi_s$ and actual law $\nu=\Omega^2dU$.
Choose the forest in that note more specifically. Write each fine
vertex uniquely as $x=\ell X+r$, with
$r\in\{0,\ldots,\ell-1\}^3$. For $r\ne0$, its parent is obtained
by decreasing the first nonzero component of $r$ by one.
The stored positive link from that parent to $x$ belongs to
$\mathcal F$. Each component is one cell, contains $\ell^3$
vertices, and has maximal root-path length $3(\ell-1)$.
It includes the first $\ell-1$ links of every sampled chain and
excludes every last-chain link in $\mathcal B$.
The counts and exact coordinates of V33 therefore apply unchanged.

Use the original fine product metric and left-multiplication frames.
A gauge parameter $\xi=(\xi_x)$ vanishes at every coarse root.
Its infinitesimal gauge field and Gram matrix are
\[
 (D_U\xi)_{xy}=\xi_x-\operatorname{Ad}_{U_{xy}}\xi_y,
 \qquad K_U=D_U^*D_U,\qquad
 C_\ell=3(\ell-1)(\ell^3-1).
 \tag{V34.1}\label{c18v34:pinned-def}
\]
Stars denote adjoints for the unweighted Euclidean vertex/link
metrics. In particular $K_U$ is not $L_g$, $\bar L_g$ or a physical
Hamiltonian.

\begin{lemma}[Pinned gauge inverse, at every configuration]
\label{c18v34:pinned}
Uniformly in the volume, coupling and fine configuration,
\[
 C_\ell^{-1}I\preceq K_U\preceq12I,\qquad
 \|K_U^{-1}\|\le C_\ell,\qquad
 \|K_U^{-1}D_U^*\|\le\sqrt{C_\ell}.
 \tag{V34.2}\label{c18v34:pinned-bound}
\]
\end{lemma}
\begin{proof}
Along the forest path from a root to $x$, solve
$\xi_a-\operatorname{Ad}_{U_{ab}}\xi_b=(D_U\xi)_{ab}$
successively for $\xi_b$. The root value is zero and every adjoint
transport is orthogonal. Cauchy--Schwarz gives
\[
 |\xi_x|^2\le d(x)\sum_{e\ {\rm on\ the\ path}}|(D_U\xi)_e|^2 .
\]
Sum first within a cell: there are $\ell^3-1$ nonroot vertices and
$d(x)\le3(\ell-1)$. Sum over the disjoint forest components.
This yields $\|\xi\|^2\le C_\ell\|D_U\xi\|^2$.
Conversely each vertex has six incident links, and
$|\xi_x-\operatorname{Ad}_{U_{xy}}\xi_y|^2
\le2|\xi_x|^2+2|\xi_y|^2$ gives the upper bound.
Finally
$(K_U^{-1}D_U^*)(K_U^{-1}D_U^*)^*=K_U^{-1}$.
No flatness, vacuum estimate or probabilistic mixing enters.
\end{proof}

Define the original-metric orthogonal projection
\[
 P_U=I-D_UK_U^{-1}D_U^*.
 \tag{V34.3}\label{c18v34:projection}
\]
It is globally smooth and gauge equivariant. The pinning condition
is preserved by conjugation under arbitrary fine gauge transformations,
so this assertion includes transformations at the coarse roots.
The identities $P_U^2=P_U=P_U^*$ and $P_UD_U=0$ are direct.
Since root-fixed gauges preserve $c$, also
\[
 Dc(U)D_U=0,\qquad Dc(U)P_U=Dc(U).
 \tag{V34.4}\label{c18v34:output}
\]
We have not taken a quotient by the possibly nonfree coarse-root
action, and have not used a changing-rank pseudoinverse.

\subsection{Projection preserves the full score exactly}

\begin{theorem}[Score-preserving internal horizontalization]
\label{c18v34:score}
Let $V$ be a smooth vector field on the full fine configuration
space, equivariant under the root-fixed group $\mathcal H_0$.
Then
\[
 \operatorname{div}_{\nu}(PV)=\operatorname{div}_{\nu}V .
 \tag{V34.5}\label{c18v34:div}
\]
Consequently, for any smooth equivariant full lift $A$ with
$Dc\,A=I$, the lift $A^{\rm h}=PA$ satisfies
\[
 \begin{aligned}
 Dc\,A^{\rm h}&=I,& S^{A^{\rm h}}&=S^A,& r^{A^{\rm h}}&=r^A,\\
 (A^{\rm h})^*A^{\rm h}
 &=A^*A-(D_U^*A)^*K_U^{-1}(D_U^*A)& &\preceq A^*A .
 \end{aligned}
 \tag{V34.6}\label{c18v34:lift}
\]
Scores and residuals use the actual conditional centering of V33.
For every smooth coarse bank $v(g)$, the full divergences of
$PA\,v(c)$ and $A\,v(c)$ also agree, with no omitted derivative of $v$.
\end{theorem}
\begin{proof}
The difference $W=(I-P)V$ is a smooth equivariant field tangent to
the internal-gauge orbit. In the global coordinates $(k,g,y)$ of
V33, its only components are in $k$. The root-fixed action sends
$k$ to $kh^{-1}$ and leaves $g,y$ fixed. Equivariance therefore
makes $W$ a linear combination of the fields
\[
 \left.\frac{d}{dt}\right|_{t=0}\exp(tE_{x,a})k
\]
with coefficients depending on $g,y$ but not on $k$.
These fields have zero divergence for product Haar $dk$.
The full native measure is exactly
$\rho(g)p_g(y)\,dg\,dk\,dy$, whose density is independent of $k$.
The weighted divergence of $W$ is therefore zero. This computation
is a measure-divergence identity in exact coordinates; it does not
replace the fine metric by a product metric.

Apply this identity to each fixed coarse column. Equation
\eqref{c18v34:output} gives the output identity. Equality of the
scores gives equality of their conditional means and residuals.
The Gram formula follows by expanding $A^*PA$.
Finally $A\,v(c)$ is still internally equivariant, since $c$ is
internally invariant, so the same proof applies to this vector field.
Equivalently the extra product-rule term is unchanged because
$Dc\,PA=Dc\,A$.
\end{proof}

The equivariance hypothesis is essential: the Haar divergence of
an arbitrary orbit-tangent field with $k$-dependent coefficients
need not vanish. The identity pays the derivatives of $P$ in this
particular divergence exactly; it is not a bound on every derivative
of $P$ or a substitute for projection-energy stability.
V30's variable-bank score estimates are preserved, not discarded.
For V28's weighted-optimal lift only the unweighted fine-metric
norm improvement is asserted; optimality in its original
$Q$-weighted objective is not claimed for $PA$.
Nor is the norm decrease claimed to be strict at every configuration.

\subsection{Optimal orbit compensation and the induced kinetic form}

Write $U^0=\Phi_{-s}(W^0(g,y))$ and use left-multiplication frames
in the $k$ variables. Gauge equivariance of the flow gives
$\chi_g(k,y)=k^{-1}\cdot U^0$. At $k=I$ its differential is
\[
 D_{(k,y)}\chi_g=(-D_{U^0},B_s),\qquad B_s=D_y\chi_g(I,y).
 \tag{V34.7}\label{c18v34:chart}
\]
Thus the orbit metric is exactly $K_{U^0}$, not a matrix whose
inverse requires an exponential inverse-flow estimate. At other
$k$ the displayed frame is pushed forward by the gauge isometry
$k^{-1}$. Hence every metric block in these compatible frames,
and in particular the $yy$ block of the inverse metric, is
independent of $k$. Set
\[
 M_s=
 \begin{pmatrix}
 K_{U^0}&-D_{U^0}^*B_s\\
 -B_s^*D_{U^0}&B_s^*B_s
 \end{pmatrix},
 \qquad a_{g,s}=[M_s^{-1}]_{yy}.
 \tag{V34.8}\label{c18v34:metric}
\]
This $a_{g,s}$ is exactly V33's averaged coefficient, not a new
conditional law.

\begin{proposition}[Least fine-metric cost at a fixed chord velocity]
\label{c18v34:minimum}
For each chord velocity $q$, minimization over orbit velocities
$\xi$ gives
\[
 \begin{aligned}
 \xi_*&=K_{U^0}^{-1}D_{U^0}^*B_sq,&
 \|\xi_*\|&\le\sqrt{C_\ell}e^{16s}\|q\|,\\
 -D_{U^0}\xi_*+B_sq&=P_{U^0}B_sq,&
 a_{g,s}^{-1}&=B_s^*P_{U^0}B_s,\\
 \min_\xi\|-D_{U^0}\xi+B_sq\|^2
 &=q^*a_{g,s}^{-1}q.&
 \end{aligned}
 \tag{V34.9}\label{c18v34:minimum-formula}
\]
\end{proposition}
\begin{proof}
Complete the square using the positive matrix $K_{U^0}$.
The remainder is the Schur complement of its block in $M_s$,
which is positive because $\chi_g$ is an immersion.
The bound uses \eqref{c18v34:pinned-bound} and
$\|B_s\|\le e^{16s}$ from V33.
\end{proof}

For the already constructed flux
$q_t=\bar QF_t(\bar T_g)\bar r_g$, write $q_{t,v}=q_tv$.
Project the entire V33 lift:
\[
 A_t^{\rm h}=P(A+Z_t^{\rm red})
            =A^{\rm h}+PZ_t^{\rm red}.
\]
Theorem~\ref{c18v34:score}, applied to V33's equivariant correction,
and \eqref{c18v34:minimum-formula} give
\[
 \begin{aligned}
 Dc\,A_t^{\rm h}&=I,\qquad
 r_t^{\rm h}(\chi_g(k,y))=e^{-t\bar T_g}\bar r_g(y),\\
 \mathbb E_{\nu_g}|PZ_t^{\rm red}v|^2
 &=\int q_{t,v}^*a_{g,s}^{-1}q_{t,v}\,p_g\,dy\\
 &\le \frac{t}{2\alpha_s^{\rm red}}\,
                         \mathbb E_{\nu_g}|v\cdot r|^2 .
 \end{aligned}
 \tag{V34.10}\label{c18v34:corrected}
\]
The last inequality follows either from projection contraction
and V33 or from its local-flux bound. This construction uses only
the already paid local inverses and $K_U^{-1}$.
It does not invert $\bar L_g$ or $\bar T_g$ on an arbitrary source.

\subsection{An explicit current--inverse--current kinetic gain}

Let $W^0=W^0(g,y)$ be the unflowed representative and restrict its
gauge differential to the three edge classes. Define
\[
 H_\ell=D_{\mathcal F}^*D_{\mathcal F}
                     +D_{\mathcal B}^*D_{\mathcal B},
 \qquad J_y=D_{\mathcal C},\qquad
 K_{W^0}=H_\ell+J_y^*J_y.
 \tag{V34.11}\label{c18v34:split}
\]
Each forest link equals $I$. Every last-chain link has a coarse
root at its head, whose parameter vanishes; its differential is
just the parameter at its tail. Consequently $H_\ell$ is a
configuration-independent, scalar rooted graph matrix tensored
with the three-dimensional colour identity. It is block diagonal
over cells, consisting of the forest energy and the extra
last-chain tail penalties. In particular
\[
 C_\ell^{-1}I\preceq H_\ell\preceq12I,\qquad \|J_y\|^2\le12.
 \tag{V34.12}\label{c18v34:H}
\]

\begin{proposition}[Exact unflowed coefficient and comparison after flow]
\label{c18v34:current}
At zero inverse-flow time the exact reduced kinetic coefficient is
\[
 \begin{aligned}
 a_{g,0}&=I+J_yH_\ell^{-1}J_y^*,\\
 a_{g,0}^{-1}
 &=I-J_y(H_\ell+J_y^*J_y)^{-1}J_y^*,\\
 I&\preceq a_{g,0}\preceq(1+12C_\ell)I .
 \end{aligned}
 \tag{V34.13}\label{c18v34:woodbury}
\]
For every $s\ge0$, at the same $(g,y)$,
\[
 e^{-32s}a_{g,0}\preceq a_{g,s}\preceq e^{32s}a_{g,0}.
 \tag{V34.14}\label{c18v34:flow}
\]
\end{proposition}
\begin{proof}
At $s=0$, the chord differential inserts $q$ only on
$\mathcal C$, isometrically, so $B_0^*B_0=I$ and
$D_{W^0}^*B_0=J_y^*$. Thus
\[
 M_0=\begin{pmatrix}
 H_\ell+J_y^*J_y&-J_y^*\\-J_y&I
 \end{pmatrix}.
\]
The Schur formula gives the second line of
\eqref{c18v34:woodbury}. Its inverse is the first line:
multiply the two matrices and use
$(H_\ell+J_y^*J_y)H_\ell^{-1}
=I+J_y^*J_yH_\ell^{-1}$. No matrices are assumed to commute.
The bounds follow from \eqref{c18v34:H}.
Both $D\Phi_s$ and $D\Phi_{-s}$ have norm at most $e^{16s}$.
Applied to every chart tangent vector, this proves
$e^{-32s}M_0\preceq M_s\preceq e^{32s}M_0$.
Invert these inequalities and restrict to the $yy$ principal block.
\end{proof}

For clarity write $p_{g,s}$ for the actual density $p_g$ of V33
at the chosen flow time, including its inverse-flow Jacobian.
Define the comparison form using \emph{this same density}:
\[
 \begin{aligned}
 \mathcal E^\circ_{g,s}(f,f)
 &=\int\left\{|\nabla_y f|^2+
     (J_y^*\nabla_yf)^*H_\ell^{-1}(J_y^*\nabla_yf)\right\}
                                  p_{g,s}\,dy,\\
 e^{-32s}\mathcal E^\circ_{g,s}
 &\preceq\overline{\mathcal E}_{g,s}
                 \preceq e^{32s}\mathcal E^\circ_{g,s}.
 \end{aligned}
 \tag{V34.15}\label{c18v34:current-energy}
\]
Replacing $p_{g,s}$ by $p_{g,0}$ in this comparison is not justified.
The quantity $J_y^*\nabla_yf$ is the chord gauge current; the second
term is its positive inverse-$H_\ell$ energy.
Dually, for a correction velocity,
\[
 q^*a_{g,s}^{-1}q
 \le e^{32s}\left\{|q|^2-
 (J_y^*q)^*(H_\ell+J_y^*J_y)^{-1}(J_y^*q)\right\}.
 \tag{V34.16}\label{c18v34:saving}
\]
This is an explicit nonnegative saving in the previous
$e^{32s}|q|^2$ bound, without a claim that the saving is uniformly
positive on the actual source.

\subsection{A native physical direction with exactly zero gain}

Put $m=|V_{\rm f}|-|V_{\rm c}|$.
Since $|\mathcal C|=2m$, $J_y$ has $3m$ columns and $6m$ rows.
Therefore $a_{g,0}-I$ has rank at most $3m$ and a kernel of dimension
at least $3m$ at every configuration. A rank count alone does not
establish that the kernel contains gradients of physical functions.
For the chosen forest, however, an explicit one does.

In any cell take the plaquette based at residual coordinate
$r=(1,0,0)$ in directions $2,3$. Each of its four vertices is
internal. Each of its four links belongs to $\mathcal C$:
a forest edge in direction $2$ requires residual coordinate
$r_1=0$, and one in direction $3$ requires $r_1=r_2=0$;
none of these four links is a last-chain link.
Let
\[
 f(y)=\tfrac12\operatorname{Tr}
 \left(y_{x,2}y_{x+e_2,3}y_{x+e_3,2}^{-1}y_{x,3}^{-1}\right).
 \tag{V34.17}\label{c18v34:plaquette}
\]
All four chord coordinates are free even when $g$ is fixed.

\begin{proposition}[Physical zero-current test]
\label{c18v34:zero}
For every $g$, $s$ and coupling, this smooth nonconstant function has
positive variance under $p_{g,s}dy$, and
\[
 \begin{aligned}
 J_y^*\nabla_y f&=0,\qquad a_{g,0}\nabla_yf=\nabla_yf,\\
 \mathcal E^\circ_{g,s}(f,f)
 &=\int|\nabla_yf|^2p_{g,s}\,dy
   =4\int(1-f^2)p_{g,s}\,dy .
 \end{aligned}
 \tag{V34.18}\label{c18v34:zero-energy}
\]
As a function of the original fine configuration, this is the
gauge-invariant Wilson trace of that plaquette evaluated on
$\Phi_sU$.
\end{proposition}
\begin{proof}
Apply arbitrary gauge rotations at the internal endpoints to the
chord variables. In the displayed loop all intermediate rotations
cancel and its product changes only by conjugation. Differentiating
gives $\langle\nabla_yf,J_y\xi\rangle=0$ for every pinned parameter
$\xi$, proving the first identity. Then use
\eqref{c18v34:woodbury} and
\eqref{c18v34:current-energy}.
For a single $SU(2)$ factor with the round normalization used
throughout these notes, the half trace $z$ has
$|\nabla z|^2=1-z^2$. Fixing the other three links makes the loop
a left/right translate, possibly composed with inversion, of that
factor. All these maps are isometries. Each of the four distinct
links contributes $1-f^2$, proving the factor four.
The density is smooth positive on the full chord product, while
varying one of the four links makes $f$ nonconstant, so its variance
is strictly positive.

Finally $W^0=k\cdot\Phi_sU$. The loop trace is invariant under
this transformation and under all fine gauge transformations.
Thus the test is a genuine physical observable, not an engineered
non-Wilson probability model.
\end{proof}

This excludes a strict uniform improvement
$\mathcal E^\circ_{g,s}\ge(1+\delta)
\int|\nabla_y f|^2p_{g,s}dy$ with $\delta>0$ on all nonconstant
physical tests. It does \emph{not} say that the native physical
gap vanishes, that this test has low energy, or that it carries
any particular amount of the lift's source spectral weight.
At $s>0$ we claim only \eqref{c18v34:flow}, not
$a_{g,s}\nabla f=\nabla f$. The paid gauge inverse and the positive
current term alone cannot close the physical low-frequency step.

\subsection{Verification and the next upstream obligation}

The accompanying diagnostic checks the cell forests on actual
periodic graphs, the pinned-path bounds on numerical nonabelian
fields, and exact rational nonabelian Schur, Woodbury, projection,
full-output and Wilson-loop identities. A variable-metric
circle chart checks the divergence cancellation algebra separately;
it is labelled a fixture, not the native vacuum.
None of these finite checks replaces the analytic proofs or
computes the many-link Wilson ground state.

The native progress is a score-preserving, full-output reduction of
the original lift's norm, with an explicitly paid gauge inverse
and an exact current correction to the kinetic form.
Crucially $r^{A_t^{\rm h}}=r^{A_t^{\rm red}}$, so V33's remaining
low-frequency inverse moment is unchanged, not newly bounded.
The next upstream task must use the actual source and Wilson law
to control physical directions beyond this gauge current.
Projection-energy stability, conditional/coarse variance,
physical return and cofinal scale transport remain unpaid.
The interacting four-dimensional continuum construction and
positive physical Yang--Mills mass gap remain \emph{unproved}.
The companion-program checkpoint is V35; the next broad
primary-literature sweep is V40.
% END CYCLE 18 VERSION 34 APPEND
% BEGIN CYCLE 18 VERSION 35 APPEND
\clearpage
\section[V35: the native leading score and its exact conditional inverse]{V35: the native leading score and its exact conditional inverse}
\label{c18v35:context}

V34 finds genuine physical gradients on which its extra current
energy vanishes. That does not determine the spectral weight of
an actual lift score. We now compute such a score from the
\emph{native Wilson ground state}, rather than from an arbitrary
test density. For an explicitly selected full lift at spatial
flow time $s=0$, its leading electric-limit coefficient belongs
to a single eigenspace of the original conditional generator:
the eigenvalue is $12-6/\ell$. Its full coarse covariance and
inverse cost can therefore be evaluated exactly. An explicit
vertical correction cancels this leading score.

The result is a fixed-regulator expansion about $\beta=0$ in
$-\Delta+\beta E_{\rm sp}$. This is the electric-dominated,
physically strong-coupling expansion, not the weak-coupling
continuum regime. The coefficient bounds below are uniform in
spatial volume; the perturbative remainders are \emph{not} proved
uniform. Nor is $s=0$ silently substituted for the positive
spatial-flow time required by another stage of the programme.
The full continuum and mass-gap objectives remain unchanged.

\subsection{Companion checkpoint, provenance and choice of lift}

The scheduled V35 review rechecked inventories, SHA--256 hashes,
and terminal result/status sections of NS V26, SM--Einstein V72,
shell-mixing V16, parallel YM V5, the supplied ESM V64,
source-selection V52, native-stationarity V54, exact-action V45,
and the newer ESM V69. All are unchanged since V30.
Their finite source, fixed time, restricted chart and remainder
qualifications still do not provide a native YM conditional
spectral estimate. Their useful direction here is to calculate
the \emph{actual source coefficient} before assigning it to an
obstructing or controlled sector. This is not a recertification
of their complete historical proofs. The next companion
checkpoint is V40.

The ordinary first electric perturbation of a Wilson vacuum is
not new. A targeted primary-abstract check of M.~Beccaria,
\emph{Wave functions for Hamiltonian Lattice Gauge Theory},
\href{https://arxiv.org/abs/hep-lat/0003016}{arXiv:hep-lat/0003016},
confirms the historical use of strong-coupling plaquette wave
functions and improved trial states. No trial state, numerical
gap or Monte Carlo conclusion is imported here. The native
coefficient is derived from its eigenvalue equation below.
The new calculation is its rooted-fibre score, exact conditional
eigenvalue, full-output covariance and native score correction.
Archived f17's temporal pinned-Green coefficient and f5/V19's
calibrated one-link Fourier score cancellation concern different
operators and sources; they are not repeated as these results.
The next broad primary-literature sweep remains V40.

Keep the cellwise forest $\mathcal F$, last-chain edges
$\mathcal B$, and free chords $\mathcal C$ of V34 on
$N=\ell n$, $\ell\ge2$, $n\ge5$. In this note set $s=0$,
so $c=\mathfrak b_\ell$ and
$U=\chi(k,g,y)=k^{-1}\cdot W^0(g,y)$.
Let $A^{\rm tr}$ differentiate $g$ at fixed $(k,y)$ in the
coarse left-multiplication frames. Each column inserts a
variation on one distinct edge of $\mathcal B$, followed by
a gauge isometry. Thus
\[
 Dc\,A^{\rm tr}=I,\qquad (A^{\rm tr})^*A^{\rm tr}=I,\qquad
 A^{\rm h}=PA^{\rm tr},\quad (A^{\rm h})^*A^{\rm h}\preceq I .
 \tag{V35.1}\label{c18v35:lift}
\]
Both lifts are smooth and gauge equivariant. Under a fixed gauge
transformation the chart variables undergo fixed left/right
translations associated with their roots; these translations
intertwine $D_g\chi$. V34 applies to the second lift and preserves
its score exactly. This is a declared change of lift from V28's
weighted-optimal one. No formula below for the new score is
asserted to be the score of that different lift.

\subsection{The first coefficient of the actual vacuum}

Use the established round $SU(2)$ metric, normalized product Haar,
and no physical-time factor in the conditional generator. Put
\[
 T=-\sum_{e\in E_{\rm f}}\Delta_e,\qquad
 W(U)=\sum_{p\in\mathcal P_N}w_p(U),\qquad
 w_p=\tfrac12\operatorname{Tr}U_p .
\]
The additive constant in $\beta E_{\rm sp}
=\beta|\mathcal P_N|-\beta W$ does not change the ground state.
Let $\Omega_\beta$ be the normalized positive ground state of
$T-\beta W$. For fixed $N$, its eigenvalue and eigenvector are
real analytic near zero. To see the isolation directly, $T$
has a simple zero eigenvalue and next eigenvalue three on the
full link space, while $\|W\|_\infty\le M=|\mathcal P_N|$.
On $|z|=3/2$ the unperturbed resolvent has norm at most $2/3$.
The Neumann series for the bounded perturbation converges when
$|\beta|M<3/2$ and gives an analytic rank-one Riesz projection.
For real $\beta$ in this interval it selects the isolated bottom
state. Elliptic regularity and positive normalization give the
corresponding expansion in every fixed smooth norm on the compact
configuration space. This sufficient radius already depends on
the volume and is not a cofinal estimate.

Each plaquette contains four distinct fine links, so
$Tw_p=12w_p$ and $\int w_p\,dU=0$. Differentiating the normalized
ground-state equation at zero therefore gives
\[
 \Omega_\beta=1+\frac{\beta}{12}W+O_N(\beta^2),\qquad
 \Omega_\beta^2=1+\frac{\beta}{6}W+O_N(\beta^2).
 \tag{V35.2}\label{c18v35:ground}
\]
These are coefficients of the true vacuum, not an ansatz.

Write the exact conditional law and its coarse density as
\[
 p_{g,\beta}(y)=\frac{\Omega_\beta(W^0(g,y))^2}{\rho_\beta(g)},
 \qquad
 \rho_\beta(g)=\int\Omega_\beta(W^0(g,y))^2\,dy .
 \tag{V35.3}\label{c18v35:law}
\]
Every elementary plaquette contains a free chord. Indeed
$\mathcal F\cup\mathcal B$ consists of the coarse cubic graph
with its edges subdivided into $\ell$ links, together with
attached trees. Its shortest cycle has length $4\ell>4$.
Integrating a chord occurring once in $w_p(W^0)$ makes its
Haar integral zero. Consequently, with $W_g(y)=W(W^0(g,y))$,
\[
 \rho_\beta(g)=1+O_{N,\ell}(\beta^2),\qquad
 p_{g,\beta}=1+\frac{\beta}{6}W_g+O_{N,\ell}(\beta^2).
 \tag{V35.4}\label{c18v35:conditional}
\]
All remainders in this note are uniform in $g$ at fixed $N,\ell$,
in each fixed smooth norm, unless a different norm is specified.

Let $D_{e,a}^{g}$ be a coarse left derivative and let
$r_\beta$ be the conditionally centred score of either lift in
\eqref{c18v35:lift}. The exact Haar chart gives
$S_{e,a}=D_{e,a}^{g}\log\Omega_\beta(W^0)^2$ and hence
$r_{\beta,e,a}=D_{e,a}^{g}\log p_{g,\beta}$.
In particular
\[
 r_\beta=\beta r_1+O_{N,\ell}(\beta^2),\qquad
 (r_1)_{e,a}=\frac16D_{e,a}^{g}W_g
     =\frac16\sum_{p\ni b_e}D_{e,a}^{g}w_p(W^0),
 \tag{V35.5}\label{c18v35:source}
\]
where $b_e\in\mathcal B$ is the last link of coarse edge $e$.
This explicitly retains every coarse column.

\subsection{The exact rooted Green diagonal}

Let $H_\ell$ be V34's scalar cell matrix, with the colour identity
suppressed. Its inverse is zero between different cells.
For a nonroot residual vertex $r$, let
$j=\max\{i:r_i\ne0\}$. Then its diagonal entry is
\[
 h(r):=(H_\ell^{-1})_{rr}
       =\sum_{i<j}r_i+\frac{r_j(\ell-r_j)}{\ell},
 \qquad \tau_\ell=\frac{\ell-1}{\ell}.
 \tag{V35.6}\label{c18v35:green}
\]
Here $h(e_i)=h((\ell-1)e_i)=\tau_\ell$.

For completeness, the cell forest has three coordinate-axis
backbones, each with a grounded endpoint at the root and one
additional grounding edge at its last-chain tail. All other
branches attach to these backbones as trees. Eliminate their
leaves from the quadratic energy. A branch carrying no source
has constant minimizing value and adds no conductance.
A unit source at $r$ traverses a unique branch of length
$\sum_{i<j}r_i$ to $r_je_j$. From there the two paths to ground
have lengths $r_j$ and $\ell-r_j$ in parallel.
The minimum flow energy is their parallel resistance plus the
branch length, giving \eqref{c18v35:green}. Equivalently, solve
$H_\ell u=\delta_r$ and use
$u_r=u^*H_\ell u$ for the same unit flow; this identifies that
energy with the inverse diagonal without a probability-law
inverse.

Let $X_{x,a}$ be the chord gauge vector field at internal vertex
$x$ and colour $a$; its fine-link coefficient is the corresponding
column of $J_y$. It is divergence-free for chord Haar. The
unperturbed native reduced generator from V34 is exactly
\[
 \bar L_0
 =-\sum_{f\in\mathcal C}\Delta_f
  -\sum_{x,z}\sum_{a=1}^3
          (H_\ell^{-1})_{xz}X_{x,a}X_{z,a}.
 \tag{V35.7}\label{c18v35:generator}
\]
It is independent of $g$. Its form dominates the product-Haar
form and therefore has centred gap at least three. Define also
the unweighted current Casimir
\[
 \mathcal G=-\sum_x\sum_{a=1}^3X_{x,a}^2 .
 \tag{V35.8}\label{c18v35:current}
\]
Neither $\mathcal G$ nor $H_\ell$ is substituted for $\bar L_0$.

\subsection{A single native eigenvalue for the complete leading score}

Fix $b_e$ in direction $j$, with tail
$x=R+(\ell-1)e_j$ and head $R'=R+\ell e_j$.
Its four adjacent plaquettes are specified by a perpendicular
direction $i\ne j$ and a positive or negative displacement in that
direction. The following list gives the free-chord part. In the
last row the same plaquette is shared by two incoming last-chain
edges and is counted only once.
\[
\begin{array}{c|c|c|c}
 \text{type}&|\mathcal B\cap p|&|\mathcal C\cap p|&
                       h(a)+h(b)\\ \hline
 +i,\ i<j&1&1&1+2\tau_\ell\\
 +i,\ i>j&1&2&2\tau_\ell\\
 -i&2&2&2\tau_\ell
\end{array}
 \tag{V35.9}\label{c18v35:types}
\]
In every case the free chords form one open path with endpoints
$a,b$ in different cells.

Here is a direct verification. For $+i$ the link from $R'$ to
$R'+e_i$ is in $\mathcal F$. If $i<j$, the link from $x$ to
$x+e_i$ is also in $\mathcal F$, leaving the single chord from
$x+e_i$ to $R'+e_i$. Its endpoints have diagonals
$1+\tau_\ell,\tau_\ell$. If $i>j$, those two links from $x$ through
$x+e_i$ to $R'+e_i$ are both chords; their endpoints have diagonals
$\tau_\ell,\tau_\ell$. For $-i$, the two nonchord links are the
incoming last-chain links ending at $R'$. The remaining open path
joins their two tails, which are in distinct neighbouring cells
and each have diagonal $\tau_\ell$.
These observations also prove that different listed plaquettes
have different chord supports. A two-link right-angle path
determines its elementary square; a one-link case is determined
by its direction and its single earlier residual displacement.

\begin{theorem}[Exact leading-source eigenspace]
\label{c18v35:eigen}
For every $g$, every full coarse column of $r_1$ obeys
\[
 \mathcal G r_1=6r_1,\qquad
 \bar L_0 r_1=\lambda_\ell r_1,\qquad
 \lambda_\ell=12-\frac6\ell\ge9 .
 \tag{V35.10}\label{c18v35:eigenvalue}
\]
\end{theorem}
\begin{proof}
A listed plaquette, or any of its coarse derivatives, is a
fundamental matrix coefficient of its open chord-path product.
The other factor depends only on $g$; differentiating it may make
it a nonunit matrix, which does not change its representation.
Each of the $k=1$ or $2$ chord Casimirs contributes three.
Chord gauge rotations cancel at each internal path vertex; at
each of the two endpoints the fundamental Casimir is three.
Thus $\mathcal G$ contributes six.

In \eqref{c18v35:generator} only the endpoint gauge fields remain.
Vector fields at different vertices commute, so cross terms
with an intermediate vertex also vanish. The two endpoint cells
are distinct, giving $(H_\ell^{-1})_{ab}=0$.
The current part is therefore the scalar
$3h(a)+3h(b)$. The total eigenvalue is
$3k+3h(a)+3h(b)$, which every row of
\eqref{c18v35:types} makes equal to
$6+6\tau_\ell=12-6/\ell$.
Sum the actual plaquette derivatives in
\eqref{c18v35:source}.
\end{proof}

This theorem is a native conditional calculation, not the
unconditional plaquette eigenvalue twelve. It includes the
induced metric's current term, the internal orbit elimination,
and the complete coarse output.

\subsection{The full coefficient covariance and inverse cost}

Let $\Sigma_1(g)=\int r_1r_1^*\,dy$.
For a coarse bank $v=(v_e)$ in the fixed left frames, write
\[
 \widetilde v_e=\operatorname{Ad}_{g_e^{-1}}v_e .
\]
At a coarse vertex $z$, denote the three incoming coarse links
by $e_1(z),e_2(z),e_3(z)$.

\begin{theorem}[Exact covariance on all coarse columns]
\label{c18v35:covariance}
The full matrix is characterized by
\[
 v^*\Sigma_1(g)v
 =\frac1{144}\left\{
 2\sum_{e\in E_{\rm c}}|v_e|^2+
 \sum_{z\in V_{\rm c}}\sum_{i<j}
       |\widetilde v_{e_i(z)}-\widetilde v_{e_j(z)}|^2
              \right\}.
 \tag{V35.11}\label{c18v35:covariance-form}
\]
After the displayed orthogonal endpoint-frame rotations, each
coarse vertex contributes the block
\[
 \frac1{144}(5I_3-\mathbf1\mathbf1^*)\otimes I_3.
\]
Consequently its eigenvalues and the leading inverse cost are
\[
 \begin{aligned}
 \operatorname{spec}\Sigma_1
  &=\left\{\frac1{72}\ [3|V_{\rm c}|],\
                   \frac5{144}\ [6|V_{\rm c}|]\right\},\\
 R_1(g)&:=\int r_1\,\bar L_0^{-1}r_1^*\,dy
                  =\lambda_\ell^{-1}\Sigma_1(g).
 \end{aligned}
 \tag{V35.12}\label{c18v35:inverse}
\]
The bracketed quantities are multiplicities.
\end{theorem}
\begin{proof}
A Haar $SU(2)$ matrix is a unit quaternion with second moment
$I_4/4$. Hence the variance of a half-trace matrix coefficient
with coefficient quaternion $m$ is $|m|^2/4$.
An open product of independent Haar chords is Haar.

Different plaquettes in \eqref{c18v35:types} have different chord
supports. Integrate a chord belonging to just one of two
different supports; it occurs linearly in that coefficient,
so every cross-plaquette moment is zero.
There are two one-$\mathcal B$ plaquettes per last-chain link.
The directional derivative on each has variance $|v_e|^2/4$.
For a two-$\mathcal B$ plaquette with incoming edges $e,f$ at the
same coarse root, the fixed part of its word is, up to inverses
and cyclic rearrangement, $M=g_eg_f^{-1}$. Its derivative is
$v_eM-Mv_f$, whose quaternion norm is
$|\operatorname{Ad}_{g_e^{-1}}v_e-
  \operatorname{Ad}_{g_f^{-1}}v_f|$.
Its variance is the square of this norm divided by four.
There is one such plaquette for each incoming pair.
Finally the source factor $1/6$ is squared, giving $1/144$.

The three-vertex complete-graph Laplacian is
$3I_3-\mathbf1\mathbf1^*$. Adding the two unary contributions
gives the displayed block and its eigenvalues two and five.
Apply Theorem~\ref{c18v35:eigen} for the inverse formula.
No coarse column has been deleted to obtain these constants.
\end{proof}

Let $R_\beta(g)$ be the original-metric native conditional inverse
cost of $r_\beta$, for the selected lift. At fixed $N,\ell$,
\[
 \begin{aligned}
 \Sigma_\beta(g)&=\beta^2\Sigma_1(g)+O_{N,\ell}(\beta^3),\\
 R_\beta(g)&=\frac{\beta^2}{\lambda_\ell}\Sigma_1(g)
                                     +O_{N,\ell}(\beta^3).
 \end{aligned}
 \tag{V35.13}\label{c18v35:cost-expansion}
\]
The remainders are matrix operator-norm statements, uniform in $g$,
not estimates independent of $N$.

Here is a justification of the inverse expansion that does not
assume an interacting uniform gap. For fixed regulator,
$p_{g,\beta}=1+O(\beta)$ in smooth norm uniformly in $g$.
On the Haar-mean-zero $H^1$ space, the form
$\int(\nabla f)^*a_0\nabla f\,p_{g,\beta}dy$ remains coercive,
using the Haar gap three and V34's $a_0\succeq I$.
Its coefficients and the mean-zero right-hand side
$r_\beta p_{g,\beta}$ have Taylor expansions in bounded form norms.
A Neumann expansion of this coercive variational problem yields
the solution
$u_\beta=\beta r_1/\lambda_\ell+O_{N,\ell}(\beta^2)$ in $H^1$,
up to an irrelevant constant. Pair with $r_\beta p_{g,\beta}$
to obtain \eqref{c18v35:cost-expansion}.

\subsection{Zero-current selection and an explicit native correction}

Define the fixed closed subspace
\[
 \mathcal Z=\{f\in L^2(dy):X_{x,a}f=0
                      \text{ in distributions for all }x,a\}.
\]
It includes V34's physical free-chord Wilson plaquette.
Theorem~\ref{c18v35:eigen} and Haar integration by parts give
$\int f\,r_1\,dy=0$ for every $f\in\mathcal Z$.
This says something about the actual leading source, not only
about the existence of a zero-current test.

Let $\Pi_{\mathcal Z,\beta}$ be the orthogonal projection onto
this same subspace in the equivalent norm $L^2(p_{g,\beta}dy)$.
For each bank $v$,
\[
 \|\Pi_{\mathcal Z,\beta}(v\cdot r_\beta)\|_{L^2(p_{g,\beta})}
        =O_{N,\ell}(\beta^2)|v|.
 \tag{V35.14}\label{c18v35:selection}
\]
Indeed $r_\beta p_{g,\beta}=\beta r_1+O_{N,\ell}(\beta^2)$
in Haar $L^2$. Pair with an arbitrary unit vector of
$\mathcal Z$ in the weighted norm, use the zero Haar pairing
of its leading term, and use uniform equivalence of the two norms.
The squared projected weight is thus $O_{N,\ell}(\beta^4)$.
At nonzero $\beta$ this projection is not asserted to commute
with the native generator; it is not a low-frequency spectral
projection.

There is also a direct constructive consequence. Let
$\nabla_{\rm fib}$ denote the original fine induced gradient on
$c^{-1}(g)$, with no artificial time factor, and set
\[
 Z_\beta=\frac{\beta}{\lambda_\ell}\nabla_{\rm fib}r_1,
 \qquad
 \widetilde A_\beta=A^{\rm h}+Z_\beta .
 \tag{V35.15}\label{c18v35:correction}
\]
The vector-valued gradient is taken column by column after
pulling $r_1(g,y)$ back to the original fibre. This is a smooth
equivariant vertical field. In particular it is internally
horizontal, since its potential is internally invariant.

\begin{theorem}[Explicit first-source cancellation]
\label{c18v35:corrected}
For the actual native law, at fixed $N,\ell$ and $s=0$,
\[
 \begin{aligned}
 Dc\,\widetilde A_\beta&=I,\\
 \widetilde r_\beta
  &=r_\beta-\frac{\beta}{\lambda_\ell}\bar L_\beta r_1
                                    =O_{N,\ell}(\beta^2),\\
 \mathbb E_{\nu_{g,\beta}}|Z_\beta v|^2
  &=\frac{\beta^2}{\lambda_\ell}v^*\Sigma_1(g)v
                                 +O_{N,\ell}(\beta^3)|v|^2,\\
 R^{\widetilde A_\beta}(g)&=O_{N,\ell}(\beta^4).
 \end{aligned}
 \tag{V35.16}\label{c18v35:corrected-cost}
\]
The error statements are uniform in $g$. The construction uses
the explicit eigenfunction, not a whole-fibre interacting inverse.
\end{theorem}
\begin{proof}
Verticality proves the full-output identity. The full weighted
divergence of a vertical field equals its conditional divergence,
and its conditional mean is zero. Therefore the new residual is
exactly the second line. Since
$\bar L_\beta r_1=\bar L_0r_1+O_{N,\ell}(\beta)
=\lambda_\ell r_1+O_{N,\ell}(\beta)$, its linear term cancels.

At zero coupling, integration by parts gives
$\int|\nabla_{\rm fib}(v\cdot r_1)|^2d\nu_{g,0}
=\lambda_\ell v^*\Sigma_1v$. Replace the density by
$p_{g,\beta}=1+O_{N,\ell}(\beta)$ to obtain the norm expansion.
Finally fixed-regulator density comparison with Haar and
$a_0\succeq I$ bounds the centred inverse of $\bar L_\beta$
for sufficiently small $\beta$. Apply that bound to the
$O_{N,\ell}(\beta^2)$ corrected residual. This proves the last
line without claiming a cofinal interacting spectral floor.
\end{proof}

The full lift, actual density and actual metric have all been
retained. Existing variable-bank derivative and projection-energy
obligations are not removed by this coefficient calculation.

\subsection{Verification, advancement and the unpaid remainder}

The diagnostic enumerates all source plaquette types on fine
tori of sides $10,15,20$ with $\ell=2,3,4$ and coarse side five.
It checks all exact cell Green diagonals, unique chord supports,
different endpoint cells, and the common eigenvalues
$9,10,21/2$. Rational quaternion second derivatives independently
check both path Casimirs. Exact matrix contractions check the
complete nine-component block at a coarse vertex, its frame
rotations, covariance eigenvalues and inverse-cost coefficient.
These finite regressions support the algebra; the all-volume
coefficient statements rest on the proofs. The many-link native
vacuum and the higher-order remainder have not been numerically
computed, and no proof-assistant certification is asserted.

V34's zero-current physical examples are therefore not leading
source obstructions for this explicitly chosen native lift.
The leading source has a paid conditional inverse, and the
explicit correction removes it with a volume-independent
leading cost. This is stronger than assigning a generic spectral
floor to an unspecified source.

The next upstream obligation is the \emph{actual corrected
higher-order source}, including any transfer into physical
zero-current directions, followed by a regulator-uniform
remainder or a genuinely nonperturbative estimate.
The displayed analytic radius and error constants do not supply
that estimate, and a finite electric-limit expansion cannot be
continued to the continuum coupling merely by formal iteration.
Positive-flow transport, projection-energy stability,
conditional/coarse variance and physical scale calibration
remain separate obligations. The interacting four-dimensional
continuum construction and positive physical Yang--Mills mass
gap remain \emph{unproved}.
% END CYCLE 18 VERSION 35 APPEND
% BEGIN CYCLE 18 VERSION 36 APPEND
\clearpage
\section[V36: the second native score, adjoint cost and local correction]{V36: the second native score, adjoint cost and local correction}
\label{c18v36:context}

V35 cancels the first coefficient of an actual Wilson-vacuum
score. Its remainder is not thereby controlled at the continuum
coupling. This note computes the next coefficient, proves that
it is nonzero, and constructs its local, unperturbed inverse.
The second residual has an explicitly paid adjoint component,
but still no Haar zero-current component. The entire
second inverse-cost coefficient admits a volume-independent
bound at fixed blocking factor. These are statements about
coefficients, not convergence of the expansion.

Throughout, retain V35's round $SU(2)$ metric, three-dimensional
spatial torus $N=\ell n$, $\ell\ge2$, $n\ge5$, full coarse
output, and the declared tree lift $A^{\rm h}$ at spatial flow
time $s=0$. The Hamiltonian is $T-\beta W$, with its exact
positive normalized vacuum. Write $\lambda=12-6/\ell$,
$h=H_\ell^{-1}$ and
\[
 \Gamma_0(f,q)=(\nabla_{\mathcal C}f)^*a_0\nabla_{\mathcal C}q
 =\sum_{f'\in\mathcal C}\nabla_{f'}f\cdot\nabla_{f'}q
   +\sum_{x,z,a}h_{xz}(X_{x,a}f)(X_{z,a}q).
 \tag{V36.1}\label{c18v36:gamma}
\]
Here $a_0$ is the exact, $\beta$- and $g$-independent reduced
inverse metric from V34, not a replacement metric.
Every error term below is in a fixed smooth norm or the
specified matrix norm, uniformly in $g$ at fixed $N,\ell$.
No volume-independent error constant is implicit.

\subsection{New companion information and the precise transfer}

Five companion files appeared during this iteration. Their
terminal result/status sections were reviewed and their source
hashes recorded without editing them. ESM V80 has advanced to a
two-loop, weight-six bosonic source construction on its stated
perturbative domain. Its hard-annulus remainder and rooted
kernel estimates use fixed source arity, profiles and inverse
margins. They do not construct the present native YM vacuum
or estimate its continuum-coupling conditional inverse.
The useful transfer is to organize a coefficient by
\emph{connected, rooted local terms} and to preserve the
physical quotient rather than discard it.

Source-selection V65's finite matrix-unit defect gap is
conditional on represented source actions; it is not a YM
spectral gap. Stationarity V61 identifies an off-constraint
leading defect and an initialized-path cancellation, but its
state-only source remains unevaluated. Exact-action V55
separates a longer-time normalization from an unevaluated
same-root slow contraction. Arithmetic-loading V17 likewise
requires represented projectors before its finite carrier
tests apply. These qualifications reinforce the present
choice to calculate the native source itself. This targeted
transfer review does not recertify those manuscripts in full.
The next scheduled companion and broad-literature checkpoints
remain V40.

Connected logarithmic vacuum expansions are not claimed as a
new method. A targeted primary-abstract check of A.~Wichmann,
D.~Schuette, B.~C.~Metsch and V.~Wethkamp,
\emph{The Coupled Cluster Method in Hamiltonian Lattice Field
Theory: SU(2) Glueballs},
\href{https://arxiv.org/abs/hep-lat/0112015}{hep-lat/0112015},
concerns a truncated closed-path Hamiltonian basis in two
spatial dimensions. No truncation error or mass prediction
from that work is imported. Earlier notes already treat
ordinary connected effective-Hamiltonian coefficients.
The calculation here concerns the different, rooted-fibre
score left by V35's particular native correction.

\subsection{The connected second coefficient of the actual vacuum}

Let $p\sim q$ mean an unordered pair of distinct elementary
plaquettes sharing one fine edge. Orient their words at that
edge as $UA$ and $UB$, reversing a plaquette if needed;
its half-trace is unchanged by reversal. Put
\[
 \chi_p=4w_p^2-1,\qquad R_{pq}=\tfrac12\operatorname{Tr}(AB^{-1}).
\]
The word in $R_{pq}$ is the simple six-edge outer loop of the
two plaquettes, including when they lie in different planes.
All these statements are on the full fine-link space until
the pullback to $W^0(g,y)$ is explicitly taken.

\begin{lemma}[Native logarithmic coefficient]
\label{c18v36:loglemma}
With $M=|\mathcal P_N|$, define
\[
 F_2=-\frac1{2304}\sum_p\chi_p+
       \sum_{p\sim q}\left(\frac{R_{pq}}{702}
                                -\frac{w_pw_q}{936}\right).
 \tag{V36.2}\label{c18v36:F2}
\]
The normalized ground state satisfies
\[
 \log\Omega_\beta^2
   =\frac{\beta}{6}W+
        \beta^2\left(F_2-\frac{M}{288}\right)+O_N(\beta^3).
 \tag{V36.3}\label{c18v36:logvac}
\]
\end{lemma}
\begin{proof}
Use the fixed-volume analytic eigenpair established in V35.
For $u=\log\Omega_\beta$, its eigenvalue equation is
$Tu=\beta W+E_\beta+|\nabla u|^2$.
Its first coefficient is $W/12$, up to a constant, and
$E_2=-\int W T^{-1}W\,dU=-M/48$. The last identity uses
$Tw_p=12w_p$, $\int w_p^2=1/4$, and zero cross-plaquette
moments. Therefore the nonconstant coefficient of $\beta^2$
in $\log\Omega_\beta^2$ solves
\[
 TF_2=\frac1{72}|\nabla W|^2-\frac{M}{24}.
 \tag{V36.4}\label{c18v36:logeq}
\]

For a single link, the half-trace is a linear coordinate on
the unit quaternion sphere, so
$|\nabla w_p|^2$ on that link is $1-w_p^2$.
Its four-link sum is $4(1-w_p^2)$.
The nonconstant self term in \eqref{c18v36:logeq} is thus
$-\chi_p/72$. The spin-one character $\chi_p$ has link
Casimir eight on each of its four links, giving
$T\chi_p=32\chi_p$.

For a shared link, the quaternion coordinate identity gives
\[
 \nabla_Uw(UA)\cdot\nabla_Uw(UB)=R_{pq}-w_pw_q.
\]
The product rule for the positive Laplacian then gives
\[
 T(w_pw_q)=26w_pw_q-2R_{pq},\qquad TR_{pq}=18R_{pq}.
\]
Consequently
\[
 T^{-1}(R_{pq}-w_pw_q)
       =\frac2{39}R_{pq}-\frac1{26}w_pw_q .
\]
The cross term in \eqref{c18v36:logeq} has coefficient $1/36$,
which gives exactly \eqref{c18v36:F2}. Disjoint link supports
have zero gradient contraction, so there are no other terms.
Every summand in $F_2$ has zero fine Haar mean. Finally
$\int(W/6)^2\,dU=M/144$; expanding
$\int\exp(\log\Omega_\beta^2)\,dU=1$ fixes the constant
as $-M/288$.
\end{proof}

\begin{lemma}[No second-order coarse normalization]
\label{c18v36:normalization}
Different elementary plaquettes have different nonempty
$\mathcal C$-supports. On the rooted chart,
\[
 \int F_2(W^0(g,y))\,dy=0,\quad
 \int W_g^2\,dy=\frac M4,\quad
 \rho_\beta(g)=1+O_{N,\ell}(\beta^3).
 \tag{V36.5}\label{c18v36:rho}
\]
\end{lemma}
\begin{proof}
Nonemptiness follows from the girth $4\ell\ge8$ of
$\mathcal F\cup\mathcal B$ after attached trees are ignored.
Two different elementary plaquettes share at most one fine
edge. Equal chord supports of size at least two are therefore
impossible. Equal singleton supports would leave their
six-edge outer cycle entirely in
$\mathcal F\cup\mathcal B$, also impossible.
Integrating a chord in the symmetric difference now shows
$\int w_pw_q\,dy=0$ for $p\ne q$.
Integrating any chord in $p$ gives
$\int w_p^2\,dy=1/4$ and $\int\chi_p\,dy=0$.
Every six-edge outer loop contains a chord occurring once,
so its integral vanishes. These prove the first two identities.
In the expansion of the conditional integral of
\eqref{c18v36:logvac}, $-M/288$ cancels
$\frac12\int(W_g/6)^2dy$. This proves the last identity.
\end{proof}

Below all $w_p$, $\chi_p$, $R_{pq}$ and $F_2$ denote their
pullbacks to the rooted chart. In particular the actual
uncorrected score is
\[
 r_\beta=\beta r_1+\beta^2D_gF_2+O_{N,\ell}(\beta^3),
 \qquad r_1=\frac16D_gW_g.
 \tag{V36.6}\label{c18v36:r2}
\]
The conditional normalizer has not been dropped by assumption.

\subsection{The actual source remaining after the first correction}

V35's correction has potential $\phi_1=r_1/\lambda$.
The native weighted generator satisfies, on fixed smooth
functions,
\[
 \bar L_\beta f=\bar L_0f-
             \frac{\beta}{6}\Gamma_0(W_g,f)+O_{N,\ell}(\beta^2).
\]
This is the divergence-form identity
$\bar L_\beta=\bar L_0-\Gamma_0(\log p_{g,\beta},\,\cdot\,)$.
It uses the original induced metric.
Thus the residual of $A^{[1]}=A^{\rm h}+
\beta\nabla_{\rm fib}\phi_1$ is
\[
 \begin{aligned}
 \widetilde r_\beta&=\beta^2t_2+O_{N,\ell}(\beta^3),\\
 t_2&=D_gF_2+\frac1{36\lambda}
                           \Gamma_0(W_g,D_gW_g)\\
    &=D_g\left(F_2+\frac1{72\lambda}\Gamma_0(W_g,W_g)\right).
 \end{aligned}
 \tag{V36.7}\label{c18v36:t2}
\]
Here $D_g$ commutes with the metric coefficients and the chord
derivatives. Formula \eqref{c18v36:t2} evaluates the source;
it is not a notation for an unknown second remainder.

\subsection{The zero-current component still vanishes}

Let $\mathcal S=\{p:p\cap\mathcal B\ne\varnothing\}$,
$W_B=\sum_{p\in\mathcal S}w_p$ and $W_F=W_g-W_B$.
The subscript on $W_F$ means absence of $\mathcal B$, not
absence of free chords. Recall
$\mathcal Z=\bigcap_{x,a}\ker X_{x,a}$ and its Haar orthogonal
projection $\Pi_{\mathcal Z}$.

\begin{lemma}[Endpoint selection and the cell subgroup]
\label{c18v36:endpoints}
Different plaquettes in $\mathcal S$ have different unordered
endpoint pairs for their open chord paths. Moreover
\[
 \Pi_{\mathcal Z}W_B^2=\frac{|\mathcal S|}{4},\qquad
 \Pi_{\mathcal Z}(W_FW_B)=0,\qquad
 \Pi_{\mathcal Z}(W_B\bar L_0W_F)=0.
 \tag{V36.8}\label{c18v36:selection-products}
\]
\end{lemma}
\begin{proof}
Use V35's three source types. For a one-$\mathcal B$
plaquette the endpoint cells are nearest neighbours and
determine the last-chain edge; its displacement specifies
the endpoint vertices uniquely. Its two different positive
perpendicular displacements have different endpoints.
For a two-$\mathcal B$ plaquette the endpoint cells differ
in two coordinates; the endpoint vertices are the two
last-chain tails and determine the incoming pair and its
common coarse head. Such a pair cannot equal a
one-$\mathcal B$ pair. Periodicity introduces no ambiguity
when $n\ge5$.
An unmatched endpoint of two different open paths carries
one fundamental representation, whose gauge Haar average
is zero. For the square of one path, averaging an endpoint
gives $1/4$. This proves the first identity.

For the other identities let $K$ be the subgroup of chord
gauge transformations that has the same value at every
internal vertex of each cell. There are no chords incident
to a coarse root: its outgoing axis links are forest links
and its incoming axis links are last-chain links.
One may therefore extend a cell value to its root without
changing its action on chords. All forest links stay the
identity, and every plaquette without $\mathcal B$ is
invariant under $K$. Thus $W_F$ is $K$-invariant.
Every source path has endpoints in two distinct cells, so
its $K$ average is zero.

The operator $\bar L_0$ commutes with $K$: its Haar link
Laplacians are bi-invariant, and each scalar product
$\sum_aX_{x,a}X_{z,a}$ with $h_{xz}\ne0$ has both vertices
in the same cell. Hence $\bar L_0W_F$ is also
$K$-invariant. First averaging under $K$, and then under
the full chord gauge group, proves the other two identities.
Commutation with independent rotations at all individual
vertices is neither required nor asserted.
\end{proof}

\begin{theorem}[Second-order zero-current exclusion]
\label{c18v36:zero}
For every $g$,
\[
 \Pi_{\mathcal Z}t_2=0.
 \tag{V36.9}\label{c18v36:zeroselection}
\]
For V35's corrected residual, the weighted projection onto
the same function space satisfies
\[
 \|\Pi_{\mathcal Z,\beta}(v\cdot\widetilde r_\beta)\|
                  _{L^2(p_{g,\beta})}
       =O_{N,\ell}(\beta^3)|v|.
 \tag{V36.10}\label{c18v36:weighted}
\]
\end{theorem}
\begin{proof}
First consider $D_gF_2$. A differentiated self term is an
open-path adjoint coefficient, with nontrivial endpoint
charges. A product of two source plaquettes has an
unmatched endpoint unless they coincide; the pair sum in
$F_2$ excludes coincidence. A source plaquette times a
non-source plaquette has zero $K$ average.
A fused six-edge loop that depends on $g$ contains a
last-chain edge and at least one chord. Its chord set is
a proper, nonempty subset of a simple cycle and therefore
has an endpoint with exactly one chord. Averaging at that
internal vertex kills its fundamental coefficient.
All terms with no $\mathcal B$ are independent of $g$.
Since $D_g$ commutes with these fixed group averages,
$\Pi_{\mathcal Z}D_gF_2=0$.

The subspace $\mathcal Z$ reduces $\bar L_0$. Indeed all its
$X$ fields vanish, while the product link Laplacian
preserves chord gauge invariance. Equivalently, in the
closed quadratic form the cross term between
$\mathcal Z$ and its orthogonal complement is zero.
Thus $\Pi_{\mathcal Z}\bar L_0=\bar L_0\Pi_{\mathcal Z}$
on smooth functions. Use the identity
\[
 2\Gamma_0(f,q)=f\bar L_0q+q\bar L_0f-\bar L_0(fq)
\]
and V35's $\bar L_0W_B=\lambda W_B$.
Lemma~\ref{c18v36:endpoints} gives
\[
 \Pi_{\mathcal Z}\Gamma_0(W_F,W_B)=0,\qquad
 \Pi_{\mathcal Z}\Gamma_0(W_B,W_B)
                              =\lambda|\mathcal S|/4 .
\]
The remaining $\Gamma_0(W_F,W_F)$ is independent of $g$.
Its derivative, and that of the displayed constant, vanish.
Apply \eqref{c18v36:t2} to prove \eqref{c18v36:zeroselection}.

Finally $\widetilde r_\beta p_{g,\beta}
=\beta^2t_2+O_{N,\ell}(\beta^3)$ in Haar $L^2$.
Pair against weighted unit vectors of $\mathcal Z$ and use
the fixed-regulator equivalence of Haar and weighted norms.
This proves \eqref{c18v36:weighted}.
\end{proof}

In particular $\int t_2dy=0$. The projection in
\eqref{c18v36:weighted} is not a low-frequency spectral
projection for $\bar L_\beta$. No exclusion at third order,
or at fixed nonzero continuum coupling, follows.

\subsection{A nonzero adjoint component with exact inverse cost}

Let $P_{\rm ev}$ average the independent sign changes
$y_f\mapsto\epsilon_f y_f$, $\epsilon_f\in\{1,-1\}$, on
every free chord. This is distinct from $\Pi_{\mathcal Z}$.
It commutes with $\bar L_0$ because it commutes with all
left and right link derivatives.

\begin{theorem}[Exact even-chord component]
\label{c18v36:adjoint}
The even component of the second native source is
\[
 t_{2,{\rm ev}}=P_{\rm ev}t_2
       =-\frac{11}{6912}D_g\sum_{p\in\mathcal S}\chi_p .
 \tag{V36.11}\label{c18v36:even}
\]
It obeys
\[
 \bar L_0t_{2,{\rm ev}}=\lambda_{\rm ad}t_{2,{\rm ev}},
 \quad \lambda_{\rm ad}=\frac83\lambda=32-\frac{16}{\ell},
 \quad \mathcal Gt_{2,{\rm ev}}=16t_{2,{\rm ev}} .
 \tag{V36.12}\label{c18v36:adeigen}
\]
Writing $\Sigma_1$ for V35's full coarse covariance,
\[
 \int t_{2,{\rm ev}}t_{2,{\rm ev}}^*dy
                  =\frac{121}{124416}\Sigma_1(g).
 \tag{V36.13}\label{c18v36:adcov}
\]
\end{theorem}
\begin{proof}
Each fundamental plaquette coefficient is odd on every
chord in its support. Distinct supports leave an odd chord
in any product of two different plaquette coefficients;
left and right derivatives preserve this parity.
Each fused outer loop likewise has an odd chord.
Thus only the self terms survive $P_{\rm ev}$ in
\eqref{c18v36:t2}.

For a source plaquette its $k$-chord open path has endpoint
cells distinct. Internal path rotations cancel, and the
two endpoint Green cross coefficients are zero. The
round-sphere coordinate identity therefore gives
\[
 \Gamma_0(w_p,w_p)
       =(k+h(a)+h(b))(1-w_p^2)
       =\frac{\lambda}{3}(1-w_p^2).
\]
Its coarse derivative gives
$\Gamma_0(w_p,D_gw_p)=-(\lambda/3)w_pD_gw_p$.
The self contribution to $t_2$ is consequently
\[
 -\left(\frac1{288}+\frac1{108}\right)w_pD_gw_p
       =-\frac{11}{6912}D_g\chi_p.
\]
This proves \eqref{c18v36:even}.
For an adjoint path coefficient each link Casimir is eight,
and each endpoint Casimir is eight. The same calculation
as V35 gives $8(k+h(a)+h(b))=(8/3)\lambda$ and current
Casimir sixteen, proving \eqref{c18v36:adeigen}.

For a unit quaternion with scalar coordinate $z$, a unit
tangent derivative of $\chi=4z^2-1$ is $8zx$ up to sign,
where $x$ is an orthogonal coordinate.
The Haar fourth moment is $\int z^2x^2=1/24$, so its
variance is $8/3$. The corresponding variance for
$(D w)/6$ is $1/144$. Their ratio is $384$.
For a bank on a two-$\mathcal B$ plaquette the tangent
coefficient has the same endpoint-frame difference as in
V35; the preceding calculation multiplies its squared
norm. Distinct plaquette adjoint coefficients are
orthogonal by integrating a chord in their symmetric
difference. Hence the complete unary and incoming-pair
covariance is $384\Sigma_1$ before multiplying by
$(11/6912)^2$. This gives \eqref{c18v36:adcov}.
\end{proof}

Define, on Haar-mean-zero functions,
\[
 \Sigma_2(g)=\int t_2t_2^*dy,\qquad
 R_2(g)=\int t_2\bar L_0^{-1}t_2^*dy .
\]
Since $P_{\rm ev}$ is an orthogonal reducing projection,
the even and odd inverse energies add without cross terms.
Thus
\[
 R_2(g)\succeq
       \frac{121}{124416\,\lambda_{\rm ad}}\Sigma_1(g)
       \succ0.
 \tag{V36.14}\label{c18v36:lower}
\]
In particular V35's cancellation does not extend accidentally
to second order. Its leading remaining inverse cost is
\[
 R^{A^{[1]}}_\beta(g)
                 =\beta^4R_2(g)+O_{N,\ell}(\beta^5).
 \tag{V36.15}\label{c18v36:secondcost}
\]
The fixed-regulator variational inverse expansion used in
V35 justifies this identity. The known part of the source is
an adjoint charged sector, not the zero-current sector.

\subsection{An explicit volume-independent bound for the whole coefficient}

The lower bound just found does not estimate all of $t_2$.
Here is a deliberately nonoptimal upper bound for its
complete coarse operator norm. Put
\[
 C_\ell=3(\ell-1)(\ell^3-1),\qquad
 B_\ell=\frac{173}{936}
                 +\frac{8/9+(8/3)C_\ell}{\lambda}.
 \tag{V36.16}\label{c18v36:B}
\]

\begin{theorem}[Full second-coefficient upper bound]
\label{c18v36:upper}
For all $n\ge5$ and every $g$,
\[
 \Sigma_2(g)\preceq6561B_\ell^2 I,\qquad
 R_2(g)\preceq2187B_\ell^2 I .
 \tag{V36.17}\label{c18v36:upperbound}
\]
These are independent of $n$, not of $\ell$.
\end{theorem}
\begin{proof}
Fix one unit coarse frame column $i=(e,a)$.
There are four source plaquettes incident to its last-chain
link and at most twelve distinct edge-sharing neighbours
of each. Thus at most forty-eight unordered pairs can
contribute to $D_iF_2$.
Use $|D_i\chi_p|\le8$,
$|D_i(w_pw_q)|\le2$ and $|D_iR_{pq}|\le1$.
The last bound holds because that link occurs once on the
outer loop, or cancels when it is the shared edge.
It follows that
\[
 |D_iF_2|\le\frac{4\cdot8}{2304}
       +48\left(\frac1{702}+\frac2{936}\right)
       =\frac{173}{936}.
\]

The finite Dirichlet matrix $H_\ell$ has nonpositive
off-diagonal entries and is positive definite. Its
maximum principle gives $h_{xz}\ge0$.
The Green diagonal from V35 is at most $3(\ell-1)$.
Positive definiteness gives
$|h_{xz}|\le\sqrt{h_{xx}h_{zz}}$, whence each cell row
sum is at most $C_\ell$.
For every chord $f$ and internal vertex $x$,
\[
 |\nabla_f W_g|\le4,\qquad |X_xW_g|\le12.
\]
There are four plaquettes at a link and twelve at a
vertex. If both incident edges of a plaquette are
chords, their gauge variations cancel at that vertex;
if exactly one is a chord its contribution has norm at
most one. These observations prove the bounds.
V35's source paths have at most two chords and exactly
two current endpoints. Their mixed coarse/chord
derivatives have norm at most one, as follows directly
from their quaternion matrix coefficients. Therefore
\[
 \sum_f|\nabla_f(r_1)_i|\le\frac43,\qquad
 \sum_x|X_x(r_1)_i|\le\frac43.
\]
Applying the preceding row bound in
\eqref{c18v36:gamma} gives
$|\Gamma_0(W_g,(r_1)_i)|\le16/3+16C_\ell$.
Since the drift term in $t_{2,i}$ is
$\Gamma_0(W_g,(r_1)_i)/(6\lambda)$,
we obtain $|t_{2,i}|\le B_\ell$.

It remains to avoid a volume factor when converting this
column estimate into a matrix estimate.
Assign each chord to the cell of its positively oriented
tail. The variables of $t_{2,i}$ all lie in cells at
periodic coarse $\ell^\infty$ distance at most two from
the head of $e$. For $D_iF_2$ this follows by adding one
edge-sharing plaquette to a plaquette incident to the
last-chain link. For the link-gradient part of $\Gamma_0$
only plaquettes meeting a source chord contribute.
For its current part, the source has only its two
endpoint currents; $h_{xz}$ couples an endpoint only to
vertices in that same cell. A plaquette incident to such
a vertex has its chord tails at most one further cell
away. The endpoints of a source lie at most one cell
from the coarse head. These facts establish the claimed
radius-two support.

Every column has Haar mean zero by
\eqref{c18v36:zeroselection}. Columns with disjoint
radius-two cell boxes depend on disjoint Haar variables,
so their covariance is zero. A given box can overlap
only boxes whose heads are within periodic radius four.
There are at most $9^3$ such heads and nine coarse
columns per head. Each covariance row therefore has at
most $9^4=6561$ potentially nonzero entries, each of
absolute value at most $B_\ell^2$. The symmetric Schur
bound proves the first inequality in
\eqref{c18v36:upperbound}.
Finally $\bar L_0\ge3$ on centred functions, as in V35,
so its inverse quadratic form is at most one third of
the covariance form. This proves the second inequality.
\end{proof}

Combining \eqref{c18v36:lower} and
\eqref{c18v36:upperbound} sandwiches the entire second
inverse-cost coefficient between explicit positive
operators. No bound for a remainder in
\eqref{c18v36:secondcost} has been inferred from this.

\subsection{A finite local inverse and the second native correction}

There is a constructive version of the preceding use of
$\bar L_0^{-1}$. Expand \eqref{c18v36:t2} into
coarse derivatives of the self and pair terms of $F_2$
and into terms $\Gamma_0(w_q,D_iw_p)/(36\lambda)$,
where $p$ contains the selected last-chain link.
Keep each such full contraction as one term $\eta$.

\begin{lemma}[Local Peter--Weyl inversion]
\label{c18v36:finite}
Every nonzero $\eta$ is Haar centred, depends on at most
six free chords, and has spin at most one on each chord.
On its chord support $K$, let $\mathcal V_K$ be the
tensor product of the matrix-coefficient spaces of spins
$0,\frac12,1$, with the global constant removed.
Then
\[
 \dim\mathcal V_K\le14^6-1,\qquad
 \bar L_0\mathcal V_K\subseteq\mathcal V_K,\qquad
 \bar L_0|_{\mathcal V_K}\succeq3I.
 \tag{V36.18}\label{c18v36:finiteblock}
\]
Each local inverse is consequently a positive finite
matrix inverse with known entries from $h$ and the
$SU(2)$ representation generators.
\end{lemma}
\begin{proof}
A source plaquette has at most two chords, and any other
plaquette has at most four. Their union has size at most
six, also for a fused-loop term.
Left and right derivatives preserve each link spin,
while a product of two fundamental coefficients has
only spins zero and one on a common link.
Self and fused terms obey the same upper limit.

Each differentiated $F_2$ summand has zero mean by the
individual integrals in Lemma~\ref{c18v36:normalization}.
For a contraction term, integration by parts and V35
give
\[
 \int\Gamma_0(w_q,D_iw_p)\,dy
       =\lambda\int w_qD_iw_p\,dy=0.
\]
For $p\ne q$ integrate a chord in the symmetric
difference; for $p=q$ differentiate the constant
integral $\int w_p^2dy=1/4$.

Link Laplacians and link left/right generators preserve
matrix-coefficient spaces and do not introduce new
link variables. The coefficients $h_{xz}$ are fixed
geometric scalars. Therefore $\bar L_0$ preserves the
finite cylinder space, and its self-adjointness
preserves its centred subspace.
The dimension of the spin $0,\frac12,1$ space on a
link is $1+4+9=14$.
Finally its form dominates the product-Haar form,
whose centred gap is three. This proves all claims.
\end{proof}

For each column define the smooth function
\[
 (\phi_2)_i=\sum_{\eta\ {\rm in}\ t_{2,i}}
               (\bar L_0|_{\mathcal V_{K(\eta)}})^{-1}\eta .
 \tag{V36.19}\label{c18v36:phi2}
\]
The finite sum solves
$\bar L_0\phi_2=t_2$ and $\int\phi_2dy=0$.
It is the unique centred global solution, but
\eqref{c18v36:phi2} constructs it from local finite
matrices rather than an unknown interacting inverse.
No assertion is made that all these matrices have
been numerically inverted. The formula and their
strict positivity define an exact algebraic
construction. Each term preserves its variable
support. Gauge equivariance follows either by
intertwining these local spaces or by uniqueness of
the centred global solution and equivariance of the
operator and source.

\begin{theorem}[Second explicit-order native correction]
\label{c18v36:correction}
Set
\[
 A^{[2]}=A^{\rm h}+
          \beta\nabla_{\rm fib}\phi_1+
          \beta^2\nabla_{\rm fib}\phi_2 .
 \tag{V36.20}\label{c18v36:A2}
\]
It retains all coarse outputs and satisfies
\[
 \begin{aligned}
 Dc\,A^{[2]}&=I,\\
 r^{A^{[2]}}_\beta&=O_{N,\ell}(\beta^3),\\
 R^{A^{[2]}}_\beta(g)&=O_{N,\ell}(\beta^6),\\
 \mathbb E_{\nu_{g,\beta}}
       |\beta^2\nabla_{\rm fib}(\phi_2\cdot v)|^2
       &=\beta^4v^*R_2(g)v+O_{N,\ell}(\beta^5)|v|^2 .
 \end{aligned}
 \tag{V36.21}\label{c18v36:corrected}
\]
\end{theorem}
\begin{proof}
Both gradient corrections are vertical and smooth in
the original fibre metric. Their conditional
divergences are $-\beta\bar L_\beta\phi_1$ and
$-\beta^2\bar L_\beta\phi_2$ and have zero conditional
mean. Thus the full output is unchanged, and
\eqref{c18v36:t2} together with
$\bar L_0\phi_2=t_2$ cancels the second coefficient.
Fixed-regulator density comparison with Haar bounds
the remaining conditional inverse for sufficiently
small $\beta$, yielding the order-six cost.
At $\beta=0$, integration by parts gives
$\int|\nabla_{\rm fib}(\phi_2\cdot v)|^2d\nu_{g,0}
=v^*R_2v$. Smooth density comparison supplies the last
line. This line concerns the added second correction;
it is not a claim that cross terms in the total lift
norm vanish.
\end{proof}

\subsection{Verification, scope and the next upstream test}

The exact diagnostic enumerates all plaquette chord
supports on fine tori of sides $10,15,20$ at
$\ell=2,3,4$, verifies all $1125$ source endpoint
pairs per torus, and checks representative local
pair supports in each coordinate direction.
Translation covers the other coarse heads.
Rational quaternion second derivatives independently
verify the shared-link identity, the eigenvalues
eighteen and thirty-two, and the fundamental and
adjoint open-path Casimirs. Exact Haar moments and
the full nine-component coarse bank reproduce
\eqref{c18v36:adcov}. Four sixteen-dimensional
unperturbed blocks are inverted exactly, including
a two-chord same-cell example with nonzero Green
cross terms. The latter is not replaced by a
diagonal-current approximation. All checks pass.
These regressions complement the proofs; they do
not compute the many-link native vacuum, certify
all finite blocks, or bound the Taylor remainder.

The advancement over V35 is the actual second
source \eqref{c18v36:t2}, its zero-current
cancellation, its nonzero adjoint inverse cost,
a bound for its entire coefficient, and a
constructive local second correction.
Merely repeating this at successively higher
orders would not address convergence. The next
upstream test should instead seek a local norm
in which the connected vacuum/source recursion
and conditional correction have a radius and
remainder independent of spatial volume.
Even a successful small-$\beta$ result would
remain an electric/strong-coupling theorem; a
bridge to the required weak-coupling continuum
regime would still need proof.

Positive spatial-flow transport, variable-bank
derivative costs, projection-energy stability,
native conditional/coarse variance and physical
scale calibration remain separate obligations.
No companion finite defect gap, unperturbed
Casimir, or formal perturbative cancellation is
identified with the physical gap.
The interacting four-dimensional continuum
construction and positive physical Yang--Mills
mass gap remain \emph{unproved}.
% END CYCLE 18 VERSION 36 APPEND
% BEGIN CYCLE 18 VERSION 37 APPEND
\clearpage
\section[V37: a uniform native log-vacuum and conditional remainder]{V37: a uniform native log-vacuum and conditional remainder}
\label{c18v37:context}

V36 controls two native score coefficients, but its error
constants still depend on volume. We now address that
dependence rather than compute a third coefficient.
An energy-resolved Fourier interaction norm closes the
quadratic equation for the \emph{actual vacuum logarithm}.
The electric denominators pay both the derivatives and
the number of overlapping links. This yields a
volume-independent analytic disk and a pointwise Hessian
bound, including on the exact rooted section.
The latter gives a uniform gap for the native
\emph{reduced conditional generator}. Combining it with
V36's local correction gives a genuinely volume-uniform
order-six inverse-cost remainder.

The new conclusion is restricted to an explicit small
electric coupling:
\[
 \beta_*=\frac{3}{512\mathrm e},\qquad
 \rho_* =2-\frac3{\mathrm e}>0.
 \tag{V37.1}\label{c18v37:constants}
\]
Here $\mathrm e=\exp(1)$, not a lattice edge.
It is not a passage to continuum weak coupling.
The full interacting four-dimensional mass-gap
objective remains open.

\subsection{What is inherited, and what is not being repeated}

The current V36 source, its preservation check and its
mathematical diagnostics were revalidated.
V5 already proves a small-coupling lattice construction
using energy-weighted creation operators; V12 dresses
the physical plaquette band. Those results are not
repackaged as a new qualitative lattice gap.
Their Hilbert-space coefficient norm does not by itself
give the pointwise conditional log-density Hessian
needed here. Archived f1/V29 also uses a trace-class
Fourier norm, but for a different $SU(3)$ transfer
operator and a stated fusion-cumulant gate. We do not
assume or reprove that gate. The present continuous
Hamiltonian logarithmic equation has a quadratic
gradient nonlinearity, whose bound is proved below.

For the last curvature-to-gap step, the primary source
is D.~Bakry and M.~\'Emery,
\emph{Diffusions hypercontractives},
\href{https://numdam.org/item/SPS_1985__19__177_0/}
{S\'eminaire de probabilit\'es XIX (1985), 177--206},
Proposition 3 and the integrated spectral observation
on p.~203. Their weighted Bochner identity is applied
only after the native Hessian has been bounded.
We include the short compact-manifold spectral argument.
No curvature assumption is imported as a conclusion.
This is a targeted check, not the scheduled broad
literature sweep. The next companion and broad
literature checkpoints remain V40; the five new
companions reviewed in V36 remain unchanged.

Use the same $SU(2)$ round metric, $T=-\sum_e\Delta_e$,
$W=\sum_pw_p$, $N=\ell n$, $\ell\ge2$, $n\ge5$,
and spatial flow time $s=0$. The coupling is
$\beta=b/\kappa$ after dividing the original
Hamiltonian by $\kappa$. No physical-time factor is
inserted into the conditional generator.

\subsection{Fourier blocks and electric denominators}

For a finite set of fine links, irreducible product
representations are indexed by nonnegative integers
$\mathbf s=(s_e)$, with spin $s_e/2$ and active set
$I(\mathbf s)=\{e:s_e>0\}$. Write
\[
 f_{\mathbf s}(U)=\Tr(A_{\mathbf s}\pi_{\mathbf s}(U)),
 \qquad
 \|f\|_{\rm A}=\sum_{\mathbf s}\|A_{\mathbf s}\|_1,
 \qquad
 E_{\mathbf s}=\sum_e s_e(s_e+2).
 \tag{V37.2}\label{c18v37:Fourier}
\]
The matrix trace norm absorbs the usual Fourier
dimension factor. Thus $Tf_{\mathbf s}=
E_{\mathbf s}f_{\mathbf s}$. A constant block has
$E_{\mathbf0}=0$ and will be omitted whenever
$T^{-1}$ is used.

\begin{lemma}[Product and derivative estimates]
\label{c18v37:Fourier-lemma}
The norm $\|\cdot\|_{\rm A}$ dominates the uniform norm
and is submultiplicative. Removing the constant
Fourier block is contractive. For a unit left
derivative $D_{e,v}$ on link $e$,
\[
 \|D_{e,v}f_{\mathbf s}\|_{\rm A}
      \le s_e\|A_{\mathbf s}\|_1,\qquad
 \|D_{e,v}D_{f,w}f_{\mathbf s}\|_{\rm A}
      \le s_es_f\|A_{\mathbf s}\|_1 .
 \tag{V37.3}\label{c18v37:derivatives}
\]
For every nonconstant block, with $m=|I(\mathbf s)|$,
\[
 \frac{s_e}{E_{\mathbf s}}\le\frac13,\qquad
 \frac{\sum_fs_f}{E_{\mathbf s}}\le\frac13,\qquad
 \frac{s_e\sum_fs_f}{E_{\mathbf s}}\le m.
 \tag{V37.4}\label{c18v37:denominators}
\]
\end{lemma}
\begin{proof}
For unitary $\pi$, $|\Tr(A\pi)|\le\|A\|_1$.
A product of coefficients is represented by
$A\otimes B$ in $\pi\otimes\sigma$.
Decompose this representation unitarily into
irreducibles with multiplicity spaces, pinch to
its diagonal isotypic blocks, and take a partial
trace over each multiplicity space.
Pinching is trace-norm contractive, and
$\|\Tr_{\rm mult}X\|_1\le\|X\|_1$ by trace-norm
duality. Summing the blocks proves
submultiplicativity without a dimension multiplier.
Deleting a summand of the defining norm proves the
constant-projection assertion.

In spin $s/2$, a unit round-metric generator has
operator norm $s$ and Casimir $s(s+2)$.
Multiplying the coefficient matrix by one or two
such generators proves \eqref{c18v37:derivatives}.
Since $s(s+2)\ge3s$ for integer $s\ge1$,
the first two denominator bounds follow.
For the last use
$2s_es_f\le s_e^2+s_f^2$ and
$E_{\mathbf s}\ge\sum_fs_f^2$:
the ratio is at most $(m+1)/2\le m$.
\end{proof}

These facts apply to arbitrarily high spins;
no electric cutoff is introduced.

\subsection{The connected family norm and a quadratic contraction}

A label $Z$ is a nonempty connected set of plaquettes,
where sharing a fine edge defines adjacency. Retain
coefficients $q_{Z,\mathbf s}$ with
$\varnothing\ne I(\mathbf s)\subseteq E(Z)$.
The cover $Z$ need not be the minimal support.
For $a\ge0$ set
\[
 \|q\|_a=\sup_e
       \sum_{\substack{Z,\mathbf s\\ e\in E(Z)}}
          e^{a|Z|}\|A_{Z,\mathbf s}\|_1,\qquad
 \mathcal Lq=\sum_{Z,\mathbf s}q_{Z,\mathbf s}.
 \tag{V37.5}\label{c18v37:norm}
\]
At each finite volume this is a Banach family space;
the finite set of covers is combined with a weighted
$\ell^1$ sum of trace-class Fourier blocks. Its
physical lift can have an extensive norm. No
volume-independent bound for that extensive norm
is asserted.

Let $\Gamma_{\rm f}(u,v)=
\sum_{e,a}(D_{e,a}u)(D_{e,a}v)$ denote the
\emph{fine product-metric} gradient contraction,
not V36's reduced $\Gamma_0$.
Define a bilinear family map $\mathcal B(q,r)$ by
expanding
$\Gamma_{\rm f}(T^{-1}\mathcal Lq,
T^{-1}\mathcal Lr)$, keeping cover $Z\cup Y$
for each input pair, resolving the product into
Fourier blocks, and deleting its constant block.
A term vanishes unless the active sets share a
link; every remaining cover union is connected.

\begin{theorem}[No-loss native quadratic estimate]
\label{c18v37:bilinear}
For all finite volumes and every $a\ge0$,
\[
       \|\mathcal B(q,r)\|_a
                      \le\frac23\|q\|_a\|r\|_a .
 \tag{V37.6}\label{c18v37:Bbound}
\]
No spatial weight is spent at an iteration.
\end{theorem}
\begin{proof}
The contribution of one pair of blocks is bounded by
\[
 3\sum_{f\in I(\mathbf s)\cap I(\mathbf t)}
      \frac{s_f}{E_{\mathbf s}}
      \frac{t_f}{E_{\mathbf t}}
      \|A_{Z,\mathbf s}\|_1\|A_{Y,\mathbf t}\|_1 .
\]
The factor three counts the orthonormal colour
directions. Fix an anchor edge $e$ of the output
cover. It belongs to $E(Z)$ or $E(Y)$; counting
both possibilities is an upper bound.
In the first case fix $(Z,\mathbf s)$ with
$e\in E(Z)$. For each common active link $f$,
the anchored sum over $(Y,\mathbf t)$ is bounded
by $\|r\|_a/3$, because
$t_f/E_{\mathbf t}\le1/3$.
The remaining sum $\sum_fs_f/E_{\mathbf s}$ is
at most $1/3$. Hence this case costs at most
$\|q\|_a\|r\|_a/3$.
The other case gives the same bound.
Finally
$e^{a|Z\cup Y|}\le e^{a|Z|}e^{a|Y|}$.
These estimates prove \eqref{c18v37:Bbound},
including covers larger than their active supports.
\end{proof}

Represent $W$ by its one-plaquette family $\mathsf W$.
Expanding $\frac12\Tr(U_1U_2U_3U_4)$ gives sixteen
products of matrix entries with coefficient $1/2$,
so $\|w_p\|_{\rm A}\le8$. Independent link inversions
preserve this norm. Each edge belongs to four
plaquettes; consequently
\[
                 \|\mathsf W\|_a\le J_a:=32e^a .
 \tag{V37.7}\label{c18v37:Wnorm}
\]

\begin{theorem}[A volume-uniform analytic native logarithm]
\label{c18v37:log}
Let
\[
 R_0=\frac38,\qquad
 \beta_a=\frac{3}{512e^a}.
\]
For complex $|\beta|\le\beta_a$ the map
\[
       q\longmapsto\beta\mathsf W+\mathcal B(q,q)
 \tag{V37.8}\label{c18v37:fixedpoint}
\]
has a unique fixed point $q_\beta$ in the
$\|\cdot\|_a$ ball of radius $R_0$.
It is holomorphic in the open disk, and
$\|q_\beta\|_a\le2J_a|\beta|$.
For real $\beta$ in this disk, put
$u_\beta=T^{-1}\mathcal Lq_\beta$. Then
\[
 \Omega_\beta(U)=
       \frac{\exp u_\beta(U)}
       {(\int\exp(2u_\beta)\,dU)^{1/2}}
 \tag{V37.9}\label{c18v37:vacuum}
\]
is the actual positive normalized ground state of
$T-\beta W$.
\end{theorem}
\begin{proof}
At the edge of the disk,
$|\beta|J_a\le3/16$. The ball is mapped into the
ball of radius
$3/16+(2/3)(3/8)^2=9/32<3/8$.
The Lipschitz constant on it is at most
$(4/3)R_0=1/2$. Banach contraction gives existence
and uniqueness. Uniform Picard convergence on
complex subdisks gives holomorphy.
The smaller ball of radius $2J_a|\beta|$ is also
invariant and contractive, giving the claimed
bound. All constants are independent of volume.

The fixed point preserves reality for real
$\beta$ and full fine-gauge invariance, because
the initial Wilson family, Fourier projections,
electric inverse and gradient contraction
intertwine those actions. Its lift obeys
\[
 Tu_\beta=\beta W+\Gamma_{\rm f}(u_\beta,u_\beta)
            -\int\Gamma_{\rm f}(u_\beta,u_\beta)\,dU.
 \tag{V37.10}\label{c18v37:logPDE}
\]
The Fourier estimates give convergent first and
second derivatives at fixed volume. Thus
$u_\beta$ is $C^2$; elliptic regularity applied
to \eqref{c18v37:logPDE} then bootstraps to
smoothness. The exponential is positive and
solves the eigenvalue equation, with
$E_\beta=-\int\Gamma_{\rm f}(u_\beta,u_\beta)dU$.
For any smooth $f$, its ground-state transform gives
\[
 \langle e^{u_\beta}f,(T-\beta W-E_\beta)
                         e^{u_\beta}f\rangle
       =\int e^{2u_\beta}|\nabla f|^2\,dU\ge0.
\]
Connectedness of the compact configuration space
makes the kernel one-dimensional. This identifies
the eigenfunction with the native ground state,
without analytic continuation from a
volume-dependent isolation radius.
\end{proof}

The construction fixes $\int u_\beta\,dU=0$.
The extensive scalar needed to normalize
$\Omega_\beta$ is retained in
\eqref{c18v37:vacuum}, not assigned a uniform bound.
Only its derivatives along the configuration
variables disappear.

\subsection{Uniform derivative tails and the actual conditional gap}

For a family $q$ let $v=T^{-1}\mathcal Lq$.
For $a>0$, Lemma~\ref{c18v37:Fourier-lemma} gives
\[
 \begin{aligned}
 \sup_{e,b}\|D_{e,b}v\|_\infty&\le\frac13\|q\|_a,\\
 \sup_e\sum_f\|\operatorname{Hess}_{ef}v\|_{\rm op,\infty}
                         &\le\frac4{\mathrm e\,a}\|q\|_a,\\
 \| (D_{e,b}D_{f,c}v)_{(e,b),(f,c)}\|_{\rm op,\infty}
                         &\le\frac{12}{\mathrm e\,a}\|q\|_a .
 \end{aligned}
 \tag{V37.11}\label{c18v37:C2}
\]
In the last line it suffices to use the scalar
row-and-column Schur bounds; derivatives need
not commute on the same link.
For the Hessian, same-link blocks are the
symmetrized second derivatives in a
bi-invariant orthonormal frame.
Their norm is bounded by $s_e^2$, and distinct-link
blocks by $s_es_f$. Sum
\eqref{c18v37:denominators}, use
$|I(\mathbf s)|\le4|Z|$, and then
$|Z|e^{-a|Z|}\le1/(\mathrm e a)$.
This proves the second line; summing the three
scalar colour columns proves the third.
The first line follows directly from
$s_e/E_{\mathbf s}\le1/3$.
All bounds also hold after freezing any subset
of links to arbitrary group values.

Write $q_\beta=\sum_{k\ge1}\beta^kq_k$ and
$u_k=T^{-1}\mathcal Lq_k$. Cauchy's inequality
and the uniform ball bound give, for
$t=|\beta|/\beta_a<1$,
\[
 \left\|q_\beta-\sum_{k=1}^m\beta^kq_k\right\|_a
       \le Q_m(t):=\frac38\,\frac{t^{m+1}}{1-t},
       \qquad m\ge0.
 \tag{V37.12}\label{c18v37:tails}
\]
For $m=0$ the sum is empty. Equation
\eqref{c18v37:C2} applies to each of these tails.
The coefficients of the \emph{lifted} logarithm
agree with V35--V36:
$u_1=W/12$ and $u_2=F_2/2$.
The scalar normalization in V36 is absent here
only because $u_\beta$ is Haar centred.

Now fix $a=1$ and real $|\beta|\le\beta_*$. The
exact rooted section freezes the forest links
to the identity and the last-chain links to
$g$, leaving the independent chord variables
$y$. It is a totally geodesic product section
of the fine product manifold. Put
$u_\beta^g(y)=u_\beta(W^0(g,y))$. Its conditional
law is exactly
\[
 p_{g,\beta}(y)=
     \frac{\exp(2u_\beta^g(y))}
          {\int\exp(2u_\beta^g)\,dy}.
 \tag{V37.13}\label{c18v37:conditional}
\]
The full fine normalization cancels from this
ratio, not from the native vacuum itself.

\begin{theorem}[Uniform reduced conditional coercivity]
\label{c18v37:gap}
For every $g$, every $n\ge5$, every $\ell\ge2$
and every real $|\beta|\le\beta_*$,
\[
 \operatorname{Var}_{p_{g,\beta}}f
       \le\frac1{\rho_*}\int|\nabla_{\mathcal C}f|^2p_{g,\beta}dy
       \le\frac1{\rho_*}
                 \int\Gamma_0(f,f)p_{g,\beta}dy.
 \tag{V37.14}\label{c18v37:Poincare}
\]
Thus $\bar L_\beta\ge\rho_* I$ on centred
reduced functions.
\end{theorem}
\begin{proof}
Each round $SU(2)$ factor has Ricci tensor $2I$.
By \eqref{c18v37:C2} and $\|q_\beta\|_1\le3/8$,
the operator norm of the conditional
log-density Hessian is at most $3/\mathrm e$.
Consequently
$\operatorname{Ric}-\operatorname{Hess}\log p_{g,\beta}
\succeq\rho_*I$ on the product of chords.

For the product-metric positive weighted
Laplacian $L_{\rm prod}$, integrated weighted
Bochner gives
$\|L_{\rm prod}f\|_2^2
\ge\rho_*\int|\nabla_{\mathcal C}f|^2p$.
Apply this to a nonconstant eigenfunction
on the compact connected product: its eigenvalue
$\mu$ satisfies $\mu^2\ge\rho_*\mu$ and is
therefore at least $\rho_*$.
Spectral expansion proves the first inequality.
V34's exact induced inverse metric satisfies
$a_0\succeq I$, proving the second.
\end{proof}

This is a theorem about the exact \emph{reduced}
conditional law and original induced metric.
It does not assert the same floor for all
unreduced gauge-orbit functions. V33 identifies
the reduced inverse with the native full-fibre
inverse on the invariant scores used here.
Nor is the product curvature gap silently
identified with a physically normalized
continuum Hamiltonian gap.

\subsection{Quantitative local banks for the two corrections}

To turn the log-density remainder into a full
coarse matrix estimate, metric derivatives must
also be paid. Retain $C_\ell,B_\ell$ from
\eqref{c18v36:B} and set
\[
 \begin{aligned}
 A_\ell&=1+12C_\ell,&
 P_1&=\frac4{3\lambda},&
 P_2&=\frac{14^3}{3}B_\ell,\\
 K_j&=104976P_j
        \left[A_\ell(1+\mathrm e^{-1})+4C_\ell\right],
       &&j=1,2.
 \end{aligned}
 \tag{V37.15}\label{c18v37:bankconstants}
\]
These deliberately large constants are finite
and explicit functions of $\ell$, not of $n$
or $g$. Let $\phi_1=r_1/\lambda$ and let
$\phi_2$ be V36's finite local inverse.

\begin{lemma}[Original-metric bank derivative bound]
\label{c18v37:bank}
For $v=T^{-1}\mathcal Lq$ and its rooted
restriction $v^g$,
\[
 \|D_y\Gamma_0(v^g,\phi_j)\|_{\rm op,\infty}
                         \le K_j\|q\|_1,\qquad j=1,2.
 \tag{V37.16}\label{c18v37:bankbound}
\]
The matrix has every chord derivative row
and every coarse-output column.
\end{lemma}
\begin{proof}
We give the constants to avoid assuming a
uniform local-bank estimate.
Each column of $\phi_1$ is a sum of four
fundamental paths of at most two chords.
The Fourier norm of such a path coefficient
or its unit coarse derivative is at most two.
Thus its total termwise Fourier norm is at
most $P_1$.
For $\phi_2$, V36 decomposes each source
column into centred terms $\eta$ of at most
six chords and spin at most one.
The termwise estimates in the proof of
Theorem~\ref{c18v36:upper}, taken before summing
over the incident plaquettes, give
$\sum_\eta\|\eta\|_\infty\le B_\ell$.
On such a Peter--Weyl space, Cauchy--Schwarz
and Plancherel give
$\|f\|_{\rm A}\le14^3\|f\|_2$:
the sum of squared representation dimensions
is at most $14^6$.
The centred inverse of $\bar L_0$ has
$L^2$ norm at most $1/3$ and preserves that
space. Hence the total termwise Fourier norm
of $\phi_2$ is at most $P_2$.
Both banks have chord-tail support in the
radius-two cell box of their coarse head.

Use scalar orthonormal derivative indices
$\alpha,\delta,\gamma$ on chords and a
scalar coarse index $i$. For either bank
write
$b_{\alpha i}=\sum_\delta(a_0)_{\alpha\delta}
D_\delta\phi_i$.
The induced metric is $a_0=I+JhJ^*$.
A chord has at most two endpoints; a vertex
has at most six incident chords. Since $h$
has nonnegative entries and row sums at most
$C_\ell$, the block row and column sums of
$a_0$ are at most $A_\ell$.
The scalar sums are at most $3A_\ell$.
Each nonconstant block of $J$ is a signed
$\operatorname{Ad}_{y_f}$, with first
unit-derivative norm at most two.
Differentiating either factor of $JhJ^*$,
summing three derivative directions, and
then three scalar column entries gives
\[
 \sup_\delta\sum_{\gamma,\alpha}
       \|D_\gamma(a_0)_{\alpha\delta}\|_\infty
 \le432C_\ell ,
 \tag{V37.17}\label{c18v37:metricderivative}
\]
and the analogous row estimate.
Indeed the undifferentiated block incidence
sum is at most $12C_\ell$; the two factors
and three directions contribute at most
$12$, giving $144C_\ell$ before the scalar
factor three.

An individual bank term has at most eighteen
scalar derivative directions, with first
derivative factor at most two and second
derivative factor at most four. For total
termwise norm $P$ this gives
\[
 \begin{aligned}
 \sup_i\sum_\alpha\|b_{\alpha i}\|_\infty
                         &\le108A_\ell P,\\
 \sup_i\sum_{\gamma,\alpha}
             \|D_\gamma b_{\alpha i}\|_\infty
                         &\le(15552C_\ell+3888A_\ell)P.
 \end{aligned}
 \tag{V37.18}\label{c18v37:bcolumns}
\]
The factors are respectively
$3A_\ell\cdot36$ and
$432C_\ell\cdot36+3A_\ell\cdot1296$.

Multiplication by $a_0$ increases the
chord-tail support radius to at most four:
from a radius-two chord one goes to an
endpoint cell, stays in that cell through
$h$, then goes to an incident chord tail.
A derivative of $a_0$ introduces only one
of those same chords. Therefore, for a
fixed derivative row, at most $9^4=6561$
coarse scalar columns can contribute.
Schur bounds give operator norms at most
$81$ times the respective column bounds
in \eqref{c18v37:bcolumns}. The second bound
is interpreted as a bound on the matrix
with entries $\sum_\alpha
\|D_\gamma b_{\alpha i}\|_\infty$.

Finally differentiate
$\Gamma_0(v^g,\phi_i)=
\sum_\alpha b_{\alpha i}D_\alpha v^g$.
The term with two derivatives of $v^g$
is bounded by
$(12/\mathrm e)\|q\|_1$ times the first
bank norm. The other is bounded by
$\|q\|_1/3$ times the second.
Using \eqref{c18v37:bcolumns}, their sum is
\[
 81P\left(\frac{1296A_\ell}{\mathrm e}
                      +5184C_\ell+1296A_\ell\right)\|q\|_1.
\]
This is \eqref{c18v37:bankbound}.
\end{proof}

The finite local inverses have not all been
numerically evaluated. Their positivity,
dimension bound and exact support preservation
are what make the displayed constants
computable and independent of volume.

\subsection{The volume-uniform native corrected remainder}

For a bank $F(g,y)$ let
$\mathsf C_\beta F=F-\int Fp_{g,\beta}dy$.
This centring does not change its $y$ derivatives.
The native identities are exactly
\[
 r_\beta=\mathsf C_\beta(2D_gu_\beta^g),
 \qquad
 \bar L_\beta=\bar L_0-
                         2\Gamma_0(u_\beta^g,\,\cdot\,).
 \tag{V37.19}\label{c18v37:exact}
\]
Define
\[
 D_\ell=
 \frac{18}{\mathrm e\,\beta_*^3}
   +\frac{3K_1}{2\beta_*^2}
   +\frac{3K_2}{2\beta_*}.
 \tag{V37.20}\label{c18v37:D}
\]

\begin{theorem}[Uniform order-six corrected inverse cost]
\label{c18v37:remainder}
For V36's original-metric lift $A^{[2]}$,
all $n\ge5$, $\ell\ge2$, all $g$, and real
$|\beta|\le\beta_*/2$,
\[
 \begin{aligned}
 Dc\,A^{[2]}&=I,\\
 \int r^{A^{[2]}}_\beta
          (r^{A^{[2]}}_\beta)^*p_{g,\beta}dy
                  &\preceq\frac{D_\ell^2}{\rho_*}\beta^6 I,\\
 R^{A^{[2]}}_\beta(g)
                  &\preceq\frac{D_\ell^2}{\rho_*^2}\beta^6 I .
 \end{aligned}
 \tag{V37.21}\label{c18v37:uniformremainder}
\]
The constants are independent of spatial
volume and coarse configuration; no column
of the coarse bank has been removed.
\end{theorem}
\begin{proof}
Write
$v_m=u_\beta^g-\sum_{k=1}^m\beta^ku_k^g$.
V36 gives
$\bar L_0\phi_1=r_1=2D_gu_1^g$ and
\[
 \bar L_0\phi_2=t_2
        =2D_gu_2^g+2\Gamma_0(u_1^g,\phi_1).
\]
Substitute these into the exact residual
$r_\beta-\beta\bar L_\beta\phi_1
-\beta^2\bar L_\beta\phi_2$.
Conditional divergences have zero mean,
so the result is exactly
\[
 r^{A^{[2]}}_\beta=\mathsf C_\beta
 \left[2D_gv_2+
        2\beta\Gamma_0(v_1,\phi_1)+
        2\beta^2\Gamma_0(u_\beta^g,\phi_2)\right].
 \tag{V37.22}\label{c18v37:residual}
\]
This identity retains the actual normalizer
and drift, rather than comparing extensive
densities with Haar.

For $t=|\beta|/\beta_*\le1/2$,
\eqref{c18v37:tails} gives
$Q_m(t)\le(3/4)t^{m+1}$.
The mixed chord/coarse derivative of the
first term in \eqref{c18v37:residual}
is a submatrix of twice the fine second
derivative matrix, bounded by
$(24/\mathrm e)Q_2(t)$.
Lemma~\ref{c18v37:bank} bounds the other
two derivative matrices by
$2|\beta|K_1Q_1(t)$ and
$2\beta^2K_2Q_0(t)$.
Thus their sum has operator norm at most
$D_\ell|\beta|^3$, uniformly in $g,y,N$.

For each fixed coarse bank $v$ apply the
first inequality in
\eqref{c18v37:Poincare} to the centred
scalar $v\cdot r^{A^{[2]}}_\beta$.
Its chord gradient has norm at most
$D_\ell|\beta|^3|v|$, proving the
covariance bound. The centred inverse
of the native reduced generator has
norm at most $1/\rho_*$, proving the
inverse-cost bound. V33 transfers that
cost to the full native invariant score.
Verticality of the two corrections
preserves $Dc\,A^{[2]}=I$, as in V36.
\end{proof}

This upgrades the \emph{scope} of V36's
order-six estimate on a stated coupling
interval. It does not change V36's exact
second coefficient or claim that its
adjoint part is the entire second source.

\subsection{Checks, advancement and the next genuine obstruction}

The diagnostic checks $5460$ multi-spin
denominator cases, exact generator and
Casimir matrices through spin three,
the sixteen-dimensional plaquette
coefficient trace norm, and three
connected-cover overlap fixtures.
It checks the contraction ball, ten
scalar-majorant coefficients, the
metric derivative constants, the
explicit cubic-remainder algebra and
the inherited V36 pair coefficient.
The overlap fixtures test combinatorial
majorants, not a many-link native
vacuum. The all-spin, all-volume,
analytic and conditional spectral
claims rest on the proofs above.
No proof-assistant certification is
asserted.

The previous missing small-coupling
uniform remainder is now paid in the
specific native score channel:
the vacuum logarithm converges in a
local norm, the exact reduced
conditional inverse is uniformly
bounded, and the explicit second
correction has uniform order-six cost.
The constants are intentionally
conservative; for example
$\beta_*\simeq0.00215554$ and
$\rho_*\simeq0.896362$.
Their positivity, not numerical
sharpness, is the present conclusion.

The next obstruction is no longer a
volume factor in this electric-limit
remainder. It is the mismatch between
this proved small-$\beta$ domain and
the required physical weak-coupling
continuum trajectory, together with
positive spatial-flow transport and
the remaining projection-energy and
coarse-variance estimates.
Another small-coupling coefficient
would not bridge those regimes.
The next upstream test must preserve
the native source and metric while
examining a scale or coupling
transport that can leave this disk.
No such transport is proved here.
The full interacting four-dimensional
continuum construction and positive
physical Yang--Mills mass gap remain
\emph{unproved}.
% END CYCLE 18 VERSION 37 APPEND
% BEGIN CYCLE 18 VERSION 38 APPEND
\clearpage
\section[V38: positive spatial flow and the curvature barrier]{V38: positive spatial flow and the curvature barrier}
\label{c18v38:context}

V37 pays a volume-uniform native conditional inverse at zero spatial
flow and small electric coupling. The next issue is transport, not
another zero-flow Taylor coefficient. We prove here that the actual
flowed conditional law and the original fine metric retain a uniform
reduced gap on an explicit positive interval. The associated
full-output score has a volume-uniform inverse cost, and an exact
Poisson correction has a bounded conditional Gram matrix.

The interval is small. More importantly, an explicit native
configuration shows that the same product-metric curvature criterion
is genuinely false at a required block-flow time, already at zero
coupling. This is an obstruction to this proof route, not a
counterexample to a conditional spectral gap or the Yang--Mills
mass gap. It directs the next step toward an integrated or
source-weighted estimate rather than repeated worst-case curvature
bounds.

The Jacobian, invariant source reduction, internal-gauge projection
and all-time metric comparison are inherited from V31--V34.
Archived f16.2/V74 already differentiates the exact transported density
once, and f16.2/V69 already identifies center-sensitive flow curvature
and an unstable homogeneous mode. We do not claim those facts again.
The additional estimates below are the locally summed second
variation, the resulting actual conditional inverse at positive flow,
and the exact pushed-density curvature obstruction. A targeted
search of the archived YM sources did not locate this last
combination; this is not a claim of literature novelty or a complete
equivalence audit.

\subsection{The exact law and the norms to be transported}

Retain $N=\ell n$, $\ell\ge2$, $n\ge5$, round $SU(2)$ links,
$T=-\sum_e\Delta_e$, $W=\sum_pw_p$, and the actual normalized
positive ground state $\Omega_\beta$ of $T-\beta W$.
Let $\Phi_s$ be the spatial flow with velocity
$F=\nabla W=-\nabla E_{\rm sp}$, with no physical-time factor.
Let $J_{-s}$ denote its inverse-flow Haar Jacobian, not its
differential matrix. The density in flowed links is exactly
\[
 \begin{aligned}
 h_{\beta,s}(V)&=\Omega_\beta(\Phi_{-s}V)^2J_{-s}(V),\\
 \log J_{-s}(V)&=12\int_0^s W(\Phi_{-t}V)\,dt,\\
 p_{g,\beta,s}(y)&=
 \frac{h_{\beta,s}(W^0(g,y))}
 {\int h_{\beta,s}(W^0(g,z))\,dz}.
 \end{aligned}
 \tag{V38.1}\label{c18v38:law}
\]
Here $W^0$ sets the rooted forest links $\mathcal F$ to $I$,
the last-chain links $\mathcal B$ to $g$, and the free chords
$\mathcal C$ to $y$. These are V33's exact Haar coordinates.
No induced-metric determinant is to be multiplied into this law.
The metric enters through V34's reduced kinetic coefficient
$a_{g,s}$:
\[
 e^{-32s}a_0\preceq a_{g,s}\preceq e^{32s}a_0,\qquad
 a_0\succeq I.
 \tag{V38.2}\label{c18v38:metric}
\]
All gradients and Hessians on the full link product below use its
round product metric. For a real scalar function $\psi$ write
\[
 G(\psi)=\sup_{V,e}|\nabla_e\psi(V)|,\qquad
 H(\psi)=\sup_V\|\operatorname{Hess}\psi(V)\|_{2\to2}.
 \tag{V38.3}\label{c18v38:norms}
\]
The first is a supremum of \emph{one-link} norms, not the norm of
the extensive gradient. With $u_\beta=\log\Omega_\beta$ up to
an additive constant, V37 gives, for real $|\beta|\le\beta_*$,
\[
 \begin{gathered}
 \beta_*=\frac{3}{512e},\qquad R_\beta=64e|\beta|\le\frac38,\\
 G(u_\beta)\le\frac{R_\beta}{3},\qquad
 H(u_\beta)\le\frac{4R_\beta}{e},\qquad
 G(W)\le4,\quad H(W)\le16 .
 \end{gathered}
 \tag{V38.4}\label{c18v38:input}
\]
The vector one-link bound follows directly from V37's
spin-generator norm for every unit direction, not by estimating
three components separately. The physical range is $\beta\ge0$;
the symmetric real interval is only useful for the analytic bound.

\subsection{A second variation that can contract an extensive local gradient}

For an array of linear link blocks $B_{ef}$, denote by
$\|B\|_{\rm S}$ the larger of the maximal block-norm row and
column sums. V31's local majorant implies, for either direction
of flow and for every propagator over an interval of length $t$,
\[
 \|\mathsf J(t,r)\|_{\rm S}\le e^{16|t-r|}.
 \tag{V38.5}\label{c18v38:SchurJ}
\]
The bound is proved in linkwise parallel frames; replacing these
by invariant frames multiplies each block on both sides by
orthogonal matrices and leaves the bound unchanged.

Use the invariant fields $X_{e,v}$ with tangent $vV_e$, and
the coefficient matrix $\mathsf J_t=D\Phi_{\pm t}$. In these frames
the forward variational generator is
\[
 (\mathsf M v)_e=\sum_f X_{f,v_f}F_e+[F_e,v_e].
\]
For backward time its sign changes. For three links $a,i,j$
let $\mathsf T_{aij}$ be the bilinear block mapping a derivative
direction at $a$ and an input at $j$ to the corresponding
coefficient derivative of $\mathsf M$ at output $i$.
Its norm is the supremum over the two unit input vectors.

\begin{lemma}[Three anchored tensor sums]
\label{c18v38:tensor}
The single-plaquette contribution satisfies
\[
 \|\mathsf T^p_{aij}\|\le
 3\,\mathbf1_{\{a,i,j\in E(p)\}}.
 \tag{V38.6}\label{c18v38:Tblock}
\]
Consequently each of the three anchored sums is at most $192$:
\[
 \sup_a\sum_{i,j}\|\mathsf T_{aij}\|\le192,\qquad
 \sup_i\sum_{a,j}\|\mathsf T_{aij}\|\le192,\qquad
 \sup_j\sum_{a,i}\|\mathsf T_{aij}\|\le192.
 \tag{V38.7}\label{c18v38:Tsum}
\]
\end{lemma}
\begin{proof}
A derivative of a plaquette trace in a unit imaginary quaternion
direction inserts that unit quaternion at its unique forward link
occurrence, or its negative at its inverse occurrence. An ordered
derivative of order at most three is still the scalar part of
one product of unit quaternions, including when links repeat.
It has absolute value at most one. This holds for arbitrary unit
directions, not just coordinate axes.

Pairing the differentiated force-gradient term with any unit
output gives such a third derivative, so its bilinear block norm
is at most one. In the same-link commutator term,
$|[z,v]|\le2|z||v|$, while the derivative of the single-plaquette
force has norm at most one by duality with its second directional
derivative. This contributes at most two. Both terms are supported
on the four links of that plaquette. This proves
\eqref{c18v38:Tblock}, with no color-counting multiplier.
An anchored link belongs to four plaquettes, with at most
$4^2$ ordered choices for the other two links. Thus every sum
is at most $4\cdot4^2\cdot3=192$.
\end{proof}

Let $\mathsf K_{aij,t}$ be the bilinear block of the ordinary
invariant coefficient derivative of $\mathsf J_t$.
In particular $\mathsf K_0=0$; this is not the covariant Hessian
of an arbitrarily framed matrix field.

\begin{lemma}[Locally summed second flow variation]
\label{c18v38:second}
For $t\ge0$, in either time direction,
\[
 \begin{gathered}
 \sup_a\sum_{i,j}\|\mathsf K_{aij,t}\|\le k(t),\qquad
 \sup_j\sum_{a,i}\|\mathsf K_{aij,t}\|\le k(t),\\
 k(t)=12e^{16t}(e^{16t}-1).
 \end{gathered}
 \tag{V38.8}\label{c18v38:Ksum}
\]
For every smooth real $\psi$,
\[
 \begin{aligned}
 G(\psi\circ\Phi_{\pm t})&\le e^{16t}G(\psi),\\
 H(\psi\circ\Phi_{\pm t})&\le
 e^{32t}H(\psi)+k(t)G(\psi).
 \end{aligned}
 \tag{V38.9}\label{c18v38:composition}
\]
\end{lemma}
\begin{proof}
Differentiating the coefficient variational equation, including
the frame commutator already in $\mathsf M$, gives V29's
variation-of-constants formula. In block notation its norm
majorant is
\[
 \|\mathsf K_{aij,t}\|
 \le\int_0^t\sum_{k,l,m}
 \|\mathsf J(t,r)_{ik}\|\,
 \|\mathsf T_{lkm}(\Phi_{\pm r}V)\|\,
 \|\mathsf J_{r,la}\|\,\|\mathsf J_{r,mj}\|\,dr .
 \tag{V38.10}\label{c18v38:Duhamel}
\]
For fixed $a$, sum over $i,j$. The propagator column sum
and the last differential row sum cost
$e^{16(t-r)}e^{16r}$. The tensor sum with $l$ fixed costs
$192$, and the remaining differential column sum costs
$e^{16r}$. This gives
$192\int_0^t e^{16(t-r)}e^{32r}dr=k(t)$.
With $j$ fixed, use the tensor sum with $m$ fixed and the
same row/column bounds. This proves both inequalities.

The first composition estimate uses a column sum of $\mathsf J_t$.
For the second, take a constant invariant vector $v=(v_e)$
and the initial product geodesic $\exp(\tau v)V$.
At the image point its velocity is $\mathsf J_tv$.
Differentiating the scalar function twice along this curve gives
the Hessian contribution
$\operatorname{Hess}\psi(\mathsf J_tv,\mathsf J_tv)$ and
\[
 \sum_{a,i,j}
 \big\langle\nabla_i\psi,\mathsf K_{aij,t}(v_a,v_j)\big\rangle.
\]
No connection term survives the diagonal evaluation: the
Levi--Civita connection of a constant invariant vector with itself
is zero on each round group factor. Equivalently the ordered
second scalar derivative evaluated on the same vector is the
Riemannian Hessian. Bound the displayed term by
$G(\psi)\sum_{a,j}B_{aj}|v_a||v_j|$, where
$B_{aj}=\sum_i\|\mathsf K_{aij,t}\|$.
The two sums in \eqref{c18v38:Ksum} and the Schur bound yield
$G(\psi)k(t)\|v\|^2$. Finally
$\|\mathsf J_t\|_{2\to2}\le e^{16t}$.
\end{proof}

This estimate has a different use from V29's packed
$\ell^2$ derivative estimate. Contracting the latter directly
with the full gradient of an extensive $\psi$ can introduce
a volume factor. The anchored sums above pay that contraction
with $G(\psi)$ instead.

\subsection{A paid positive-flow conditional gap}

Put $E_s=e^{16s}$ and define
\[
 \begin{aligned}
 \mathcal H_\beta(s)
  &=\frac{8R_\beta}{e}E_s^2
    +8R_\beta E_s(E_s-1)
    +6(E_s^2-1)+18(E_s-1)^2,\\
 \rho_\beta(s)&=2-\mathcal H_\beta(s).
 \end{aligned}
 \tag{V38.11}\label{c18v38:H}
\]
All constants here are independent of $n,\ell,g$.

\begin{theorem}[Actual flowed law, original reduced metric]
\label{c18v38:gap}
For real $|\beta|\le\beta_*$ and $s\ge0$ with
$\mathcal H_\beta(s)<2$, the exact reduced conditional generator
obeys
\[
 \operatorname{gap}\bar L_{g,\beta,s}
 \ge e^{-32s}\rho_\beta(s)
 \quad\hbox{for every }g,n\ge5,\ell\ge2.
 \tag{V38.12}\label{c18v38:gapbound}
\]
In particular, uniformly for
\[
 |\beta|\le\frac{3}{512e},\qquad
 0\le s\le\frac1{1024},
\]
one has
\[
 \mathcal H_\beta(s)\le\frac{620}{441},\qquad
 \operatorname{gap}\bar L_{g,\beta,s}
 \ge\frac{4061}{7056}>\frac12 .
 \tag{V38.13}\label{c18v38:rectangle}
\]
This is the invariant reduced generator, with no physical
factor $\kappa$ inserted.
\end{theorem}
\begin{proof}
Apply \eqref{c18v38:composition} to $2u_\beta\circ\Phi_{-s}$
and to the exact integral in \eqref{c18v38:law}. It gives
\[
 \begin{aligned}
 H(\log h_{\beta,s})
 &\le\frac{8R_\beta}{e}E_s^2
      +\frac{2R_\beta}{3}k(s)
      +12\int_0^s[16e^{32t}+4k(t)]\,dt\\
 &=\mathcal H_\beta(s).
 \end{aligned}
 \tag{V38.14}\label{c18v38:logH}
\]
In particular the last two terms of \eqref{c18v38:H} are a
genuine Jacobian cost, not a magnetic coupling replacement.
Freezing distinct fine links to $I$ and $g$ is a totally
geodesic product section. The normalizing logarithm in
\eqref{c18v38:law} is constant in $y$. Therefore the
$y$-Hessian of $\log p_{g,\beta,s}$ has norm at most
$\mathcal H_\beta(s)$.

The product round $S^3$ Ricci tensor is $2I$.
For the product-metric diffusion
$\mathcal A=\Delta_y+\nabla_y\log p\cdot\nabla_y$, the integrated
Bochner identity yields
\[
 \int(\mathcal A f)^2p\,dy
 =\int\bigl[
   |\operatorname{Hess}_y f|^2+
   (2I-\operatorname{Hess}_y\log p)(\nabla_y f,\nabla_y f)
  \bigr]p\,dy .
 \tag{V38.15}\label{c18v38:Bochner}
\]
Apply it to each nonconstant eigenfunction of the compact
elliptic operator $-\mathcal A$. If its eigenvalue is
$\lambda>0$, then $\lambda^2\|f\|^2\ge
\rho_\beta(s)\lambda\|f\|^2$. Spectral expansion gives
$\operatorname{Var}_p(f)\le\rho_\beta(s)^{-1}
\int|\nabla_y f|^2p$. Comparison with
\eqref{c18v38:metric} proves \eqref{c18v38:gapbound}.
This restates the needed Bakry--Emery step explicitly; it
does not assert that $a_{g,s}$ is a product metric.

For the rectangle, use
$E_s\le e^{1/64}\le64/63$, $R_\beta\le3/8$, and
$e^{-1}<3/8$. Every term in \eqref{c18v38:H} is increasing in
$R_\beta$ and $E_s\ge1$. Substitution gives
$\mathcal H_\beta(s)\le620/441$ and
$\rho_\beta(s)\ge262/441$.
Also $e^{-32s}\ge1-1/32=31/32$, proving the displayed floor.
\end{proof}

For primary context the Bochner criterion is the one used in
D.~Bakry and M.~\'Emery, \emph{Diffusions hypercontractives},
\href{https://numdam.org/item/SPS_1985__19__177_0/}
{S\'eminaire de probabilit\'es 19 (1985), 177--206};
V37 records the relevant original passages. Its bibliographic
page was rechecked here. M.~L\"uscher's
\href{https://arxiv.org/abs/1006.4518}
{\emph{Properties and uses of the Wilson flow in lattice QCD}}
provides broader flow context. Only that paper's abstract was
rechecked in this iteration; no continuum or mixing conclusion
from it is imported. The broad literature checkpoint remains V40.

\subsection{Every coarse column, its inverse cost and an exact correction}

Let $A_s^{\rm tr}$ differentiate the exact chart
$\chi_s(k,g,y)=k^{-1}\cdot\Phi_{-s}W^0(g,y)$ in $g$
at fixed $k,y$. Let $A_s^{\rm h}=P A_s^{\rm tr}$, with
V34's \emph{internal-gauge} projection $P$.
These are full coarse-output lifts, not one selected observable.
Their common centred score is
\[
 r_{g,\beta,s}=
 \mathcal C_{g,\beta,s}
      D_g\log h_{\beta,s}(W^0(g,y)),\qquad
 \mathcal C_{g,\beta,s}f=f-\int f p_{g,\beta,s}\,dy .
 \tag{V38.16}\label{c18v38:score}
\]
The formula follows from the exact full chart density
$h_{\beta,s}(W^0)\,dk\,dg\,dy$ and the zero Haar divergence
of each coarse invariant frame. V34 preserves that divergence
under $P$. Moreover
\[
 Dc\,A_s^{\rm h}=I,\qquad
 (A_s^{\rm h})^*A_s^{\rm h}
 \preceq(A_s^{\rm tr})^*A_s^{\rm tr}\preceq E_s^2 I .
 \tag{V38.17}\label{c18v38:lift}
\]

\begin{theorem}[Uniform source inverse and score-free full lift]
\label{c18v38:source}
In the domain of Theorem \ref{c18v38:gap}, write
$H=\mathcal H_\beta(s)$ and $\rho=2-H$.
The full coarse covariance and native invariant inverse-cost matrices
satisfy
\[
 \Sigma(g)\preceq\frac{H^2}{\rho}I,\qquad
 R(g)\preceq E_s^2\frac{H^2}{\rho^2}I .
 \tag{V38.18}\label{c18v38:sourcebound}
\]
Let $\varphi_i=\bar L_{g,\beta,s}^{-1}r_i$ be centred and pulled
back to the original fibre, and set, column by column,
\[
 A_s^\sharp=A_s^{\rm h}+\nabla_{\rm fib}\varphi .
 \tag{V38.19}\label{c18v38:sharp}
\]
Then
\[
 \begin{aligned}
 Dc\,A_s^\sharp&=I,& r^{A_s^\sharp}&=0,\\
 \mathbb E_{\nu_g}\big[(A_s^\sharp)^*A_s^\sharp\big]
 &\preceq\frac{4E_s^2}{\rho^2}I.
 \end{aligned}
 \tag{V38.20}\label{c18v38:sharpbound}
\]
On the rectangle \eqref{c18v38:rectangle}, the last upper bound
is less than $12I$.
\end{theorem}
\begin{proof}
For a fixed full coarse bank $v$, the $y$-gradient of $v\cdot r$
is a mixed $\mathcal C,\mathcal B$ Hessian block of $\log h$.
The two link sets are disjoint; there is no same-link connection
term. Its norm is at most $H|v|$ by
\eqref{c18v38:logH}. Conditional centering subtracts a function
constant in $y$. The product Poincare estimate in the preceding
proof gives the covariance bound. The native reduced gap is
at least $E_s^{-2}\rho$, so its centred inverse is at most
$E_s^2/\rho$. V33 identifies this inverse cost exactly with
that of the internally invariant source in the original fibre.
This proves the second bound, without replacing the conditional
law or dropping any output column.

For each finite regulator these elliptic centred Poisson equations
have unique smooth solutions. Smooth dependence on $g$ follows,
for example, by inverting $\bar L_g+\Pi_g$ on a fixed compact
function space, where $\Pi_g f=\int f p_g$, and using the gap
and elliptic regularity locally in $g$. This is existence and
smoothness, not a uniform bound on differentiated solution banks.
Gauge transformations intertwine the data and generator;
uniqueness gives the equivariant column transformation.

The correction is vertical and its full weighted divergence is
its conditional divergence, namely $-\bar L\varphi=-r$.
Its mean divergence is zero. This proves both identities in
\eqref{c18v38:sharpbound}. Its conditional squared norm on $v$
equals $v^*Rv$ by integration by parts. The triangle inequality
in the conditional $L^2$ space of tangent fields gives
\[
 \bigl(\mathbb E_{\nu_g}|A_s^\sharp v|^2\bigr)^{1/2}
 \le E_s|v|+\frac{E_sH}{\rho}|v|
 =\frac{2E_s}{\rho}|v|.
\]
We do \emph{not} set the cross term to zero:
$A_s^{\rm h}$ is horizontal only to internal gauge orbits,
not to the whole fibre. Finally $E_s^2\le32/31$ and
$\rho\ge262/441$ bound its squared coefficient by
$6223392/531991<12$.
\end{proof}

The Poisson correction principle itself is inherited, not new.
What is now paid is its volume-uniform cost for this actual
positive-flow law. The correction need not be finite range.
For a smooth variable bank $v(g)$ its additional divergence term
is a deterministic coarse derivative because $Dc\,A^\sharp=I$;
it disappears under conditional centering. This proves score
cancellation for such banks but does not bound their
projection energies or derivatives.

There is a useful normalization check at the origin. At fixed
$N,\ell$, in every fixed smooth norm, put
$\epsilon=|\beta|+s$. The exact law gives
\[
 \log h_{\beta,s}
   =\left(\frac{\beta}{6}+12s\right)W+
       \text{a constant}+O_{N,\ell}(\epsilon^2),\qquad
 r_{g,\beta,s}=(\beta+72s)r_1+O_{N,\ell}(\epsilon^2).
 \tag{V38.21}\label{c18v38:tangent}
\]
Indeed $u_\beta=\beta W/12+O(\beta^2)$,
$\Phi_{-s}=\mathrm{id}+O(s)$ and
$\log J_{-s}=12sW+O(s^2)$. Each plaquette has a free chord,
so its Haar conditional mean vanishes. Conditional centering
therefore changes this first term only at second order.
With V35's $\lambda_\ell=12-6/\ell$ and covariance $\Sigma_1$,
analyticity of the fixed-regulator inverse further gives
\[
 R(g)=\frac{(\beta+72s)^2}{\lambda_\ell}\Sigma_1(g)
       +O_{N,\ell}(\epsilon^3)
 \tag{V38.22}\label{c18v38:tangentcost}
\]
in coarse operator norm at fixed regulator. These Taylor
remainders are not asserted to be uniform in volume.
In particular a correction that only cancels powers of $\beta$
at $s=0$ cannot simply be declared to have the same order at
$s>0$. The number $\beta+72s$ is a first-order density tangent,
not an exact effective Wilson coupling or a renormalization map.

\subsection{A native failure of uniform product curvature at block time}

The exponential majorant in \eqref{c18v38:H} is not by itself
evidence of failure of a sharper curvature argument. We therefore
test the \emph{actual} density on an admissible configuration.

\begin{proposition}[Negative product curvature at $s=\ell^2$]
\label{c18v38:obstruction}
For $\beta=0$, $\ell=2$, $n=5$, hence $N=10$, there are
coarse coordinates $g_*$ and chord coordinates $y_*$
such that at $s=4$ every unit single-chord color direction $v$
has
\[
 \operatorname{Hess}_y\log p_{g_*,0,4}(y_*)[v,v]
 >\frac{2997}{1000}>2 .
 \tag{V38.23}\label{c18v38:negative}
\]
Thus $2I-\operatorname{Hess}_y\log p_{g_*,0,4}$ is not
positive semidefinite. A product-metric Bakry--Emery lower
bound valid at all configurations cannot extend the small-flow
argument to all required block times.
\end{proposition}
\begin{proof}
On any even fine torus define center-valued stored links by
\[
 U^*_{x,1}=I,\qquad U^*_{x,2}=(-1)^{x_1}I,\qquad
 U^*_{x,3}=(-1)^{x_1+x_2}I .
 \tag{V38.24}\label{c18v38:center}
\]
The signs are periodic and each of the three plaquette types
has holonomy $-I$. All first trace derivatives vanish.
Writing $d_1$ for the ordinary real periodic lattice curl,
the Hessian of $W$ at this configuration is
\[
 A=d_1^*d_1\otimes I_3 .
 \tag{V38.25}\label{c18v38:curl}
\]
To see the normalization, vary the four plaquette links along
$\exp(\tau v_e)U^*_e$. Their center factors commute.
The second derivative of the scalar trace is
$|v_{x,i}+v_{x+\hat i,j}-v_{x+\hat j,i}-v_{x,j}|^2$.
Commutators contribute zero scalar part at this order.
Sum over all plaquettes. This gives \eqref{c18v38:curl}
with no omitted factor two.

The configuration is a stationary point of the spatial flow.
At that point the backward differential is $e^{-tA}$.
At $\beta=0$ the original vacuum is exactly constant, so
only the inverse Jacobian remains. Since $\nabla W=0$ there,
the second-flow-variation term in the composition Hessian vanishes.
Consequently the \emph{exact} Hessian is
\[
 \operatorname{Hess}\log h_{0,s}(U^*)
   =12\int_0^s e^{-tA}Ae^{-tA}\,dt
   =6(I-e^{-2sA}).
 \tag{V38.26}\label{c18v38:exactcenterH}
\]

For one color the Fourier curl Gram has two eigenvalues
\[
 \lambda(k)=4\sum_{j=1}^3\sin^2(\pi k_j/N)
\]
and one zero eigenvalue at each nonzero momentum; at zero
momentum all three vanish. Indeed its symbol is
$\lambda(k)I-d(k)d(k)^*$ with
$d_j(k)=e^{2\pi i k_j/N}-1$.
Let $Q$ project onto the nonzero curl subspace.
Translation and cubic symmetry make its diagonal constant;
its rank is $2(N^3-1)$ per color. Thus any unit coordinate
vector has
\[
 \langle v,Qv\rangle=\frac{2(N^3-1)}{3N^3}.
 \tag{V38.27}\label{c18v38:Qdiag}
\]
The least positive eigenvalue is
$\lambda_{\min}=4\sin^2(\pi/N)$. Spectral calculus gives
$I-e^{-2sA}\succeq(1-e^{-2s\lambda_{\min}})Q$.
At $N=10$, $\lambda_{\min}=(3-\sqrt5)/2>1/3$.
At $s=4$, $e^{-2s\lambda_{\min}}<1/4$, for instance
by the first three terms of the positive exponential series.
Equations \eqref{c18v38:exactcenterH}--\eqref{c18v38:Qdiag}
therefore give the lower bound
\[
 6\cdot\frac34\cdot\frac{2(10^3-1)}{3\cdot10^3}
   =\frac{2997}{1000}.
\]

It remains to put this point in the declared conditional
section, rather than silently changing the fibre. Starting at
each coarse root with gauge value $I$, choose center gauge
values successively along the forest to make its links $I$.
All remaining links are still central. Denote their
$\mathcal B,\mathcal C$ values by $g_*,y_*$.
This is an internal-gauge isometry, its tangent adjoint action
is the identity, and $\log h_{0,s}$ is gauge invariant.
Thus the computed Hessian transfers to $W^0(g_*,y_*)$.
There are $2(N^3-n^3)>0$ free chords. Varying any single
one with $g_*$ fixed is a geodesic in the product section.
The conditional normalizer is constant on that section, so
the same lower bound applies to its conditional Hessian.
This proves \eqref{c18v38:negative}.
\end{proof}

This is a density-curvature statement in the comparison product
metric. It does not calculate the Bakry--Emery tensor of the
nonproduct native metric, nor prove a small eigenvalue.
The density is positive at the displayed point; continuity
also gives a neighborhood of negative product curvature.
No estimate of that neighborhood's native probability or
source spectral weight has been made here. In particular
an exponentially rare bad region could be compatible with
a useful integrated inequality.

\subsection{Verification and the resulting change of direction}

The diagnostic
\texttt{scripts/verify\_ym\_c18\_v38\_flowed\_conditional\_gap.py}
checks the exact scalar integrals and rational constants,
8,000 unit-direction quaternion third-derivative words
(including repeated links, reversed links and non-axis directions),
and V29's ordered local tensor with its frame sign.
It enumerates all 3,000 plaquettes of the native $N=10$ center
configuration, performs the prescribed forest gauge fixing,
and checks its exact $12$-by-$12$ single-plaquette Hessian.
The 1,000 momentum symbols give an additional floating-point
regression of the curl identity and the density-Hessian diagonal.
The latter is approximately $3.99475965$; the proof uses instead
the exact lower bound $2997/1000$.
V37's scalar contraction constants are revalidated.
These tests do not compute a many-link interacting vacuum,
integrate a nonlinear flow trajectory, or certify a proof assistant.

V38 replaces V37, which is exactly recoverable by removing this
appendix and restoring its title. The other 52 sources are unchanged.
Only \texttt{.tex} files are in \texttt{latex\_notes};
builds and diagnostics remain outside. Companion and literature
checkpoints remain V40, the last opportunistic companion review
was V36, and the archive/restart remains V100.

The volume-uniform positive-flow conditional inverse and full-output
correction are paid only on the stated small-coupling rectangle.
The bound $s\le1/1024$ does not reach block time $s=\ell^2\ge4$
or transport the weak-coupling continuum trajectory.
The native example excludes the proposed everywhere-positive
product-curvature route at that time.

Next, retain the signed Jacobian Hessian and control its bad-region
contribution to an integrated conditional or source-weighted inverse,
or prove another metric-faithful scale transport.
Uniform solution-bank derivatives, projection energy, coarse variance,
physical return and scale calibration remain unpaid.
The interacting four-dimensional continuum construction and positive
physical Yang--Mills mass gap remain \emph{unproved}.
% END CYCLE 18 VERSION 38 APPEND
% BEGIN CYCLE 18 VERSION 39 APPEND
\clearpage
\section[V39: native coarea geometry and the remaining curvature cost]{V39: native coarea geometry and the remaining curvature cost}
\label{c18v39:context}

V38 proves a positive-flow conditional estimate in a small rectangle,
but its counterexample concerns the \emph{product comparison metric}.
Before estimating the probability of that comparison metric's bad
curvature, we go upstream to the actual fibre geometry.
The spatial-flow Jacobian is exactly a Riemannian volume factor
when the transported metric is retained. After conditioning and
internal-gauge reduction, its contribution separates into the
normal coarea Jacobian, orbit volume and quotient metric volume.
The score still contains their full coordinate derivative.

We prove that separation for the declared native model, derive its
weighted curvature with all extrinsic terms, and compute the exact
quotient metric on V38's stationary center configuration. There
the native kinetic coefficient increases with flow time, despite
the failure of the product curvature test. This does not prove
positivity of native curvature: it proves why the earlier negative
tensor is not the native tensor to insert into a new estimate.
The remaining integrated estimate is formulated for that correct
tensor and the same source.

Archived f17/V28 already warns that orbit volume is part of a
quotient probability; Shell Mixing V16's inherited V4 already
combines a constrained determinant with coarea. Current V33--V34
provide the exact rooted coordinates and quotient kinetic matrix.
Those principles are not new. The new application combines
both successive reductions for the full flowed sampled fibre,
retains their score derivative, and evaluates the stationary
native metric. The review was targeted, not an exhaustive
novelty certificate.

\subsection{The two reductions and their exact determinants}

Let $\mathcal M=SU(2)^{E_{\rm f}}$ have its original product metric
$\mathfrak g$, and retain $N=\ell n$, $\ell\ge2$, $n\ge5$,
$c=\mathfrak b_\ell\circ\Phi_s$, $s\ge0$, and the actual
law $\nu=\Omega_\beta^2dU$. In this appendix the geometric
identities require no smallness of $\beta$ or $s$.
For the physical Hamiltonian take $\beta\ge0$.
The chain differential has full row rank and spatial flow is
a diffeomorphism, so $c$ is a smooth submersion everywhere.
Write
\[
 C_U=Dc(U),\qquad G_U=C_UC_U^*,\qquad
 \mathcal J_c(U)=\sqrt{\det G_U}.
 \tag{V39.1}\label{c18v39:normal}
\]
Here $\mathcal J_c$ is the normal Jacobian of the sampled map.
It is not the inverse spatial-flow Jacobian $J_{-s}$.

The compact connected fibre $\mathcal F_g=c^{-1}(g)$ has the
original induced metric. The root-fixed group $\mathcal H_0$
acts freely and isometrically on it. Its quotient
$\mathcal Q_g=\mathcal F_g/\mathcal H_0$ is the smooth compact
manifold with V33's global chord coordinates $y$.
At $U^0=\Phi_{-s}W^0(g,y)$ let
\[
 K=D_{U^0}^*D_{U^0},\qquad
 a=a_{g,s},\qquad q=a^{-1}.
 \tag{V39.2}\label{c18v39:quotient}
\]
$q$ is the covariant quotient metric, expressed in product
Haar orthonormal chord frames; $a$ is its contravariant
kinetic coefficient. The determinants of $K,G,a$ are taken
in their respective original orthonormal frames, including
all color components. The pinned gauge matrix $K$ is strictly
positive by V34.

\begin{theorem}[Exact geometric density, without a second Jacobian]
\label{c18v39:density}
On the rooted section, up to the fixed normalization of
Riemannian volumes versus normalized Haar measures,
\[
 J_{-s}(W^0)^2=\frac{\det K}{\det G\,\det a}.
 \tag{V39.3}\label{c18v39:det}
\]
The actual reduced conditional probability can be written either
as $p_g(y)\,dy$, with V38's exact $p_g$, or as
$\omega_g(y)\,d\operatorname{vol}_{q}(y)$, where
\[
 \begin{aligned}
 \omega_g(y)&=Z_g^{-1}
   \Omega_\beta(U^0)^2
       \frac{\sqrt{\det K(U^0)}}{\sqrt{\det G(U^0)}},\\
 d\operatorname{vol}_{q}(y)&=
                  (\det a(y))^{-1/2}\,dy .
 \end{aligned}
 \tag{V39.4}\label{c18v39:weight}
\]
The harmless fixed volume factors are absorbed into $Z_g$;
in the normalized invariant-frame determinant identity
\eqref{c18v39:det} no such factor is needed.
\end{theorem}
\begin{proof}
Use the full chart $\chi_s(g,k,y)=k^{-1}\cdot\Phi_{-s}W^0(g,y)$
and let $\mathsf M=(D\chi_s)^*D\chi_s$ be its full metric
matrix in product Haar orthonormal frames. The unflowed
tree change of variables has Haar Jacobian one. Thus
$\det\mathsf M=J_{-s}^2$.
Partition it into the $g$ coordinates and $z=(k,y)$.
Since $c\circ\chi_s=g$, the $gg$ principal block of
$\mathsf M^{-1}$ is $G$.
The Schur determinant identity therefore gives
\[
 \det\mathsf M=\frac{\det\mathsf M_{zz}}{\det G}.
 \tag{V39.5}\label{c18v39:Schur}
\]
The $kk$ block of the fibre metric $\mathsf M_{zz}$ is $K$.
Its Schur complement is precisely the quotient metric $q$,
by minimizing over pinned gauge tangents as in V34. Hence
$\det\mathsf M_{zz}=\det K\,\det q=\det K/\det a$.
This proves \eqref{c18v39:det}.

Alternatively, ordinary coarea first gives the native conditional
law, relative to the induced fibre volume, as
\[
 d\nu_g=Z_g^{-1}\frac{\Omega_\beta^2}{\mathcal J_c}
                       \,d\operatorname{vol}_{\mathcal F_g}.
 \tag{V39.6}\label{c18v39:fibrelaw}
\]
Both factors in its density are invariant under $\mathcal H_0$.
The metric volume of that orbit is a fixed Haar normalization
times $\sqrt{\det K}$; the action is free and isometric.
Integrate over the orbit to obtain \eqref{c18v39:weight}.
Multiplying its two lines recovers
$Z_g^{-1}\Omega_\beta(U^0)^2J_{-s}(W^0)\,dy$
by \eqref{c18v39:det}, namely V38's law.
Thus the two descriptions agree, rather than prescribing
an additional measure factor.
\end{proof}

In particular the scalar probability potential for the native
quotient metric is
\[
 V_g^{\rm geom}
 =-2\log\Omega_\beta(U^0)-\tfrac12\log\det K(U^0)
                         +\tfrac12\log\det G(U^0)
                         +\text{a fibre constant}.
 \tag{V39.7}\label{c18v39:potential}
\]
It is not $-\log p_g$ paired with $q$.
The latter is the potential only when the reference volume
and metric are the product Haar ones.

\subsection{What cancels geometrically, and what stays in the source}

There is a useful ambient check, before either reduction.
In flowed link variables give $\mathcal M$ the metric
$\widetilde{\mathfrak g}_s=(\Phi_{-s})^*\mathfrak g$.
Then
\[
 d\operatorname{vol}_{\widetilde{\mathfrak g}_s}
       =J_{-s}\,dV,\qquad
 h_{\beta,s}\,dV
       =(\Omega_\beta^2\circ\Phi_{-s})
                       d\operatorname{vol}_{\widetilde{\mathfrak g}_s}.
 \tag{V39.8}\label{c18v39:ambient}
\]
The map $\Phi_{-s}$ is an isometry for these two declared
metrics. Naturality of Ricci curvature and the covariant Hessian
therefore gives
\[
 \operatorname{Ric}_{\widetilde{\mathfrak g}_s}
 -2\operatorname{Hess}_{\widetilde{\mathfrak g}_s}
                    (\log\Omega_\beta\circ\Phi_{-s})
 =(\Phi_{-s})^*
          (2\mathfrak g-2\operatorname{Hess}_{\mathfrak g}
                                      \log\Omega_\beta).
 \tag{V39.9}\label{c18v39:ambientRic}
\]
At $\beta=0$ this is exactly $2\widetilde{\mathfrak g}_s$
for every $s$. At small $\beta$ the analogous ambient bound
is already a consequence of V37; no new physical gap is claimed.
Conditioning on $c$ is the additional step that this ambient
isometry does \emph{not} perform.

Nor does \eqref{c18v39:potential} allow deletion of the flow
Jacobian from the coarse score.
Let $u=\log\Omega_\beta$ and let $D_g$ be the full fixed
coarse frame derivative in V38's tree lift. Its exact residual is
\[
 r_g=\mathcal C_gD_g
 \left[
  2u(U^0)+\tfrac12\log\det K(U^0)
           -\tfrac12\log\det G(U^0)
           -\tfrac12\log\det a_{g,s}(y)
 \right].
 \tag{V39.10}\label{c18v39:source}
\]
Indeed the bracket is $\log h_{\beta,s}(W^0)$ by
\eqref{c18v39:det}; $D_g\log Z_g$ is removed by centering.
When a conditional metric depends on $g$, its volume derivative
is part of the density derivative relative to the fixed
Haar chart. The last term is therefore essential.
All derivatives include the dependence of $U^0$ on $g$.
In particular V38's first-order coefficient $\beta+72s$
is unchanged. No score or coarse-output component is removed.

\subsection{The native quotient curvature and its nonnegative gauge gain}

Let $\nabla^q,\operatorname{Ric}_q$ refer to the quotient
metric. Its actual weighted Ricci tensor is
\[
 \mathcal R_Q=
 \operatorname{Ric}_q-2\operatorname{Hess}_q(u|_{\mathcal Q_g})
 -\tfrac12\operatorname{Hess}_q\log\det K
 +\tfrac12\operatorname{Hess}_q\log\det G .
 \tag{V39.11}\label{c18v39:RicQ}
\]
Here restrictions are interpreted through the rooted section,
not by differentiating an ambient scalar while ignoring
the section's acceleration.

Write $\pi:\mathcal F_g\to\mathcal Q_g$ for the Riemannian
quotient map. On $\mathcal F_g$ put
$\alpha=\Omega_\beta^2/\mathcal J_c$ and
$\mathcal R_F=\operatorname{Ric}_{\mathcal F_g}
                    -\operatorname{Hess}_{\mathcal F_g}\log\alpha$.
Let $X$ be the horizontal lift of a tangent $v\in T\mathcal Q_g$.
At a point choose orthonormal horizontal vectors $E_i$ and
orbit-tangent vectors $U_a$ in $T\mathcal F_g$.
For horizontal lifts of local base fields define
\[
 \begin{aligned}
 \mathcal A(E_i,X)&=(\nabla^{F}_{E_i}X)_{\rm orb},\\
 \mathcal T(U_a,U_b)&=(\nabla^{F}_{U_a}U_b)_{\rm hor},\\
 \mathcal S_{\rm orb}(X)
 &=2\sum_{i,a}|\langle\mathcal A(E_i,X),U_a\rangle|^2
       +\sum_{a,b}|\langle X,\mathcal T(U_a,U_b)\rangle|^2 .
 \end{aligned}
 \tag{V39.12}\label{c18v39:gaugegain}
\]
These expressions are independent of the chosen orthonormal
bases, and $\mathcal S_{\rm orb}\ge0$.

\begin{theorem}[Exact weighted curvature under internal-gauge reduction]
\label{c18v39:submersion}
With precisely the above measure and metric,
\[
 \mathcal R_Q(v,v)=\mathcal R_F(X,X)+\mathcal S_{\rm orb}(X).
 \tag{V39.13}\label{c18v39:Ricgain}
\]
\end{theorem}
\begin{proof}
The weighted fibre diffusion
$\mathcal A_F=\Delta_F+\nabla^F\log\alpha\cdot\nabla^F$
commutes with the isometric group action, since $\alpha$ is
invariant. It therefore sends basic functions to basic functions.
For smooth $f,h$ on $\mathcal Q_g$, disintegrating its Dirichlet
form gives
\[
 -\int(h\circ\pi)\mathcal A_F(f\circ\pi)\,d\nu_g
 =\int\langle\nabla^q h,\nabla^q f\rangle\,p_gdy .
\]
This identifies $\mathcal A_F(f\circ\pi)$ with
$(\mathcal A_Qf)\circ\pi$, where
$\mathcal A_Q=\Delta_q+\nabla^q\log\omega_g\cdot\nabla^q
             =-\bar L_g$.
The pointwise squared gradient is also preserved. Applying
these two intertwining identities in the definition of
$\Gamma_2$ proves pointwise intertwining of $\Gamma_2$.

For $F=f\circ\pi$ the gradient is horizontal. Its horizontal
Hessian block equals $\operatorname{Hess}_q f$ lifted by $\pi$.
Its mixed block is
$\operatorname{Hess}_F F(E_i,U_a)
 =\langle\mathcal A(E_i,\nabla^F F),U_a\rangle$.
Its orbit-orbit block is
$\operatorname{Hess}_F F(U_a,U_b)
 =-\langle\nabla^F F,\mathcal T(U_a,U_b)\rangle$.
The last equality follows by differentiating $U_bF=0$;
the mixed equality is the gradient definition of the Hessian.
The horizontal block follows from the horizontal part of the
Levi--Civita connection of a Riemannian submersion, which is
also obtained directly from the Koszul formula on horizontal
lifts.

Squaring the blocks gives the quotient Hessian square plus
\eqref{c18v39:gaugegain}; the mixed block occurs twice.
Apply the weighted Bochner identity on both manifolds and
cancel the common quotient Hessian square. At a point choose
$f$ with prescribed gradient $v$. This proves
\eqref{c18v39:Ricgain}.
\end{proof}

This is the free-isometric-action case of weighted submersion
geometry, not a new general theorem. For primary context,
J.~Lott, \emph{Some geometric properties of the Bakry--Emery--Ricci
tensor}, \href{https://arxiv.org/abs/math/0211065}{arXiv:math/0211065},
Theorem 2, equations (3.1)--(3.11), and the group-action example
at the end of Section 3 were inspected. His measure-transport
hypothesis matters; here the group action and the displayed
Dirichlet argument verify the needed intertwining.
We do not apply that theorem directly to
$c:\mathcal M\to SU(2)^{E_{\rm c}}$ with its product coarse metric:
that map is generally not a Riemannian submersion for that metric.

\subsection{The remaining cost is extrinsic to the sampled fibre}

We can make $\mathcal R_F$ in \eqref{c18v39:Ricgain} entirely
explicit in ambient geometric quantities. Let $\mathrm{II}$
be the second fundamental form of
$\mathcal F_g\subset(\mathcal M,\mathfrak g)$, with convention
$\mathrm{II}(v,w)=(\nabla_v^{\mathcal M}w)_{\perp}$.
Its mean curvature is the unnormalized trace
$\mathsf H=\sum_j\mathrm{II}(e_j,e_j)$ over a fibre
orthonormal basis. Let $n_a$ be an orthonormal normal basis.
Set $\lambda=\log\mathcal J_c=\tfrac12\log\det G$.
For any $X\in T\mathcal F_g$ define the signed quadratic form
\[
 \begin{aligned}
 \mathcal D_c(X)={}&
 \sum_a\langle R^{\mathcal M}(X,n_a)n_a,X\rangle
 +\sum_j|\mathrm{II}(X,e_j)|^2\\
 &-\operatorname{Hess}_{\mathcal M}\lambda(X,X)
 -\langle\mathsf H-2\nabla^{\mathcal M}u
                  +\nabla^{\mathcal M}\lambda,
                         \mathrm{II}(X,X)\rangle .
 \end{aligned}
 \tag{V39.14}\label{c18v39:cost}
\]
The Riemann tensor convention has sectional curvature $+1$
on a unit round sphere.

\begin{proposition}[Native curvature with all coarea and shape terms]
\label{c18v39:Gauss}
For the horizontal lift $X$ of $v$,
\[
 \mathcal R_Q(v,v)=
 2|X|^2-2\operatorname{Hess}_{\mathcal M}u(X,X)
        -\mathcal D_c(X)+\mathcal S_{\rm orb}(X).
 \tag{V39.15}\label{c18v39:nativeRic}
\]
Moreover the shape tensor is determined by the original
sampled map:
\[
 \mathrm{II}(v,w)
       =-C^*G^{-1}(\nabla dc)(v,w)
                  \quad(v,w\in T\mathcal F_g).
 \tag{V39.16}\label{c18v39:II}
\]
\end{proposition}
\begin{proof}
Contract the Gauss equation over an orthonormal fibre basis.
Since the ambient Ricci tensor is $2\mathfrak g$, it gives
\[
 \operatorname{Ric}_{F}(X,X)
 =2|X|^2-\sum_a\langle R^{\mathcal M}(X,n_a)n_a,X\rangle
 +\langle\mathsf H,\mathrm{II}(X,X)\rangle
 -\sum_j|\mathrm{II}(X,e_j)|^2 .
\]
For any ambient scalar $\psi$ the restriction Hessian satisfies
\[
 \operatorname{Hess}_{F}(\psi|_F)(X,X)
 =\operatorname{Hess}_{\mathcal M}\psi(X,X)
                  +\langle\nabla^{\mathcal M}\psi,
                                      \mathrm{II}(X,X)\rangle .
\]
Substitute $\log\alpha=2u-\lambda$, then use
\eqref{c18v39:Ricgain}. Every term has the sign in
\eqref{c18v39:cost}. For \eqref{c18v39:II}, differentiate
$dc(w)=0$ covariantly in a fibre-tangent direction $v$.
The target-connection term is zero because both input
derivatives of $c$ vanish. Thus
$C\,\mathrm{II}(v,w)=-(\nabla dc)(v,w)$.
The inverse of $C$ on the normal space is $C^*G^{-1}$.
\end{proof}

The normal-curvature sum is between $0$ and $2|X|^2$,
since all ambient product sectional curvatures are nonnegative.
There is no corresponding sign assertion for the remaining
terms of $\mathcal D_c$. In particular its mean-curvature
trace may not be bounded by a dimension-free operator norm
without paying the trace contraction.
For computing its coarea Hessian one has, in compatible
orthonormal target frames with covariant derivatives,
\[
 \operatorname{Hess}_{\mathcal M}\lambda(X,X)
 =\tfrac12\operatorname{Tr}\!
 \left[G^{-1}(\nabla^2G)(X,X)
       -G^{-1}(\nabla_XG)G^{-1}(\nabla_XG)\right].
 \tag{V39.17}\label{c18v39:logdetH}
\]
This follows by two differentiations along an ambient geodesic,
parallel transporting target frames. It identifies a specific
third-map-derivative and trace cost, not a bound for it.
Ignoring the $\mathrm{II}$ term when changing an ambient
Hessian into a fibre Hessian would be incorrect.

\subsection{Exact metric on the native center obstruction}

There is an all-time calculation at the actual stationary
configuration used in V38. Take even $\ell$, so $N=\ell n$
is even, and put
$U^*_{x,1}=I$, $U^*_{x,2}=(-1)^{x_1}I$,
$U^*_{x,3}=(-1)^{x_1+x_2}I$.
Every sampled straight chain has product $I$:
its center sign is constant along its own direction and
is raised to the even power $\ell$. Thus this example lies
in the coarse fibre $g=I$, not merely in an unspecified
exceptional coarse fibre.
After root-fixed forest gauge fixing it remains central.

Use full color tangent spaces, and let
$A=d_1^*d_1\otimes I_3$ as in V38.
Let $D$ be the pinned ordinary gradient with the orientation
of V34, let $P_D=D(D^*D)^{-1}D^*$, and let $I_C$ insert
the free chord link coordinates. Let $B$ sum tangents along
each sampled chain. At this central point all adjoint transports
are identities, so
\[
 AD=0,\qquad BD=0,\qquad BB^*=\ell I.
 \tag{V39.18}\label{c18v39:centralmaps}
\]
The first identity is curl of a gradient; the second telescopes
between pinned roots.

\begin{theorem}[Opposite native metric responses at the two stationary backgrounds]
\label{c18v39:stationary}
At the all-minus center background, for every finite $s\ge0$,
\[
 \begin{aligned}
 G_-(s)&=B e^{2sA}B^*,&
 q_-(s)&=I_C^*(e^{-2sA}-P_D)I_C,\\
 q_-'(s)&=-2I_C^*Ae^{-2sA}I_C\preceq0,&
 a_-'(s)&=2a_-(s)I_C^*Ae^{-2sA}I_Ca_-(s)\succeq0 .
 \end{aligned}
 \tag{V39.19}\label{c18v39:minus}
\]
Both $q_-(s)$ and $G_-(s)$ are strictly positive.
At the flat identity background the exact formulas instead are
\[
 G_+(s)=B e^{-2sA}B^*,\qquad
 q_+(s)=I_C^*(e^{2sA}-P_D)I_C,\qquad
 q_+'(s)\succeq0,\quad a_+'(s)\preceq0.
 \tag{V39.20}\label{c18v39:plus}
\]
All statements are in the original induced and quotient metrics.
\end{theorem}
\begin{proof}
At the all-minus point the force vanishes and its derivative
is $A$; at flatness it is $-A$. Hence the forward differentials
are $e^{sA}$ and $e^{-sA}$ respectively.
Compose with $B$ to obtain the two normal Grams.
The inverse-flow chord insertions are $e^{-sA}I_C$ and
$e^{sA}I_C$. V34's exact quotient metric is their Gram
after projection orthogonal to $D$.
By $AD=0$, either exponential acts as the identity on
$\operatorname{Ran}D$ and commutes with $P_D$.
This gives both formulas for $q$ and, by differentiation,
their signs and the inverse-matrix derivatives.

For strict positivity, if $e^{-sA}I_Cv$ is a pinned gauge
tangent, applying $e^{sA}$ shows $I_Cv\in\operatorname{Ran}D$.
Its forest coordinates vanish. A pinned gradient vanishing
on the forest has every vertex value zero, hence $v=0$.
The same argument applies at flatness. Full row rank of $B$
and invertibility of the exponentials prove positivity of $G$.
\end{proof}

Writing $M=|\mathcal P_N|=3N^3$, the all-minus determinant
ledger at the point is
\[
 \log J_{-s}=-12sM,\qquad
 \log\det a_-(s)=\log\det K-\log\det G_-(s)+24sM.
 \tag{V39.21}\label{c18v39:centraldet}
\]
It follows from V38's exact Liouville identity and
\eqref{c18v39:det}. Here $K=D^*D$ at the stationary point
does not vary with $s$. It does vary when the configuration
is varied; no configuration derivative of $K$ is set to zero.

Thus V38's positive product Hessian
$6(I-e^{-2sA})$ occurs together with an increasing native
kinetic coefficient. It is not an isolated negative
potential term for that coefficient. Conversely the favorable
product Hessian at flatness comes with the opposite metric
change. Neither pointwise matrix monotonicity proves a
conditional spectral gap or a sign for \eqref{c18v39:nativeRic}.
Those require configuration derivatives of the metric,
not only its values along these stationary backgrounds.

\subsection{The integrated and source-weighted obligations in the correct geometry}

For a smooth scalar $f$ on the quotient, the exact integrated
identity is
\[
 \|\bar L_g f\|_{L^2(p_g)}^2
  =\int\left(
       |\operatorname{Hess}_q f|^2+
       \mathcal R_Q(\nabla^q f,\nabla^q f)
       \right)p_g\,dy .
 \tag{V39.22}\label{c18v39:integrated}
\]
This is the weighted Bochner identity on a compact manifold
without boundary, using \eqref{c18v39:weight}.
A volume-uniform lower bound by
$\delta\int|\nabla^q f|^2p_gdy$ would give the required
reduced conditional gap by applying it to eigenfunctions.
V39 does not prove that lower bound.
The term that must be controlled is the full signed
\eqref{c18v39:nativeRic}, including its nonnegative orbit
gain, not $2I-\operatorname{Hess}_{\rm product}\log p_g$.
Native probabilities alone do not bound the weighted gradient
integral of its negative part for arbitrary low-energy states.

For completeness, the source target also has an exact
one-form formulation. Let $\delta_\omega$ be the adjoint of
$d$ for $\omega_gd\operatorname{vol}_q$ and let
$L_1=d\delta_\omega+\delta_\omega d$ on one-forms.
On the closure of exact forms it intertwines $d\bar L_g=L_1d$.
For any centred smooth real source $r_v=v\cdot r_g$,
\[
 \begin{aligned}
 \operatorname{Var}_{p_g}(r_v)
    &=\langle dr_v,L_1^{-1}dr_v\rangle,\\
 \langle r_v,\bar L_g^{-1}r_v\rangle
    &=\langle dr_v,L_1^{-2}dr_v\rangle .
 \end{aligned}
 \tag{V39.23}\label{c18v39:oneform}
\]
The inverses here are restricted to exact forms.
To prove these identities, expand $r_v$ in nonconstant
orthonormal scalar eigenfunctions $f_j$ with eigenvalues
$\lambda_j>0$. Their differentials satisfy
$L_1df_j=\lambda_jdf_j$ and
$\langle df_i,df_j\rangle=\lambda_j\delta_{ij}$.
Both displayed equalities follow term by term and then by
spectral closure. The compact fixed-regulator gap suffices
for this identity; uniform inverse bounds are not assumed.
The source remains \eqref{c18v39:source}. In particular
controlling only its first inverse one-form moment would
control its variance, not yet its native inverse cost.
This is a target specification, not a new paid inverse estimate.

\subsection{Verification, ownership and the next calculation}

The diagnostic
\texttt{scripts/verify\_ym\_c18\_v39\_native\_curvature\_geometry.py}
checks three exact nonorthogonal determinant/Schur fixtures,
the signs in a compact warped-torus quotient, and the
ambient-to-sphere Hessian correction. These are explicitly
algebraic/geometric fixtures, not substitute YM measures.
On the actual $N=10,\ell=2,n=5$ cubic torus it builds the
3,000-link one-color curl, 875-column pinned gradient and
375-row full sampler. It verifies curl--gradient and
sampling--gradient annihilation, all identity coarse outputs,
and $BB^*=2I$. Fourier exponentiation tests the exact
stationary metric formulas on four chord vectors at six times,
including $s=4$. The color factor is exactly $I_3$.
No native Ricci tensor, interacting vacuum or unknown
spectral gap is numerically certified.

Only V39 replaces V38; removing this appendix and restoring
the title recovers V38 exactly. All 52 other sources are
unchanged. Only \texttt{.tex} files remain in
\texttt{latex\_notes}; builds and diagnostics are outside.
V40 is due for both the companion checkpoint and the broad
primary-literature sweep; V100 remains the archive/restart.

The preceding turn was progress, and this iteration adds
the exact native geometric target and a stationary evaluation
that changes the next calculation. It does not pay the
block-time conditional gap. At V40, the useful questions
are whether the signed coarea/shape contractions in
\eqref{c18v39:cost} admit local integrated bounds without a
coarse-dimension trace loss, and whether the source in
\eqref{c18v39:oneform} avoids any uncontrolled low modes.
One must not declare the flow Jacobian absent from the score,
or count its entire product Hessian again as an independent
native potential after \eqref{c18v39:weight}.
Weak-coupling scale transport, physical return, reconstruction
and a positive physical Yang--Mills mass gap remain unproved.
% END CYCLE 18 VERSION 39 APPEND
% BEGIN CYCLE 18 VERSION 40 APPEND
\clearpage
\section[V40: stationary score cancellation and a native derivative barrier]{V40: stationary score cancellation and a native derivative barrier}
\label{c18v40:context}

V39 identifies the correct native curvature tensor but does not
bound its negative part in the conditional law. The scheduled
companion and literature review suggests testing the actual
source before placing a worst-direction curvature estimate
around it. We do that at two exact native stationary backgrounds,
using the complete chain-averaged lift already constructed in
V25. No coarse column, fine harmonic mode or conditional measure
is replaced.

There is a useful cancellation: the forward-flow growth in the
score derivative is normal to the sampled fibre and disappears
under the \emph{full fibre} projection. At the all-minus background
this gives an all-time, volume-independent bound, including the
actual interacting vacuum in V37's coupling disk. At flatness the
remaining term instead has an exponentially large native derivative
at the required time $s=\ell^2$. Both conclusions are exact
pointwise statements for the original model. They neither prove
nor disprove the conditional inverse estimate.

The previous iteration was progress. This iteration changes the
next source estimate: a global polynomial bound on this explicit
lift's native score gradient is false, even though its gradient is
well controlled at V38's product-curvature obstruction. A
source-weighted inverse must retain the operator rather than
substituting that false pointwise envelope.

\subsection{Scheduled checkpoint: what transfers and what does not}

The current companion files were rehashed and their relevant
terminal results reread: NS V26, SM--Einstein V72, source-selection
V65, stationarity V61, exact-action Einstein V55, perturbative
ESM V80, arithmetic-loading V17, shell-mixing V16 and parallel
YM V5. They are unchanged since the most recent reviews.
The useful lessons are precise but are not new YM assumptions.
NS V26 contracts on the actual input class; none of its numerical
kernel constants is imported. Selection V65's finite positive
defect witness concerns a specified represented source, not an
arbitrary spectral subspace. Stationarity V61 retains its original
constraint defect and still lacks a state-only source calculation.
Einstein V55 still has an uncomputed same-root tangency coefficient.
ESM V80 is a fixed two-loop, weight-six bosonic result, not a
nonperturbative conditional inequality. SM V72's auxiliary
many-copy limit, arithmetic V17's represented-projector requirement,
shell V16's measured one-loop calculation, and parallel YM V5's
extensive-charge warning retain their stated scopes.

The broad primary-literature checkpoint covered the following
directions. Reading levels and rejected transfers are recorded
so that an abstract is not treated as a proved input.
\begin{enumerate}
\item Bauerschmidt and Bodineau.\par
\emph{Log-Sobolev inequality for the continuum sine-Gordon model},
\href{https://arxiv.org/abs/1907.12308}{arXiv:1907.12308}.
The Gaussian setup, Theorems 1.2 and 2.5, and the displayed
time-dependent gradient calculation were inspected. Their
criterion combines the evolving covariance with the renormalized
Hessian; negative instantaneous contributions need not be
estimated separately. Our compact conditional law has no proved
Polchinski representation satisfying their hypotheses.
\item Bauerschmidt--Bodineau--Dagallier,
\emph{Stochastic dynamics and the Polchinski equation: an introduction},
\href{https://arxiv.org/abs/2307.07619}{arXiv:2307.07619}.
The primary abstract was checked for its links to stochastic
localization, transport and multiscale functional inequalities.
No theorem is imported from the survey.
\item Leli\`evre--Zhang,
\emph{Pathwise estimates for effective dynamics: the case of nonlinear
vectorial reaction coordinates},
\href{https://arxiv.org/abs/1805.01928}{arXiv:1805.01928}.
The coarea law, Section 2.2 and the pathwise theorem were reread.
Their projected derivative budgets are relevant, but Assumption 4
requires the uniform conditional Poincar\'e inequality we still lack.
This remains a method comparison, not a way to obtain that premise.
\item Rose--Stollmann,
\emph{Manifolds with Ricci curvature in the Kato class: heat kernel
bounds and applications},
\href{https://arxiv.org/abs/1804.04094}{arXiv:1804.04094}.
The resolvent/heat Kato quantities, Section 2's one-form domination
and Theorem 3.2 were inspected. Smallness is measured relative to
the native heat operator, not just by a bad-set probability.
Their unweighted geometric theorem is not applied to our weighted,
growing-dimensional quotient. In particular its explicit
dimension-dependent smallness threshold cannot be suppressed.
\item Shen--Zhu--Zhu,
\emph{A stochastic analysis approach to lattice Yang--Mills at strong
coupling}, \href{https://arxiv.org/abs/2204.12737}{arXiv:2204.12737}.
The primary abstract's strong-coupling regime and Wilson Gibbs
measure were rechecked. They are not the weak-coupling quantum
ground-state conditional law in this programme.
\item Chevyrev--Shen,
\emph{Invariant measure and universality of the 2D Yang--Mills
Langevin dynamic},
\href{https://arxiv.org/abs/2302.12160}{arXiv:2302.12160},
and Chevyrev--Qu--Shen,
\emph{Scaling limit of the 3D abelian Yang--Mills Langevin dynamics},
\href{https://arxiv.org/abs/2608.27828}{arXiv:2608.27828v2},
revised 17 September 2026. The primary abstracts were checked.
The former's invariant-measure identification and the latter's
abelian weak-coupling scaling are useful model-ownership examples;
neither is a four-dimensional nonabelian quantum mass-gap theorem.
\end{enumerate}

This is a broad targeted sweep, not an exhaustive literature review
or an independent validation of every cited proof. The calculation
below is proved directly. It uses V25's existing lift and V39's
existing geometry; neither construction is claimed as new.

\subsection{The full averaged lift and the source being differentiated}

Retain $\mathcal M=SU(2)^{E_{\rm f}}$, $N=\ell n$,
$n\ge5$, $\ell\ge2$, the unit-round product metric,
$\mathcal W=\sum_p w_p$, and spatial flow generated by
$\nabla\mathcal W$. Let $b=\mathfrak b_\ell$ and
$c_s=b\circ\Phi_s$. In flowed variables $W$ set
\[
 B_W=Db(W),\qquad B_WB_W^*=\ell I,\qquad
 Y_v(W)=\ell^{-1}B_W^*v,\qquad
 \mathcal B_{s,v}(U)=D\Phi_{-s}(\Phi_sU)Y_v(\Phi_sU).
 \tag{V40.1}\label{c18v40:lift}
\]
$v$ is in the full fixed coarse invariant frame.
Disjoint straight chains and their orthogonal prefix adjoints
give $B_WB_W^*=\ell I$. V25 proves
$\operatorname{div}_{dW}Y_v=0$ and
$Dc_s\mathcal B_{s,v}=v$. Thus this is the same complete
averaged bank, not a selected mask.

Put $u_\beta=\log\Omega_\beta$ and
\[
 h_{\beta,s}(W)=
       \Omega_\beta(\Phi_{-s}W)^2J_{-s}(W),\qquad
 S^{\rm av}_{s,v}=(Y_v\log h_{\beta,s})\circ\Phi_s,\qquad
 r^{\rm av}_{s,v}=S^{\rm av}_{s,v}
                       -\mathbb E[S^{\rm av}_{s,v}\mid c_s].
 \tag{V40.2}\label{c18v40:score}
\]
The divergence change-of-variables formula gives
$S^{\rm av}_{s,v}=\operatorname{div}_{\Omega_\beta^2dU}
\mathcal B_{s,v}$. In particular the Jacobian is retained.

The averaged source is not declared equal to V38's tree source.
If $\mathcal A^{\rm tr}_{s,v}$ denotes that full tree lift, then
$V_v=\mathcal A^{\rm tr}_{s,v}-\mathcal B_{s,v}$ is vertical and
\[
 r^{\rm tr}_{s,v}-r^{\rm av}_{s,v}
       =\operatorname{div}_{\nu_g}V_v,\qquad
 \|r^{\rm tr}_{s,v}-r^{\rm av}_{s,v}\|_{-1,g}
       \le\|V_v\|_{L^2(\nu_g)} .
 \tag{V40.3}\label{c18v40:change}
\]
Here the inverse norm is for the original induced fibre energy,
restricted to internal-gauge-invariant scalars, equivalently
V33's reduced inverse. For the inequality one may first
project $V_v$ off internal gauge orbits as in V34.
To verify \eqref{c18v40:change}, disintegrate integration by
parts for a vertical field. Its conditional divergence has
mean zero. Pair it with a centred $f$ and apply Cauchy--Schwarz
to $-\int df(V_v)\,d\nu_g$, then use the variational dual
definition of the inverse norm. This standard comparison does
not provide a new uniform bound on $V_v$.
The rest of V40 concerns precisely $r^{\rm av}$.

\subsection{A stationary normal cancellation with the vacuum term retained}

At a stationary point $U_*$ of $\mathcal W$ suppose also that
$\nabla u_\beta(U_*)=0$. All central configurations satisfy
both conditions. Indeed diagonal global gauge conjugation fixes
such a configuration; invariance of either scalar forces every
one-link color gradient to be an adjoint-invariant vector in
$\mathfrak{su}(2)$, hence zero. The ground scalar is gauge
invariant by uniqueness of the positive finite-regulator ground
state.

In original tangent frames at $U_*$ write
\[
 H_*=\operatorname{Hess}\mathcal W(U_*),\qquad
 M_\beta=\operatorname{Hess}u_\beta(U_*),\qquad
 B=Db(U_*),\qquad E_s=e^{sH_*}.
 \tag{V40.4}\label{c18v40:stationarydata}
\]
These are ordinary self-adjoint tangent endomorphisms where
appropriate. No commutation between $M_\beta$ and $H_*$ is
assumed. Define
\[
 G_s=BE_s^2B^*,\qquad
 P_s=I-E_sB^*G_s^{-1}BE_s .
 \tag{V40.5}\label{c18v40:P}
\]
$P_s$ is the original orthogonal projection onto
$T_{U_*}c_s^{-1}(b(U_*))=\ker(BE_s)$.
It is the \emph{full fibre} projection, not V34's projection
only off internal gauge orbits.

Let $\mathcal J_{\beta,s}$ send a coarse vector $v$ to the
native fibre gradient of $r^{\rm av}_{s,v}$ at $U_*$.
For internal-gauge-invariant sources this gradient is horizontal
for $\mathcal F_g\to\mathcal Q_g$, so its norm is the true
quotient-gradient norm.

\begin{theorem}[Exact full-bank stationary score derivative]
\label{c18v40:cancellation}
For every finite $s\ge0$, under the two stationary conditions above,
\[
 \boxed{\displaystyle
 \mathcal J_{\beta,s}
       =\frac{2}{\ell}P_s(M_\beta-3I)e^{-sH_*}B^*.}
 \tag{V40.6}\label{c18v40:J}
\]
At a central configuration the centred source itself is zero:
$r^{\rm av}_{s,v}(U_*)=0$ for every $v$.
\end{theorem}
\begin{proof}
At a critical point the flow fixes $U_*$ and
$D\Phi_s(U_*)=E_s$. The exact Liouville formula in V38 is
\[
 \log J_{-s}(W)=12\int_0^s\mathcal W(\Phi_{-t}W)\,dt .
\]
Its gradient at $U_*$ vanishes. Because the inner scalar
gradient is zero, the Hessian chain rule has no second-flow
acceleration term. Consequently
\[
 \begin{aligned}
 \operatorname{Hess}\log J_{-s}(U_*)
 &=12\int_0^s e^{-tH_*}H_*e^{-tH_*}\,dt
       =6(I-e^{-2sH_*}),\\
 \operatorname{Hess}\log h_{\beta,s}(U_*)
 &=2e^{-sH_*}M_\beta e^{-sH_*}
                         +6(I-e^{-2sH_*}) .
 \end{aligned}
 \tag{V40.7}\label{c18v40:Hess}
\]
These identities hold for zero or negative eigenvalues of $H_*$
as well, by finite-dimensional functional calculus.

Since $\nabla\log h_{\beta,s}(U_*)=0$, differentiating
$Y_v\log h_{\beta,s}$ contributes no derivative of $Y_v$.
Pulling back by $\Phi_s$ and projecting tangentially gives
\[
 \mathcal J_{\beta,s}
 =\frac1\ell P_s
       \{2M_\beta e^{-sH_*}
                    +6(E_s-e^{-sH_*})\}B^* .
 \tag{V40.8}\label{c18v40:raw}
\]
The conditional mean in \eqref{c18v40:score} is constant along
the fibre, so its tangential derivative is zero.
The decisive exact identity is $P_sE_sB^*=0$, from
\eqref{c18v40:P}. Cancel that term in \eqref{c18v40:raw}
to obtain \eqref{c18v40:J}. Notice that moving
$M_\beta$ through $e^{-sH_*}$ would be unjustified.

For internal gauge tangents $D\xi$, gauge invariance and
criticality give $H_*D=0$, $M_\beta D=0$, while $BD=0$.
Thus the result is orthogonal to those orbit tangents,
as also follows from invariance of the source.

Finally $S^{\rm av}_{s,v}(U_*)=0$. At a central coarse output
the smooth positive marginal $\rho_s(g)$ is invariant under
diagonal conjugation and its coarse gradient is zero by the
same adjoint argument. Integration by parts for the full lift
gives $\mathbb E[S^{\rm av}_{s,v}\mid c_s=g]
=X_v\log\rho_s(g)$, since the coarse invariant frame has zero
Haar divergence. This proves the assertion about centering.
\end{proof}

\subsection{The exact Gram and an all-time bound at the all-minus background}

At $\beta=0$ the actual electric ground scalar is one, so
$M_0=0$, not a Gaussian approximation. The preceding theorem
gives the full coarse derivative Gram
\[
 \boxed{\displaystyle
 \mathcal J_{0,s}^*\mathcal J_{0,s}
   =36\left\{\ell^{-2}Be^{-2sH_*}B^*
                         -(Be^{2sH_*}B^*)^{-1}\right\}\succeq0 .}
 \tag{V40.9}\label{c18v40:Gram}
\]
Indeed insert $P_s$ between the two factors
$e^{-sH_*}B^*$ in the squared norm in \eqref{c18v40:J}
and use $BB^*=\ell I$. Positivity is a squared norm, not an
assumed comparison of two noncommuting matrices.
Writing $Q=B/\sqrt\ell$ identifies the bracket with
$\ell^{-1}\{Qe^{-2sH_*}Q^*-(Qe^{2sH_*}Q^*)^{-1}\}$,
the usual compression/inverse defect.

For a fixed regulator its first nonzero small-time coefficient is
\[
 \mathcal J_{0,s}^*\mathcal J_{0,s}
  =\frac{144s^2}{\ell^2}BH_*P_0H_*B^*+O_{N,\ell}(s^3),
 \qquad P_0=I-\ell^{-1}B^*B .
 \tag{V40.10}\label{c18v40:small}
\]
To check this, expand both exponentials to second order and
the inverse of $\ell I+2sBH_*B^*+2s^2BH_*^2B^*$.
The constant and linear terms cancel, and the quadratic
term is the displayed expression. If $\operatorname{Ran}B^*$
is $H_*$-invariant, \eqref{c18v40:Gram} vanishes for every
$s$, since the exponential then preserves both blocks.
Thus normal--tangent mixing, rather than the mere size of
$H_*$, produces this stationary source derivative.

\begin{corollary}[Native all-minus source derivative, including the interacting vacuum]
\label{c18v40:minus}
For even $\ell$, take V38--V39's central configuration
\[
 U^*_{x,1}=I,\qquad U^*_{x,2}=(-1)^{x_1}I,\qquad
 U^*_{x,3}=(-1)^{x_1+x_2}I .
\]
Then $g=I$, $H_*=A=d_1^*d_1\otimes I_3\succeq0$ and,
for every $s\ge0$,
\[
 \|\mathcal J_{\beta,s}\|_{\rm op}
       \le\frac{6+2\|M_\beta\|_{\rm op}}{\sqrt\ell}.
 \tag{V40.11}\label{c18v40:minusbound}
\]
In particular, with V37's
$\beta_*=3/(512e)$ and $R_\beta=64e|\beta|\le3/8$,
\[
 \|\mathcal J_{\beta,s}\|_{\rm op}
 \le\frac{6+8R_\beta/e}{\sqrt\ell}
 \le\frac{6+3/e}{\sqrt\ell}
       \quad(0\le\beta\le\beta_*).
 \tag{V40.12}\label{c18v40:interacting}
\]
At $\beta=0$, the full Gram is bounded above by
$36\ell^{-1}I$. Every bound is independent of $n$ and $s$.
\end{corollary}
\begin{proof}
Use $\|P_s\|\le1$, $\|e^{-sA}\|\le1$ and
$\|B^*\|=\sqrt\ell$ in \eqref{c18v40:J}.
V37 bounds the actual log-vacuum Hessian by $4R_\beta/e$,
giving \eqref{c18v40:interacting}. The color and every coarse
column are included in these operator norms.
\end{proof}

This is at the native point which defeats V38's global
product-curvature criterion. It does not assert a uniform
neighborhood radius, a probability estimate for that neighborhood,
or a bound on the source inverse integrated over the fibre.

\subsection{A flat native derivative barrier at the required block time}

At flat identity $H_*=-A$. There is no all-time replacement of
$e^{sA}$ by a contraction in \eqref{c18v40:J}. We can compute
a complete nine-dimensional coarse source sector exactly:
the three constant coarse link directions, with all three
colors. Normalize each constant coarse vector to have norm one.
The fine harmonic alias is \emph{included}, and no coarse
longitudinal direction is removed.

Use unitary fine and coarse Fourier transforms as in V17.
At coarse momentum zero the aliases are
$k=2\pi m/\ell$, $m\in\{0,\ldots,\ell-1\}^3$.
For direction $i$, the straight-chain multiplier is
$\ell^{-1/2}$ when $m_i=0$ and zero otherwise.
If $m_i=0$, the $i$th component of $d(k)$ is zero and
$A(k)e_i=\lambda(k)e_i$, with
$\lambda(k)=4\sum_j\sin^2(\pi m_j/\ell)$.
All off-diagonal sampled entries vanish. Define
\[
 \sigma_\ell^\pm(s)
       =\sum_{j=0}^{\ell-1}
                      \exp\{\pm8s\sin^2(\pi j/\ell)\}.
 \tag{V40.13}\label{c18v40:sigma}
\]
Summing the two remaining alias indices proves
\[
 \left.Be^{\pm2sA}B^*\right|_{q=0}
       =\ell^{-1}\big(\sigma_\ell^\pm(s)\big)^2 I_9 .
 \tag{V40.14}\label{c18v40:alias}
\]
In particular the $m=0$ harmonic contribution is
$\ell^{-1}I_9$; discarding it would change the inverse.

\begin{theorem}[Exact stationary zero-momentum source Gram]
\label{c18v40:qzero}
At $\beta=0$, the restrictions of the native derivative Grams
to this full constant-source sector are $\tau_-(s)I_9$ at the
all-minus background and $\tau_+(s)I_9$ at flatness, where
\[
 \begin{aligned}
 \tau_-(s)&=\frac{36}{\ell^3}\big(\sigma_\ell^-(s)\big)^2
                  -\frac{36\ell}{\big(\sigma_\ell^+(s)\big)^2},\\
 \tau_+(s)&=\frac{36}{\ell^3}\big(\sigma_\ell^+(s)\big)^2
                  -\frac{36\ell}{\big(\sigma_\ell^-(s)\big)^2}.
 \end{aligned}
 \tag{V40.15}\label{c18v40:tau}
\]
For every even $\ell\ge2$, at the actual block time $s=\ell^2$,
\[
 \boxed{\displaystyle
 \tau_+(\ell^2)\ge
             18\ell^{-3}e^{16\ell^2},\qquad
 \|\mathcal J_{0,\ell^2}\|_{\rm op}
             \ge3\sqrt2\,\ell^{-3/2}e^{8\ell^2}.}
 \tag{V40.16}\label{c18v40:barrier}
\]
These bounds hold for every coarse side $n\ge5$.
\end{theorem}
\begin{proof}
The two background signs in \eqref{c18v40:Gram} and the
literal alias sum \eqref{c18v40:alias} give
\eqref{c18v40:tau}. For even $\ell$, the term $j=\ell/2$
gives $\sigma_\ell^+(s)\ge e^{8s}$, while the harmonic term
gives $\sigma_\ell^-(s)\ge1$. Therefore
\[
 \tau_+(s)\ge36\ell^{-3}e^{16s}-36\ell .
\]
At $s=\ell^2$, $e^{16\ell^2}\ge2\ell^4$, for example from
the quadratic term of its exponential series.
The subtracted term is at most half the first term.
Taking square roots proves the norm bound. No estimate of
an unknown vacuum or a spectral gap enters this argument.
\end{proof}

For $\ell=2$, set $z=e^{-8s}$. The exact values simplify to
\[
 \tau_-(s)=\frac92
      \frac{(1-z)^2(1+6z+z^2)}{(1+z)^2},\qquad
 \tau_+(s)=e^{16s}\tau_-(s).
 \tag{V40.17}\label{c18v40:ell2}
\]
For example $\tau_-(4)$ is approximately $4.5$, whereas
$\tau_+(4)$ is approximately $2.806\,10^{28}$.
The identities, not these floating values, prove the barrier.
Both centred sources are zero at their central backgrounds;
zero source value does not imply zero source derivative.

The barrier is for the declared averaged bank, not every possible
lift. It rules out a volume-and-block-scale-uniform polynomial
pointwise bound for that bank's native score derivative.
It is not a lower bound for its inverse cost: a growing source
derivative may be compensated by a correspondingly stiff
conditional operator. Even on a fixed circle,
$r_k(x)=\sin(kx)$ has gradient norm of order $k$ but
$\langle r_k,(-\partial_x^2)^{-1}r_k\rangle=1/(2k^2)$
under probability Haar measure. This elementary comparison is
only a warning against that invalid inference, not a YM model.

\subsection{Verification and the next source-weighted task}

The new diagnostic
\texttt{scripts/verify\_ym\_c18\_v40\_stationary\_score.py}
checks the full matrix cancellation with a noncommuting rational
vacuum-Hessian fixture, the compression/inverse identity and
the second-order coefficient. It separately proves
\eqref{c18v40:ell2} by exact symbolic algebra on the actual four
surviving zero-momentum aliases with eigenvalues $0,4,4,8$.
That native alias calculation is distinct from the abstract
rational fixture.

On the $N=10,\ell=2,n=5$ native lattice it rechecks the central
geometry, reconstructs all 125 complete 24-dimensional
one-color alias symbols and all three coarse rows, and tests
the projection identity at four times. An independent
real-space straight-chain Fourier calculation checks
\eqref{c18v40:alias}. The block times for $\ell=2,4,6$ and
the full $\ell=2,s=4$ projection are checked with 100-digit
arithmetic. These are regression checks of identities:
the all-volume estimates are proved above, while no interacting
vacuum Hessian, native neighborhood probability or conditional
inverse is numerically certified.

Only V40 replaces V39. Removing this appendix and restoring the
title recovers V39 byte for byte; all 52 other sources are
unchanged. Builds stay outside \texttt{latex\_notes}, which contains
only \texttt{.tex} sources. The next companion checkpoint is V45,
the next broad literature checkpoint V50, and the archive/restart
remains V100.

V39's exact one-form identity applies to the declared averaged
source as well:
\[
 \langle r^{\rm av}_{s,v},\bar L_g^{-1}r^{\rm av}_{s,v}\rangle
 =\langle dr^{\rm av}_{s,v},L_{1,g}^{-2}
                                      dr^{\rm av}_{s,v}\rangle .
 \tag{V40.18}\label{c18v40:target}
\]
It must now be estimated jointly. The flat derivative barrier
shows why replacing $dr$ by a global polynomial envelope
before using $L_{1,g}^{-2}$ cannot work for this bank.
Conversely, the cancellation at the all-minus background shows
why its bad product curvature cannot by itself demonstrate a
large native source inverse. A useful next calculation must
either retain this joint source/operator structure on native
neighborhoods with controlled conditional weight, or build and
pay for a different full lift. Another isolated point value,
a bare bad-set probability, or another Gaussian reference gap
would not discharge that task.

The required block-time conditional inverse, weak-coupling scale
transport, physical return, interacting continuum reconstruction
and positive physical Yang--Mills mass gap remain unproved.
% END CYCLE 18 VERSION 40 APPEND
% BEGIN CYCLE 18 VERSION 41 APPEND
\clearpage
\section[V41: a normal lift on a polynomial flat neighborhood]{V41: a normal lift on a polynomial flat neighborhood}
\label{c18v41:context}

V40's exponential derivative barrier concerns a specified averaged
lift, not the conditional inverse itself. We now test the already
available full normal lift rather than transferring that barrier
to every possible lift. The new results are an exact identification
of the averaged lift's stationary excess energy, the complete normal
score formula, and explicit polynomial derivative bounds on a
nonlinear native neighborhood through the required time $s=\ell^2$.
The neighborhood is not asserted to be typical. The scalar
divergence bound still contains the fine dimension.

This uses V24's full normal lift, V26--V27's minimum-transport
distinction, V29's invariant-frame differential calculus, and V40's
stationary identity. Those constructions are not new here.
The previous turn was progress. The current V40 source and all
52 other note sources were checked before editing; the companion
and broad-literature checkpoints remain V45 and V50.
For primary context, the conditional coarea density in
Leli\`evre--Zhang,
\href{https://arxiv.org/html/1805.01928v1}
{\emph{Pathwise estimates for effective dynamics: the case of
nonlinear vectorial reaction coordinates}}, equations (7)--(8),
was rechecked. It contains the inverse square-root Gram determinant.
No conditional Poincar\'e premise from that work is imported.
The Wilson-flow
\href{https://arxiv.org/abs/1006.4518}{primary abstract of L\"uscher}
was also rechecked for context only. All estimates below are derived
in our spatial, unit-round normalization.

\subsection{The full normal score, with its trace and coarea terms}

Keep $\mathcal M=SU(2)^{E_{\rm f}}$, $N=\ell n$, $n\ge5$,
$\ell\ge2$, $\mathcal W=\sum_p w_p$, and
$c_s=b_\ell\circ\Phi_s$ with $\dot\Phi=\nabla\mathcal W$.
Let $X_i$ be the fine invariant orthonormal frame, and use the
full coarse invariant frame. Matrix derivatives below are
derivatives of coefficients in these frames. Put
\[
 C=Dc_s,\quad G=CC^*,\quad Z=G^{-1},\quad
 R=C^*Z,\quad P=I-RC,\quad u=\log\Omega_\beta .
 \tag{V41.1}\label{c18v41:normal}
\]
V17 proves that $c_s$ is a submersion at every finite time.
Consequently $R$ is a globally smooth full-output lift, $CR=I$,
and $P$ is the original orthogonal fibre projection.
The lift is equivariant and orthogonal to every fibre tangent,
including internal gauge directions. It is not a Poisson-corrected
or score-free lift.

For a fixed coarse coefficient vector $v$, define the actual scores
\[
 S^R_v=\operatorname{div}_{\Omega_\beta^2dU}(Rv),\qquad
 r^R_v=S^R_v-\mathbb E[S^R_v\mid c_s].
 \tag{V41.2}\label{c18v41:score}
\]
Write $C_i=X_iC$, $G_i=C_iC^*+CC_i^*$, and
$\tau_c=\sum_i C_i e_i$. Since invariant frame fields have
zero Haar divergence, direct differentiation gives
\[
 S^R_v
 =\langle Z\tau_c,v\rangle
  -\sum_i\langle R^*e_i,G_iZv\rangle
  +2\langle\nabla u,Rv\rangle .
 \tag{V41.3}\label{c18v41:trace}
\]
Indeed $X_iR=C_i^*Z-C^*ZG_iZ$; take its $i$th fine
coefficient and sum. The middle sum is not deleted or
replaced by an operator norm without paying its contraction.
In these product invariant frames $\tau_c$ is also the map
tension: the domain diagonal connection terms vanish, and
the target connection term in each diagonal pair is a bracket
of a vector with itself.

There is an equivalent intrinsic formula. Let
$\lambda=\frac12\log\det G$, and let
$\mathsf H_{\rm fib}$ be the unnormalized mean-curvature
vector of the original fibre, with the sign convention of V39.
Then
\[
 S^R_v
 =\langle 2\nabla u-\nabla\lambda-\mathsf H_{\rm fib},Rv\rangle .
 \tag{V41.4}\label{c18v41:geometry}
\]
To prove it, flow the fibre by the projectable normal field $Rv$.
The coarse field $v$ preserves coarse Haar volume. The induced
fibre-volume variation is
$-\langle\mathsf H_{\rm fib},Rv\rangle$, by differentiating the
tangent Gram determinant. The coarea density relative to that
volume is $\Omega_\beta^2e^{-\lambda}$; its logarithmic variation
is $\langle2\nabla u-\nabla\lambda,Rv\rangle$.
Their sum is the variation of the disintegrated unnormalized
ambient measure, proving \eqref{c18v41:geometry}. Conditional
normalization then subtracts the mean in \eqref{c18v41:score}.
In particular there is no removal of the coarea determinant.
The original-coordinate formula avoids an inverse-flow
Jacobian representation, not its geometric content.

\subsection{What the stationary averaged-source barrier measures}

At a stationary point of $\mathcal W$ write $H=\operatorname{Hess}
\mathcal W$, $E=e^{sH}$ and $B=Db_\ell$, as in V40. At this point
the averaged lift and normal lift are
\[
 A^{\rm av}=\ell^{-1}E^{-1}B^*,\qquad
 R=EB^*(BE^2B^*)^{-1},\qquad CR=CA^{\rm av}=I .
 \tag{V41.5}\label{c18v41:stationarylifts}
\]
Thus $PA^{\rm av}=A^{\rm av}-R$. At $\beta=0$, V40.6 yields
the exact identity
\[
 \boxed{\mathcal J^{\rm av}_{0,s}=-6(A^{\rm av}-R),\qquad
 (A^{\rm av}-R)^*(A^{\rm av}-R)
       =\frac1{36}(\mathcal J^{\rm av}_{0,s})^*
                         \mathcal J^{\rm av}_{0,s}.}
 \tag{V41.6}\label{c18v41:excess}
\]
Both sides use the same original tangent metric. The proof is
just $BB^*=\ell I$, the displayed normal projection, and V40's
proved stationary score identity. Orthogonality also gives
$(A^{\rm av})^*A^{\rm av}=G^{-1}+
(A^{\rm av}-R)^*(A^{\rm av}-R)$.

At flat identity, in V40's complete nine-dimensional constant
coarse sector, this reads
\[
 \left.R^*R\right|_{q=0}
   =\frac{\ell}{(\sigma_\ell^-(s))^2}I_9\preceq\ell I_9,\qquad
 \left.(A^{\rm av}-R)^*(A^{\rm av}-R)\right|_{q=0}
   =\frac{\tau_+(s)}{36}I_9 .
 \tag{V41.7}\label{c18v41:excessflat}
\]
The fine harmonic alias is retained. Hence V40's exponentially
large derivative is exactly six times an unnecessary vertical
part of that particular lift at the stationary point. This
does not compute the derivative of $S^R$, or a conditional
integral of the lift difference. We compute the normal score's
own derivatives next; we do not transfer an inverse estimate
from the averaged bank.

\subsection{Local coefficient bounds without a fine-volume factor}

Use V29's moving-frame generator
\[
 \dot{\mathsf J}_t=\mathsf M(\Phi_tU)\mathsf J_t,\qquad
 (\mathsf M v)_e=\sum_jv_jX_jF_e+[F_e,v_e],\qquad
 F_e=(\nabla\mathcal W)_e .
\]
Let $D_z=\sum_i z_iX_i$ for a constant fine coefficient vector.
Ordered $D_wD_z$ derivatives are not symmetrized.
Retain V29's two-sided packed first-derivative bound, with
\[
 b_1=288\sqrt3,\qquad
 \|D_z\mathsf M\|\le b_1|z|.
\]
The following additional bound will be used:
\[
 \|D_wD_z\mathsf M\|\le b_2|w||z|,\qquad b_2=1728 .
 \tag{V41.8}\label{c18v41:Msecond}
\]
Here and below matrix norms are operator norms.
For a proof, decompose $\mathsf M=\sum_p\mathsf M^p$ as in V29.
Every entry of $X_jX_i\mathsf M^p$ has magnitude at most three
and vanishes unless its four indices lie in the twelve
link/color slots $I_p$. Ordered derivatives through order four
of a plaquette trace insert unit quaternions in its word and
have magnitude at most one. The differentiated commutator
contributes at most two; no number of terms depending on
the lattice occurs.

Contracting the four-index tensor for one plaquette with
$w,z,v$ bounds the output norm by
$3\cdot12^2|w_{I_p}||z_{I_p}||v_{I_p}|=432
|w_{I_p}||z_{I_p}||v_{I_p}|$.
Each output slot belongs to four plaquettes. Squaring their
sum, bounding $|w_{I_p}|\le|w|$, $|z_{I_p}|\le|z|$,
and using $\sum_p|v_{I_p}|^2=4|v|^2$ gives
$16\cdot432^2|w|^2|z|^2|v|^2$, which proves the claim.

For the straight-chain sampler $\mathsf B=Db_\ell$ put
\[
 a_{\rm B}=2\sqrt3\,\ell,\qquad b_{\rm B}=36\ell^{3/2}.
\]
V29 supplies $\|\mathsf B\|=\sqrt\ell$ and both packed
first-derivative bounds with constant $a_{\rm B}^2$. In addition,
\[
 \|D_wD_z\mathsf B\|\le b_{\rm B}|w||z| .
 \tag{V41.9}\label{c18v41:Bsecond}
\]
Indeed a block is an adjoint prefix. Two ordered invariant
derivatives insert two Lie-algebra representation generators,
each of norm at most two, so every scalar entry is at most four.
For one chain the tensor has at most
$3(3\ell)^3$ scalar entries, hence Hilbert--Schmidt norm at
most $36\ell^{3/2}$. The chain supports of all three input
slots are disjoint between coarse rows. Applying the local
bound and summing squared row norms proves the displayed
global trilinear bound. These deliberately loose constants
are independent of $n$.

\subsection{A quantitative nonlinear neighborhood through block time}

Set $T=\ell^2$ and define constants
\[
 \begin{aligned}
 k_T&=b_1Te^3,&
 l_T&=\tfrac32 b_1^2T^2e^5+b_2Te^4,\\
 p_1&=a_{\rm B}e^2+\sqrt\ell\,k_T,&
 p_2&=b_{\rm B}e^3+3a_{\rm B}e\,k_T+\sqrt\ell\,l_T,\\
 m_{\rm f}&=\frac2\pi e^{-3\pi^2},&
 \delta_\ell&=\min\left\{\frac18,\frac1{2b_1T},
                   \frac{m_{\rm f}}{2\sqrt\ell\,p_1}\right\}.
 \end{aligned}
 \tag{V41.10}\label{c18v41:constants}
\]
The symbol $l_T$ is a scalar majorant, not a block length.
Let $\mathcal U_\ell$ be the original product-metric ball
$d(U,I)<\delta_\ell$. No mode or link has been removed in
defining this neighborhood.

\begin{theorem}[Stable flow and full normal Gram on a native neighborhood]
\label{c18v41:tube}
For $U\in\mathcal U_\ell$ and $0\le t\le T$, the flow stays
within distance $2\delta_\ell$ of identity. Its forward
variational propagators have norm at most $e$.
For every $0\le s\le T$, in original invariant frames,
\[
 \begin{gathered}
 \|C_s\|\le e\sqrt\ell,\qquad
 \sum_i C_{i,s}^*C_{i,s}\preceq p_1^2I,\qquad
 \sum_i C_{i,s}C_{i,s}^*\preceq p_1^2I,\\
 \|D_wD_zC_s\|\le p_2|w||z|,\qquad
 G_s(U)\succeq\frac{m_{\rm f}^2}{4\ell}I .
 \end{gathered}
 \tag{V41.11}\label{c18v41:tubebounds}
\]
All coarse outputs, fine harmonics and longitudinal directions
are included. The constants and radius are independent of $n$.
\end{theorem}
\begin{proof}
At identity $\mathsf M(I)=-A$ with
$A=d_1^*d_1\otimes I_3\succeq0$, where $d_1$ is the lattice curl.
The coefficient bound for $D\mathsf M$ gives
$\|\mathsf M(U)-\mathsf M(I)\|\le b_1r$ at distance $r<1/4$,
by integration along the radial product geodesic.
The symmetric part is $\operatorname{Hess}\mathcal W$.

Write $U=\exp x$ linkwise, with $|x|=r$.
Along $\exp(\theta x)$ the radial invariant vector is $x$.
Integrating the upper Hessian bound along this geodesic,
and using $\nabla\mathcal W(I)=0$, gives
$\langle F(U),x\rangle\le b_1r^3/2$.
Consequently a nonzero radial distance along the actual flow
satisfies
\[
 r'(t)\le\tfrac12 b_1r(t)^2,\qquad
 r(t)\le\frac{r(0)}{1-b_1r(0)t/2}.
 \tag{V41.12}\label{c18v41:radial}
\]
A first-exit argument proves this inequality throughout $[0,T]$:
our choice gives $b_1\delta_\ell T\le1/2$, so the right side
is at most $4r(0)/3<2\delta_\ell\le1/4$.
There is no cut-locus issue in this ball. The zero initial
distance is the fixed identity trajectory.
The symmetric variational generator is bounded above by
$2b_1\delta_\ell I$, hence every forward propagator between
times $r$ and $t$ has norm at most
$e^{2b_1\delta_\ell(t-r)}\le e$.
No contraction of the inverse propagator is asserted.

For completeness let
$K_z(t)=D_z\mathsf J_t$ and
$L_{wz}(t)=D_wD_z\mathsf J_t$.
Differentiating the exact coefficient equation gives
\[
 \begin{aligned}
 \dot K_z={}&\mathsf M K_z+
                (D_{\mathsf J_tz}\mathsf M)\mathsf J_t,\\
 \dot L_{wz}={}&\mathsf M L_{wz}
 +(D_{\mathsf J_tw}\mathsf M)K_z
 +(D_{\mathsf J_tz}\mathsf M)K_w\\
 &+(D_{\mathsf J_tw}D_{\mathsf J_tz}\mathsf M)\mathsf J_t
 +(D_{K_wz}\mathsf M)\mathsf J_t .
 \end{aligned}
 \tag{V41.13}\label{c18v41:jets}
\]
In the iterated $D$ term its coefficient directions are frozen
while differentiating $\mathsf M$; their variation is precisely
the last term. Both jets start at zero. The frame commutators
are already inside $\mathsf M$.
Duhamel's formula and the forward propagator bound yield,
at $t\le T$,
\[
 \|K_z(t)\|\le b_1te^3|z|,\qquad
 \|L_{wz}(t)\|\le
       \left(\tfrac32b_1^2t^2e^5+b_2te^4\right)|w||z|.
 \tag{V41.14}\label{c18v41:jetbounds}
\]
There are three first-generator/first-jet terms in the
second equation; their integral is
$\frac32b_1^2t^2e^5$. The second-generator term costs
$b_2te^4$. V29's derivative-index contraction proof gives
both packed $K$ bounds with constant $(b_1te^3)^2$.

Now $C_s=\mathsf B(\Phi_sU)\mathsf J_s$.
Its first derivative is V29.13. The two packed bounds
follow by the same direct-sum triangle inequality, using
$a_{\rm B}e^2+\sqrt\ell k_T=p_1$.
Its second derivative has the five terms
\[
 \begin{aligned}
 D_wD_zC_s={}&
 (D_{\mathsf J_sw}D_{\mathsf J_sz}\mathsf B)\mathsf J_s
 +(D_{K_wz}\mathsf B)\mathsf J_s\\
 &+(D_{\mathsf J_sz}\mathsf B)K_w
 +(D_{\mathsf J_sw}\mathsf B)K_z+\mathsf B L_{wz}.
 \end{aligned}
 \tag{V41.15}\label{c18v41:Csecond}
\]
They give $p_2$ exactly as in \eqref{c18v41:constants}.

It remains to prove the full Gram floor, rather than a
transverse or selected-output compression.
At flatness, at coarse momentum $q\in[-\pi,\pi)^3$ take
the primary fine alias $k=q/\ell$. In unitary Fourier
normalization its sampler is diagonal with entries
$\ell^{-3/2}\sum_{j=0}^{\ell-1}e^{ijk_a}$.
Each magnitude is at least $2/(\pi\sqrt\ell)$, including
the limiting value when $q_a=0$.
The full heat matrix $e^{-sA(k)}$ has smallest eigenvalue
at least $e^{-s\lambda(k)}$, and
$\lambda(k)=4\sum_a\sin^2(k_a/2)\le3\pi^2/\ell^2$.
The primary block alone therefore has least singular value
at least $m_{\rm f}/\sqrt\ell$ for $s\le T$.
Summing the positive Grams of every alias gives
$G_s(I)\succeq m_{\rm f}^2\ell^{-1}I$.
This argument retains the fine primary harmonic at $q=0$.

Along the radial path from $I$ to $U$ the first derivative
bound proves $\|C_s(U)-C_s(I)\|\le p_1d(U,I)$.
Thus for every coarse $v$,
\[
 |C_s(U)^*v|
 \ge\left(\frac{m_{\rm f}}{\sqrt\ell}
                       -p_1\delta_\ell\right)|v|
 \ge\frac{m_{\rm f}}{2\sqrt\ell}|v|.
\]
This proves the last part of \eqref{c18v41:tubebounds}.
\end{proof}

The radius is polynomial rather than exponentially small
in block length. Indeed, with
$a_1=2\sqrt3e^2+b_1e^3$ and
$c_\delta=\min\{1/8,(2b_1)^{-1},m_{\rm f}/(2a_1)\}>0$,
\[
 p_1\le a_1\ell^{5/2},\qquad
 \delta_\ell\ge c_\delta\ell^{-3}.
 \tag{V41.16}\label{c18v41:radius}
\]
The constant $m_{\rm f}$ is very small; these are rigorous
existence bounds, not practical numerical radii.
The conclusions also hold on the union of gauge translates
of $\mathcal U_\ell$. Gauge actions are isometries, intertwine
the flow and sampler, and rotate all fixed invariant frame
indices by constant orthogonal matrices. This saturates a
neighborhood of the gauge orbit of identity; it does not
cover all flat toron sectors or assert that their complement
has small probability.

\subsection{Derivatives and the actual normal score on this neighborhood}

Define
\[
 K=\frac{2\sqrt\ell}{m_{\rm f}},\qquad
 K_1=2K^2p_1,\qquad
 K_2=10K^3p_1^2+2K^2p_2 .
 \tag{V41.17}\label{c18v41:Rconstants}
\]
These bounds concern the globally defined $R$ evaluated on the
neighborhood, not a lift with a sharp cutoff.

\begin{theorem}[Polynomial normal-lift jets and the paid trace cost]
\label{c18v41:normaljets}
On the neighborhood of Theorem \ref{c18v41:tube}, for all
$s\le\ell^2$,
\[
 \begin{gathered}
 \|R\|\le K,\qquad
 \sum_i\|(X_iR)v\|^2\le K_1^2|v|^2,\qquad
 \|D_wD_zR\|\le K_2|w||z| .
 \end{gathered}
 \tag{V41.18}\label{c18v41:Rjets}
\]
Let $D_{\rm f}=3|E_{\rm f}|=9N^3$.
For $0\le\beta\le\beta_*=3/(512e)$ put
$\mu_\beta=4R_\beta/e$ with $R_\beta=64e\beta$, as in V37.
Then the actual full normal score satisfies
\[
 \begin{aligned}
 |S^R_v(U)|
 &\le\left[\sqrt{D_{\rm f}}K_1
                 +2\mu_\beta\delta_\ell K\right]|v|,\\
 |\nabla_{\rm fib}r^R_v(U)|
 &\le\left[D_{\rm f}K_2
       +2\mu_\beta\{(1+\delta_\ell)K+
                                      \delta_\ell K_1\}\right]|v|.
 \end{aligned}
 \tag{V41.19}\label{c18v41:scorebound}
\]
The fibre gradient uses the original induced metric.
The same norm is the internally reduced quotient-gradient
norm. The explicit $D_{\rm f}$ factors have not been
replaced by volume-independent constants.
\end{theorem}
\begin{proof}
The Gram floor gives $\|R\|\le K$. Differentiate the full
right inverse without choosing an eigenbasis:
\[
 D_zR=P(D_zC)^*Z-R(D_zC)R .
 \tag{V41.20}\label{c18v41:Rfirst}
\]
Since $\|P\|\le1$, $\|Z\|\le K^2$, this gives
$\|D_zR\|\le2K^2p_1|z|$.
For the packed assertion use the two packed $C$ bounds in
\eqref{c18v41:tubebounds} and the direct-sum triangle
inequality in \eqref{c18v41:Rfirst}. Each summand costs
at most $K^2p_1|v|$.

Useful differentiated identities are
\[
 D_wP=-P(D_wC)^*R^*-R(D_wC)P,\qquad
 D_wZ=(D_wR)^*R+R^*(D_wR).
\]
They imply $\|D_wP\|\le2Kp_1|w|$ and
$\|D_wZ\|\le4K^3p_1|w|$.
Differentiating \eqref{c18v41:Rfirst} once more gives six
terms. In their displayed order their bounds are
\[
 2K^3p_1^2,\quad K^2p_2,\quad4K^3p_1^2,\quad
 2K^3p_1^2,\quad K^2p_2,\quad2K^3p_1^2,
\]
times $|w||z|$. Their sum is $K_2|w||z|$.
No commutation of derivatives or matrix factors was used.

For the score use the exact original-frame formula
\[
 S^R_v=\sum_i\langle e_i,(X_iR)v\rangle
                         +2\langle\nabla u,Rv\rangle .
\]
Cauchy--Schwarz and the packed bound pay the first contraction
by $\sqrt{D_{\rm f}}K_1|v|$.
V37 gives $\|\operatorname{Hess}u\|\le\mu_\beta$ everywhere.
Global conjugation invariance gives $\nabla u(I)=0$,
so integrating the covariant Hessian along the radial
geodesic yields $|\nabla u(U)|\le\mu_\beta\delta_\ell$.
This proves the first line of \eqref{c18v41:scorebound}.

For a unit constant coefficient vector $w$, the differentiated
divergence is
$\sum_i\langle e_i,D_wX_iR\,v\rangle$.
The second jet bound pays each of its $D_{\rm f}$ terms by
$K_2|v|$. This is the asserted, still extensive trace estimate.
For the ground term, the coefficient derivative of $\nabla u$
differs from its covariant derivative by the invariant-frame
connection. On unit-round $SU(2)$ that connection has bilinear
norm at most one; the product has the same bound. Hence
$\|D_w(\nabla u)\|\le\mu_\beta(1+\delta_\ell)|w|$.
The product rule bounds its contribution by
$2\mu_\beta\{(1+\delta_\ell)K+\delta_\ell K_1\}|v|$.
Taking the supremum over $w$ bounds the ambient scalar gradient.
The subtracted conditional mean is constant along the fibre,
and orthogonal projection cannot increase the norm.
Equivariance makes the scalar source invariant under the
internal gauge group, so its fibre gradient is horizontal
for that quotient. This proves every assertion.
\end{proof}

In particular $p_2=O(\ell^{9/2})$, $K_1=O(\ell^{7/2})$,
$K_2=O(\ell^{13/2})$, with constants independent of $n$.
At $\beta=0$ the bound for the normal score derivative is
$O(n^3\ell^{19/2})$ on this native neighborhood, including
flat identity at $s=\ell^2$. This polynomial bound contrasts
with V40's $e^{8\ell^2}\ell^{-3/2}$ lower bound for the
averaged lift at that same point. It proves that the latter
growth is not an unavoidable stationary derivative obstruction
for all full lifts. It does not prove a global polynomial bound
for $R$, a volume-uniform bound for its divergence, or a mass gap.

\subsection{The remaining quantitative task and verification scope}

The ball in \eqref{c18v41:constants} bounds the total
product distance, not an average distance per link. Its
volume-independent radius therefore does not make it a
thermodynamically large set. No lower bound on its actual
conditional weight is proved. The score $r^R$ is centered
using the entire fibre, not the part lying in the ball;
the local bound on $S^R$ is not silently promoted to a
local bound on $r^R$.

The candidate changes, but the required quantity is still
\[
 R^R_{\gamma\eta}(g)=
       \langle r^R_\gamma,\bar L_g^{-1}r^R_\eta\rangle_{\nu_g}.
 \tag{V41.21}\label{c18v41:inverse}
\]
Nothing above bounds this matrix. The next useful step is
to localize the normal-score trace or estimate this joint
inverse directly in the actual conditional law, with the
cost of the region outside the neighborhood included.
The minimum-lift identity in V26 and the change-of-lift
certificate in V27 remain in force. A pointwise saving
in lift norm cannot be substituted for their conditional
source transport cost. Neither finite-volume polynomial
growth nor a missing conditional weight pays the
thermodynamic and weak-coupling limits.

The diagnostic
\texttt{scripts/verify\_ym\_c18\_v41\_normal\_lift\_tube.py}
checks exact rational pseudoinverse first and second jets,
the full divergence formula on a nonlinear finite-dimensional
fixture, and the stationary excess identity. It separately
checks the actual native flat alias matrices with high
precision, retaining all fine and coarse components, and
checks the radius and polynomial majorants. The abstract
fixture is not a native vacuum. No nonlinear conditional
law, neighborhood probability, interacting vacuum or
Poisson inverse is numerically certified. The displayed
all-volume neighborhood and derivative bounds are proved
by the local incidence and differential inequalities above,
not by sampling or a proof assistant.

V41 replaces only V40. Removing this appendix and restoring
the prior title recovers V40 byte for byte. All other 52
sources are unchanged. Only \texttt{.tex} sources remain in
\texttt{latex\_notes}; build and inspection products are
outside it. The next companion checkpoint is V45, the next
broad literature checkpoint V50, and the archive/restart V100.
The required conditional inverse, weak-coupling scale
transport, physical return, interacting four-dimensional
continuum reconstruction and positive physical Yang--Mills
mass gap remain unproved.
% END CYCLE 18 VERSION 41 APPEND
% BEGIN CYCLE 18 VERSION 42 APPEND
\clearpage
% Keep the new contents entry compact without changing inherited text.
\addtocontents{toc}{\protect\setcounter{tocdepth}{1}}
\section[V42: spatially summed normal scores without a volume loss]{V42: spatially summed normal scores without a volume loss}
\label{c18v42:context}

V41 bounds the actual normal lift on a nonlinear neighborhood,
but contracts its second derivative by the entire fine dimension.
This appendix removes that volume loss on a slightly smaller,
explicit neighborhood through $s=\ell^2$. It keeps the full
normal Gram inverse and all coarse output columns. The price is
a very large polynomial in block length, not a conditional
Poincar\'e theorem or a claim that the neighborhood is typical.

V38 already uses three anchored tensor sums for a flow variation.
We extend that bookkeeping to spatial weights, all slots of the
second coefficient derivative, the full inverse Gram, and the
diagonal trace defining the normal score. V41's normal lift and
its coarea formula are inherited. The new quantitative target is
its full coarse-bank source derivative with a constant independent
of the fine and coarse volumes.

The previous turn was progress, and V41 plus all 52 other sources
were revalidated. For primary context, the introduction and
Section 9 of Benzi--Boito,
\href{https://arxiv.org/html/1307.7119v1}
{\emph{Decay properties for functions of matrices over
$C^*$-algebras}}, were inspected, as was the primary abstract of
Benzi--Simoncini,
\href{https://arxiv.org/abs/1501.07376}
{\emph{Decay bounds for functions of matrices with banded or
Kronecker structure}}. Their matrix-decay context is relevant;
our flowed Gram is not declared banded. The exponential-conjugation
argument below is proved directly from its weighted row and
column sums. No general novelty for that inverse-decay method
is claimed. This targeted check does not replace the scheduled
broad sweep at V50; the companion checkpoint remains V45.

\subsection{Weighted anchored sums and the trace that they control}

Retain $N=\ell n$, $n\ge5$, $\ell\ge2$, and the original fine
product metric. Every fine scalar index $(x,i,a)$ is assigned
its link-tail position $x$ on the fine periodic cubic lattice.
A coarse scalar index $(z,i,a)$ is assigned position $\ell z$
on that same lattice. Let $d$ be periodic Manhattan distance
between the assigned positions. Distinct orientations, colors,
or fine/coarse indices can have distance zero.

For a scalar-entry tensor $T$ with at least one slot define
\[
 \mathfrak a_\theta(T)
 =\max_{k}\sup_{i_k}
   \sum_{(i_j)_{j\ne k}}
       |T_{i_1\cdots i_r}|
           e^{\theta\operatorname{diam}(i_1,\ldots,i_r)},
 \qquad \theta\ge0.
 \tag{V42.1}\label{c18v42:anchor}
\]
The diameter is computed from $d$. For two slots this is the
larger weighted absolute row or column sum, denoted also by
$\|\cdot\|_{{\rm S},\theta}$; it bounds the $\ell^2$ operator
norm. We use scalar color coefficients here, not V38's smaller
block norms. Slot permutations preserve \eqref{c18v42:anchor}.

\begin{lemma}[Connected contraction and diagonal trace]
\label{c18v42:contract}
Contracting two tensors along one or more common index slots
has anchored norm at most the product of their anchored norms,
provided a free slot remains. Identifying and summing two slots
of one tensor also cannot increase its anchored norm if at least
one free slot remains.
\end{lemma}
\begin{proof}
Take absolute values first. When the tensors share an index,
the diameter of their combined external indices is at most the
sum of the two original diameters, by the triangle inequality
through that shared position. The exponential weight therefore
splits into the product of the two weights. Anchor a remaining
external slot in the first tensor. For fixed shared indices,
sum the other tensor using one shared slot as its anchor;
summing a subset of its unanchored slots is bounded by its
full anchored sum. Now sum the remaining slots of the first
tensor. This proves the product bound. If the chosen external
slot is in the second tensor, reverse the argument.
For a diagonal identification, the restricted sum is a subset
of the unrestricted nonnegative sum with the same external anchor,
and deleting slots cannot increase diameter.
\end{proof}

All networks below are connected and retain external indices.
A disconnected closed trace would not be covered by this lemma;
it could introduce a volume factor.
In particular, for $R=C^*(CC^*)^{-1}$ set
\[
 (\partial^2R)_{j,i;k,\alpha}=X_jX_iR_{k\alpha},
 \qquad
 \mathsf T^R_{j\alpha}=\sum_iX_jX_iR_{i\alpha}.
\]
The diagonal-trace part of the lemma gives exactly
\[
 \|\mathsf T^R\|_{{\rm S},\theta}
       \le\mathfrak a_\theta(\partial^2R).
 \tag{V42.2}\label{c18v42:trace}
\]
It bounds both sums of the remaining fine/coarse matrix.
This is stronger than estimating its individual summands
and multiplying by $D_{\rm f}$.

\subsection{A full flat heat identity, without deleting zero modes}

Let $A=d_1^*d_1\otimes I_3$ be the flat spatial curl Gram on
all fine links and colors. Let $\Delta$ be the scalar nearest
neighbor positive lattice Laplacian, repeated on link directions
and colors. In Fourier variables
\[
 A(k)=\lambda(k)I-d(k)d(k)^*,\qquad
 \Delta(k)=\lambda(k)I,\qquad A^2=\Delta A.
\]
Here the $d(k)=0$ case has $A(k)=0$.
Finite-dimensional functional calculus, including the zero
eigenspaces, gives
\[
 \boxed{e^{-tA}=I-\int_0^t A e^{-r\Delta}\,dr.}
 \tag{V42.3}\label{c18v42:heat}
\]
Thus longitudinal and harmonic inputs remain exactly unchanged;
there is no singular inverse of $\Delta$ in this formula.

The scalar kernel $e^{-t\Delta}$ is the continuous-time
nearest-neighbor random-walk kernel, with total jump rate six.
Each jump changes the assigned position by distance at most one.
The number of jumps is Poisson with mean $6t$, so
\[
 \|e^{-t\Delta}\|_{{\rm S},\theta}
       \le e^{6t(e^\theta-1)} .
 \tag{V42.4}\label{c18v42:walk}
\]
This proof applies on every periodic volume; repeated orientations
and colors are diagonal copies. The plaquette incidence gives
$\|A\|_{{\rm S},0}\le16$, and its nonzero entries have
tail distance at most two. Therefore
\[
 \|e^{-tA}\|_{{\rm S},\theta}
 \le1+16e^{2\theta}t\,e^{6t(e^\theta-1)} .
 \tag{V42.5}\label{c18v42:flatSchur}
\]
These are spatial differential matrices, not the conditional
diffusion or the physical quantum Hamiltonian.

\subsection{A polynomial spatial bound on the actual nonlinear flow}

Keep $T=\ell^2$, $b_1=288\sqrt3$, $m_{\rm f}$ and
$\delta_\ell$ from V41. Define
\[
 \begin{aligned}
 \theta_0&=\frac1{24(1+T)},&
 H_0&=1+16e^2T,& H&=eH_0,\\
 \rho_\ell&=\min\left\{\delta_\ell,
                    \frac1{2e b_1 H_0T}\right\},&
 \mathcal V_\ell&=\{U:d_{\rm prod}(U,I)<\rho_\ell\}.
 \end{aligned}
 \tag{V42.6}\label{c18v42:domain}
\]
In particular $\rho_\ell\ge c_\rho\ell^{-4}$ with
\[
 c_\rho=\min\{c_\delta,[2e b_1(1+16e^2)]^{-1}\}>0,
\]
where $c_\delta$ is V41.16. This is still a total product-distance
neighborhood, not a per-link typicality condition.

Use the actual invariant-frame generator $\mathsf M$ from V29
and V41. Its coefficient difference from $\mathsf M(I)=-A$
obeys, for $W$ in the radius-$1/4$ product ball,
\[
 \|\mathsf M(W)+A\|_{{\rm S},\theta_0}
             \le e b_1\,d_{\rm prod}(W,I).
 \tag{V42.7}\label{c18v42:perturb}
\]
To check the scalar row/column bound needed here, fix an output
index in a plaquette contribution to $D_z\mathsf M$. Its
entries have magnitude at most three, as in V29.
There are at most four incident plaquettes and twelve input
slots in each. Summing the derivative direction over twelve
slots gives
$4\cdot12\cdot3\sqrt{12}|z|=b_1|z|$.
Fixing an input index gives the same bound.
The spatial weight costs at most $e^{2\theta_0}\le e$.
Integrate this coefficient derivative along the radial product
geodesic to prove \eqref{c18v42:perturb}.

\begin{lemma}[Actual spatial variational propagators]
\label{c18v42:propagator}
For $U\in\mathcal V_\ell$ and $0\le r\le t\le T$, the forward
invariant-frame propagator $\mathsf V(t,r;U)$ obeys
\[
 \|\mathsf V(t,r;U)\|_{{\rm S},\theta_0}\le H .
 \tag{V42.8}\label{c18v42:propbound}
\]
\end{lemma}
\begin{proof}
V41's radial proof, applied with initial radius $\rho_\ell$,
keeps the trajectory within $2\rho_\ell$.
Thus $Q_t=\mathsf M(\Phi_tU)+A$ has weighted norm at most
$\epsilon=2e b_1\rho_\ell$.
Since $\theta_0<1$ and $e^{\theta_0}-1\le2\theta_0$,
$6T(e^{\theta_0}-1)\le1/2$.
Equation \eqref{c18v42:flatSchur} is consequently at most $H_0$
through time $T$.
The exact Duhamel equation is
\[
 \mathsf V(t,r)=e^{-(t-r)A}
       +\int_r^t e^{-(t-u)A}Q_u\mathsf V(u,r)\,du .
\]
The weighted Schur product inequality and Gronwall give
$\|\mathsf V(t,r)\|_{{\rm S},\theta_0}
\le H_0e^{\epsilon H_0(t-r)}\le eH_0$ by
\eqref{c18v42:domain}. No inverse propagator bound is used.
\end{proof}

Here and below all derivatives are ordinary ordered coefficient
derivatives. Their moving-frame terms are already in $\mathsf M$.
Introduce scalar anchored majorants
\[
 \begin{gathered}
 t_1=1728e,\qquad t_2=20736e,\qquad
 b_0=3e\ell,\quad b_{\rm s1}=18e\ell^2,\quad
 b_{\rm s2}=108e\ell^3,\\
 k=t_1TH^3,\qquad
 l=\tfrac32t_1^2T^2H^5+t_2TH^4,\\
 c_0=b_0H,\qquad
 c_1=b_{\rm s1}H^2+b_0k,\qquad
 c_2=b_{\rm s2}H^3+3b_{\rm s1}Hk+b_0l .
 \end{gathered}
 \tag{V42.9}\label{c18v42:jetconstants}
\]
These symbols are majorants, not couplings or operator eigenvalues.

\begin{proposition}[All-slot spatial sums for the full map]
\label{c18v42:mapjets}
For $C_s=D(b_\ell\circ\Phi_s)$, $U\in\mathcal V_\ell$ and
$0\le s\le T$,
\[
 \mathfrak a_{\theta_0}(C_s)\le c_0,\qquad
 \mathfrak a_{\theta_0}(\partial C_s)\le c_1,\qquad
 \mathfrak a_{\theta_0}(\partial^2 C_s)\le c_2 .
 \tag{V42.10}\label{c18v42:Cjets}
\]
\end{proposition}
\begin{proof}
For the three-index tensor $\partial\mathsf M$, fixing any
scalar index leaves at most four plaquettes, $12^2$ choices
for the other slots, and entry bound three. Its unweighted
anchored sum is at most $1728$. For
$\partial^2\mathsf M$, it is at most
$4\cdot12^3\cdot3=20736$, using V41's fourth ordered
plaquette derivative bound. All positions are in one plaquette,
so $e^{2\theta_0}\le e$ proves the $t_1,t_2$ bounds.
Unlike the trilinear operator norm in V41, this sums every
unanchored scalar index.

For the sampler $\mathsf B=Db_\ell$, each row has $\ell$
rotation blocks. Absolute scalar sums are at most $3\ell$.
First and second coefficient derivatives of a prefix have
scalar entries bounded by two and four. Fixing a coarse
output leaves at most $(3\ell)^2$ or $(3\ell)^3$ entries.
Fixing any fine slot gives no larger count, since that slot
belongs to at most one coarse chain. These give
$3\ell,18\ell^2,108\ell^3$.
All these positions lie in one length-$\ell$ chain, including
the coarse tail. Since $\theta_0\ell\le1$, its weight costs
at most $e$, proving the sampler bounds.

The first and second variational equations are precisely
V41.13, with zero initial derivative tensors. Apply
Lemma \ref{c18v42:contract} to their Duhamel formulas and use
\eqref{c18v42:propbound}. The first jet costs $t_1sH^3$.
The three first-generator/first-jet terms in the second
equation cost $\frac32t_1^2s^2H^5$ after integration;
the second-generator term costs $t_2sH^4$.
This proves the anchored bounds $k,l$ through time $T$.
Finally $C_s=\mathsf B(\Phi_sU)\mathsf J_s$.
Its first derivative has two connected terms, and its
ordered second derivative has V41.15's five connected terms.
Their products give exactly $c_0,c_1,c_2$.
All possible external anchors were retained at every step.
\end{proof}

\subsection{Spatially summing the full normal Gram inverse}

This subsection concerns the finite coarse matrix $G=CC^*$,
not the still unbounded conditional Poisson inverse.
V41's full-output lower bound remains valid on $\mathcal V_\ell$:
\[
 G\succeq g_0 I,\qquad
 g_0=\frac{m_{\rm f}^2}{4\ell},\qquad
 \|G\|_{{\rm S},\theta_0}\le g_{\rm S}:=c_0^2 .
\]
Define
\[
 \eta=\frac{\theta_0g_0}{4g_{\rm S}},\qquad
 \zeta=\frac\eta2,\qquad
 z_0=\frac{18}{g_0}\left(1+\frac4{\eta\ell}\right)^3 .
 \tag{V42.11}\label{c18v42:inverseconstants}
\]
Our explicit constants give $g_0\le g_{\rm S}$, hence
$\eta\le\theta_0/4$.

\begin{lemma}[Full inverse decay and its weighted sum]
\label{c18v42:inverse}
Every scalar entry of the full inverse $Z=G^{-1}$ satisfies
\[
 |Z_{\alpha\gamma}|\le\frac2{g_0}e^{-\eta d(\alpha,\gamma)},
 \qquad
 \|Z\|_{{\rm S},\zeta}\le z_0 .
 \tag{V42.12}\label{c18v42:Zbound}
\]
The bounds are independent of $n$.
\end{lemma}
\begin{proof}
Fix a coarse index $\gamma$ and conjugate by
$D_{\alpha\alpha}=e^{\eta d(\alpha,\gamma)}$.
The triangle inequality gives
\[
 |(DGD^{-1}-G)_{\alpha\beta}|
 \le |G_{\alpha\beta}|\,\eta d(\alpha,\beta)
                       e^{\eta d(\alpha,\beta)}.
\]
For $u\ge0$,
$u e^{-(\theta_0-\eta)u}\le(\theta_0-\eta)^{-1}
\le2/\theta_0$. Both weighted row/column bounds of $G$
therefore imply
\[
 \|DGD^{-1}-G\|_{2\to2}
 \le\frac{2\eta g_{\rm S}}{\theta_0}=\frac{g_0}{2}.
\]
The Neumann-series inverse around $G$ has norm at most $2/g_0$.
Its $(\alpha,\gamma)$ entry equals
$e^{\eta d(\alpha,\gamma)}Z_{\alpha\gamma}$, proving the
entry estimate. Every conjugation is finite at each regulator;
the resulting constants do not depend on its size.

There are nine coarse scalar indices per coarse vertex, and
coarse-tail distances are $\ell$ times periodic coarse
Manhattan distances. Choose nearest integer representatives
and bound their sum by the sum on $\mathbb Z^3$. Thus
\[
 \sup_\gamma\sum_\alpha e^{-\eta d(\alpha,\gamma)/2}
 \le9\left(1+\frac2{e^{\eta\ell/2}-1}\right)^3
 \le9\left(1+\frac4{\eta\ell}\right)^3 .
\]
Multiply by $2/g_0$. Symmetry gives the other sum, proving
\eqref{c18v42:Zbound}. No coarse direction has been projected out.
\end{proof}

\subsection{The normal-lift trace now has no volume factor}

All subsequent anchored norms use $\zeta\le\theta_0$.
Set
\[
 \begin{aligned}
 g_1&=2c_0c_1,&g_2&=2c_0c_2+2c_1^2,\\
 z_1&=z_0^2g_1,&z_2&=2z_0^3g_1^2+z_0^2g_2,\\
 r_0&=c_0z_0,&r_1&=c_1z_0+c_0z_1,&
 r_2&=c_2z_0+2c_1z_1+c_0z_2 .
 \end{aligned}
 \tag{V42.13}\label{c18v42:Rconstants}
\]
Then on the declared neighborhood
\[
 \mathfrak a_\zeta(\partial^jG)\le g_j,\quad
 \mathfrak a_\zeta(\partial^jZ)\le z_j\quad(j=1,2),
 \qquad
 \mathfrak a_\zeta(\partial^jR)\le r_j\quad(j=0,1,2).
 \tag{V42.14}\label{c18v42:Rjets}
\]
To verify these statements with no hidden commutation, put
$G_i=X_iG$, $G_{ji}=X_jX_iG$. The exact identities are
\[
 \begin{aligned}
 X_iZ&=-ZG_iZ,\\
 X_jX_iZ&=ZG_jZG_iZ+ZG_iZG_jZ-ZG_{ji}Z,\\
 X_jX_iR&=(X_jX_iC)^*Z+(X_iC)^*(X_jZ)
                  +(X_jC)^*(X_iZ)+C^*(X_jX_iZ).
 \end{aligned}
 \tag{V42.15}\label{c18v42:inversejets}
\]
Differentiate $G=CC^*$ for the two $g_j$ bounds.
Every term in these identities is a connected contraction,
so Lemma \ref{c18v42:contract}, \eqref{c18v42:Cjets} and
\eqref{c18v42:Zbound} give \eqref{c18v42:Rjets}.
In particular \eqref{c18v42:trace} now gives
\[
 \boxed{\left\|\left(\sum_iX_jX_iR_{i\alpha}\right)_
                                  {j,\alpha}\right\|_{2\to2}
                         \le r_2.}
 \tag{V42.16}\label{c18v42:paidtrace}
\]
This is the previously unpaid trace contraction, not merely
one more bound on individual coefficient derivatives.

Let $S^R_v,r^R_v$ be exactly V41.2's scores for this normal lift
and the actual ground-state density. Keep the V41 constants
$K,K_1$ and $\mu_\beta=4R_\beta/e$, for
$0\le\beta\le3/(512e)$. Define
\[
 L_{\ell,\beta}
 =r_2+2\mu_\beta\bigl[(1+\rho_\ell)K+\rho_\ell K_1\bigr].
 \tag{V42.17}\label{c18v42:L}
\]

\begin{theorem}[Volume-independent native normal-score derivative]
\label{c18v42:score}
For every $n\ge5$, $\ell\ge2$, $0\le s\le\ell^2$,
$U\in\mathcal V_\ell$, and every full coarse coefficient vector
$v$,
\[
 \boxed{\begin{aligned}
 |\nabla S^R_v(U)|&\le L_{\ell,\beta}|v|,\\
 |\nabla_{\rm fib}r^R_v(U)|&\le L_{\ell,\beta}|v|,\\
 |S^R_v(U)|&\le\rho_\ell L_{\ell,\beta}|v| .
 \end{aligned}}
 \tag{V42.18}\label{c18v42:scorebounds}
\]
The constants are independent of $n$. The gradient and
fibre-gradient norms are the original ones, and the latter
equals the norm after internal-gauge reduction.
\end{theorem}
\begin{proof}
The exact invariant-frame expression is
\[
 S^R_v=\sum_iX_i(Rv)_i+2\langle\nabla u,Rv\rangle.
\]
The derivative of its first term is the fine/coarse matrix
in \eqref{c18v42:paidtrace}, contracted with $v$; Schur's
bound gives $r_2|v|$ without a dimension factor.
As in V41, $\|\operatorname{Hess}u\|\le\mu_\beta$ and
$\nabla u(I)=0$ imply $|\nabla u(U)|\le\mu_\beta\rho_\ell$.
The invariant-frame connection costs at most
$|\nabla u(U)|$ in the coefficient derivative of this gradient.
Together with the inherited $\|R\|\le K$ and
$\|D_wR\|\le K_1|w|$, the derivative of the second term
costs at most
$2\mu_\beta[(1+\rho_\ell)K+\rho_\ell K_1]|v|$.
This proves the ambient gradient bound.

The subtracted conditional mean is constant along the entire
fibre, so its tangential derivative is zero. Projecting the
ambient scalar gradient orthogonally proves the second bound.
The actual measure, lift and map are equivariant; the source
is internally gauge invariant. Its gradient is therefore
horizontal for the internal quotient, giving the asserted norm.

Finally the coarse score vector $S^R(I)$ vanishes.
Indeed simultaneous global conjugation fixes $I$ and rotates
each coarse color block by the adjoint representation.
Equivariance of the normal lift and invariance of the actual
ground density make each block of $S^R(I)$ invariant under
all such rotations, hence zero. Integrate the ambient gradient
bound along the radial product geodesic from $I$ to $U$.
It stays in $\mathcal V_\ell$ and has length at most $\rho_\ell$,
proving the last assertion.
\end{proof}

The final operator-norm and scalar conclusions also hold on the
union of gauge translates of $\mathcal V_\ell$, by the isometries
and the corresponding orthogonal coarse-frame changes.
The intermediate scalar-entry weighted sums were proved in the
identity-centered frame; no invariance of an entrywise absolute
norm under color rotations is assumed.

The displayed constants are polynomial in $\ell$, although
very large. Directly from their formulas,
\[
 \begin{gathered}
 H=O(\ell^2),\quad c_0=O(\ell^3),\quad
 c_1=O(\ell^9),\quad c_2=O(\ell^{15}),\\
 \eta^{-1}=O(\ell^9),\quad z_0=O(\ell^{25}),\quad
 g_1=O(\ell^{12}),\quad g_2=O(\ell^{18}),\\
 z_1=O(\ell^{62}),\quad z_2=O(\ell^{99}),\quad
 r_2=O(\ell^{102}),\qquad
 L_{\ell,\beta}=O(\ell^{102}) .
\end{gathered}
 \tag{V42.19}\label{c18v42:powers}
\]
The last bound is uniform over the stated small-electric-coupling
interval. These powers and constants are not claimed to be
sharp or practically useful. The point is that neither a factor
of $n$ nor an exponential in $\ell^2$ is needed for this local
full-bank trace bound.
For scale, the displayed formulas give
$\log_{10}r_2\simeq476.07$ at $\ell=2$.
This is a conservative majorant, not a computed native score.

\clearpage
\subsection{The conditional-weight and inverse obligations are unchanged}

Theorem \ref{c18v42:score} replaces V41's extensive local
derivative bound by a volume-independent one. It does not turn
the total-distance neighborhood into a high-probability set.
No estimate of its actual conditional weight or of the
source on its complement has been obtained. In particular
the last bound in \eqref{c18v42:scorebounds} is for the raw
score $S^R$, not the globally centered $r^R$.

The full conditional inverse remains
\[
 \langle r^R_v,\bar L_g^{-1}r^R_v\rangle_{\nu_g}
   =\langle dr^R_v,L_{1,g}^{-2}dr^R_v\rangle_{\nu_g}.
 \tag{V42.20}\label{c18v42:target}
\]
The inverse Gram $Z$ estimated above is not either inverse
operator in this identity. A global polynomial source-gradient
bound, even if it were proved, would still need an appropriate
conditional inverse or joint spectral estimate. Here only
a precisely declared neighborhood is covered.

The next useful step is to replace the total-distance criterion
by a spatially local, native configuration criterion with its
source-weighted exceptional cost paid, or estimate the joint
inverse directly including the complement. Another refinement
of the polynomial degree alone would not settle this issue.
Weak-coupling transport, the physical return, cofinal scale
calibration, interacting four-dimensional reconstruction and
the positive physical Yang--Mills mass gap remain unproved.

The diagnostic
\texttt{scripts/verify\_ym\_c18\_v42\_spatial\_normal\_score.py}
separately checks weighted anchored contractions, inverse jets
and their diagonal trace on finite matrix fixtures; the exact
native cubic curl/Laplacian identity; native flat heat kernels
without harmonic or longitudinal removal; and all stated
constant recurrences. The fixtures do not simulate a YM vacuum.
The prior V41 mathematical regression is replayed.
These checks support the algebra and normalizations; the
uniform theorem is the written proof above, not sampling or
proof-assistant certification.

V42 replaces only V41. Removing this appendix and restoring
the title recovers V41 byte for byte. The other 52 note
sources are unchanged. The \texttt{latex\_notes} folder
contains only \texttt{.tex} files, with all builds and QA
outside it. The next companion checkpoint is V45, the next
broad literature checkpoint V50, and archive/restart remains V100.
\addtocontents{toc}{\protect\setcounter{tocdepth}{2}}
% END CYCLE 18 VERSION 42 APPEND
% BEGIN CYCLE 18 VERSION 43 APPEND
\clearpage
\addtocontents{toc}{\protect\setcounter{tocdepth}{1}}
\section[V43: conditional anti-concentration of the flat neighborhood]{V43: conditional anti-concentration of the flat neighborhood}
\label{c18v43:context}

V42 removes the volume factor in a local normal-score derivative
bound, but leaves the weight of its neighborhood unpaid. We now
test that weight in the actual conditional law. The result changes
the next step: at fixed block length and coupling, even the full
gauge saturation of that neighborhood has superexponentially small
conditional probability as the volume grows. Enlarging it to a
sub-root-volume total-distance tube does not fix this issue.

This is an upper probability bound, not a source-weighted inverse
estimate. It neither contradicts V42 nor disproves a Yang--Mills gap.
V19's local, weak-coupling Wilson-loop metric estimate concerns a
different event and a different direction; it is not being redone.
V32's finite-replacement density bound and V33's exact forest
disintegration are inherited. The new calculation combines them
with many disjoint elementary holonomies and the total-distance
constraint. It also gives a conditional defect-count estimate
without declaring the interacting chords independent.

V42 and the 52 other sources were revalidated. A targeted archive
search found earlier small-ball estimates for covariant spatial
spectra, notably f16.2/V74, but not this conditional
configuration-tube conclusion. This is not a claim of general
literature novelty. The primary abstract of Ligterink--Walet--Bishop,
\href{https://arxiv.org/abs/hep-lat/0001028}
{\emph{Towards a Many-Body Treatment of Hamiltonian Lattice
SU(N) Gauge Theory}}, was rechecked for maximal-tree context only.
No spectral theorem is imported from it. The next companion and
broad-literature checkpoints remain V45 and V50.

\subsection{The actual conditional law and two density costs}

Use the same $SU(2)$ round metric, fine side $N=\ell n$,
$\ell\ge2$, $n\ge5$, and actual ground law $\nu=\Omega_\beta^2dU$.
In this appendix $\beta\ge0$ is finite and $s\ge0$ is finite.
The smaller V42 coupling range will only be needed when applying
its derivative theorem. Put $V=\Phi_sU$ and
$c_s=\mathfrak b_\ell\circ\Phi_s$.
The reduced conditional law at $c_s(U)=g$ is exactly
\[
 p_{g,\beta,s}(y)\,dy
 =\frac{h_{\beta,s}(W^0(g,y))}
        {\int h_{\beta,s}(W^0(g,z))\,dz}\,dy,
 \quad
 h_{\beta,s}(V)=\Omega_\beta(\Phi_{-s}V)^2J_{-s}(V).
 \tag{V43.1}\label{c18v43:law}
\]
The inverse-flow Jacobian satisfies
$\log J_{-s}=12\int_0^s W(\Phi_{-t}V)\,dt$.
Gauge coordinates have normalized Haar law and integrate out
on the events considered below. No extra metric determinant
is inserted into \eqref{c18v43:law}.

Recall V32's constants
\[
 \begin{gathered}
 a_\beta=\min\{2\beta,\pi\sqrt{8\beta}\},\\
 F(0)=1,\qquad
 F(t)=\min\{e^{16t},\,3(80t+1)^3+1\}\quad(t>0),\\
 D_s=\pi\left(2a_\beta F(s)+48\int_0^sF(t)\,dt\right),
 \qquad Q_s=6\pi a_\beta+72s .
 \end{gathered}
 \tag{V43.2}\label{c18v43:costs}
\]
The letter $a_\beta$ here is a scalar score bound, not the
reduced kinetic matrix $a_{g,s}$.

For any set $J$ of $k$ chord coordinates, its conditional
density given $g$ and all other chords obeys
\[
 e^{-kD_s}\le p_{J\mid J^c,g}(y_J)\le e^{kD_s}.
 \tag{V43.3}\label{c18v43:localdensity}
\]
Indeed changing one chord changes one link of $W^0$; V32 bounds
the log-density change by $D_s$. Telescope $k$ replacements.
A positive function of log oscillation at most $kD_s$, after
normalization against probability Haar measure, lies between
$e^{-kD_s}$ and $e^{kD_s}$. The normalizer is retained in this
argument. In particular every one-chord conditional has the
same lower bound $e^{-D_s}$, uniformly in its exterior.

A second, extensive bound is useful when all the chords are
integrated:
\[
        e^{-Q_sN^3}\le p_{g,\beta,s}(y)\le e^{Q_sN^3}.
 \tag{V43.4}\label{c18v43:globaldensity}
\]
To verify it, $|\nabla_e\log\Omega_\beta|\le a_\beta$ gives
oscillation at most $2\pi a_\beta|E_{\rm f}|$ for
$2\log\Omega_\beta$, including after composition with $\Phi_{-s}$.
Since $|E_{\rm f}|=|\mathcal P_N|=3N^3$ and each $w_p\in[-1,1]$,
the Jacobian log oscillation is at most $24s|\mathcal P_N|$.
Their sum is $Q_sN^3$. Restriction to the section and subsequent
normalization cannot increase this oscillation.
This bound is extensive and is not a uniform comparison of
the full density with Haar.

\subsection{Many disjoint holonomies with exact Haar reference}

Two elementary plaquettes are adjacent here if they share a
stored fine link. Each plaquette has at most twelve neighbors:
four links, with at most three other plaquettes through each.
A greedy independent set therefore provides
\[
 \mathcal P_*=\{p_1,\ldots,p_K\},\qquad
 K\ge\left\lceil\frac{3N^3}{13}\right\rceil ,
 \tag{V43.5}\label{c18v43:packing}
\]
with pairwise disjoint link supports. The set is fixed
independently of the configuration and coupling.

Every elementary plaquette contains a free chord of V33's
forest. For completeness, $\mathcal F\cup\mathcal B$ consists
of the subdivided coarse cubic graph with trees attached:
a forest extension never joins two previously assigned
components. Its shortest cycle has length $4\ell>4$.
Thus a four-edge elementary cycle cannot be contained in
$\mathcal F\cup\mathcal B$.
Choose one chord $f_i$ of each $p_i$ and set $J=\{f_i\}_{i=1}^K$.
These selected chords are distinct.

Let $H_i(V)$ be the oriented plaquette product and
$\theta_i(V)=d_{S^3}(H_i(V),I)\in[0,\pi]$.
After fixing $g$ and all unselected chords, each
\[
                  H_i=A_i y_{f_i}^{\sigma_i}B_i,
             \qquad \sigma_i\in\{1,-1\},
 \tag{V43.6}\label{c18v43:holonomy}
\]
has fixed $A_i,B_i$ and depends on no other selected chord.
Left/right multiplication and inversion preserve Haar measure.
Consequently under the \emph{reference} product Haar measure on
$y_J$, the $H_i$ are independent Haar $SU(2)$ variables.
The actual conditional density is still present in
\eqref{c18v43:localdensity}; no independence under $p_g$ follows.

The one-holonomy cap probability is
\[
 b(r)=\frac2\pi\int_0^r\sin^2t\,dt
      =\frac{r-\frac12\sin(2r)}{\pi},\quad 0\le r\le\pi .
 \tag{V43.7}\label{c18v43:cap}
\]
A stronger bound for the joint radial constraint is
\[
 \mathbb P_{\rm Haar}\left\{\sum_{i=1}^K\theta_i^2\le R^2\right\}
 \le \frac{\bigl(R^3/(2\sqrt\pi)\bigr)^K}
             {\Gamma(3K/2+1)} .
 \tag{V43.8}\label{c18v43:ball}
\]
In exponential coordinates $x\in\mathbb R^3$, $|x|<\pi$,
normalized Haar measure is
$(2\pi^2)^{-1}(\sin|x|/|x|)^2\,dx$.
Its density is at most $(2\pi^2)^{-1}$.
The cut locus is Haar null. Enlarge the resulting integration
domain to the Euclidean $3K$-ball of radius $R$ and use its
volume $\pi^{3K/2}R^{3K}/\Gamma(3K/2+1)$.
This proves \eqref{c18v43:ball}, also when its right side exceeds one.

Combining the two density bounds separately with the same
reference calculation yields, uniformly in $g$,
\[
 \begin{split}
 &\nu_g\left\{\sum_i\theta_i(\Phi_sU)^2\le R^2\right\}\\
 &\qquad\le
 \min\left\{1,\,
  e^{A_{K,N,s}}
  \frac{\bigl(R^3/(2\sqrt\pi)\bigr)^K}
       {\Gamma(3K/2+1)}\right\},\\
 &A_{K,N,s}=\min\{KD_s,Q_sN^3\}.
 \end{split}
 \tag{V43.9}\label{c18v43:radialprob}
\]
For the first bound condition on unselected chords before
integrating; for the second use \eqref{c18v43:globaldensity}
and integrate all chords. The Haar radial probability in
both cases is independent of the frozen exterior and of $g$.
This justifies taking the smaller exponent. It does not
claim the global bound conditional on every exterior.

\subsection{The gauge-saturated total-distance tube is rare}

Write $\mathcal G=SU(2)^{V_{\rm f}}$ for the full fine gauge
group and define distance to the pure-gauge orbit by
\[
 d_{\mathcal G}(U)=\inf_{h\in\mathcal G}d_{\rm prod}(h\cdot U,I),
 \qquad \mathcal T(r)=\{U:d_{\mathcal G}(U)<r\}.
 \tag{V43.10}\label{c18v43:tube}
\]
The infimum is attained since $\mathcal G$ is compact.
This includes every gauge translate of the identity-centered
ball. It is not the set of all flat toron sectors.

For any configuration $V$,
\[
          \sum_{i=1}^K\theta_i(V)^2\le4d_{\mathcal G}(V)^2.
 \tag{V43.11}\label{c18v43:geometry}
\]
For a fixed gauge, the bi-invariant triangle inequality bounds
a plaquette angle by the sum of its four link distances to $I$.
Square and use Cauchy--Schwarz with four terms. The supports
are disjoint, so summation uses each link at most once.
Angles are invariant under conjugation of the plaquette
holonomy. Taking the gauge infimum proves the assertion.

The full spatial-flow differential has operator norm at most
$e^{16s}$ by the existing all-configuration bound.
Gauge equivariance and $\Phi_s(I)=I$ therefore give
\[
        d_{\mathcal G}(\Phi_sU)\le e^{16s}d_{\mathcal G}(U).
 \tag{V43.12}\label{c18v43:flowdistance}
\]
Apply the Lipschitz bound along a minimizing product geodesic
after choosing a minimizing gauge. No path in that argument
is required to stay in one coarse fibre.

\begin{theorem}[Conditional tube anti-concentration]
\label{c18v43:antic}
For every finite $\beta,s$, every $g$, and every $r>0$,
\[
 \nu_g(\mathcal T(r))
 \le\min\left\{1,\,
 e^{A_{K,N,s}}
 \frac{\bigl((2e^{16s}r)^3/(2\sqrt\pi)\bigr)^K}
      {\Gamma(3K/2+1)}\right\}.
 \tag{V43.13}\label{c18v43:general}
\]
For V42's radius $\rho_\ell$ and $0\le s\le\ell^2$, the sharper
bound is
\[
 \boxed{\displaystyle
 \nu_g(\mathcal T(\rho_\ell))
 \le\min\left\{1,\,
 e^{A_{K,N,s}}
 \frac{\bigl((4\rho_\ell)^3/(2\sqrt\pi)\bigr)^K}
      {\Gamma(3K/2+1)}\right\}.}
 \tag{V43.14}\label{c18v43:v42}
\]
Here the tube is the full gauge saturation of the V42 domain.
The probability bounds alone hold for every finite $\beta$;
the score estimate on that tube still requires V42's
small-electric-coupling interval.
\end{theorem}
\begin{proof}
Equations \eqref{c18v43:geometry}--\eqref{c18v43:flowdistance}
place $\mathcal T(r)$ inside the event in
\eqref{c18v43:radialprob} with $R=2e^{16s}r$.
This proves the first statement.
For the second, V41--V42's nonlinear stability proof keeps
each trajectory starting in its ball within $2\rho_\ell$
of identity through time $\ell^2$.
Equivariance applies the same assertion on each gauge translate.
Thus $d_{\mathcal G}(\Phi_sU)\le2\rho_\ell$ on the tube.
Use $R=4\rho_\ell$ in \eqref{c18v43:radialprob}.
\end{proof}

The estimate includes all gauge translates at once: no union
bound over gauge transformations, or unbounded gauge volume,
has been used. It is uniform even in coarse values for which
the tube does not meet the fibre.

For fixed $\ell,\beta$ it implies, uniformly in $g$ and
$0\le s\le\ell^2$,
\[
 \log\nu_g(\mathcal T(\rho_\ell))
       \le-\tfrac32K\log K+O_{\ell,\beta}(K)
       \quad(n\longrightarrow\infty).
 \tag{V43.15}\label{c18v43:rate}
\]
If the probability is zero its logarithm is $-\infty$.
One elementary explicit version is obtained by setting
\[
 \begin{split}
 Q_T&=6\pi a_\beta+72\ell^2,\\
 B_{\ell,\beta}
   &=e^{13Q_T/3}\frac{(4\rho_\ell)^3}{2\sqrt\pi}
                     \left(\frac{2e}{3}\right)^{3/2}.
 \end{split}
\]
Then $\nu_g(\mathcal T(\rho_\ell))
\le\min\{1,e(B_{\ell,\beta}/K^{3/2})^K\}$.
Indeed $N^3\le13K/3$, and integration over $[z,z+1]$ gives
$\Gamma(z+1)\ge e^{-1}(z/e)^z$ for $z>0$.
The constants can make the upper bound vacuous at accessible
finite volumes; no useful onset volume is asserted.
For example, using the guaranteed packing size at
$\ell=2$, $n=5$, $\beta=\beta_*$ and $s=4$, the uncapped
log upper bound in \eqref{c18v43:v42} is about $257928$,
so the displayed probability bound there is only one.

More generally at fixed $\ell,\beta,s$, if
$r_N/\sqrt{N^3}\to0$, \eqref{c18v43:general} gives
$\nu_g(\mathcal T(r_N))\to0$ uniformly in $g$.
After the same gamma estimate its log upper bound per $K$
is $3\log(r_N/\sqrt K)+O_{\ell,\beta,s}(1)$.
This conclusion is not uniform in a coupling or flow time
that changes with $N$.

\subsection{A defect-count bound without interacting independence}

A total-distance tube is not the only possible local criterion.
To identify what an all-plaquette small-field criterion would cost,
fix $0<r<\pi$ and define
\[
 Z_r(U)=\sum_{i=1}^K
              \mathbf1_{\{\theta_i(\Phi_sU)>r\}},\qquad
 \alpha_r=e^{-D_s}[1-b(r)]>0 .
 \tag{V43.16}\label{c18v43:defects}
\]
These are defects relative to this chosen cap, not a claim
that such plaquettes have an unstable physical mode.

\begin{theorem}[Native conditional lower-tail estimate]
\label{c18v43:count}
For every $t\ge0$, uniformly in $g$,
\[
 \begin{aligned}
 \mathbb E_{\nu_g}e^{-tZ_r}
       &\le(1-\alpha_r+\alpha_re^{-t})^K,\\
 \mathbb E_{\nu_g}Z_r&\ge\alpha_rK,\\
 \nu_g\{Z_r\le\alpha_rK/2\}&\le e^{-\alpha_rK/8},\\
 \nu_g\{Z_r=0\}&\le(1-\alpha_r)^K .
 \end{aligned}
 \tag{V43.17}\label{c18v43:countbounds}
\]
\end{theorem}
\begin{proof}
Condition on all chords except $y_{f_i}$. The cap complement
has Haar measure $1-b(r)$ by \eqref{c18v43:holonomy}.
The one-coordinate density lower bound therefore makes its
conditional probability at least $\alpha_r$.
Every other selected defect indicator is measurable in that
exterior, because the selected plaquettes have disjoint supports.
Writing $I_i$ for this indicator gives
$\mathbb E[e^{-tI_i}\mid y_{f_i}^c,g]
\le1-\alpha_r+\alpha_re^{-t}$.
Remove these factors successively by the tower property.
This proves the exponential-moment bound without factoring
the law. Summing the individual probability floors gives
the expectation statement. Taking $t\to\infty$ proves the
zero-defect bound by bounded convergence.

For the third statement Markov's inequality gives
\[
 \nu_g\{Z_r\le\alpha_rK/2\}
 \le\exp\{t\alpha_rK/2-\alpha_rK(1-e^{-t})\}.
\]
Choose $t=1/2$ and use $1-e^{-1/2}>3/8$.
The exponent is at most $-\alpha_rK/8$.
\end{proof}

This also yields a geometric consequence on the original
unflowed configurations:
\[
 \begin{aligned}
 \mathbb E_{\nu_g}d_{\mathcal G}(U)^2
     &\ge\frac{e^{-32s}r^2\alpha_rK}{4},\\
 \nu_g\left\{d_{\mathcal G}(U)<
             e^{-16s}r\sqrt{\alpha_rK/8}\right\}
     &\le e^{-\alpha_rK/8}.
 \end{aligned}
 \tag{V43.18}\label{c18v43:width}
\]
Indeed $r^2Z_r\le\sum_i\theta_i^2
\le4e^{32s}d_{\mathcal G}(U)^2$.
Thus distance from the pure-gauge orbit is at least of
root-volume order with a positive, possibly extremely small,
parameter-dependent constant.

In fact the general tube and width bounds extend to the
entire elementary-flat set
\[
 \mathcal Z_{\rm flat}=\{Z:H_p(Z)=I\text{ for every }p\},
 \qquad d_{\rm flat}(U)=\inf_{Z\in\mathcal Z_{\rm flat}}
                                       d_{\rm prod}(U,Z).
 \tag{V43.19}\label{c18v43:allflat}
\]
This compact set includes nontrivial flat toron sectors.
For each $Z$ in it, bi-invariant telescoping gives
$\theta_i(V)\le\sum_{e\in p_i}d(V_e,Z_e)$, since $H_i(Z)=I$.
Thus $\sum_i\theta_i^2\le4d_{\rm flat}(V)^2$.
Every elementary-flat $Z$ is stationary for $\Phi_s$,
because each plaquette trace is at its maximum there.
The same Lipschitz proof gives
$d_{\rm flat}(\Phi_sU)\le e^{16s}d_{\rm flat}(U)$.
Consequently \eqref{c18v43:general} and \eqref{c18v43:width}
hold with distance to this whole set in place of
$d_{\mathcal G}$. The sub-root-volume conclusion does too.
No union bound over flat sectors is needed.
The sharper $4\rho_\ell$ radius in \eqref{c18v43:v42} and
the V42 derivative theorem still apply only to the stated
pure-gauge tube; they have not been proved on all toron sectors.

At $\beta=s=0$ the reference law is the exact conditional law.
The selected holonomies are then independent Haar even after
conditioning on $g$; hence $Z_r$ is exactly binomial with
parameters $K$ and $1-b(r)$.
For $r=\pi/2$ half of the selected plaquettes are outside
the cap on average. For positive $\beta$ or $s$ the proofs
above retain their correlated conditional density instead.
The constant $\alpha_r$ can be extremely small at $s=\ell^2$;
the theorem asserts no efficient numerical defect estimate
at that time.

\clearpage
\subsection{What this changes, and what it does not pay}

The neighborhood in V42 cannot be promoted to a typical set
in the thermodynamic limit at fixed parameters, even after
full gauge saturation. Its complement has conditional
probability tending to one. Replacing the total-distance
condition by the requirement that every selected plaquette
lie in a fixed proper cap also has vanishing probability.
A local method must allow defects and pay their contributions,
not demand a globally defect-free configuration.

These are unweighted probability statements. They do not
bound the normal score on the complement, the conditional
mean used in centering, or the spectral weight in
$\langle r^R,\bar L_g^{-1}r^R\rangle_{\nu_g}$.
Conversely a positive density of the cap defects does not
imply a small conditional gap: at $\beta=s=0$ the reduced
conditional gap is already paid in V37.
There is no spectral inference from this counting argument.

The next upstream target is therefore a source-weighted,
spatially local treatment with defects allowed, or a direct
conditional inverse estimate that does not rely on the
single flat tube. V19's weak-coupling local geometry may
supply a starting local chart, but its one-loop-covector
conclusion must not be upgraded to full-bank score control.
The parameter dependence in \eqref{c18v43:costs} and
\eqref{c18v43:defects} must be tracked before any simultaneous
coupling, block-size and volume limit is claimed.
The full interacting four-dimensional continuum, physical
return and positive physical Yang--Mills mass gap remain unproved.

The diagnostic
\texttt{scripts/verify\_ym\_c18\_v43\_conditional\_tube\_weight.py}
checks native forest and plaquette packing incidence, quaternion
holonomy isometries and the geometric inequality, exact cap
normalization, the radial gamma majorant, and the defect
moment argument on a separately labelled correlated finite
fixture. It also evaluates costs in logarithms and replays
V42's mathematics. It does not sample the native interacting
vacuum or certify an inverse. The uniform statements are
the written proofs, not a proof-assistant certificate.

During verification the companion ESM programme was externally
restarted: V80 is now V80\_f1 with the same byte hash, and
a new V14 appeared. Its restart introduction and latest
audit/status discussion were inspected, not all its proofs.
Its controlled projection-component and complement-cancellation
criteria do not supply the native conditional inverse required
here. No ESM theorem is imported, and the V45 companion
checkpoint is unchanged.

V43 replaces only V42. Removing this appendix and restoring
the title recovers V42 byte for byte. The contents of the
other 52 original sources are preserved: 51 remain at their
original paths and ESM V80 is the hash-identical external
archive just described. The new ESM V14 is also untouched.
Only \texttt{.tex} sources remain in
\texttt{latex\_notes}. All build and QA products stay outside.
Companion review remains due at V45, the broad sweep at V50,
and archive/restart at V100.
\addtocontents{toc}{\protect\setcounter{tocdepth}{2}}
% END CYCLE 18 VERSION 43 APPEND
% BEGIN CYCLE 18 VERSION 44 APPEND
\clearpage
\addtocontents{toc}{\protect\setcounter{tocdepth}{1}}
\section[V44: global normal-score gradients and the low spectrum]{V44: global normal-score gradients and the low spectrum}
\label{c18v44:context}

V43 proves that V42's total-distance flat neighborhood is rare
at fixed block parameters as volume increases. We therefore do
not further refine its polynomial constants as a route to
typical configurations. Here we remove the neighborhood
restriction altogether, at the explicit price of exponential
block-time constants. The new result is an all-configuration,
all-output, volume-independent bound on the derivative of the
actual normal score. It gives a uniform first spectral moment
and pays the high-frequency part of its conditional inverse.
It does not pay the remaining low-frequency part.

V30 already proves a full-bank mean-square estimate for its
weighted lift; we do not present another averaged variance
bound as new. V42 already supplies the anchored tensor algebra,
inverse-decay lemma and normal-lift identities. The work here
is to supply their global inputs, contract the actual vacuum
gradient without an extensive norm, and apply the resulting
estimate on every entire native fibre. No gauge, longitudinal,
harmonic, toron or coarse-output sector is discarded.
The vacuum estimates still require the small-electric-coupling
interval of V37, not the continuum weak-coupling regime.

The V43 source and other 53 note sources were rechecked.
The primary matrix-decay context in the introduction and
Section 9 of
\href{https://arxiv.org/html/1307.7119v1}{Benzi--Boito,
\emph{Decay properties for functions of matrices over
$C^*$-algebras}} was revisited. The argument below uses
V42's proved weighted conjugation, not a claim that the
flowed Gram is banded. This is a targeted check; the next
broad literature checkpoint remains V50.

\subsection{Global spatial flow and map derivatives}

Keep the native spatial torus of side $N=\ell n$, $n\ge5$,
$\ell\ge2$, and all scalar link/color indices in their original
invariant frames. Retain the tail positions, periodic Manhattan
distance, and all-slot norm $\mathfrak a_\theta$ of
\eqref{c18v42:anchor}. For a two-slot matrix this is the larger
weighted absolute row or column sum $\|\cdot\|_{{\rm S},\theta}$.
In this appendix define new majorants
\[
 T=\ell^2,\qquad \theta=\frac1{2\ell},\qquad
 \omega=4+12e^{2\theta}<24,\qquad H=3e^{\omega T}.
 \tag{V44.1}\label{c18v44:H}
\]
These replace, and do not improve, V42's local constants.

\begin{lemma}[Global weighted variational propagator]
\label{c18v44:propagator}
For every fine configuration $U$ and $0\le r\le t\le T$,
the actual forward variational propagator satisfies
\[
 \|\mathsf V(t,r;U)\|_{{\rm S},\theta}
                 \le3e^{\omega(t-r)}\le H .
 \tag{V44.2}\label{c18v44:prop}
\]
\end{lemma}
\begin{proof}
Let $\mathsf N_{ef}$ count plaquettes shared by two distinct
fine links, with zero diagonal. Its row and column sums are
twelve; nonzero entries have tail distance at most two.
As in V24 and V31, in linkwise parallel frames the Hessian
generator has diagonal block norm at most four and
off-diagonal block norm at most $\mathsf N_{ef}$.
The time-ordered integral series therefore bounds each
propagator block by the corresponding entry of
$\exp((t-r)(4I+\mathsf N))$.
The weighted row and column sums of $4I+\mathsf N$ are at
most $4+12e^{2\theta}$. Submultiplicativity and the exponential
series bound those of its exponential by $e^{\omega(t-r)}$.
Changing to invariant frames rotates individual blocks on
both sides, preserving their operator norms and positions.
For each scalar row or column, a $3$-by-$3$ block contributes
at most three times its operator norm. This proves the result.
No nearness to a stationary point or positivity of the
Hessian is used.
\end{proof}

For the moving-frame generator $\mathsf M$ and chain sampler
$\mathsf B=Db_\ell$, the all-configuration incidence bounds are
\[
\begin{gathered}
 \mathfrak a_\theta(\partial\mathsf M)\le t_1=1728e,\qquad
 \mathfrak a_\theta(\partial^2\mathsf M)\le t_2=20736e,\\
 \mathfrak a_\theta(\mathsf B)\le b_0=3e\ell,\quad
 \mathfrak a_\theta(\partial\mathsf B)\le b_1=18e\ell^2,\quad
 \mathfrak a_\theta(\partial^2\mathsf B)\le b_2=108e\ell^3 .
\end{gathered}
 \tag{V44.3}\label{c18v44:localinputs}
\]
Here $b_1,b_2$ are sampler majorants, not V41's generator
operator-norm constants. To verify global validity, each
ordered plaquette derivative through order four is bounded
by unit-quaternion insertion. An entry of either differentiated
generator costs at most three, including its commutator term.
An anchored slot meets four plaquettes; the other slots have
twelve choices each. The counts are $4\cdot12^2\cdot3$
and $4\cdot12^3\cdot3$, with weight at most
$e^{2\theta}\le e$. Sampler prefixes are orthogonal adjoint
matrices at every configuration. Their first two derivatives
have entry bounds two and four, and the chain counts are
$3\ell,18\ell^2,108\ell^3$. Their weight is at most
$e^{\theta\ell}\le e$. These are the V42 incidence proofs;
neither one used its neighborhood restriction.

Set
\[
\begin{aligned}
 k&=t_1TH^3,&
 l&=\tfrac32t_1^2T^2H^5+t_2TH^4,\\
 c_0&=b_0H,&
 c_1&=b_1H^2+b_0k,\\
 c_2&=b_2H^3+3b_1Hk+b_0l .
\end{aligned}
 \tag{V44.4}\label{c18v44:Cconstants}
\]
The letter $l$ denotes a scalar, not a block length.
For $C_s=D(b_\ell\circ\Phi_s)$, globally in $U$ and $s\le T$,
\[
 \mathfrak a_\theta(C_s)\le c_0,\qquad
 \mathfrak a_\theta(\partial C_s)\le c_1,\qquad
 \mathfrak a_\theta(\partial^2C_s)\le c_2.
 \tag{V44.5}\label{c18v44:Cjets}
\]
Indeed apply the connected-contraction lemma
\ref{c18v42:contract} to the exact ordered variational
equations \eqref{c18v41:jets}, with zero initial jets.
One first-generator term gives $t_1sH^3$; the three
first-generator/first-jet terms integrate to
$\frac32t_1^2s^2H^5$, and the second-generator term gives
$t_2sH^4$. The two first chain-rule terms and the five
second chain-rule terms in \eqref{c18v41:Csecond} then
give \eqref{c18v44:Cjets}. Every tensor network retains
an external anchor. No closed extensive trace is bounded
by this argument.

\subsection{The global full Gram and its normal right inverse}

Let $G_s=C_sC_s^*$, $Z_s=G_s^{-1}$ and $R_s=C_s^*Z_s$.
The full-output global floor is
\[
 G_s\succeq \ell e^{-32s}I\succeq g_0I,\qquad
 g_0=\ell e^{-32T},\qquad
 \|G_s\|_{{\rm S},\theta}\le g_{\rm S}=c_0^2 .
 \tag{V44.6}\label{c18v44:Gram}
\]
For clarity, $\mathsf B\mathsf B^*=\ell I$ exactly, since
the straight chains are edge-disjoint. The inverse
variational propagator has operator norm at most $e^{16s}$.
Consequently
$|\mathsf J_s^*\mathsf B^*v|\ge e^{-16s}\sqrt\ell\,|v|$
for every full coarse $v$, proving the floor without a
Fourier compression or a special background.

Define
\[
 \eta=\frac{\theta g_0}{4g_{\rm S}},\qquad
 \zeta=\frac\eta2,\qquad
 z_0=\frac{18}{g_0}\left(1+\frac4{\eta\ell}\right)^3.
 \tag{V44.7}\label{c18v44:Zconstants}
\]
Then $\eta\le\theta/4$, and Lemma \ref{c18v42:inverse}
applies with these global inputs:
\[
 |Z_{\alpha\gamma}|\le\frac2{g_0}
              e^{-\eta d(\alpha,\gamma)},\qquad
 \mathfrak a_\zeta(Z)\le z_0 .
 \tag{V44.8}\label{c18v44:Z}
\]
Its proof conjugates $G$ by the distance exponential,
bounds the perturbation by $g_0/2$, and sums the resulting
entry bound over nine scalar coarse indices per vertex.
The sum is bounded by the corresponding $\mathbb Z^3$ sum,
uniformly in $n$. This is the inverse of a finite normal
Gram, not the conditional kinetic inverse.

Now put
\[
\begin{aligned}
 g_1&=2c_0c_1,& g_2&=2c_0c_2+2c_1^2,\\
 z_1&=z_0^2g_1,& z_2&=2z_0^3g_1^2+z_0^2g_2,\\
 r_0&=c_0z_0,& r_1&=c_1z_0+c_0z_1,\\
 r_2&=c_2z_0+2c_1z_1+c_0z_2 .
\end{aligned}
 \tag{V44.9}\label{c18v44:Rconstants}
\]
The exact ordered inverse identities
\eqref{c18v42:inversejets} give
\[
 \mathfrak a_\zeta(\partial^jR_s)\le r_j\quad(j=0,1,2),
 \qquad
 \left\|\left(\sum_iX_jX_iR_{i\alpha}\right)_{j,\alpha}
                                   \right\|_{{\rm S},0}\le r_2 .
 \tag{V44.10}\label{c18v44:Rjets}
\]
All inputs to that argument have now been proved globally.
The diagonal identification in its last bound retains the
external fine derivative and coarse output, so the trace
cost has no volume factor.

\subsection{Contracting the extensive vacuum term correctly}

Let $u=\log\Omega_\beta$ up to its normalization constant.
For $0\le\beta\le\beta_*=3/(512e)$, write
\[
 R_\beta=64e\beta,\qquad
 a_\beta=R_\beta/3,\qquad h_\beta=12R_\beta/e.
 \tag{V44.11}\label{c18v44:vacuumconstants}
\]
The actual native vacuum construction and scalar
row-and-column proof of \eqref{c18v37:C2} give, everywhere,
\[
 \sup_i|X_i u|\le a_\beta,\qquad
 \|(X_jX_i u)_{j,i}\|_{{\rm S},0}\le h_\beta .
 \tag{V44.12}\label{c18v44:vacuum}
\]
The second assertion is for ordered coefficient derivatives,
not just the covariant Hessian. We need no bound on the
full extensive norm $|\nabla u|$.

For the same normal lift and density as in V41, define
\[
 S^R_\alpha=\sum_iX_iR_{i\alpha}
                  +2\sum_i(X_i u)R_{i\alpha},\qquad
 r^R_\alpha=S^R_\alpha-
              \mathbb E_{\nu_g}[S^R_\alpha]\quad\hbox{on }c_s^{-1}(g).
 \tag{V44.13}\label{c18v44:scoredefinition}
\]
Both the conditional expectation and $\nu_g$ use the
entire actual fibre, with its coarea density retained.
Set
\[
 b=r_1+2a_\beta r_0,\qquad
 L=r_2+2h_\beta r_0+2a_\beta r_1 .
 \tag{V44.14}\label{c18v44:L}
\]

\begin{theorem}[Global full-bank normal-score gradient]
\label{c18v44:score}
For every $U$, every coarse value $g$, $0\le s\le\ell^2$,
and constant full coarse vector $v$, in the stated coupling
interval,
\[
\begin{gathered}
 \|(X_jS^R_\alpha)_{j,\alpha}\|_{{\rm S},0}\le L,\qquad
 |\nabla S^R_v(U)|\le L|v|,\\
 |\nabla_{\rm fib}r^R_v(U)|\le L|v|
       \quad(U\in c_s^{-1}(g)),\\
 |S^R_\alpha(U)|\le b,\qquad |r^R_\alpha(U)|\le2b .
\end{gathered}
 \tag{V44.15}\label{c18v44:scorebounds}
\]
The constants do not depend on $n$ or $g$. The fibre norm
is the original induced norm and equals the internally
reduced quotient-gradient norm for these invariant sources.
\end{theorem}
\begin{proof}
Differentiate \eqref{c18v44:scoredefinition} in the original
ordered invariant frame:
\[
 X_jS^R_\alpha
 =\sum_iX_jX_iR_{i\alpha}
  +2\sum_i(X_jX_i u)R_{i\alpha}
  +2\sum_i(X_i u)(X_jR_{i\alpha}).
 \tag{V44.16}\label{c18v44:scoregradient}
\]
The three matrices have Schur bounds $r_2$,
$2h_\beta r_0$ and $2a_\beta r_1$, respectively.
For the last one, anchor either $j$ or $\alpha$ in
$\partial R$, bound $|X_i u|$ by $a_\beta$, and sum
the other two slots. This explains why the extensive
$\ell^2$ norm of the vacuum gradient is unnecessary.
There is no additional connection term: all derivatives
in this exact scalar identity are ordered coefficient
derivatives, and \eqref{c18v44:vacuum} has that convention.
Schur's test gives the ambient gradient bound for all $v$.

The conditional mean is constant along its fibre.
Orthogonal projection of the ambient gradient bounds
the intrinsic fibre gradient. Internal gauge equivariance
makes each source invariant under the based internal
group; the fibre gradient is horizontal and its quotient
norm is unchanged. Finally anchor $\alpha$ and identify
the two fine slots of $\partial R$ to bound the divergence
by $r_1$. The other raw-score term is at most
$2a_\beta r_0$. Centering doubles this component bound,
since it is global on the whole fibre.
\end{proof}

The last line is a componentwise estimate, not a
dimension-free pointwise bound on the full score vector.
It gives $|S^R_v|\le b\|v\|_1$, and does not justify
replacing $\|v\|_1$ by $|v|$. The full-bank gradient
bound, in contrast, really uses $|v|$.

The price of covering all configurations must not be hidden.
With absolute implicit constants, uniformly in the stated
coupling interval, the displayed recurrences imply
\[
\begin{gathered}
 c_0=O(\ell e^{24T}),\quad
 c_1=O(\ell^3e^{72T}),\quad c_2=O(\ell^5e^{120T}),\\
 z_0=O(\ell^2e^{272T}),\quad
 z_1=O(\ell^8e^{640T}),\quad z_2=O(\ell^{14}e^{1008T}),\\
 r_0=O(\ell^3e^{296T}),\quad
 r_1=O(\ell^9e^{664T}),\quad
 L=O(\ell^{15}e^{1032T}),\qquad T=\ell^2 .
\end{gathered}
 \tag{V44.17}\label{c18v44:growth}
\]
For example $\eta^{-1}=O(\ell^2e^{80T})$, which explains
the large spatial summation cost. These are conservative
majorants, not lower bounds on native scores. The result
pays the complement of V42's neighborhood for the first
source derivative; it neither gives a global polynomial
bound nor pays an inverse-weighted complementary integral.
At $\ell=2$ and $\beta=\beta_*$ the displayed formulas
give $\log_{10}L\simeq1846.59$. This is a particularly
large proof constant, not a measured physical scale.

\subsection{A paid high-frequency inverse and an unpaid low-frequency inverse}

Let $\bar L_g\ge0$ be the actual reduced conditional
kinetic operator, with original induced metric, and let
$E_g$ be its spectral resolution. At each finite regulator
it is the positive weighted elliptic operator on the
compact connected reduced fibre. Its kernel is the
constants, and the scores in \eqref{c18v44:scoredefinition}
are centered. No uniform positive eigenvalue is assumed.
Define the positive matrix-valued source measure
\[
 M_g(B)_{\alpha\gamma}
   =\langle r^R_\alpha,E_g(B)r^R_\gamma\rangle_{\nu_g}.
 \tag{V44.18}\label{c18v44:measure}
\]
Smoothness and the Dirichlet-form identity give
\[
 \boxed{\int_{(0,\infty)}\lambda\,dM_g(\lambda)
   =\left(\int\langle dr^R_\alpha,dr^R_\gamma\rangle
                                     d\nu_g\right)_{\alpha\gamma}
   \preceq L^2 I .}
 \tag{V44.19}\label{c18v44:moment}
\]
To prove the matrix inequality, contract with an arbitrary
$v$ and integrate the squared fibre-gradient bound
\eqref{c18v44:scorebounds} against the normalized actual
law. No replacement of that law by product Haar is made.

Consequently, for every $\varepsilon>0$,
\[
\begin{aligned}
 \int_{[\varepsilon,\infty)}\lambda^{-1}\,dM_g(\lambda)
                  &\preceq \varepsilon^{-2}L^2 I,\\
 \int_{[\varepsilon,\infty)}
       e^{-2t\lambda}\lambda^{-1}\,dM_g(\lambda)
                  &\preceq e^{-2t\varepsilon}
                              \varepsilon^{-2}L^2I\quad(t\ge0).
\end{aligned}
 \tag{V44.20}\label{c18v44:high}
\]
These follow from $\lambda^{-1}\le\varepsilon^{-2}\lambda$
on the indicated spectral band. They hold on every fibre
with the same $L$, at the required block time as well as
earlier times. The scalar spectral comparison is elementary;
the input newly paid here is its uniform native first moment.

The desired full inverse is still
\[
 \mathcal R^R(g)=
       \int_{(0,\varepsilon)}\lambda^{-1}\,dM_g(\lambda)
       +\int_{[\varepsilon,\infty)}
                               \lambda^{-1}\,dM_g(\lambda).
 \tag{V44.21}\label{c18v44:unpaid}
\]
Only the second term has been bounded uniformly here.
A first moment cannot bound the first one. Even a
two-state reversible chain with centered eigenfunction
$r=\pm1$ and eigenvalue $\delta$ has
$\|r\|_\infty=1$, $\langle r,Lr\rangle=\delta\le1$,
but $\langle r,L^{-1}r\rangle=\delta^{-1}$.
With jump rate $\delta/2$ its pointwise energy is also
$\delta$. This is a logical counterexample to that
inference, not a model of YM or a counterexample to its gap.

For specificity, one sufficient genuinely additional input
would be, uniformly in $g,n$ and $0<t\le\varepsilon$,
\[
 M_g((0,t])\preceq A t^{1+q}I,\qquad q>0.
 \tag{V44.22}\label{c18v44:lowtarget}
\]
Integration by parts of the scalar spectral measure after
contraction with $v$ would give
\[
 \int_{(0,\varepsilon]}\lambda^{-1}dM_g(\lambda)
             \preceq A(1+q^{-1})\varepsilon^q I .
 \tag{V44.23}\label{c18v44:lowpayment}
\]
The boundary at zero vanishes because $q>0$.
Equation \eqref{c18v44:lowtarget} is a declared target,
not a consequence of \eqref{c18v44:moment}.
Alternatively a conditional spectral floor would suffice,
with its full block-time and coupling cost paid.

\subsection{Scope, verification and next direction}

The geometric derivative obstruction outside a tiny
neighborhood has now been paid with exponential constants.
The upstream issue is therefore the source's low conditional
spectrum, not the conditional probability of another small
flat ball. A useful next calculation should test actual
low-mode coupling or a native conditional coercivity route,
including the source and its metric together. Optimizing
the exponent in \eqref{c18v44:growth} alone cannot settle
\eqref{c18v44:unpaid}. V38's short-flow conditional theorem
and V39's coupled density--metric identity remain available;
neither already supplies a block-time gap by itself.

The diagnostic
\texttt{scripts/verify\_ym\_c18\_v44\_global\_normal\_score.py}
checks native plaquette incidence and weighted positive
propagator sums, full nonflat quaternion chain samplers,
normal-Gram algebra, ordered source contractions on a
separately labelled finite matrix fixture, high-frequency
spectral inequalities and explicit constant recurrences.
The fixture and two-state warning are not native vacuum
simulations. V43's mathematical regression is replayed.
The all-volume claims rest on the proofs above, not on
finite samples or proof-assistant certification.

V44 replaces only V43. Removing this appendix and restoring
the title recovers V43 byte for byte. The other 53 source
files are untouched, and only \texttt{.tex} files remain
in \texttt{latex\_notes}; build and QA products are outside.
The next companion review is V45, broad sweep V50, and
archive/restart V100. The conditional low-spectrum input,
weak-coupling transport, physical return, cofinal scale
calibration, interacting four-dimensional reconstruction
and positive physical Yang--Mills mass gap remain unproved.
\addtocontents{toc}{\protect\setcounter{tocdepth}{2}}
% END CYCLE 18 VERSION 44 APPEND
% BEGIN CYCLE 18 VERSION 45 APPEND
\clearpage
\addtocontents{toc}{\protect\setcounter{tocdepth}{1}}
\section[V45: block-two geometry and normal-score cancellation]{V45: block-two geometry and normal-score cancellation}
\label{c18v45:context}

V44 controls the normal source derivative globally but does not
control its low conditional spectrum. Rather than optimize its
exponential majorant, we compute an actual source cancellation.
At block length two the unflowed sampled fibres have an especially
simple native geometry. This allows an exact first-flow calculation
for the normal lift, including its changing vertical correction.
Its sign differs from the tree-lift sign in V38.

The result is not a block-time or continuum theorem. For $\ell=2$,
the first normal source is proportional to $\beta-24s$ and vanishes
on the short-flow curve $s=\beta/24$. We also calculate its full
conditional covariance and native eigenvalue. Remainders below
are explicitly fixed-regulator remainders. The required block
time is $s=4$, not $s=\beta/24$.

\subsection{Scheduled companion review and transfer boundary}

The previous V44 turn made progress, and its source and preservation
checks pass. The scheduled five-version review covered the current
SM V72 and NS V26 terminal results, and the related Source Selection
V65, Stationarity V61, Einstein Bridge V55, ESM restart V14, Shell
Mixing V16, parallel YM V5 and Arithmetic Loading V17 scopes.
Except for the already recorded ESM restart after V40, the reviewed
files are unchanged. This is a relevant-results review, not an
independent recertification of all their accumulated proofs.

ESM V14 distinguishes equivalence of protected zero-mode ranges,
extension to the complementary carrier, and an identity-component
unitary path. Its protected spectral floor and coefficientwise,
positive-reference domain do not prove a nonperturbative YM floor.
Source Selection's finite Choi-defect gap likewise requires the
actual represented source. Stationarity and Einstein Bridge retain
uncomputed same-source or same-root coefficients. SM V72's
many-copy, bounded-time limit is not a physical finite-species
continuum limit. NS V26's input-specific contraction calculation
suggests testing the actual tensor contraction, not importing its
numerical constants or grid claims. Shell Mixing keeps its
coarea/source factors and lacks a common nonperturbative remainder;
parallel YM excludes an extensive global-charge tightness shortcut.
Arithmetic Loading still requires represented source projectors.

The transfer is therefore methodological: compute the normal source
on its actual carrier and keep the density, metric and lift together.
No companion gap or cancellation hypothesis is promoted to a YM
conclusion. For literature context, the primary abstract of
M.~L\"uscher,
\href{https://arxiv.org/abs/0907.5491}
{\emph{Trivializing maps, the Wilson flow and the HMC algorithm}},
was checked. Approximate trivialization by a flow is not a new
general idea. Our calculation concerns the native Hamiltonian
ground density, this block-two sampler and its normal score;
no trivializing-map theorem or Wilson-Gibbs law is substituted
for those objects. The broad literature sweep remains due at V50.

\subsection{Block-two fibres are totally geodesic products}

For the remainder of this appendix fix $\ell=2$, $N=2n$, $n\ge5$.
Let $\mathcal E_{\rm ch}$ be the union of the sampled two-link
chains, and let $\mathcal E_{\rm free}$ be the other fine links.
This is a decomposition before internal gauge reduction, not
the rooted forest/chord decomposition of V33.
For a coarse link $e$ the two fine variables on its chain,
conditioned on $b_2(U)_e=g_e$, can be written
\[
             (U_{e,1},U_{e,2})=(x_e,x_e^{-1}g_e).
 \tag{V45.1}\label{c18v45:pair}
\]
Chains are edge-disjoint. Inversion and left/right multiplication
are isometries of the unit-round bi-invariant $SU(2)$ metric.
The graph in \eqref{c18v45:pair} is therefore totally geodesic
in the product, with induced metric twice the unit-round metric.
Indeed a pair of corresponding geodesics remains in that graph,
which proves the vanishing of its second fundamental form.

Consequently the full fibre $\mathcal F_{g,0}=b_2^{-1}(g)$ has
product metric
\[
 \mathfrak g_{F,0}
   =\bigoplus_{e\in E_{\rm c}}2\mathfrak g_{SU(2)}
       \ \oplus\!
       \bigoplus_{a\in\mathcal E_{\rm free}}\mathfrak g_{SU(2)},
 \qquad \mathrm{II}_0=0 .
 \tag{V45.2}\label{c18v45:product}
\]
At $\beta=0$, its normalized conditional law is product Haar
in these variables. The normal Jacobian is constant since
$BB^*=2I$. Its positive weighted Laplacian is
\[
 L_{F,0}=\frac12\sum_{e\in E_{\rm c}}T_{x_e}
                   +\sum_{a\in\mathcal E_{\rm free}}T_{U_a}.
 \tag{V45.3}\label{c18v45:L0}
\]
In particular its full, unreduced gap is exactly $3/2$.
This statement includes internal gauge functions; it is not
V37's larger comparison floor on the reduced invariant space.
The Ricci tensor of this product is at least $\mathfrak g_{F,0}$.

Let $R_0=B^*/2$ and let $R_{0,e,v}$ denote one coarse-link
column in direction $v\in\mathfrak{su}(2)$.
Its ambient flow on this pair, starting on the $g_e$ fibre, is
\[
 x_e(t)=e^{tv/2}x_e,\qquad
 g_e(t)=e^{tv}g_e,\qquad
 U_{e,2}(t)=x_e^{-1}e^{tv/2}g_e .
 \tag{V45.4}\label{c18v45:isometry}
\]
The other pair and free variables are fixed. The first formula
is a round isometry between the two fibre coordinate spaces.
Thus these horizontal flows identify fibres isometrically,
preserve their normalized Haar laws, and intertwine $L_{F,0}$.
Also $\operatorname{div}R_{0,\alpha}=0$, either by this
volume statement and constant normal Jacobian, or because
each coefficient has zero derivative in its own fine link.

\subsection{The actual boundary plaquettes and their eigenvalue}

Every elementary plaquette contains either zero or two sampled
fine links. To check this, for a plaquette in the $(i,j)$ plane,
its third coordinate must be even for any edge to be sampled.
If it is even, exactly one of the two consecutive $i$-edge
transverse coordinates is even, and likewise for the $j$ edges.
The two captured edges belong to different chains and meet
at one coarse root. Write $\mathcal S$ for these plaquettes and
\[
 W_{\mathcal S}=\sum_{p\in\mathcal S}w_p,\qquad
 f_\alpha=R_{0,\alpha}W=R_{0,\alpha}W_{\mathcal S}.
 \tag{V45.5}\label{c18v45:f}
\]
At each coarse root there are twelve such plaquettes:
choose two distinct axes and a sign on each axis.
Hence $|\mathcal S|=12n^3$. The other $12n^3$ elementary
plaquettes do not enter $f$.

A term $w_p$ with $p\in\mathcal S$ has fundamental spin in
two pair coordinates and two free links. Thus
\[
 L_{F,0}w_p=(\tfrac32+\tfrac32+3+3)w_p=9w_p,\qquad
 \mathbb E_{\nu_{g,0}}w_p=0 .
 \tag{V45.6}\label{c18v45:Wmode}
\]
The zero mean follows already by integrating one of its free
links. The isometry \eqref{c18v45:isometry} commutes with
the fibre Laplacian and conditional integration, giving
\[
        L_{F,0}f_\alpha=9f_\alpha,\qquad
        \mathbb E_{\nu_{g,0}}f_\alpha=0 .
 \tag{V45.7}\label{c18v45:fmode}
\]
The map, vacuum and normal lift are internally gauge equivariant;
$f_\alpha$ is invariant under the root-fixed internal group.
Thus \eqref{c18v45:fmode} holds also for the exact reduced
operator $\bar L_{g,0}$. No gauge or coarse column is deleted.
The eigenvalue agrees numerically with V35 at $\ell=2$,
but $f$ is a different, normal-lift bank, not its last-edge bank.

\subsection{The first native curvature variation, with bulk cancellation}

This subsection concerns the full fibre before the internal
quotient. Put $F=\nabla W$ and
$Q_s=R_sC_s$. Locally in time identify $\mathcal F_{g,0}$
with $\mathcal F_{g,s}$ by the normal velocity
\[
                    V_s=-Q_sF .
 \tag{V45.8}\label{c18v45:normalvelocity}
\]
Indeed $\partial_sc_s=C_sF$, since $c_s=b_2\Phi_s$ and the
flow is autonomous, so $\partial_sc_s+C_sV_s=0$.
All statements in this subsection are derivatives at $s=0$,
where this smooth normal transport exists at each finite
regulator.

Let $\lambda_s=\frac12\log\det(C_sC_s^*)$.
Since $BB^*=2I$ identically, differentiation of the Gram,
or its determinant along the forward flow, gives
\[
 \left.\partial_s\lambda_s\right|_0
     =\operatorname{Tr}(Q_0\operatorname{Hess}W)
     =\Delta_{\mathcal M}W-\Delta_{F,0}W
     =-3W_{\mathcal S}.
 \tag{V45.9}\label{c18v45:normaltrace}
\]
In the middle equality $\mathrm{II}_0=0$ is essential.
Plaquettes not meeting the chains have full fibre eigenvalue
twelve, so cancel exactly; \eqref{c18v45:Wmode} treats the others.
Because $\lambda_0$ is a global constant, the transport of its
argument adds no derivative.

Under normal transport, the metric derivative is
$-2\langle\mathrm{II}_0,V_0\rangle=0$.
At $\beta=0$, the density relative to induced fibre volume
is proportional to $e^{-\lambda_s}$ by V39. Thus, pulling
both metric and probability back by this transport,
\[
 \dot{\mathfrak g}_{F,0}=0,\qquad
 \left.\partial_s\log\omega_{F,g,s}\right|_0=3W_{\mathcal S},
 \qquad
 \left.\partial_s\mathcal R_{F,g,s}\right|_0
       =-3\operatorname{Hess}_{F,0}W_{\mathcal S}.
 \tag{V45.10}\label{c18v45:Ricvariation}
\]
The normalization derivative vanishes by
\eqref{c18v45:Wmode}. Since the metric derivative vanishes
identically as a tensor field on the fibre, its Ricci derivative
vanishes; the last formula follows by differentiating the
weighted Hessian. These are intrinsic native identities,
not the product comparison curvature of V38.
After internal reduction V39's orbit-gain tensor must still
be included; its derivative has not been set to zero here.
Neither a first derivative nor the initial Ricci floor proves
a curvature floor at $s=4$.

\subsection{Differentiating the changing normal lift}

All divergences in the next calculation are for the original
ambient product Haar law at $\beta=0$, or, when marked by $F$,
its actual conditional law on the unflowed fibre.
Let
\[
 A_s=(\Phi_{-s})_*R_0\circ\Phi_s,\qquad
 \mathcal V_s=A_s-R_s=P_sA_s,\qquad P_s=I-Q_s .
\]
These are the averaged transported lift and its vertical excess,
now differentiated at zero flow rather than just at a stationary
background. At zero, $A_0=R_0$ and $\mathcal V_0=0$.

\begin{lemma}[First excess and its exact divergence]
\label{c18v45:excess}
For each coarse column, with $\mathsf H_W=\operatorname{Hess}W$,
\[
 \dot A_0=[F,R_0],\qquad
 \dot{\mathcal V}_0=-2P_0\mathsf H_WR_0,\qquad
 \operatorname{div}\dot A_{0,\alpha}=12f_\alpha .
 \tag{V45.11}\label{c18v45:excessjet}
\]
\end{lemma}
\begin{proof}
Differentiating the pullback of a vector field gives its Lie
bracket with $F$. To compute its tangential component, pull
back a fibre tangent as well. Their initial inner product
is zero. The derivative of the pulled-back ambient metric
is $-2\operatorname{Hess}W$. Equivalently, at fixed $U$,
the derivative of their pairing at $\Phi_sU$ vanishes, while
the two inverse differentials contribute
$-2\langle\mathsf H_WR_0,T\rangle$.
The normal lift remains orthogonal to these moving tangents,
so their excess has the second derivative in
\eqref{c18v45:excessjet}. This argument includes the change
of the horizontal projection.
Finally
$\operatorname{div}[F,R_0]
=F(\operatorname{div}R_0)-R_0(\operatorname{div}F)$.
Use $\operatorname{div}R_0=0$ and $\Delta W=-12W$.
\end{proof}

The remaining contraction is not replaced by an operator norm.
For one column put
\[
 K_\alpha=P_0(\nabla R_{0,\alpha})^*F .
\]
Differentiating $f_\alpha=\langle F,R_{0,\alpha}\rangle$ gives
\[
 P_0\mathsf H_WR_{0,\alpha}
       =\nabla_F f_\alpha-K_\alpha .
 \tag{V45.12}\label{c18v45:K}
\]
We compute $\operatorname{div}_FK_\alpha$ explicitly.

On the pair \eqref{c18v45:pair}, write a vertical tangent as
the first-factor velocity $\xi x$. Its norm squared is
$2|\xi|^2$. For constant coarse $v$, the bi-invariant connection
gives
\[
 \nabla_{\xi}^{\perp}R_{0,v}=R_{0,-[\xi,v]/2}.
 \tag{V45.13}\label{c18v45:connection}
\]
For example in right-invariant fine frames the vertical
coefficients are $(\xi,-\operatorname{Ad}_{x^{-1}}\xi)$
and the normal coefficients are
$(v/2,\operatorname{Ad}_{x^{-1}}v/2)$.
Differentiate these coefficients and use
$\nabla_{X_a}X_b=-\frac12X_{[a,b]}$ to obtain
\eqref{c18v45:connection}. The derivative is normal because
the fibre is totally geodesic.

For one plaquette let
$\mathbf f=(R_{0,a}w_p)_{a=1}^3$ on this coarse chain.
Its contribution to $K_v$, expressed as a first-factor
velocity, is
\[
                         k_v=-\tfrac14[v,\mathbf f].
 \tag{V45.14}\label{c18v45:k}
\]
Indeed the normal force is $2\sum_a f_aR_{0,a}$ and
$\langle R_{0,a},R_{0,b}\rangle=\delta_{ab}/2$;
pair \eqref{c18v45:connection} with that force and use the
vertical metric factor two.

If the plaquette uses the first chain edge, then
$f_a=\frac12X_aw_p$ in the fibre $x$ coordinate.
If it uses the second, then $f_a=-\frac12X_aw_p$.
There is never more than one edge from this chain in a
plaquette. Choose unit imaginary quaternions $\tau_a$, so
$[\tau_a,\tau_b]=2\epsilon_{abc}\tau_c$, while the
right-invariant derivative fields obey
$[X_a,X_b]=-2\epsilon_{abc}X_c$.
For $\sigma=1$ or $-1$ accordingly, \eqref{c18v45:k} yields
\[
 \operatorname{div}_FK_v
 =-\frac{\sigma}{4}\sum_{a,b}\epsilon_{cba}X_aX_bw_p
 =-\frac{\sigma}{2}X_cw_p=-f_c\quad(v=\tau_c).
 \tag{V45.15}\label{c18v45:divK}
\]
Here antisymmetrizing $X_aX_b$ gives
$\sum_{a,b}\epsilon_{cba}X_aX_b=2X_c$.
The constant metric factor two changes the gradient-vector
coefficients but not the Haar divergence in $x$.
Other pairs contribute nothing to this column's connection.
Summing the plaquettes, then using \eqref{c18v45:fmode},
proves
\[
 \operatorname{div}_F(P_0\mathsf H_WR_{0,\alpha})
       =-9f_\alpha+f_\alpha=-8f_\alpha .
 \tag{V45.16}\label{c18v45:contracted}
\]
Since $\dot{\mathcal V}_0$ is tangent and the original normal
Jacobian is constant, its ambient divergence equals this
conditional fibre divergence. Thus no extra coarea term
has been omitted in applying \eqref{c18v45:contracted}.

\begin{theorem}[Exact first normal-score cancellation]
\label{c18v45:cancellation}
For the actual normal score of V41, at $(\beta,s)=(0,0)$,
\[
 S^R_{0,0}=r^R_{0,0}=0,\qquad
 \left.\partial_\beta r^R\right|_{0,0}=\frac16 f,\qquad
 \left.\partial_s r^R\right|_{0,0}=-4f .
 \tag{V45.17}\label{c18v45:first}
\]
All coarse columns and every coarse value $g$ are included.
The derivatives may be pulled back by either the exact tree
chart or the normal fibre transport; their first values agree.
\end{theorem}
\begin{proof}
At zero coupling the raw score is $\operatorname{div}R_s$.
From \eqref{c18v45:excessjet} and
\eqref{c18v45:contracted},
\[
 \operatorname{div}\dot R_{0,\alpha}
 =12f_\alpha
      +2\operatorname{div}_F(P_0\mathsf H_WR_{0,\alpha})
 =-4f_\alpha .
\]
For the coupling derivative use the native ground-state
coefficient $u_\beta=\beta W/12+O_n(\beta^2)$ from V35--V37.
The vacuum term $2R_0u_\beta$ contributes $f/6$.
The base raw score vanishes identically; changing the fibre
chart or differentiating its centering measure therefore adds
no first-order term. The first coefficients already have
conditional mean zero by \eqref{c18v45:fmode}.
\end{proof}

The $12f$ transported-lift term alone has the opposite sign.
Subtracting the changing vertical excess is essential.
Likewise \eqref{c18v45:Ricvariation}'s density derivative
cannot be differentiated in a fictitiously fixed coarse
connection to recover the normal score. Both computations
refer to their specified transports.

\subsection{The whole leading covariance and its retained kernel}

For a coarse vertex $z$ let its six incident oriented legs
be labelled by $(i,+)$ and $(i,-)$. Define the six-vector
of outgoing root velocities $K_zv$ by
\[
 (K_zv)_{(i,+)}=v_{(z,i)},\qquad
 (K_zv)_{(i,-)}
       =-\operatorname{Ad}_{g_{(z-\hat i,i)}^{-1}}
                                      v_{(z-\hat i,i)} .
 \tag{V45.18}\label{c18v45:Kstar}
\]
Let $L_{\rm oct}$ be the graph Laplacian on these six legs,
joining precisely pairs from different axes, with each
edge weight one. It is the octahedral graph $K_{2,2,2}$.
Its eigenvalues are $0,4,6$ with multiplicities $1,3,2$.
For instance the three within-axis differences have
eigenvalue four, the two axis-constant differences have
eigenvalue six, and the constant vector has eigenvalue zero.

\begin{theorem}[Exact all-output leading covariance]
\label{c18v45:covariance}
For every coarse configuration $g$,
\[
 \Sigma_f(g):=\mathbb E_{\nu_{g,0}}[ff^*]
 =\frac1{16}\sum_z K_z^*(L_{\rm oct}\otimes I_3)K_z
 \preceq\frac34 I,\qquad
 \langle f,\bar L_{g,0}^{-1}f\rangle=\frac19\Sigma_f(g).
 \tag{V45.19}\label{c18v45:Sigma}
\]
The full covariance kernel consists precisely of vectors
for which some vertex colors $\zeta_z$ satisfy
\[
 v_{(z,i)}=\zeta_z,\qquad
 \zeta_{z+\hat i}
       =-\operatorname{Ad}_{g_{(z,i)}^{-1}}\zeta_z .
 \tag{V45.20}\label{c18v45:kernel}
\]
\end{theorem}
\begin{proof}
Base a plaquette of $\mathcal S$ at its unique coarse root.
The normal coarse variation gives outgoing velocities
$(K_zv)_a/2$ and $(K_zv)_b/2$ on its two captured edges.
Gauge invariance at that root makes its derivative depend
on their difference. Its remaining two fine links are free.
Integrating either one makes the plaquette holonomy Haar.
For a Haar unit quaternion its imaginary vector has second
moment $I_3/4$. Therefore this plaquette contributes
$|(K_zv)_a-(K_zv)_b|^2/16$.

Different plaquettes have different unordered pairs of
unsampled edges: two perpendicular edges at the opposite
corner determine the elementary square uniquely. Thus
for two distinct terms at least one unsampled edge occurs
in only one of them. Its fundamental Haar integral kills
the cross term. Sum the twelve pairs at each root to obtain
the equality in \eqref{c18v45:Sigma}. This uses the exact
unflowed zero-coupling law only, not independence at
positive coupling or flow.
Since $L_{\rm oct}\preceq6I$ and every coarse link contributes
isometrically at its two endpoints,
$\sum_z|K_zv|^2=2|v|^2$, proving the upper bound.
The inverse equality follows from \eqref{c18v45:fmode}.
Finally zero quadratic form requires the six root velocities
to agree at every vertex because the octahedral graph is
connected. This is exactly \eqref{c18v45:kernel}.
\end{proof}

At $g=I$, a unitary coarse Fourier transform gives, per color,
with $d_i(k)=1-e^{-ik_i}$,
\[
 \Sigma_f(k)=\frac1{16}
       \left[\operatorname{diag}(8+|d_i(k)|^2)
                               -\overline{d(k)}\,d(k)^{\mathsf T}\right].
 \tag{V45.21}\label{c18v45:Fourier}
\]
In particular $\Sigma_f(0)=I/2$. On an even coarse torus
the checkerboard mode $k=(\pi,\pi,\pi)$ has one zero
orientation eigenvalue and two eigenvalues $3/4$, each
repeated over color. On an odd coarse torus the kernel
at $g=I$ is zero, since alternating vertex colors cannot
close around an odd cycle. These are covariance eigenvalues,
not small eigenvalues of the conditional kinetic operator.
All directions, including the kernel, remain in the formula.

\subsection{What the cancellation buys, and the next actual coefficient}

At each fixed $n$, smooth elliptic perturbation on the compact
fibres and the native vacuum expansion give, after a fixed
smooth fibre identification, for any fixed derivative order,
\[
 r^R_{g,\beta,s}
   =\left(\frac{\beta}{6}-4s\right)f
                 +O_n((|\beta|+|s|)^2).
 \tag{V45.22}\label{c18v45:expansion}
\]
The constants may be chosen uniformly in $g$ at this fixed
regulator by compactness. They have not been bounded uniformly
in $n$. The base operator has a positive gap and all the
coefficient families are smooth; identifying the varying
constant-orthogonal spaces by their smooth densities therefore
also gives the fixed-regulator inverse expansion
\[
 \mathcal R^R(g;\beta,s)
   =\frac{(\beta-24s)^2}{324}\Sigma_f(g)
                       +O_n((|\beta|+|s|)^3).
 \tag{V45.23}\label{c18v45:inverseexpansion}
\]
This is the normal bank's coefficient, not V35's tree-bank
coefficient with a coupling renamed.

Along $s=\beta/24$, for $\beta\ge0$ sufficiently small,
\[
       r^R_{g,\beta,\beta/24}=O_n(\beta^2),\qquad
       \mathcal R^R(g;\beta,\beta/24)=O_n(\beta^4).
 \tag{V45.24}\label{c18v45:tuned}
\]
The second statement follows from the first and a bounded
fixed-regulator centered inverse, not merely from cancelling
the quadratic term in \eqref{c18v45:inverseexpansion}.
In fact V38's already proved uniform conditional gap applies
on the stated small-coupling interval along this curve:
$\beta_*/24=1/(4096e)<1/1024$.
It does not remove the unproved volume dependence of the
second source remainder. No new gap theorem is inferred.

This is a first-order cancellation for the actual native
normal lift, without adding a Poisson correction. It is not
an exact score-free map, an RG coupling law, or control
through $s=\ell^2=4$. In particular the geometric Gram
inverse and the conditional inverse remain distinct.
The spectral mass condition in V44 is still unpaid at the
required block time.

The next useful source calculation is now explicit:
determine
$q_g=\lim_{\beta\to0}\beta^{-2}
 r^R_{g,\beta,\beta/24}$ with its complete normal-lift,
flow, density and centering terms. A second scalar time
adjustment $s=\beta/24+a\beta^2$ changes it only by
$-4a f$. Hence components outside the $f$ bank cannot
be removed by choosing $a$. They must be computed, corrected
on their actual carrier, or retained as an obstruction.
V36's zero-flow second coefficient can inform this calculation
but does not already contain these changing-map terms.
This test should precede any assertion of a volume-uniform
or iterated trivialization.

The diagnostic
\texttt{scripts/verify\_ym\_c18\_v45\_first\_normal\_flow.py}
checks the literal plaquette word with both chain positions,
its ordered derivatives, normal right-inverse time jet and
divergence. One populated rational quaternion test differentiates
through third word order and gives six exact zero residuals
for the coefficient $-4f$; numerical differences separately
check all four endpoint orientations. It also checks degree-two
Haar moments, native plaquette incidence on sides ten and twelve,
the exact octahedral spectrum, and full flat Fourier matrices.
These are supporting regressions, not an interacting conditional
simulation or proof-assistant certificate. V44's mathematical
regression is replayed.

V45 replaces only V44, whose text is recovered byte for byte
by removing this appendix and restoring its title. All other
53 source files are preserved. Only \texttt{.tex} files remain
in \texttt{latex\_notes}, with build products outside.
The next companion review and broad literature sweep are V50;
archive/restart remains V100. Uniform nonlinear remainders,
block-time low-spectrum control, weak-coupling transport,
physical scale calibration, interacting continuum reconstruction
and the positive physical YM mass gap remain unproved.
\addtocontents{toc}{\protect\setcounter{tocdepth}{2}}
% END CYCLE 18 VERSION 45 APPEND
% BEGIN CYCLE 18 VERSION 46 APPEND
\clearpage
\addtocontents{toc}{\protect\enlargethispage{2\baselineskip}}
\addtocontents{toc}{\protect\setcounter{tocdepth}{1}}
\section[V46: the second tuned source and its adjoint obstruction]{V46: the second tuned source and its adjoint obstruction}
\label{c18v46:context}

V45 cancelled the first normal source on $s=\beta/24$.
The present note computes the complete second coefficient as
a finite local differential expression, evaluates its entire
centre-even sector, and proves that no second scalar adjustment
of flow time can cancel that sector. Its conditional inverse
cost is strictly positive on each individual coarse color
column. We also identify a different, adjoint-character flow
which cancels this particular sector at second order. The
mixed-plaquette sectors are retained, not declared cancelled.

The setting remains $SU(2)$, block length two, $N=2n$, $n\ge5$,
with the native Hamiltonian $T-\beta W$, its exact vacuum, and
the full normal lift. All statements about Taylor remainders
are fixed-regulator statements unless expressly stated otherwise.
This is not an estimate at the required block time $s=4$ or
at continuum weak coupling. The physical YM mass gap remains
unproved.

The previous V45 turn was progress; its source, preservation
check and mathematical regressions have been revalidated.
No companion changed at the opening inventory. The scheduled
companion review was V45 and the next companion and broad
literature reviews are V50. A targeted recheck of the primary
\href{https://arxiv.org/abs/0907.5491}{L\"uscher abstract}
retains the limited context recorded in V45: systematic
flow-based trivialization is an existing idea. No result for
a Wilson-Gibbs action is substituted for this Hamiltonian
ground density and this normal connection.

\subsection{Pulling the actual normal connection to a fixed fibre}

Use V45's pair/free product chart for $\mathcal F_{g,0}$.
On the full fine-link manifold put
\[
 F=\nabla W,\quad H=\operatorname{Hess}W,\quad
 B=Db_2,\quad R_0=B^*/2,\quad P=I-R_0B.
 \tag{V46.1}\label{c18v46:base}
\]
Here $H$ is also viewed as the metric-raised endomorphism.
Let $\Phi_s$ be the $F$ flow. At $Y=\Phi_s(U)$ push the
normal lift forward, and pull the original ambient metric
back, as follows:
\[
 \widehat R_s(Y)=D\Phi_s(U)R_s(U),\qquad
 m_s=(\Phi_{-s})^*\mathfrak g.
 \tag{V46.2}\label{c18v46:fixed}
\]
Then $B\widehat R_s=I$, and $\widehat R_s$ is orthogonal
to $\ker B$ for $m_s$, not generally for $\mathfrak g$.
Consequently its vertical correction has the expansion
\[
 \widehat R_s=R_0+sV_1+s^2V_2+O_n(s^3),
 \quad
 \begin{cases}
 V_1=2PHR_0,\\
 V_2=P\{4HPH-2H^2-\nabla_FH\}R_0.
 \end{cases}
 \tag{V46.3}\label{c18v46:Vjets}
\]
Both displayed corrections are tangent to the unflowed fibre.
They are not the jets of V45's original-coordinate vertical
excess, which has a different sign and transport convention.

\begin{proof}
Since $\mathcal L_F\mathfrak g=2H$, the inverse-flow pullback
has metric expansion
\[
 m_s=\mathfrak g-2sH+s^2(\nabla_FH+2H^2)+O_n(s^3).
 \tag{V46.4}\label{c18v46:metricjet}
\]
For a symmetric two-tensor, $\mathcal L_FH=\nabla_FH+2H^2$;
this follows directly by differentiating its two tangent
arguments and using $\nabla F=H$. Write
$\widehat R_s=R_0+sV_1+s^2V_2+\cdots$, with $PV_j=V_j$.
The identity $Pm_s\widehat R_s=0$ gives successively
$V_1=2PHR_0$ and
$V_2=2PHV_1-P(\nabla_FH+2H^2)R_0$.
These are exactly \eqref{c18v46:Vjets}. Thus both the metric
change and the orthogonal-lift change have been included.
\end{proof}

All entries in these jets are explicit ordered derivatives
of elementary plaquette words. For clarity, in right-invariant
orthonormal fine frames set
\[
 D_{ji}=X_iX_jW,\qquad A=\operatorname{diag}_e
          (\operatorname{ad}_{F_e}).
\]
Their coefficient formulas are
\[
 H=D+\tfrac12A,\qquad
 \nabla_FH=\sum_kF_kX_kH-\tfrac12AH+\tfrac12HA.
 \tag{V46.5}\label{c18v46:frame}
\]
The sign uses $\nabla_{X_a}X_b=-X_{[a,b]}/2$.
This supplies a finite recipe through fourth word order
when the divergence of $V_2$ is taken. It does not require
the unknown interacting conditional inverse.

\subsection{The complete second source, including density and centering}

Let $F_2$ be the already proved native coefficient (V36.2),
recalled here only to specify the input:
\[
 F_2=-\frac1{2304}\sum_p(4w_p^2-1)
       +\sum_{p\sim q}
           \left(\frac{R_{pq}}{702}-\frac{w_pw_q}{936}\right).
 \tag{V46.6}\label{c18v46:F2}
\]
The second sum is over distinct fine-edge-sharing plaquettes;
$R_{pq}$ is their six-edge outer-loop half-trace as defined
in V36. This coefficient, not a Gibbs replacement, occurs
in $\log\Omega_\beta^2=\beta W/6+\beta^2F_2+
\text{constant}+O_n(\beta^3)$.

The ambient measure pushed to $Y$ has logarithmic density
relative to original product Haar measure
\[
 \begin{split}
 \ell_{\beta,s}(Y)
  &=\log\Omega_\beta^2(\Phi_{-s}Y)
                       +\log\operatorname{Jac}(\Phi_{-s})(Y)\\
  &=(\beta/6+12s)W+\beta^2F_2
       -(\beta s/6+6s^2)|F|^2
       +\text{constant}+O_n((|\beta|+|s|)^3).
 \end{split}
 \tag{V46.7}\label{c18v46:density}
\]
Indeed $\operatorname{div}F=-12W$, so the second logarithm
is $12\int_0^s W(\Phi_{-t}Y)\,dt$.
This proves both the $12sW$ and $-6s^2|F|^2$ terms.
Pulling the vacuum coefficient back supplies the other
mixed term. The inverse-flow Jacobian has not been omitted.

Write $\Pi_g a=a-\mathbb E_{\nu_{g,0}}a$, where
$\nu_{g,0}$ is the exact zero-coupling product Haar
conditional law in V45's pair/free coordinates.
Let $f=R_0W$. Define the following complete local expression,
column by column:
\[
 Q=R_0F_2-\frac5{288}R_0|F|^2
                  +\frac1{36}V_1W+\frac1{576}\operatorname{div}V_2,
 \qquad q_g=\Pi_g Q.
 \tag{V46.8}\label{c18v46:Q}
\]
Here $R_0F_2$ denotes its bank of directional derivatives,
and the divergence sums the original fine invariant
directions in each vector-field column. Equations
\eqref{c18v46:Vjets}, \eqref{c18v46:frame} and
\eqref{c18v46:F2} specify all its inputs explicitly.
Mixed plaquettes are present in this formula even though
their individual irreducible components are not tabulated below.

\begin{theorem}[Second tuned coefficient]
\label{c18v46:second}
For $s=\beta/24+a\beta^2$, with fixed scalar $a$, the
actual centered normal source, identified with
$\mathcal F_{g,0}$, satisfies
\[
 r^R_{g,\beta,s}=\beta^2(q_g-4af)+O_n(\beta^3).
 \tag{V46.9}\label{c18v46:tuned}
\]
The remainder is uniform in $g$ at fixed $n$ in any fixed
smooth norm. No uniform-in-$n$ remainder follows here.
\end{theorem}
\begin{proof}
The pushed raw score is
$\operatorname{div}\widehat R_s+\widehat R_s\ell_{\beta,s}$.
At $s=\beta/24$, \eqref{c18v46:density} becomes
$2\beta W/3+\beta^2(F_2-5|F|^2/288)+\text{constant}
+O_n(\beta^3)$.
By V45, $\operatorname{div}V_1=-16f$.
The coefficient of $\beta$ is consequently
$2f/3-16f/24=0$ identically on the full fine-link space.
The second coefficient is precisely \eqref{c18v46:Q}.
The change $a\beta^2$ in $s$ adds
$a(12f+\operatorname{div}V_1)=-4af$.

The base score and its entire tuned first coefficient vanish
pointwise. Thus differentiating either the fibre identification
or the conditional centering measure creates no additional
second-order term. Centering the second coefficient uses
only $\nu_{g,0}$; also $\Pi_gf=f$.
Smooth compact-fibre perturbation and the native vacuum
expansion give the stated fixed-regulator remainder.
All functions are internally gauge invariant, so the same
identities hold on the exact internally reduced carrier.
\end{proof}

\subsection{A centre projection that isolates every self-plaquette term}

On $\mathcal F_{g,0}$ independently change the sign of each
pair coordinate $x_e$, and of each free-link coordinate.
Call this finite product group $\mathcal C_{\rm cen}$, and let
$\mathsf E$ be its normalized averaging projection.
For a pair, $x_e\mapsto-x_e$ changes both fine links' signs
and leaves their product $g_e$ fixed. These transformations
are isometries preserving product Haar measure; they commute
with the internal gauge action and with $\bar L_{g,0}$.
They do not preserve the positive-coupling measure, and are
used only on the actual Taylor coefficients at the base law.

Every plaquette has two or four distinct free links and
defines a nontrivial character $\varepsilon_p$ of
$\mathcal C_{\rm cen}$. Distinct plaquettes have different characters.
In fact their sets of free links cannot agree: each has
at least two such links, whereas distinct elementary
plaquettes share at most one fine edge when $n\ge5$.
The finer character including pair signs is therefore
also distinct. In particular
\[
 \mathsf E f=0,\qquad
 \mathsf E(\hbox{a two-plaquette term of types }p,q)=0
                       \quad(p\ne q).
 \tag{V46.10}\label{c18v46:parity}
\]
To justify the latter statement also for geometric terms,
the right-invariant derivative fields commute with centre
translations. The coefficients of $B,R_0,P$ depend on
adjoint matrices and are centre-even. Thus a differentiated
$p$ term still has character $\varepsilon_p$, and a
quadratic $p,q$ term has character $\varepsilon_p\varepsilon_q$.
The six-edge term $R_{pq}$ has the same product character.

It remains to evaluate, rather than assume, the self terms.
For a single captured plaquette put
$w=w_p$, $F_p=\nabla w_p$, $H_p=\operatorname{Hess}w_p$,
$f_p=R_0w_p$, and let $V_{1,p},V_{2,pp}$ be
\eqref{c18v46:Vjets} with $W=w_p$.
Then the exact contractions are
\[
 \begin{array}{c|c}
 \text{quantity}&\text{value}\\ \hline
 R_0|F_p|^2&-8w_pf_p\\
 V_{1,p}w_p&-6w_pf_p\\
 \operatorname{div}V_{2,pp}&4w_pf_p.
 \end{array}
 \tag{V46.11}\label{c18v46:selftable}
\]

Here is a finite algebraic proof of the contractions, including
how the constant in the last line is determined. The plaquette
and the unused partner of each captured edge form a connected
six-edge graph with one cycle. Gauge-transform five tree
edges to the identity, leaving a single free holonomy
$q=(w,\mathbf t)\in S^3$.
The displayed quantities transform covariantly in each
coarse color vector. In this gauge all three are polynomials
of degree at most two in $q$, because each is quadratic in
the single plaquette's derivatives and all $B,P$ coefficients
are independent of this free holonomy. The divergence does
not change that degree. They are even under $q\mapsto-q$.
Residual common conjugation rotates $\mathbf t$.
An even, degree-at-most-two $SO(3)$-equivariant vector
polynomial in $(w,\mathbf t)$ is a constant times
$w\mathbf t$: a constant invariant vector does not exist,
and the antisymmetric vector in $\mathbf t\otimes\mathbf t$
vanishes. Thus one nonzero rational $w\mathbf t$ determines
each column's coefficient, not merely a sampled value of
an arbitrary function.

Take both captured edges outgoing. In the literal word
$U_0U_4U_5^{-1}U_2^{-1}$, with chains $(U_0,U_1)$ and
$(U_2,U_3)$, set $U_4=(3/5,0,0,4/5)$ and all other
$U_j=1$. For the first chain's third color the exact
arithmetic using \eqref{c18v46:frame} gives
\[
 w=\frac35,\quad f_{p,3}=-\frac25,\quad
 R_{0,3}|F_p|^2=\frac{48}{25},\quad
 (V_{1,p}w_p)_3=\frac{36}{25},\quad
 (\operatorname{div}V_{2,pp})_3=-\frac{24}{25}.
 \tag{V46.12}\label{c18v46:canonical}
\]
The second chain gives the opposite third-color values;
the other color values vanish. These equalities can be
checked solely with
$(a_0,\mathbf a)(b_0,\mathbf b)
=(a_0b_0-\mathbf a\cdot\mathbf b,
a_0\mathbf b+b_0\mathbf a+\mathbf a\times\mathbf b)$:
an ordered derivative of a positive link inserts a unit
imaginary quaternion on the left, and a derivative of an
inverse link inserts its negative on the right.
Using this rule through fourth order in
\eqref{c18v46:frame} and \eqref{c18v46:Vjets} proves
\eqref{c18v46:canonical} by rational arithmetic.
For an explicit check on the final divergence contraction,
at this canonical configuration its sum over the three
colors of fine link $U_0$ is
$(0,0,-24/25;0,0,0)$, its sum over those of $U_2$ is
$(0,0,0;0,0,24/25)$, and each of the four other fine-link
color sums is zero. The semicolon separates the two
coarse chains. These are the entries of $\sum_iX_i(V_{2,pp})_i$
before combining the fine-link sums.
The first entry also follows directly from
$|F_p|^2=4(1-w_p^2)$.
The covariance and degree argument then proves
\eqref{c18v46:selftable} for every configuration.
Reversing an incoming two-link chain is a fibre and base
isometry; the vertical divergence transforms covariantly.
This covers all four captured-edge orientations.
Thus the argument does not infer an all-configuration
identity from a numerical sample.

A plaquette with no captured edge has identically zero
self contribution: its single-plaquette flow fixes $b_2$
and acts only on free factors, so its normal lift remains
$R_0$ and $R_0w_p=0$.

\begin{theorem}[The complete centre-even second coefficient]
\label{c18v46:even}
For all $g$, with $\mathcal S$ as in V45, put
\[
 h=\sum_{p\in\mathcal S}w_pf_p,
 \qquad J=\sum_{p\in\mathcal P_N}\chi_p,
 \quad \chi_p=4w_p^2-1.
\]
Then
\[
 \boxed{\mathsf E q_g=-\frac7{288}h
                          =-\frac7{2304}R_0J,\qquad
         \mathsf E(q_g-4af)=\mathsf E q_g.}
 \tag{V46.13}\label{c18v46:evenanswer}
\]
\end{theorem}
\begin{proof}
By \eqref{c18v46:parity}, only self terms survive in
\eqref{c18v46:Q}. Their coefficients, in units of $w_pf_p$,
are respectively
\[
 -\frac1{288}+\frac{40}{288}
                      -\frac6{36}+\frac4{576}=-\frac7{288}.
\]
The first term uses $R_0\chi_p=8w_pf_p$ in
\eqref{c18v46:F2}. Each $w_pf_p$ has conditional mean zero:
integrate a free link and differentiate
$\mathbb E(w_p^2)=1/4$ using the horizontal isometry of V45.
All cross terms have nontrivial character and hence mean
zero after centre averaging. Therefore centering $Q$ does
not change its displayed even part. Finally
$\mathsf E f=0$ proves the second assertion.
In fact the same calculation proves
$\mathbb E_{\nu_{g,0}}Q=0$, so $q_g$ is the restriction
of $Q$ to this fibre. This vanishing is proved, not an
assumption that the conditional normalizer stays fixed.
\end{proof}

\subsection{Native eigenvalue, covariance and unavoidable quartic cost}

The character $\chi_p$ has spin one on each of its four
pair/free coordinates. Thus for a captured plaquette
\[
 \bar L_{g,0}\chi_p=(4+4+8+8)\chi_p=24\chi_p,
 \qquad \bar L_{g,0}h=24h.
 \tag{V46.14}\label{c18v46:adjointmode}
\]
The second equality uses the intertwining horizontal
isometry and $w_pf_p=R_0\chi_p/8$.
It is an identity for the reduced invariant functions,
not a new assertion that the whole reduced gap is 24.

Let $\Sigma_f(g)$ be the full octahedral covariance matrix
(V45.19). Then
\[
 \mathbb E_{\nu_{g,0}}[hh^*]=\frac16\Sigma_f(g),
 \qquad
 \langle h,\bar L_{g,0}^{-1}h\rangle=\frac1{144}\Sigma_f(g).
 \tag{V46.15}\label{c18v46:hcovariance}
\]
For a proof, the root velocity difference for one plaquette
is paired with $\mathbf t/2$ as in V45. On a Haar unit
quaternion,
$\mathbb E[w^2t_at_b]=\delta_{ab}/24$ instead of
$\mathbb E[t_at_b]=\delta_{ab}/4$.
This gives the factor $1/6$. Distinct terms again have a
free edge occurring in only one. On that edge $w_pf_p$
has zero Haar mean by the derivative-of-$1/4$ identity,
so their cross covariance vanishes. Sum over the roots
and use \eqref{c18v46:adjointmode}.

\begin{theorem}[Scalar retiming cannot remove the second source]
\label{c18v46:obstruction}
For every fixed scalar $a$, let
$\mathcal Q_a(g)=\langle q_g-4af,
\bar L_{g,0}^{-1}(q_g-4af)\rangle$, a full-coarse-output
matrix. Then
\[
 \boxed{\mathcal Q_a(g)\succeq
                   \frac{49}{11943936}\Sigma_f(g).}
 \tag{V46.16}\label{c18v46:lower}
\]
In particular, for any one unit coarse-link color column,
its diagonal entry is at least $49/23887872$, independently
of $n,g,a$. Along $s=\beta/24+a\beta^2$ the native
inverse-source matrix has fixed-regulator expansion
\[
 \mathcal R^R(g;\beta,s)=\beta^4\mathcal Q_a(g)+O_n(\beta^5).
 \tag{V46.17}\label{c18v46:quartic}
\]
No $C^2$ scalar retiming with derivative $1/24$ at zero
can make the whole source $o(\beta^2)$.
\end{theorem}
\begin{proof}
The orthogonal projection $\mathsf E$ commutes with the
base conditional generator and its centered inverse.
The even and complementary inverse quadratic forms
therefore add without cross terms. Keeping the even
one gives, by \eqref{c18v46:evenanswer} and
\eqref{c18v46:hcovariance},
$(7/288)^2\Sigma_f/144$, which is
\eqref{c18v46:lower}. No complementary sector was deleted
from the actual source; it contributes a positive
semidefinite matrix to this coefficient.

Each single color column has $\Sigma_f$ diagonal $1/2$:
its two endpoints each have degree four in the
octahedral star and the prefactor is $1/16$.
There is no self-loop for $n\ge5$.
The fixed-regulator centered inverse is smooth near zero,
and \eqref{c18v46:tuned} gives \eqref{c18v46:quartic}.
For a merely $C^2$ time curve the corresponding remainders
are little-oh rather than the displayed cubic/quintic
ones. Its second coefficient still has the nonzero even
part \eqref{c18v46:evenanswer}, proving the last assertion.
\end{proof}

The semidefinite kernel of $\Sigma_f$ remains the full
alternating adjoint-parallel kernel found in V45. The
column lower bounds do not turn \eqref{c18v46:lower} into
a positive multiple of the identity on every combined
coarse direction. These are source-cost coefficients,
not physical or conditional spectral-gap lower bounds.

\subsection{Coefficient locality is uniform; the nonlinear remainder is not yet}

One can prove a uniform upper bound for this second
coefficient without solving the positive-coupling problem.
Make a graph whose vertices are V45's pair/free factors,
joining two factors when they meet in an elementary
plaquette. Each pair factor meets at most eight plaquettes,
and each free factor at most four. A plaquette uses four
distinct factors, so the graph degree is at most 24.

Every component $Q_\alpha$ in \eqref{c18v46:Q} depends only
on the radius-two factor ball about its coarse chain.
Indeed the vacuum term is a differentiated self plaquette
or fine-edge-connected pair. In $HPH$, the right Hessian
starts on the specified chain and the next Hessian can
meet only one of its factors; $P$ joins no different
pair/free factors. In $\nabla_FH$, the force plaquette
must meet a differentiated link of the Hessian plaquette.
The remaining norm and divergence terms obey the same
two-plaquette rule. Differentiating $B,R_0,P$ adds only
the coordinate of its own pair, not a remote plaquette.
Thus no extensive divergence estimate is being used.

There are at most $1+24+24^2=601$ factors in each such
ball. The finite local formulas, whose only denominators
are the displayed numerical constants, imply a finite
bound $|q_{g,\alpha}-4af_\alpha|\le K_a$ independent
of $n,g$. To see that this is a genuine uniform constant,
take the maximum of twice the uncentered suprema over
the finitely many rooted, oriented neighbourhood types
of that bounded size. Each supremum is over compact
products of $SU(2)$ variables and has polynomial
coefficient functions. Periodic identifications only
identify variables in this finite list. Conditional
centering at the product law preserves locality and
at most doubles a component supremum.

Two centered components with disjoint balls are independent
under $\nu_{g,0}$. Their covariance vanishes. For a fixed
column, at most
\[
 D_4=3\sum_{j=0}^4 24^j=1038603
 \tag{V46.18}\label{c18v46:degree}
\]
other color columns can have an intersecting ball.
The scalar row/column sum bound and V45's full unreduced
gap $3/2$ therefore give
\[
 \mathbb E[(q_g-4af)(q_g-4af)^*]\preceq D_4K_a^2I,
 \qquad \mathcal Q_a(g)\preceq\tfrac23D_4K_a^2I.
 \tag{V46.19}\label{c18v46:upper}
\]
The invariant restriction can only improve this comparison
floor. This proves existence of a volume-uniform
\emph{coefficient} bound. It does not bound the
$O_n(\beta^3)$ or $O_n(\beta^5)$ constants, justify
interchanging the thermodynamic and small-coupling
limits, or reach $s=4$.

\subsection{An adjoint map correction and what it leaves unresolved}

The obstruction identifies a direction other than scalar
retiming. Let $\Psi_t$ be the flow of $\nabla J$, with
$J=\sum_p\chi_p$ as above. For a single captured
character put $f_\chi=R_0\chi_p$. Its zero-density,
zero-flow normal-score response is
\[
 \left.\partial_t r^{R,\chi}_{g,0,t}\right|_0=-14f_\chi.
 \tag{V46.20}\label{c18v46:charresponse}
\]
To prove this, V45's changing-excess calculation applies
linearly to any gradient generator. The transported
divergence is $32f_\chi$ because
$\Delta_{\mathcal M}\chi_p=-32\chi_p$.
The same pair connection (V45.13)--(V45.15) gives
$\operatorname{div}_FK_\chi=-f_\chi$: on a used
chain edge $f_{\chi,a}=\pm X_a\chi_p/2$, and the
commutator computation is independent of the spin.
Since $\bar L_{g,0}f_\chi=24f_\chi$,
\[
 \operatorname{div}_F(P\operatorname{Hess}\chi_pR_0)
       =-24f_\chi+f_\chi=-23f_\chi.
\]
Adding the twice-vertical correction gives
$32f_\chi-46f_\chi=-14f_\chi$.
An uncaptured character has zero response, as its flow
fixes the sampler. Every term has zero conditional mean.

Consequently the different, gauge-equivariant sampler
\[
 c_\beta^{\rm adj}
    =b_2\circ\Psi_{-\beta^2/4608}\circ\Phi_{\beta/24}
 \tag{V46.21}\label{c18v46:adjmap}
\]
has the same first cancellation and second normal-source
coefficient
\[
 q_g^{\rm adj}=q_g+\frac7{2304}R_0J,
 \qquad \mathsf E q_g^{\rm adj}=0.
 \tag{V46.22}\label{c18v46:adjcancel}
\]
At this order the composition order of the two small
flows changes nothing, since the second duration is
$O(\beta^2)$. Equation \eqref{c18v46:charresponse} gives
the added term; its sign is positive because this small
adjoint flow is backwards. Both flows are smooth
diffeomorphisms on each compact regulated manifold.
Their sampler remains a full-output submersion. Its
conditional measure and its normal lift must be those
of the new map; they have not been kept frozen.

This is a concrete second-order map correction, not a
proof that the mixed-plaquette source vanishes. In
particular \eqref{c18v46:upper} does not upgrade it to
an exact or convergent trivialization. The next useful
calculation is the retained component
$(I-\mathsf E)q_g$, using the full expression
\eqref{c18v46:Q}. One should determine whether it lies
in the actual normal-score response of local
two-plaquette gradient generators, retaining any
range obstruction and the native metric. This comes
before claims of iteration, volume-uniform remainder,
block-time low-spectrum control or weak-coupling transport.

\subsection{Verification and continuation record}

The script
\texttt{scripts/verify\_ym\_c18\_v46\_tuned\_second\_source.py}
implements literal quaternion word derivatives, the
covariant Hessian derivative and the vertical metric jets.
Exact rational dual-number differentiation at
\eqref{c18v46:canonical} verifies all six divergence
entries through fourth word order. Numerical complex-step
checks independently cover all four captured-edge
orientations and the three contractions in
\eqref{c18v46:selftable}. The written covariance/degree
reduction, not numerical sampling, extends the rational
calculation to all configurations.
Native incidence checks on $n=5,6$ verify distinct
nonempty free-link signatures and the factor degree bound.
The Haar fourth moment, quartic lower coefficient and
adjoint correction constant are checked exactly.
The 24-point rational sphere rule is checked on all 70
monomials through degree four, and tests the actual local
self covariance. The character's $-14f_\chi$ normal response
is separately checked in all four orientations. An
independent short backward-flow variational integration
checks both pushed normal-lift jets against the exact
weighted Gram right inverse; it is a finite local regression,
not a conditional-vacuum calculation.
V45's mathematical regressions are replayed.
These checks are not a proof-assistant certificate or an
interacting conditional simulation.

V46 replaces V45 only; removing this appendix and restoring
its title recovers the full V45 source byte for byte.
The other 53 source files are unchanged. Only
\texttt{.tex} files are kept in \texttt{latex\_notes};
build and diagnostic products are outside that folder.
The next scheduled companion review and broad literature
sweep are V50, and the archive/restart point remains V100.
The mixed second source, a volume-uniform nonlinear
remainder, required block-time low-spectrum control,
physical weak-coupling transport and interacting
continuum reconstruction remain unresolved. Nothing
here establishes the full positive physical YM mass gap.
\addtocontents{toc}{\protect\setcounter{tocdepth}{2}}
% END CYCLE 18 VERSION 46 APPEND
% BEGIN CYCLE 18 VERSION 47 APPEND
\clearpage
\addtocontents{toc}{\protect\vspace{-0.75\baselineskip}}
\addtocontents{toc}{\protect\setcounter{tocdepth}{1}}
\section[V47: chain-mediated sources and a non-gradient correction]{V47: chain-mediated sources and a non-gradient correction}
\addtocontents{toc}{\protect\vspace{-\baselineskip}}
\label{c18v47:context}

V46 removed the entire centre-even second source by an
adjoint-character flow, leaving the distinct-plaquette
terms of its complete normal-source formula. Here we
evaluate and cancel one whole remaining family: two
plaquettes with disjoint fine edges which use opposite
halves of a sampled chain. Their interaction is created
by the normal geometry of the sampler, not by an ambient
second-order vacuum interaction between these plaquettes.
A scalar-gradient response on the local patch cannot
cancel it. A different, explicitly weighted and
gauge-equivariant vector field does cancel it. This
reduces the unpaid second-order calculation to plaquettes
sharing a fine edge.

All statements retain V45--V46's native $SU(2)$ model,
$N=2n$, $n\ge5$, block length two, exact Hamiltonian
vacuum and all coarse color columns. The expansion is
at $\beta=0$ on each fixed regulator. In particular,
the correction below is not a convergent trivialization,
a weak-coupling result, or a proof of the physical
four-dimensional YM mass gap.

The previous V46 turn was progress. Its preservation
and mathematical regressions have been replayed.
The opening inventory contains the same other 53
sources. No new companion hypothesis is needed here.
The primary abstract of Boyle et al.,
\href{https://arxiv.org/abs/2212.11387}{\emph{Use of
Schwinger--Dyson equation in constructing an approximate
trivializing map}} (2022), provides limited methodological
context: one must test the response of an actual chosen
flow-kernel basis. Its approximate lattice construction
is not a theorem about our native normal connection.
No coefficient or convergence claim is imported from it.
The next scheduled companion review and broad literature
sweep remain V50.

\subsection{A general first variation of the normal source}

Write $B=Db_2$, $R_0=B^*/2$ and $P_0=I-R_0B$.
Divergence without a subscript is for ambient product
Haar measure. A divergence of a bank is taken columnwise.
Let $\Pi_g$ subtract the base conditional mean on
$\mathcal F_{g,0}$. The notation $R_0\phi$ means the
bank whose $\alpha$ component is $R_{0,\alpha}\phi$,
not multiplication of a vector field by $\phi$.

\begin{lemma}[Native response to an arbitrary smooth flow]
\label{c18v47:response}
Let $\Xi_t^G$ be the flow of a smooth vector field $G$
on the fine-link manifold, with the sampler changed to
$b_2\circ\Xi_t^G$. At the base density $\Omega_0=1$,
the first centred normal-source response is
\[
 \mathcal D(G)=\Pi_g\left[
   \operatorname{div}\{P_0S_GR_0\}
                  -R_0\operatorname{div}G\right],
 \qquad S_G=\nabla G+(\nabla G)^* .
 \tag{V47.1}\label{c18v47:D}
\]
In particular
\[
 \mathcal D(\nabla\phi)
  =\Pi_g\left[2\operatorname{div}
       (P_0\operatorname{Hess}\phi R_0)+R_0T\phi\right].
 \tag{V47.2}\label{c18v47:Dgrad}
\]
If $BG=0$ identically, then $\mathcal D(G)=0$.
\end{lemma}
\begin{proof}
Use $Y=\Xi_t^G(U)$ to identify the moving fibre with
the fixed fibre of $b_2$. The inverse-flow metric and
pushed log density have first derivatives $-S_G$ and
$-\operatorname{div}G$. Differentiating the orthogonality
equation for the right inverse of $B$ gives
$\widehat R_t=R_0+tP_0S_GR_0+O(t^2)$.
The score in this chart is the divergence of
$\widehat R_t$ plus its application to the pushed log
density. Its base value is
$\operatorname{div}R_0=0$. Consequently differentiation
of the centering measure and change of chart add no
term at first order. This proves \eqref{c18v47:D}.
For $G=\nabla\phi$, use $S_G=2\operatorname{Hess}\phi$
and $\operatorname{div}G=-T\phi$.
If $BG=0$, then $b_2\circ\Xi_t^G=b_2$ exactly, so its
normal source does not change.
\end{proof}

For checking the formula directly in fine invariant
frames, let $M_G=DG+\operatorname{ad}G$ be the
moving-frame differential of the flow, and let $B_G$
denote the directional derivative of the matrix $B$.
Then
\[
 C_1=B_G+BM_G,\qquad
 R_1=\tfrac12P_0C_1^*-R_0C_1R_0 .
 \tag{V47.3}\label{c18v47:originaljet}
\]
Indeed, differentiate $C_t^*(C_tC_t^*)^{-1}$ at
$C_0=B$, using $BB^*=2I$. In original coordinates
$\mathcal D(G)=\Pi_g\operatorname{div}R_1$.
This is a second computational expression for the same
response, not a replacement of the metric by a frozen
one. Its $R_1$ is not the pushed vertical jet in the
proof of the lemma.

\subsection{The native ten-link patch and its complete mixed source}

Call an unordered pair $\{p,q\}$ \emph{chain-mediated}
if its fine edge sets are disjoint but it uses opposite
halves of one sampled chain. Both plaquettes are
captured. Their common chain is denoted $A$, and their
other captured chains $B$ and $C$. The completed local
patch consists of these three two-link chains and four
free links. One representative has sampled products
\[
 g_A=U_0U_1,\qquad g_B=U_2U_3,\qquad g_C=U_4U_5,
\]
and plaquette words
\[
 w_p=\tfrac12\Tr(U_0U_6U_7^{-1}U_2^{-1}),\qquad
 w_q=\tfrac12\Tr(U_1^{-1}U_9^{-1}U_8U_5).
 \tag{V47.4}\label{c18v47:words}
\]
For example, on an $xy$ plane take $U_0,U_1$ on the
$x$ chain from $(0,0)$ to $(2,0)$, $U_2,U_3$ on the
$y$ chain from $(0,0)$ to $(0,2)$, and $U_4,U_5$ on
the $y$ chain from $(2,-2)$ to $(2,0)$.
The remaining links are
$U_6:(1,0)\to(1,1)$,
$U_7:(0,1)\to(1,1)$,
$U_8:(1,-1)\to(2,-1)$ and
$U_9:(1,-1)\to(1,0)$.
The same graph also describes the pairs in different
coordinate planes, with the corresponding edge
relabelings and possible reversals.

Base the two loops at the common midpoint and choose
their orientations so that their last edges are the
two halves pointing towards that midpoint:
\[
 A_p=U_6U_7^{-1}U_2^{-1}U_0,\qquad
 A_q=U_9^{-1}U_8U_5U_1^{-1}.
\]
Define the two gauge-invariant potentials
\[
 \mathcal P=w_pw_q,\qquad
 \mathcal Z=\tfrac12\Tr(A_pA_q^{-1}),\qquad
 u=R_0\mathcal P,\quad v=R_0\mathcal Z .
 \tag{V47.5}\label{c18v47:potentials}
\]
Interchanging $p$ and $q$ leaves both potentials
unchanged. All three coarse triples are retained in
$u,v$. Their other coarse components vanish.

For this paragraph only set
$F_p=\nabla w_p$, $H_p=\operatorname{Hess}w_p$,
$E_p=H_pR_0$, $N_p=BE_p$ and $f_p=R_0w_p$;
define the analogous quantities for $q$. Put
\[
 Y=P_0(H_pP_0H_q+H_qP_0H_p)R_0
   =-P_0(E_pN_q+E_qN_p).
 \tag{V47.6}\label{c18v47:Y}
\]
The equality follows from $H_pH_q=H_qH_p=0$, because
the two fine edge sets are disjoint.

\begin{proposition}[Complete chain-mediated coefficient]
\label{c18v47:mixed}
The contribution of the unordered pair $\{p,q\}$ to
the complete tuned second coefficient of V46 is
\[
 q^{pq}=-\tfrac1{18}(N_pf_q+N_qf_p)
                          +\tfrac1{144}\operatorname{div}Y.
 \tag{V47.7}\label{c18v47:mixedformula}
\]
It has conditional mean zero, and exactly
\[
 q_A^{pq}=\frac{u_A}{576},\qquad
 q_B^{pq}=\frac{u_B}{192},\qquad
 q_C^{pq}=\frac{u_C}{192}.
 \tag{V47.8}\label{c18v47:mixedvalue}
\]
\end{proposition}
\begin{proof}
Polarize \eqref{c18v46:Q} with $W=w_p+w_q$ and retain
one factor from each plaquette. The vacuum coefficient
$F_2$ has no interaction for fine-edge-disjoint
plaquettes. Also $\langle F_p,F_q\rangle=0$ and
$\nabla_{F_p}H_q=\nabla_{F_q}H_p=0$.
The cross part of $V_2$ is $4Y$. The cross part of
$V_1W$ is
$-2(N_pf_q+N_qf_p)$: in $2P_0H_pR_0$ paired
with $F_q$, only $P_0=I-R_0B$ contributes.
These facts give \eqref{c18v47:mixedformula} with
the displayed unordered-pair normalization.

Every remaining expression is fundamental in each of
the four free links. Haar integration in any one of
them gives zero, even after restriction to any base
fibre. Thus centering adds nothing. The full evaluation
\eqref{c18v47:mixedvalue} is included in the exact
finite calculation below, together with a degree and
covariance argument extending it to every configuration.
\end{proof}

\subsection{A response table and its all-configuration proof}

For a scalar function on this patch let
$\nabla_A$ retain only the two components on links
$0,1$, and let $\nabla_E$ retain only the components
on links $2,5$. All other link components are zero.
These projections are orthogonal fine-link projections,
not modifications of the metric. They preserve gauge
equivariance. The unused partners $3,4$ have zero
derivative for both potentials. The free-link part of
either full gradient is vertical, so
\[
 \mathcal D(\nabla\phi)
   =\mathcal D(\nabla_A\phi)+\mathcal D(\nabla_E\phi),
 \qquad \phi\in\{\mathcal P,\mathcal Z\}.
 \tag{V47.9}\label{c18v47:split}
\]

\begin{lemma}[Exact retained-bank response table]
\label{c18v47:tablelemma}
For every configuration and every coarse value,
the response table is
\[
\begin{array}{c|cc}
 G & \mathcal D(G)_A & \mathcal D(G)_e,\ e=B,C\\ \hline
 \nabla_A\mathcal P & -13u_A+v_A & 2u_e+v_e\\
 \nabla_E\mathcal P & 3u_A       & -14u_e+v_e\\
 \nabla_A\mathcal Z & -9v_A      & 6v_e\\
 \nabla_E\mathcal Z & 3v_A       & -10v_e
\end{array}
 \tag{V47.10}\label{c18v47:table}
\]
\end{lemma}
\begin{proof}
Here are both a finite rational certificate and the
reason it determines the full identities. The ten-link
graph has nine vertices and two independent cycles.
A fine gauge transformation puts a spanning tree of
eight links at the identity, leaving $U_6=A$ and $U_8=D$.
This tree contains all six captured links. Write the
two remaining unit quaternions as $A=(a_0,a)$ and
$D=(d_0,d)$. Residual fine gauge transformations act by
common conjugation, that is, by common $SO(3)$ rotation
on $a,d$ and on each coarse color triple.

All expressions in \eqref{c18v47:mixedformula} and
\eqref{c18v47:table} are homogeneous of degree one in
each remaining quaternion. For the mixed expression
this follows because every term contains exactly one
word derivative of each plaquette; for the response
it follows because \eqref{c18v47:originaljet} is linear
in the potential and all its coefficients depend only
on captured links. Word differentiation preserves
this bidegree. Every equivariant vector polynomial of
this bidegree is therefore a linear combination of
\[
 d_0a,\qquad a_0d,\qquad a\times d.
 \tag{V47.11}\label{c18v47:basis}
\]
Indeed the scalar--vector and vector--scalar tensor
products each have one vector summand, and the
vector--vector tensor product has the single vector
summand given by the cross product. The
scalar--scalar product has none.

It suffices to evaluate at
\[
 U_e=1\ (e\ne6,8),\qquad
 U_6=(3/5,4/5,0,0),\quad
 U_8=(5/13,0,12/13,0).
 \tag{V47.12}\label{c18v47:canonical}
\]
The three vectors in \eqref{c18v47:basis} are nonzero
and linearly independent there. Thus their
coefficients in each of the three output triples
are determined by this one exact nine-component
evaluation, not by a numerical extrapolation.

For completeness the finite evaluation is as follows,
with triples separated by semicolons:
\[
 \begin{split}
 65u&=(-10,18,0;\ 10,0,0;\ 0,-18,0),\\
 65v&=(-20,36,-48;\ 10,-18,24;\ 10,-18,24),\\
 37440q^{pq}&=(-10,18,0;\ 30,0,0;\ 0,-54,0).
 \end{split}
 \tag{V47.13}\label{c18v47:certificate1}
\]
The four response columns in the order of the table
are $M/65$, where
\[
 M=
 \begin{pmatrix}
 110&-30&180&-60\\
 -198&54&-324&108\\
 -48&0&432&-144\\
 30&-130&60&-100\\
 -18&-18&-108&180\\
 24&24&144&-240\\
 10&10&60&-100\\
 -54&234&-108&180\\
 24&24&144&-240
 \end{pmatrix}.
 \tag{V47.14}\label{c18v47:certificate2}
\]
One can verify these rational entries directly as
follows. In the fine frame $X_{e,a}U_e=\mathbf i_aU_e$,
replace a positive word factor by
$\mathbf i_aU_e$, and an inverse factor by
$-U_e^{-1}\mathbf i_a$; iterate with the product rule.
Form the component derivative matrix
$(DG)_{ji}=X_iG_j$, add the block
$\operatorname{ad}G$ to obtain $M_G$, and use
\eqref{c18v47:originaljet}. Its divergence is
$\sum_iX_i(R_1)_{i\alpha}$.
For \eqref{c18v47:mixedformula} form the covariant
Hessians
$H_p=DF_p+\frac12\operatorname{ad}F_p$ and
$H_q=DF_q+\frac12\operatorname{ad}F_q$, then
\eqref{c18v47:Y}, and again sum the diagonal
directional derivatives. Here the adjoint block is
twice the usual cross-product matrix. These rules
use rational arithmetic and word derivatives only
through order three; the disjoint fine supports
removed the fourth-order term. They produce
\eqref{c18v47:certificate1}--\eqref{c18v47:certificate2}.
Substituting the first two rows of
\eqref{c18v47:certificate1} in the claimed table
reproduces every entry of $M/65$.

The equivariant bidegree argument then proves both
\eqref{c18v47:mixedvalue} and \eqref{c18v47:table}
on the gauge-fixed graph and hence on all configurations.
This use of fine gauge covariance does not restrict
the coarse value $g$ to the identity: the identity tree
is a device for verifying a covariant polynomial
identity on the ambient configuration space.
Link relabeling and reversal are product-metric
isometries and give all native patch orientations,
including the induced inversion of a coarse product.
\end{proof}

\subsection{Why a local scalar potential is insufficient}

\begin{proposition}[A local scalar-gradient obstruction]
\label{c18v47:obstruction}
No linear combination of $\nabla\mathcal P$ and
$\nabla\mathcal Z$ has response $-q^{pq}$ as a full
coarse bank. More strongly, no smooth
full-fine-gauge-invariant scalar potential supported
on this ten-link graph has that response under its
ordinary gradient. This is a local-patch statement,
not a prohibition on larger-support or nonlocal
scalar constructions.
\end{proposition}
\begin{proof}
For $\phi=a\mathcal P+b\mathcal Z$, the two responses
from \eqref{c18v47:table} are
\[
 \mathcal D(\nabla\phi)_A
       =-10a\,u_A+(a-6b)v_A,\qquad
 \mathcal D(\nabla\phi)_e
       =-12a\,u_e+(2a-4b)v_e .
 \tag{V47.15}\label{c18v47:scalar}
\]
At \eqref{c18v47:canonical}, $u_A,v_A$ are independent,
as are $u_B,v_B$. Cancellation would require
$a-6b=2a-4b=0$, hence $a=b=0$, whereas $q^{pq}\ne0$.
Equivalently, the two scalar response columns have
rank two and adjoining the source raises the rank
to three.

To extend this conclusion to a smooth scalar on the
same graph, project it onto spin $1/2$ on each of its
four free links. The operator
$\phi\mapsto\mathcal D(\nabla\phi)$ commutes with each
such Peter--Weyl projection: its coefficients involve
only captured links, and the remaining differentiations
are invariant derivatives, which commute with the
free-link Casimir. Conditional centering also commutes
with these projections; the mean is independent of
free links and the projected sector has zero mean.
The source is already in this joint fundamental sector.
Thus any proposed cancellation would persist after
projection of the potential. Projection preserves
fine gauge invariance.

Fix the same spanning tree. A gauge-invariant function
on this graph is a function of its two holonomies,
invariant under common conjugation. Fundamental
dependence on a free edge of each loop restricts the
function to fundamental matrix coefficients in each
holonomy. Under common conjugation the fundamental
matrix-coefficient space is $0\oplus1$. Its tensor
square has precisely two scalar summands, namely
$a_0d_0$ and $a\cdot d$. These are spanned by
$\mathcal P$ and $\mathcal Z$. The projected potential
therefore has exactly the form already excluded.
\end{proof}

This obstruction concerns the actual normal-score
response. It does not follow merely from orthogonality
of distinct plaquette words, and it does not assert
that the Hamiltonian has an obstruction to a mass gap.

\subsection{An explicit non-gradient repair}

\begin{theorem}[Cancellation of a chain-mediated pair]
\label{c18v47:repair}
Put $D_*=3649536$ and define the gauge-equivariant
vector field
\[
 G_{pq}=\frac1{D_*}\left(
   828\nabla_A\mathcal P+1476\nabla_E\mathcal P
       +211\nabla_A\mathcal Z+357\nabla_E\mathcal Z
                      \right).
 \tag{V47.16}\label{c18v47:generator}
\]
Then $\mathcal D(G_{pq})=-q^{pq}$ on every base fibre,
with every coarse column retained.
\end{theorem}
\begin{proof}
The coefficients of $u_A,u_e,v_A,v_e$ in the numerator
of the response, respectively, are
\[
 \begin{gathered}
 -13(828)+3(1476)=-6336=-D_*/576,\\
 2(828)-14(1476)=-19008=-D_*/192,\\
 828-9(211)+3(357)=0,\\
 828+1476+6(211)-10(357)=0.
 \end{gathered}
\]
Apply \eqref{c18v47:table}. Each link projection
commutes with gauge transformations, so the vector
field is gauge equivariant. No conditional inverse
or unproved Poisson estimate is used.
\end{proof}

The unequal link weights are essential to this repair.
In particular \eqref{c18v47:generator} has no
gauge-invariant scalar-gradient representative modulo
a vertical field supported on this patch, by
Proposition~\ref{c18v47:obstruction} and
$\mathcal D(G_{\rm vertical})=0$. This is the precise
sense in which the local direction of the programme
has changed.

There is also a simple two-mode description of the
base conditional dependence. In pair coordinates
$(U_0,U_1)=(x,x^{-1}g_A)$, the word $\mathcal Z$
does not depend on $x$ and
$\mathbb E_x\mathcal P=\mathcal Z/4$.
The last identity follows from
$\mathbb E_{x\in S^3}x_ix_j=\delta_{ij}/4$, or
fundamental matrix-coefficient orthogonality.
The other two pair factors contribute $3$ in total
to the conditional Casimir, and the four free factors
contribute $12$. The remaining shared-pair tensor
product has spins zero and one, with eigenvalues
zero and four in its doubled metric. Hence
\[
 L_0\mathcal Z=15\mathcal Z,\qquad
 L_0\mathcal P=19\mathcal P-\mathcal Z .
 \tag{V47.17}\label{c18v47:modes}
\]
The isometric $R_0$ transports commute with $L_0$, so
\[
 L_0u=19u-v,\quad L_0v=15v,\quad
 L_0^{-1}u=\frac{u}{19}+\frac{v}{285}.
 \tag{V47.18}\label{c18v47:inverse}
\]
These are identities in this finite source sector.
The inverse here is the mean-zero base inverse, not a
positive-coupling conditional gap estimate. They are
not needed to implement the cancellation.

\subsection{All-volume superposition and the remaining pair family}

\begin{lemma}[Native incidence and exhaustive second-order split]
\label{c18v47:incidence}
There are exactly $36n^3$ chain-mediated unordered
pairs. Each belongs to a unique common sampled chain,
and each captured plaquette participates in six.
Each fine link occurs in the fine support of at most
24 such pairs; each captured fine link is active in
exactly 24 of their generators. After the self terms,
every possible second-order mixed term in
\eqref{c18v46:Q} is either fine-edge-sharing or
chain-mediated.
\end{lemma}
\begin{proof}
Each half of a sampled chain belongs to four
plaquettes. Among the $4\cdot4$ choices across its
two halves, four use the same transverse direction
and share the transverse fine edge at the midpoint.
The remaining twelve are fine-edge-disjoint.
The native block-two parity pattern gives exactly
one other captured chain per plaquette. For these
twelve choices the other chains are distinct.
Thus the common chain is unique and its completion
has three two-link chains and four free links.
There are $3n^3$ sampled chains, proving the count.
A captured plaquette has two captured edges; across
either corresponding chain it has three disjoint
partners. This proves the six-partner count.
Every fine link lies in four plaquettes, giving the
upper incidence bound $4\cdot6=24$. If the link is
captured, all four plaquettes are captured, and all
these pair memberships are distinct and active.

Finally the fibre projection $P_0$ couples only the
two links of one sampled chain; its free blocks are
identities. Covariant word derivatives have support
only on their plaquette fine edges. Therefore a
mixed contraction involving two different
plaquettes vanishes unless they share a fine edge
or a pair factor. The latter case, when there is no
common fine edge, is precisely the chain-mediated
case. The native vacuum $F_2$ itself has only
fine-edge-sharing cross terms. This exhausts the
second coefficient and does not omit distant pairs.
\end{proof}

Let $G_{\rm ch}=\sum_{\{p,q\}\ {\rm chain-mediated}}G_{pq}$
and denote its flow by $\Xi_t$.
Keep V46's $J=\sum_p(4w_p^2-1)$, with flow
$\Psi_t^J$, and $W=\sum_pw_p$, with flow $\Phi_s^W$.
Define the new sampler
\[
 c_\beta^\sharp
   =b_2\circ\Xi_{\beta^2}\circ
          \Psi^J_{-\beta^2/4608}\circ\Phi^W_{\beta/24}.
 \tag{V47.19}\label{c18v47:globalmap}
\]
All these flows are complete on every finite compact
regulated manifold. By linearity of the first response,
the two order-$\beta^2$ corrections add at this order;
their interaction with the order-$\beta$ flow affects
only higher orders. Theorem~\ref{c18v47:repair} and
V46's adjoint cancellation therefore give
\[
 r^\sharp_{g,\beta}
   =\beta^2\!\!\sum_{\substack{\{p,q\}\\
                   p,q\ {\rm share\ a\ fine\ edge}}}
                 q_g^{pq}+O_n(\beta^3).
 \tag{V47.20}\label{c18v47:remaining}
\]
Here $q_g^{pq}$ is the corresponding polarized
coefficient of \eqref{c18v46:Q}; it includes the native
vacuum contribution for that pair. The conditional
measure, centering and normal lift in the left side
are those of the new sampler. Formula
\eqref{c18v47:remaining} is not a claim that this last
pair family vanishes, or that its expression is the
same as \eqref{c18v47:mixedformula}.

\subsection{A uniform geometric bound, with its limitation}

\begin{proposition}[Full-output Gram stability of the corrected map]
\label{c18v47:gram}
In the native product metrics, for all configurations
and all finite $n\ge5$,
\[
 2e^{-2\Lambda_\beta}I\ \preceq\
 Dc_\beta^\sharp(Dc_\beta^\sharp)^*
       \ \preceq\ 2e^{2\Lambda_\beta}I,\qquad
 \Lambda_\beta=\frac83|\beta|+\frac{749}{1408}\beta^2.
 \tag{V47.21}\label{c18v47:grambound}
\]
In particular the normal right inverse has norm at
most $e^{\Lambda_\beta}/\sqrt2$.
\end{proposition}
\begin{proof}
For $\mathcal P$ or $\mathcal Z$, any ordered
invariant derivative has absolute value at most one.
For $\mathcal Z$ this is a half-trace word with eight
distinct fine links and unit quaternion insertions.
For $\mathcal P$, each fine link belongs to only one
of its two factors, so each ordered derivative is
again a product of two such bounded word derivatives.

Set
\[
 \gamma=\frac{1833}{3649536}
        =\frac{611}{1216512}.
\]
The absolute value of each component of one generator,
and of each of its component derivatives, is at most
$\gamma$, since the larger sum of its two relevant
weights is $1476+357=1833$.
A component of $G_{\rm ch}$ is active in at most 24
patches. A patch depends on eight links (24 fine
color directions) and acts on four links (12 directions).
The row and column absolute sums of $DG_{\rm ch}$
are consequently at most $576\gamma$ and $288\gamma$.
The adjoint block has row and column sums at most
$96\gamma$: its two off-diagonal entries have the
factor two from the $SU(2)$ bracket. It follows that
\[
 \|M_{G_{\rm ch}}\|_2
  \le\sqrt{672\gamma\cdot384\gamma}
  \le672\gamma=\frac{4277}{12672}.
 \tag{V47.22}\label{c18v47:MG}
\]
All incidence constants are independent of volume.

Conservative bounds for the other two moving-frame
differentials suffice. Each link lies in four
plaquettes, each with 12 color slots, so
$\|M_{\nabla W}\|_2\le48+16=64$.
For $\chi_p=4w_p^2-1$ the first and second ordered
derivatives are bounded by 8 and 16. This gives
$\|M_{\nabla J}\|_2\le768+128=896$.
The variational equation for each flow, and for its
inverse, therefore bounds the norms of the differential
and its inverse by the exponential of the duration
times the corresponding constant. The composed fine
diffeomorphism in \eqref{c18v47:globalmap} has both
norms bounded by
\[
 \exp\left\{\frac{64}{24}|\beta|
       +\left(\frac{896}{4608}
                       +\frac{4277}{12672}\right)\beta^2\right\}
       =e^{\Lambda_\beta}.
\]
At its final point $B B^*=2I$. Apply the two-sided
singular-value bound to $Dc_\beta^\sharp=B\,D\mathcal F_\beta$
to obtain \eqref{c18v47:grambound}.
\end{proof}

This is a genuine volume-independent estimate for the
corrected sampler's geometry and full-output rank.
It is not a bound for the conditional vacuum density,
a uniform nonlinear remainder in
\eqref{c18v47:remaining}, or a low-frequency conditional
source inverse. The short map deformation has still
not controlled required block time four, continuum
scale transport, or a positive physical mass.

\subsection{Verification record and the next upstream obligation}

The reproducible diagnostic is
\begin{center}\small
\path{scripts/verify_ym_c18_v47_chain_mixed_source.py}.
\end{center}
It differentiates the literal ten-link words and the
native normal right inverse. Exact rational dual
arithmetic checks all nine source components and all
four grouped response columns at
\eqref{c18v47:canonical}, with nine exactly zero
correction residuals. Exact linear algebra checks the
weighted response rank four, scalar rank two and
augmented scalar rank three. Separate complex-step
regressions check the full formulas at the canonical
point and three nonflat random configurations; these
are regressions, not the all-configuration proof.
Native incidence enumerations at $n=5,6$ give 4500
and 7776 pairs, twelve per chain, completed patch size
ten and link load at most 24. The written incidence
argument establishes the statement for all $n\ge5$.
V46's mathematical diagnostics were replayed before
this appendix was written. None of these checks is
a proof-assistant certificate or an interacting
continuum simulation.

The completed work should not be repeated in V48:
the centre-even self source and the chain-mediated
source now have explicit native map corrections.
The next unpaid local task is the fine-edge-sharing
family in \eqref{c18v47:remaining}, including its
actual $F_2$ contribution and all coarse columns.
The present scalar obstruction is a reason to test
the native response of link-weighted local fields
there, not to assume that the same weights or
coefficients will work. A volume-uniform nonlinear
remainder remains a distinct subsequent obligation.
Then come required block-time conditional low
frequencies, weak-coupling transport and interacting
continuum reconstruction.

V47 replaces V46 only. Removing this appendix and
restoring the title recovers the complete V46 source
byte for byte; the other 53 sources remain unchanged.
Only \texttt{.tex} files are kept in
\texttt{latex\_notes}. Build and diagnostic products
remain outside it. The next companion and broad
literature checkpoints are V50 and the archive/restart
point is V100. The full positive physical Yang--Mills
mass-gap goal remains open.
\addtocontents{toc}{\protect\setcounter{tocdepth}{2}}
% END CYCLE 18 VERSION 47 APPEND
% BEGIN CYCLE 18 VERSION 48 APPEND
\clearpage
\addtocontents{toc}{\protect\enlargethispage{\baselineskip}}
\addtocontents{toc}{\protect\setcounter{tocdepth}{1}}
\section[V48: cancellation of the complete second normal source]{V48: cancellation of the complete second normal source}
\addtocontents{toc}{\protect\vspace{-1.25\baselineskip}}
\label{c18v48:context}

V47 left precisely the fine-edge-sharing terms of V46's
complete second normal source. This note evaluates all of
them, including the native vacuum contribution, and supplies
explicit local map corrections. Together with the already
proved self and chain-mediated corrections, they cancel
the entire second coefficient. The corrected normal source
is $O_n(\beta^3)$ and its conditional inverse quadratic form
is $O_n(\beta^6)$ at each fixed regulator. A separate
full-output geometric estimate is uniform in volume.
The nonlinear remainder and physical continuum mass gap
are not thereby proved uniform or closed.

We retain $SU(2)$, $N=2n$, $n\ge5$, block length two,
the exact Hamiltonian vacuum and every coarse color column.
Here $T=-\sum_e\Delta_e$ is the dimensionless electric
operator and $\beta=b/\kappa$; $T$ is not a duration symbol
in this appendix. Normal lifts and conditional operators
are those of the actual changing sampler in the native
metric. No new Hamiltonian or substitute Gibbs law is used.

The previous V47 turn was progress; its preservation and
mathematical checks have been replayed. All 53 other
sources are unchanged at the opening inventory.
V36's vacuum coefficient is reused, not claimed as new;
its fixed-map correction differs from the present normal
connection. A targeted recheck of the primary abstracts of
\href{https://arxiv.org/abs/0907.5491}{L\"uscher (2009)}
and \href{https://arxiv.org/abs/2212.11387}{Boyle et al.\ (2022)}
retains the existing flow-construction context only.
No coefficient or exact trivialization theorem is imported.
The next companion review and broad literature sweep are V50.

\subsection{A curvature identity removes the fourth word derivative}

Keep $B=Db_2$, $R_0=B^*/2$, $P_0=I-R_0B$ and V47's
native response $\mathcal D$. Let an unordered pair
$p\ne q$ share exactly one fine edge $e$. As in V36,
orient its two words at that edge as $UA$ and $UD$.
Set
\[
 \mathcal P=w_pw_q,\quad
 \mathcal Z=\tfrac12\Tr(AD^{-1}),\quad
 C=\langle F_p,F_q\rangle=\mathcal Z-\mathcal P,\quad
 F_p=\nabla w_p,\quad H_p=\operatorname{Hess}w_p .
 \tag{V48.1}\label{c18v48:C}
\]
The six-edge word $\mathcal Z$ is precisely V36's $R_{pq}$.
A plaquette can be reversed before this orientation choice,
since its $SU(2)$ half-trace is unchanged.
Let $H_C=\operatorname{Hess}C$ and let $K_{pq}$ vanish
off the common link, with block there
\[
 (K_{pq})_e=2CI_3-F_{p,e}\otimes F_{q,e}
                            -F_{q,e}\otimes F_{p,e}.
 \tag{V48.2}\label{c18v48:K}
\]

\begin{lemma}[Covariant Hessian product identity]
\label{c18v48:curvature}
On the product of unit three-spheres,
\[
 \nabla_{F_p}H_q+\nabla_{F_q}H_p
       =H_C-H_pH_q-H_qH_p-K_{pq}.
 \tag{V48.3}\label{c18v48:codazzi}
\]
The polarized cross part of V46's pushed second normal
jet is consequently
\[
 Z_2^{pq}=P_0\left[
 4(H_pP_0H_q+H_qP_0H_p)-H_pH_q-H_qH_p-H_C+K_{pq}
                          \right]R_0 .
 \tag{V48.4}\label{c18v48:Z2}
\]
\end{lemma}
\begin{proof}
The gradient of $C$ is $H_pF_q+H_qF_p$.
The Hessian commutation identity gives, for tangent $v$,
\[
 (\nabla_vH_p)F_q=(\nabla_{F_q}H_p)v+
                                      \mathcal R(v,F_q)F_p .
\]
Differentiate the gradient of $C$ and use the corresponding
identity with $p,q$ exchanged. Our sphere-curvature
convention is
$\mathcal R(v,x)y=\langle x,y\rangle v-\langle v,y\rangle x$.
Product curvature acts separately on fine links, and both
gradients are nonzero together only on $e$. Its two terms
give exactly $K_{pq}v$. This proves
\eqref{c18v48:codazzi}. Polarize
$P_0(4HP_0H-2H^2-\nabla_FH)R_0$ and substitute that
identity to obtain \eqref{c18v48:Z2}.
\end{proof}

This retains sphere curvature, not a flat-field-space
approximation. The divergence of \eqref{c18v48:Z2}
now requires word derivatives only through order three.
Put
\[
 u=R_0\mathcal P,\quad v=R_0\mathcal Z,\quad
 f_p=R_0w_p,\quad N_p=BH_pR_0,\quad
 \mathcal N=N_pf_q+N_qf_p .
\]
As before, $R_0\phi$ denotes its bank of directional
derivatives.

\begin{lemma}[Complete shared-edge pair formula]
\label{c18v48:pairformula}
The unordered pair contributes
\[
 q^{pq}=\frac{125}{5616}v-\frac{41}{1872}u
            -\frac{\mathcal N}{18}
                 +\frac1{576}\operatorname{div}Z_2^{pq}
 \tag{V48.5}\label{c18v48:Q}
\]
to the complete tuned second source. Each nonzero case
below has zero base conditional mean for every $g$.
\end{lemma}
\begin{proof}
The native $F_2$ contribution is $v/702-u/936$.
The cross part of $R_0|\nabla W|^2$ is $2(v-u)$.
The cross part of $V_1W$, with $V_1=2P_0HR_0$, is
$2(v-u)-2\mathcal N$. In \eqref{c18v46:Q} their combined
coefficient of $v-u$ is $-5/144+1/18=1/48$.
Adding \eqref{c18v48:Z2} proves \eqref{c18v48:Q}
with the actual vacuum input and unordered normalization.
In each nonzero case there is a free fine link belonging
to only one plaquette. Every expression is fundamental
in that free link, while normal-matrix coefficients use
only captured links. Haar integration in this variable
annihilates it, so conditional centering adds nothing.
\end{proof}

\subsection{An exhaustive native classification}

A fine link is captured exactly when both transverse
coordinates are even. A plaquette is captured exactly
when its coordinate perpendicular to its plane is even;
it then has two captured fine links, and otherwise none.
The shared-edge pairs are
\[
\begin{array}{c|c|c|c}
\text{type}&\text{shared link}&\text{captured plaquettes}
   &(\text{completed chains},\text{fine links})\\ \hline
{\rm I}&\text{captured}&2&(3,10)\\
{\rm II}&\text{free}&2&(3,9)\\
{\rm III}&\text{free}&1&(2,9)\\
{\rm 0}&\text{free}&0&(0,7).
\end{array}
 \tag{V48.6}\label{c18v48:types}
\]
Completion includes both halves of each touched sampled
chain. Type 0 gives zero: both potentials and all word
derivatives have only free support, so
$H_pR_0=H_qR_0=H_CR_0=K_{pq}R_0=0$, and its entire
normal bank vanishes.

Here are representatives. In I and II the sampled
products are $U_0U_1,U_2U_3,U_4U_5$; in III they are
$U_0U_1,U_2U_3$. The displayed words have half-traces
$w_p,w_q$; all other links are free:
\[
\begin{array}{c|cc}
 &p&q\\ \hline
{\rm I}&U_0U_6U_7^{-1}U_2^{-1}
        &U_0U_8U_9^{-1}U_4^{-1}\\
{\rm II}&U_6U_7^{-1}U_2^{-1}U_0
         &U_6U_8U_4^{-1}U_1^{-1}\\
{\rm III}&U_4U_5^{-1}U_2^{-1}U_0
          &U_4U_6U_7^{-1}U_8^{-1}.
\end{array}
 \tag{V48.7}\label{c18v48:words}
\]
In II, the two halves of the first sampled chain occur
in different plaquettes. Its coarse triple is $A$;
the others are $B,C$, collectively $E$.
The shared fine edge is free, not that sampled chain.
This distinguishes II from V47's fine-edge-disjoint
chain-mediated pair.

\begin{proposition}[Evaluated complete pair sources]
\label{c18v48:evaluated}
For all configurations, coarse values and retained columns,
\[
\begin{array}{c|l}
\text{case}&q^{pq}\\ \hline
{\rm I}&-\dfrac{133}{14976}u+\dfrac{149}{22464}v\\[10pt]
{\rm II},A&-\dfrac{29}{14976}u_A+\dfrac{649}{44928}v_A\\[10pt]
{\rm II},e=B,C&\dfrac{101}{14976}u_e+\dfrac{383}{22464}v_e\\[10pt]
{\rm III}&\dfrac{19}{2496}u-\dfrac{23}{11232}v .
\end{array}
 \tag{V48.8}\label{c18v48:sources}
\]
\end{proposition}
\begin{proof}
The exact contractions are
\[
\begin{array}{c|cc}
\text{case}&\mathcal N&\operatorname{div}Z_2^{pq}\\ \hline
{\rm I}&(v-u)/2&-17u/2+7v\\
{\rm II},A&-u_A/2&-9(u_A+v_A)/2\\
{\rm II},e=B,C&(u_e-v_e)/2&65u_e/2-19v_e\\
{\rm III}&0&17u-14v .
\end{array}
 \tag{V48.9}\label{c18v48:contractions}
\]
Their finite all-configuration proof follows below.
Insert this table in \eqref{c18v48:Q}. For example,
the III coefficients are
$-41/1872+17/576=19/2496$ and
$125/5616-14/576=-23/11232$.
\end{proof}

\subsection{The actual response range and its finite proof}

Let $\nabla_{\rm cap}\phi$ retain only captured fine
gradient components in the pair's support.
The omitted free part is vertical and has zero
$\mathcal D$ response. In I and III,
\[
\begin{array}{c|cc}
 &\mathcal D(\nabla_{\rm cap}\mathcal P)
 &\mathcal D(\nabla_{\rm cap}\mathcal Z)\\ \hline
{\rm I}&-10u&-10v\\
{\rm III}&-18u+2v&-10v.
\end{array}
 \tag{V48.10}\label{c18v48:scalarresponses}
\]
These are also the responses of the full scalar gradients.

For II let $\nabla_A$ retain links $0,1$, and
$\nabla_E$ links $2,4$. The partners $3,5$ have zero
derivatives for both potentials. The table is
\[
\begin{array}{c|cc}
G&\mathcal D(G)_A&\mathcal D(G)_e,\ e=B,C\\ \hline
\nabla_A\mathcal P&-15u_A+3v_A&2u_e+v_e\\
\nabla_E\mathcal P&3u_A&-16u_e+3v_e\\
\nabla_A\mathcal Z&-3v_A&6v_e\\
\nabla_E\mathcal Z&3v_A&-4v_e .
\end{array}
 \tag{V48.11}\label{c18v48:IIresponse}
\]
This is not V47's response table: the common free
edge changes both the source and response values.

\begin{lemma}[Finite proof of the contraction and response tables]
\label{c18v48:finiteproof}
Equations \eqref{c18v48:contractions}--\eqref{c18v48:IIresponse}
hold on all native embeddings of the three graphs.
\end{lemma}
\begin{proof}
The completed graphs have nine, eight and eight vertices,
respectively, and two independent cycles.
Their captured subgraphs have no cycles. Complete
each to a spanning tree and gauge-fix that tree to
the identity. In I leave $U_6=A,U_8=D$; in II leave
$U_7=A^{-1},U_8=D$; in III leave
$U_5=A^{-1},U_7=D^{-1}$. Every other link, including
every captured link, is the identity. The plaquette
holonomies are then $A,D$.

All expressions in the proposed tables are homogeneous
of degree one in each remaining quaternion.
This is immediate for the potentials, for $F_pF_q$
and its differentiated contractions, and for
\eqref{c18v48:Z2}. The normal-response operator is
linear in the potential with matrix coefficients
depending only on captured links. It preserves this
bidegree. Residual fine gauge transformations act
by common conjugation on $A,D$ and by the adjoint
action on each output triple. With
$A=(a_0,a)$, $D=(d_0,d)$, the equivariant vector
bidegree space is exactly
\[
 \operatorname{span}\{d_0a,\ a_0d,\ a\times d\}.
 \tag{V48.12}\label{c18v48:degree}
\]
These are the scalar--vector, vector--scalar and
antisymmetric vector--vector summands; no other vector
summand exists. Evaluate at
\[
 A=(3/5,4/5,0,0),\qquad D=(5/13,0,12/13,0),
 \tag{V48.13}\label{c18v48:canonical}
\]
where the three vectors in \eqref{c18v48:degree}
are linearly independent. The banks, separated
into coarse triples by semicolons, are
\[
\begin{array}{c|l}
 &65u\\ \hline
{\rm I}&(-10,-18,0;\ 10,0,0;\ 0,18,0)\\
{\rm II}&(-10,18,0;\ 10,0,0;\ 0,18,0)\\
{\rm III}&(-10,0,0;\ 10,0,0)
\end{array}
 \tag{V48.14}\label{c18v48:ucert}
\]
and
\[
\begin{array}{c|l}
 &65v\\ \hline
{\rm I}&(0,0,0;\ 10,-18,24;\ -10,18,-24)\\
{\rm II}&(-20,36,-48;\ 10,-18,24;\ -10,18,-24)\\
{\rm III}&(-10,18,-24;\ 10,-18,24).
\end{array}
 \tag{V48.15}\label{c18v48:vcert}
\]
Literal differentiation gives
\[
\begin{array}{c|l}
 &65\operatorname{div}Z_2^{pq}\\ \hline
{\rm I}&(85,153,0;\ -15,-126,168;\ -70,-27,-168)\\
{\rm II}&(135,-243,216;\ 135,342,-456;\ 190,243,456)\\
{\rm III}&(-30,-252,336;\ 30,252,-336).
\end{array}
 \tag{V48.16}\label{c18v48:divcert}
\]
The first and second word derivatives in
$N_pf_q+N_qf_p$ give the middle column of
\eqref{c18v48:contractions} at these same points.
For II the four response columns in
\eqref{c18v48:IIresponse} are $1/65$ times
\[
 \begin{pmatrix}
 90&-30&60&-60\\
 -162&54&-108&108\\
 -144&0&144&-144\\
 30&-130&60&-40\\
 -18&-54&-108&72\\
 24&72&144&-96\\
 -10&-30&-60&40\\
 54&-234&108&-72\\
 -24&-72&-144&96
 \end{pmatrix}.
 \tag{V48.17}\label{c18v48:responsecert}
\]
For I the summed scalar response columns are
$-10u,-10v$ at \eqref{c18v48:ucert}--\eqref{c18v48:vcert};
for III they are $-18u+2v,-10v$.

Here are explicit rules for checking the finite arithmetic.
In the fine frame $X_{e,a}U_e=\mathbf i_aU_e$,
differentiate a positive word factor by inserting
$\mathbf i_a$ on its left and an inverse factor by
inserting $-\mathbf i_a$ on its right. For a gradient
$F$, form $H=DF+\frac12\operatorname{ad}F$,
where $(DF)_{ji}=X_iF_j$ and the adjoint block is
twice the cross-product matrix. Use \eqref{c18v48:Z2}
and sum $X_i(Z_2^{pq})_{i\alpha}$ for its divergence.
For a link-projected gradient $G$, use
\[
 C_1=B_G+B(DG+\operatorname{ad}G),\qquad
 R_1=\tfrac12P_0C_1^*-R_0C_1R_0,\qquad
 \mathcal D(G)_\alpha=\sum_iX_i(R_1)_{i\alpha}.
\]
This is \eqref{c18v47:originaljet} with zero centering
in the present sector. Differentiating the literal
words \eqref{c18v48:words} by these rules gives
\eqref{c18v48:divcert}--\eqref{c18v48:responsecert}
and the scalar columns just stated. Every operand
at \eqref{c18v48:canonical} is rational.
Independent covariant differentiation of $H_p,H_q$
also verifies \eqref{c18v48:codazzi} there.

Equality at this exact point determines all coefficients
in \eqref{c18v48:degree}, proving the identities for
arbitrary $A,D$. Fine gauge covariance proves them
before tree fixing, for every coarse value, not only
the identity coarse fibre. Product-metric link
reversals and relabelings give every native embedding.
The argument is a finite polynomial proof, not an
extrapolation from numerical samples.
\end{proof}

\subsection{Explicit corrections and a local scalar obstruction}

\begin{theorem}[Cancellation for every shared-edge pair]
\label{c18v48:corrections}
The following smooth gauge-equivariant fields satisfy
$\mathcal D(G_{pq})=-q^{pq}$:
\[
 \begin{split}
 G_{pq}^{\rm I}
  &=-\frac{133}{149760}\nabla_{\rm cap}\mathcal P
             +\frac{149}{224640}\nabla_{\rm cap}\mathcal Z,\\
 G_{pq}^{\rm III}
  &=\frac{19}{44928}\nabla_{\rm cap}\mathcal P
             -\frac1{8320}\nabla_{\rm cap}\mathcal Z,\\
 G_{pq}^{\rm II}
  &=\frac{-161\nabla_A\mathcal P+1457\nabla_E\mathcal P
        -65405\nabla_A\mathcal Z-82118\nabla_E\mathcal Z}
           {3504384}.
 \end{split}
 \tag{V48.18}\label{c18v48:generators}
\]
For type 0 take $G_{pq}=0$.
\end{theorem}
\begin{proof}
I and III follow from \eqref{c18v48:scalarresponses}.
For III the response is
$-(19/2496)u+(23/11232)v$; for I it is
$(133/14976)u-(149/22464)v$.
For II put $D_*=3504384$. The numerator's coefficients
of $u_A,u_e,v_A,v_e$, respectively, are
\[
 \begin{gathered}
 -15(-161)+3(1457)=6786=D_*\,29/14976,\\
 2(-161)-16(1457)=-23634=-D_*\,101/14976,\\
 3(-161)-3(-65405)+3(-82118)=-50622=-D_*\,649/44928,\\
 -161+3(1457)+6(-65405)-4(-82118)
                         =-59748=-D_*\,383/22464.
 \end{gathered}
\]
Use \eqref{c18v48:IIresponse} and compare
\eqref{c18v48:sources}. Fine-link projections commute
with gauge transformations. The potentials and selected
link groups are intrinsic to the unordered embedded
pair, so the field is independent of word orientation
conventions.
\end{proof}

I and III could instead use full scalar gradients
for this first response, since the omitted free parts
are vertical. II cannot.

\begin{proposition}[Type II is not a local scalar-gradient response]
\label{c18v48:scalarobstruction}
No smooth full-fine-gauge-invariant scalar potential
supported on the completed type II graph has ordinary
gradient response $-q^{pq}$ as the full retained bank.
\end{proposition}
\begin{proof}
For $\phi=a\mathcal P+b\mathcal Z$, summing the
response table gives
\[
 \mathcal D(\nabla\phi)_A=a(-12u_A+3v_A),\qquad
 \mathcal D(\nabla\phi)_e=-14a\,u_e+(4a+2b)v_e .
\]
At \eqref{c18v48:canonical}, $u_A,v_A$ are independent.
The $A$ triple alone would require
\[
 -12a=\frac{29}{14976},\qquad
 3a=-\frac{649}{44928},
\]
which are incompatible.

For a general scalar on this graph, project onto
fundamental matrix coefficients in the unshared free
chord links $7,8$. The response operator commutes
with these Peter--Weyl projections: its coefficients
depend only on captured links, and its differentiations
are invariant derivatives. Conditional centering
also commutes with them. The source already lies
in this sector, and projection preserves fine gauge
invariance. After the same tree fixing, the projected
potential is a common-conjugation-invariant fundamental
matrix coefficient in each holonomy $A,D$.
The two scalar summands in
$(0\oplus1)\otimes(0\oplus1)$ are $a_0d_0$ and $a\cdot d$,
spanned by $\mathcal P,\mathcal Z$.
This reduces to the excluded case.
\end{proof}

This is a distinct local response obstruction, not
a no-go for larger-support or nonlocal scalar
potentials and not an obstruction theorem for the
YM mass gap. It explains the unequal type II link
weights without deleting its obstructing coarse column.

\subsection{Native incidence and full second-order cancellation}

\begin{lemma}[Counts and support loads]
\label{c18v48:incidence}
The numbers of unordered pairs of types I, II, III
and 0 are
\[
 36n^3,\qquad12n^3,\qquad48n^3,\qquad48n^3 .
 \tag{V48.19}\label{c18v48:counts}
\]
For a fixed fine link, the numbers of pairs of
types I, II and III containing it in their fine
support are at most $18,8,16$, respectively.
These also bound active-link incidences.
\end{lemma}
\begin{proof}
There are $6n^3$ captured fine links, each incident
to four captured plaquettes and hence six I pairs.
A free link with exactly one odd transverse
coordinate has two captured and two uncaptured
incident plaquettes; there are $12n^3$ such links.
Each gives one II, four III and one type 0 pair.
The $6n^3$ free links with two odd transverse
coordinates each give six type 0 pairs.
Distinct plaquettes share at most one fine edge at
these torus sizes, so this counts each pair once.

Each captured plaquette has six I partners,
two II partners and four III partners.
An uncaptured plaquette has four III partners:
exactly two of its fine edges have one odd
transverse coordinate, each with two captured
neighbours.
For a captured fine link, the four incident
plaquettes give $24$ I memberships, but the six
pairs sharing that link were each counted twice;
its I load is $18$. A free link has at most two
incident captured plaquettes, so its I load is at
most $12$. The II load is at most $4\cdot2=8$.
Every incident plaquette, captured or not, has
four III partners, so the III load is at most
$4\cdot4=16$. Completion adds no dependence or
active component on an unused partner.
\end{proof}

Let $G_{\rm ed}$ sum \eqref{c18v48:generators} over
all fine-edge-sharing pairs, with flow
$\Xi_t^{\rm ed}$. Keep V47's chain-mediated flow
$\Xi_t^{\rm ch}$, V46's adjoint flow $\Psi_t^J$
and V45's Wilson flow $\Phi_s^W$. Define
\[
 c_\beta^{**}
  =b_2\circ\Xi_{\beta^2}^{\rm ed}
       \circ\Xi_{\beta^2}^{\rm ch}
       \circ\Psi_{-\beta^2/4608}^{J}
       \circ\Phi_{\beta/24}^{W}.
 \tag{V48.20}\label{c18v48:map}
\]

\begin{theorem}[The entire second coefficient is removed]
\label{c18v48:secondclosed}
At every fixed regulator, the actual centred normal
source for \eqref{c18v48:map}, expressed on the base
fibre, satisfies
\[
 r^{**}_{g,\beta}=\beta^3h_g+O_n(\beta^4).
 \tag{V48.21}\label{c18v48:cubic}
\]
This is uniform in $g$ in each fixed smooth norm
for fixed $n$; $h_g$ has base conditional mean zero.
For its actual mean-zero conditional generator
$L^{**}_{g,\beta}$, the inverse-source matrix has
expansion
\[
 \left(\left\langle r^{**}_\alpha,
       (L^{**}_{g,\beta})^{-1}r^{**}_\delta
       \right\rangle_{\nu^{**}_{g,\beta}}\right)_{\alpha\delta}
 =
 \beta^6\left(\left\langle h_\alpha,L_{g,0}^{-1}h_\delta
                 \right\rangle_{\nu_{g,0}}\right)_{\alpha\delta}
       +O_n(\beta^7).
 \tag{V48.22}\label{c18v48:sextic}
\]
The third coefficient $h_g$ is not evaluated or
asserted to be nonzero or uniformly bounded.
\end{theorem}
\begin{proof}
V47's corrected map has identically zero first
coefficient and second coefficient equal to the
sum of all fine-edge-sharing pair sources.
The new flow has duration $\beta^2$, so at second
order it adds the sum of base responses
$\mathcal D(G_{pq})$. Its interaction with an
order-$\beta$ term starts at order three.
Theorem~\ref{c18v48:corrections} cancels precisely
that sum. The self, chain-mediated and
fine-edge-sharing terms exhaust V46's coefficient.
The exact Hamiltonian vacuum remains unchanged.

At fixed $n$, the vacuum eigenpair is analytic
near zero coupling and all fields and flows are
smooth. The sampler is a submersion with compact
connected fibres. Pullback by the composed
diffeomorphism gives smooth elliptic families on
a fixed fibre. Taylor expansion gives
\eqref{c18v48:cubic}. Its lower coefficients
vanish pointwise, so differentiating the
conditional mean-zero identity shows that $h_g$
has zero base mean.

The base unreduced pair/free fibre has mean-zero
gap $3/2$. Identify the measures by their positive
square roots and add the zero-eigenspace projection
to the conditional generator. The resulting
operator is invertible at zero coupling.
An elliptic resolvent identity, or a Neumann
argument from $H^2$ to $L^2$, makes its inverse
and the mean-zero projection smooth for small
$\beta$ at fixed regulator. Compactness of the
coarse parameter space makes this uniform in $g$
for that $n$. Insert \eqref{c18v48:cubic} and the
$O_n(\beta)$ changes of inverse and measure to
obtain \eqref{c18v48:sextic}. Restriction to
internal-gauge-invariant functions preserves
the identities without a smooth quotient chart.
\end{proof}

\subsection{A separate uniform geometric estimate}

Define
\[
 \gamma_{\rm I}=\frac{947}{224640},\qquad
 \gamma_{\rm II}=\frac{43973}{1752192},\qquad
 \gamma_{\rm III}=\frac{407}{224640},\qquad
 A_*=18\gamma_{\rm I}+8\gamma_{\rm II}+16\gamma_{\rm III}
       =\frac{1338841}{4380480}.
 \tag{V48.23}\label{c18v48:gamma}
\]
For potential weights $a,b$, $4|a|+|b|$ bounds
each component and each component derivative of
the projected gradient. Indeed $\mathcal P$
contains the common fine link twice: a first
ordered derivative has at most two terms and
a second at most four, each bounded by one.
The six-edge $\mathcal Z$ uses distinct fine
links, so its ordered derivatives are bounded
by one. For II use the largest absolute
$\mathcal P$ and $\mathcal Z$ weights across
the groups. These give exactly
\eqref{c18v48:gamma}.

Each pair depends on seven fine links (21 color
directions) and acts on at most four captured
links (12 directions). The support-load lemma
bounds row and column absolute sums of
$DG_{\rm ed}$ by $21A_*$ and $12A_*$.
Its adjoint block adds at most $4A_*$ to each.
Consequently
\[
 \|M_{G_{\rm ed}}\|_2
    \le\sqrt{25A_*\cdot16A_*}
    \le25A_*=\frac{6694205}{876096}.
 \tag{V48.24}\label{c18v48:edbound}
\]
Combining with V47's bounds for the other flows yields
\[
 \begin{gathered}
 2e^{-2\Lambda_\beta^{**}}I
 \ \preceq\ Dc_\beta^{**}(Dc_\beta^{**})^*
 \ \preceq\ 2e^{2\Lambda_\beta^{**}}I,\\
 \Lambda_\beta^{**}
 =\frac83|\beta|+
       \left(\frac{749}{1408}+\frac{6694205}{876096}\right)\beta^2
 =\frac83|\beta|+\frac{157525571}{19274112}\beta^2 .
 \end{gathered}
 \tag{V48.25}\label{c18v48:gram}
\]
To see this directly, integrate the moving-frame
variational equations for each flow and its inverse.
Both differential norms are bounded by the
exponential of duration times the corresponding
differential bound. At the final point $BB^*=2I$.
The composed fine diffeomorphism's two-sided
singular-value bounds give \eqref{c18v48:gram}
for the rectangular coarse derivative. The normal
right inverse has norm at most
$e^{\Lambda_\beta^{**}}/\sqrt2$.

This geometric estimate holds for every real $\beta$
and all finite $n\ge5$. It does not make
\eqref{c18v48:cubic} or \eqref{c18v48:sextic}
uniform in volume, or imply convergence of the
vacuum expansion for arbitrary $\beta$.
A metric/rank estimate is not a low-frequency
conditional estimate.

\subsection{Verification and the next upstream obligation}

The reproducible diagnostic is
\begin{center}\small
\path{scripts/verify_ym_c18_v48_shared_edge_source.py}.
\end{center}
It checks the complete pair formula, contractions,
native response tables and fixed correction constants.
The three rational canonical calculations verify
24 source components and their exactly zero
correction residuals, retaining all coarse output
columns. The individual captured-link response
columns are also evaluated. An independent exact
covariant-Hessian check gives 2358 zero matrix
entries across the three curvature tests.
Exact operands remain rational or rational-dual;
numeric-to-symbolic mixing in an intermediate
constant check was caught and removed before
accepting the tests.

Complex-step regressions cover the canonical point
and three nonflat configurations per type.
Separate covariant-Hessian regressions use three
nonflat configurations per type. Native incidence
checks at $n=5,6$ verify classification, patch
sizes, counts and support loads. The written
parity and bidegree arguments, not these finite
samples, establish the all-volume and
all-configuration claims. V47's mathematical
diagnostics were replayed. These checks are
not a proof-assistant certificate or an
interacting continuum simulation.

V48 replaces V47 only. Removing this appendix
and restoring the title recovers the complete
V47 source byte for byte; the other 53 sources
are unchanged. Only \texttt{.tex} files remain
in \texttt{latex\_notes}; build products and
diagnostics are outside that directory.
The next companion and broad primary-literature
checkpoints are V50, and archive/restart is V100.

The local second-order cancellation problem is
now closed in its stated native block-two regime.
The next step should go upstream to nonlinear
estimates: combine the actual vacuum's spatially
summable derivatives with the full corrected
normal connection to prove, or disprove, a
volume-uniform remainder and conditional inverse
bound. Another formal Taylor coefficient would
not pay that obligation. Even a uniform
small-$\beta$ result would still require a
justified route to required block time and
physical weak coupling. It would not replace
scale transport, physical return and
interacting continuum reconstruction.
The full positive physical YM mass gap remains unproved.
\addtocontents{toc}{\protect\setcounter{tocdepth}{2}}
% END CYCLE 18 VERSION 48 APPEND
% BEGIN CYCLE 18 VERSION 49 APPEND
\clearpage
\addtocontents{toc}{\protect\vspace{1.25\baselineskip}}
\addtocontents{toc}{\protect\setcounter{tocdepth}{1}}
\section[V49: a uniform cubic source and conditional inverse]{V49: a uniform cubic source and conditional inverse}
\label{c18v49:context}
\addtocontents{toc}{\protect\clearpage}

V48 cancels the first two normal-source coefficients but only
states a fixed-volume remainder. This note pays that
particular uniformity obligation. The main device is to
differentiate the raw score in the original configuration
and project onto the actual fibre \emph{before} taking its
coupling Taylor remainder. Conditional centering then
disappears exactly. We need only the already constructed
vacuum's first two spatial derivatives, not a new
high-spatial-derivative vacuum expansion.

The resulting estimate holds on a very conservative,
volume-independent electric-coupling interval. It includes
the actual conditional density and a gap for the entire
unreduced conditional generator. It is not a proof at
the required Wilson block time four, at continuum weak
coupling, or of a positive physical YM mass gap.

\subsection{Fixed data, nonduplication and the claimed scope}

Work throughout with round $\mathrm{SU}(2)$, $N=2n$,
$n\ge5$, $T=-\sum_e\Delta_e$, and the actual normalized
vacuum $\Omega_\beta$ of $T-\beta W$. Define
\[
 \begin{split}
 \mathcal F_\beta
  &=\Xi^{\rm ed}_{\beta^2}\circ\Xi^{\rm ch}_{\beta^2}
       \circ\Psi^J_{-\beta^2/4608}\circ\Phi^W_{\beta/24},\\
 c_\beta&=b_2\circ\mathcal F_\beta,\qquad
 \psi_\beta=\mathcal F_\beta^{-1}.
 \end{split}
 \tag{V49.1}\label{c18v49:map}
\]
The fields, order of composition, all coarse outputs,
and fine metric are exactly V48's. We do not alter the
Hamiltonian vacuum or optimize over a substitute law.

V37 already proves a uniform native log-vacuum and a
uniform remainder for a \emph{different, fixed-map
vertical lift correction}. V38 proves a short-Wilson-flow
conditional gap. V42 and V44 pay connected spatial tensor
sums for the normal lift and score. None of those results
alone bounds the nonlinear remainder of
\eqref{c18v49:map}. We reuse their proved ingredients,
rather than count their lattice gaps as new results.
V48's exact pair calculation is also not repeated.

Let
\[
 K=10^9,\qquad
 \varepsilon=\frac1{10000K}=10^{-13},\qquad
 \beta_*=\frac3{512e},\qquad \delta=\frac{\beta_*}{2}.
 \tag{V49.2}\label{c18v49:interval}
\]
The tiny value of $\varepsilon$ is a sufficient proof
radius, not a physical threshold or an optimized estimate.
All estimates below are uniform in $n$, in the coarse value
$g$, and in configurations. The finite constant in the
cubic bound is defined by explicit finite-jet majorants
below; its numerical size is not evaluated.

\subsection{Local inputs for the composed map}

Use scalar invariant fine-link/color indices and the
unweighted all-slot norm $\mathfrak a_0$ of
\eqref{c18v42:anchor}. For a one-slot vector it is the
component supremum, not an extensive Euclidean norm.
For matrices it is the maximum absolute row or column
sum. Every tensor contraction used below is connected
and retains an external slot.

For each of the four individual vector fields
\[
 Y\in\{\nabla W,\nabla J,G_{\rm ch},G_{\rm ed}\},
\]
write $Y_i$ for its invariant scalar coefficient; $X_i$
continues to denote a fine invariant derivative. The
following uniform bounds suffice for every integer $m\ge0$:
\[
 \mathfrak a_0(\partial^m Y)\le A_m:=800\,48^m.
 \tag{V49.3}\label{c18v49:local}
\]
Here $\partial$ denotes ordered fine invariant coefficient
derivatives. These are bounds on each entire field, not
on only one of its patch summands.

\begin{proof}
Each patch uses at most eight fine links, hence at most
24 scalar directions. V47's chain-mediated support load
is at most 24; V48's three nonzero shared-edge loads sum
to at most $18+8+16=42$. The one-plaquette fields have
load four. A bound of 100 therefore covers each family,
whether the anchored slot is an output or a derivative.

Every ordered derivative of a trace word with distinct
links is bounded by one. For a product of two plaquette
traces, allocating each derivative between its two
factors gives the bound $2^r$ at order $r$. This also
bounds shared-link products. The adjoint potential
$4w_p^2-1$ costs at most $4\,2^r$ for $r\ge1$.
The absolute sums of the two weights on each active link
in the V47 and V48 fields are less than one. Captured-link
projections are fixed selections of link components;
they do not add a configuration-dependent derivative.
Thus a component derivative of order $m$ of a patch
field costs at most $2^{m+3}$. For an anchored slot,
there are at most $24^m$ choices of the other scalar
slots. Multiplication gives \eqref{c18v49:local}.
Unused completing links do not enlarge the dependency
support or generate additional derivatives.
\end{proof}

The invariant-frame variational generator is
$M_Y=DY+\operatorname{ad}Y$. Its bracket block has
absolute row and column sums at most four times the
component bound, and the same incidence argument gives
\[
 \mathfrak a_0(\partial^mM_Y)
       \le A_{m+1}+4A_m.
 \tag{V49.4}\label{c18v49:Mlocal}
\]
In particular, $K$ exceeds the right side for $m=0,1,2$
and exceeds $A_0,A_1,A_2,A_3$. Identifying the output and
one derivative slot shows that for
$d_Y=\operatorname{div}Y$,
\[
 \sup_i|X_i d_Y|\le A_2,\qquad
 \|(X_jX_i d_Y)_{j,i}\|_{\rm S}\le A_3 .
 \tag{V49.5}\label{c18v49:divlocal}
\]
The scalar value of $d_Y$ itself may be extensive.
It is not assigned a uniform bound.

Set
\[
 \begin{gathered}
 d_\beta=|\beta|/24+(2+1/4608)\beta^2\le|\beta|
       \quad (|\beta|\le1/4),\\
 E_\beta=e^{Kd_\beta},\qquad
 k_\beta=Kd_\beta e^{3Kd_\beta}.
 \end{gathered}
 \tag{V49.6}\label{c18v49:Ek}
\]
Concatenating the signed autonomous flows, in either
direction and in either composition order, gives
\[
 \|D\mathcal F_\beta\|_{\rm S},
 \ \|D\psi_\beta\|_{\rm S}\le E_\beta,\qquad
 \mathfrak a_0(\partial D\mathcal F_\beta),
 \ \mathfrak a_0(\partial D\psi_\beta)\le k_\beta .
 \tag{V49.7}\label{c18v49:flowbounds}
\]
These bounds also hold for every partial concatenation.

\begin{proof}
Parametrize each segment by its nonnegative absolute
duration and change the sign of its field when necessary.
The first variational equation has generator bounded
by $K$ in the scalar Schur norm, so its propagators are
bounded by $e^{K(t-r)}$. Differentiate once. The
inhomogeneous term is a connected contraction of
$\partial M_Y$ with two first variations. With any of
its three spatial slots anchored its norm is at most
$K e^{2Kt}$. Variation of constants bounds the second
variation by $Kt e^{3Kt}$. This argument applies to a
piecewise chosen generator through the concatenation
times: variations are continuous at a switch. The
invariant-frame bracket is already included in
\eqref{c18v49:Mlocal}. This is the same coefficient
convention as V38's second variation.
\end{proof}

\subsection{Finite mixed jets without a spatial volume factor}

We record carefully the regularity input used in the
Taylor estimate. The coupling variable is a parameter;
it is never a summed spatial slot or a spatial anchor.

\begin{lemma}[Uniform mixed map jets]
\label{c18v49:jets}
For each fixed pair $k,j$ with $k+j\le6$, the ordered
invariant coefficient jets of $\mathcal F_\beta$ and
$\psi_\beta$, with $k$ coupling derivatives and $j$
fine derivatives, have finite all-spatial-slot bounds
on $|\beta|\le\varepsilon$, independent of $n$.
For jets without a fine input slot, use the output
component supremum. In particular, writing
\[
 C_\beta=Dc_\beta ,
\]
there are finite constants $c_{k,j}$ bounding
$\mathfrak a_0(\partial_\beta^k\partial^j C_\beta)$
for $0\le k\le3$, $0\le j\le2$.
\end{lemma}

\begin{proof}
Here is a direct finite-order justification, including
the parameter derivatives. Realize every link as a
unit quaternion. Each field's ambient velocity is its
imaginary coefficient multiplied by that quaternion.
Inversions in its words become quaternion conjugations.
Thus each velocity is a finite local real polynomial,
of degree at most nine, with bounded coefficients and
the incidence established above. On the product of
radius-two link balls its ordinary derivatives through
any fixed order have finite anchored bounds. There
is no sum over the full set of links: a chosen anchor
meets only the patch load in \eqref{c18v49:local}.

Rescale each of the four flow segments to a unit interval.
Its equation is
\[
 \dot U=\tau(\beta)P_X(U),\qquad
 \tau(\beta)\in\{\beta/24,-\beta^2/4608,\beta^2\}.
 \tag{V49.8}\label{c18v49:parameterODE}
\]
The inverse has reversed order and signs. The coefficients
$\tau,\tau',\tau''$ are bounded on the stated interval
and higher derivatives vanish. The real trajectories
remain on the unit spheres. Differentiate the equation
in any ordered collection of coupling and fine input
directions. At total order $m$, the only unknown
order-$m$ variation on the right is the linear term
$\tau P_X'(U)Y_m$. Every other term is a finite
Leibniz/partition product of lower variations and a
derivative of $P_X$, multiplied by a derivative of
$\tau$. Each product is a connected rooted tensor
network. Parameter variations are one-slot vectors
bounded by a component supremum; contracting such a
vector into a spatial tensor does not close a trace.
The connected-contraction inequality therefore
bounds every spatial anchoring of each term.

For an explicit way to select finite majorants, take
$B\ge1$ to dominate all local ambient polynomial and
frame derivatives needed through order eight, all
$\tau$ derivatives, and the initial coordinate jets.
There are only four field types and finitely many
local polynomial patterns, so $B$ is obtained by
finite coefficient sums on radius-two link balls,
using the bound 100 on patch load. Enlarge it to
cover the at most four quaternion coordinates per
link. A deliberately generous recursive majorant
for orders $1\le m\le6$ is
\[
 \begin{split}
 z_m(t)=e^{Bt}\bigg[
 B^m m!+
 B\,2^{m+2}(m+1)!
 \int_0^t e^{-Bs}
       \Big(1+\sum_{q=1}^{m-1}z_q(s)\Big)^m ds
                    \bigg],\quad 0\le t\le4 .
 \end{split}
 \tag{V49.9}\label{c18v49:recursion}
\]
Enlarging $B$ also bounds the homogeneous coefficient.
The displayed partition multiplier overcounts the
finite ordered chain and product rules. Induction
and Gronwall bound the actual jets by these finite
numbers. No high-order variation occurs in the
inhomogeneity defining $z_m$. Finite products of
these numbers cover the initial/final local frame
changes. Invariant derivatives of the initial
quaternion are local jets and obey the same bounds.
This proves uniform finite mixed jets in the
invariant convention, not a bound on a global
configuration displacement in Euclidean norm.

Finally $C=B_2(\mathcal F_\beta U)D\mathcal F_\beta$,
where $B_2=Db_2$ has fixed two-link incidence.
Two fine derivatives of $C$ and three coupling
derivatives need at most six total derivatives of
the fine flow. The finite chain rule and connected
contractions give the constants $c_{k,j}$.
\end{proof}

This lemma is a finite smooth-jet result. It does not
postulate a holomorphic map on a growing-dimensional
configuration tube, uniform high spatial derivatives
of the interacting vacuum, or convergence of all
orders of a normal-source correction.

\subsection{The full normal inverse in a summable matrix norm}

Put
\[
 G_\beta=C_\beta C_\beta^*,\quad Z_\beta=G_\beta^{-1},
 \quad R_\beta=C_\beta^*Z_\beta,\quad
 P_\beta=I-R_\beta C_\beta .
 \tag{V49.10}\label{c18v49:normal}
\]
These are respectively the full coarse Gram, its inverse,
the normal lift and the orthogonal fibre projection.
The notation $R_\beta$ here is a lift, not V37's
scalar bound on its Fourier family.

The two-link sampler obeys
\[
 B_2B_2^*=2I,\qquad
 \mathfrak a_0(B_2)\le6,\qquad
 \mathfrak a_0(\partial B_2)\le72 .
\]
The constants follow by counting its two adjoint-prefix
blocks and their first derivatives. Along the concatenated
fine flow, the component velocity is at most $K$. Hence
\[
 \begin{gathered}
 \|C_\beta-B_2(U)\|_{\rm S}
 \le \Delta_\beta:=72Kd_\beta E_\beta+6(E_\beta-1),\\
 \|G_\beta-2I\|_{\rm S}
       \le\Delta_\beta(12+\Delta_\beta)<\frac18,
       \qquad |\beta|\le\varepsilon .
 \end{gathered}
 \tag{V49.11}\label{c18v49:Gramclose}
\]
Indeed the variation of $B_2(\mathcal F_\beta U)$ costs
$72Kd_\beta$, and the invariant differential differs
from identity by at most $E_\beta-1$. Since
$Kd_\beta\le1/10000$,
\[
 E_\beta\le\frac{10000}{9999},\qquad
 \Delta_\beta\le\frac{78}{9999},
\]
which proves the strict last inequality.

The Neumann series about $2I$ therefore converges in
the \emph{full} Schur norm and gives
\[
 \|Z_\beta\|_{\rm S}\le\frac8{15}.
 \tag{V49.12}\label{c18v49:Z}
\]
No output projection, selected Fourier band or assumed
conditional spectral gap is involved in this finite
matrix inverse.

Differentiating $GZ=I$ in the ordered fine directions
and in $\beta$ recursively expresses each derivative
of $Z$ as a sum of connected products of $Z$ and
derivatives of $G$. For example,
\[
 \partial Z=-Z(\partial G)Z,\qquad
 \partial_\beta Z=-Z(\partial_\beta G)Z .
\]
Every further product-rule term retains the order of
the fine derivatives; their commutation is not assumed.
Lemma~\ref{c18v49:jets}, \eqref{c18v49:Z}, and the
all-slot contraction bound thus define finite constants
\[
 \begin{aligned}
 \mathfrak a_0(\partial_\beta^k\partial^jR_\beta)
      &\le r_{k,j} &&(0\le k\le3,\ 0\le j\le2),\\
 \mathfrak a_0(\partial_\beta^kP_\beta)&\le p_k
                                      &&(0\le k\le3).
 \end{aligned}
 \tag{V49.13}\label{c18v49:Rjets}
\]
One may choose these constants by the finite derivative
rules just specified and the majorants in
\eqref{c18v49:recursion}. Their definition does not
take a supremum over unknown volume-dependent
inverses. It involves only finite scalar arithmetic
and integrations of previously bounded functions.

\subsection{The projected raw-score argument}

Let $u_\beta=\log\Omega_\beta$ up to its scalar
normalization. V37's analytic Fourier-family construction
has radius $\beta_*$ in the $a=1$ norm and norm at most
$3/8$ there. Cauchy estimates on disks of radius $\delta$
centered in $|\beta|\le\delta$, followed by its proved
first- and second-derivative bounds, give
\[
 \begin{gathered}
 \sup_i|\partial_\beta^k X_i u_\beta|\le a_k,
 \qquad
 \|(\partial_\beta^k X_jX_i u_\beta)_{j,i}\|_{\rm S}
             \le h_k,\\
 a_k=\frac{k!}{8\delta^k},\qquad
 h_k=\frac{9k!}{2e\delta^k},\qquad 0\le k\le3 .
 \end{gathered}
 \tag{V49.14}\label{c18v49:vacuumjets}
\]
The estimates concern derivatives at the original
fine configuration $U$. They do not differentiate
$u_\beta(\psi_\beta V)$ three times in $\beta$.
This distinction avoids an unsupported fifth
spatial derivative bound. The normalization constant
drops out before applying Cauchy's estimate.

Define the uncentered normal score and its ambient
fine derivative by
\[
 \begin{gathered}
 S_{\beta,\alpha}
   =\sum_iX_i(R_\beta)_{i\alpha}
       +2\sum_i(X_i u_\beta)(R_\beta)_{i\alpha},\\
 (D S_\beta)_{j\alpha}
   =\sum_iX_jX_i(R_\beta)_{i\alpha}
      +2\sum_i(X_jX_i u_\beta)(R_\beta)_{i\alpha}
      +2\sum_i(X_i u_\beta)X_j(R_\beta)_{i\alpha},\\
 \mathcal K_\beta=P_\beta D S_\beta .
 \end{gathered}
 \tag{V49.15}\label{c18v49:K}
\]
The fine index of $D S$ is its row index.
For a constant full coarse bank $v$, the exact identity is
\[
 \nabla_{\mathrm{fib}} r_{\beta,v}
       =\mathcal K_\beta v,\qquad
 r_{\beta,\alpha}
       =S_{\beta,\alpha}-\mathbb E_{\nu_{g,\beta}}
                                     S_{\beta,\alpha}.
 \tag{V49.16}\label{c18v49:fibreidentity}
\]
The conditional expectation is constant \emph{on that
fibre}. Its ambient dependence through $g=c_\beta(U)$
is annihilated by $P_\beta$. Thus there is no derivative
of a conditional partition function in
\eqref{c18v49:fibreidentity}.

\begin{lemma}[Vanishing jets transfer to the projected gradient]
\label{c18v49:vanishing}
At each fixed regulator, V48 implies the identities
\[
 \mathcal K_0=0,\qquad
 \partial_\beta\mathcal K_0=0,\qquad
 \partial_\beta^2\mathcal K_0=0
 \quad\hbox{on the entire fine configuration space}.
 \tag{V49.17}\label{c18v49:zerojets}
\]
These are pointwise identities, not averaged estimates.
\end{lemma}

\begin{proof}
Identify $c_\beta^{-1}(g)$ with $b_2^{-1}(g)$ by
$U=\psi_\beta(V)$. V48 gives
$r_{g,\beta}(\psi_\beta V)=O_n(\beta^3)$ in every
fixed smooth norm. Differentiate along any tangent
vector to $b_2^{-1}(g)$. Its image by $D\psi_\beta$
is tangent to $c_\beta^{-1}(g)$, and these images
span that tangent space. The metric-dual transformation
between these fibre covectors and
$\mathcal K_\beta(\psi_\beta V)$ is smooth and
invertible. It follows at fixed $n$ that this last
vector bank is $O_n(\beta^3)$ in a smooth norm.
Compose back by $\mathcal F_\beta$ to obtain
$\mathcal K_\beta(U)=O_n(\beta^3)$ on the full
space. Differentiating at zero proves
\eqref{c18v49:zerojets}.

Only vanishing of finite jets, not their numerical
remainder constants, is imported in this argument.
There is no assertion that the raw score's second
coefficient vanishes before centering; a coarse-only
coefficient would also be annihilated here.
\end{proof}

Define the finite constants
\[
 \begin{split}
 D_q&=r_{q,2}
   +2\sum_{a=0}^q\binom qa
             (h_a r_{q-a,0}+a_a r_{q-a,1}),\\
 M&=\sum_{q=0}^3\binom3q p_{3-q}D_q,
       \qquad C_3=M/6 .
 \end{split}
 \tag{V49.18}\label{c18v49:C3}
\]
A diagonal contraction of $\partial^2R$ retains the
external fine and coarse slots. V42's trace bound
therefore applies, including after coupling
differentiation. The last term of $D S$ is bounded
by a component supremum of $\nabla u$, not its
extensive $\ell^2$ norm. Product rules in
\eqref{c18v49:K} now give
\[
 \|\partial_\beta^3\mathcal K_\beta\|_{\rm S}\le M .
\]
Together with \eqref{c18v49:zerojets}, the integral
Taylor formula proves the \emph{uniform} estimate
\[
 \|\mathcal K_\beta\|_{\rm S}\le C_3|\beta|^3,
 \qquad
 |\nabla_{\rm fib}r_{\beta,v}|
                  \le C_3|\beta|^3|v|.
 \tag{V49.19}\label{c18v49:cubicgradient}
\]
For negative $\beta$ integrate over the oriented
segment and take its absolute length. Smooth
finite-volume identities and a uniform third-jet
bound suffice; no volume-dependent Taylor radius
has been promoted to a uniform one.

\subsection{The actual conditional law and its uniform inverse}

The ambient density in $V=\mathcal F_\beta(U)$ is exactly
\[
 h_\beta(V)=\Omega_\beta(\psi_\beta V)^2
                       \operatorname{Jac}\psi_\beta(V).
 \tag{V49.20}\label{c18v49:density}
\]
On $b_2^{-1}(g)$ use V45's pair/free product parametrization
$(U_1,U_2)=(x,x^{-1}g)$ and the unused fine links.
Its base metric $m_0$ has factor $2g_{\rm round}$
on each pair variable and $g_{\rm round}$ on each
free link. The embedding is totally geodesic.
Product Haar is its normalized volume, denoted
$\nu_{g,0}$. Since $B_2B_2^*=2I$, the base coarea
Jacobian is constant, so the conditional law pulled
back from the original fibre is exactly
\[
 d\widetilde\nu_{g,\beta}(V)
   =\frac{h_\beta(V)}{\int_{b_2^{-1}(g)}h_\beta\,d\nu_{g,0}}
                                  d\nu_{g,0}(V).
 \tag{V49.21}\label{c18v49:conditionallaw}
\]
There is no additional induced-volume determinant
to multiply into this formula. The metric of the
original fibre after pullback is
$m_\beta=(\psi_\beta)^*\mathfrak g|_{b_2^{-1}(g)}$.

Put $\mathcal R_\beta=64e|\beta|$ and
\[
 H_\beta=
 \frac{8\mathcal R_\beta}{e}E_\beta^2
       +\frac{2\mathcal R_\beta}{3}k_\beta
       +Kd_\beta(E_\beta^2+k_\beta).
 \tag{V49.22}\label{c18v49:densityH}
\]
Then
\[
 \|\operatorname{Hess}\log h_\beta\|_{2\to2}
                  \le H_\beta<\frac18
           \quad (|\beta|\le\varepsilon).
 \tag{V49.23}\label{c18v49:Hsmall}
\]

\begin{proof}
V37 gives the one-link gradient bound
$\mathcal R_\beta/3$ and covariant Hessian bound
$4\mathcal R_\beta/e$ for $u_\beta$.
The second-variation chain rule, with the same
anchored contraction as V38, bounds the Hessian
of $2u_\beta\circ\psi_\beta$ by the first two
terms of \eqref{c18v49:densityH}.

For each signed inverse-flow segment, Liouville's
identity expresses its logarithmic Jacobian as
an integral of the signed divergence of its field
along a partial flow. For a composition, the
logarithms add after the corresponding partial
pullbacks. Every such partial concatenation has
duration at most $d_\beta$ and bounds
\eqref{c18v49:flowbounds}. By
\eqref{c18v49:divlocal}, the divergence has
component gradient bound $K$ and Hessian norm
bound $K$. Its pullback Hessian is thus at most
$K(E_\beta^2+k_\beta)$. Summing the absolute
segment durations proves the last term of
\eqref{c18v49:densityH}, without bounding the
extensive undifferentiated logarithmic Jacobian.

Finally $Kd_\beta\le1/10000$,
$E_\beta\le10000/9999$ and
$k_\beta\le1/9997$. Substitution, with $e<3$
and $|\beta|\le10^{-13}$, gives
\eqref{c18v49:Hsmall}. This last inequality is
very conservative.
\end{proof}

\begin{theorem}[Uniform actual conditional gap]
\label{c18v49:gap}
For all $n\ge5$, all $g$ and $|\beta|\le\varepsilon$,
the mean-zero generator $L_{g,\beta}$ for the
original induced fibre metric and actual conditional
law satisfies
\[
 \operatorname{gap}L_{g,\beta}
     \ge e^{-2\Lambda_\beta}(1-H_\beta)>\frac12,
 \qquad
 \Lambda_\beta=\frac83|\beta|
             +\frac{157525571}{19274112}\beta^2 .
 \tag{V49.24}\label{c18v49:gapbound}
\]
The same floor holds upon restriction to the
internal-gauge-invariant subspace.
\end{theorem}

\begin{proof}
The Ricci endomorphism of a round pair factor
with metric $2g_{\rm round}$ is $I$; for a free
factor it is $2I$. Hence
$\operatorname{Ric}(m_0)\succeq m_0$.
Total geodesy makes the fibre Hessian of
$\log h_\beta$ the restriction of its ambient
Hessian. Normalizing \eqref{c18v49:conditionallaw}
does not change that Hessian in fibre directions.

For the $m_0$ diffusion with density
\eqref{c18v49:conditionallaw}, integrated weighted
Bochner gives
\[
 \int(\mathcal A f)^2d\widetilde\nu
  =\int\bigl(
    |\operatorname{Hess}_{m_0}f|^2+
    (\operatorname{Ric}(m_0)-\operatorname{Hess}_{m_0}
           \log h_\beta)(\nabla f,\nabla f)
           \bigr)d\widetilde\nu .
\]
Applying this to its nonconstant eigenfunctions
proves the Poincare constant $(1-H_\beta)^{-1}$.
The space is compact and connected, so the
zero eigenspace consists only of constants.

V48's two-sided fine differential estimate gives
$e^{-2\Lambda_\beta}m_0\preceq m_\beta
\preceq e^{2\Lambda_\beta}m_0$ on the fixed fibre.
Its dual metric therefore bounds the actual
Dirichlet form below by $e^{-2\Lambda_\beta}$
times this reference-metric form, using the
\emph{same} measure \eqref{c18v49:conditionallaw}.
This proves the first inequality. On the stated
interval $\Lambda_\beta<1/8$ and $H_\beta<1/8$;
$(7/8)e^{-1/4}>1/2$ proves the numerical floor.

All fields and the law intertwine the internal
gauge isometries. Restricting the Poincare
inequality to invariant functions cannot lower
its floor. This does not require a smooth
quotient or removal of singular orbit strata.
\end{proof}

The weighted Bochner step is included above. For primary
context, the bibliographic page of Bakry and Emery's
\href{https://numdam.org/item/SPS_1985__19__177_0/}
{\emph{Diffusions hypercontractives} (1985)}
was rechecked. The flow-construction context of
\href{https://arxiv.org/abs/0907.5491}{Luscher's
\emph{Trivializing maps, the Wilson flow and the HMC algorithm}}
was also rechecked at abstract level. No Wilson-Gibbs
trivialization result is substituted for the native
Hamiltonian or for the estimates proved here.
This targeted check is not the broad V50 sweep.

\subsection{Uniform sextic cost and the limit of what is closed}

\begin{theorem}[Full-bank uniform nonlinear source cost]
\label{c18v49:cost}
Let $\Sigma_{g,\beta}$ be the covariance matrix of
the actual centered normal source, and let
$\mathcal Q_{g,\beta}$ be its inverse-source matrix.
Then, for every $n,g$ and $|\beta|\le\varepsilon$,
\[
 \begin{gathered}
 \Sigma_{g,\beta}\preceq2C_3^2|\beta|^6 I,\qquad
 \mathcal Q_{g,\beta}:=
  \bigl(\langle r_\alpha,L_{g,\beta}^{-1}r_\gamma
                                     \rangle_{\nu_{g,\beta}}\bigr)
       \preceq4C_3^2|\beta|^6 I .
 \end{gathered}
 \tag{V49.25}\label{c18v49:sextic}
\]
In particular, the normal source is $O(|\beta|^3)$
as an operator from the full coarse $\ell^2$ bank
to the actual conditional $L^2$ space, uniformly
in volume and coarse configuration.
\end{theorem}

\begin{proof}
For a constant full bank $v$, its score has actual
conditional mean zero. Apply
\eqref{c18v49:gapbound} to this scalar score and
then \eqref{c18v49:cubicgradient}:
\[
 \|r_v\|_{L^2(\nu)}^2
 \le2\int|\nabla_{\rm fib}r_v|^2d\nu
 \le2C_3^2|\beta|^6|v|^2 .
\]
The mean-zero inverse has norm at most two by the
same conditional gap. This proves both matrix
inequalities without estimating individual
columns and multiplying by the number of outputs.
\end{proof}

The inherited Poisson-correction principle now applies
with a uniformly small cost for this \emph{new map}.
At fixed regulator let $\phi_\alpha=L^{-1}r_\alpha$
be centered and set $A^\sharp=R+\nabla_{\rm fib}\phi$.
Smooth ellipticity gives the solution; uniqueness
preserves its internal gauge covariance. Its
vertical weighted divergence is $-L\phi=-r$, so
\[
 \begin{gathered}
 Dc_\beta A^\sharp=I,\qquad r^{A^\sharp}=0,\\
 \mathbb E_\nu[(A^\sharp)^*A^\sharp]
 \preceq\left(\frac12 e^{2\Lambda_\beta}
                    +4C_3^2|\beta|^6\right)I .
 \end{gathered}
 \tag{V49.26}\label{c18v49:sharp}
\]
Here the mixed term is genuinely zero: $R$ is
normal to the \emph{entire} fibre, and the added
gradient is tangent to it, pointwise. This is
not the merely internal-gauge-horizontal lift
for which V38 correctly retained mixed terms.
The new correction is not claimed to be local,
nor are its coarse derivatives bounded here.

The distinction from V48 is precise. Its
$O_n(\beta^3)$ centered-source remainder now
has a volume-uniform conditional $L^2$ and
fibre-gradient bound, and its inverse cost has
the uniform sextic upper bound
\eqref{c18v49:sextic}. This does not prove a
volume-uniform order-seven \emph{asymptotic
expansion} of that inverse matrix, an evaluated
or nonzero cubic coefficient, or an all-orders
normal trivialization.

\subsection{Verification and the next upstream task}

The diagnostic
\begin{center}\small
\path{scripts/verify_ym_c18_v49_uniform_normal_remainder.py}
\end{center}
checks the local numerical majorants, exact
Schur/Neumann arithmetic, the gap-radius
inequalities, the factor $1/6$ in the Taylor
remainder and its Leibniz constants. A moving-fibre
symbolic test includes a nonzero coarse-only raw
coefficient and verifies that its projected
gradient vanishes, whereas a derivative along
the wrong fixed fibre need not vanish.
Tensor regressions check connected diagonal
trace contractions without a volume multiplier.
V48's full mathematical regression is replayed.
These checks support the written uniform
arguments; finite examples are not treated as
a proof for all volumes or as an interacting
conditional simulation.

V49 replaces V48 only. Removing this appendix
and restoring the title recovers the preceding
source byte for byte. All other source files
are preserved and only \texttt{.tex} files are
kept in \texttt{latex\_notes}. The previous goal
turn was progress; this turn pays the stated
nonlinear uniformity obligation. The next
companion review and broad primary-literature
sweep are V50; archive/restart remains V100.

The next task should not be another small-coupling
coefficient or a reproof of a fixed-spacing
electric gap. Determine what quantitative
estimate can transport this controlled native
normal connection toward the required block
deformation. The first Wilson-flow duration
here is $\beta/24$, not four. The actual
conditional low spectrum at long block time,
coarse variance and return, stability of
variable banks and projection energies,
weak-coupling transport, and physical scale
calibration still require proofs. No rescaling
of the coupling or time in this note reaches
that regime. A nontrivial interacting
four-dimensional continuum reconstruction and
its positive physical Yang--Mills mass gap
remain unproved.
\addtocontents{toc}{\protect\setcounter{tocdepth}{2}}
% END CYCLE 18 VERSION 49 APPEND
% BEGIN CYCLE 18 VERSION 50 APPEND
\clearpage
\addtocontents{toc}{\protect\setcounter{tocdepth}{2}}
\section[V50: temporal transport and native metric strain]{V50: temporal transport and native metric strain}
\label{c18v50:context}

\paragraph{Context and change of direction (21 September 2026).}
V49 pays a small-electric-coupling, volume-uniform cubic
normal-source bound and a conditional inverse bound for its
corrected block-two map. It does not transport those estimates
to Wilson block time four. This scheduled fifth-version
companion review and tenth-version literature sweep therefore
turns to the \emph{temporal} conditional source and the
\emph{material change of the actual fibre metric}. A small
normal source is not a small temporal source. Below the latter
is computed exactly at the electric origin, together with its
Poisson solution and the resulting metric strain.

The new native identity cancels the entire uncaptured
plaquette sector of the initial strain. The surviving
captured sector is uniformly bounded but has both signs.
This is an initial-time theorem, not an integrated bound
through time four and not a physical Yang--Mills mass gap.

\subsection{Companion review, literature sweep and nonduplication}

The current companions in \path{latex_notes} were checked
against the V45 inventory; all nine source hashes are
unchanged. Their terminal results and relevant hypotheses
were reread, not independently recertified line by line.
They are later continuations of the supplied related notes.
The following transfer decisions retain their scope.

The supplied Downloads versions were also checked at their
terminal ledgers: perturbative ESM V64 has fixed-order
common-action/Hessian compatibility but an unpopulated
critical source array; source-selection V52 optimizes a
declared one-query class without supplying the historical
preparation bank; stationarity V54 pays a linear covector
response with two fast scales but not the nonlinear native
remainder; exact-action V45 resolves frozen intermediate
frequencies without a moving nonlinear energy estimate.
These distinctions persist in the later companions and
support the same actual-source, actual-path discipline.

\begin{itemize}
\item Exact-action Einstein bridge V55 retains a five-row
  tangency obstruction at fixed cutoff; stationarity V61
  distinguishes the initialized path from ambient tests
  and still owes its actual state-only source. We therefore
  evaluate the actual YM temporal source on its own path.
\item Source-selection V65's finite naturality/Choi gap and
  arithmetic V17's projector reconstruction require actual
  represented sources. Neither is a conditional YM gap.
\item Perturbative ESM V14, after its archived first cycle,
  requires controlled partial isometries, complements and
  the correct path component. Its positive-reference,
  coefficientwise estimates do not supply our nonlinear
  conditional inverse. SM--Einstein V72's bounded-time,
  fixed-mode many-copy limit is not a finite-species continuum.
\item NS V26 motivates evaluating contracted sources before
  taking absolute bounds; its grid-dependent global
  inequalities are not imported. Shell-mixing V16 retains
  coarea and normal-return terms but lacks a nonperturbative
  mixed-curvature remainder. YM-parallel V5's extensive-charge
  obstruction reinforces the distinction between global
  scalar norms and local operator norms.
\end{itemize}

\paragraph{What the wider search supplies.}
The following primary sources were checked on 21 September
2026 across transport, multiscale functional inequalities,
effective dynamics and rigorous lattice/continuum YM.
They motivate a direction, not an automatic transfer theorem.

\begin{enumerate}
\item Moser, \emph{On the volume elements on a manifold}
  (1965), \href{https://doi.org/10.1090/S0002-9947-1965-0182927-5}{Theorem 2},
  constructs a diffeomorphism from a prescribed smooth
  positive volume-form path by solving a linear divergence
  equation and integrating its vector field. Our f17 V36
  already uses a conditional Poisson/Moser correction for
  a fixed six-spin Wilson block. We reuse that principle;
  we do not claim it as new.
\item Fathi--Mikulincer--Shenfeld,
  \href{https://arxiv.org/abs/2305.03786v3}{\emph{Transportation onto log-Lipschitz perturbations}}
  (v3, 24 August 2026), treats global log-Lipschitz
  perturbations with curvature assumptions. Its Theorems
  1--3 depend on a \emph{global} gradient bound, not a
  one-link component bound. The extensive temporal energies
  computed below rule out substituting the latter for
  the former in a uniform first-order estimate.
\item Bauerschmidt--Bodineau--Dagallier,
  \href{https://arxiv.org/abs/2307.07619v2}{\emph{Stochastic dynamics and the Polchinski equation}},
  Section 5.3, and Shenfeld,
  \href{https://arxiv.org/abs/2205.01642v4}{\emph{Exact renormalization groups and transportation of measures}},
  connect renormalized Hessian bounds to Lipschitz transport
  and functional inequalities. Their Gaussian-reference
  flow and multiscale lower bounds have not been identified
  with the present conditioned compact-group flow.
\item Conforti--Eichinger,
  \href{https://arxiv.org/abs/2502.01353}{\emph{A coupling approach to Lipschitz transport maps}},
  offers a Euclidean coupling route under weaker convexity
  assumptions. Leli\`evre--Zhang,
  \href{https://arxiv.org/abs/1805.01928}{\emph{Pathwise estimates for effective dynamics}},
  concerns nonlinear/vectorial reaction coordinates.
  Neither abstract-level transfer identifies our actual
  time-dependent conditional inverse; the latter method's
  conditional-Poincar\'e input was already audited in V40.
\item Shen--Zhu--Zhu,
  \href{https://arxiv.org/abs/2204.12737}{\emph{A stochastic analysis approach to lattice Yang--Mills at strong coupling}},
  proves functional inequalities and correlation decay for
  a lattice Gibbs law in a strong-coupling regime. That
  law is not our Hamiltonian vacuum conditional law.
  Chevyrev--Shen's
  \href{https://arxiv.org/abs/2302.12160v3}{two-dimensional invariant-measure and universality theorem}
  and Chevyrev--Qu--Shen's
  \href{https://arxiv.org/abs/2608.27828v2}{three-dimensional abelian scaling-limit theorem}
  are important model-specific advances, not a four-dimensional
  nonabelian continuum mass-gap proof.
\end{enumerate}

The resulting adjustment is to test \emph{density transport
and metric deformation together}. The archive already has
fixed-block co-moving transport, and V45--V49 already have
the pair geometry, local Hessian bounds and short-time gap
inputs. The new calculation below is the native full-fibre
\emph{temporal} Poisson solution and its material strain,
including a sign witness and exact volume scaling. Searches
of the current note and YM archives did not locate this
identity; this is a targeted nonduplication check, not a
recertification of the entire historical corpus.

\subsection{The actual Wilson-time path and its conditional source}

Keep $\mathrm{SU}(2)$ with the round metric, the spatial
torus $N=2n$, $n\ge5$, $T=-\sum_e\Delta_e$, and the
normalized positive vacuum $\Omega_\beta$ of $T-\beta W$.
Set $F=\nabla W$, and let $\Phi_t$ be its complete flow
on the compact fine configuration manifold $Q$.
There is no rescaling of $t$ or $\beta$ in this note.
For the initial map choose either
\[
 S_\beta=I\quad\hbox{or}\quad S_\beta=\mathcal F_\beta
       \text{ of \eqref{c18v49:map}}.
 \tag{V50.1}\label{c18v50:initialmap}
\]
These are two explicitly distinguished paths, not equal
maps at nonzero coupling. Define
\[
 \begin{aligned}
 F_{\beta,t}&=\Phi_t\circ S_\beta,&
 c_{\beta,t}&=b_2\circ F_{\beta,t},\\
 \psi_{\beta,t}&=S_\beta^{-1}\circ\Phi_{-t},&
 \overline m_{\beta,t}&=\psi_{\beta,t}^{*}\mathfrak g.
 \end{aligned}
 \tag{V50.2}\label{c18v50:path}
\]
For $S_\beta=I$, $t=4$ is the literal required Wilson
block map $b_2\Phi_4$. For $S_\beta=\mathcal F_\beta$,
the starting point is V49's corrected map, but the endpoint
is $b_2\Phi_4\mathcal F_\beta$, not $b_2\Phi_4$.
No endpoint equivalence is assumed.

Fix any coarse value $g$. Work on $M_g=b_2^{-1}(g)$,
identified by the pair/free coordinates of V45 and V49.
Let $m_0=\mathfrak g|_{M_g}$ and $\nu_{g,0}$ be its
normalized volume: a product of round Haar laws, with
metric $2g_{\rm round}$ on each pair and $g_{\rm round}$
on each free edge. The metric and probability density
actually transported from the original fibre are
\[
 \begin{split}
 m_{\beta,t}&=\overline m_{\beta,t}|_{TM_g},\\
 h_{\beta,t}(V)&=\Omega_\beta(\psi_{\beta,t}V)^2
                         \operatorname{Jac}\psi_{\beta,t}(V),\\
 p_{\beta,t}&=h_{\beta,t}/Z_{\beta,t}(g),\qquad
 Z_{\beta,t}(g)=\int_{M_g}h_{\beta,t}\,d\nu_{g,0}.
 \end{split}
 \tag{V50.3}\label{c18v50:law}
\]
Here $h$ is first defined on the \emph{ambient} manifold.
The conditioned law is $p_{\beta,t}\nu_{g,0}$, not $p$
times the moving metric volume. The metric is used for
the Dirichlet energy. As in V49 there is no additional
induced-volume determinant in \eqref{c18v50:law}.

\begin{lemma}[Exact temporal source]
Suppress $\beta,g$ temporarily, and let $\mathbb E_t$
refer to $p_t\nu_{g,0}$. Then
\[
 \begin{split}
 a_t:=\partial_t\log p_t&=A_t-\mathbb E_t A_t,\\
 A_t&=12W-F\log h_t\\
    &=\big(12W-F\log h_0\big)\circ\Phi_{-t}.
 \end{split}
 \tag{V50.4}\label{c18v50:source}
\]
The $F$ derivative in this formula is ambient, not a
derivative of the restricted conditional density.
\end{lemma}

\begin{proof}
Each plaquette character $w_p=\tfrac12\Tr U_p$ satisfies
$\sum_e\Delta_e w_p=-12w_p$, hence
$\operatorname{div}_{\mathfrak g}F=-12W$.
Composition of densities gives
$h_t=(h_0\circ\Phi_{-t})\operatorname{Jac}\Phi_{-t}$.
Differentiating at fixed $V$ gives the last expression
for $A_t=\partial_t\log h_t$. Equivalently the ambient
continuity equation is
$\partial_t h_t=-\operatorname{div}_{\mathfrak g}(h_tF)$,
which gives the middle expression. Restrict to $M_g$
only after this derivative. Differentiation under the
compact-fibre integral yields
$\partial_t\log Z_t=\mathbb E_t A_t$. This proves the
normalization term as well as the sign in
\eqref{c18v50:source}.
\end{proof}

\subsection{A transport criterion retaining the moving metric}

For a tangent vector field $Y$ on $M_g$, define
$\operatorname{div}_{p_t\nu_{g,0}}Y$ by its integration
by parts against the \emph{probability measure}.
Let
\[
 L_t f=-\operatorname{div}_{p_t\nu_{g,0}}
                       (\nabla_{m_t}f),\qquad
 \phi_t=L_t^{-1}a_t,\qquad Y_t=\nabla_{m_t}\phi_t,
 \qquad \mathbb E_t\phi_t=0.
 \tag{V50.5}\label{c18v50:poisson}
\]
At fixed finite volume, smoothness, positivity and
connectedness give a smooth centered solution along any
compact time interval. This finite-volume elliptic fact
does \emph{not} assert a lower eigenvalue bound uniform
in volume, coarse value or time.

\begin{proposition}[Material metric criterion]
Let $\Theta_0=I$ and $\dot\Theta_t=Y_t\circ\Theta_t$.
Then $(\Theta_t)_*(p_0\nu_{g,0})=p_t\nu_{g,0}$.
Here time subscripts suppress the fixed $\beta$; write
$m_{\rm ini}=m_{\beta,0}$ for the initial metric, which
need not equal the fixed pair/free reference $m_0$.
Define the symmetric material strain tensor
\[
 \mathcal D_t:=\partial_t m_t+\mathcal L_{Y_t}m_t
      =-\left.\mathcal L_F\overline m_t\right|_{TM_g}
                         +2\operatorname{Hess}_{m_t}\phi_t.
 \tag{V50.6}\label{c18v50:strain}
\]
If an integrable scalar $\sigma(t)$ satisfies
\[
 \mathcal D_t\le\sigma(t)m_t
       \quad\hbox{at every point of }M_g,
 \tag{V50.7}\label{c18v50:gate}
\]
and the initial Poincar\'e gap is at least $\lambda_0$,
then
\[
 \lambda_t\ge\lambda_0
              \exp\left(-\int_0^t\sigma(s)\,ds\right).
 \tag{V50.8}\label{c18v50:gap}
\]
The same conclusion is uniform in $n,g$ only when
$\lambda_0$ and the integral bound are uniform in them.
\end{proposition}

\begin{proof}
Equation \eqref{c18v50:poisson} gives
$\operatorname{div}_{p_t\nu_{g,0}}Y_t=-a_t$.
Together with $\partial_t p_t=a_tp_t$, this is precisely
the continuity equation, proving the pushforward
identity by differentiating $\Theta_t^*(p_t\nu_{g,0})$.
Compactness gives the finite-volume flow. This is the
already-used Poisson/Moser mechanism.

Since $\overline m_t=(\Phi_{-t})^*\overline m_0$,
$\partial_t\overline m_t=-\mathcal L_F\overline m_t$.
Restriction is taken \emph{after} the ambient Lie
derivative, since $F$ need not be tangent to $M_g$.
The identity $\mathcal L_{\nabla_m\phi}m=2\operatorname{Hess}_m\phi$
proves \eqref{c18v50:strain}.

Put $B(t)=\int_0^t\sigma(s)\,ds$. For each initial
tangent vector $v$, differentiation and
\eqref{c18v50:gate} give
\[
 \frac{d}{dt}(\Theta_t^*m_t)(v,v)
       =(\Theta_t^*\mathcal D_t)(v,v)
       \le\sigma(t)(\Theta_t^*m_t)(v,v).
\]
Thus $\Theta_t^*m_t\le e^{B(t)}m_{\rm ini}$. The dual operator
inequality gives
\[
 |d(f\circ\Theta_t)|_{m_{\rm ini}^{-1}}^2
       \le e^{B(t)}|df|_{m_t^{-1}}^2\circ\Theta_t.
\]
Apply the initial Poincar\'e inequality to
$f\circ\Theta_t$, and use the pushforward identity on
both variance and energy. This gives
\eqref{c18v50:gap}. No comparison of unnormalized
densities by their extensive oscillation is used.
\end{proof}

For $|\beta|\le10^{-13}$ the corrected choice
$S_\beta=\mathcal F_\beta$ has $\lambda_0>1/2$ by
V49. For $S_\beta=I$, V37 gives
$\|\operatorname{Hess}_{m_0}\log\Omega_\beta^2\|\le512|\beta|$;
total geodesy and $\operatorname{Ric}_{m_0}\ge m_0$
give $\lambda_0\ge1-512|\beta|>1/2$. These are reused
short-regime inputs. The unproved new premise is the
\emph{integrated actual strain} in
\eqref{c18v50:gate}, not existence of a finite-volume
Poisson solution. An inverse-energy bound alone does not
control its Hessian, and assuming a uniform conditional
gap to prove that premise would require a separate
noncircular argument.

\subsection{Exact native temporal inversion at the electric origin}

Both choices in \eqref{c18v50:initialmap} are $I$ at
$\beta=0$. In the rest of this subsection all quantities
are evaluated at $(\beta,t)=(0,0)$, so $p_0=1$ and
$m_0$ is the pair/free product metric.

A fine edge in direction $i$ is captured precisely when
both transverse coordinates are even. A plaquette in
the $ij$ plane is captured when its remaining coordinate
is even. Such a plaquette has two captured edges in
\emph{distinct} coarse chains and two free edges; an
uncaptured plaquette has four free edges. Write
\[
 W=W_{\rm cap}+W_{\rm unc},\qquad
 W_{\rm cap}=\sum_{p\ {
m captured}}w_p.
 \tag{V50.9}\label{c18v50:split}
\]
Each class has exactly $12n^3$ plaquettes.

\begin{lemma}[Conditional plaquette spectrum and orthogonality]
For every coarse $g$ and every plaquette,
\[
 L_0w_p=\begin{cases}9w_p,&p\ {
m captured},\\
                     12w_p,&p\ {
m uncaptured},\end{cases}
 \qquad
 \int w_p\,d\nu_{g,0}=0,
 \qquad
 \int w_pw_q\,d\nu_{g,0}=\frac14\,\mathbf1_{p=q}.
 \tag{V50.10}\label{c18v50:characters}
\]
\end{lemma}

\begin{proof}
In pair coordinates $(U_1,U_2)=(x,x^{-1}g_\alpha)$,
each plaquette uses at most one half of each chain.
Dependence on any appearing factor is therefore a
fundamental matrix coefficient, whether $x$ or $x^{-1}$
occurs. Its positive round Laplacian eigenvalue is $3$;
for the pair metric it is $3/2$. A captured plaquette
has eigenvalue $2(3/2)+2(3)=9$, and an uncaptured one
has eigenvalue $4(3)=12$.

Each plaquette has a free Haar edge appearing once.
Integrating it proves the zero mean and, for a square,
the character moment $\int w_p^2=1/4$, independently
of the other factors and of $g$.
For distinct plaquettes there is a free edge belonging
to just one of them. To check this last geometric fact,
an uncaptured plaquette has all four edges free, and
two distinct unit plaquettes share at most one edge.
A captured plaquette's two free edges form the adjacent
interior corner of its coarse face; those two directed
geometric edges determine its unit plaquette uniquely.
Consequently two distinct captured plaquettes cannot
have the same free-edge pair, and a free edge in their
symmetric difference integrates the product to zero.
The same single-edge argument covers a mixed pair.
The size assumption excludes periodic identifications
of these local edges. This proves every assertion for
arbitrary $g$, without a flat-coarse specialization.
\end{proof}

\begin{theorem}[Native initial temporal transport and strain]
At $(\beta,t)=(0,0)$, for every $n\ge5$ and every $g$,
\[
 \begin{split}
 a_0&=12(W_{\rm cap}+W_{\rm unc}),\\
 \phi_0=L_0^{-1}a_0&=\frac43W_{\rm cap}+W_{\rm unc},\\
 Y_0&=\nabla_{m_0}\left(\frac43W_{\rm cap}+W_{\rm unc}\right),\\
 \mathcal D_0&=\frac23\operatorname{Hess}_{m_0}W_{\rm cap}.
 \end{split}
 \tag{V50.11}\label{c18v50:main}
\]
In particular,
\[
 -\frac{32}{3}m_0\le\mathcal D_0\le\frac{32}{3}m_0.
 \tag{V50.12}\label{c18v50:uniform}
\]
The entire uncaptured sector cancels from the strain.
This does not say that its temporal source vanishes.
\end{theorem}

\begin{proof}
At the origin $h_0=1$. Equations
\eqref{c18v50:source} and \eqref{c18v50:characters}
give $a_0=12W$ and the unique centered solution
displayed in \eqref{c18v50:main}.

The graph $x\mapsto(x,x^{-1}g_\alpha)$ is totally
geodesic: inversion and translation are isometries of
the round group, so a geodesic of its metric
$2g_{\rm round}$ maps to a product geodesic. The full
pair/free product has the same property. Therefore
\[
 \left.\operatorname{Hess}_{\mathfrak g}W\right|_{TM_g}
     =\operatorname{Hess}_{m_0}(W|_{M_g}),\qquad
 \partial_t m_t|_0=-2\operatorname{Hess}_{m_0}W.
\]
Adding $2\operatorname{Hess}_{m_0}\phi_0$ cancels the
uncaptured coefficient $-2+2=0$ and leaves captured
coefficient $-2+8/3=2/3$.

For completeness, along an ambient product geodesic
with fine-link velocity $v_e$, twice differentiating
the four-factor normalized trace gives
\[
 |\operatorname{Hess}_{\mathfrak g}w_p(v,v)|
      \le\left(\sum_{e\in p}|v_e|\right)^2
      \le4\sum_{e\in p}|v_e|^2.
\]
This follows also by inserting the Lie algebra
generators in the word: their operator norms are
$|v_e|$, including inverse factors. Each fine edge
belongs to at most four captured plaquettes. Sum the
bound and restrict to the totally geodesic fibre to
obtain $|\operatorname{Hess}_{m_0}W_{\rm cap}(v,v)|
\le16m_0(v,v)$. Multiplication by $2/3$ proves
\eqref{c18v50:uniform}. No volume factor occurs in
this quadratic-form estimate.
\end{proof}

\subsection{A native sign obstruction and extensive scalar costs}

\begin{proposition}[The residual is not pointwise contractive]
On the single coarse fibre $g_\alpha=I$ for every
coarse edge, there are configurations and unit tangent
vectors with $\mathcal D_0(v,v)=-4/3$ and $+4/3$.
Thus the choice $\sigma(0)=0$ in
\eqref{c18v50:gate} is false for this Poisson transport.
\end{proposition}

\begin{proof}
The all-identity fine field has every $w_p=1$. An
all-negative-plaquette central field is
\[
 U_1(x)=(-1)^{x_2+x_3}I,\qquad
 U_2(x)=(-1)^{x_3}I,\qquad U_3(x)=I.
 \tag{V50.13}\label{c18v50:centre}
\]
It is periodic since $N$ is even. Direct multiplication
gives every $U_p=-I$ and every captured two-edge product
$I$. Hence both fields belong to the same coarse fibre.

Vary only the free edge in direction $1$ at $(0,1,0)$
by a unit-speed round geodesic. Exactly two captured
plaquettes contain it. If their initial central sign
is $\epsilon=\pm1$, each has $w_p(s)=\epsilon\cos s$.
This is a fibre geodesic, and its velocity is a unit
vector of $m_0$. Thus the captured Hessian on it is
$-2\epsilon$, and \eqref{c18v50:main} gives
$\mathcal D_0(v,v)=-4\epsilon/3$. This proves both
signs using actual native configurations, not a toy
potential. It excludes this particular everywhere
contractive transport claim; it does not exclude an
integrable positive strain bound or another transport.
\end{proof}

\begin{proposition}[Exact scalar volume scaling]
At the electric origin, for every coarse $g$,
\[
 \begin{aligned}
 \mathbb E_0 a_0^2&=864n^3,&
 \mathbb E_0|\nabla_{m_0}W|^2&=63n^3,\\
 \mathbb E_0|Y_0|^2
       =\langle a_0,L_0^{-1}a_0\rangle&=84n^3,&
 \mathbb E_0|\nabla_{m_0}a_0|^2&=9072n^3.
 \end{aligned}
 \tag{V50.14}\label{c18v50:extensive}
\]
In particular $\|Y_0\|_\infty\ge\sqrt{84}\,n^{3/2}$
and $\|\nabla a_0\|_\infty\ge\sqrt{9072}\,n^{3/2}$.
Neither contradicts the uniform strain bound.
\end{proposition}

\begin{proof}
There are $12n^3$ orthogonal characters of each type,
each with squared norm $1/4$. Thus
$\|W\|_2^2=6n^3$ and
$\langle W,L_0W\rangle=(9+12)12n^3/4=63n^3$.
Using the Poisson coefficients $4/3,1$ gives
\[
 \langle a_0,L_0^{-1}a_0\rangle
   =12\left(\frac43+1\right)\frac{12n^3}{4}
   =84n^3.
\]
The other two identities follow by multiplying by
$12^2$. Integration by parts identifies these inner
products with their gradient energies. Supremum norms
dominate probability-space $L^2$ norms.
\end{proof}

One precise consequence for the literature transfer is
that no constants $C,t_0>0$, independent of $n$, can
obey
\[
 \|\nabla_{m_0}\log p_{0,t}\|_\infty\le Ct
       \quad(0\le t\le t_0)
 \tag{V50.15}\label{c18v50:noglobal}
\]
for this path. Indeed, divide by $t$ and send $t$ to
zero at each fixed finite volume. Smoothness gives
$\|\nabla a_0\|_\infty\le C$, contradicting
\eqref{c18v50:extensive} as $n\to\infty$.
This rules out a \emph{uniform linear global
log-Lipschitz budget}, not every possible
dimension-free transport theorem. Local derivative
bounds and cancellations in a metric strain remain
compatible with extensive scalar transport energy.
Also, $a_0=12W$ is not V49's coarse-direction normal
source: the latter vanishes at the origin, whereas
the temporal source does not.

\subsection{Verification, preserved scope and the next unpaid estimate}

The accompanying diagnostic enumerates the native
periodic lattices $n=5,6$, checks captured/free counts,
distinct pair factors, free-edge orthogonality witnesses,
and the all-negative central configuration. It checks
the eigenvalues $9,12$ by quaternion finite differences
for every plane/parity type with genuinely noncommuting
random coarse values. Exact rational checks recover
$4/3,1$, the residual $2/3,0$, all four extensive
coefficients, and the direction of the metric dual
comparison. These are regression checks; the proofs
above establish the all-$n$, all-$g$ statements.
No interacting conditional law is numerically simulated,
and no proof-assistant certification is claimed.

V49 is recovered byte for byte from this source by
removing the V50 append and restoring the old title.
The other 53 source files are hash-protected. Only the
active YM version is replaced; all archive cycles and
companion files remain untouched. Build products and
verification records stay outside \path{latex_notes}.
The next companion checkpoint is V55, the next broad
literature checkpoint V60, and archive/restart remains
V100.

\paragraph{What has been paid.}
The exact temporal conditional source is identified
on the literal Wilson path as well as on the path
starting at V49's corrected map. Its native initial
Poisson solution is explicit on every coarse fibre.
The resulting material metric strain loses the entire
uncaptured sector, has a uniform initial operator
bound, and has an actual positive-sign witness.
This is a concrete new path calculation rather than
another small-$\beta$ normal-source coefficient.

\paragraph{What remains unpaid.}
The next upstream estimate is control of
\[
 -\left.\mathcal L_F\overline m_{\beta,t}\right|_{TM_g}
       +2\operatorname{Hess}_{m_{\beta,t}}
              L_{\beta,t}^{-1}a_{\beta,t}
 \tag{V50.16}\label{c18v50:next}
\]
relative to $m_{\beta,t}$, with a volume/coarse-value
uniform integrable upper bound over the required
time interval. The calculation at $t=0$ does not
bound this tensor at $t>0$, and its smooth
finite-volume dependence alone supplies no uniform
continuation radius. In particular one must not
insert the value $32/3$ into \eqref{c18v50:gap}
for the whole interval $[0,4]$ without a new proof.
The f17 long-time curvature and tube obstructions
are not removed by the initial cancellation.

The appropriate next attack is a spatially localized
Hessian estimate for the \emph{actual temporal}
Poisson solution, or a controlled alternative
solution of the same continuity equation. It must
be obtained without presupposing the sought
long-time conditional gap. If that step succeeds,
the coarse variance/return, variable-bank stability,
weak-coupling scale transport and physical time
calibration still need proofs. Neither branch of
\eqref{c18v50:initialmap} has yet supplied a
nontrivial interacting four-dimensional continuum
construction or its positive physical Yang--Mills
mass gap. The full goal remains open.
\addtocontents{toc}{\protect\setcounter{tocdepth}{2}}
\addtocontents{toc}{\protect\clearpage}
% END CYCLE 18 VERSION 50 APPEND
% BEGIN CYCLE 18 VERSION 51 APPEND
\clearpage
\section[V51: normal transport and unavoidable initial expansion]{V51: normal transport and unavoidable initial expansion}
\label{c18v51:context}

V50 isolates a temporal Poisson Hessian whose control would
transport a conditional gap. We now change coordinates on
the \emph{same literal Wilson path} to remove its tangential
geometric drift before asking for a density correction.
The full normal bank already controlled in V44 constructs
this geometric transport through the entire interval
$0\le t\le4$, uniformly in volume at the level of derivative
majorants. What it does not construct with uniform strain
is the subsequent probability-preserving correction.

This separation reveals a source term not captured by the
centered scores of constant coarse columns. It also upgrades
V50's sign witness: on the full unreduced native fibre,
\emph{every} smooth exact conditional-probability transport
has an initial expansion eigenvalue at least $8/7$ at the
all-negative central field. The bound is sharp at that one
point if arbitrary smooth divergence-free corrections are
allowed. It is not a lower bound on a physical energy gap,
nor a claim of globally optimal transport.

The normal velocity and its initial density derivative were
already identified in V45. The global Gram and lift jets
are V44 results, and V50 supplies the plaquette spectrum.
We reuse them explicitly rather than reprove them as new.
The generic Poisson/Moser mechanism, including the circularity
of assuming its uniform inverse on arbitrary sources, is
already recorded in archive f7 V19 and f17 V36. The new
work concerns this native normal-chart source, the all-time
geometric bound, and the transport-independent trace.

\subsection{An exact normal chart on the literal path}

Let the fine torus have side $N=2n$, $n\ge5$, with the
product unit-round $\mathrm{SU}(2)$ metric $\mathfrak g$.
Write $F=\nabla W$, let $\Phi_t$ be its autonomous forward
flow, and set $c_t=b_2\circ\Phi_t$. The vacuum measure
$\nu=\Omega_\beta^2dU$ is fixed as $t$ changes. In the
original fine and coarse invariant orthonormal frames put
\[
\begin{aligned}
 C_t&=Dc_t,&G_t&=C_tC_t^*,&R_t&=C_t^*G_t^{-1},\\
 Q_t&=R_tC_t,&P_t&=I-Q_t,&q_t&=C_tF,\qquad
 \mathcal N_t=R_tq_t=Q_tF .
\end{aligned}
 \tag{V51.1}\label{c18v51:normal}
\]
Here $q_t$ consists of coarse-frame velocity coefficients,
which can depend on the whole fine configuration; it is
not presumed to descend to a vector field of $c_t(U)$ alone.
$Q_t$ is the orthogonal normal projection to the fibre.
Define the ambient time-dependent diffeomorphism $\eta_t$ by
\[
 \eta_0=I,\qquad \dot\eta_t=-\mathcal N_t\circ\eta_t.
 \tag{V51.2}\label{c18v51:eta}
\]

\begin{proposition}[Global normal identification]
For every finite $n$, $\eta_t$ exists smoothly throughout
$[0,4]$ and satisfies $c_t\circ\eta_t=b_2$. Its restriction
maps $M_g=b_2^{-1}(g)$ diffeomorphically onto $c_t^{-1}(g)$.
No correction of the endpoint $b_2\Phi_4$ is made.
\end{proposition}

\begin{proof}
The sampler is a submersion and $\Phi_t$ a diffeomorphism.
Thus $G_t$ is positive definite; V44 gives the quantitative
global floor in \eqref{c18v44:Gram}. All fields in
\eqref{c18v51:normal} are smooth on a compact finite-dimensional
manifold and compact time interval, so their time-dependent
flow and inverse exist there. Autonomy gives
$D\Phi_tF=F\circ\Phi_t$, hence $\partial_tc_t=C_tF=q_t$.
Since $C_tR_t=I$, differentiation yields
$\partial_t(c_t\eta_t)=q_t-C_t\mathcal N_t=0$.
The equality to $b_2$ follows at time zero, and the inverse
diffeomorphism proves equality, not just inclusion, of fibres.
\end{proof}

\subsection{Paying the normal geometric strain through time four}

Use $c_j,r_j$ from \eqref{c18v44:Cconstants} and
\eqref{c18v44:Rconstants}, with $\ell=2$, $T=4$.
For coefficient derivatives in the same invariant frames
and V42's all-slot norm at weight zero, define
\[
\begin{gathered}
 Q_0^\sharp=r_0c_0,\qquad Q_1^\sharp=r_1c_0+r_0c_1,\\
 Q_2^\sharp=r_2c_0+2r_1c_1+r_0c_2,\\
 n_0=4Q_0^\sharp,\qquad n_1=4Q_1^\sharp+48Q_0^\sharp,\\
 n_2=4Q_2^\sharp+96Q_1^\sharp+576Q_0^\sharp,\qquad
 K_{\rm geom}=n_1+4n_0.
\end{gathered}
 \tag{V51.3}\label{c18v51:constants}
\]
The superscript distinguishes these scalar majorants from
the projection $Q_t$ and its values at different times.

\begin{proposition}[Volume-independent geometric majorants]
Globally in configuration and $0\le t\le4$,
\[
 \mathfrak a_0(\partial^j\mathcal N_t)\le n_j\quad(0\le j\le2),
 \qquad \|D\eta_t\|_{{\rm S},0},\ 
         \|(D\eta_t)^{-1}\|_{{\rm S},0}\le e^{K_{\rm geom}t}.
 \tag{V51.4}\label{c18v51:jets}
\]
Let $m_0=\mathfrak g|_{TM_g}$ and
$\widehat m_t=\eta_t^*\mathfrak g|_{TM_g}$. Then
\[
\begin{gathered}
 e^{-2K_{\rm geom}t}m_0\le\widehat m_t\le e^{2K_{\rm geom}t}m_0,\\
 \partial_t\widehat m_t
   =2\eta_t^*\langle\mathcal N_t,\mathrm{II}_t\rangle,
 \qquad
 -2K_{\rm geom}\widehat m_t\le\partial_t\widehat m_t
                            \le2K_{\rm geom}\widehat m_t .
\end{gathered}
 \tag{V51.5}\label{c18v51:metric}
\]
Here $\mathrm{II}_t(v,w)=(\nabla_vw)^\perp$ on the actual
moving fibre. These constants do not involve a conditional
Poincar\'e inequality.
\end{proposition}

\begin{proof}
Each fine link meets four plaquettes and each plaquette
has twelve scalar link/color slots. Ordered derivatives
of its fundamental trace, with unit generators inserted,
are bounded by one. Anchoring any slot therefore gives
$\mathfrak a_0(\partial^jF)\le4\cdot12^j$ for $j=0,1,2$.
For a one-slot vector this is a component supremum, not
the extensive Euclidean norm of the vector.
Differentiate $Q=RC$ and then $\mathcal N=QF$ in ordered
frames. The connected contractions of V42 give exactly
the three $Q_j^\sharp$ and three $n_j$ above. The two
first-derivative product terms at second order account
for the coefficient $96$; no extensive closed trace is used.

In invariant frames the variational generator for
$V_t=-\mathcal N_t$ consists of its coefficient derivative
and a bracket matrix. The unit-round $\mathfrak{su}(2)$
structure constants have absolute value two; each row
and column of this block-diagonal bracket matrix has at
most two nonzero entries. Its Schur norm is at most $4n_0$.
The full generator has Schur norm at most $K_{\rm geom}$.
The time-ordered integral series bounds both forward
and backward variational propagators by $e^{K_{\rm geom}t}$.
Schur's test gives the corresponding operator bounds
and hence the first line of \eqref{c18v51:metric}.

In a bi-invariant frame the covariant derivative has
half the bracket term, so
$\|\nabla\mathcal N_t\|_{{\rm S},0}\le n_1+2n_0\le K_{\rm geom}$.
For vectors tangent to a fibre, normality gives
$\langle\nabla_v\mathcal N_t,w\rangle
=-\langle\mathcal N_t,\mathrm{II}_t(v,w)\rangle$.
Now differentiate the pullback metric using
$\partial_t\eta_t^*\mathfrak g=\eta_t^*\mathcal L_{-\mathcal N_t}\mathfrak g$.
This proves its exact shape formula and the two-sided
bound. In particular we did not multiply an extensive
$|\mathcal N_t|$ by a separate curvature estimate.
\end{proof}

The majorants are extremely large, because they inherit
V44's full-configuration inverse-Gram constants. Their
volume independence is a mathematical statement, not a
claim of small distortion, a useful physical scale, or
polynomial dependence on block time. Most importantly,
$\eta_t$ alone need not transport the conditional probability.

\subsection{The actual normal-chart source includes a velocity trace}

Denote the fixed pair/free product Haar probability on
$M_g$ by $\mu_g$; it is normalized $m_0$-volume. Define
\[
\begin{aligned}
 k_t(U)&=\Omega_\beta(\eta_tU)^2\operatorname{Jac}\eta_t(U),\\
 \widehat p_t(U)&=\frac{k_t(U)}{\int_{M_g}k_t\,d\mu_g},
 &d\widehat\nu_t&=\widehat p_t\,d\mu_g .
\end{aligned}
 \tag{V51.6}\label{c18v51:density}
\]
The Jacobian is ambient relative to product Haar. Since
$c_t\eta_t=b_2$, disintegration of $\eta_t^*\nu$ under
$b_2$ shows that $\widehat\nu_t$ is precisely the pullback
of the actual conditional law on $c_t^{-1}(g)$.
The coarea Jacobian of $b_2$ is constant and is absorbed
in normalization. There is no extra moving fibre-volume
determinant to multiply into \eqref{c18v51:density};
$\widehat m_t$ is retained separately in the Dirichlet form.

\begin{proposition}[Exact normalized normal-flux source]
With ambient weighted divergence for the fixed $\nu$, put
\[
 d_t=\operatorname{div}_{\nu}\mathcal N_t,
 \qquad \widehat b_t=\partial_t\log\widehat p_t.
\]
Then
\[
\begin{aligned}
 \widehat b_t
   &=-\bigl[d_t\circ\eta_t-
                \mathbb E_{\widehat\nu_t}(d_t\circ\eta_t)\bigr],\\
 d_t&=\sum_\alpha\bigl(q_{t,\alpha}S^R_{t,\alpha}
                         +R_{t,\alpha}q_{t,\alpha}\bigr),
 \qquad S^R_{t,\alpha}=\operatorname{div}_{\nu}R_{t,\alpha}.
\end{aligned}
 \tag{V51.7}\label{c18v51:flux}
\]
The second sum is a derivative of the coarse velocity
coefficients and cannot in general be omitted.
\end{proposition}

\begin{proof}
The Lie derivative of the fixed measure $\nu$ under
$-\mathcal N_t$ is $-d_t\nu$. Consequently
$\partial_t\log k_t=-d_t\circ\eta_t$.
Differentiating the normalizing integral gives its
$\widehat\nu_t$-mean, proving the first identity.
The second is the weighted divergence product rule
applied to $\mathcal N_t=\sum_\alpha q_{t,\alpha}R_{t,\alpha}$.
\end{proof}

There is one useful special case. If
$q_{t,\alpha}=\bar q_{t,\alpha}\circ c_t$ for all $\alpha$,
then $R_{t,\alpha}q_{t,\alpha}
=(E^c_\alpha\bar q_{t,\alpha})\circ c_t$ is fibre-constant.
Centering \eqref{c18v51:flux} then leaves
$\sum_\alpha\bar q_{t,\alpha}(g)r^R_{t,\alpha}$ before
the minus sign and pullback. Without this basicness
property even zero constant-column scores do not imply
zero temporal source. At the electric origin the native
example below has every $S^R_{0,\alpha}=0$ but
$\widehat b_0=3W_{\rm cap}\ne0$.

The derivative of the actual raw flux does nevertheless
have a useful volume-independent \emph{component} bound.
For $0\le\beta\le\beta_*=3/(512e)$ use $a_\beta,h_\beta$
from \eqref{c18v44:vacuumconstants}. Then
\[
 \sup_j|X_jd_t|\le
 B_\beta:=n_2+2h_\beta n_0+2a_\beta n_1 .
 \tag{V51.8}\label{c18v51:fluxgradient}
\]
Indeed, with $u=\log\Omega_\beta$,
\[
 X_jd_t=\sum_iX_jX_i\mathcal N_{t,i}
       +2\sum_i(X_jX_i u)\mathcal N_{t,i}
       +2\sum_i(X_i u)(X_j\mathcal N_{t,i}).
\]
The first diagonal contraction retains the external
$j$ anchor in $\partial^2\mathcal N_t$; the other two
use respectively the Hessian Schur bound and gradient
component bound of \eqref{c18v44:vacuum}. This proves
\eqref{c18v51:fluxgradient}. It is not a uniform bound
on $|\nabla d_t|$ or on the extensive scalar $d_t$.
Neither it nor the geometric metric comparison supplies
a conditional inverse at low frequency.

\subsection{The remaining density correction and its exact initial cost}

On the fixed $M_g$ use the actual pulled-back measure
and metric, and set
\[
 \widehat L_t=-\operatorname{div}_{\widehat\nu_t}
                    \nabla_{\widehat m_t},\qquad
 \chi_t=\widehat L_t^{-1}\widehat b_t,
 \qquad Z_t=\nabla_{\widehat m_t}\chi_t,
 \qquad\mathbb E_{\widehat\nu_t}\chi_t=0.
 \tag{V51.9}\label{c18v51:correction}
\]
At each finite regulator the centred inverse exists
smoothly on this compact connected fibre. Let $\theta_t$
be the flow of $Z_t$ starting at $I$. The continuity
equation gives $(\theta_t)_*\widehat\nu_0=\widehat\nu_t$,
so $\Xi_t=\eta_t\theta_t$ is an exact conditional-measure
transport onto the moving native fibre. Its material
strain in the normal chart is
\[
 \widehat{\mathcal D}_t
  =2\eta_t^*\langle\mathcal N_t,\mathrm{II}_t\rangle
                +2\operatorname{Hess}_{\widehat m_t}\chi_t .
 \tag{V51.10}\label{c18v51:strain}
\]
This follows from \eqref{c18v51:metric} and
$\mathcal L_{\nabla_m\chi}m=2\operatorname{Hess}_m\chi$.
V50's already-proved comparison applies to an integrable
uniform upper bound for \eqref{c18v51:strain}. Only its
first term has been bounded above here through time four.
The finite-volume solution in \eqref{c18v51:correction}
does not give a volume-uniform Hessian estimate.

At $(\beta,t)=(0,0)$ there is an exact reduction:
\[
\begin{aligned}
 \partial_t\widehat m_t|_0&=0,&
 d_0&=-3W_{\rm cap},&\widehat b_0&=3W_{\rm cap},\\
 \chi_0&=\tfrac13W_{\rm cap},&
 \mathbb E_{\mu_g}\widehat b_0^2&=27n^3,&
 \mathbb E_{\mu_g}|Z_0|_{m_0}^2&=3n^3 .
\end{aligned}
 \tag{V51.11}\label{c18v51:initial}
\]

\begin{proof}
The initial fibre is totally geodesic, so its normal
metric derivative is zero. V45 already proves the initial
coarea density derivative $3W_{\rm cap}$ in
\eqref{c18v45:Ricvariation}. To check the raw divergence
as well, on a chain $(A,B)$ its normal column in direction
$v$ has the flow
\[
 (A,B)\longmapsto
 \bigl(e^{hv/2}A,\ A^{-1}e^{hv/2}AB\bigr),
 \qquad AB\longmapsto e^{hv}AB .
 \tag{V51.12}\label{c18v51:pairflow}
\]
This is indeed its autonomous flow: its coarse product
evolves as shown and its second component differentiates
as $A(h)^{-1}vA(h)B(h)/2$. It preserves product Haar
by its triangular change of variables. Therefore every
$S^R_{0,\alpha}$ vanishes. Since $G_0=2I$,
$q_{0,\alpha}=2R_{0,\alpha}W$, and
$d_0=2\sum_\alpha R_{0,\alpha}^2W$.
A captured plaquette meets two distinct chains and at
most one half of each. Inserting each unit generator
at half speed gives $\sum_{a=1}^3R_{\alpha,a}^2w_p
=-3w_p/4$ on each such chain. Uncaptured plaquettes
have no such contribution. Thus $d_0=-3W_{\rm cap}$
pointwise, agreeing with V45's coarea calculation.

V50's \eqref{c18v50:characters} gives zero conditional
mean, $L_0W_{\rm cap}=9W_{\rm cap}$, and
$\|W_{\rm cap}\|_2^2=12n^3/4=3n^3$ for every $g$.
It proves all remaining identities in
\eqref{c18v51:initial} by inversion and integration by
parts. In particular the variance is $9(3n^3)$ and
the inverse energy is that variance divided by nine.
\end{proof}

This example also proves that $q_0$ is not basic: if it
were, the special case following \eqref{c18v51:flux}
and the zero column scores would force $\widehat b_0=0$,
contradicting its positive variance. The scalar energy
$3n^3$ differs from V50's $84n^3$ because the coordinate
motion and hence the correction source have changed.
It is not a factor-28 improvement in the same optimization
problem. The resulting initial Poisson material strain is
still exactly $(2/3)\operatorname{Hess}_{m_0}W_{\rm cap}$.

\subsection{An expansion obstruction for every exact transport}

\begin{theorem}[Universal initial trace and the $8/7$ bound]
\label{c18v51:universal}
Work at $\beta=0$. Let $\Xi_t:M_g\to c_t^{-1}(g)$
be any smooth family of diffeomorphisms with $\Xi_0$
the inclusion, transporting $\mu_g$ to the actual
conditional probability. Its initial material metric
strain $\mathcal D_0^\Xi=\partial_t(\Xi_t^*\mathfrak g)|_0$
satisfies the pointwise identity
\[
 \operatorname{tr}_{m_0}\mathcal D_0^\Xi=-6W_{\rm cap}.
 \tag{V51.13}\label{c18v51:trace}
\]
On the fibre $g_\alpha=I$, at the all-negative central
configuration $p_-$ of \eqref{c18v50:centre},
\[
 \lambda_{\max}\bigl(m_0^{-1}\mathcal D_0^\Xi(p_-)\bigr)
                         \ge\frac87 .
 \tag{V51.14}\label{c18v51:expansion}
\]
No gradient-velocity, minimization or gauge-equivariance
assumption is imposed on $\Xi_t$.
\end{theorem}

\begin{proof}
Compose with the inverse normal chart:
$\theta_t=\eta_t^{-1}\Xi_t$ is a diffeomorphism of $M_g$
transporting $\widehat\nu_0$ to $\widehat\nu_t$.
Its initial velocity $Z$ is tangent. Equivalently,
differentiating the fibre constraint forces the normal
part of $\dot\Xi_0$ to equal $-\mathcal N_0$, so
$Z=\dot\Xi_0+\mathcal N_0$. The continuity equation
and \eqref{c18v51:initial} give
$\operatorname{div}_{\mu_g}Z=-3W_{\rm cap}$.
Because $\mu_g$ is constant times $m_0$-volume and
the normal metric derivative vanishes identically,
$\mathcal D_0^\Xi=\mathcal L_Zm_0$ and
$\operatorname{tr}_{m_0}\mathcal D_0^\Xi
=2\operatorname{div}_{m_0}Z=-6W_{\rm cap}$.

The fine graph has $24n^3$ links and the coarse graph
$3n^3$ links. Fixing each disjoint two-link product
removes one $\mathrm{SU}(2)$ factor. Thus $M_g$ has
$21n^3$ independent group factors and real dimension
$d=63n^3$. At $p_-$ all $12n^3$ captured plaquette
characters equal $-1$, giving trace $72n^3$.
The maximum eigenvalue of a symmetric endomorphism is
at least its trace divided by its dimension, which is
$72/63=8/7$. This proves the bound for every such transport.
\end{proof}

This forbids an everywhere initially nonexpanding exact
transport on the \emph{full unreduced} fibre. It does
not forbid an integrable positive expansion budget,
and therefore is not a no-go theorem for the mass gap.
Nor does this dimension/trace argument automatically
give the same number after internal gauge reduction;
the orbit-volume and quotient-metric terms would have
to be included. The obstruction concerns physical
conditional measures but not yet a physical time scale.

\subsection{Sharpness at one point by solenoidal jet adjustment}

The lower bound is not just a nonoptimal average forced
by our selected Poisson velocity. The following local
divergence-free construction specifies exactly its scope.

\begin{lemma}[A prescribed traceless symmetric jet]
On a smooth Riemannian manifold of dimension $d\ge2$,
with measure a constant multiple of volume, let $p$
be a point and $B:T_pM\to T_pM$ a symmetric traceless
endomorphism. There is a smooth vector field $K$, supported
in a coordinate neighborhood of $p$, such that
$\operatorname{div}K=0$, $K(p)=0$, and $\nabla K(p)=B$.
\end{lemma}

\begin{proof}
Use orthonormal Riemann normal coordinates $x$ at $p$,
writing volume as $\rho(x)dx$ with $\rho(0)=1$.
Choose a smooth compactly supported cutoff $\zeta=1$
near zero inside the chart. Put
\[
 A_{ij}(x)=\frac{(Bx)_i x_j-(Bx)_j x_i}{d},\qquad
 K^i(x)=\rho(x)^{-1}\sum_j\partial_j(\zeta A_{ij})(x).
 \tag{V51.15}\label{c18v51:jet}
\]
Antisymmetry gives
$\sum_i\partial_i(\rho K^i)=
\sum_{i,j}\partial_i\partial_j(\zeta A_{ij})=0$.
Since $\operatorname{tr}B=0$, direct differentiation
gives $\sum_j\partial_j A_{ij}=(Bx)_i$.
Near zero $K=\rho^{-1}Bx$, so $K(0)=0$ and
$DK(0)=B$; its covariant derivative there is also $B$.
Extension by zero is smooth because of the compact
support of the cutoff. This proves the lemma.
\end{proof}

\begin{proposition}[Pointwise sharpness on the native fibre]
For fixed finite $n$ on $g=I$ at $p_-$, there exists
a smooth exact conditional transport, starting at the
inclusion, with
\[
             \mathcal D_0^\Xi(p_-)=\frac87m_0(p_-).
 \tag{V51.16}\label{c18v51:sharp}
\]
Thus the infimum of the largest initial strain eigenvalue
at this one point, over all these transports, is $8/7$.
\end{proposition}

\begin{proof}
Start with the Poisson solution in the normal chart.
Identify its initial strain at $p_-$ with the symmetric
endomorphism $D_P=(2/3)m_0^{-1}\operatorname{Hess}_{m_0}W_{\rm cap}$.
Theorem \ref{c18v51:universal} gives
$\operatorname{tr}D_P=(8/7)d$. Therefore
$B=\tfrac12[(8/7)I-D_P]$ is symmetric and traceless.
Apply the lemma to obtain $K_0$ with this jet and
$\operatorname{div}_{\mu_g}K_0=0$.

To realize it by an \emph{exact} time-dependent transport,
not just a formal first jet, set
$K_t=(\widehat p_0/\widehat p_t)K_0=K_0/\widehat p_t$
at $\beta=0$ on the fixed normal chart. Then
$\operatorname{div}_{\widehat\nu_t}K_t=0$ at every time.
Adding $K_t$ to the smooth Poisson velocity in
\eqref{c18v51:correction} leaves its continuity equation
unchanged. At each fixed finite volume this velocity
has a smooth flow on $[0,4]$. Compose that flow with
$\eta_t$ to obtain the required exact transport.
At time zero its metric strain gains
$\mathcal L_{K_0}m_0(p_-)=2m_0B$, proving
\eqref{c18v51:sharp}; the universal lower bound proves
the assertion about the infimum.
\end{proof}

This local construction asserts neither a globally
optimal maximum strain nor a uniform bound on $K_t$
as volume grows. It is not claimed to be gauge-equivariant
or smooth as a selection over all coarse values. It
shows that the native pointwise obstruction persists
after allowing arbitrary solenoidal corrections, and
identifies its sharp value in that unrestricted class.

\subsection{Transfer audit, verification and the next unpaid estimate}

The primary context was checked against
\href{https://arxiv.org/abs/math/0211065}{Lott,
\emph{Some geometric properties of the Bakry--Emery--Ricci tensor}}:
its fibrewise measure-transport condition is a substantive
hypothesis, not a consequence of having a normal right inverse.
We do not assert that $b_2\Phi_t$ is a Riemannian submersion
to the original product coarse metric. The Euclidean
heat-flow results of
\href{https://arxiv.org/abs/2404.15205}{Brigati--Pedrotti,
\emph{Heat flow, log-concavity, and Lipschitz transport maps}}
likewise do not supply a Hessian estimate for this compact
conditional Poisson equation. No theorem from those
papers is imported with unverified hypotheses.

The new diagnostic enumerates the native $n=5,6$
incidences and the same-fibre central witness. On all
24 plane/parity types it checks the full-pair normal
flow and $2\sum R_\alpha^2w_p$ by quaternion finite
differences. Symbolic tests check the antisymmetric
potential, weighted divergence and jet in dimensions
$3,4,6$, and rational tests check $8/7$, $27$, and $3$.
These are regression fixtures; the displayed proofs,
not finite-size testing, give the all-volume statements.
There is no interacting conditional simulation or
proof-assistant certification.

Removing this appendix and restoring the prior title
recovers V50 byte for byte. The other 53 sources are
hash-protected and unchanged. Only the latest active
YM version is retained, with only \texttt{.tex} files
in \path{latex_notes}; build products stay elsewhere.
The last companion and broad literature checkpoint was
V50; the next are V55 and V60 respectively. Archive and
restart remain scheduled at V100.

\paragraph{Direction after this calculation.}
The actual geometry of the normal identification is
now controlled through the required block time, with
poor but volume-independent constants. The remaining
transport estimate is source-specific control of
\[
 \operatorname{Hess}_{\widehat m_t}
       \widehat L_t^{-1}
       \bigl[-\operatorname{center}_{\widehat\nu_t}
          ((\operatorname{div}_{\nu}(Q_tF))\circ\eta_t)\bigr]
 \tag{V51.17}\label{c18v51:next}
\]
or of a solenoidal modification of the same continuity
equation. One must retain the velocity-derivative trace
in \eqref{c18v51:flux}. Replacing it by a constant-column
score, or assuming a uniform inverse on arbitrary
centered sources, would bypass rather than solve the
remaining low-spectrum problem. The $8/7$ obstruction
also rules out pursuing zero expansion as the needed
bound; a positive integrable upper budget would suffice.

No uniform bound for \eqref{c18v51:next} on $[0,4]$ has
been proved here. Even such a conditional bound would
leave coarse variance/return, variable-bank stability,
weak-coupling scale transport, physical time calibration,
and interacting continuum reconstruction to be proved.
The full four-dimensional Yang--Mills mass gap remains
open in this programme.
\addtocontents{toc}{\protect\setcounter{tocdepth}{2}}
\addtocontents{toc}{\protect\clearpage}
% END CYCLE 18 VERSION 51 APPEND
% BEGIN CYCLE 18 VERSION 52 APPEND
\clearpage
\section[V52: low-band temporal selection and the connected return]{V52: low-band temporal selection and the connected return}
\label{c18v52:context}

V51 proves that every exact probability transport must
expand at some initial native configuration. That is a
pointwise statement, not the derivative of a spectral
gap. We now test the actual conditional operator on its
lowest electric eigenspace. Its first variation vanishes
on an entire low-energy band, even though the pointwise
transport obstruction is nonzero. The first resonant
energy and the complete second-order virtual return from
the lowest eigenspace can be calculated exactly.

This moves the next calculation from a worst-configuration
transport bound to the actual low-energy spectral block.
It does not replace the requested long-time and continuum
theorems with a perturbative one. In particular the direct
second variation of the conditional metric and density
is kept as an unpaid term, and no regulator-uniform
remainder is inferred from the exact coefficients.

All results in this appendix concern the literal
$c_t=b_2\Phi_t$ path at $\beta=0$, $N=2n$, $n\ge5$,
and every fixed coarse value $g$. The fibre is the full
unreduced one. V45's product metric, V50's captured
plaquette character calculation, and V51's normal chart
are reused. The earlier connected-return arguments in
this cycle concern different Hamiltonian perturbations;
the operator here is the \emph{flow-time conditional}
kinetic operator, not $T-\beta W$ with $t$ renamed
as a coupling.

\subsection{The actual first operator variation in a fixed Hilbert space}

Use $M_g$, $\mu_g$, $\widehat m_t$, $\widehat p_t$ and
$\widehat L_t$ from V51. Multiplication by
$\widehat p_t^{1/2}$ is a unitary map from
$L^2(\widehat\nu_t)$ to $L^2(\mu_g)$. Define
\[
 \mathcal K_t=\widehat p_t^{1/2}\widehat L_t
                            \widehat p_t^{-1/2},\qquad
 \mathcal K_t=L_0+tV+t^2A_2+O_{n,g}(t^3).
 \tag{V52.1}\label{c18v52:operator}
\]
The Taylor convention makes $A_2=\mathcal K''_0/2$.
At fixed regulator these are smooth elliptic operators
with common domain $H^2(M_g)$ in the fixed Hilbert space;
the remainder is in the norm $H^2\to L^2$. For this
appendix constants in such Taylor remainders are allowed
to depend on the regulator. No uniformity in $n$ is
hidden in this notation.

\begin{proposition}[A multiplication first variation]
For every coarse value,
\[
                 V=-\frac{27}{2}W_{\rm cap}.
 \tag{V52.2}\label{c18v52:V}
\]
There is no first-order differential part in this
normal-chart, fixed-Hilbert-space variation.
\end{proposition}

\begin{proof}
By \eqref{c18v51:initial},
$\widehat m'_0=0$ and
$b:=\partial_t\log\widehat p_t|_0=3W_{\rm cap}$.
Consequently $\widehat L'_0 f=-\langle db,df\rangle_{m_0^{-1}}$.
Differentiating the two square-root multiplication factors
in \eqref{c18v52:operator} gives
\[
 \mathcal K'_0f
 =\tfrac12 bL_0f-\langle db,df\rangle
                         -\tfrac12L_0(bf)
 =-\tfrac12(L_0b)f .
\]
Here the positive Laplacian product identity is
$L_0(bf)=bL_0f+fL_0b-2\langle db,df\rangle$.
Since $L_0W_{\rm cap}=9W_{\rm cap}$ by
\eqref{c18v50:characters}, the coefficient is $-27/2$.
The vanishing metric derivative is essential to this
simple formula. It does not hold for an arbitrary
choice of fibre coordinates without the associated
unitary coordinate correction.
\end{proof}

For every $g$, $\|V\|\le162n^3$; equality holds on
$g=I$ since $W_{\rm cap}=\pm12n^3$ at the all-identity
and all-negative configurations. An extensive full
operator norm therefore supplies no volume-independent
smallness argument.

\subsection{A sharp low-energy selection rule}

The independent pair coordinates have metric
$2g_{\rm round}$, and the free coordinates have
$g_{\rm round}$. For a tensor Peter--Weyl vector with
degrees $\ell_j=0,1,\ldots$, its electric energy is
\[
 E(\boldsymbol\ell)
   =\frac12\sum_{j\ {\rm pair}}\ell_j(\ell_j+2)
        +\sum_{j\ {\rm free}}\ell_j(\ell_j+2).
 \tag{V52.3}\label{c18v52:energy}
\]
Let $\Pi_E$ be the spectral projection of $L_0$ for $E$.

\begin{theorem}[No matrix elements below energy sum nine]
For any two electric energies,
\[
 E+E'<9\quad\Longrightarrow\quad
                    \Pi_{E'}V\Pi_E=0 .
 \tag{V52.4}\label{c18v52:selection}
\]
Thus $\Pi_{[0,\Lambda]}V\Pi_{[0,\Lambda]}=0$
for $\Lambda<9/2$. The threshold is sharp: the
$E=9/2$ compression has both positive and negative
Rayleigh values, of magnitude at least $27/8$.
\end{theorem}

\begin{proof}
A captured plaquette uses two distinct pair factors
and two distinct free factors, once each. As a function
of any one appearing factor it is a fundamental matrix
coefficient, possibly with inverse or fixed coarse
matrices inserted. Multiplication sends degree $\ell$
only to $\ell'=\ell+1$ or $\ell-1$; this is
$\ell\otimes1=(\ell+1)\oplus(\ell-1)$ in degree notation,
with a negative degree omitted.
For each allowed nonnegative pair,
$\ell(\ell+2)+\ell'(\ell'+2)\ge3$.
The sum of the two state energies over the four
appearing factors is therefore at least
$3/2+3/2+3+3=9$. Other factors have nonnegative energy.
This proves the rule term by term, for arbitrary
coarse matrices and all representation degrees.

For sharpness choose the plaquette in the $12$ plane
at $(0,0,0)$ and write its oriented word
$AB C^{-1}D^{-1}$. Its first and last edges $A,D$
are first halves of distinct captured chains; $B,C$
are free. Set
\[
 f=\Tr(AB),\qquad h=\Tr(DC),\qquad
 w_p=\tfrac12\Tr\bigl((AB)(DC)^{-1}\bigr).
 \tag{V52.5}\label{c18v52:witness}
\]
Both functions have norm one, mean zero and energy
$3/2+3=9/2$, and they are orthogonal. The products
$AB$ and $DC$ are independent Haar variables, even
when the unused second chain halves are constrained
by arbitrary $g$. Fundamental character orthogonality
gives $\langle h,w_pf\rangle=1/4$.

No other captured plaquette contributes to
$\langle h,W_{\rm cap}f\rangle$. Its two free edges
would have to be exactly $B,C$; those two free edges
determine the captured plaquette uniquely, as in V50.
Otherwise a free fundamental variable appears only once
and integrates to zero. Every diagonal contribution
to $\langle f,W_{\rm cap}f\rangle$ or
$\langle h,W_{\rm cap}h\rangle$ also vanishes: each
test depends on only one free edge whereas each
captured plaquette needs two. Thus the compression
of $V$ to the span of $f,h$ is
\[
             \begin{pmatrix}0&-27/8\\-27/8&0\end{pmatrix}.
 \tag{V52.6}\label{c18v52:resonance}
\]
The unit vectors $(f+h)/\sqrt2$ and $(f-h)/\sqrt2$
give the asserted Rayleigh values. This two-dimensional
span is not asserted to be invariant under the entire
$E=9/2$ compressed operator.
\end{proof}

\subsection{The lowest gap is stationary to first order}

Write $E_*=3/2$ and $\Pi=\Pi_{E_*}$.
This eigenspace is the orthogonal direct sum of the
four-dimensional fundamental coefficient spaces
$\mathcal H_i$ of the $3n^3$ pair factors, with every
other factor constant. Its dimension is $12n^3$.
The next distinct positive eigenvalue is three.

\begin{corollary}[Actual conditional gap at the electric origin]
At each fixed finite $n$ and coarse $g$, the full
unreduced conditional gap satisfies
\[
        \lambda_{\rm fib}(t,g)=\frac32+O_{n,g}(t^2).
 \tag{V52.7}\label{c18v52:gap}
\]
In particular its right derivative at zero is zero.
This is not a regulator-uniform continuation estimate.
\end{corollary}

\begin{proof}
On a small circle enclosing $E_*$ but no other
eigenvalue of $L_0$, the resolvent identity and the
$H^2\to L^2$ Taylor expansion in
\eqref{c18v52:operator} give a smooth Riesz projection
$\Pi(t)$ with the same finite rank. For sufficiently
small $t$ at this fixed regulator, the polar part of
$\Pi(t)\Pi+(I-\Pi(t))(I-\Pi)$ provides a smooth unitary
identification of the two spectral subspaces.
The compressed operator in this identification has
first derivative $\Pi V\Pi$; differentiating the
unitary only adds a commutator with $L_0$, whose
compression to the scalar $E_*$ block is zero.
Equation \eqref{c18v52:selection} makes that derivative
zero. Finite-dimensional spectral stability therefore
puts every eigenvalue in this cluster at $E_*+O(t^2)$.

The zero eigenvalue stays exactly zero and simple:
its normalized vector in the fixed Hilbert space is
$\widehat p_t^{1/2}$. Compactness, positive density,
and ellipticity at fixed regulator imply norm-resolvent
continuity and preserve the number of eigenvalues
below an isolating cutoff between $3/2$ and three.
Hence this cluster gives the first positive eigenvalue
for small $t$. These facts give \eqref{c18v52:gap},
not a lower bound at $t=4$.
\end{proof}

V51's $8/7$ pointwise expansion obstruction and
\eqref{c18v52:gap} therefore coexist on the same
native model. The strongest pointwise transport
criterion can lose information relevant to the
lowest spectral block. Conversely the zero first
variation says nothing about the sign of its
second variation.

\subsection{The complete virtual return from the lowest block}

Let $\Pi^\perp=I-\Pi$, and define the reduced resolvent
on that complement and the return
\[
 R_*=\Pi^\perp(L_0-E_*)^{-1}\Pi^\perp,\qquad
 \mathcal S_*=\Pi VR_*V\Pi .
 \tag{V52.8}\label{c18v52:return}
\]
The inverse has a negative ground-state denominator,
but $V\Pi$ has no ground-state component by
\eqref{c18v52:selection}. All denominators that
actually contribute below are positive.

\begin{theorem}[Exact extensive part and connected remainder]
On the entire degenerate lowest block, for every $g$,
\[
 \mathcal S_*=
       \left(\frac{243}{4}n^3+\frac{81}{40}\right)\Pi .
 \tag{V52.9}\label{c18v52:mainreturn}
\]
The corresponding ground-state return is
\[
 s_{\rm vac}:=
   \langle V1,L_0^{-1}V1\rangle=\frac{243}{4}n^3.
 \tag{V52.10}\label{c18v52:vacuum}
\]
Consequently $\mathcal S_*-s_{\rm vac}\Pi=(81/40)\Pi$.
This is a return coefficient, not yet the whole
second-order spectral shift.
\end{theorem}

\begin{proof}
Put $a=27/2$ so that $V=-a\sum_p w_p$, with the sum
over captured plaquettes. Distinct plaquettes do not
cross in $\Pi w_qR_*w_p\Pi$: their free-edge pairs
are distinct, leaving a free factor of degree one
on one side and degree zero on the other.
The product resolvent preserves each factor's
representation degree, so it cannot remove this
orthogonality. This argument applies to every
initial and final pair factor.

For $p=q$, the operator $w_pR_*w_p$ preserves
the sign parity of every pair coordinate. Each
appearance of $w_p$ flips the parity on its two
pair factors, and the resolvent commutes with all
individual sign changes. Thus the return cannot
move $\mathcal H_i$ to a different $\mathcal H_j$.
It remains to compute its quadratic form for
$f_i\in\mathcal H_i$ of norm one.

If $p$ does not use factor $i$, $w_pf_i$ is an
eigenfunction of energy $E_*+9$, and has squared
norm $1/4$. Its return without $a^2$ is $1/36$.
If $p$ does use $i$, write that unit quaternion
as $x$ and $f_i(x)=2\sum_{\mu=0}^3c_\mu x_\mu$,
where $\sum|c_\mu|^2=1$. With the other three
factors collected, $w_p=x\cdot z$. Fixed coarse
rotations or inversion merely rotate these real
coordinates. The vector $z$ is Haar distributed
on $S^3$, and each of its coefficients belongs
to degree one in each of those other three factors.

The degree-zero projection in $x$ is
$(1/2)\sum_\mu c_\mu z_\mu$. Since
$\int x_\mu x_\nu\,dH=\delta_{\mu\nu}/4$,
its squared norm is $1/16$.
The total norm of $w_pf_i$ is $1/4$, obtained
by first integrating a free Haar factor.
The remaining degree-two component therefore
has norm squared $3/16$. These statements hold
for complex coefficients by the same Hermitian
moment calculation.

The other pair contributes energy $3/2$ and
the two free factors contribute six. The two
intermediate energies and denominators are thus
\[
 \begin{array}{c|c|c|c}
 \text{degree on }i&\text{squared norm}&
               \text{intermediate energy}&E'-E_*\\ \hline
 0&1/16&15/2&6\\
 2&3/16&23/2&10
 \end{array}
 \qquad
 \frac{1/16}{6}+\frac{3/16}{10}=\frac7{240}.
 \tag{V52.11}\label{c18v52:channels}
\]
The value is independent of $f_i$, so polarization
makes the return scalar on $\mathcal H_i$.
Each pair contains two fine links, each in four
captured plaquettes; no unit plaquette uses both
halves. Exactly eight captured plaquettes use $i$.
There are $12n^3$ captured plaquettes altogether.
Consequently the return on every $\mathcal H_i$ is
\[
 \begin{aligned}
 a^2\left(\frac{12n^3-8}{36}
                         +8\frac7{240}\right)
 &=\frac{243}{4}n^3+\frac{81}{40},\\
 \frac7{240}-\frac1{36}&=\frac1{720}.
 \end{aligned}
 \tag{V52.12}\label{c18v52:sum}
\]
This proves \eqref{c18v52:mainreturn} on the whole
block, not just for an isolated test function.
Finally $V1=-aW_{\rm cap}$, whose squared norm is
$a^2(3n^3)$ and whose electric energy is nine.
Division by nine proves \eqref{c18v52:vacuum}.
\end{proof}

There is a useful warning about global spectral
subspace perturbation. Squaring the denominators in
the same proof gives
\[
 \Pi V R_*^2V\Pi
      =\left(\frac{27}{4}n^3+\frac{153}{200}\right)\Pi .
 \tag{V52.13}\label{c18v52:dressing}
\]
Indeed the disjoint term is $1/324$, the incident
term is $1/576+3/1600=13/3600$, and their difference
is $17/32400$. Multiplication by $8a^2$ gives
$153/200$. Differentiating the spectral projection
shows $\Pi^\perp\Pi'(0)\Pi=-R_*V\Pi$.
Thus the norm of this off-diagonal derivative grows
like $n^{3/2}$ even though the connected return
\eqref{c18v52:mainreturn} is volume-independent after
the vacuum subtraction. This forbids a uniform
bound on that \emph{global projection derivative};
it does not forbid uniform spectral-gap bounds,
local dressing estimates, or a different comparison.

\subsection{What the exact second-order coefficient still contains}

The direct coefficient $A_2$ in
\eqref{c18v52:operator} must not be set to zero.
In fact the exact zero eigenvalue already forces
its extensive vacuum expectation:
\[
                  \langle1,A_2 1\rangle=s_{\rm vac}.
 \tag{V52.14}\label{c18v52:calibration}
\]
To see this, write the normalized ground vector
$\widehat p_t^{1/2}=1+t\,b/2+O(t^2)$.
The first-order equation gives
$b/2=-L_0^{-1}V1=(3/2)W_{\rm cap}$.
Take the inner product of the second-order ground
equation with $1$ to obtain
$\langle1,A_2 1\rangle
=-\langle1,V(b/2)\rangle=s_{\rm vac}$.
This is forced by the actual conditional operator,
not an adjustable subtraction of physical energy.

Define the vacuum-subtracted direct block
\[
 \mathcal A_{\rm conn}
       :=\Pi A_2\Pi-\langle1,A_2 1\rangle\Pi .
 \tag{V52.15}\label{c18v52:direct}
\]
The exact effective second coefficient at $E_*$ is
\[
 \mathcal B_*=\mathcal A_{\rm conn}-\frac{81}{40}\Pi .
 \tag{V52.16}\label{c18v52:effective}
\]
For completeness, solve the complementary first-order
eigenvector equation with initial vector $f\in\Pi$:
its orthogonal correction is $-R_*Vf$. Substitution
into the projected second-order equation yields
$\Pi A_2\Pi-\Pi VR_*V\Pi$. Equations
\eqref{c18v52:mainreturn} and
\eqref{c18v52:calibration} give
\eqref{c18v52:effective}. Equivalently one may use
the smoothly transported Riesz subspace of the
previous corollary; its second effective coefficient
is the same because $\Pi V\Pi=0$.
In finite dimension its eigenvalues determine the
$t^2$ coefficients of the cluster, with a fixed-regulator
$O(t^3)$ remainder. In particular
\[
 \lambda_{\rm fib}(t,g)=\frac32+
     t^2\lambda_{\min}(\mathcal B_*)+O_{n,g}(t^3).
 \tag{V52.17}\label{c18v52:secondgap}
\]
The sign cannot be read off from $-81/40$ alone:
$\mathcal A_{\rm conn}$ is not yet evaluated.

Here is a concrete expression for the missing term
that retains both the density and the moving metric.
Expand the inverse metric and log density on $M_g$ as
\[
 \widehat m_t^{-1}=m_0^{-1}+t^2a_2+O(t^3),
 \qquad \log\widehat p_t=t b+t^2c+O(t^3),
 \qquad b=3W_{\rm cap}.
 \tag{V52.18}\label{c18v52:seconddata}
\]
For a real smooth test $f$, the exact fixed-space
quadratic form is
\[
 \langle f,\mathcal K_tf\rangle
 =\int\widehat m_t^{-1}
   \left(df-\tfrac f2d\log\widehat p_t,\,
         df-\tfrac f2d\log\widehat p_t\right)d\mu_g .
\]
Consequently
\[
 \langle f,A_2f\rangle
 =\int\left[
       a_2(df,df)-f\langle df,dc\rangle_{m_0^{-1}}
                       +\frac{f^2}{4}|db|_{m_0^{-1}}^2
        \right]d\mu_g .
 \tag{V52.19}\label{c18v52:directform}
\]
Polarization gives the complete bilinear form and its
complex extension. The tensor $a_2$ and scalar $c$
come from the actual normal chart, not a replacement
Gibbs density. Omitting either of their contributions
would compute a different operator.

\subsection{Audit and direction of the next calculation}

The relevant primary perturbation context was checked
against \href{https://arxiv.org/abs/1903.00785}
{Greenbaum--Li--Overton, \emph{First-order Perturbation
Theory for Eigenvalues and Eigenvectors}}. Its main
simple-eigenvalue setting is not substituted for the
$12n^3$-fold block here; the projection and complementary
equations needed for that block were given above.
The usual representation basis is described in
\href{https://arxiv.org/abs/2307.11829}
{Bauer--D'Andrea--Freytsis--Grabowska,
\emph{A new basis for Hamiltonian SU(2) simulations}}.
We use exact Peter--Weyl selection, not a truncation
as a proxy for the continuum weak-coupling limit.
This is a targeted check, not the scheduled broad
literature sweep.

The new diagnostic checks the eight-plaquette load of
each actual coarse pair at $n=5,6$, free-pair uniqueness,
and the native resonant plaquette. Exact $S^3$ moments
check the full four-dimensional channel Gram matrices
$I/4$, $I/16$, and $3I/16$ and the character matrix
element $1/4$. Rational tests check both returns and
the projection-derivative coefficient. A finite
representation enumeration supplements, but does not
replace, the all-degree proof of the energy-sum rule.
No interacting conditional law or long-time spectral
gap is numerically simulated.

V51 is recovered byte for byte by removing this append
and restoring its title. The other 53 sources remain
hash-protected and unchanged. Only the latest active
YM source is retained in \path{latex_notes}, with all
build products outside it. Companion review remains
due at V55, broad literature review at V60, and archive
and restart at V100.

\paragraph{New evidence and the next unpaid term.}
The native lowest conditional spectral block does not
see the first-order perturbation. There is nevertheless
a genuine resonance at $9/2$, and a nonzero connected
virtual return at second order. Thus neither worst-point
initial expansion nor a blanket assertion that all low
modes decouple adequately describes this path.
The immediate upstream calculation is the actual
vacuum-subtracted direct form
\eqref{c18v52:directform} on $\Pi$, followed by a
regulator-uniform remainder or a nonperturbative
low-block comparison. Another generic finite-volume
perturbation theorem would not pay that remainder.

Even a successful temporal conditional-gap continuation
would leave the coarse variance/return and variable-bank
stability estimates, weak-coupling scale transport,
physical time calibration and interacting continuum
construction. The full Yang--Mills mass gap is still
unproved. The results here concern the unreduced
conditional kinetic spectrum at the electric origin,
not a newly constructed physical continuum Hamiltonian.
\addtocontents{toc}{\protect\setcounter{tocdepth}{2}}
\addtocontents{toc}{\protect\clearpage}
% END CYCLE 18 VERSION 52 APPEND
% BEGIN CYCLE 18 VERSION 53 APPEND
\clearpage
\section{V53: The complete initial quadratic conditional-gap shift}
\label{c18v53:section}

\subsection{Context, inherited input and the new coefficient}

This continuation evaluates the direct second coefficient left
open in V52. We keep exactly its native setting: the spatial
torus has side $N=2n$, $n\ge5$, the group is unit-round
$\mathrm{SU}(2)$, $\beta=0$, and $c_t=b_2\Phi_t$ is the literal
positive Wilson-gradient flow followed by block-two multiplication.
The full unreduced conditional fibre is used throughout.
No physical Hamiltonian is identified with this conditional operator.

The base fibre $M_g$ is the totally geodesic product of $3n^3$
pair factors
$(A_i,B_i)=(x_i,x_i^{-1}g_i)$, with metric twice the round
metric, and $18n^3$ free round factors. Write $m=m_0$ and
$\mu=\mu_g$ for its metric and normalized product Haar measure.
The normal chart of V51 has velocity $-N_t=-Q_tF$ and
$\widehat m'_0=0$. Its fixed-space conditional operator is
$\mathcal K_t=L_0+tV+t^2A_2+O_{n,g}(t^3)$, with
$V=-(27/2)W_{\rm cap}$, as proved in
\eqref{c18v52:operator}. We use the entire lowest positive block
\[
 E_*=\frac32,\qquad
 \operatorname{ran}\Pi=\bigoplus_i\mathcal H_i,\qquad
 \mathcal H_i=\operatorname{span}_{\mathbb C}
                  \{\,2(x_i)_\mu:0\le\mu\le3\,\}.
 \tag{V53.1}\label{c18v53:block}
\]
In particular $\dim_{\mathbb C}\operatorname{ran}\Pi=12n^3$.

\begin{theorem}[Complete initial second coefficient]
\label{c18v53:main}
For every coarse value $g$ and every $n\ge5$, the actual
vacuum-subtracted direct block and effective second coefficient are
\[
 \mathcal A_{\rm conn}=-\frac32\Pi,\qquad
 \mathcal B_*=-\frac{141}{40}\Pi .
 \tag{V53.2}\label{c18v53:answer}
\]
Consequently the full unreduced conditional gap satisfies
\[
 \lambda_{\rm fib}(t,g)
       =\frac32-\frac{141}{40}t^2+O_{n,g}(t^3)
       \qquad(t\longrightarrow0).
 \tag{V53.3}\label{c18v53:gap}
\]
This is a fixed-regulator Taylor statement. The displayed
coefficient is regulator-independent, but a regulator-independent
remainder or interval of validity is not asserted.
\end{theorem}

The proof below pays the previously unknown metric contribution.
The density terms are not omitted: internal gauge averaging shows
that their compression equals exactly the vacuum contribution.
The normal acceleration need not be computed independently,
because total geodesicity and normality determine its only
contribution to the second metric variation.

\subsection{Internal gauge averaging of the scalar terms}

Let $\mathcal G_{\rm mid}$ be the product of the gauge groups at
the midpoints of the coarse two-link chains, with all other
vertex transformations set to the identity. These midpoints are
distinct: their coordinate parity has exactly one odd entry.
For the midpoint belonging to pair $i$, the action is
\[
       x_i\longmapsto x_i k^{-1},\qquad
       B_i\longmapsto kB_i,\qquad k\in\mathrm{SU}(2).
 \tag{V53.4}\label{c18v53:gauge}
\]
It also changes the incident transverse free links, but no
other pair coordinate. Both coarse endpoints are fixed, so the
action preserves each $g$, not merely its conjugacy class.

\begin{lemma}[Scalar compression by actual internal symmetry]
\label{c18v53:scalarlemma}
If a smooth scalar $S$ on $M_g$ is invariant under
$\mathcal G_{\rm mid}$, then
\[
                  \Pi S\Pi=\left(\int S\,d\mu\right)\Pi.
 \tag{V53.5}\label{c18v53:scalar}
\]
The actual pulled metric and density are invariant under this
group, at every chart time. Hence the coefficients $b,c$ in
\eqref{c18v52:seconddata}, $|db|_m^2$, and $L_0c$ are invariant.
\end{lemma}

\begin{proof}
The base action preserves product Haar measure. For $f,h$
depending only on pair $i$, average $\overline f h$ over its
midpoint group. Since $x_i k^{-1}$ is Haar distributed for fixed
$x_i$, the result is the constant $\langle f,h\rangle$.
Invariance of $S$ and $\mu$ proves the same-block assertion.
For $f\in\mathcal H_i$, $h\in\mathcal H_j$, $i\ne j$, average
over the midpoint of $i$. It leaves $h$ unchanged and annihilates
$f$, whose Haar mean is zero. This proves the off-diagonal assertion.
No independence between $S$ and the free links is assumed.

Gauge transformations are ambient isometries preserving the
Wilson action. Its gradient flow is therefore equivariant.
The block map is unchanged under a gauge transformation that
fixes all coarse vertices. Thus $c_t$ is unchanged, its tangent
kernel and orthogonal normal projector transform equivariantly,
and so does $N_t=Q_tF$. Uniqueness for the normal-chart ordinary
differential equation implies that the chart commutes with the
action. Its pulled metric, volume Jacobian and normalized
conditional density are invariant. Taking Taylor coefficients
and applying the invariant base metric and Laplacian proves
the final claims.
\end{proof}

For a real test function, integrate the $c$ term of
\eqref{c18v52:directform} by parts with the positive Laplacian:
\[
 \begin{aligned}
 \langle f,A_2f\rangle
       &=\int a_2(df,df)\,d\mu+\int f^2 S\,d\mu,\\
 S&=-\frac12 L_0c+\frac14|db|_m^2,\\
 \int S\,d\mu
       &=\frac14\langle b,L_0b\rangle
         =\frac{243}{4}n^3=s_{\rm vac}.
 \end{aligned}
 \tag{V53.6}\label{c18v53:density}
\]
Here $b=3W_{\rm cap}$, $L_0W_{\rm cap}=9W_{\rm cap}$ and
$\int W_{\rm cap}^2\,d\mu=3n^3$, all inherited native identities.
The integral of $L_0c$ vanishes on the compact fibre.
Lemma~\ref{c18v53:scalarlemma} and polarization now show
\[
       \langle f,\mathcal A_{\rm conn}h\rangle
              =\int a_2(d\overline f,dh)\,d\mu,
       \qquad f,h\in\operatorname{ran}\Pi .
 \tag{V53.7}\label{c18v53:metriconly}
\]
Thus the unknown density coefficient is eliminated from this
particular compressed difference by a proved symmetry identity,
not by declaring it negligible or constant.

\subsection{The normal second metric variation, with its signs}

In this subsection $N=N_0$ is a normal field along the base fibre,
and $\nabla^\perp$ is the ambient Levi--Civita normal connection.
It is not the coarse-lift connection used in earlier notes.
Use curvature convention
$R(X,Y)Z=\nabla_X\nabla_YZ-\nabla_Y\nabla_XZ-\nabla_{[X,Y]}Z$;
the unit sphere has
$R(X,Y)Z=\langle Y,Z\rangle X-\langle X,Z\rangle Y$.

\begin{lemma}[Normal-chart metric two-jet at a totally geodesic fibre]
\label{c18v53:variationlemma}
For the actual chart, and tangent vectors $v,w$ to $M_g$,
\[
 \begin{aligned}
 \widehat m''_0(v,w)
    ={}&2\langle\nabla^\perp_vN,\nabla^\perp_wN\rangle
          -2\langle R(v,N)N,w\rangle\\
       &-\operatorname{Hess}_{M_g}|N|^2(v,w).
 \end{aligned}
 \tag{V53.8}\label{c18v53:secondmetric}
\]
Only the initial normal field occurs in this formula.
\end{lemma}

\begin{proof}
Let $V_t=\partial_t\eta_t=-N_t\circ\eta_t$ and
$A=\nabla_tV_t|_0$ be the acceleration along the variation.
For a parameter tangent $v_t=d\eta_t(v)$, normality gives
$\langle V_t,v_t\rangle=0$. Its derivative at zero, using
$\nabla_tv_t=\nabla_vV_0$, gives
\[
       \langle A,v\rangle
         =-\langle V_0,\nabla_vV_0\rangle
         =-\frac12v(|N|^2),\qquad
       A^\top=-\frac12\operatorname{grad}_{M_g}|N|^2 .
 \tag{V53.9}\label{c18v53:acceleration}
\]
Differentiate $\langle v_t,w_t\rangle$ twice and commute
covariant derivatives along the variation. At zero this gives
\[
 \begin{aligned}
 \widehat m''_0(v,w)
  ={}&2\langle\nabla_vV_0,\nabla_wV_0\rangle
       +2\langle R(V_0,v)V_0,w\rangle\\
     &+\langle\nabla_v A,w\rangle+\langle\nabla_w A,v\rangle.
 \end{aligned}
 \tag{V53.10}\label{c18v53:unreducedvariation}
\]
The two curvature terms from differentiation agree by the
curvature symmetries. Since the base second fundamental form
is zero, $\nabla_vN$ is purely normal and the normal component
of $A$ contributes nothing to the last line. Substitute
\eqref{c18v53:acceleration} and $V_0=-N$.
The last line is minus the Hessian of $|N|^2$ and
$R(V_0,v)V_0=-R(v,N)N$. This proves the formula.
\end{proof}

Since $\widehat m'_0=0$, inverse differentiation gives
$a_2=-\tfrac12m^{-1}\widehat m''_0m^{-1}$.
Invariance of this tensor and $\mu$ implies that the real
quadratic form in \eqref{c18v53:metriconly} is scalar on every
real $\mathcal H_i$. Indeed right multiplication by
$\mathrm{SU}(2)$ acts transitively on the unit sphere of its
four real coefficient vectors. Between distinct $i,j$, the
central midpoint transformation $k=-1$ changes the sign of one
argument and not the other, so the mixed form vanishes.
Complexification proves the same assertion on the complex block.
This argument uses the real form and transitivity; it does not
incorrectly assume complex irreducibility of the four-dimensional
right-regular coefficient space.

Call its scalar on $\mathcal H_i$ $d_i$. The real orthonormal
basis $f_{i,\mu}=2(x_i)_\mu$ obeys
$\sum_\mu df_{i,\mu}\otimes df_{i,\mu}=4g_{\rm round}=2m_i$.
Taking its average and contracting gives
\[
             d_i=-\frac14\int
                       \operatorname{tr}_i\widehat m''_0\,d\mu .
 \tag{V53.11}\label{c18v53:trace}
\]
Here the trace uses the three unit tangent vectors of pair $i$,
not its four scalar eigenfunctions. The pairwise trace of the
Hessian term in \eqref{c18v53:secondmetric} integrates to zero:
it is the pair Laplacian with negative sign, and pairwise Haar
integration annihilates it. Consequently, with an orthonormal
pair frame $T_{i,a}$,
\[
 \begin{aligned}
 d_i&=-\frac12(D_i-C_i),\\
 D_i&=\int\sum_{a=1}^3
                  |\nabla^\perp_{T_{i,a}}N|^2\,d\mu,\\
 C_i&=\int\sum_{a=1}^3
               \langle R(T_{i,a},N)N,T_{i,a}\rangle\,d\mu .
 \end{aligned}
 \tag{V53.12}\label{c18v53:DC}
\]

\subsection{Native normal frames and the exact local contractions}

Let $e_1,e_2,e_3$ be imaginary unit quaternions, with
$[e_a,e_b]=2\sum_c\epsilon_{abc}e_c$. On pair
$(A,B)=(x,x^{-1}g)$ the tangent and normal orthonormal frames are
\[
       T_a=\frac1{\sqrt2}(xe_a,-e_aB),\qquad
       H_a=\frac1{\sqrt2}(xe_a,+e_aB).
 \tag{V53.13}\label{c18v53:frames}
\]
All normal directions to the whole base fibre are the $H_{i,a}$;
the free factors supply no additional normal directions.
The bi-invariant ambient connection gives
\[
       \nabla^\perp_{T_a}H_b
          =\frac1{\sqrt2}\sum_c\epsilon_{abc}H_c .
 \tag{V53.14}\label{c18v53:connection}
\]
For clarity, $\nabla_{Ae_a}(Ae_b)=A[e_a,e_b]/2$,
whereas $\nabla_{e_aB}(e_bB)=-[e_a,e_b]B/2$.
The opposite sign in the second component of $T_a$ makes the
two components of $\nabla_{T_a}H_b$ have the same sign.
Equation~\eqref{c18v53:connection} follows, and its coefficients
are exactly the tangent connection coefficients on the pair
with metric twice the round metric.

Write the contribution of captured plaquette $p$ to normal
component $j$ as
$N_{j,p}=\sum_b(H_{j,b}w_p)H_{j,b}$.
Uncaptured plaquettes have no normal derivative. A captured
plaquette uses one half of each of two different pairs; hence
\[
        H_{j,b}w_p=\sigma_{j,p}T_{j,b}w_p,\qquad
        \sigma_{j,p}=
           \begin{cases}
              +1,&\text{first half},\\
              -1,&\text{second half}.
           \end{cases}
 \tag{V53.15}\label{c18v53:halfsign}
\]
This equality includes inversely oriented occurrences, because
the two directional derivatives on that same link still agree
or differ by the indicated sign. The normal vector $N_{j,p}$
is therefore the isometric normal copy of
$\sigma_{j,p}\operatorname{grad}_j w_p$.
The connection equality in \eqref{c18v53:connection} shows that
its covariant derivative in pair $i$ is the corresponding copy
of $\sigma_{j,p}\operatorname{Hess}_{i,j}w_p$.
For $i\ne j$ this is the mixed derivative on the product;
normal frames in pair $j$ do not depend on pair $i$.

The ambient curvature is the product curvature of unit spheres.
If $N_i=\sum_b n_bH_{i,b}$, its contribution along $T_{i,a}$
is $(|n|^2-n_a^2)/2$. Summing over $a$ gives
\[
       \sum_a\langle R(T_{i,a},N)N,T_{i,a}\rangle
                          =|N_i|^2 .
 \tag{V53.16}\label{c18v53:curvature}
\]
All other normal components have zero mixed ambient curvature
with $T_i$. The factor $1/2$ before summing comes from the two
ambient sphere components of the normalized frames, not from
changing the original metric convention.

\begin{lemma}[Local Haar contractions]
\label{c18v53:locallemma}
For a captured plaquette $p$ using distinct pairs $i,j$,
\[
 \begin{aligned}
 \int|\operatorname{Hess}_{i}w_p|^2\,d\mu&=\frac3{16},\\
 \int|\operatorname{Hess}_{i,j}w_p|^2\,d\mu&=\frac9{16},\\
 \int|N_{i,p}|^2\,d\mu&=\frac38 .
 \end{aligned}
 \tag{V53.17}\label{c18v53:moments}
\]
Cross terms belonging to distinct captured plaquettes vanish
in both integrals defining $D_i,C_i$.
\end{lemma}

\begin{proof}
For fixed other factors, $w_p$ is a real degree-one spherical
coordinate on pair $i$. On its metric $2g_{\rm round}$,
\[
       \operatorname{Hess}_{i}w_p=-\frac12w_pm_i,\qquad
       L_iw_p=\frac32w_p.
 \tag{V53.18}\label{c18v53:hessian}
\]
There are three tangent directions. Also
$\int w_p^2\,d\mu=1/4$, by integrating either free Haar link.
This proves the first identity. Integration by parts in the
two independent pair factors gives
\[
 \int|\operatorname{Hess}_{i,j}w_p|^2\,d\mu
       =\langle w_p,L_iL_jw_p\rangle
       =\left(\frac32\right)^2\frac14=\frac9{16}.
 \tag{V53.19}\label{c18v53:mixed}
\]
The normal-copy identity and the pair eigenvalue give the
third moment as $(3/2)(1/4)=3/8$.

Two distinct captured plaquettes have distinct pairs of free
links, as proved in the native incidence analysis of V50 and
used in V52. There is a free link present in only one of them.
Taking pair derivatives or the normal-frame connection terms
introduces no new free-link dependence. In a cross inner
product that free factor consequently still has degree one,
and its Haar integral is zero. Frames contract only within a
single normal pair. This proves all the claimed cross-term
cancellations, including those with covariant derivatives.
\end{proof}

Only plaquettes using $i$ can contribute to $D_i$.
Each such plaquette has one normal component in $i$ and one in
its other pair. Equations~\eqref{c18v53:halfsign} and
\eqref{c18v53:moments} give $3/16+9/16=3/4$ for its
contribution to $D_i$, and $3/8$ for its contribution to $C_i$.
There are exactly eight such plaquettes per pair, so
\[
                   D_i=6,\qquad C_i=3,\qquad d_i=-\frac32 .
 \tag{V53.20}\label{c18v53:totals}
\]
These are actual contractions of the native normal field,
not bounds obtained from its maximum norm.

\begin{proof}[Completion of Theorem~\ref{c18v53:main}]
The symmetry reduction, trace formula and
\eqref{c18v53:totals} show
$\mathcal A_{\rm conn}=-(3/2)\Pi$ on the whole degenerate block.
Insert this in the inherited exact virtual-return formula
\eqref{c18v52:effective}:
$-(3/2)-81/40=-141/40$.
The fixed-regulator perturbation conclusion
\eqref{c18v52:secondgap} proves \eqref{c18v53:gap}.
In particular the gap is strictly less than $3/2$ for all
sufficiently small positive $t$, at each fixed regulator.
This does not determine its value or even a quantitative
positive lower bound at the required block time.
\end{proof}

\subsection{The lowest branch is not an internal gauge singlet}

Let $P_{\rm mid}$ denote Haar averaging over
$\mathcal G_{\rm mid}$ in the fixed Hilbert space.
It is an orthogonal projection. Equations
\eqref{c18v53:block} and \eqref{c18v53:gauge} imply
\[
                       P_{\rm mid}\Pi=\Pi P_{\rm mid}=0 .
 \tag{V53.21}\label{c18v53:nonsinglet}
\]
Indeed averaging at the midpoint of the single occupied pair
annihilates every basis function in that block.
The absence of a singlet persists for the entire continued
cluster, not merely its initial eigenvectors.

\begin{corollary}[No physical-gap identification from this branch]
\label{c18v53:singletcor}
For sufficiently small $|t|$ at fixed $n,g$, the spectral
projection $\Pi(t)$ continuing the $3/2$ cluster satisfies
$P_{\rm mid}\Pi(t)=0$.
\end{corollary}

\begin{proof}
The fixed-space operator commutes with the group: the metric
and density are invariant, and multiplication by their density
square root also commutes with it. Its spectral projection
therefore commutes with $P_{\rm mid}$. By fixed-regulator
norm continuity choose $t$ so that $\|\Pi(t)-\Pi\|<1$.
If $P_{\rm mid}\Pi(t)$ were a nonzero projection, its norm
would be one. But
\[
        \|P_{\rm mid}\Pi(t)\|
             =\|P_{\rm mid}(\Pi(t)-\Pi)\|<1 ,
\]
a contradiction.
\end{proof}

Any physical gauge-invariant subspace must, in particular, be
invariant under these midpoint transformations. Thus the
coefficient computed here is not the coefficient of its lowest
positive eigenvalue. A uniform Poincare bound on the full fibre
could still control an invariant subspace by restriction;
an equality of their first positive eigenvalues does not follow.
This distinction now has an explicit spectral proof.

\subsection{Audit and the next upstream obligations}

The new verification script checks native incidence at $n=5,6$
and all twelve captured plane/parity types with noncommuting
coarse values. It differentiates actual quaternion plaquette
words to check the first/second-half signs, the pair Hessian
and the normal-connection correction. Exact symbolic Haar
moments and rational arithmetic check the three contractions
and the coefficient $-141/40$. Two separately labelled
non-native geometric fixtures check the acceleration/Hessian
sign on an initially straight normal variation and the
curvature sign on a unit-sphere latitude. They are sign checks,
not replacements for the native proof. No interacting
conditional-law simulation or long-time spectral certification
is claimed.

A targeted primary-source check considered
\href{https://arxiv.org/abs/2412.00969v2}
{Zawadzki, \emph{Variations of metric that preserve a Riemannian
submersion and geometry of its fibers}}.
Its stated variations change the ambient metric while fixing
the induced fibre metric. That is different from the present
fixed ambient metric and moving native fibre, so no second
variation theorem from that setting is imported. The formula
needed here was derived directly in
Lemma~\ref{c18v53:variationlemma}.
This targeted scope check does not replace the broad literature
checkpoint scheduled at V60.

Removing this append and restoring the title recovers V52
byte for byte. The other 53 sources are hash-protected and
unchanged. The prior active filename is replaced, while the
archives and companion notes are preserved. Only \texttt{.tex}
sources are retained in \path{latex_notes}; compilation and
inspection products remain outside that folder.
The next companion checkpoint is V55 and the next archive and
restart is V100.

\paragraph{What changed mathematically.}
The unknown direct coefficient has now been evaluated rather
than relabelled as a new assumption. Density normalization,
normal connection, ambient curvature and tangential
acceleration have all been accounted for, yielding a scalar,
negative, complete quadratic shift. It is not legitimate to
extrapolate that Taylor polynomial to time four: V52 already
exhibited the volume growth of a global spectral-projector
derivative, and no uniform Taylor remainder has been paid.

\paragraph{Next target and limits.}
There are two genuinely new upstream tasks. First, identify
the bottom of the electric spectrum invariant under the
\emph{full based internal gauge group}, acting at all fine
vertices except the coarse vertices, and its temporal couplings.
Invariance under the midpoint subgroup alone is only a necessary
condition. The $3/2$ block cannot stand in for this spectrum.
Second, derive a connected/local
remainder or a nonperturbative conditional comparison on the
required time interval. Neither task is supplied by another
finite-volume perturbation theorem or by the three constants
in \eqref{c18v53:moments}.
Coarse variance/return, variable-bank stability, weak-coupling
scale transport, physical-time calibration and interacting
continuum reconstruction remain separate unpaid obligations.
The full Yang--Mills mass gap remains unproved.
\addtocontents{toc}{\protect\setcounter{tocdepth}{2}}
\addtocontents{toc}{\protect\clearpage}
% END CYCLE 18 VERSION 53 APPEND
% BEGIN CYCLE 18 VERSION 54 APPEND
\clearpage
\section{V54: The first internal-singlet band and its native linear shift}
\label{c18v54:section}

\subsection{Why the invariant calculation changes the target}

V53 computed the complete initial quadratic shift of the
\emph{unreduced} conditional gap and proved that its $3/2$
cluster contains no midpoint-gauge singlet. V33 had already
shown that the actual invariant source has zero spectral weight
on noninvariant modes. We now enforce the full based internal
gauge group, not just the midpoint subgroup, and determine
its entire first positive electric block and its actual
first-order temporal compression. No new gauge-fixing metric
or replacement conditional density is introduced.

Keep the literal $c_t=b_2\Phi_t$ path, $\beta=0$, $N=2n$,
$n\ge5$, unit-round $\mathrm{SU}(2)$ and every fixed coarse
value $g$. Let
\[
 \mathcal G_{\rm int}
  =\{h:h_v=1\text{ at every coarse vertex}\},\qquad
 \mathcal H_{\rm int}=L^2(M_g,\mu_g)^{\mathcal G_{\rm int}} .
 \tag{V54.1}\label{c18v54:group}
\]
Thus transformations at \emph{all} noncoarse fine vertices
are averaged. Coarse-root transformations are not quotiented.
In V51's equivariant normal chart, $\mathcal H_{\rm int}$ is
a fixed reducing subspace of V52's operator
$\mathcal K_t=L_0+tV+O_{n,g}(t^2)$, where
$V=-(27/2)W_{\rm cap}$. Let $\lambda_{\rm int}(t,g)$ be its
first positive eigenvalue on this subspace.

\begin{theorem}[Complete invariant first band and right derivative]
\label{c18v54:main}
The first positive eigenvalue of $L_0$ on
$\mathcal H_{\rm int}$ is exactly nine. Its eigenspace has
complex dimension $72n^3$. It is the orthogonal sum of
six four-dimensional path channels at each of $3n^3$
fine vertices having exactly two odd coordinates.
Four of those six channels are captured-plaquette channels
and two are straight-path channels, described below.

For every coarse value, the full compression to this eigenspace
has extremal eigenvalues
\[
 \lambda_{\max}(P_9W_{\rm cap}P_9)=\frac1{\sqrt2},\qquad
 \lambda_{\min}(P_9W_{\rm cap}P_9)=-\frac1{\sqrt2}.
 \tag{V54.2}\label{c18v54:extremes}
\]
Consequently
\[
       \lambda_{\rm int}(t,g)
          =9-\frac{27}{2\sqrt2}t+O_{n,g}(t^2),
                       \qquad t\downarrow0 .
 \tag{V54.3}\label{c18v54:gap}
\]
The coefficient is independent of volume and coarse holonomy.
The remainder and the time interval are only asserted at
fixed regulator. This is a conditional invariant spectrum,
not the full reconstructed physical Yang--Mills Hamiltonian.
\end{theorem}

\subsection{A cut graph that retains the original electric metric}

Let $A$ denote the first fine half of every sampled two-link
chain and $B$ its second half. On $M_g$, retain the independent
coordinates $x_e=U_{A_e}$ and all free fine links, and eliminate
$U_{B_e}=x_e^{-1}g_e$ exactly. The remaining independent
variables are edges of the graph obtained by deleting $B$,
with the original fine vertices and the coarse vertices
marked as roots. Call it $\Gamma_{\rm cut}$.

This is only a representation of the existing coordinates,
not a new metric. The electric energy of a product
Peter--Weyl degree assignment is exactly
\[
 E(\boldsymbol\ell)
     =\frac12\sum_{e\in A}\ell_e(\ell_e+2)
         +\sum_{e\in E_{\rm free}}\ell_e(\ell_e+2).
 \tag{V54.4}\label{c18v54:energy}
\]
Each nontrivial prefix factor costs at least $3/2$;
each nontrivial free factor costs at least three.
Because every coarse-root transformation is the identity,
the group action on this cut graph is precisely ordinary
left/right endpoint multiplication. The deterministic
$B$ links still transform correctly, since the coarse
$g_e$ is fixed. In particular the group action on the
independent coordinates contains no hidden $g$ dependence.

For each degree assignment, the tensor product of matrix
coefficient spaces is invariant under the group.
Averaging at each nonroot vertex requires an invariant
tensor among its incident representation indices.
A vertex with exactly one nontrivial incident edge has
no such tensor: a nontrivial irreducible representation
has no fixed vector. At a degree-two vertex with two
fundamental edges the invariant contraction is unique,
up to scale. This follows from the one-dimensional
intertwiner space between two equivalent irreducibles,
with duals supplied by the edge orientations.
No contraction is imposed at a root.

These elementary consequences of Peter--Weyl decomposition
apply to every degree, not to a chosen truncation.
They give the following complete low-energy classification.

\begin{lemma}[All nonconstant invariant modes of energy at most nine]
\label{c18v54:classification}
Each nonzero degree component is fundamental on two prefix
factors and two free factors. Its two free links meet at a vertex
with exactly two odd coordinates and end at two different
midpoints. The prefix factor attached to each midpoint
completes a four-edge path between coarse roots, which
may be the same root. There are no other nonconstant
invariant modes at energy at most nine.
\end{lemma}

\begin{proof}
Write the parity count of a vertex for its number of odd
coordinates. A root has count zero, a midpoint count one.
Each midpoint is incident to exactly one retained prefix
edge; its other captured edge was in $B$. A free link
joins parity counts $1,2$ or $2,3$, and no free link
is incident to a root. These assertions follow directly
from the definition of captured links.

Consider the support of one nonzero degree assignment.
With no occupied free edge, any occupied prefix edge
has a degree-one midpoint, so there is no invariant.
With one occupied free edge, an endpoint of parity count
at least two has no available prefix edge and is a
degree-one nonroot vertex. This is again impossible.
At least two free edges are therefore necessary.

If there are three occupied free edges and energy at most
nine, they are fundamental and there is no occupied
prefix edge. The free graph is simple and bipartite;
three edges cannot contain a cycle. A nonempty forest
has a leaf, here necessarily a nonroot, which is impossible.
Four occupied free edges already cost at least twelve.

Thus there must be exactly two occupied free edges.
Neither can have degree two or higher in Peter--Weyl
label, since $8+3>9$. Their cost is six, leaving at most
three for prefix factors. There can only be at most
two fundamental prefix factors.
The two free edges must share their endpoint of parity
count at least two: otherwise one such endpoint is a
degree-one nonroot with no prefix available.
Their other endpoints must both be midpoints; otherwise
another vertex without a prefix is a leaf.
Their common endpoint then has parity count two.
Both midpoints require their unique prefix edge.
All four edges are fundamental and the total cost is
$3/2+3+3+3/2=9$. This proves necessity.

Conversely, this four-edge path has unique fundamental
contractions at its three nonroot vertices. Its holonomy
matrix coefficients give nonzero invariant functions,
with energy nine by \eqref{c18v54:energy}.
Decomposing an arbitrary eigenfunction into orthogonal
degree assignments proves completeness of the list.
\end{proof}

\begin{corollary}[The next invariant energy and its resolvent buffer]
\label{c18v54:buffer}
The next distinct invariant electric energy after nine is twelve.
For $Q_9=I-P_9$ on $\mathcal H_{\rm int}$,
the reduced inverse $Q_9(L_0-9)^{-1}Q_9$ has norm $1/3$.
\end{corollary}
\begin{proof}
Below energy twelve there are at most three occupied free
edges. With two, the preceding support argument still forces
the same four-edge path. Its degree-two internal contractions
require all four representation labels to agree. Its energies
are $3\ell(\ell+2)$, namely nine and then at least twenty-four.
With three free edges, all are fundamental below twelve.
Their support is a nonempty forest with at least two leaves,
each requiring a different prefix edge. The cost is at least
$9+2(3/2)=12$. Zero or one free edge remains impossible.
An uncaptured fine plaquette has four free links, an invariant
nonconstant character and exact energy twelve, so the bound
is attained. The only lower energy outside $P_9$ is the
constant ground of energy zero. The distances from nine
are therefore nine below and at least three above, with
three attained. Functional calculus proves the norm.
\end{proof}

\subsection{The six channels and their exact Haar normalization}

At a vertex $y$ with exactly two odd coordinates, there
are four adjacent midpoints, in a square plane.
Number them cyclically $0,1,2,3$. There are $3n^3$
such vertices. Let $Z_r$ be the path holonomy from
the root of midpoint $r$ to $y$, first along its
prefix edge and then along its free spoke, using
an inverse for a negatively traversed stored edge.
The four $Z_r$ are independent Haar quaternions:
their four prefixes and four free spokes are distinct
independent factors, and Haar multiplication preserves
the marginal law. Independence here is local to one
centre; adjacent centres can share underlying factors.

Define, for $r<s$,
\[
       X_{rs}=Z_rZ_s^{-1},\qquad
       f_{rs,\mu}=2(X_{rs})_\mu,\quad 0\le\mu\le3 .
 \tag{V54.5}\label{c18v54:paths}
\]
Every $X_{rs}$ is invariant at both midpoints and at
$y$. Its endpoints are roots where no invariance is
imposed. Each $X_{rs}$ is Haar distributed, so its
four displayed real functions are orthonormal.
They also form a complex basis of its four-dimensional
fundamental matrix-coefficient channel.

Distinct pairs $\{r,s\}$ have different pairs of free
links; a unique free factor of degree one makes their
inner product zero. The same argument proves
orthogonality between channels at different centres.
Lemma~\ref{c18v54:classification} therefore gives
\[
  \operatorname{ran}P_9
     =\bigoplus_y\ \bigoplus_{0\le r<s\le3}
          \operatorname{span}_{\mathbb C}
                    \{f_{rs,\mu}:0\le\mu\le3\},\qquad
  \dim\operatorname{ran}P_9=(3n^3)\,6\,4=72n^3 .
 \tag{V54.6}\label{c18v54:dimension}
\]
The four adjacent pairs
$01,12,23,03$ are the four captured plaquettes around
$y$, with any missing second halves restored by fixed
coarse factors. We call these the corner channels.
The opposite pairs $02,13$ give two collinear free
spokes and are the straight channels. They are equally
low in electric energy, even though their full paths
are not elementary fine plaquette loops.

There are thus $12n^3$ corner channels and $6n^3$
straight channels. A captured Wilson character is only
one particular real coefficient vector in its
four-dimensional corner channel; it is not that entire
channel, much less the entire eigenspace.

\subsection{The actual coarse holonomy remains as a local twist}

At one centre the captured action has the form
\[
       W_y=\sum_{rs\in\{01,12,23,03\}}
                        \operatorname{Re}(K_{rs}X_{rs}),
          \qquad K_{rs}\in\mathrm{SU}(2).
 \tag{V54.7}\label{c18v54:Wy}
\]
The $K_{rs}$ are the actual coarse link products from
the eliminated halves. This is not a flat-coarse-field
assumption. To check orientations concretely, let
$b,l,u,r$ be the bottom, left, top and right positive
coarse links of the enclosing square. With midpoint
order south, west, north, east, direct substitution
$B=x^{-1}g$ gives
\[
       K_{01}=1,\quad K_{12}=l^{-1},\quad
       K_{23}=ru^{-1},\quad K_{03}=b^{-1}.
 \tag{V54.8}\label{c18v54:coarsecoeff}
\]
This formula holds in each of the three spatial planes.
Translations and reversal of the square only change
the conventions for its coarse holonomy.

For fixed $g$, replace $Z_r$ by $\widetilde Z_r=a_rZ_r$.
This induces orthogonal changes of the real channel
bases and preserves the local Haar integrals.
The coefficient becomes
$\widetilde K_{rs}=a_sK_{rs}a_r^{-1}$.
Choose $a_0=1$ and successively
$a_{s}=a_rK_{rs}^{-1}$ along $01,12,23$.
Those three coefficients become one and the fourth is
\[
       H=\widetilde K_{03}=lur^{-1}b^{-1}.
 \tag{V54.9}\label{c18v54:twist}
\]
It is the holonomy around the actual enclosing coarse
plaquette, with the indicated orientation.
Thus, in these channel bases,
\[
 W_y=\operatorname{Re}X_{01}
        +\operatorname{Re}X_{12}
        +\operatorname{Re}X_{23}
        +\operatorname{Re}(H X_{03}).
 \tag{V54.10}\label{c18v54:normalizedWy}
\]
Tildes have been dropped after the basis change.
Independent choices of channel bases at different
centres are permitted on the orthogonal sum
\eqref{c18v54:dimension}. We do not assert an independent
global change of all overlapping underlying link variables.

\subsection{A complete local compression rather than a plaquette truncation}

\begin{lemma}[Locality and the exact $24$-dimensional matrix]
\label{c18v54:matrixlemma}
The compression $P_9W_{\rm cap}P_9$ is a direct sum
over centres of real symmetric $24$-dimensional matrices.
In corner order $01,12,23,03$ and straight order $02,13$,
the matrix at coarse twist $H$ is
\[
 M_H=\frac14
       \begin{pmatrix}0_{16}&B_H\\ B_H^{\mathsf T}&0_8\end{pmatrix},
 \qquad
 B_H=\begin{pmatrix}
                I&R_H\\ I&I\\ R_H^{\mathsf T}&I\\ I&I
             \end{pmatrix},\qquad
 R_Hq=\overline H\,\overline q .
 \tag{V54.11}\label{c18v54:matrix}
\]
Each entry of $B_H$ is a real $4$ by $4$ block and
$I=I_4$. Quaternion conjugation is denoted by a bar.
\end{lemma}

\begin{proof}
A basis function in a path channel has degree one
on exactly its two free spokes. In a matrix element
with a captured plaquette, centre parity on each free
link requires every free factor to occur an even
number of times. Free spokes of these paths have a
unique endpoint of parity count two. Hence all three
channels in a nonzero matrix element must have the
same centre.

At that centre, let $\gamma,\delta$ be the two path
pairs and $p$ the plaquette pair. The necessary parity
condition is
$\gamma\mathbin{\triangle}\delta=p$,
where the sets are spoke indices. It forces a
triangle on three different spokes. Since $p$ is an
adjacent pair of the four-cycle, exactly one of
$\gamma,\delta$ is a corner and the other is straight.
This proves the two zero diagonal blocks and leaves
the eight displayed off-diagonal blocks.

All their integrals follow from the exact identity
\[
       \int q_\alpha q_\beta\,dH(q)=\frac{\delta_{\alpha\beta}}4,
       \qquad
       \int q\,\operatorname{Re}(q a^{-1})\,dH(q)=\frac a4 .
 \tag{V54.12}\label{c18v54:Haar}
\]
For example, in the row $01$, column $02$, the
plaquette is $12$. Integrating its shared spoke with
$f_{01}$ replaces $Z_1$ by $Z_2$ with factor $1/4$,
giving the block $I/4$.
In the row $01$, column $13$, the plaquette is $03$.
Integrating $Z_0$ in
$2Z_0Z_1^{-1}\operatorname{Re}(H Z_0Z_3^{-1})$
gives
$\tfrac12\overline H Z_3Z_1^{-1}
 =\tfrac12\overline H\,\overline X_{13}$.
In normalized coordinates this gives $R_H/4$.
Applying the same identity to the remaining six
positions gives, row by row,
$(I,I)$, $(R_H^{\mathsf T},I)$ and $(I,I)$,
in addition to the first row $(I,R_H)$.
The transpose accounts for the reversed path
orientation. Orthogonality of the four channel
coordinates completes \eqref{c18v54:matrix}.
\end{proof}

\begin{lemma}[A sharp norm independent of coarse frustration]
\label{c18v54:normlemma}
For every unit quaternion $H$,
$\lambda_{\max}(M_H)=1/\sqrt2$ and
$\lambda_{\min}(M_H)=-1/\sqrt2$.
\end{lemma}

\begin{proof}
$R_H$ is orthogonal, so direct block multiplication gives
\[
       B_H^{\mathsf T}B_H
       =\begin{pmatrix}
          4I&2(I+R_H)\\
          2(I+R_H^{\mathsf T})&4I
        \end{pmatrix}.
 \tag{V54.13}\label{c18v54:Gram}
\]
Its largest eigenvalue is
$4+2\|I+R_H\|$, as follows by singular-value
decomposition of the off-diagonal block.
The norm is at most eight.
It is exactly eight: choose any unit quaternion
$q$ with $q^2=\overline H$. Such a square root
exists also when $H=-1$. Then
$R_Hq=\overline H\,\overline q=q^2q^{-1}=q$,
and therefore $\|I+R_H\|=2$.
Equivalently $(q,q)$ is an eigenvector of
\eqref{c18v54:Gram} of eigenvalue eight.
The nonzero eigenvalues of the symmetric
off-diagonal matrix in \eqref{c18v54:matrix}
are plus and minus the singular values of $B_H$,
divided by four. This proves both extrema.
No smooth choice of the square root over all $g$
is needed for this spectral statement.
\end{proof}

As a further exact check, for $H=1$ the local spectrum is
\[
 \begin{array}{c|ccccc}
  \text{eigenvalue}&-1/\sqrt2&-1/2&0&1/2&1/\sqrt2\\ \hline
  \text{multiplicity}&1&6&10&6&1 .
 \end{array}
 \tag{V54.14}\label{c18v54:flat}
\]
Indeed $R_1$ is quaternion conjugation: it is $+1$
on the scalar component and $-1$ on the three imaginary
components. In \eqref{c18v54:Gram} this gives one
eigenvalue eight, six eigenvalues four and one zero.
The rectangular off-diagonal construction supplies
the remaining zero multiplicities. The extreme
eigenvalues, but not all these multiplicities,
persist for arbitrary $H$.

\subsection{The actual gap shifts linearly, and what that does not prove}

\begin{proof}[Completion of Theorem~\ref{c18v54:main}]
The support classification and exact channel basis
give the electric gap and multiplicity. The locality
lemma and sharp local norm give
\eqref{c18v54:extremes} for every centre and hence
on their entire orthogonal sum.

The full based gauge group preserves the native
metric, density and normal chart by the same
equivariance argument as in V33 and V53. Its fixed
invariant subspace reduces $\mathcal K_t$.
At fixed $n,g$, the compact elliptic operator has
an isolated finite-dimensional energy-nine cluster.
The fixed-space perturbation argument of V52
applies to this reducing subspace with the original
metric and density. Its first cluster matrix is
$P_9VP_9=-(27/2)P_9W_{\rm cap}P_9$.
The exact zero ground state is invariant and
remains zero. For positive $t$, the lowest member
of the positive cluster therefore has first
coefficient $-27/(2\sqrt2)$, with fixed-regulator
$O(t^2)$ remainder. This proves \eqref{c18v54:gap}.
\end{proof}

The direction of the limit matters. The first
compressed spectrum is symmetric, so on a
two-sided extension the minimum behaves as
$9-(27/(2\sqrt2))|t|+O_{n,g}(t^2)$.
Equation~\eqref{c18v54:gap} asserts its right
derivative along the positive flow, not a
two-sided derivative of the minimum.

There is also a concrete truncation warning.
Let $P_{\rm cor}$ and $P_{\rm str}$ project onto
the corner and straight channels within $P_9$.
Then
\[
 P_{\rm cor}VP_{\rm cor}=0,\qquad
 P_{\rm str}VP_{\rm str}=0,\qquad
 \|P_{\rm cor}VP_{\rm str}\|=\frac{27}{2\sqrt2}.
 \tag{V54.15}\label{c18v54:missed}
\]
In particular a calculation retaining only the
captured plaquette channels would predict zero
first-order shift of its own compressed matrix,
while the true first invariant gap has a
strictly negative first-order shift.
Individual plaquette Rayleigh quotients are
insufficient to see this degeneracy splitting.
The straight channels cannot be integrated out
using a positive electric separation: they have
exactly the same electric energy.

The new invariant origin is larger than V53's
unreduced origin, but it is not stationary under
the flow. Neither fact proves continuation of a
positive uniform gap to block time four. In
particular substituting $t=4$ in the Taylor
polynomial would be outside any proved error
control. The based internal restriction also
does not impose the remaining coarse gauge
constraints, identify physical time or construct
an interacting continuum measure.

\subsection{Verification, nonduplication and the next estimate}

The native diagnostic checks the cut-graph
parities, the single prefix at each midpoint,
the six paths and four/two corner/straight split
at all relevant vertices for $n=5,6$.
It verifies the twelve captured plaquette
orientations in the three planes with genuinely
noncommuting coarse values and checks the exact
coarse-square twist after the tree rotations.
Symbolic quaternion second moments verify every
entry of \eqref{c18v54:matrix} for an arbitrary
quaternion $H$, as well as
\eqref{c18v54:Gram} and the flat multiplicities.
Numerical tests on 24 additional unit twists
supplement the symbolic calculation; the proof
for all twists is Lemma~\ref{c18v54:normlemma}.
There is no interacting-law simulation,
continuum certification or proof-assistant check.

For primary literature context, the abstract of
\href{https://arxiv.org/abs/hep-lat/9909067}
{Burgio--De Pietri--Morales-Tecotl--Urrutia--Vergara,
\emph{Hamiltonian LGT in the complete Fourier
analysis basis}} was checked. Harmonic analysis
and vertex intertwiners are established tools,
not a new method claimed here. The new results
are the complete low-energy classification for
this literal conditioned pair/free metric and
its native coarse-twisted temporal compression.
The support and matrix arguments above supply
the needed proofs; no continuum conclusion is
imported from that abstract.

V53 is preserved byte for byte after restoring
its title and removing this append. The other
53 sources remain hash-protected and unchanged.
Only the latest active YM source is retained
under its current filename; archives remain
untouched and all compilation products stay
outside \path{latex_notes}.
The next companion review is V55, the next
broad literature sweep V60, and archive/restart
remains scheduled at V100.

\paragraph{The next unpaid estimate.}
The first invariant block is now known, including
the channels missed by a plaquette-only
compression and the full coarse holonomy.
The next spectral calculation must control the
return from the \emph{entire} energy-nine
subspace to its complement, together with the
actual direct higher metric/density terms.
The local $24$-dimensional first matrix is not
an invariant subspace of the full finite-time
operator. A connected, regulator-uniform
remainder or nonperturbative conditional
comparison is still required. Source overlap
with the split channels must also be tested
for the actual represented source; the norm
alone does not determine an inverse-source cost.
Coarse return, weak-coupling scale transport,
physical calibration and interacting continuum
reconstruction remain open. The full
Yang--Mills mass gap is still unproved.
\addtocontents{toc}{\protect\setcounter{tocdepth}{2}}
\addtocontents{toc}{\protect\clearpage}
% END CYCLE 18 VERSION 54 APPEND
% BEGIN CYCLE 18 VERSION 55 APPEND
\clearpage
\section{V55: Actual source visibility and native flow reversal}
\label{c18v55:section}

\subsection{Companion checkpoint and the source-resolved target}

The scheduled checkpoint reread the terminal mathematical
ledgers of Source Selection V65, Stationarity V61,
Einstein Bridge V55, restarted Perturbative ESM V14,
SM--Einstein V72, NS V26, Shell Mixing V16,
parallel YM V5 and Arithmetic Loading V17.
All nine active files are unchanged since V50.
The four supplied downloads, ESM V64, Source Selection V52,
Stationarity V54 and Einstein Bridge V45, were also
hash-checked and their terminal ledgers reread; they
remain unchanged.

The useful transfer is source-specific. Source Selection
and Arithmetic Loading require the actual represented
source or projector; a finite spectrum does not supply
its occurrence. Stationarity's principal control leaves
a state-only source unpaid, while Einstein Bridge retains
the moving five-row tangency test. Perturbative ESM keeps
its controlled-complement and coefficientwise scope.
NS's rank-one calculation reinforces testing the actual
input. SM's fixed-mode many-copy limit, Shell Mixing's
measured one-loop calculation and parallel YM's global
charge obstruction are not imported as uniform
conditional or physical mass-gap estimates.

Accordingly, we now test the \emph{actual temporal
normal-chart source}, not an arbitrary vector that
maximizes V54's first matrix. Throughout, retain
$\beta=0$, $N=2n$, $n\ge5$, every coarse value $g$,
and the literal normal chart of V51. Put
\[
 a_t=\partial_t\log\widehat p_t,\qquad
 v_t=\widehat p_t^{1/2}a_t,\qquad
 \psi_t=\widehat p_t^{1/2},\qquad
 a_0=v_0=b=3W_{\rm cap}.
 \tag{V55.1}\label{c18v55:source}
\]
Normalization gives $\langle\psi_t,v_t\rangle=0$.
This is V51's temporal source, including its normal
velocity trace. It is not silently identified with
the centered score of a fixed coarse lift column.
V54's entire energy-nine invariant band is used.

\subsection{An exact three-atom source measure in the first split}

At one two-odd centre use V54's normalized corner/straight
basis, coarse twist $H$, and matrix $M_H$.
Write $c=\operatorname{Re}H$, let $e=1\in\mathbb H$,
and identify real quaternion coordinates with $\mathbb R^4$.
The coefficient vector of the actual local action $W_y$ is
\[
        s_H=\frac12(e,e,e,\overline H;\,0,0),\qquad
        \|s_H\|^2=1 .
 \tag{V55.2}\label{c18v55:s}
\]
The factor $1/2$ comes from the normalized channel
functions $2(X_{rs})_\mu$. The semicolon separates
the four corner from the two straight channels.

\begin{theorem}[The source sees only three first-order spectral values]
\label{c18v55:measurethm}
The spectral measure of $M_H$ in the unit vector $s_H$ is
\[
   \rho_H=\frac{1+c}{4}\delta_{1/\sqrt2}
                +\frac{1+c}{4}\delta_{-1/\sqrt2}
                +\frac{1-c}{2}\delta_0 .
 \tag{V55.3}\label{c18v55:measure}
\]
This holds for every $H\in\mathrm{SU}(2)$, including
$H=-1$, where $M_Hs_H=0$ although
$\|M_H\|=1/\sqrt2$. Other eigenvalues of the complete
$24$-dimensional matrix have zero weight in this source.
\end{theorem}

\begin{proof}
Let $k=e+\overline H$. V54's orthogonal map
$R_Hq=\overline H\,\overline q$ satisfies
$R_Hk=R_H^{\mathsf T}k=k$,
$R_He=\overline H$ and
$R_H^{\mathsf T}e=\overline H$.
For $s_{\rm cor}=(e,e,e,\overline H)/2$,
the explicit $B_H$ of \eqref{c18v54:matrix} gives
\[
 B_H^{\mathsf T}s_{\rm cor}=(k,k),\qquad
 B_H(k,k)=(2k,2k,2k,2k).
 \tag{V55.4}\label{c18v55:Bsource}
\]
Consequently
\[
 \begin{aligned}
 M_Hs_H&=(0;\,k/4,k/4),\\
 M_H^2s_H&=(k/8,k/8,k/8,k/8;\,0,0),\\
 M_H^3s_H&=\tfrac12M_Hs_H,\qquad
 \langle s_H,M_Hs_H\rangle=0,\qquad
 \|M_Hs_H\|^2=\frac{1+c}{4}.
 \end{aligned}
 \tag{V55.5}\label{c18v55:cyclic}
\]
The polynomial $z(z^2-1/2)$ annihilates this cyclic
source vector. For a self-adjoint matrix, its spectral
measure is therefore supported on the three roots.
Their weights sum to one. The vanishing first moment
makes the two nonzero weights equal; the second moment
in \eqref{c18v55:cyclic} determines each as $(1+c)/4$.
The remaining weight is $(1-c)/2$.
Equivalently the projected source vectors are
\[
 \begin{aligned}
 s_+&=M_H^2s_H+\frac1{\sqrt2}M_Hs_H,&
 s_-&=M_H^2s_H-\frac1{\sqrt2}M_Hs_H,\\
 s_0&=s_H-2M_H^2s_H .
 \end{aligned}
 \tag{V55.6}\label{c18v55:projections}
\]
They are pairwise orthogonal eigenvectors with the
claimed squared norms. These polynomial formulas
are asserted on the source cyclic space; $2M_H^2$
is not declared a projection on the whole band.
\end{proof}

Let $\mathcal W_{\rm c}(g)=\sum_y\operatorname{Re}H_y$,
the coarse Wilson sum with one plaquette per two-odd
centre. V54's orthogonal decomposition and the factor
three in \eqref{c18v55:source} give the spectral weights
of $b$ for the \emph{first splitting matrix}
$P_9VP_9=-(27/2)P_9W_{\rm cap}P_9$:
\[
 \begin{array}{c|c}
 \text{first temporal slope}&\text{squared source weight}\\ \hline
 -27/(2\sqrt2)&(9/4)\bigl(3n^3+\mathcal W_{\rm c}(g)\bigr)\\[4pt]
 0&(9/2)\bigl(3n^3-\mathcal W_{\rm c}(g)\bigr)\\[4pt]
 +27/(2\sqrt2)&(9/4)\bigl(3n^3+\mathcal W_{\rm c}(g)\bigr).
 \end{array}
 \tag{V55.7}\label{c18v55:globalweights}
\]
Their sum is the inherited $\|b\|^2=27n^3$.
They are leading split weights, not an assertion
that these three spectral values remain the exact
finite-time spectrum.

\subsection{A native failure of uniform leading-source visibility}

\begin{proposition}[The lowest split band can be invisible to the source]
\label{c18v55:counter}
For every even coarse side $n\ge6$ there is a legitimate
periodic coarse field for which the leading source
weight at slope $-27/(2\sqrt2)$ is zero, although the
source norm is $27n^3$ and that slope is present.
No positive lower bound on its normalized leading
weight can hold uniformly over $n,g$.
\end{proposition}

\begin{proof}
On coarse vertices $z\in(\mathbb Z/n\mathbb Z)^3$ set
\[
           g_{z,i}=(-1)^{\sum_{j>i}z_j}I .
 \tag{V55.8}\label{c18v55:coarseminus}
\]
Even $n$ makes this periodic. For a plaquette in
directions $i<j$, shifting the $i$-link in direction
$j$ changes its sign once; shifting the $j$-link
in direction $i$ does not. Its product is $-I$.
Thus every $H_y=-I$ and
$\mathcal W_{\rm c}(g)=-3n^3$.
Equation~\eqref{c18v55:globalweights} gives zero
weight in both extreme split bands and full weight
in the zero-slope subspace. V54 proves that the
extreme slopes are nevertheless present.
\end{proof}

This is not merely a measure-zero qualification to
a uniform estimate: by continuity, arbitrarily small
leading weights occur in open neighbourhoods of this
field, each of positive coarse Haar measure at fixed
even $n$. It does not disprove an averaged visibility
estimate. At the actual $\beta=0,t=0$ coarse product
Haar reference, each coarse plaquette has zero mean,
so the averaged fractions in the three rows of
\eqref{c18v55:globalweights} are respectively
$1/4,1/2,1/4$. No such average is inferred at
nonzero coupling or from a different coarse law.

There is also a finite-time consequence with the
proper fixed-regulator scope. At the field
\eqref{c18v55:coarseminus}, let
$\Pi_{\rm low}(t)$ project onto the whole subcluster
with first slope $-27/(2\sqrt2)$, for small positive $t$.
Then the \emph{actual} source satisfies
\[
                  \|\Pi_{\rm low}(t)v_t\|^2=O_{n,g}(t^2).
 \tag{V55.9}\label{c18v55:smallweight}
\]
To justify this, first transport the isolated energy-nine
cluster smoothly to its initial space, as in V52.
Subtract nine and divide its matrix by $t$.
The resulting matrix extends smoothly to $P_9VP_9$.
Its lowest first-slope eigenspace is isolated from the
other first slopes at this fixed coarse field.
Its projection therefore equals its limiting projection
plus $O(t)$. Undoing the smooth cluster transport
preserves that estimate. Since $v_t=b+O(t)$ and the
limiting projection annihilates $b$, the norm is $O(t)$
and its square is $O(t^2)$.
This does not assert exact finite-time annihilation,
nor a uniform constant as $n$ or $g$ varies.

\subsection{A block-preserving link-sign involution reverses the flow}

For a stored fine link $e=(x,i)$, define
\[
   \epsilon_{x,i}=(-1)^{\sum_{j>i}x_j},\qquad
   (\sigma U)_{x,i}=\epsilon_{x,i}U_{x,i}.
 \tag{V55.10}\label{c18v55:sigma}
\]
This uses the even fine side $2n$ and is valid for
\emph{every} $n\ge5$, unlike the separate all-negative
\emph{coarse} example, which used even $n$.
It is a Haar-preserving product isometry and an
involution. It is not an internal gauge transformation.

\begin{lemma}[Exact native reversal identities]
\label{c18v55:reversal}
The link map fixes every captured fine link, preserves
$b_2$, and satisfies
\[
 \begin{gathered}
 W(\sigma U)=-W(U),\qquad
 \Phi_t\sigma=\sigma\Phi_{-t},\qquad
 c_t\sigma=c_{-t},\\
 N_t(\sigma U)=-D\sigma\,N_{-t}(U),\qquad
 \eta_t\sigma=\sigma\eta_{-t}.
 \end{gathered}
 \tag{V55.11}\label{c18v55:flow}
\]
Here $N_t=Q_tF$ is the actual V51 normal field,
not a substituted tangent transport.
\end{lemma}

\begin{proof}
On a captured link, every transverse coordinate is
even, so $\epsilon_{x,i}=1$. Thus $b_2\sigma=b_2$.
In an $i,j$ fine plaquette with $i<j$, the same
sign calculation as in \eqref{c18v55:coarseminus}
gives a product of signs equal to minus one.
Every plaquette character, captured or not, changes sign.

An isometry reversing $W$ sends its gradient to its
negative: $F(\sigma U)=-D\sigma F(U)$. Uniqueness
for the autonomous flow gives
$\Phi_t\sigma=\sigma\Phi_{-t}$. Applying $b_2$
gives the third identity.
Differentiation yields
$Dc_t(\sigma U)D\sigma=Dc_{-t}(U)$.
Since $D\sigma$ is orthogonal, the corresponding
orthogonal normal projectors obey
$Q_t(\sigma U)=D\sigma Q_{-t}(U)(D\sigma)^{-1}$.
Combining with the gradient identity proves the
normal-field identity.
Finally differentiating $\sigma\eta_{-t}\sigma$
shows that it solves
$\dot\eta_t=-N_t\circ\eta_t$ with initial identity.
Uniqueness proves the last formula.
All these are finite-dimensional smooth identities
on any common finite time interval. In particular
V51's chart on $[0,4]$ acquires its negative-time
counterpart by this reflection.
\end{proof}

On $M_g$, $\sigma$ preserves the base metric and
product Haar measure. At $\beta=0$ the exact density
is V51's normalized ambient Jacobian, so
\[
        \sigma^*\widehat m_t=\widehat m_{-t},\qquad
        \widehat p_t\circ\sigma=\widehat p_{-t},\qquad
        a_t\circ\sigma=-a_{-t}.
 \tag{V55.12}\label{c18v55:law}
\]
For the density assertion, use
$\eta_t\sigma=\sigma\eta_{-t}$ and
$\operatorname{Jac}\sigma=1$. Normalizing on the
same base fibre preserves the identity, since
$\sigma$ preserves that fibre's Haar measure.
The last identity follows by differentiation.

Let $Tf=f\circ\sigma$ on the fixed Hilbert space.
The full metric and density relations imply
\[
        T\mathcal K_tT=\mathcal K_{-t},\qquad
        T\psi_t=\psi_{-t},\qquad Tv_t=-v_{-t}.
 \tag{V55.13}\label{c18v55:Ksymmetry}
\]
The map also commutes with the based internal gauge
action, since its link factors are central.
Thus these identities restrict to
$\mathcal H_{\rm int}$ without discarding any
invariant source mode.

The old all-negative centre configuration in V50
is therefore being used here as an \emph{operation}
on every field, not only as a point witness.
On V54's first band, $T$ is minus the identity
on all corner channels and plus the identity on
all straight channels. For a corner, its two free
signs multiply to minus one; for a straight, the
two collinear signs are equal and their product
is plus one. This gives the exact grading
$TM_HT=-M_H$ at first order, consistently with
its off-diagonal form.

\subsection{The actual inverse-source cost is even in flow time}

Define the native temporal inverse cost, with its
original moving metric and density, by
\[
 \mathcal C(t,g)
   =\langle a_t,\widehat L_t^{-1}a_t\rangle_{\widehat\nu_t}
   =\langle v_t,\mathcal K_t^+v_t\rangle_{\mu_g}.
 \tag{V55.14}\label{c18v55:cost}
\]
The inverse on the left is centered. On the right,
the pseudoinverse is zero on $\psi_t$.
Gauge invariance makes its full-space and
internally invariant evaluations agree.

\begin{theorem}[A native parity constraint, not a frozen-source surrogate]
\label{c18v55:costthm}
At each fixed regulator and coarse value,
\[
      \mathcal C(t,g)=\mathcal C(-t,g),\qquad
      \mathcal C(t,g)=3n^3+O_{n,g}(t^2).
 \tag{V55.15}\label{c18v55:evencost}
\]
In particular $\partial_t\mathcal C(0,g)=0$,
and every odd Taylor coefficient at zero vanishes.
The coefficient of $t^2$ and a uniform remainder
are not evaluated here.
\end{theorem}

\begin{proof}
Functional calculus and \eqref{c18v55:Ksymmetry}
give $T\mathcal K_t^+T=\mathcal K_{-t}^+$.
The two source signs cancel in the quadratic form,
proving exact evenness.
For smoothness near zero, replace the
pseudoinverse by
$(\mathcal K_t+|\psi_t\rangle\langle\psi_t|)^{-1}$
inside this source quadratic form.
The added ground projection makes the compact
elliptic operator invertible; it changes nothing
because $v_t\perp\psi_t$.
At fixed regulator its inverse and the source
depend smoothly on $t$, so the even function has
the stated Taylor parity. The value at zero is
the inherited V51 cost:
$L_0b=9b$, $\|b\|^2=27n^3$, hence
$\mathcal C(0,g)=3n^3$.
\end{proof}

This result is compatible with V54's linear
decrease of the minimum spectral gap. The minimum
has an absolute-value cusp when both time signs
are considered, while the smooth source quadratic
form is even. It is also stronger than a
calculation with $a_t$ replaced by $a_0$:
the identity retains the actual varying source.
It uses the Haar vacuum at $\beta=0$; no invariance
of a positive-coupling ground density under
$\sigma$ has been assumed.

\subsection{An explicit diagnostic model and the remaining native term}

For clarity about the limit of first-order data,
one may freeze only the first band and its initial
source. This defines, rather than identifies with
\eqref{c18v55:cost}, the model
\[
 \mathcal C_{\rm lin}(t,g)
     =\left\langle b,
       \left(9P_9+tP_9VP_9\right)^{-1}b\right\rangle,
                         \qquad |t|<\frac{2\sqrt2}{3}.
 \tag{V55.16}\label{c18v55:model}
\]
The inverse is taken on $\operatorname{ran}P_9$.
The exact source measure immediately gives
\[
 \begin{aligned}
 \mathcal C_{\rm lin}(t,g)
   &=\sum_y\left[
      \frac{1-c_y}{2}
       +\frac{1+c_y}{2}\frac1{1-9t^2/8}\right],\\
   &=3n^3+\frac9{16}
          \bigl(3n^3+\mathcal W_{\rm c}(g)\bigr)t^2+O_{n,g}(t^4).
 \end{aligned}
 \tag{V55.17}\label{c18v55:modelanswer}
\]
Here $c_y=\operatorname{Re}H_y$.
At the all-negative coarse field this model is
constant even though its full matrix still has
descending extreme eigenvalues.
Most importantly,
$\mathcal C_{\rm lin}$ is \emph{not} asserted to
equal $\mathcal C$, even to second order.
The actual source derivatives, direct
metric/density coefficients and complementary
spectral return are absent from this model.
The rational pole is not a physical singularity
or a proved breakdown time of the native path.

\subsection{Audit and the next noncircular calculation}

The diagnostic reconstructs the complete V54
matrix from exact quaternion Haar moments for
an arbitrary unit $H$, then checks the source
cubic identity, the three orthogonal projected
vectors and their weights. It independently
verifies the frozen model resolvent equation.
Native lattice checks at $n=5,6$ confirm every
plaquette sign, fixed captured links, and all
corner/straight grading signs. An actual
periodic coarse field at $n=6$ verifies the
zero-leading-visibility counterexample.
No normal-flow trajectory or interacting
conditional law is numerically simulated;
flow reversal and cost evenness use the
written smooth identities.

The targeted primary abstract of
\href{https://arxiv.org/abs/hep-lat/0410029}
{Li--Meurice, \emph{Lattice gluodynamics at
negative $g^2$}}, was checked. Relating positive
and negative Wilson couplings by centre-valued
link changes in $\mathrm{SU}(2)$ is established
context. Here no negative-coupling Gibbs law is
substituted for the vacuum: the needed
block-preserving operation and its actual
flow, normal-chart, density and source
consequences at $\beta=0$ were proved explicitly.
This targeted check does not replace V60's
scheduled broad literature sweep.

Removing this append and restoring its title
recovers V54 byte for byte. The other 53
sources remain hash-protected and unchanged;
only the newest active YM filename is retained.
All build and inspection products remain
outside \path{latex_notes}. The completed
companion checkpoint is recorded at V55;
the next companion and broad literature
checkpoints are both V60. Archive/restart
remains V100.

\paragraph{What the new evidence changes.}
A universal first spectral slope does not
imply universal visibility to the native
source. The source has an exact coarse-dependent
three-atom first-split measure and a genuine
leading-visibility counterexample. Yet its
actual inverse cost has no linear term, by a
native symmetry that retains the complete
source and operator.
The next calculation is the actual quadratic
coefficient of \eqref{c18v55:cost}, with both
source derivatives and the full complementary
return retained, followed by a regulator-uniform
local remainder or finite-time comparison.
Neither the model coefficient in
\eqref{c18v55:modelanswer} nor another gap-only
bound supplies it.
Coarse-source scores, coarse variance/return,
weak-coupling scale transport, physical-time
calibration and interacting continuum
reconstruction remain separate obligations.
The full Yang--Mills mass gap is still unproved.
\addtocontents{toc}{\protect\setcounter{tocdepth}{2}}
\addtocontents{toc}{\protect\clearpage}
% END CYCLE 18 VERSION 55 APPEND
% BEGIN CYCLE 18 VERSION 56 APPEND
\clearpage
\section{V56: Connected temporal cost, native density jets and a local cumulant}
\label{c18v56:section}

\subsection{The actual cost, not the frozen first-band model}

V55 determined the leading spectral weights of the actual
temporal source and proved that its inverse cost is even.
It did not determine that cost's quadratic coefficient.
This continuation derives a connected formula for that
coefficient, proves a bound proportional to the number
of coarse sites, and evaluates its fourth-cumulant term
exactly. The remaining local contractions are specified
but not assigned a numerical value or sign.

Keep the literal normal chart, $\beta=0$, fine side
$2n$, $n\ge5$, and any fixed coarse field $g$.
All expectations below are under the original
pair/free product Haar probability $\mu=\mu_g$ on $M_g$.
Write $L=L_0$, $\Gamma(f,h)=m_0^{-1}(df,dh)$,
$\Gamma(f)=\Gamma(f,f)$, and $R=L^+$, zero on constants.
Thus
\[
 b=3W_{\rm cap},\qquad Lb=9b,\qquad
 \mathbb E b=0,\qquad B_2:=\mathbb E b^2=27n^3 .
 \tag{V56.1}\label{c18v56:base}
\]
The full based-internal invariant electric gap is nine
by V54. The full unreduced product gap is $3/2$.
These are conditional electric operators, not an
identification with a physical continuum Hamiltonian.

To keep normalization visible, use the \emph{raw}
normal-chart log Jacobian
\[
 u_t=\log\operatorname{Jac}\eta_t,\qquad
 j=\partial_t^2u_t|_0,\qquad k=\partial_t^3u_t|_0,
 \qquad
 \widehat m_t^{-1}=m_0^{-1}+t^2A+O_{n,g}(t^3).
 \tag{V56.2}\label{c18v56:jets}
\]
Here $u'_0=b$ pointwise. The tensor $A$ is precisely
V53's actual inverse-metric coefficient $a_2$, not an
adjustable counterterm. In particular neither $j$ nor
$k$ is a derivative of an assumed exponential tilt.

\subsection{Raw jets obtained from the actual normal projector}

The following finite-jet recipe fixes all native inputs
in \eqref{c18v56:jets}. Use the original invariant fine
and coarse frames, in which adjoints use fixed Euclidean
inner products. Set
\[
 C_r=\partial_t^r(Dc_t)|_0,\quad
 G_r=\partial_t^r(C_tC_t^*)|_0,\quad
 S_r=\partial_t^r(G_t^{-1})|_0,\qquad r=0,1,2.
\]
There is no inverse of a large variable matrix to assume
at this origin: $G_0=2I$. Product and inverse differentiation give
\[
\begin{aligned}
 G_1&=C_1C_0^*+C_0C_1^*,\\
 G_2&=C_2C_0^*+2C_1C_1^*+C_0C_2^*,\\
 S_0&=\tfrac12 I,\qquad S_1=-\tfrac14G_1,\qquad
 S_2=\tfrac14(G_1^2-G_2),\\
 Q_1&=C_1^*S_0C_0+C_0^*S_0C_1+C_0^*S_1C_0,\\
 Q_2&=C_2^*S_0C_0+C_0^*S_0C_2+2C_1^*S_0C_1\\
    &\hspace{5mm}+2C_1^*S_1C_0+2C_0^*S_1C_1+C_0^*S_2C_0 .
\end{aligned}
 \tag{V56.3}\label{c18v56:projector}
\]
Here $Q_r=\partial_t^rQ_t|_0$, so the subscript two
denotes a derivative, not the coefficient of $t^2$.
The $C_r$ are derivatives of $c_t=b_2\Phi_t$ itself.
They may be obtained by differentiating the Wilson
gradient ordinary differential equation twice and
then the block map; no replacement coarse motion is used.

Put $N_r=Q_rF$ and $d_r=\operatorname{div}N_r$, with
$N=N_0$ and ambient Haar divergence. In expressions
such as $Nd_1$, a vector field acts as a derivation.
Then the exact raw density derivatives are
\[
 \begin{aligned}
 j&=-d_1+Nd_0,\\
 k&=-d_2+2Nd_1+N_1d_0-N(Nd_0).
 \end{aligned}
 \tag{V56.4}\label{c18v56:raw}
\]
Indeed V51 gives $u'_t=-d_t\circ\eta_t$.
Differentiation along the chart applies the material
operator $\partial_t-N_t$ to the scalar before pullback.
Applying it once and twice proves \eqref{c18v56:raw},
including the $N_1d_0$ term and the last minus sign.
In the invariant ambient frame,
$\operatorname{div}N_r=\sum_iX_i(N_r^i)$:
the Haar divergence of each frame vector vanishes.

V55's link-sign operation also applies before density
normalization: $u_t\circ\sigma=u_{-t}$.
Consequently $b,k$ are odd under $\sigma$, $j$ is even,
and
\[
 \mathbb E k=\mathbb E(bj)=\mathbb E b^3=0.
 \tag{V56.5}\label{c18v56:parity}
\]
This parity is not a claim that $j$ vanishes.

\subsection{The first corrected Poisson source is local after centering}

Define $Z_2=\mathbb E(j+b^2)$.

\begin{lemma}[The actual coarse marginal has zero second derivative]
\label{c18v56:normalizer}
For every $g$,
\[
       Z_2=0,\qquad \mathbb E j=-27n^3 .
 \tag{V56.5a}\label{c18v56:Zzero}
\]
This is a second-jet assertion, not a statement that
the coarse law remains Haar at nonzero time.
\end{lemma}

\begin{proof}
Let $\rho_t$ be the density of $(\Phi_t)_*dU$
relative to ambient product Haar. The original
fine Wilson sum, including uncaptured plaquettes,
satisfies $\operatorname{div}F=-12W$. The continuity
equation therefore gives
\[
 \rho'_0=12W,\qquad
 \rho''_0=144W^2-12|\nabla W|_{\mathfrak g}^2 .
\]
The density of $(c_t)_*dU$ relative to coarse
product Haar is both $\mathbb E_\mu\rho_t$ and
$\mathbb E_\mu e^{u_t}$. The second equality follows
from $c_t\eta_t=b_2$ and ambient change of variables.

There are $24n^3$ fine plaquettes. Conditional on
$b_2=g$, each character has second moment $1/4$.
For distinct plaquettes, a free factor occurring
once annihilates their cross moment, as in V50.
The same argument annihilates the cross terms
in the ambient gradient square: a nonzero gradient
contraction differentiates a common fine link,
and a free factor belonging to only one plaquette
remains odd under its central sign. For a single
plaquette,
$|\nabla w_p|_{\mathfrak g}^2=4(1-w_p^2)$,
because each of its four link variables sees a
unit degree-one spherical coordinate.
Consequently
\[
 \mathbb E_\mu W^2=6n^3,\qquad
 \mathbb E_\mu|\nabla W|_{\mathfrak g}^2=72n^3.
\]
Thus $\mathbb E_\mu\rho''_0=144(6n^3)-12(72n^3)=0$.
On the other hand differentiating $e^{u_t}$ twice
at zero gives $j+b^2$. This proves the first identity;
\eqref{c18v56:base} proves the second.
\end{proof}

Expansion of $e^{u_t}/\mathbb E e^{u_t}$ gives
\[
\begin{aligned}
 \widehat p_t&=1+tb+t^2r+t^3s+O_{n,g}(t^4),\\
 r&=\tfrac12(j+b^2-Z_2),\\
 s&=\tfrac16(k+3bj+b^3-3Z_2b).
\end{aligned}
 \tag{V56.6}\label{c18v56:density}
\]
The missing scalar cubic normalizer is zero by
\eqref{c18v56:parity}; $\mathbb E r=\mathbb E s=0$.
Although $Z_2=0$ for this native path, it is retained
in the algebra below to display the normalization
cancellations explicitly.

Let $q_t$ solve the actual weak Poisson equation, but
choose its harmless additive constant by
$\mathbb E_\mu q_t=0$, instead of V51's moving-law mean.
For real tests $f$,
\[
 \mathcal E_t(q_t,f)=\ell_t(f),\quad
 \mathcal E_t(f,h)=
  \mathbb E\bigl[\widehat p_t\widehat m_t^{-1}(df,dh)\bigr],
 \quad
 \ell_t(f)=\mathbb E[f\,\partial_t\widehat p_t].
 \tag{V56.7}\label{c18v56:weak}
\]
This is the original equation
$\widehat L_tq_t=\partial_t\log\widehat p_t$.
Its solution differs from $\chi_t$ only by a constant,
and $\mathcal C(t,g)=\ell_t(q_t)=\mathcal E_t(q_t,q_t)$.

\begin{proposition}[Actual first corrected source]
\label{c18v56:corrector}
The fixed-Haar-centered solution satisfies
\[
 q_0=b/9,\qquad q'_0=Rh,\qquad
 h=\operatorname{center}_{\mu}
                \left(j+\frac{\Gamma(b)}9\right).
 \tag{V56.8}\label{c18v56:h}
\]
For the original moving-law-centered solution,
$\chi'_0=Rh-3n^3$. That constant has no gradient or
Hessian and does not alter the inverse cost.
\end{proposition}

\begin{proof}
The coefficient of $t$ in the weak form is
$\mathcal E_0(q'_0,f)=2\mathbb E(rf)-\mathcal E_1(q_0,f)$,
where $\mathcal E_1(f,h)=\mathbb E[b\Gamma(f,h)]$.
Its strong right hand side is
\[
 2r-\bigl[-\operatorname{div}_{\mu}(b\nabla q_0)\bigr]
 =2r-b^2+\Gamma(b)/9=j+\Gamma(b)/9-Z_2.
\]
Since $\mathbb E\Gamma(b)=\mathbb E(bLb)=9B_2$,
its last constant is exactly the mean of the preceding
two terms. This proves \eqref{c18v56:h}.
Finally $\chi_t=q_t-\mathbb E_{\widehat\nu_t}q_t$,
so $\chi'_0=q'_0-\mathbb E(bq_0)=Rh-B_2/9$.
\end{proof}

In particular the extensive product $b^2$ in the
normalized density cancels \emph{before} inversion.
A bound on that product separately would create an
artificial disconnected volume factor.
For the native path, \eqref{c18v56:Zzero} also gives
$h=j+\Gamma(b)/9$ without any further global constant.

\subsection{A connected formula for the complete quadratic cost}

For bounded random variables let $\kappa$ denote their
joint cumulant under $\mu$, defined by the finite
moment--partition formula
\[
 \kappa(X_1,\ldots,X_q)=
 \sum_{\pi} (|\pi|-1)!(-1)^{|\pi|-1}
                 \prod_{B\in\pi}\mathbb E\prod_{i\in B}X_i .
 \tag{V56.9}\label{c18v56:cumulant}
\]
In particular $\kappa(b,b,j)=\mathbb E(b^2j)-B_2\mathbb E j$,
$\kappa(b,b,Lj)=\mathbb E(b^2Lj)$ and
$\kappa_4(b)=\mathbb E b^4-3B_2^2$.

\begin{theorem}[The native connected quadratic coefficient]
\label{c18v56:costthm}
At each fixed regulator and coarse value,
\[
 \mathcal C(t,g)=3n^3+t^2\mathcal C_2(g)+O_{n,g}(t^4),
 \tag{V56.10}\label{c18v56:expansion}
\]
where all terms below refer to the actual normal chart:
\[
\begin{aligned}
 \mathcal C_2(g)={}&\mathbb E[hRh]
       +\frac19\mathbb E[bk]
       +\frac5{18}\kappa(b,b,j)
       +\frac5{54}\kappa_4(b)\\
       &+\frac1{324}\kappa(b,b,Lj)
       -\frac1{81}\mathbb E[A(db,db)] .
\end{aligned}
 \tag{V56.11}\label{c18v56:connected}
\]
The first term uses the entire centered electric
inverse, not its compression to the nine-band.
\end{theorem}

\begin{proof}
Write $q_t=q_0+tq_1+t^2q_2+O(t^3)$.
The variational value
$2\ell_t(q)-\mathcal E_t(q,q)$ is maximized at $q_t$,
on smooth functions modulo constants.
Alternatively substitute the weak equations at orders
zero, one and two. In either way all occurrences of
$q_2$ cancel, yielding
\[
 \mathcal C_2
  =\mathbb E[hRh]+6\mathbb E[sq_0]
       -\mathbb E[r\Gamma(q_0)]
       -\mathbb E[A(dq_0,dq_0)] .
 \tag{V56.12}\label{c18v56:variational}
\]
Thus the moving source and the moving metric have
both been retained. Substituting \eqref{c18v56:density}
and $q_0=b/9$ into the last three terms gives
\[
\begin{aligned}
 &\frac19\mathbb E[bk]+\frac13\mathbb E[b^2j]
       +\frac19\mathbb E b^4-\frac{Z_2}{3}B_2\\
 &\hspace{5mm}
       -\frac1{162}\mathbb E[(j+b^2-Z_2)\Gamma(b)]
       -\frac1{81}\mathbb E[A(db,db)].
\end{aligned}
\]
The positive-Laplacian product rule gives
$L(b^2)=18b^2-2\Gamma(b)$. Integration by parts also gives
\[
 \mathbb E[b^2\Gamma(b)]=3\mathbb E b^4,\qquad
 \mathbb E[j\Gamma(b)]
       =9\mathbb E[jb^2]-\tfrac12\mathbb E[b^2Lj].
 \tag{V56.13}\label{c18v56:ibp}
\]
Use these identities, $\mathbb E\Gamma(b)=9B_2$,
and $Z_2=\mathbb E j+B_2$. The disconnected products
cancel into the cumulants in \eqref{c18v56:connected}.
Smoothness at each finite regulator follows as in V55.
Its exact evenness removes the cubic coefficient and
gives the fixed-regulator fourth-order remainder in
\eqref{c18v56:expansion}.
\end{proof}

Equation~\eqref{c18v56:connected} is an exact evaluation
recipe and a cancellation identity, not a numerical
evaluation of all six contractions. In particular it
is not V55's frozen-model coefficient.

\subsection{Independent scalar action cells and an exact fourth cumulant}

There is a special simplification for $W_{\rm cap}$,
although not for $j,k$ or the full operator.

\begin{lemma}[A global Haar statement with a metric caveat]
\label{c18v56:independence}
For fixed $g$ under $\mu_g$, the scalar variables
$W_y$ of \eqref{c18v54:normalizedWy}, indexed by the
$3n^3$ two-odd centers, are independent. Their laws
depend on their own coarse holonomies $H_y$.
This assertion does not make the electric operator
a product of operators on these scalar variables.
\end{lemma}

\begin{proof}
Condition first on all pair prefixes $x_i$ and on
the unused free links. A free spoke between a one-odd
midpoint and a two-odd center has exactly one such
center. Hence the four free spokes used at different
centers form disjoint families of independent Haar
variables. Orient a spoke from its midpoint to its
center, call its value $U_{y,r}$, and set
$Z_{y,r}=x_{i(y,r)}U_{y,r}$, as in V54.
Conditional on every prefix these variables are
independent Haar, with a joint law independent of
those prefixes. V54's additional rotations by $a_r(g)$
are deterministic Haar-preserving multiplications.
The $W_y$ are now functions of disjoint four-variable
families, with the only coefficients $H_y(g)$.
Their conditional joint law therefore factorizes
and is independent of the conditioning variables.
Integrating those variables proves the assertion.

Equivalently, the change from free spokes to $Z$,
while retaining the prefixes, is a triangular
Haar-preserving change of variables. It is not a
product isometry: a derivative in a prefix changes
several $Z$ variables. This is why the scalar
factorization neither contradicts V54's metric
warning nor supplies a new spectral product rule.
\end{proof}

\begin{proposition}[An evaluated native fourth-cumulant contribution]
\label{c18v56:fourththm}
With $c_y=\operatorname{Re}H_y$,
\[
 \mathbb E W_y^2=1,\qquad
 \mathbb E W_y^4=\frac{11}{4}+\frac38c_y,\qquad
 \kappa_4(W_y)=-\frac14+\frac38c_y.
 \tag{V56.14}\label{c18v56:fourthlocal}
\]
Consequently, for the actual $b=3W_{\rm cap}$,
\[
 \begin{aligned}
 \kappa_4(b)&=-\frac{243}{4}n^3+
                     \frac{243}{8}\mathcal W_{\rm c}(g),\\
 \frac5{54}\kappa_4(b)
            &=-\frac{45}{8}n^3+
                     \frac{45}{16}\mathcal W_{\rm c}(g).
 \end{aligned}
 \tag{V56.15}\label{c18v56:fourthglobal}
\]
This evaluates one term of the actual coefficient,
not the whole coefficient.
\end{proposition}

\begin{proof}
Represent the four independent Haar quaternions by
unit vectors $z_0,\ldots,z_3$ in $\mathbb R^4$.
The first three corner characters are their adjacent
dot products, and the fourth is
$z_0^{\mathsf T}O_Hz_3$, where $O_H$ is an orthogonal
matrix with $\operatorname{tr}O_H=4\operatorname{Re}H$.
Changing $z_r$ to $-z_r$ shows that a moment vanishes
unless the edge multiplicities give even degree at
every vertex of the four-cycle.

At order two only a doubled single edge survives,
with moment $1/4$. At order four the surviving patterns
are: four copies of one edge; two copies of each of
two edges; and one copy of every edge.
Their moments are respectively $1/8$, $1/16$,
and
\[
 \frac{\operatorname{tr}O_H}{4^4}
                      =\frac{\operatorname{Re}H}{64}.
\]
For the middle case, integrate the outer endpoints
if the two edges meet; each integration gives $1/4$.
Disjoint edges factor directly. The last expression
contracts one second moment at each of the four vertices.
There are four single-edge choices, six unordered
two-edge choices with multinomial coefficient six,
and the cycle has coefficient $24$. Hence
\[
 \mathbb E W_y^4
   =4(1/8)+6\cdot6(1/16)+24(c_y/64)
   =11/4+3c_y/8.
\]
The second-moment calculation gives one. Subtract
three times its square and use the preceding
independence lemma. Fourth cumulants add for
independent sums, and scaling by three multiplies
a fourth cumulant by $81$. This proves both lines
of \eqref{c18v56:fourthglobal}.
\end{proof}

In particular this term can have either sign.
At $H_y=I$ for every $y$ it equals $45n^3/16$;
at the all-negative coarse configuration of V55 it
equals $-225n^3/16$. Other terms in
\eqref{c18v56:connected} remain and can change the
sign of the complete answer.

\subsection{Uniform size of the actual coefficient from finite locality}

Here \emph{local} refers to the original independent
pair/free factors, not to the $Z$ variables in the
last subsection. Form their graph by joining two
distinct factors when they occur in one fine plaquette.
There are $m=21n^3$ vertices. A pair occurs in eight
plaquettes, a free link in four, with at most three
other factors per plaquette. Thus the degree is at
most $24$, uniformly also at periodic boundaries.

\begin{lemma}[Finite native jets have bounded connected supports]
\label{c18v56:locality}
There are decompositions, indexed by the $m$ factors,
\[
 b=\sum_x b_x,\quad j=\sum_x j_x,\quad k=\sum_x k_x,
 \quad h=\sum_x h_x,\quad
 A(db,db)=\sum_x a_x,\qquad \mathbb E h_x=0,
 \tag{V56.16}\label{c18v56:localdecomp}
\]
in which every summand is supported within graph
distance three of $x$. The same holds for $Lj_x$.
There is a finite number $B\ge1$, independent of $n,g$,
bounding the suprema of
$b_x,j_x,k_x,h_x,Lj_x,a_x$ in these decompositions.
The construction uses only fixed-order derivatives,
finite incidence counts, and $G_0^{-1}=I/2$.
It assumes no nonzero-time conditional estimate.
\end{lemma}

\begin{proof}
Assign each captured plaquette to one of its factors,
for example the first in a fixed ordering. This gives
$b_x$ supported at distance one, with uniformly
bounded multiplicity. The gradient component $F_i$
depends only on plaquettes incident to its fine link.
After pair identification these factors lie at
distance one of its factor.

For precision about transposes and matrix products,
retain individual differentiated plaquette monomials.
At order zero, $C_0$ and $Q_0$ couple only the two
halves of a single pair, and their coefficients depend
only on that pair. A time differentiation of the
Wilson-flow jet attaches one incident differentiated
plaquette at a factor already present; a derivative
whose factor is absent is zero. Hence a nonzero term
of $C_r$ contains at most $r$ connected plaquette
attachments, with both its row pair and its column
factor incident to that connected support. The
invariant-frame coefficients at a pair use only that
pair and do not enlarge support. This proves the
assertion for $r=1,2$ directly by differentiating once
and twice. A transpose exchanges the two external
anchors without disconnecting that support.
A contracted matrix index in \eqref{c18v56:projector}
joins supports at the same factor. Thus each term
of $Q_r$ still has at most $r$ attachments and is
within distance $r$ of either external factor.
Multiplying by $F$ gives $N_r$ at distance $r+1$
of its output factor, for $r=0,1,2$.

Ambient differentiation does not enlarge a
coefficient's dependence. Group the three color
components and, for a pair, its two fine-link halves.
The divergence then gives a local decomposition
of $d_r$ of radius $r+1$. In $Nd_0$ a nonzero
differentiated factor lies in the support of the
particular $d_0$ term; the coefficient of $N$ adds
at most one adjacent layer. Similarly the four terms
in $k$ in \eqref{c18v56:raw} add at most three layers
in total. This constructs $j_x,k_x$ of radii two
and three. The product base metric makes $\Gamma(b)$
a sum of squared derivatives at each factor, each
supported at distance one. Center each summand of
$j+\Gamma(b)/9$ separately to obtain $h_x$.
Applying the product Laplacian to $j_x$ preserves
its factor support.

For the metric term use V53's exact second variation
\eqref{c18v53:secondmetric}, not a scalar comparison.
The field $N_i$ has radius one. A derivative of it
in a tangent factor outside that neighborhood is zero.
The two differentiated $N$ factors in the normal
gradient square share their normal-component anchor;
the curvature term is confined to one ambient pair;
and each term of the Hessian of $\sum_i|N_i|^2$ has
the same radius-one anchor. Contracting the two tangent
slots with $db$ adds at most one further layer,
because each gradient component of $b$ has radius one.
This gives $a_x$ of radius at most two.

Finally every displayed coefficient is a finite sum
of fixed-order derivatives of unit-quaternion
products, fixed pair-frame matrices and their
contractions. The only inverse entering these jets
has been replaced by the finite expressions
\eqref{c18v56:projector}. The number of factors,
terms and derivatives per anchor is bounded using
the degree $24$, independently of total volume.
Derivatives of these fixed-order group expressions
are uniformly bounded for all unit fine and coarse
quaternions. Taking a maximum of these finite bounds,
and doubling where a mean is subtracted, supplies
$B$. Equivalently it is a maximum over the finitely
many bounded-size local incidence patterns; periodic
identifications only identify entries in such a
pattern and do not introduce an unbounded family.
No spectral inverse is used in constructing $B$.
\end{proof}

The constants in this lemma are not asserted small.
An explicit numerical optimization of $B$ is not
needed for the following all-volume conclusion.
Set
\[
 D=\sum_{r=0}^{6}24^r,\qquad
 C_*=
 21\left(\frac{DB^2}{3}
         +\frac{91}{18}D^2B^3+\frac{B}{81}\right)
                  +\frac{225}{16}.
 \tag{V56.17}\label{c18v56:boundconstant}
\]

\begin{theorem}[No disconnected volume square in the actual quadratic cost]
\label{c18v56:uniform}
For every $n\ge5$ and every $g$,
\[
 \mathbb E[hRh]\le \frac{21DB^2}{9}n^3,\qquad
                 |\mathcal C_2(g)|\le C_*n^3.
 \tag{V56.18}\label{c18v56:uniformbound}
\]
The second assertion bounds the \emph{complete}
coefficient in \eqref{c18v56:connected}, not just
the frozen first-band term. It does not bound the
$O_{n,g}(t^4)$ remainder uniformly.
\end{theorem}

\begin{proof}
Two local supports can meet only if their anchors
are within distance six. Each anchor has at most
$D$ such partners. Disjoint factor families are
independent under the original product Haar measure.
Thus $\mathbb E(h_xh_y)=0$ for disjoint supports,
and
$\|h\|_2^2\le mDB^2$. The full $h$ is internally
gauge invariant: $j,b$ and the base metric are, by
equivariance of the actual normal chart. V54's
invariant electric gap therefore gives
$\mathbb E[hRh]\le\|h\|_2^2/9$.

The inverse term itself can also be written as a
sum over intersecting supports. The product heat
semigroup preserves cylinder functions, so
$Rh_y=\int_0^\infty e^{-sL}h_y\,ds$ depends only
on the factors of $h_y$ and has mean zero.
The integral converges by the full product gap.
Hence $\mathbb E(h_xRh_y)=0$ for disjoint supports.
This support statement does not assert that the
coefficients of the finite-time inverse are local.

For the remaining terms use multilinearity of
cumulants. If their factor-intersection graph is
disconnected, the variables split into independent
families, and the joint cumulant vanishes.
One proof is to differentiate
$\log\mathbb E\exp(\sum_i t_iX_i)$ at zero:
independence splits that logarithm into two sums
with no joint derivative. Boundedness justifies
these finite derivatives and reproduces
\eqref{c18v56:cumulant}.

For two variables bounded by $B$, the
moment--partition bound is $2B^2$; for three it
is $6B^3$. There are at most $mD$ connected
ordered pairs and $3mD^2$ connected ordered
triples. For the latter, choose a spanning tree
on three labeled positions, its root anchor,
and at most $D$ choices for each succeeding
anchor; repeated anchors are included.
Consequently
\[
 |\mathbb E(bk)|\le2mDB^2,\quad
 |\kappa(b,b,j)|,\ |\kappa(b,b,Lj)|
                         \le18mD^2B^3 .
\]
Here $\mathbb E(bk)=\kappa(b,k)$ because $\mathbb E b=0$.
Also $|\mathbb E A(db,db)|\le mB$.
The exact expression \eqref{c18v56:fourthglobal}
bounds the fourth-cumulant contribution by
$225n^3/16$, since $|\mathcal W_{\rm c}|\le3n^3$.
Insert these bounds into \eqref{c18v56:connected}.
The two $DB^2$ contributions sum to $DB^2/3$,
and the triple contributions to
$(5+1/18)D^2B^3=(91/18)D^2B^3$.
Since $m=21n^3$, this is exactly
\eqref{c18v56:boundconstant}.
\end{proof}

All local jets here are finite polynomials in the
original group-coordinate functions. Thus each centered
$h_x$ belongs to a finite Peter--Weyl sum on a bounded
set of factors. The action of $R$ on it can be computed
by decomposing that finite sum and dividing each
nonconstant mode by its exact sum of pair/free
Casimir energies. This is an exact finite calculation,
not a cutoff approximation to a long-time interacting
inverse. It specifies a way to evaluate the remaining
local contractions without ever forming a matrix
whose dimension grows with the total lattice volume.

\subsection{Audit, new obligation and scope}

The diagnostic independently checks the weak-form
quadratic coefficient against an exactly solved
Fourier Poisson two-jet on a circle with nonzero
raw $j,k$ and inverse-metric variation. This is a
clearly non-native algebra fixture, not a YM
conditional trajectory. A separate one-dimensional
Jacobian fixture checks every material-derivative
sign in \eqref{c18v56:raw}; an exact rectangular
matrix checks the Gram inverse and projector
two-jets. Native incidence checks at $n=5,6$
verify disjoint spoke ownership and the factor
graph degree (the observed degree is $20$, within
the proved bound $24$). The actual noncommuting
coarse-twist check from V54 is replayed.
Native free-factor parity and exact single-plaquette
moments check the two conditional expectations used
in \eqref{c18v56:Zzero}.
Exact sphere moments, cycle contractions and
edge-multiplicity enumeration check
\eqref{c18v56:fourthlocal}.
The all-volume bound is the written locality
and cumulant proof, not an inference from these
two finite lattices.

The primary abstract of
\href{https://arxiv.org/abs/math/0612562}
{Lott, \emph{Some geometric calculations on
Wasserstein space}}, was checked for context.
It concerns the Wasserstein geometry of a smooth
compact manifold. We do not import a fixed-metric
curvature formula for our moving fibre metric:
the metric term in \eqref{c18v56:connected} was
retained and derived directly from the actual
weak Poisson problem. This is a targeted check;
the next broad literature sweep remains V60.
The scheduled companion review was completed
at V55 and is next due at V60.

Removing this append and restoring the title
recovers V55 byte for byte. The other 53 sources
are hash-protected and unchanged. Only the new
active YM filename is kept; all build and
inspection products remain outside
\path{latex_notes}. Archive/restart remains V100.

\clearpage
\paragraph{What has advanced.}
The actual first Poisson correction is now a
centered local source, the complete cost coefficient
has a connected formula with all native density
and metric terms present, and its magnitude is
proved to grow at most like $n^3$ uniformly in $g$.
One nontrivial term is evaluated exactly through
a newly justified scalar Haar factorization.
This is stronger than merely observing an even
Taylor series or bounding a frozen matrix.

\paragraph{What is not closed.}
No sign or fully evaluated numerical formula
for $\mathcal C_2(g)$ is claimed. The remaining
native $j,k,A$ contractions and the local electric
inverse in \eqref{c18v56:connected} should now be
evaluated on their connected factor supports,
or controlled directly in a finite-time local norm.
The coefficient bound alone gives neither a
regulator-uniform Taylor radius nor a bound at
block time four. A useful source-cost estimate
would still not by itself control the Hessian
of the Poisson corrector required for V51's
transport comparison. Coarse variance/return,
weak-coupling scale transport, physical-time
calibration and interacting continuum
reconstruction remain separate obligations.
The full Yang--Mills mass gap is still unproved.
\addtocontents{toc}{\protect\setcounter{tocdepth}{2}}
\addtocontents{toc}{\protect\clearpage}
% END CYCLE 18 VERSION 56 APPEND
% BEGIN CYCLE 18 VERSION 57 APPEND
\clearpage
\section{V57: the complete native first-corrector spectrum}
\label{c18v57:section}
\addtocontents{toc}{\protect\setcounter{tocdepth}{1}}

\subsection{Context and the specific debt paid}

V56 identified the actual first Poisson correction but left
its source in terms of a differentiated normal projector.
This append removes that derivative and computes the
\emph{entire} electric inverse energy of this source.
The result is not a restriction to V54's first band and
not the frozen resolvent in V55.

Keep the literal native sampler, normal chart and probability
law of V51--V56. In particular the fine spatial side is
$2n$, $n\ge5$, the group is unit-round $\mathrm{SU}(2)$,
the coupling here is $\beta=0$, and the coarse field $g$
is arbitrary and fixed. All expectations below are under
the original product-Haar conditional law $\mu_g$.
Write
\[
 \begin{gathered}
 M=M_g,\quad L=L_0,\quad R=L^+,\quad
 C=W_{\rm cap},\quad U=W_{\rm uncap},\quad W=C+U,\\
 LC=9C,\quad LU=12U,\quad
 \Gamma(f,h)=\langle df,dh\rangle_{m_0^{-1}},\quad
 \Gamma(f)=\Gamma(f,f).
 \end{gathered}
 \tag{V57.1}\label{c18v57:base}
\]
The letter $C$ in this append denotes the captured
\emph{action}, not the full inverse cost $\mathcal C$.
For the ambient gradient $F=\nabla W$, let
$P=P_0$ be the tangent projector of $b_2$ and $Q=I-P$.
Put
\[
       T=PF=\nabla_M W,\qquad N=QF,\qquad S=|N|^2 .
 \tag{V57.2}\label{c18v57:normal}
\]
V56 uses $u_t=\log\operatorname{Jac}_{\rm amb}\eta_t$,
$b=u'_0=3C$, $j=u''_0$, and the fixed-Haar-centered
Poisson representative $q_t$. Its exact first correction is
\[
           q'_0=Rh,\qquad h=j+\Gamma(b)/9=j+\Gamma(C).
 \tag{V57.3}\label{c18v57:target}
\]
There is no centering term missing in this formula:
V56 proved $\mathbb E j=-27n^3$ and
$\mathbb E\Gamma(C)=27n^3$.
The original moving-law-centered derivative differs
from $q'_0$ by the constant $-3n^3$.

The supplied ESM V64, source-selection V52, stationarity
V54 and exact-action V45 files have unchanged hashes
relative to the V55 review. Their useful constraint is
retained: identify the actual source, chart and spectral
weights before importing a finite-spectrum conclusion.
No companion's finite-mode or fixed-cutoff statement is
being promoted to a YM continuum theorem.

\subsection{Eliminating the differentiated projector}

\begin{proposition}[An exact scalar formula for the native density two-jet]
\label{c18v57:jprop}
On every base fibre,
\[
 \boxed{\quad
       j=(L-12)S+3\Gamma(W,C),\qquad
       h=(L-12)S+4\Gamma(C)+3\Gamma(U,C).
       \quad}
 \tag{V57.4}\label{c18v57:j}
\]
The formula uses the original normal field $N$, not a
chosen tangential replacement for the normal chart.
\end{proposition}

\begin{proof}
Let $\Phi_t$ be the ambient flow of $F$ and set
$\theta_t=\Phi_t\circ\eta_t$. Since
$b_2\Phi_t\eta_t=b_2$, this is an ambient diffeomorphism
preserving each base fibre. Its initial velocity along
$M$ is $T$. Let $Y_t=\dot\theta_t\circ\theta_t^{-1}$
on $M$ and let $H=\nabla F$ be the ambient Hessian
endomorphism.

We first compute its Eulerian velocity derivative.
The base fibre is totally geodesic, so along tangent
directions the tangent and normal subbundles are
parallel. Consequently
\[
             \nabla_M S=2PHN .
 \tag{V57.5}\label{c18v57:gradS}
\]
If $A_\eta=\nabla_t\dot\eta_t|_0$, differentiating
the normality of $\dot\eta_t$ to $d\eta_t(TM)$ gives
\[
 \langle A_\eta,v\rangle
    =-\langle-N,\nabla_v(-N)\rangle
    =-\tfrac12\,vS,\qquad v\in TM .
\]
This is the acceleration identity used in
Lemma~\ref{c18v53:variationlemma}.
The covariant chain rule for $\theta_t=\Phi_t\eta_t$
therefore gives
\[
 \begin{aligned}
 P A_\theta
    &=P A_\eta-2PHN+PHF
      =PHT-2PHN,\\
 Y'_0&=A_\theta-\nabla_T^M T
      =-2PHN=-\nabla_M S .
 \end{aligned}
 \tag{V57.6}\label{c18v57:Y}
\]
Here $A_\theta$ is tangent because $\theta_t$ stays
in the totally geodesic $M$, and
$\nabla_T^M T=PHT$. Thus no normal acceleration
component has been guessed.

Put $\rho_t=\operatorname{Jac}_{\rm amb}\Phi_{-t}$.
The ambient divergence satisfies $\operatorname{div}F=-12W$:
each fine plaquette is fundamental on its four
unit-round link factors. Hence, at a fixed ambient point,
\[
 \left.\partial_t\log\rho_t\right|_0=12W,\qquad
 \left.\partial_t^2\log\rho_t\right|_0=-12|F|^2 .
 \tag{V57.7}\label{c18v57:rho}
\]
The initial block map has constant coarea Jacobian,
since its Gram matrix is $2I$ in coarse invariant frames.
For a diffeomorphism preserving $b_2$, the coarea
change-of-variables formula therefore identifies its
ambient Jacobian with its fibre Jacobian. Apply this to
$\theta_t$ and use $\eta_t=\Phi_{-t}\theta_t$:
\[
 u_t=\log\rho_t\circ\theta_t
          +\log\operatorname{Jac}_{M}\theta_t .
 \tag{V57.8}\label{c18v57:factorjac}
\]
This identity is for unnormalized Jacobians; it is not
an assumed factorization of the moving conditional law.

The second derivative of the first term is
$-12|F|^2+24TW$. The second derivative of the second is
$\operatorname{div}_M Y'_0+T\operatorname{div}_M T$.
Using $|F|^2=S+\Gamma(W)$,
$\operatorname{div}_M T=-LW$ and
\eqref{c18v57:Y}, their sum is
\[
 \begin{aligned}
 j&=-12S+12\Gamma(W)+LS-\Gamma(W,LW)\\
  &=(L-12)S+3\Gamma(W,C),
 \end{aligned}
\]
because $LW=12W-3C$. Adding $\Gamma(C)$ proves
the second identity. This also checks the first
Jacobian derivative: $12W-LW=3C=b$.
\end{proof}

\subsection{The normal-copy identity and the diagonal sector}

For a captured pair factor $i$, let $\Gamma_i$ denote
its contribution to $\Gamma$; its metric is twice the
round metric. For a captured plaquette $p$, put
$w_p=\operatorname{Re}U_p$, and let
$\sigma_{ip}$ be $+1$ if it uses the first half of
pair $i$ and $-1$ if it uses the second. Set it to
zero if $p$ misses the pair. These signs describe
which half is used, not the orientation of traversal.

The tangent and normal frames in
\eqref{c18v53:frames}--\eqref{c18v53:halfsign}
give the pointwise identity
\[
     S=\sum_i\ \sum_{p,q\ {\rm captured}}
               \sigma_{ip}\sigma_{iq}\Gamma_i(w_p,w_q).
 \tag{V57.9}\label{c18v57:S}
\]
Indeed, the normal gradient contribution of a plaquette
on a pair is the isometric normal copy of its pair
tangent gradient, with the displayed sign. An uncaptured
plaquette has no captured link and gives no normal
contribution. Squaring the sum on each orthogonal
pair normal factor proves the formula.

\begin{lemma}[The entire free-link-even source sector]
\label{c18v57:diagonal}
Let $\Pi_{\rm even}$ average independently over the
central sign change on every original free link. Then
\[
 \begin{gathered}
  \xi_p=w_p^2-\tfrac14,\qquad
  L\xi_p=24\xi_p,\qquad
  \mathbb E\xi_p^2=\tfrac1{16},\\
  \Pi_{\rm even}h=-24\sum_{p\ {\rm captured}}\xi_p,\qquad
  \Pi_{\rm even}q'_0=-\sum_{p\ {\rm captured}}\xi_p,\\
  \|\Pi_{\rm even}h\|_2^2=432n^3,\qquad
  \langle\Pi_{\rm even}h,R\Pi_{\rm even}h\rangle=18n^3 .
 \end{gathered}
 \tag{V57.10}\label{c18v57:self}
\]
The distinct $\xi_p$ in these sums are orthogonal.
\end{lemma}

\begin{proof}
A captured plaquette has two distinct pair factors and
two distinct free factors. Its diagonal contribution to
$S$ is $1-w_p^2$, and to $\Gamma(C)$ is
$3(1-w_p^2)$. Thus its diagonal contribution to
\eqref{c18v57:j} is
\[
       (L-12)(1-w_p^2)+12(1-w_p^2)
              =L(1-w_p^2)=-24\xi_p .
\]
The square of a fundamental coordinate minus its Haar
mean is degree two on each of these four factors.
The degree-two energies are four per pair and eight
per free factor, giving twenty-four.
The sphere moments $\mathbb E w_p^2=1/4$,
$\mathbb E w_p^4=1/8$ give the norm.

Different fine plaquettes have different free-link sets.
For a mixed product and any of its differentiated
contractions, a free factor in their symmetric difference
remains odd. The Laplacian preserves this parity.
Thus $\Pi_{\rm even}$ annihilates every distinct-plaquette
term in $h$. For two distinct $\xi_p,\xi_q$, integrate
a private free factor of $p$ first: the conditional mean
of $\xi_p$ is zero. This proves orthogonality.
There are $12n^3$ captured plaquettes, proving the last
line and the inverse formula. This argument does not
require independence of the kinetic variables of
different action cells.
\end{proof}

\subsection{All mixed motifs and their exact two-channel decomposition}

Consider the product $f=w_pw_q$ in a nonzero
distinct-plaquette contraction. If the plaquettes share
one original pair or free factor, let $f_0$ be its
degree-zero projection and $f_2=f-f_0$.
There is one special case: adjacent corners around the
same two-odd vertex share both a pair prefix and a free
spoke. Their common variable is the path holonomy
$Z=xU$ (with fixed translations or inverses according
to the native orientation). In this case project in
that common $Z$. The scalar part removes \emph{both}
shared factors; its complement has degree two on both.
This defines $f_0,f_2$ in every row of the following table.

\begin{lemma}[Exhaustive native motif table]
\label{c18v57:motifs}
The contribution of one motif to $h$ is
$a_0f_0+a_2f_2$ with the following data.
Counts in the table are multiples of $n^3$.
CC pairs are unordered; UC pairs have specified types.
\[
\begin{array}{@{}lrrr@{}}
 \toprule
 \text{motif} & \text{count}/n^3 &(E_0,E_2)&(a_0,a_2)\\
 \midrule
 \text{CC, adjacent at one centre} &12 &(9,21)&(45,-3)\\
 \text{CC, distinct centres, same half} &36 &(15,19)&(21,-11)\\
 \text{CC, distinct centres, opposite halves} &36 &(15,19)&(3,3)\\
 \text{UC, one shared free spoke} &48 &(15,23)&(9,-3)\\
 \bottomrule
\end{array}
 \tag{V57.11}\label{c18v57:table}
\]
In every row $Lf_r=E_rf_r$, and
\[
       \|f_0\|_2^2=\frac1{64},\qquad
       \|f_2\|_2^2=\frac3{64},\qquad
       \langle f_0,f_2\rangle=0 .
 \tag{V57.12}\label{c18v57:fusionnorm}
\]
There are no other mixed contributions to $h$.
\end{lemma}

\begin{proof}
First consider a single common Haar quaternion $e$.
Conditionally on the other factors the two characters
are linear coordinates $e\cdot a$, $e\cdot b$ for
unit vectors $a,b\in\mathbb R^4$. Each plaquette has
a private free factor, so the remaining Haar integrals
make $a,b$ independent and uniform, conditionally on
all factors they otherwise share. Fixed coarse
multiplications are orthogonal maps and do not change
this statement. Thus
\[
 \begin{aligned}
 f_0&=(a\cdot b)/4,\\
 \mathbb E_e f^2&=\bigl(1+2(a\cdot b)^2\bigr)/24,\\
 \mathbb E(a\cdot b)^2&=1/4 .
 \end{aligned}
 \tag{V57.13}\label{c18v57:fusion}
\]
These identities give $\|f\|_2^2=1/16$ and
\eqref{c18v57:fusionnorm}. The polynomial
$(e\cdot a)(e\cdot b)-(a\cdot b)|e|^2/4$
is homogeneous harmonic of degree two. This proves
the claimed scalar/degree-two decomposition.

For adjacent corners at one centre, use the common
path variable in place of $e$. The two exterior path
variables have distinct private free factors, so the
same norm proof applies. The degree-two component
in the common path has Casimir degree two separately
in its prefix and spoke, by invariance under left
and right multiplication. In particular it has
combined shared energy $4+8=12$; this assertion
does not identify the global kinetic operator with
one in independent path coordinates.

For the coefficients, the positive-Laplacian product
rule on one shared factor gives
\[
  \Gamma_i(w_p,w_q)
       =\frac{(L_iw_p)w_q+w_p(L_iw_q)-L_i(w_pw_q)}2 .
 \tag{V57.14}\label{c18v57:productrule}
\]
For a common pair its coefficients on $(f_0,f_2)$
are $(3/2,-1/2)$; for a common free factor they
are $(3,-1)$. A same-centre adjacent pair has both,
so its full $\Gamma(w_p,w_q)$ has coefficients
$(9/2,-3/2)$. The two plaquettes use opposite
halves of the shared pair, and their exterior
factors have energy nine. Its contribution is
\[
        -2(L-12)\Gamma_i(w_p,w_q)
                         +8\Gamma(w_p,w_q).
\]
Inserting energies nine and twenty-one gives
$(45,-3)$.

A CC pair at distinct centres shares just one
pair factor. With
$\sigma=\sigma_{ip}\sigma_{iq}$ its contribution is
\[
        2\{\sigma(L-12)+4\}\Gamma_i(w_p,w_q).
 \tag{V57.15}\label{c18v57:CC}
\]
Its exterior factors have energy fifteen and the
common degree-two pair adds four. This gives
$(21,-11)$ if $\sigma=1$ and $(3,3)$ if $\sigma=-1$.
A UC pair has no normal cross term and occurs only
in $3\Gamma(U,C)$. The common free factor has
coefficients $(3,-1)$, exterior energy fifteen and
degree-two increment eight. Its contribution
is therefore $(9,-3)$ at energies fifteen and
twenty-three.

For completeness, the lattice counts and exhaustion
are local, not an extrapolation from a few volumes.
Each captured pair has eight incident captured
plaquettes, four using each half. Among its
$\binom82=28$ unordered pairs, twelve use the same
half and sixteen use opposite halves. Four of
the latter meet at the same two-odd centre and
share its transverse free spoke. The other twelve
have distinct centres. Distinct captured plaquettes
cannot share two pair factors: their two pair
chains determine their square uniquely. Also any
shared captured free spoke forces the same centre
and its unique adjacent pair prefix. Consequently
there is no double counting and no additional
CC overlap. Multiplication by the $3n^3$ pair
factors gives the first three counts.

There are $12n^3$ free spokes joining one-odd to
two-odd vertices. Each lies in two captured and
two uncaptured plaquettes, producing four UC
pairs. Such a pair shares only that one free
edge: the uncaptured square's two edges incident
to its three-odd vertex are never captured
spokes, and its two remaining spokes meet at
a one-odd vertex, whereas the captured square's
two free spokes meet at a two-odd vertex.
Sharing both would force two different meeting
vertices for the same adjacent-edge pair.
Hence the UC count is $48n^3$.
Uncaptured plaquettes have no pair factor,
and there is no $\Gamma(U)$ term in $h$.
Disjoint-factor pairs have zero contraction.
These observations exhaust \eqref{c18v57:j}.
The assumption $n\ge5$ excludes periodic short
identifications in these local configurations.
\end{proof}

\subsection{Why almost all motifs are orthogonal}

It would be incorrect to add all entries of
\eqref{c18v57:table} in squared norm independently.
Two routes at a common centre can interfere.
We now isolate exactly that exception.

\begin{lemma}[Parity separation and the sole coherent return]
\label{c18v57:orthogonality}
Distinct motifs in \eqref{c18v57:table} are orthogonal
under every spectral function of $L$, except possibly
for their energy-nine components at the same centre.
The latter are exactly the two adjacent-corner routes
joining the same pair of opposite spokes.
\end{lemma}

\begin{proof}
Associate to a motif the free-link sign signature
$\mathcal F(p)\mathbin{\triangle}\mathcal F(q)$.
Every $f_r$ has that signature. The independent
central sign changes are commuting Haar isometries
commuting with $L$, so different signatures are
orthogonal under spectral functions of $L$.

For CC motifs at distinct centres the signature
contains four spokes, two at each centre.
The owner of a captured spoke is its unique
two-odd endpoint. At either centre the two
spokes determine the captured corner square,
so the signature determines the unordered motif.

A UC signature contains the uncaptured square's
two edges incident to its three-odd vertex.
These are distinguishable from captured spokes
by their endpoint parities. Together they
determine that uncaptured square uniquely.
Taking its free set's symmetric difference
with the signature then recovers the captured
square's free pair, which determines the
captured square. Thus UC signatures are unique,
and cannot coincide with CC signatures.

For CC motifs at one centre the signature is
the two unshared, opposite spokes. Exactly
two adjacent-corner routes give it, one through
each of the other two spokes of the four-cycle.
The energy-twenty-one part of each route has
degree two on its shared spoke. The other route
does not contain that spoke at all. Integrating
it annihilates the cross product, even if other
pair prefixes are shared. Energy-nine and
energy-twenty-one components are orthogonal
by their distinct eigenvalues. Only the two
energy-nine route components can interfere.
The diagonal sector has empty signature and
is already handled by Lemma~\ref{c18v57:diagonal}.
\end{proof}

Let $P_9$ be the full invariant energy-nine projection
from V54, including both corner and straight channels.
The product rule and the motif classification imply
\[
 \begin{gathered}
    P_9\Gamma(C)=\tfrac92 P_9C^2,\qquad
    P_9S=-\tfrac32 P_9C^2,\qquad
    P_9\Gamma(U,C)=0,\\
    P_9h=\tfrac{45}{2}P_9C^2 .
 \end{gathered}
 \tag{V57.16}\label{c18v57:P9}
\]
For the second identity, an unordered adjacent
CC pair contributes $-2(3/2)f_0$ to $S$ and
$2f_0$ to $C^2$; all other pieces miss energy nine.
The first identity also follows directly from
$L(C^2)=18C^2-2\Gamma(C)$.

Recall the actual coarse holonomy $H_y$ and write
\[
                    W_c(g)=\sum_y\operatorname{Re}H_y .
 \tag{V57.17}\label{c18v57:Wc}
\]
V55's source calculation gives
$\|M_{H_y}s_{H_y}\|^2=(1+\operatorname{Re}H_y)/4$,
where $s_{H_y}$ is the coefficient vector of $W_y$.
Since the original $C$ lies in $\operatorname{ran}P_9$,
\[
       \|P_9C^2\|_2^2
        =\|P_9 C P_9 C\|_2^2
        =\frac{3n^3+W_c(g)}4 .
 \tag{V57.18}\label{c18v57:C2nine}
\]
This reuses V54--V55's \emph{complete} local compression
and its arbitrary-coarse-field twist, not a flat-field
or independent-route approximation.

\clearpage
\subsection{The complete source measure and inverse energy}

\begin{theorem}[An exact six-energy formula for the actual first correction]
\label{c18v57:spectrumthm}
Let $P_E=\mathbf1_{\{E\}}(L)$. The native source in
\eqref{c18v57:target} has spectral support contained in
$\{9,15,19,21,23,24\}$. Its complete unnormalized
spectral measure is
\[
\begin{array}{@{}cr@{}}
 \toprule
 E & \|P_Eh\|_2^2\\[3pt]
 \midrule
 9  & \dfrac{2025}{16}\bigl(3n^3+W_c(g)\bigr)\\[8pt]
 15 & \dfrac{2511}{8}n^3\\[8pt]
 19 & \dfrac{1755}{8}n^3\\[8pt]
 21 & \dfrac{81}{16}n^3\\[8pt]
 23 & \dfrac{81}{4}n^3\\[8pt]
 24 & 432n^3\\
 \bottomrule
\end{array}
 \tag{V57.19}\label{c18v57:measure}
\]
The energy-nine weight can vanish; no omitted
energy carries additional source weight.
In particular
\[
 \boxed{\quad
  \mathbb E[hRh]
      =K n^3+\frac{225}{16}\bigl(3n^3+W_c(g)\bigr),
  \qquad
  K=18+\frac{540}{19}+\frac{27}{112}+\frac{567}{115}.
       \quad}
 \tag{V57.20}\label{c18v57:inverse}
\]
Equivalently $K=12625731/244720$. This is the entire
first-corrector inverse energy in
\eqref{c18v56:connected}, not the entire coefficient
$\mathcal C_2(g)$.
\end{theorem}

\begin{proof}
The complete expansion is the diagonal sector plus
the four rows of \eqref{c18v57:table}.
Thus its only possible energies are the six listed.
Lemmas~\ref{c18v57:motifs} and
\ref{c18v57:orthogonality} justify adding all
non-nine squared norms separately. Per unit $n^3$,
the energy-fifteen weight is
\[
       \frac{36\cdot21^2+36\cdot3^2+48\cdot9^2}{64}
                   =\frac{2511}{8},
\]
and the energy-nineteen weight is
$3(36\cdot11^2+36\cdot3^2)/64=1755/8$.
The weights at twenty-one and twenty-three are,
respectively, $3(12\cdot3^2)/64=81/16$ and
$3(48\cdot3^2)/64=81/4$.
The diagonal weight is $432n^3$.
Equations~\eqref{c18v57:P9} and
\eqref{c18v57:C2nine} give the coherent
energy-nine weight without separating its routes.

Since $\mathbb E h=0$, the spectral theorem yields
$\mathbb E[hRh]=\sum_{E>0}\|P_Eh\|_2^2/E$.
The non-nine sum is
\[
       18+\frac{837}{40}+\frac{1755}{152}
                        +\frac{27}{112}+\frac{81}{92}=K.
\]
The nine term gives the second term of
\eqref{c18v57:inverse}. Every spectral component
of this particular source has been retained.
\end{proof}

Several checks and consequences are worth recording.
First, directly from \eqref{c18v57:S},
$\mathbb ES=9n^3$, whereas
$\mathbb E\Gamma(W,C)=\langle W,LC\rangle=27n^3$.
Thus \eqref{c18v57:j} independently recovers
$\mathbb E j=-27n^3$, agreeing with V56's
coarse-marginal calibration. Second,
\[
 \begin{gathered}
 \|h\|_2^2=\frac{5481}{4}n^3+\frac{2025}{16}W_c(g),\\
 q'_0=\sum_{E\in\{9,15,19,21,23,24\}}E^{-1}P_Eh,
       \qquad P_9q'_0=\tfrac52 P_9C^2 .
 \end{gathered}
 \tag{V57.21}\label{c18v57:q}
\]
All these sources are invariant under the based
internal gauge group. Hence the full-space and
invariant-space inverses agree on them.

Since $-3n^3\le W_c(g)\le3n^3$, the exact uniform
source-energy bounds are
\[
       Kn^3\le\mathbb E[hRh]
                     \le\left(K+\frac{675}{8}\right)n^3 .
 \tag{V57.22}\label{c18v57:range}
\]
The upper endpoint is attained at a flat coarse field.
For even $n$, V55's all-negative coarse plaquette
example attains the lower endpoint: its entire
energy-nine correction vanishes, while the other
five energies still give $Kn^3>0$.
This is a source calculation at the electric
reference, not a positivity statement about
the curvature of the full temporal cost.

\subsection{Verification, literature boundary and the next upstream task}

The reproducible diagnostic is
\begin{center}\small
\path{scripts/verify_ym_c18_v57_native_corrector_spectrum.py}.
\end{center}
It enumerates native incidence on $n=5,6$, checks
all four motif counts and the uniqueness of the
free-link signatures, and identifies the two
same-centre routes as the sole repeated signatures.
It verifies the harmonic trace subtraction and
all rational spectral weights, and replays V55's
arbitrary-unit-quaternion source matrix. The
all-volume claims rest on the local classification
and Haar proofs above, not numerical extrapolation.

A separate, explicitly \emph{non-native} sign fixture
uses the flat two-torus metric
$dx^2+\tfrac13dy^2$, fibres $y=\text{constant}$ and
$W=\cos(3x)\cos y$. Its ambient eigenvalue is twelve
and fibre eigenvalue nine. Differentiating the
actual sampler $c_t=y\circ\Phi_t$ gives its
projector derivative directly; symbolic
calculation verifies that
$-\operatorname{div}(Q'_0F)+N\operatorname{div}N$
equals \eqref{c18v57:j}.
This checks signs and material derivatives, not
a native lattice trajectory or continuum limit.

For context, the primary abstract of Bauer,
D'Andrea, Freytsis and Grabowska,
\emph{A new basis for Hamiltonian SU(2) simulations}
(2023), \url{https://arxiv.org/abs/2307.11829},
was rechecked. It discusses representation-basis
Hamiltonian truncation and its coupling dependence.
The distinction here is important: the finite
support in \eqref{c18v57:measure} is an exact
property of one polynomial source for the original
electric operator. It is \emph{not} a finite-dimensional
replacement of the interacting Hamiltonian.
No weak-coupling approximation or gap estimate
is imported from that work. This is a targeted
context check, not the next broad sweep.
The companion review and broad literature sweep
remain due at V60; archive/restart remains V100.

Removing this append and restoring the title
recovers V56 byte for byte. The other 53
sources are protected and unchanged. Only V57
is retained as the active YM filename.
Compilation and inspection products remain
outside \path{latex_notes}.

\clearpage
\paragraph{What has advanced.}
The unevaluated projector derivative in the
native source has been eliminated. Its exact
two-plaquette decomposition has been resolved
over all six electric energies, including the
coherent coarse-holonomy-dependent return.
Thus the entire term $\mathbb E[hRh]$ in
V56's actual cost formula is now evaluated.
No unknown source modes were discarded to obtain it.

\paragraph{What is still open.}
The mixed $j,k,A$ terms in
\eqref{c18v56:connected} have not all been
evaluated. The explicit $j$ formula now makes
the two $j$ cumulants a finite local Haar
calculation; the cubic density contraction
$\mathbb E[bk]$ and inverse-metric contraction
$\mathbb E[A(db,db)]$ remain distinct obligations.
A positive inverse energy cannot determine
the sign of their full sum.

The next useful step is to resolve those actual
local contractions, or use the explicit corrector
to obtain a volume-uniform local derivative
estimate. Merely iterating another abstract
resolvent bound would leave the real debt
unchanged. Even a fully evaluated $\mathcal C_2$
would not give a uniform Taylor remainder,
control through block time four, or the
Poisson-Hessian estimate required by V51.
Coarse variance/return, weak-coupling scale
transport, physical-time normalization and
interacting continuum reconstruction are still
unpaid. The full Yang--Mills mass gap remains
unproved.

\addtocontents{toc}{\protect\setcounter{tocdepth}{2}}
\addtocontents{toc}{\protect\clearpage}
% END CYCLE 18 VERSION 57 APPEND
% BEGIN CYCLE 18 VERSION 58 APPEND
\clearpage
\section{V58: density cumulants and the actual metric contribution}
\label{c18v58:section}
\addtocontents{toc}{\protect\setcounter{tocdepth}{1}}

\subsection{Context: finishing three contractions, not claiming a gap}

V57 resolved the whole first-corrector inverse energy.
The next debt in V56's actual temporal coefficient
is not another spectral truncation: it is the
remaining density and metric contractions.
This append evaluates both mixed $j$ cumulants
and the \emph{entire} inverse-metric contribution.
After these calculations, only one scalar
cubic-density contraction remains unevaluated in
that particular coefficient.

Keep the native unit-round $\mathrm{SU}(2)$ model
at $\beta=0$, fine side $2n$, $n\ge5$, and arbitrary
fixed coarse field $g$. All expectations are
under the original product-Haar fibre law.
Use V57's notation
\[
 \begin{gathered}
 C=W_{\rm cap},\quad U=W_{\rm uncap},\quad W=C+U,\quad
 L=L_0,\quad R=L^+,\quad b=3C,\\
 D=\nabla_M C,\quad N=Q_0\nabla W=Q_0\nabla C,\quad
 S=|N|^2,\quad W_c=\sum_y\operatorname{Re}H_y .
 \end{gathered}
 \tag{V58.1}\label{c18v58:base}
\]
The equality in the definition of $N$ follows because
$U$ uses no captured link. Let $A$ continue to mean
the actual coefficient in
$\widehat m_t^{-1}=m_0^{-1}+t^2A+O(t^3)$.
It is not the chart acceleration, nor a metric
chosen for convenience. The raw density coefficients
remain $j=u''_0$, $k=u'''_0$ for
$u_t=\log\operatorname{Jac}_{\rm amb}\eta_t$.

The full interacting continuum construction and
positive physical YM mass gap remain open.
An exact second-order coefficient at the electric
reference does not pay a finite-time or continuum
estimate.

\subsection{The complete mixed scalar contractions}

Let $\xi_p=w_p^2-1/4$ for a captured plaquette $p$.
V57 proved $L\xi_p=24\xi_p$ and
$\mathbb E\xi_p^2=1/16$, with different $\xi_p$
orthogonal. It also defined $f_0,f_2$ for each
mixed motif, with squared norms $1/64,3/64$.
Subtracting $\Gamma(C)$ from V57's $h$ gives
\[
 \begin{gathered}
 j_{\rm diag}=-27n^3-21\sum_p\xi_p,\qquad
 S_{\rm diag}=9n^3-\sum_p\xi_p,\\
\begin{array}{@{}lrrr@{}}
 \toprule
 \text{motif}&(E_0,E_2)&j\text{ coefficients}&S\text{ coefficients}\\
 \midrule
 \text{CC, adjacent}&(9,21)&(36,0)&(-3,1)\\
 \text{CC, same half}&(15,19)&(18,-10)&(3,-1)\\
 \text{CC, opposite halves}&(15,19)&(0,4)&(-3,1)\\
 \text{UC}&(15,23)&(9,-3)&(0,0)\\
 \bottomrule
\end{array}
 \end{gathered}
 \tag{V58.2}\label{c18v58:scalarrows}
\]
In particular $P_9j=18P_9C^2$ and
$P_9S=-\tfrac32P_9C^2$.
The adjacent, same-half and opposite-half counts
are respectively $12n^3,36n^3,36n^3$.

\begin{proposition}[Two exact cumulants and a metric Hessian input]
\label{c18v58:scalarprop}
Uniformly for all coarse $g$,
\[
 \begin{aligned}
 \kappa(b,b,j)&=-\frac{81}{4}n^3+\frac{81}{2}W_c,\\
 \kappa(b,b,Lj)&=-\frac{6075}{2}n^3+\frac{729}{2}W_c,\\
 \mathbb E[S L^2(C^2)]&=-27n^3-\frac{243}{8}W_c .
 \end{aligned}
 \tag{V58.3}\label{c18v58:scalars}
\]
The last expression is not asserted positive;
$S$ and $L^2(C^2)$ are different functions.
\end{proposition}

\begin{proof}
We must first check the full $C^2$, including products
of plaquettes that do not share a factor. Its diagonal
centered part is $\sum_p\xi_p$.
Its unordered distinct-plaquette terms are $2w_pw_q$.

For distinct centres, the free-link signature has
two spokes at each centre, which recover the two
captured squares uniquely. It cannot match a
contributing $j$ or $S$ motif unless it is that
very pair of squares. At a common centre an adjacent
pair has a two-spoke signature; its only repetition
is the other route between those opposite spokes.
An opposite-corner pair has a four-spoke signature
at one centre and matches none of the contributing
rows. A UC row contains edges incident to a
three-odd vertex and cannot match any product
of two captured plaquettes. These are exactly
V57's free-parity mechanisms, now applied to
\emph{all} of $C^2$, not only its overlapping pairs.
Thus no disconnected $C^2$ term has been silently
dropped.

Write $F_j=\mathbb E[C^2(j-\mathbb Ej)]$ and
$G_j=\mathbb E[C^2Lj]$. The diagonal contributions
are $-63n^3/4$ and $-378n^3$.
The energy-nine contributions, with both routes
retained, are
\[
 18\,\frac{3n^3+W_c}{4},\qquad
 162\,\frac{3n^3+W_c}{4}.
\]
For one different-centre same-half motif,
the contributions are
\[
       \frac{2(18-30)}{64}=-\frac38,\qquad
       \frac{2(15\cdot18-19\cdot30)}{64}=-\frac{75}{8}.
\]
For one opposite-half motif they are
$3/8$ and $57/8$. The two types have equal
count $36n^3$. There is no adjacent energy-twenty-one
contribution because its $j$ coefficient is zero.
Therefore
\[
 F_j=-\frac94n^3+\frac92W_c,\qquad
 G_j=-\frac{675}{2}n^3+\frac{81}{2}W_c .
\]
Multiplication by nine proves the two cumulants,
using $\mathbb Eb=0$, $\mathbb Ej=-27n^3$ and
$\mathbb E Lj=0$.

For the last expression in \eqref{c18v58:scalars},
the diagonal contribution is $-432n^3$.
The two different-centre types cancel, since
their $S$ coefficients are opposite and their
counts agree. The adjacent energy-nine contribution
is $-(243/8)(3n^3+W_c)$, using
$\|P_9C^2\|^2=(3n^3+W_c)/4$.
Its energy-twenty-one contribution is
\[
       (12n^3)\,\frac{2\cdot3}{64}\,21^2
                     =\frac{3969}{8}n^3 .
\]
Adding gives the stated value.
\end{proof}

\subsection{A compact eigenfunction identity for the metric term}

\begin{lemma}[Integrated Hessian against an eigenfunction gradient]
\label{c18v58:hesslemma}
On a compact Riemannian manifold without boundary,
use the positive Laplacian $L=-\operatorname{div}\nabla$
and its normalized volume measure.
If $LC=\lambda C$ and $D=\nabla C$, then for
every smooth real $S$,
\[
       \mathbb E[\operatorname{Hess}S(D,D)]
                         =\frac14\mathbb E[S L^2(C^2)] .
 \tag{V58.4}\label{c18v58:hessibp}
\]
\end{lemma}

\begin{proof}
Since $\operatorname{div}D=-\lambda C$ and
$\nabla_DD=\tfrac12\nabla\Gamma(C)$, integration
of the derivative $D(DS)$ gives
\[
 \begin{aligned}
 \mathbb E[\operatorname{Hess}S(D,D)]
  &=\lambda\mathbb E[C\Gamma(C,S)]
          -\tfrac12\mathbb E[\Gamma(\Gamma(C),S)]\\
  &=\frac{\lambda}{2}\mathbb E[S L(C^2)]
          -\tfrac12\mathbb E[S L\Gamma(C)] .
 \end{aligned}
\]
The positive-Laplacian product rule is
$\Gamma(C)=\lambda C^2-\tfrac12L(C^2)$.
Substitution cancels the two first-order Laplacian
terms and proves the formula.
No curvature inequality or Hessian bound has been
assumed.
\end{proof}

The actual normal-chart metric two-jet in
\eqref{c18v53:secondmetric} now reduces its integrated
contribution to three explicit scalar quantities.
Define
\[
 \mathcal D=\mathbb E|\nabla_D^\perp N|^2,\qquad
 \mathcal R=\mathbb E\langle R_{\rm amb}(D,N)N,D\rangle .
 \tag{V58.5}\label{c18v58:DR}
\]
The curvature convention is V53's positive-sphere
convention. Since $m'_0=0$, inverse differentiation
and \eqref{c18v58:hessibp} give
\[
       \mathbb E[A(dC,dC)]
            =\mathcal R-\mathcal D
                        +\frac18\mathbb E[S L^2(C^2)] .
 \tag{V58.6}\label{c18v58:metricreduction}
\]
Every object here belongs to the original normal
chart. In particular, discarding the Hessian term
would change the answer.

\subsection{The ambient curvature contraction is exactly local}

\begin{lemma}[The complete curvature contraction]
\label{c18v58:curvlemma}
For the actual native fibre,
\[
                         \mathcal R=9n^3 .
 \tag{V58.7}\label{c18v58:curvature}
\]
\end{lemma}

\begin{proof}
Use the parallel tangent-to-normal copy on pair $i$
from \eqref{c18v53:frames}--\eqref{c18v53:connection}.
Let $v_p=\nabla_i w_p$ in its three-dimensional
unit tangent coordinates. Split the captured
plaquettes incident to this pair by their half:
\[
       a_i=\sum_{\sigma_{ip}=1}v_p,\qquad
       c_i=\sum_{\sigma_{ip}=-1}v_p .
\]
(The vector $c_i$ is not the sampling map.)
The pair components of $D$ and the copied $N$
are $a_i+c_i$ and $a_i-c_i$.
The two ambient unit spheres give
\[
 \langle R_{\rm amb}(D_i,N_i)N_i,D_i\rangle
  =\frac12\bigl(|D_i|^2|N_i|^2
                  -\langle D_i,N_i^{\rm copy}\rangle^2\bigr)
  =2|a_i\times c_i|^2 .
 \tag{V58.8}\label{c18v58:wedge}
\]
Different pair normal components and all free
factors give no further ambient curvature term.

To integrate its quartic expansion, each original
free link must have even central-sign parity.
At a two-odd centre the four captured squares
are the four edges of a four-cycle whose vertices
are its free spokes. A four-plaquette product
has even free parity only through repeated
plaquette pairs or the four different edges
of one such cycle. But a fixed captured pair
is incident to only two of the four squares
at any one centre. The latter possibility
therefore cannot occur in \eqref{c18v58:wedge}.
Products with exactly two odd-multiplicity
distinct plaquettes also have nonempty free
signature. Thus only the identical first-half
pair and identical second-half pair survive.

For a first-half square $p$ and a second-half
square $q$, their private free factors make
$v_p,v_q$ independent isotropic tangent vectors
conditionally on the common factors, with
\[
          \mathbb E(v_p)_a(v_p)_b
                    =\mathbb E(v_q)_a(v_q)_b
                    =\frac{\delta_{ab}}8 .
\]
This remains true when they share a free spoke:
integrate the private spokes first.
Consequently $2\mathbb E|v_p\times v_q|^2=3/16$.
There are sixteen first/second-half pairs at
each of the $3n^3$ captured pair factors.
This gives $48n^3(3/16)=9n^3$.
\end{proof}

\subsection{The full normal derivative, including its coherent return}

\begin{proposition}[A source-weighted normal derivative contraction]
\label{c18v58:normalprop}
The other geometric term in
\eqref{c18v58:metricreduction} is
\[
            \mathcal D=\frac{225}{4}n^3-\frac{243}{64}W_c .
 \tag{V58.9}\label{c18v58:normalvalue}
\]
Its coarse-holonomy term cannot be obtained by
treating the adjacent routes as independent.
\end{proposition}

\begin{proof}
The normal-copy map on each pair is an isometry
parallel for its tangent and normal connections.
It is constant along the other pair and free factors.
Thus, in copied tangent coordinates,
\[
 J_i:=(\nabla_D^\perp N)_i^{\rm copy}
       =\sum_{p,q}\sigma_{ip}
                    \operatorname{Hess}_{i,*}w_p(\nabla w_q).
 \tag{V58.10}\label{c18v58:J}
\]
The sum is over captured $p,q$. Terms missing a
shared differentiation factor vanish.
For an unordered distinct pair, its contribution is
\[
 J_{i;pq}=\sigma_{ip}\operatorname{Hess}_{i,*}w_p(\nabla w_q)
          +\sigma_{iq}\operatorname{Hess}_{i,*}w_q(\nabla w_p).
 \tag{V58.11}\label{c18v58:Jpq}
\]
A sign is zero when the corresponding plaquette
misses $i$. There is no factor of two on a
diagonal $p=q$ term.

All these vector components have the same
free-link signatures as their scalar motifs:
the pair frames are independent of the free
links, and differentiation preserves their
central parities. V57's classification therefore
again applies. There are diagonal terms and
the three CC types, but no UC term because
$D=\nabla C$. The only possible cross return
between distinct motifs is between the two
adjacent routes of one straight signature.

\paragraph{Diagonal terms.}
For one captured square,
$\Gamma(w_p)=3(1-w_p^2)$, so
$\operatorname{Hess}w_p(\nabla w_p)=-3w_p\nabla w_p$.
Summing its two pair normal components gives
\[
        \sum_i|J_{i;pp}|^2=9w_p^2(1-w_p^2),\qquad
        \mathbb E\sum_i|J_{i;pp}|^2=\frac98 .
 \tag{V58.12}\label{c18v58:Jself}
\]
Different diagonal terms are orthogonal:
on a private free factor,
$\mathbb E[w_p\nabla_iw_p]
=\tfrac12\nabla_i\mathbb E[w_p^2]=0$.
They are also orthogonal to all nonempty
mixed signatures. Their full contribution is
$12n^3(9/8)=27n^3/2$.

\paragraph{Different-centre motifs.}
Such a motif has one common pair and two exterior
pairs. At an exterior pair only one square occurs.
Differentiating the scalar contraction therefore gives
\[
          J_{i;pq}=\sigma_{ip}\nabla_i\Gamma(w_p,w_q)
\]
when $p$ is the incident square, and similarly for $q$.
Both fusion components have exterior pair energy three,
and $\Gamma(w_p,w_q)=\tfrac32f_0-\tfrac12f_2$.
Hence the summed exterior squared norm is
\[
                3\left(\frac{9}{4\cdot64}
                      +\frac{3}{4\cdot64}\right)=\frac9{64}.
 \tag{V58.13}\label{c18v58:Jouter}
\]

For the common pair, rotate its value to the identity
for the Haar moment computation.
The two characters are then linear coordinates
with independent uniform coefficient quaternions
$a=(a_0,\mathbf a)$ and $c=(c_0,\mathbf c)$.
The unit pair gradients are $\mathbf a/\sqrt2$,
$\mathbf c/\sqrt2$, and the pair Hessians are
$-a_0I/2$, $-c_0I/2$.
For equal and opposite half signs, respectively,
the middle components, up to an overall sign, are
\[
 -\frac{a_0\mathbf c+c_0\mathbf a}{2\sqrt2},
       \qquad
  \frac{c_0\mathbf a-a_0\mathbf c}{2\sqrt2}.
 \tag{V58.14}\label{c18v58:Jmiddle}
\]
Each has squared mean $3/64$ by the sphere
second moments. Thus either different-centre
motif has full squared mean $12/64=3/16$.
Their total is $72n^3(3/16)=27n^3/2$.

\paragraph{Adjacent motifs at one centre.}
There is now a shared pair and a shared free
spoke, and the shared pair half signs are opposite.
The exterior argument still applies, but its
scalar contraction is
$\Gamma(w_p,w_q)=\tfrac92f_0-\tfrac32f_2$.
Its summed exterior squared norm is therefore
\[
              3\left(\frac{81}{4\cdot64}
                       +\frac{27}{4\cdot64}\right)=\frac{81}{64}.
 \tag{V58.15}\label{c18v58:Jadjouter}
\]
The middle component also contains a mixed
pair/free Hessian. It is essential not to
replace it by a second derivative in one
independent path coordinate.

Here is the explicit local computation in the
original product metric. Write the common path
as $Z=xV$, where $x$ is its pair prefix and
$V$ its free spoke, absorbing any fixed
translations and inverses in orthogonal frames.
For the moment computation rotate their values
to $x=V=1$; the two exterior coefficient
quaternions remain independent Haar.
For $w_p=a\cdot(xV)$, and imaginary quaternion
basis $e_\alpha e_\beta
=-\delta_{\alpha\beta}
+\sum_\gamma\epsilon_{\alpha\beta\gamma}e_\gamma$,
the Hessian blocks are
\[
 \begin{aligned}
 (\operatorname{Hess}_{i,i}w_p)_{\alpha\beta}
       &=-\frac{a_0}{2}\delta_{\alpha\beta},\\
 (\operatorname{Hess}_{i,V}w_p)_{\alpha\beta}
       &=\frac1{\sqrt2}
          \left(-a_0\delta_{\alpha\beta}
                    +\sum_\gamma\epsilon_{\alpha\beta\gamma}a_\gamma\right).
 \end{aligned}
 \tag{V58.16}\label{c18v58:mixedH}
\]
The first line includes the pair Levi--Civita
connection. The second has no cross-factor
connection term and is the direct derivative
of $x e_\alpha V e_\beta$.
These formulas also show why the product
metric cannot be replaced by a scalar
path-metric prescription.

Substitution into \eqref{c18v58:Jpq} gives its
middle component, up to a common sign,
\[
        \frac1{\sqrt2}
        \left\{\frac32(c_0\mathbf a-a_0\mathbf c)
                              -2\mathbf a\times\mathbf c\right\}.
 \tag{V58.17}\label{c18v58:Jadjmiddle}
\]
The two vectors in braces are pointwise orthogonal.
Their squared means before the displayed coefficients
are both $3/8$. Hence the squared mean of
\eqref{c18v58:Jadjmiddle} is $75/64$.
Together with \eqref{c18v58:Jadjouter} this is
$156/64=39/16$ for one adjacent motif.
The independent-motif contribution is
$12n^3(39/16)=117n^3/4$.

All local moment reductions used private free
factors and Haar invariance. They are valid for
arbitrary coarse coefficients, not only a flat
coarse field.

\paragraph{The sign of the coherent return.}
The preceding three totals sum to $225n^3/4$,
but the adjacent-route cross terms remain.
For one straight signature, its two routes have
different middle pairs, so their middle components
cannot pair. They have the same two exterior
pairs. At each exterior pair their half signs
are \emph{opposite}: the two corners incident
to that spoke use the two different halves.

The $f_2$ component of a route has a private
degree-two shared free spoke and integrates to
zero against the other route. For the two $f_0$
components, denoted $f_0^{(1)},f_0^{(2)}$, the
exact four-cycle Haar integral is
\[
            \mathbb E[f_0^{(1)}f_0^{(2)}]
                         =\frac{\operatorname{Re}H_y}{64}.
 \tag{V58.18}\label{c18v58:routeoverlap}
\]
Indeed, integrate each route's shared spoke first;
their product then has the same integral as
the four different captured characters around
the centre. Three successive fundamental Haar
contractions give the indicated coarse character,
as in V56's fourth-cumulant calculation.

The two exterior pair Laplacians have total
energy three on either $f_0$. Integration by
parts and the reversed signs therefore make
the \emph{cross term in the squared norm}
for one straight signature
\[
        -2\cdot3\left(\frac92\right)^2
                   \frac{\operatorname{Re}H_y}{64}
                      =-\frac{243}{128}\operatorname{Re}H_y .
 \tag{V58.19}\label{c18v58:Jcross}
\]
There are two straight signatures at each centre.
Summing proves the correction $-243W_c/64$ and
\eqref{c18v58:normalvalue}. No other signature
can contribute a cross term.
\end{proof}

\subsection{The exact metric cost and the remaining scalar debt}

\begin{theorem}[The complete actual inverse-metric contraction]
\label{c18v58:metricthm}
For every $n\ge5$ and every fixed coarse field,
\[
 \boxed{\quad
     \mathbb E[A(dC,dC)]=-\frac{405}{8}n^3,\qquad
     \mathbb E[A(db,db)]=-\frac{3645}{8}n^3 .
       \quad}
 \tag{V58.20}\label{c18v58:metric}
\]
Consequently its contribution to the actual
quadratic temporal cost is
\[
                 -\frac1{81}\mathbb E[A(db,db)]
                                      =\frac{45}{8}n^3 .
 \tag{V58.21}\label{c18v58:metriccost}
\]
\end{theorem}

\begin{proof}
Insert \eqref{c18v58:scalars},
\eqref{c18v58:curvature} and
\eqref{c18v58:normalvalue} into
\eqref{c18v58:metricreduction}:
\[
 \begin{aligned}
 \mathbb E[A(dC,dC)]
   &=9n^3-\left(\frac{225}{4}n^3-\frac{243}{64}W_c\right)
                 +\frac18\left(-27n^3-\frac{243}{8}W_c\right)\\
   &=-\frac{405}{8}n^3 .
 \end{aligned}
\]
The two coarse-holonomy contributions cancel.
Since $db=3\,dC$, the other formulas follow.
This cancellation uses the actual normal derivative
and the integrated metric Hessian, not a symmetry
assumption that the answer must be constant.
\end{proof}

\begin{corollary}[Only the cubic-density contraction is unpaid in this coefficient]
\label{c18v58:costcor}
With V57's exact constant
$K=18+540/19+27/112+567/115$, the actual coefficient is
\[
 \begin{aligned}
 \mathcal C_2(g)
    &=\mathbb E[hRh]+\frac19\mathbb E[bk]
                                -15n^3+\frac{243}{16}W_c\\
    &=\left(K+\frac{435}{16}\right)n^3
                          +\frac{117}{4}W_c+\frac19\mathbb E[bk].
 \end{aligned}
 \tag{V58.22}\label{c18v58:remainder}
\]
This is not yet a numerical evaluation or a sign
determination of $\mathcal C_2(g)$.
\end{corollary}

\begin{proof}
The two $j$ terms contribute
$-45n^3/8+45W_c/4$ and
$-75n^3/8+9W_c/8$.
V56's fourth-cumulant term is
$-45n^3/8+45W_c/16$.
Its constant part cancels
\eqref{c18v58:metriccost}, leaving
$-15n^3+243W_c/16$.
Use V57's complete inverse energy for the
second line of \eqref{c18v58:remainder}.
\end{proof}

Thus the next finite-jet calculation has a precise
target, the actual $\mathbb E[bk]$ with
$k=u'''_0$. A local cubic-density formula or an
integration-by-parts reduction must retain the
normal chart's time dependence. Setting $k=0$,
using a frozen tilt, or transferring a companion
finite-spectrum statement would alter the coefficient.

\subsection{Verification, literature context and scope}

The exact diagnostic is
\begin{center}\small
\path{scripts/verify_ym_c18_v58_density_metric_cost.py}.
\end{center}
On the native $n=5,6$ lattices it checks all
$1{,}124{,}250$ and $3{,}357{,}936$ distinct
captured-plaquette pairs in $C^2$, respectively.
The only matches to the contributing scalar
signatures are their actual native motifs.
It also checks both reversed exterior half signs
for every coherent route pair, and that four
different captured squares incident to one pair
never have empty free signature.

The local metric calculation is checked by exact
quaternion multiplication and rational sphere
moments, including the pair connection subtraction,
the mixed pair/free Hessian, all middle and exterior
norms, and the curvature wedge.
V55's arbitrary-unit-quaternion source matrix
is replayed to check the coherent spectral input.
An independent circle calculation checks
\eqref{c18v58:hessibp} with the positive Laplacian.
That circle is explicitly a non-native sign
fixture. Neither it nor the finite incidence
enumeration replaces the all-volume proofs above.

An additional numerical regression uses 64 independent
Haar samples at each of a flat $n=5$ coarse field
and an $n=6$ field with every coarse plaquette
holonomy equal to minus the identity. It differentiates
the original pair/free plaquette variables, including
the covariant same-factor Hessians, and directly
assembles $\nabla_D^\perp N$ and the ambient curvature.
Midpoint gauge fixing is used only for evaluating
these invariant scalar integrals, not to replace
their differential operator. This is a sampling
cross-check of the exact formulas, not a proof
or a finite-time normal-trajectory simulation.

A targeted primary-source check consulted Fabrice
Baudoin's \emph{Lecture 13: The Bochner's formula}
(10 September 2013),
\url{https://fabricebaudoin.blog/2013/09/10/lecture-13-the-bochners-formula/}.
Its vector-field derivation highlights the
connection terms behind Riemannian Hessians.
Here the needed integrated identity is proved
directly in \eqref{c18v58:hessibp}; the mixed
Hessian retains the quaternion commutator term.
The source's diffusion-sign Laplacian is opposite
to this note's positive $L$. No global curvature,
finite-time Hessian or spectral-gap estimate
is imported. This is not the scheduled broad
sweep, which remains due at V60, together with
the next companion checkpoint.

Removing this append and restoring the title
recovers V57 byte for byte. The other 53 sources
are protected and unchanged. Only V58 is kept
as the active YM source, with all build and
inspection products outside \path{latex_notes}.
Archive/restart remains V100.

\clearpage
\paragraph{What has advanced.}
Both actual mixed density cumulants are evaluated,
and the complete native inverse-metric contraction
is computed with its connection, curvature and
coherent return included. A nontrivial coarse
dependence cancels in the full metric contribution.
Together with V57, this reduces the formerly
multi-term unknown coefficient to the single
actual cubic-density contraction in
\eqref{c18v58:remainder}.

\paragraph{What has not advanced to a theorem.}
No value or sign of $\mathbb E[bk]$ has been
established here. Nor has a regulator-uniform
Taylor remainder, block-time-four comparison,
or the source-specific Poisson-Hessian bound
required by V51 been obtained. Even completing
this coefficient would leave those finite-time
obligations, coarse variance and return,
weak-coupling scale transport, physical-time
calibration and interacting continuum
reconstruction unpaid.
The full physical Yang--Mills mass gap remains
unproved and the goal remains active.

\addtocontents{toc}{\protect\setcounter{tocdepth}{2}}
\addtocontents{toc}{\protect\clearpage}
% END CYCLE 18 VERSION 58 APPEND
% BEGIN CYCLE 18 VERSION 59 APPEND
\clearpage
\section{V59: the actual cubic density and the completed quadratic cost}
\label{c18v59:section}
\addtocontents{toc}{\protect\setcounter{tocdepth}{1}}

\subsection{Context: finishing the native coefficient, not the mass-gap proof}

V58 left exactly one contraction unpaid in V56's quadratic
temporal coefficient, namely $\mathbb E[bk]$ with the
\emph{actual} raw cubic density jet $k=u'''_0$. This append
evaluates it. The calculation first returns upstream to
the ambient metric to remove the second differentiated
normal projector. Covariant integration by parts then
removes all third spatial derivatives. Finally a complete
free-link parity classification reduces the remaining
quartic integral to explicitly evaluated native motifs.

The result closes this finite-jet calculation, not the
Yang--Mills programme. We still work at $\beta=0$ on a
finite spatial torus of side $2n$, $n\ge5$, with group
$\mathrm{SU}(2)$, arbitrary fixed coarse field $g$, and
the original block-two sampler and normal chart.
Expectations are under the conditional product Haar law
$\mu_g$ on the original fibre $M=M_g$.
The original physical continuum mass-gap goal is unchanged.

Retain $C=W_{\rm cap}$, $U=W_{\rm uncap}$, $W=C+U$,
the positive electric $L$, $LC=9C$, $LW=12W-3C$,
and $b=3C$. Set
\[
 \begin{gathered}
 F=\nabla_{\rm amb}W,\quad H=\nabla_{\rm amb}F,\quad
 P=P_0,\quad Q=I-P,\\
 T=PF=\nabla_MW,\quad D=\nabla_MC,\quad
 N=QF,\quad S=|N|^2,\quad Z=\Gamma(W,C),\quad G=\Gamma(C),\\
 \mathcal B=C\,LW-Z,\qquad
 \mathcal H=\langle\operatorname{Hess}_MW,
                         \operatorname{Hess}_MC\rangle
                    +\operatorname{Ric}_M(T,D).
 \end{gathered}
 \tag{V59.1}\label{c18v59:notation}
\]
Here $H$ is an ambient endomorphism; $\mathcal H$ is a scalar.
The symbol $D$ is the source gradient, not $T$.
Use the curvature convention
$\mathrm{Rm}(X,Y)Z=\nabla_X\nabla_YZ-\nabla_Y\nabla_XZ
-\nabla_{[X,Y]}Z$, for which a unit sphere has
$\mathrm{Rm}(X,Y)Z=\langle Y,Z\rangle X-\langle X,Z\rangle Y$.
All identities below also hold with real plaquette weights
in $W$ and the corresponding captured sub-sum $C$.
This permits exact polarization, not a change of sampler.

\subsection{The metric equation for the actual tangential velocity}

Put $\theta_t=\Phi_t\eta_t$ as in V57 and
$Y_t=\dot\theta_t\theta_t^{-1}$ on $M$.
Let $\gamma_t=\Phi_{-t}^*g_{\rm amb}$, a metric on the
ambient space. Superscripts on $Y$ below denote fixed-point
Eulerian time derivatives, not material accelerations.

\begin{proposition}[The second velocity jet without $Q_2$]
\label{c18v59:velocity}
The actual velocity is determined by
\[
 Y_t\in TM,\qquad \gamma_t(Y_t-F,v)=0\quad(v\in TM).
 \tag{V59.2}\label{c18v59:normality}
\]
At the origin,
\[
 \begin{aligned}
 \gamma'_0&=-2H,\qquad
 \gamma''_0=2\nabla_FH+4H^2,\\
 Y_0&=T,\qquad Y^{(1)}=-2PHN=-\nabla_MS,\\
 Y^{(2)}&=-8PHPHN+2P(\nabla_FH)N+4PH^2N .
 \end{aligned}
 \tag{V59.3}\label{c18v59:Y2}
\]
In the first line the endomorphisms are lowered by the
ambient metric to give symmetric two-tensors.
\end{proposition}

\begin{proof}
Since $\eta_t=\Phi_{-t}\theta_t$, at a point of the
original fibre its velocity is
$D\Phi_{-t}(Y_t-F)$.
The moving fibre tangent is $D\Phi_{-t}(TM)$.
The defining normality of the chart velocity gives
\eqref{c18v59:normality}. Conversely the restriction of
$\gamma_t$ to $TM$ is positive definite, so that equation
determines its tangential solution uniquely. Thus we have
not substituted a freely chosen tangential transport.

Differentiation of a pullback by the autonomous flow gives
$\gamma'_0=-\mathcal L_Fg_{\rm amb}=-2H$.
For a covariant Hessian,
$\mathcal L_FH=\nabla_FH+2H^2$, proving the second derivative.
Differentiate \eqref{c18v59:normality} once and twice,
using $Y_0-F=-N$. Its first derivative is
$\langle Y^{(1)},v\rangle+2H(N,v)=0$.
The second is
\[
 \langle Y^{(2)},v\rangle
 -4H(Y^{(1)},v)
 -(2\nabla_FH+4H^2)(N,v)=0.
\]
Substitution proves \eqref{c18v59:Y2}.
The identity $\nabla_MS=2PHN$ was proved in V57 from
total geodesy, and remains valid with weighted plaquettes.
\end{proof}

\begin{proposition}[The actual raw cubic density]
\label{c18v59:rawk}
On the original fibre,
\[
 \begin{aligned}
 k={}&24H(N,N)-72H(T,N)-6H(N,D)\\
    &+3\operatorname{Hess}_{\rm amb}C(T,T)+3H(D,T)\\
    &+\operatorname{div}_M Y^{(2)}+2T(LS).
 \end{aligned}
 \tag{V59.4}\label{c18v59:k}
\]
There is no derivative of an assumed Gibbs tilt in this formula.
\end{proposition}

\begin{proof}
V57's unnormalized Jacobian identity is
$u_t=\alpha_t\circ\theta_t+\log\operatorname{Jac}_M\theta_t$,
where $\alpha_t=\log\operatorname{Jac}_{\rm amb}\Phi_{-t}$.
Since $\operatorname{div}_{\rm amb}F=-12W$,
\[
 \alpha'_0=12W,\quad \alpha''_0=-12|F|^2,\quad
 \alpha'''_0=24H(F,F).
\]
The third derivative of the composed term is
$\alpha'''_0+3T\alpha''_0+
3(Y^{(1)}+T^2)\alpha'_0$.
The fibre Jacobian contributes
\[
 \operatorname{div}_MY^{(2)}
 +2T\operatorname{div}_MY^{(1)}
 +Y^{(1)}\operatorname{div}_MT
 +T^2\operatorname{div}_MT.
\]
Use $\operatorname{div}_MT=-LW$, $\operatorname{div}_MY^{(1)}=LS$,
$LW=12W-3C$, $\nabla_T^MT=PHT$, and
$Y^{(1)}=-2PHN$. The $H(T,T)$ terms cancel;
the remaining terms are exactly \eqref{c18v59:k}.
In particular the two occurrences of the material
derivative in $2T\operatorname{div}_MY^{(1)}$ are both retained.
\end{proof}

\subsection{A quartic formula containing only first and second derivatives}

Define
\[
 \begin{aligned}
 \mathfrak F(W,C)={}&6C H(N,N)
 +S\{-36CW-72C^2-9Z+12G+\mathcal H\}\\
 &+3Z\mathcal B
 +2\langle QHD,QHT\rangle
 -2\langle\mathrm{Rm}(F,D)F,N\rangle .
 \end{aligned}
 \tag{V59.5}\label{c18v59:F}
\]
Every term is homogeneous of degree four in the plaquette
weights. Only derivatives of order at most two occur.

\begin{theorem}[Integrated cubic reduction]
\label{c18v59:reduction}
For every native weighted potential described above,
\[
                 \mathbb E[Ck]=\mathbb E\mathfrak F(W,C).
 \tag{V59.6}\label{c18v59:integral}
\]
For the native pair/free geometry the last curvature
contraction equals
$\langle\mathrm{Rm}(D,N)N,D\rangle$ pointwise.
\end{theorem}

\begin{proof}
First integrate the divergence in \eqref{c18v59:k}:
$\mathbb E[C\operatorname{div}Y^{(2)}]=-\mathbb E\langle D,Y^{(2)}\rangle$.
Commuting the covariant Hessian gives
\[
 (\nabla_FH)(N,D)
   =(\nabla_DH)(N,F)+\langle\mathrm{Rm}(F,D)F,N\rangle.
 \tag{V59.7}\label{c18v59:commutation}
\]
Along $M$, $\nabla_DN=QHD$ and $\nabla_DF=HD$.
Consequently, if $r=\langle\mathrm{Rm}(F,D)F,N\rangle$,
\[
 (\nabla_FH)(N,D)
 =D[H(N,F)]-\langle QHD,HF\rangle-\langle HN,HD\rangle+r.
\]
Use $\mathbb E[Df]=9\mathbb E[Cf]$ and \eqref{c18v59:Y2}.
After separating the tangent and normal parts of $HD$,
\[
 -\mathbb E\langle D,Y^{(2)}\rangle
 =\mathbb E\{6\langle PHD,PHN\rangle
       +2\langle QHD,QHT\rangle
       -18C H(N,F)-2r\}.
 \tag{V59.8}\label{c18v59:Y2integral}
\]
This is where the third derivatives disappear.

Here are the remaining scalar integrations, displayed
to fix the signs of the positive Laplacian:
\[
 \begin{aligned}
 \mathbb E[C\,T f]&=\mathbb E[\mathcal B f],&
 \mathbb E[C\,D f]&=\mathbb E[(9C^2-G)f],\\
 2\mathbb E[C\,T(LS)]&=2\mathbb E[S L\mathcal B],&
 T Z&=\operatorname{Hess}_MC(T,T)+H(D,T),\\
 L\mathcal B&=C(252W-90C)-45Z+9G+2\mathcal H,\\
 6\mathbb E\langle PHD,PHN\rangle
   &=3\mathbb E[S(12Z-3G-\mathcal H)].
 \end{aligned}
 \tag{V59.9}\label{c18v59:ibp}
\]
For the last line, $2PHN=\nabla S$ and
\[
 \operatorname{div}_M[(\operatorname{Hess}_MW)D]
       =-\Gamma(C,LW)+\mathcal H.
\]
The line for $L\mathcal B$ follows from the product rule
$L(fh)=fLh+hLf-2\Gamma(f,h)$ and the polarized Bochner
identity
$LZ=\Gamma(LW,C)+\Gamma(W,LC)-2\mathcal H$.
These identities can also be obtained at a normal frame
by commuting the two covariant derivatives of a gradient;
their curvature sign is the convention stated above.

Substitute \eqref{c18v59:Y2integral} in \eqref{c18v59:k}.
Use $H(T,N)=TS/2$ and $H(D,N)=DS/2$.
The terms multiplied by $S$ become
$-36CW-72C^2-9Z+12G+\mathcal H$; the two terms with
$\operatorname{Hess}C(T,T)$ and $H(D,T)$ become
$3Z\mathcal B$. This proves \eqref{c18v59:integral}.

For the curvature assertion, uncaptured plaquettes have
no pair gradient, so $T_i=D_i$ on every pair factor.
In orthonormal tangent/normal copies write these as $d,n$.
The ambient two-sphere coordinates are
\[
 F_i=((d+n)/\sqrt2,(-d+n)/\sqrt2),\quad
 D_i=(d/\sqrt2,-d/\sqrt2),\quad
 N_i=(n/\sqrt2,n/\sqrt2).
\]
The unit-sphere curvature formula gives in both
contractions
$\tfrac12(|d|^2|n|^2-\langle d,n\rangle^2)$.
Free factors have $N_i=0$, so the identity holds after summation.
\end{proof}

\subsection{The complete quartic parity classification}

For independent weights $z_p$, put
$W_z=\sum_p z_pw_p$ and $C_z=\sum_{p\ {\rm cap}}z_pw_p$.
Write $J(z)=\mathbb E[C_zk_z]$.
This is a homogeneous quartic polynomial.
A free link central sign changes every differentiated
plaquette containing that link by minus one.
The covariant derivatives used above preserve that sign.

\begin{lemma}[No unlisted quartic support survives]
\label{c18v59:parity}
The only possibly nonzero monomials of $J$ are:
a fourth power of one captured weight;
$z_p^2z_q^2$ for an overlapping captured--captured or
captured--uncaptured pair; and the product of the four
distinct captured corner weights of one coarse plaquette.
Here overlap means sharing an original pair/free factor.
\end{lemma}

\begin{proof}
Different fine plaquettes have different sets of free links.
Therefore parity kills $z_p^3z_q$ and $z_p^2z_qz_r$
for distinct indicated plaquettes. A surviving fourth
power must be captured, since $C_z$ vanishes when only
uncaptured weights are present.

Consider four distinct plaquettes with at least one captured.
Every uncaptured plaquette contains exactly two edges incident
to one three-odd vertex. At that vertex its two incident
directions are not opposite. These pairs of incident edges
are distinct for different uncaptured plaquettes.
No captured plaquette contains a three-odd edge.
One uncaptured plaquette cannot have even parity there;
two distinct ones cannot have equal two-edge signatures.
For three uncaptured plaquettes, even parity requires the
three two-edge sets to form a triangle at the same three-odd
vertex. This is one of the eight choices of an octant.
Their other six free edges are distinct and remain odd.
A single captured plaquette has only two free edges, and
cannot cancel them. Thus no four-distinct surviving
monomial involving an uncaptured plaquette exists.

The two free edges of a captured plaquette meet at its
unique two-odd centre. At each such centre the four
captured corner plaquettes form the four-cycle on its
four free spokes. No proper nonempty subset of that cycle
has even incidence. Free edges belong to only one such
centre. Four distinct captured plaquettes with even free
incidence must therefore be precisely this entire cycle.

It remains to exclude a repeated pair with disjoint
pair/free support. For the potential $z_pw_p+z_qw_q$
on such supports, the ambient flows and normal-chart
equations separate, including all captured half pairs.
Thus $k_z=z_p^3k_p+z_q^3k_q$; multiplying by $C_z$
has no $z_p^2z_q^2$ term. This argument uses the actual
chart, not cancellation of only a selected summand
of \eqref{c18v59:F}. It also covers separated supports
whose unintegrated expression in \eqref{c18v59:F}
has apparent disconnected terms.
\end{proof}

The overlapping classes and counts were identified in V57:
$12n^3$ adjacent captured pairs,
$36n^3$ different-centre same-half pairs,
$36n^3$ different-centre opposite-half pairs, and
$48n^3$ captured--uncaptured pairs sharing a free edge.
There are $12n^3$ self motifs and $3n^3$ four-corner cycles.
The $n\ge5$ restriction excludes short periodic
identifications in these local support arguments.

\subsection{Exact finite Haar evaluation of every surviving motif}

We give both the local differential recipe and its finite
rational integration rule. The table below is a
computer-assisted finite arithmetic calculation with
an explicit reproducing recipe, not a numerical fit or
a proposed two-parameter ansatz.

Use unnormalized quaternion frames with flow
$q\mapsto q\exp(se_a)$ on each pair/free factor.
Let $\omega_i=1/2$ on pairs and
$\omega_i=1$ on free factors; these are inverse metric
weights. Write $g_{p,i}$ for the three gradient covector
components and $h_{p,ij}$ for the $3$ by $3$ covariant
Hessian block of $w_p$ in these frames.
The pair half sign is V57's $\sigma_{ip}$.
The differential data determine all entries of
\eqref{c18v59:F} by
\[
 \begin{gathered}
 D_i=\omega_i\sum_{p\ {\rm cap}}z_pg_{p,i},\quad
 T_i=\omega_i\sum_pz_pg_{p,i},\quad
 N_i=\omega_i\sum_pz_p\sigma_{ip}g_{p,i},\\
 S=\sum_{i\ {\rm pair}}2|N_i|^2,\quad
 Z=\sum_i\left(\sum_{p\ {\rm cap}}z_pg_{p,i}\right)\cdot T_i,
 \quad G=\sum_i\omega_i^{-1}|D_i|^2,\\
 H(N,N)=\sum_{p,i,j}z_p\sigma_{ip}\sigma_{jp}
                         N_i^{\mathsf T}h_{p,ij}N_j,\\
 (J_D)_i=\sum_{p,j}z_p\sigma_{ip}h_{p,ij}D_j,\qquad
 (J_T)_i=\sum_{p,j}z_p\sigma_{ip}h_{p,ij}T_j,\\
 \langle QHD,QHT\rangle=\sum_i\omega_i(J_D)_i\cdot(J_T)_i,\\
 \mathcal H=\sum_{i,j}\omega_i\omega_j
       \langle h_{W,ij},h_{C,ij}\rangle_{\rm F}
                  +2\sum_iT_i\cdot D_i,\\
 \langle\mathrm{Rm}(F,D)F,N\rangle
   =2\sum_{i\ {\rm pair}}
       (|D_i|^2|N_i|^2-(D_i\cdot N_i)^2).
 \end{gathered}
 \tag{V59.10}\label{c18v59:localrecipe}
\]
Zero half signs restrict the normal terms to pairs.
The Frobenius pairing is indicated by the subscript.
The Ricci tensor of each radius-scaled three-sphere
has components $2I$ in the unnormalized frame, explaining
the factor in $\mathcal H$.
The normal-copy connection is the parallel copy proved in V58.
For $i\ne j$ the normal--normal Hessian of a plaquette is
$\sigma_{ip}\sigma_{jp}h_{p,ij}$; on the same pair it is
$-w_pI$, since a captured plaquette uses exactly one half.
Thus the formula for $H(N,N)$ includes its diagonal blocks.

For clarity the data $g_p,h_p$ require no numerical
differentiation. Set all pair prefixes to the identity
using internal midpoint gauge transformations. This is
only an integration gauge choice, not a change of metric.
A positive letter $q$ has derivative $q e_a$; an inverse
letter $q^{-1}$ has derivative $-e_aq^{-1}$.
For a second captured half $x^{-1}g$, use these same
insertions before setting $x=1$.
Different-factor Hessian blocks are obtained by two
insertions in the ordered four-letter Wilson word.
Every same-factor covariant block is
\[
                         h_{p,ii}=-w_pI_3 .
 \tag{V59.11}\label{c18v59:diaghess}
\]
Indeed a degree-one spherical coordinate has this
covariant Hessian; the antisymmetric connection term
in two invariant derivatives has been subtracted.
All metric factors are in $\omega_i$, not dropped
by treating the frames as orthonormal.

\begin{lemma}[The exact small cubature rules]
\label{c18v59:cubature}
For centrally even polynomials on the unit quaternion
sphere of degree at most two, normalized Haar integration
equals the equally weighted average on
\[
                   \mathcal A=\{1,\mathbf i,\mathbf j,\mathbf k\}.
\]
For centrally even polynomials of degree at most four it
equals the equally weighted average on
\[
 \mathcal Q=\mathcal A\ \cup\
 \{(1+\epsilon_1\mathbf i+\epsilon_2\mathbf j+
                  \epsilon_3\mathbf k)/2:\epsilon_a=\pm1\}.
 \tag{V59.12}\label{c18v59:rule}
\]
Neither rule is claimed for odd polynomials.
\end{lemma}

\begin{proof}
The second moments are $\delta_{ab}/4$ for both rules.
On $\mathcal Q$, each fourth coordinate moment is
$(1+8/16)/12=1/8$, and each mixed squared-coordinate
moment is $(8/16)/12=1/24$. All other even-total-degree
monomials of degree at most four with an odd coordinate
exponent average to zero by the displayed signs.
These are exactly the sphere's second and fourth moments.
Constants agree as well; linearity proves the assertions.
\end{proof}

Apply the rules \emph{after} extracting the desired
plaquette coefficient. A repeated pair has degree two
on each private free factor and degree four on a shared
free factor. A self motif has degree four on each of
its two free factors. A four-corner cycle has degree
two on each of its four free factors.
Thus the quadrature sizes are respectively
$256$ for either different-centre captured pair,
$192$ for an adjacent pair, $3072$ for a
captured--uncaptured pair, $144$ for a self motif,
and $256$ for the four-corner cycle.

Define the six summands of $\mathfrak F$ in the order
\[
 \begin{aligned}
 F_1&=6C H(N,N),&F_2&=S(-36CW-72C^2),\\
 F_3&=S(-9Z+12G),&F_4&=S\mathcal H,\\
 F_5&=3Z\mathcal B,&
 F_6&=2\langle QHD,QHT\rangle
                  -2\langle\mathrm{Rm}(F,D)F,N\rangle .
 \end{aligned}
 \tag{V59.13}\label{c18v59:six}
\]
For a two-plaquette support the exact coefficient extraction
on any homogeneous quartic function $f(z_1,z_2)$ is
\[
 [z_1^2z_2^2]f
       =\frac{f(1,1)+f(1,-1)}2-f(1,0)-f(0,1).
\]
For the four-corner support it is
\[
 [z_1z_2z_3z_4]f
       =2^{-4}\sum_{\epsilon\in\{\pm1\}^4}
                       \epsilon_1\epsilon_2\epsilon_3\epsilon_4 f(\epsilon).
 \tag{V59.14}\label{c18v59:polarization}
\]
A single self coefficient is $f(1)$.
These are identities for each $F_a$, not approximate
finite differences.

The exact resulting coefficients at trivial coarse
holonomy are:
\[
\renewcommand{\arraystretch}{1.35}
\setlength{\arraycolsep}{3pt}
\begin{array}{c|rrrrrr|r}
 &F_1&F_2&F_3&F_4&F_5&F_6&\text{sum}\\ \hline
 \text{self}&-3/4&-27/2&45/8&117/8&-27/4&9/4&3/2\\
 \text{same half}&-9/16&-81/2&171/16&489/16&-9/16&3/8&0\\
 \text{opposite half}&-9/16&-81/2&153/16&483/16&-9/16&0&-15/8\\
 \text{adjacent}&-9/16&-81/2&135/16&459/16&-81/16&9/2&-9/2\\
\text{C--U}&0&0&0&0&-9/16&0&-9/16\\
 \text{four corners}&27/16&81/2&-81/16&-243/32&243/16&-243/32&297/8
\end{array}
 \tag{V59.15}\label{c18v59:table}
\]
For reproducing the finite sums, let $p(x;i,j)$ mean the
standard positively oriented fine plaquette at $x$.
Use $p((0,0,0);0,1)$ for the first plaquette.
The respective second plaquettes for the same-half,
opposite-half, adjacent, and captured--uncaptured rows are
\[
 p((0,0,0);0,2),\quad p((0,1,0);1,2),\quad
 p((0,1,0);0,1),\quad p((0,1,0);0,2).
\]
The four-corner row uses $p((a,b,0);0,1)$,
$a,b\in\{0,1\}$.
These coordinates, the four-letter word rule,
\eqref{c18v59:localrecipe}, \eqref{c18v59:rule},
and \eqref{c18v59:polarization} completely specify
every rational entry of the table.

Two rows have short independent hand evaluations.
For a self motif $W=C=w$, put $r=1-w^2$. Then
\[
 \begin{gathered}
 S=r,\quad Z=G=3r,\quad H(N,N)=-wr,\quad
 \mathcal H=18+9w^2,\\
 |QHD|^2=9w^2r,\quad \langle\mathrm{Rm}(F,D)F,N\rangle=0.
 \end{gathered}
\]
Substitution gives $\mathfrak F=12w^2(1-w^2)$,
whose mean is $3/2$.
For a captured--uncaptured pair, only $F_5$ has the
required two-plus-two weight. Its coefficient is
$3\Gamma(p,q)(12pq-\Gamma(p,q))$.
On their common free factor write $pq=f_0+f_2$.
Then $\Gamma(p,q)=3f_0-f_2$ and the orthogonal
squared norms are $1/64,3/64$.
Thus $\mathbb E[pq\,\Gamma(p,q)]=0$ and
$\mathbb E\Gamma(p,q)^2=3/16$, giving $-9/16$.

The other rows are exact finite sums by the displayed
recipe. The accompanying verifier uses integer quaternion
numerators with each input coordinate represented over
denominator two. Each $w_p,g_p,h_p$ then has denominator
$16$; all six scalar summands have common denominator
$2^{19}$. Polarization and the finite average are performed
with integers and exact rational division, not floating
point. There are at most four plaquettes and sixteen
independent factors. Bounds $|w_p|,|g_{p,i}^a|,
|h_{p,ij}^{ab}|\le1$ and $\omega_i\le1$ bound every
scaled intermediate, including the sixteen-sign sum,
strictly below $2^{60}$; signed 64-bit arithmetic cannot
overflow here. The largest numerator actually found
in the polarized six-term tables is $1811939328$.
The elementary moment rule and these finite sums are
the arithmetic part of the proof, rather than empirical
checks of a conjectured formula.

\begin{lemma}[Coarse dependence of the motif coefficients]
\label{c18v59:coarse}
All self and repeated-pair entries in \eqref{c18v59:table}
are independent of $g$.
A four-corner coefficient is its displayed value times
$\operatorname{Re}H_y$, where $H_y$ is that coarse plaquette.
\end{lemma}

\begin{proof}
The coarse links present in a self or overlapping
two-plaquette support form a tree: at most three
distinct coarse links occur and there is no short
periodic identification. Coarse-vertex gauge changes
set their values to the identity; these are product
isometries preserving the Haar integral and every
scalar contraction used above. The remaining word
incidences are the shared-factor types specified in
the table. Relabeling factors, reversing a quaternion
word by an isometry, and translating its private Haar
variables leave the value unchanged. The only retained
pair distinction is its half-sign product; an adjacent
pair additionally shares its common free link.
This gives precisely the three captured-pair types
and the single captured--uncaptured type already listed.

For four corners, fix three coarse links to the identity
and leave the fourth equal to $H_y$.
In the coefficient with each corner weight used once,
that remaining link occurs in exactly one of the four
plaquette words, once and in the fundamental representation.
The derivative insertions preserve this real linear
dependence on the four coordinates of $H_y$.
The integrated scalar is invariant under conjugation
of $H_y$, so its linear dependence is a multiple of
$\operatorname{Re}H_y$. Evaluation at $H_y=1$
determines the multiplier in \eqref{c18v59:table}.
This proves the dependence; it was not inferred by
fitting the two central holonomies.
\end{proof}

\subsection{The full cubic contraction and a sign-changing temporal coefficient}

\begin{theorem}[The previously unpaid native contraction]
\label{c18v59:answer}
With $\mathcal W_{\rm c}(g)=\sum_y\operatorname{Re}H_y$,
\[
 \boxed{\begin{aligned}
 \mathbb E[Ck]&=-\frac{261}{2}n^3+\frac{297}{8}\mathcal W_{\rm c}(g),\\
 \mathbb E[bk]&=-\frac{783}{2}n^3+\frac{891}{8}\mathcal W_{\rm c}(g).
 \end{aligned}}
 \tag{V59.16}\label{c18v59:bk}
\]
These are raw normal-chart density contractions.
\end{theorem}

\begin{proof}
The support classification is complete, so sum the table:
\[
 \begin{aligned}
 \mathbb E[Ck]
 &=12n^3(3/2)+36n^3(0)+36n^3(-15/8)\\
 &\quad+12n^3(-9/2)+48n^3(-9/16)
                         +(297/8)\mathcal W_{\rm c}(g).
 \end{aligned}
\]
This simplifies to the first expression. The second
uses $b=3C$. No normalization derivative has been
inserted into $k$.
\end{proof}

\begin{corollary}[Completed quadratic coefficient, with fixed-regulator scope]
\label{c18v59:cost}
The actual cost of V55--V58 satisfies
\[
 \begin{aligned}
 \mathcal C(t,g)&=3n^3+t^2\mathcal C_2(g)+O_{n,g}(t^4),\\
 \mathcal C_2(g)&=\frac{1079217}{30590}\,n^3
                           +\frac{333}{8}\mathcal W_{\rm c}(g).
 \end{aligned}
 \tag{V59.17}\label{c18v59:C2}
\]
The error estimate here is not uniform in volume or
coarse field.
\end{corollary}

\begin{proof}
V58 gives
$\mathcal C_2=(K+435/16)n^3+(117/4)\mathcal W_{\rm c}
+\mathbb E[bk]/9$, with $K=12625731/244720$.
Substitute \eqref{c18v59:bk}. The volume coefficient
is $K-261/16=1079217/30590$, and the holonomy
coefficient is $117/4+99/8=333/8$.
V55's exact time-reversal symmetry and fixed-regulator
smoothness give the fourth-order remainder.
\end{proof}

At the flat coarse field and V55's legitimate all-negative
coarse-plaquette field, respectively, the coefficient is
\[
 \mathcal C_2(g_{\rm flat})=\frac{19596573}{122360}n^3>0,\qquad
 \mathcal C_2(g_{\rm minus})=-\frac{10962837}{122360}n^3<0.
 \tag{V59.18}\label{c18v59:sign}
\]
The second field exists for every even $n\ge6$.
Consequently neither an everywhere increasing nor an
everywhere decreasing small-positive-time cost can hold
uniformly over all coarse fields. For each fixed such
field the strict sign holds for sufficiently small
positive time by \eqref{c18v59:C2}; no common time radius
is claimed. A negative quadratic coefficient is fully
compatible with the nonnegative cost near its positive
initial value $3n^3$.

Under the \emph{reference coarse Haar law},
$\mathbb E_g\mathcal W_{\rm c}=0$, so the mean coefficient
is $1079217n^3/30590$.
This is the mean of this coefficient at the reference
law, not a derivative after silently holding fixed a
different evolving coarse marginal.
For every coarse value the exact coefficient bound is
\[
 -\frac{10962837}{122360}n^3\le\mathcal C_2(g)
                  \le\frac{19596573}{122360}n^3.
 \tag{V59.19}\label{c18v59:bounds}
\]
The upper endpoint is attained for all $n$, and the
lower endpoint for the even $n$ just described.

\subsection{Audit, direction adjustment and the unpaid continuum obligations}

The geometric verifier independently differentiates
the actual sampler projector twice on a flat
three-torus with metric $\operatorname{diag}(1,1/3,1/8)$,
potential $\cos(3x)\cos y+\cos(2x)\cos z$, and fibre
$y$ fixed. Exact Laurent-polynomial algebra compares
V56's raw $k$ against \eqref{c18v59:k}, and its integral
against \eqref{c18v59:integral}, retaining the arbitrary
coarse coordinate. This is a non-native sign fixture,
not a simulated YM trajectory. Separate sphere and
native pair-frame calculations check the curvature
commutation sign.

The integer motif verifier checks the Haar rules,
the three-odd and two-odd incidence arguments at $n=5,6$,
and all $82$ shared-slot/half-sign types present at
$n=5$: $30$ same-half, $23$ opposite-half, $4$ adjacent,
and $25$ captured--uncaptured types. Their exact cubatures
agree with the four pair rows.
It also evaluates the four-corner coefficient at
$H_y=-1,\mathbf i,\mathbf j,\mathbf k$:
the first reverses every table entry and the last
three give zero. These are independent arithmetic
checks of the proved linear conjugation-invariant
dependence, not its justification.

The reproducing scripts are
\begin{center}\small
\path{scripts/verify_ym_c18_v59_cubic_density.py}\\
\path{scripts/verify_ym_c18_v59_quartic_motifs.py}.
\end{center}
The written geometric and support proofs extend to
every stated $n,g$; the finite-lattice checks do not
replace that quantifier. The table is supported by
the explicit exact finite computation above, not
by proof-assistant certification of the historical corpus.

The targeted primary recheck of
\href{https://fabricebaudoin.blog/2013/09/10/lecture-13-the-bochners-formula/}
{Baudoin's \emph{Lecture 13: The Bochner's formula}}
was used only to audit connection and Ricci terms;
its diffusion-sign operator has the opposite sign
to our positive $L$. The needed integrations and
commutation signs are supplied in this append.
No global curvature-to-gap conclusion is imported.
The next broad literature sweep and companion
checkpoint are both due at V60, not replaced by this
targeted check.

The direction now changes from evaluating zero-time
jets to controlling the actual nonzero-time source
and Poisson Hessian. The full first-corrector inverse,
all mixed density and metric contractions, and the
cubic contraction are paid. Recomputing them or using
the frozen first-band resolvent is not a substitute
for the remaining estimate.
The sign change in \eqref{c18v59:sign} specifically
excludes a monotonicity shortcut for the actual cost.

Still unpaid are a regulator-uniform finite-time
remainder or nonlinear comparison through block time
four, the source-specific Poisson-Hessian control,
coarse variance and return, variable-bank stability,
weak-coupling scale transport with physical-time
calibration, and interacting continuum reconstruction.
None follows from the exact electric finite-jet
coefficient just obtained. The full physical
Yang--Mills mass gap remains unproved; the goal is active.

Removing this append and restoring the title recovers
V58 byte for byte. Companion work was left untouched.
During this append the separate perturbative ESM source
was reorganized externally into general and phenomenological
files; that change is recorded for the V60 checkpoint,
not overwritten or attributed to this YM calculation.
Only V59 remains the active filename; builds and
inspection products are outside \path{latex_notes}.
Archive/restart remains due at V100.

\addtocontents{toc}{\protect\setcounter{tocdepth}{2}}
\addtocontents{toc}{\protect\clearpage}
% END CYCLE 18 VERSION 59 APPEND
% BEGIN CYCLE 18 VERSION 60 APPEND
\clearpage
\addtocontents{toc}{\protect\setcounter{tocdepth}{1}}
\section[V60: finite-time source locality and the Poisson-Hessian obstruction]{V60: finite-time source locality and the Poisson-Hessian obstruction}
\label{c18v60:context}

V59 completes the quadratic temporal-cost coefficient. Its sign
depends on the coarse field, so it does not provide a monotonicity
argument. We now leave that Taylor calculation and pay a different
input: a spatial Hessian row bound for the \emph{actual temporal
source}, throughout the entire block interval $0\le t\le4$.
The bound is uniform in volume and coarse value, but very large.
It is proved at $\beta=0$, not on a continuum weak-coupling path.

The accompanying analytic results identify what this input does and
does not buy. The Haar Poisson inverse preserves a Hessian row norm
with constant $4/3$. Conversely, a gauge-invariant source on the
native fibre can have uniformly bounded single-factor gradients,
bounded energy per volume, and even bounded square-summed Hessian
rows of its solution, while the solution's Hessian operator norm
diverges. That example is not the Wilson temporal source.
It excludes an inference from weaker norms, not the native programme.

We use V42's anchored tensor algebra, V44's full Gram inverse, and
V51's exact normal chart. V38 already supplies a small-time
conditional gap; another proof of that fact is not the advance here.
The new finite-time estimate concerns the differentiated temporal
divergence. The remaining nonlinear, source-specific Poisson
estimate is kept separate throughout.

\subsection{The V60 companion and primary-literature checkpoint}

The supplied ESM V64, SM-source V52, stationarity V54 and Einstein
bridge V45 are hash-identical to those reviewed at V55. Their terminal
scope remains finite-order or fixed-cutoff, with actual source or
continuum acquisition still open. We also rechecked the current
bridge V55, stationarity V61, SM-source V65, SM-continuum V72,
NS-stability V26, shell-mixing V16 and YM-parallel V5.
This is a targeted transfer audit, not an independent certification
of every proof in these large sources.

Three externally added or reorganized files were examined more
closely. ESM-general V44 has a fixed-order, fixed-panel soft-forest
analysis, not an all-orders finite-coupling estimate. ESM-pheno V40
warns that a finite-volume coboundary identity need not survive the
thermodynamic limit. Arithmetic-loading V41 still requires the
actual matching-core channel occurrence. These support keeping the
literal source, the represented operator and the order of limits
together. None supplies the conditional YM Poisson Hessian.

The broad primary-source sweep covered the following distinct routes.
\begin{enumerate}
\item
\href{https://arxiv.org/html/2001.09910}{Thompson,
\emph{Approximation of Riemannian measures by Stein's method}},
Proposition 6.2 and Theorem 12.2, bounds Poisson/semigroup Hessians
using a global source-gradient norm and curvature-derivative inputs.
A bounded gradient on each factor is not that norm. The example
below makes this distinction quantitative on our fibre.
No convention for the drift or the factor $1/2$ in that paper is
silently substituted for our positive $L$.
\item
\href{https://arxiv.org/abs/2110.05131}{Chen--Cheng--Thalmaier,
\emph{Bismut--Stroock Hessian formulas and local Hessian estimates}}
was screened at primary abstract level. Localization there is
geometric localization; an estimate in our spatial row norm requires
its own argument. It is a potential method, not an imported bound.
\item
\href{https://arxiv.org/html/1309.0862}{Menz,
\emph{The approach of Otto--Reznikoff revisited}}, Theorem 1.7,
requires uniform correlation decay over conditioned finite
subsystems, along with its structural interaction hypotheses.
Those hypotheses are not established here for the moving YM law.
The paper's real-spin model is not identified with a compact
gauge fibre.
\item
\href{https://arxiv.org/html/2503.24372}{Bauerschmidt--Dagallier,
\emph{A criterion on the free energy for log-Sobolev inequalities
in mean-field particle systems}}, includes free-energy and
dense-graph criteria. Theorem 1.5, in particular, has a
growing-degree hypothesis. The fixed-degree cubic lattice and its
conditioned Wilson flow do not satisfy that hypothesis merely
because their volume grows.
\item
\href{https://arxiv.org/html/0711.3621}{van Enter--Ruszel,
\emph{Gibbsianness versus Non-Gibbsianness of time-evolved planar
rotor models}}, Theorem 4.1, demonstrates a conditioning-induced
obstruction for a low-temperature planar-rotor model under
independent Brownian evolution. It warns against inferring
conditional mixing from a well-defined local evolution.
It is not a counterexample to our deterministic Wilson flow.
\item
\href{https://arxiv.org/abs/2204.12737}{Shen--Zhu--Zhu,
\emph{A stochastic analysis approach to lattice Yang--Mills at
strong coupling}}, was rechecked at primary abstract level.
Its lattice Gibbs functional inequalities are not an estimate for
the native Hamiltonian-vacuum conditional law and do not by
themselves perform the physical continuum limit.
\end{enumerate}
The sweep changes the target norm rather than asserting a literature
shortcut. The proofs below are supplied directly. General heat-kernel,
tensorization and differentiation methods are not claimed as new
mathematics. The next companion checkpoint is V65 and the next
broad sweep V70.

\subsection{One more spatial derivative of the actual normal field}

Keep $T=4$, $\ell=2$, and all constants $H,t_1,t_2,b_0,b_1,b_2$,
$k,l,c_0,c_1,c_2,z_0,z_1,z_2,r_0,r_1,r_2$ from V44.
They concern the full sampler, not a selected transverse bank.
The weight $\zeta>0$ is the one in \eqref{c18v44:Zconstants}.
Ordered coefficient derivatives are in the original invariant
frames. In particular, they are not symmetrized covariant jets.

Define the additional positive majorants
\[
\begin{aligned}
 t_3&=248832e,& b_3&=648e\ell^4,\\
 m&=TH\bigl(4t_1Hl+3t_1k^2+6t_2H^2k+t_3H^4\bigr),\\
 c_3&=b_3H^4+6b_2H^2k+3b_1k^2+4b_1Hl+b_0m,\\
 g_3&=2c_0c_3+6c_1c_2,\\
 z_3&=z_0^2g_3+6z_0^3g_1g_2+6z_0^4g_1^3,\\
 r_3&=c_3z_0+3c_2z_1+3c_1z_2+c_0z_3 .
\end{aligned}
\tag{V60.1}\label{c18v60:thirdconstants}
\]
Here $m$ is a scalar majorant, not a metric. For $0\le j\le3$
put
\[
 f_j=4e\,12^j,\qquad
 v_j=\sum_{\alpha+\gamma+\delta=j}
       \frac{j!}{\alpha!\gamma!\delta!}\,
                  r_\alpha c_\gamma f_\delta .
\tag{V60.2}\label{c18v60:v}
\]
The deliberately enlarged $v_j$ need not equal V51's smaller $n_j$.

\begin{lemma}[Uniform third ordered normal-field jet]
\label{c18v60:normaljet}
For every $n\ge5$, every fine configuration, and $0\le t\le T$,
\[
 \mathfrak a_\zeta(\partial^3 C_t)\le c_3,\qquad
 \mathfrak a_\zeta(\partial^3 R_t)\le r_3,\qquad
 \mathfrak a_\zeta(\partial^j\mathcal N_t)\le v_j
                         \quad(0\le j\le3).
\tag{V60.3}\label{c18v60:normaljetbound}
\]
No conditional kinetic inverse is used.
\end{lemma}

\begin{proof}
An ordered derivative of a plaquette word, including repeated
derivatives in one edge, inserts unit quaternions and has magnitude
at most one. In the moving-frame generator $\mathsf M$ the
differentiated commutator term costs at most two in addition to the
ordinary derivative. Thus an entry of $\partial^3\mathsf M$ is at
most three per incident plaquette. Fixing any one slot leaves four
plaquettes and at most $12^4$ assignments of the other slots.
Their diameter is at most two. This proves
$\mathfrak a_\theta(\partial^3\mathsf M)\le4\cdot12^4\cdot3e=t_3$.
It is the next instance of V44's incidence proof, now through
five derivatives of the plaquette trace.

Three derivatives of an adjoint sampler prefix have entry bound
$2^3$. With one of its five slots anchored, at most
$(3\ell)^4$ assignments remain; the chain weight is at most $e$.
This gives $\mathfrak a_\theta(\partial^3\mathsf B)\le b_3$.
The same counting for the coefficient field $F$ gives
$\mathfrak a_\zeta(\partial^jF)\le f_j$.
For $j=0$ the bound $4e$ simply enlarges the component bound four.

Write $\mathsf J=D\Phi_t$, with its first and second ordered
coefficient jets bounded by $k,l$ in V44. Differentiating
$\dot{\mathsf J}=(\mathsf M\circ\Phi_t)\mathsf J$ three times,
and moving $(\mathsf M\circ\Phi_t)\partial^3\mathsf J$ to the
homogeneous part, gives forcing norm at most
\[
 4t_1Hl+3t_1k^2+6t_2H^2k+t_3H^4 .
\]
Indeed the first three pullback derivatives of $\mathsf M$
have respective bounds
$t_1H$, $t_2H^2+t_1k$, and
$t_3H^3+3t_2Hk+t_1l$. The ordered product rule gives
three first-derivative/second-jet terms, three
second-derivative/first-jet terms, and one third-derivative
term. Expanding gives precisely the displayed coefficients.
Directions in an iterated derivative of $\mathsf M$ are frozen;
their derivatives are the $k,l$ terms, as in
\eqref{c18v41:jets}.
The third jet starts at zero. Duhamel's formula and the propagator
bound $H$ prove the bound $m$.

Apply the same ordered product rule to
$C_t=(\mathsf B\circ\Phi_t)\mathsf J_t$. It gives $c_3$ in
\eqref{c18v60:thirdconstants}. All contractions retain an external
slot, so V42's anchored contraction lemma applies without a
volume factor. The bounds can be decreased from weight $\theta$
to $\zeta$.

Next $G=CC^*$ gives $g_3$. Differentiating $ZG=I$ three times
gives one $ZG^{(3)}Z$ term, six ordered terms involving one first
and one second Gram derivative, and six terms with three first
Gram derivatives. Their norms give $z_3$, using V44's
\emph{full} weighted inverse bound $z_0$. Finally $R=C^*Z$
gives $r_3$, and $\mathcal N=RCF$ gives the multinomial
formula \eqref{c18v60:v}. Slot order affects the tensors but not
these upper bounds.
\end{proof}

\subsection{A finite-time Hessian row bound for the native temporal source}

Let $\mu_g$ and $m_0$ be the product Haar law and pair/free metric
on $M_g=b_2^{-1}(g)$. Denote its product factors by $i,j$.
For a scalar function $h$ define
\[
 \mathcal H(h)=
       \sup_{x\in M_g}\max_i\sum_j
       \|(\operatorname{Hess}_{m_0}h)_{ij}(x)\|_{\rm op}.
\tag{V60.4}\label{c18v60:Hnorm}
\]
The supremum is outside the sum. This is a pointwise block row norm,
not a sum of separate block suprema and not a square-summed row norm.
By symmetry and the Schur test it bounds the Hessian operator norm.

Set
\[
\begin{aligned}
 K_1&=v_1+4v_0,&K_2&=v_2+4v_1,\\
 A_1&=e^{TK_1},&A_2&=TK_2A_1^3,\\
 B_2&=v_3A_1^2+v_2A_2+2v_2A_1,&K_*&=12B_2 .
\end{aligned}
\tag{V60.5}\label{c18v60:budget}
\]
These constants are finite and independent of $n,g$.
They are not asserted to be numerically useful.

\begin{theorem}[The actual electric temporal source has a local Hessian budget]
\label{c18v60:actualsource}
At $\beta=0$, throughout $0\le t\le4$, the actual normal-chart
log density $u_t=\log\widehat p_t$ and temporal score
$a_t=\partial_tu_t=\widehat b_t$ satisfy
\[
          \mathcal H(a_t)\le K_*,
       \qquad \mathcal H(u_t)\le tK_* .
\tag{V60.6}\label{c18v60:sourcebound}
\]
The conclusion is uniform over $n\ge5$ and all coarse $g$.
It controls the source and density, not
$\operatorname{Hess}_{\widehat m_t}\widehat L_t^{-1}a_t$.
\end{theorem}

\begin{proof}
At $\beta=0$, V51's raw temporal flux is the ambient Haar divergence
\[
 d_t=\sum_\alpha X_\alpha\mathcal N_{t,\alpha},
 \qquad
 a_t=-d_t\circ\eta_t+
                   \mathbb E_{\widehat\nu_t}(d_t\circ\eta_t).
\tag{V60.7}\label{c18v60:raw}
\]
Each invariant $X_\alpha$ has Haar divergence zero. In
$X_jd_t$ and $X_kX_jd_t$ the output index and one derivative
index of $\mathcal N_t$ are identified and summed.
An external derivative anchor remains. Consequently
\[
 \sup_j|X_jd_t|\le v_2,\qquad
      \mathfrak a_0(\partial^2d_t)\le v_3 .
\tag{V60.8}\label{c18v60:rawjets}
\]
This is the differentiated trace estimate, not a bound on the
extensive undifferentiated divergence.

The coefficient generator for $D\eta_t$ has Schur norm at most
$K_1$; its first spatial derivative has anchored norm at most
$K_2$. The commutator coefficients on a unit-round group contribute
at most four times the corresponding component or anchored norm.
The same ordered variational equations used in V51 therefore give
\[
       \mathfrak a_0(D\eta_t)\le A_1,\qquad
       \mathfrak a_0(\partial D\eta_t)\le A_2 .
\tag{V60.9}\label{c18v60:etajets}
\]
One may see the second bound directly by integrating the forcing
$(\partial\mathsf M_{\mathcal N})[D\eta_t,D\eta_t]$
against the first variational propagator.

The ordered scalar pullback rule now gives bounds
$v_2A_1$ for each first derivative of $d_t\circ\eta_t$ and
$v_3A_1^2+v_2A_2$ for its ordered second-derivative Schur norm.
In invariant orthonormal fine frames the connection is
$\nabla_{E_a}E_b=\sum_c\epsilon_{abc}E_c$ on each link.
For a fixed scalar row, subtracting the connection term costs
at most twice the gradient component bound. Thus its ambient
covariant Hessian Schur norm is at most $B_2$.

The inclusion of $M_g$ is totally geodesic. In pair coordinates
its derivative has a block $I/\sqrt2$ on one half and an
orthogonal block divided by $\sqrt2$ on the other; a free factor
has an identity block. Both scalar absolute row and column sums
are less than or equal to two. Restricting an ambient Hessian
therefore costs at most four in scalar Schur norm. Passing from
scalar rows to $3$-by-$3$ block operator rows costs at most three.
This proves the factor $12$ in $K_*$, uniformly in $g$.
The conditional mean in \eqref{c18v60:raw} is constant on the
fibre and has no spatial derivative.

Finally $\widehat p_0=1$ and $\partial_tu_t=a_t$ on the fixed
normal chart. Integrating its spatial Hessian in time proves
the second inequality in \eqref{c18v60:sourcebound}.
No derivative of a conditional expectation and no conditional
spectral gap was needed.
\end{proof}

This extends V51's component-gradient estimate by a genuine
summation over the remaining spatial index. The new input
is $\partial^3\mathcal N_t$, not another finite-time Taylor
coefficient. At $\beta\ne0$, differentiating
$2\,d\log\Omega_\beta(\mathcal N_t)$ twice additionally requires
a suitable anchored third derivative of $\log\Omega_\beta$.
That input is not supplied by V44's first and second derivative
bounds alone, so \eqref{c18v60:sourcebound} is not extended to
that case here.

\subsection{A local Hessian estimate for the Haar Poisson inverse}

Let $L_0=-\Delta_{m_0}$ and $P_s=e^{-sL_0}$ on the same fibre.
The following is an analytic reference estimate for arbitrary
smooth sources, not another evaluation of $a_0$.

\begin{lemma}[A one-factor, Banach-valued mixing bound]
\label{c18v60:heat}
On a round $\mathrm{SU}(2)$ factor with metric
$\kappa^{-1}g_{\rm round}$, $\kappa\in\{1/2,1\}$, let $h$ be a
continuous function with values in any finite-dimensional normed
space, with Haar mean zero. Then
\[
 \|P_s^{(\kappa)}h\|_\infty\le D(s)\|h\|_\infty,\qquad
 D(s)=
 \begin{cases}
 1,&0\le s\le1,\\
 \tfrac12 e^{-3(s-1)/2},&s>1,
 \end{cases}
 \qquad \int_0^\infty D(s)\,ds=\frac43 .
\tag{V60.10}\label{c18v60:D}
\]
The jump in the chosen majorant at $s=1$ is harmless.
\end{lemma}

\begin{proof}
For $s\le1$ use Markov contraction. For $s\ge1$, subtract
the constant kernel and use its scalar $L^1$ norm. Cauchy--Schwarz
and the round-sphere spectrum give
\[
 \|k_s(x,\cdot)-1\|_1^2
 \le\|k_s(x,\cdot)-1\|_2^2
 =\sum_{l\ge1}(l+1)^2e^{-2\kappa l(l+2)s}.
\]
Homogeneity makes the diagonal spectral multiplicity
$(l+1)^2$ independent of $x$. Since
$2\kappa l(l+2)\ge3l$, $(l+1)^2\le4^l$, and $e^3>20$,
the sum at $s=1$ is less than
$\sum_{l\ge1}(4e^{-3})^l<1/4$.
For later times factor out $e^{-3(s-1)}$.
The scalar kernel bound multiplies the norm of a vector-valued
integrand just as it does a scalar one, proving the assertion.
\end{proof}

\begin{theorem}[Haar inverse preserves the Hessian row norm]
\label{c18v60:haar}
For smooth $a$ on $M_g$, put
$\psi=L_0^{-1}(a-\mu_g a)$ with $\mu_g\psi=0$. Then
\[
              \mathcal H(\psi)\le\frac43\mathcal H(a).
\tag{V60.11}\label{c18v60:haarbound}
\]
The constant does not depend on the number of factors, the coarse
field, or the support of $a$.
\end{theorem}

\begin{proof}
Choose invariant orthonormal frames $E_{ia}$ on each factor.
The Casimir commutes with each $E_{ia}$. In these frames
\[
 (\operatorname{Hess}h)_{ia,jb}
   =E_{ia}E_{jb}h-
        \mathbf1_{\{i=j\}}\sqrt{\kappa_i}
                           \sum_c\epsilon_{abc}E_{ic}h .
\tag{V60.12}\label{c18v60:frameH}
\]
Thus every Hessian coefficient commutes with $P_s$.

Fix $i$. Every coefficient in the entire $i$th block row of
$\operatorname{Hess}a$ has mean zero upon integration in factor
$i$, with the other variables held fixed. This follows by
integration of an invariant derivative; it also holds for the
connection term. Regard this row as one vector-valued function,
with norm $\sum_j\|H_{ij}\|_{\rm op}$.
The other-factor heat operators preserve its factor-$i$ zero
mean and contract its supremum norm. Applying
Lemma \ref{c18v60:heat} in factor $i$ therefore gives
\[
        \mathcal H(P_sa)\le D(s)\mathcal H(a).
\]
The representation
$\psi=\int_0^\infty P_s(a-\mu_ga)\,ds$ holds on this compact
connected product. The displayed integrable Hessian bound
justifies taking two spatial derivatives under the integral;
alternatively use finite spectral sums and smooth approximation.
Integration proves \eqref{c18v60:haarbound}.
It is essential that the same factor-$i$ kernel acts on the
whole row before its norm is estimated.
\end{proof}

In particular, \eqref{c18v60:sourcebound} yields, through time four,
\[
 \mathcal H\!\left(L_0^{-1}
                (a_t-\mu_ga_t)\right)\le\frac43K_* .
\tag{V60.13}\label{c18v60:frozen}
\]
This is a bound for the actual source under a \emph{reference}
inverse. It is not the actual density correction:
generally $\mu_ga_t\ne0$, $\widehat p_t\ne1$, and
$\widehat m_t\ne m_0$. Equation \eqref{c18v60:frozen} cannot
be inserted into V51's strain formula.

\subsection{A gauge-invariant collective-source obstruction on the native fibre}

We now prove why the row norm matters. This comparison does
not alter the Wilson potential or claim that the following
source occurs in its normal chart.

In each coarse cell $z\in(\mathbb Z/n\mathbb Z)^3$, choose the
fine plaquette in directions $1,2$ based at $2z+e_3$.
These $M=n^3$ plaquettes are vertex-disjoint. Every one of
their four edges is free because its transverse third coordinate
is odd. Write $w_i$ for the real normalized holonomy trace.
They are gauge-invariant functions, independent under $\mu_g$,
and available on every coarse fibre. With the native product
metric,
\[
 L_0w_i=12w_i,\qquad
 |\nabla w_i|^2=4(1-w_i^2),\qquad
        \langle\nabla w_i,\nabla w_j\rangle=0\quad(i\ne j).
\tag{V60.14}\label{c18v60:loops}
\]
These identities follow from the four free round factors;
they are not replacements by independent auxiliary spins.

Define
\[
 S=M^{-1/2}\sum_{i=1}^M w_i,\quad
 Q=M^{-1}\sum_{i=1}^M w_i^2,\quad
 \varphi(s)=\sqrt{1+s^2},\quad
 \psi_M=\sqrt M\,\bigl(\varphi(S)-\mu_g\varphi(S)\bigr).
\tag{V60.15}\label{c18v60:psi}
\]
Its centered Poisson source is exactly
\[
 a_M=L_0\psi_M
   =4\sqrt M\,\bigl(3S\varphi'(S)-(1-Q)\varphi''(S)\bigr),
       \qquad \mu_ga_M=0 .
\tag{V60.16}\label{c18v60:asource}
\]

\begin{theorem}[Per-factor control does not bound pointwise strain]
\label{c18v60:obstruction}
The smooth gauge-invariant pair $(a_M,\psi_M)$ satisfies
\[
\begin{gathered}
 \sup_e\|\nabla_e a_M\|_\infty\le44,\qquad
 \sup_e\|\nabla_e\psi_M\|_\infty\le1,\\
 \langle a_M,L_0^{-1}a_M\rangle_{\mu_g}
                  =\int|\nabla\psi_M|^2d\mu_g\le4M,\\
 \sup_x\max_e\sum_f
       \|(\operatorname{Hess}\psi_M)_{ef}(x)\|_{\rm HS}^2
                  \le32 .
\end{gathered}
\tag{V60.17}\label{c18v60:small}
\]
Unused pair/free factors contribute zero.
Nevertheless, at any configuration with all selected $w_i=0$,
\[
 \|\operatorname{Hess}\psi_M\|_{\rm op}=4\sqrt M,\qquad
 \mathcal H(\psi_M)\ge4\sqrt M,\qquad
 \mathcal H(a_M)\ge144\sqrt M+\frac{32}{\sqrt M}.
\tag{V60.18}\label{c18v60:large}
\]
The underlying Haar generator has a volume-independent gap.
\end{theorem}

\begin{proof}
The chain rule and \eqref{c18v60:loops} give
\eqref{c18v60:asource}. For a selected loop variable,
\[
 \partial_{w_i}a_M
  =4\left[
    3\varphi'(S)+3S\varphi''(S)
       -(1-Q)\varphi'''(S)
       +\frac{2w_i}{\sqrt M}\varphi''(S)\right].
\tag{V60.19}\label{c18v60:agrad}
\]
Here $|\varphi'|\le1$, $|\varphi''|\le1$,
$|s\varphi''(s)|\le1$, $|\varphi'''(s)|\le3$, $0\le Q\le1$.
Each chosen fine edge belongs to one selected loop and has
$|\nabla_e w_i|\le1$. This proves $44$.
Similarly $\nabla_e\psi_M=\varphi'(S)\nabla_e w_i$ gives
the component bound and total energy bound $4M$.
Integration by parts proves the cost identity.

Let $R=\sum_iw_i$. Pointwise,
\[
 \operatorname{Hess}\psi_M
   =\frac{\varphi''(S)}{\sqrt M}\,dR\otimes dR
                              +\varphi'(S)\operatorname{Hess}R .
\tag{V60.20}\label{c18v60:Hpsi}
\]
For a fixed selected edge $e$, the squared Hilbert--Schmidt
block row sum of the first term is at most
$M^{-1}|\nabla_eR|^2|\nabla R|^2\le4$.
In the second term only the four edges of that plaquette occur.
Every diagonal or mixed Hessian block of a single plaquette
has operator norm at most one: a diagonal block is $-w_iI$,
and a mixed block is a two-insertion quaternion trace.
Its Hilbert--Schmidt norm squared is at most three.
The second term's row sum is therefore at most twelve.
Using $|A+B|^2\le2|A|^2+2|B|^2$ gives $32$.

All $w_i=0$ can be realized simultaneously using the disjoint
free edges without changing $g$. There
$\varphi'(0)=0$, $\varphi''(0)=1$, and $|dR|^2=4M$,
so \eqref{c18v60:Hpsi} is a rank-one tensor with eigenvalue
$4\sqrt M$. Each of its $4M$ nonzero blocks in a chosen edge row
has operator norm $1/\sqrt M$.

At the same configuration $\partial_{w_i}a_M=0$ and
\[
   \partial_{w_i}\partial_{w_j}a_M
                 =\frac{4}{\sqrt M}(9+2\delta_{ij}).
\tag{V60.21}\label{c18v60:aH}
\]
Hence no Hessian-of-$w_i$ term remains in its chain rule.
For an edge in loop $i$, the four blocks in loop $i$ each
have norm $44/\sqrt M$, and the $4(M-1)$ blocks in other
loops each have norm $36/\sqrt M$.
Their sum is the last quantity in \eqref{c18v60:large}.

Finally the full product's smallest nonzero Casimir is $3/2$
on a pair factor, or $3$ on a free factor. Tensorization
therefore gives the asserted uniform gap. Restricting the
Poincare inequality to gauge-invariant functions cannot
decrease its lower bound; our source and solution are already
invariant under the gauge transformations preserving the fibre.
\end{proof}

Multiplying both $a_M$ and $\psi_M$ by $M^{-1/4}$ makes the
source-gradient component bound tend to zero and makes the
cost divided by $M$ tend to zero. The square-summed Hessian
row bound tends to zero as well. The Hessian operator norm
still grows as $4M^{1/4}$. Thus even these strengthened
smallness statements do not imply a volume-uniform pointwise
Poisson-strain estimate. This does not contradict
Theorem \ref{c18v60:haar}: the source's \emph{row-summed}
Hessian norm diverges. It also does not contradict
Theorem \ref{c18v60:actualsource}, which excludes this particular
collective scaling for the actual temporal source.

\subsection{The remaining nonlinear passage is explicit}

Even after paying \eqref{c18v60:sourcebound}, Haar tensorization
cannot be used for the actual inverse without an additional
argument. To exhibit that argument's missing input, first
retain $m_0$ but put the \emph{actual} density $e^u$ on it:
\[
 {\cal L}_u=-\Delta_{m_0}-\Gamma(u,\cdot),\qquad
             \Gamma(f,h)=\langle\nabla f,\nabla h\rangle_{m_0}.
\tag{V60.22}\label{c18v60:Lu}
\]
For invariant Killing fields $X,Y$ on product factors,
and any solution $\mathcal L_u\psi=a$, the exact identities are
\[
\begin{aligned}
 {\cal L}_u(X\psi)
   &=Xa+\Gamma(Xu,\psi),\\
 {\cal L}_u(XY\psi)
   &=XYa+\Gamma(XYu,\psi)
               +\Gamma(Yu,X\psi)+\Gamma(Xu,Y\psi).
\end{aligned}
\tag{V60.23}\label{c18v60:commutators}
\]
Indeed $X$ commutes with $\Delta_{m_0}$ and, being Killing,
$X\Gamma(f,h)=\Gamma(Xf,h)+\Gamma(f,Xh)$.
Applying this once and twice proves the identities with the
displayed signs. They are not equations with the Haar operator
on the left. They involve a third spatial derivative of $u$
and coupled derivatives of $\psi$ on the right.

The factorwise Haar cancellation used in
\eqref{c18v60:haarbound} is not preserved by an arbitrary
interacting drift. Moreover $\mathcal L_{u_t}$ still differs
from $\widehat L_t$ when $\widehat m_t\ne m_0$.
Solving with $\mathcal L_{u_t}$ would give a nonminimal velocity
for the same continuity equation, but its material strain
would have to be computed with the actual metric. Neither
replacement removes an estimate.

The next upstream calculation is therefore the nonlinear
local response: either control these coupled differentiated
equations (including the third density derivative and
anchored metric coefficients), or construct another velocity
with a proved row-summed differential bound for the same
continuity equation. The finite constants in
\eqref{c18v60:budget} establish locality, not a small
interaction condition through time four. Their magnitude does
not certify a long-time Dobrushin or curvature inequality.
No additional zero-time cost coefficient is needed for this step.

\subsection{Verification and scope after V60}

The diagnostic
\begin{center}\small
\path{scripts/verify_ym_c18_v60_local_poisson.py}
\end{center}
checks the ordered third-jet product multiplicities, inverse
jet majorant, differentiated collective source, its exact
Hessian at zero, the rescaled exponents, the one-factor
heat-kernel majorant, and the native disjoint free-plaquette
packing at $n=5,6$. A noncommuting two-factor group fixture
checks \eqref{c18v60:commutators}; it is not a simulated
conditional YM trajectory. The general-volume assertions
are supplied by the written incidence, contraction and
analytic proofs, not by finite-size tests or proof-assistant
certification.

The native temporal source now has a spatial Hessian row
budget on the full electric block interval. The comparison
counterexample explains why the previously available
per-factor gradient and scalar inverse costs were insufficient.
The reference Haar inverse supplies a usable local norm
estimate, but not the required nonlinear continuation.

Still unpaid are the actual source-specific Poisson-Hessian
or alternative-velocity strain bound through time four,
the interacting-vacuum extension, coarse variance/return,
variable-bank stability, weak-coupling scale transport,
physical-time calibration, and interacting continuum
reconstruction. In particular, a positive physical
four-dimensional Yang--Mills mass gap remains unproved.

Removing this append and restoring the title recovers V59
byte for byte. Other sources are preserved; only the active
YM filename is replaced. The PDF skill is used for internal
layout verification, with all build products outside
\path{latex_notes}; that folder contains only \texttt{.tex}
sources. Archive/restart remains due at V100.
\addtocontents{toc}{\protect\setcounter{tocdepth}{2}}
\addtocontents{toc}{\protect\clearpage}
% END CYCLE 18 VERSION 60 APPEND
% BEGIN CYCLE 18 VERSION 61 APPEND
\clearpage
\section{V61: Anchored entropy decoupling and nonlinear probability transport}
\label{c18v61:context}

V60 paid a spatial Hessian row budget for the actual electric
temporal source throughout the block interval, but its Haar
inverse estimate did not apply to the moving density.
This append makes that nonlinear passage on an explicit,
volume-uniform \emph{short} interval. It constructs a velocity
for the actual conditional probability law and bounds its
strain in the actual moving metric. It does not extend the
interval to time four, and it is not a continuum mass-gap proof.

The new input is one-factor drift removal. The total change
of drift is controlled by an anchored Hessian row, even when
the norm of the full drift is extensive. A path-entropy
comparison then transfers a one-factor Haar mixing step to
the density-weighted diffusion. Ordered derivative rows have
the required Haar cancellation each time the argument is
restarted. A third density derivative supplies the remaining
inhomogeneous term. We first prove this analytic statement
on arbitrary finite products, and then pay its inputs on the
native electric normal chart. This is not another Taylor
coefficient or a replacement of the moving law by Haar.

\subsection{Norms, centering, and the nonlinear analytic statement}

Let $M$ be a finite product of copies of $\mathrm{SU}(2)$ with
bi-invariant metrics $\kappa_i^{-1}m_{\rm round}$,
$\kappa_i\in\{1/2,1\}$. Let $\mu$ be product normalized Haar
measure and $m_0$ the product metric. Write $E_\alpha=E_{ia}$
for its invariant orthonormal frame; there are three colors
per factor. All scalar derivatives below are ordered, with
no commutation of consecutive $E_\alpha$ understood.
For a scalar matrix $A$, set
\[
 \|A\|_{\rm S}=
 \max\left\{\max_\alpha\sum_\beta|A_{\alpha\beta}|,
                  \max_\beta\sum_\alpha|A_{\alpha\beta}|\right\}.
\]
We use the following configuration-uniform norms:
\begin{align*}
 G(f)&=\sup_x\max_i|\nabla_i f(x)|,\\
 J(f)&=\sup_x\| (E_\alpha E_\beta f(x))_{\alpha\beta}\|_{\rm S},\\
 T_3(f)&=\sup_x\mathfrak a_0
       ((E_\gamma E_\alpha E_\beta f(x))_{\gamma\alpha\beta}),\\
 H(f)&=\sup_x\max_i\sum_j
             \|(\operatorname{Hess}_{m_0}f)_{ij}(x)\|_{\rm op}.
\end{align*}
Here $\mathfrak a_0$ is V42's all-slot norm: fix any one
scalar slot, sum the absolute values over every other slot,
and take the maximum over the fixed slot and its position.
The supremum over configurations is \emph{outside} the sums.
No sum of separate configuration suprema is assumed.

Put $D=\pi\sqrt2<5$, a diameter bound for every factor.
The constant connection coefficients on factor $i$ are
$\Gamma_{ab}^{c}=\sqrt{\kappa_i}\epsilon_{abc}$, up to the
consistent orientation choice of the frame. Therefore
\begin{equation}
 G(f)\le D H(f),\qquad
 J(f)\le3H(f)+2G(f),\qquad
 \|\operatorname{Hess}_{m_0}f\|_{\infty,\rm op}\le J(f)+G(f).
 \tag{V61.1}\label{c18v61:norms}
\end{equation}
For the first assertion, fix every coordinate except $i$,
choose a maximum of that one-factor function, and integrate
its gradient along a minimizing geodesic of length at most
$D$. For the second, convert each $3$ by $3$ Hessian block
to scalar absolute row/column sums, with factor at most $3$.
The connection correction has scalar Schur norm at most
$2G(f)$. For the third, its block is a cross-product matrix
of operator norm at most $|\nabla_i f|$, and scalar Schur
norm bounds the operator norm of the ordered matrix.

For smooth real $u$, let $\nu=Z^{-1}e^u\mu$ and
\[
 \mathcal L_u=-\Delta_{m_0}-\Gamma(u,\cdot),\qquad
 P_s^u=e^{-s\mathcal L_u},\qquad
 \Gamma(f,h)=\langle\nabla f,\nabla h\rangle_{m_0}.
\]
The variable $s$ is auxiliary diffusion time, not Wilson
block time and not physical time.

\begin{theorem}[Anchored nonlinear Poisson estimate near Haar]
\label{c18v61:poisson}
Suppose $q=H(u)\le1/256$, and put $c=T_3(u)$.
For every smooth $a$ with $\nu(a)=0$, let
$\psi=\mathcal L_u^{-1}a$ be the solution normalized by
$\nu(\psi)=0$. Then
\begin{align}
 G(\psi)&\le2G(a),\nonumber\\
 J(\psi)&\le5J(a)+20cG(a),\nonumber\\
 \|\operatorname{Hess}_{m_0}\psi\|_{\infty,\rm op}
       &\le5J(a)+(20c+2)G(a).
 \tag{V61.2}\label{c18v61:poissonbounds}
\end{align}
The constants do not depend on the number of factors.
No smallness of $c$ is required.
\end{theorem}

\subsection{One-factor removal and the entropy comparison}

Denote averaging only in factor $i$ by $\mu_i$ and set
$u^{(i)}=\mu_i u$. This potential is independent of that
factor, but may retain all interactions among the others.
We claim
\begin{equation}
 \|\nabla u-\nabla u^{(i)}\|_{\infty,m_0}\le2Dq.
 \tag{V61.3}\label{c18v61:driftdifference}
\end{equation}
The $i$ component of the difference has norm at most $Dq$
by \eqref{c18v61:norms}. For $j\ne i$, compare the $j$
gradients at $x_i$ and $y_i$, holding all other coordinates
fixed. Product parallel transport in the $j$ tangent space
is the identity along this path. Integration along an $i$
geodesic gives
\[
 \sum_{j\ne i}|\nabla_j u(x)-\nabla_j u^{(i)}(x)|
 \le \int \int_0^{d(x_i,y_i)}
            \sum_{j\ne i}\|\operatorname{Hess}_{ji}u\|_{\rm op}
                         \,dr\,\mu_i(dy_i)\le Dq.
\]
The Hessian is symmetric, so the column used here has the
same bound as the defining row. Combining these bounds in
the full product tangent norm proves
\eqref{c18v61:driftdifference}. In particular, this difference
is not estimated by the extensive norm of either drift.

Here is the precise probabilistic estimate used below.
For two smooth diffusions on this compact product, both
with generator $\Delta_{m_0}+$ drift and with the same
deterministic initial point, suppose their drift difference
has norm at most $\delta$. Their path laws up to time $s$
satisfy
\[
 {\cal H}(\mathbf P\mid\mathbf R)
      =\frac14\,\mathbf E_{\mathbf P}
            \int_0^s|b-\widetilde b|^2(X_r)\,dr
      \le\frac{s\delta^2}{4}.
\]
Indeed the diffusion can be written in the global invariant
frame with noise $\sqrt2 E_\alpha$. The drift-change control
is $(b-\widetilde b)/\sqrt2$. Its boundedness gives the
exponential-martingale condition; the logarithm of the
likelihood ratio, averaged under the changed law, is one
half of the squared control integral. Smoothness and
compactness give uniqueness of both stochastic equations.
Thus this is an entropy \emph{equality}, not an unjustified
use of the opposite-direction general entropy inequality.

Pinsker's inequality for variation norm,
$\|P-R\|_{\rm var}\le\sqrt{2{\cal H}(P\mid R)}$, and
pushforward to the terminal point imply, also for
finite-dimensional Banach-valued $h$,
\begin{equation}
 \|(P_s^u-P_s^{u^{(i)}})h\|_\infty
     \le \delta\sqrt{s/2}\,\|h\|_\infty.
 \tag{V61.4}\label{c18v61:entropy}
\end{equation}
For the Banach assertion, integrate the pointwise norm
against the total variation measure; no dimension-dependent
coordinate summation is needed.

At time one the normalized heat kernel of either allowed
single factor satisfies
$\|k_1-1\|_{L^1(\mu_i)}<1/2$, as proved in
Lemma \ref{c18v60:heat}. The sharper endpoint value follows
directly from that lemma's kernel calculation, even though
its convenient piecewise majorant uses $1$ at the endpoint.
Since $u^{(i)}$ does not depend on $i$, its diffusion
factorizes into this factor heat semigroup and a Markov
semigroup on the remaining factors. Hence $\mu_i h=0$
pointwise in the other coordinates gives
\begin{equation}
 \|P_1^u h\|_\infty\le\alpha\|h\|_\infty,\qquad
 \alpha=\frac12+\sqrt2Dq=\frac12+2\pi q\le\frac{17}{32}.
 \tag{V61.5}\label{c18v61:seed}
\end{equation}
We do \emph{not} claim that $P_s^u$ preserves $\mu_i$
centering. The next argument supplies a freshly centered
derivative row or column at each integer time.

The literature input is the bounded-control case of the
path-entropy identity, with its normalization and uniqueness
hypotheses made explicit above. For the more general
finite-entropy formulation see C.~L\'eonard,
\emph{Girsanov theory under a finite entropy condition},
Theorem 2.3,
\url{https://arxiv.org/html/1101.3958v2}.
Its uniqueness condition is essential to the equality used
here. No entropy estimate for the unknown continuum YM
measure is imported from this reference.

\subsection{Derivative recurrence and proof of the Poisson theorem}

Put $f_s=P_s^u a$. The component Bochner equation and
Kato inequality give
\begin{equation}
 G(f_s)\le e^{-(1-q)s}G(a)\le e^{-s/2}G(a).
 \tag{V61.6}\label{c18v61:gradient}
\end{equation}
For completeness, the equation for the $i$ component of
$df_s$ has diffusion and drift covariant derivatives,
a term $-\operatorname{Ric}_i df_{s,i}$, and forcing
$\sum_j(\operatorname{Hess}u)_{ij}df_{s,j}$.
Here $\operatorname{Ric}_i=2\kappa_i I\ge I$.
At a maximum of the largest component norm the last term
is at most $q\max_j|df_{s,j}|$. Applying the maximum
principle to regularized component norms and then letting
the regularization tend to zero proves the first bound.
This uses the row sum at the same configuration, not the
sum of separately maximized coefficients.

Let $K_{\alpha\beta}(f)=E_\alpha E_\beta f$.
The Killing identities of V60, or direct differentiation,
give the exact coefficient equation
\begin{align}
 \partial_s K_{\alpha\beta}(f_s)
 &= (\Delta+\Gamma(u,\cdot))K_{\alpha\beta}(f_s)
                                    +F_{\alpha\beta}(s),\nonumber\\
 F_{\alpha\beta}(s)
 &=\sum_\gamma(E_\gamma E_\alpha E_\beta u)(E_\gamma f_s)
   +\sum_\gamma(E_\gamma E_\beta u)
                         (E_\gamma E_\alpha f_s)\nonumber\\
 &\hspace{12mm}
   +\sum_\gamma(E_\gamma E_\alpha u)
                         (E_\gamma E_\beta f_s).
 \tag{V61.7}\label{c18v61:derivativePDE}
\end{align}
In particular these are ordered, possibly noncommuting
derivatives. The two matrix products are
$K(f_s)^T K(u)+K(u)^T K(f_s)$.

Write $R(s)=J(f_s)$ and $h=J(u)\le(3+2D)q<13q$.
All-slot contraction and Schur submultiplicativity give
\begin{equation}
 \sup_x\|F(s,x)\|_{\rm S}
       \le cG(a)e^{-s/2}+2hR(s).
 \tag{V61.8}\label{c18v61:forcing}
\end{equation}
A scalar Markov operator acting coefficientwise contracts
the uniform matrix Schur norm. Consequently Duhamel's
formula and Gronwall yield, for $0\le r\le1$,
\[
 R(s+r)\le e^{2hr}\bigl(R(s)+r cG(a)e^{-s/2}\bigr).
\]

Every row of $K(f)$ has zero Haar average in the factor of
its first index, by integration of its outer derivative.
Every column has zero Haar average in the factor of its
second index: if the two factors differ, commute the two
derivatives; if they agree, integrate the outer derivative
as before. Regard a row or column as an $\ell^1$-valued
function. Applying \eqref{c18v61:seed} separately to these
anchored functions in the initial term of Duhamel's formula,
while using \eqref{c18v61:forcing} in its integral term,
gives
\[
 R(s+1)\le(\alpha+e^{2h}-1)R(s)
                  +e^{2h}cG(a)e^{-s/2}.
\]
For example the scalar convolution upper bound solves
$y'=2hy+cG(a)e^{-s/2}$ with the last exponential held at its
left-endpoint value; its integrated forcing is at most
$e^{2h}cG(a)e^{-s/2}$. Only the initial term uses anchoring.

Since $2h<26/256<1/8$,
\[
 e^{2h}\le\frac87,\qquad
 \alpha+e^{2h}-1\le\frac{17}{32}+\frac17
                      =\frac{151}{224}<\frac34.
\]
It follows that for integers $k\ge0$ and $0\le r\le1$,
\begin{align}
 R(k+1)&\le\frac34R(k)+\frac87cG(a)e^{-k/2},\nonumber\\
 R(k+r)&\le\frac87\bigl(R(k)+cG(a)e^{-k/2}\bigr).
 \tag{V61.9}\label{c18v61:recurrence}
\end{align}
Summing the first inequality, and then integrating the
second, proves
\[
 \int_0^\infty R(s)\,ds
 \le\frac{32}{7}R(0)
       +\frac{312}{49}\frac{cG(a)}{1-e^{-1/2}}
 \le5J(a)+20cG(a).
\]
For the last step, $e^{1/2}>3/2$ implies
$(1-e^{-1/2})^{-1}<3$, and $936/49<20$.

At each fixed finite volume, compact elliptic theory gives
the centered inverse and the representation
$\psi=\int_0^\infty P_s^u a\,ds$. Any finite-volume scalar
mixing constants needed for that representation are not
used in the derivative estimates. Equations
\eqref{c18v61:gradient} and \eqref{c18v61:recurrence}
give absolutely integrable first and second derivative
bounds. They justify differentiation of the integral and
prove the first two assertions of
Theorem \ref{c18v61:poisson}. Equation
\eqref{c18v61:norms} proves the last assertion.
Equivalently one may first integrate to a finite upper
limit and pass to the normalized elliptic solution.
This completes the proof of the theorem.

\subsection{Paying the native third density derivative}

We now return to $\ell=2$, $T=4$, $\beta=0$, $n\ge5$,
and arbitrary coarse configuration $g$. All notation for
the normal chart $\eta_t$, the normal field $\mathcal N_t$,
the actual law $\widehat\nu_t=e^{u_t}\mu_g$, and the source
$a_t=\partial_tu_t$ is as in V60.
In particular $\widehat\nu_t(a_t)=0$ and $u_0=0$.
Let $K_*$ be exactly V60's constant in
\eqref{c18v60:budget}; thus
\begin{equation}
 H(a_t)\le K_*,\qquad H(u_t)\le tK_*
       \quad(0\le t\le T).
 \tag{V61.10}\label{c18v61:oldinput}
\end{equation}
Here is a finite fourth-jet recursion supplying the
additional input, rather than assuming a volume-uniform
third density derivative.

Use V44's weight, propagation bound $H$, inverse bound
$z_0$, and retained weight $\zeta$. For $1\le r\le4$ put
$t_r=144\mathrm e\,12^r$; for $0\le r\le4$ put
\[
 b_r=3\mathrm e\,\ell(6\ell)^r,\qquad
 f_r=4\mathrm e\,12^r .
\]
These extend the local bounds used in V60.
Define the exponential partial Bell polynomials by
\[
 {\sf B}_{p,k}(x_1,\ldots,x_{p-k+1})
 =\sum_{\substack{\sum m_j=k\\\sum jm_j=p}}
          p!\prod_j\frac{(x_j/j!)^{m_j}}{m_j!}.
\]
The following are \emph{new, possibly larger} majorants;
they do not redefine any sharper constants in V44 or V60.
Set $j_0=H$ and successively for $1\le r\le4$ define
\begin{align*}
 M_p^\#&=\sum_{k=1}^p t_k
                  {\sf B}_{p,k}(j_0,\ldots,j_{p-k})
                      \quad(1\le p\le r),\\
 j_r&=TH\sum_{p=1}^r {r\choose p}M_p^\#j_{r-p}.
\end{align*}
This is triangular: the right side for $j_r$ involves only
$j_0,\ldots,j_{r-1}$. Next put $B_0^\#=b_0$ and
\[
 B_p^\#=\sum_{k=1}^p b_k
                       {\sf B}_{p,k}(j_0,\ldots,j_{p-k})
                       \quad(1\le p\le4).
\]
For $0\le r\le4$, define
\begin{align*}
 C_r^\#&=\sum_{p=0}^r{r\choose p}B_p^\#j_{r-p},&
 G_r^\#&=\sum_{p=0}^r{r\choose p}C_p^\#C_{r-p}^\#,\\
 Z_0^\#&=z_0,&
 Z_r^\#&=z_0\sum_{p=1}^r{r\choose p}G_p^\#Z_{r-p}^\#
                                  \quad(r\ge1),\\
 R_r^\#&=\sum_{p=0}^r{r\choose p}C_p^\#Z_{r-p}^\#,&
 v_r^\#&=\sum_{\alpha+\gamma+\delta=r}
          \frac{r!}{\alpha!\gamma!\delta!}
                    R_\alpha^\#C_\gamma^\#f_\delta .
\end{align*}
In particular $C_0^\#=b_0H$ is unchanged, so the inherited
$z_0,\zeta$ from the full Gram inverse remain applicable.

\begin{lemma}[Fourth normal jet with retained anchors]
\label{c18v61:fourthjet}
Uniformly in $n,g$ and $0\le t\le T$,
\[
 \mathfrak a_\zeta(\partial^r\mathcal N_t)\le v_r^\#
                        \quad(0\le r\le4).
\]
All these constants are finite and explicit finite
recursions in the previously displayed geometric constants.
\end{lemma}
\begin{proof}
The variational generator $M$ of the Wilson flow has
$\mathfrak a_\theta(\partial^r M)\le t_r$ for $1\le r\le4$.
Indeed a fixed differentiated slot belongs to at most four
plaquettes; each remaining slot has at most twelve scalar
choices, and the differentiated local coefficient is at
most three. The connected diameter weight is at most
$\mathrm e$. The same counting for the sampler gives
$\mathfrak a_\theta(\partial^r Db_\ell)\le b_r$:
there are at most $3\ell$ choices initially and at most
$6\ell$ for each additional derivative. The bound for
$\partial^rF$ is $f_r$, as in V60.

Let $J_t=D\Phi_t$. Differentiate
$\dot J_t=(M\circ\Phi_t)J_t$ in $r$ ordered spatial
directions. The undifferentiated propagator is bounded by
$H$. The chain-rule partitions for
$\partial^p(M\circ\Phi_t)$ are bounded by $M_p^\#$.
Their differentiated flow factors of orders $1,2,\ldots$
are bounded by $j_0,j_1,\ldots$. The product rule then
gives the displayed Duhamel majorant for $j_r$, because
$\partial^rJ_0=0$ for $r\ge1$.
Ordered words are kept distinct in the tensor formulas;
the positive scalar partition counts are the ordinary
Bell coefficients. No interchange of invariant derivatives
is used.

Apply the same partition argument to
$C_t=(Db_\ell\circ\Phi_t)J_t$ to obtain $C_r^\#$.
The product $G_t=C_tC_t^*$ gives $G_r^\#$.
Differentiating $G_tZ_t=I$ in ordered slots and isolating
the undifferentiated $G_t$ gives the stated inverse
recursion with its binomial upper counts.
The full inverse bound is $z_0$; this is not an inverse
of a selected finite block. The product $R_t=C_t^*Z_t$
gives $R_r^\#$, and $\mathcal N_t=R_tC_tF$ gives
$v_r^\#$.

Every contraction in these formulas is connected and has
an external slot. V42's weighted all-slot contraction
inequality therefore applies throughout. Differentiated
divergence traces below retain their derivative slots.
There is no unanchored sum proportional to volume in
this argument.
\end{proof}

For clarity, the subsequent constants use only the enlarged
$v_r^\#$:
\begin{align*}
 k_1&=v_1^\#+4v_0^\#,&
 k_2&=v_2^\#+4v_1^\#,&
 k_3&=v_3^\#+4v_2^\#,\\
 E_0&=e^{Tk_1},&
 E_1&=Tk_2E_0^3,&
 E_2&=TE_0(3k_2E_0E_1+k_3E_0^3),\\
 d_1&=v_2^\#E_0,&
 d_2&=v_3^\#E_0^2+v_2^\#E_1,&
 d_3&=v_4^\#E_0^3+3v_3^\#E_0E_1+v_2^\#E_2.
\end{align*}
The moving-frame variational equation for $\eta_t$ gives
bounds $E_0,E_1,E_2$ for $D\eta_t,\partial D\eta_t,
\partial^2D\eta_t$ respectively. To see this directly,
the generator for $D\eta_t$ is bounded by $k_1$, its first
spatial derivative by $k_2$, and its second by $k_3$.
These include the bracket coefficients, whose all-slot
norm is at most four. One and two differentiations of its
Duhamel equation give exactly the displayed majorants.

Write $d_t=\operatorname{div}_{\rm Haar}\mathcal N_t$.
A differentiated divergence trace has
$\partial d_t,\partial^2d_t,\partial^3d_t$ bounded by
$v_2^\#,v_3^\#,v_4^\#$ in their anchored norms.
Composition with $\eta_t$ gives $d_1,d_2,d_3$ for the
corresponding ordered derivatives. The un-differentiated
extensive trace has not been bounded by these constants.

The initial fibre inclusion $i_g:M_g\longrightarrow
{\cal A}_{2n}$ has invariant-frame derivative matrix $E$
with
\[
 \mathfrak a_0(E)\le2,\qquad
 \mathfrak a_0(\partial E)\le36,\qquad
 \mathfrak a_0(\partial^2 E)\le324.
\]
Here free factors are constant inclusions. On a pair
factor, the two half-edge derivatives are a constant
isometry component and an orthogonal adjoint matrix,
each scaled by $1/\sqrt2$. Each derivative of an adjoint
entry is at most two, and each second derivative at most
four. Only three colors in the same factor can occur.
Counting the remaining slots gives the deliberately loose
bounds $36,324$, independently of $g$.

Since $a_t=-d_t\circ\eta_t\circ i_g+
\widehat\nu_t(d_t\circ\eta_t\circ i_g)$, its spatially
constant mean disappears after differentiation.
The three chain-rule types in its third derivative give
\begin{equation}
 \begin{split}
 K_3&=8d_3+216d_2+324d_1,\qquad B=4d_1,\qquad
 A=3K_*+2B,\\
 T_3(a_t)&\le K_3,\quad T_3(u_t)\le tK_3,\quad
 G(a_t)\le B,\quad J(a_t)\le A .
 \end{split}
 \tag{V61.11}\label{c18v61:paid}
\end{equation}
The coefficient $216$ is $3\cdot2\cdot36$.
For the gradient, each scalar component is at most $2d_1$
and there are three components per factor, so $4d_1$
suffices. The final bound follows from
\eqref{c18v61:norms} and \eqref{c18v61:oldinput}.
Finally $u_t=\int_0^t a_r\,dr$ proves the third-density
bound. These arguments apply at $\beta=0$; no higher
interacting-vacuum derivative has silently been assumed.

\subsection{A velocity for the actual law, with the actual metric}

Define the strictly positive, volume-uniform time
\begin{equation}
 \tau=\min\{4,(256K_*)^{-1}\}.
 \tag{V61.12}\label{c18v61:tau}
\end{equation}
For $0\le t\le\tau$, solve in the \emph{fixed product metric}
\[
 \mathcal L_{u_t}\psi_t=a_t,\qquad
 \widehat\nu_t(\psi_t)=0,\qquad V_t=\nabla_{m_0}\psi_t .
\]
The measure in this operator is the \emph{actual} pulled
conditional law. This is a nonminimal choice of velocity;
it is not asserted to equal
$\nabla_{\widehat m_t}\widehat L_t^{-1}a_t$.

By Theorem \ref{c18v61:poisson} and
\eqref{c18v61:paid}, put
\[
 P_*=2B,\qquad H_*=5A+(80K_3+2)B .
\]
Then $G(\psi_t)\le P_*$ and
$\|\nabla^0 V_t\|_{\infty,\rm op}\le H_*$ on this
interval. The first is a \emph{component} bound on $V_t$,
not a bound on its possibly extensive full length.

We check its material strain rather than discard the
moving metric. Let
$\widehat m_t=(\eta_t\circ i_g)^*m_{\rm ambient}$ and let
$\overline J_t=(D\eta_t\circ i_g)E$ be its frame matrix.
The preceding jet bounds give
\[
 \mathfrak a_0(\overline J_t)\le2E_0,\qquad
 \mathfrak a_0(\partial\overline J_t)\le4E_1+36E_0,
 \qquad
 \mathfrak a_0(\partial\widehat m_t)
                 \le16E_0E_1+144E_0^2.
\]
The base connection has all-slot norm at most two.
Correcting the two metric slots therefore gives
\begin{equation}
 \mathfrak a_0(\nabla^0\widehat m_t)\le M_1,\qquad
 M_1=16E_0E_1+160E_0^2 .
 \tag{V61.13}\label{c18v61:metricjet}
\end{equation}
All the displayed tensor norms include a uniform
configuration supremum. The full ambient flow and its
inverse also satisfy the operator bounds $E_0$; the initial
inclusion is an isometry in operator norm. Consequently
\[
 E_0^{-2}m_0\le\widehat m_t\le E_0^2m_0,\qquad
 |\partial_t\widehat m_t|\le2k_1\widehat m_t.
\]
The last estimate is the normal-flow variational estimate
of V51 with the enlarged constants above.

For any vector field $V$ the exact identity is
\[
 {\cal L}_V\widehat m
   =\nabla^0_V\widehat m
       +\widehat m(\nabla^0 V\,\cdot,\cdot)
       +\widehat m(\cdot,\nabla^0 V\,\cdot).
\]
Contracting the first tensor in
\eqref{c18v61:metricjet} with the components of $V_t$
gives operator norm at most $M_1P_*$, with no factor
equal to the square root of volume. The other two terms
are bounded by $2E_0^2H_*$ relative to $m_0$. Thus
\begin{equation}
 \partial_t\widehat m_t+{\cal L}_{V_t}\widehat m_t
     \le \sigma\widehat m_t,\qquad
 \sigma=2k_1+E_0^2(M_1P_*+2E_0^2H_*).
 \tag{V61.14}\label{c18v61:strain}
\end{equation}
This is a bound on the actual material metric, although the
velocity was found using a different metric.

\begin{theorem}[Native short-time conditional probability transport]
\label{c18v61:transport}
At $\beta=0$, for every $n\ge5$ and coarse $g$, there is
a smooth conditional transport $\Xi_t:M_g\to c_t^{-1}(g)$,
$0\le t\le\tau$, taking the initial conditional Haar law
to the actual conditional law on the moving fibre. It
satisfies the volume-uniform bound
\[
 \Xi_t^*m_{\rm ambient}\le e^{\sigma t}m_0 .
\]
It is equivariant under the fibre-preserving internal
gauge action. The time $\tau$ and constant $\sigma$ are
the explicit finite constants above.
\end{theorem}
\begin{proof}
Finite-volume elliptic regularity and the centered
normalization make $\psi_t$, hence $V_t$, smooth in space
and time. Let $\vartheta_t$ be its flow starting at the
identity. The Poisson equation has the sign
\[
 \operatorname{div}_{\mu_g}(e^{u_t}V_t)
       =-a_te^{u_t}=-\partial_t e^{u_t}.
\]
Uniqueness of the smooth continuity equation on the compact
manifold therefore gives
$(\vartheta_t)_*\mu_g=\widehat\nu_t$.
Take $\Xi_t=\eta_t\circ i_g\circ\vartheta_t$.
The normal-chart identity $c_t\eta_t i_g=g$ puts its image
in the correct fibre and proves the probability assertion.
Pulling back \eqref{c18v61:strain} by $\vartheta_t$ and
integrating gives the metric inequality.

For a fibre-preserving gauge transformation, the base
metric and Haar law are invariant, as are $u_t,a_t$.
Uniqueness of the centered elliptic solution makes
$\psi_t$ invariant. Its gradient flow and the normal flow
are equivariant, proving the last assertion.
\end{proof}

As a consequence, the actual conditional Poincare gap on
this interval is at least $(3/2)e^{-\sigma t}$:
pull a function back by $\Xi_t$, use the product Haar gap
$3/2$, and use its differential bound in the energy.
A volume-uniform short-time conditional gap was already
obtained in V38. The new result here is the explicit
nonlinear probability correction and its actual material
strain, not the bare existence of another short-time gap.

\subsection{What this resolves, and the next upstream obstruction}

The Haar inverse of V60 is no longer substituted for the
actual density-weighted inverse: the drift is retained in
\eqref{c18v61:derivativePDE} and its effect is bounded by
\eqref{c18v61:entropy}--\eqref{c18v61:recurrence}.
Nor is the base metric substituted for the material
metric: the latter appears in
\eqref{c18v61:metricjet}--\eqref{c18v61:strain}.
These are genuine nonlinear finite-time statements, not
formal power-series remainders.

However, \eqref{c18v61:tau} is only a positive short
interval; the constants are extremely large. The bound
$H(u_t)\le tK_*$ does not verify $H(u_t)\le1/256$ at
$t=4$. Repeated auxiliary diffusion steps in
\eqref{c18v61:recurrence} do not extend Wilson block time.
Restarting a Wilson interval does not reset the accumulated
density to Haar or erase its interactions. Thus successive
applications with a fictitious fresh density would be
circular.

The next upstream target is a local response estimate
beyond this small-interaction neighborhood, for the native
density rather than for an arbitrary smooth potential.
One possible route is a quantitatively mixing one-factor
conditional reference retaining its diagonal potential,
with the off-diagonal response treated separately.
That route still owes a uniform conditional mixing bound
and a summable influence estimate. Additional finite jets
alone would not pay either obligation.

The diagnostic
\begin{center}\small
\path{scripts/verify_ym_c18_v61_anchored_decoupling.py}
\end{center}
checks the rational recurrence bounds, the noise/entropy
normalization, Bell and inverse-jet coefficients, and the
ordered derivative equation on noncommuting two-factor
group fixtures. Those fixtures are not native YM
trajectories. The analytic, all-volume claims rest on the
written proofs, not on finite symbolic tests or a formal
proof-assistant certificate.

V60 was the scheduled companion and broad-literature
checkpoint; the next such reviews remain V65 and V70.
This targeted Girsanov import does not replace those
reviews. The interacting-vacuum extension, strain through
time four, coarse return/variance, variable-bank stability,
weak-coupling scale transport, physical-time calibration,
and interacting continuum reconstruction remain unpaid.
A full positive physical four-dimensional Yang--Mills
mass gap is still unproved.

Removing this append and restoring the title recovers
V60 byte for byte. Only the active YM filename is replaced;
other sources are preserved. All compilation and visual
verification products stay outside \path{latex_notes},
which contains only \texttt{.tex} sources. The next append
is V62; archive/restart remains due at V100.
\addtocontents{toc}{\protect\setcounter{tocdepth}{2}}
\addtocontents{toc}{\protect\clearpage}
% END CYCLE 18 VERSION 61 APPEND
% BEGIN CYCLE 18 VERSION 62 APPEND
\clearpage
\section{V62: Conditional diagonal coercivity and a nonlinear transport obstruction}
\label{c18v62:context}

V61 constructs an exact probability correction near Haar. Its
smallness condition cannot be removed merely by replacing Haar's
one-factor gap with a positive conditional gap: the responses of
different coordinates still interact. This append identifies that
interaction in a compact-manifold Bochner identity and proves a
comparison obstruction adapted to V61's norms.

The one-coordinate gap itself is not a new missing lemma.
V31--V33 already proved such bounds for the actual fibre law in
inverse-Wilson coordinates, with its metric comparison retained.
We reuse their density-comparison argument in the normal chart,
without equating the two coordinate systems. The new analytic
distinction is between conditional diagonal coercivity, the
\emph{signed} mixed-Hessian coupling, and an absolute-value
matrix sufficient condition. A positive matrix condition is
not inferred from finite entries.

The second result is a smooth gauge-invariant probability path
on the same native product fibre. Every fixed-order unweighted
anchored density derivative, the source Hessian row norm, and
every one-coordinate inverse remain uniformly bounded, but the
integrated strain of any exact transport grows at least linearly
with volume. This is not the actual Wilson density. It shows why
the nonlinear continuation needs more than those particular
inputs. In particular, unlike V60's comparison, its temporal
source has a uniformly bounded Hessian row norm.

\subsection{Conditional diagonal coercivity on the compact product}

Use the product $(M_g,m_0,\mu_g)$ of V61 and temporarily suppress
$g$. Write $\nu=Z^{-1}e^u\mu$, with $u$ smooth. Let $\nu_i$
be its conditional law in factor $i$, with the other factors
fixed, and let
\[
 L=-\Delta_{m_0}-\langle\nabla u,\nabla\cdot\rangle,\qquad
 L_i=-\Delta_i-\langle\nabla_i u,\nabla_i\cdot\rangle.
\]
Suppose $\lambda_i>0$ is a lower bound on the scalar
Poincare gap of $\nu_i$, uniform in the exterior.
It need not be a lower bound on
$\operatorname{Ric}_i-\operatorname{Hess}_{ii}u$ pointwise.

For every smooth $\varphi$, the conditional integrated Bochner
identity and the scalar spectral theorem give
\begin{equation}
 \begin{split}
 &\int\bigl(\|\nabla_i d_i\varphi\|_{\rm HS}^2+
       \langle d_i\varphi,
       (\operatorname{Ric}_i-\operatorname{Hess}_{ii}u)
                         d_i\varphi\rangle\bigr)\,d\nu_i\\
 &\hspace{15mm}
       =\int(L_i\varphi)^2\,d\nu_i
       \ge\lambda_i\int|d_i\varphi|^2\,d\nu_i .
 \end{split}
 \tag{V62.1}\label{c18v62:conditionalBochner}
\end{equation}
Indeed the local operator is self-adjoint and nonnegative;
on each nonconstant spectral component, $\ell^2\ge\lambda_i\ell$.
The constant component contributes zero. This proves the last
inequality without convexity of the one-factor potential.
There is no boundary term on the group.

Let $a=L\psi$ and set $v_i=d_i\psi$. Product geometry and
integration by parts yield the directional identity
\begin{equation}
 \begin{split}
 \int\langle v_i,d_i a\rangle\,d\nu
  =\int\bigg[
   &\|\nabla_i v_i\|_{\rm HS}^2+
      \langle v_i,(\operatorname{Ric}_i
                          -\operatorname{Hess}_{ii}u)v_i\rangle\\
   &+\sum_{j\ne i}\|\nabla_jv_i\|_{\rm HS}^2
      -\sum_{j\ne i}
              \langle v_i,(\operatorname{Hess}_{ij}u)v_j\rangle
                         \bigg]\,d\nu .
 \end{split}
 \tag{V62.2}\label{c18v62:directional}
\end{equation}
One way to verify every term is to differentiate the scalar
equation covariantly:
$d_iL\psi$ is the weighted rough Laplacian on $d_i\psi$,
plus $\operatorname{Ric}_i d_i\psi$, minus
$\sum_j(\operatorname{Hess}_{ij}u)d_j\psi$.
Taking its inner product with $v_i$ gives
\eqref{c18v62:directional}. The Ricci term is
$2\kappa_i I$, not zero. Terms with different factors have no
mixed curvature because the metric is a product.
Integrating \eqref{c18v62:conditionalBochner} in the exterior
controls the entire first line of \eqref{c18v62:directional}.

This use of scalar conditional gaps is valid here because
$v_i=d_i\psi$ is exact in its own factor. It does not prove
the same diagonal floor on arbitrary one-forms. In particular,
no unstated one-form spectral gap is substituted for a
scalar conditional gap.

\subsection{Two sufficient coupling tests, with different losses}

Define the symmetric tangent-bundle matrix $\mathsf S_u(x)$ by
\[
 (\mathsf S_u)_{ii}=\lambda_i I,\qquad
 (\mathsf S_u)_{ij}=-\operatorname{Hess}_{ij}u(x)\quad(i\ne j).
\]
The entries are evaluated at the same configuration.
Summing \eqref{c18v62:directional}, retaining the mixed
derivative squares as nonnegative terms, gives
\begin{equation}
 \|L\psi\|_{L^2(\nu)}^2
       \ge\int\langle d\psi,\mathsf S_u(x)d\psi\rangle\,d\nu .
 \tag{V62.3}\label{c18v62:signed}
\end{equation}
Consequently the sufficient condition
\begin{equation}
       \mathsf S_u(x)\succeq\rho I
                \quad\hbox{for every }x,\qquad \rho>0,
 \tag{V62.4}\label{c18v62:signedtest}
\end{equation}
implies $\operatorname{gap}(L)\ge\rho$.
To prove this last assertion, test \eqref{c18v62:signed} on
each nonconstant eigenfunction of $L$: its eigenvalue
$\ell>0$ satisfies $\ell^2\ge\rho\ell$.
Compact ellipticity supplies a complete discrete spectral
resolution at each finite volume.
This test can retain matrix signs; it does not require
replacing each mixed block by its absolute norm.

There is also a directional, but potentially weaker, test.
Choose numbers
\[
 k_{ij}=k_{ji}\ge
   \sup_x\|\operatorname{Hess}_{ij}u(x)\|_{\rm op}
           \quad(i\ne j),\qquad
 \mathsf A_{ii}=\lambda_i,\quad\mathsf A_{ij}=-k_{ij}.
\]
The supremum is now taken \emph{entry by entry}. Thus a bound
on $\sup_x\sum_j\|\operatorname{Hess}_{ij}u(x)\|$ does not,
by itself, bound $\sum_j k_{ij}$.

\begin{proposition}[Compact-product directional response]
\label{c18v62:response}
If $\mathsf A$ is positive definite, its inverse is
entrywise nonnegative. For every centered smooth $a$ and
$\psi=L^{-1}a$, put
\[
 r_i=\|d_i\psi\|_{L^2(\nu)},\qquad
 b_i=\|d_i a\|_{L^2(\nu)} .
\]
Then
\begin{equation}
 r\le\mathsf A^{-1}b\quad\hbox{entrywise},\qquad
 |\operatorname{Cov}_\nu(a,h)|
       \le\sum_{i,j}(\mathsf A^{-1})_{ij}
                   b_j\|d_i h\|_{L^2(\nu)}.
 \tag{V62.5}\label{c18v62:responsebound}
\end{equation}
In particular $\operatorname{gap}(L)\ge\lambda_{\min}(\mathsf A)$.
\end{proposition}
\begin{proof}
Equations \eqref{c18v62:conditionalBochner}--\eqref{c18v62:directional}
and Cauchy--Schwarz imply
$\lambda_i r_i^2\le b_i r_i+\sum_{j\ne i}k_{ij}r_i r_j$.
Division when $r_i>0$ gives $\mathsf A r\le b$.
When $r_i=0$ that row inequality holds directly.

Choose $c\ge\max_i\lambda_i$. The matrix
$\mathsf B=cI-\mathsf A$ is entrywise nonnegative.
Its spectral radius equals its largest eigenvalue,
$c-\lambda_{\min}(\mathsf A)<c$. Hence
$\mathsf A^{-1}=c^{-1}\sum_{m\ge0}(\mathsf B/c)^m\ge0$.
Multiplication proves the first bound.
Finally $\operatorname{Cov}_\nu(a,h)=\int\langle d\psi,dh\rangle\,d\nu$
proves the covariance bound. Taking $h=a$ and using the
operator norm of $\mathsf A^{-1}$ proves the gap assertion.
\end{proof}

The directional covariance method is an established one;
see G.~Menz, \emph{A Brascamp--Lieb type covariance estimate},
Theorems 2.3 and 2.7,
\url{https://arxiv.org/html/1402.5160}.
The import here is its diagonal-versus-mixed-response
organization. Equations \eqref{c18v62:conditionalBochner}
and \eqref{c18v62:directional} spell out the compact
curved-product implementation, including the Ricci and
exact-gradient qualifications. Neither the cited Euclidean
result nor our implementation verifies positivity of the
native YM coupling matrix.

\subsection{Application to the normal-chart law, without a false closure}

At $\beta=0$ and $0\le t\le4$, V60--V61 give
$H(u_t)\le tK_*$ and $G(u_t)\le DtK_*$, where $D=\pi\sqrt2$.
A one-factor geodesic therefore gives
$\operatorname{osc}_i u_t\le D^2tK_*$.
The already used Haar density-comparison argument allows the
following choice of normal-chart conditional floor:
\begin{equation}
 \lambda_i(t)= 3\kappa_i e^{-D^2tK_*}
               \ge\lambda_*(t):=\frac32e^{-D^2tK_*}.
 \tag{V62.6}\label{c18v62:nativefloor}
\end{equation}
This merely instantiates the old local-gap argument in the
current coordinates. In particular it is not a new proof
of a full-fibre gap at time four.

For this law the signed test uses exactly
$-\operatorname{Hess}_{ij}u_t$, with no vacuum or metric
replacement. A sufficient, cruder consequence of the paid
row bound is
\[
 \mathsf S_{u_t}(x)\succeq
       \bigl(\lambda_*(t)-tK_*\bigr)I .
\]
This again helps only while its right side is positive.
The finite upper bound $tK_*$ does not become smaller
than $\lambda_*(t)$ by being independent of volume.

If either coupling test is verified with a uniform
positive floor $\rho$, then the actual metric comparison
of V61 gives
\begin{equation}
 \widehat L_t\succeq E_0^{-2}\mathcal L_{u_t},
 \qquad \operatorname{gap}(\widehat L_t)\ge E_0^{-2}\rho.
 \tag{V62.7}\label{c18v62:actualmetric}
\end{equation}
Both forms use the actual probability $\widehat\nu_t$.
For a centered source $a$, the directional test would
also give
\[
 \langle a,\widehat L_t^{-1}a\rangle_{\widehat\nu_t}
 \le E_0^2\langle a,\mathcal L_{u_t}^{-1}a\rangle_{\widehat\nu_t}
 =E_0^2\sum_i r_i^2
 \le E_0^2\|\mathsf A^{-1}b\|_2^2 .
\]
The inverse inequality follows from the variational
formula on centered functions. It is an integrated cost,
not a pointwise Hessian or material-strain bound.
V60 already distinguishes these conclusions.

The diagonal potential has now been absorbed
nonperturbatively into a scalar conditional gap, rather
than discarded. What remains is a native signed coupling
estimate or an equally strong alternative. No positivity
of $\mathsf S_{u_4}$ or $\mathsf A$ is asserted here.

\subsection{A comparison path with all unweighted local budgets bounded}

Use V60's $M=n^3$ vertex-disjoint uncaptured plaquettes,
based at $2z+e_3$ in the $12$ plane. Their four edges are
free fibre factors. Let $w_r$ be the normalized real
holonomy trace of plaquette $r$ and put
\[
 W=\sum_{r=1}^M w_r,\qquad m=W/M .
\]
The disjointness proof is already in V60 and is not
repeated. Under $\mu_g$ the $w_r$ are independent, with
the same distribution as the real coordinate of Haar
$\mathrm{SU}(2)$:
\[
 d\gamma(w)=\frac2\pi\sqrt{1-w^2}\,dw,\qquad -1\le w\le1.
\]
Also $\Gamma(w_r)=4(1-w_r^2)$ in the product metric.
These statements are independent of $g$.

For $0\le s\le4$, define a genuine probability path
\begin{equation}
 U_s=\frac{8s}{M}W^2,\qquad
 Z_s=\int e^{U_s}\,d\mu_g,\qquad
 \nu_s^{\rm cmp}=Z_s^{-1}e^{U_s}\mu_g .
 \tag{V62.8}\label{c18v62:path}
\end{equation}
Write $u_s^{\rm cmp}=U_s-\log Z_s$. Its actual temporal
source, not an independently chosen forcing, is
\begin{equation}
 a_s^{\rm cmp}=\partial_su_s^{\rm cmp}
      =\frac8M W^2-\nu_s^{\rm cmp}\left(\frac8M W^2\right).
 \tag{V62.9}\label{c18v62:source}
\end{equation}
Both the law and the source are invariant under the
internal gauge group, since each $w_r$ is a closed-loop
trace. The comparison starts at Haar.
Its time $s$ is a parameter of this comparison, not a
statement about the Wilson normal flow.

\begin{lemma}[Uniform unweighted anchored jets and conditional gaps]
\label{c18v62:comparisoninputs}
For every fixed integer $r\ge1$, the invariant ordered
derivatives satisfy
\begin{align}
 \sup_{0\le s\le4}\mathfrak a_0(\partial^r u_s^{\rm cmp})
          &\le32\,2^r12^{r-1},\nonumber\\
 \sup_{0\le s\le4}\mathfrak a_0(\partial^r a_s^{\rm cmp})
          &\le 8\,2^r12^{r-1}.
 \tag{V62.10}\label{c18v62:comparisonjets}
\end{align}
The norms include the configuration supremum.
In particular, uniformly in $M,g,s$,
\begin{equation}
 \begin{split}
 G(u_s^{\rm cmp})&\le64,\quad J(u_s^{\rm cmp})\le1536,\quad
 H(u_s^{\rm cmp})\le4672,\\
 G(a_s^{\rm cmp})&\le16,\quad J(a_s^{\rm cmp})\le384,\quad
 H(a_s^{\rm cmp})\le1168 .
 \end{split}
 \tag{V62.11}\label{c18v62:comparisonnorms}
\end{equation}
Every one-factor conditional has Poincare gap at least
$3e^{-128}$ throughout the path.
\end{lemma}
\begin{proof}
A nonzero ordered derivative of $w_r$ uses only its twelve
scalar edge directions. Quaternion multiplication by a unit
imaginary generator has operator norm one, so every such
derivative has absolute value at most one. Fixing a scalar
slot fixes its plaquette. Consequently, for $p\ge1$,
\[
 \mathfrak a_0(\partial^p W)\le12^{p-1},\qquad
 \sum_{\text{all }p\text{ slots}}|\partial^pW|\le M12^p,
 \qquad |W|\le M .
\]
These are unweighted bounds. Apply the ordered product rule
to $cW^2/M$. The two terms with an undifferentiated $W$
contribute at most $2c12^{r-1}$. In each remaining term,
the factor containing the fixed slot is bounded in its
anchored norm, and the other factor in its full summed
norm. Their product is at most $M12^{r-1}$.
The sum of all product multiplicities is $2^r$.
Taking $c\le32$ and $c=8$ proves
\eqref{c18v62:comparisonjets}; scalar normalizers disappear.

On a participating free edge,
$\nabla_e(cW^2/M)=2c\,m\,\nabla_e w_r$, with
$|\nabla_e w_r|\le1$. Thus its component gradient is at most
$2c$; all other factors have zero gradient.
For any scalar $f$, conversion of the ordered matrix into
Hessian blocks gives $H(f)\le3J(f)+G(f)$: sum the absolute
scalar entries over each block row, then use the
block-diagonal cross-product connection term of norm at
most $G(f)$. Together with the $r=2$ bound this gives
\eqref{c18v62:comparisonnorms}.

When one participating edge varies, only one $w_r$ varies.
The derivative of $c(S+w)^2/M$ in $w\in[-1,1]$ is at
most $2c$, since $|S+w|\le M$. Its oscillation is therefore
at most $4c\le128$. The conditional gap is at least
$3e^{-128}$ by the free-factor Haar comparison.
Unused factors stay Haar; pair factors have gap $3/2$.
This proves the common floor. It is positive but not
claimed numerically useful.
\end{proof}

\subsection{Exponential full-gap collapse at the endpoint}

\begin{theorem}[A collective bottleneck despite the paid local inputs]
\label{c18v62:bottleneck}
For the endpoint measure in \eqref{c18v62:path},
\begin{equation}
 \operatorname{gap}(-\operatorname{div}_{\nu_4^{\rm cmp}}
                                  \nabla_{m_0})
       \le \frac{128}{M}e^{-3M}.
 \tag{V62.12}\label{c18v62:smallgap}
\end{equation}
The same upper bound applies after restricting to
gauge-invariant test functions.
\end{theorem}
\begin{proof}
The elementary event $w\in[1/2,3/4]$ has Haar probability
\[
 \int_{1/2}^{3/4}\frac2\pi\sqrt{1-w^2}\,dw
       \ge\frac{\sqrt7}{8\pi}>\frac1{16}.
\]
For the last strict inequality, $\pi<4<2\sqrt7$ suffices.
If every $w_r\ge1/2$, then $m\ge1/2$. Since
$U_4=32Mm^2$, independence gives
\[
 Z_4\ge16^{-M}e^{8M}.
\]
On the central strip $\{|m|\le1/4\}$, $U_4\le2M$.
Its unnormalized probability is at most $e^{2M}$, and hence
\begin{equation}
 \nu_4^{\rm cmp}(|m|\le1/4)
       \le e^{-(6-\log16)M}<e^{-3M}.
 \tag{V62.13}\label{c18v62:strip}
\end{equation}
Here $\log16<3$, for example by V60's elementary
$e^3>20$ bound.

The measure is invariant under simultaneous
$w_r\mapsto-w_r$: multiply one free link of each disjoint
plaquette by the central element $-I$. This preserves
Haar, the fixed sampled coarse configuration, and $W^2$.
Let $F=\operatorname{clip}(4m,-1,1)$. It is odd under this
symmetry, so its mean is zero, and
\[
 \operatorname{Var}_{\nu_4^{\rm cmp}}F
      \ge1-e^{-3M}>\frac12 .
\]
Moreover $\Gamma(m)=4M^{-2}\sum_r(1-w_r^2)\le4/M$.
The weak derivative of the clipping function is zero off
the central strip and has magnitude four inside. Thus
\[
 \int|\nabla F|^2\,d\nu_4^{\rm cmp}
       \le\frac{64}{M}e^{-3M}.
\]
This Lipschitz function belongs to the form domain; smooth
approximations converge in that domain on the compact
product with smooth positive density. Its Rayleigh
quotient proves \eqref{c18v62:smallgap}.
Since $F$ is a function of gauge-invariant traces, it also
proves the restricted assertion.
\end{proof}

There is no deteriorating one-coordinate well parameter
in this example: the time interval and coefficient $32$
are fixed as $M$ grows. The exponentially small gap is
collective. This is distinct from the earlier
one-coordinate double-well comparisons.

\subsection{The same path forces large exact probability strain}

Consider any smooth flow $\vartheta_s$ on $M_g$ starting
at the identity and satisfying
$(\vartheta_s)_*\mu_g=\nu_s^{\rm cmp}$.
Suppose its velocity $V_s$ has an upper material-strain
bound
\[
 {\cal L}_{V_s}m_0\preceq\sigma_M(s)m_0,\qquad
 S_M=\int_0^4\sigma_M(s)\,ds<\infty .
\]
The metric here is fixed; the time dependence is all in
the probability. Pullback and Gronwall imply
$\vartheta_4^*m_0\preceq e^{S_M}m_0$.
Haar's product gap $3/2$ therefore gives an endpoint gap
at least $(3/2)e^{-S_M}$.
Together with \eqref{c18v62:smallgap}, this proves
\begin{equation}
           S_M\ge 3M+\log M+\log(3/256).
 \tag{V62.14}\label{c18v62:strainlower}
\end{equation}
Thus no volume-uniform integrated strain bound can hold
for this probability path, for \emph{any} choice of exact
transport velocity.

In particular, solve the actual centered Poisson equation
\[
 -\operatorname{div}_{\nu_s^{\rm cmp}}\nabla_{m_0}\psi_s
             =a_s^{\rm cmp}.
\]
Finite-volume compact elliptic theory gives a smooth family.
Its gradient flow transports the law by the same continuity
equation argument as V61. Since
$\mathcal L_{\nabla\psi_s}m_0=2\operatorname{Hess}_{m_0}\psi_s$,
\begin{equation}
 2\int_0^4\sup_x\lambda_{\max}
                 (\operatorname{Hess}_{m_0}\psi_s(x))\,ds
       \ge3M+\log M+\log(3/256).
 \tag{V62.15}\label{c18v62:Poissonstrain}
\end{equation}
No claim that the source has overlap with the slowest odd
spectral mode was used. The argument is a transport
distortion bound for the actual even temporal source
\eqref{c18v62:source}, not an inverse-norm claim for an
arbitrary forcing.

The conditional coupling test detects an obstruction in
this comparison. Choose a configuration with all selected
$w_r=0$, and write $p_e=\nabla_e w_r$ for each of its $4M$
free edges. Then $|p_e|=1$. For distinct such factors,
\[
 (\operatorname{Hess}_{ef}U_4)
                   =\frac{64}{M}p_e\otimes p_f .
\]
The same formula holds within one plaquette because the
term proportional to $W$ vanishes. At an exterior with
all other $w_r=0$, the conditional distribution of the
varying trace is proportional to $e^{32w^2/M}d\gamma(w)$.
It has mean zero and second moment at least $1/4$:
the derivative of $\mathbf E_\theta(w^2)$ for the tilt
$e^{\theta w^2}d\gamma$ is
$\operatorname{Var}_\theta(w^2)\ge0$, and its value at
$\theta=0$ is $1/4$. The Rayleigh quotient of
$w$ for this one-edge law is
\[
 \frac{1-\mathbf E(w^2)}{\mathbf E(w^2)}\le3.
\]
Consequently every valid exterior-uniform conditional
floor for a participating edge is at most three.
Testing $\mathsf S_{U_4}$ on the block vector with
components $p_e$, zero on other factors, gives
\begin{equation}
 \frac{\langle p,\mathsf S_{U_4}p\rangle}{|p|^2}
       \le 3-\frac{64}{M}(4M-1)<0 .
 \tag{V62.16}\label{c18v62:matrixfailure}
\end{equation}
The absolute matrix test cannot improve this value.
This failure is consistent with the collective
bottleneck; it is not evidence that the signed test is
necessary for every gapped probability law.

\subsection{What is and is not transferred to YM}

The comparison in \eqref{c18v62:path} obeys the compact
native geometry, gauge invariance, initial Haar law,
uniform one-factor floors, and every fixed-order
\emph{unweighted} anchored derivative bound.
Its true temporal source also has the bounded Hessian
row norm absent in V60's comparison. Nonetheless
\eqref{c18v62:Poissonstrain} rules out extending V61 on the
strength of these properties alone.

It does \emph{not} obey a fixed positive spatial
exponential weight uniformly in volume. At the
configuration used in \eqref{c18v62:matrixfailure}, mixed
entries between distant selected plaquettes are of order
$1/M$, regardless of their separation. For plaquettes a
distance proportional to $n$ apart, even a single entry
weighted by $e^{\zeta\,\mathrm{distance}}$ is unbounded
as $M=n^3\to\infty$, for every fixed $\zeta>0$.
One may realize the nonzero scalar components by setting
one edge on each plaquette to a fixed imaginary unit
quaternion and the other three to the identity.
Thus this comparison does not contradict the stronger
weighted native jet estimates, nor does it establish a
transition in the actual Wilson conditional law.

The resulting direction is sharper: retain the native
spatial structure and estimate its \emph{signed mixed
response relative to the conditional diagonal energies},
or supply a different native phase-exclusion argument.
Replacing all this structure by the unweighted constants
in V61 loses information that may be decisive.
Conversely, finite weighted locality alone has not been
proved here to imply a uniform coupling margin.
We must actually evaluate the native signed form or a
successful multiscale substitute, rather than name such
a margin as an assumption and count it as paid.

The diagnostic
\begin{center}\small
\path{scripts/verify_ym_c18_v62_conditional_coupling.py}
\end{center}
checks the Bochner curvature and sign on exact sphere
fixtures, noncommuting directional identities at Haar,
the ordered product multiplicities, bottleneck exponents,
and the matrix test. The finite fixtures do not certify
a YM trajectory or a continuum theorem; the uniform
comparison statements are proved analytically above.

The next companion review remains V65 and the next broad
literature sweep V70; this was a targeted primary-source
import, not a replacement for those checkpoints.
The native time-four coupling margin and pointwise
probability strain remain unproved. Interacting-vacuum
control, coarse return/variance, weak-coupling scale
transport, physical-time calibration and interacting
continuum reconstruction are still outstanding.
There is still no proof of the full four-dimensional
physical Yang--Mills mass gap.

Only the active filename changes from V61 to V62.
Removing this append and restoring the title recovers
V61 byte for byte. All other sources are preserved,
and \path{latex_notes} contains only \texttt{.tex}
files. The PDF skill is used for internal layout checks;
its build products remain outside that folder.
Archive/restart remains due at V100.
\addtocontents{toc}{\protect\setcounter{tocdepth}{2}}
\addtocontents{toc}{\protect\clearpage}
% END CYCLE 18 VERSION 62 APPEND
% BEGIN CYCLE 18 VERSION 63 APPEND
\clearpage
\section[V63: native spatial tails and approximate boundary separation]{V63: native spatial tails and approximate boundary separation}
\label{c18v63:context}

V62 separates conditional diagonal coercivity from the signed
mixed response that still has to be controlled. Its comparison
path rules out a continuation based only on unweighted local
norms, but does not have the native spatial exponential bounds.
We now retain that spatial information. At electric coupling
$\beta=0$, the actual normal-chart density has uniformly
exponentially summable Hessian rows throughout Wilson smoothing
time $0\le t\le4$. This gives a quantitative approximate boundary
separation for its \emph{normalized} block laws, with an error
independent of the number of distant exterior factors.

These are bounds for the existing Wilson construction, not for a
new comparison density. They do not establish positivity of V62's
near-field signed matrix. In particular, a conditional boundary
estimate with the intervening collar pinned is not a mixing
estimate after that collar is integrated out.

\subsection{Companion transfer, exact scope, and non-duplication}

The newly supplied files were compared with the last review by
SHA-256. ESM general V44, Einstein Bridge V55, Stationarity V61,
and Source Selection V65 are byte-identical to the previously
reviewed copies; their terminal scope was rechecked rather than
described as a new full audit. The new ESM phenomenology V41
and Arithmetic V42 terminal developments were read. Their
same-record incidence and classical/coherent distinctions are
useful safeguards, but neither supplies a native conditional
coercivity constant for the YM flow. The general ESM statements
remain fixed-order perturbative statements. The initialized
stationarity and bridge results retain their unpaid source or
tangency rows; the source-selection result retains its actual
occurrence requirement.

The Grand Tensor V079 was reviewed in its relevant YM boundary,
source-normalization, and continuum-scope sections, not certified
line by line. Its \texttt{YM-Wilson-blanket} result proves exact
conditional independence for the finite-plaquette Wilson law in
the original link coordinates. That exact separation is not
silently transferred to our flowed normal-chart law. Here we
prove the quantitative replacement from the latter law's actual
mixed derivatives, and retain its boundary-mediated covariance.
This is the substantive change of direction prompted by this
companion review.

For non-duplication, f16.2 V31 already constructs a summable
interaction and proves normalized conditional truncation and the
oscillation-to-TV inequality on its restricted retained-bridge
window. We reuse that elementary normalization inequality, not
its unproved extension to another model. No new interaction
reconstruction theorem is claimed. Current V42 supplies weighted
contractions and V60 the native third normal-field jet; the new
step is retaining their positive weight through the pullback and
extracting unrestricted compact-fibre boundary and score bounds.

For primary context, the abstract of
\href{https://arxiv.org/abs/2001.03880}{Barbieri--G\'omez--Marcus--Meyerovitch--Taati,
\emph{Gibbsian representations of continuous specifications}}
was rechecked. Its symbolic-configuration representation results
are not invoked for $\mathrm{SU}(2)$ variables.
V62's \href{https://arxiv.org/html/1402.5160}{Menz covariance
reference} remains the context for its separate positive
coupling-matrix requirement. The present proofs do not import
a positivity margin from either paper. This targeted review
does not replace the next broad sweep at V70. The next scheduled
companion checkpoint remains V65.

\subsection{The positive spatial weight survives the native pullback}

Keep $N=2n$, $n\ge5$, a fixed arbitrary coarse configuration $g$,
the compact product $M_g=b_2^{-1}(g)$ and its pair/free metric
$m_0$. The pair factors have metric twice the unit-round metric;
the free factors have unit-round metric. Thus every factor has
diameter at most
\[
                         D=\pi\sqrt2.
\]
Give each free factor its fine-link tail as anchor. Give each
pair factor the tail of its first captured half-link. The
second half-link has distance one from that anchor.
Use periodic fine-lattice Manhattan distance $d(i,j)$ between
anchors. At most six factor labels have one anchor (a harmless
overcount of at most three free and three pair labels).

All $v_j$, $K_1,K_2,A_1,A_2,B_2,K_*$ below are exactly those of
\eqref{c18v60:v} and \eqref{c18v60:budget}, not the enlarged
fourth-jet constants of V61. Keep the strictly positive
$\zeta$ from \eqref{c18v44:Zconstants}. Define
\[
 \begin{split}
 \mathcal H_\zeta(h)
   &=\sup_{x\in M_g}\max_i\sum_j e^{\zeta d(i,j)}
            \|(\operatorname{Hess}_{m_0}h)_{ij}(x)\|_{\rm op},\\
 K_\zeta&=12e^{2\zeta}B_2=e^{2\zeta}K_* .
 \end{split}
 \tag{V63.1}\label{c18v63:Hweight}
\]
The supremum is still \emph{outside} the row sum. All statements
are uniform in $g$ but compare configurations on the same
fixed fibre; no coarse derivative is asserted.

\begin{theorem}[Native weighted density and temporal-score Hessians]
\label{c18v63:weighted}
For the actual electric normal-chart probability
$d\widehat\nu_t=\widehat p_t\,d\mu_g$, put
$u_t=\log\widehat p_t$ and $a_t=\partial_tu_t$. Then
\[
 \boxed{\quad
       \mathcal H_\zeta(a_t)\le K_\zeta,\qquad
       \mathcal H_\zeta(u_t)\le tK_\zeta
                   \quad(0\le t\le4).\quad}
 \tag{V63.2}\label{c18v63:native}
\]
No conditional gap, typical-set restriction, or smallness of
$tK_\zeta$ is assumed.
\end{theorem}

\begin{proof}
V60 proves
$\mathfrak a_\zeta(\partial^j\mathcal N_t)\le v_j$
for $0\le j\le3$ on the complete ambient fine product.
Let $d_t=\sum_\alpha X_\alpha\mathcal N_{t,\alpha}$.
The differentiated trace retains an external derivative anchor.
V42's diagonal-trace lemma therefore gives, now with the weight
retained,
\[
 \sup_j|X_jd_t|\le v_2,\qquad
              \mathfrak a_\zeta(\partial^2d_t)\le v_3.
 \tag{V63.3}\label{c18v63:trace}
\]
It makes no assertion about the extensive undifferentiated
trace.

The coefficient generator of $D\eta_t$ is the ordered derivative
of $-\mathcal N_t$ plus the moving-frame bracket. A bracket
coefficient connects only colors on the same link, so it costs
no spatial weight. Its norm and first coefficient derivative
are bounded by $K_1=v_1+4v_0$ and $K_2=v_2+4v_1$ in the
same weighted norms. The connected variational equations and
their initial data $D\eta_0=I$, $\partial D\eta_0=0$ give
\[
 \mathfrak a_\zeta(D\eta_t)\le A_1,\qquad
 \mathfrak a_\zeta(\partial D\eta_t)\le A_2.
 \tag{V63.4}\label{c18v63:flow}
\]
For the second inequality its forcing is bounded by
$K_2 A_1^2$ and its propagator by $A_1$, giving
$TK_2A_1^3=A_2$. Configuration composition changes coefficients,
not the assigned lattice positions; the coefficient bounds
already take their supremum over all configurations.

The scalar pullback rule gives gradient bound $v_2A_1$ and
ordered second-derivative weighted Schur bound
$v_3A_1^2+v_2A_2$ for $d_t\circ\eta_t$.
Subtracting the invariant-frame connection adds at most
$2v_2A_1$. This term too has all its color indices on one link.
Consequently the ambient covariant Hessian has weighted Schur
norm at most $B_2$.

The inclusion $i_g:M_g\longrightarrow\mathcal U_N$ is totally
geodesic. Its scalar derivative has unweighted row and column
sums at most two, as in V60. Every nonzero entry connects
anchors at distance at most one, hence
\[
                    \mathfrak a_\zeta(Di_g)\le2e^\zeta.
\]
The covariant restriction formula has no second fundamental
form term. Contracting with $Di_g$ on its two remaining slots
costs at most $4e^{2\zeta}$. Conversion of scalar rows into
three-by-three operator-block rows costs at most three.
This proves the constant $12e^{2\zeta}B_2$.

Finally the actual coarea identity is
\[
 a_t=-d_t\circ\eta_t\circ i_g
             +\mathbb E_{\widehat\nu_t}
                         (d_t\circ\eta_t\circ i_g).
 \tag{V63.5}\label{c18v63:score}
\]
The last term is spatially constant on this fibre.
Thus it contributes no Hessian. Since $u_0=0$ and
$\partial_tu_t=a_t$ in the fixed chart, time integration proves
the density estimate. The normalized probability and its
normalizer have not been replaced.
\end{proof}

\subsection{Separate entry suprema and a paid signed spatial tail}

Write
$k_{ij}(t)=\sup_x\|(\operatorname{Hess}_{m_0}u_t)_{ij}(x)\|_{\rm op}$.
For $0\le\xi<\zeta$ set
\[
 Q(s)=\left(\frac{1+e^{-s}}{1-e^{-s}}\right)^3\quad(s>0).
\]
Then
\[
 k_{ij}(t)\le tK_\zeta e^{-\zeta d(i,j)},\qquad
 \sup_i\sum_j e^{\xi d(i,j)}k_{ij}(t)
                 \le6tK_\zeta Q(\zeta-\xi).
 \tag{V63.6}\label{c18v63:envelope}
\]
Indeed each entry follows from \eqref{c18v63:native} before
taking its separate supremum. Representatives of a periodic
coordinate displacement can be chosen injectively in
$\mathbb Z^3$, preserving its Manhattan length. The full
integer-lattice sum of $e^{-s|x|_1}$ is $Q(s)$.
The anchor multiplicity proves the displayed row bound.
This is a justified interchange at a \emph{smaller} weight,
not the invalid assertion
$\sum_j\sup_x\|H_{ij}(x)\|\le\sup_x\sum_j\|H_{ij}(x)\|$.

There is a sharper tail statement for V62's pointwise signed
matrix $S_t(x)$. Let $S_t^{[R]}(x)$ keep its same diagonal
conditional-gap floors and keep its off-diagonal blocks only
when $d(i,j)<R$, where $R>0$. Then
\[
 \boxed{\quad
 \sup_x\|S_t(x)-S_t^{[R]}(x)\|_{\rm op}
                 \le tK_\zeta e^{-\zeta R}.\quad}
 \tag{V63.7}\label{c18v63:signedtail}
\]
The difference is a symmetric block matrix, and its block
row sums are bounded directly by the weighted pointwise row
bound. The Schur test proves \eqref{c18v63:signedtail} without
a lattice-count factor.

Consequently a \emph{verified} inequality
$S_t^{[R]}(x)\succeq\rho_R I$ for all $x$ would give
$S_t(x)\succeq(\rho_R-tK_\zeta e^{-\zeta R})I$.
The tail is now paid, but no positive $\rho_R$ is assigned.
The retained entries still belong to the original global
normal-chart density. Cutting matrix entries does not identify
them with entries computed on a smaller periodic lattice, nor
establish a finite-window spectral certificate.

\subsection{Mixed finite differences without an exterior-volume factor}

For disjoint factor sets $A,C$ define
$d(A,C)=\min_{i\in A,j\in C}d(i,j)$, with the empty-set case
interpreted as infinite distance. Let all remaining coordinates
be fixed. For a smooth scalar $h$, suppose
$\mathcal H_\zeta(h)\le K_h$. If $x_A,y_A$ are two interior
values and $z_C,z'_C$ two exterior values, then
\[
 \begin{split}
 \big|h(x_A,z_C)-h(y_A,z_C)
       -h(x_A,z'_C)+h(y_A,z'_C)\big|
       \le |A|D^2K_h e^{-\zeta d(A,C)}.
 \end{split}
 \tag{V63.8}\label{c18v63:rectangle}
\]
To prove this, join each changed coordinate by a minimizing
constant-speed factor geodesic on $[0,1]$, of speed at most
$D$. Change all coordinates of $A$ simultaneously along one
parameter and all of $C$ along another. Because the two sets
are disjoint, the mixed derivative of the resulting scalar
rectangle is exactly the sum of its cross Hessian blocks
paired with these velocities. At every point its absolute
value is at most
\[
 D^2\sum_{i\in A}\sum_{j\in C}\|(\operatorname{Hess}h)_{ij}\|
       \le |A|D^2K_h e^{-\zeta d(A,C)}.
\]
Double integration proves the claim. It is valid at cut-locus
endpoints since a minimizing geodesic can be chosen; no
differentiation of its choice is used.

In particular, changing \emph{all} far coordinates together
does not require summing independent configuration suprema.
For $h=u_t$ take $K_h=tK_\zeta$; for $h=a_t$ take
$K_h=K_\zeta$.

\subsection{Normalized block conditionals and their temporal scores}

For a nonempty set $A$ and an exterior value $z$ define the
everywhere positive smooth conditional version
\[
 d\nu_{t,A}^{\,z}(x_A)
  =\frac{e^{u_t(x_A,z)}\,d\mu_A(x_A)}
         {\int e^{u_t(y_A,z)}\,d\mu_A(y_A)}.
 \tag{V63.9}\label{c18v63:kernel}
\]
These are conditionals of the \emph{actual} $\widehat\nu_t$,
not conditionals of a truncated Wilson action. Let $z,z'$
agree on every exterior factor of distance less than $R$
from $A$, and put
\[
 \varepsilon_A(R)=|A|D^2K_\zeta e^{-\zeta R},
 \qquad \delta_{t,A}(R)=t\varepsilon_A(R),\qquad
 \theta_{t,A}(R)=\tanh\!\left(\frac{\delta_{t,A}(R)}4\right).
 \tag{V63.10}\label{c18v63:errors}
\]

\begin{theorem}[Full-time electric conditional boundary control]
\label{c18v63:conditionals}
Uniformly over these exterior values and $0\le t\le4$,
\[
 \begin{gathered}
 \operatorname{osc}_{x_A}
                [u_t(x_A,z)-u_t(x_A,z')]\le\delta_{t,A}(R),\\
 e^{-\delta_{t,A}(R)}
      \le\frac{d\nu_{t,A}^{\,z}}{d\nu_{t,A}^{\,z'}}\le
                                     e^{\delta_{t,A}(R)},\\
 d_{\rm TV}(\nu_{t,A}^{\,z},\nu_{t,A}^{\,z'})
                                    \le\theta_{t,A}(R).
 \end{gathered}
 \tag{V63.11}\label{c18v63:TV}
\]
The bound depends on $|A|$, but not on the size of the changed
exterior, the ambient volume, or $g$.
\end{theorem}

\begin{proof}
Apply \eqref{c18v63:rectangle} to the exterior coordinates
which differ. If the potential difference has minimum $m$
and maximum $M$, its normalizing integral lies between
$e^m$ and $e^M$ against the second conditional probability.
This proves the likelihood-ratio bounds.

For completeness, the sharp elementary TV conversion already
used in f16.2 V31 retains both normalizers. If a probability
likelihood ratio $r$ ranges from $l$ to $L$ with
$\mathbb E r=1$ and $L/l\le e^\delta$, its positive-part
chord bound gives
\[
 \tfrac12\mathbb E|r-1|
      \le\frac{(L-1)(1-l)}{L-l}
      \le\frac{e^{\delta/2}-1}{e^{\delta/2}+1}.
\]
The last expression follows by maximizing over $l,L$ at the
fixed ratio bound. The constant-ratio case is zero.
This proves the last line of \eqref{c18v63:TV}.
\end{proof}

Let $b_{t,A}^{\,z}=\partial_t\log(d\nu_{t,A}^{\,z}/d\mu_A)$,
where the exterior is held fixed during time differentiation.
Then
\[
 b_{t,A}^{\,z}(x_A)
       =a_t(x_A,z)-\mathbb E_{\nu_{t,A}^{\,z}}a_t(\cdot,z).
 \tag{V63.12}\label{c18v63:localscore}
\]
In particular these are the actual centered conditional scores.

\begin{proposition}[Locality of the normalized temporal source]
\label{c18v63:scorelocal}
For the same $z,z'$,
\[
 \|b_{t,A}^{\,z}-b_{t,A}^{\,z'}\|_\infty
    \le \varepsilon_A(R)
           +|A|D^2K_\zeta\,\theta_{t,A}(R).
 \tag{V63.13}\label{c18v63:scorebound}
\]
\end{proposition}

\begin{proof}
For any scalar $h$ on one fixed factor, its factor gradient
vanishes at a minimum. Parallel transport along a minimizing
factor geodesic therefore bounds that gradient everywhere by
$D\sup\|(\operatorname{Hess}h)_{ii}\|$.
Apply this to $a_t$ with all other coordinates fixed.
Changing the coordinates in $A$ one at a time yields
$\operatorname{osc}_{A}a_t(\cdot,z')\le|A|D^2K_\zeta$.

Put $h=a_t(\cdot,z)-a_t(\cdot,z')$, and abbreviate the two
probabilities by $P,Q$. Their centered-score difference is
\[
       h-\mathbb E_P h
           +(\mathbb E_Q-\mathbb E_P)a_t(\cdot,z').
\]
Its first term is bounded by $\operatorname{osc}h\le
\varepsilon_A(R)$ using \eqref{c18v63:rectangle}.
The second is bounded by
\[
       \operatorname{osc}_{A}a_t(\cdot,z')\,d_{\rm TV}(P,Q).
\]
This proves \eqref{c18v63:scorebound}. In particular the
derivative of a conditional normalizer was not discarded.
\end{proof}

One immediate implementation is to keep the exterior inside
the radius-$R$ collar and set the remaining factor coordinates
to any fixed reference values on this same fibre.
Equation \eqref{c18v63:TV} controls the resulting block kernel
and \eqref{c18v63:scorebound} its score. This construction
evaluates the original $u_t$ on the original torus.
The separately defined frozen-collar kernels need not form
one mutually compatible joint law.

There is also a gap comparison on the \emph{fixed} product
metric. If $\lambda_A^z$ is the scalar conditional diffusion
gap of \eqref{c18v63:kernel} using $m_0|_A$, then
\[
 e^{-\delta_{t,A}(R)}\lambda_A^{z'}
       \le\lambda_A^z
       \le e^{\delta_{t,A}(R)}\lambda_A^{z'}.
 \tag{V63.14}\label{c18v63:gapcompare}
\]
Indeed if $l\le dP/dQ\le L$, their Dirichlet energies and
variances, expressed as infima over constant centers, compare
by $l,L$. The Rayleigh quotient loss is at most $L/l$,
which here is at most $e^\delta$, not $e^{2\delta}$.
This comparison does not provide a positive reference-block
margin uniform in its size, and is not a comparison of the
different actual metrics $\widehat m_t$ at the two exteriors.

\subsection{Approximate Markov blanket with the boundary term retained}

Partition the factor labels into $A,B,C$, with
$d(A,C)\ge R>0$. Fix any $B$ value $b$, and write
$\pi_b=\widehat\nu_t(dx_A,dx_C\mid x_B=b)$ and
$\pi_{A,b},\pi_{C,b}$ for its marginals. Define
$\delta=\delta_{t,A}(R)$ and $\theta=\theta_{t,A}(R)$.

\begin{theorem}[Native conditional boundary separation]
\label{c18v63:blanket}
The actual joint conditional obeys
\[
 \boxed{\quad
 d_{\rm TV}(\pi_b,\pi_{A,b}\otimes\pi_{C,b})\le\theta,
 \qquad
 \frac{d\pi_b}{d(\pi_{A,b}\otimes\pi_{C,b})}\ge e^{-\delta}.
 \quad}
 \tag{V63.15}\label{c18v63:product}
\]
For real square-integrable $F=F(x_A)$ and $G=G(x_C)$,
\[
 |\operatorname{Cov}_{\pi_b}(F,G)|
     \le(1-e^{-\delta})
           \sqrt{\operatorname{Var}_{\pi_{A,b}}F\,
                 \operatorname{Var}_{\pi_{C,b}}G}.
 \tag{V63.16}\label{c18v63:L2}
\]
For bounded real writers one also has
\[
 |\operatorname{Cov}_{\pi_b}(F,G)|
             \le\tfrac14\theta\,
                   \operatorname{osc}F\,\operatorname{osc}G .
 \tag{V63.17}\label{c18v63:bounded}
\]
All three statements hold with the original boundary law and
normalizers, including atypical fixed boundary values.
\end{theorem}

\begin{proof}
Disintegrate $\pi_b$ over $C$, writing its $A$-kernel as $P_c$.
Every pair $P_c,P_{c'}$ satisfies \eqref{c18v63:TV}, because
all changed coordinates lie in $C$. Averaging over $c'$ gives
$d_{\rm TV}(P_c,\pi_{A,b})\le\theta$. Integrating over $c$
proves the joint TV claim. The pairwise likelihood ratio
$P_c/P_{c'}$ is at least $e^{-\delta}$ pointwise; averaging
the denominator comparison shows $P_c\ge
e^{-\delta}\pi_{A,b}$. This proves the joint minorization.

If $\delta>0$, set $\alpha=e^{-\delta}$ and write
\[
 \pi_b=\alpha\,\pi_{A,b}\otimes\pi_{C,b}
                         +(1-\alpha)\widetilde\pi_b .
\]
The remainder is a probability with the same two marginals.
For marginally centered $F,G$ the product term is zero,
and Cauchy--Schwarz under $\widetilde\pi_b$ proves
\eqref{c18v63:L2}. At $\delta=0$ the law is the exact product.

For the bounded claim,
$m(c)=\mathbb E_{P_c}F$ has oscillation at most
$\theta\operatorname{osc}F$, by the pairwise TV estimate.
The covariance equals $\operatorname{Cov}_{\pi_{C,b}}(m,G)$.
Each variance is at most one quarter of its squared range;
Cauchy--Schwarz proves \eqref{c18v63:bounded}.
\end{proof}

The corresponding full-law identity is
\[
 \begin{split}
 \operatorname{Cov}_{\widehat\nu_t}(F,G)
 &=\operatorname{Cov}_{\widehat\nu_t}
       \bigl(\mathbb E[F\mid B],\mathbb E[G\mid B]\bigr)
                           +\mathcal E_{A,C\mid B},\\
 |\mathcal E_{A,C\mid B}|
 &\le(1-e^{-\delta})
       \sqrt{\mathbb E\operatorname{Var}(F\mid B)\,
             \mathbb E\operatorname{Var}(G\mid B)} .
 \end{split}
 \tag{V63.18}\label{c18v63:mediation}
\]
This is total covariance followed by
\eqref{c18v63:L2} and Cauchy--Schwarz in $B$.
It is the controlled-error counterpart of the Grand Tensor's
exact Wilson boundary mediation. The first covariance on the
right is \emph{not} declared small.

For example, the elementary three-spin probability
$P(s_1,s_2,s_3)\propto e^{J(s_1s_2+s_2s_3)}$ has exact
conditional independence of $s_1,s_3$ given $s_2$, but
\[
       \mathbb E[s_1s_3]=\tanh^2 J .
\]
Indeed $\mathbb E[s_1\mid s_2]=\mathbb E[s_3\mid s_2]
=s_2\tanh J$ and $s_2$ is symmetric. The correlation is
entirely the retained boundary term.
This elementary comparison is not the native Wilson law;
it explains why even zero conditional-separation error
cannot justify deleting that term.

\subsection{What is paid and the next unresolved native calculation}

V63 pays the actual spatial Hessian tail, normalized
conditional boundary error, conditional temporal-score error,
and a boundary-conditioned $L^2$ separation estimate through
smoothing time four at $\beta=0$. No small-time continuation
of the nonlinear Poisson inverse is used in these proofs.
The constants are explicit but extremely large; mere finiteness
is not a useful numerical coupling margin.

The next load-bearing calculation remains the \emph{near}
native signed response relative to conditional diagonal
energies, or a quantitative control of the retained boundary
message covariance in \eqref{c18v63:mediation}.
The tail bound \eqref{c18v63:signedtail} removes one spatial
error from that calculation. It does not perform it.
There is no fresh Haar restart, no integration of the collar
followed by an unjustified reuse of its pinned estimate,
and no claim of thermodynamic consistency from separately
frozen exterior kernels.

The verifier
\begin{center}\small
\path{verify_ym_c18_v63_native_boundary.py}
\end{center}
checks periodic anchor counts, weighted summation and restriction
bookkeeping, exact finite-probability normalizer and score
identities, the minorization argument, and the boundary-mediated
three-spin control. These are regression fixtures for the proof
algebra, not enclosures of a native Wilson trajectory or a
proof-assistant verification. The inherited weighted normal-field
certificates are used in their stated scopes and are not
independently recertified in full here.

The build audit also corrects an earlier verification record:
the V62 log does contain repeated equation hyperlink destinations
and an inherited \texttt{amsmath} warning for old binomial syntax.
They persist in the preserved body; the earlier zero-warning
description was inaccurate. They are not new mathematical
estimates or a reason to alter an inherited proof silently.
The V63 appendix's line overflow was repaired before inspection.

Removing this appendix and restoring V62's title recovers V62
byte for byte. Only the active YM filename is replaced; all
other sources and supplied downloads remain untouched.
Build and inspection products remain outside
\texttt{latex\_notes}. The next companion review is V65,
the next broad literature sweep V70, and archive/restart V100.

The actual probability-transport strain through time four,
its interacting-vacuum extension, weak-coupling scale transport,
physical return and clock calibration, nontrivial
reflection-positive continuum construction, and full positive
physical Yang--Mills mass gap remain unproved.
% END CYCLE 18 VERSION 63 APPEND
% BEGIN CYCLE 18 VERSION 64 APPEND
\clearpage
\section{V64: Gauge-constrained native coercivity and a pure-gauge comparison loss}
\label{c18v64:context}

\subsection{Context, literature intake, and the nonduplicative target}

This continuation addresses the actual electric normal-chart law, not
a replacement Gibbs model. V63 paid an exponentially small spatial
tail for the signed matrix of V62, but did not establish positivity
of its near part. The new user-supplied literature suggestions
motivate a sharper question: which directions does the actual
invariant temporal source use? We prove a native midpoint Ward
inequality, spatial decay for the already constructed gauge
projection, and an exact first-order loss in the unrestricted
comparison matrix. That loss has a pure-gauge minimizing direction.
This is a reason to sharpen the comparison, not a mass-gap proof.

The attachment's proposed contraction $\mathbb E[b\,u'''_0]$
refers to the V58 frontier. V59 already evaluated it and closed the
stated second-order scalar cost coefficient. We do not reopen that
calculation. Likewise V33 already restricted the scalar inverse by
gauge symmetry, and V34 already constructed the pinned gauge
projection and proved its uniform Gram inverse bound. Neither
construction is counted as a new result here. The present
contribution is their quantitative use in V62's \emph{native signed}
test, including a calculation of a loss that the unrestricted test
actually incurs.

The following is a targeted primary-source screen, not the
scheduled broad V70 sweep. Abstracts or publication records were
checked for the cited works; additionally we inspected the
introductory structure and the relevant hypothesis statements in
the two most immediately relevant recent preprints. No cited
theorem is applied without verifying its model assumptions.
\begin{itemize}
\item Abdelghani--Nissim,
\href{https://arxiv.org/abs/2309.07399}{\emph{Geometric Derivation
of the Finite N Master Loop Equation}}, supplies compact-group
integration-by-parts context. It does not make the completed V59
contraction a new outstanding task.
\item Buchholz--Cotar--Schweiger,
\href{https://arxiv.org/abs/2608.14526}{\emph{Gradient Gibbs measures
with non-convex potentials and the universality class of the
Gaussian Free Field}}, concerns real gradient fields. Its
Helffer--Sjostrand construction uses a log-mixture representation;
Section 7 requires positive bounded conditional second derivatives.
No such representation for our compact Wilson law is established.
We adopt the question of a source-restricted inverse, not that
paper's coercivity conclusion.
\item Otto--Reznikoff,
\href{https://doi.org/10.1016/j.jfa.2006.10.002}{\emph{A new
criterion for the logarithmic Sobolev inequality and two
applications}}, and Caputo--Parisi,
\href{https://arxiv.org/abs/2004.10574}{\emph{Block factorization
of the relative entropy via spatial mixing}}, organize conditional
coercivity and coupling/mixing inputs. V62--V63 do not yet pay the
native coupling or strong-mixing input. Stroock--Zegarlinski's
\href{https://doi.org/10.1016/0022-1236(92)90003-2}{continuous-spin
LSI criterion} likewise has a local-specification hypothesis.
The attachment's conservative-spin citation is specifically
Landim--Panizo--Yau,
\href{https://www.numdam.org/item/AIHPB_2002__38_5_739_0/}
{\emph{Spectral gap and logarithmic Sobolev inequality for unbounded
conservative spin systems}}, not an identification of our
physical evolution with a Lu--Yau dynamics.
\item Qin,
\href{https://arxiv.org/abs/2504.01247}{\emph{On spectral gap
decomposition for Markov chains}}, v2, makes the actual operator
factorization or Dirichlet-form comparison essential: Section 4.1
specifies a contraction and reversible conditional kernels.
A physical Doob transform alone does not verify that comparison
or bound the idealized coarse chain. V63's boundary-message
covariance must remain in any two-scale argument.
\item Bony--Popoff,
\href{https://arxiv.org/abs/1802.02882}{\emph{Low-lying eigenvalues
of semiclassical Schr\"odinger operator with degenerate wells}},
and Bauerschmidt--Bodineau,
\href{https://arxiv.org/abs/1907.12308}{\emph{Log-Sobolev inequality
for the continuum sine-Gordon model}}, suggest localization and
multiscale potential estimates. A native bad-region energy bound,
or control of the actual running potential, is still needed.
Baudel--Berglund's
\href{https://www.idpoisson.fr/berglund/spectralTheoryForRandomPoincareMaps.pdf}
{\emph{Spectral theory for random Poincar\'e maps}} was screened
only at its abstract: weak-noise periodic-orbit metastability
does not identify our physical transfer kernel.
\item Conrady--Khavkine,
\href{https://arxiv.org/abs/0706.3423}{\emph{An exact string
representation of 3d SU(2) lattice Yang--Mills theory}}, is an
algebraic backup, not a four-dimensional positive gap theorem.
Bourgain--Gamburd's
\href{https://arxiv.org/abs/1108.6264}{\emph{A Spectral Gap Theorem
in SU(d)}} requires the specified algebraic generators of a
dense subgroup; our spatially coupled law is not that convolution
operator. Bach--Ballesteros--Geisler,
\href{https://arxiv.org/abs/2508.07643}{\emph{The Spectral
Renormalization Flow Based on the Smooth Feshbach--Schur Map:
Introduction of the Semi-Group Property}}, does not supply a
missing complementary inverse.
\end{itemize}
The ScienceDirect full-text requests were inaccessible; indexed
publisher abstracts and the Numdam publication record support
only the screening statements just made. No full-paper audit of
all thirteen references is claimed.

\subsection{Native geometry and the exact invariant subspace}

Throughout the calculations below, $N=2n$, $n\ge5$, $\beta=0$,
$g$ is fixed, and $0\le t\le4$ denotes Wilson smoothing time.
Retain the original product fibre
\[
 M_g=\mathfrak b_2^{-1}(g),\qquad
 \widehat\nu_t=e^{u_t}\mu_g,\qquad a_t=\partial_tu_t,\qquad
 L_t=-\Delta_{m_0}-\langle\nabla u_t,\nabla\cdot\rangle .
\]
The density is normalized. There are $3n^3$ pair factors
$(A,A^{-1}g_e)$ with metric twice the round metric and $18n^3$
free-link factors with round metric. Put $\kappa_i=1/2$ on
pairs and $\kappa_i=1$ on free factors. All adjoints, gradients
and tangent norms in this note use $m_0$, unless expressly stated.

Let $\mathcal H_0$ be the fine gauge group with identity value
at every coarse vertex $2z$. Its action preserves $M_g$, $m_0$
and $\mu_g$. Equivariance of the native normal field and its flow
implies that $u_t$ and $a_t$ are invariant. Thus $L_t$ commutes
with the action. Averaging over $\mathcal H_0$ is an orthogonal
projection commuting with its resolvent. In particular the
unique centered solution of $L_t\psi=a_t$ is invariant.
Existence at each fixed finite volume and time follows from
compact ellipticity; no uniform inverse bound is implicit in
this statement.

Write $D_g(x)$ for the infinitesimal internal gauge action,
in orthonormal product-factor frames at $x\in M_g$. Via the
isometric inclusion into the fine link product it is exactly
the restriction of V34's $D_U$, at $U=i_g(x)$. Define
\[
 K_g=D_g^*D_g,\qquad P_g=I-D_gK_g^{-1}D_g^*.
 \tag{V64.1}\label{c18v64:projection}
\]
V34, with $\ell=2$, already gives
$21^{-1}I\preceq K_g\preceq12I$. Thus $P_g$ is the smooth
orthogonal projection onto $\ker D_g^*$, at every configuration.
For invariant $f$, $P_g\nabla f=\nabla f$. We call this tangent
subspace horizontal. This terminology refers only to the
pinned internal gauge action, not to a new physical Hamiltonian.

\begin{proposition}[Native midpoint Ward budget]
\label{c18v64:ward}
Decompose a horizontal vector $v=(p,f)$ into its pair and free
components. Uniformly in $g$, $x$ and $n$,
\[
 \|p\|^2\le2\|f\|^2,\qquad
 \left\langle v,\operatorname{diag}(3\kappa_i I)v\right\rangle
       \ge2\|v\|^2 .
 \tag{V64.2}\label{c18v64:wardbudget}
\]
\end{proposition}
\begin{proof}
A sampled-chain midpoint $w=2z+e_i$ has exactly four incident
free links in the two transverse coordinate directions.
Its other two links are the two halves of its pair. A gauge
parameter supported at $w$ has pair component of magnitude
$\sqrt2$ in the normalized pair frame; each incident free
component has magnitude one. Hence the midpoint equation
$D_g^*v=0$ has the form
\[
 \sqrt2\,O_wp_{z,i}
       +\sum_{\substack{e\ {\rm free}\\ e\sim w}}
                       \epsilon_{we}O_{we}f_e=0 ,
 \tag{V64.3}\label{c18v64:midpoint}
\]
where the $O$'s are orthogonal adjoint transports and
$\epsilon_{we}\in\{-1,1\}$. Cauchy--Schwarz gives
$2|p_{z,i}|^2\le4\sum_{e\sim w,\ {\rm free}}|f_e|^2$.
Every midpoint has exactly one odd coordinate. A free nearest
neighbor edge has at most one such endpoint: its endpoint
parity classes differ in one coordinate, and an edge joining
a coarse vertex to a midpoint is captured, not free.
Thus no free edge is counted twice when summing this inequality
over midpoints. This proves the first assertion. Finally
\[
 \tfrac32\|p\|^2+3\|f\|^2-2(\|p\|^2+\|f\|^2)
       =\|f\|^2-\tfrac12\|p\|^2\ge0 .
\]
\end{proof}
The number $2$ is a floor for this \emph{diagonal comparison
form}, not a computation of the true invariant Haar gap.
V54 already obtained the larger invariant scalar Haar gap $9$.
No improvement over that exact result is claimed.

\subsection{A spatial estimate for V34's existing projection}

Use fine torus graph distance $d$. For factor indices use V63's
anchor convention: the root of a pair's first half, or the
tail of a free edge. For rectangular block matrices define
$\|\cdot\|_{{\rm Sch},\xi}$ to be the maximum of the weighted
row and column sums of block operator norms, with weight
$e^{\xi d}$. The same definition applies between factor and
internal-vertex blocks. Distances are measured between their
respective anchors. This is a block norm, not a count of
individual scalar frame entries.

\begin{proposition}[Quasilocality of the pinned projection]
\label{c18v64:locality}
Set $q=251/252$, $\alpha=\log(252/251)$ and
$Q(s)=((1+e^{-s})/(1-e^{-s}))^3$. Then
\[
 \|(K_g^{-1})_{vw}\|\le21e^{-\alpha d(v,w)},\qquad
 \|P_g\|_{{\rm Sch},\xi}
       \le1+756e^{2\xi}Q(\alpha-\xi)\quad(0<\xi<\alpha).
 \tag{V64.4}\label{c18v64:quasilocal}
\]
These bounds are uniform in volume, coarse configuration and
fibre point.
\end{proposition}
\begin{proof}
The Gram matrix has range one on the internal fine vertices.
Let $B=I-K_g/12$. V34's bounds give $0\preceq B\preceq qI$,
so
\[
 K_g^{-1}=\frac1{12}\sum_{k=0}^{\infty}B^k .
\]
The $(v,w)$ block of $B^k$ vanishes for $k<d(v,w)$ and otherwise
has norm at most $q^k$. Summing the geometric series gives the
first bound, since $[12(1-q)]^{-1}=21$.
For any fine vertex the weighted sum over internal vertices is
at most the sum over all $\mathbb Z^3$ displacement vectors,
namely $Q(\alpha-\xi)$. Thus
$\|K_g^{-1}\|_{{\rm Sch},\xi}\le21Q(\alpha-\xi)$.

A free-factor row of $D_g$ has at most two blocks of norm one,
each within distance one. A pair row has one block of norm
$\sqrt2$ at its midpoint, also at distance one. Each internal
vertex column has either six unit free blocks, or four such
blocks and one pair block of norm $\sqrt2$. Consequently
$\|D_g\|_{{\rm Sch},\xi}\le6e^\xi$, and the same holds for its
adjoint. The triangle inequality for distance proves
submultiplicativity of this weighted block Schur norm.
Apply it to \eqref{c18v64:projection}.
\end{proof}
This proves decay for a matrix already known to exist. It is
not an inverse estimate for $L_t$ or for the physical transfer.
No bound on derivatives of $P_g$ is asserted here.

\subsection{The correctly restricted signed test and inverse}

Let $\lambda_i(t)>0$ be any exterior-uniform scalar conditional
gap floors for the actual law. As in V62 define
\[
 (\mathsf S_t)_{ii}=\lambda_i(t)I,\qquad
 (\mathsf S_t)_{ij}=-\operatorname{Hess}_{ij}u_t\quad(i\ne j).
\]
Write $\operatorname{gap}_{\rm inv}(L_t)$ for the scalar gap
on $\mathcal H_0$-invariant functions orthogonal to constants.

\begin{theorem}[Gauge-constrained conditional comparison]
\label{c18v64:comparison}
If, pointwise on $M_g$,
\[
 P_g\mathsf S_tP_g\succeq\rho P_g,\qquad \rho>0,
 \tag{V64.5}\label{c18v64:test}
\]
then $\operatorname{gap}_{\rm inv}(L_t)\ge\rho$. For every
centered invariant scalar source $a$,
\[
 \int|\nabla L_t^{-1}a|^2\,d\widehat\nu_t
       \le\rho^{-1}\|a\|_{L^2(\widehat\nu_t)}^2 .
 \tag{V64.6}\label{c18v64:cost}
\]
Furthermore, if $\mathsf S_t^{[R]}$ is V63's off-diagonal
spatial truncation and
$P_g\mathsf S_t^{[R]}P_g\succeq\rho_RP_g$ everywhere, then
\[
 \operatorname{gap}_{\rm inv}(L_t)
       \ge\rho_R-tK_\zeta e^{-\zeta R}
 \quad\hbox{whenever the right side is positive}.
 \tag{V64.7}\label{c18v64:tail}
\]
\end{theorem}
\begin{proof}
V62 proves, for scalar $\psi$,
\[
 \|L_t\psi\|_2^2
   \ge\int\langle\nabla\psi,\mathsf S_t\nabla\psi\rangle
                            \,d\widehat\nu_t .
\]
For invariant $\psi$ its gradient is horizontal. Apply
\eqref{c18v64:test} to each nonconstant invariant eigenfunction.
If its eigenvalue is $\ell$, then $\ell^2\ge\rho\ell$.
The invariant subspace is reducing and has a complete scalar
eigenbasis. Expanding $a$ in that basis gives
$\langle a,L_t^{-1}a\rangle\le\rho^{-1}\|a\|_2^2$,
which is \eqref{c18v64:cost}.
Finally $\|P_g\|_{\rm op}=1$, so compression does not enlarge
V63's pointwise tail bound
$\|\mathsf S_t-\mathsf S_t^{[R]}\|
 \le tK_\zeta e^{-\zeta R}$. Apply the first part.
\end{proof}
There is no weighted-Schur constant in the error in
\eqref{c18v64:tail}. Conversely, $P_g\mathsf S_t^{[R]}P_g$
is not strictly finite range; \eqref{c18v64:quasilocal} describes
the projection's spatial spread. Freezing a collar still
requires V63's actual conditional normalization, not replacement
by a smaller periodic lattice.

For orientation, V62's crude floors
$\lambda_i(t)=3\kappa_i e^{-D^2tK_*}$, $D=\pi\sqrt2$,
and the bound on the off-diagonal Hessian row sums give
\[
 P_g\mathsf S_tP_g
 \succeq\bigl(2e^{-D^2tK_*}-tK_*\bigr)P_g .
 \tag{V64.8}\label{c18v64:crudefloor}
\]
Indeed the diagonal term uses \eqref{c18v64:wardbudget};
the off-diagonal symmetric operator has norm at most $tK_*$
by its block row bound. This improves the $3/2$ diagonal
budget in V62's unrestricted test. It is neither an optimized
small-time theorem nor a positive estimate at $t=4$.
In particular it is not a replacement for V54's exact Haar
spectral calculation or V61's actual short-time strain bound.

The inherited comparison
$E_0^{-2}m_0\preceq\widehat m_t\preceq E_0^2m_0$ implies,
under \eqref{c18v64:test}, an invariant scalar gap
$E_0^{-2}\rho$ for the actual pulled-back metric, by comparison
of Dirichlet energies. Its integrated source cost is at most
$E_0^2\rho^{-1}\|a_t\|_2^2$. These are scalar energy estimates;
they do not bound the pointwise Hessian or probability-transport
strain of the Poisson solution.

To make the source-projected Witten idea precise, let
$\mathfrak H_1=dd_{\widehat\nu_t}^*
                 +d_{\widehat\nu_t}^*d$
be the nonnegative weighted Hodge Laplacian on one-forms.
Let $\mathcal E_{\rm inv}$ be the $L^2$ closure of
$\{df:f\text{ smooth and invariant}\}$.
The standard identity $\mathfrak H_1d=dL_t$ follows by
$d^2=0$ and $d_{\widehat\nu_t}^*d=L_t$ on scalars.
More explicitly, if $f_j$ is an invariant normalized
nonconstant eigenfunction with eigenvalue $\ell_j$,
then $df_j/\sqrt{\ell_j}$ is an orthonormal one-form
eigenfamily spanning $\mathcal E_{\rm inv}$.
Thus this closed subspace reduces $\mathfrak H_1$, and
\[
 dL_t^{-1}a
   =(\mathfrak H_1|_{\mathcal E_{\rm inv}})^{-1}da
 \quad\hbox{for smooth centered invariant }a .
 \tag{V64.9}\label{c18v64:witten}
\]
This proves the identity without assuming a gap on arbitrary
one-forms. At fixed finite volume the inverse on this subspace
exists; a uniform bound still requires a quantitative estimate.

Pointwise horizontality is necessary for membership in
$\mathcal E_{\rm inv}$, not sufficient. General horizontal
forms need not be invariant or exact. Multiplication by the
configuration-dependent $P_g$ is not asserted to commute
with $\mathfrak H_1$; differentiating it introduces derivatives
of $P_g$. Accordingly we do \emph{not} replace
$P_g\mathfrak H_1^{-1}P_g$ by the inverse of a compressed
differential operator. A still narrower physical source family
would require its own justified invariant subspace or resolvent
estimate. The preceding theorem supplies neither automatically.

\subsection{An exact native first-order loss, even with optimal conditional floors}

We now evaluate the signed matrix itself, rather than a generic
comparison example. Set $g=I$ and let $x=I$ denote the point
where every product factor is the identity. For each factor,
define the optimal exterior-uniform conditional floor
\[
 \lambda_i^{\rm opt}(t)
   =\inf_z\operatorname{gap}
       \bigl(-\Delta_i-\langle\nabla_i u_t(\cdot,z),\nabla_i\cdot\rangle\bigr).
\]
The infimum is over that factor's full compact exterior.
Use these floors in $\mathsf S_t^{\rm opt}$.

\begin{lemma}[No linear loss in the native one-factor gap]
\label{c18v64:localgap}
For each fixed $n,g$, uniformly in the conditional exterior,
\[
 \operatorname{gap}(L_{i,t}^{\,z})
       =3\kappa_i+O_{n,g}(t^2),\qquad
 \lambda_i^{\rm opt}(t)=3\kappa_i+O_{n,g}(t^2)
 \quad(t\longrightarrow0).
 \tag{V64.10}\label{c18v64:localexpansion}
\]
No volume-uniform quadratic remainder is claimed.
\end{lemma}
\begin{proof}
The inherited identity is $u'_0=b=3W_{\rm cap}$.
For fixed exterior, every captured plaquette containing a
given factor contains it exactly once, possibly inverted.
A quaternion and its inverse have entries linear in its
four real coordinates. Thus $b(\cdot,z)=c(z)+\ell_z$, where
$\ell_z$ is a degree-one spherical harmonic on that $S^3$.
For the metric in factor $i$,
$\Delta_i\ell_z=-3\kappa_i\ell_z$.

The unitary multiplication by the square root of the conditional
density transforms $L_{i,t}^{\,z}$ to the fixed Haar space as
\[
 H_{i,t}^{\,z}
  =-\Delta_i+\tfrac12\Delta_i u_t(\cdot,z)
                   +\tfrac14|\nabla_i u_t(\cdot,z)|^2 .
 \tag{V64.11}\label{c18v64:ground}
\]
The sign here corresponds to density $e^{u_t}$.
Consequently its first perturbation is multiplication by
$-(3\kappa_i/2)\ell_z$.
The first nonzero Haar eigenspace has eigenvalue $3\kappa_i$
and is spanned by the four quaternion coordinate functions.
The compression of this perturbation to that eigenspace is
zero: each matrix entry is an integral of three linear
coordinates, odd under $q\mapsto-q$.

For completeness the quadratic remainder is uniform in $z$
at fixed $n,g$. Smoothness of the actual $u_t$ on a compact
time neighborhood and compact fibre gives a uniform
operator-norm expansion of the bounded multiplication
potential in \eqref{c18v64:ground}. The Haar spectral
cluster at $3\kappa_i$ is separated from $0$ and
$8\kappa_i$. On a fixed contour around that cluster, the
resolvent Neumann expansion is uniform in $z$ for small $t$.
The resulting spectral projection equals the Haar projection
plus $O_{n,g}(t)$ and has rank four. Transport it unitarily
to the Haar subspace; its restricted matrix is
$3\kappa_i I+tP H'_0P+O_{n,g}(t^2)$.
Here $PH'_0P=0$. The zero eigenvalue remains simple
(the positive conditional ground state), and the next
cluster stays separated. Hence the first positive eigenvalue
has the stated expansion, uniformly in $z$. Taking the
infimum preserves it.
\end{proof}

\begin{theorem}[A pure-gauge minimizing direction of the initial comparison loss]
\label{c18v64:loss}
For each fixed $n\ge5$,
\[
 \lambda_{\min}\bigl(\mathsf S_t^{\rm opt}(I)\bigr)
       =\tfrac32-12t+O_n(t^2)\qquad(t\downarrow0).
 \tag{V64.12}\label{c18v64:slope}
\]
A minimizing direction of the first-order pair block is
constant on all pair factors and zero on all free factors.
It is a pure internal gauge direction, so $P_g(I)$ kills it.
\end{theorem}
\begin{proof}
At the identity, the Hessian of a normalized Wilson plaquette
trace is minus the square of its linearized oriented curl.
For each of the three Lie-algebra colors independently, put
$\mathsf K=-\operatorname{Hess}W_{\rm cap}(I)$.
The normalized coefficient of a captured first half is
$+1/\sqrt2$ times its pair coordinate and that of its second
half is $-1/\sqrt2$; free-link coefficients have magnitude one.
Since $b=3W_{\rm cap}$, the off-diagonal derivative of the
signed matrix is $3\mathsf K_{\rm off}$.

Here is an explicit enumeration of its pair block. For each
coarse plaquette in directions $i<j$ based at $z$, the four
captured fine corner plaquettes have pair parts of their
linearized curls, up to an irrelevant overall sign,
\[
 \begin{array}{ll}
 (p_i(z)-p_j(z))/\sqrt2,&
 (-p_i(z)+p_j(z+e_i))/\sqrt2,\\
 (p_j(z)-p_i(z+e_j))/\sqrt2,&
 (p_i(z+e_j)-p_j(z+e_i))/\sqrt2 .
 \end{array}
 \tag{V64.13}\label{c18v64:corners}
\]
These are all captured plaquettes. Every pair factor appears
eight times, giving diagonal $4$ in $\mathsf K_{pp}$.
Each has eight distinct pair neighbors, each with off-diagonal
entry $-1/2$. Thus the off-diagonal pair block has row sum
$-4$ and absolute row sum $4$. Its least eigenvalue is
exactly $-4$, attained by the constant pair vector.

At $t=0$, $\mathsf S_0^{\rm opt}$ has eigenvalue $3/2$ on
the pair subspace and $3$ on the free subspace. By
\eqref{c18v64:localexpansion} its diagonal first derivative
is zero. Finite-dimensional spectral perturbation of the
isolated pair cluster therefore has first-order matrix
$3(\mathsf K_{pp})_{\rm off}$, whose minimum is $-12$.
This proves \eqref{c18v64:slope}; mixing with the separated
free cluster enters only the fixed-volume quadratic remainder.

Finally choose a gauge parameter equal to a fixed color
$\xi$ at \emph{every} noncoarse fine vertex, and zero at
coarse vertices. At $U=I$, every free edge joins noncoarse
vertices, so its gauge variation is zero. The first half
of every captured pair has variation $-\xi$, and the second
half $+\xi$. Its orthonormal pair component is therefore
$-\sqrt2\,\xi$, constant over all pairs. This realizes the
claimed minimizing vector as $D_g\xi$, and
$P_gD_g\xi=0$.
\end{proof}

This is a loss of a \emph{sufficient comparison matrix}, not
the derivative of the native scalar spectral gap. It occurs
even if the diagonal floors are the exact optimal conditional
ones, so it cannot be blamed only on the crude oscillation
comparison. It does not show that the matrix becomes negative
at $t=1/8$: the remainder has not been controlled there.
Nor does it show that every dangerous later mode is pure
gauge, or that the horizontal compression is positive at
time four. The result identifies a concrete direction that
the source-restricted estimate need not pay.

\subsection{Verification and the remaining native obligation}

The new exact regression script
\begin{center}\small
\path{verify_ym_c18_v64_horizontal.py}
\end{center}
enumerates the captured plaquettes and midpoint incidence
for $n=5,6,7$, checks the rational pair-block entries and
constant gauge vector, the antipodal first-harmonic matrix
elements, and the Schur/Gram constants. The proofs above
supply the arbitrary-volume statements. The checks are not
a rigorous enclosure of a finite-time Wilson trajectory,
nor a proof-assistant certification of inherited results.

The next near-field calculation should estimate
$P_g\mathsf S_t^{[R]}P_g$ on the actual original fibre, retaining
both the Ward constraint and V63's paid tail. If this route
still loses the physical response, one must keep the positive
mixed-derivative terms discarded in V62, or prove a bound
on the exact invariant source subspace; one must not infer
an inverse by merely projecting a differential operator.
The complementary two-scale route still has to control
V63's boundary-message covariance under the actual law.
These are distinct unpaid estimates, not interchangeable
names for a solved bottleneck.

Only the active YM source is replaced. Removing this appendix
and restoring V63's title recovers V63 byte for byte; supplied
documents and other programmes are not altered. Build and
inspection files stay outside \texttt{latex\_notes}.
The next companion checkpoint remains V65, the next broad
primary-literature sweep V70, and archive/restart V100.

The actual probability-transport strain through time four,
the interacting-vacuum extension, weak-coupling scale transport,
coarse physical return, physical clock calibration, a
nontrivial reflection-positive continuum construction, and
the full positive physical Yang--Mills mass gap remain unproved.
% END CYCLE 18 VERSION 64 APPEND
% BEGIN CYCLE 18 VERSION 65 APPEND
\clearpage
\section{V65: Flat-point horizontal squares and a positive native initial slope}
\label{c18v65:context}

\subsection{The companion checkpoint and the calculation it selects}

The scheduled V65 checkpoint rehashed the eleven active companion
files recorded at V60; all eleven hashes are unchanged. The terminal
NS V26 and SM--Einstein V72 discussions were reread, along with the
current scope statements in the other branches. This is a targeted
frontier review, not a recertification of their entire cumulative
manuscripts. NS distinguishes a norm on the actual rank-one inputs
from an unrestricted tensor norm, while retaining its stated
verification limitations. SM V72 distinguishes a bounded-time,
fixed-mode many-copy limit from a physical finite-species continuum.
Neither supplies a Yang--Mills coercivity constant.

The remaining checked files were Exact-Action Einstein Bridge V55,
Native Regenerative Stationarity V61, SM Source Selection V65,
general perturbative ESM V44, phenomenology ESM V40, arithmetic
V41 and its Renewal V17 branch, Shell Mixing V16, and YM Parallel
V5. Their native tangency, source-occurrence, physical-writer,
noncommuting-curvature and extensive-charge restrictions remain
in force. In particular Shell Mixing V16 expressly warns that a
flat-holonomy calculation cannot replace a mixed-curvature estimate.
The more recent supplied arithmetic V42 and phenomenology V41
copies retain their V63 review; the older in-folder counterparts
are not treated as superseding them.

The resulting choice is concrete: evaluate the \emph{constrained}
flat-point matrix of V64, including the whole horizontal space
rather than one selected vector. V64 found an unrestricted initial
loss $3/2-12t+O_n(t^2)$ and removed a pure-gauge minimizer, but
did not compute the restricted minimum. Here an exact native
sum-of-squares identity does that. We also prove a
volume-uniform bound for the corresponding linear matrix pencil.
That pencil is not silently substituted for the nonlinear density.

The previous goal turn made progress by adding V64's native Ward
inequality and pure-gauge loss calculation. Its source hash and
current mathematical body were checked before this continuation.
The following results are new relative to that calculation; the
gauge projection, one-factor gap expansion and actual spatial tail
remain inherited results, with their original scopes.

\subsection{Flat-point notation and an orthogonal Ward remainder}

Work at $g=I$, $x=I$ on the original product fibre $M_I$ of V64.
Let $N=2n$, $n\ge5$, and retain $\beta=0$. All quantities below
are evaluated in the orthonormal pair/free frames of $m_0$.
Let
\[
 \mathcal H=\ker D_g(I)^*,\qquad
 \Lambda=\operatorname{diag}(3\kappa_i I),\qquad
 \mathsf K=-\operatorname{Hess}W_{\rm cap}(I).
\]
The space $\mathcal H$ is a tangent space, not the Hilbert space
of physical states. A vector in it is denoted $v=(p,f)$.
Pair components $p_i(z)$ are assigned to the midpoint $2z+e_i$.

Partition the free edges into $\mathcal F_1$, joining a
one-odd-coordinate midpoint to a two-odd-coordinate vertex,
and $\mathcal F_0$, joining a two-odd vertex to a three-odd
vertex. There are $12n^3$ and $6n^3$ such edges, respectively.
These exhaust the free edges. Let $D_f$ be the ordinary
oriented incidence map from internal vertex vectors to free
edge vectors at the identity: $(D_f\xi)_{xy}=\xi_x-\xi_y$.
Its adjoint is the divergence. Extend $p$ to the internal
vertices by $\widetilde p=p$ at the midpoints and zero elsewhere,
and define
\[
        Gp=\frac1{2\sqrt2}D_f\widetilde p,\qquad r=f-Gp .
 \tag{V65.1}\label{c18v65:remainder}
\]
Thus $Gp$ is supported on $\mathcal F_1$. The letter $G$ in
this appendix denotes this explicit map, not the normal-map
Gram matrix used earlier.

\begin{lemma}[Exact Ward orthogonality]
\label{c18v65:ward}
For every $v=(p,f)\in\mathcal H$,
\[
 \begin{aligned}
 (D_f^*r)_{\rm midpoint}&=0,&
 \langle Gp,r\rangle&=0,& \|Gp\|^2&=\tfrac12\|p\|^2,\\
 \|v\|^2&=\tfrac32\|p\|^2+\|r\|^2,&
 \langle v,\Lambda v\rangle&=2\|v\|^2+\|r\|^2 .
 \end{aligned}
 \tag{V65.2}\label{c18v65:orthogonal}
\]
\end{lemma}
\begin{proof}
At the identity the pair component of a gauge parameter at a
midpoint is $-\sqrt2\,\xi$. Hence horizontality there says
$(D_f^*f)_{\rm midpoint}=\sqrt2\,p$.
Every midpoint has four free neighbors, all of parity two,
and no two midpoints share a free edge. Therefore
$(D_f^*Gp)_{\rm midpoint}=\sqrt2\,p$ and
$\|D_f\widetilde p\|^2=4\|p\|^2$.
The divergence assertion follows. Also
$\langle Gp,r\rangle
=(2\sqrt2)^{-1}\langle\widetilde p,D_f^*r\rangle=0$.
Consequently $\|f\|^2=\|p\|^2/2+\|r\|^2$.
Insert this into $\|v\|^2=\|p\|^2+\|f\|^2$ and
$\langle v,\Lambda v\rangle=3\|p\|^2/2+3\|f\|^2$.
\end{proof}

For the two-odd vertex $w=2z+e_i+e_j$, $i<j$, abbreviate
\[
 a=p_i(z),\quad b=p_j(z),\quad
 c=p_i(z+e_j),\quad d=p_j(z+e_i).
\]
Define vertex arrays
\[
       (\mathcal Cp)_w=a+b+c+d,\qquad
       (\mathcal Ep)_w=a+c-b-d .
 \tag{V65.3}\label{c18v65:CE}
\]
The two-odd Ward equation gives
\[
       \mathcal Cp=2\sqrt2\,(D_f^*r)_{\rm two\text{-}odd},
       \qquad \|\mathcal Cp\|^2\le48\|r\|^2 .
 \tag{V65.4}\label{c18v65:Cbound}
\]
Indeed $(D_f^*Gp)_w=-(\mathcal Cp)_w/(2\sqrt2)$ and there
is no pair component at $w$. Each two-odd vertex has six
incident free edges, and each free edge has exactly one
two-odd endpoint. Thus Cauchy--Schwarz followed by summation
gives
$\|(D_f^*r)_{\rm two\text{-}odd}\|^2\le6\|r\|^2$.
This proves both assertions without a volume factor.

\subsection{The native curvature squares}

Let $J_p p$ denote the pair part of the linearized oriented
curl on all captured plaquettes, and $J_f f$ its free part.
At the flat identity,
\[
 \langle v,\mathsf K v\rangle=\|J_p p+J_f f\|^2,\qquad
          \mathsf K_{pp}=J_p^*J_p .
\]
These are Hessians of the actual normalized Wilson traces.
They are not an assumed Gaussian measure or a surrogate action.

\begin{proposition}[Corner identity and the coupled curl]
\label{c18v65:squares}
For arbitrary pair vectors $p$,
\[
 \langle p,\mathsf K_{pp}p\rangle
    =4\|p\|^2+\tfrac14\|\mathcal Ep\|^2
                    -\tfrac14\|\mathcal Cp\|^2,\qquad
 \|J_p p\|^2\le8\|p\|^2 .
 \tag{V65.5}\label{c18v65:corner}
\]
For $v=(p,f)\in\mathcal H$, with $r$ from
\eqref{c18v65:remainder},
\[
 \begin{aligned}
 J_p p+J_f f&=\tfrac32J_p p+J_f r,&
             \|J_f r\|^2&\le4\|r\|^2,\\
 \langle v,\operatorname{diag}(\mathsf K)v\rangle
             &=5\|p\|^2+2\|r_{\mathcal F_1}\|^2 .
 \end{aligned}
 \tag{V65.6}\label{c18v65:curl}
\]
\end{proposition}
\begin{proof}
The four pair curls around $w$ are, in a consistent orientation,
$(a-b)/\sqrt2$, $(-a+d)/\sqrt2$, $(b-c)/\sqrt2$ and
$(c-d)/\sqrt2$, as enumerated in V64. Their squared sum is
\[
 |a|^2+|b|^2+|c|^2+|d|^2
       +\tfrac14|a+c-b-d|^2-\tfrac14|a+b+c+d|^2 .
\]
Each midpoint occurs at four two-odd centers. Summing proves
the identity in \eqref{c18v65:corner}. The upper bound follows
from V64's pair matrix: diagonal $4$, eight off-diagonal
entries of size $1/2$ per row, so its operator norm is at most
$8$.

On every captured corner, the free curl of $Gp$ is exactly
one half of its pair curl. To check the signs, at the corner
with pair curl $(a-b)/\sqrt2$, the two free contributions are
$a/(2\sqrt2)$ and $-b/(2\sqrt2)$. At the next corner they
are $-a/(2\sqrt2)$ and $d/(2\sqrt2)$; the other two follow
the displayed orientation. Hence $J_fGp=J_pp/2$.

Every captured plaquette has two free edges, both in
$\mathcal F_1$, and every such free edge occurs in two captured
plaquettes. Therefore
$\sum_{\rm cap}|(J_fr)_q|^2
\le2\sum_{\rm cap}\sum_{e\in q,\ {\rm free}}|r_e|^2
=4\|r_{\mathcal F_1}\|^2\le4\|r\|^2$.
The diagonal of $\mathsf K$ equals $4$ on pairs, $2$ on
$\mathcal F_1$, and $0$ on $\mathcal F_0$.
Because $Gp$ vanishes on $\mathcal F_0$,
$\langle Gp,r_{\mathcal F_1}\rangle=0$ by
\eqref{c18v65:orthogonal}. Thus its diagonal quadratic form is
\[
 4\|p\|^2+2\|f_{\mathcal F_1}\|^2
       =5\|p\|^2+2\|r_{\mathcal F_1}\|^2 .
\]
\end{proof}

The equality space for the diagonal floor is consequently
\[
 \mathcal E_2
  =\{v\in\mathcal H:\langle v,\Lambda v\rangle=2\|v\|^2\}
  =\{(p,Gp):\mathcal Cp=0\}.
 \tag{V65.7}\label{c18v65:equality}
\]
The first equality follows from \eqref{c18v65:orthogonal};
the second uses all the remaining Ward equations.
At three-odd vertices $Gp$ vanishes on every incident edge.
At two-odd vertices its divergence vanishes precisely when
$\mathcal Cp=0$. In particular this is the \emph{full} internal
constraint, not midpoint invariance alone. On this space
\[
 \langle p,\mathsf K_{pp}p\rangle
      =4\|p\|^2+\tfrac14\|\mathcal Ep\|^2\ge4\|p\|^2 .
 \tag{V65.8}\label{c18v65:positive}
\]

\subsection{A uniform bound for the literal first-order pencil}

Define the symmetric matrix pencil
\[
       \mathsf A_t=\Lambda+3t\,\mathsf K_{\rm off},
       \qquad \mathsf K_{\rm off}
                   =\mathsf K-\operatorname{diag}(\mathsf K).
 \tag{V65.9}\label{c18v65:pencil}
\]
V64 proves that this is the first-order expansion of the
optimal-conditional-diagonal signed matrix at $g=I,x=I$:
\[
             \mathsf S_t^{\rm opt}(I)
                      =\mathsf A_t+O_n(t^2).
 \tag{V65.10}\label{c18v65:nativeremainder}
\]
The remainder in this display is an operator-norm statement
at each fixed volume. It is \emph{not} the uniform quadratic
constant proved for the pencil below.

\begin{theorem}[Flat horizontal pencil, uniformly in volume]
\label{c18v65:penciltheorem}
For every $n\ge5$, every $v\in\mathcal H$ and
$0\le t\le1/190$,
\[
       \langle v,\mathsf A_t v\rangle
           \ge(2+8t-864t^2)\|v\|^2 .
 \tag{V65.11}\label{c18v65:uniform}
\]
For even $n\ge6$ there is a nonzero horizontal vector with
Rayleigh quotient exactly $2+8t$, for every real $t$.
Thus, on those tori,
\[
 2+8t-864t^2
   \le\min_{\substack{v\in\mathcal H\\\|v\|=1}}
                  \langle v,\mathsf A_t v\rangle
   \le2+8t\qquad(0\le t\le1/190).
 \tag{V65.12}\label{c18v65:sandwich}
\]
\end{theorem}
\begin{proof}
Put $P=\|p\|$ and $R=\|r\|$ for this proof only.
Equations \eqref{c18v65:Cbound}--\eqref{c18v65:curl} give
\[
 \begin{split}
 \langle v,\mathsf K v\rangle
 &=\tfrac94\|J_pp\|^2
              +3\operatorname{Re}\langle J_pp,J_fr\rangle
              +\|J_fr\|^2\\
 &\ge9P^2-27R^2-12\sqrt2\,PR .
 \end{split}
\]
Here $\|J_pp\|^2\ge4P^2-12R^2$ and
$\|J_pp\|\|J_fr\|\le\sqrt8\,P\cdot2R$.
Also $\langle v,\Lambda v\rangle=3P^2+3R^2$.
Subtract the diagonal term of \eqref{c18v65:curl}, using
$\|r_{\mathcal F_1}\|\le R$, to obtain
\[
 \langle v,\mathsf A_t v\rangle
       \ge(3+12t)P^2+(3-87t)R^2-36\sqrt2\,tPR .
\]
Since $\|v\|^2=3P^2/2+R^2$, subtracting
$(2+8t)\|v\|^2$ leaves a lower bound
\[
 (1-95t)R^2-36\sqrt2\,tPR
       \ge-\frac{648t^2}{1-95t}P^2
       \ge-864t^2\|v\|^2 .
\]
The last two steps use completion of the square,
$1-95t\ge1/2$ and $P^2\le2\|v\|^2/3$.
This proves \eqref{c18v65:uniform} with no spectral-separation
constant that could deteriorate with $n$.

If $n$ is even, choose fixed Lie-algebra vectors $q_i$, not
all zero, and set
\[
           p_i(z)=(-1)^{z_1+z_2+z_3}q_i,\qquad f=Gp .
 \tag{V65.13}\label{c18v65:checkerboard}
\]
Periodicity uses the parity assumption. In every corner,
$a+c=0$ and $b+d=0$, so $\mathcal Cp=\mathcal Ep=0$.
Thus $v\in\mathcal H$, $R=0$,
$\|J_pp\|^2=4P^2$, $\langle v,\mathsf K v\rangle=9P^2$,
and the diagonal curvature contribution is $5P^2$.
It follows that
$\langle v,\mathsf A_tv\rangle=(3+12t)P^2$ and
$\|v\|^2=3P^2/2$. This proves the exact Rayleigh quotient
and \eqref{c18v65:sandwich}.
\end{proof}
The upper witness need not remain an eigenvector of the pencil
for $t\ne0$. Only its Rayleigh quotient was used. No analogous
checkerboard witness is asserted on an odd coarse torus.

\begin{corollary}[Actual flat-point initial slope after gauge reduction]
\label{c18v65:initial}
For each fixed even $n\ge6$,
\[
 \min_{\substack{v\in\mathcal H\\\|v\|=1}}
       \langle v,\mathsf S_t^{\rm opt}(I)v\rangle
                 =2+8t+O_n(t^2)\qquad(t\downarrow0).
 \tag{V65.14}\label{c18v65:initialslope}
\]
\end{corollary}
\begin{proof}
The minimum Rayleigh quotient changes by at most the
operator-norm change of the matrix.
Apply \eqref{c18v65:nativeremainder} to
\eqref{c18v65:sandwich}. V64's remainder includes the true
optimal exterior-uniform one-factor gaps and the actual
density Hessian, not just the first density jet.
\end{proof}
Equivalently this is the least eigenvalue of
$P_g(I)\mathsf S_t^{\rm opt}(I)P_g(I)$ \emph{on its horizontal
range}. Its full ambient matrix also has vertical zero
eigenvalues and must not be assigned the value in
\eqref{c18v65:initialslope}. Gauge equivariance transports the
same conclusion to the based internal gauge orbit of this
flat representative.

The contrast with V64 is exact: the unrestricted comparison
has the initial minimum $3/2-12t+O_n(t^2)$, whereas the
horizontal one has $2+8t+O_n(t^2)$ on the stated even tori.
Neither quantity is the invariant scalar gap of the
differential operator, whose Haar value was already $9$,
nor the physical Hamiltonian gap.

\subsection{A native central-flux configuration reverses the favorable sign}

The flat restriction is substantive, not merely a qualification
about the size of a Taylor remainder. Reuse V50's central
configuration, with the coordinate order reversed, on the
\emph{same} fibre $M_I$.
For a fine vertex $x\in(\mathbb Z/N\mathbb Z)^3$, set
\[
       U_i^\pi(x)=(-1)^{\sum_{j<i}x_j}I .
 \tag{V65.15}\label{c18v65:piflux}
\]
This is periodic because $N=2n$ is even. Every captured half
link is $I$: its base is $2z$ or $2z+e_i$, and all coordinates
with index less than $i$ are even there. Hence
$\mathfrak b_2(U^\pi)=I$ exactly.
The existence of this field and V55's associated flow reversal
are inherited. The new assertion below is its \emph{horizontal
signed-comparison} response, not another construction of a
negative-plaquette field.

\begin{proposition}[Opposite native initial Rayleigh response]
\label{c18v65:pi}
Every fine plaquette of $U^\pi$ equals $-I$.
In left-multiplication frames its internal gauge matrix is
the same as at $U=I$, but
\[
 \operatorname{Hess}W_{\rm cap}(U^\pi)=+\mathsf K,\qquad
 \mathsf S_t^{\rm opt}(U^\pi)
                  =\Lambda-3t\mathsf K_{\rm off}+O_n(t^2).
 \tag{V65.16}\label{c18v65:piresponse}
\]
For every fixed even $n\ge6$ the checkerboard vector
\eqref{c18v65:checkerboard}, placed in these frames, is
horizontal and has native signed Rayleigh quotient
\[
 \frac{\langle v,\mathsf S_t^{\rm opt}(U^\pi)v\rangle}
      {\|v\|^2}=2-8t+O_n(t^2).
 \tag{V65.17}\label{c18v65:pirayleigh}
\]
In particular the least horizontal quotient at $U^\pi$ is
at most this expression. Equality for the least quotient
is not asserted.
\end{proposition}
\begin{proof}
For $i<j$, increasing $x_i$ changes the sign of $U_j^\pi$,
whereas increasing $x_j$ leaves the sign of $U_i^\pi$
unchanged. The four central factors in the oriented
plaquette word therefore multiply to $-I$. All their
adjoint transports are the identity, so the gauge
incidence matrix in the chosen frames is unchanged.

In a nearby plaquette word the central signs commute
through the exponential variations. Its normalized trace
is minus the corresponding trace near the identity.
Its quadratic Hessian is consequently the negative of
the flat-identity Hessian. The pair inclusion has the
same first variation and the plaquette's first derivative
is zero at either central value, so restriction to $M_I$
preserves this Hessian sign reversal. Summing captured
plaquettes proves the first identity in
\eqref{c18v65:piresponse}.

Use $u'_0=3W_{\rm cap}$ and V64's zero linear variation
of the exterior-uniform optimal conditional floors.
This proves the second identity. The gauge matrix is the
same, so the checkerboard vector remains horizontal.
Its diagonal Haar quotient is $2$ and its quotient for
$3\mathsf K_{\rm off}$ is $8$, by the explicit computation
in Theorem~\ref{c18v65:penciltheorem}. This proves
\eqref{c18v65:pirayleigh}.
\end{proof}
Thus even at the same coarse value the horizontal test has
a decreasing native Rayleigh direction elsewhere in the
configuration space. The favorable flat-point initial slope
is not uniform in that space. This does not prove loss of
positivity at any specified positive time, and says nothing
by itself about the probability or physical energy of a
neighborhood of $U^\pi$ under the required interacting law.

\subsection{What this pays and what it does not}

This calculation resolves the flat-point sign question left
open in V64 and shows that its Ward floor is sharp on even
tori. The positive first-order coefficient is not an inference
from deleting one unfavorable vector: it follows from the
full horizontal-space lower bound and an actual attaining
Rayleigh witness. The constant $864$ controls the spectral
minimum of the \emph{linear pencil} uniformly in volume.
It does not bound the native nonlinear Taylor remainder.

The flat gauge orbit is only a subset of the configuration
space. No concentration of the actual electric law near it
has been proved or assumed. No conclusion at all other
configurations, all coarse $g$, smoothing time four, or
interacting vacuum follows by continuity without a quantified
neighborhood and a corresponding control of its complement.
The fixed-volume $O_n(t^2)$ in \eqref{c18v65:initialslope}
therefore remains distinct from V63's all-configuration
spatial tail and V61's actual short-time strain estimate.

The next useful step is not to optimize another flat-point
constant. One must control the signed response at nonflat
native configurations, including the central-flux witness,
with the exact Ward constraint retained,
or replace the pointwise test by an integrated exact-gradient
estimate under the actual law. A nonlinear temporal remainder
and its spatial/volume uniformity remain necessary for a
finite-time argument. The alternative boundary route must
still pay V63's retained conditional-mean covariance.
The companion checkpoint supplies no shortcut through these
obligations.

The verifier
\begin{center}\small
\path{verify_ym_c18_v65_flat_squares.py}
\end{center}
constructs horizontal vectors from actual fine plaquette curls
on $n=5,6,8$, checks all incidence and square identities by
rational arithmetic, and evaluates the checkerboard witnesses
on $n=6,8$. It also checks the displayed pencil bounds at
specified rational times, the literal central plaquette
signs, and the reversed checkerboard Rayleigh response.
The arbitrary-vector and
arbitrary-volume bounds are the analytic proofs above, not
claims extrapolated from those fixtures. No nonlinear
native trajectory or physical continuum gap is certified.

Removing this appendix and restoring V64's title recovers
V64 byte for byte. Only the active YM source is replaced;
other programmes and supplied downloads remain untouched.
Only \texttt{.tex} files are kept in \texttt{latex\_notes}.
The next companion review and broad literature checkpoint
are V70; archive/restart remains V100.

The full interacting four-dimensional Yang--Mills continuum
construction and positive physical mass gap remain unproved.
The goal is unchanged and remains active.
% END CYCLE 18 VERSION 65 APPEND
% BEGIN CYCLE 18 VERSION 66 APPEND
\clearpage
\section{V66: Gauge-stabilizer mixed energy at arbitrary native configurations}
\label{c18v66:context}

V65's favorable flat-point comparison slope has an opposite-sign
native nonflat witness. We therefore go upstream of the comparison
matrix: V62 discarded mixed derivative squares of the actual
invariant scalar response. These squares are not independent of
its gradient. Differentiating gauge invariance forces a definite
amount of them at every configuration, with no flatness or
probability-concentration assumption. This append proves that
quantitative statement and inserts it into both the pointwise
Bochner identity and V62's integrated conditional test.

The result is a native finite-volume geometric estimate, uniform
in volume and coarse data. It is not a construction of the
interacting continuum theory. The literature directions screened
in V64 still motivate retaining signed response and exact gradients;
no unverified external mixing or convexity hypothesis is used here.
The previous gauge projection, first-order Ward identity, conditional
Bochner identity and exact invariant Haar gap are inherited results,
not counted again as new.

\subsection{Native product geometry and the energy that was dropped}

Keep $N=2n$, $n\ge5$, and the fibre $(M_g,m_0)$ of V64.
A pair coordinate is $(A,A^{-1}g_e)$ with metric twice the
round metric on $SU(2)$; a free-link coordinate has the round
metric. Write $\mathcal P$ for pair factors and split free factors
as in V65:
\[
 \mathcal F_1=\{\hbox{midpoint--two-odd edges}\},\qquad
 \mathcal F_0=\{\hbox{two-odd--three-odd edges}\}.
\]
Here the parity class of a vertex is the number of its odd
coordinates. There are $12n^3$ factors in $\mathcal F_1$ and
$6n^3$ in $\mathcal F_0$. Both endpoints of every free edge
are internal, so both admit independent gauge parameters in
the group $\mathcal H_0$ pinned at all coarse roots.

For a smooth real $\mathcal H_0$-invariant scalar $f$, put
$v=\nabla f$, $H=\operatorname{Hess}_{m_0}f$, and
\[
 P_f=\sum_{i\in\mathcal P}|v_i|^2,\qquad
 F_{1,f}=\sum_{e\in\mathcal F_1}|v_e|^2,\qquad
 F_{0,f}=\sum_{e\in\mathcal F_0}|v_e|^2,\qquad
 \mathcal M(f)=\sum_{i\ne j}\|H_{ij}\|_{\rm HS}^2 .
 \tag{V66.1}\label{c18v66:energy}
\]
All these quantities are pointwise. The sum defining $\mathcal M$
is ordered: each off-diagonal symmetric pair is counted twice.
The proof of V64's midpoint Ward budget actually gives the
slightly more informative form
\[
              P_f\le2F_{1,f}.
 \tag{V66.2}\label{c18v66:retainedward}
\]
Indeed only midpoint-incident free edges occur in that proof;
they are exactly $\mathcal F_1$, and none is counted twice.
We reuse this observation, not a stronger unproved gauge constraint.

\subsection{A frozen two-endpoint stabilizer}

\begin{lemma}[Native free-factor stabilizer identity]
\label{c18v66:stabilizer}
Fix a configuration $U=i_g(x)$ and a free edge $e=(a,b)$
with stored orientation $a\to b$. For $\xi\in\mathfrak{su}(2)$
choose gauge parameters
\[
 \eta_a=\xi,\qquad \eta_b=\operatorname{Ad}_{U_e^{-1}}\xi,
 \qquad \eta_z=0\quad(z\notin\{a,b\}).
 \tag{V66.3}\label{c18v66:parameters}
\]
Let $Q_\xi$ be the fundamental gauge vector field for these
parameters, frozen at their displayed values. At $x$ its
$e$-component vanishes. Let
$B_e\xi=(Q_{\xi,j}(x))_{j\ne e}$ in orthonormal factor frames.
Then
\[
 H_{e,-e}B_e\xi=[\xi,v_e],\qquad B_e^*B_e=10I_3.
 \tag{V66.4}\label{c18v66:identity}
\]
The bracket uses the round-sphere normalization
$[\xi,z]=2\xi\times z$. If $\mathsf W$ has value $2I$ on
pair columns and $I$ on free columns, then
\[
 B_e^*\mathsf W^{-1}B_e=
 \begin{cases}
 9I_3,&e\in\mathcal F_1,\\
 10I_3,&e\in\mathcal F_0.
 \end{cases}
 \tag{V66.5}\label{c18v66:weightedgram}
\]
These identities hold for every $g,x,n$ in the stated regime.
\end{lemma}
\begin{proof}
In left-multiplication frames a free edge has gauge velocity
$\eta_a-\operatorname{Ad}_{U_e}\eta_b$, which is zero for
\eqref{c18v66:parameters}. Along
$U_e(s)=\exp(sw)U_e$, with all other factor coordinates fixed,
the derivative of this velocity at zero is
$-[w,\xi]=[\xi,w]$. There is no connection correction at that
point because $Q_{\xi,e}(x)=0$. The other components of
$Q_\xi$ depend only on their own factor coordinates and on the
frozen parameters, so their covariant derivatives in direction
$w_e$ vanish.

Gauge invariance gives $df(Q_\xi)=0$ as an identity of functions.
Differentiating it at $x$ yields
\[
 0=H(w_e,Q_\xi)+\langle v_e,[\xi,w]\rangle
   =H(w_e,Q_\xi)-\langle[\xi,v_e],w\rangle .
\]
Since $Q_{\xi,e}(x)=0$, this proves the first part of
\eqref{c18v66:identity}. The dependence of the choice
\eqref{c18v66:parameters} on the base configuration has
\emph{not} been differentiated: at each base point it selects
one fixed global infinitesimal gauge action. Differentiating
a newly selected parameter field instead would not prove
the displayed identity.

Each endpoint has six incident fine links. Apart from their
common edge $e$, the remaining ten links are distinct, and
no other factor couples both endpoints. A single-endpoint
action on a free factor is a signed orthogonal adjoint map
from $\xi$, hence contributes $I_3$ to its Gram matrix.
A midpoint action on a pair is right multiplication of its
first-half coordinate, so its normalized pair velocity is
$\sqrt2$ times an orthogonal adjoint map. It contributes
$2I_3$, accounting exactly for the two captured half-links.

For $e\in\mathcal F_1$, the midpoint endpoint has one pair
and three other free edges; the two-odd endpoint has five
other free edges. Thus the nonzero columns of $B_e$ consist
of one pair and eight free factors. For $e\in\mathcal F_0$,
there are ten free factors and no pairs. The parity description
and $n\ge5$ exclude wraparound identifications of these local
stars. Orthogonality of every adjoint map now gives
$2+8=10$ or $10$, respectively. Inserting $\mathsf W^{-1}$
reduces the pair contribution from $2$ to $1$, proving
\eqref{c18v66:weightedgram}. No holonomy is required to be central.
\end{proof}

The word stabilizer here means a gauge transformation fixing
the selected \emph{factor value}; it need not fix the whole
configuration. No singular gauge quotient or global stabilizer
stratum is invoked.

\begin{theorem}[Forced mixed energy for invariant scalars]
\label{c18v66:mixed}
At every native configuration,
\[
       \mathcal M(f)\ge
          \frac89 F_{1,f}+\frac45 F_{0,f}
          \qquad(f\in C^\infty(M_g)^{\mathcal H_0}).
 \tag{V66.6}\label{c18v66:mixedbound}
\]
\end{theorem}
\begin{proof}
For an orthonormal color basis $(\xi_\alpha)_{\alpha=1}^3$,
\[
 \sum_\alpha|[\xi_\alpha,v_e]|^2=8|v_e|^2.
\]
By the lemma and the Hilbert--Schmidt/operator norm inequality,
\[
 \begin{split}
 8|v_e|^2
 &=\|H_{e,-e}B_e\|_{\rm HS}^2\\
 &\le \|H_{e,-e}\mathsf W^{1/2}\|_{\rm HS}^2
           \|\mathsf W^{-1/2}B_e\|_{\rm op}^2 .
 \end{split}
\]
Consequently
\[
 2\sum_{j\in\mathcal P}\|H_{ej}\|_{\rm HS}^2
 +\sum_{\substack{j\ {\rm free}\\j\ne e}}\|H_{ej}\|_{\rm HS}^2
 \ge
 \begin{cases}
 (8/9)|v_e|^2,&e\in\mathcal F_1,\\
 (4/5)|v_e|^2,&e\in\mathcal F_0.
 \end{cases}
 \tag{V66.7}\label{c18v66:row}
\]
Columns outside the star can be retained on the left because
they add nonnegative terms. Sum over all free $e$.
The free--free blocks have precisely their ordered multiplicity
in $\mathcal M$. Since $H_{je}=H_{ej}^*$, the coefficient two
on free--pair blocks accounts for both orders. The only
remaining terms of $\mathcal M$ are pair--pair squares, which
are nonnegative. This proves \eqref{c18v66:mixedbound}.
\end{proof}

A pointwise horizontal vector by itself has no prescribed
Hessian. Thus \eqref{c18v66:mixedbound} is a statement about
derivatives of invariant scalar functions, not a bound on all
horizontal vector fields or on arbitrary one-forms. In particular
it is not an additional positive matrix term in V65's bare
pointwise pencil unless that pencil is placed back inside
the scalar Bochner argument.

\subsection{An invariant pointwise Bochner estimate}

For any smooth invariant $u$, let
$\nu=Z^{-1}e^u\mu_g$ and
$A_u=\Delta_{m_0}+\langle\nabla u,\nabla\cdot\rangle=-L_u$.
Write
\[
 \Gamma(f)=|\nabla f|^2,\qquad
 \Gamma_{2,u}(f)=\tfrac12A_u\Gamma(f)
                       -\langle\nabla f,\nabla A_u f\rangle .
\]
The product Ricci tensor is $I$ on pairs and $2I$ on free
factors. The Bochner identity and the new mixed estimate give
\[
 \Gamma_{2,u}(f)\ge
 P_f+\frac{26}{9}F_{1,f}+\frac{14}{5}F_{0,f}
             -\operatorname{Hess}u(\nabla f,\nabla f).
 \tag{V66.8}\label{c18v66:bochner}
\]
Indeed the diagonal Hessian squares may be dropped, while
the mixed squares contribute \eqref{c18v66:mixedbound}.
Moreover, by \eqref{c18v66:retainedward},
\[
 \begin{split}
 P_f+\frac{26}{9}F_{1,f}+\frac{14}{5}F_{0,f}
       -\frac{44}{27}(P_f+F_{1,f}+F_{0,f})
   =\frac{17}{27}(2F_{1,f}-P_f)+\frac{158}{135}F_{0,f}
   \ge0 .
 \end{split}
 \tag{V66.9}\label{c18v66:floor}
\]
In particular, if $h$ is a uniform upper bound for the quadratic
form of $\operatorname{Hess}u$ on horizontal unit vectors, then
\[
       \Gamma_{2,u}(f)\ge\rho\Gamma(f),
       \qquad \rho=\frac{44}{27}-h ,
       \qquad f\ {\rm invariant}.
 \tag{V66.10}\label{c18v66:rho}
\]
The sufficient choice $h=\sup_x\|\operatorname{Hess}u(x)\|_{\rm op}$
is permitted, but discards signs. Neither $44/27$ nor the mixed
coefficients are asserted optimal.

\begin{corollary}[Auxiliary invariant semigroup and the native application]
\label{c18v66:semigroup}
For the above fixed law and $P_s=e^{sA_u}$,
\[
       |\nabla P_s f|^2
           \le e^{-2\rho s}P_s|\nabla f|^2
       \quad(s\ge0,\ f\ {\rm smooth\ and\ invariant}).
 \tag{V66.11}\label{c18v66:gradient}
\]
If $\rho>0$, the scalar Poincare gap restricted to invariant
functions is at least $\rho$.

For the actual normal-chart law at $\beta=0$ and $0\le t\le4$,
V60's constant $K_*$ permits
\[
       \rho_t=\frac{44}{27}-tK_* .
 \tag{V66.12}\label{c18v66:native}
\]
Thus \eqref{c18v66:gradient} holds for each fixed $t$, and the
invariant gap conclusion is positive when $tK_*<44/27$.
These statements are uniform in $n,g$ in the declared regime.
\end{corollary}
\begin{proof}
The gauge action is isometric and preserves $u$, so $A_u$
and its semigroup commute with it. Hence every $P_s f$ is
invariant when $f$ is. Fix $s$ and set
$F(r)=P_r\Gamma(P_{s-r}f)$ for $0\le r\le s$.
Differentiation on the compact smooth manifold gives
\[
       F'(r)=2P_r\Gamma_{2,u}(P_{s-r}f)\ge2\rho F(r).
\]
Positivity of $P_r$ and \eqref{c18v66:rho} were used only
on invariant scalar functions. Integrating in $r$ proves
\eqref{c18v66:gradient}; an unrestricted curvature lower
bound is not being assumed.

Alternatively, integration of \eqref{c18v66:rho} gives
$\|L_u f\|_2^2\ge\rho\langle f,L_u f\rangle$.
On the invariant subspace the compact elliptic operator has
a discrete spectral resolution, with constants its only
zero mode. Testing nonconstant invariant eigenfunctions proves
the gap assertion.

Finally V60 bounds the symmetric Hessian block row sum of
$u_t$ by $tK_*$. Symmetry gives the same column sum, so the
block Schur inequality bounds its full operator norm by
$tK_*$. This proves \eqref{c18v66:native} without converting
a sum of pointwise block norms into a sum of separate suprema.
\end{proof}

Here $s$ is auxiliary heat time; $t$ is Wilson smoothing time.
Neither is physical time. V54's exact invariant Haar gap $9$
is much larger than $44/27$; the latter is only a pointwise
Bochner floor obtained after discarding other nonnegative
information. Small-time positivity was already known.
The new contribution is the explicit, configuration-uniform
mixed energy which remains available even away from Haar,
not a claim to have newly established the existence of a
short interval.

If $\rho_t>0$, V61's actual metric comparison gives the scalar
invariant form consequence
$\operatorname{gap}_{\rm inv}(\widehat L_t)\ge E_0^{-2}\rho_t$.
Both sides use the same actual law $\widehat\nu_t$.
This does not transfer \eqref{c18v66:gradient} to the nonproduct
metric, nor bound a pointwise Poisson Hessian or transport strain.
The published majorant $K_*$ does not verify positivity at $t=4$.

\subsection{Putting the energy back into the signed conditional test}

Keep V62's scalar conditional gap floors $\lambda_i$,
uniform in the exterior, and its signed matrix $\mathsf S_u$.
Define the diagonal bundle endomorphism
\[
 (\mathsf C)_{ii}=
 \begin{cases}
 0,&i\in\mathcal P,\\
 (8/9)I,&i\in\mathcal F_1,\\
 (4/5)I,&i\in\mathcal F_0 .
 \end{cases}
\]
For every smooth invariant scalar $\psi$,
\[
 \|L_u\psi\|_2^2
 \ge\int\langle\nabla\psi,(\mathsf S_u+\mathsf C)
                           \nabla\psi\rangle\,d\nu .
 \tag{V66.13}\label{c18v66:conditional}
\]
To prove this, sum V62's directional identity before dropping
its mixed squares. Its conditional integrated Bochner inequality
controls the diagonal terms because $d_i\psi$ is an exact local
gradient. What remains is
$\int\langle\nabla\psi,\mathsf S_u\nabla\psi\rangle+\int\mathcal M(\psi)$.
Apply \eqref{c18v66:mixedbound} pointwise. Thus no diagonal
conditional curvature hypothesis has replaced a scalar gap.

In particular
\[
 P_g(\mathsf S_u+\mathsf C)P_g\succeq\rho P_g
                 \quad\hbox{at every configuration},\quad \rho>0
 \tag{V66.14}\label{c18v66:signedtest}
\]
is sufficient for an invariant scalar gap $\rho$. This uses
$P_g\nabla\psi=\nabla\psi$ followed by the scalar spectral argument;
it does not invert an arbitrary compressed one-form operator.
V63's truncation error for signed off-diagonal blocks is unchanged,
since $\mathsf C$ is diagonal and $\|P_g\|_{\rm op}=1$.

For normalization only, at Haar the diagonal form in
\eqref{c18v66:conditional} has the horizontal floor $62/27$:
\[
 \tfrac32P_f+\tfrac{35}{9}F_{1,f}+\tfrac{19}{5}F_{0,f}
       \ge \tfrac{62}{27}(P_f+F_{1,f}+F_{0,f}).
\]
Subtract the right side and use $P_f\le2F_{1,f}$.
This improves V64's diagonal comparison floor $2$, not the
actual invariant Haar gap $9$. No native signed margin through
time four is inferred from this normalization check.

\subsection{Verification, limitations and the next upstream estimate}

The exact diagnostic
\begin{center}\small
\path{scripts/verify_ym_c18_v66_stabilizer_energy.py}
\end{center}
enumerates every free-edge star on $n=5,6,8$ tori and checks
the pair/free multiplicities. On all 24 parity/orientation
plaquette types it uses noncommuting rational unit quaternions,
with the actual captured-pair coordinate derivatives. It verifies
72 exact stabilizer Hessian rows and 72 transported Gram identities,
the bracket normalization, the midpoint Ward budget, and the
mixed-energy and Ricci bounds. These are native observable
fixtures, not finite-difference approximations. The uniform
statements for all invariant functions rest on the analytic
proofs above; finite fixtures are not proof-assistant certification.

The all-configuration assertion includes V65's central-flux
configuration and its neighborhoods, without estimating their
probability. It does not reverse or invalidate V65's negative
initial signed-comparison witness: it supplies an additional
derivative energy which that pointwise matrix had discarded.
The next useful calculation is whether this retained energy
can absorb the \emph{actual} signed density Hessian on the
invariant response through $t=4$, or whether an actual-law
integrated bad-region estimate is needed. A new upper bound
on its absolute norm alone does not pay that absorption.

The nonlinear normal-chart source and density are still only
the $\beta=0$ construction studied here. The interacting-vacuum
extension, source-specific pointwise Poisson Hessian and strain
through the block time, coarse return/variance, scale transport,
physical-clock calibration, nontrivial reflection-positive
continuum construction and positive physical YM mass gap remain
unproved. The goal remains active.

Removing this append and restoring V65's title recovers V65
byte for byte. Only the active YM source filename is replaced;
all other sources and supplied attachments are preserved.
Only \texttt{.tex} files are kept in \path{latex_notes}.
The companion and broad primary-literature checkpoints remain
V70, and archive/restart remains V100.
% END CYCLE 18 VERSION 66 APPEND
% BEGIN CYCLE 18 VERSION 67 APPEND
\clearpage
\section{V67: Finite-time native central-flux probability and source load}
\label{c18v67:context}

V66 retained a positive mixed-derivative energy, but did not
control the actual density curvature through Wilson time four.
V65's central-flux witness therefore cannot simply be declared
irrelevant on grounds of rarity. This append pays an actual
probability estimate for a specified adverse family: small
gauge-invariant neighborhoods of central configurations with
a sufficiently negative captured Wilson sum. The proof uses
the literal normal-chart Jacobian, not an exponential-tilt model
or a presumed mixing law.

The new estimate holds through the full smoothing interval,
uniformly in volume, at the electric origin $\beta=0$ and
coarse identity $g=I$. It also bounds the temporal source
carried by those neighborhoods. It does not identify them
with the entire set of negative curvature, control the
Poisson inverse on them, or construct the interacting continuum.
The distinction between a small source load and a controlled
resolvent remains essential.

\subsection{An exact stationary-point Jacobian on the normal path}

Use V51's ambient fine product metric $\mathfrak g$, Wilson
gradient $F=\nabla W$, flow $\Phi_t$, coarse map
$c_t=b_2\Phi_t$, and normal identification $\eta_t$:
\[
 C_t=Dc_t,\quad Q_t=C_t^*(C_tC_t^*)^{-1}C_t,\quad
 \dot\eta_t=-Q_tF\circ\eta_t,\qquad c_t\eta_t=b_2.
\]
All Jacobians below are ambient Haar Jacobians. At $\beta=0$,
the actual density on the fixed fibre is
$\widehat p_t=k_t/Z_t$, where $k_t=\operatorname{Jac}\eta_t$
and $Z_t=\int k_t\,d\mu_g$. No new fibre-volume factor is inserted.

\begin{lemma}[Stationary native determinant and tangent isometry]
\label{c18v67:stationary}
Let $p$ be an ambient stationary point of $F$, put
$A=\nabla F(p)$ and $B=Db_2(p)$, and use the induced tangent
inner products at the fixed point and its coarse image.
For $0\le t\le4$,
\[
 \eta_t(p)=\Phi_t(p)=p,\qquad C_t(p)=Be^{tA},
 \qquad
 k_t(p)=
 \left(\frac{\det(Be^{2tA}B^*)}{\det(BB^*)}\right)^{-1/2}.
 \tag{V67.1}\label{c18v67:det}
\]
Moreover, on the initial fibre tangent $\ker B$,
\[
             (D\eta_t(p))^*D\eta_t(p)=I .
 \tag{V67.2}\label{c18v67:isometry}
\]
Thus $\widehat m_t(p)=m_0(p)$ at such a point. This is
a value of the metric, not a bound on its spatial derivatives.
\end{lemma}
\begin{proof}
The point is fixed by $\Phi_t$ since $F(p)=0$, and by $\eta_t$
since $Q_t(p)F(p)=0$. The linearized autonomous gradient flow
is $D\Phi_t(p)=e^{tA}$, with $A=A^*$, proving the formula for
$C_t$. Write $J_t=D\eta_t(p)$. Since the derivative of $Q_t$
is multiplied by $F(p)=0$, its variational equation is
\[
                 J'_t=-Q_t(p)A J_t,\qquad J_0=I .
\]
Consequently
\[
 \partial_t\log\det J_t=-\operatorname{tr}(Q_tA)
       =-\tfrac12\partial_t\log\det(Be^{2tA}B^*).
\]
Integration from zero proves \eqref{c18v67:det}. Positivity
of the Jacobian follows from the flow starting at the identity.

Differentiating $c_t\eta_t=b_2$ gives $C_t(p)J_t=B$.
If $v,w\in\ker B$, then $J_tv,J_tw\in\ker C_t(p)$, whereas
$Q_t$ projects onto its orthogonal complement. Hence
\[
 \partial_t\langle J_tv,J_tw\rangle
 =-\langle Q_tAJ_tv,J_tw\rangle
  -\langle J_tv,Q_tAJ_tw\rangle=0 .
\]
Its initial value is $\langle v,w\rangle$, proving
\eqref{c18v67:isometry}. This does not say that $J_t$ is
an ambient isometry; its normal determinant generally changes.
\end{proof}

We will use the following elementary finite-matrix fact.
If $A=A^*$ and $V^*V=I$, set
\[
 \phi_A(t)=\log\det(V^*e^{2tA}V),\qquad
 \Pi_t=e^{tA}V(V^*e^{2tA}V)^{-1}V^*e^{tA}.
\]
Direct differentiation gives
\[
 \begin{aligned}
 \Pi'_t&=(I-\Pi_t)A\Pi_t+\Pi_tA(I-\Pi_t),\\
 \phi'_A(t)&=2\operatorname{tr}(\Pi_tA),\\
 \phi''_A(t)&=4\|(I-\Pi_t)A\Pi_t\|_{\rm HS}^2\ge0 .
 \end{aligned}
 \tag{V67.3}\label{c18v67:convex}
\]
For completeness, write $E_t=e^{tA}V$, differentiate
$\Pi_t=E_t(E_t^*E_t)^{-1}E_t^*$, and use $E'_t=AE_t$.
Taking the trace of $2\Pi'_tA$ gives the last equality.
Thus
\[
       \phi_A(t)\ge2t\,\operatorname{tr}(V^*AV)
       \qquad(t\ge0).
 \tag{V67.4}\label{c18v67:tangent}
\]
No positivity of $A$ is required. This permits mixed central
fluxes, whose Wilson Hessians need not have one sign.

\subsection{Central configurations and their actual density ratios}

Keep $N=2n$, $n\ge5$, and $g=I$. Let $\mathcal Z$ consist
of the points of $M_I$ at which every fine link is in
$\{I,-I\}$. In pair/free coordinates it has $21n^3$
independent binary variables: the two captured halves
of a pair have the same sign because their product is $I$.
Write $p_+$ for the all-identity point and
\[
 C(z)=W_{\rm cap}(z)=\sum_{p\ {\rm captured}}s_p(z),
 \qquad s_p(z)\in\{-1,1\},\qquad z\in\mathcal Z .
\]
There are $12n^3$ captured plaquettes. Let $d_p$ be the
oriented real fine-link curl row around plaquette $p$.

\begin{proposition}[Finite-time central likelihood bound]
\label{c18v67:likelihood}
Every $z\in\mathcal Z$ is stationary for the native Wilson
flow. In the fine orthonormal color frames,
\[
 A_z=-\sum_p s_p(z)d_p^*d_p\otimes I_3,\quad
 BB^*=2I_{9n^3},\quad V=B^*/\sqrt2,\quad
 \operatorname{tr}(V^*A_zV)=-3C(z).
 \tag{V67.5}\label{c18v67:trace}
\]
Here $B$ is the same pair-sum matrix at all these central points.
For the actual normalized normal-chart density,
\[
       \frac{\widehat p_t(z)}{\widehat p_t(p_+)}
                    \le e^{\,3tC(z)},\qquad 0\le t\le4 .
 \tag{V67.6}\label{c18v67:likelihoodbound}
\]
\end{proposition}
\begin{proof}
A first trace derivative of a central plaquette vanishes.
Its scalar second derivative is $-s_p(z)$ times the squared
oriented sum of the link color variations. Commutator terms
have zero scalar part at this order. This gives the first
identity, also used for the all-positive and all-negative
fields in V38 and V65.

Central adjoint transports are identity, so $B$ sums the
two variations in each captured pair. The normalized normal
column has $1/\sqrt2$ on each half, proving $BB^*=2I$
and $V^*V=I$. An uncaptured plaquette has no such columns.
A captured plaquette uses two distinct captured pairs, one
half from each. Its curl Gram therefore has trace
$1/2+1/2=1$ on these normal columns, per color. Summing
the signed contributions over three colors proves
$\operatorname{tr}(V^*A_zV)=-3C(z)$.
No independence of plaquette signs is assumed.

By the stationary determinant formula and
\eqref{c18v67:tangent},
\[
              \log k_t(z)\le3tC(z).
\]
At $p_+$, $A_{p_+}=-\mathcal L$, where
$\mathcal L=d_1^*d_1\otimes I_3\succeq0$ is the full
ordinary fine curl Gram. Hence
$V^*e^{-2t\mathcal L}V\preceq I$, and
$\log k_t(p_+)\ge0$. Subtraction gives
\eqref{c18v67:likelihoodbound}, since the same conditional
normalizer $Z_t$ cancels. No bound on that coarse marginal
or its free energy is needed.
\end{proof}

In particular, the inherited all-negative field
\[
 (p_-)_{x,i}=(-1)^{\sum_{j<i}x_j}I
 \tag{V67.7}\label{c18v67:minus}
\]
is periodic on every even fine torus, has both captured
halves $I$, and has every fine plaquette $-I$.
Thus it belongs to the same fibre $M_I$, and
\[
           \widehat p_t(p_-)/\widehat p_t(p_+)
                    \le e^{-36tn^3}.
 \tag{V67.8}\label{c18v67:minusratiobound}
\]
This field itself is old; the finite-time likelihood bound
and the probability consequences below are the new work.
The estimate is not a probability assigned to a point.

\subsection{From point values to gauge-invariant neighborhoods}

Let $K_*$ be V60's actual Hessian constant, without any
improvement to its very large numerical value. V60 and the
symmetric block Schur inequality give, throughout $[0,4]$,
\[
       \|\operatorname{Hess}_{m_0}u_t\|_{\infty,\rm op}
                      \le tK_*,\qquad u_t=\log\widehat p_t.
 \tag{V67.9}\label{c18v67:H}
\]
At every $z\in\mathcal Z$ one has $\nabla u_t(z)=0$.
Indeed the entire construction is equivariant under
simultaneous global conjugation, which fixes coarse identity
and every central field. Its tangent action is the adjoint
three-dimensional color rotation on each product factor.
There is no nonzero vector fixed by all those rotations,
so an invariant scalar has zero gradient there. This
argument uses a symmetry fixing the point, not a claim
that the pinned internal gauge action has stabilizers.

For $\varepsilon>0$, form a product geodesic box
\[
 B_z(\varepsilon)
   =\{x\in M_I:\ d_i(x_i,z_i)<\varepsilon
                         \text{ for every product factor }i\},
 \qquad
 \mathcal U_z(\varepsilon)=\mathcal H_0 B_z(\varepsilon).
\]
The latter is open and gauge invariant. Product-factor distances
and Haar measure here use the pair/free metric $m_0$.
The fixed product has $m=21n^3$ factors.

Multiplication by the central factor signs of $z$ defines an
involutive product isometry $\sigma_z:M_I\to M_I$ sending
$p_+$ to $z$. It preserves $\mu_I$, commutes with
$\mathcal H_0$, and sends $B_{p_+}(\varepsilon)$ and its
gauge saturation onto their counterparts at $z$.
This does not assert that $\sigma_z$ preserves the time-$t$
density or commutes with the normal flow.

\begin{theorem}[Actual-law suppression of central neighborhoods]
\label{c18v67:neighborhood}
Fix $0<\eta\le12$ and choose the volume-independent radius
\[
 \varepsilon_\eta=
     \min\left\{\frac{\pi}{4},
                  \sqrt{\frac{\eta}{14K_*}}\right\}.
 \tag{V67.10}\label{c18v67:radius}
\]
If $z\in\mathcal Z$ and $C(z)\le-\eta n^3$, then
\[
 \widehat\nu_t(\mathcal U_z(\varepsilon_\eta))
       \le e^{-(3/2)\eta tn^3}
                   \widehat\nu_t(\mathcal U_{p_+}(\varepsilon_\eta))
       \le e^{-(3/2)\eta tn^3}.
 \tag{V67.11}\label{c18v67:singlebox}
\]
Define the gauge-invariant adverse neighborhood family
\[
 \mathcal E_\eta=
    \bigcup_{\substack{z\in\mathcal Z\\C(z)\le-\eta n^3}}
                       \mathcal U_z(\varepsilon_\eta).
\]
Then, for every $n\ge5$ and $0\le t\le4$,
\[
 \widehat\nu_t(\mathcal E_\eta)\le
 q_\eta(t,n):=
 \min\{1,\exp[(14\log2-\tfrac32\eta t)n^3]\}.
 \tag{V67.12}\label{c18v67:union}
\]
In particular $\eta=6$ means that at least three quarters
of the captured plaquettes of the central configuration
are negative, and
\[
        \widehat\nu_4(\mathcal E_6)
          \le e^{-(36-14\log2)n^3}<e^{-26n^3}.
 \tag{V67.13}\label{c18v67:time4}
\]
\end{theorem}
\begin{proof}
For $x\in B_{p_+}(\varepsilon_\eta)$, the product minimizing
geodesic and \eqref{c18v67:H}, together with zero gradients
at the two central endpoints, give
\[
 \begin{aligned}
 u_t(\sigma_zx)-u_t(x)
 &\le u_t(z)-u_t(p_+)
             +tK_*\sum_{i=1}^{m}d_i(x_i,(p_+)_i)^2\\
 &\le -3\eta tn^3+21tK_*\varepsilon_\eta^2n^3
 \le-\tfrac32\eta tn^3 .
 \end{aligned}
 \tag{V67.14}\label{c18v67:paired}
\]
Each of the two Taylor bounds costs half of the displayed
quadratic error. The choice $\varepsilon_\eta\le\pi/4$
stays within the injectivity radius of every factor.
For $x=h y$ in $\mathcal U_{p_+}$, gauge invariance of $u_t$
and $\sigma_z h=h\sigma_z$ transfer the same inequality
from $y$ to $x$. Changing variables by the measure-preserving
$\sigma_z$ proves \eqref{c18v67:singlebox}. No estimate of
orbit volumes is inserted.

It remains to count the union without pretending that
the plaquettes are independent. The pinned center-valued
gauge group has $2^{7n^3}$ elements, because there are
$8n^3-n^3$ internal vertices. Its action on the
$2^{21n^3}$ central product configurations is free.
Indeed a gauge transformation fixing all fine links must
have equal center values across every adjacent edge; the
connected fine lattice and pinned coarse roots force all
of them to be $I$. Thus there are $2^{14n^3}$ center-gauge
orbits.

These are also exactly the intersections of full pinned
gauge orbits with $\mathcal Z$: if $h$ sends one central
configuration to another, adjacent vertex values of $h$
differ by a center sign. Starting at a pinned root forces
every vertex value of $h$ to be central. Gauge-related
central points have the same captured Wilson sum, and
their saturated boxes coincide, since the action is a
product-factor isometry. Taking one representative per
eligible orbit and using the union bound proves
\eqref{c18v67:union}, irrespective of overlaps.

For $\eta=6$, $C=12n^3-2N_-$ gives $N_-\ge9n^3$.
Substitute $t=4$ in \eqref{c18v67:union}.
The elementary bound $\log2<5/7$ gives the last inequality
in \eqref{c18v67:time4}; for example
$1+5/7+(5/7)^2/2+(5/7)^3/6>2$ proves it.
\end{proof}

This is a probability bound for all gauge transforms of the
specified boxes, not just a thin coordinate slice through
an orbit. Nevertheless the radius is extremely small because
$K_*$ is huge. The set $\mathcal E_6$ is \emph{not} the full
event $W_{\rm cap}(x)\le-6n^3$, nor is it proved to contain
every adverse-curvature configuration. The bound is specific
to coarse identity and $\beta=0$. No uniform statement over
coarse fields or interacting vacua is inferred.

\subsection{The temporal source carried by the suppressed family}

The sets just defined do not depend on $t$. Let
$a_t=\partial_tu_t$, so $\int a_t\,d\widehat\nu_t=0$.
V60's $H(a_t)\le K_*$ and V61's one-factor gradient estimate,
with $D=\pi\sqrt2$, imply
\[
 \|\nabla_i a_t\|_\infty\le DK_*,
 \qquad \|a_t\|_\infty\le mD^2K_*,
 \qquad \|\nabla a_t\|_\infty^2\le mD^2K_*^2.
 \tag{V67.15}\label{c18v67:sourceuniform}
\]
For the middle bound, connect any two product points one
factor at a time. The oscillation is at most
$\sum_i D\|\nabla_i a_t\|_\infty$; the mean-zero condition
places zero between the minimum and maximum. These are
extensive bounds, not falsely claimed anchored inverse estimates.

Multiplication by \eqref{c18v67:union} gives the actual
source-load estimates
\[
 \begin{aligned}
 \int_{\mathcal E_\eta} a_t^2\,d\widehat\nu_t
       &\le m^2D^4K_*^2 q_\eta(t,n),\\
 \int_{\mathcal E_\eta}|\nabla a_t|^2\,d\widehat\nu_t
       &\le mD^2K_*^2 q_\eta(t,n),\\
 \left|\partial_t\widehat\nu_t(\mathcal E_\eta)\right|
       &=\left|\int_{\mathcal E_\eta}a_t\,d\widehat\nu_t\right|
         \le mD^2K_* q_\eta(t,n).
 \end{aligned}
 \tag{V67.16}\label{c18v67:sourceload}
\]
The derivative follows by differentiating a smooth density
over a fixed measurable set at finite volume. At time four
and $\eta=6$, all three right sides have the exponential
factor in \eqref{c18v67:time4}, with the displayed polynomial
volume prefactors. Their constants can still be enormous.

These bounds do not yet control
$\int_{\mathcal E_6}|\nabla L_t^{-1}a_t|^2\,d\widehat\nu_t$.
Small probability alone cannot control the fraction of gradient
energy of an arbitrary invariant scalar supported there.
To see this in the actual geometry, choose a nonzero
nonnegative smooth bump supported inside the nonempty open
$\mathcal E_6$ and average it over $\mathcal H_0$.
Its support remains a compact subset of $\mathcal E_6$;
it is nonzero and nonconstant, since \eqref{c18v67:time4}
makes $\mathcal E_6$ proper. All its gradient energy is in
$\mathcal E_6$, so the ratio of that energy to its total
energy is one, however small the probability bound.
Subtracting its mean does not change this observation.
The actual Poisson response is not an arbitrary such bump;
a source-to-region resolvent or capacity estimate is still owed.

\subsection{Finite-mode verification and the next nonperturbative task}

The stationary density is also directly computable without
integrating a nonlinear trajectory. For the all-negative point,
$\mathcal L=d_1^*d_1\otimes I_3$ is translation invariant.
For one color and coarse momentum
$k\in\{0,\ldots,n-1\}^3$, fold the eight fine momenta
$\theta_i=\pi k_i/n+\pi r_i$, $r\in\{0,1\}^3$, and put
\[
 d_i=e^{i\theta_i}-1,\quad
 c_i=\frac{1+e^{-i\theta_i}}4,\quad
 D_c=\operatorname{diag}(c_i),\quad
 \lambda=\sum_i|d_i|^2,\quad
 M_\theta=D_c^*\left(I-\frac{dd^*}{\lambda}\right)D_c .
\]
For $\lambda=0$ set $M_\theta=0$ in the next formula.
Normalized coarse-to-fine Fourier transformation of a pair
column gives the coefficient $c_i$, and the fine curl symbol
is $\lambda I-dd^*$. Therefore
\[
 G_k(s)=I+\sum_{r\in\{0,1\}^3}
                  (e^{2s\lambda}-1)M_\theta,\qquad
 \phi_{\mathcal L}(s)=3\sum_k\log\det G_k(s).
 \tag{V67.17}\label{c18v67:bloch}
\]
The identity term follows from
$\sum_rD_c^*D_c=I$. The factor three is color multiplicity.
The exact central log likelihood ratio is
$\tfrac12[\phi_{\mathcal L}(-t)-\phi_{\mathcal L}(t)]$.
This is a finite-time formula, not a Taylor truncation.

The diagnostic
\begin{center}\small
\path{scripts/verify_ym_c18_v67_native_flux_probability.py}
\end{center}
enumerates native $n=5,6,8$ incidences, verifies the captured
normal trace for uniform and mixed center signs, and checks
the pinned center-gauge action rank over $\mathbb F_2$.
A small symbolic matrix fixture checks the projection,
log-determinant and tangent-isometry identities; it is
explicitly not a native lattice configuration.
The native Bloch regression uses 100-digit arithmetic for
$n=5,6$ at $t=0,1/32,1/4,1,4$. It checks the determinant
and trace formulas, with the all-volume inequalities supplied
by the analytic proofs, not by floating-point tolerances.
No interacting trajectory or proof-assistant certificate is claimed.

This advances beyond the previous absence of an actual
probability estimate for the central witness. The next task
is to control the remaining noncentral adverse regions and
the actual invariant source response, without replacing
source load by an assumed inverse bound. One possible extension
is a local central-flux comparison with a boundary cost;
no such contour or boundary comparison has been proved here.
V66's mixed energy remains available for an eventual integrated
absorption estimate, but that absorption is not supplied by
\eqref{c18v67:sourceload} alone.

The source-specific strain through time four, interacting
vacuum extension, coarse return, weak-coupling scale transport,
physical-clock calibration, nontrivial reflection-positive
continuum construction and positive physical Yang--Mills
mass gap remain unproved. The goal remains active.

Removing this append and restoring V66's title recovers
V66 byte for byte. Other sources are preserved and only
the active YM filename is replaced. Only \texttt{.tex}
files are kept in \path{latex_notes}; build and diagnostic
products stay outside. The next companion and broad primary
literature checkpoints remain V70, and archive/restart V100.
% END CYCLE 18 VERSION 67 APPEND
% BEGIN CYCLE 18 VERSION 68 APPEND
\clearpage
\section{V68: Local native flux droplets with surface cost and arbitrary exterior}
\label{c18v68:context}

V67 proved rarity of specified globally near-central configurations.
It did not give a local droplet estimate with an uncontrolled
exterior. We now localize the stationary determinant comparison.
Dropping plaquettes crossing a block boundary costs surface area,
not the volume of the surrounding torus. A boundary-layer sum
using V63's weighted Hessian bound then allows the entire exterior
to be arbitrary and noncentral, at another explicit surface cost.
The resulting statement concerns actual normalized conditional
probabilities and their temporal scores throughout $0\le t\le4$.

The scope remains $\beta=0$, coarse identity $g=I$, fine side
$N=2n$, $n\ge5$. The prescribed near-central condition is only
inside the block; the constants are very large. Neither this
local suppression nor its conditional form is a mixing theorem
or a proof of the physical mass gap.
The stationary formula and weighted spatial bounds are inherited
from V67 and V63; their local combination is the new step.

\subsection{A native cut that does not split captured pairs}

Fix a coarse cube $R=\{0,\ldots,\ell-1\}^3$, with $2\le\ell<n$;
coarse translates are allowed. Let $D_R$ be its fine vertex cube
of side $L=2\ell$. Assign a stored fine edge to the interior if
its \emph{tail} is in $D_R$. Write $\mathsf E_R$ for these edges.
This is an edge-index partition, not an assertion that every
assigned edge has both endpoints in $D_R$.

Both halves of each captured pair belong to the same part:
their tails lie in the same coarse cell. Thus the pair/free
factor set splits as $\mathcal I_R\sqcup\mathcal I_R^c$, with
\[
                  m_R:=|\mathcal I_R|=21\ell^3.
\]
The normalized coarse normal inclusion $V=B^*/\sqrt2$ of V67
also splits into interior and exterior columns.

Call a plaquette a boundary plaquette if its four stored edges
are not all in one part, and let $b_R$ be their number. The
exact counts are
\[
 b_R=48\ell^2-6\ell,\qquad
 N_{{\rm cap},R}=12\ell^3-12\ell^2+3\ell ,
 \tag{V68.1}\label{c18v68:counts}
\]
where the second number counts captured plaquettes wholly inside
$\mathsf E_R$. To verify them, a plaquette in plane $i,j$ has
tails $x,x+e_i,x+e_j$. There are $2L^2-L$ mixed triples with
$x$ inside and $2L^2$ with $x$ outside, giving $4L^2-L$ per
plane. A wholly interior plaquette has $L-1$ choices for each
of $x_i,x_j$ and $L$ choices for the remaining coordinate.
It is captured exactly when that remaining coordinate is even.
There are $\ell(L-1)^2$ such plaquettes per plane. Summing the
three planes proves \eqref{c18v68:counts}. The condition $\ell<n$
keeps opposite boundary layers distinct, including after translation.

For a central configuration $z$, retain V67's actual stationary
Wilson Hessian
\[
       A_z=-\sum_p s_p(z)d_p^*d_p\otimes I_3 .
\]
Delete only boundary plaquette terms to obtain the auxiliary
matrix $A_z^\circ=A_{z,R}\oplus A_{z,R^c}$. It is used as a
matrix comparison, not as a substitute native probability law.
One has
\[
 \|A_z\|_{\rm op},\ \|A_z^\circ\|_{\rm op}\le12,\qquad
 \|A_z-A_z^\circ\|_1\le12b_R .
 \tag{V68.2}\label{c18v68:cutnorm}
\]
Indeed the unsigned fine curl Gram has Fourier symbol
$\lambda I-dd^*$, with $\lambda=\sum_i|e^{i\theta_i}-1|^2\le12$.
Inserting signs, zeros, or coefficients of absolute value at
most one between the curl and its adjoint preserves this norm
bound. Each removed plaquette has a rank-one Gram of trace four
per color, hence trace norm twelve over three colors. The same
operator-norm bound holds along the straight interpolation to
the cut matrix.

\subsection{A noncommuting determinant boundary estimate}

For an isometry $V$ and a self-adjoint matrix $A$, define
\[
       \mathcal F_t(A)
          =-\tfrac12\log\det(V^*e^{2tA}V).
\]
At central configurations this is exactly $\log k_t(z)$,
the raw normal-chart log density, by V67. No monotonicity of
the matrix exponential in the Loewner order is assumed.

\begin{lemma}[Trace-norm continuity of the compressed log determinant]
\label{c18v68:logdet}
If $A,B$ are self-adjoint and $\|A\|,\|B\|\le12$, then
\[
 |\mathcal F_t(A)-\mathcal F_t(B)|
       \le c(t)\|A-B\|_1,\qquad
       c(t)=\frac{e^{24t}-1}{24},\qquad t\ge0.
 \tag{V68.3}\label{c18v68:modulus}
\]
\end{lemma}
\begin{proof}
Along a self-adjoint perturbation $E$, Duhamel differentiation
and cyclicity of trace give
\[
 D\mathcal F_t(A)[E]
 =-\tfrac12\int_{-t}^{t}
      \operatorname{tr}\bigl(e^{rA}\Pi_t e^{-rA}E\bigr)\,dr,
 \qquad
 \Pi_t=e^{tA}V(V^*e^{2tA}V)^{-1}V^*e^{tA}.
 \tag{V68.4}\label{c18v68:derivative}
\]
For example, differentiate $e^{2tA}$ as
$\int_0^{2t}e^{(2t-s)A}Ee^{sA}\,ds$ inside the log determinant
and set $r=s-t$ after moving factors around the trace.
The matrix $\Pi_t$ is an orthogonal projection. Thus
\[
 \|e^{rA}\Pi_t e^{-rA}\|_{\rm op}
       \le e^{|r|(\lambda_{\max}(A)-\lambda_{\min}(A))}
       \le e^{24|r|}.
\]
The trace-norm/operator-norm inequality and integration of
$\tfrac12e^{24|r|}$ give
$|D\mathcal F_t(A)[E]|\le c(t)\|E\|_1$.
Integrate along $(1-s)A+sB$, whose norm is at most twelve.
\end{proof}

A central interior pattern is a choice
$w\in\{I,-I\}^{\mathcal I_R}$. Its captured halves inherit the
same sign, as required by $g=I$. Let
\[
 C_R(w)=\sum_{\substack{p\ {\rm captured}\\
                             {\rm edges}(p)\subset\mathsf E_R}}s_p(w).
\]
For any central exterior $\xi$, write $z^w=(w,\xi)$ and
$z^+=(I,\xi)$. The normal columns and cut exponential split
over the edge partition. The exterior block is identical
for $z^w,z^+$, so its determinant cancels. On the interior,
\[
 \operatorname{tr}(V_R^*A_{z^w,R}V_R)=-3C_R(w).
\]
This is V67's normal-column count restricted to wholly interior
plaquettes. Its compressed-exponential convexity therefore gives
$\mathcal F_t(A_{z^w,R})\le3tC_R(w)$.
The repaired interior has negative semidefinite Hessian, hence
$\mathcal F_t(A_{z^+,R})\ge0$. Applying
\eqref{c18v68:modulus} once to each full matrix gives
\[
 u_t(z^w)-u_t(z^+)
 \le 3tC_R(w)+\mathfrak b_t(R),\qquad
 \mathfrak b_t(R)=b_R(e^{24t}-1).
 \tag{V68.5}\label{c18v68:center}
\]
Here $u_t=\log\widehat p_t$ is the actual normalized density.
Its common normalizer cancels. The coefficient of $b_R$ is
$2\cdot12c(t)=e^{24t}-1$, including all three colors.
The estimate is uniform over every central exterior and the
size of the surrounding torus.

\subsection{Finite sign changes have spatially summable density response}

Keep $D=\pi\sqrt2$ and V60--V63's paid constants
\[
 H(u_t)\le tK_*,\qquad
 \mathcal H_\zeta(u_t)\le tK_\zeta,\qquad
 K_\zeta=e^{2\zeta}K_*,\qquad \zeta>0.
\]
Distances use V63's fine-lattice anchors: pair first-half
tails and free-link tails. No new density regularity is assumed.

Let $\sigma$ multiply selected product factors by center signs,
and let $S$ be the set of changed factors, $k=|S|$.
It is a product isometry. Put
$\Psi_t(x)=u_t(\sigma x)-u_t(x)$. Then, pointwise,
\[
 \sum_j|\nabla_j\Psi_t(x)|\le3kDtK_*.
 \tag{V68.6}\label{c18v68:flipgradient}
\]
If $\mathcal T$ is a set of unchanged factors whose anchors
have distance greater than $r$ from every factor of $S$, then
\[
 \sum_{j\in\mathcal T}|\nabla_j\Psi_t(x)|
                     \le kDtK_\zeta e^{-\zeta r}.
 \tag{V68.7}\label{c18v68:fliptail}
\]

Here is a proof retaining the placement of the supremum.
For $j\in S$, the isometry formula for the gradient and
$G(u_t)\le DH(u_t)$ bound its contribution by $2DtK_*$.
For $j\notin S$, the $j$-coordinates of $x,\sigma x$ coincide.
Connect those points by simultaneous minimizing geodesics
in the changed factors only, each of length at most $D$.
Parallel transport in the unchanged factor is identity, and
\[
 \nabla_j\Psi_t
       =\int_0^1\sum_{i\in S}
             (\operatorname{Hess}u_t)_{ji}(x(s))\,\dot x_i(s)\,ds .
\]
Sum the norms over $j\notin S$, interchange finite sums and
the integral, and use symmetry of the Hessian blocks and the
row bound at the same configuration $x(s)$. This costs
$kDtK_*$. Together with the changed rows it proves
\eqref{c18v68:flipgradient}. Restricting the sum to $\mathcal T$
and using the weighted row bound instead gives
\eqref{c18v68:fliptail}. In particular we did not replace a
supremum of row sums by a sum of separate entry suprema.

\subsection{Paying for an entirely arbitrary exterior}

For a factor anchored inside the fine cube, let $\rho_i$ be
the fine vertex distance from its anchor to the complement
of $D_R$. There are at most three factor labels at any anchor:
pair first halves and free links form disjoint subsets of
the three stored edge directions. A six-face estimate gives
\[
 \sum_{i\in\mathcal I_R}e^{-\zeta\rho_i}
        \le\frac{72\ell^2}{e^\zeta-1}.
 \tag{V68.8}\label{c18v68:layers}
\]
Indeed, for $L=2\ell$, the exponential of minus the minimum
distance to the six faces is at most the sum of their six
exponentials. Each face contributes at most
$L^2\sum_{q\ge1}e^{-\zeta q}=L^2/(e^\zeta-1)$.
Multiplication by the anchor multiplicity three proves the
bound. It remains valid for translated cubes on the torus.

Apply the proof of \eqref{c18v68:fliptail} separately to each
changed interior factor, retaining its own distance to the
exterior. For any sign change supported in $\mathcal I_R$,
\[
 \sum_{j\notin\mathcal I_R}|\nabla_j\Psi_t|
   \le DtK_\zeta
             \sum_{i\in\mathcal I_R}e^{-\zeta\rho_i}
   \le\frac{72DtK_\zeta}{e^\zeta-1}\ell^2 .
 \tag{V68.8a}\label{c18v68:surfacegradient}
\]
Consequently changing \emph{all} exterior coordinates arbitrarily,
with the interior fixed, changes $\Psi_t$ by at most
$tB_\zeta\ell^2$, where
\[
 B_\zeta=\frac{72D^2K_\zeta}{e^\zeta-1},\qquad
 \mathfrak s_t(R)=\mathfrak b_t(R)+tB_\zeta\ell^2 .
 \tag{V68.8b}\label{c18v68:surfacecost}
\]
Use simultaneous product geodesics of length at most $D$
to see this directly from \eqref{c18v68:surfacegradient}.
This is a surface penalty, not a restriction on an intermediate
collar. Its constant is finite but extremely large.

Fix $0<\eta\le12$ and choose
\[
 \epsilon_\eta=
       \min\left\{\frac{\pi}{4},\frac{\eta}{42DK_*}\right\},
 \qquad
 \mathcal A_+=\{x:\ d_i(x_i,I)<\epsilon_\eta
                                    \ (i\in\mathcal I_R)\}.
 \tag{V68.8c}\label{c18v68:radius}
\]
Every exterior coordinate is unrestricted. For a central interior
pattern $w$, let $\sigma_w$ multiply only interior factor
coordinates by its signs, and set
\[
 \mathcal A_w=\sigma_w\mathcal A_+,\qquad
 \mathcal E_+=\mathcal H_0\mathcal A_+,\qquad
 \mathcal E_w=\mathcal H_0\mathcal A_w=\sigma_w\mathcal E_+.
\]
The last equality uses that central sign multiplication
commutes with the pinned gauge action. Because the action
is factorwise, these saturated events still depend only
on factors in $\mathcal I_R$. No exterior variable is
silently required to be near a center element.

\begin{theorem}[Native local repair comparison]
\label{c18v68:repair}
If $C_R(w)\le-\eta\ell^3$, then for every $x\in\mathcal E_+$,
\[
 u_t(\sigma_wx)-u_t(x)
       \le-\kappa_t,\qquad
 \kappa_t=\tfrac32\eta t\ell^3-\mathfrak s_t(R),
       \qquad 0\le t\le4 .
 \tag{V68.9}\label{c18v68:repairpoint}
\]
Let $\widehat\nu_t^R(\cdot\mid\xi)$ be the canonical conditional
law in factors $\mathcal I_R$, with every other factor fixed
to $\xi$. For every exterior $\xi$,
\[
 \widehat\nu_t^R(\mathcal E_w\mid\xi)
      \le e^{-\kappa_t}
                   \widehat\nu_t^R(\mathcal E_+\mid\xi).
 \tag{V68.10}\label{c18v68:conditional}
\]
The same comparison holds without conditioning.
\end{theorem}
\begin{proof}
First take $x\in\mathcal A_+$. Choose a full central point
$z^+$ equal to $I$ in the interior and to any center signs
outside. Put $\Psi_t=u_t\circ\sigma_w-u_t$ and
$k=|\operatorname{supp}\sigma_w|\le21\ell^3$.
Equation \eqref{c18v68:center} gives
$\Psi_t(z^+)\le-3\eta t\ell^3+\mathfrak b_t(R)$.

Change only the interior coordinates of $z^+$ to those of
$x$ along simultaneous product geodesics, each of length
at most $\epsilon_\eta$. Equation \eqref{c18v68:flipgradient}
bounds the change of $\Psi_t$ by
\[
          3kDtK_*\epsilon_\eta
                       \le\tfrac32\eta t\ell^3 .
\]
Then change all exterior coordinates to their arbitrary
values in $x$, each along a geodesic of length at most $D$.
The paid surface estimate \eqref{c18v68:surfacegradient}
bounds this second change by $tB_\zeta\ell^2$.
Both estimates use gradient sums at the current configuration,
not sums of separate suprema. Adding them proves
\eqref{c18v68:repairpoint} on $\mathcal A_+$.

The density is pinned-gauge invariant and $\sigma_w$ commutes
with that action. Thus $\Psi_t$ is invariant, extending the
same bound to $\mathcal E_+$.
Finally $\sigma_w$ fixes all exterior factors, preserves
interior product Haar measure, and maps each exterior slice
of $\mathcal E_+$ bijectively onto that of $\mathcal E_w$.
Change variables in the conditional density
\[
 \frac{e^{u_t(y,\xi)}\,d\mu_R(y)}
           {\int e^{u_t(y',\xi)}\,d\mu_R(y')}.
\]
Its normalizer is identical on the two sides. This proves
\eqref{c18v68:conditional} for the positive-density version
defined at every exterior, including empty slices. Integration
over the exterior proves the unconditioned version.
\end{proof}

Thus the repair comparison allows every exterior, including
arbitrary noncentral configurations immediately outside the
block. It has not replaced their law by a central approximation:
their effect was bounded by \eqref{c18v68:surfacegradient}.
The \emph{interior} event remains a specified small neighborhood
of central patterns and their gauge transforms.

\subsection{Pattern entropy and a finite-time droplet bound}

Let $\mathcal W_{\eta,R}$ be all interior central patterns
with $C_R(w)\le-\eta\ell^3$ and let
$\mathcal E_{\eta,R}=\bigcup_{w\in\mathcal W_{\eta,R}}\mathcal E_w$.
There are at most $2^{21\ell^3}$ patterns, irrespective of
their plaquette constraints or gauge equivalences.
Using this deliberately crude local count avoids subtracting
gauge degrees of freedom whose boundary values have been fixed.
A union bound in \eqref{c18v68:conditional} yields, for every
exterior $\xi$,
\[
 \widehat\nu_t^R(\mathcal E_{\eta,R}\mid\xi)
 \le Q_{\eta,R}(t):=
 \min\left\{1,\exp\left[
   (21\log2-\tfrac32\eta t)\ell^3+\mathfrak s_t(R)\right]\right\}.
 \tag{V68.11}\label{c18v68:droplet}
\]
No independence between plaquettes, patterns, or droplets
is used. The same upper bound holds for
$\widehat\nu_t(\mathcal E_{\eta,R})$.

For a concrete full-block-time consequence choose $\eta=6$.
The all-negative interior sign field qualifies for every
$\ell\ge2$, since $N_{{\rm cap},R}\ge6\ell^3$.
Put $\delta=36-21\log2>21$. At $t=4$, if
\[
            \ell\ge\frac{96(e^{96}-1)+8B_\zeta}{\delta},
 \tag{V68.12}\label{c18v68:size}
\]
then $b_R\le48\ell^2$ and \eqref{c18v68:droplet} imply
\[
 \widehat\nu_4^R(\mathcal E_{6,R}\mid\xi)
       \le e^{-\delta\ell^3/2}\qquad\hbox{for every }\xi .
 \tag{V68.13}\label{c18v68:t4}
\]
This is a finite, volume-independent size threshold, but an
astronomically large one. It is not a useful numerical block
scale or a weak-coupling continuum estimate. The radius
$\epsilon_\eta$ is also very small. The improvement over V67
is local support and a completely arbitrary exterior with
a paid surface penalty, not better constants or larger
neighborhoods.

\subsection{A local conditional temporal-source load}

The preceding events are independent of $t$. Reuse V63's
normalized conditional-score identity, now in this block:
\[
 s_t^{R,\xi}(y)
  :=\partial_t\log\frac{d\widehat\nu_t^R}{d\mu_R}(y\mid\xi)
   =a_t(y,\xi)-
        \mathbb E_{\widehat\nu_t^R(\cdot\mid\xi)}a_t(\cdot,\xi).
 \tag{V68.14}\label{c18v68:score}
\]
V60--V61 give $\|\nabla_i a_t\|_\infty\le DK_*$.
The interior oscillation and conditional centering therefore imply
\[
 \|s_t^{R,\xi}\|_\infty\le m_RD^2K_*,
 \qquad
 \|\nabla_R s_t^{R,\xi}\|_\infty^2\le m_RD^2K_*^2,
 \qquad m_R=21\ell^3 .
\]
Combining with \eqref{c18v68:droplet} proves
\[
 \begin{aligned}
 \int_{\mathcal E_{\eta,R}}|s_t^{R,\xi}|^2\,d\widehat\nu_t^R
       &\le m_R^2D^4K_*^2Q_{\eta,R}(t),\\
 \int_{\mathcal E_{\eta,R}}|\nabla_Rs_t^{R,\xi}|^2\,d\widehat\nu_t^R
       &\le m_RD^2K_*^2Q_{\eta,R}(t),\\
 \left|\partial_t\widehat\nu_t^R(\mathcal E_{\eta,R}\mid\xi)\right|
       &\le m_RD^2K_*Q_{\eta,R}(t).
 \end{aligned}
 \tag{V68.15}\label{c18v68:localload}
\]
All integrals use the exterior slice and its actual conditional
law. The last identity follows by differentiating that normalized
density on a fixed measurable event. The polynomial prefactors
depend on the block volume, not on the surrounding lattice.
This is a conditional temporal score; replacing the global
Poisson source $a_t$ by $s_t^{R,\xi}$ would change the equation
and is not licensed by these bounds.

\subsection{Verification and the remaining response problem}

The diagnostic
\begin{center}\small
\path{scripts/verify_ym_c18_v68_local_flux_boundary.py}
\end{center}
enumerates six native translated cuts on $n=5,6,8$ tori,
including wrapping cubes. It checks that normal columns do
not straddle the cut, the exact boundary and captured counts,
the signed interior normal trace, exterior cancellation, and
eligibility of the all-negative pattern, and the weighted
boundary-layer sum with exact rational weights. An exact symbolic
two-dimensional matrix calculation checks the Duhamel trace
formula and a rank-one plaquette calculation checks the
three-color trace norm. A noncommuting four-dimensional
matrix fixture tests the determinant modulus through $t=4$
at 90-digit precision. The latter is explicitly a matrix
regression, not a simulated native nonlinear density.
The uniform native inequalities are the written proofs,
not claims extrapolated from those fixtures.

The new estimate combines a paid native central bulk term,
a determinant boundary cost and a surface-summed spatial tail.
It removes any near-central restriction on the exterior,
but does not supply suppression of all possible contours
or all adverse-curvature configurations: the specified
interior neighborhoods remain part of the event. It also
does not control the fraction of energy of $L_t^{-1}a_t$
there; V67's localized invariant-bump warning still applies.
The next upstream issue is the actual source-to-region
response or the remaining noncentral interior configurations,
rather than another flat-point Taylor coefficient.

The all-coarse and interacting-vacuum extensions, transport
strain through time four, global mixing/coercivity, coarse
return, scale transport, physical clock, interacting continuum
construction and positive physical Yang--Mills mass gap remain
unproved. The goal remains active.

Removing this append and restoring V67's title recovers V67
byte for byte. Other sources are preserved; only the active
YM source filename is replaced. Only \texttt{.tex} files are
kept in \path{latex_notes}. All build and diagnostic products
stay outside. Companion and broad primary-literature reviews
remain due at V70, with archive/restart at V100.
% END CYCLE 18 VERSION 68 APPEND
% BEGIN CYCLE 18 VERSION 69 APPEND
\clearpage
\section{V69: Rare-source inverse energy and the conditional-mean remainder}
\label{c18v69:context}

V68 proves a local rare-droplet probability bound for every
exterior, but that is not an estimate on the inverse equation.
This append turns it into an \emph{unshifted} inverse-energy
estimate for an explicitly represented part of the native
temporal source. It retains the actual moving metric and
allows arbitrary square-integrable coefficients depending
on the whole exterior of the block. No global mixing
constant is assumed. The exact source representation is
essential: the full temporal source, its Hessian response,
and its discarded conditional means are not controlled here.

The local conditional solver is inherited machinery from
V31, not a new general method. V62 already gives the needed
one-factor floors in the present normal chart. The new
calculation combines them with V68's exterior-uniform
probability bound, without a block-volume exponential
comparison or a fictitious global gap. V20's positive-shift
physical resolvent and V31's finite-time lift correction
are different statements; neither supplies the estimate
below for this particular rare source.

The user's literature suggestions remain useful in their
source-specific inverse and conditional-mean directions.
Their proposed V58 cubic contraction has already been
completed in V59 and is not repeated. We rechecked
Caputo--Menz--Tetali,
\href{https://www.numdam.org/item/AFST_2015_6_24_4_691_0.pdf}
{\emph{Approximate tensorization of entropy at high temperature}},
Section 1.1 and Theorem 2.1: conditional expectations define
the rate-one heat-bath form, while global tensorization
requires an additional weak-dependence hypothesis.
Only that standard form is used here, with a direct proof;
its global inequality is not imported into YM.
This targeted check does not replace the broad V70 review.

\subsection{The actual law, local projections and paid constants}

Fix \(\beta=0\), coarse configuration \(g=I\), \(N=2n\),
\(n\ge5\) and \(0\le t\le4\). On the full normal-chart fibre
write
\[
 \nu=\widehat\nu_t=e^{u_t}\mu_I,\qquad a=a_t=\partial_tu_t,
 \qquad \nu(a)=0,\qquad
 \widehat L=-\operatorname{div}_{\nu}\nabla_{\widehat m_t}.
\]
Divergence here is defined by the probability volume \(\nu\);
it is not replaced by Riemannian volume.
The product reference metric is \(m_0\), with factor metrics
\(\kappa_i^{-1}m_{\rm round}\), \(\kappa_i\in\{1/2,1\}\).
Keep the full-interval constant \(E_0\) of V61, so that
\[
 E_0^{-2}m_0\preceq\widehat m_t\preceq E_0^2m_0.
 \tag{V69.1}\label{c18v69:metric}
\]
This comparison holds through time four independently
of V61's much shorter probability-transport interval.

Let \(P_i\) be expectation in factor \(i\) under its actual
conditional law, with all other factors fixed, and put
\(D_i=I-P_i\). These are orthogonal projections in \(L^2(\nu)\),
not derivatives. Write \(\nabla_i\) for the product-metric
gradient in that factor. With \(D=\pi\sqrt2\), define
\[
 S=D^2K_*,\qquad
 \lambda_i(t)=3\kappa_i e^{-D^2tK_*},\qquad
 \lambda_*(t)=\tfrac32 e^{-D^2tK_*}.
 \tag{V69.2}\label{c18v69:constants}
\]
The inherited bounds imply
\[
 |D_i a|\le S,\qquad
 \lambda_i\|D_i f\|_{L^2(\nu)}^2
        \le\int|\nabla_i f|^2\,d\nu .
 \tag{V69.3}\label{c18v69:local}
\]
Indeed \(|\nabla_i a|\le DK_*\) bounds its one-factor
oscillation by \(S\). Similarly the oscillation of \(u_t\)
in that factor is at most \(D^2tK_*\). The conditional
density ratio and the Haar gap \(3\kappa_i\) give the
second inequality. Conditional normalizers cancel.
These assertions hold for every exterior; no smallness
of \(tK_*\) is required.

Use the block \(\mathcal I_R\), event
\(E=\mathcal E_{\eta,R}\) and number \(Q=Q_{\eta,R}(t)\)
from V68, with \(2\le\ell<n\) and \(0<\eta\le12\):
\[
 m_R=|\mathcal I_R|=21\ell^3,\qquad
       \nu(E\mid x_{\mathcal I_R^c}=\xi)\le Q
              \quad\hbox{for every }\xi .
 \tag{V69.4}\label{c18v69:rare}
\]
The event depends only on the interior factors and is
pinned-gauge invariant. Its near-central interior
restriction is retained; its entire exterior is arbitrary.

\subsection{A represented-source inverse estimate}

At each fixed finite regulator the compact connected
manifold and positive smooth density give an inverse
of \(\widehat L\) on centered \(L^2\) sources. We use its
energy norm
\[
 \|h\|_{-1,\widehat L}^2
    :=\langle h,\widehat L^{-1}h\rangle_\nu
     =\sup_f\left\{2\langle h,f\rangle_\nu
                 -\int|df|_{\widehat m_t^{-1}}^2d\nu\right\}.
 \tag{V69.5}\label{c18v69:dual}
\]
The supremum is over real form-domain functions.
Complete the square with the centered solution of
\(\widehat L\psi=h\) to prove the equality.
The qualitative finite-volume inverse is not assigned
a uniform spectral norm.

\begin{lemma}[Local centering pays an unshifted inverse]
\label{c18v69:represented}
For any finite factor set \(J\) and \(F_i\in L^2(\nu)\),
put \(h=\sum_{i\in J}D_iF_i\). Then \(\nu(h)=0\) and
\[
 \|h\|_{-1,\widehat L}^2
       \le E_0^2\sum_{i\in J}\lambda_i^{-1}
                                      \|D_iF_i\|_2^2
       \le E_0^2\sum_{i\in J}\lambda_i^{-1}\|F_i\|_2^2 .
 \tag{V69.6}\label{c18v69:inverse}
\]
No commutation between distinct \(P_i\), independence
of factors under \(\nu\), or global Poincare inequality
with a uniform constant is required.
\end{lemma}
\begin{proof}
Orthogonality and \eqref{c18v69:local} give
\[
 \begin{split}
 |\langle h,f\rangle_\nu|
 &=\left|\sum_i\langle D_iF_i,D_if\rangle_\nu\right|\\
 &\le\left(\sum_i\lambda_i^{-1}\|D_iF_i\|_2^2\right)^{1/2}
          \left(\sum_i\lambda_i\|D_if\|_2^2\right)^{1/2}\\
 &\le\left(\sum_i\lambda_i^{-1}\|D_iF_i\|_2^2\right)^{1/2}
          \left(\int|df|_{m_0^{-1}}^2d\nu\right)^{1/2}.
 \end{split}
\]
By \eqref{c18v69:metric}, the last energy is at most
\(E_0^2\) times the actual energy in \eqref{c18v69:dual}.
Maximize \(2E_0\sqrt{AB}-B\) over \(B\ge0\) to obtain
\(E_0^2A\). Each \(D_i\) is a contraction.
There is no extra \(|J|\) from squaring a sum of sources.
\end{proof}

This estimate also has an explicit divergence representation.
Let \(L_i=-\Delta_i-\langle\nabla_i u_t,\nabla_i\cdot\rangle\)
be the conditional one-factor generator and solve
\(L_i\phi_i=D_iF_i\) with \(P_i\phi_i=0\).
The proved local floor gives
\[
 \begin{gathered}
 Y_i=\nabla_i\phi_i,\qquad
 -\operatorname{div}_\nu\!\left(\sum_iY_i\right)=h,\\
 \int\left|\sum_iY_i\right|_{\widehat m_t}^2d\nu
          \le E_0^2\sum_i\lambda_i^{-1}\|D_iF_i\|_2^2 .
 \end{gathered}
 \tag{V69.7}\label{c18v69:current}
\]
Divergence is in the weak sense: integrate by parts
only in coordinate \(i\), then integrate its exterior.
The \(Y_i\) are orthogonal in \(m_0\), though not necessarily
in \(\widehat m_t\); \eqref{c18v69:metric} pays that distinction.
For measurable \(F_i\), approximate by smooth functions,
apply the uniform conditional floor, and pass in the
gradient Hilbert space to the conditional variational
solutions. No derivative of an event indicator is taken.

The gradient \(\nabla_{\widehat m_t}\widehat L^{-1}h\)
has the least actual \(L^2\) norm among currents with
this divergence. The difference from any such current
has zero weighted divergence and is orthogonal to
that gradient by integration by parts. Thus this is
a bound for the actual inverse, not a product-metric
inverse substituted into its equation.

\subsection{Droplet fluctuations, exterior banks and actual inverse energy}

Define the one-factor temporal scores \(a_i=D_i a\) and
the represented rare source
\[
 h_{R,E}=\sum_{i\in\mathcal I_R}D_i({\bf1}_E a_i).
 \tag{V69.8}\label{c18v69:source}
\]
The second centering is deliberate:
\({\bf1}_E a_i\) is not conditionally centered in general.

\begin{theorem}[Native rare-source response through time four]
\label{c18v69:rareinverse}
Let \(\bar\nu_R\) be the exterior marginal of the actual
law and \(c\in L^2(\bar\nu_R)\) any real function of
the entire exterior \(\xi=x_{\mathcal I_R^c}\).
Then \(c\,h_{R,E}\) is centered and
\[
 \begin{split}
 \|c\,h_{R,E}\|_{-1,\widehat L}^2
 &\le E_0^2 S^2Q
            \left(\sum_{i\in\mathcal I_R}\lambda_i^{-1}\right)
                         \|c\|_{L^2(\bar\nu_R)}^2\\
 &\le E_0^2\lambda_*^{-1}m_RS^2Q
                         \|c\|_{L^2(\bar\nu_R)}^2 .
 \end{split}
 \tag{V69.9}\label{c18v69:rareenergy}
\]
Thus \(c\mapsto c\,h_{R,E}\) is bounded from exterior
\(L^2\) into the actual negative energy space, with
the displayed squared operator-norm bound.
No exterior derivative of \(c\) is required.
\end{theorem}
\begin{proof}
For \(i\in\mathcal I_R\), \(c\) is independent of coordinate
\(i\), so \(D_i(c{\bf1}_E a_i)=cD_i({\bf1}_E a_i)\).
Apply Lemma \ref{c18v69:represented} with
\(F_i=c{\bf1}_E a_i\). Disintegration in the \emph{whole}
exterior, \eqref{c18v69:rare}, and \(|a_i|\le S\) give
\[
 \|F_i\|_2^2
  \le S^2\int c(\xi)^2\,\nu(E\mid\xi)\,d\bar\nu_R(\xi)
  \le S^2Q\|c\|_2^2 .
\]
This proves integrability for unbounded \(c\).
Exterior \(L^2\) truncation passes the bound to this
class; conditional centering gives zero integral.
For the current construction, multiplication by \(c\)
creates no divergence term, because its \(i\) derivative
vanishes in the distributional sense. Other derivatives
may be uncontrolled. This is not a strain or Hessian
estimate for variable banks.
\end{proof}

If \(c\) is pinned-gauge invariant, so are the sources
and their global Poisson solutions. The gauge action
maps each factor by a Haar-preserving isometry without
mixing indices. Invariance of \(\nu\) makes each \(P_i\)
commute with that action. The event and \(a\) are
invariant, and the actual metric is equivariant.
Uniqueness of the centered solution proves the claim.
No horizontal projection is asserted to commute with
a one-form inverse.

For \(\eta=6\), \(t=4\), and the explicit size condition
\eqref{c18v68:size}, let \(\delta=36-21\log2>21\). Then
\[
 \|c\,h_{R,E}\|_{-1,\widehat L}^2
 \le 14E_0^2D^4K_*^2 e^{4D^2K_*}
           \ell^3e^{-\delta\ell^3/2}\|c\|_2^2 .
 \tag{V69.10}\label{c18v69:t4}
\]
Here \(14=21(2/3)\); heat-bath time is not rescaled.
Constants are independent of the surrounding lattice
and exterior, but extremely large. The bound holds
for every admissible \(n>\ell\), not at a fixed torus
while \(\ell\) tends to infinity. No useful physical
scale is thereby identified.

The left side is the \emph{total} energy of the
response to this represented source, throughout
the fibre. It is not the energy of
\(\widehat L^{-1}a\) restricted to \(E\).
Neither the source nor the current need be supported
on \(E\): conditional averaging spreads a function
along a coordinate.

\subsection{The conditional mean is an explicit unsolved source}

Here is precisely what the second centering removes.
For each interior factor put
\[
 \begin{gathered}
 p_i=P_i{\bf1}_E,\qquad q_i=P_i({\bf1}_E a_i),\\
 z_i={\bf1}_E a_i-\nu({\bf1}_E a_i),\qquad
 r_i=q_i-\nu({\bf1}_E a_i).
 \end{gathered}
\]
All projections use the law at the displayed time.
Since the event is time independent, differentiating
its normalized conditional density gives
\[
 \partial_t p_i=q_i,\qquad
 z_i=D_i({\bf1}_E a_i)+r_i,\qquad
 \nu(r_i)=0,\qquad \|r_i\|_2^2\le S^2Q .
 \tag{V69.11}\label{c18v69:remainder}
\]
The conditional score is \(a-P_i a=a_i\).
The last inequality follows from
\[
 \operatorname{Var}_\nu(q_i)\le\|q_i\|_2^2
 \le\|{\bf1}_Ea_i\|_2^2\le S^2\nu(E).
\]
Thus the remainder is the centered time derivative
of an actual conditional event probability, not a
normalization constant that may be discarded.
It need not vanish.

Since \(r_i\) is independent of coordinate \(i\),
\(D_ir_i=0\). The \(i\)-direction estimate therefore
gives no inverse bound on it. Its small \(L^2\) norm
does not by itself bound its negative energy norm.
A uniform estimate on the latter, or a further exact
divergence representation, is still required for \(z_i\).

There is a second, separate source distinction. Put
\[
 T_R=\sum_{i\in\mathcal I_R}D_i,\qquad
 T_{R,E}=\sum_{i\in\mathcal I_R}D_i{\bf1}_E D_i .
\]
These are bounded self-adjoint positive operators;
the indicator in the second formula is multiplication.
Directly,
\[
 \begin{gathered}
 h_{R,E}=T_{R,E}a,\qquad
 T_{R,E}+T_{R,E^c}=T_R,\\
 \langle f,T_{R,E}f\rangle_\nu
       =\sum_{i\in\mathcal I_R}\int_E|D_if|^2d\nu .
 \end{gathered}
 \tag{V69.12}\label{c18v69:partition}
\]
The rare and complementary represented sources sum
to \(T_Ra\), \emph{not} to \(a\).
A full-coordinate heat-bath inverse, or a martingale
decomposition of \(a\) with controlled coefficients,
is an additional obligation. The \(P_i\) generally
do not commute; they cannot be cancelled or reordered
as independent Haar averages.

The paid result also gives the form inequality
\[
 T_{R,E}\widehat L^{-1}T_{R,E}
               \preceq E_0^2\lambda_*^{-1}T_{R,E}.
 \tag{V69.13}\label{c18v69:sandwich}
\]
Apply \eqref{c18v69:inverse} with
\(F_i={\bf1}_ED_if\) and
\(\|D_iF_i\|_2^2\le\int_E|D_if|^2d\nu\) to prove it.
The range is centered, so the inverse is well defined.
This is a represented-source sandwich, not a
spectral floor for the full fibre.

The next useful calculation is now sharper: bound
the native remainders \(r_i\), or assemble a controlled
representation of the actual \(a_t\). Repeating the
local conditional floor or the probability of \(E\)
does not address either input. The complement still
includes uncontrolled noncentral configurations.

\subsection{Verification, limits and version handoff}

The diagnostic
\begin{center}\small
\path{scripts/verify_ym_c18_v69_rare_source_inverse.py}
\end{center}
checks conditional projections for a strictly positive,
correlated three-bit rational law. It verifies their
noncommutation, the centered-source and probability
derivative identities, the positive sandwich, an
exterior-dependent bank, and an anisotropic generator
comparison using exact rational arithmetic.
This is an algebraic regression, \emph{not} a lattice
YM law. The native constants and the all-exterior
probability input come from the written proofs and
V68; its independent native-incidence regression is
also replayed. No proof assistant or numerical
mass-gap claim is used.

The advance is an actual unshifted inverse-energy
bound for a specified rare-source family, including
exterior \(L^2\) banks and the actual metric. It does
not close the full temporal-source inverse, pointwise
Poisson Hessian or transport strain. All-coarse and
interacting-vacuum extensions, global mixing,
coarse return, physical scale transport and a
nontrivial interacting continuum with a positive
physical Yang--Mills mass gap remain unproved.

V68 is recoverable byte for byte after removing this
append and restoring its title. Only the active
source filename is replaced; other sources are
preserved. The PDF skill is used solely for internal
typesetting and visual checks, outside
\path{latex_notes}, which contains only \texttt{.tex}
files. Both the companion checkpoint and the broad
primary-literature review are due at V70.
Archive/restart remains V100. The goal remains active.
% END CYCLE 18 VERSION 69 APPEND
% BEGIN CYCLE 18 VERSION 70 APPEND
\clearpage
\section{V70: Rare conditional-mean messages, collars and distant inverse energy}
\label{c18v70:context}

V69 controlled a locally recentered rare source, but left the
conditional mean \(r_i\) in \eqref{c18v69:remainder}. This append
examines that actual remainder rather than replacing it by
a different law. Averaging it over its block gains a second
rare-probability factor in squared norm. Its derivatives in
distant exterior directions admit a surface-summed bound.
These estimates pay the actual unshifted inverse energy of
the corresponding distant-resampling source. They do not
pay the inverse of the entire remainder: the near-boundary
message, and the assembly of the full temporal source,
remain separate obligations.

V63 already differentiated normalized conditional densities
and obtained collar locality; V68 already proved the rare
event bound for every entire exterior. Neither is repeated
as a new general theorem. The new point is to retain that
rare factor in the differentiated V69 message and then use
the paid V69 source representation. All statements below
retain its scope: \(\beta=0\), coarse identity, \(N=2n\),
\(n\ge5\), and \(0\le t\le4\). In particular this is not
an interacting four-dimensional continuum construction.

\subsection{Scheduled companion and primary-literature checkpoint}

The V70 checkpoint reread the terminal scope and obstruction
passages of NS V26, SM--Einstein V72, source selection V65,
stationarity V61, Einstein bridge V55, shell mixing V16,
parallel YM V5, ESM general V44, ESM phenomenology V40,
arithmetic loading V41 and renewal arithmetic V17.
Their hashes are unchanged since the V65 checkpoint.
The ESM V80\_f1 terminal scope was also checked. This is
a targeted review, not recertification of all cumulative
proofs or of numerical certification grids. No source was
modified and no companion gap constant is imported.

The pertinent lessons are concrete. Source selection and
arithmetic distinguish a represented operator from its
occurrence in the actual equation; stationarity keeps the
uncomputed residual. Shell mixing retains the mixed-curvature
calculation beyond the measured one-loop assembly.
NS's pointwise contraction formulas do not certify an
uncomputed integrated bound. The ESM and SM results keep
their fixed-order, fixed-mode or auxiliary many-copy limits;
the Einstein slow gate does not supply a physical clock.
Accordingly, the conditional mean below is retained as an
actual source, and an inverse applied to a resampling
operator is not cancelled to claim an inverse on that mean.

The scheduled broad sweep screened primary literature in
spatial mixing, entropy factorization, spectral decomposition,
metastable capacities, nonconvex gradient fields and lattice
gauge dynamics. The following are possible mechanisms, not
imported YM theorems.
\begin{enumerate}
\item
\href{https://arxiv.org/abs/2004.10574}
{Caputo--Parisi, \emph{Block factorization of the relative entropy
via spatial mixing}}, assumes a spatial-mixing input.
The classical Lu--Yau
\href{https://deepblue.lib.umich.edu/bitstream/handle/2027.42/46486/220_2005_Article_BF02098489.pdf?sequence=1}
{\emph{Spectral gap and logarithmic Sobolev inequality for Kawasaki
and Glauber dynamics}} also requires mixing hypotheses.
Their primary abstracts were screened; no full proof is
claimed to have been rechecked and no mixing assumption is
silently supplied.
\item
\href{https://arxiv.org/html/2503.19419v1}
{Caputo--Chen--Parisi, \emph{Factorizations of Relative Entropy
Using Stochastic Localization}}, was checked at its
Theorem 2.2 and component extensions. Its interaction
representation and high-temperature spectral margin are
additional inputs. Finiteness of our \(K_\zeta\), which
can be huge, is not that smallness condition.
\item
\href{https://arxiv.org/html/2504.01247v2}
{Qin, \emph{On spectral gap decomposition for Markov chains}}
(revision 10 September 2026), was checked at the abstract
and introduction. Its local-gap decomposition retains a
benchmark-chain gap and a Markov sandwich representation.
V69's source-form sandwich alone is not that representation
or a proof of the benchmark gap.
\item
The primary abstracts of
\href{https://arxiv.org/abs/2608.14526}
{Buchholz--Cotar--Schweiger, \emph{Gradient Gibbs measures with
non-convex potentials and the universality class of the
Gaussian free field}}, and
\href{https://ems.press/content/serial-article-files/31531?nt=1}
{Bovier--Gayrard--Klein, \emph{Metastability in reversible
diffusion processes II}}, were screened. The former concerns
real gradient interfaces with special potential hypotheses;
the latter uses a small-noise metastable-well setting.
Neither structure has been verified for the native
compact normal-chart law uniformly in scale or volume.
\item
\href{https://arxiv.org/abs/2204.12737}
{Shen--Zhu--Zhu, \emph{A stochastic analysis approach to
lattice Yang--Mills at strong coupling}}, and
\href{https://arxiv.org/abs/2608.27828}
{Chevyrev--Qu--Shen, \emph{Scaling limit of the 3D abelian
Yang--Mills Langevin dynamics}}, were screened at their
primary abstracts. Strong-coupling lattice Gibbs results
and a three-dimensional abelian dynamical scaling limit
do not supply the four-dimensional nonabelian physical
mass gap or the presently missing native-source estimate.
\end{enumerate}
The direction is therefore to quantify the actual message,
not assume the coarse mixing which these possible global
mechanisms would still need. Detailed review scopes and
source links are recorded outside the notes folder.

\subsection{Geometry and the inherited analytic inputs}

Keep \(\nu,u,a,\widehat L,P_i,D_i,D,S,\lambda_i,\lambda_*,E_0\)
from V69. Thus \(D_i=I-P_i\) is a conditional projection
difference, while \(\nabla_i\) is a product-metric derivative.
Write \(P_R\) for expectation over all factors
\(\mathcal I_R\), conditional on the entire exterior.
For \(i\in\mathcal I_R\), the tower property gives
\(P_RP_i=P_R\); no commutation of different single-factor
projections is used. Let
\[
 a_i=D_i a,\qquad E=\mathcal E_{\eta,R},\qquad
 Q=Q_{\eta,R}(t),\qquad C_\zeta=\frac{72}{e^\zeta-1}.
\]
In particular \(|a_i|\le S\) and \(P_R{\bf1}_E\le Q\)
pointwise in the entire exterior. The event is a fixed
measurable set of interior variables, independent of
time and of every exterior coordinate.

Use V68's fine cube \(D_R\) of side \(2\ell\), and the
fine torus graph distance \(d\) of the factor anchors.
For an interior factor let
\(\rho_i=d(\operatorname{anchor}(i),D_R^c)\ge1\).
For an integer \(r\ge0\), partition the exterior as
\[
 \begin{split}
 F_r&=\{j\notin\mathcal I_R:
               d(\operatorname{anchor}(j),D_R)\ge r+1\},\\
 C_r&=\mathcal I_R^c\setminus F_r .
 \end{split}
 \tag{V70.1}\label{c18v70:collar}
\]
Empty far sets give zero sums below. In particular
\(F_0=\mathcal I_R^c\). For \(i\in\mathcal I_R,j\in F_r\),
\[
 d(i,j)\ge \rho_i+r,\qquad
 \sum_{i\in\mathcal I_R}e^{-\zeta\rho_i}
                         \le C_\zeta\ell^2 .
 \tag{V70.2}\label{c18v70:depth}
\]
For the distance inequality take a shortest path from
the interior anchor to the exterior one. Before its
first exit it uses at least \(\rho_i\) edges. The last
interior vertex before that exit lies in \(D_R\);
the remaining path after the first outside vertex
therefore uses at least \(r\) edges. This also works
for translated or wrapping cubes on the torus.
The second inequality is V68's six-face layer estimate,
with at most three factor anchors per fine vertex.

The weighted Hessian bounds inherited from V63 are
\[
 \mathcal H_\zeta(a)\le K_\zeta,\qquad
 \mathcal H_\zeta(u)\le tK_\zeta,\qquad
 \mathcal H_\zeta(f)=
   \sup_x\max_i\sum_j e^{\zeta d(i,j)}
                      \|(\operatorname{Hess}f)_{ij}(x)\| .
 \tag{V70.3}\label{c18v70:jets}
\]
The supremum is outside each row sum. We keep this order
throughout; summing separate entrywise suprema would not
be justified by this input.

\subsection{Differentiating the actual one-factor remainder}

For \(j\ne i\), put \(b_j=\nabla_j u\) and
\(B_{ij}=b_j-P_i b_j\). Vectors in factor \(j\) are
canonically identified while only factor \(i\) varies.
Differentiating the numerator and denominator of the
actual conditional expectation gives, for smooth \(f\),
\[
 \nabla_j P_i f
 =P_i\nabla_j f+P_i\big[(f-P_i f)B_{ij}\big].
 \tag{V70.4}\label{c18v70:diff}
\]
There is no omitted conditional normalizer. In particular
\[
 \nabla_j a_i=D_i\nabla_j a-P_i(a_iB_{ij}).
 \tag{V70.5}\label{c18v70:ai}
\]
The same differentiation is valid for functions only
measurable in interior variables and smooth in the
exterior coordinates being differentiated, with bounded
derivatives. All applications below have that property
on a fixed finite compact fibre.

\begin{lemma}[Rare-factor bound for the distant raw message]
\label{c18v70:rawlemma}
For \(i\in\mathcal I_R\), retain V69's
\(p_i=P_i{\bf1}_E\), \(q_i=P_i({\bf1}_E a_i)\) and
\(r_i=q_i-\nu({\bf1}_E a_i)\). Then
\[
 \begin{split}
 \sum_{j\in F_r}|\nabla_j a_i|
       &\le D K_\zeta(1+tS)e^{-\zeta(\rho_i+r)},\\
 \sum_{j\in F_r}|\nabla_j r_i|
       &\le p_iD K_\zeta(1+2tS)e^{-\zeta(\rho_i+r)}.
 \end{split}
 \tag{V70.6}\label{c18v70:raw}
\]
Consequently,
\[
 \int\left(\sum_{j\in F_r}|\nabla_j r_i|\right)^2d\nu
 \le D^2K_\zeta^2(1+2tS)^2
                     e^{-2\zeta(\rho_i+r)}Q .
 \tag{V70.7}\label{c18v70:rawL2}
\]
\end{lemma}
\begin{proof}
Connect two values of factor \(i\) by a minimizing
geodesic of length at most \(D\), leaving all other
coordinates fixed. Integrating its mixed Hessian,
summing over \(F_r\) at the same point of the path,
and using \eqref{c18v70:depth}--\eqref{c18v70:jets}
bounds the summed change of \(\nabla_j a\) by
\(DK_\zeta e^{-\zeta(\rho_i+r)}\), and that of \(b_j\)
by \(DtK_\zeta e^{-\zeta(\rho_i+r)}\).
Averaging the second endpoint under the actual
\(i\)-conditional law gives these same bounds on
\(\sum_{F_r}|D_i\nabla_j a|\) and
\(\sum_{F_r}|B_{ij}|\). Equation \eqref{c18v70:ai}
and \(|a_i|\le S\) prove the first assertion.

Since \({\bf1}_E\) is independent of \(x_j\) for
\(j\in F_r\), differentiation using \eqref{c18v70:diff}
and \eqref{c18v70:ai} gives the exact three-term identity
\[
 \nabla_j q_i
   =P_i({\bf1}_E D_i\nabla_j a)
       -p_iP_i(a_iB_{ij})
       +P_i({\bf1}_E a_iB_{ij}).
 \tag{V70.8}\label{c18v70:three}
\]
Each expectation containing the indicator is bounded
by \(p_i\) times the uniform summed bound on its other
factor. The other covariance term already has the
factor \(p_i\). This gives \(1+2tS\), not \(1+tS\).
The subtracted \(\nu({\bf1}_Ea_i)\) is a spatial
constant, so the same derivative bound holds for \(r_i\).
Finally \(0\le p_i\le1\) implies
\(\int p_i^2d\nu\le\int p_i d\nu=\nu(E)\le Q\).
No derivative of the sharp indicator was taken.
\end{proof}

This lemma concerns exterior partial derivatives only.
It does not assert interior smoothness of \(r_i\), nor
an inverse-energy estimate on \(r_i\) itself.

\subsection{The block-averaged mean: a rare surface message}

Define an actual exterior function
\[
 M_i=P_Rq_i=P_R({\bf1}_E a_i),\qquad
             P_Rr_i=M_i-\nu(M_i).
 \tag{V70.9}\label{c18v70:mean}
\]
The equality is the tower property. Introduce
\[
 A_i(r)=D K_\zeta e^{-\zeta r}
       \left[(1+tS)e^{-\zeta\rho_i}
                         +tS C_\zeta\ell^2\right].
 \tag{V70.10}\label{c18v70:A}
\]

\begin{theorem}[Uniformly rare, surface-summed exterior derivative]
\label{c18v70:meanthm}
For every entire exterior, \(M_i\) is smooth and
\[
 |M_i|\le SQ,\qquad
 \|P_Rr_i\|_2^2\le S^2Q^2,\qquad
 \sum_{j\in F_r}|\nabla_j M_i|\le Q A_i(r).
 \tag{V70.11}\label{c18v70:meanbounds}
\]
The constants are independent of the volume of the
surrounding lattice. The block polynomial in \(A_i\)
is a surface term, not an ambient-volume term.
\end{theorem}
\begin{proof}
The first bound follows from \(P_R{\bf1}_E\le Q\).
Centering decreases squared \(L^2\) norm, proving
the second bound. Smoothness follows by differentiation
under the integral over the fixed measurable set \(E\):
on each finite compact fibre the density and \(a_i\)
and their exterior derivatives are bounded, and the
conditional normalizer is positive.

For \(j\) exterior, the exact normalized derivative is
\[
 \nabla_j M_i
 =P_R({\bf1}_E\nabla_j a_i)
   +P_R\big[{\bf1}_Ea_i(b_j-P_Rb_j)\big].
 \tag{V70.12}\label{c18v70:meanidentity}
\]
To bound the second term, compare \(b_j\) at two
interior configurations with the same exterior.
Change the interior factors one by one along
geodesics of length at most \(D\). The weighted
row bound, at each point of each path, gives
\[
 \sum_{j\in F_r}|b_j-P_Rb_j|
 \le DtK_\zeta e^{-\zeta r}
               \sum_{k\in\mathcal I_R}e^{-\zeta\rho_k}
 \le DtK_\zeta C_\zeta\ell^2e^{-\zeta r}.
 \tag{V70.13}\label{c18v70:boundaryscore}
\]
Averaging the comparison configuration under \(P_R\)
does not change the uniform bound. The first term of
\eqref{c18v70:meanidentity} is bounded by the first line
of \eqref{c18v70:raw}, multiplied by \(Q\).
The second is bounded by \eqref{c18v70:boundaryscore}
times \(SQ\). Their sum is \(Q A_i(r)\).
\end{proof}

The factor \(Q^2\) in squared norm is stronger than
V69's \(Q\) for the raw \(r_i\), but averaging has
removed its within-block fluctuation. Explicitly,
\[
 r_i=(r_i-P_Rr_i)+(M_i-\nu M_i).
 \tag{V70.14}\label{c18v70:twopieces}
\]
Neither term is declared negligible in inverse norm.

There is also a collar approximation in the actual law.
Let \(\bar\nu_R\) be its exterior marginal and define
\[
 M_i^{[r]}=
   \mathbb E_{\bar\nu_R}[M_i\mid x_{C_r}].
\]
This is a conditional expectation, not a frozen-exterior
replacement density. Two exteriors agreeing on \(C_r\)
can be joined by simultaneous geodesics in their far
factors, each of length at most \(D\).
Equation \eqref{c18v70:meanbounds} therefore gives
\[
 \|M_i-M_i^{[r]}\|_\infty\le DQ A_i(r),
 \qquad \nu(M_i-M_i^{[r]})=0 .
 \tag{V70.15}\label{c18v70:approx}
\]
The norm may be read as an essential supremum; the
positive smooth marginal supplies canonical conditional
versions. This also bounds the \(L^2\) error.
It is not an inverse-energy bound for that error:
locality of a function alone does not control a small
eigenvalue on which it might have nonzero projection.

\subsection{An actual inverse estimate for the distant-resampling part}

Now use the original full-law projections \(D_j\), not
single-factor projections in the exterior marginal.
Define the centered source
\[
 h_{i,r}^{\rm far}=\sum_{j\in F_r}D_jM_i .
 \tag{V70.16}\label{c18v70:farsource}
\]
Although \(M_i\) depends only on the exterior, \(P_jM_i\)
for an exterior \(j\) can depend on interior coordinates
through its full conditional law. This dependence is
retained in the source and in its inverse.

\begin{theorem}[Distant message resampling has a paid unshifted inverse]
\label{c18v70:inverse}
For the actual metric and actual full law,
\[
 \begin{split}
 \|h_{i,r}^{\rm far}\|_{-1,\widehat L}^2
 &\le E_0^2\sum_{j\in F_r}
                      \lambda_j^{-1}\|D_jM_i\|_2^2\\
 &\le E_0^2\lambda_*^{-2}
                  \int\sum_{j\in F_r}|\nabla_jM_i|^2d\nu\\
 &\le E_0^2\lambda_*^{-2}Q^2 A_i(r)^2 .
 \end{split}
 \tag{V70.17}\label{c18v70:farinverse}
\]
No factor \(|F_r|\) or uniform global spectral gap is used.
\end{theorem}
\begin{proof}
The first line is Lemma \ref{c18v69:represented}, with
\(F_j=M_i\). Apply the conditional Poincare bound
\eqref{c18v69:local} once more to each \(M_i\);
this accounts for the second power of \(\lambda_*^{-1}\).
Finally
\(\sum_j|\nabla_j M_i|^2
 \le(\sum_j|\nabla_j M_i|)^2\),
and apply \eqref{c18v70:meanbounds}. This is a direct
bound on the actual inverse, not an exterior-marginal
inverse or a positive-shift resolvent.
\end{proof}

For \(\eta=6,t=4\) and the V68 size threshold
\eqref{c18v68:size}, let \(\delta=36-21\log2\).
The explicit consequence is
\[
 \begin{split}
 \|h_{i,r}^{\rm far}\|_{-1,\widehat L}^2
 \le {}&E_0^2\lambda_*(4)^{-2}D^2K_\zeta^2
                   e^{-\delta\ell^3-2\zeta r}\\
 &{}\times
       \left[(1+4S)e^{-\zeta\rho_i}
                         +4S C_\zeta\ell^2\right]^2 .
 \end{split}
 \tag{V70.18}\label{c18v70:t4}
\]
Its constants and threshold remain extremely large.
The result holds for every admissible \(n>\ell\);
it neither identifies a physical block scale nor
asserts uniform weak-coupling control.

The distinction between this estimate and the desired
source bound is exact. If \(T=\sum_jD_j\), then
\[
 T M_i=\sum_{j\in C_r}D_jM_i+h_{i,r}^{\rm far},
 \tag{V70.19}\label{c18v70:split}
\]
since \(D_jM_i=0\) for interior \(j\).
At \(r=0\) the collar is empty, so
\eqref{c18v70:farinverse} bounds \(T M_i\).
It still does not bound \(M_i-\nu M_i\) in negative
energy norm: cancelling \(T\) would require a further
inverse or controlled representation which has not
been paid. For \(r>0\) the near-collar part also
remains. Thus an exponential far tail has not been
promoted to a global mixing theorem.

\subsection{Verification and the next upstream obligation}

The diagnostic
\begin{center}\small
\path{scripts/verify_ym_c18_v70_boundary_message.py}
\end{center}
checks the normalized derivative, its three terms,
the block tower identity, and a nonzero normalization
correction in a strictly positive correlated rational
three-bit family with a smooth external parameter.
It also checks the represented inverse budget for
the resulting block message in an anisotropic
finite-state generator. That law is not native YM;
the finite-state check is not a gradient estimate
on a compact Lie group. Separate native anchor
enumerations check depth separation and layer sums
on translated and wrapping tori. The uniform native
analytic inequalities above are written proofs,
not extrapolations from these diagnostics. V69's
independent exact inverse regression is replayed.

The next source-level task is to control the near-boundary
mean or a literal representation of
\(M_i-\nu M_i\), and the within-block piece in
\eqref{c18v70:twopieces}, with coefficients whose cost
does not reintroduce an uncontrolled global inverse.
Repeating a general decomposition criterion would not
pay this input. Noncentral adverse configurations,
the full temporal-source inverse, Poisson Hessian,
transport strain through time four, all-coarse and
interacting-vacuum extensions, coarse return,
physical-scale transport and a nontrivial interacting
continuum with positive physical Yang--Mills mass gap
remain unproved.

Removing this append and restoring V69's title recovers
V69 byte for byte. Only the active source filename is
replaced; other sources and legacy archives are preserved.
The PDF skill is used for internal layout checks with
all build products outside \path{latex_notes}; that
folder retains only \texttt{.tex} files.
The next companion checkpoint is V75 and the next
broad literature sweep V80; archive/restart is V100.
The goal remains active.
% END CYCLE 18 VERSION 70 APPEND
% BEGIN CYCLE 18 VERSION 71 APPEND
\clearpage
\section{V71: Surface-cost block relaxation and the within-block rare remainder}
\label{c18v71:context}

V70 left two different parts of the actual remainder:
its within-block fluctuation and its exterior mean.
This append pays the inverse energy of the first part.
The step upstream is a conditional block Poincare bound
whose inverse costs an exponential of surface area,
rather than an exponential of block volume. That cost
can be absorbed by V68's rare-volume suppression at an
explicit, extremely large block size. There is no assumed
uniform global gap, no cancellation of the heat-bath
operator \(T\), and no estimate here for the remaining
exterior mean itself.

The method is the classical coordinate-replacement,
or canonical-path, comparison. For primary context we
checked Berger--Kenyon--Mossel--Peres,
\href{https://arxiv.org/abs/math/0308284}
{\emph{Glauber Dynamics on Trees and Hyperbolic Graphs}},
Proposition 1.1 and its Section 2.1 proof.
Their cut-width argument motivates the ordering and
hybrid-configuration comparison, not a new principle.
We give the continuous-product argument directly below:
our density is the actual normal-chart density, not
an Ising or finite-range pair-interaction density.
The needed crossing cost is proved from V63's weighted
Hessian bound, including its long spatial tails.

The scope remains V68--V70's \(\beta=0\), coarse identity,
\(N=2n\), \(n\ge5\), \(0\le t\le4\), with a block
\(2\le\ell<n\). This is an auxiliary finite-block
conditional estimate; it is not a physical spectral
gap or a four-dimensional continuum construction.
The companion checkpoint remains V75 and the next
broad literature checkpoint V80; this targeted method
check supplements the completed V70 sweep.

\subsection{A weighted crossing bound for every ordered cut}

Use the actual law \(\nu=e^u\mu_I\), product metric \(m_0\),
and the constants \(D,S,K_\zeta,E_0,\lambda_*\) of
V69--V70. In particular
\[
 S=D^2K_*,\quad
 \operatorname{osc}_{x_i}u\le tS,\quad
 \mathcal H_\zeta(u)\le tK_\zeta,\quad
 \lambda_*=\tfrac32 e^{-tS}.
 \tag{V71.1}\label{c18v71:inputs}
\]
Fix the entire exterior \(\xi\) and write
\(\pi_\xi=\nu(\,\cdot\mid x_{\mathcal I_R^c}=\xi)\).
Its reference measure is the product of the normalized
factor Haar measures. Every estimate below holds for
each exterior, not only almost every exterior.

Order the \(m=m_R=21\ell^3\) interior factors first
lexicographically by their local anchor coordinates
in the fine cube of side \(L=2\ell\), then by direction
at a common anchor. Local coordinates are used even
when the translated cube wraps around the torus.
Let \(A\) be a prefix of this order and \(B\) its
remaining interior factors. Define
\[
 C_{\rm cut}=\frac{216}{e^\zeta-1},\qquad
 W_R(t)=tD^2K_\zeta(C_{\rm cut}\ell^2+3).
 \tag{V71.2}\label{c18v71:W}
\]

\begin{lemma}[Surface-summed crossing rows]
\label{c18v71:cut}
If \(A,B\) are nonempty, then
\[
 \sum_{i\in A}
   e^{-\zeta d(\operatorname{anchor}(i),
                    \{\operatorname{anchor}(j):j\in B\})}
                  \le C_{\rm cut}\ell^2+3 .
 \tag{V71.3}\label{c18v71:rows}
\]
Empty cuts have zero crossing interaction.
Consequently, for any two interior configurations \(x,y\),
with the same arbitrary exterior \(\xi\), let
\(z=(y_A,x_B)\) and \(w=(x_A,y_B)\). Then
\[
 |u(x,\xi)+u(y,\xi)-u(z,\xi)-u(w,\xi)|\le W_R(t).
 \tag{V71.4}\label{c18v71:cross}
\]
\end{lemma}
\begin{proof}
Let \(P\) be the set of completed anchor vertices
strictly preceding the possibly partially processed
anchor \((p,q,s)\). This vertex prefix is a disjoint
union of at most three rectangular boxes:
\[
 \begin{gathered}
 [0,p)\times[0,L)\times[0,L),\\
 \{p\}\times[0,q)\times[0,L),\qquad
 \{p\}\times\{q\}\times[0,s).
 \end{gathered}
\]
Empty boxes are omitted. At a completed anchor,
every factor in \(B\) has its anchor outside \(P\).
For a rectangular box \(H\) with all side lengths
at most \(L<N\), the number of vertices at any
fixed positive depth from \(H^c\) is at most \(6L^2\).
Thus, in periodic graph distance,
\[
 \sum_{v\in H}e^{-\zeta d(v,H^c)}
                   \le \frac{6L^2}{e^\zeta-1}.
\]
For \(v\in H\subset P\), one has
\(d(v,P^c)\ge d(v,H^c)\); the exponential inequality
therefore has the direction needed for an upper bound.
There are at most three factor labels at each vertex.
Summing the preceding bound over the three disjoint
boxes and their factor multiplicities costs
\(54L^2/(e^\zeta-1)=C_{\rm cut}\ell^2\).
The partial anchor adds at most three factors, each
of weight at most one. This proves \eqref{c18v71:rows}.
The six-face count is unchanged under coarse translation.

For \eqref{c18v71:cross}, join \(x_A\) to \(y_A\) and
\(x_B\) to \(y_B\) by independent simultaneous factor
geodesics, parametrized by \(v,w\in[0,1]\), each factor
speed at most \(D\). The mixed derivative of \(u\)
on this two-parameter rectangle is bounded by
\[
 D^2\sum_{i\in A,j\in B}
          \|(\operatorname{Hess}_{m_0}u)_{ij}\|
 \le D^2tK_\zeta
       \sum_{i\in A}e^{-\zeta d(i,B)}.
\]
At each point, the weighted row sum is evaluated at
that same configuration, before taking its supremum.
Integrate over the unit square and use
\eqref{c18v71:rows}. No sum of separate entrywise
suprema is introduced. The exterior does not vary.
\end{proof}

If \(p_\xi\) denotes the density of \(\pi_\xi\) relative
to product Haar, all its normalizers cancel in the
four-point ratio, giving
\[
             p_\xi(x)p_\xi(y)
       \le e^{W_R(t)}p_\xi(z)p_\xi(w).
 \tag{V71.5}\label{c18v71:ratio}
\]
In particular the crossing estimate concerns the
full normalized conditional law, not a cut-off
Hamiltonian or a density with its tails discarded.

\subsection{Continuous coordinate replacement and a block spectral floor}

For this fixed exterior, let \(P_i^\xi\) average
factor \(i\) under its actual conditional law.
These are the restrictions of the full-law \(P_i\)
for \(i\in\mathcal I_R\). Define
\[
 C_R(t)=m\exp(W_R(t)+2tS),\qquad
 \gamma_R(t)=\frac{\lambda_*(t)}{C_R(t)}
       =\frac{3}{2m}\exp(-W_R(t)-3tS).
 \tag{V71.6}\label{c18v71:gamma}
\]

\begin{theorem}[Uniform-exterior block floor with surface cost]
\label{c18v71:floor}
For every entire exterior \(\xi\),
\[
 \begin{split}
 \operatorname{Var}_{\pi_\xi}(f)
 &\le C_R(t)\sum_{i\in\mathcal I_R}
                   \|f-P_i^\xi f\|_{L^2(\pi_\xi)}^2,\\
 \operatorname{Var}_{\pi_\xi}(f)
 &\le \gamma_R(t)^{-1}
                   \int|\nabla_R f|_{m_0}^2\,d\pi_\xi .
 \end{split}
 \tag{V71.7}\label{c18v71:poincare}
\]
The first holds for \(L^2\) functions and the second
on the gradient form domain. The heat-bath generator
has rate one per factor; it is not divided by \(m\).
\end{theorem}
\begin{proof}
We supply the comparison for a positive density on
a compact product, with the paid inputs
\eqref{c18v71:inputs} and \eqref{c18v71:ratio}.
For \(x,y\), let \(z^{k}\) have the first \(k\)
coordinates of \(y\) and the remaining coordinates
of \(x\). Telescoping and Cauchy--Schwarz imply
\[
 \operatorname{Var}_{\pi_\xi}(f)
 \le \frac m2\sum_{k=1}^m
   \iint|f(z^k)-f(z^{k-1})|^2
                          \,d\pi_\xi(x)d\pi_\xi(y).
 \tag{V71.8}\label{c18v71:path}
\]
Fix \(k\), and exchange the prefix \(A=\{1,\ldots,k-1\}\)
between \(x,y\). This is a measure-preserving bijection
for the two copies of product Haar. Write
\(z=z^{k-1}=(y_A,x_B)\), \(w=(x_A,y_B)\).
The increment is \(f(z^{k\leftarrow w_k})-f(z)\).
By \eqref{c18v71:ratio}, its integral in
\eqref{c18v71:path} is at most
\[
 e^{W_R(t)}
   \int d\pi_\xi(z)\int d\pi_\xi^{(k)}(v)
              |f(z^{k\leftarrow v})-f(z)|^2 ,
 \tag{V71.9}\label{c18v71:marginalstep}
\]
where \(\pi_\xi^{(k)}\) is the \(k\)-coordinate
\emph{marginal}, not its conditional law.

To make that distinction quantitative, the
one-coordinate log oscillation at most \(tS\)
implies
\[
 e^{-tS}\le
 \frac{d\pi_\xi(\,\cdot\mid z_{-k})}{d\mu_k}(v)
                  \le e^{tS}.
\]
Averaging the conditional density over \(z_{-k}\)
gives the same two bounds on the marginal density.
Hence
\[
 d\pi_\xi^{(k)}(v)
       \le e^{2tS}d\pi_\xi(v\mid z_{-k}).
 \tag{V71.10}\label{c18v71:marginalratio}
\]
The integral in \eqref{c18v71:marginalstep} is therefore
at most \(2e^{W_R(t)+2tS}\|f-P_k^\xi f\|_2^2\).
Substitution in \eqref{c18v71:path} proves the first
line of \eqref{c18v71:poincare}.
This explicitly pays the conversion of a marginal
draw to an actual conditional draw; treating them
as equal would be an error.

V69's one-factor conditional Poincare bound gives
\[
 \sum_i\|f-P_i^\xi f\|_2^2
       \le\lambda_*^{-1}
                  \int|\nabla_R f|_{m_0}^2\,d\pi_\xi .
\]
This proves the second line. First work with bounded
smooth functions; truncation proves the first inequality
for \(L^2\), and form-domain density proves the second.
Every bound used is uniform in \(\xi\).
\end{proof}

The loss \(m\) is a non-sharp path-length loss, already
present for product measure in this proof. The exponent
\(W_R(t)\) grows like \(\ell^2\), not \(\ell^3\).
In particular this floor is allowed to tend to zero
as the block grows. It is not approximate tensorization
with a volume-independent constant and is not a proof
of global mixing through time four.

\subsection{Paying a genuinely block-centered source in the actual inverse}

Retain the actual moving-metric generator \(\widehat L\)
and \(E_0^{-2}m_0\preceq\widehat m_t\preceq E_0^2m_0\).
Let \(P_R\) be full-block expectation as in V70.

\begin{lemma}[Conditional block centering removes the ambient inverse]
\label{c18v71:blockinverse}
For any \(h\in L^2(\nu)\) satisfying \(P_Rh=0\),
\[
       \|h\|_{-1,\widehat L}^2
          \le E_0^2\gamma_R(t)^{-1}\|h\|_2^2 .
 \tag{V71.11}\label{c18v71:inverse}
\]
Only the finite block floor just proved is used.
\end{lemma}
\begin{proof}
For any real form-domain \(f\),
\[
 \begin{split}
 |\langle h,f\rangle_\nu|
 &=|\langle h,f-P_Rf\rangle_\nu|\\
 &\le \|h\|_2\gamma_R^{-1/2}
                \left(\int|\nabla_Rf|_{m_0}^2d\nu\right)^{1/2}\\
 &\le E_0\|h\|_2\gamma_R^{-1/2}
                \left(\int|df|_{\widehat m_t^{-1}}^2d\nu\right)^{1/2}.
 \end{split}
\]
The second line is \eqref{c18v71:poincare} integrated
over the exterior. Maximizing the dual formula
\eqref{c18v69:dual} gives \eqref{c18v71:inverse}.
No derivative of \(P_Rf\) or of an exterior coefficient
is required, and no global spectral norm is inserted.
\end{proof}

Thus the use of a block solver here has an explicit
paid surface cost. Unlike applying an unspecified
conditional-gap criterion, the proof does not move
the missing estimate into an assumption.

\subsection{The V70 within-block remainder is now controlled}

Use \(E,Q,a_i,p_i,q_i,r_i,M_i\) from V69--V70, and set
\[
 F_i={\bf1}_Ea_i-M_i,\qquad
 G_i=q_i-M_i=r_i-P_Rr_i .
 \tag{V71.12}\label{c18v71:sources}
\]
Both are genuinely block centered. For every exterior,
\[
 \begin{split}
 P_RF_i^2&\le P_R({\bf1}_Ea_i^2)\le S^2Q,\\
 P_RG_i^2&\le P_Rq_i^2
        \le S^2P_Rp_i^2
        \le S^2P_Rp_i
         =S^2P_R{\bf1}_E\le S^2Q.
 \end{split}
 \tag{V71.13}\label{c18v71:rareload}
\]
Here \(|q_i|\le Sp_i\), \(0\le p_i\le1\),
and \(P_RP_i=P_R\). The sharp event is never
differentiated in an interior direction.

\begin{theorem}[Actual inverse of the within-block rare remainder]
\label{c18v71:rare}
For any real \(c\in L^2(\bar\nu_R)\) depending on the
entire exterior, each \(H_i\in\{F_i,G_i\}\) satisfies
\[
 \begin{split}
 \|cH_i\|_{-1,\widehat L}^2
 &\le E_0^2\gamma_R(t)^{-1}S^2Q\|c\|_2^2\\
 &=14E_0^2S^2\ell^3
                   e^{W_R(t)+3tS}Q\|c\|_2^2 .
 \end{split}
 \tag{V71.14}\label{c18v71:rareinverse}
\]
These are bounds on \(F_i,G_i\) themselves, not on
\(T_RF_i\) or \(T_RG_i\).
\end{theorem}
\begin{proof}
Because \(c\) is exterior measurable, \(P_R(cH_i)=0\).
Disintegrating \eqref{c18v71:rareload} gives
\(\|cH_i\|_2^2\le S^2Q\|c\|_2^2\).
Use \eqref{c18v71:inverse} and
\(\gamma_R^{-1}=(2m/3)e^{W_R+3tS}\).
Finally \(m=21\ell^3\). Unbounded \(c\) is admitted
by this \(L^2\) estimate, or by exterior truncation.
No exterior smoothness or transport strain is asserted.
\end{proof}

For \(\eta=6,t=4\), keep
\(\delta=36-21\log2>21\) and V68's size threshold.
If in addition
\[
       \ell\ge\frac{16D^2K_\zeta C_{\rm cut}}{\delta},
 \tag{V71.15}\label{c18v71:size}
\]
then \(4D^2K_\zeta C_{\rm cut}\ell^2\le
\delta\ell^3/4\). Since \(Q\le e^{-\delta\ell^3/2}\),
\[
 \|cH_i\|_{-1,\widehat L}^2
 \le 14E_0^2S^2e^{12D^2K_\zeta+12S}
            \ell^3e^{-\delta\ell^3/4}\|c\|_2^2 .
 \tag{V71.16}\label{c18v71:t4}
\]
This is the useful volume-versus-surface gain: despite
a deteriorating block gap, the actual inverse energy
of these rare centered sources decays with \(\ell\).
All thresholds and prefactors are finite but extremely
large. The assertion applies to every admissible
\(n>\ell\), not to \(\ell\to\infty\) on a fixed torus.

For completeness the same argument applies to V68's
actual conditional temporal score
\(s_R=a-P_Ra\), for which \(|s_R|\le mS\).
Let
\[
 B_R=P_R({\bf1}_Es_R)=\partial_tP_R{\bf1}_E,\qquad
 Y_R={\bf1}_Es_R-B_R .
 \tag{V71.17}\label{c18v71:score}
\]
The derivative includes the full block normalizer,
as in \eqref{c18v68:score}. Then \(P_RY_R=0\) and
\[
 \|cY_R\|_{-1,\widehat L}^2
       \le E_0^2\gamma_R^{-1}m^2S^2Q\|c\|_2^2 .
 \tag{V71.18}\label{c18v71:blockscore}
\]
Under \eqref{c18v71:size} and V68's threshold its
upper bound is \eqref{c18v71:t4} multiplied by \(m^2\),
which is still only a polynomial in \(\ell\).
This is not permission to replace the full temporal
source \(a\) by \(s_R\) in its global inverse equation.

\subsection{Exactly what is left at the exterior}

The V69 centered raw source has the exact decomposition
\[
 \begin{split}
 z_i&={\bf1}_Ea_i-\nu({\bf1}_Ea_i)
       =F_i+(M_i-\nu M_i),\\
 r_i&=G_i+(M_i-\nu M_i).
 \end{split}
 \tag{V71.19}\label{c18v71:remaining}
\]
The inverse energies of \(F_i\) and \(G_i\) are paid
by \eqref{c18v71:rareinverse}; the inverse energy of
\(M_i-\nu M_i\) is not. V70 still supplies its rare
amplitude, distant derivatives and the inverse of
its far-resampling action, not its own inverse.
For the block score the parallel identity is
\[
 {\bf1}_Es_R-\nu({\bf1}_Es_R)
       =Y_R+(B_R-\nu B_R).
 \tag{V71.20}\label{c18v71:remainingblock}
\]
Both remaining means are actual exterior functions.

Enlarging the block cannot simply be asserted to close
this gap. Although V70's collar approximation has an
exponentially small \(L^2\) error in collar radius,
our inverse block cost grows as an exponential of
surface area. Multiplying these two estimates does
not give a summable large-radius inverse bound.
Nor does block centering remove the marginal source:
\(P_R(M_i-\nu M_i)=M_i-\nu M_i\).
The next upstream obligation is a native bound on that
marginal dynamics/source, or an exact source representation
with independently controlled coefficients, not another
application of the now-paid interior solver.

\subsection{Verification, preservation and scope}

The diagnostic
\begin{center}\small
\path{scripts/verify_ym_c18_v71_surface_block_inverse.py}
\end{center}
checks native pair/free anchor orderings and every
prefix's three-box decomposition on translated and
wrapping cubes. Exact rational weighted sums verify
the crossing-row geometry. A separate correlated
finite-state fixture enumerates the hybrid bijection,
density cross ratios, marginal-to-conditional comparison,
and the Poincare form inequality. It checks the
block-centered rare source and its actual anisotropic
inverse budget, including exterior-dependent coefficients.
The latter fixture is not a YM law or a compact-Lie-group
gradient calculation. The all-volume native theorem is
the written proof above, not a numerical extrapolation.
V70's independent derivative regression is replayed.

This step removes the previously unpaid within-block
inverse response in \eqref{c18v70:twopieces}. It does not
remove the exterior mean, prove a uniform global gap,
control every noncentral adverse configuration, or close
the full temporal inverse, Poisson Hessian or transport
strain through time four. All-coarse and interacting-vacuum
extensions, coarse return, physical-scale transport,
a nontrivial reflection-positive interacting continuum
and the positive physical Yang--Mills mass gap remain
unproved. The full goal remains active.

Removing this append and restoring V70's title recovers
V70 byte for byte. Only the active YM filename is replaced;
other sources and legacy archives remain unchanged.
All typesetting, PDF-skill visual checks and diagnostics
are outside \path{latex_notes}, which contains only
\texttt{.tex} files. Next companion review: V75;
next broad literature sweep: V80; archive/restart: V100.
% END CYCLE 18 VERSION 71 APPEND
% BEGIN CYCLE 18 VERSION 72 APPEND
\clearpage
\section[V72: Whole-torus bounds and complete rare-source control]{V72: Whole-torus surface bounds and complete mesoscopic rare-source control}
\label{c18v72:context}

V71 paid the within-block response but retained the
exterior mean. Here that mean is bounded by a proved
whole-torus inverse estimate, with its ambient-volume
dependence fully displayed. Its cost is exponential
in the torus surface scale \(n^2\). A sufficiently
large mesoscopic droplet, of side comparable to
\(n^{2/3}\), has enough rare-volume suppression to
overcome this cost. The resulting estimates apply
to complete centered rare sources, not just their
single-factor or block-recentered components.

There is also a source-level consequence for the
actual \(a_t=\partial_tu_t\): the inverse energy
of its centered restriction to the union of all
translated specified mesoscopic droplets tends to
zero as \(n\) grows. No conditional-mean remainder
is discarded in that statement. The complement
remains in the equation and is not proved coercive.
All estimates concern the existing \(\beta=0\),
coarse-identity normal-chart law. The bound below
is a finite-volume auxiliary spectral floor which
deteriorates with volume, not a physical mass gap.

The canonical-path method and its primary-source
attribution are inherited from V71. The new geometry
is the whole periodic torus, which V71's hypothesis
\(\ell<n\) did not cover. The new analytic use is
retaining the mean, and balancing its ambient inverse
cost against rare volume. We do not replace this
with a hypothetical volume-independent marginal gap.
The scheduled reviews remain V75 and V80.

\subsection{The periodic whole-torus crossing estimate}

Let \(\mathcal I_n\) be all pair/free factors on the
fine torus of side \(N=2n\). Their number is
\[
                         m_n=21n^3 .
\]
Keep \(D,S,K_\zeta,C_{\rm cut},E_0\) from V69--V71.
At most three factors have any one fine anchor.
Order all factors lexicographically by fine coordinates
in \(\{0,\ldots,N-1\}^3\), then by direction.
Distances are periodic graph distances, not distances
in an artificially opened cube. Put
\[
 W_n(t)=tD^2K_\zeta(C_{\rm cut}n^2+3),\qquad
 \Gamma_n(t)=\frac{1}{14n^3}
                       e^{-W_n(t)-3tS}.
 \tag{V72.1}\label{c18v72:constants}
\]

\begin{lemma}[A proper periodic prefix has surface crossing cost]
\label{c18v72:geometry}
For any nonempty proper factor prefix \(A\), with
\(B=\mathcal I_n\setminus A\),
\[
           \sum_{i\in A}e^{-\zeta d(i,B)}
                         \le C_{\rm cut}n^2+3 .
 \tag{V72.2}\label{c18v72:cut}
\]
The corresponding four-point log-density difference
has absolute value at most \(W_n(t)\).
\end{lemma}
\begin{proof}
As in V71, the completed anchor vertices before
the current anchor \((p,q,s)\) are the disjoint
union of
\[
 \begin{gathered}
 [0,p)\times[0,N)\times[0,N),\\
 \{p\}\times[0,q)\times[0,N),\qquad
 \{p\}\times\{q\}\times[0,s).
 \end{gathered}
\]
Each nonempty box is a proper subset of the torus.
Some directions now run around a complete periodic
cycle. Such a direction has no boundary face and
does not enter the distance to the box complement.
In a proper direction the two end faces give at
most \(2N^2\) vertices at each fixed depth. Therefore
every positive-depth layer has at most \(6N^2\)
vertices, even if two of the directions are full
cycles. Summing the geometric series gives
\(6N^2/(e^\zeta-1)\) for each box.

Distance to the completed prefix complement is
at least distance to the complement of the box
containing the anchor. Sum over at most three boxes,
at most three labels per anchor, and at most three
partially processed labels. The result is
\(54N^2/(e^\zeta-1)+3=C_{\rm cut}n^2+3\).
This proves \eqref{c18v72:cut} without assuming
that a periodic face is an exterior boundary.

The mixed-geodesic rectangle calculation in the
proof of Lemma \ref{c18v71:cut} uses only the
weighted row bound \(\mathcal H_\zeta(u_t)\le tK_\zeta\)
and the displayed crossing sum. It gives the
four-point bound on the full native density.
All configurations, including noncentral ones,
are allowed in this comparison.
\end{proof}

\begin{theorem}[Explicit whole-torus auxiliary spectral floor]
\label{c18v72:global}
For \(0\le t\le4\), the actual law satisfies
\[
 \operatorname{Var}_{\nu_t}(f)
       \le\Gamma_n(t)^{-1}
                       \int|df|_{m_0^{-1}}^2d\nu_t.
 \tag{V72.3}\label{c18v72:poincare}
\]
Consequently, for any centered \(h\in L^2(\nu_t)\),
\[
       \|h\|_{-1,\widehat L_t}^2
             \le E_0^2\Gamma_n(t)^{-1}\|h\|_2^2.
 \tag{V72.4}\label{c18v72:inverse}
\]
\end{theorem}
\begin{proof}
Use V71's continuous coordinate-replacement proof
on all \(m_n\) factors, now with the four-point
cost \(W_n(t)\) supplied by the preceding lemma.
The marginal-to-conditional draw conversion still
costs \(e^{2tS}\), not one. Thus
\[
 \operatorname{Var}_{\nu_t}(f)
       \le m_ne^{W_n(t)+2tS}
                         \sum_{i\in\mathcal I_n}\|D_if\|_2^2 .
\]
The one-factor floor
\(\lambda_*=\tfrac32e^{-tS}\) gives
\eqref{c18v72:poincare}, since \(2m_n/3=14n^3\).
Heat-bath updates still have rate one per factor.
Moving-metric comparison and the centered inverse
variational formula give \eqref{c18v72:inverse}.
No smallness of \(tK_\zeta\), thermodynamic mixing,
or physical spectral comparison is used.
\end{proof}

The lower bound for the actual auxiliary gap is
\(E_0^{-2}\Gamma_n(t)\). Its exponential dependence
on \(n^2\) is essential to what follows. In particular,
one cannot set \(n=\infty\) in this estimate to
obtain a positive spectral floor. We have paid an
explicit finite ambient inverse, not removed its
volume dependence.

\subsection{A complete rare source, with its mean retained}

Fix a V68 block \(R\), \(2\le\ell<n\), its event \(E\),
and the all-exterior conditional bound \(Q\).
Let \(b\) be a real measurable function on the entire
fibre with \(|b|\le\mathfrak B\). Set
\[
 J={\bf1}_E b,\qquad M=P_RJ,\qquad U=J-M .
\]
Alternatively replace \(J\) by \(P_i({\bf1}_E b)\)
for any interior factor \(i\); the same \(M\)
is obtained by the tower property. In both cases,
\[
 |M|\le\mathfrak B Q,\qquad
 P_RU=0,\qquad P_RU^2\le\mathfrak B^2Q .
 \tag{V72.5}\label{c18v72:rare}
\]
For the projected choice, conditional Jensen and
\(P_RP_i=P_R\) prove the last bound. This permits
sharp events without taking their derivatives.

For a real \(c\in L^2(\bar\nu_R)\) of the entire
exterior, define the complete centered bank source
\[
 \begin{split}
 Z(c)&=cJ-\nu_t(cJ)\\
     &=cU+\{cM-\nu_t(cM)\}.
 \end{split}
 \tag{V72.6}\label{c18v72:split}
\]
The braces contain the actual remaining mean.
For a variable bank it is not generally correct
to replace them by \(c(M-\nu_tM)\).

\begin{theorem}[Full rare-bank inverse with explicit ambient cost]
\label{c18v72:rareinverse}
For either choice of \(J\),
\[
 \begin{split}
 \|cM-\nu_t(cM)\|_{-1,\widehat L_t}^2
       &\le E_0^2\mathfrak B^2
                          \Gamma_n(t)^{-1}Q^2\|c\|_2^2,\\
 \|Z(c)\|_{-1,\widehat L_t}^2
       &\le E_0^2\mathfrak B^2
          \left(\gamma_R(t)^{-1/2}\sqrt Q
                 +\Gamma_n(t)^{-1/2}Q\right)^2\|c\|_2^2 .
 \end{split}
 \tag{V72.7}\label{c18v72:bank}
\]
Here \(\gamma_R(t)\) is V71's proved block floor.
\end{theorem}
\begin{proof}
Equation \eqref{c18v72:rare} and disintegration give
\[
 \|cU\|_2^2\le\mathfrak B^2Q\|c\|_2^2,\qquad
 \|cM-\nu_t(cM)\|_2^2\le\mathfrak B^2Q^2\|c\|_2^2.
\]
The first source is block centered, so V71's
\eqref{c18v71:inverse} applies. The second is
globally centered, so use \eqref{c18v72:inverse}.
Apply the triangle inequality in the actual negative
energy norm to \eqref{c18v72:split}, and square.
Although the two pieces are orthogonal in ordinary
\(L^2\), no orthogonality for \(\widehat L_t^{-1}\)
is asserted. Exterior truncation covers unbounded
\(L^2\) banks; no derivative of \(c\) is required.
\end{proof}

With \(b=a_i\) and \(\mathfrak B=S\), the choices of
\(J\) give the complete \(z_i\) and \(r_i\) from V69
when \(c=1\). Thus their conditional means are now
included, at the explicit ambient cost. With
\(b=s_R=a-P_Ra\), take \(\mathfrak B=m_RS\);
this includes V71's mean \(B_R\). A fixed \(\ell\)
still leaves a bound growing exponentially in
\(n^2\). The useful uniform decay statement below
therefore requires a growing, mesoscopic block,
not a fixed local observable.

\subsection{A late-time mesoscopic regime with all constants displayed}

The rare-volume estimate can be made uniform on a
late interval, not only at \(t=4\). Choose
\[
 \frac{7\log2}{3}<\tau\le4,\qquad
 \delta_\tau=9\tau-21\log2>0 .
 \tag{V72.8}\label{c18v72:late}
\]
For example, \(\tau=2\) gives \(\delta_\tau>3\).
Let \(B_\zeta\) be V68's exterior surface constant
and define
\[
 \begin{gathered}
 L_{\rm prob}(\tau)
       =\frac{96(e^{96}-1)+8B_\zeta}{\delta_\tau},\\
 \alpha=4D^2K_\zeta C_{\rm cut},\qquad
 b_*=12D^2K_\zeta+12S,\qquad
 L_*(\tau)=\max\{2,L_{\rm prob}(\tau),4\alpha/\delta_\tau\}.
 \end{gathered}
 \tag{V72.9}\label{c18v72:thresholds}
\]
For \(\eta=6\), every translated block and every
entire exterior satisfy
\[
 Q_{6,R}(t)\le e^{-\delta_\tau\ell^3/2}
       \quad(\tau\le t\le4,\ \ell\ge L_{\rm prob}(\tau)).
 \tag{V72.10}\label{c18v72:Q}
\]
Indeed V68's exponent is at most
\[
 -\delta_\tau\ell^3+
             \{48(e^{96}-1)+4B_\zeta\}\ell^2.
\]
This uses the unchanged event and all-exterior bound
from V68. No new conditional mixing assertion enters.

If
\[
 \ell\ge L_*(\tau),\qquad
             \ell^3\ge\frac{2\alpha}{\delta_\tau}n^2,
 \tag{V72.11}\label{c18v72:meanregime}
\]
then, uniformly for \(t\in[\tau,4]\), the two squared
contributions in the proof of \eqref{c18v72:bank}
are bounded respectively by
\[
 \begin{split}
 A_{\rm int}&=14E_0^2\mathfrak B^2 e^{b_*}
                       \ell^3e^{-\delta_\tau\ell^3/4}\|c\|_2^2,\\
 A_{\rm mean}&=14E_0^2\mathfrak B^2 e^{b_*}
                       n^3e^{-\delta_\tau\ell^3/2}\|c\|_2^2 .
 \end{split}
 \tag{V72.12}\label{c18v72:two}
\]
For the first, \(W_R(t)+3tS\le\alpha\ell^2+b_*\)
and \(\alpha\ell^2\le\delta_\tau\ell^3/4\).
For the second, \(W_n(t)+3tS\le\alpha n^2+b_*\),
the rare squared factor is \(Q^2\), and
\(\alpha n^2\le\delta_\tau\ell^3/2\).
Thus
\[
 \|Z(c)\|_{-1,\widehat L_t}^2
 \le 28E_0^2\mathfrak B^2e^{b_*}
   \left(\ell^3e^{-\delta_\tau\ell^3/4}
             +n^3e^{-\delta_\tau\ell^3/2}\right)\|c\|_2^2 .
 \tag{V72.13}\label{c18v72:complete}
\]
We used \((\sqrt{A_{\rm int}}+\sqrt{A_{\rm mean}})^2
\le2(A_{\rm int}+A_{\rm mean})\); the inverse cross
term was not dropped. For a family of banks the bound
is in operator norm from the actual exterior \(L^2\)
space at each time.

\subsection{The actual temporal source on all translated droplets}

The global score itself has a paid, if extensive,
amplitude bound:
\[
                     |a_t|\le m_nS=21n^3S .
 \tag{V72.14}\label{c18v72:amplitude}
\]
Each one-coordinate oscillation of \(a_t\) is at most
\(S\), by V60--V69. Connect any two configurations by
changing their \(m_n\) factors one at a time, then
average the second endpoint under \(\nu_t\).
Since \(\nu_t(a_t)=0\), this proves
\eqref{c18v72:amplitude} without replacing \(a_t\)
by a local conditional score. It could already be
used in \eqref{c18v72:complete} with
\(\mathfrak B=m_nS\).

A direct estimate controls a whole rare sector at once.
Let \(\mathcal R_{n,\ell}\) be all \(n^3\) coarse
translations of the side-\(\ell\) cube, and put
\[
 \begin{split}
 \mathcal U_{n,\ell}
       &=\bigcup_{R\in\mathcal R_{n,\ell}}\mathcal E_{6,R},\\
 h_{\rm drop}(t)
       &={\bf1}_{\mathcal U_{n,\ell}}a_t
                    -\nu_t({\bf1}_{\mathcal U_{n,\ell}}a_t).
 \end{split}
 \tag{V72.15}\label{c18v72:union}
\]
Overlaps are allowed. No independence between droplets
is asserted, and no sum of overlapping indicators is
substituted for the union indicator.

\begin{theorem}[Complete mesoscopic droplet-source response]
\label{c18v72:actualsource}
For \(\ell\ge L_{\rm prob}(\tau)\) and
\(t\in[\tau,4]\),
\[
 \begin{split}
 \nu_t(\mathcal U_{n,\ell})
                   &\le n^3e^{-\delta_\tau\ell^3/2},\\
 \|h_{\rm drop}(t)\|_{-1,\widehat L_t}^2
 &\le 6174E_0^2S^2e^{b_*}n^{12}
                      e^{\alpha n^2-\delta_\tau\ell^3/2}.
 \end{split}
 \tag{V72.16}\label{c18v72:unionenergy}
\]
In particular, if
\[
                  \ell^3\ge\frac{4\alpha}{\delta_\tau}n^2,
 \tag{V72.17}\label{c18v72:unionregime}
\]
then
\[
 \|h_{\rm drop}(t)\|_{-1,\widehat L_t}^2
       \le6174E_0^2S^2e^{b_*}
                    n^{12}e^{-\delta_\tau\ell^3/4}.
 \tag{V72.18}\label{c18v72:uniondecay}
\]
\end{theorem}
\begin{proof}
The first line of \eqref{c18v72:unionenergy} is the
union bound and \eqref{c18v72:Q}. Centering decreases
ordinary squared \(L^2\) norm, so
\[
 \|h_{\rm drop}(t)\|_2^2
   \le m_n^2S^2\nu_t(\mathcal U_{n,\ell}).
\]
Apply the actual inverse bound \eqref{c18v72:inverse}
and \(W_n(t)+3tS\le\alpha n^2+b_*\).
The coefficient is \(14\cdot21^2=6174\);
the powers of \(n\) are \(3+6+3=12\), from
the inverse prefactor, score amplitude squared
and number of translates. Finally
\eqref{c18v72:unionregime} absorbs
\(\alpha n^2\) into \(\delta_\tau\ell^3/4\).
All conditional and global centerings have been
retained; no remainder is omitted.
\end{proof}

To exhibit an admissible scale, set
\[
 d_*=\max\{1,(4\alpha/\delta_\tau)^{1/3}\},
 \qquad
                  \ell_n=\lceil d_*n^{2/3}\rceil .
 \tag{V72.19}\label{c18v72:scale}
\]
For example, every integer \(n\) larger than
\[
 \max\{5,(d_*+1)^3,(L_*(\tau)/d_*)^{3/2}\}
\]
has \(L_*(\tau)\le\ell_n<n\) and satisfies
\eqref{c18v72:unionregime}. Hence the bound
\[
 \sup_{\tau\le t\le4}
 \|h_{\rm drop}(t)\|_{-1,\widehat L_t}^2
 \le6174E_0^2S^2e^{b_*}
                  n^{12}e^{-\delta_\tau d_*^3 n^2/4}
                         \longrightarrow0
 \tag{V72.20}\label{c18v72:vanish}
\]
is explicit. Also \(\ell_n/n\to0\). The scale is
mesoscopic in this finite-lattice thermodynamic
sense only; \(d_*\) and the onset volume are enormous.
It does not identify a physical continuum scale,
and the fibre Hilbert spaces for different \(n\)
have not been identified as a continuum field space.

\subsection{The retained complement is the next source problem}

Let
\[
 h_{\rm rest}(t)
   ={\bf1}_{\mathcal U_{n,\ell}^c}a_t
                   -\nu_t({\bf1}_{\mathcal U_{n,\ell}^c}a_t).
\]
The exact centered identity and its actual inverse are
\[
 \begin{gathered}
 a_t=h_{\rm drop}(t)+h_{\rm rest}(t),\\
 \widehat L_t^{-1}a_t-\widehat L_t^{-1}h_{\rm rest}(t)
                      =\widehat L_t^{-1}h_{\rm drop}(t).
 \end{gathered}
 \tag{V72.21}\label{c18v72:actualsplit}
\]
Thus \eqref{c18v72:vanish} bounds the total moving-metric
Dirichlet energy of the difference between these two
Poisson responses. It is not the energy of
\(\widehat L_t^{-1}a_t\) spatially restricted to the
rare event. Its response may spread throughout
the fibre. Nor is it a Hessian or transport-strain
estimate; that stronger norm remains unproved.

The union only covers V68's specified near-central,
high-negative-flux droplets at a growing scale.
Its complement contains all sorts of noncentral
configurations and is not proved a favorable-curvature
region. Small droplets at fixed support are not
uniformly handled by this estimate. In particular,
the exponentially worsening ambient inverse cannot
be used to infer a uniform inverse on a fixed-block
mean. The remaining \(h_{\rm rest}\), smaller-scale
mean response, and source-complete mixing/transport
remain the next upstream obligations. A negligible
source sector is not a positive spectral bottom for
all physical excitations.

\subsection{Verification and version handoff}

The diagnostic
\begin{center}\small
\path{scripts/verify_ym_c18_v72_torus_rare_source.py}
\end{center}
checks whole-torus native pair/free orderings, including
periodic full-cycle directions in the prefix boxes.
Exact rational layer sums and independent periodic
distance checks verify the surface crossing input.
A correlated finite-state regression tests the global
Poincare comparison, variable-bank centering, the
nonorthogonal inverse splitting, and the full centered
source on overlapping-event unions. It also checks the
late-time and mesoscopic exponent bookkeeping exactly.
The fixture is not native YM; the native analytic
claims are proved above, not inferred numerically.
V71's independent block-floor regression is replayed.

The progress is complete rare-source inverse control,
including conditional means, in a stated volume-dependent
mesoscopic regime, and a vanishing actual temporal-source
droplet contribution on a late interval. The full mass
gap remains unproved. Uniform small-scale mean control,
the complementary temporal-source inverse, Poisson
Hessian, strain, all-coarse and interacting-vacuum
extensions, coarse return, physical clocks/scales and
a nontrivial interacting reflection-positive continuum
remain open.

V71 is recoverable byte for byte by removing this append
and restoring its title. Only the active YM filename
is replaced; other sources and archives are preserved.
Only \texttt{.tex} files are kept in \path{latex_notes};
builds, diagnostics and PDF-skill visual checks stay
outside. Next companion checkpoint: V75; broad primary
literature sweep: V80; archive/restart: V100.
The full goal remains active.
% END CYCLE 18 VERSION 72 APPEND
% BEGIN CYCLE 18 VERSION 73 APPEND
\clearpage
\section[V73: Sparse conditional coercivity and exact source feedback]{V73: Sparse conditional coercivity through time four and exact source feedback}
\label{c18v73:context}

V72 controls complete specified mesoscopic rare sources, but
does not give a useful inverse bound for their nonrare
complement. Applying its whole-torus comparison to that
complement would reintroduce the exponential surface cost.
We therefore go upstream to the source decomposition, not
to a further estimate of the same droplet probability.

This append proves a volume-independent conditional gap
on sufficiently separated sets of native factors throughout
\(0\le t\le4\). Such a set can contain a fixed positive
fraction of all factors; the fraction and the gap can be
extremely small. The separation makes the already-paid
spatial tail smaller than the already-paid one-factor floor.
Unlike the full-fibre test in V62, this restricted margin
is actually verified. It pays a nonrare, conditionally
centered component of the full temporal source.

We then eliminate this fast component in the actual
moving-metric Dirichlet form. The exact remaining source
contains a feedback term, not just the marginal temporal
score. Its feedback and the effective metric comparison
are bounded explicitly. The remaining marginal source
inverse and a full physical mass gap are not claimed.
The scope remains the electric
\(\beta=0\), coarse-identity normal-chart law, with
\(N=2n\), \(n\ge5\).

\subsection{A separated native set with a populated margin}

Use the notation and constants of V61--V72:
\[
 \nu_t=e^{u_t}\mu_I,\quad a_t=\partial_tu_t,\quad
 S=D^2K_*,\quad D=\pi\sqrt2,\quad K_\zeta=e^{2\zeta}K_* .
\]
The inputs, already proved on the actual law for \(0\le t\le4\), are
\[
 \begin{gathered}
 \mathcal H_\zeta(u_t)\le tK_\zeta,\qquad
 \|\nabla_i a_t\|_\infty\le DK_*,
 \qquad \operatorname{gap}(\nu_{t,i}^{\,z})\ge
                  \tfrac32e^{-tS},\\
 E_0^{-2}m_0\preceq\widehat m_t\preceq E_0^2m_0 .
 \end{gathered}
 \tag{V73.1}\label{c18v73:inputs}
\]
The last bound is the geometric comparison through time four,
not an extension of V61's short-time probability-strain result.
Set
\[
 \begin{gathered}
 \lambda_4=\tfrac32e^{-4S},\qquad \rho=\lambda_4/2,\\
 R_*=\max\left\{1,\left\lceil
                  \zeta^{-1}\log(8K_\zeta/\lambda_4)
                              \right\rceil\right\}.
 \end{gathered}
 \tag{V73.2}\label{c18v73:separation}
\]
All these constants are independent of \(n\). A set \(A\)
of pair/free factors is called separated if distinct
\(i,j\in A\) have periodic anchor distance \(d(i,j)\ge R_*\).
In particular two labels with the same anchor cannot both
belong to \(A\).

\begin{lemma}[A positive-density separated class]
\label{c18v73:colour}
The \(m_n=21n^3\) native factors can be partitioned into at most
\[
                       J_*=3(2R_*-1)^3
 \tag{V73.3}\label{c18v73:colours}
\]
separated classes. One class has at least \(m_n/J_*\) factors.
This statement holds without divisibility of \(N\) by \(R_*\).
\end{lemma}
\begin{proof}
Join distinct labels when their periodic anchor distance is
less than \(R_*\). A periodic ball of that radius is contained
in the image of an integer cube with \((2R_*-1)^3\) points.
There are at most three native labels per anchor, including
pair labels only at their first half. Thus each graph degree
is at most \(J_*-1\). Greedy vertex colouring uses at most
\(J_*\) colours; each colour is a separated class.
The largest class has the stated size.
If \(R_*\) exceeds the torus diameter, classes may be singletons;
the assertion is still valid. No false wrap-around separation
from a coordinate residue colouring is used.
\end{proof}

Let \(P_A\) denote integration over the entire set \(A\) under
its actual conditional law, with every factor in \(A^c\)
frozen; write \(Q_A=I-P_A\).
These are orthogonal conditional projections in \(L^2(\nu_t)\),
not products of the one-factor projections.

\begin{theorem}[Uniform sparse conditional gap through time four]
\label{c18v73:sparse}
For every separated \(A\), every entire exterior \(z\),
and every \(0\le t\le4\),
\[
 \operatorname{Var}_{\nu_{t,A}^{\,z}} f
     \le \rho^{-1}
       \int\sum_{i\in A}|\nabla_i f|_{m_0}^2\,d\nu_{t,A}^{\,z}.
 \tag{V73.4}\label{c18v73:conditionalgap}
\]
Consequently, for the actual full Dirichlet form
\(\mathcal E_t(f)=\int|df|_{\widehat m_t^{-1}}^2d\nu_t\),
\[
        \|Q_Af\|_2^2
          \le\rho^{-1}\sum_{i\in A}\|\nabla_i f\|_2^2
          \le E_0^2\rho^{-1}\mathcal E_t(f).
 \tag{V73.5}\label{c18v73:fastform}
\]
\end{theorem}
\begin{proof}
At each fixed configuration, the weighted row estimate gives
\[
 \sum_{\substack{j\in A\\j\ne i}}
       \|(\operatorname{Hess}u_t)_{ij}\|_{\rm op}
       \le tK_\zeta e^{-\zeta R_*}
       \le\lambda_4/2,\qquad i\in A .
\]
Here the supremum remains outside the row sum. We have not
summed separate configuration suprema. The conditional
one-factor laws inside \(A\) are the same native one-factor
laws with all their other coordinates fixed, so their scalar
floors are at least \(\lambda_4\).

Apply the conditional integrated Bochner identity (V62.1)
and sum (V62.2) only over factors in \(A\), with exterior \(z\)
fixed. The mixed derivative squares are nonnegative.
The signed mixed Hessian matrix is symmetric, and its
off-diagonal block row sums are at most \(\lambda_4/2\).
The elementary symmetric block Schur bound therefore gives
\[
      \|L_A f\|_{2,z}^2
         \ge(\lambda_4/2)\sum_{i\in A}\|\nabla_i f\|_{2,z}^2.
\]
The conditional product is compact and connected, with smooth
positive density. Apply this inequality to its nonconstant
elliptic eigenfunctions to obtain (V73.4).
Disintegrate in \(z\) and use the inverse-metric comparison
in (V73.1) to obtain (V73.5).
The argument requires scalar exact gradients, not a gap for
arbitrary one-forms.
\end{proof}

The one-factor-versus-mixed-Hessian mechanism is standard;
the primary context is G.~Menz,
\emph{A Brascamp--Lieb type covariance estimate},
Theorems 2.3 and 2.7,
\url{https://arxiv.org/html/1402.5160}.
We use V62's compact-product proof, not the Euclidean theorem
without its hypotheses. The new input here is the native
separation that verifies the margin at the full smoothing time.
It proves neither independence of the separated variables
nor mixing between different classes.

\subsection{A nonrare part of the actual source is controlled}

For centered \(h\in L^2(\nu_t)\), define the actual inverse
energy by its usual centered variational formula. Theorem
\ref{c18v73:sparse} implies
\[
 \|Q_Ah\|_{-1,\widehat L_t}^2
       \le E_0^2\rho^{-1}\|Q_Ah\|_2^2 .
 \tag{V73.6}\label{c18v73:fastinverse}
\]
Indeed, \(\langle Q_Ah,f\rangle=\langle Q_Ah,Q_Af\rangle\);
Cauchy--Schwarz and (V73.5), followed by optimization, prove
(V73.6). Thus, for the actual full temporal score,
\[
 \begin{split}
 \|Q_Aa_t\|_2^2&\le \rho^{-1}|A|D^2K_*^2,\\
 \|Q_Aa_t\|_{-1,\widehat L_t}^2
             &\le E_0^2\rho^{-2}|A|D^2K_*^2 .
 \end{split}
 \tag{V73.7}\label{c18v73:source}
\]
This is a volume-extensive bound with a volume-independent
coefficient, valid at every configuration and exterior and
through time four. No rare-event indicator is needed.
It is not the full source estimate:
\[
                         a_t=Q_Aa_t+P_Aa_t .
 \tag{V73.8}\label{c18v73:sourcesplit}
\]
The second term remains globally centered, but need not vanish
or have a small norm. A large cardinality of \(A\) does not
prove that a fixed fraction of the source energy is captured.

The same projection sharpens what can be said about V72's
complement. Write its exact split as
\(a_t=h_{\rm drop}+h_{\rm rest}\), on the same fixed torus.
By \(L^2\) contraction, (V72.14)--(V72.16), and (V73.6),
\[
 \begin{split}
 \|Q_Ah_{\rm drop}\|_{-1,\widehat L_t}^2
 &\le E_0^2\rho^{-1}(21n^3S)^2
                       \nu_t(\mathcal U_{n,\ell})\\
 &\le441E_0^2\rho^{-1}S^2n^9
                         e^{-\delta_\tau\ell^3/2},
                  \qquad \tau\le t\le4 .
 \end{split}
 \tag{V73.9}\label{c18v73:sparsedrop}
\]
As before, \(\ell\ge L_{\rm prob}(\tau)\), \(2\le\ell<n\),
and \(\tau>7\log2/3\). There is no ambient surface exponential.
For any \(\epsilon>0\), choosing
\[
 \ell_n=\left\lceil\max\left\{2,L_{\rm prob}(\tau),
       \left(\frac{2(9+\epsilon)\log n}{\delta_\tau}\right)^{1/3}
                              \right\}\right\rceil
 \tag{V73.10}\label{c18v73:logscale}
\]
gives \(\ell_n<n\) for all sufficiently large \(n\) and makes
(V73.9) at most \(441E_0^2\rho^{-1}S^2n^{-\epsilon}\).
It bounds the total Dirichlet energy of
\[
 \widehat L_t^{-1}Q_Aa_t-
       \widehat L_t^{-1}Q_Ah_{\rm rest}
                  =\widehat L_t^{-1}Q_Ah_{\rm drop}.
\]
This logarithmic-volume droplet scale only controls the
\(Q_A\)-centered source. It does not replace V72's larger
scale for the complete rare source, and it does not bound
the complementary conditional mean. The sharp indicator is
never differentiated.

\subsection{Eliminate only the fast operator whose inverse is paid}

An exact form reduction identifies what must be solved next.
Fix \(t,n,A\), suppress their subscripts, and work in real
mean-zero spaces. Let
\[
 \mathcal V=H^1(M_I,\nu)\cap\{\nu f=0\},\qquad
 \mathcal F=\{v\in\mathcal V:P_Av=0\},\qquad
 \mathcal Y=\{y\in\mathcal V:P_Ay=y\}.
\]
We use the actual \(\mathcal E\) throughout, not a product
metric substituted into the final inverse.

The projection is bounded on the form domain, with a
volume-independent gradient constant. For an exterior derivative,
\[
 \nabla_jP_Af=P_A\nabla_j f+
             P_A[(f-P_Af)\nabla_j u].
 \tag{V73.11}\label{c18v73:projectionderivative}
\]
Put \(b_{\rm pr}(t)=tK_*/\rho\),
\(C_{\rm pr}(t)=1+b_{\rm pr}(t)^2\), and let
\(\bar\nu_t\) be the actual \(A^c\) marginal. Its product
metric form is denoted by
\(\bar{\mathcal E}_0(y)=\int|\nabla_{A^c}y|^2d\bar\nu_t\).
Then
\[
 \bar{\mathcal E}_0(P_Af)
       \le C_{\rm pr}(t)\int|\nabla f|_{m_0}^2d\nu_t
       \le E_0^2C_{\rm pr}(t)\mathcal E_t(f).
 \tag{V73.11a}\label{c18v73:projectionenergy}
\]
Here is the proof, including the normalization covariance.
Fix an exterior tangent vector \(v\) and write
\(F_v=\langle\nabla_{A^c}u_t,v\rangle\). Conditional
Poincare and Cauchy--Schwarz give
\[
 \begin{split}
 |\operatorname{Cov}_A(f,F_v)|
 &\le \rho^{-1}
       \bigl(P_A|\nabla_A f|^2\bigr)^{1/2}
       \bigl(P_A|(\operatorname{Hess}u_t)_{A,A^c}v|^2\bigr)^{1/2}\\
 &\le b_{\rm pr}(t)|v|\bigl(P_A|\nabla_A f|^2\bigr)^{1/2}.
 \end{split}
\]
The last step uses the full operator bound
\(\|\operatorname{Hess}u_t\|_{\rm op}\le tK_*\), not a sum
over exterior coordinates. Consequently (V73.11) yields
\[
 |\nabla_{A^c}P_Af|
   \le\bigl(P_A|\nabla_{A^c}f|^2\bigr)^{1/2}
             +b_{\rm pr}(t)\bigl(P_A|\nabla_A f|^2\bigr)^{1/2}.
\]
Square, use \((r+b_{\rm pr}(t)s)^2\le C_{\rm pr}(t)(r^2+s^2)\),
and integrate.
This proves (V73.11a), initially for smooth functions and
then by approximation for \(H^1\).
The finite-volume elliptic gap, or (V72.3) with metric
comparison, makes \(\mathcal V\) a Hilbert space with
inner product \(\mathcal E\); \(\mathcal F\) is closed.

For any centered \(h\in L^2(\nu)\), there is a unique
\(w_h\in\mathcal F\) such that
\[
          \mathcal E(w_h,v)=\langle h,v\rangle
                         \quad(v\in\mathcal F).
 \tag{V73.12}\label{c18v73:w}
\]
The fast coercivity (V73.5) and Riesz representation give
\[
 \begin{split}
 \mathcal E(w_h)&\le E_0^2\rho^{-1}\|Q_Ah\|_2^2,\\
 \mathcal E(w_{a_t})&\le E_0^2\rho^{-2}|A|D^2K_*^2.
 \end{split}
 \tag{V73.13}\label{c18v73:paidfast}
\]
In operator language, the form compression of \(\widehat L\)
to \(\ker P_A\) has floor \(E_0^{-2}\rho\).
This is a populated fast-complement inverse, not an assumed
gap for the whole operator. Generally
\(w_h\ne\widehat L^{-1}Q_Ah\).

For \(y\in\mathcal Y\), let \(k_y\in\mathcal F\) be the
energy projection determined by
\[
 \mathcal E(k_y,v)=\mathcal E(y,v)\quad(v\in\mathcal F).
\]
Define the harmonic lift, effective form, and effective source
functional by
\[
 \begin{split}
 Hy&=y-k_y,\qquad
 \mathfrak s(y,z)=\mathcal E(Hy,Hz),\\
 \ell_h(y)&=\langle P_Ah,y\rangle-\mathcal E(w_h,y).
 \end{split}
 \tag{V73.14}\label{c18v73:effective}
\]
The last expression is a functional on the exterior form
domain; an \(L^2\) density for it is not needed or asserted.
All operators and projections here depend on \(t\).

\begin{theorem}[Exact source-aware energy reduction]
\label{c18v73:schur}
The full actual inverse energy satisfies
\[
 \boxed{\quad
 \|h\|_{-1,\widehat L}^2
     =\mathcal E(w_h)+
          \sup_{y\in\mathcal Y}
                 \{2\ell_h(y)-\mathfrak s(y,y)\}.
       \quad}
 \tag{V73.15}\label{c18v73:identity}
\]
Moreover,
\[
 \mathfrak s(y,y)
       =\inf_{v\in\mathcal F}\mathcal E(y+v)
       =\mathcal E(y,y)-\mathcal E(k_y,k_y).
 \tag{V73.16}\label{c18v73:schurform}
\]
The fast term has a volume-uniform coercivity input.
The remaining dual supremum is retained, not set to zero.
\end{theorem}
\begin{proof}
Energy projection gives \(\mathcal E(Hy,v)=0\) for
\(v\in\mathcal F\), and \(P_AHy=y\).
Every \(f\in\mathcal V\) has the unique decomposition
\[
                  f=x+Hy,\qquad y=P_Af,\quad x\in\mathcal F .
\]
Pythagoras proves (V73.16) and
\(\mathcal E(f)=\mathcal E(x)+\mathfrak s(y,y)\).
Also
\[
 \begin{split}
 \langle h,Hy\rangle
   &=\langle P_Ah,y\rangle-\langle Q_Ah,k_y\rangle\\
   &=\langle P_Ah,y\rangle-\mathcal E(w_h,k_y)
     =\ell_h(y).
 \end{split}
\]
For the last equality use \(w_h\in\mathcal F\) in the defining
equation of \(k_y\). Thus the full variational supremum splits
into independent suprema over \(x\) and \(y\).
The former is attained at \(w_h\), with value
\(\mathcal E(w_h)\), proving (V73.15).
At fixed volume the latter supremum is finite because the
full compact elliptic inverse exists; this existence statement
does not supply a uniform bound on it.
\end{proof}

This is the usual Schur completion of squares, proved here
at form level to avoid illicit products of unbounded blocks.
For a finite block matrix it reads
\[
 A_{\rm slow}^{\rm eff}=C-B^*A_{\rm fast}^{-1}B,\qquad
 h_{\rm slow}^{\rm eff}=h_{\rm slow}
                         -B^*A_{\rm fast}^{-1}h_{\rm fast}.
\]
Both corrections are retained. In particular the reduced
form is no larger than the compression \(\mathcal E(y,y)\).
A lower bound on that compression would not automatically
be a lower bound on \(\mathfrak s\). Nor may the feedback
\(\mathcal E(w_h,y)\) be discarded when \(P_Ah\) is small.
These are not assertions of a physical transfer-operator
Schur gap.

\begin{proposition}[Paid effective metric and feedback comparison]
\label{c18v73:slowcomparison}
For every \(y\in\mathcal Y\),
\[
 \frac{\bar{\mathcal E}_0(y)}{E_0^2C_{\rm pr}(t)}
        \le\mathfrak s(y,y)
        \le E_0^2\bar{\mathcal E}_0(y).
 \tag{V73.16a}\label{c18v73:effectivecomparison}
\]
For the feedback functional \(b_h(y)=-\mathcal E(w_h,y)\),
its dual energy norm satisfies
\[
 \|b_h\|_{-1,\mathfrak s}^2
       :=\sup_{y\in\mathcal Y}\{2b_h(y)-\mathfrak s(y,y)\}
       \le E_0^2\rho^{-1}\|Q_Ah\|_2^2 .
 \tag{V73.16b}\label{c18v73:feedbackinverse}
\]
\end{proposition}
\begin{proof}
For the lower form bound use (V73.11a) on \(Hy\), since
\(P_AHy=y\). For the upper bound take \(v=0\) in
(V73.16), and then use metric comparison.
For the feedback, \(Q_AHy=-k_y\), so (V73.5) gives
\(\|k_y\|_2^2\le E_0^2\rho^{-1}\mathfrak s(y,y)\).
The identity
\(\mathcal E(w_h,y)=\langle Q_Ah,k_y\rangle\)
then proves (V73.16b) by Cauchy--Schwarz and optimization.
No lower bound on the full slow spectral gap is used.
\end{proof}

There is a convenient consequence directly in source norms.
Let \(\bar L_0\) denote the actual marginal's product-metric
weighted generator. Its finite-volume centered inverse exists,
but no volume-uniform gap is presumed. For every centered \(h\),
\[
 \|h\|_{-1,\widehat L_t}
   \le E_0\rho^{-1/2}\|Q_Ah\|_2
            +E_0\sqrt{C_{\rm pr}(t)}\|P_Ah\|_{-1,\bar L_0}.
 \tag{V73.16c}\label{c18v73:twosource}
\]
Indeed
\(\langle h,f\rangle=\langle Q_Ah,Q_Af\rangle+
\langle P_Ah,P_Af\rangle\).
Bound the first pairing by (V73.5), and the second by
marginal inverse-energy duality and (V73.11a).
Taking the supremum over unit-energy \(f\) proves the result.
In particular,
\[
 \|a_t\|_{-1,\widehat L_t}
    \le E_0\rho^{-1}\sqrt{|A|}\,DK_*
           +E_0\sqrt{C_{\rm pr}(t)}\|P_Aa_t\|_{-1,\bar L_0}.
 \tag{V73.16d}\label{c18v73:fullsource}
\]
Thus the fast source, the change of form, and the source
feedback have explicit costs. The actual marginal temporal
source inverse on the right remains to be controlled.
This is a two-scale estimate with its fast hypotheses paid,
not a proof that the slow quantity is small.

The sparse choice is stable under fixed pinned gauge
transformations: they act factor by factor, preserve the
law and the actual metric, and commute with \(P_A\).
The constructions consequently preserve the invariant
subspace for invariant sources. A particular colour
choice need not preserve lattice translations; no such
symmetry of an individual fast correction is required.

\subsection{The marginal score is not the effective source}

Write \(x\) for the factors in \(A\), \(y\) for those in \(A^c\),
and define the actual normalized marginal density by
\[
 \bar u_t(y)=\log\int e^{u_t(x,y)}\,d\mu_A(x).
\]
Because the factor set is independent of \(t\), differentiation
at fixed exterior gives
\[
                         \partial_t\bar u_t=P_Aa_t .
 \tag{V73.17}\label{c18v73:marginalscore}
\]
This is the marginal temporal score, with zero marginal mean.
It is not the effective forcing \(\ell_{a_t}\) in (V73.14).

There is a second obstruction to repeating the sparse gap
argument on the marginal without new estimates. For exterior
slots \(j,k\), differentiation of the normalized integral gives
\[
 \operatorname{Hess}_{jk}\bar u_t
    =P_A(\operatorname{Hess}_{jk}u_t)
          +\operatorname{Cov}_A(\nabla_j u_t,\nabla_k u_t).
 \tag{V73.18}\label{c18v73:marginalhessian}
\]
The covariance is a positive semidefinite block matrix on
exterior tangent vectors. For a fixed such vector \(v\), apply
(V73.4) to
\(F_v=\langle\nabla_{A^c}u_t,v\rangle\), with \(y,v\) fixed.
Since \(\nabla_AF_v=(\operatorname{Hess}u_t)_{A,A^c}v\),
\[
 \begin{split}
 0\preceq\operatorname{Cov}_A(\nabla_{A^c}u_t)
   &\preceq \rho^{-1}
          P_A\!\left[(\operatorname{Hess}u_t)_{A,A^c}^*
                     (\operatorname{Hess}u_t)_{A,A^c}\right]\\
   &\preceq \rho^{-1}(tK_*)^2I .
 \end{split}
 \tag{V73.19}\label{c18v73:feedbackbound}
\]
All inequalities are quadratic-form inequalities at the fixed
exterior. Hence a valid, but potentially much worse, bound is
\[
       \|\operatorname{Hess}\bar u_t\|_{\rm op}
                      \le tK_*+\rho^{-1}(tK_*)^2 .
 \tag{V73.20}\label{c18v73:marginalbound}
\]
This is an operator bound, not an anchored weighted-row bound.
The original native constants are not automatically inherited
after marginalization. The positive covariance contribution
can enlarge positive log-density curvature. The actual
effective Dirichlet form has its additional correction
(V73.16); its comparison with the marginal product-metric
form is now paid by (V73.16a), with the displayed large
constant. What is not paid is the marginal source inverse
itself, or stability of the required spatial bounds under
further elimination.

\subsection{Verification and the next upstream problem}

The diagnostic
\begin{center}\small
\path{scripts/verify_ym_c18_v73_sparse_source_reduction.py}
\end{center}
checks native periodic conflict-graph colourings, separation,
and cardinality bounds at sample integer separations; it does
not numerically instantiate the enormous native \(R_*\).
Exact rational matrix regressions check
the conditional projection, fast coercivity, nonzero feedback,
the corrected slow form, the full inverse-energy identity,
and the feedback and two-scale source inequalities.
A separate exactly solvable normalized integral checks the
positive marginal-Hessian covariance. These are algebraic and
geometric regressions, not simulations or certificates of an
interacting YM vacuum. The analytic proofs above supply the
all-volume claims. V72's math regression is replayed.

The advance is a paid, volume-uniform fast sector at time four
for the actual normal-chart law, together with its nonrare
temporal-source bound and exact remaining marginal-source problem.
The sparse fraction can be astronomically small. No uniform
inverse of the retained marginal or effective form, source-specific
pointwise Hessian, transport strain, or global mixing is
proved. The feedback functional and effective metric now
have the explicit bounds (V73.16a)--(V73.16d).
The next useful work is control of \(P_Aa_t\) in the actual
marginal inverse, with a stable spatial structure under
elimination, not another application of the same one-factor
gap or a reset of the marginal to Haar.
V72's complement has not been declared coercive.

All-coarse and interacting-vacuum extensions, coarse return,
physical-scale and clock control, nontrivial interacting
reflection-positive continuum reconstruction and the positive
physical YM mass gap remain open. The full goal stays active.

V72 is recoverable by removing this append and restoring
its title; only its active filename is replaced. Other sources
and archives are unchanged. Only \texttt{.tex} files are stored
in \path{latex_notes}; verification and build files stay outside.
The next companion checkpoint is V75, the next broad primary
literature sweep V80, and archive/restart V100.
% END CYCLE 18 VERSION 73 APPEND
% BEGIN CYCLE 18 VERSION 74 APPEND
\clearpage
\section[V74: Finite-step marginal locality and source propagation]{V74: Finite-step marginal locality, exact source propagation and the cost of iteration}
\label{c18v74:context}

V73 paid a sparse conditional inverse and the metric and
source-feedback costs of one elimination. Its remaining
quantity was the actual marginal temporal-source inverse.
A repeated argument could not simply reuse the original
native Hessian constants. Here we derive the missing spatial
bounds for the marginal itself, including the normalization
covariance, and propagate a localized representation of the
full temporal source.

The outcome is a nonperturbative, finite-step marginalization
theorem through smoothing time four. One-factor oscillation
bounds survive every marginal exactly; exponential spatial
decay survives each fixed number of suitably sparse
eliminations, with displayed losses. These facts do not give
a contraction under arbitrarily many eliminations. We compute
the growth of the resulting scalar majorants and retain the
terminal marginal source. This rules out declaring the
iteration closed merely because each individual step exists.

The scope remains \(N=2n\), \(n\ge5\), \(\beta=0\), coarse
identity, the actual normal-chart law
\(\nu_t=e^{u_t}\mu_I\), and \(0\le t\le4\).
All retained coordinates keep their original pair/free
product metrics and periodic anchor distances. Marginalizing
them is not a physical change of lattice spacing or clock.

\subsection{One-factor floors survive arbitrary marginalization}

Write \(I_0=\mathcal I_n\), and for any retained set
\(J\subset I_0\) define the actual marginal density by
\[
 e^{u_t^J(x_J)}
       =\int e^{u_t(x_J,x_{J^c})}\,d\mu_{J^c}(x_{J^c}).
 \tag{V74.1}\label{c18v74:marginal}
\]
There is no additional normalization: both sides give the
normalized marginal of \(\nu_t\). Retain
\[
 S=D^2K_*,\qquad
 \lambda=\tfrac32e^{-4S},\qquad \rho=\lambda/2 .
 \tag{V74.2}\label{c18v74:floor}
\]

\begin{lemma}[An unchanged marginal one-factor oscillation bound]
\label{c18v74:oscillation}
For every \(J\), every \(i\in J\), every exterior in \(J\setminus\{i\}\),
and \(0\le t\le4\),
\[
     \operatorname{osc}_i u_t^J\le tS,\qquad
     \operatorname{gap}((\nu_t^J)_i^{\,z})
                  \ge3\kappa_i e^{-tS}\ge\lambda .
 \tag{V74.3}\label{c18v74:stablefloor}
\]
Here \(\kappa_i=1/2\) for a pair and \(1\) for a free factor.
\end{lemma}
\begin{proof}
For two values \(x_i,x'_i\), the original pointwise
one-factor bound gives
\[
 e^{-tS}e^{u_t(x'_i,x_{-i})}
   \le e^{u_t(x_i,x_{-i})}
   \le e^{tS}e^{u_t(x'_i,x_{-i})}
\]
for every eliminated configuration as well as the fixed
retained exterior. Integrating the same inequalities over
\(J^c\) proves the oscillation bound for the logarithm of
(V74.1). The one-factor conditional normalizer cancels
in the density ratio. Comparing its variance and gradient
form with the one-factor Haar law costs \(e^{tS}\);
the Haar gap is \(3\kappa_i\). This proves (V74.3).
No independence of retained factors follows.
\end{proof}

Thus the diagonal floor need not be degraded after each
elimination. Spatial mixed response, rather than the existence
of one-factor inverses, is still the substantive issue.

\subsection{An entrywise spatial envelope and a weighted sparse response}

For a nonnegative matrix \(K\) indexed by factor labels,
including rectangular matrices, use
\[
 \|K\|_\xi
 =\max\left\{\sup_i\sum_j e^{\xi d(i,j)}K_{ij},
             \sup_j\sum_i e^{\xi d(i,j)}K_{ij}\right\}.
 \tag{V74.4}\label{c18v74:matrixnorm}
\]
Triangle inequality for the original periodic anchor distance
implies \(\|KL\|_\xi\le\|K\|_\xi\|L\|_\xi\).
This remains true for labels sharing an anchor; a
pseudometric is enough. All matrices below dominate
separate configuration suprema, not a falsely interchanged
supremum and row sum.

For \(\eta>0\), put
\[
 V_\eta=3\left(\frac{1+e^{-\eta}}{1-e^{-\eta}}\right)^3 .
 \tag{V74.5}\label{c18v74:volume}
\]
At most three original factor labels have one anchor.
Injecting periodic shortest displacements into \(\mathbb Z^3\)
therefore gives
\(\sum_i e^{-\eta d(i,j)}\le V_\eta\), uniformly in \(n\)
and also after removing any labels.

V63's weighted Hessian estimate supplies the explicit
time-uniform initial envelope and norm bound
\[
 \begin{gathered}
 k^{(0)}_{ij}=4K_\zeta e^{-\zeta d(i,j)},\qquad
 \xi_0=\zeta/2,\\
 \sup_x\|(\operatorname{Hess}u_t)_{ij}(x)\|_{\rm op}
                     \le k^{(0)}_{ij},\\
 K_0=\max\{\lambda,4K_\zeta V_{\zeta/2}\},\qquad
                    \|k^{(0)}\|_{\xi_0}\le K_0 .
 \end{gathered}
 \tag{V74.6}\label{c18v74:initial}
\]
The loss from \(\zeta\) to \(\xi_0\) pays the entrywise
suprema. This is V63's envelope mechanism with the exact
native anchor multiplicity, not a new claim of a smaller
unweighted Hessian constant.

Now suppose a current marginal on \(J\) has the floor
(V74.3), and its full Hessian blocks, including diagonal
blocks, are dominated by a symmetric nonnegative
configuration-independent matrix \(k\) with
\(\|k\|_\xi\le K\), \(K\ge\lambda\).
Choose \(0<\eta<\xi\), and a factor set \(A\subset J\)
such that the anchors of any two distinct factor labels
in \(A\) are at distance at least
\[
 R=\max\left\{1,\left\lceil
             \frac{\log(2K/\lambda)}{\xi-\eta}\right\rceil\right\}.
 \tag{V74.7}\label{c18v74:R}
\]
Such sets with positive factor density are provided by V73's
conflict-graph colouring, now on \(J\).
Set \(B=J\setminus A\), let \(k_{AA}^{\circ}\) have zero
diagonal, and define the numerical matrix
\[
                   G_A=(\lambda I-k_{AA}^{\circ})^{-1}.
 \tag{V74.8}\label{c18v74:G}
\]

\begin{lemma}[A paid weighted conditional response]
\label{c18v74:response}
The inverse in (V74.8) exists, is entrywise nonnegative, and
\[
       \|G_A\|_\eta\le\rho^{-1}.
 \tag{V74.9}\label{c18v74:Gbound}
\]
For the actual conditional law on \(A\), with every \(B\)
coordinate fixed, and smooth real \(f,h\),
\[
 |\operatorname{Cov}_A(f,h)|
       \le\sum_{i,j\in A}(G_A)_{ij}
                   \|\nabla_i f\|_{2,A}\|\nabla_j h\|_{2,A}.
 \tag{V74.10}\label{c18v74:covariance}
\]
Its scalar product-metric gap is at least \(\rho\).
All assertions are uniform in the entire retained exterior.
\end{lemma}
\begin{proof}
Separation and (V74.7) give
\[
 \|k_{AA}^{\circ}\|_\eta
       \le e^{-(\xi-\eta)R}K\le\lambda/2=\rho .
\]
Consequently
\[
 G_A=\lambda^{-1}\sum_{m\ge0}
                        (k_{AA}^{\circ}/\lambda)^m
\]
converges in the weighted matrix norm, is nonnegative, and
has norm at most \(1/(\lambda-\rho)=1/\rho\).
The unweighted symmetric matrix
\(\lambda I-k_{AA}^{\circ}\) has floor at least \(\rho\).

For clarity, the analytic input is V62's compact-product
directional calculation, not a Gaussian replacement.
For \(\psi=L_A^{-1}(f-P_Af)\), its exact-gradient
conditional Bochner identity gives
\[
 \lambda r_i\le b_i+\sum_{\substack{j\in A\\j\ne i}}k_{ij}r_j,
 \qquad
 r_i=\|\nabla_i\psi\|_{2,A},\quad b_i=\|\nabla_i f\|_{2,A}.
\]
It uses the genuine one-factor scalar gap for the complete
diagonal term; no pointwise convexity of that term is needed.
The case \(r_i=0\) satisfies the same row inequality.
Multiplication by \(G_A\), followed by
\(\operatorname{Cov}_A(f,h)=
\sum_i\langle\nabla_i\psi,\nabla_i h\rangle_A\),
proves (V74.10), using symmetry of \(G_A\).
The same integrated identity, summed and tested on
eigenfunctions, proves the conditional gap.
Compact ellipticity ensures existence at fixed regulator.
\end{proof}

This is the directional covariance mechanism of
\href{https://arxiv.org/html/1402.5160}{G.~Menz,
\emph{A Brascamp--Lieb type covariance estimate}},
Theorems 2.3 and 2.7, in the compact implementation already
proved in V62. What is supplied here is its weighted
native margin after actual marginalization. The Euclidean
result is not imported without its hypotheses.

\subsection{The full marginal Hessian and projected source profile}

Let \(P_A\) be actual conditional expectation on \(A\) for
the current marginal, and
\(u^B=\log\int e^{u^J}\,d\mu_A\).
Then the next marginal Hessian has the entrywise envelope
\[
              k^+=k_{BB}+k_{BA}G_Ak_{AB},\qquad
              \|k^+\|_\eta\le K+K^2/\rho .
 \tag{V74.11}\label{c18v74:envelopeupdate}
\]
In particular
\(\sup_x\|(\operatorname{Hess}u^B)_{bc}\|_{\rm op}
\le k^+_{bc}\) for every \(b,c\in B\).
This includes its diagonal blocks.

\begin{proof}
The normalized marginal identity from V73 is
\[
 \operatorname{Hess}_{bc}u^B
    =P_A(\operatorname{Hess}_{bc}u^J)
            +\operatorname{Cov}_A(\nabla_bu^J,\nabla_cu^J).
\]
At a fixed exterior, contract the last block against unit
tangent vectors in factors \(b,c\). The two scalar functions
on \(A\) have gradient profiles bounded respectively by
\(k_{ib}\) and \(k_{jc}\). Equation (V74.10) bounds the
covariance block by
\(\sum_{i,j\in A}k_{bi}(G_A)_{ij}k_{jc}\).
The direct term is bounded by \(k_{bc}\), proving the
entrywise assertion. The weighted matrix norm and (V74.9)
prove its stated bound. The covariance is retained, not
discarded by centering or a claim of conditional independence.
\end{proof}

More generally, if a smooth \(f\) on the current coordinates
has pointwise derivative bounds \(\|\nabla_i f\|_\infty\le F_i\),
then
\[
 \|\nabla_bP_Af\|_\infty
       \le F_b+\sum_{i,j\in A}k_{bj}(G_A)_{ji}F_i,
                    \qquad b\in B .
 \tag{V74.12}\label{c18v74:profilemap}
\]
Indeed differentiate the normalized integral:
\(\nabla_bP_Af=P_A\nabla_bf+
\operatorname{Cov}_A(f,\nabla_bu^J)\), then apply (V74.10).

If \(F_i\le L e^{-\xi d(i,o)}\) for any original anchor \(o\),
the triangle inequality and the weighted column/row bounds give
\[
       \|\nabla_bP_Af\|_\infty
          \le L(1+K/\rho)e^{-\eta d(b,o)} .
 \tag{V74.13}\label{c18v74:localizedprofile}
\]
To verify the summation explicitly, multiply the feedback
sum in (V74.12) by \(e^{\eta d(b,o)}\), use
\(d(b,o)\le d(b,j)+d(j,i)+d(i,o)\), and bound
\(F_i e^{\eta d(i,o)}\le L\).
The resulting row sum of \(k_{BA}G_A\) is at most \(K/\rho\).
The direct term costs at most \(L\). No factor \(|A|\)
or exterior volume is introduced.

These are spatial derivative bounds, not merely an
unweighted operator estimate. In particular \(k^+\) is
a deterministic envelope for the actual normalized marginal.
It need not have the original finite-range Wilson form,
and it need not have a small interaction norm.

\subsection{A localized representation of the actual full temporal score}

A source profile centered at one anchor is useful only if
it occurs in the actual source. Here it does. Fix an ordering
\(1,\ldots,m_n\) of all original factors and a reference point
\(o\in M_I\), independent of \(t\). Write \(x^{[i]}\) for
the configuration using \(x_1,\ldots,x_i\) and reference
values on all later factors, and put
\[
       \phi_{i,t}(x)=a_t(x^{[i]})-a_t(x^{[i-1]}).
 \tag{V74.14}\label{c18v74:telescope}
\]
Then
\[
 \begin{gathered}
 a_t=\sum_{i\in I_0}(\phi_{i,t}-\nu_t\phi_{i,t}),\qquad
 |\phi_{i,t}|\le S,\\
 \|\nabla_j\phi_{i,t}\|_\infty
                 \le D K_\zeta e^{-\zeta d(i,j)}
                         \quad\hbox{for every }i,j .
 \end{gathered}
 \tag{V74.15}\label{c18v74:nativesource}
\]

\begin{proof}
The differences telescope to \(a_t(x)-a_t(o)\).
Subtract its \(\nu_t\)-mean and use \(\nu_ta_t=0\).
The one-coordinate oscillation bound gives \(|\phi_{i,t}|\le S\).
For \(j>i\) its derivative is zero. For \(j=i\), it is
bounded by \(DK_*\le DK_\zeta\).
For \(j<i\), differentiate the difference along factor \(j\),
then integrate the mixed Hessian of \(a_t\) on a minimizing
factor-\(i\) geodesic of length at most \(D\).
V63 bounds each such mixed block by
\(K_\zeta e^{-\zeta d(i,j)}\), proving the last assertion.
The factor-\(j\) tangent space is the same at the two
endpoints of this path because \(i\ne j\).
No differentiability of a chosen minimizing geodesic in
its endpoints is needed: only the fixed-endpoint estimate
of the already smooth difference is used.
\end{proof}

If invariant summands are desired, replace \(\phi_{i,t}\)
by its average \(\gamma_{i,t}\) under the compact pinned
gauge group. The action is factorwise isometric and the law
and \(a_t\) are invariant. Averaging therefore preserves
every displayed bound and the centered sum identity.
The chosen reference and ordering need not themselves
be invariant or translation covariant. In what follows use
these averaged summands, and set \(L_0=DK_\zeta\).
They have exponentially decaying derivative profiles, not
strict finite support.

For any retained \(J\) let
\[
 \gamma_{i,t}^{\,J}
       =\mathbb E_{\nu_t}[\gamma_{i,t}\mid x_J].
\]
The actual marginal temporal score, for a set \(J\) chosen
independently of time, is
\[
 a_t^J=\partial_tu_t^J
      =\mathbb E_{\nu_t}[a_t\mid x_J]
      =\sum_{i\in I_0}
                   (\gamma_{i,t}^{\,J}-\nu_t\gamma_{i,t}).
 \tag{V74.16}\label{c18v74:scoreidentity}
\]
The first equality after the derivative follows directly
by differentiating (V74.1). The tower property proves the
last equality. Every original source label stays in this
sum, including labels of variables that have been eliminated.
Deleting those labels would not be the actual marginal score.

\subsection{A finite-depth theorem with every cost retained}

Start with \(I_0,k^{(0)},\xi_0,K_0\) from (V74.6).
For \(r\ge0\), recursively set
\[
 \begin{gathered}
 \xi_{r+1}=\xi_r/2,\qquad
 K_{r+1}=K_r+K_r^2/\rho,\qquad
 L_{r+1}=L_r(1+K_r/\rho),\\
 R_r=\max\left\{1,\left\lceil
                  \frac{2}{\xi_r}\log(2K_r/\lambda)
                                 \right\rceil\right\},
 \qquad J_r=3(2R_r-1)^3 .
 \end{gathered}
 \tag{V74.17}\label{c18v74:recursion}
\]
Colour the conflict graph on the current \(I_r\), choose a
largest \(R_r\)-separated colour class \(A_r\), and set
\(I_{r+1}=I_r\setminus A_r\).
All choices use geometry and the time-uniform constants,
not \(t\) or the current field configuration.
Define \(k^{(r+1)}\) by the matrix update (V74.11).
Stop if no variables remain; the final centered score is
then zero.

\begin{theorem}[Actual marginal locality at every fixed depth]
\label{c18v74:finite}
For each finite depth \(r\), before termination, uniformly
in \(n\), \(0\le t\le4\), and all retained configurations:
\[
 \begin{gathered}
 \operatorname{osc}_i u_t^{I_r}\le tS,\qquad
 \sup_x\|(\operatorname{Hess}u_t^{I_r})_{ij}(x)\|_{\rm op}
                                      \le k^{(r)}_{ij},\\
 \|k^{(r)}\|_{\xi_r}\le K_r,\qquad
 |\gamma_{i,t}^{\,I_r}|\le S,\qquad
 \|\nabla_j\gamma_{i,t}^{\,I_r}\|_\infty
                       \le L_r e^{-\xi_r d(i,j)},\\
 L_r=L_0K_r/K_0,\qquad
 \|\nabla_j a_t^{I_r}\|_\infty\le G_r:=L_rV_{\xi_r}.
 \end{gathered}
 \tag{V74.18}\label{c18v74:induction}
\]
Each \(A_r\)-conditional law has scalar gap at least \(\rho\).
\end{theorem}
\begin{proof}
The initial envelope is (V74.6), and (V74.15) gives the
initial source profile, also at the smaller weight \(\xi_0\).
All marginals have the floor (V74.3).
Apply (V74.7)--(V74.13) with
\((\xi,\eta,K)=(\xi_r,\xi_{r+1},K_r)\).
This proves the next Hessian envelope and source profile,
and supplies the next conditional gap without assuming a
gap for the full marginal.
The composition of successive normalized conditional
expectations equals the original conditional expectation
onto \(I_{r+1}\), by the tower property.
Thus the propagated functions are precisely those in
(V74.16), not arbitrary local test functions.
Conditional expectation preserves the amplitude bound.

Both \(L_r\) and \(K_r\) multiply by \(1+K_r/\rho\)
at each step, giving \(L_r/L_0=K_r/K_0\).
Finally differentiate (V74.16) and sum its profiles over
all original anchors, using (V74.5). This yields \(G_r\),
with no original-volume factor.
All constants depend on the fixed depth \(r\); uniformity
in \(r\) is not part of this theorem.
\end{proof}

To keep the actual inverse problem visible, let \(L_{t,r}\)
be the positive weighted generator of the marginal
\(\nu_t^{I_r}\) in its retained product metric.
Here \(L_{t,r}\) is an operator, whereas the scalar \(L_r\)
above is a source-profile constant.
Set
\[
 C_r=1+(K_r/\rho)^2,\qquad
 \Lambda_0=1,\quad \Lambda_r=\prod_{j<r}\sqrt{C_j}.
\]
V73's two-scale source proof applied to each current
marginal gives, for every available finite depth \(r\),
\[
 \begin{split}
 \|a_t\|_{-1,\widehat L_t}
 \le E_0\bigg\{
    &\sum_{j<r}\Lambda_j\rho^{-1}\sqrt{|A_j|}\,G_j\\
    &+\Lambda_r
               \|a_t^{I_r}\|_{-1,L_{t,r}}\bigg\}.
 \end{split}
 \tag{V74.19}\label{c18v74:sourceiteration}
\]
Indeed at stage \(j\),
\(\|Q_{A_j}a_t^{I_j}\|_2
 \le\rho^{-1/2}\sqrt{|A_j|}G_j\),
and the conditional-projection energy constant is \(C_j\).
The latter uses
\(\|\operatorname{Hess}u_t^{I_j}\|_{\rm op}\le K_j\).
Iterate the norm inequality (V73.16c).
The factor \(E_0\) appears only once, in the initial comparison
between the actual moving metric and the original product
metric; intermediate marginals use their product metrics.
No terminal inverse, conditional mean, or source feedback
has been omitted. A finite-volume terminal zero does not
make the constants uniform in volume.

\subsection{Why these scalar bounds do not close the infinite iteration}

The recurrence is not a contraction estimate. Define
\(x_r=K_r/\rho\); then \(x_0\ge2\) and
\[
 x_{r+1}=x_r+x_r^2,\qquad
 x_r\ge2^{\,2^r},\qquad
 \Lambda_r\ge2^{\,2^r-1}.
 \tag{V74.20}\label{c18v74:growth}
\]
The first lower bound follows inductively from
\(x_{r+1}\ge x_r^2\). Since \(\sqrt{1+x_r^2}\ge x_r\),
the product bound follows by summing \(2^j\) for \(j<r\).
These are lower bounds on the \emph{chosen scalar majorants},
not on the actual interactions, inverse energies, or physical
correlation lengths.

The guaranteed colour densities also deteriorate:
\[
 R_r\ge\frac{2\log2}{\xi_0}\,4^r,\qquad
             \sum_{r\ge0}J_r^{-1}<\infty,\qquad
             \prod_{r\ge0}(1-J_r^{-1})>0.
 \tag{V74.21}\label{c18v74:densityloss}
\]
Indeed \(\xi_r=\xi_0/2^r\) and
\(\log(2K_r/\lambda)=\log x_r\ge2^r\log2\).
Since \(R_r\ge1\),
\(J_r=3(2R_r-1)^3\ge3R_r^3\), giving a summable geometric
upper bound on \(J_r^{-1}\).
Also \(J_r\ge3\), so convergence of that series makes
the infinite product strictly positive.

Choosing a largest class only guarantees
\[
       |I_r|\le m_n\prod_{j<r}(1-J_j^{-1}).
 \tag{V74.22}\label{c18v74:remainingcount}
\]
The right-hand fraction has a positive limit.
This bound therefore does not certify exhaustion at a
volume-uniform rate. It does not prove that the actual
algorithm leaves a positive fraction: every nonempty
finite set loses at least one factor, and eventually
terminates at finite volume. Nor does (V74.20) rule out
better matrix, signed, or native estimates.

What has been established is a faithful finite-step spatial
and source calculus for the actual normal-chart marginals.
What has not been established is stable contraction of that
calculus as depth grows. The next attack must retain more
native information in the covariance feedback, exploit
source-specific cancellation, or supply a genuine mixing
or renormalization mechanism. Repeating this scalar
majorant recurrence is not a route to a proved uniform
marginal source inverse.

\subsection{Verification and handoff}

The diagnostic
\begin{center}\small
\path{scripts/verify_ym_c18_v74_marginal_locality.py}
\end{center}
tests weighted matrix convolution, positive sparse response,
and sequential envelope/profile updates with exact rational
arithmetic. Separate finite-state regressions check marginal
density-ratio preservation and the normalized temporal-score
tower identities, with every original source summand retained.
The tests also verify the scalar recurrence and its
noncontractive growth. They are not native probability
simulations or a proof-assistant certification.
The all-volume analytic statements are proved above;
V73's exact math regression is replayed.

For non-duplication, f16.2 V79's colour argument concerned
independent original Wilson-history links and explicitly did
not identify a marginalized colour with the same Wilson law.
No such independence is asserted here. V63's separate-supremum
envelope and V73's fast/source-form calculations are inputs,
not counted again as new results. The present advance is
the spatially weighted actual marginal update and exact
propagation of the full native temporal-source representation.

The next scheduled companion review is V75, and the next
broad primary-literature sweep is V80. The feedback/iteration
loss identified here should inform that checkpoint.
The actual marginal source inverse, pointwise Poisson
Hessian and probability strain, interacting-vacuum and
all-coarse extensions, coarse return, physical scale/clock,
nontrivial interacting reflection-positive continuum and
full positive physical YM mass gap remain unproved.
The complete goal remains active.

V73 is recoverable byte for byte by removing this append and
restoring its title. Only the active filename is replaced;
other sources and archives are preserved. Only \texttt{.tex}
files remain in \path{latex_notes}; all checks and build
products are outside that folder. Archive/restart remains V100.
% END CYCLE 18 VERSION 74 APPEND
% BEGIN CYCLE 18 VERSION 75 APPEND
\clearpage
\section[V75: Native marginal source regeneration]{V75: Native marginal source regeneration and the first straight-path interaction}
\label{c18v75:context}

V74 established spatial control at every fixed elimination depth,
but its absolute-value majorants grew rather than contracted.
A proposed improvement by signed cancellation must be tested on
the actual density, including its own time dependence, before
it is used in an iteration. This appendix computes a complete
local marginal two-jet and then the first nonzero source in a
straight-path marginal. The native density term and the
normalization covariance reinforce each other in this channel.
Neither term may be removed.

The scope is the actual normal-chart probability
\(\nu_t=e^{u_t}\mu_I\), at \(\beta=0\), coarse identity,
\(N=2n\), \(n\ge5\). Exact symmetry statements below hold
through \(0\le t\le4\). Taylor expansions are at \(t=0\)
at a fixed regulator. Their displayed coefficients hold
for every \(n\), but no volume-uniform Taylor remainder
or positive time interval is asserted. This is not a
replacement Gibbs measure \(e^{3tW_{\rm cap}}\mu_I\).

\subsection{Scheduled companion checkpoint and retained inputs}

The V75 checkpoint rehashed twelve active companion sources
against V70 and reread their terminal mathematical scope.
All twelve hashes are unchanged. The core SM V72 and NS V26
notes respectively concern a fixed-time, fixed-mode auxiliary
many-copy limit and pointwise rank-one kernel norms.
The shell-mixing V16 note retains its complete flat one-loop
coefficient but leaves the mixed noncommuting curvature
coefficient and a nonperturbative remainder unpaid.
The source-selection V65, stationarity V61 and exact-action
V55 notes still retain their historical-occurrence,
state-only-source and nonlinear slow-source gates.
The arithmetic and perturbative ESM companions do not supply
the present native marginal inverse either.

The useful transfer is methodological: keep the actual
source and its complete normalization terms. No companion
gap, physical clock or uniform nonperturbative estimate
is imported. The full checkpoint is recorded in
\begin{center}\small
\path{artifacts/ym_c18_v75_companion_review.json}.
\end{center}

For context, exact link integration and effective-variable
representations are established lattice-gauge methods; see
H.~Vairinhos and P.~de Forcrand,
\href{https://arxiv.org/abs/1409.8442}
{\emph{Lattice gauge theory without link variables}}.
Only its primary abstract was checked in this targeted review.
Its auxiliary-field representation for the Wilson action
is not identified with our normal-chart law, and no
theorem from it is used below. The scheduled broad sweep
remains V80.

The inherited native identities needed here are V54's
four-path geometry, V55's time-reflection operation,
V56's vanishing second normalization derivative, and
V57's scalar formula for the density two-jet. Write
\[
 C=W_{\rm cap},\quad U=W_{\rm uncap},\quad W=C+U,\quad
 L=L_0,\quad b=\left.\partial_tu_t\right|_0=3C,\quad
 j=\left.\partial_t^2u_t\right|_0 .
\]
Then
\[
 LC=9C,\quad LW=12W-3C,\quad
 S=|Q_0\nabla W|^2,\qquad
 j=(L-12)S+3\Gamma(W,C).
 \tag{V75.1}\label{c18v75:native}
\]
V56 proves \(\mu_I(j+b^2)=0\); thus the raw and normalized
logarithmic second derivatives agree here. These identities
are inputs, not new results of V75.

\subsection{Four paths at one centre, with the original kinetic metric}

Fix a fine vertex \(y\) with exactly two odd coordinates.
There are \(3n^3\) such centres. Number its four neighbouring
midpoints cyclically \(0,1,2,3\), and let \(Z_r\) be V54's
path from the corresponding coarse root to \(y\): a prefix
half-link followed by a free spoke, with inverses where required.
At coarse identity the four captured plaquettes at \(y\) are
\[
 c_0=w_{01},\quad c_1=w_{12},\quad
 c_2=w_{23},\quad c_3=w_{30},\qquad
 w_{rs}=\operatorname{Re}(Z_rZ_s^{-1})=Z_r\cdot Z_s .
\]
Put
\[
 C_y=\sum_{r=0}^3c_r,\quad
 T_y=w_{02}+w_{13},\quad
 \Xi_y=\sum_{r=0}^3(c_r^2-\tfrac14),\quad
 A_y=\sum_{r=0}^3c_rc_{r+1},
 \qquad c_4=c_0 .
 \tag{V75.2}\label{c18v75:local}
\]
The dot product is the real quaternion dot product.
Let \(\mathsf P_y\) denote conditional expectation under
the reference \(\mu_I\) onto these four paths, not under
\(\nu_t\). Their reference law is product Haar.

Here is a useful way to justify the conditional integrations.
Hold all prefix factors fixed and replace the four distinct
free spokes by \(Z_0,\ldots,Z_3\). These are separate Haar
translations or inversions. All other free links are still
integrated with their original Haar law. Averaging the
prefixes afterwards does not change the resulting identities.
This does not assert independence under \(\nu_t\), or
independence of the electric derivatives at different centres.

On a function of these four paths alone, the original
electric generator is exactly
\[
 L F(Z_0,\ldots,Z_3)=L_yF,\qquad
 L_y=\tfrac32\sum_{r=0}^3(-\Delta_{Z_r}).
 \tag{V75.3}\label{c18v75:localL}
\]
Indeed each path uses its own distinct prefix, with
kinetic weight \(1/2\), and its own free spoke, with weight
one. Left/right Haar multiplication transfers both round
Laplacians to that path. There is no shared factor between
two paths at this one centre. Consequently this subspace
is reducing for \(L\), and on smooth functions
\(\mathsf P_y L=L_y\mathsf P_y\).
For example this follows first on the product
Peter--Weyl core, then by the self-adjoint heat semigroup.
It does not diagonalize \(L\) jointly over all centres.

\begin{lemma}[The exterior integrations needed for the complete two-jet]
\label{c18v75:projection}
With \(q=n^3\),
\[
\begin{aligned}
 \mathsf P_y C&=C_y,&
 \mathsf P_y C^2&=C_y^2+3q-1,\\
 \mathsf P_y(UC)&=0,&
 \mathsf P_y S&=9q-\Xi_y+A_y-2T_y,\\
 \mathsf P_y\Gamma(W,C)
      &=\Gamma_y(C_y)+9(3q-1).
\end{aligned}
 \tag{V75.4}\label{c18v75:projections}
\]
Here \(\Gamma_y\) is the carré du champ of \(L_y\).
\end{lemma}
\begin{proof}
Every captured square outside \(y\) has two free spokes
owned by its own two-odd centre. A free factor occurring
once integrates to zero. Distinct captured squares have
distinct free-link pairs. Thus the only exterior terms
surviving in \(C^2\) are the diagonal squares, each of
mean \(1/4\). There are \(12q-4\) such squares.
Every product \(UC\) has a private free link incident to
a three-odd vertex, so it vanishes under this projection.
These arguments remain valid with the prefixes fixed.

For \(S\), use V57's normal-copy identity (V57.9).
The diagonal part is \(9q-\sum_p(w_p^2-1/4)\);
only the four terms at \(y\) retain their centered part.
A distinct-centre CC contraction has a private exterior
free factor and vanishes. An adjacent pair of squares
at \(y\) shares one prefix and its spoke, and uses opposite
halves of that prefix. If \(f=c_rc_{r+1}\) and \(f_0\)
is its shared-path Haar projection, its contribution to
\(S\) is
\[
 -2\Gamma_{\rm pair}(w_p,w_{p'})
          =-3f_0+f_2=f-4f_0,\qquad f_2=f-f_0 .
\]
The scalar/degree-two pair energies give the coefficients
\(3/2,-1/2\) for \(\Gamma_{\rm pair}\), as in V57.
The Haar identity
\(\int(z\cdot a)(z\cdot b)\,dz=(a\cdot b)/4\)
shows that \(4f_0\) is the opposite-spoke character.
Each opposite pair has two routes. Summing the four
adjacent motifs therefore gives \(A_y-2T_y\).
There is no UC normal term.

Finally the positive-Laplacian product rule and (V75.1)
give
\[
 \Gamma(W,C)=\tfrac12\{21WC-3C^2-L(WC)\}.
\]
Project, use \(WC=C^2+UC\), and commute the projection
with \(L\). Since \(L_yC_y=9C_y\), the result is
\(\Gamma_y(C_y)+9(3q-1)\).
This accounts for every exterior contribution, including
its constant part.
\end{proof}

\subsection{The complete native four-path marginal two-jet}

Let
\[
 p_{y,t}(Z)=\mathsf P_y(e^{u_t}),\qquad
 v_{y,t}=\log p_{y,t}.
\]
This is the exact normalized four-path marginal density.
The normalization covariance identity at zero reads
\[
 \partial_t v_{y,0}=\mathsf P_yb,\qquad
 \partial_t^2v_{y,0}
       =\mathsf P_yj+\mathsf P_yb^2-(\mathsf P_yb)^2 .
 \tag{V75.5}\label{c18v75:logmarg}
\]

\begin{theorem}[A full native local marginal, not a plaquette-only ansatz]
\label{c18v75:four}
For every centre and every \(n\ge5\),
\[
 \partial_t v_{y,0}=3C_y,\qquad
 J_y:=\partial_t^2v_{y,0}
       =-9-21\Xi_y+18T_y
       =12-21\sum_{r=0}^3c_r^2+18T_y .
 \tag{V75.6}\label{c18v75:fourjet}
\]
Accordingly, for every fixed derivative order \(k\),
\[
 v_{y,t}=3tC_y+\tfrac12t^2J_y+O_{n,C^k}(t^3).
 \tag{V75.7}\label{c18v75:fourTaylor}
\]
The remainder is not asserted uniform in \(n\).
\end{theorem}
\begin{proof}
The elementary sphere product rule gives
\[
\begin{aligned}
 L_y\Xi_y&=24\Xi_y,& L_yT_y&=9T_y,\\
 L_yA_y&=21A_y-6T_y,&
 \Gamma_y(C_y)&=9-3\Xi_y+6T_y-3A_y .
\end{aligned}
 \tag{V75.8}\label{c18v75:localalgebra}
\]
For example, on one adjacent product its degree-zero
shared-path part is \(w_{\rm opposite}/4\), of energy nine,
and its degree-two complement has energy twenty-one.
This yields \(21f-3w_{\rm opposite}\).
Summing gives the third identity. The diagonal terms of
\(\Gamma_y(C_y)\) are \(3(1-c_r^2)\); its four mixed
adjacent terms are \(3w_{\rm opposite}-3c_rc_{r+1}\).
This proves the last identity without a numerical fit.

Now substitute (V75.4) into (V75.1), and use (V75.8).
All \(A_y\) terms cancel, giving
\[
                   \mathsf P_yj=-27q-21\Xi_y+18T_y .
\]
The variance term in (V75.5) is exactly \(9(3q-1)\).
It is not zero merely because the exterior first moment
vanishes. Adding it gives (V75.6).
Smooth finite-dimensional flow, a smooth positive density
and integration on a compact fibre justify the fixed-regulator
Taylor expansion in each \(C^k\) norm.
As a normalization check, the Haar means are
\(\int J_y=-9\) and \(\int(3C_y)^2=9\).
Thus the second derivative of \(\int p_{y,t}\) is zero.
\end{proof}

\subsection{Exactly Haar single paths, but a regenerated physical source}

The based gauge group acts at \(y\) by simultaneous right
multiplication of all four \(Z_r\). The normal-chart law
is invariant under this action. Therefore each individual
\(Z_r\) has exactly normalized Haar law for every time in
the chart. Its marginal temporal score is identically zero.
This consequence of V33/V51's gauge invariance is recalled
to distinguish a gauge marginal from the following observable;
it is not counted as a new gauge-reduction theorem.

Retain two opposite paths, say \(Z_1,Z_3\). Their exact
joint density is invariant under simultaneous right
multiplication, so it has the form
\[
                  q_t(X),\qquad X=Z_1Z_3^{-1}.
 \tag{V75.9}\label{c18v75:relative}
\]
The change \((Z_1,Z_3)\leftrightarrow(X,Z_3)\) preserves
product Haar. Thus \(q_t\) is also the normalized density
of the internally gauge-invariant straight holonomy \(X\).
No coarse-root orbit has been quotiented.

\begin{theorem}[The first generated straight-path coefficient]
\label{c18v75:straight}
Let \(w(X)=\operatorname{Re}X\). At coarse identity,
\[
 \partial_t\log q_t\big|_0=0,\qquad
 \partial_t^2\log q_t\big|_0=27w(X).
 \tag{V75.10}\label{c18v75:straightjet}
\]
The marginal is even in \(t\). Consequently, at fixed regulator,
\[
\begin{aligned}
 q_t(X)&=1+\tfrac{27}{2}t^2w(X)+O_{n,C^k}(t^4),\\
 \log q_t(X)&=\tfrac{27}{2}t^2w(X)+O_{n,C^k}(t^4),\\
 a_t^X(X):=\partial_t\log q_t(X)
           &=27t\,w(X)+O_{n,C^k}(t^3),\\
 \nu_t(w(X))&=\tfrac{27}{8}t^2+O_n(t^4).
\end{aligned}
 \tag{V75.11}\label{c18v75:straightTaylor}
\]
Moreover \(a_t^X=\mathbb E_{\nu_t}[a_t\mid X]\) exactly.
\end{theorem}
\begin{proof}
Integrate the four-path density, not its logarithm, over
\(Z_0,Z_2\). Set \(\mathsf E_{02}\) for this product-Haar
integration with \(Z_1,Z_3\) fixed. Then
\[
\begin{aligned}
 \mathsf E_{02}(3C_y)&=0,\\
 \mathsf E_{02}J_y&=-9+18w_{13},\\
 \mathsf E_{02}(3C_y)^2&=9+9w_{13}.
\end{aligned}
 \tag{V75.12}\label{c18v75:twoterms}
\]
Indeed each corner square has mean \(1/4\),
\(\mathsf E_{02}w_{02}=0\), and
\[
 C_y=(Z_0+Z_2)\cdot(Z_1+Z_3),\qquad
 \mathsf E_{02}C_y^2
      =\tfrac12|Z_1+Z_3|^2=1+w_{13}.
\]
The second density derivative is
\(\mathsf E_{02}(J_y+9C_y^2)=27w_{13}\).
The first derivative is zero, so the logarithmic derivative
has the same second value. This proves (V75.10).

V55's central sign operation \(\sigma\) preserves the fibre
and obeys \(\widehat p_t\circ\sigma=\widehat p_{-t}\).
It changes the signs of adjacent-spoke characters and fixes
opposite-spoke holonomies as quaternions, not just their real
parts: the signs on opposite paths are equal. Hence
\(X\circ\sigma=X\) and \(q_t=q_{-t}\).
This gives the even remainders in (V75.11).
Haar's \(\int w=0,\ \int w^2=1/4\) gives the moment coefficient.
Finally differentiating the exact pushforward density under
the time-independent observable \(X\) gives
\(a_t^X=\mathbb E_{\nu_t}[a_t\mid X]\), with its full conditional
normalizer. In particular the displayed leading source has
a higher-order correction that enforces its exact mean zero;
it is not a finite-time replacement source.
\end{proof}

The sign is informative. The direct native second derivative
contributes \(18w_{13}\), while the conditional variance
contributes another \(9w_{13}\); their constants cancel.
Keeping only \(e^{3tC_y}\) would yield \(9w_{13}\) instead of
\(27w_{13}\). Gauge averaging does not remove this channel,
and the zero single-path score does not imply a zero
two-path score. This is a populated noncancellation test,
not a general assertion that signed improvements are impossible.

\subsection{The straight bank is visible to the actual inverse problem}

For every centre and each of its two opposite pairs, retain
the four normalized real coefficients
\(F_{\alpha,\mu}=2(X_\alpha)_\mu\), \(0\le\mu\le3\).
There are \(6n^3\) quaternion channels and \(24n^3\)
real functions. Let \(\mathcal V_{\rm st}\) be their real
linear span. These are V54's straight channels, not newly
postulated observables. Define their variational response
using the original moving-metric form:
\[
 \mathcal C_{\rm st}(t)
  :=\sup_{f\in\mathcal V_{\rm st}}
       \{2\nu_t(a_tf)-\mathcal E_t(f,f)\}.
 \tag{V75.13}\label{c18v75:bank}
\]
Constants in \(f\) do not affect this expression because
\(\nu_t(a_t)=0\). At every fixed regulator,
\[
 0\le\mathcal C_{\rm st}(t)
       \le\langle a_t,\widehat L_t^{-1}a_t\rangle_{\nu_t},
 \qquad
 \mathcal C_{\rm st}(t)
       =\tfrac{243}{2}n^3t^2+O_n(t^4).
 \tag{V75.14}\label{c18v75:bankresult}
\]

\begin{proof}
The inequality is the restriction of the centered inverse's
Dirichlet variational principle to a finite test space.
It does not assert that this space reduces the positive-time
generator. At zero, V54's distinct free-link signatures and
Haar moments give
\[
 \mu_I(F_{\alpha,\mu}F_{\beta,\nu})
       =\delta_{\alpha\beta}\delta_{\mu\nu},\qquad
 L F_{\alpha,\mu}=9F_{\alpha,\mu}.
\]
Thus the energy Gram matrix \(M(0)\) of these tests is \(9I\).
It is positive definite at every fixed \(t\): zero energy
would make their linear combination constant on the connected
fibre, which is excluded by their zero-time orthogonality
to constants and one another.

Let \(d_{\alpha,\mu}(t)=\nu_t(a_tF_{\alpha,\mu})\).
Differentiation of the expectation and (V75.11) give
\[
 d_{\alpha,0}(t)=\tfrac{27}{2}t+O_n(t^3),\qquad
 d_{\alpha,\mu}(t)=O_n(t^3)\quad(\mu>0).
\]
For the vector coordinates this follows from
\(\int X_\mu\,\operatorname{Re}X=0\).
Every test is fixed by \(\sigma\), whereas the source changes
sign. Hence \(d(t)\) is odd. The complete moving metric
also obeys V55's reflection identity, so \(M(t)\) is even
and \(M(t)=9I+O_n(t^2)\).
Completing the finite-dimensional square gives
\(\mathcal C_{\rm st}=d^{\mathsf T}M^{-1}d\).
There is one nonzero leading component per quaternion
channel. Therefore its quadratic coefficient is
\[
             (6n^3)\frac{(27/2)^2}{9}
                       =\frac{243}{2}n^3 ,
\]
as claimed. All expansions concern finite matrices at fixed
\(n\); no growing-dimensional inverse remainder is assumed.
\end{proof}

This is a lower bound on the full inverse cost, not its
missing upper bound. It is not an additional term to be
added to V59's full cost expansion: corner and straight
tests need not be energy-orthogonal at positive time.
The coefficient is also not a physical spectral gap.

\subsection{Verification, direction and limits}

The exact-arithmetic diagnostic
\begin{center}\small
\path{scripts/verify_ym_c18_v75_native_marginal_source.py}
\end{center}
applies the original weighted sphere Laplacian to the local
polynomials, integrates sphere monomials rationally, verifies
both normalization terms, and checks all four coordinates
of the relative quaternion. Periodic native incidence on
\(n=5,6\) checks the centre four-cycles, distinct prefixes,
mixed-motif exclusions and sign grading. V74's math
regression is replayed. The all-volume conclusions follow
from the geometric and Haar proofs above, not from sampling
those two tori. No finite-time YM simulation or proof-assistant
certification is claimed.

V54 already identified the straight channels; V57--V59
already computed the required unprojected density/source
jets and full temporal-cost coefficients. Those are not
recounted as new results. The new calculation is their
complete normalized four-path marginal, its two-path
density and score, and the actual straight-bank variational
response. The old gauge disintegration of V33 is not redone.

The next estimate must control the regenerated physical
path interactions and their source together. A cancellation
argument that deletes the first straight channel is now
excluded by the native coefficient \(27\). An estimate
that resums or bounds it remains possible. In particular,
one should seek a uniform connected-marginal remainder or
a signed response bound that retains these paths, rather
than claim closure from vanishing one-path marginals or
iterate V74's worsening scalar bounds.

The full time-four marginal source inverse, pointwise
Poisson Hessian and probability strain, interacting-vacuum
and all-coarse extensions, physical scale/clock, nontrivial
reflection-positive interacting continuum and positive
physical Yang--Mills mass gap remain unproved.
The goal remains active. The next companion and broad
literature checkpoints are both V80; archive/restart is V100.

Removing this append and restoring the title recovers V74
byte for byte. Only the immediate active predecessor is
replaced, and only \texttt{.tex} files are kept in
\path{latex_notes}. All builds and verification records
remain outside that folder.
% END CYCLE 18 VERSION 75 APPEND
% BEGIN CYCLE 18 VERSION 76 APPEND
\clearpage
\section[V76: Uniform additive straight-holonomy response]{V76: Uniform additive straight-holonomy response and a local marginal-energy solver}
\label{c18v76:context}

V75 showed that the first straight-path source survives the
complete native normalization. This appendix retains that
interaction, at every order in each straight holonomy, instead
of computing another Taylor coefficient. Its new input is a
geometric fact about the whole straight family: each channel
has two private free factors, whereas each prefix factor is
used by only four channels. This gives a populated, uniformly
conditioned additive response problem throughout the chart.

The result is not a full marginal or full-source inverse
estimate. \emph{Additive} means a sum of arbitrary functions of
one straight holonomy at a time, not an arbitrary function of
all holonomies jointly. The exact one-holonomy marginal sources
remain on the right-hand side. Controlling their aggregate load,
and then the nonadditive residual, is still necessary.
The setting is the actual electric normal-chart law
\(\nu_t=e^{u_t}\mu_I\), coarse identity, \(N=2n\), \(n\ge5\),
and \(0\le t\le4\). There is no expansion in \(t\) below.

\subsection{Private factors, shared prefixes and the actual marginals}

Let \(\mathcal A_n\) be the \(6n^3\) straight quaternion channels
of V54/V75. Each \(X_\alpha\) is the relative holonomy of two
opposite paths at a two-odd centre \(y_\alpha\).
It uses two free spokes and two distinct pair/prefix factors.
Different channels have disjoint free-spoke pairs, including
the two channels at the same centre.

For completeness, a free spoke joining a one-odd midpoint to
a two-odd centre has that unique centre as owner. At this centre
the two opposite pairs partition its four spokes. A prefix
factor has a unique one-odd midpoint; its four transverse
neighbouring two-odd centres each use that prefix in one
straight channel. Thus every prefix is used by exactly four
channels, and every channel uses exactly two prefixes.
These statements are about original independent pair/free
coordinates, not about independence under \(\nu_t\).

Join two distinct channels if they share a prefix. Write
\(\mathcal G_n\) for this graph and \(d_{\mathcal G}\) for its
graph distance, with value \(+\infty\) between components.
Every degree is at most six. Adjacent centres have periodic
fine-lattice Manhattan distance at most two, so
\[
 d_{\mathcal G}(\alpha,\beta)
       \ge \left\lceil d(y_\alpha,y_\beta)/2\right\rceil
 \quad\hbox{when the left side is finite}.
 \tag{V76.1}\label{c18v76:graph}
\]
No probabilistic separation between graph components is asserted.

For a smooth function \(h\) on the unit-round \(S^3=\mathrm{SU}(2)\),
varying any one of the four factors in \(h(X_\alpha)\), with
the others fixed, is an isometry of that sphere: a composition
of left/right multiplication and possibly inversion.
Consequently the squared derivative norm is
\(|\nabla h|^2(X_\alpha)\) for each free factor, and
\(\tfrac12|\nabla h|^2(X_\alpha)\) for each prefix in the
original product metric \(m_0\). Shared-prefix derivatives
are not declared orthogonal.

Let \(q_{\alpha,t}\,dX\) be the exact marginal law of \(X_\alpha\).
It is smooth and positive at every finite regulator.
The inherited one-factor oscillation bound
\(\operatorname{osc}_i u_t\le tS\), \(S=D^2K_*\), gives
\[
 \operatorname{osc}\log q_{\alpha,t}\le tS,\qquad
 e^{-tS}\le q_{\alpha,t}\le e^{tS}.
 \tag{V76.2}\label{c18v76:qbound}
\]
Indeed condition on all factors except one of this channel's
free spokes. Its normalized conditional density has pointwise
ratios at most \(e^{tS}\). The isometric change from that spoke
to \(X_\alpha\) preserves Haar measure. Average over the actual
law of the frozen factors: for any \(X,X'\), the same inequality
holds before each averaging, hence after averaging. This proves
the ratio bound for the marginal, not just for a pinned law.

Use the round marginal Dirichlet form
\[
 \mathcal D_{\alpha,t}(h,k)
       =\int \langle\nabla h,\nabla k\rangle q_{\alpha,t}\,dX,
 \qquad
 \mathcal L_{\alpha,t}
       =-\Delta-\langle\nabla\log q_{\alpha,t},\nabla\cdot\rangle .
 \tag{V76.3}\label{c18v76:localform}
\]
Its scalar gap is at least \(3e^{-tS}\).
To verify the constant, the Haar gap is three;
\(\operatorname{Var}_{q}h\le(\max q)\operatorname{Var}_{dX}h\)
and \(\mathcal D_q(h)\ge(\min q)\mathcal D_{dX}(h)\).
Their ratio is at most \(e^{tS}\).
This is a one-group marginal statement, not a gap for their joint law.

\subsection{A density-adapted additive energy space}

Fix \(t,n\). Let \(\mathcal H_{\alpha,t}\) be \(H^1(S^3)\)
modulo constants, with inner product \(\mathcal D_{\alpha,t}\),
using the representative of \(q_{\alpha,t}\)-mean zero.
The preceding marginal gap makes this a Hilbert space.
Set
\[
 \begin{gathered}
 \mathcal H_t=\bigoplus_{\alpha\in\mathcal A_n}\mathcal H_{\alpha,t},
 \qquad
 \|h\|_{\mathcal H_t}^2=\sum_\alpha\mathcal D_{\alpha,t}(h_\alpha),\\
 \mathsf T h=\sum_\alpha h_\alpha(X_\alpha).
 \end{gathered}
 \tag{V76.4}\label{c18v76:H}
\]
In particular \(\nu_t(\mathsf T h)=0\). Define
\[
 \mathcal B_t(h,k)
      =\int\langle\nabla_{m_0}\mathsf T h,
                       \nabla_{m_0}\mathsf T k\rangle\,d\nu_t,
 \qquad
 \mathcal A_t(h,k)=\mathcal E_t(\mathsf T h,\mathsf T k).
 \tag{V76.5}\label{c18v76:forms}
\]
The first uses the original product metric but the
\emph{actual} density; it is not the Haar form.

\begin{theorem}[Private-spoke coercivity at all harmonic orders]
\label{c18v76:coercivity}
For every \(h\in\mathcal H_t\),
\[
 \begin{gathered}
 2\|h\|_{\mathcal H_t}^2
       \le\mathcal B_t(h,h)\le6\|h\|_{\mathcal H_t}^2,\\
 2E_0^{-2}\|h\|_{\mathcal H_t}^2
       \le\mathcal A_t(h,h)\le6E_0^2\|h\|_{\mathcal H_t}^2.
 \end{gathered}
 \tag{V76.6}\label{c18v76:coercive}
\]
Here \(E_0\) is V61's geometric comparison constant through time
four. The constants two and six require no smallness of the
density, no Taylor truncation, and no bound on the number or
harmonic degree of the channel functions.
\end{theorem}
\begin{proof}
First take smooth \(h_\alpha\).
Each free spoke belongs to only one summand. The total
contribution of these private derivatives is exactly
\[
             2\sum_\alpha\mathcal D_{\alpha,t}(h_\alpha).
\]
All other product-metric terms are nonnegative, proving the
lower bound. For a prefix \(i\), let \(v_{i\alpha}\) be the
round-frame derivative of \(h_\alpha(X_\alpha)\).
At most four are nonzero, and its original metric weight
is \(1/2\). Pointwise Cauchy--Schwarz gives
\[
 \tfrac12\left|\sum_{\alpha\ni i}v_{i\alpha}\right|^2
             \le2\sum_{\alpha\ni i}|v_{i\alpha}|^2 .
\]
Each channel has two prefixes, each derivative with norm
\(|\nabla h_\alpha|(X_\alpha)\). Summing and averaging bounds
the prefix contribution by
\(4\sum_\alpha\mathcal D_{\alpha,t}(h_\alpha)\).
Together with the private part this is six.
No product-measure factorization has been used.
V61's comparison
\(E_0^{-2}m_0\preceq\widehat m_t\preceq E_0^2m_0\)
gives the same two-sided comparison for the inverse-metric
forms. Finally smooth functions are dense in each weighted
\(H^1\) space, and the bounds extend the maps and forms
continuously. The lower bound also proves that \(\mathsf T\)
is injective modulo constants.
\end{proof}

The adaptation to the actual one-channel energies matters.
In a fixed Haar norm an oscillation factor would appear.
In (V76.6), the normalized product-metric condition number
is at most three even if the marginal is strongly non-Haar.
This controls an additive subspace, not all scalar functions.

\subsection{A genuinely local solver in the channel index}

Let \(\mathsf B_t,\mathsf A_t\) be the bounded self-adjoint
operators on \(\mathcal H_t\) representing the two forms.
Then
\[
 2I\preceq\mathsf B_t\preceq6I,\qquad
 E_0^{-2}\mathsf B_t\preceq\mathsf A_t\preceq E_0^2\mathsf B_t .
 \tag{V76.7}\label{c18v76:operators}
\]
If distinct channels share no prefix, their product-metric
derivatives have disjoint supports, pointwise. Their block
of \(\mathsf B_t\) is therefore zero, regardless of any
probabilistic correlations. With \(\Pi_\alpha\) the
orthogonal channel projection, set
\[
 \mathsf R_t=I-\mathsf B_t/6,\qquad
 \mathsf P_K=\tfrac16\sum_{j=0}^K\mathsf R_t^j .
\]
Spectral calculus gives
\[
 \begin{gathered}
 0\preceq\mathsf R_t\preceq\tfrac23I,\qquad
 \mathsf B_t^{-1}=\tfrac16\sum_{j=0}^\infty\mathsf R_t^j,\\
 \|\Pi_\alpha\mathsf B_t^{-1}\Pi_\beta\|
       \le\tfrac12(2/3)^{d_{\mathcal G}(\alpha,\beta)}.
 \end{gathered}
 \tag{V76.8}\label{c18v76:inverse}
\]
The norm in the last expression is between the corresponding
marginal-energy Hilbert spaces, not their unweighted \(L^2\)
norms. When the graph distance is infinite the block is zero.

Here is the proof, including its infinite harmonic scope.
A product of \(j\) block-range-one operators has range at most
\(j\) in \(\mathcal G_n\). Hence all summands before
\(j=d_{\mathcal G}(\alpha,\beta)\) vanish in this block.
The remaining operator-norm sum is at most
\((1/6)\sum_{j\ge d}(2/3)^j=(1/2)(2/3)^d\).
This uses only bounded operators on the direct-sum Hilbert
space; it does not require finitely many harmonics.
For every \(r\in\mathcal H_t\), put
\[
 G=\langle r,\mathsf B_t^{-1}r\rangle,\qquad
 G_K=2\langle r,\mathsf P_Kr\rangle
              -\langle\mathsf P_Kr,\mathsf B_t\mathsf P_Kr\rangle .
\]
Then
\[
 \begin{aligned}
 \|r-\mathsf B_t\mathsf P_Kr\|_{\mathcal H_t}
       &\le(2/3)^{K+1}\|r\|_{\mathcal H_t},\\
 0\le G-G_K
       &\le\tfrac12(2/3)^{2K+2}\|r\|_{\mathcal H_t}^2 .
 \end{aligned}
 \tag{V76.9}\label{c18v76:truncation}
\]
The exact residual is \(\mathsf R_t^{K+1}r\).
Completing the square makes the exact error in the second
line
\(\langle\mathsf R_t^{K+1}r,
          \mathsf B_t^{-1}\mathsf R_t^{K+1}r\rangle\);
use \(\|\mathsf B_t^{-1}\|\le1/2\).
If \(r\) is supported on a set of channels,
\(\mathsf P_Kr\) is supported within graph distance \(K\).

The inverse-decay mechanism is classical and already occurs
in archive f16.2 V11. For primary context see
S.~Demko, W.~F.~Moss and P.~W.~Smith,
\href{https://doi.org/10.1090/S0025-5718-1984-0758197-9}
{\emph{Decay rates for inverses of band matrices}},
Math.\ Comp.\ 43 (1984), 491--499.
This targeted check confirmed the indexed primary abstract;
direct full-text retrieval was unavailable. No unverified
theorem is imported: (V76.8)--(V76.9) are proved here.
The new input is the native, all-time, all-harmonic interval
\([2,6]\), not a new inverse-decay theorem.

Locality here is precise but limited. It concerns a solver
on channel indices with the exact marginal forms and exact
blocks supplied. Neither \(q_{\alpha,t}\) nor a nonzero
block of \(\mathsf B_t\) is asserted to be determined by
a finite exterior collar. They are expectations under the
full law. Moreover the actual metric can couple distinct
factor supports: no finite-range or decay statement for
\(\mathsf A_t^{-1}\) follows just from its form comparison.

\subsection{The complete one-holonomy sources, with no normalization dropped}

Let \(a_t=\partial_tu_t\), and define
\[
 a_{\alpha,t}(X)
       =\mathbb E_{\nu_t}[a_t\mid X_\alpha=X]
       =\partial_t\log q_{\alpha,t}(X).
\]
These are the exact scores, not their first nonzero jets.
Let
\[
 r_\alpha=\mathcal L_{\alpha,t}^{-1}a_{\alpha,t},
 \qquad
 I_\alpha(t)=\mathcal D_{\alpha,t}(r_\alpha)
       =\langle a_{\alpha,t},
                  \mathcal L_{\alpha,t}^{-1}a_{\alpha,t}\rangle_{q_{\alpha,t}}.
 \tag{V76.10}\label{c18v76:loads}
\]
All inverses are centered one-group inverses. The one-group
gap already proved ensures their existence and gives
\[
 I_\alpha(t)\le\tfrac13 e^{tS}
                       \|a_{\alpha,t}\|_{L^2(q_{\alpha,t})}^2 .
 \tag{V76.11}\label{c18v76:singleinverse}
\]
For each fixed test \(h\), conditional expectation and
integration by parts give the exact source functional
\[
 \begin{gathered}
        \nu_t(a_t\mathsf T h)
             =\sum_\alpha\int a_{\alpha,t}h_\alpha q_{\alpha,t}\,dX
             =\langle r,h\rangle_{\mathcal H_t},\\
 \|r\|_{\mathcal H_t}^2=\sum_\alpha I_\alpha(t).
 \end{gathered}
 \tag{V76.12}\label{c18v76:source}
\]
No derivative of a time-dependent choice of centered
representatives is being taken in this identity.

\begin{theorem}[All-order additive response in the actual metric]
\label{c18v76:response}
Define
\[
 C_{\rm add}(t)
   =\sup_{h\in\mathcal H_t}
       \{2\nu_t(a_t\mathsf T h)-\mathcal E_t(\mathsf T h)\},
 \qquad
 G_{\rm add}(t)=\langle r,\mathsf B_t^{-1}r\rangle_{\mathcal H_t}.
\]
Then, at every declared \(t,n\),
\[
 \begin{gathered}
 C_{\rm add}(t)=\langle r,\mathsf A_t^{-1}r\rangle_{\mathcal H_t},
 \qquad
 E_0^{-2}G_{\rm add}\le C_{\rm add}\le E_0^2G_{\rm add},\\
 \frac{1}{6E_0^2}\sum_\alpha I_\alpha(t)
      \le C_{\rm add}(t)
      \le\frac{E_0^2}{2}\sum_\alpha I_\alpha(t),\\
 C_{\rm st}(t)\le C_{\rm add}(t)
      \le \langle a_t,\widehat L_t^{-1}a_t\rangle_{\nu_t}.
 \end{gathered}
 \tag{V76.13}\label{c18v76:responsebounds}
\]
For \(h_K=\mathsf P_Kr\), the actual trial
\(f_K=E_0^{-2}\mathsf T h_K\) satisfies
\[
 2\nu_t(a_tf_K)-\mathcal E_t(f_K)
 \ge E_0^{-2}\{G_{\rm add}
           -\tfrac12(2/3)^{2K+2}\sum_\alpha I_\alpha(t)\}.
 \tag{V76.14}\label{c18v76:trial}
\]
\end{theorem}
\begin{proof}
Coercivity and the Hilbert-space completion of the square
give the first equality. Inverting the positive operator
comparisons (V76.7) gives its two-sided comparison with
\(G_{\rm add}\). Since
\(\tfrac16I\preceq\mathsf B_t^{-1}\preceq\tfrac12I\),
(V76.12) gives the second line.
The first-harmonic straight bank is contained in the additive
space modulo constants, and the latter is a subspace of
the full Dirichlet space; restriction of the same variational
principle proves the last line.
Finally \(\mathcal E_t(\mathsf T h_K)\le E_0^2\mathcal B_t(h_K)\);
substitute the scaled trial and apply (V76.9).
\end{proof}

This is a finite-time upper bound for the \emph{additive}
response in terms of exact local source loads, not an upper
bound for the full inverse. A uniform gap on each one-group
marginal controls its inverse operator; it does not bound
its input \(a_{\alpha,t}\). That score is a conditional
expectation of the full extensive score, so long-range
response has not disappeared.

\subsection{A concrete finite-bank corollary}

For V75's four real tests \(F_{\alpha,\mu}=2(X_\alpha)_\mu\),
subtract their \(q_{\alpha,t}\)-means when selecting representatives.
Write the \(4\) by \(4\) marginal-energy block as
\[
 \begin{gathered}
 D_{\alpha,t}
       =4\left(I_4-\int XX^{\mathsf T}q_{\alpha,t}(X)\,dX\right),
 \qquad D_t=\bigoplus_\alpha D_{\alpha,t},\\
 3e^{-tS}I_4\preceq D_{\alpha,t}\preceq4I_4,
 \end{gathered}
 \tag{V76.15}\label{c18v76:finiteD}
\]
The formula follows from
\(\nabla X_\mu\cdot\nabla X_\nu=\delta_{\mu\nu}-X_\mu X_\nu\).
The lower bound integrates against \(q\ge e^{-tS}\) and uses
the Haar second moment \(I_4/4\); the upper bound is pointwise.
Let \(B_t,M_t\) be the energy Grams of these tests in \(m_0\)
and \(\widehat m_t\), both integrated against \(\nu_t\).
Restriction of (V76.6) gives
\[
 \begin{gathered}
 2D_t\preceq B_t\preceq6D_t,\qquad
 2E_0^{-2}D_t\preceq M_t\preceq6E_0^2D_t,\\
 6e^{-4S}E_0^{-2}I\preceq M_t\preceq24E_0^2 I .
 \end{gathered}
 \tag{V76.16}\label{c18v76:finitebound}
\]
Thus V75's growing energy matrix has a volume-uniform lower
bound at every time in the chart. The block-diagonally
normalized matrix \(D_t^{-1/2}B_tD_t^{-1/2}\) has spectrum
in \([2,6]\) and the same channel graph. Its inverse obeys
(V76.8), in these normalized coordinates.
At zero \(D_0=3I\) and \(B_0=M_0=9I\), consistently with V75.
This uniform conditioning does not make V75's Taylor
remainder uniform: time derivatives of the exact source
pairings and of the marginal blocks still need control.

\subsection{Verification, non-duplication and the next unpaid step}

The diagnostic
\begin{center}\small
\path{scripts/verify_ym_c18_v76_additive_straight_response.py}
\end{center}
checks native incidence on three periodic tori: private-spoke
ownership, the fourfold prefix load, the two prefixes per
channel and the channel graph. Exact noncommutative quaternion
algebra checks the private derivative identity and the shared
prefix energy accounting. An exact rational positive-matrix
fixture checks inverse decay, finite propagation and the
variational truncation identity; that fixture is not the
native finite-time law. V75's exact algebra is replayed.
The general bounds follow from the proofs above, not from
finite-size tests or proof-assistant certification.

V54 already supplied the channels and their Haar spectrum;
V75 supplied their native marginal jets and the first-harmonic
variational response. The new statements are all-time
conditioning on the whole additive, all-harmonic space,
its explicitly populated local product-metric solver, and
its comparison to the actual moving-metric response.
The generic finite-range inverse argument is inherited.

The next source estimate is now sharply separated from
kinetic conditioning: bound
\(\sum_\alpha I_\alpha(t)\) by a volume-extensive quantity
with a uniform coefficient for the actual chart, while
retaining V75's regenerated interaction. Even that would
control only \(C_{\rm add}\); mixed functions of several
holonomies and other source channels remain outside this
space. No assertion of statistical independence, global
mixing, a depth-stable full marginal response, or a full
Poisson-Hessian bound has been made.

The interacting-vacuum and all-coarse extensions, physical
scale and clock, nontrivial reflection-positive continuum
construction, and positive physical Yang--Mills mass gap
remain unproved. The previous turn was progress, and the
goal remains active. Companion and broad literature
checkpoints are next due at V80; archive/restart is V100.

Removing this append and restoring the title recovers V75
byte for byte. Only the immediate active filename is replaced.
Only \texttt{.tex} files remain in \path{latex_notes};
all diagnostics and build products stay outside that folder.
% END CYCLE 18 VERSION 76 APPEND
% BEGIN CYCLE 18 VERSION 77 APPEND
\clearpage
\section[V77: Native radial source flux and distant correlations]{V77: Native radial source flux, collective symmetry and the distant-correlation remainder}
\label{c18v77:context}

V76 paid the additive kinetic inverse, but left its exact
one-holonomy source loads unbounded through time four.
We now reduce that remaining native source question, rather
than impose another energy criterion. At coarse identity
all straight channels have the same central marginal.
Its entire temporal inverse load is an explicit weighted
norm of a single radial distribution flux. Every fixed-radius
part of the actual source has a uniform contribution to
that norm. The remaining term is a signed sum of distant
source--cap covariances, not a derivative tail that has
already been estimated.

The scope remains \(\beta=0\), \(g=I\), \(N=2n\), \(n\ge5\),
and \(0\le t\le4\). The radial reduction and the local-source
bound hold on this whole interval. A further bound on the
\emph{complete} radial flux is verified only on V61's
previously established short interval. Neither that interval
nor a full physical mass gap is extended here.

\subsection{The actual symmetry reduces all straight marginal sources}

Choose the straight channels with a consistent orientation.
A channel can be labeled by
\[
       (z;i,j),\qquad z\in(\mathbb Z/n\mathbb Z)^3,\quad i\ne j:
\]
its centre is \(2z+e_i+e_j\), its two midpoint paths have
prefix direction \(j\), and their displacement is \(2e_i\).
Translations by \(2z\) and permutations of the three axes
are transitive on these \(6n^3\) labels. Changing an
orientation replaces \(X\) by \(X^{-1}\) and does not
change \(s=\operatorname{Re}X\).

The Wilson action and ambient product metric are invariant
under these transformations. The block map intertwines them
with coarse coordinate permutations, fixing coarse identity.
Consequently its flowed differential and orthogonal normal
projector intertwine them as well. Uniqueness for the
normal-chart ODE makes \(\eta_t\) equivariant. Its ambient
Jacobian and the normalized law \(\nu_t\) are invariant.
The same argument applies to simultaneous conjugation of
every link by one \(h\in\mathrm{SU}(2)\); it acts on every
\(X_\alpha\) by \(X_\alpha\mapsto hX_\alpha h^{-1}\).
These are symmetries of the actual law, not of a
plaquette-only surrogate.

It follows that there is one function \(Q_{n,t}\) such that
\[
 q_{\alpha,t}(X)=Q_{n,t}(s),\quad s=\operatorname{Re}X,\qquad
 p_{n,t}(s)=\omega(s)Q_{n,t}(s),\quad
 \omega(s)=\frac2\pi\sqrt{1-s^2},\quad -1<s<1 .
 \tag{V77.1}\label{c18v77:radiallaw}
\]
Here \(q_{\alpha,t}\) is V76's normalized group marginal,
and \(p_{n,t}\) is the Lebesgue density of its real part.
Smooth conjugation-invariant functions on \(S^3\) are smooth
functions of \(s\) on the closed interval. Near either pole,
the smooth radial-function lemma, obtained by successive
Hadamard factorizations at the origin, writes the function
as a smooth function of the squared imaginary radius.
That radius is \(1-s^2\), with a fixed sign of \(s\) in the
chosen pole chart. Thus \(Q_{n,t}\) is smooth
and strictly positive, including the endpoints, at fixed
regulator. V76 gives
\[
 e^{-tS}\le Q_{n,t}\le e^{tS},\qquad
 \operatorname{osc}\log Q_{n,t}\le tS.
 \tag{V77.2}\label{c18v77:Qbounds}
\]
No uniform bound on \(\partial_t Q_{n,t}\) follows from this.

Write
\[
 A_{n,t}(s)=\partial_t\log Q_{n,t}(s),\qquad
 F_{n,t}(s)=\int_{-1}^s p_{n,t}(r)\,dr,\qquad
 J_{n,t}(s)=\partial_tF_{n,t}(s).
\]
The marginal source and the distribution flux are exactly
\[
 \begin{split}
 A_{n,t}(s)&=\mathbb E_{\nu_t}[a_t\mid
                            \operatorname{Re}X_\alpha=s],\\
 J_{n,t}(s)&=\nu_t\!\left(a_t\,{\bf1}_{\{\operatorname{Re}X_\alpha\le s\}}\right),
 \qquad J_{n,t}(-1)=J_{n,t}(1)=0 .
 \end{split}
 \tag{V77.3}\label{c18v77:nativeflux}
\]
The indicator is independent of time. Differentiation under
the compact integral is valid at fixed \(n\), and
\(\nu_t a_t=0\) retains the complete normalizer.
The flux is a covariance with the full native source.

\subsection{An exact radial inverse with its endpoint conditions}

Suppress \(n,t\) temporarily and put \(D(s)=1-s^2\).
On central functions the round group operator of V76 is
\[
 \mathcal L_Q h
    =-\frac1p\frac{d}{ds}\left(Dp\,h'\right)
    =-D h''+\left(3s-D(\log Q)'\right)h' .
 \tag{V77.4}\label{c18v77:radialoperator}
\]
Indeed the radial round Laplacian is \(Dh''-3sh'\),
and \((D\omega)'=-3s\omega\). The factor here is one,
not the four-link factor in archive f17 V8.

\begin{theorem}[The complete marginal load is its radial distribution flux]
\label{c18v77:fluxinverse}
Let \(r=\mathcal L_Q^{-1}A\), with \(Q\,dX\)-mean zero.
Then \(r\) is central and
\[
 \begin{gathered}
 r'(s)=-\frac{J(s)}{D(s)p(s)},\\
 I_n(t):=\langle A,\mathcal L_Q^{-1}A\rangle_{Q\,dX}
       =\int_{-1}^1\frac{|J_{n,t}(s)|^2}{D(s)p_{n,t}(s)}\,ds,
 \qquad
 \sum_{\alpha\in\mathcal A_n} I_\alpha(t)=6n^3 I_n(t).
 \end{gathered}
 \tag{V77.5}\label{c18v77:exactload}
\]
All formulas include the entire source, not only its first harmonic.
\end{theorem}
\begin{proof}
The weighted elliptic operator commutes with conjugation,
and its centered inverse is unique. Its solution for the
central source \(A\) is therefore central. Equivalently,
averaging any variational test over conjugation preserves
the source pairing and cannot increase its energy.

The equation is \(-(Dp\,r')'=\partial_t p\).
Integration from \(-1\) gives \(Dp\,r'=-J+c\).
At either endpoint \(Dp\) is a positive smooth factor
times \((1-s^2)^{3/2}\). A nonzero \(c\) would have
infinite energy \(\int |c-J|^2/(Dp)\,ds\), since
\(J(-1)=0\). Hence \(c=0\).
Smoothness of \(\partial_tQ\) and its zero total mass give
\[
 J(s)=O((1+s)^{3/2})\quad(s\downarrow-1),\qquad
 J(s)=O((1-s)^{3/2})\quad(s\uparrow1).
\]
These bounds make the displayed energy finite and the
boundary flux zero. More precisely, integrating the smooth
endpoint expansions shows that \(J/(Dp)\) is smooth there;
the resulting potential is the smooth group solution up to
its centering constant. Integration by parts yields
\(\langle A,r\rangle=\int D|r'|^2p\), proving the integral.
Finally all \(6n^3\) channels have the same marginal
and source by (V77.1)--(V77.3).
\end{proof}

There is also a spectral form suitable for moment estimates.
Let \(\chi_\ell(s)=U_\ell(s)\) be the \(\mathrm{SU}(2)\)
character, with \(U_0=1\), \(U_1=2s\) and
\(U_{\ell+1}=2sU_\ell-U_{\ell-1}\).
These form an orthonormal basis in \(L^2(\omega\,ds)\)
and have radial Haar energy \(\ell(\ell+2)\).
Indeed \(s=\cos\theta\) and
\(U_\ell(\cos\theta)\sin\theta=\sin((\ell+1)\theta)\)
reduce orthogonality and completeness to the sine basis
on \((0,\pi)\); differentiating this identity gives
\(-D U_\ell''+3sU_\ell'=\ell(\ell+2)U_\ell\).
For the actual moments \(m_\ell(t)=\nu_t(\chi_\ell(s_\alpha))\),
define
\[
 \mathcal J_n(t)
     =\frac{\pi}{2}\int_{-1}^1
                \frac{|J_{n,t}(s)|^2}{(1-s^2)^{3/2}}\,ds
     =\sum_{\ell\ge1}
          \frac{|\partial_t m_\ell(t)|^2}{\ell(\ell+2)} .
 \tag{V77.6}\label{c18v77:characters}
\]
For the equality, solve the same one-dimensional divergence
equation with Haar weight \(\omega\) and source
\(\partial_tQ\), whose Haar mean is zero. Its spectral
expansion and its flux integral are the two sides.
Smoothness at fixed \(n\) justifies the expansion and limit.
The density comparison gives
\[
             e^{-tS}\mathcal J_n(t)\le I_n(t)
                           \le e^{tS}\mathcal J_n(t).
 \tag{V77.7}\label{c18v77:fluxcomparison}
\]
The differentiated moments on the right are those of the
actual full law. No uniform truncation of this series is
asserted. Bounding density moments without bounding their
time derivatives would not establish (V77.7)'s needed input.

\subsection{Collective symmetry in the complete additive variational problem}

Symmetry also acts on V76's additive energy space by
permuting channels and conjugating their group arguments.
The source functional and both quadratic forms are invariant.
Averaging a test over this compact symmetry group preserves
its source pairing and decreases its quadratic energy.
The transitivity and centrality proved above imply that the
averaged test has the form
\[
                  f_h=\sum_{\alpha\in\mathcal A_n}h(s_\alpha)
 \tag{V77.8}\label{c18v77:collective}
\]
for a single real function \(h\), centered under \(p_{n,t}\).
Thus the suprema defining both \(C_{\rm add}\) and
\(G_{\rm add}\) can be taken over these collective radial
tests. This is not a restriction to a finite character cutoff.

For a more explicit fixed-product-metric description, write
\(\mathcal D_p(h,k)=\int D h'k'p\,ds\) and \(N_c=6n^3\).
If distinct channels \(\alpha,\beta\) share prefix \(i\), put
\[
 \kappa_{\alpha\beta}(x)
     =\left\langle\nabla_i^{\rm round}s_\alpha,
                           \nabla_i^{\rm round}s_\beta\right\rangle .
\]
They share at most one prefix for \(n\ge5\). Then the
product-metric energy per channel is exactly
\[
 \begin{split}
 \mathfrak b_{n,t}(h,k)
 &=3\mathcal D_p(h,k)
   +\frac{1}{2N_c}\sum_{\{\alpha,\beta\}\in E(\mathcal G_n)}
       \nu_t\!\left(\kappa_{\alpha\beta}
        [h'(s_\alpha)k'(s_\beta)+h'(s_\beta)k'(s_\alpha)]\right),\\
 2\mathcal D_p(h)&\le\mathfrak b_{n,t}(h,h)
                                      \le6\mathcal D_p(h).
 \end{split}
 \tag{V77.9}\label{c18v77:collectiveform}
\]
The diagonal coefficient is two free derivatives plus two
prefix derivatives of weight \(1/2\), hence three.
Each unordered shared-prefix pair contributes its two
cross terms, each of weight \(1/2\). This proves the
identity directly; V76 proves the inequalities.
The expectations in the kernel remain full-law expectations,
not products of one-channel marginals.

In particular,
\[
 \begin{gathered}
 \frac{G_{\rm add}(t)}{N_c}
      =\sup_h\left\{2\int A_{n,t}h\,p_{n,t}\,ds
                                      -\mathfrak b_{n,t}(h,h)\right\},\\
 \frac{n^3}{E_0^2}I_n(t)\le C_{\rm add}(t)
                               \le3E_0^2 n^3 I_n(t).
 \end{gathered}
 \tag{V77.10}\label{c18v77:collectiveinverse}
\]
The second line is V76 with its now identified common source
load. Therefore a uniform volume-extensive additive response
is equivalent, up to these fixed geometric constants, to a
uniform bound on this one scalar flux norm. It is still not
equivalent to a full Yang--Mills mass gap.

\subsection{Every fixed-radius native source contribution is bounded}

Use V74's actual gauge-averaged source representation,
with its fixed ordering and reference configuration:
\[
 a_t=\sum_{i\in I_0}(\gamma_{i,t}-\nu_t\gamma_{i,t}),\qquad
 |\gamma_{i,t}|\le S,\qquad
 \|\nabla_j\gamma_{i,t}\|_\infty
                         \le DK_\zeta e^{-\zeta d(i,j)} .
 \tag{V77.11}\label{c18v77:sourcepieces}
\]
The summands have derivative tails, not strict finite support.
No new source decomposition is claimed.

Fix a canonical channel \(\alpha\), and an integer \(R\ge2\).
Let \(B_R\) contain source labels whose anchors have distance
at most \(R\) from \(y_\alpha\), and set
\[
 M_R=3(2R+1)^3,\qquad
 g_R=\sum_{i\in B_R}\gamma_{i,t},\qquad
 b_R(s)=\mathbb E_{\nu_t}[g_R-\nu_tg_R\mid s_\alpha=s],
 \qquad b_R^{\rm far}=A_{n,t}-b_R .
\]
There are at most three original factor labels per anchor,
so \(|B_R|\le M_R\), also on a wrapped torus. Define the
two radial source loads using the same exact marginal \(p\):
\[
 I_R^{\rm near}=\langle b_R,\mathcal L_Q^{-1}b_R\rangle_p,
 \qquad
 I_R^{\rm far}=\langle b_R^{\rm far},
                                \mathcal L_Q^{-1}b_R^{\rm far}\rangle_p .
\]

\begin{theorem}[A paid local contribution and an exact distant remainder]
\label{c18v77:nearfar}
For every \(n\ge5\) and \(0\le t\le4\),
\[
 \begin{gathered}
 I_R^{\rm near}\le\frac{e^{tS}}{3}M_R^2S^2,\qquad
 \left|\sqrt{I_n(t)}-\sqrt{I_R^{\rm far}}\right|
                                  \le\frac{e^{tS/2}}{\sqrt3}M_RS,\\
 I_R^{\rm far}
     =\int_{-1}^1\frac{|J_R^{\rm far}(s)|^2}{D(s)p_{n,t}(s)}\,ds,\\
 J_R^{\rm far}(s)
     =\sum_{i\notin B_R}
        \operatorname{Cov}_{\nu_t}
            \left(\gamma_{i,t},{\bf1}_{\{s_\alpha\le s\}}\right).
 \end{gathered}
 \tag{V77.12}\label{c18v77:nearfarbound}
\]
For any fixed \(R\), uniformly bounding \(I_n(t)\) over
\(n,t\) is equivalent to uniformly bounding this distant
signed covariance norm.
\end{theorem}
\begin{proof}
The range of \(g_R\) has length at most \(2M_RS\).
Any real random variable with range length \(L\) has variance
at most \(L^2/4\): if its endpoints are \(a,b\), the inequality
\((Y-a)(b-Y)\ge0\) yields
\(\operatorname{Var}Y\le(\mathbb EY-a)(b-\mathbb EY)\le L^2/4\).
Conditional expectation contracts centered \(L^2\) norm,
so \(\|b_R\|_{L^2(p)}^2\le M_R^2S^2\).
Use V76's round marginal gap \(3e^{-tS}\).

The inverse form defines a Hilbert norm on centered sources.
Its triangle and reverse-triangle inequalities, applied to
\(A=b_R+b_R^{\rm far}\), give the first line.
The flux proof of (V77.5) applies to any smooth centered
radial source \(b\), with flux \(\int_{-1}^s bp\).
For \(b=b_R^{\rm far}\), conditional expectation and
(V77.11) identify this flux with the last line of (V77.12).
All source labels outside \(B_R\) remain in the sum.
Finally \(e^{tS/2}M_RS/\sqrt3\le e^{2S}M_RS/\sqrt3\)
is independent of \(n,t\) for fixed \(R\), proving the
equivalence.
\end{proof}

The derivative tail in (V77.11) is not a bound on the
unconditioned covariances in (V77.12). Marginalization
can transmit a distant response through the retained law.
We have not replaced this signed sum by a sum of absolute
values, assumed independence, or removed the conditional
normalizers. The new bound pays its fixed-radius part;
further calculations confined to that part cannot settle
the remaining uniform question.

\subsection{A complete flux bound on the already-paid short interval}

There is a nonempty interval on which the whole flux is
already controlled by a native input, rather than left
as an unevaluated condition. Recall V61's exact probability
velocity \(V_t\), with
\[
 \partial_t\nu_t+\operatorname{div}_{m_0}(\nu_t V_t)=0,\qquad
 \sup_i|V_{t,i}|_{m_0}\le P_*,
 \qquad
 0\le t\le\tau:=\min\{4,(256K_*)^{-1}\}.
\]
Put \(K_{\rm path}=(2+\sqrt2)P_*\). Each straight map has
two free factors of speed at most \(P_*\) and two prefix
factors of round speed at most \(P_*/\sqrt2\).
The triangle inequality in its target tangent space gives
\[
                |dX_\alpha(V_t)|\le K_{\rm path}.
\]
Push the continuity equation to \(s_\alpha=\operatorname{Re}X_\alpha\).
Its scalar velocity is
\(v_t(s)=\mathbb E[ds_\alpha(V_t)\mid s_\alpha=s]\), with
\[
 |v_t(s)|\le K_{\rm path}\sqrt{1-s^2},\qquad
 \partial_t p+\partial_s(pv_t)=0,\qquad
 J_{n,t}(s)=-p_{n,t}(s)v_t(s).
\]
These identities first follow against smooth scalar tests.
The bound makes \(pv_t\) tend to zero at both endpoints,
so integrating determines the last identity with no
boundary constant. Thus the populated flux estimate is
\[
 \begin{gathered}
 |J_{n,t}(s)|\le K_{\rm path}p_{n,t}(s)\sqrt{1-s^2},\\
 I_n(t)\le K_{\rm path}^2,\qquad
 C_{\rm add}(t)\le3E_0^2K_{\rm path}^2 n^3
                  \quad(0\le t\le\tau).
 \end{gathered}
 \tag{V77.13}\label{c18v77:shorttime}
\]
This uses V61's existing full-source transport result; it
is not a new continuation theorem beyond \(\tau\).
For the distant term it also gives
\(\sqrt{I_R^{\rm far}}\le K_{\rm path}
                         +e^{tS/2}M_RS/\sqrt3\) on that interval.
No such bound has been proved here for \(\tau<t\le4\).

As a normalization regression, V75's already computed jet
gives, at fixed \(n\),
\[
 \begin{gathered}
 \partial_tQ_{n,t}(s)=27ts+O_n(t^3),\qquad
 J_{n,t}(s)=-\frac{18}{\pi}t(1-s^2)^{3/2}+O_n(t^3)
                        \ \hbox{at each fixed }s,\\
 r'(s)=9t+O_n(t^3),\qquad
 I_n(t)=\frac{243}{4}t^2+O_n(t^4).
 \end{gathered}
 \tag{V77.14}\label{c18v77:regression}
\]
The endpoint-weighted remainder needed in the last identity
follows from the smooth fixed-regulator Taylor expansion
and zero total mass, as in the proof of (V77.5).
This reproduces V75's straight response coefficient after
the zero-time factor three in the additive energy.
It is a check, not a new source estimate or a uniform
Taylor remainder.

\subsection{Literature, verification and direction}

Radial changes of variables are not new to this programme.
Archive f4 V71 already treats an exact one-link Wilson
quantile map; f17 V8 treats a radial score for the full
Wilson history law. Their densities and sources are not
the present normal-chart temporal marginal. No exponential
Wilson ansatz or theorem about its radial transport is
substituted for \(Q_{n,t}\).

For primary context the abstract of R.~Peyre,
\href{https://arxiv.org/abs/1104.4631}
{\emph{Comparison between \(W_2\) distance and
\(\dot H^{-1}\) norm, and localisation of Wasserstein distance}},
was checked. It concerns weighted negative Sobolev and
transport comparisons. We use only that perspective:
the present compact radial flux identity and all constants
are derived above. No endpoint transport distance is used
as a bound on the squared speed along this prescribed
time-dependent law.

The diagnostic
\begin{center}\small
\path{scripts/verify_ym_c18_v77_radial_source_flux.py}
\end{center}
checks transitivity of the native channel labels, the
radial operator, character normalizations and flux signs
with exact arithmetic. A nonconstant positive one-group
comparison density tests the weighted Poisson formula;
it is not identified with the native finite-time density.
Native incidence and V75's source algebra are replayed.
These tests supplement the proofs and are not a
proof-assistant certificate.

The new reduction isolates a precise upstream task:
control the distant source--cap covariance norm in
(V77.12), or equivalently the weighted derivative-moment
sum (V77.6), for the actual law after the short-time interval.
The near-source estimate is now paid for every fixed \(R\).
A bound on \(Q\) alone, additional local Taylor coefficients,
or another inversion of V76's conditioned additive matrix
does not settle this target. Even a complete additive
source bound would leave nonadditive and other-channel
parts of the full inverse problem.

The interacting-vacuum extension, all-coarse control,
full probability strain, physical scaling and clock,
nontrivial reflection-positive continuum construction,
and positive physical Yang--Mills mass gap remain unproved.
The previous turn was progress and the goal remains active.
The next companion and broad literature checkpoints are
V80; archive/restart remains V100.

Removing this append and restoring the title recovers V76
byte for byte. Only its active filename is replaced.
Only \texttt{.tex} files are stored in \path{latex_notes};
checks and build products are outside that folder.
% END CYCLE 18 VERSION 77 APPEND
% BEGIN CYCLE 18 VERSION 78 APPEND
\clearpage
\section[V78: A finite-range obstruction to source-flux closure]{V78: A finite-range obstruction to source-flux closure in the straight-holonomy coordinates}
\label{c18v78:context}

V77 isolates an actual distant source--cap covariance norm.
Before trying to bound it using only the estimates already paid,
we test whether those estimates can imply the desired conclusion.
The answer is negative even with strictly finite-range interactions,
not merely V62's long-range mean-field comparison. We embed a
classical ordered spin law into the same straight-holonomy
coordinates. Its one-channel marginal is exactly Haar at one
time, while its \emph{actual temporal marginal score} has inverse
load of order \(n^6\).

This is an obstruction to a class of proposed arguments, not
a counterexample to the Wilson normal-chart law. The comparison
starts from the same product Haar law and even has the same
initial score \(3W_{\rm cap}\), but its later density is not
the normal Jacobian in V51. No native estimate beyond V61's
interval, nor a physical mass gap, is proved in this append.

\subsection{A simultaneous Haar coordinate change}

Put \(V=n^3\), \(n\ge5\), and write a straight channel as
\(\alpha=(z,\kappa)\), where
\(z\in(\mathbb Z/n\mathbb Z)^3\) and
\(\kappa=(i,j)\), \(i\ne j\), ranges over six types.
Use precisely the group variables \(X_{z,\kappa}\) of V76--V77.
Each is the ratio of two coarse-root-to-centre midpoint paths.
Its two free spokes are private to that channel; its two
prefix factors may be shared.

Conditional on every prefix, replace the first free spoke in
each pair by \(X_{z,\kappa}\), retaining the other free spoke.
This is a separate left/right Haar translation or inversion
in every replaced variable. The map is smoothly invertible,
and all such replacements can be made simultaneously because
the free pairs are disjoint. Hence
\[
 \mu_I=\left(\prod_{z,\kappa}dX_{z,\kappa}\right)d\mu_{\rm sp},
 \qquad X_{z,\kappa}\in\mathrm{SU}(2)\simeq S^3 ,
 \tag{V78.1}\label{c18v78:product}
\]
with \(15V\) spectator factors in \(\mu_{\rm sp}\).
This is a statement about probability coordinates, not a
product isometry for the original kinetic metric \(m_0\).

Pinned gauge transformations act trivially on \(X_{z,\kappa}\):
both paths acquire the same right multiplier at their common
centre and have fixed coarse-root gauge values. Simultaneous
global conjugation conjugates every \(X\). Thus functions of
their real parts and quaternion dot products below are
pinned-gauge invariant and globally conjugation invariant.

\subsection{An explicit smooth comparison path, with its real source}

For each type define nearest-neighbour edges on the coarse torus,
counting each undirected edge once. Set
\[
 C=W_{\rm cap},\qquad
 H=\sum_\kappa\sum_z\sum_{r=1}^3
           X_{z,\kappa}\cdot X_{z+e_r,\kappa},\qquad
 M=\sum_{\kappa,z}s_{z,\kappa},\quad
 s_{z,\kappa}=\operatorname{Re}X_{z,\kappa}.
\]
Here \(X\cdot Y=\operatorname{Re}(XY^{-1})\). Define
\[
 \begin{gathered}
 \alpha(t)=3t(1-t^2/9)^2,\qquad
 j(t)=8t^2/9,\qquad
 h(t)=t^2(t^2-9)/54,\\
 U_t=\alpha(t)C+j(t)H+h(t)M,\qquad
 d\widetilde\nu_t=Z_t^{-1}e^{U_t}\,d\mu_I,\qquad
 \widetilde a_t=\dot U_t-\widetilde\nu_t(\dot U_t).
 \end{gathered}
 \tag{V78.2}\label{c18v78:path}
\]
The tilde distinguishes this law from the native \(\nu_t\).
It is positive and smooth on \(0\le t\le4\), and
\[
 \widetilde\nu_0=\mu_I,\qquad \widetilde a_0=3C,\qquad
 \alpha(3)=\alpha'(3)=h(3)=0,\quad j(3)=8,\quad h'(3)=1 .
 \tag{V78.3}\label{c18v78:times}
\]
The equality of initial scores uses \(\mu_I C=0\).
We do not identify higher jets with V75's native jets.

At \(t=3\), (V78.1) shows that the \(X\)'s form six
independent nearest-neighbour \(O(4)\) spin systems, each with
density proportional to
\[
       \exp\left(8\sum_{\langle z,w\rangle}X_z\cdot X_w\right)
                         \prod_zdX_z .
 \tag{V78.4}\label{c18v78:spinlaw}
\]
The spectators remain independent Haar variables in these
coordinates. The complete temporal score is
\[
 \widetilde a_3=\frac{16}{3}
                      (H-\widetilde\nu_3H)+M .
 \tag{V78.5}\label{c18v78:fullscore}
\]
In particular the interaction derivative has not been
silently deleted from the source.

\subsection{The comparison has the spatial estimates at issue}

The terms in \(U_t\), and also those in \(\dot U_t\), depend
on at most eight original pair/free factors. Their anchor
diameter is at most six in fine-lattice distance. Every
original factor occurs in at most eight \(C\)-terms,
twenty-four \(H\)-terms and four \(M\)-terms:
a prefix occurs in four channels, each channel has six
neighbours, and each fine captured half-link meets at most
four plaquettes. Private free factors have smaller loads.

For an explicit ordered-derivative bound let
\[
 \begin{split}
 A(t)&=\max\{|\alpha(t)|,|j(t)|,|h(t)|\},\\
 A_1(t)&=\max\{|\alpha'(t)|,|j'(t)|,|h'(t)|\},\\
 B_r(\zeta)&=36e^{6\zeta}8^r24^{r-1}\qquad(r\ge1).
 \end{split}
\]
In the original invariant orthonormal pair/free frames,
\[
 \mathfrak a_\zeta(\partial^r\log\widetilde p_t)
       \le B_r(\zeta)A(t),\qquad
 \mathfrak a_\zeta(\partial^r\widetilde a_t)
       \le B_r(\zeta)A_1(t).
 \tag{V78.6}\label{c18v78:jets}
\]
Indeed a trace word has at most eight group occurrences;
each ordered derivative inserts a unit generator in one
of those occurrences, giving the bound \(8^r\).
After fixing one scalar derivative slot there are at most
\(24^{r-1}\) choices for the remaining factor/color slots
of one term. At most thirty-six terms meet the fixed slot.
Their spatial diameter contributes at most \(e^{6\zeta}\).
Pair-frame scaling only reduces insertion norms.
Normalizing constants have no spatial derivatives.
The same counting works with any slot fixed, so it is an
all-slot estimate with the configuration supremum outside
the sums. Covariant Hessian rows are bounded as well:
the connection correction gives, for either scalar,
at most three times its ordered second-derivative bound
plus six times its ordered first-derivative bound.
All constants are independent of \(n\).

Changing one original factor changes the potential by an
oscillation of at most
\[
 B_{\rm fac}(t)=16|\alpha(t)|+48j(t)+8|h(t)| .
\]
Consequently every one-factor conditional product-metric gap
is at least \((3/2)e^{-B_{\rm fac}(t)}\), by the elementary
Haar density comparison already used in V62.
For a private free spoke the sharper bound
\[
 B_{\rm ch}(t)=4|\alpha(t)|+12j(t)+2|h(t)|
\]
applies: that spoke meets two captured plaquettes and one
channel. Conditional on everything else, replacing it by
\(X_\alpha\) preserves Haar measure. Integrating the normalized
conditional densities therefore proves, for every channel,
\[
 e^{-B_{\rm ch}(t)}\le\widetilde q_{\alpha,t}(X)
                  \le e^{B_{\rm ch}(t)},\qquad
 \operatorname{osc}\log\widetilde q_{\alpha,t}
                  \le B_{\rm ch}(t).
 \tag{V78.7}\label{c18v78:marginalbounds}
\]
Translations and positive axis permutations preserve the
path and are transitive on the channel labels. Conjugation
makes these common marginals central. Thus locality,
the relevant symmetries, bounded conditional oscillations,
and all fixed-order weighted spatial score jets hold
through time four. None is a claim about the native
normal Jacobian.

\subsection{A finite-volume infrared estimate and its normalization}

We use the established finite-periodic Gaussian-domination
theorem of Fr\"ohlich--Simon--Spencer,
\href{https://math.caltech.edu/SimonPapers/65.pdf}
{\emph{Infrared Bounds, Phase Transitions and Continuous
Symmetry Breaking}}, Theorem 2.3 and its proof in Section 2.
For the law proportional to \(e^{j\sum_{\langle z,w\rangle}
X_z\cdot X_w}\prod_zdX_z\), \(X_z\in S^3\), it gives
\[
 \mathbb E|\widehat X(k)|^2
       \le\frac4{j\lambda(k)}\quad(k\ne0),\qquad
 \widehat X(k)=V^{-1/2}\sum_z e^{-2\pi i k\cdot z/n}X_z,\quad
 \lambda(k)=4\sum_{r=1}^3\sin^2(\pi k_r/n).
 \tag{V78.8}\label{c18v78:IR}
\]
This is a cited theorem, not a newly proved Gaussian
domination result. Its hypotheses here are a cubic periodic
box, positive nearest-neighbour coupling, and a common
single-spin measure with all quadratic exponential moments;
normalized surface measure on \(S^3\) satisfies them.
Apply the theorem to real sine and cosine modes in each
of four components to obtain the displayed normalization.
No infinite-volume phase theorem is needed below.

Here is an elementary uniform bound for the resulting
finite lattice sum. Choose signed representatives with
\(|k_r|\le n/2\), and group nonzero points by
\(q=\max_r|k_r|\). Concavity of sine on \([0,\pi/2]\)
gives \(\lambda(k)\ge16q^2/n^2\).
A shell has at most
\((2q+1)^3-(2q-1)^3=24q^2+2\le26q^2\) points.
Hence
\[
 \frac1V\sum_{k\ne0}\lambda(k)^{-1}
        \le\frac{26}{16n}\lfloor n/2\rfloor\le\frac{13}{16}.
 \tag{V78.9}\label{c18v78:latticesum}
\]
Parseval and \(|X_z|=1\) now imply, for \(T=\sum_zX_z\),
\[
 \frac{\mathbb E|T|^2}{V^2}
     \ge1-\frac{13}{4j}.
 \qquad\text{In particular at }j=8,\quad
 \frac{\mathbb E|T|^2}{V^2}\ge\frac{19}{32}>\frac12 .
 \tag{V78.10}\label{c18v78:order}
\]

\subsection{Haar marginals with a volume-growing exact radial source}

Fix a channel type and its origin, suppressing the type.
At time three, global \(O(4)\) rotations in this one copy
imply that \(X_0\) is exactly uniform on \(S^3\).
They also imply
\[
 \mathbb E[X_z\mid X_0=x]=c_zx,\qquad
 c_z=\mathbb E(X_0\cdot X_z),\qquad |c_z|\le1.
\]
To prove the first identity, the conditional vector must
be fixed by the \(O(3)\) stabilizer of \(x\); equivariance
under the transitive \(O(4)\) action makes its scalar
coefficient independent of \(x\). Taking the dot product
with \(x\) and integrating identifies \(c_z\).
Translation invariance gives
\[
 b_n:=\sum_zc_z=\frac{\mathbb E|T|^2}{V}\ge V/2.
 \tag{V78.11}\label{c18v78:b}
\]

Every bond energy is \(O(4)\)-invariant. Its conditional
expectation given \(X_0=x\) is therefore constant in \(x\),
equal to its unconditional expectation. The other five
copies are independent. Thus the centered \(H\)-term in
(V78.5) has zero one-channel projection, while
the \(M\)-term projects to \(b_n\operatorname{Re}x\).
This proves the following exact statement.
\begin{theorem}[The source-flux bound is not forced by the paid local structures]
\label{c18v78:obstruction}
For the comparison path (V78.2), all channels at \(t=3\)
have
\[
 \widetilde Q_{n,3}(s)=1,\qquad
 \widetilde A_{n,3}(s)=b_ns,\qquad
 \widetilde J_{n,3}(s)
       =-\frac{2b_n}{3\pi}(1-s^2)^{3/2},
\]
and their exact round marginal inverse loads obey
\[
 \widetilde I_n(3)=\frac{b_n^2}{12}\ge\frac{V^2}{48}.
 \tag{V78.12}\label{c18v78:load}
\]
\end{theorem}
\begin{proof}
Differentiating the normalized marginal gives the conditional
expectation of the complete score (V78.5), evaluated above.
With \(\omega=(2/\pi)\sqrt{1-s^2}\), direct integration gives
\(\int_{-1}^s r\omega(r)\,dr
=-(2/3\pi)(1-s^2)^{3/2}\).
The round Haar generator sends \(s\) to \(3s\), and
\(\int s^2\omega(s)\,ds=1/4\).
Therefore its centered inverse sends \(b_ns\) to \(b_ns/3\),
with energy \(b_n^2/12\).
This also follows from V77's flux formula.
\end{proof}

The distant term itself can be evaluated without discarding
its signs. At time three use the exact local source pieces
\[
 \gamma_{z,\kappa}
    =\frac{16}{3}\sum_{r=1}^3
          X_{z,\kappa}\cdot X_{z+e_r,\kappa}
                        +s_{z,\kappa},\qquad |\gamma_{z,\kappa}|\le17.
\]
Their centered sum is \(\widetilde a_3\).
For a fixed-radius set \(B_R\) of sites in the observed copy,
remove the terms with these labels; other copies have
zero cap covariance. Every individual bond term also has
zero cap covariance by the same \(O(4)\) argument. Hence
\[
 \begin{split}
 \widetilde J_R^{\rm far}(s)
   &=-\frac{2}{3\pi}(1-s^2)^{3/2}
                       \left(b_n-\sum_{z\in B_R}c_z\right),\\
 \widetilde I_R^{\rm far}
   &=\frac1{12}\left(b_n-\sum_{z\in B_R}c_z\right)^2
     \ge\frac1{12}\bigl(V/2-|B_R|\bigr)_+^2 .
 \end{split}
 \tag{V78.13}\label{c18v78:far}
\]
For example a coarse maximum-distance ball has
\(|B_R|\le(2R+1)^3\), independently of \(n\).
The fixed-radius local bound is fully compatible with
this divergent distant load. No positivity assumption
on the individual \(c_z\) was used.

V76's additive kinetic inequalities are geometric
pointwise inequalities, valid also under this comparison
law with the original \(m_0\). Their product-metric
variational response therefore satisfies
\[
             \widetilde G_{\rm add}(3)
                   \ge V\widetilde I_n(3)\ge V^3/48 .
 \tag{V78.14}\label{c18v78:additive}
\]
Thus this is not merely an artifact of a badly conditioned
additive kinetic matrix or of a rough indicator function.

\subsection{What this rules out, and the next native obligation}

The classical infrared bound is an imported theorem.
The new programme-specific point is its simultaneous
embedding in the actual straight-channel factor incidence
and the explicit temporal-source/cap calculation.
V62's earlier obstruction lacked uniform positive spatial
weights. This one has finite range. Archive f16.2 V61's
\(O(4)\) alignment comparison concerns different Wilson
ribbon integrals and did not supply this source obstruction.

The comparison deliberately has no claim to satisfy the
native identity
\[
 k_t=\operatorname{Jac}\eta_t,\qquad
 \dot\eta_t=-Q_t\nabla W\circ\eta_t
                      \quad(\beta=0)
\]
or V75's later native jets. Consequently it does not prove
that the Wilson path orders, or that its V77 load diverges.
It proves that volume-independent local jets, gauge symmetry,
conditional gaps, Haar-comparable marginals, and even the
correct initial score cannot alone rule that phenomenon out.

Further work on V77 must use a constraint specific to the
Wilson Jacobian/source relation, or a native quantitative
phase-exclusion estimate that controls its overlap with
the relevant collective modes. Another local density floor,
finite-radius inverse, or unspecialized spatial derivative
bound will not provide the missing implication.
The complete interacting continuum construction, physical
scale/clock and positive physical YM gap remain unproved.

The diagnostic
\begin{center}\small
\path{scripts/verify_ym_c18_v78_finite_range_obstruction.py}
\end{center}
checks the original factor incidence and term supports,
the polynomial path and score coefficients, the exact
radial identities, and the lattice-shell counting.
It does not numerically certify Gaussian domination or
simulate a native finite-time law. The classical theorem
is used with the hypotheses stated above.

The previous goal turn was progress. This turn supplies
a verified obstruction that changes the permissible next
argument, not a new native upper bound.
Scheduled companion and broad literature reviews remain
V80; archive/restart remains V100. Removing this append
and restoring the title recovers V77 byte for byte.
Only the active version is replaced, and only
\texttt{.tex} files are stored in \path{latex_notes}.
% END CYCLE 18 VERSION 78 APPEND
% BEGIN CYCLE 18 VERSION 79 APPEND
\clearpage
\section[V79: Native flat bias and straight anti-alignment droplets]{V79: Native flat bias, a noncentral stationary test and straight anti-alignment droplets}
\label{c18v79:context}

V78 shows why local regularity cannot by itself settle V77's
distant source response. We now use information specific to
the actual Wilson normal Jacobian. Its value at the flat
configuration grows at a rate bounded below by a positive
constant times volume, throughout the smoothing interval.
Combining this with the stationary determinant identity
gives a persistent native likelihood and source bias against
specified anti-aligned stationary configurations. It also
pays a new local droplet probability with arbitrary exterior.

The scope is \(\beta=0\), coarse identity, \(N=2n\),
\(n\ge5\), \(0\le t\le4\). Write \(\mathsf v=n^3\).
All densities in this append are the native
\(\nu_t=k_t\mu_I/Z_t\), \(k_t=\operatorname{Jac}\eta_t\),
and \(a_t=\partial_t\log(k_t/Z_t)\).
No comparison Gibbs density is substituted. The full
source inverse and physical mass gap remain unproved.

\subsection{The flat normal Jacobian has an extensive positive rate}

At the all-identity point \(p_+\), let
\(\mathcal L=d_1^*d_1\otimes I_3=-\operatorname{Hess}W(p_+)\),
\(B=Db_2(p_+)\), and \(\mathsf V=B^*/\sqrt2\).
Thus \(\mathsf V^*\mathsf V=I\). V67 gives
\[
 \begin{gathered}
 k_t(p_+)=\det(\mathsf V^*e^{-2t\mathcal L}\mathsf V)^{-1/2},
 \qquad
 \Pi_t=\operatorname{proj}\operatorname{Ran}
                                  (e^{-t\mathcal L}\mathsf V),\\
 r_+(t):=\partial_t\log k_t(p_+)
                          =\operatorname{tr}(\Pi_t\mathcal L).
 \end{gathered}
 \tag{V79.1}\label{c18v79:flatrate}
\]
These are ambient normal determinants, not tangent-volume
determinants or probability-preserving transport Jacobians.

Take principal coarse Fourier frequencies
\(p=2\pi k/n\), \(-\pi\le p_i\le\pi\), and put
\[
 \lambda_*(p)=4\sum_{i=1}^3\sin^2(p_i/4),\qquad
 R_n=6\sum_p\lambda_*(p).
\]
\begin{theorem}[A native flat density and source floor]
\label{c18v79:flat}
For the actual normal chart,
\[
 R_n\le r_+(t)\le36\mathsf v,\qquad
 \log k_t(p_+)\ge tR_n,\qquad R_n\ge12\mathsf v .
 \tag{V79.2}\label{c18v79:flatfloor}
\]
\end{theorem}
\begin{proof}
It suffices first to work with one color. Let \(d_{\rm f}\)
and \(d_{\rm c}\) be the fine and coarse vertex gradients,
and let \(\mathcal R\) restrict a fine vertex field to
the coarse roots. Summing two consecutive link differences
gives the exact identity
\[
                     Bd_{\rm f}=d_{\rm c}\mathcal R .
\]
Therefore \(\mathsf V\) sends
\(\ker d_{\rm c}(p)^*\) into \(\ker d_{\rm f}^*\).
This coarse subspace has dimension two for \(p\ne0\)
and dimension three at \(p=0\). It is not a deletion
of physical or longitudinal modes from the determinant.

The coarse Bloch space has the eight fine aliases
\(\theta=p/2+\pi\varepsilon\), \(\varepsilon\in\{0,1\}^3\).
On fine divergence-free vectors the curl Gram acts with
eigenvalue
\[
       \lambda(\theta)=4\sum_i\sin^2(\theta_i/2).
\]
Each alias has \(\lambda(\theta)\ge\lambda_*(p)\).
This follows componentwise from
\(\sin^2(p_i/4)\le\cos^2(p_i/4)\) for \(|p_i|\le\pi\).

The space \(e^{-t\mathcal L}\mathsf V\ker d_{\rm c}(p)^*\)
still lies in the same fine divergence-free Bloch space.
It has dimension at least two and lies in
\(\operatorname{Ran}\Pi_t\). Choose two orthonormal vectors
there and extend them to an orthonormal basis of the
three-dimensional normal Bloch space. The first two
quadratic forms of \(\mathcal L\) are at least
\(\lambda_*(p)\); the last is nonnegative.
Thus the trace contribution is at least \(2\lambda_*(p)\).
Summing all frequencies and three colors proves \(r_+\ge R_n\).
Integration, with \(k_0=1\), proves the logarithmic bound.
V67's compressed-exponential convexity says that \(r_+\)
is nonincreasing. Its initial value is
\(-\operatorname{tr}(\mathsf V^*\operatorname{Hess}W(p_+)
\mathsf V)=36\mathsf v\), giving the upper bound.

For the last estimate define
\[
 c_n=\frac1n\sum_{k\ {\rm principal}}\cos(\pi k/n)
 =\begin{cases}
    [n\sin(\pi/(2n))]^{-1},&n\ \hbox{odd},\\
    n^{-1}\cot(\pi/(2n)),&n\ \hbox{even}.
   \end{cases}
\]
The finite geometric sum proves these identities and gives
\(R_n=36\mathsf v(1-c_n)\).
Concavity of sine on \([0,\pi/6]\) yields
\(n\sin(\pi/(2n))\ge3/2\) for \(n\ge3\);
also \(\cot x\le\csc x\) on this interval.
Hence \(c_n\le2/3\), and \(R_n\ge12\mathsf v\).
\end{proof}

The proof used two test directions to lower-bound a
\emph{full} normal trace; it did not replace that trace
by a transverse determinant. In particular it does not
assume that the moving projection commutes with
\(\mathcal L\). V67 only used \(k_t(p_+)\ge1\);
the extensive lower bound here is additional information.

\subsection{All stationary points: a persistent relative source bias}

Let \(p\in M_I\) be any ambient Wilson stationary point,
not necessarily central, and put \(C(p)=W_{\rm cap}(p)\).
For \(A_p=\operatorname{Hess}W(p)\) and
\(\mathsf V_p=Db_2(p)^*/\sqrt2\), one has
\[
                \operatorname{tr}(\mathsf V_p^*A_p\mathsf V_p)
                              =-3C(p).
 \tag{V79.3}\label{c18v79:stationarytrace}
\]
Indeed a plaquette trace, as a function of one unit-round
link, is a height function on \(S^3\), so its diagonal
link Hessian is minus its trace times the round metric.
The two consecutive captured halves of one pair occur
in no common plaquette. Their off-diagonal Hessian block
is therefore zero. A normalized normal column has
squared norm \(1/2\) on each half, with isometric adjoint
transports. Taking its three-color trace and counting
two captured half-links per captured plaquette proves
(V79.3). This trace identity in fact holds at any point;
stationarity is needed for V67's time-dependent determinant.

That determinant and its convex logarithm imply
\[
 \partial_t\log k_t(p)\le3C(p).
\]
Subtract (V79.2). Since \(p,p_+\) are fixed points of
\(\eta_t\), and the same normalizer cancels, we obtain
\[
 \begin{gathered}
 a_t(p)-a_t(p_+)
       \le3C(p)-R_n\le3C(p)-12\mathsf v,\\
 \log\frac{p_t(p)}{p_t(p_+)}
       \le t\{3C(p)-R_n\}
       \le t\{3C(p)-12\mathsf v\}.
 \end{gathered}
 \tag{V79.4}\label{c18v79:bias}
\]
Here \(p_t=k_t/Z_t\).
This is a relative source and likelihood statement.
It does not assert that \(a_t(p_+)\) itself is positive,
or that a point has positive probability.

There is a useful noncentral test of precisely the
anti-alignment channel used in V78. Suppose \(n\) is even.
Fix a unit imaginary quaternion \(q\), so \(q^2=-I\),
and set
\[
             U_{x,i}=q^{\,\sum_{j<i}x_j}.
 \tag{V79.5}\label{c18v79:abelian}
\]
The fine side \(2n\) is divisible by four, so this field
is periodic. Each ordered plaquette in a plane \(i<j\)
has holonomy \(q\). On each captured pair the other
coordinates are even, making the product of its two
halves \(I\). Thus the field lies in \(M_I\) and
\(C(p)=0\).

It is stationary: derivatives in the two color directions
orthogonal to \(q\) have zero trace. In the \(q\)-direction,
the two adjacent plaquettes in each transverse direction
have equal sine and opposite incidence, and cancel.
Every straight midpoint-path ratio is \(q^2\) or \(q^{-2}\),
hence \(-I\). This verifies a noncentral stationary field
with all \(6\mathsf v\) straight holonomies anti-aligned.
Equation (V79.4) gives, throughout the interval,
\[
   p_t(p)/p_t(p_+)\le e^{-12t\mathsf v},\qquad
   a_t(p)-a_t(p_+)\le-12\mathsf v .
 \tag{V79.6}\label{c18v79:noncentralbias}
\]
For comparison, V78's artificial law assigns these two
configurations equal density at time three: both have
the same aligned \(H\), while its \(C\) and field
coefficients vanish then. The native law does not.
This is an exact distinction between the two laws,
not yet a probability estimate for a noncentral phase.

\subsection{A global central straight-defect neighborhood}

At a central configuration write the four path signs
at a two-odd centre as \(z_0,z_1,z_2,z_3\).
The two straight signs and the captured sum there satisfy
\[
 C_y=(z_0+z_2)(z_1+z_3),\qquad
 d_y={\bf1}_{z_0z_2=-1}+{\bf1}_{z_1z_3=-1},\qquad
 C_y\le4-2d_y .
 \tag{V79.7}\label{c18v79:signs}
\]
If either straight sign is negative then \(C_y=0\);
if neither is negative then \(C_y\le4\).
This proves the inequality in all cases.
For the total number \(D_{\rm st}\) of negative straight
channels,
\[
               C(z)\le12\mathsf v-2D_{\rm st}(z).
\]
In particular \(D_{\rm st}\ge(3/4)6\mathsf v\) implies
\(C(z)\le3\mathsf v\). Equation (V79.4) then bounds its
central likelihood ratio by \(e^{-3t\mathsf v}\).
Unlike V67's earlier family, this permits positive \(C(z)\).

Use V67's product-metric boxes and full pinned-gauge
saturations, now with radius
\[
       \varepsilon_{\rm g}
          =\min\{\pi/4,(84K_*)^{-1/2}\}.
\]
The gradients of \(u_t=\log p_t\) vanish at central
configurations, and \(H(u_t)\le tK_*\).
The paired Taylor comparison of a central sign translate
with the flat box costs at most
\(21tK_*\varepsilon_{\rm g}^2\mathsf v\le t\mathsf v/4\).
It follows that, for the union \(\mathcal E_{\rm st}\)
of these saturated boxes with \(D_{\rm st}\ge9\mathsf v/2\),
\[
 \nu_t(\mathcal E_{\rm st})
    \le\min\{1,\exp[(14\log2-11t/4)\mathsf v]\}.
 \tag{V79.8}\label{c18v79:globalprob}
\]
The \(2^{14\mathsf v}\) center-gauge orbit count is exactly
V67's count; no independence of the straight signs is
assumed. At time four the exponent is
\(-\delta\mathsf v\), where
\[
                         \delta=11-14\log2>1 .
 \tag{V79.9}\label{c18v79:delta}
\]
For example the first four terms of \(e^{5/7}\) exceed
two, proving \(\log2<5/7\).

\subsection{Local flat repair retains the extensive gain}

We next retain arbitrary exterior configurations, rather
than use (V79.8) only on a whole-torus neighborhood.
Take V68's cube \(R\) of coarse side \(\ell\), now with
\(5\le\ell<n\). Its \(21\ell^3\) interior factors have
stored fine-link tails in a fine cube of side \(2\ell\).
The cut counts already proved in V68 are
\[
 b_R=48\ell^2-6\ell,\qquad
 N_{{\rm cap},R}=12\ell^3-12\ell^2+3\ell .
\]
The flat interior cut matrix differs from the flat
periodic matrix on side \(2\ell\) by deletion of
\[
             w_R=24\ell^2-6\ell
 \tag{V79.10}\label{c18v79:wrap}
\]
plaquette terms. Indeed the latter has \(24\ell^3\)
plaquettes and the former has \(3(2\ell)(2\ell-1)^2\).
Both act on the same interior stored edges and normal
columns after their natural periodic identification.

Use V68's trace-norm modulus
\[
 |\mathcal F_t(A)-\mathcal F_t(B)|
       \le\frac{e^{24t}-1}{24}\|A-B\|_1,\qquad
 \mathcal F_t(A)=-\tfrac12\log\det(\mathsf V^*e^{2tA}\mathsf V).
\]
Each deleted plaquette has trace norm twelve, including
three colors. The new flat lower bound on the periodic
interior therefore gives
\[
 \mathcal F_t(A_{{\rm flat},R})
       \ge12t\ell^3-\tfrac12w_R(e^{24t}-1).
\]
The signed pattern cut still has
\(\mathcal F_t(A_{w,R})\le3tC_R(w)\).
Cutting both full matrices costs
\(b_R(e^{24t}-1)\), as in V68, and their exterior
determinants cancel. Thus for every central exterior,
\[
 \begin{split}
 u_t(w,\xi)-u_t(I,\xi)
   &\le3tC_R(w)-12t\ell^3
                      +(b_R+w_R/2)(e^{24t}-1)\\
   &\le3tC_R(w)-12t\ell^3
                      +60(e^{24t}-1)\ell^2 .
 \end{split}
 \tag{V79.11}\label{c18v79:localcenter}
\]
No operator monotonicity of the exponential is used;
the cut matrices remain comparison matrices, not new laws.

There are \(3\ell(\ell-1)^2\) two-odd centres for which
all four captured plaquettes are wholly interior.
For a centre in plane \(i,j\), its coarse coordinates
in those directions range from \(0\) to \(\ell-2\);
the third ranges from \(0\) to \(\ell-1\).
Count only their \(6\ell(\ell-1)^2\) straight channels,
and denote their negative-sign count by \(D_R^{\rm st}\).
Other wholly interior captured plaquettes contribute
at most one each. Applying (V79.7) at complete centres yields
\[
 \begin{split}
 C_R(w)&\le N_{{\rm cap},R}-2D_R^{\rm st}(w),\\
 D_R^{\rm st}(w)\ge\tfrac34\,6\ell(\ell-1)^2
   &\ \Longrightarrow\
 C_R(w)\le3\ell^3+6\ell^2-6\ell .
 \end{split}
 \tag{V79.12}\label{c18v79:localdefects}
\]
Consequently (V79.11) is at most
\(-3t\ell^3+[60(e^{24t}-1)+18t]\ell^2\)
for every eligible central pattern.

\subsection{Boundary-safe gauge counting and the actual conditional bound}

Let
\[
 \varepsilon_{\rm loc}
       =\min\{\pi/4,(252DK_*)^{-1}\},\qquad D=\pi\sqrt2,
 \quad B_\zeta=\frac{72D^2K_\zeta}{e^\zeta-1}.
\]
Form V68's interior boxes of radius \(\varepsilon_{\rm loc}\),
with every exterior coordinate unrestricted, and their
full pinned-gauge saturations \(\mathcal E_w\).
The central sign change \(\sigma_w\) fixes the exterior
and commutes with the gauge action.

For \(\Psi_t=u_t\circ\sigma_w-u_t\), V68 proves
\(\sum_i|\nabla_i\Psi_t|\le3kDtK_*\), where
\(k\le21\ell^3\), and bounds the total arbitrary-exterior
change of \(\Psi_t\) by \(tB_\zeta\ell^2\).
Moving the interior from flat to its radius box thus
costs at most \(63DtK_*\varepsilon_{\rm loc}\ell^3
\le t\ell^3/4\).
Adding these already proved errors to (V79.11)--(V79.12)
gives, on the gauge-saturated flat event,
\[
 \Psi_t\le-\tfrac{11}{4}t\ell^3+S_t\ell^2,\qquad
 S_t=60(e^{24t}-1)+(18+B_\zeta)t .
 \tag{V79.13}\label{c18v79:repair}
\]
Changing variables by \(\sigma_w\) in the actual conditional
density proves this same exponential probability comparison
for every entire exterior \(\xi\), with the flat conditional
event on the right. Its normalizer is retained.

We can count patterns more sharply here than V68's
deliberately crude \(2^{21\ell^3}\).
Consider only center-valued gauge transformations supported
at the noncoarse vertices of
\(\{1,\ldots,2\ell-1\}^3\); all other gauge values are \(I\).
Their six incident stored edges lie inside \(R\), so
they fix every exterior factor. The number of available
vertices is
\[
 g_R=(2\ell-1)^3-(\ell-1)^3
                       =7\ell^3-9\ell^2+3\ell .
\]
Their action on central interior configurations is free.
If all interior edges were fixed, adjacent gauge values
would agree; connectivity to the fixed boundary or a
pinned coarse root forces every value to be \(I\).
The eligibility count is gauge invariant. Patterns in
one such orbit have identical saturated events.
Thus at most
\[
                  2^{21\ell^3-g_R}
                        =2^{14\ell^3+9\ell^2-3\ell}
 \tag{V79.14}\label{c18v79:localorbits}
\]
representatives are needed. No gauge value on a fixed
exterior has been subtracted in this count.

\begin{theorem}[Native local straight anti-alignment suppression]
\label{c18v79:localprobthm}
Let \(\mathcal E_R^{\rm st}\) be the union of the above
saturated interior boxes satisfying (V79.12).
For every entire exterior configuration \(\xi\),
\[
 \nu_t^R(\mathcal E_R^{\rm st}\mid\xi)
 \le Q_R(t):=\min\!\left\{1,
   \exp\left[(14\log2-11t/4)\ell^3+
                    (9\log2+S_t)\ell^2\right]\right\}.
 \tag{V79.15}\label{c18v79:localprob}
\]
The same bound holds unconditionally.
At \(t=4\), put
\[
 A_4=9\log2+60(e^{96}-1)+72+4B_\zeta .
\]
If \(5\le\ell<n\) and \(\ell\ge2A_4/\delta\), then
\[
          \nu_4^R(\mathcal E_R^{\rm st}\mid\xi)
                       \le e^{-\delta\ell^3/2}
                       \quad\hbox{for every }\xi ,
 \tag{V79.16}\label{c18v79:time4}
\]
with \(\delta>1\) from (V79.9).
\end{theorem}
\begin{proof}
Use (V79.13)'s conditional comparison for one representative
of each orbit in (V79.14), sum, and bound the flat event's
conditional probability by one. Dropping the favorable
term \(-3\ell\log2\) gives (V79.15). Integrating over the
exterior proves the unconditional bound. At time four
the exponent is \(-\delta\ell^3+A_4\ell^2\), which is
at most \(-\delta\ell^3/2\) under the stated threshold.
\end{proof}

This proves a nonempty new family: private free-spoke
choices can make both straight signs negative at every
complete centre. It is not necessary for the captured
sum to be strongly negative, as it was in V68.
The radius is extremely small and the threshold
astronomically large. Neither is a physical continuum scale.

\subsection{Verification, non-duplication and the remaining source problem}

The new flat trace floor, the stationary noncentral
anti-alignment test, and the straight-defect criterion
are native information absent from V78's counterexample.
The exact local gauge rank also supplies the entropy
saving needed for the weaker bulk bias. V67's determinant,
V68's noncommuting cut and arbitrary-exterior bounds,
and V76's channel geometry are reused, not reproved as
new general principles. No new external theorem is
imported in this append.

The diagnostic
\begin{center}\small
\path{scripts/verify_ym_c18_v79_native_flat_bias.py}
\end{center}
checks the local cut and complete-centre counts, the
boundary-safe gauge rank over \(\mathbb F_2\), and the
stationary Abelian phases exactly. Eight-alias native
flat determinants and their source rates are evaluated
at high precision as normalization regressions.
The written Bloch-space argument, not these finite samples,
proves the all-volume floor. These checks are not
nonlinear probability simulations or a proof assistant.

The noncentral stationary field in (V79.5) has a pointwise
likelihood bound, not a gauge-saturated neighborhood
probability bound from (V79.15). That theorem retains a
near-central assumption on its \emph{interior}; only the
exterior is arbitrary. The union does not exhaust
noncentral collective alignment or all configurations
carrying V77's distant source--cap covariance.
Small event weight alone still does not bound its
unshifted source inverse.

Thus the next substantive task is a native estimate on
the remaining noncentral collective/source sector, or
a direct bound on V77's complete derivative-moment sum.
The full probability strain, interacting-vacuum and
all-coarse extensions, physical scaling/clock,
nontrivial reflection-positive continuum construction,
and positive physical YM mass gap remain unproved.
The previous goal turn and this turn both made progress;
the complete goal remains active.

V80 is the next scheduled companion review and broad
primary-literature checkpoint; the new native bias and
the uncovered noncentral source sector should guide it.
Archive/restart remains V100. Removing this append and
restoring the title recovers V78 byte for byte. Only
its active filename is replaced. The \path{latex_notes}
folder contains only \texttt{.tex}; diagnostics and
build products stay outside it.
% END CYCLE 18 VERSION 79 APPEND
% BEGIN CYCLE 18 VERSION 80 APPEND
\clearpage
\section[V80: Noncentral neighborhoods and their actual source response]{V80: Noncentral neighborhoods and their actual source response}
\label{c18v80:context}

V79 separates the native law from V78's comparison model at a
noncentral stationary configuration, but a density at one point
does not control a probability. We now pass to a genuine
gauge-saturated neighborhood of that configuration, including
every color orientation. We then bound the contribution of the
\emph{actual centered source restricted to this event} to the
complete additive straight-holonomy inverse. No indicator is
silently substituted for the source and no mass term is added
to the inverse.

This is not a bound for the complementary source or for arbitrary
nonadditive functions of all holonomies. The continuum construction
and positive physical mass gap remain open. Throughout, the law is
the native electric normal-chart law
\(\nu_t=e^{u_t}\mu_I=k_t\mu_I/Z_t\), with
\(a_t=\partial_tu_t\), coarse identity, \(N=2n\), and
\(0\le t\le4\). General statements use \(n\ge5\);
the Abelian family below uses even \(n\ge6\).
Write \(\mathsf v=n^3\), \(D=\pi\sqrt2\),
\(S=D^2K_*\), and \(\overline S=\max\{1,S\}\).

\subsection{The scheduled companion and primary-literature checkpoint}

Fresh hashes show that the twelve companion sources reviewed at
V75 are unchanged. Their terminal mathematical scope statements
were revisited, not their entire cumulative proofs recertified.
In particular, SM V72 is a fixed-mode, fixed-time auxiliary
many-copy limit; NS V26 retains the full regularity problem;
shell-mixing V16 still leaves the noncommuting mixed-curvature
coefficient open. The perturbative ESM branches retain fixed
order/source scope. The Einstein-bridge, regenerative and
source-selection branches still distinguish an actual source
from a formally available inverse or invariant subspace.
No companion supplies the missing YM source estimate.

The scheduled broad sweep screened primary abstracts across
the following directions. These are applicability assessments,
not imports of unverified theorems.
\begin{itemize}
\item Lattice dynamics:
Shen--Zhu--Zhu,
\href{https://arxiv.org/abs/2204.12737}{arXiv:2204.12737},
prove strong-coupling lattice results. The normal-chart law here
is not that Wilson Gibbs law, and a physical weak-coupling
continuum conclusion cannot be transferred.
\item Nonconvex covariance representations:
Buchholz--Cotar--Schweiger,
\href{https://arxiv.org/abs/2608.14526}{arXiv:2608.14526},
treat real gradient interfaces through convex-mixture
representations. Such a representation has not been established
for this compact nonabelian determinant density.
\item Stochastic localization and entropy factorization:
Caputo--Chen--Parisi,
\href{https://arxiv.org/abs/2503.19419}{arXiv:2503.19419},
retain a high-temperature or weak-interaction input.
V78 explains why bounded local jets alone cannot replace it here.
\item Multiscale functional inequalities:
Bauerschmidt--Bodineau,
\href{https://arxiv.org/abs/1907.12308}{arXiv:1907.12308},
use control of the actual Polchinski evolution.
A finite density jet does not supply that control.
\item Markov spectral decomposition:
Qin,
\href{https://arxiv.org/abs/2504.01247v2}{arXiv:2504.01247v2},
revised 10 September 2026, retains benchmark and component gaps.
Neither an exact physical-vacuum Doob kernel nor its required
uniform comparison follows from our auxiliary chart.
\item Loop-equation/positivity methods:
Liu--Yang,
\href{https://arxiv.org/abs/2601.04316v2}{arXiv:2601.04316v2},
organize direct and indirect finite loop relations.
The suggested V58 cubic contraction was already completed in
V59; repeating it would not bound the full temporal remainder.
\item Trivializing flows and dynamical scaling:
L\"uscher,
\href{https://arxiv.org/abs/0907.5491}{arXiv:0907.5491},
distinguishes trivializing constructions from Wilson-flow
approximations. Chevyrev--Qu--Shen,
\href{https://arxiv.org/abs/2608.27828v2}{arXiv:2608.27828v2},
revised 17 September 2026, treats three-dimensional \(U(1)\)
dynamics, not four-dimensional interacting nonabelian YM.
\end{itemize}
The resulting direction is to retain exact probability
normalization and the source projection, using V33's already
proved Haar factorization. No external gap theorem is invoked
in the proofs below. Detailed file hashes and screening scope
are recorded outside \path{latex_notes}.

\subsection{A fixed forest and equal-Haar quotient boxes}

Specialize V33's forest to block size two. For a noncoarse
vertex \(x\), choose its parent by subtracting one in the first
odd coordinate. Each edge points from that parent to \(x\).
Every tree has one coarse root and depth at most three.
It contains the first half of every captured pair. Let
\(\mathcal F,\mathcal B,\mathcal C\) denote its forest edges,
the last captured halves, and the remaining chords. Then
\[
 |\mathcal F|=7\mathsf v,\qquad
 |\mathcal B|=3\mathsf v,\qquad
 |\mathcal C|=14\mathsf v.
 \tag{V80.1}\label{c18v80:counts}
\]
V33's exact coordinate change, disintegrated at coarse identity,
gives
\[
 M_I\longleftrightarrow
       \mathcal H_0\times \mathrm{SU}(2)^{\mathcal C},
 \qquad
 \mu_I=dk\,dy,\qquad
 \nu_t=dk\,e^{\bar u_t(y)}dy,\quad
 \bar u_t(y)=u_t(W^0(I,y)).
 \tag{V80.2}\label{c18v80:quotient}
\]
All forest and last-half links in \(W^0(I,y)\) equal \(I\).
This reuses the geometric Haar change of variables, with the
present density \(k_t/Z_t\); it does not identify that density
with V33's differently defined pushed density.
The native \(u_t,a_t\) are \(\mathcal H_0\)-invariant.
No orbit-volume factor is added, and no product quotient
metric is substituted for the actual Dirichlet form.

Every chord is an original free factor. Thus the inherited
gradient bounds and \(u_0=0\) give
\[
 \|\nabla_j\bar u_t\|_\infty\le DtK_*,
 \qquad
 \operatorname{osc}_j\bar u_t\le tS .
 \tag{V80.3}\label{c18v80:grad}
\]
For \(y^p\) the forest coordinates of \(p\in M_I\), define
\[
 B_r(y^p)=\{y:\ d_{S^3}(y_j,y_j^p)\le r
                         \text{ for every }j\in\mathcal C\},
 \qquad E_p(r)=\mathcal H_0\times B_r(y^p),
 \tag{V80.4}\label{c18v80:boxes}
\]
where the last set is interpreted through (V80.2).
These are full gauge-saturated neighborhoods, not individual
gauge-slice points. Their probability can be computed in \(dy\).
The map \(y_j\mapsto y_j^p y_j\) sends the flat box onto
\(B_r(y^p)\) and preserves product Haar. It need not commute
with the original gauge action: it is used only after the
exact disintegration.

\begin{lemma}[Stationary comparison of noncentral boxes]
\label{c18v80:comparison}
If \(p\) is an ambient Wilson stationary point and
\(C(p)=W_{\rm cap}(p)\), then
\[
 \nu_t(E_p(r))
 \le
 \exp\{t[3C(p)-12\mathsf v+28DK_*r\mathsf v]\}
                              \nu_t(E_I(r)).
 \tag{V80.5}\label{c18v80:boxcompare}
\]
For \(0<r<\pi\), writing
\[
 h(r)=\frac{r-\sin r\cos r}{\pi},
\]
one also has
\[
 \nu_t(E_I(r))
       \le \min\{1,[e^{tS}h(r)]^{14\mathsf v}\}.
 \tag{V80.6}\label{c18v80:flatbox}
\]
\end{lemma}
\begin{proof}
Move the \(14\mathsf v\) chord coordinates, each by distance
at most \(r\). Equation (V80.3) bounds each of the two errors
\(\bar u_t(y^p y)-\bar u_t(y^p)\) and
\(\bar u_t(y)-\bar u_t(I)\) in absolute value by
\(14\mathsf v DtK_*r\).
V79's stationary relative-density estimate bounds their
center difference by \(t[3C(p)-12\mathsf v]\).
Exponentiation and the equal-Haar change of variables prove
(V80.5), with the actual common normalizer retained.

The radial Haar density on unit-round \(S^3\) is
\(2\sin^2\theta/\pi\). Its integral from zero to \(r\)
is \(h(r)\).
By (V80.3), every one-chord conditional density with all
other chords fixed is at most \(e^{tS}\) relative to Haar.
Its ball probability is therefore at most \(e^{tS}h(r)\).
Conditioning on the other coordinates, successively removing
one ball indicator at a time, proves the product bound.
This argument does not assert independence under \(\nu_t\).
\end{proof}

\subsection{All color orientations of a noncentral stationary family}

For even \(n\), let \(p(q)\) be V79's field
\[
 U_{x,i}=q^{\,\sum_{j<i}x_j},\qquad
 q\in S^2\subset\operatorname{Im}\mathbb H,\quad q^2=-I.
\]
Every such field is stationary, has \(C(p(q))=0\), and
has all straight holonomies equal to \(-I\).
Along a forest, gauge transport is a product of powers of \(q\).
Consequently each \(y_j^{p(q)}\) is one of
\(I,-I,q,-q\), with its choice independent of the orientation
of \(q\). In particular,
\[
 d_{S^3}(y_j^{p(q)},y_j^{p(q')})
                        \le d_{S^2}(q,q').
 \tag{V80.7}\label{c18v80:orientation}
\]
Set the volume-independent constants
\[
 \begin{gathered}
 \varepsilon=\min\{\pi/8,(56DK_*)^{-1}\},\qquad
 r=2\varepsilon,\qquad L=-\log h(r)>0,\\
 M_\varepsilon=(1+\pi/\varepsilon)^3,\qquad
 \kappa_{\rm nc}=\frac{11L}{\overline S}>0,
 \qquad
 \mathcal N_n=\bigcup_{q\in S^2}E_{p(q)}(\varepsilon).
 \end{gathered}
 \tag{V80.8}\label{c18v80:constants}
\]
The union is compact and hence measurable. It is invariant
under the pinned gauge group and under global color conjugation.

\begin{theorem}[A uniform native noncentral neighborhood probability]
\label{c18v80:probthm}
For every even \(n\ge6\) and every \(0\le t\le4\),
\[
 \begin{split}
 \nu_t(\mathcal N_n)
 &\le \min\left\{1,\
 M_\varepsilon e^{-11t\mathsf v}
       \min\{1,\exp[14\mathsf v(tS-L)]\}\right\}\\
 &\le \min\{1,M_\varepsilon e^{-\kappa_{\rm nc}\mathsf v}\}.
 \end{split}
 \tag{V80.9}\label{c18v80:prob}
\]
\end{theorem}
\begin{proof}
Take a maximal Euclidean \(2\varepsilon/\pi\)-separated
set on the unit \(S^2\subset\mathbb R^3\).
The disjoint ambient balls of radius \(\varepsilon/\pi\)
lie in the ball of radius \(1+\varepsilon/\pi\);
there are at most \(M_\varepsilon\) centers.
Maximality gives a Euclidean cover. Since the spherical
distance is at most \(\pi/2\) times chord distance, it
is a spherical \(\varepsilon\)-net.
By (V80.7), each \(E_{p(q)}(\varepsilon)\) lies in
\(E_{p(q_j)}(2\varepsilon)\) for some net point \(q_j\).

For these enlarged boxes, (V80.5) has exponent at most
\(-11t\mathsf v\), since \(C=0\) and
\(28DK_*r=56DK_*\varepsilon\le1\).
Combine with (V80.6) and sum over the net.
For the second bound, replace \(S\) by \(\overline S\).
The exponent per volume is bounded by
\[
 \min\{-11t,\ -14L+(14\overline S-11)t\}.
\]
The decreasing and increasing affine functions cross at
\(t=L/\overline S\), where both equal
\(-11L/\overline S\). Their minimum never exceeds that
value, whether or not the crossing lies in \([0,4]\).
This proves the uniform bound, including \(t=0\).
\end{proof}

This uses both the native stationary bias at later times
and small Haar volume at earlier times. It is not an assertion
that the flat neighborhood is typical.
The radius is independent of \(n\), but extremely small.
The theorem covers a noncentral family that was only a
pointwise test in V79. It does not cover all configurations
whose straight holonomies are close to \(-I\).

\subsection{The actual restricted source, not only its event weight}

For a measurable event \(E\), define its centered source piece
\[
 b_{E,t}={\bf1}_E a_t-\nu_t({\bf1}_E a_t).
 \tag{V80.10}\label{c18v80:sourcepiece}
\]
It is generally nonzero outside \(E\), because its scalar
centering must be retained. It is not the temporal score
of the law conditioned on \(E\).

There are \(21\mathsf v\) original independent pair/free
factors. The inherited \(\|\nabla_i a_t\|_\infty\le DK_*\)
and factor diameter at most \(D\) give
\(\operatorname{osc}_i a_t\le S\).
Telescoping through all factors, and using \(\nu_t(a_t)=0\),
therefore yields
\[
 \|a_t\|_\infty\le21\mathsf v S,\qquad
 \|b_{E,t}\|_2^2
   =\operatorname{Var}_{\nu_t}({\bf1}_E a_t)
   \le441\mathsf v^2S^2\nu_t(E).
 \tag{V80.11}\label{c18v80:sourceL2}
\]
These are bounds for the complete actual source, including
its normalization term, not just its value at stationary points.

Let \(\mathcal V_{\rm add}\) be V76's space of all sums
\(\sum_{\alpha=1}^{6\mathsf v}f_\alpha(X_\alpha)\), completed
in its actual energy norm modulo constants. For centered
\(b\in L^2(\nu_t)\), write
\[
 \mathcal G_{\rm add,t}(b)
   =\sup_{f\in\mathcal V_{\rm add}}
             \{2\nu_t(bf)-\mathcal E_t(f,f)\}.
 \tag{V80.12}\label{c18v80:add}
\]
No restriction to the first character or to radial test
functions is made.

\begin{theorem}[Complete additive response of a rare source piece]
\label{c18v80:sourcebound}
For every measurable \(E\),
\[
 \mathcal G_{\rm add,t}(b_{E,t})
        \le441E_0^2e^{tS}S^2\mathsf v^3\nu_t(E).
 \tag{V80.13}\label{c18v80:addbound}
\]
In particular, for \(E=\mathcal N_n\),
\[
 \begin{split}
 \mathcal G_{\rm add,t}(b_{\mathcal N_n,t})
 &\le441E_0^2e^{4S}S^2M_\varepsilon
            \mathsf v^3e^{-\kappa_{\rm nc}\mathsf v}\\
 &\le441E_0^2e^{4S}S^2M_\varepsilon
             \left(\frac{3}{e\kappa_{\rm nc}}\right)^3 .
 \end{split}
 \tag{V80.14}\label{c18v80:uniforminverse}
\]
This is uniform in all even \(n\ge6\) and \(0\le t\le4\).
\end{theorem}
\begin{proof}
Let \(h_\alpha=\mathbb E_{\nu_t}[b_{E,t}\mid X_\alpha]\).
It is centered in the actual marginal \(q_{\alpha,t}dX\).
Conditional expectation contracts \(L^2\), and V76 gives
the round one-channel gap \(3e^{-tS}\). Hence
\[
 I_\alpha^E:=
 \langle h_\alpha,\mathcal L_{\alpha,t}^{-1}h_\alpha\rangle
 \le \frac{e^{tS}}{3}\|b_{E,t}\|_2^2
 \le147e^{tS}S^2\mathsf v^2\nu_t(E).
 \tag{V80.15}\label{c18v80:oneinverse}
\]
This inverse is unshifted on the marginal mean-zero space.

For \(f=\sum_\alpha f_\alpha(X_\alpha)\), its source
pairing is \(\sum_\alpha\langle h_\alpha,f_\alpha\rangle\).
V76's private-spoke lower bound gives
\(\mathcal E_t(f,f)\ge
2E_0^{-2}\sum_\alpha\mathcal D_{\alpha,t}(f_\alpha)\).
Taking the scalar variational supremum separately in each
channel therefore proves
\[
 \mathcal G_{\rm add,t}(b_{E,t})
              \le\frac{E_0^2}{2}\sum_{\alpha=1}^{6\mathsf v}I_\alpha^E.
\]
Together with (V80.15), this is (V80.13).
The marginal Riesz representatives exist by their positive
gaps, so the same argument applies to the \(L^2\), possibly
nonsmooth source (V80.10).
Finally insert (V80.9) and maximize
\(x^3e^{-\kappa_{\rm nc}x}\) for \(x>0\);
its maximum is \((3/(e\kappa_{\rm nc}))^3\).
No gap for the joint law was used.
\end{proof}

\subsection{An exact reduction of the remaining radial and additive source}

Because \(\mathcal N_n\) is globally conjugation invariant,
its one-channel projected source is radial.
Let \(s_\alpha=\operatorname{Re}X_\alpha\),
let \(p_{\alpha,t}(s)\,ds\) be its actual radial law,
and put \(F_{\alpha,t}(s)=\nu_t(s_\alpha\le s)\).
The radial flux of this source piece is exactly
\[
 J_\alpha^{\mathcal N}(s)=
 \nu_t\!\left[
 {\bf1}_{\mathcal N_n}a_t
      \bigl({\bf1}_{\{s_\alpha\le s\}}-F_{\alpha,t}(s)\bigr)
 \right].
 \tag{V80.16}\label{c18v80:flux}
\]
V77's radial integration argument applies to every centered
radial source, not only a temporal score. It gives
\[
 I_\alpha^{\mathcal N}
  =\int_{-1}^1
       \frac{|J_\alpha^{\mathcal N}(s)|^2}
            {(1-s^2)p_{\alpha,t}(s)}\,ds
  \le147e^{4S}S^2M_\varepsilon
                  \mathsf v^2e^{-\kappa_{\rm nc}\mathsf v}.
 \tag{V80.17}\label{c18v80:radialcost}
\]
The source is centered, so both endpoint fluxes vanish.
The right side controls the entire one-holonomy inverse
load, including all its characters, for this source piece.

For the complementary piece, retain its full scalar centering:
\[
 a_t=b_{\mathcal N_n,t}+b_{\mathcal N_n^c,t}.
 \tag{V80.18}\label{c18v80:split}
\]
The fluxes split with the same sign, and the square root of
(V80.12) is a dual seminorm. Thus
\[
 \left|
 \sqrt{\mathcal G_{\rm add,t}(a_t)}
 -\sqrt{\mathcal G_{\rm add,t}(b_{\mathcal N_n^c,t})}
 \right|
 \le\sqrt{\mathcal G_{\rm add,t}(b_{\mathcal N_n,t})}.
 \tag{V80.19}\label{c18v80:reduction}
\]
There is an identical triangle inequality for the square
roots of the one-channel costs in (V80.17).
In particular the error in the additive dual norm divided
by \(\sqrt{\mathsf v}\) tends to zero exponentially with
volume. This isolates a genuine noncentral sector that
cannot carry a volume-growing additive-source obstruction.
It does not set the complementary source to zero.

\subsection{Verification, non-duplication and next obligation}

The diagnostic
\begin{center}\small
\path{scripts/verify_ym_c18_v80_noncentral_source.py}
\end{center}
checks the chosen forest counts and depth, its coordinate
inverse and gauge invariance on noncommuting quaternion-group
configurations, the Abelian chord phases, the Haar-ball
normalization, exponent crossing, inverse constants, and
the centered source split. Finite group tests are word
regressions, not proofs of the Haar disintegration; that
proof is already V33's. The measure and all-volume inverse
bounds above rest on their written analytic proofs, not
Monte Carlo or a proof-assistant certificate.

The new results are the noncentral quotient-box comparison,
the continuous-orientation family with an all-time probability
bound, and its complete actual additive/radial source cost.
V33's forest disintegration, V79's stationary bias and V76's
kinetic comparison are reused. V67--V72's central droplet
estimates are not recounted as new noncentral results.

The remaining problem is not another isolated small-box
probability estimate. One must control the complementary
source in (V80.18), or directly bound V77's complete
distant source--cap covariance. Our neighborhoods constrain
all \(14\mathsf v\) quotient coordinates; their complement
contains almost all other noncentral geometries, including
configurations with coherent straight holonomies. A local
noncentral patch with arbitrary exterior is generally
\emph{not stationary}, so V79's stationary determinant
cannot be applied to such a patched field without a new
interface argument. The full nonadditive inverse, pointwise
Poisson Hessian/strain, all-coarse and interacting-vacuum
extensions, physical scale, and nontrivial reflection-positive
continuum with positive physical mass gap remain unproved.

The preceding goal turn and this one both make progress;
the full goal remains active. The next companion checkpoint
is V85, the next broad literature checkpoint V90, and the
archive/restart remains V100. Removing this append and
restoring the title recovers V79 byte for byte. Only the
active filename is replaced; \path{latex_notes} contains
only \texttt{.tex}, with checks and build products elsewhere.
% END CYCLE 18 VERSION 80 APPEND
% BEGIN CYCLE 18 VERSION 81 APPEND
\clearpage
\section[V81: Noncentral interfaces and conditional source bias]{V81: Noncentral interfaces, conditional source bias and rare block labels}
\label{c18v81:context}

V80 controls a neighborhood constraining the whole torus.
This does not yet say that a noncentral patch is suppressed
when everything outside the patch is arbitrary. The obstruction
is real: patching a stationary configuration generally destroys
stationarity, so the stationary determinant formula cannot be
applied to the patched field.

We resolve this interface step using the already proved
mixed-Hessian bounds, not a new frozen determinant.
Periodicity distributes the native global bias over a block;
a simultaneous interpolation bounds the discrepancy by its
boundary area. The same proof applies to the actual temporal
source, not only the log density. We then prove a local
probability bound conditional on the entire original exterior,
and a relative-source bound for a fixed quotient exterior.
The two conditioning statements are deliberately distinguished.

These results still concern a specified noncentral family,
not the whole complementary source of V80. The full interacting
continuum and physical YM mass gap remain unproved.
The setting is \(\beta=0\), coarse identity, \(N=2n\),
even \(n\ge6\), and \(0\le t\le4\).
Keep V80's depth-three forest, \(14n^3\) free chords,
\(\bar u_t(y)=u_t(W^0(I,y))\), and
\(\bar a_t(y)=a_t(W^0(I,y))\).

\subsection{The retained Hessian bounds and an even-block identity}

Fix the positive spatial weight \(\zeta\) of V63.
Its supremum remains outside the row sum:
\[
 \mathcal H_\zeta(h)
  =\sup_x\max_i\sum_j e^{\zeta d(i,j)}
                     \|(\operatorname{Hess}_{m_0}h)_{ij}(x)\|.
\]
On the forest section, varying chords varies only original
free factors. Restriction therefore preserves the bounds
\[
 \begin{array}{c|c|c|c}
 h&\theta_h&\mathcal H_\zeta(\bar h)
                       &\sup_j\|\nabla_j\bar h\|_\infty\\ \hline
 u_t&t&tK_\zeta&DtK_*\\
 a_t&1&K_\zeta&DK_*
 \end{array}
 \qquad
 B_\zeta=\frac{72D^2K_\zeta}{e^\zeta-1},\quad D=\pi\sqrt2 .
 \tag{V81.1}\label{c18v81:inputs}
\]
The chord metric here is merely a way to read the restricted
derivatives; no quotient product metric replaces the native
Dirichlet form.

Let \(p(q)\) be V79's stationary Abelian configuration, for a
unit imaginary quaternion \(q\). Its forest chord values are
\(y_j^q=q^{m_j}\), \(m_j\in\{0,1,2,3\}\).
Choose
\[
 A_j(q)=\alpha_{m_j}q,\qquad
 (\alpha_0,\alpha_1,\alpha_2,\alpha_3)
                       =(0,\pi/2,\pi,-\pi/2),\qquad
 \gamma_j(s)=\exp(sA_j(q)).
 \tag{V81.2}\label{c18v81:path}
\]
Every speed is at most \(\pi\le D\).
The forest is covariant under coarse translations.
The field and this chosen path have coarse period two:
shifting a fine coordinate by four leaves all their phases
unchanged. Both native scalars \(\bar u_t,\bar a_t\) are
coarse-translation invariant.

Let \(R\) be any translated coarse cube of even side
\(2\le\ell<n\), including cubes crossing a torus seam.
Let \(\mathcal C_R\) consist of chords whose stored tails
lie in its fine cube of side \(2\ell\).
Then \(|\mathcal C_R|=14\ell^3\).
For either scalar in (V81.1), define
\[
 T_R(h)=\sum_{j\in\mathcal C_R}
       \int_0^1
        \langle\nabla_j\bar h(\gamma(s)),\dot\gamma_j(s)\rangle\,ds .
\]
Periodicity gives the exact finite-volume identity
\[
 T_R(h)=\frac{\ell^3}{n^3}
                   [\bar h(y^q)-\bar h(I)]
                     \le-12\theta_h\ell^3 .
 \tag{V81.3}\label{c18v81:bulk}
\]
Indeed a chord is specified by one of fourteen within-cell
types and one of eight coarse parity classes. Every class
has \(\ell^3/n^3\) of its torus occurrences in \(R\).
The integrand is constant within each such class at each
\(s\), by translation covariance and (V81.2).
Summing proves the equality. The inequality is precisely
V79's global native density/source bias.
No identification with a smaller-volume density is made.

\subsection{A surface comparison without stationarity of the patch}

For an interior chord, let \(\rho_j\) be the distance of
its fine tail from the complement of the fine cube.
V68's six-face count, restricted to the chord subset, gives
\[
 \sum_{j\in\mathcal C_R}e^{-\zeta\rho_j}
                      \le\frac{72\ell^2}{e^\zeta-1}.
 \tag{V81.4}\label{c18v81:layers}
\]
The actual chord distance to an exterior chord is at least
\(\rho_j\). Thus at \emph{each common configuration},
\[
 \sum_{i\in\mathcal C_R}\sum_{j\notin\mathcal C_R}
          \|(\operatorname{Hess}\bar h)_{ij}\|
 \le \theta_h K_\zeta
                  \frac{72\ell^2}{e^\zeta-1}.
 \tag{V81.5}\label{c18v81:cutrows}
\]
Separate entry suprema are not summed.

\begin{lemma}[Native noncentral center comparison with arbitrary exterior]
\label{c18v81:interface}
For every quotient exterior \(\xi=y_{\mathcal C_R^c}\),
\[
 \bar h(y_R^q,\xi)-\bar h(I_R,\xi)
 \le \theta_h\left[-12\ell^3+\frac32B_\zeta\ell^2\right],
 \qquad h=u_t\text{ or }a_t .
 \tag{V81.6}\label{c18v81:center}
\]
\end{lemma}
\begin{proof}
First keep the exterior flat. In the derivative of the
interior-only path, compare its gradient to the same
interior gradient on the global path \(\gamma(s)\).
Connect all exterior coordinates simultaneously from
\(I\) to \(\gamma_{\mathcal C_R^c}(s)\), with parameter
\(r\mapsto\gamma_{\mathcal C_R^c}(rs)\). Each exterior
speed is at most \(sD\); the interior path speed is at
most \(D\). The mixed derivative uses only the cross
Hessian blocks in (V81.5), evaluated at the same point.
Integrating over \(r\in[0,1]\) and \(s\in[0,1]\) costs
\[
       \theta_h B_\zeta\ell^2\int_0^1s\,ds
                            =\tfrac12\theta_h B_\zeta\ell^2.
\]
Together with (V81.3), this proves (V81.6) for flat exterior
with coefficient \(1/2\) instead of \(3/2\).

Now compare
\(\bar h(y_R^q,\xi)-\bar h(I_R,\xi)\) to its flat-exterior
value. Use the two-parameter rectangle whose interior
path is (V81.2) and whose exterior coordinates travel
simultaneously along chosen minimizing geodesics from
\(I\) to \(\xi\). Both speeds are bounded by \(D\).
The cross-Hessian sum (V81.5) bounds this rectangle by
\(\theta_h B_\zeta\ell^2\). Adding the two errors proves
the claim. These are simultaneous-path estimates; the
placement of the supremum in (V81.1) is respected.
No patched point was declared Wilson stationary.
\end{proof}

Use V80's
\[
 \varepsilon=\min\{\pi/8,(56DK_*)^{-1}\},\qquad r=2\varepsilon .
\]
For the quotient boxes
\(B_R(z,r)=\{y_R:d(y_j,z_j)\le r,\ j\in\mathcal C_R\}\),
the gradient bound in (V81.1) and (V81.6) imply the
stronger two-point estimate
\[
 \sup_{y_R\in B_R(y_R^q,r)}\bar h(y_R,\xi)
 -\inf_{z_R\in B_R(I_R,r)}\bar h(z_R,\xi)
 \le \theta_h[-11\ell^3+\tfrac32B_\zeta\ell^2].
 \tag{V81.7}\label{c18v81:twopoint}
\]
Each box costs at most \(14\ell^3rD\theta_h K_*\);
the sum costs at most \(\theta_h\ell^3\).
Henceforth take \(\ell\) even with
\[
                    \max\{2,3B_\zeta\}\le\ell<n .
 \tag{V81.8}\label{c18v81:size}
\]
Then the right side of (V81.7) is at most
\(-10\theta_h\ell^3\). The threshold is finite but
extremely large; it is not a physical length scale.

\subsection{The repair fixes the original exterior and preserves its Haar law}

The preceding argument used quotient coordinates.
We now verify that its probability comparison is also
valid after conditioning on the entire original exterior.

Define \(T_{q,R}\) at fixed forest gauge coordinate \(k\)
by left multiplying \(y_e\) by \(y_e^q\) for
\(e\in\mathcal C_R\), and changing no other chord.
Recall V33's inverse formula
\[
                    U_e=k_x^{-1}W_e^0k_z,\qquad e=(x,z).
\]
On a changed chord the repair is therefore exactly
\[
                 U_e\longmapsto
                  k_x^{-1}y_e^q k_x\,U_e .
 \tag{V81.9}\label{c18v81:originalrepair}
\]
All forest coordinates depend only on forest links,
not on any chord. Those forest links are unchanged.
The captured last halves are unchanged as well.
Every unselected original chord is unchanged, so in
particular every original factor with tail outside
the fine cube is fixed.

Conditional on all forest factors, (V81.9) is a product
of Haar translations in the selected original free
factors. Fubini proves Haar preservation, also on
\emph{every} original-exterior conditional fibre.
The inverse uses \((y_e^q)^{-1}\).
The map commutes with pinned gauge transformations:
\(k_x\) changes to \(k_xh_x^{-1}\), and (V81.9) is
conjugated accordingly. This is a gauge-adapted
noncentral repair, not a literal constant multiplication
of the original links.

Let \(E_{q,R}(r)\) be the inverse image of
\(B_R(y_R^q,r)\) with every other quotient and gauge
coordinate unrestricted; let \(E_{0,R}(r)\) use the flat
box. The map \(T_{q,R}\) takes the latter event onto
the former. Equation (V81.7) for \(u_t\) and conditional
Haar preservation show, for every original exterior
\(\omega\),
\[
 \nu_t(E_{q,R}(r)\mid U_{\rm ext}=\omega)
 \le e^{-10t\ell^3}
          \nu_t(E_{0,R}(r)\mid U_{\rm ext}=\omega).
 \tag{V81.10}\label{c18v81:originalratio}
\]
The same statement holds conditional on every quotient
exterior. All normalizers here are those of the actual
conditionals and cancel only in this measure-preserving
comparison.

To estimate the flat probability for an original exterior,
condition additionally on all interior forest factors.
Now all \(k\)'s are fixed and the \(14\ell^3\) selected
original free factors are related to \(y_R\) by independent
left/right Haar translations. Every one-chord conditional
has log oscillation at most \(tS\), with \(S=D^2K_*\).
The same sequential conditioning used in V80 yields
\[
 \nu_t(E_{0,R}(r)\mid U_{\rm ext}=\omega)
      \le\min\{1,[e^{tS}h(r)]^{14\ell^3}\},\qquad
 h(r)=\frac{r-\sin r\cos r}{\pi}.
 \tag{V81.11}\label{c18v81:originalflat}
\]
This bound is uniform in the extra forest data and
therefore survives averaging them out. No assertion
of conditional product independence is used.

\subsection{A local noncentral probability and its actual conditional source}

Define the local continuous-orientation event
\[
 \mathcal N_R=\bigcup_{q\in S^2}E_{q,R}(\varepsilon),\qquad
 L=-\log h(2\varepsilon),\quad
 M_\varepsilon=(1+\pi/\varepsilon)^3,\quad
 \overline S=\max\{1,S\},\quad
 \kappa_{\rm loc}=\frac{10L}{\overline S}.
 \tag{V81.12}\label{c18v81:constants}
\]
Unlike V80's whole-torus event, \(\mathcal N_R\) imposes
conditions on only \(14\ell^3\) chord variables.
It is gauge invariant and leaves the entire exterior
unrestricted.

\begin{theorem}[Arbitrary-exterior native noncentral block suppression]
\label{c18v81:probability}
Under (V81.8), for every original exterior \(\omega\),
\[
 \begin{split}
 \nu_t(\mathcal N_R\mid U_{\rm ext}=\omega)
 &\le \min\{1,M_\varepsilon e^{-10t\ell^3}
                       \min(1,e^{14\ell^3(tS-L)})\}\\
 &\le \rho_\ell:=
           \min\{1,M_\varepsilon e^{-\kappa_{\rm loc}\ell^3}\}.
 \end{split}
 \tag{V81.13}\label{c18v81:prob}
\]
These bounds hold for every quotient exterior as well,
and hence unconditionally, uniformly in \(n\) and
\(0\le t\le4\).
\end{theorem}
\begin{proof}
V80's fixed \(\varepsilon\)-net of \(S^2\) has at most
\(M_\varepsilon\) points. Each chord center is \(I,-I,q\)
or \(-q\), so the radius-\(\varepsilon\) union is covered
by their radius-\(2\varepsilon\) boxes.
Sum (V81.10)--(V81.11) over the net.
For the all-time bound use
\[
 \min\{-10t,\,-14L+(14\overline S-10)t\}
                       \le-\frac{10L}{\overline S}.
\]
The two affine functions cross at \(t=L/\overline S\).
This is a bound on their minimum for all \(t\ge0\),
even if the crossing lies outside the chart interval.
\end{proof}

There is also a direct statement about the actual source.
Here the conditioning variable is explicitly
\(\xi=y_{\mathcal C_R^c}\); let \(\bar\nu_t^R(\cdot\mid\xi)\)
be that normalized quotient conditional.
Write \(A_R=B_R(I_R,2\varepsilon)\).

\begin{proposition}[A native relative conditional-source rate]
\label{c18v81:source}
For every fixed quotient exterior \(\xi\),
\[
 \begin{split}
 &\mathbb E_{\bar\nu_t^R}[\bar a_t\mid\mathcal N_R,\xi]
       -\mathbb E_{\bar\nu_t^R}[\bar a_t\mid A_R,\xi]
                          \le-10\ell^3,\\
 &\partial_t\log
   \frac{\bar\nu_t^R(\mathcal N_R\mid\xi)}
        {\bar\nu_t^R(A_R\mid\xi)}
                          \le-10\ell^3 .
 \end{split}
 \tag{V81.14}\label{c18v81:sourcerate}
\]
\end{proposition}
\begin{proof}
Apply (V81.7) to \(h=a_t\).
It compares every point in each noncentral box with
every point of \(A_R\), at the same exterior, and is
uniform over \(q\). Thus it applies to their union.
Averaging over the two events proves the first line.

The events have positive probability and are independent
of \(t\). Smooth finite-volume density on a compact
space permits differentiation under their integrals.
The actual conditional score is
\(\bar a_t-\mathbb E[\bar a_t\mid\xi]\).
Consequently the derivative of the log event-probability
ratio is the difference of event means in the first line;
the conditional normalizer cancels. This proves the
second line without replacing the source by an indicator.
\end{proof}

The proposition is a relative rate, not a claim that
either individual conditional event probability decreases
at that rate. Nor is its derivative inequality asserted
after averaging the forest variables while fixing an
original exterior. Those mixing weights depend on time.
The probability bound (V81.13) survives that averaging
by uniformity; a derivative bound need not.

\subsection{Simultaneous rare labels and their cluster sizes}

The quotient-exterior version of (V81.13) can be assembled
without assuming probabilistic independence.
Take disjoint chord blocks \(R_1,\ldots,R_m\) satisfying
the same size condition and put
\(Z_j={\bf1}_{\mathcal N_{R_j}}\).
Each \(Z_j\) is a function only of its own quotient
chords. Conditioning on earlier labels is weaker than
conditioning on the entire exterior of the next block,
so
\[
 \mathbb P(Z_j=1\mid Z_1,\ldots,Z_{j-1})\le\rho_\ell .
 \tag{V81.15}\label{c18v81:sequential}
\]
For example one can sample the labels successively
using independent uniform random variables. Comparing
each conditional success threshold to \(\rho_\ell\)
couples them below independent Bernoulli\((\rho_\ell)\)
variables. This is one-sided stochastic domination,
not independence of the native labels.

In particular, for every subset \(J\) and \(\lambda\ge0\),
\[
 \mathbb P(Z_j=1\ \forall j\in J)\le\rho_\ell^{|J|},
 \qquad
 \mathbb E\exp\!\left(\lambda\sum_{j=1}^mZ_j\right)
          \le[1+\rho_\ell(e^\lambda-1)]^m.
 \tag{V81.16}\label{c18v81:moments}
\]
The first also follows by successive conditional
expectations; expansion of the exponential product
then proves the second.

For a concrete spatial statement, suppose additionally
\(\ell\mid n\), and partition the torus into these blocks.
Join face-neighboring blocks, a graph of degree at most six.
For a fixed block \(o\), if its bad-label component has
at least \(k\) vertices, there is a connected \(k\)-set
of bad blocks containing \(o\). A deterministic
depth-first traversal of a spanning tree gives a walk
of length \(2(k-1)\) whose visited set is that \(k\)-set.
There are at most \(6^{2(k-1)}=36^{k-1}\) such walks.
Hence, for \(1\le k\le(n/\ell)^3\),
\[
 \mathbb P(|\mathcal K_o|\ge k)
                       \le \rho_\ell(36\rho_\ell)^{k-1},
 \tag{V81.17}\label{c18v81:cluster}
\]
where \(\mathcal K_o\) is empty if \(Z_o=0\).
An even block size satisfying both (V81.8) and
\[
                \ell^3\ge
                 \frac{\log(72M_\varepsilon)}{\kappa_{\rm loc}}
 \tag{V81.18}\label{c18v81:rarity}
\]
has \(\rho_\ell\le1/72\), so (V81.17) is at most
\(\rho_\ell\,2^{-(k-1)}\). All constants are independent
of the number of blocks and of chart time.

This is a bound for clusters of the specified rare
labels. It is not exponential clustering of arbitrary
YM observables in their complementary region.
No global total-defect tightness is asserted: the
number of rare labels can still be extensive, as the
parallel sector notes warned.

\subsection{Verification and the remaining source-coverage problem}

The exact diagnostic
\begin{center}\small
\path{scripts/verify_ym_c18_v81_noncentral_interface.py}
\end{center}
checks the 112 periodic chord classes and their block
fractions on even tori, including seam-crossing blocks.
It verifies the forest-adapted repair and its inverse,
pinned-gauge equivariance, and unchanged original exterior
using noncommuting quaternion-group configurations.
It also checks the six-face count, the \(1/2+1=3/2\)
surface coefficient, radius budget, exponent crossing,
and elementary rare-label moment algebra.

An exact native Wilson-force regression further checks
the failed shortcut: on \(n=6,\ell=2\), the patched
forest-section field has 207 nonzero Cartan derivative
rows, of maximum absolute value four, whereas the
global periodic field is stationary. The interface
proof above does not require the patch to be stationary.
These finite tests supplement the analytic proofs;
they neither simulate the nonlinear native probability
nor provide proof-assistant certification.

The new input is the periodic global-to-local bias
identity combined with the native mixed-Hessian
surface comparison. The ensuing conditional source
rate and many-block bounds were not available from
V80's whole-torus neighborhood theorem.
V33's Haar factorization, V63's simultaneous rectangles,
V68's layer count, V79's global stationary bias and
V80's orientation net are reused, not reproved as new
general principles. No new external theorem is imported.

The next step must address source coverage or the
complementary conditional-mean return. Small bad-label
clusters alone do not imply a source inverse or a
positive gap on their complement. The events still
constrain every chord of a block to a small neighborhood
of a specific periodic family; they do not exhaust
noncentral coherent configurations. One cannot simply
sum (V81.14) into an unconditional source estimate.
A full distant source--cap bound, nonadditive inverse,
probability strain, interacting-vacuum/all-coarse
extension, physical scaling, and nontrivial
reflection-positive continuum with positive physical
mass gap remain unproved.

The preceding goal turn and this one both made progress;
the full goal remains active. The completed V80 companion
and broad literature checkpoint is retained; the next
checkpoints are V85 and V90 respectively. Archive/restart
remains V100. Removing this append and restoring the
title recovers V80 byte for byte. Only the active
filename is replaced, and only \texttt{.tex} files
remain in \path{latex_notes}.
% END CYCLE 18 VERSION 81 APPEND
% BEGIN CYCLE 18 VERSION 82 APPEND
\clearpage
\section[V82: Source-preserving removal of noncentral patches]{V82: Source-preserving removal of noncentral patches from the complete additive response}
\label{c18v82:context}

V81 bounds specified noncentral patches, even with the whole
exterior fixed. It does not show that their complement has a
gap. We now use that result to make a precise reduction of the
remaining source problem: condition on the absence of all
translated patches at a logarithmic-volume block scale, and
compare the \emph{entire} additive inverse response before
and after this conditioning. Both the temporal source and
the actual Dirichlet form change. Neither change is dropped.

The essential new estimate is rarity conditional on the value
of any straight holonomy, not just on a spatial exterior.
A private forest chord gives a Haar coordinate for that
holonomy. Saturating the patch event along this chord has a
fixed cost, independent of the ambient volume. This controls
the event-weighted energy of every additive test simultaneously,
and also the derivative of the conditioned marginal density.

This is stronger than replacing a centered rare source while
keeping the original law, as in V72 and V80. It is not a
coercivity theorem on the hard-conditioned full configuration
space. The no-patch source remains to be estimated; the full
interacting continuum and physical YM mass gap remain unproved.

\subsection{The actual patch union and a private quotient chord}

Work at \(\beta=0\), coarse identity, \(N=2n\), even \(n\ge6\),
and \(0\le t\le4\). Put \(v=n^3\). Retain V80's fixed forest
and V81's events \(\mathcal N_R\), constants
\(\varepsilon,M_\varepsilon,\kappa_{\rm loc},B_\zeta\).
Choose an even side satisfying
\[
                 \max\{2,3B_\zeta\}\le\ell<n .
\]
Let \(\mathcal R_{n,\ell}\) be the \(v\) coarse translations
of the side-\(\ell\) cube, with the same periodic centers
\(y^q\) as in V81. Define the time-independent events
\[
 \mathcal U=\bigcup_{R\in\mathcal R_{n,\ell}}\mathcal N_R,
 \qquad \mathsf G=\mathcal U^c,\qquad
 w_t=\nu_t(\mathcal U).
 \tag{V82.1}\label{c18v82:event}
\]
Overlaps are allowed. V81 immediately supplies the input
\[
                   w_t\le vM_\varepsilon
                                      e^{-\kappa_{\rm loc}\ell^3}.
 \tag{V82.2}\label{c18v82:union}
\]
The union bound itself is not claimed as a new method.

Every V76 straight channel \(X_\alpha\) has a private chord
in this particular forest. To see this directly, its
two-odd center has four free spokes and exactly one
incoming forest spoke, obtained by subtracting the first
odd coordinate. The two opposite pairs defining the
channels partition those four spokes. Each pair therefore
contains at least one chord. Captured prefixes are fixed
to identity on the forest section. Varying the chosen
chord, with all other quotient coordinates fixed, changes
\(X_\alpha\) by left/right multiplication and possibly
inversion. It is a Haar-preserving isometry of
\(\mathrm{SU}(2)\). There are \(3v\) one-chord and \(3v\)
two-chord channels. A chord can be chosen separately for
each; the chosen private chords are distinct.

The same statement may be checked before gauge fixing:
pinned gauge factors cancel in the relative holonomy,
whose endpoints are coarse roots. We use the product
quotient Haar measure only for disintegration, never
as a replacement for the moving Dirichlet metric.

\subsection{Rarity survives conditioning on an arbitrary straight value}

Write \(\bar\nu_t=e^{\bar u_t(y)}\,dy\) for the quotient
law. Recall \(S=D^2K_*\) and the one-chord oscillation
bound \(\operatorname{osc}_j\bar u_t\le tS\).
Set
\[
 h_\varepsilon=\frac{\varepsilon-\sin\varepsilon
                                  \cos\varepsilon}{\pi},
 \quad C_{\rm obs}=\frac{e^{12S}}{h_\varepsilon},
 \quad
 \delta_{n,\ell}
       =C_{\rm obs}M_\varepsilon v
                                  e^{-\kappa_{\rm loc}\ell^3}.
 \tag{V82.3}\label{c18v82:delta}
\]
Here \(h_\varepsilon\) uses radius \(\varepsilon\), not
V81's doubled-radius covering boxes.

\begin{lemma}[One-chord saturation and observable-conditional rarity]
\label{c18v82:saturation}
For every straight channel, its actual conditional probability
has the bound
\[
 e_{\alpha,t}(X):=\nu_t(\mathcal U\mid X_\alpha=X)
       \le \frac{e^{3tS}}{h_\varepsilon}\,w_t
       \le\delta_{n,\ell}.
 \tag{V82.4}\label{c18v82:conditional}
\]
It holds for a canonical disintegration, in particular
almost everywhere under the marginal. No independence
between a patch and the channel is assumed.
\end{lemma}
\begin{proof}
Choose a private chord \(j\) for the channel and let
\(\operatorname{Sat}_j\mathcal U\) be the cylinder obtained
by allowing arbitrary values of that chord whenever its
fiber meets \(\mathcal U\). All events here are measurable:
the individual continuous-orientation events and their
finite union are compact; projection to the remaining
compact coordinates is compact.

A nonempty chord fiber of \(\mathcal U\) contains a Haar
ball of radius \(\varepsilon\). Indeed choose a patch and
an orientation witnessing nonemptiness. If \(j\) is in
that patch, keeping its other constraints leaves precisely
such a permitted ball in coordinate \(j\); if \(j\) is
outside it, the whole chord fiber is permitted.
This argument does not require a measurable choice of
orientation.

Let \(z=y_{j^c}\), with actual marginal \(m_t(dz)\).
The conditional chord density \(r_t(y_j\mid z)\) satisfies
\[
                        e^{-tS}\le r_t\le e^{tS}.
\]
Consequently
\[
 m_t(\operatorname{Sat}_j\mathcal U)
                   \le e^{tS}h_\varepsilon^{-1}w_t .
 \tag{V82.5}\label{c18v82:satcost}
\]
For fixed \(z\), change from \(y_j\) to \(X_\alpha\)
by the Haar isometry just proved. If \(q_{\alpha,t}\)
is the group marginal density and \(q_{\alpha,t}^{\,U}\)
is the subprobability density restricted to \(\mathcal U\),
then, for that disintegration,
\[
 q_{\alpha,t}(X)\ge e^{-tS},\qquad
 q_{\alpha,t}^{\,U}(X)
      \le e^{tS}m_t(\operatorname{Sat}_j\mathcal U).
\]
Divide, apply (V82.5), and then (V82.2). The three
oscillation costs are displayed; none is a conditional
independence assertion.
\end{proof}

\subsection{The conditioned source includes its changing normalization}

Suppose \(\delta=\delta_{n,\ell}<1\), and define the
auxiliary conditioned law and its actual score by
\[
 \nu_t^{\,G}=\frac{{\bf1}_{\mathsf G}\nu_t}{1-w_t},
 \qquad
 a_t^{\,G}=a_t-\nu_t^{\,G}(a_t)
                 =a_t+\frac{\dot w_t}{1-w_t},
 \qquad
 \dot w_t=\nu_t({\bf1}_{\mathcal U}a_t).
 \tag{V82.6}\label{c18v82:score}
\]
The signs use \(\nu_t(a_t)=0\). In particular, the score
of the conditioned law is not merely \({\bf1}_{\mathsf G}a_t\).
The fixed event permits differentiation under the compact
integrals. All identities below concern this prescribed
time-dependent conditioned law, not a new dynamics.

Let \(q_\alpha^G\) be its exact group marginal. Then
\[
 q_\alpha^G=q_\alpha\frac{1-e_\alpha}{1-w_t},
 \qquad
 (1-\delta)q_\alpha\le q_\alpha^G
                                  \le(1-\delta)^{-1}q_\alpha .
 \tag{V82.7}\label{c18v82:qG}
\]
We suppress \(t\) in some subscripts. Marginals of the
hard-conditioned law need only be measurable; we do not
assume smoothness across the event boundary.

Use the inherited amplitude bound
\[
                  \mathfrak B=21vS,\qquad |a_t|\le\mathfrak B .
\]
Writing \(A_\alpha=\mathbb E[a_t\mid X_\alpha]\), direct
differentiation gives
\[
 \dot e_\alpha
   =\mathbb E[{\bf1}_{\mathcal U}a_t\mid X_\alpha]
                                      -e_\alpha A_\alpha,
 \quad |\dot e_\alpha|\le2\mathfrak B e_\alpha,
 \quad |\dot w_t|\le\mathfrak B w_t .
 \tag{V82.8}\label{c18v82:edot}
\]
Thus the actual marginal scores satisfy
\[
 \partial_t\log q_\alpha^G-\partial_t\log q_\alpha
          =-\frac{\dot e_\alpha}{1-e_\alpha}
                               +\frac{\dot w_t}{1-w_t}.
 \tag{V82.9}\label{c18v82:margscore}
\]
This is the normalization term that a naive averaging of
V81's relative-source inequality would miss.

For the variational comparison we need a centered density
derivative, not a difference of scores under different laws.
Set
\[
 b_\alpha=\frac{\partial_tq_\alpha^G-\partial_tq_\alpha}
                                                        {q_\alpha},
 \qquad
 T=\frac{3\mathfrak B\delta}{(1-\delta)^2}.
\]
Then
\[
                \int b_\alpha q_\alpha\,dX=0,
                         \qquad |b_\alpha|\le T .
 \tag{V82.10}\label{c18v82:b}
\]
Indeed let
\(m_\alpha=\mathbb E[{\bf1}_{\mathcal U}a_t\mid X_\alpha]\).
Differentiating \(q_\alpha^G=(q_\alpha-q_\alpha^U)/(1-w_t)\)
and using \(\partial_tq_\alpha=q_\alpha A_\alpha\),
\(\partial_tq_\alpha^U=q_\alpha m_\alpha\), gives
\[
 b_\alpha
    =\frac{w_tA_\alpha-m_\alpha}{1-w_t}
                   +\frac{(1-e_\alpha)\dot w_t}{(1-w_t)^2}.
\]
Its absolute value is bounded by
\(\mathfrak B[2\delta/(1-\delta)+\delta/(1-\delta)^2]\le T\).
Its mean vanishes because both marginal densities integrate
to one. This proof retains the time-dependent mixing weights.

\subsection{Simultaneous comparison of all additive energies and sources}

Let \(\mathcal F_{\rm add}\) be the real additive test space
of V76, initially smooth functions
\[
             f=\sum_{\alpha=1}^{6v}h_\alpha(X_\alpha).
\]
Every group harmonic is allowed. Write
\[
 \begin{array}{ll}
 \mathcal E_t(f)=\displaystyle\int|df|_{\widehat m_t^{-1}}^2d\nu_t,
 &L_t(f)=\nu_t(a_tf),\\[2pt]
 \mathcal E_t^G(f)=\displaystyle\int|df|_{\widehat m_t^{-1}}^2d\nu_t^G,
 &L_t^G(f)=\nu_t^G(a_t^Gf).
 \end{array}
 \tag{V82.11}\label{c18v82:forms}
\]
The metric is the same actual moving metric in both rows.
No derivative is taken of the hard-event indicator.
Set
\[
 \mathcal D_t(h)=
       \sum_\alpha\int|\nabla h_\alpha|^2q_{\alpha,t}\,dX .
\]
V76's private-spoke proof gives, pointwise before integration,
\[
 2\sum_\alpha|\nabla h_\alpha(X_\alpha)|^2
 \le |df|_{m_0^{-1}}^2
 \le6\sum_\alpha|\nabla h_\alpha(X_\alpha)|^2.
 \tag{V82.12}\label{c18v82:pointwise}
\]
The upper bound pays the fourfold prefix overlap; the lower
bound uses only private free spokes. In particular
\(\mathcal E_t(f)\ge2E_0^{-2}\mathcal D_t(h)\).

\begin{proposition}[All-test additive form and source stability]
\label{c18v82:stability}
Put \(c=3E_0^4\delta\), and assume \(c<1\). Then
\[
 \frac{1-c}{1-w_t}\mathcal E_t(f)
       \le\mathcal E_t^G(f)\le
                   \frac{1}{1-w_t}\mathcal E_t(f),
 \tag{V82.13}\label{c18v82:energy}
\]
and
\[
 |L_t^G(f)-L_t(f)|
       \le d_t\sqrt{\mathcal E_t(f)},\qquad
 d_t=E_0e^{tS/2}\sqrt v\,T .
 \tag{V82.14}\label{c18v82:sourcebound}
\]
These inequalities extend to the additive energy completion.
\end{proposition}
\begin{proof}
By (V82.4), conditioning separately on each \(X_\alpha\)
gives
\[
 \begin{split}
 \int_{\mathcal U}|df|_{\widehat m_t^{-1}}^2d\nu_t
 &\le6E_0^2\sum_\alpha
       \nu_t({\bf1}_{\mathcal U}
                         |\nabla h_\alpha(X_\alpha)|^2)\\
 &\le6E_0^2\delta\mathcal D_t(h)
 \le3E_0^4\delta\,\mathcal E_t(f).
 \end{split}
\]
Subtract this event energy and divide by \(1-w_t\) to
prove (V82.13). A small unconditioned event probability
alone would not justify the second line.

For each test choose its additive constants so that
\(\int h_\alpha q_\alpha=0\) at the displayed time.
Centering does not change its gradient or the pairing
with (V82.10). The actual marginal gap \(3e^{-tS}\) from
V76, Cauchy--Schwarz, and (V82.10) show
\[
 \begin{split}
 |L_t^G(f)-L_t(f)|
 &=\left|\sum_\alpha\int h_\alpha b_\alpha q_\alpha\,dX\right|\\
 &\le\frac{Te^{tS/2}}{\sqrt3}
           \sqrt{6v}\sqrt{\mathcal D_t(h)}
 \le E_0e^{tS/2}\sqrt v\,T\sqrt{\mathcal E_t(f)} .
 \end{split}
\]
The identities \(L_t(f)=\partial_t\nu_t(f)\) and
\(L_t^G(f)=\partial_t\nu_t^G(f)\) are used only for
time-independent tests; the displayed centering is a
constant chosen at that time. No derivative of a
time-varying test is silently omitted.

The two energies are equivalent on this space by
(V82.13), and the source difference is continuous
there. Passing to its completion proves the last claim.
There is no corresponding assertion for the full
nonadditive form domain.
\end{proof}

Define the complete additive variational responses
\[
 C_t=\sup_{f\in\mathcal F_{\rm add}}
                       \{2L_t(f)-\mathcal E_t(f)\},
 \qquad
 C_t^G=\sup_{f\in\mathcal F_{\rm add}}
                       \{2L_t^G(f)-\mathcal E_t^G(f)\}.
 \tag{V82.15}\label{c18v82:responses}
\]
Here \(C_t\) is V76's actual \(C_{\rm add}\). The second
is defined by a variational form; we do not assert that
a globally connected hard-conditioned manifold or a
uniform full Poisson inverse exists.

The square roots are dual norms of the respective
source functionals. Proposition \ref{c18v82:stability}
therefore gives
\[
 \sqrt{1-w_t}\,(\sqrt{C_t}-d_t)_+
 \le\sqrt{C_t^G}\le
     \sqrt{\frac{1-w_t}{1-c}}\,(\sqrt{C_t}+d_t).
 \tag{V82.16}\label{c18v82:dual}
\]
For clarity, first measure both \(L_t,L_t^G\) in the
\(\mathcal E_t\) dual norm: their distance is at most
\(d_t\), so apply the triangle and reverse-triangle
inequalities. Then compare the dual norms using
(V82.13). This proves (V82.16) without equating the
two inverse operators.

\subsection{An unconditional normalized-response error at logarithmic scale}

No uniform bound on \(C_t/v\) has yet been proved.
Nevertheless the preceding comparison gives an
\emph{absolute}, not merely relative, error after
normalization. The crude amplitude estimate already yields
\[
                      \sqrt{C_t/v}
                            \le E_0e^{tS/2}\mathfrak B .
 \tag{V82.17}\label{c18v82:crude}
\]
Indeed \(|\mathbb E[a_t\mid X_\alpha]|\le\mathfrak B\);
repeat the marginal-gap and private-spoke argument
in (V82.14), now with \(\mathfrak B\) in place of \(T\).
This estimate can grow linearly in \(v\), and is used
only to pay the small form perturbation.

\begin{theorem}[Source-preserving no-patch reduction]
\label{c18v82:reduction}
Take
\[
 \delta\le\delta_*:=\min\{1/2,(6E_0^4)^{-1}\},\qquad
 K_{82}=21E_0e^{2S}S(6E_0^4+18).
\]
Then, uniformly for \(0\le t\le4\),
\[
             \left|\sqrt{C_t^G/v}-\sqrt{C_t/v}\right|
                                      \le K_{82}\,v\delta .
 \tag{V82.18}\label{c18v82:error}
\]
For any fixed \(\sigma>0\), choose \(\ell_n\) to be the
least even integer at least
\[
 \max\left\{2,3B_\zeta,\
 \left[
  \frac{\log(C_{\rm obs}M_\varepsilon/\delta_*)
                         +(2+\sigma)\log v}{\kappa_{\rm loc}}
 \right]^{1/3}\right\}.
 \tag{V82.19}\label{c18v82:scale}
\]
For all sufficiently large even \(n\), \(\ell_n<n\),
\(\ell_n/n\to0\), and
\[
 \delta_{n,\ell_n}\le\delta_*v^{-1-\sigma},\qquad
 \sup_{0\le t\le4}
       \left|\sqrt{C_t^G/v}-\sqrt{C_t/v}\right|
                     \le K_{82}\delta_*v^{-\sigma}\longrightarrow0 .
 \tag{V82.20}\label{c18v82:limit}
\]
Thus a volume-uniform extensive additive-response bound
for the native law is equivalent to that bound for this
actual no-patch conditioned law on the same asymptotic
family. Neither bound is established by the reduction.
\end{theorem}
\begin{proof}
Put \(D_t=\sqrt{C_t/v}\), \(D_t^G=\sqrt{C_t^G/v}\),
and \(\epsilon_t=d_t/\sqrt v\).
Since \(c\le1/2\),
\((1-c)^{-1/2}\le1+2c\) and is at most \(\sqrt2\);
also \(\sqrt{1-w_t}\ge1-w_t\).
Equation (V82.16) implies
\[
           |D_t^G-D_t|\le2cD_t+\sqrt2\,\epsilon_t .
\]
For the lower-side error use \(w_t\le\delta\le2c\);
if \(D_t<\epsilon_t\), the same bound follows from
\(D_t^G\ge0\). Now (V82.17), \(t\le4\), and
\(T\le12\mathfrak B\delta\) give
\[
 |D_t^G-D_t|
 \le E_0e^{2S}\mathfrak B\delta
                              (6E_0^4+12\sqrt2)
 \le K_{82}v\delta .
\]
Here \(E_0\ge1\), as in the retained metric comparison.
This proves (V82.18) without presuming the sought
volume-uniform bound.

Substitute (V82.19) in (V82.3) to obtain
\(\delta\le\delta_*v^{-1-\sigma}\). Rounding upward to
an even integer preserves that inequality. The threshold
constants are fixed; hence \(\ell_n=O(1+(\log v)^{1/3})\)
and \(\ell_n/n\to0\) along even \(n\). Eventually all
V81 size conditions hold. Equation (V82.20) follows.
The vanishing difference proves the equivalence of the
two uniform bounds. A fixed block size would not give
this result as the number of translates grows.
\end{proof}

The event and its complement are globally conjugation
invariant, but the fixed forest and periodic pattern need
not preserve every channel-permuting symmetry. No equality
of all conditioned marginals was used. All \(6v\) channels,
and all their harmonics, are included in (V82.18).
The logarithmic block scale is an auxiliary thermodynamic
scale with extremely large constants, not a physical
continuum length.

\subsection{What has changed, and what has not}

This append supplies a genuine source-level reduction:
the source and metric effects of conditioning away
\emph{all} translated specified mesoscopic patches have
vanishing normalized additive inverse cost. V81's
unconditioned probability estimate alone did not give
the energy bound (V82.13), and its fixed-quotient relative
source rate did not give (V82.9). The private-chord
saturation estimate pays both missing distinctions.

The remaining law is not an independently sampled
good-block model. Its normalized weight is exactly
\({\bf1}_{\mathsf G}e^{u_t}/(1-w_t)\), and its source is
(V82.6). Hard exclusion may introduce nonsmooth boundaries
or disconnected regions. Integration-by-parts arguments
on that restricted space must account for any such
boundaries; no global spectral gap is inherited from
the one-holonomy marginal gap.

The next useful estimate is now on the actual conditioned
source: its complete distant source--cap covariance, or
an equivalent full additive source bound. The mere
absence of these particular periodic patches does not
supply favorable curvature or exhaust possible coherent
configurations. Further rare-family estimates are useful
only if they help prove that remaining estimate.
Nonadditive inverse control, probability strain, an
interacting-vacuum/all-coarse extension, physical scaling
and clock, and a nontrivial reflection-positive continuum
with positive physical mass gap all remain unproved.

The exact diagnostic
\begin{center}\small
\path{scripts/verify_ym_c18_v82_conditioned_response.py}
\end{center}
checks the native private-chord incidence on even tori,
including wrapped spokes, and the Haar-coordinate
bijection using noncommuting quaternion-group values.
It replays V76's pointwise energy algebra. A strictly
positive 512-state comparison law checks saturation,
observable conditioning, changing normalizers, and the
density-derivative bound using exact rational arithmetic.
The omitted normalizer changes the answer in that test.
A positive matrix fixture checks the dual variational
comparison. These are diagnostics, not simulations of
the nonlinear native probability, all-volume proof
certificates, or proof-assistant verification.

No new external theorem is imported. The V80 broad
literature and companion reviews remain the latest
scheduled checkpoints; the next are V90 and V85
respectively. The previous goal turn and this turn
are progress, while the full goal remains active.
Removing this append and restoring the title recovers
V81 byte for byte. Only the active filename is replaced;
only \texttt{.tex} files remain in \path{latex_notes}.
Archive/restart remains V100.
% END CYCLE 18 VERSION 82 APPEND
% BEGIN CYCLE 18 VERSION 83 APPEND
\clearpage
\section[V83: Surviving coherence and a native joint source]{V83: Surviving coherence and a native joint source}
\label{c18v83:context}

V82 paid the change of the complete \emph{additive} source
problem when V81's noncentral patches are excluded. It did
not prove that the remaining event has favourable curvature,
or control nonadditive observables. We now examine those two
limitations directly. First, the no-patch event still projects
onto \emph{every} assignment of the straight holonomies.
Second, the actual joint marginal has a nonadditive source
already at its first nonzero order. Both statements concern
the native normal-chart law, not an invented comparison law.

Throughout, \(\beta=0\), the coarse field is \(I\), the fine
side is \(2n\), and \(v=n^3\). The coverage statement uses
even \(n\ge6\) and V82's event. The joint-source identities
hold for every \(n\ge5\). All Taylor estimates below are
at a fixed regulator; their remainders are \emph{not}
volume-uniform. The full interacting continuum mass gap
remains unproved.

\subsection{The retained inputs and the actual new question}

We reuse V75's four-path notation at each two-odd centre
\(y\), and orient its two straight channels as
\[
 X_{y,1}=Z_0Z_2^{-1},\qquad X_{y,2}=Z_1Z_3^{-1},
 \qquad s_{y,r}=\operatorname{Re}X_{y,r}.
 \tag{V83.1}\label{c18v83:bank}
\]
There are \(3v\) centres and \(6v\) channels. Write \(X\)
for the entire bank and \(\mathsf P\) for conditional
expectation onto \(X\) under the \emph{reference} Haar
law \(\mu_I\). V78's Haar factorization and V82's private
chord construction are retained, not proved again as new
results.

To avoid conflict with V80's oscillation constant, put
\[
 S_{\rm osc}=D^2K_*,\qquad
 C=W_{\rm cap},\quad U=W_{\rm uncap},\quad W=C+U,\quad
 \mathcal N=|Q_0\nabla W|^2.
\]
Let \(u_t\) be the \emph{normalized} native log density,
\(L=L_0\), \(b=u'_0=3C\), and \(j=u''_0\).
V56--V57 and \eqref{c18v75:native} give
\[
 LC=9C,\quad LW=12W-3C,\qquad
 j=(L-12)\mathcal N+3\Gamma(W,C).
 \tag{V83.2}\label{c18v83:inputs}
\]
The vanishing second derivative of the native log
normalizer is part of this input. A second derivative
of an assumed tilt \(e^{3tC}\) would omit \(j\) and
give a different answer.

\subsection{A hidden chord that defeats every excluded patch}

Use exactly V80's first-odd-coordinate forest. Work with
axes numbered \(0,1,2\), and define
\[
 e_z=(2z+(1,1,0),\,2),\qquad
 z\in(\mathbb Z/n\mathbb Z)^3 .
 \tag{V83.3}\label{c18v83:hidden}
\]
These \(v\) distinct edges join two-odd vertices to
three-odd vertices. They are \(F_0\) chords, not forest
edges, and do not occur in any straight channel.

\begin{lemma}[A central quotient value independent of orientation]
\label{c18v83:reference}
For every unit imaginary quaternion \(q\), the quotient
chord of V79's field \(p(q)\) at every \(e_z\) is \(-I\).
\end{lemma}
\begin{proof}
The stored fine field is
\(U_{x,i}=q^{\sum_{j<i}x_j}\), with \(q^2=-I\).
For \(x=2z+(1,1,0)\) and \(w=x+(0,0,1)\), the forest
products from their coarse root are
\[
 k_x=(-1)^{z_0},\qquad k_w=(-1)^{z_1},\qquad
 U_{x,2}=(-1)^{z_0+z_1+1}.
\]
For example the path to \(w\) first traverses direction
\(2\), then \(1\), then \(0\), contributing exponents
\(2z_0+2z_1\), \(2z_0\), and zero. Thus
\(k_xU_{x,2}k_w^{-1}=-I\).
Even \(n\) makes these phase formulas consistent with
the periodic identifications. The result is independent
of both \(z\) and the axis \(q\).
\end{proof}

Let \(\mathcal U\) be V82's union of all translated
V81 events \(\mathcal N_R\), and \(G=\mathcal U^c\).
The event \(\mathcal N_R\) requires every chord in its
cube to be within \(\varepsilon\le\pi/8\) of some
\(p(q)\) quotient value. For \(0<\eta<\pi-\varepsilon\),
set
\[
 T_\eta=\{d(y_{e_z},I)<\eta\text{ for every }z\},
 \qquad h(\eta)=\frac{\eta-\sin\eta\cos\eta}{\pi}.
 \tag{V83.4}\label{c18v83:cylinder}
\]

\begin{theorem}[The no-patch event has full straight projection]
\label{c18v83:coverage}
One has \(T_\eta\subset G\), and the projection
\(G\longrightarrow\mathrm{SU}(2)^{6v}\) onto \(X\) is
surjective. In the canonical smooth disintegration,
for every \(x\) and \(0\le t\le4\),
\[
 \mu_I(T_\eta\mid X=x)=h(\eta)^v,\qquad
 \nu_t(G\mid X=x)\ge
       [e^{-tS_{\rm osc}}h(\eta)]^v>0 .
 \tag{V83.5}\label{c18v83:conditionalcoverage}
\]
In particular all channels can equal \(-I\), or any
other prescribed collection, on the surviving event.
\end{theorem}
\begin{proof}
Every nonempty translated coarse cube contains an
\(e_z\). Its reference value is \(-I\) by the lemma,
whereas points of \(T_\eta\) have distance greater than
\(\pi-\eta>\varepsilon\) from \(-I\) at that chord.
Thus no \(\mathcal N_R\) can occur, including cubes
crossing a torus seam.

In the forest section a straight channel is the product
of its two stored opposite spokes. At most one is a
forest edge. Choosing one private chord in every pair
and replacing it by that product is a Haar-preserving
bijection, with all remaining chords fixed. Different
channels use different selected chords. The \(14v\)
quotient variables therefore become \(6v\) independent
Haar variables \(X\) and \(8v\) Haar spectators.
There are \(5v\) \(F_0\) chords among those spectators,
including all \(e_z\). This is V78's factorization
after the \(7v\) internal gauge coordinates are removed.

Choose all distinguished chords equal to \(I\) and
solve the private product equation for each prescribed
\(X_\alpha=x_\alpha\). This proves surjectivity.
Product Haar gives the first equality in (V83.5).

At fixed \(X\), changing \(y_{e_z}\) changes no other
chord or observable. V80's single-chord oscillation
bound is therefore still \(tS_{\rm osc}\). The
conditional density of this chord, with all other
coordinates fixed, lies between \(e^{-tS_{\rm osc}}\)
and \(e^{tS_{\rm osc}}\) relative to Haar. Its conditional
probability of the radius-\(\eta\) ball is at least
\(e^{-tS_{\rm osc}}h(\eta)\). Integrating the other
spectators and conditioning successively on the
preceding distinguished balls preserves this bound.
Multiplying proves (V83.5). Positivity and smoothness
of the original density give a canonical disintegration
for every \(x\).
\end{proof}

The lower bound is exponentially small in \(v\).
It is a support statement, not a uniform probability
floor or a spectral inequality. No completed field is
declared Wilson stationary. It follows in particular
that exclusion of these patches alone cannot imply
a pointwise restriction on the straight bank that
fails at some bank assignment.

\subsection{The joint electric generator retains shared prefixes}

Let \(p\) range over the \(3v\) captured pairs, with
first-half variable \(A_p\). A left infinitesimal
rotation of \(A_p\) induces simultaneous left or
inverse-right rotations on its four incident channels.
Let \(Y_{p,a}\) be this vector field on the entire bank,
for an orthonormal imaginary-quaternion basis indexed
by \(a=1,2,3\). This definition fixes its signs using
the original path orientations.

\begin{lemma}[An exact reducing joint algebra]
\label{c18v83:generator}
The subspace of functions of \(X\) is reducing for
\(L\), and its positive generator is
\[
 \mathcal L_X=
  2\sum_{\alpha=1}^{6v}(-\Delta_{X_\alpha})
       -\frac12\sum_{p,a}Y_{p,a}^{\,2},\qquad
 \mathsf P L=\mathcal L_X\mathsf P .
 \tag{V83.6}\label{c18v83:L}
\]
For every smooth joint function \(F\),
\[
 2\sum_\alpha|\nabla_\alpha F|^2
 \le |d(F\circ X)|_{m_0^{-1}}^2
 \le 6\sum_\alpha|\nabla_\alpha F|^2 .
 \tag{V83.7}\label{c18v83:elliptic}
\]
\end{lemma}
\begin{proof}
Every free spoke occurs in only one channel, whose
two free-spoke Casimirs each give \(-\Delta_{X_\alpha}\).
With prefixes fixed, the adjoint rotations in their
derivatives preserve that Casimir and its energy.
The first-half factor has metric \(2g_{\rm round}\);
its inverse kinetic weight is \(1/2\). Its action on
every incident channel is precisely \(Y_{p,a}\).
These fields are Haar-divergence-free. This proves
(V83.6) on smooth cylinder functions, including all
cross derivatives generated by shared prefixes.

The private free-spoke energy gives the lower bound
in (V83.7). Each prefix has four incident channels;
Cauchy--Schwarz bounds the square of their sum by
four times the sum of their squares. Each channel
has two different prefixes. The prefix contribution
is consequently at most
\((1/2)\cdot4\cdot2\sum_\alpha|\nabla_\alpha F|^2\).
Adding the private contribution proves the upper bound.

The joint operator is symmetric and elliptic on the
compact product group. Uniqueness for its heat
equation shows that the full self-adjoint heat
semigroup preserves this subspace. Self-adjointness
then preserves its orthogonal complement as well,
which proves the projection identity.
\end{proof}

This is not a product Laplacian: the prefix term
couples different centres. At one centre, however,
the two opposite channels have four distinct prefixes
and four distinct free spokes. Consequently
\[
 \begin{gathered}
 H=\sum_{y,r}s_{y,r},\qquad K=\sum_y s_{y,1}s_{y,2},\\
 \mathcal L_X H=9H,\qquad \mathcal L_X K=18K .
 \end{gathered}
 \tag{V83.8}\label{c18v83:HK}
\]
The eigenvalue \(9\) is the two-spoke contribution
\(6\) plus the two half-weight prefix contributions
\(3\). The product at a single centre has no shared
factor between its channels, hence eigenvalue \(18\).
No independence of the electric derivatives at
different centres is assumed.

\subsection{Joint conditional moments, including the normal field}

\begin{lemma}[Complete conditional input at order two]
\label{c18v83:projection}
Under the reference projection onto the \emph{entire}
bank,
\[
 \begin{aligned}
 \mathsf PC&=0,&
 \mathsf PC^2&=3v+H+K,&
 \mathsf P(UC)&=0,\\
 \mathsf P\mathcal N&=9v-\tfrac32H,&
 \mathsf P\Gamma(W,C)&=27v+\tfrac92H,&
 \mathsf Pj&=-27v+18H .
 \end{aligned}
 \tag{V83.9}\label{c18v83:projections}
\]
\end{lemma}
\begin{proof}
Fix the prefixes first. The four paths at each centre
are independent Haar, and centres have separate free
spokes. Conditional on the two opposite products
\(A=X_{y,1}\) and \(B=X_{y,2}\), a common right
multiplication reduces the four paths, for invariant
scalar calculations, to
\[
 (Z_0,Z_1,Z_2,Z_3)=(AY,B,Y,I),\qquad Y\text{ Haar}.
\]
The remaining relative variables \(Y\) are independent
between centres. With \(c_r,T_y,\Xi_y,A_y\) as in V75,
\[
 \begin{aligned}
 C_y&=\operatorname{Re}[(B^{-1}+I)(A+I)Y],\\
 \mathsf P C_y&=0,&
 \mathsf P C_y^2&=(1+s_{y,1})(1+s_{y,2}),\\
 \mathsf P\Xi_y&=0,&
 \mathsf P A_y&=\tfrac12(s_{y,1}+s_{y,2}).
 \end{aligned}
 \tag{V83.10}\label{c18v83:localmoments}
\]
Indeed \(\int\operatorname{Re}(MY)^2\,dY=|M|^2/4\)
and quaternion norms multiply. Also
\(|I+A|^2=2(1+\operatorname{Re}A)\).
Each \(c_r^2\) has conditional mean \(1/4\).
The four adjacent products each give one quarter
of an opposite character, with each opposite
character occurring twice. These prove (V83.10).
Different-centre products \(C_yC_{y'}\) average to
zero, giving the first two identities in (V83.9).
Every uncaptured plaquette has free links incident
to a three-odd vertex. Those links are absent from
the whole straight bank and from \(C\); integration
of a private such link proves \(\mathsf P(UC)=0\).

For clarity, a projection onto just one centre would
not justify the claimed \(\mathsf P\mathcal N\).
Use instead V57's normal-copy decomposition itself.
Its diagonal and same-centre part, at centre \(y\),
is exactly
\[
                  3-\Xi_y+A_y-2T_y .
 \tag{V83.11}\label{c18v83:normalcentre}
\]
The remaining terms are contractions of captured
normal forces from different centres sharing a
captured pair. Fixing prefixes, each component of
the force from one centre is linear in that centre's
relative Haar variable \(Y\), so has conditional
mean zero even with both straight products fixed.
Relative variables at different centres are
independent. Thus every such cross contraction
vanishes under the full-bank projection. Equivalently,
one can set the prefixes to \(I\) by midpoint gauge
transformations: the normal energy is one half
the sum, over captured pairs, of the squared sum
of their two ambient half-link forces. Each
captured plaquette force contains one relative
\(Y\), and the preceding cancellation applies
componentwise. No UC normal term occurs.
Using (V83.10) in (V83.11) gives
\(3-\tfrac32T_y\), hence the fourth identity in (V83.9).

The positive-generator product rule gives
\[
 \Gamma(W,C)=\tfrac12\{21WC-3C^2-L(WC)\}.
\]
Apply (V83.6), \(\mathsf P(UC)=0\), and (V83.8):
\[
 \mathsf P\Gamma(W,C)
 =\tfrac12(18-\mathcal L_X)(3v+H+K)
 =27v+\tfrac92H .
\]
Finally insert the projected normal energy and this
identity into (V83.2). It gives \(-27v+18H\).
All constant contributions have been retained.
\end{proof}

\subsection{A native source invisible to all single-channel marginals}

Let \(q_{X,t}=\mathsf P(e^{u_t})\) be the actual normalized
joint density relative to product Haar \(dX\).
Its score is \(A_{X,t}=\partial_t\log q_{X,t}
=\nu_t(a_t\mid X)\).

\begin{theorem}[The first nonadditive joint source]
\label{c18v83:jointsource}
At zero,
\[
 \partial_t\log q_{X,0}=0,\qquad
 J_X:=\partial_t^2\log q_{X,0}=27H+9K .
 \tag{V83.12}\label{c18v83:J}
\]
At each fixed \(n\), for every fixed differentiability
order \(k\),
\[
 \log q_{X,t}=\tfrac12t^2(27H+9K)+O_{n,C^k}(t^4),
 \qquad A_{X,t}=t(27H+9K)+O_{n,C^k}(t^3).
 \tag{V83.13}\label{c18v83:Taylor}
\]
The projection of \(9K\) onto every individual
\(X_\alpha\) is zero.
\end{theorem}
\begin{proof}
Differentiate the normalized marginal, not a raw
exponential model. The two log derivatives are
\[
 \mathsf Pb,\qquad
 \mathsf Pj+\mathsf Pb^2-(\mathsf Pb)^2 .
\]
By (V83.9), the first is zero and the second is
\((-27v+18H)+9(3v+H+K)=27H+9K\).
Its Haar mean is zero, as normalization requires.

V55's sign involution fixes the prefixes and gives
the same sign to the two spokes in each opposite
pair, hence fixes every \(X_\alpha\). The exact
identity \(\widehat p_t\circ\sigma=\widehat p_{-t}\)
therefore makes \(q_{X,t}\) even. Smooth dependence
at the fixed compact regulator proves (V83.13).
Finally the other character in each product in
\(K\) is independent Haar with mean zero. Integrating
onto any single channel annihilates every term of
\(K\).
\end{proof}

Thus V75's marginal coefficient \(27s_\alpha\) is
recovered, not corrected. The new information is
that the native joint coefficient also contains
\(9K\). A collection of all one-channel marginal
identities cannot detect it. The coverage theorem
shows that neither \(K\) nor its coherent
configurations are removed as pointwise functions
by V82's exclusion. It does \emph{not} assert that
the conditioned law has the same coefficient
(V83.12); that law has its own normalizer and
conditional moments.

\subsection{Its exact leading joint variational response}

Define the actual full-bank response by
\[
 \mathcal C_X(t)=\sup_{F\in C^\infty(\mathrm{SU}(2)^{6v})}
 \left\{2\nu_t[a_tF(X)]
          -\nu_t[|d(F\circ X)|_{\widehat m_t^{-1}}^2]\right\}.
 \tag{V83.14}\label{c18v83:response}
\]
This admits all joint functions, not only additive
ones. It is a restriction of the full native
variational source inverse, and therefore a lower
bound for that full response.

\begin{proposition}[A positive mixed contribution]
\label{c18v83:responseprop}
At each fixed regulator,
\[
 \mathcal C_X(t)
 =\left(\frac{243}{2}+\frac{27}{32}\right)vt^2
       +O_n(t^4)
 =\frac{3915}{32}vt^2+O_n(t^4).
 \tag{V83.15}\label{c18v83:coefficient}
\]
The first summand is V75's inherited additive
coefficient. The second is the contribution of
\(9K\), and is strictly positive.
\end{proposition}
\begin{proof}
Push the actual quadratic form in (V83.14) onto
\(X\). Its coefficient tensor is the conditional
expectation of the contracted native inverse
metric. V51's metric comparison and (V83.7)
bound it between \(2E_0^{-2}I\) and \(6E_0^2I\).
At each fixed \(n,t\), its coefficients and the
positive density \(q_{X,t}\) are smooth. The
fixed-regulator mean-zero elliptic inverse thus
exists. Smooth perturbation of this inverse near
zero, or the corresponding coercive variational
equation on functions modulo constants, gives
\[
 \mathcal C_X(t)
  =t^2\langle J_X,\mathcal L_X^{-1}J_X\rangle_{dX}
      +O_n(t^3).
\]
This statement uses only a neighbourhood of zero
at each fixed \(n\). The involution also satisfies
\(\sigma^*\widehat m_t=\widehat m_{-t}\), so the
pushed form is even and the source functional
is odd. Its inverse response is even; smoothness
improves the remainder to \(O_n(t^4)\).

Independent product Haar and \(\int s^2=1/4\) give
\[
 \|H\|_2^2=\frac{6v}{4},\qquad
 \|K\|_2^2=\frac{3v}{16},\qquad
 \langle H,K\rangle=0 .
\]
Using (V83.8),
\[
 \mathcal L_X^{-1}J_X=3H+\tfrac12K,\qquad
 \langle J_X,\mathcal L_X^{-1}J_X\rangle
 =\frac{27^2}{9}\frac{6v}{4}
       +\frac{9^2}{18}\frac{3v}{16}.
\]
This is (V83.15). Since \(J_X\) consists entirely
of these eigenfunctions, no other joint harmonic
contributes at this leading order.
\end{proof}

The extra \(27v/32\) must not be added to V59's
already computed \emph{full} inverse expansion:
it is a decomposition of the joint-bank response,
which the full inverse already contains through
its own variational problem. Uniform ellipticity
of the kinetic tensor also does not give a
volume-uniform gap for its nonproduct probability.

\subsection{Checks, interpretation and the next upstream obligation}

The independent diagnostic
\begin{center}\small
\path{scripts/verify_ym_c18_v83_joint_source_coverage.py}
\end{center}
checks the forest phase, chord partition, arbitrary
prescribed quaternion bank completions, and prefix
incidence on native periodic tori \(n=6,8\).
It checks the ambient half-link normal forces at
24 centres in each of \(n=5,6\), including all
planes and boundary representatives, with five
rational unit-quaternion pairs. The eight quaternion
basis points give an exact Haar quadrature for the
degree-at-most-two relative polynomials used here.
All force means vanish and their squared averages
give (V83.11) after projection. Exact rational
algebra separately checks (V83.12) and (V83.15).
These are regression checks supporting the written
all-volume proofs, not a simulation of the native
finite-time law or a proof-assistant certificate.

Removing this append and restoring its title
recovers V82 byte for byte. The other 55 sources
remain unchanged. Only the newest active YM
filename is retained, and all build products stay
outside \path{latex_notes}. The preceding corpus
is preserved, not recertified line by line.
No new external theorem is imported. The last
scheduled companion and broad-literature reviews
are V80; the next are V85 and V90 respectively.
Archive/restart remains V100.

The direction changes in a specific way. Hard
exclusion of the previously studied noncentral
patches does not remove any straight-bank
configuration. Furthermore, the native source
on that surviving bank is not additive, even
at its first nonzero order. The next useful
estimate must retain this mixed source and its
normalization, and bound an actual finite-time
covariance or inverse response; another isolated
Taylor coefficient or another rare-family
exclusion would not supply that estimate.
V77's distant source-cap term, volume-uniform
source control through time four, Poisson Hessian
and probability strain, extension beyond coarse
identity and \(\beta=0\), physical scaling, and
nontrivial reflection-positive interacting
continuum reconstruction remain unpaid.
These finite-regulator results do not establish
the physical Yang--Mills mass gap.
% END CYCLE 18 VERSION 83 APPEND
% BEGIN CYCLE 18 VERSION 84 APPEND
\clearpage
\section[V84: Native Fisher budgets and temporal concentration]{V84: Native Fisher budgets and temporal concentration}
\label{c18v84:context}

V83 found an actual mixed source that all single-channel
marginals miss. A further centrewise Taylor expansion does
not control its finite-time remainder. We instead return
to the unnormalized ambient Jacobian and retain the
\emph{entire} temporal score. This yields two nonperturbative
budgets: a time-integrated Fisher bound on every fixed
coarse fibre, and a coarse-probability-averaged Fisher
bound at every chart time. Neither needs a conditional
spectral gap. The normalizer is the source of the exact
variance identity, not a term to discard.

The setting is the native \(\beta=0\) normal chart,
fine side \(2n\), \(n\ge5\), \(v=n^3\), and \(0\le t\le T=4\).
Unlike V83's coverage result, the estimates here hold
for \emph{every} fixed coarse field \(g\).
They do not prove pointwise Fisher control at \(g=I,t=4\),
an unshifted inverse estimate, or the physical YM mass gap.
We identify explicitly how concentration can remain.

\subsection{The full ambient law and its fixed-fibre normalization}

Use V51's fields on the fine configuration manifold
\(\mathcal M\), with its unit-round product metric and
normalized Haar measure \(\mu_{\rm f}\):
\[
 \begin{gathered}
 F=\nabla W,\quad c_t=b_2\Phi_t,\quad C_t=Dc_t,\quad
 R_t=C_t^*(C_tC_t^*)^{-1},\\
 Q_t=R_tC_t,\quad P_t=I-Q_t,\quad
 N_t=Q_tF,\qquad \dot\eta_t=-N_t\circ\eta_t .
 \end{gathered}
 \tag{V84.1}\label{c18v84:normal}
\]
Here \(N_t\) denotes the normal vector field, not the
fine side length. V51 proves that \(\eta_t\) is a global
ambient diffeomorphism with \(c_t\eta_t=b_2\).

Let
\[
 \begin{gathered}
 k_t=\operatorname{Jac}_{\mu_{\rm f}}\eta_t,\qquad
 r_t=\log k_t,\qquad d_t=\operatorname{div}_{\mu_{\rm f}}N_t,\\
 Z_g(t)=\int_{M_g}k_t\,d\mu_g,\qquad
 \nu_{g,t}=\frac{k_t}{Z_g(t)}\mu_g,\qquad
 a_{g,t}=\dot r_t-\partial_t\log Z_g(t).
 \end{gathered}
 \tag{V84.2}\label{c18v84:laws}
\]
The restriction of \(r_t\) to \(M_g\) is understood in
the second line. These are precisely the native
conditional probability and score, not a Wilson
exponential ansatz. We keep \(r_t\) raw and reserve
\(r_t-\log Z_g\) for the normalized log density.

Product Haar disintegrates as \(d\mu_{\rm f}=d\mu_g\,dg\)
under \(b_2\): each coarse link is the product of a
separate fine pair, whose first half remains Haar
after conditioning on that product. Therefore
\[
 \bar\nu_t=k_t\mu_{\rm f},\quad
 (\eta_t)_*\bar\nu_t=\mu_{\rm f},\quad
 \int Z_g(t)\,dg=1,\quad Z_g(0)=1 .
 \tag{V84.3}\label{c18v84:ambient}
\]
The coarse probability in this identity is \(Z_g(t)\,dg\),
not Haar at positive time. The raw temporal identities are
\[
 \dot r_t=-d_t\circ\eta_t,\qquad
 \ddot r_t=-(\partial_t d_t-N_td_t)\circ\eta_t .
 \tag{V84.4}\label{c18v84:raw}
\]

\subsection{A time derivative of the normal field from spatial geometry}

The following identity avoids assuming a new conditional
inverse or a new temporal regularity estimate.
All covariant derivatives in this subsection are ambient.

\begin{lemma}[Exact normal-projector transport]
\label{c18v84:transport}
Put \(H=\nabla F=\operatorname{Hess}W\). Then
\[
 \partial_tQ_t-\nabla_FQ_t=P_tHQ_t+Q_tHP_t,\qquad
 \partial_tN_t-\nabla_FN_t=(I-2Q_t)HN_t .
 \tag{V84.5}\label{c18v84:transportidentity}
\]
\end{lemma}
\begin{proof}
Autonomy gives \(\partial_tc_t=C_tF\).
Covariantly differentiating this map identity in a
domain direction \(V\), and using symmetry of the
second covariant differential of \(c_t\), gives
\[
 D_tC_t(V)=(\nabla_FC_t)(V)+C_t(\nabla_VF).
\]
The target covariant derivative is used on the left.
Target orthonormal-frame rotations do not change
\(Q_t\). The derivative of the orthogonal projector
associated with any full-row-rank \(C\), in direction
\(\dot C\), is
\[
 \dot Q=P\dot C^*R^*+R\dot C P .
\]
It follows either by differentiating
\(Q=C^*(CC^*)^{-1}C\), including the inverse derivative,
or directly from \(Q^2=Q\) and its off-diagonal blocks.
Apply it first to \(D_tC\), then to \(\nabla_FC\).
Their difference is \(CH\); since \(H=H^*\) and \(RC=Q\),
the first identity follows. Finally
\(\nabla_F(QF)=(\nabla_FQ)F+QHF\).
Subtract this from \((\partial_tQ)F\), use \(P+Q=I\),
and collect \(PHN-QHN=(I-2Q)HN\).
\end{proof}

We now pay a bound for (V84.4) with existing constants.
Let \(c_j,r_j,Q_j^\sharp,n_j\) be exactly V51's constants
in \eqref{c18v51:constants}. Rename the numerical
majorants for clarity:
\[
 q_0=Q_0^\sharp,\quad q_1=Q_1^\sharp,\quad
 N_0=n_0,\quad N_1=n_1,\quad N_2=n_2 .
\]
Thus \(\mathfrak a_0(\partial^jN_t)\le N_j\) for \(j\le2\),
and \(\mathfrak a_0(Q_t)\le q_0\),
\(\mathfrak a_0(\partial Q_t)\le q_1\).
These are scalar anchored all-slot norms, with the
configuration supremum outside their sums. For a vector
with just one slot the norm is its component supremum.
The elementary native plaquette bounds are
\(\mathfrak a_0(\partial^jF)\le4\cdot12^j\), \(j=0,1,2\).

Define the following deliberately loose constants:
\[
 \begin{aligned}
 k_1&=N_1+4N_0,&
 k_2&=N_2+16N_1+64N_0,\\
 h_0&=48+4\cdot4=64,&
 h_1&=576+16\cdot48+64\cdot4=1600,\\
 q_1^{\rm cov}&=q_1+8q_0,& j_0&=1+2q_0,\\
 W_1&=4k_2+k_1h_0+2q_1^{\rm cov}h_0N_0
                +j_0h_1N_0+j_0h_0k_1 .
 \end{aligned}
 \tag{V84.6}\label{c18v84:constants}
\]
They are uniform in \(n,g,t\), though extremely large.

\begin{lemma}[Paid raw temporal derivatives]
\label{c18v84:rawbounds}
With
\[
 A=72N_1,\qquad B=72(W_1+N_0N_2),
 \tag{V84.7}\label{c18v84:AB}
\]
one has globally
\[
             |\dot r_t|\le Av,\qquad
             |\ddot r_t|\le Bv .
 \tag{V84.8}\label{c18v84:rawbound}
\]
\end{lemma}
\begin{proof}
In invariant unit-round frames the connection coefficients
are constant and their all-slot norm is at most four
(in fact two suffices). Passing from an ordered first
coefficient derivative of a vector to its covariant
derivative costs at most four times the vector norm.
For the second derivative the three first-order
connection terms and two quadratic terms give the
bound \(N_2+12N_1+32N_0\), which is at most \(k_2\).
The same calculation for \(F\) gives the bounds \(h_0\)
for \(H\) and \(h_1\) for \(\nabla H\). A covariant
derivative of the two-slot tensor \(Q\) costs at most
eight times its zero-order norm, giving \(q_1^{\rm cov}\).

Differentiate the second identity in (V84.5) once
in space. Its five terms are bounded respectively by
\[
 k_2\cdot4,\quad k_1h_0,\quad
 2q_1^{\rm cov}h_0N_0,\quad j_0h_1N_0,\quad j_0h_0k_1 .
\]
These are the derivative of \(\nabla_FN\) in its
two arguments, and the three product derivatives of
\((I-2Q)HN\). Anchored contraction retains an external
slot in every term, so there is no volume factor in
\(\mathfrak a_0(\nabla\partial_tN)\le W_1\).

There are \(24v\) fine links and \(72v\) scalar
fine-frame directions. Haar divergence is
\(d_t=\sum_\alpha X_\alpha N_{t,\alpha}\), so
\(|d_t|\le72vN_1\).
The complete trace of \(\nabla\partial_tN\) is at most
\(72vW_1\). Also each differentiated divergence satisfies
\(|X_\beta d_t|\le N_2\): its trace contraction still
has the external derivative anchor \(\beta\).
Consequently \(|N_td_t|\le72vN_0N_2\).
Use (V84.4). The trace dimension is \(72v\), not the
\(63v\) dimension of a fixed coarse fibre.
\end{proof}

\subsection{A complete Fisher budget on every fixed coarse fibre}

Define the full temporal Fisher information and the
conditional log normalizer per coarse volume by
\[
 \mathcal I_{n,g}(t)=\nu_{g,t}(a_{g,t}^2),\qquad
 \phi_{n,g}(t)=v^{-1}\log Z_g(t).
 \tag{V84.9}\label{c18v84:fisher}
\]

\begin{theorem}[Native time-integrated source control]
\label{c18v84:time}
For every \(g,n\), with no conditional-gap hypothesis,
\[
 \begin{gathered}
 |\phi'_{n,g}(t)|\le A,\qquad
 \phi''_{n,g}(t)
    =v^{-1}\nu_{g,t}(\ddot r_t)
                   +v^{-1}\mathcal I_{n,g}(t)\ge-B,\\
 \int_0^T\mathcal I_{n,g}(t)\,dt
                      \le vK_T,\qquad K_T=A+BT,\quad T=4 .
 \end{gathered}
 \tag{V84.10}\label{c18v84:timebudget}
\]
More generally, for \(0\le s<t\le T\),
\(\int_s^t\mathcal I_{n,g}(r)\,dr
 \le v[2A+B(t-s)]\).
\end{theorem}
\begin{proof}
Differentiate the normalized integral defining \(Z_g\):
\[
 \partial_t\log Z_g=\nu_{g,t}(\dot r_t),\qquad
 \partial_t^2\log Z_g
       =\nu_{g,t}(\ddot r_t)
             +\operatorname{Var}_{\nu_{g,t}}(\dot r_t).
\]
The variance is precisely \(\mathcal I_{n,g}\); subtracting
the mean in (V84.2) is essential. Equation (V84.8)
gives the derivative bounds and semiconvexity.

At zero, V51--V56's raw identity is
\(\dot r_0=3W_{\rm cap}\). Each captured plaquette has
a private free Haar factor at every coarse field, hence
\(\phi'_{n,g}(0)=0\). Integrating the displayed second
derivative identity from zero to \(T\), bounding the
endpoint derivative by \(A\) and the expected raw second
derivative below by \(-Bv\), gives \(K_T\).
For a general interval bound the two endpoint derivatives
separately by \(A\).
\end{proof}

This controls the full score, not a sum of one-factor
conditional scores and not only its rare part. At zero,
the inherited checks are
\[
 \mathcal I_{n,g}(0)=27v,\qquad
 \nu_{g,0}(\ddot r_0)=-27v,\qquad \phi''_{n,g}(0)=0 .
\]
They illustrate why the raw second derivative cannot
be dropped. They are consistency checks, not the proof
of the all-time integral.

\subsection{A complementary bound at every time, averaged over coarse fields}

\begin{theorem}[Exact coarse/conditional Fisher decomposition]
\label{c18v84:coarse}
At every \(0\le t\le4\),
\[
 \begin{split}
 &\int Z_g(t)
       \left[\mathcal I_{n,g}(t)
                   +|\partial_t\log Z_g(t)|^2\right]dg\\
 &\hspace{12mm}=\int_{\mathcal M}d_t^2\,d\mu_{\rm f}
                       \le vK_{\rm amb},\\
 &\hspace{12mm}K_{\rm amb}=72(k_1^2+2N_0^2).
 \end{split}
 \tag{V84.11}\label{c18v84:coarsebudget}
\]
In particular the coarse-probability average of the
complete conditional Fisher information is bounded
by \(vK_{\rm amb}\).
\end{theorem}
\begin{proof}
The full ambient score is \(\dot r_t\), of mean zero
under \(\bar\nu_t\). Conditional on \(b_2=g\), its mean
is \(\partial_t\log Z_g\) and its centered part is
\(a_{g,t}\). The conditional variance identity and
(V84.3)--(V84.4) give exactly the first equality:
\[
 \int(\dot r_t)^2\,d\bar\nu_t
       =\int d_t^2\,d\mu_{\rm f}.
\]
This change of variables is ambient. It does not
replace a fixed-fibre conditional probability by Haar.

For any smooth vector field \(N\) on a closed Riemannian
manifold with its normalized volume measure, integration
by parts and the Ricci commutator give
\[
 \int(\operatorname{div}N)^2
 =\int\left[
    \sum_{\alpha,\beta}
       (\nabla_\alpha N^\beta)(\nabla_\beta N^\alpha)
                   +\operatorname{Ric}(N,N)\right].
 \tag{V84.12}\label{c18v84:dividentity}
\]
Indeed integrate \(-N^\alpha\nabla_\alpha\operatorname{div}N\),
commute the two covariant derivatives, then integrate
the derivative in the other index. This is the same
signed divergence identity used in V30, now applied to
the actual temporal normal field under ambient Haar.
The derivative contraction is not generally a squared
Hilbert--Schmidt norm; it is bounded above by that norm.
On the fine product \(\operatorname{Ric}=2\mathfrak g\).
The anchored bound gives
\(\|\nabla N_t\|_{\rm HS}^2\le72v k_1^2\) and
\(|N_t|^2\le72vN_0^2\). Substitute these in (V84.12).
\end{proof}

The probability weight \(Z_g(t)\) cannot be removed.
In particular (V84.11) does not bound the value at the
single prescribed coarse field \(I\). It bounds both
the coarse score and the conditional score only after
the displayed coarse averaging. This differs from
V30's spatial/coarse source-bank Fisher estimate.

\subsection{Joint sources and shifted inverses are included}

Let \(X\) be any time-independent measurable observable
on \(M_g\), including the \emph{entire} V83 straight
bank at \(g=I\). Its exact marginal score is
\(A_{X,t}=\nu_{g,t}(a_{g,t}\mid X)\), in the
Radon--Nikodym sense if no smooth marginal coordinates
are selected. Conditional Jensen gives
\[
 \mathcal I_{X,n,g}(t)
     :=\nu_{g,t}(A_{X,t}^2)\le\mathcal I_{n,g}(t).
 \tag{V84.13}\label{c18v84:processing}
\]
Differentiation against bounded time-independent tests
proves the score identity at each finite regulator.
Thus both preceding averaged budgets retain every
joint component, including V83's mixed term.
This does not bound a \emph{sum} of overlapping
individual Fisher projections by one copy of
\(\mathcal I_{n,g}\).

For \(\lambda>0\), retain the actual moving-metric
generator \(\widehat L_{g,t}\) and define
\[
 \mathcal C_\lambda(n,g,t)
   =\langle a_{g,t},
          (\lambda+\widehat L_{g,t})^{-1}a_{g,t}\rangle_{\nu_{g,t}} .
\]
Nonnegativity of its quadratic form, without any
uniform positive gap, gives
\[
 \begin{gathered}
 0\le\mathcal C_\lambda\le\lambda^{-1}\mathcal I_{n,g},\\
 \int_0^T\mathcal C_\lambda\,dt\le vK_T/\lambda,\qquad
 \int Z_g(t)\mathcal C_\lambda\,dg
                             \le vK_{\rm amb}/\lambda .
 \end{gathered}
 \tag{V84.14}\label{c18v84:shifted}
\]
For example maximize
\(2\nu(af)-\mathcal E(f)-\lambda\nu(f^2)\)
and discard only the nonnegative energy to prove
the first bound. Restricting the same variational
problem to arbitrary joint \(F(X)\) preserves it.

The shift is in the auxiliary generator's units.
It is not a physical mass gap. If \(d\rho_{n,g,t}(E)\)
is the actual source spectral measure, the still
unpaid difference from the unshifted response is
\[
 \mathcal C_0-\mathcal C_\lambda
 =\int_{(0,\infty)}
        \frac{\lambda}{E(E+\lambda)}\,d\rho_{n,g,t}(E).
 \tag{V84.15}\label{c18v84:lowenergy}
\]
At fixed regulator this is finite by compact ellipticity.
The Fisher bounds control the measure's total mass,
not this inverse-weighted low-energy mass uniformly.
No \(\lambda\downarrow0\) limit of (V84.14) is bounded here.

\subsection{The remaining temporal concentration has a precise form}

The time budget immediately implies
\[
 \left|\{t\in[0,T]:\mathcal I_{n,g}(t)>Rv\}\right|
                         \le K_T/R,\qquad R>0 .
 \tag{V84.16}\label{c18v84:exceptional}
\]
It gives no regulator-uniform pointwise bound at a
specified time. A limit statement locates the possible
singular part more precisely.

\begin{proposition}[Source-concentration measures and the normalizer]
\label{c18v84:concentration}
For any sequence of regulators and coarse fields,
a subsequence has
\[
 \phi_{n,g}\longrightarrow\phi\ \hbox{uniformly on }[0,T],
 \qquad v^{-1}\mathcal I_{n,g}(t)\,dt
                         \ \rightharpoonup\ \sigma ,
\]
where \(\sigma\) is a positive finite measure of mass
at most \(K_T\), and \(\phi+Bt^2/2\) is convex.
There is a function \(e\in L^\infty(0,T)\),
\(|e|\le B\), such that, in distributions on \((0,T)\),
\[
                     \sigma=\phi''-e(t)\,dt .
 \tag{V84.17}\label{c18v84:defect}
\]
The singular part of \(\sigma\) is exactly the singular
part of \(\phi''\). At any interior \(t_0\),
\[
             \sigma(\{t_0\})
                       =\phi'_+(t_0)-\phi'_-(t_0).
 \tag{V84.18}\label{c18v84:atom}
\]
\end{proposition}
\begin{proof}
The pressure functions vanish at zero and are uniformly
\(A\)-Lipschitz by (V84.10), so a subsequence converges
uniformly. Semiconvexity passes to the limit.
The positive Fisher measures have uniformly bounded
mass on a compact interval, so extract a weakly
convergent subsequence. Finally
\(e_{n,g}=v^{-1}\nu_{g,t}(\ddot r_t)\) is bounded in
\(L^\infty\); extract a weak-star subsequence as well.
Against every smooth compactly supported temporal
test, pass to the limit in
\(v^{-1}\mathcal I_{n,g}=\phi''_{n,g}-e_{n,g}\).
This proves (V84.17).
The second distributional derivative of a convex
function is the measure of increments of its
monotone first derivative. Subtracting the smooth
quadratic and the absolutely continuous \(e\,dt\)
does not affect its singular part, and a first
derivative jump has mass (V84.18).
\end{proof}

No unique thermodynamic limit, absolute continuity
of \(\sigma\), or convergence of Fisher values at
each time has been proved. A singular continuous
part is also possible at this level of control.
The formula for atoms concerns interior times;
boundary concentration at \(T=4\) is not ruled out
by that distributional identity.

Here is an explicit comparison illustrating this
logical limitation, not a native YM law. On two
states \(\epsilon=\pm1\), with uniform reference,
let \(m\) be a volume parameter and put
\[
 r_m(t,\epsilon)=m h(t)\epsilon,\qquad
 h(t)=t^2(t^2-1),\qquad
 \phi_m(t)=m^{-1}\log\cosh(mh(t)).
\]
Its raw first and second derivatives are bounded
by constants times \(m\) on \([0,4]\), its initial
density is uniform, and \(\phi_m'(0)=0\). Nevertheless
\[
 \frac{\mathcal I_m(t)}m
       =m h'(t)^2\operatorname{sech}^2(mh(t)),\qquad
 \frac{\mathcal I_m(1)}m=4m .
 \tag{V84.19}\label{c18v84:comparison}
\]
Uniformly \(\phi_m\to|h|\). Near \(t=1\), the change
of variable \(z=mh(t)\), where \(h'(1)=2\), shows
that the Fisher measure converges to \(4\delta_1\).
The peak therefore grows while its integral stays
bounded, precisely as the derivative jump of
\(|h|\) predicts. This example does not assert a
native phase transition or the ambient identity
(V84.11) for this comparison law.

\subsection{Verification, literature scope and direction}

The diagnostic
\begin{center}\small
\path{scripts/verify_ym_c18_v84_fisher_budgets.py}
\end{center}
checks the differentiated full projector and normal
transport identity with exact rational noncommuting
matrix data. Exact \(S^3\) integration checks the
divergence identity, including a Killing-field test
that distinguishes its signed contraction from a
Hilbert--Schmidt square. A positive finite correlated
law checks the complete normalization derivatives,
coarse/conditional Fisher chain rule and data processing.
Native periodic incidence at \(n=5,6\) checks the
fine/fibre dimensions and the zero-time source
normalization. High-precision integration checks the
two-state concentration example. The fixtures are
explicitly not simulations of the native finite-time
probability; the all-volume bounds are proved above.

For primary context, G.~E.~Crooks,
\href{https://arxiv.org/abs/0706.0559}
{\emph{Measuring thermodynamic length}},
connects thermodynamic length and Fisher information.
Only the primary abstract was checked in this
targeted review. No theorem about dissipation,
mixing or YM is imported. The present temporal
variance identity, geometric constants and source
bounds are derived directly. This targeted check
does not replace the next scheduled broad sweep.

The previous goal turn was progress; this turn
adds control of the full native source throughout
the chart under the stated averages. It does not
claim a new long-time pointwise source inverse.
The next useful task is to prevent the actual
source from concentrating at the required time
and coarse field, or to control its low-energy
spectral weight directly in (V84.15). The normalizer's
semiconvex limit identifies one concentration
mechanism without excluding it. These are separate
from the physical scaling and clock, all-coupling
vacuum control, and nontrivial reflection-positive
interacting continuum construction still required
for the full YM mass gap.

Removing this append and restoring the title
recovers V83 byte for byte. The other 55 sources
are unchanged; only the newest active YM filename
is retained. The old corpus is preserved, not
recertified line by line. All build products remain
outside \path{latex_notes}, which contains only
\texttt{.tex}. The last scheduled companion and
broad reviews were V80; the next are V85 and V90.
Archive/restart remains V100. The goal remains active.
% END CYCLE 18 VERSION 84 APPEND
% BEGIN CYCLE 18 VERSION 85 APPEND
\clearpage
\section[V85: Exponential source budgets and conditioned inverse tails]{V85: Exponential source budgets and conditioned inverse tails}
\label{c18v85:context}

V84 controls the full native score in mean square under two
different averages. We now keep the exact ambient Haar
representation and obtain an exponential square moment.
This estimates all amplitudes of the actual source, including
its joint components and its conditional normalizer, rather
than selecting another geometric droplet family. The new
moment bound can absorb V72's already-paid surface-cost
inverse on a precisely specified large-amplitude tail.

The setting remains the native \(\beta=0\) normal chart,
\(n\ge5\), \(v=n^3\), and \(0\le t\le4\). The exponential
estimate is averaged over the \emph{actual} coarse probability.
It is not a pointwise assertion at \(g=I\). The unshifted
inverse of the typical source, its volume-uniform control,
and the physical continuum mass gap remain open.
Here an \emph{uncontrolled inverse} means the required
volume-uniform bound: V72 already supplies a deteriorating
finite-volume bound and is not redone here.

\subsection{Scheduled companion review and inherited inputs}

The twelve active companion sources listed in the V80 review
were inventoried and hashed again. None has changed.
Their terminal mathematical scope and outstanding source
obligations were rechecked, not recertified cumulatively.
The SM source-selection and arithmetic notes retain their
actual-occurrence requirements; the stationarity and
Einstein-bridge notes retain residual/source and slow-tangency
obligations. The ESM notes retain their fixed-order and
order-of-limits restrictions. NS supplies no new YM mixing
constant. The shell-mixing note still needs its mixed-curvature
assembly; its flat potential does not supply a probability gap.

The useful transfer is to preserve the actual probability
denominator and quantify the information cost of further
conditioning. No companion spectral margin or physical clock
is imported. The review also located V72's whole-torus
inverse, preventing a duplicate canonical-path argument.
V63's weighted Hessian bounds, V51's metric comparison,
and V84's ambient Jacobian identity are the inputs below.
The next companion and broad-literature checkpoints are both
V90; archive/restart remains V100.

Write \(k_t,r_t,Z_g,\nu_{g,t},a_{g,t},d_t,\eta_t\) exactly as in
\eqref{c18v84:laws}. Define the actual coarse probability
\[
                  d\zeta_{n,t}(g)=Z_g(t)\,dg .
\]
For a fibre quantity \(F_g\), an integral against
\(\zeta_{n,t}\) will always be displayed explicitly.
The ambient probability is \(\bar\nu_t=k_t\mu_{\rm f}\);
\((\eta_t)_*\bar\nu_t=\mu_{\rm f}\) and
\(\dot r_t=-d_t\circ\eta_t\).
Use V84's \(N_2\), so that
\[
 \sup_{\alpha,U,t}|X_\alpha d_t(U)|\le N_2,\qquad
 \int d_t\,d\mu_{\rm f}=0 .
 \tag{V85.1}\label{c18v85:rawgradient}
\]
The first follows from the anchored differentiated trace;
the second is integration of a divergence on the closed
fine product. Put
\[
          \Lambda=18\pi^2N_2^2,\qquad
          \Theta=16\Lambda=288\pi^2N_2^2 .
 \tag{V85.2}\label{c18v85:constants}
\]
These are positive volume-independent majorants, not
claims of small numerical constants.

\subsection{Product concentration before disintegration}

\begin{lemma}[A native ambient exponential square moment]
\label{c18v85:ambient}
For every real \(s\), every \(n\), and every \(t\in[0,4]\),
\[
 \int e^{s d_t}\,d\mu_{\rm f}
       \le e^{\Lambda v s^2/2},\qquad
 \int \exp\!\left(\frac{d_t^2}{4\Lambda v}\right)d\mu_{\rm f}
       \le\sqrt2 .
 \tag{V85.3}\label{c18v85:ambientmoment}
\]
The same inequalities hold for \(\dot r_t\) under
\(\bar\nu_t\).
\end{lemma}
\begin{proof}
There are \(24v\) independent fine links under Haar,
each a unit-round three-sphere of diameter \(\pi\).
By \eqref{c18v85:rawgradient} its three-color gradient
has norm at most \(\sqrt3N_2\).
Changing one whole link changes \(d_t\) by at most
\(c=\pi\sqrt3N_2\). Reveal the links successively.
The Doob martingale increments have conditional mean
zero and conditional range length at most \(c\):
averaging unrevealed independent links preserves the
one-link oscillation bound.

For completeness, if a centered variable \(Y\) lies
in an interval of length \(c\), set
\(q(s)=\log\mathbb E e^{sY}\).
Every exponentially tilted law has the same support,
so \(q''(s)=\operatorname{Var}_s(Y)\le c^2/4\);
the latter follows by comparison with the midpoint
of the support interval. Since \(q(0)=q'(0)=0\),
integration gives \(q(s)\le c^2s^2/8\) for both signs
of \(s\). Apply this conditional inequality to each
martingale increment and iterate conditional expectations.
The result is
\[
 \log\int e^{s d_t}\,d\mu_{\rm f}
       \le (24v)c^2s^2/8
       =9\pi^2N_2^2vs^2=\Lambda vs^2/2 .
\]
Centering uses the second identity in (V85.1), not an
assumption on a fixed conditional fibre.

Let \(G\) be an independent standard real Gaussian.
Its elementary Gaussian integral gives
\(\mathbb E_G e^{Gx/\sqrt{2\Lambda v}}
=e^{x^2/(4\Lambda v)}\).
Tonelli and the first inequality then bound the second
integral in (V85.3) by \(\mathbb E e^{G^2/4}=\sqrt2\).
Finally change variables by \(\eta_t\).
Haar concentration is used before conditioning; no
independence of the conditional fibre coordinates
at positive time has been asserted.
\end{proof}

This is the classical bounded-differences argument,
with its native oscillation constant paid above.
For attribution the primary publisher record of
C.~McDiarmid,
\href{https://www.cambridge.org/core/books/abs/surveys-in-combinatorics-1989/on-the-method-of-bounded-differences/AABA597B562BDA7D89C6077E302694FB}
{\emph{On the method of bounded differences}},
in \emph{Surveys in Combinatorics, 1989}, was checked.
The proof needed here is included in full; no uninspected
chapter theorem or conditional mixing claim is invoked.
This targeted check is not the scheduled V90 broad sweep.

\subsection{Centering costs a constant, not an unknown inverse}

\begin{theorem}[Exponential control of the complete conditional score]
\label{c18v85:conditional}
At every fixed \(t\in[0,4]\),
\[
 \begin{split}
 &\int d\zeta_{n,t}(g)\,
       \nu_{g,t}\!\left[
          \exp\!\left(\frac{a_{g,t}^2}{\Theta v}\right)\right]
             \le\sqrt2,\\
 &\int \exp\!\left(\frac{\mathcal I_{n,g}(t)}
                              {\Theta v}\right)d\zeta_{n,t}(g)
             \le\sqrt2 .
 \end{split}
 \tag{V85.4}\label{c18v85:scoremoment}
\]
In particular, for \(R>0\) and every integer \(p\ge1\),
\[
 \begin{split}
 \zeta_{n,t}\{\mathcal I_{n,g}(t)>Rv\}
       &\le\sqrt2\,e^{-R/\Theta},\\
 \int \mathcal I_{n,g}(t)^p\,d\zeta_{n,t}(g)
       &\le\sqrt2\,p!(\Theta v)^p .
 \end{split}
 \tag{V85.5}\label{c18v85:fishertails}
\]
\end{theorem}
\begin{proof}
Set \(b=\dot r_t\) and \(m_g=\nu_{g,t}(b)\), so that
\(a=b-m_g\). On one fixed conditional fibre draw \(b'\)
as an independent copy of \(b\). Conditional Jensen,
\((b-b')^2\le2b^2+2(b')^2\), and Cauchy--Schwarz give
\[
 \begin{split}
 \nu_g e^{(b-m_g)^2/(16\Lambda v)}
 &\le(\nu_g\otimes\nu_g)e^{(b-b')^2/(16\Lambda v)}\\
 &\le\left(\nu_g e^{b^2/(8\Lambda v)}\right)^2
 \le\nu_g e^{b^2/(4\Lambda v)} .
 \end{split}
\]
Integrate against \(\zeta_{n,t}\) and use (V85.3) and
the exact disintegration of \(\bar\nu_t\).
The second line of (V85.4) is another conditional
Jensen inequality, since \(\mathcal I_{n,g}=\nu_g(a^2)\).
Markov's inequality and \(x^p\le p!e^x\) for \(x\ge0\)
prove (V85.5).
\end{proof}

For any fixed joint observable \(X\) on a fibre, the
exact marginal score is \(\nu_g(a\mid X)\).
The convexity of \(x\mapsto e^{x^2/(\Theta v)}\)
shows that it satisfies the first bound in (V85.4)
as well. This includes the entire V83 bank, not only
its additive channels. It does not justify adding
overlapping projection budgets.

The bound is uniform in the displayed time, but
\(\zeta_{n,t}\) changes with time. It neither supplies
one exceptional coarse set valid for all times nor
rules out V84's concentration for a prescribed
sequence of coarse fields.

\subsection{An explicit information cost for choosing coarse fields}

Let \(\rho\ll\zeta_{n,t}\) be any probability on coarse
fields and let \(H(\rho\mid\zeta_{n,t})\) denote relative
entropy, with value \(+\infty\) if necessary. Then
\[
 \int\mathcal I_{n,g}(t)\,d\rho(g)
 \le\Theta v\left[H(\rho\mid\zeta_{n,t})
                              +\tfrac12\log2\right].
 \tag{V85.6}\label{c18v85:entropy}
\]
Indeed tilt \(\zeta_{n,t}\) by
\(e^{\mathcal I/(\Theta v)}\), divide by its integral,
and use nonnegativity of the relative entropy of
\(\rho\) with respect to this tilted law.
Equation (V85.4) bounds its log normalizer.
All quantities are finite at a fixed regulator for
finite-entropy \(\rho\); truncation gives the same argument
if a general measurable version is used.

For a coarse event \(B\) of positive probability,
take \(\rho=\zeta_{n,t}(\cdot\mid B)\). Its entropy is
exactly \(\log[1/\zeta_{n,t}(B)]\).
Thus (V85.6) quantifies rather than removes the
conditioning cost. A point mass at a specified
coarse field is singular: (V85.6) gives no bound there.
Even an extensive entropy cost gives only a potentially
quadratic-in-volume Fisher bound.

\subsection{Reusing the actual surface-cost inverse at every coarse field}

For clarity record the inherited whole-torus estimate
in the form needed for the next theorem:
\[
 \|h\|_{-1,\widehat L_{g,t}}^2
       \le G_n(t)\nu_{g,t}(h^2),\qquad \nu_{g,t}(h)=0,
 \quad
 G_n(t)=14E_0^2v\,e^{W_n(t)+3tS},
 \tag{V85.7}\label{c18v85:globalinverse}
\]
where \(W_n(t)=tD^2K_\zeta(C_{\rm cut}n^2+3)\),
\(D=\pi\sqrt2\), \(S=D^2K_*\), and
\(C_{\rm cut}=216/(e^\zeta-1)\).
The squared negative norm uses the actual moving-metric
unshifted inverse. It is not a Haar inverse.

V72 states its applications at coarse identity.
Its proof of (V85.7) has the same constants for
every \(g\): the pair/free product metric and anchor
order do not change with \(g\); V63's weighted density
Hessian and V60's one-factor oscillation bounds are
uniform in \(g\); and V51's metric comparison is
global. The four-point density ratios in the
coordinate-replacement proof cancel their normalizers
on each fixed fibre. No coarse derivatives enter.
These are all the hypotheses of V71's comparison
and V72's periodic crossing calculation.
Applying those proved lemmas gives (V85.7) without
another canonical-path proof.

In particular, uniformly in \(t\in[0,4]\),
\[
 G_n(t)\le G_n^*
   :=14E_0^2v\,e^{b_*n^2+c_*},\quad
 b_*=4D^2K_\zeta C_{\rm cut},\quad
 c_*=12D^2K_\zeta+12S .
 \tag{V85.8}\label{c18v85:surface}
\]
These constants can be enormous. The lower spectral
floor \(1/G_n^*\) tends to zero; it is not the desired
uniform gap.

\subsection{The full large-amplitude source has a paid inverse tail}

At each fixed \(g,t\), and for \(R>0\), define
\[
 \begin{split}
 h_{g,t}^{R}
 &=a_{g,t}{\bf1}_{\{|a_{g,t}|>R\sqrt v\}}
      -\nu_{g,t}\!\left(a_{g,t}
                  {\bf1}_{\{|a_{g,t}|>R\sqrt v\}}\right),\\
 c_{g,t}^{R}
 &=a_{g,t}{\bf1}_{\{|a_{g,t}|\le R\sqrt v\}}
      -\nu_{g,t}\!\left(a_{g,t}
                  {\bf1}_{\{|a_{g,t}|\le R\sqrt v\}}\right).
 \end{split}
 \tag{V85.9}\label{c18v85:split}
\]
Thus \(a=h^R+c^R\) exactly and both terms are centered.
No claim is made that either term separately is
the temporal score of a new probability path.

\begin{theorem}[Conditioned, unshifted large-amplitude response]
\label{c18v85:tail}
Let \(\rho\ll\zeta_{n,t}\) and suppose
\(d\rho/d\zeta_{n,t}\le e^J\), \(J\ge0\). Then
\[
 \int\|h_{g,t}^{R}\|_{-1,\widehat L_{g,t}}^2\,d\rho(g)
 \le \sqrt2\,e^J G_n(t)v(R^2+\Theta)e^{-R^2/\Theta}.
 \tag{V85.10}\label{c18v85:tailinverse}
\]
In particular, fix \(j\ge0,\delta>0\), allow
\(J\le jn^2\), and set
\(R_n^2=\Theta(b_*+j+\delta)n^2\). Uniformly in time,
\[
 \int\|h_{g,t}^{R_n}\|_{-1,\widehat L_{g,t}}^2\,d\rho(g)
 \le14\sqrt2 E_0^2v^2(R_n^2+\Theta)
                            e^{c_*-\delta n^2}
       \longrightarrow0 .
 \tag{V85.11}\label{c18v85:vanishing}
\]
\end{theorem}
\begin{proof}
Under the joint law \(d\zeta_{n,t}(g)d\nu_{g,t}\),
put \(Y=a^2/v\). Equation (V85.4) gives
\(\mathbb P(Y>x)\le\sqrt2 e^{-x/\Theta}\).
Integration of this tail yields
\[
 \mathbb E[Y{\bf1}_{\{Y>R^2\}}]
 =R^2\mathbb P(Y>R^2)+\int_{R^2}^\infty\mathbb P(Y>x)\,dx
 \le\sqrt2(R^2+\Theta)e^{-R^2/\Theta}.
\]
Conditional centering decreases the squared norm, so
\(\nu_g[(h^R)^2]\le
\nu_g[a^2{\bf1}_{\{|a|>R\sqrt v\}}]\).
Use (V85.7) on each fibre and the displayed bound
on \(d\rho/d\zeta_{n,t}\) to prove (V85.10).
Substitute (V85.8) and the value of \(R_n\);
the exponents \(jn^2+b_*n^2\) cancel explicitly.
The remaining factor is polynomial in \(n\),
whereas \(e^{-\delta n^2}\) decays faster.
Sharp indicators cause no domain problem:
\(h^R\) is a centered \(L^2\) source for the inverse.
No derivative of the indicator is taken.
\end{proof}

The unconditioned choice is \(\rho=\zeta_{n,t}\),
\(j=0\). A coarse event with probability at least
\(e^{-jn^2}\) is another allowed choice after
normalization. Conditioning on an event with only
a volume-exponential lower probability requires
a larger \(J\); (V85.10) retains that cost, but
(V85.11) with fixed \(j\) no longer applies.

This tail is selected by the \emph{full native source
itself}, not by membership in the old central or
noncentral patches. Its threshold is of order
\(n^{5/2}\), not the typical \(\sqrt v=n^{3/2}\)
scale. Thus the theorem does not claim to remove
the typical source.

For a permitted \(\rho\), abbreviate
\[
 {\cal A}=\int\|a\|_{-1,\widehat L}^2d\rho,\qquad
 {\cal A}_{\rm core}=\int\|c^{R_n}\|_{-1,\widehat L}^2d\rho,
 \qquad
 {\cal A}_{\rm tail}=\int\|h^{R_n}\|_{-1,\widehat L}^2d\rho .
\]
The triangle inequality in the direct integral of
the fibre negative-norm Hilbert spaces gives
\[
       \left|\sqrt{\cal A}
                   -\sqrt{{\cal A}_{\rm core}}\right|
                       \le\sqrt{{\cal A}_{\rm tail}} .
 \tag{V85.12}\label{c18v85:normcomparison}
\]
Consequently a volume-uniform bound for
\({\cal A}/v\) is equivalent, under these selections,
to such a bound for \({\cal A}_{\rm core}/v\).
We have bounded their \emph{square-root} difference;
without an additional bound on the core one must
not infer that the difference of the two energies
tends to zero. The core inverse is still unpaid.

\subsection{Why exponential score moments do not finish the inverse}

Here is a separate comparison, not a native YM law.
On two equally weighted states \(\epsilon=\pm1\),
let a local probability path have raw density
\(e^{s\sqrt v\,\epsilon}\). At \(s=0\) its exact
normalized score is \(a=\sqrt v\,\epsilon\).
Choose the reversible generator at that point
\[
 (L_nf)(\epsilon)
      =\frac{e^{-n^2}}2[f(\epsilon)-f(-\epsilon)] .
\]
Then \(L_n\epsilon=e^{-n^2}\epsilon\), and
\[
 \nu e^{a^2/(\Theta v)}=e^{1/\Theta},\qquad
 \nu(a^2)/v=1,\qquad
 \frac{\langle a,L_n^{-1}a\rangle}{v}=e^{n^2}.
 \tag{V85.13}\label{c18v85:comparison}
\]
For any \(\Theta\ge2/\log2\) the first quantity is at
most \(\sqrt2\). Every fixed moment is harmless,
and a cutoff \(R_n>1\) removes nothing: the entire
large inverse remains in the core.
This example illustrates the low-energy distinction,
not a claim that this generator or path obeys the
native Wilson equations or all native local floors.

\subsection{Verification and the next load-bearing estimate}

The diagnostic
\begin{center}\small
\path{scripts/verify_ym_c18_v85_exponential_source.py}
\end{center}
checks the bounded-differences constant and Gaussian
integration normalization, exhaustive finite-product
moment regressions, and conditional symmetrization
for a nonuniform correlated finite law. Exact
finite-state calculations retain the centered tail
and core in an anisotropic reversible inverse;
the information-cost and surface-exponent balances
are checked separately. Native incidence regressions
check the fine-link and conditional-factor counts.
These tests are not simulations of the native
finite-time law and are not a proof by lattice
sampling; the uniform claims have the proofs above.

The previous goal turn was progress. This turn
adds exponential full-source control and a
vanishing unshifted tail under explicit coarse
averages/selections. It does not replace the full
mass-gap objective by that averaged problem.
Further rare-family exclusions or stronger score
moments alone would leave the central debt intact:
control the typical native source's low-energy
spectral weight, and justify any passage from
coarse ensembles to the required prescribed
coarse fields. The Poisson Hessian and transport
strain, interacting-vacuum extension, physical
scale and clock, reflection-positive nontrivial
continuum construction, and positive physical
YM gap remain unproved.

Removing this append and restoring the title
recovers V84 byte for byte. Other sources are
preserved; only the newest active YM filename
is retained. All build products stay outside
\path{latex_notes}, which contains only
\texttt{.tex} files. The prior corpus is preserved,
not independently recertified. The goal remains active.
% END CYCLE 18 VERSION 85 APPEND
% BEGIN CYCLE 18 VERSION 86 APPEND
\clearpage
\section[V86: The exact action filter and a paid spectral source sector]{V86: The exact action filter and a paid spectral source sector}
\label{c18v86:context}

V85's large-amplitude tail leaves the typical native source
in the inverse equation. We now use the special structure
of that source, rather than another moment bound. At every
chart time the full normal-chart score is exactly
\((12-\widehat L)\) applied to the centered transported
Wilson action. This pays the entire source inverse on the
auxiliary spectral sector \(E\ge6\), uniformly in volume
and at every coarse field. The unresolved part is a
precisely identified low-energy action autocorrelation.

The identities below use the actual probability, induced
moving metric and full source. They do not replace the
law by a Wilson Gibbs tilt or assume a conditional gap.
The number twelve is the ambient plaquette Casimir in
the present unit-round convention; six is a convenient
auxiliary spectral cutoff. Neither is a physical mass.
The full interacting four-dimensional continuum YM
mass-gap problem remains open.

For nonduplication, V50 already gives the exact temporal
score in the inverse-Wilson-flow chart; V51 gives the
normal-chart divergence formula, and V57 checks
\(12W-L_0W=3W_{\rm cap}\) at time zero. We reuse them.
The new step is the \emph{all-time} normal-chart
action filter, its full spectral consequence, and
its action-normalizer identities. V72's surface-cost
inverse and V85's tail comparison are not repeated.

\subsection{The actual action and its conditional generator}

Keep \(\beta=0\), \(n\ge5\), \(v=n^3\), and the native
map \(c_t=b_2\Phi_t\), normal identification \(\eta_t\),
and projection \(P_t=I-Q_t\) of V51. On a fixed \(M_g\) set
\[
 \begin{gathered}
 w_t=W\circ\eta_t,\qquad
 m_g(t)=\nu_{g,t}(w_t),\qquad
 \widetilde w_t=w_t-m_g(t),\\
 L=\widehat L_{g,t},\qquad
 {\cal V}_g(t)=\nu_{g,t}(\widetilde w_t^2),\\
 {\cal D}_g(t)=\nu_{g,t}(|dw_t|_{\widehat m_t^{-1}}^2),
 \qquad
 {\cal Q}_g(t)=\nu_{g,t}(|F|^2\circ\eta_t),\qquad F=\nabla W .
 \end{gathered}
 \tag{V86.1}\label{c18v86:data}
\]
The temporary \(L\) is the full actual fibre generator,
not its zero-time value. It is nonnegative and
self-adjoint in \(L^2(\nu_{g,t})\); at each finite
regulator its kernel consists of constants. The
probability is the pullback of ambient Haar conditioned
on \(c_t=g\). Its coarea density is retained.

We work throughout the required interval \([0,4]\).
In fact the following identities and numerical energy
bounds hold at every finite nonnegative chart time.
To justify that qualitative extension, \(F\) is smooth
on a compact fine product, so \(\Phi_t\) is complete.
The map \(c_t\) remains a submersion. Hence \(Q_t\) and
\(N_t=Q_tF\) are smooth on every compact time interval
times the fine product, and their time-dependent
flow exists there. This does not extend V51's
quantitative derivative majorants to unbounded time.

\subsection{An exact spectral filter for the complete source}

\begin{theorem}[All-time action filter, with the mean retained]
\label{c18v86:filter}
For the raw Jacobian \(r_t=\log\operatorname{Jac}\eta_t\)
and the full normalized score \(a_{g,t}\),
\[
 \boxed{\quad
 \dot r_t=12w_t-Lw_t,\qquad
 \partial_t\log Z_g(t)=12m_g(t),\qquad
 a_{g,t}=(12-L)\widetilde w_t .
 \quad}
 \tag{V86.2}\label{c18v86:filteridentity}
\]
These equalities hold on every fixed coarse fibre,
not only after coarse averaging.
\end{theorem}
\begin{proof}
At a fixed time, \(T_t=P_tF\) is a smooth ambient
vertical field for \(c_t\). Its restriction to a
moving fibre is the induced gradient of \(W\).
Ambient divergence of a vertical field equals its
conditional probability divergence. Indeed integrate
by parts against a product of a fine test and an
arbitrary smooth function of \(c_t\). The derivative
of the latter along \(T_t\) is zero. Disintegration
and localization in \(g\) give conditional integration
by parts; smoothness extends the identity to every
coarse field. This argument includes the coarea
weight and would not hold with an arbitrarily
substituted induced-volume probability.

Consequently, after pullback by \(\eta_t\),
\[
  (\operatorname{div}_{\mu_{\rm f}}T_t)\circ\eta_t=-Lw_t.
\]
Every fine plaquette is fundamental in its four
unit-round link factors, so the inherited ambient
identity is \(\operatorname{div}F=\Delta W=-12W\).
Since \(N_t=F-T_t\),
\[
  (\operatorname{div}N_t)\circ\eta_t=-12w_t+Lw_t.
\]
The normal-flow Jacobian identity
\(\dot r_t=-(\operatorname{div}N_t)\circ\eta_t\)
proves the first formula. Integrating it under the
actual \(\nu_{g,t}\) annihilates \(Lw_t\).
Differentiation of the normalizer therefore gives
the second formula. Subtracting that mean gives
the last formula, because \(L\) kills constants.
\end{proof}

Thus a cancellation at auxiliary energy twelve is
exact at every time: if an actual eigenspace has
energy twelve, the source has no component there.
This does not assert persistence of an eigenvalue
or of a zero-time eigenspace. Below twelve and
above twelve the source multiplier has opposite
signs; those signs must not be replaced by a
positive scalar before forming the inverse cost.

\subsection{Uniform energy bounds without a probability-gap hypothesis}

\begin{lemma}[Paid action energy]
\label{c18v86:energy}
At every \(g,t\),
\[
       0\le{\cal D}_g(t)\le{\cal Q}_g(t)\le384v .
 \tag{V86.3}\label{c18v86:energybound}
\]
For the actual coarse probability \(d\zeta_{n,t}=Z_g(t)\,dg\),
\[
 \begin{gathered}
 \int[{\cal V}_g(t)+m_g(t)^2]\,d\zeta_{n,t}(g)=6v,
 \qquad \int m_g(t)\,d\zeta_{n,t}(g)=0,\\
 \int{\cal Q}_g(t)\,d\zeta_{n,t}(g)=72v,\qquad
 \int{\cal D}_g(t)\,d\zeta_{n,t}(g)\le72v .
 \end{gathered}
 \tag{V86.4}\label{c18v86:ambientbudgets}
\]
\end{lemma}
\begin{proof}
The identification by \(\eta_t\) is an isometry
between its pullback metric and the induced moving
fibre metric. Therefore
\[
 |dw_t|_{\widehat m_t^{-1}}^2
                        =|P_tF|^2\circ\eta_t
                        \le |F|^2\circ\eta_t .
\]
As a function of one link, a plaquette trace is a
unit linear quaternion coordinate. Its round
gradient has norm at most one. Each fine link
occurs in four plaquettes, hence \(|F_e|\le4\).
There are \(24v\) fine links, proving the bound \(384v\).

Under ambient Haar, every plaquette has mean zero
and second moment \(1/4\). Distinct plaquettes
are orthogonal: a link in their symmetric
difference occurs once, and integrating that
Haar link gives zero. Thus
\(\mu_{\rm f}(W)=0\) and \(\mu_{\rm f}(W^2)=6v\).
Integration by parts using \(\Delta W=-12W\)
gives \(\mu_{\rm f}(|F|^2)=72v\).
Finally V84's exact ambient change of variables
and disintegration identify these Haar integrals
with the corresponding coarse averages in (V86.4).
No fixed-fibre expectation is replaced by Haar.
\end{proof}

The constants in this lemma use local plaquette
geometry, not the enormous all-time chart-jet
majorants. In particular (V86.3) is pointwise
in the coarse parameter and uniform in the
displayed time.

\subsection{A whole spectral part of the typical source is now controlled}

Let \(d\rho_{W,g,t}(E)\) be the spectral measure of
\(\widetilde w_t\) for \(L\), supported in \((0,\infty)\).
Define
\[
 \begin{gathered}
 {\cal C}_g(t)=\langle a,L^{-1}a\rangle_{\nu_{g,t}},\qquad
 {\cal S}_{W,g}(t)
        =\langle\widetilde w_t,L^{-1}\widetilde w_t\rangle_{\nu_{g,t}},\\
 \Pi_<={\bf1}_{(0,6)}(L),\qquad
 \Pi_\ge={\bf1}_{[6,\infty)}(L),\qquad
 {\cal S}_{W,<}=\int_{(0,6)}E^{-1}\,d\rho_{W,g,t}(E).
 \end{gathered}
 \tag{V86.5}\label{c18v86:spectrum}
\]
All inverse expressions are finite at a fixed regulator.
This fact is not used as a uniform bound.

\begin{theorem}[Paid full-source sector and exact low-energy debt]
\label{c18v86:sector}
The actual source spectral measure is
\[
          d\rho_{a,g,t}(E)=(12-E)^2d\rho_{W,g,t}(E).
 \tag{V86.6}\label{c18v86:measure}
\]
Its high-sector inverse cost obeys
\[
 \begin{gathered}
 {\cal C}_\ge:=\|\Pi_\ge a\|_{-1,L}^2
                \le{\cal D}_g(t)\le384v,\\
 \int{\cal C}_\ge\,d\zeta_{n,t}\le72v .
 \end{gathered}
 \tag{V86.7}\label{c18v86:high}
\]
The full cost satisfies
\[
 \boxed{\qquad
 36{\cal S}_{W,<}\le{\cal C}_g(t)
                   \le144{\cal S}_{W,<}+384v .
 \qquad}
 \tag{V86.8}\label{c18v86:low}
\]
In particular, a volume-uniform bound for the
complete source inverse is equivalent to such
a bound for this specific low-energy action
susceptibility. No global gap is assumed.
\end{theorem}
\begin{proof}
The functional calculus applied to (V86.2)
proves (V86.6). For \(E\ge6\),
\((12-E)^2/E\le E\). Hence
\[
 {\cal C}_\ge
 =\int_{[6,\infty)}\frac{(12-E)^2}{E}\,d\rho_W(E)
 \le\int E\,d\rho_W(E)
 ={\cal D}_g(t).
\]
Use (V86.3)--(V86.4).
On \(0<E<6\) the squared multiplier is between
\(36\) and \(144\). The low and high spectral
parts are orthogonal for the inverse quadratic
form, so their costs add exactly. This gives
(V86.8), not merely a triangle-inequality upper
bound on a modified source.
\end{proof}

This controls a part of the \emph{entire} typical
source and all its mixed components, on the actual
full fibre, for every \(g\). It is not the inverse
of a finite bank projection or of a source
restricted to a rare event. The energy cutoff
belongs to \(L\), not to the physical Hamiltonian.

There are also useful exact bookkeeping identities:
\[
 \begin{split}
 {\cal C}_g(t)
     &=144{\cal S}_{W,g}(t)-24{\cal V}_g(t)+{\cal D}_g(t),\\
 L^{-1}a&=12L^{-1}\widetilde w_t-\widetilde w_t,\\
 \left|\sqrt{{\cal C}_g(t)}
                       -12\sqrt{{\cal S}_{W,g}(t)}\right|
     &\le\sqrt{{\cal D}_g(t)} .
 \end{split}
 \tag{V86.9}\label{c18v86:exactcost}
\]
The first follows by expanding the spectral
multiplier, and the other two follow in the
centered inverse and its Hilbert norm.
Moreover
\({\cal S}_{W,\ge}\le{\cal D}_g(t)/36\).
Thus full action susceptibility is an equivalent
target as well, with an already-controlled
high-energy part. The negative variance term in
(V86.9) is retained; it is not an error term
whose sign may be reversed.

At zero time the inherited identities are
\(w_0=C+U\), \(LC=9C\), \(LU=12U\),
\(\nu C^2=\nu U^2=3v\), and \(\nu(CU)=0\).
Consequently
\[
 {\cal V}_g(0)=6v,\quad {\cal D}_g(0)=63v,\quad
 {\cal S}_{W,g}(0)=\tfrac7{12}v,\quad
 a_{g,0}=3C,\quad {\cal C}_g(0)=3v .
 \tag{V86.10}\label{c18v86:zero}
\]
These are regressions from V50--V57, not newly
computed Taylor coefficients. The actual source
lies entirely above the cutoff at zero time;
no claim that it stays there is made.

\subsection{The normalizer pays a static action-variance budget}

\begin{theorem}[Exact conditional action-mean evolution]
\label{c18v86:pressure}
For every fixed coarse field,
\[
 m_g'(t)=12{\cal V}_g(t)-{\cal Q}_g(t),\qquad
 \partial_t^2\log Z_g(t)
                   =144{\cal V}_g(t)-12{\cal Q}_g(t).
 \tag{V86.11}\label{c18v86:meanODE}
\]
In particular, for any finite \(T\ge0\),
\[
 \int_0^T{\cal V}_g(t)\,dt\le(2+32T)v .
 \tag{V86.12}\label{c18v86:variancebudget}
\]
At \(T=4\) this is \(130v\).
\end{theorem}
\begin{proof}
The normal chart moves with velocity \(-N_t\).
Since \(N_t=Q_tF\) is an orthogonal normal projection,
\[
             \partial_tw_t=-|N_t|^2\circ\eta_t .
\]
Differentiate the actual normalized expectation:
\[
 \begin{split}
 m_g'
 &=\nu_g(\partial_tw_t)+\nu_g(a w_t)\\
 &=-\nu_g(|N_t|^2\circ\eta_t)
             +12{\cal V}_g-\langle w_t,Lw_t\rangle_{\nu_g}\\
 &=12{\cal V}_g-{\cal Q}_g .
 \end{split}
\]
The last step uses \(|N_t|^2+|P_tF|^2=|F|^2\).
Differentiate the second identity in (V86.2)
to obtain the pressure formula.
At zero, each plaquette has a free Haar factor
after \(b_2=g\) is fixed, so \(m_g(0)=0\).
Also \(|m_g(T)|\le24v\), since \(|W|\le24v\).
Integrate the first formula and use
\({\cal Q}_g(t)\le384v\):
\[
 12\int_0^T{\cal V}_g(t)\,dt
            =m_g(T)+\int_0^T{\cal Q}_g(t)\,dt
            \le24v+384Tv .
\]
\end{proof}

For the normalized pressure \(\phi_g=v^{-1}\log Z_g\)
this also gives the explicit bounds
\[
       |\phi_g'|\le288,\qquad \phi_g''\ge-4608 .
 \tag{V86.13}\label{c18v86:pressurebounds}
\]
They improve the crude raw-geometric pressure
majorants without claiming a pointwise bound
on \({\cal V}_g\).
Equations (V86.2) and (V86.4) further give
\[
 \int|\partial_t\log Z_g|^2d\zeta_{n,t}
       =144\int m_g^2d\zeta_{n,t}\le864v .
 \tag{V86.14}\label{c18v86:coarsefisher}
\]
The conditional normalizer is an observable of
the actual action, not a removable constant.

\subsection{Static variance is not the long auxiliary-time correlation}

Fix the chart time \(t\) and coarse field \(g\).
Let \(s\ge0\) now denote the stationary diffusion
time of the \emph{frozen} generator \(L=\widehat L_{g,t}\).
Then the spectral theorem gives
\[
 \begin{split}
 {\cal S}_{W,g}(t)
   &=\int_0^\infty
       \langle\widetilde w_t,e^{-sL}\widetilde w_t\rangle_{\nu_g}\,ds,\\
 {\cal S}_{W,<}
   &=\int_0^\infty
       \langle\Pi_<\widetilde w_t,
                      e^{-sL}\Pi_<\widetilde w_t\rangle_{\nu_g}\,ds .
 \end{split}
 \tag{V86.15}\label{c18v86:correlation}
\]
Indeed integrate \(e^{-sE}\) against the positive
spectral measure and use Tonelli.
The integrands are nonnegative, even though
their individual plaquette-pair contributions
need not be. There are three distinct notions
of time: Wilson chart time \(t\), auxiliary
diffusion time \(s\), and still-uncalibrated
physical time. Integrating static variance
in \(t\) in (V86.12) does not integrate the
correlation in \(s\) in (V86.15).

For completeness, if \(X_s\) is the stationary
finite-regulator diffusion with generator \(-L\),
stationarity and reversibility imply
\[
 \lim_{S\to\infty}\frac1S
    \operatorname{Var}\!\left(\int_0^S\widetilde w_t(X_s)\,ds\right)
                         =2{\cal S}_{W,g}(t).
 \tag{V86.16}\label{c18v86:asymvariance}
\]
At finite \(S\) the left side before taking the
limit is
\(2\int_0^S(1-s/S)\langle\widetilde w_t,
e^{-sL}\widetilde w_t\rangle\,ds\).
Its nonnegative integrand increases to the
integrand in (V86.15), proving the formula.
No infinite-volume central limit theorem is
being asserted.

For primary context, the statement of Theorem 3.1,
especially its source-specific energy-dual
condition (3.7), was read in C.~Kipnis and
S.~R.~S.~Varadhan,
\href{https://www.numdam.org/article/AST_1985__132__65_0.pdf}
{\emph{Central limit theorems for additive
functionals of reversible Markov chains and
applications}}, \emph{Ast\'erisque} 132 (1985),
65--70. The present continuous-time spectral
identities are proved directly. That reference
does not supply our missing uniform energy-dual
bound, and no CLT is imported as a substitute.

\subsection{The interacting-vacuum term cannot be silently dropped}

There is an exact scope check beyond \(\beta=0\).
For a fixed smooth positive vacuum density
\(\Omega_\beta^2d\mu_{\rm f}\), put
\(u_\beta=\log\Omega_\beta\), and keep the same
geometric \(F,N_t,\eta_t\).
Use the corresponding actual conditional law
and its generator \(L_\beta\), and center all
quantities under that law. The same vertical
divergence proof gives
\[
 a_{\beta,g,t}
   =(12-L_\beta)\widetilde w_{\beta,t}
       -2\left[
         \langle\nabla u_\beta,F\rangle\circ\eta_t
        -\nu_{\beta,g,t}
           (\langle\nabla u_\beta,F\rangle\circ\eta_t)
          \right].
 \tag{V86.17}\label{c18v86:vacuum}
\]
Here \(\operatorname{div}_{\Omega_\beta^2\mu_{\rm f}}F
=-12W+2\langle\nabla u_\beta,F\rangle\).
The new centered vacuum term has no paid uniform
inverse in this append. Thus the beta-zero spectral
bound cannot simply be declared an interacting
continuum theorem. This formula identifies exactly
what an interacting extension must retain.

\subsection{Verification, progress and next obligation}

The diagnostic
\begin{center}\small
\path{scripts/verify_ym_c18_v86_action_filter.py}
\end{center}
checks the coarea divergence identity on a
nonlinear two-circle flow with exact rational
coordinate algebra; omitting its conditional
Jacobian gives a nonzero defect. Native periodic
incidence, the four-link Casimir, and noncommuting
coarse-twist regressions check the inherited
action constants and the zero-time cancellation.
Exact finite spectral calculations check the
full filter, its energy-twelve cancellation,
the energy-six inequality and all cost signs.
The mean/normalizer differentiation and the
two-time correlation integral are checked
separately. Comparison fixtures are not native
finite-time probability simulations or proofs
by sampling; the all-volume arguments are above.

The previous turn was progress. This turn pays
a whole non-rare spectral source sector at
every coarse field with explicit constants
independent of chart time and volume.
The remaining typical-source debt is now
\({\cal S}_{W,<}/v\), equivalently the long
auxiliary-time conditional action correlation,
not an unspecified score amplitude.
The useful next estimate must exploit actual
Wilson-action dynamics in that low-energy
sector. More static moments, old rare-family
exclusions, or relabeling a finite-volume
surface floor do not pay it.

Even a uniform action susceptibility would
not alone establish the Poisson Hessian and
transport strain, the interacting-vacuum
extension, physical scaling and clock,
or a nontrivial reflection-positive
four-dimensional continuum construction.
The physical YM mass gap remains unproved
and the full goal remains active.

Removing this append and restoring the title
recovers V85 byte for byte. Other sources and
legacy archives are unchanged; only the newest
active YM filename is retained. The prior corpus
is preserved, not independently recertified.
All build products stay outside \path{latex_notes},
which contains only \texttt{.tex}. The last
companion review was V85 and the last broad
review V80; both are next due at V90.
Archive/restart remains V100.
% END CYCLE 18 VERSION 86 APPEND
% BEGIN CYCLE 18 VERSION 87 APPEND
\clearpage
\section[V87: Actual action memory and the connected return]{V87: Actual action memory, failure of scalar Markov closure, and the connected return}
\label{c18v87:context}

V86 pays the full source's auxiliary sector \(E\ge6\).
The remaining action susceptibility is a dynamical question.
This append derives its actual scalar memory equation, with
a well-defined orthogonal evolution and a positive return
kernel. A native calculation shows that the total Wilson
action is not itself a Markov coordinate, even at time zero.
Consequently a diffusion fitted only to its marginal
distribution cannot silently replace the conditioned dynamics.

The new estimates identify the part of the memory amplitude
already paid by V84's Fisher budgets. They do not bound its
long-time integral uniformly in volume. The native zero-time
return fraction is \(1/49\), but the corresponding correction
to the source inverse is not negligible. A separate reversible
comparison shows why even a vanishing memory amplitude does
not establish a uniform return margin.

All statements about the native source keep \(\beta=0\),
\(n\ge5\), \(v=n^3\), the full conditional probability and
the induced moving metric of V86. The Wilson chart time \(t\)
is frozen while the auxiliary semigroup time \(s\) varies.
Neither time is identified with physical time. The full
interacting continuum construction and physical mass gap
remain unproved.

\subsection{A domain-correct action projection}

Fix \(g,t\), abbreviate \(L=\widehat L_{g,t}\) and
\(\nu=\nu_{g,t}\), and write
\[
 \begin{gathered}
 f=w_t-\nu(w_t),\qquad V=\nu(f^2),\qquad
 D=\langle f,Lf\rangle,\qquad I=\nu(a^2),\\
 a=(12-L)f,\qquad
 {\cal S}_W=\langle f,L^{-1}f\rangle,\qquad
 {\cal C}=\langle a,L^{-1}a\rangle .
 \end{gathered}
 \tag{V87.1}\label{c18v87:data}
\]
If \(V=0\), then \(f=a=0\), and their costs vanish.
Otherwise \(D>0\), and on the centered Hilbert space
\({\cal H}_0=L^2_0(\nu)\) put
\[
 \begin{gathered}
 q=f/\sqrt V,\qquad d=D/V,\qquad
 {\mathsf P}h=\langle q,h\rangle q,\qquad
 {\mathsf Q}=1-{\mathsf P},\\
 \xi={\mathsf Q}Lq=(L-d)q .
 \end{gathered}
 \tag{V87.2}\label{c18v87:projection}
\]
These Hilbert-space projections are not the geometric
normal and tangent projections of V51. Nor is
\({\mathsf P}\) conditional expectation given the action:
it retains only the one-dimensional span of \(f\).

\begin{lemma}[The orthogonal generator exists]
\label{c18v87:domain}
The restriction of the closed form of \(L\) to
\({\mathsf Q}{\cal H}_0\) defines a positive self-adjoint
operator \({\cal A}\). At a fixed regulator it is bounded
below by the positive centered gap of \(L\).
In the decomposition
\({\cal H}_0=\mathbb Rq\oplus{\mathsf Q}{\cal H}_0\),
\[
                 L=
          \begin{pmatrix}d&\xi^*\\ \xi&{\cal A}\end{pmatrix}.
 \tag{V87.3}\label{c18v87:block}
\]
The off-diagonal part is bounded. In particular
\(e^{-s{\cal A}}\) is a genuine strongly continuous
self-adjoint contraction semigroup.
\end{lemma}
\begin{proof}
The finite-regulator fibre is compact and connected,
and its metric and positive probability density are smooth.
Thus its centered elliptic generator has a positive gap,
used here only qualitatively. The vector \(q\) is smooth
and belongs to \(D(L)\). Projection onto its span preserves
the form domain, and its orthogonal complement is a closed
subspace there. The restricted form is densely defined on
that complement, closed, and has the same qualitative
lower bound.

For a vector \(h\) in the complementary form domain,
the cross form is
\(\langle Lq,h\rangle=\langle\xi,h\rangle\).
It is bounded in the Hilbert norm because \(\xi\in{\cal H}_0\).
The operator associated with the full form is therefore
the block diagonal operator \(d\oplus{\cal A}\) plus this
bounded symmetric off-diagonal operator. This proves
(V87.3), including its domain interpretation.
\end{proof}

There is no claim that the orthogonal semigroup is Markov
or positivity preserving on functions. Scalar positivity
of the kernel below follows instead from self-adjoint
spectral calculus.

\subsection{The actual autocorrelation has a positive renewal kernel}

Define
\[
 R(s)=\langle q,e^{-sL}q\rangle,\qquad
 k(s)=\langle\xi,e^{-s{\cal A}}\xi\rangle,\qquad s\ge0 .
 \tag{V87.4}\label{c18v87:kernels}
\]
Both are nonnegative and completely monotone. Set
\[
 \kappa(z)=\langle\xi,(z+{\cal A})^{-1}\xi\rangle
          =\int_0^\infty e^{-zs}k(s)\,ds,\qquad z\ge0 .
 \tag{V87.5}\label{c18v87:kappa}
\]

\begin{theorem}[Exact native action-memory equation]
\label{c18v87:memory}
At every fixed native \(g,t,n\) with \(V>0\),
\[
 \begin{gathered}
 R'(s)=-dR(s)+\int_0^s k(s-r)R(r)\,dr,\qquad R(0)=1,\\
 \int_0^\infty e^{-zs}R(s)\,ds
                  =\frac{1}{z+d-\kappa(z)},\qquad z\ge0 .
 \end{gathered}
 \tag{V87.6}\label{c18v87:volterra}
\]
In particular, putting
\[
 \delta=d-\kappa(0),\qquad \theta=\kappa(0)/d ,
 \tag{V87.7}\label{c18v87:margin}
\]
one has \(\delta>0\), \(0\le\theta<1\), and
\[
                 {\cal S}_W=\frac{V}{d(1-\theta)} .
 \tag{V87.8}\label{c18v87:susceptibility}
\]
No uniform lower bound on \(\delta\) or \(1-\theta\)
is asserted.
\end{theorem}
\begin{proof}
Decompose \(e^{-sL}q=R(s)q+y(s)\) with \(y\perp q\).
The block equations are
\[
 R'=-dR-\langle\xi,y\rangle,\qquad
 y'=-{\cal A}y-\xi R,\qquad y(0)=0 .
\]
Bounded off-diagonal coupling justifies variation of
constants:
\[
          y(s)=-\int_0^s e^{-(s-r){\cal A}}\xi R(r)\,dr .
\]
Substitution proves the first equation. The plus sign
in front of the return is essential for the dissipative
generator \(-L\).

To see strict positivity at \(z=0\), insert
\(q-{\cal A}^{-1}\xi\) in the positive quadratic form of \(L\).
Its form value is \(d-\kappa(0)\), and the vector is nonzero
because its \(q\)-component is one. The finite-regulator
gap makes this value strictly positive.
The same block elimination for \(z+L\) gives the second
identity for all \(z\ge0\). Alternatively, integrate the
Volterra equation; all integrals are finite at the
fixed regulator. At \(z=0\), multiply by \(V\).
\end{proof}

For nonnegative functions on \([0,\infty)\), let
\(*\) denote causal convolution. Define
\[
             b(s)=e^{-ds},\qquad h=b*k .
 \tag{V87.9}\label{c18v87:renewaldata}
\]
The integrated form of (V87.6) is \(R=b+h*R\), and
\(\|h\|_{L^1}=\theta\). Therefore the exact renewal
representation and its integrated remainder are
\[
 \begin{gathered}
 R=\sum_{j=0}^\infty h^{*j}*b,\qquad h^{*0}*b=b,\\
 \int_0^\infty
   \left[R-\sum_{j=0}^M h^{*j}*b\right](s)\,ds
      =\frac{\theta^{M+1}}{d(1-\theta)} .
 \end{gathered}
 \tag{V87.10}\label{c18v87:renewal}
\]
Indeed convolution by \(h\) has \(L^1\)-operator norm
\(\theta<1\). Its Neumann series converges there.
Every term is nonnegative, and its integral is
\(\theta^j/d\), so Tonelli gives the exact remainder.
This is an identity for the actual conditional diffusion,
not a postulated renewal model.

\subsection{The source cost separates into a resolved term and a return}

\begin{theorem}[Source-aware connected-return identity]
\label{c18v87:sourcecost}
The actual inverse cost has the nonnegative decomposition
\[
 \boxed{\quad
 {\cal C}
 = \frac{V(d-12)^2}{d}
        +\frac{144V}{d}\,\frac{\theta}{1-\theta}.
 \quad}
 \tag{V87.11}\label{c18v87:positivecost}
\]
The first term is exactly the Galerkin value on the trial
potential space \(\mathbb R f\). For an integer \(M\ge0\), put
\[
 {\cal C}^{[M]}
       =\frac{144V}{d}\sum_{j=0}^M\theta^j-24V+D .
\]
Then
\[
 0\le{\cal C}^{[0]}\le{\cal C}^{[M]}\le{\cal C},
 \qquad
 {\cal C}-{\cal C}^{[M]}
       =\frac{144V}{d}\,\frac{\theta^{M+1}}{1-\theta}.
 \tag{V87.12}\label{c18v87:sourceerror}
\]
\end{theorem}
\begin{proof}
Substitute (V87.8) into V86's exact identity
\({\cal C}=144{\cal S}_W-24V+D\), and use \(D=dV\).
The terms without a return combine to
\(V(d-12)^2/d\).

The variational formula for the positive inverse is
\[
 {\cal C}
 =\sup_\psi\{2\langle a,\psi\rangle-\langle\psi,L\psi\rangle\}.
\]
For \(\psi=cf\) this becomes
\(2c(12V-D)-c^2D\). Its maximum is
\((12V-D)^2/D=V(d-12)^2/d\), at \(c=(12-d)/d\).
The geometric sum and (V87.10) give (V87.12).
\end{proof}

Equivalently, for
\(\psi_0=(12-d)f/d\), the exact residual is
\[
 a-L\psi_0=-\frac{12\sqrt V}{d}\,\xi,\qquad
 {\cal C}-{\cal C}^{[0]}
       =\langle a-L\psi_0,L^{-1}(a-L\psi_0)\rangle .
 \tag{V87.13}\label{c18v87:residual}
\]
Here \(\xi\) is viewed in the full centered Hilbert space.
The equality follows by completing the variational square.
It does not replace its full inverse by \({\cal A}^{-1}\).

A small value of \(\theta\) alone is not a small relative
source error. If \(d\) is close to twelve, the resolved
term in (V87.11) is small by cancellation. At \(d=12\)
it vanishes altogether, whereas a nonzero return still
gives a positive source cost.

The exact uniform susceptibility target is
\(V/[v d(1-\theta)]\). A bound on \(1-\theta\) is not enough
if \(V/d\) has not also been controlled. For example,
\(V/d\le Av\) and \(\theta\le1-\varepsilon\) would give
\({\cal S}_W\le Av/\varepsilon\); neither hypothesis is
claimed through time four here.

\subsection{What is already paid for the genuine memory kernel}

The source decomposes orthogonally as
\[
       a=\sqrt V\,[(12-d)q-\xi].
\]
Consequently
\[
 V k(0)=I-V(12-d)^2,\qquad
               0\le V k(s)\le V k(0)\le I .
 \tag{V87.14}\label{c18v87:fisher}
\]
The middle inequality uses the spectral theorem for
\({\cal A}\). It holds for all \(s\ge0\), with chart
time still frozen.

Thus V84's already-paid constants imply, for each fixed
\(s\), every \(g\), and the required \(T=4\),
\[
 \begin{gathered}
 \int_0^T V_g(t)k_{g,t}(s)\,dt\le vK_T,\\
 \int V_g(t)k_{g,t}(s)\,d\zeta_{n,t}(g)\le vK_{\rm amb}.
 \end{gathered}
 \tag{V87.15}\label{c18v87:memorybudget}
\]
These formulas include the zero-variance case by assigning
\(Vk=0\) there. For a finite auxiliary-time window
\([0,S]\), integrating either inequality in \(s\) costs
a factor \(S\). There is no bound independent of \(S\)
from this argument. It pays the memory's amplitude,
not the integral defining \(\kappa(0)\).

\subsection{Native zero-time memory is explicit, and is not negligible}

Use the inherited native \(t=0\) decomposition
\(f=C+U\), where \(C=W_{\rm cap}\) and \(U=W_{\rm uncap}\).
At every coarse field,
\[
 LC=9C,\quad LU=12U,\quad
 \nu C^2=\nu U^2=3v,\quad \nu(CU)=0 .
\]
These facts are reused, not new Taylor coefficients.
The two orthonormal vectors
\[
             q=\frac{C+U}{\sqrt{6v}},\qquad
             p=\frac{-C+U}{\sqrt{6v}}
\]
span an invariant subspace of the full generator.
Its matrix and the resulting actual memory are
\[
 \begin{gathered}
 L\big|_{\operatorname{span}\{q,p\}}
       =\begin{pmatrix}21/2&3/2\\3/2&21/2\end{pmatrix},
 \qquad \xi=\tfrac32p,\quad {\cal A}p=\tfrac{21}{2}p,\\
 R_0(s)=\tfrac12(e^{-9s}+e^{-12s}),\qquad
 k_0(s)=\tfrac94 e^{-21s/2},\\
 d_0=\tfrac{21}{2},\qquad
 \kappa_0(0)=\tfrac3{14},\qquad
 \theta_0=\tfrac1{49},\qquad \delta_0=\tfrac{72}{7}.
 \end{gathered}
 \tag{V87.16}\label{c18v87:nativekernel}
\]
In particular, the orthogonal evolution is not obtained
by simply applying the original semigroup to \(\xi\).
The latter gives instead
\[
 \langle\xi,e^{-sL}\xi\rangle
       =\tfrac98(e^{-9s}+e^{-12s}),\qquad
 \int_0^\infty\langle\xi,e^{-sL}\xi\rangle ds=\tfrac7{32},
 \tag{V87.17}\label{c18v87:wrongkernel}
\]
which differs from \(3/14\). This is a native,
not merely abstract, distinction.

Substitution in the positive source decomposition gives
\[
 \begin{gathered}
 {\cal C}^{[0]}_0=\tfrac97v,\qquad
 {\cal C}_0-{\cal C}^{[0]}_0=\tfrac{12}{7}v,\qquad
 {\cal C}_0=3v,\\
 {\cal C}_0-{\cal C}^{[M]}_0
                   =\frac{84v}{49^{M+1}} .
 \end{gathered}
 \tag{V87.18}\label{c18v87:nativeerror}
\]
The \(1/49\) action return changes the source cost
by the fraction \(4/7\) of its exact value.
The action-line approximation would retain only
the fraction \(3/7\). This calculation warns against
discarding apparently small feedback before applying
the action-to-source filter.

\subsection{The total native action is not a Markov coordinate}

\begin{proposition}[Two equal-action configurations with unequal drift]
\label{c18v87:notmarkov}
At \(g=I,t=0\), the stationary auxiliary diffusion
projected onto the total action \(W\) is not a
time-homogeneous Markov process.
In particular there is no exact scalar Markov generator
for \(W\) obtained by replacing the fibre dynamics with
a diffusion on its action values.
\end{proposition}
\begin{proof}
Start with all fine links equal to \(I\). Consider
the two free links, both in direction one,
\[
                e_1=((0,1,0),1),\qquad
                e_0=((0,1,1),1).
\]
The direction labels here are one-based.
The first belongs to two captured and two uncaptured
plaquettes; the second belongs to four uncaptured
plaquettes. This follows directly by listing the two
plaquettes in each of the planes containing that link.
Neither belongs to a retained coarse chain.

In configuration \(U^{(j)}\), change only \(e_j\) to a
unit quaternion \(h\) with real part \(c\in(-1,1)\).
Every incident plaquette then has trace coordinate \(c\);
all others have coordinate one. Put \(b=1-c\).
Both configurations satisfy \(b_2=I\), and
\[
 \begin{array}{c|ccc}
       &W&C&U\\ \hline
 U^{(1)}&24v-4b&12v-2b&12v-2b\\
 U^{(0)}&24v-4b&12v&12v-4b
 \end{array}
\]
Since \(LW=9C+12U\),
\[
          (LW)(U^{(1)})-(LW)(U^{(0)})=6b\ne0 .
 \tag{V87.19}\label{c18v87:driftwitness}
\]
At both configurations the action is a submersion:
varying the angle of \(h\) differentiates four plaquettes,
giving derivative \(-4\sqrt{1-c^2}\ne0\).
Continuity provides two neighborhoods with overlapping
action-value intervals and separated values of \(LW\).
The positive Haar density and the local coarea theorem
then imply that \(LW\) cannot be a measurable function
of \(W\), even modulo null sets.

For completeness, let \(E\) be conditional expectation
onto the full sigma-algebra of \(W\), on \(L^2(\nu)\),
and let \(T_s=e^{-sL}\).
If the stationary projected process were Markov, its
two-point conditional transition operators \(ET_sE\)
would be a semigroup. Self-adjointness would imply
\[
 0=ET_{2s}E-(ET_sE)^2
     =ET_s(1-E)T_sE
     =[(1-E)T_sE]^*[(1-E)T_sE].
\]
Hence \(T_s\) preserves the action-measurable subspace.
Since \(W\in D(L)\), differentiating its orbit at zero
would make \(LW\) action-measurable, contradicting the
native witness.
\end{proof}

This obstruction concerns even the entire action
sigma-algebra, not only the rank-one projection
\({\mathsf P}\). It does not rule out useful enlarged
coordinates or a justified approximation with an error.
It rules out an exact memory-free replacement for the
specific native action process.

\subsection{A vanishing memory height can still have a critical return}

Here is a separate nonnative comparison, used only to
test a proposed inference. Let \(X,Y\) be independent
uniform signs, with the positive generator of independent
sign flips normalized by
\[
 LX=X,\qquad LY=\epsilon^2Y,\qquad
 L(XY)=(1+\epsilon^2)XY,\qquad 0<\epsilon\le\tfrac12 .
\]
Each sign flips at half its displayed eigenvalue;
this is a genuine reversible irreducible four-state chain.
Take
\[
 f=q=\sqrt{1-\epsilon}\,X+\sqrt\epsilon\,Y,\qquad V=1,
 \qquad a=(12-L)q .
\]
If desired, \(a\) is the actual normalized score at
\(\tau=0\) of the finite probability path proportional
to \(e^{\tau a}\). That path is not the native Wilson flow.

Put
\[
 d=1-\epsilon+\epsilon^3,\qquad
 \alpha=\epsilon+(1-\epsilon)\epsilon^2,\qquad
 B=\epsilon(1-\epsilon)(1-\epsilon^2)^2 .
\]
The complementary combination of \(X,Y\) has compressed
energy \(\alpha\); the \(XY\) mode is decoupled. Direct
two-by-two multiplication gives
\[
 \begin{gathered}
 k(s)=B e^{-\alpha s},\qquad
 k(0)=B\longrightarrow0,\qquad
 \theta=\frac{B}{d\alpha},\\
 {\cal S}_W=1-\epsilon+\epsilon^{-1},\qquad
 \delta=\frac{1}{1-\epsilon+\epsilon^{-1}}\longrightarrow0,
 \qquad \theta\longrightarrow1 .
 \end{gathered}
 \tag{V87.20}\label{c18v87:slowcomparison}
\]
Nevertheless \(d\ge1/2\), \(V/d\le2\), \(|q|\le\sqrt2\),
and \(I=\|(12-L)q\|_2^2\le144\).
Thus bounded source moments, bounded bare action
susceptibility \(V/d\), and even a vanishing memory
height do not bound the integrated return away from
its critical value. The source inverse
\(144{\cal S}_W-24+d\) diverges.
This is not evidence that the native YM kernel has this
behavior; it shows exactly which inference remains invalid.

\subsection{Literature transfer, verification and next step}

The projection/orthogonal-evolution structure is classical.
For primary context we read Sections 2--3 of H.~Mori,
\href{https://doi.org/10.1143/PTP.33.423}
{\emph{Transport, Collective Motion, and Brownian Motion}},
\emph{Progress of Theoretical Physics} 33 (1965), 423--455.
Its Hamiltonian evolution must not be used to copy the
sign of our dissipative return.
The domain distinction was also checked against
Section II.1 and the finite-rank discussion of
C.~Widder and T.~Schilling,
\href{https://arxiv.org/html/2604.20453v1}
{\emph{Generalised Langevin Dynamics: Significance and
Limitations of the Projection Operator Formalism}}
(2026 preprint).
Here bounded off-diagonal coupling, existence of the
orthogonal semigroup, and the scalar equation are
proved directly. No decay theorem or uniform memory
margin is imported from either source.

The diagnostic
\begin{center}\small
\path{scripts/verify_ym_c18_v87_action_memory.py}
\end{center}
checks exact finite spectral Schur identities, the native
two-channel convolution and source error, the periodic
equal-action/unequal-drift witness, and the reversible
vanishing-height comparison. Native noncommuting plaquette
regressions are retained. These finite tests do not certify
a native positive-time simulation or an all-volume gap;
the proofs and their scope are stated above.

The preceding turn was progress. This append supplies
actual source-specific dynamics and a native obstruction
to a tempting scalar closure. It does not count the old
zero-time eigenvalues or a general Schur formula as new
mass-gap estimates. The useful new target is a uniform
estimate on the \emph{actual integrated orthogonal return},
together with its bare action scale \(V/d\).
V84 pays a weighted memory-height budget; V86 already
pays the high-energy full-source sector. Neither permits
replacing \(e^{-s{\cal A}}\) by the original heat semigroup
or replacing a long-time integral by its initial value.
An enlarged native action/force system or a controlled
orthogonal-dynamics comparison must now estimate that return.

The low-energy source inverse, Poisson Hessian and
probability strain, interacting-vacuum term, physical
scaling and clock, and nontrivial reflection-positive
four-dimensional continuum remain unpaid.
The goal remains active, not complete.

Removing this append and restoring the title recovers
V86 byte for byte. The prior corpus is preserved, not
independently recertified. The other sources and legacy
archives are unchanged; only the newest active YM source
is retained. Build products and diagnostics stay outside
\path{latex_notes}, which contains only \texttt{.tex}.
Both scheduled companion and broad-literature reviews
are next due at V90. Archive/restart remains V100.
% END CYCLE 18 VERSION 87 APPEND
% BEGIN CYCLE 18 VERSION 88 APPEND
\clearpage
\section{V88: A native extensive action-energy floor}
\label{c18v88:main}

\subsection{Why go upstream of the return denominator?}

V87 separates the actual action susceptibility into a bare
action scale and an orthogonal return. It does not bound either
factor uniformly. We now pay a previously missing piece of
the bare scale: the native action Dirichlet energy cannot
collapse below a positive constant times volume, on the
identity coarse fibre, throughout the required chart interval.
The proof uses the actual normal flow, not a scalar Markov
replacement and not a comparison probability.

Keep \(\beta=0\), \(n\ge5\), \(v=n^3\), and \(0\le t\le4\).
Unless explicitly stated otherwise the coarse field in this
append is \(g=I\). Write
\[
 w_t=W\circ\eta_t,\quad
 {\cal V}(t)=\operatorname{Var}_{\nu_{I,t}}w_t,\quad
 {\cal D}(t)=\int |dw_t|_{\widehat m_t^{-1}}^2\,d\nu_{I,t}.
 \tag{V88.1}\label{c18v88:data}
\]
The metric and probability are exactly those of V86.
The product reference metric \(m_0\) has \(3v\) pair factors
of twice-round metric and \(18v\) free factors of unit-round
metric. The probability includes the actual coarea density.
The new result is not a Poincare inequality for arbitrary
functions: it is a lower bound on the energy of this specific,
unnormalized extensive action.

V60 and V63 already prove the required weighted normal-flow
jets and density bounds; V68 already explains how to remove
an arbitrary exterior using such spatial bounds. Those
architectures are reused, not counted again as new results.
The new native input is a persistent one-link action drop,
which supplies a local gradient witness at every chart time.
No third derivative of the pulled-back action, Taylor-time
radius, conditional inverse, or probability mixing estimate
is needed.

\subsection{A paid weighted Hessian for the pulled-back action}

Retain the strictly positive \(\zeta\), the weighted jet
constants \(A_1,A_2\) of \eqref{c18v63:flow}, and the
density constant \(K_*\) of \eqref{c18v60:budget}.
Put \(f_0=4e\), \(f_1=48e\), as in V60's Wilson-force jets,
and define
\[
 \begin{split}
 H_W&=f_1A_1^2+f_0A_2+2f_0A_1,\\
 K_W&=\max\{1,12e^{2\zeta}H_W\},\qquad
 D_*=\pi\sqrt2,\qquad S=D_*^2K_* .
 \end{split}
 \tag{V88.2}\label{c18v88:constants}
\]
All these constants are independent of \(n\) and \(g\);
they are finite but not numerically economical.

\begin{lemma}[Action locality and finite local conditional cost]
\label{c18v88:locality}
In fact on every fixed coarse fibre,
\[
 {\cal H}_\zeta(w_t)\le K_W,\qquad
 \operatorname{osc}_i u_t\le tS,\qquad
 u_t=\log\frac{d\nu_{g,t}}{d\mu_g}.
 \tag{V88.3}\label{c18v88:local}
\]
Here \({\cal H}_\zeta\) is exactly V63's weighted block-row
Hessian norm, with the configuration supremum outside the
row sum; \(\operatorname{osc}_i\) varies only factor \(i\).
For a block \(B\) of \(m\) factors, every exterior
configuration \(y\), and every measurable block event \(E\),
\[
 \nu_{g,t}(E\mid y)\ge e^{-mtS}\mu_B(E).
 \tag{V88.4}\label{c18v88:conditional}
\]
This is a finite-block comparison, not a global density ratio.
\end{lemma}

\begin{proof}
The ambient scalar gradient of \(W\) has component bound
\(f_0\), and its ordered second derivative has weighted
anchored norm at most \(f_1\). These are the first two
force-jet bounds in V60. The scalar pullback rule and
\eqref{c18v63:flow} give first-derivative bound \(f_0A_1\)
and ordered second-derivative bound
\(f_1A_1^2+f_0A_2\). Each contraction retains an external
derivative anchor; there is no undifferentiated extensive
\(W\) and no free volume sum in these estimates.
Conversion to the covariant Hessian costs at most
\(2f_0A_1\), from the same-link invariant-frame connection.
As in the proof of \eqref{c18v63:native}, the totally
geodesic pair/free inclusion costs at most
\(4e^{2\zeta}\) in scalar Schur norm, and conversion to
three-by-three operator blocks costs at most three.
This proves the first inequality.

For the second, V60 gives
\({\cal H}(u_t)\le tK_*\). With all other variables fixed,
\(u_t\) has a maximum on the compact \(i\)-th factor,
where its factor gradient vanishes. Parallel transport
along a minimizing factor geodesic of length at most
\(D_*\) gives
\(\|\nabla_i u_t\|_\infty\le D_*tK_*\).
Integrating between any two values of that factor gives
the claimed oscillation bound. Changing the \(m\) block
coordinates one at a time consequently gives
\(\operatorname{osc}_B u_t(\,\cdot\,,y)\le mtS\).
The normalized conditional density relative to \(\mu_B\)
is bounded below by the exponential of minus this
oscillation. This proves \eqref{c18v88:conditional}.
Smooth positive densities give this version for every
exterior, not merely almost every exterior.
\end{proof}

\subsection{A persistent native action drop}

Let \(p_+\) be the all-identity fine configuration.
For a free factor \(i\), let \(q_i(\alpha)\in M_I\)
have all other fine links equal to identity and its
\(i\)-th link equal to \(\exp(\alpha T)\), where \(T\)
is a unit imaginary quaternion and \(0\le\alpha\le\pi/2\).
This is a unit-speed \(m_0\) factor geodesic.

\begin{lemma}[A gradient witness at every chart time]
\label{c18v88:drop}
For every free factor \(i\) and every \(t\in[0,4]\),
there is \(\alpha_{i,t}\in(0,\pi/2)\) such that, writing
\(x_{i,t}=q_i(\alpha_{i,t})\),
\[
                 |\nabla_i w_t(x_{i,t})|\ge a_*,
                 \qquad a_*=\frac8\pi .
 \tag{V88.5}\label{c18v88:witness}
\]
The point may depend on \(i,t,n\); no measurable choice
of such points is required.
\end{lemma}

\begin{proof}
The identity configuration is stationary for the ambient
Wilson force \(F=\nabla W\), since every plaquette trace
is maximized there. Thus \(N_t(p_+)=Q_tF(p_+)=0\) at every
time, and uniqueness for the smooth normal-flow ODE gives
\(\eta_t(p_+)=p_+\). In particular \(w_t(p_+)=24v\).

A free link belongs to four distinct plaquettes. Along
\(q_i(\alpha)\) each of their traces is \(\cos\alpha\),
while every other trace is one; inverse orientations
give the same real trace. Also no captured coarse link
changes, so the entire curve stays in \(M_I\). Hence
\[
             W(q_i(\alpha))=24v-4(1-\cos\alpha).
\]
V86's exact normal-flow identity is
\[
       \partial_t w_t(x)=-|Q_tF|^2(\eta_t x)\le0.
 \tag{V88.6}\label{c18v88:decrease}
\]
Consequently \(w_t(q_i(\pi/2))\le24v-4\).
The ordinary mean-value theorem on \([0,\pi/2]\)
provides an interior point with directional derivative
at most \(-8/\pi\). The curve is unit speed in the free
factor, so its factor-gradient norm is at least \(8/\pi\).
\end{proof}

This is a finite action drop, not a small-time expansion.
The monotonicity is in normal-chart time \(t\); it is not
monotonicity of a physical correlation or of V87's
orthogonal dynamics.

\subsection{Arbitrary exterior and a volume-uniform energy floor}

Use V63's factor anchors and their periodic fine-lattice
Manhattan distance \(d(i,j)\). For a fixed integer radius
\(R\), the ball \(B_i=\{j:d(i,j)\le R\}\) has cardinality
at most \(6(2R+1)^3\). Indeed each periodic displacement
has an injectively chosen representative in the integer
box, and V63's harmless anchor-multiplicity bound is six.
The assertion also holds when the ball wraps around or
covers the entire torus.

Choose, once for the whole interval,
\[
 \begin{split}
 r&=\min\{1,a_*/(4K_W)\},\\
 R&=\max\left\{1,\left\lceil
       \zeta^{-1}\log(4D_*K_W/a_*)\right\rceil\right\},\\
 M_R&=6(2R+1)^3,\qquad
 p_r=\frac8{3\pi^3}\left(\frac r{\sqrt2}\right)^3,\\
 p_*&=\exp(-4SM_R)\,p_r^{M_R},\qquad
 c_W=18E_0^{-2}\frac{a_*^2}{4}\,p_* .
 \end{split}
 \tag{V88.7}\label{c18v88:floorconstant}
\]
Here \(E_0\ge1\) is the retained geometric metric-comparison
constant through time four:
\(E_0^{-2}m_0\preceq\widehat m_t\preceq E_0^2m_0\).
All constants in \eqref{c18v88:floorconstant} are
strictly positive and volume independent. Their nested
size is kept symbolic; floating-point underflow is not
a statement that \(c_W=0\).

\begin{theorem}[Native extensive action-energy nondegeneracy]
\label{c18v88:floor}
On the identity coarse fibre at \(\beta=0\),
\[
                  c_Wv\le{\cal D}(t)\le384v
                  \qquad(n\ge5,\ 0\le t\le4).
 \tag{V88.8}\label{c18v88:energyfloor}
\]
In particular \({\cal V}(t)>0\), so V87's normalized
action and memory decomposition are defined at all
these native parameters.
\end{theorem}

\begin{proof}
Fix \(i,t\) and the reference point \(x_{i,t}\) supplied
by Lemma \ref{c18v88:drop}. Let \(E_{i,t}\) be the block
event that every factor in \(B_i\) lies within \(m_0\)
distance \(r\) of its value at \(x_{i,t}\). There is no
restriction at all on the exterior.

Compare \(x_{i,t}\) with any full configuration in this
event by simultaneous minimizing geodesics in every
product factor. Their speeds are at most \(r\) inside
\(B_i\), and at most \(D_*\) outside. Parallel-transport
the \(i\)-th gradient along this product path.
Differentiation of that vector is the sum of the
Hessian blocks applied to the factor velocities.
At each single point of the path, the weighted row
bound \eqref{c18v88:local} therefore gives
\[
 \begin{split}
 |\nabla_i w_t(x)|
 &\ge |\nabla_i w_t(x_{i,t})|
       -rK_W-D_*e^{-\zeta R}K_W\\
 &\ge a_*/2\qquad (x\in E_{i,t}).
 \end{split}
 \tag{V88.9}\label{c18v88:gradientevent}
\]
The sum is taken before the pointwise supremum, precisely
as allowed by \({\cal H}_\zeta\). We have not replaced
it by a sum of separate block suprema or paid an
exterior-volume factor. Simultaneous geodesics remain
valid at antipodal endpoints by choosing any minimizing
geodesic; no uniqueness is needed.

For normalized Haar on the unit three-sphere, a ball of
radius \(\alpha\le\pi/2\) has measure
\[
 \frac2\pi\int_0^\alpha\sin^2s\,ds
       \ge\frac8{3\pi^3}\alpha^3,
\]
because \(\sin s\ge2s/\pi\) on that interval.
A ball of metric radius \(r\le1\) in a pair or free
factor has at least the unit-round ball measure of
radius \(r/\sqrt2\), hence at least \(p_r\).
If \(m=|B_i|\le M_R\), product Haar and
\eqref{c18v88:conditional} imply, for every exterior,
\[
 \nu_{I,t}(E_{i,t}\mid y)
      \ge e^{-mtS}p_r^m
      \ge e^{-4M_RS}p_r^{M_R}=p_* .
 \tag{V88.10}\label{c18v88:eventprobability}
\]
The same lower bound holds after integrating the exterior.
It is a local-event bound with fixed radius, not the
probability that the entire lattice is near \(p_+\).

Finally inversion of the metric comparison gives
\[
 \begin{split}
 {\cal D}(t)
 &\ge E_0^{-2}\int|\nabla_{m_0}w_t|^2\,d\nu_{I,t}\\
 &\ge E_0^{-2}\sum_{i\ {\rm free}}
             \int|\nabla_iw_t|^2\,d\nu_{I,t}\\
 &\ge18vE_0^{-2}(a_*^2/4)p_*=c_Wv .
 \end{split}
 \tag{V88.11}\label{c18v88:sum}
\]
The events for different \(i\) may overlap arbitrarily.
No independence and no disjoint-ball packing are used:
we sum distinct gradient components and individual
probability lower bounds, not joint events.
The upper bound is V86's paid energy bound.
If the variance were zero, the smooth action would be
constant under the strictly positive smooth law, hence
everywhere, contradicting the positive energy.
\end{proof}

The restriction \(g=I\) is substantive. The witness uses
a stationary all-identity configuration on this fibre.
The locality lemma holds at every \(g\), but it alone
does not supply an analogous action drop on an arbitrary
frustrated coarse fibre. No all-\(g\) energy floor is
claimed here.

\subsection{What this pays in the actual action-memory formula}

Retain V87's exact
\(d(t)={\cal D}(t)/{\cal V}(t)\) and return fraction
\(\theta(t)\in[0,1)\), and introduce the bare action scale
\[
             B(t)=\frac{{\cal V}(t)}{d(t)}
                  =\frac{{\cal V}(t)^2}{{\cal D}(t)}.
 \tag{V88.12}\label{c18v88:bare}
\]
Neither \(d(t)\) nor \(1-\theta(t)\) is asserted to
have a pointwise volume-uniform positive lower bound.

\begin{corollary}[Chart-time control of the bare scale]
\label{c18v88:budget}
For every \(0\le T\le4\),
\[
 \begin{split}
 \int_0^T\frac{dt}{d(t)}
      &\le\frac{2+32T}{c_W},\\
 \int_0^T\sqrt{\frac{B(t)}v}\,dt
      &\le\frac{2+32T}{\sqrt{c_W}} .
 \end{split}
 \tag{V88.13}\label{c18v88:timebudget}
\]
Consequently, for \(A>0\), Lebesgue chart-time measure
satisfies
\[
 \left|\{t\in[0,T]:B(t)>Av\}\right|
       \le\frac{2+32T}{\sqrt{c_WA}} .
 \tag{V88.14}\label{c18v88:exceptional}
\]
One may of course take the minimum with \(T\).
\end{corollary}

\begin{proof}
Pointwise,
\[
 \frac1d=\frac{{\cal V}}{{\cal D}}
       \le\frac{{\cal V}}{c_Wv},\qquad
 \sqrt{\frac Bv}
       =\frac{{\cal V}}{\sqrt{{\cal D}v}}
       \le\frac{{\cal V}}{\sqrt{c_W}\,v}.
\]
Integrate V86's exact variance budget
\eqref{c18v86:variancebudget}. On the exceptional set
the last square root exceeds \(\sqrt A\); integration
gives \eqref{c18v88:exceptional}.
\end{proof}

This is a weak chart-time estimate for \(B/v\), not a
pointwise estimate or an \(L^1(dt)\) bound for \(B/v\).
It does not choose a favorable time in place of the
prescribed endpoint \(t=4\), and \(t\) must not be
identified with auxiliary semigroup time \(s\) or
physical time.

For clarity, the precise remaining return requirement
can be read directly from V87:
\[
 {\cal S}_W(t)=\frac{B(t)}{1-\theta(t)},\qquad
 {\cal C}(t)=\frac{144B(t)}{1-\theta(t)}
                        -24{\cal V}(t)+{\cal D}(t).
 \tag{V88.15}\label{c18v88:remaining}
\]
If, at a particular time, \(B(t)\le Av\) and
\(1-\theta(t)\ge\varepsilon>0\), then
\[
              {\cal C}(t)\le(144A/\varepsilon+384)v .
 \tag{V88.16}\label{c18v88:conditionalcost}
\]
Thus the set where this displayed source bound fails
is contained in the union of the exceptional set in
\eqref{c18v88:exceptional} and
\(\{t:1-\theta(t)<\varepsilon\}\).
There is no estimate for the latter set in this append.
An extensive energy numerator cannot exclude a small
spectral component with a very long auxiliary-time
return; V87's comparison already prohibits that inference.

At zero time the inherited calibration is
\({\cal V}(0)=6v\), \({\cal D}(0)=63v\),
\(d(0)=21/2\), and \(B(0)=4v/7\).
These calibrate the units and the square of the variance
in \eqref{c18v88:bare}; they are not new Taylor data.

\subsection{Verification, non-duplication, and next obligation}

The companion regression
\path{scripts/verify_ym_c18_v88_action_energy_floor.py}
checks the native periodic incidence at \(n=5,6\):
there are \(18v\) free factors, every one has exactly
four incident plaquettes, and all single-free-link
action drops have the asserted formula. Exact quaternion
evaluations include both a right-angle link and a
nontrivial rational quaternion. It also checks finite
torus anchor-ball counts, the derivative sign for
orthogonal normal projections, the Haar-ball lower
bound, and the algebra of the chart-time and source
implications. These are regressions, not an all-volume
numerical proof, a finite-time native probability
simulation, or proof-assistant certification. The
all-volume argument is the one given above.

No new external gap or mixing theorem is imported in
this append. The scheduled companion and broad
primary-literature checkpoints remain V90; V87's
targeted memory references do not supply the unpaid
return margin. The unchanged historical corpus is
preserved rather than independently recertified.

The direction is now more specific: at \(g=I,\beta=0\)
the extensive action-energy denominator is paid, and
the bare scale has the explicit chart-time control
above. A useful next step must address actual action
fluctuations at the prescribed time or the integrated
orthogonal return, not repeat local density oscillation
or another abstract Schur identity. Extension of the
witness to arbitrary coarse fields is a separate
native question, not an automatic corollary.

The volume-uniform low-energy inverse, full source
Poisson Hessian and probability strain, interacting
vacuum correction, physical scaling and clock, and
nontrivial reflection-positive four-dimensional
continuum construction with positive physical mass
gap remain unproved. The full YM goal remains active.

Removing this append and restoring the title recovers
V87 byte for byte. Only V88 remains as the active
source; the other notes and legacy archives are
unchanged. Build products stay outside
\path{latex_notes}, which contains only \texttt{.tex}.
The cycle's archive/restart point remains V100.
% END CYCLE 18 VERSION 88 APPEND
% BEGIN CYCLE 18 VERSION 89 APPEND
\clearpage
\section{V89: Global compact action control and the exact inverse-flow density}
\label{c18v89:main}

\subsection{An upstream change of direction, with the old results credited}

V88 pays the action-energy denominator at \(g=I,\beta=0\),
but its local gradient events do not upper-bound endpoint
action fluctuations or distant covariances. The attempt to
promote that proof directly to such a bound is not justified.
We instead expose a global property of the \emph{actual}
conditional density, using the rooted inverse-flow coordinates
already constructed in V33--V34. No probability is replaced.

Tree--chord curvature fillings are inherited architecture.
The closure monograph's curvature-to-defect bridge, and
f14 V39's explicit acknowledgement of that bridge, were
rechecked. Also f15 V11 already proves a quadratic,
scale-correct rooted-fibre estimate for its four-dimensional
averaging map. We do not claim either general principle anew.
Here we specialize the compact filling to the current
three-dimensional block-two straight sampler, with an explicit
constant, and propagate it through the \emph{exact} finite
Wilson flow and its full Jacobian. The new conclusion is a
global potential comparison for the native conditional law,
not another flat Hessian or a selected rare-event estimate.

Keep \(n\ge5\), \(v=n^3\), fine side \(2n\), \(\beta=0\)
and \(g=I\). Put
\[
 S(U)=24v-W(U)=\sum_p[1-\operatorname{Re}U_p],\qquad
 q(V)=1-\operatorname{Re}V\quad(V\in\mathrm{SU}(2)).
 \tag{V89.1}\label{c18v89:defect}
\]
The real part is the normalized fundamental trace.
Thus \(0\le q\le2\), with \(q(V)=0\) precisely at \(V=I\).

For primary literature context, the abstract of
Shen--Zhu--Zhu,
\href{https://arxiv.org/abs/2204.12737}
{\emph{A stochastic analysis approach to lattice Yang--Mills
at strong coupling}}, explicitly assumes a small Wilson
coupling in its functional inequalities. No theorem from it
is imported here. In particular our \(\beta=0\) specifies the
ambient electric vacuum; it does not make the conditional
flowed density at chart time four Haar or put that induced
density under the paper's Wilson-coupling hypothesis.
The targeted search does not replace the broad V90 review.

\subsection{The actual rooted coordinates and a compact filling constant}

Use the particular forest of V34. Write each fine vertex
as \(x=2X+r\), \(r\in\{0,1\}^3\). Its parent reduces the
first nonzero component of \(r\). The root-to-vertex path
therefore uses coordinate directions \(3,2,1\), omitting
zero components. After the exact root-fixed gauge fixing,
all forest links and all last coarse-chain links are \(I\)
at \(g=I\). The remaining \(14v\) chord links \(y_e\)
are independent product-Haar coordinates in the reference
measure. Denote this representative by \(U^0(y)\), and put
\[
                         Q(y)=\sum_{e\in{\cal C}}q(y_e).
 \tag{V89.2}\label{c18v89:Q}
\]
This reference factorization is V33's; it is not a
factorization of the native interacting conditional density.

\begin{lemma}[Compact, noncommutative filling on the current fibre]
\label{c18v89:filling}
At every chord configuration, with no small-field restriction,
\[
                 \frac18 Q(y)\le S(U^0(y))\le16Q(y).
 \tag{V89.3}\label{c18v89:fillingbound}
\]
The constants are independent of \(n\).
\end{lemma}

\begin{proof}
Multiplication and conjugation by unit quaternions preserve
Euclidean chord distance. Since
\(q(V)=|V-I|_{\mathbb R^4}^2/2\), telescoping and
Cauchy--Schwarz give
\[
 q(V_1\cdots V_m)\le m\sum_{j=1}^m q(V_j),
 \qquad q(HVH^{-1})=q(V)=q(V^{-1}).
 \tag{V89.4}\label{c18v89:word}
\]
For the upper bound, in rooted gauge every plaquette is
a word of at most four chord links and their inverses.
Each chord is incident to four plaquettes. Summing
\eqref{c18v89:word} proves \(S\le16Q\).

For the lower bound, an edge with tail \(2X+r\) and
direction \(i\) is completed to a root-based loop by
its two forest paths and, if it crosses to the next
cell, the inverse coarse chain. That coarse chain
has product \(I\). Its loop holonomy in rooted gauge
is therefore precisely \(y_e\).

Here is an explicit plaquette filling with its entire
incidence cost. Translate \(2X\) to zero. Let \(P(r)\)
be the word of directions \(3,2,1\) with multiplicities
\(r_3,r_2,r_1\). Compare the two positive paths
\[
 \begin{array}{c|c|c}
 & A & B\\ \hline
 r_i=0 & P(r)\,i & P(r+e_i)\\
 r_i=1 & P(r)\,i & i\,i\,P(r-e_i).
 \end{array}
 \tag{V89.5}\label{c18v89:paths}
\]
They have the same endpoints; \(A B^{-1}\) is the loop
just described. Transform \(A\) into \(B\) by reading
\(B\) from left to right, at each step moving the next
matching letter of \(A\) left by adjacent swaps.
A swap \(P\,a\,b\,R\longmapsto P\,b\,a\,R\)
has ratio holonomy equal to the elementary
\((a,b)\)-plaquette based at the endpoint of \(P\),
conjugated by the holonomy along \(P\).
Telescoping all successive path ratios expresses
\(y_e\) as an ordered product of these conjugated
plaquettes. No commutativity or Abelian Stokes formula
is used.

There are \(m_e=\sum_{j<i}r_j\) swaps when \(r_i=0\),
and \(m_e=\sum_{j<i}r_j+2\sum_{j>i}r_j\) when \(r_i=1\).
Exactly the seven forest and three last-chain edge
types have \(m_e=0\); the other fourteen are the
chords, with \(1\le m_e\le4\).
All face bases remain in the half-open cell
\(2X+\{0,1\}^3\), even for a crossing chord.
For complete transparency the following table gives
the weighted incidence
\(\sum_{e:p\ {\rm in\ its\ filling}}m_e\), counting
multiplicity, at each parity and each oriented plane.
It follows by the prescribed adjacent swaps.
\[
\begin{array}{c|rrr}
 \text{face-base parity}&12&13&23\\ \hline
 000&3&7&8\\
 001&5&1&3\\
 010&1&2&5\\
 011&3&2&0\\
 100&2&6&0\\
 101&4&0&0\\
 110&0&0&0\\
 111&0&0&0
\end{array}
 \tag{V89.6}\label{c18v89:incidence}
\]
Its largest entry is eight. Each plaquette base belongs
to one such half-open cell, so no additional neighbouring-cell
multiplicity occurs. Applying \eqref{c18v89:word} to every
chord filling and summing gives \(Q\le8S\).
The paths and swaps lift to the integer lattice and
then reduce periodically; their coordinate extents
are at most two. Thus the same argument applies at
periodic seams for all the stated \(n\).
\end{proof}

The table is a native incidence specialization of the
inherited filling method. In particular it proves
that \(S=0\) on this fibre precisely on the root-fixed
gauge orbit of the identity. There is no residual
physical toron in this claim: every coarse chain is
fixed to \(I\), so products along every periodic
coarse cycle are also \(I\). This is not a claim for
unconstrained torus chords or arbitrary coarse holonomies.

\subsection{The energy deficit survives the full nonlinear flow}

Let \(F=\nabla W\) in the original fine product metric,
let \(\Phi_s\) be its forward flow, and let \(\eta_s\)
be V51's normal identification. These are different
flows and must remain distinct.

\begin{lemma}[Global deficit growth bounds, not a flat approximation]
\label{c18v89:growth}
For every fine configuration,
\[
                          |F|^2\le32S.
 \tag{V89.7}\label{c18v89:force}
\]
For either \(S_s=S(\Phi_{-s}U)\), or
\(S_s=S(\eta_s U)\), one has
\[
                S(U)\le S_s\le e^{32s}S(U)
                         \qquad(s\ge0).
 \tag{V89.8}\label{c18v89:growthbound}
\]
For the normal flow this assertion uses its smooth
extension to each finite time interval, as in V86.
It does not extend a quantitative jet bound to
unbounded time.
\end{lemma}

\begin{proof}
For one plaquette trace \(w_p\), varying any of its
four unit-round links gives a linear unit-quaternion
coordinate, so \(|\nabla_e w_p|^2=1-w_p^2\).
There are four plaquettes per link. Consequently
\[
 \begin{split}
 |F|^2
 &\le4\sum_e\sum_{p\ni e}|\nabla_e w_p|^2
   =16\sum_p(1-w_p^2)\\
 &\le32\sum_p(1-w_p)=32S .
 \end{split}
\]
Along the inverse full flow,
\(\partial_s S_s=|F|^2(\Phi_{-s}U)\).
Along the normal flow,
\(\partial_s S_s=|Q_sF|^2(\eta_s U)\), with \(Q_s\)
the orthogonal normal projection of V51.
Both derivatives lie between zero and \(32S_s\).
Integration and Gronwall prove \eqref{c18v89:growthbound}.
\end{proof}

In particular the actual normal-chart action obeys,
on the same rooted representative of its initial fibre,
\[
 \frac18 Q(y)\le24v-w_t(U^0(y))
                  \le16e^{32t}Q(y).
 \tag{V89.9}\label{c18v89:normalwell}
\]
Thus its maximum is the identity root-fixed gauge orbit,
globally, not only in a quadratic neighbourhood.
Gauge-related identity configurations are stationary
under both flows. Uniqueness here concerns the action
maximum, not the maximum of its probability density
in an arbitrary chart and not uniqueness of an
infinite-volume Gibbs state.

\subsection{A global envelope for the exact conditional potential}

Now use the \emph{inverse full-flow} rooted chart
\[
        \chi_t(k,y)=\Phi_{-t}(k^{-1}\!\cdot U^0(y))
\]
of \eqref{c18v33:law}. At \(\beta=0\) the ambient law
is Haar. V32's exact Liouville identity gives
\[
 \begin{split}
 h_t(U^0(y))
 &=\operatorname{Jac}_{\rm amb}\Phi_{-t}(U^0(y))
      =\exp\{288vt-{\cal U}_t(y)\},\\
 {\cal U}_t(y)&=12\int_0^t S(\Phi_{-s}U^0(y))\,ds,\\
 d\nu_{I,t}&=(\chi_t)_\#\left(
       dk\,\frac{e^{-{\cal U}_t(y)}}{{\cal Z}_t}\,dy\right),
 \qquad {\cal Z}_t=\int e^{-{\cal U}_t(y)}\,dy .
 \end{split}
 \tag{V89.10}\label{c18v89:actualdensity}
\]
Here \(dk,dy\) are normalized product Haar, and \(dk\)
is exactly independent. The displayed equality for
\(\nu_{I,t}\) is in the original moving fibre; it is
unitarily equivalent at each fixed time to its normal
pullback. This notation does not set \(\chi_t=\eta_t\).

\begin{theorem}[Global compact potential bounds for the native law]
\label{c18v89:potential}
For every \(t\ge0\) and every \(y\), set
\[
               c_-(t)=\frac32t,\qquad
               c_+(t)=6(e^{32t}-1).
\]
Then
\[
                    c_-(t) Q(y)\le{\cal U}_t(y)\le c_+(t) Q(y).
 \tag{V89.11}\label{c18v89:potentialbound}
\]
For \(t>0\) the product-Haar potential \({\cal U}_t\)
has the unique minimum \(y=I\); this does not assert
global convexity of that potential.
The corresponding density has its unique maximum
there in these declared coordinates.
\end{theorem}

\begin{proof}
For each \(s\), \eqref{c18v89:growthbound} gives
\[
 S(U^0(y))\le S(\Phi_{-s}U^0(y))
                       \le e^{32s}S(U^0(y)).
\]
Integrate, then apply \eqref{c18v89:fillingbound}:
the lower coefficient is \(12t/8=3t/2\);
the upper one is
\(12\cdot16(e^{32t}-1)/32=6(e^{32t}-1)\).
The potential vanishes at \(y=I\), which is stationary,
and for \(t>0\) its lower bound is positive elsewhere.
\end{proof}

This is an everywhere comparison for the native
potential, not a near-central cap estimate. It also
identifies the correct limitation of a single-well
intuition: absence of another global minimum does
not control entropic competition, secondary local
wells, a log-Sobolev constant or the complete source
inverse.

The normal-chart raw density at the flat point is
V79's compressed determinant, not \(e^{288vt}\).
The latter is the \emph{ambient inverse-flow} Jacobian
in \eqref{c18v89:actualdensity}. Their difference is
the tangential chart Jacobian. Neither their
pointwise values nor their modes may be interchanged.

\subsection{Normalization is retained, and the spectral inference is not made}

Define the one-factor Haar integral
\[
 z(c)=\int_{\mathrm{SU}(2)}e^{-cq(V)}\,dV
     =\frac2\pi\int_0^\pi
             e^{-c(1-\cos\alpha)}\sin^2\alpha\,d\alpha,
 \qquad c\ge0,\quad z(0)=1 .
 \tag{V89.12}\label{c18v89:z}
\]
Product Haar and \eqref{c18v89:potentialbound} give
the genuine normalizer bounds
\[
 z(c_+(t))^{14v}\le{\cal Z}_t\le z(c_-(t))^{14v},\qquad
 Z_I(t)=e^{288vt}{\cal Z}_t .
 \tag{V89.13}\label{c18v89:partition}
\]
The last \(Z_I\) is the actual coarse density relative
to coarse Haar, as in V86; both charts disintegrate
the same ambient probability.

For \(t>0\), \(0\le\lambda<c_-(t)\), and \(r\ge0\), the
normalized probability of a whole-fibre chord-defect
event has the explicit bound
\[
 \nu_{I,t}\{Q(y)\ge14vr\}
 \le\min\left\{1,
  \exp\left(14v\left[
       \log z(\lambda)-\log z(c_+(t))
                         -(c_-(t)-\lambda)r\right]\right)\right\}.
 \tag{V89.14}\label{c18v89:tail}
\]
The event is read in the declared \(y\) coordinates.
Indeed, on it,
\(e^{-{\cal U}_t}\le e^{-c_-(t) Q}
 \le e^{-14v(c_-(t)-\lambda)r}e^{-\lambda Q}\);
integrate the numerator and use the lower normalizer.
No independence under the native law was assumed.

The entropy term \(-14v\log z(c_+(t))\) is extensive.
In particular \(c_+(4)=6(e^{128}-1)\) is enormous,
while \(c_-(4)=6\). We do not assert that
\eqref{c18v89:tail} supplies a useful small exceptional
set at the desired endpoint. Dropping its normalizer
would manufacture such a conclusion. A bounded
quadratic envelope for an extensive potential does
not bound the oscillation of its difference from a
product potential uniformly in volume.

There is a further chart distinction relevant to
V87's target. In the present product-Haar chart set
\[
 h_t(y)=W(\Phi_{-t}U^0(y)),\qquad f_t=h_t-\nu h_t.
\]
Differentiating \eqref{c18v89:actualdensity} at fixed
\(y\) gives the normalized density score \(12f_t\).
The \emph{normal-chart} score, expressed at the same
physical point, remains \((12-L_t)f_t\) by V86.
Their difference is \(L_t f_t\), the already paid
vertical-divergence correction, not zero.
Also \(L_t\) retains the induced metric and quotient
kinetic cross terms of V33--V34, not a product Haar
Laplacian. The fixed-time distribution and variance
of the physical action are chart invariant; these
time derivatives and reference-density modes are not.

Consequently V88's energy floor still needs endpoint
control of \({\cal V}\), and V87 still needs a bound
for the integrated orthogonal return.
The present envelope alone proves neither.
It turns the next probability problem into one about
the explicitly identified native potential
\({\cal U}_4\), including its normalizer and metric;
it does not replace it by the quadratic lower
envelope \(6Q\).

\subsection{Checks and next checkpoint}

The regression
\path{scripts/verify_ym_c18_v89_global_action_well.py}
enumerates all twenty-four parity/orientation edge types,
constructs the prescribed adjacent-swap fillings and
checks the fourteen weighted face incidences in
\eqref{c18v89:incidence}. It checks exact noncommuting
quaternion path-ratio identities, including the
conjugating prefixes, and periodic \(n=5,6\) native
rooted configurations. It also checks the force-defect
constants, potential integration and partition/tail
algebra. These finite checks supplement the written
all-volume proof; they do not simulate the native
positive-time probability or certify a spectral gap.

The previous goal turn was progress: V88 proved and
verified a native positive-time energy floor.
This turn does not repeat a flat coercivity calculation
or claim the unsuccessful endpoint-variance inference.
Its new all-configuration native density comparison
exposes the remaining normalization/correlation issue.
At V90 the scheduled companion-file review and broad
primary-literature sweep must precede choosing the
next import. Relevant directions now include native
potential curvature with defects, normalized local
replacement costs, and actual source-weighted response,
with the full induced metric retained.

The interacting vacuum extension, continuum coupling
and physical clock, nontrivial reflection-positive
four-dimensional construction, and positive physical
YM mass gap remain unproved. The goal stays active.

Removing this append and restoring the title recovers
V88 byte for byte. Only V89 is the active source;
the other notes and legacy archives are unchanged.
Only \texttt{.tex} files remain in
\path{latex_notes}; verification and build products
stay outside. Archive/restart remains V100.
% END CYCLE 18 VERSION 89 APPEND
% BEGIN CYCLE 18 VERSION 90 APPEND
\clearpage
\section{V90: Boundary deletion, normalized block comparison, and source pressure}
\label{c18v90:main}

\subsection{Checkpoint and precise change of direction}

The scheduled V90 review rechecked twelve current companion files against
the V85 hashes; none changed. Their terminal results were inspected, not
silently promoted to nonperturbative YM estimates. In particular, the
general ESM V44 exceptional expansion and ESM V80 two-loop bosonic closure
are coefficientwise; source-selection V65 still needs a populated native
branch; stationarity V61 and Einstein-bridge V55 leave their actual slow
closure inputs unpaid. SM-continuum V72 is a fixed-time many-copy limit.
NS V26 supplies kernel norms, not a YM correlation inequality. Shell-mixing
V16 retains its noncommuting Born and nonperturbative remainder problems.
The arithmetic files and ESM phenomenology likewise supply no missing
native probability estimate. The parallel YM V5 warning about extensive
global labels remains relevant: use densities and finite windows, not
volume-uniform tails for an extensive total.

The broad primary-literature screening covered compact-group loop
equations, Wilson flow, classical dynamical locality, nonconvex Gibbs
inequalities, entropy factorization, orthogonal dynamics, and Markov
decomposition. These were abstract/scope checks, not independent audits
of the papers' proofs. The concrete decisions are:
\begin{itemize}
\item \href{https://arxiv.org/abs/2309.07399}{Abdelghani--Nissim's
finite-\(N\) master-loop derivation} supports intrinsic Haar integration
by parts, but the attachment's V58 cubic target is historical; it is not
the current endpoint return problem.
\item \href{https://arxiv.org/abs/2503.20457}{Widder--Zimmer--Schilling,
revised April 2026}, justify rank-one orthogonal dynamics under their
operator hypotheses. V87 already supplies the needed native rank-one
identity. Well-posed memory does not bound its time integral.
\item \href{https://arxiv.org/abs/2608.14526}{Buchholz--Cotar--Schweiger}
use a particular nonconvex gradient-interface/convex-mixture structure.
No such representation of our compact flowed law has been proved.
\item \href{https://arxiv.org/abs/2004.10574}{Caputo--Parisi} assume
strong spatial mixing; \href{https://arxiv.org/abs/2503.19419}
{Caputo, Chen and Parisi} require high temperature or weak interactions
for their stated factorizations. A surface comparison below is not
either hypothesis.
\item \href{https://arxiv.org/abs/1907.12308}{Bauerschmidt--Bodineau}
require control of the actual Polchinski evolution.
\href{https://arxiv.org/abs/2504.01247}{Qin's September 2026 revision}
retains component and comparison-chain gaps. Neither manufactures
our unpaid source inverse or the physical vacuum.
\item \href{https://arxiv.org/abs/2204.12737}{Shen--Zhu--Zhu}
assume small Wilson coupling for their lattice result.
\href{https://arxiv.org/abs/0907.5491}{L\"uscher's flow construction}
does not identify ordinary Wilson flow as an exact probability
trivialization.
\item \href{https://arxiv.org/abs/0902.0025}{Raz--Sims' classical
locality estimates} suggest comparing full and cut dynamics.
Their oscillator theorem is not imported: the compact first-order
comparison needed here is proved directly below.
\end{itemize}

V68 and V79 already cut boundary plaquettes in stationary Hessian
determinants and count the resulting faces. V63 and V71--V72 already
give spatial density/canonical-path estimates. We credit all of them.
The new object here is the \emph{full nonlinear inverse-flow potential}
of V89 at arbitrary compact configurations, together with its actual
normalizer and its Wilson-source tilt. This gives a normalized block
comparison and a thermodynamic pressure limit, but not a susceptibility
or mass-gap theorem.

\subsection{A direct all-configuration comparison of cut Wilson flows}

Keep the V89 setting: \(\mathrm{SU}(2)\), \(\beta=0\), \(g=I\),
\(n\ge5\), \(v=n^3\), unit-round fine-link metrics, and the exact rooted
representative \(U^0(y)\) with \(14v\) physical chords. For a subset
\(\mathcal P'\) of the \(24v\) fine plaquettes put
\[
 W'=\sum_{p\in\mathcal P'}w_p,\quad
 S'=\sum_{p\in\mathcal P'}(1-w_p),\quad
 \Psi_s'=\text{flow for time \(s\) of }-\nabla W'.
 \tag{V90.1}\label{c18v90:cut}
\]
The full inverse flow is \(\Psi_s=\Phi_{-s}\).
Only \(b=|\mathcal P\setminus\mathcal P'|\) terms are removed;
the fine-link configuration space is unchanged.

\begin{lemma}[Compact plaquette Hessian and nonlinear deletion cost]
\label{c18v90:flow}
For every \(s\ge0\) and every initial configuration \(U\),
\[
 \sum_e d\bigl((\Psi_sU)_e,(\Psi_s'U)_e\bigr)
       \le \frac b4(e^{16s}-1),
 \qquad
 |S(\Psi_sU)-S'(\Psi_s'U)|\le b B_s,\quad B_s=1+e^{16s}.
 \tag{V90.2}\label{c18v90:flowbound}
\]
These are uniform in the volume and the configuration.
\end{lemma}
\begin{proof}
A single normalized plaquette trace is linear in the four real
quaternion coordinates of any one of its links, including a link
occurring inversely. Its one-link gradient has norm at most one and
its diagonal covariant Hessian is \(-w_p\) times the round metric.
For two distinct links, differentiating their quaternion factors
in unit tangent directions gives a real part of a quaternion
product of norm at most one. Thus every mixed Hessian block has
operator norm at most one. Each plaquette has four distinct links
and each link is in four plaquettes. Consequently
\[
 \sup_j\sum_i\|(\nabla^2 W_\theta)_{ij}\|_{\rm op}\le16,
 \quad
 W_\theta=W'+\theta(W-W'),\quad 0\le\theta\le1.
 \tag{V90.3}\label{c18v90:Hess}
\]
The same row bound follows by symmetry. This is a covariant
block bound, not an entrywise sum with an omitted color factor.

Let \(X_\theta(s)\) solve \(\dot X_\theta=-\nabla W_\theta(X_\theta)\)
with the common initial condition. Covariant differentiation in
\(\theta\) gives
\[
 \nabla_s V_\theta
   =-\nabla^2W_\theta\,V_\theta-\nabla(W-W')(X_\theta),\qquad
 V_\theta(0)=0.
\]
Parallel frames in each round factor preserve norms. The forcing
has summed link norm at most \(4b\); (V90.3) and the upper Dini
derivative of \(\sum_e|V_{\theta,e}|\) give
\(\sum_e|V_{\theta,e}(s)|\le b(e^{16s}-1)/4\).
Integration in \(\theta\) bounds the sum of endpoint distances.
This argument also avoids any differentiation across a cut locus.

Both \(S\) and \(S'\) have link-gradient norm at most four.
At any fixed configuration \(|S-S'|\le2b\).
Apply the resulting \(4\)-Lipschitz bound in the summed distance:
\(2b+4[b(e^{16s}-1)/4]=b(1+e^{16s})\).
All finite-time flows exist by compactness.
\end{proof}

Define the exact potential and the endpoint action observable
\[
 \mathcal U_{n,t}(y)=12\int_0^t S(\Psi_sU^0(y))\,ds,\qquad
 E_{n,t}(y)=S(\Psi_tU^0(y)).
 \tag{V90.4}\label{c18v90:UE}
\]
Here \(E_{n,t}=24v-W\) at the corresponding physical configuration.
Its variance is exactly V86's \(V\), although its expression is
in the full inverse-flow chart rather than the normal chart.
Use primes for the quantities with both \(S'\) and \(\Psi'\).
For \(t\ge0\), \(h\in\mathbb R\), put
\[
 A_t=12t+\frac34(e^{16t}-1),\qquad
 C_t(h)=A_t+|h|B_t .
 \tag{V90.5}\label{c18v90:constants}
\]
Integrating (V90.2) gives the pointwise, all-configuration estimates
\[
 |\mathcal U_{n,t}-\mathcal U'_{n,t}|\le A_tb,\qquad
 |(-\mathcal U_{n,t}+hE_{n,t})
     -(-\mathcal U'_{n,t}+hE'_{n,t})|\le C_t(h)b .
 \tag{V90.6}\label{c18v90:weight}
\]
The source parameter \(h\) is not chart time, semigroup time,
a physical clock, or an already-calibrated YM coupling.

\subsection{Normalized local comparison, with every exterior retained}

Partition fine links by the coarse cell of their stored tail.
For a cube of coarse side \(\ell<n\), delete precisely the
plaquettes whose link tails lie on both sides. The count already
proved in V68 is \(b_R=48\ell^2-6\ell\).
In rooted coordinates the initial field is a product of the
interior and exterior chord assignments; all other initial links
are \(I\). After deletion, the vector field splits on the two
disjoint sets of stored links. Hence \(\mathcal U'\) and \(E'\)
are sums of an interior and an exterior quantity, at every
configuration, not only at a stationary center.

Let \(\mu_{n,t,h}\) be the normalized chord law with density
proportional to \(\exp(-\mathcal U_{n,t}+hE_{n,t})\).
At \(h=0\) this is V89's actual native law in the declared chart.
Let \(\mu_{\ell,t,h}^{\rm o}\) denote the corresponding interior
cut law, with normalized product Haar reference on its
\(14\ell^3\) chords. No gauge-volume factor is lost: the rooted
gauge variables integrate to one, and the initial forest/last
links remain the same prescribed \(I\) values.

\begin{proposition}[An exterior-uniform, normalized block comparison]
\label{c18v90:conditional}
For every exterior chord assignment \(\xi\), the smooth positive
conditional density satisfies
\[
 e^{-2C_t(h)b_R}\le
 \frac{d\mu_{n,t,h}(\,\cdot\mid\xi)}
      {d\mu_{\ell,t,h}^{\rm o}}
 \le e^{2C_t(h)b_R}.
 \tag{V90.7}\label{c18v90:conditionalbound}
\]
The conditional law on a measure-zero exterior is understood by
its continuous positive density formula.
\end{proposition}
\begin{proof}
Subtract the cut interior-plus-exterior exponent. By (V90.6)
the remainder lies between \(-C_t(h)b_R\) and \(C_t(h)b_R\).
The exterior cut exponent cancels in conditional normalization.
The integral of the remaining perturbation against the cut
interior law lies between the same two exponentials. Dividing
accounts for the factor two. This proves the assertion for all
exteriors without an independence assumption for the native law.
\end{proof}

This pays a previously explicit normalization step for arbitrary
compact interiors. It is not strong spatial mixing: its exponent
grows with the boundary area and does not decay as a boundary is
moved away from a local observable. In particular it supplies
neither the Caputo--Parisi hypothesis nor a return contraction.
At \(t=4\), \(A_4=48+\tfrac34(e^{64}-1)\) is enormous.
No useful numerical block-mixing constant is asserted.

\subsection{Open-block pressure and a full thermodynamic limit}

Define
\[
 \mathcal Z_n(t,h)=\int e^{-\mathcal U_{n,t}+hE_{n,t}}\,dy,\qquad
 p_n(t,h)=n^{-3}\log\mathcal Z_n(t,h).
 \tag{V90.8}\label{c18v90:pressure}
\]
For an isolated open cube of coarse side \(\ell\), use the
stored-link variables of its \( (2\ell)^3\) fine vertices,
the same parity-dependent rooted initial assignment, and just
the plaquette words all of whose link tails stay in the cube.
Unused outward links are retained as variables; their vector
field is zero. This defines \(\mathcal Z_\ell^{\rm o}\) and
\(p_\ell^{\rm o}=\ell^{-3}\log\mathcal Z_\ell^{\rm o}\)
independently of any containing torus.

\begin{theorem}[Native source pressure with an explicit surface error]
\label{c18v90:thermo}
For each fixed finite \(t\ge0\) and real \(h\), there exists
\[
 p(t,h)=\lim_{n\longrightarrow\infty}p_n(t,h).
\]
For every \(\ell,n\ge5\),
\[
 |p(t,h)-p_\ell^{\rm o}(t,h)|\le\frac{24C_t(h)}{\ell},
 \qquad
 |p_n(t,h)-p(t,h)|\le\frac{48C_t(h)}{n}.
 \tag{V90.9}\label{c18v90:pressureerror}
\]
The convergence is locally uniform in \(h\), and \(p(t,\cdot)\)
is convex, nondecreasing, and \(48\)-Lipschitz. At zero source,
the actual coarse-density pressure has the limit
\[
 \lim_{n\to\infty}n^{-3}\log Z_{I,n}(t)=288t+p(t,0).
 \tag{V90.10}\label{c18v90:coarsepressure}
\]
No uniqueness of an infinite-volume Gibbs state is part of this
statement.
\end{theorem}
\begin{proof}
First let \(n=k\ell\), \(k\ge2\), and cut every plaquette crossing
the tiling by \(\ell\)-cubes. The retained count per cube is
\(3(2\ell)(2\ell-1)^2\); equivalently V79's periodic/open
deletion count is \(24\ell^2-6\ell\). Thus
\[
 b=\frac{n^3}{\ell^3}(24\ell^2-6\ell)
     \le\frac{24n^3}{\ell}.
\]
The cut exponent is a sum of independent-block functions in
the reference product coordinates. Its partition function is
\((\mathcal Z_\ell^{\rm o})^{k^3}\).
Integrate the pointwise bound (V90.6), take logarithms, and divide
by \(n^3\):
\[
 |p_n(t,h)-p_\ell^{\rm o}(t,h)|
             \le24C_t(h)/\ell\qquad(\ell\mid n,\ n>\ell).
 \tag{V90.11}\label{c18v90:divisible}
\]

For general \(n>\ell\), partition each coarse coordinate into
intervals of length \(\ell\), with a final remainder if needed.
There are at most \(n/\ell+1\) boundary planes per axis.
Counting the two possible crossing directions of each fine
plaquette gives \(b\le24n^3(1/\ell+1/n)\).
The regular cubes occupy \((\ell\lfloor n/\ell\rfloor)^3\)
coarse cells, so the remainder occupies at most \(3\ell n^2\).
For any cut rectangle of \(r\) coarse cells,
\[
 0\le E'\le48r,\quad 0\le\mathcal U'\le576tr,\quad
 |\log\mathcal Z'|\le M_t(h)r,\qquad
 M_t(h)=576t+48|h|.
\]
The same bound holds for \(|p_\ell^{\rm o}|\).
Replacing the remainder contribution by that fraction of
\(p_\ell^{\rm o}\) therefore costs at most \(6M_t(h)\ell/n\).
Consequently
\[
 |p_n-p_\ell^{\rm o}|
 \le24C_t(h)(1/\ell+1/n)+6M_t(h)\ell/n.
 \tag{V90.12}\label{c18v90:remainder}
\]
The \(p_n\) are bounded. For each fixed \(\ell\), (V90.12)
places their limsup and liminf in an interval of width
\(48C_t(h)/\ell\). Letting \(\ell\to\infty\) proves existence,
and then (V90.12) proves the first bound in (V90.9).

Finally compare the periodic \(n\)-cube with its own open cube.
Removing its wrap plaquettes costs \(24n^2-6n\) terms, so
\(|p_n-p_n^{\rm o}|\le24C_t(h)/n\).
Combine this with the first bound at \(\ell=n\).
Since \(C_t(h)\) is bounded on compact \(h\)-intervals,
convergence is locally uniform. Positivity and
\(0\le E_{n,t}\le48n^3\) give convexity, monotonicity, and the
Lipschitz constant by differentiation at finite \(n\) and passage
to the limit. V89's exact identity
\(Z_{I,n}(t)=e^{288tn^3}\mathcal Z_n(t,0)\) proves (V90.10).
\end{proof}

The open-block integral is a fully specified finite compact
integral, not an evaluated number. Thus (V90.9) is a rigorous
finite-block enclosure if that integral is bounded separately;
it is not a claimed efficient computation at \(t=4\).
At \(h=0\), V89 still gives
\(14\log z(c_+(t))\le p(t,0)\le14\log z(c_-(t))\).
The new theorem replaces the need to normalize a whole arbitrarily
large torus by an open-block normalization with an explicit
surface correction; it does not assert that the correction is small
on practical block sizes.

\subsection{Actual source tails and the susceptibility obstruction}

Because the normalizer is retained, for \(h>0\) and real \(r\)
one has the exact exponential-Markov bound
\[
 \mu_{n,t,0}\{E_{n,t}/n^3\ge r\}
 \le\min\{1,\exp[n^3(p_n(t,h)-p_n(t,0)-hr)]\}.
 \tag{V90.13}\label{c18v90:tail}
\]
Using (V90.9), its exponent divided by \(n^3\) is at most
\[
 p_\ell^{\rm o}(t,h)-p_\ell^{\rm o}(t,0)-hr
 +(C_t(h)+C_t(0))(24/\ell+48/n).
 \tag{V90.14}\label{c18v90:enclosedtail}
\]
For \(\ell\mid n\), \(n>\ell\), (V90.11) gives the sharper
replacement \(24/\ell\) for the last parenthesis. The proof
is simply Markov's inequality applied to the normalized law,
followed by the displayed pressure comparisons. No selected
defect family or product approximation of the native law is used.

At each fixed regulator differentiation under the compact
integral is legitimate and gives
\[
 \partial_h p_n(t,0)=n^{-3}\mu_{n,t,0}(E_{n,t}),\qquad
 \partial_h^2p_n(t,0)=V_{I,n}(t)/n^3.
 \tag{V90.15}\label{c18v90:susceptibility}
\]
This identifies the exact remaining variance as a source
susceptibility, not a derivative of the chart-time pressure
in V86. Uniform convergence in (V90.9) cannot be differentiated
twice without an additional estimate.

\begin{proposition}[Even an analytic limit with surface accuracy is insufficient]
\label{c18v90:counterexample}
Surface-order convergence of normalized log partition functions,
even to an affine analytic function, does not imply a
volume-uniform susceptibility bound.
\end{proposition}
\begin{proof}
This is an abstract counterexample, not a replacement YM law.
On a two-point probability space of equal weights put
\(E_n=n^3\pm n^2\), \(n\ge1\). It is nonnegative and at most
\(2n^3\). Its normalized log moment generating function is
\[
 \widetilde p_n(h)=h+n^{-3}\log\cosh(hn^2),\qquad
 |\widetilde p_n(h)-h|\le |h|/n,\qquad
 \widetilde p_n''(0)=n.
 \tag{V90.16}\label{c18v90:counter}
\]
All formulas follow directly from the two-point sum and
\(\log\cosh x\le|x|\). The limiting pressure is entire and
affine, while the variance per unit volume diverges.
\end{proof}

For orientation, if one later proves that \(p(t,\cdot)\) is
differentiable at zero, convexity and exponential Markov bounds
do imply concentration of \(E_{n,t}/n^3\) about
\(\partial_hp(t,0)\) at every fixed positive tolerance.
Indeed the secant slopes for sufficiently small fixed positive
and negative \(h\) then lie inside that tolerance; local uniform
convergence transfers their strict exponential bounds to large
\(n\). This only controls macroscopic deviations. The example
(V90.16) satisfies such concentration and still has divergent
susceptibility. Differentiability at zero for the native limit
has not been established here.

\subsection{Status and next load-bearing estimate}

Newly proved are the direct nonlinear deletion estimate, the
normalized arbitrary-exterior block comparison for V89's exact
density, and the source-tilted thermodynamic pressure limit with
explicit finite-size error. They hold at all compact chord
configurations and every fixed finite chart time in the stated
\(\beta=0,g=I\) setting. The first-order covariant estimate does
not extend V60's higher-jet constants to infinite time.

The next task must improve \emph{source response}, not just
repeat free-energy convergence. A sufficient target for the
static part is a volume-uniform quadratic bound on the centered
source log moment generating function near \(h=0\), which would
bound (V90.15). Even that would leave V87's integrated orthogonal
return \(1-\theta\) to control. The induced metric, the
all-coarse/interacting-vacuum extension, probability strain,
physical scaling, and the interacting four-dimensional
continuum construction remain unpaid. No YM mass gap is claimed.

The companion checkpoint is complete; the next is V95.
The broad literature checkpoint is complete; the next is V100,
also the archive/restart point. The accompanying exact checks
test the quaternion Hessian identities, the cut geometry on
several tori and tilings, and the analytic constants. They are
not simulations of the positive-time native probability law,
nor proof-assistant verification. Removing this append,
restoring the V89 title, and appending the final document
terminator recovers the predecessor byte for byte. Only the
active V89 source is replaced; every other note and archive is
unchanged. All build products remain outside \path{latex_notes}.
% END CYCLE 18 VERSION 90 APPEND
% BEGIN CYCLE 18 VERSION 91 APPEND
\clearpage
\section[V91: Source tilts, time shifts, and action-density concentration]{V91: Finite source tilts, chart-time shifts, and native action-density concentration}
\label{c18v91:main}

\subsection{Why the surface-pressure estimate is not differentiated}

V90 controls the actual normalizer and its source pressure by open
blocks. Its surface error is not a bound on the second source
derivative. V86 already gives the infinitesimal mean/variance identity
and the chart-time variance budget; V61 already provides short-time
probability transport. Neither result is reproved as a new gap here.
Instead we exploit a finite, pointwise cancellation between an action
source and a change of Wilson chart time. The remainder is quadratic
in the source with a volume-uniform coefficient.

The new finite-volume comparison holds for every coarse field at
\(\beta=0\). For the thermodynamic step we return to \(g=I\), where
V90 proved existence of the limiting pressure. There we obtain
exponential concentration at each fixed positive tolerance for all
chart times outside a countable exceptional set. This is
\emph{macroscopic action-density concentration}, not an \(O(v)\)
variance bound, an auxiliary spectral gap, or a physical mass gap.
The prescribed endpoint \(t=4\) cannot simply be replaced by a
nearby regular time.

A targeted literature screen also considered
\href{https://arxiv.org/abs/1001.1822}{Cattiaux--Guillin's
Lyapunov approach} and
\href{https://arxiv.org/abs/0712.0235}{Cattiaux--Guillin--Wang--Wu}.
Only their primary abstracts and scope were checked. No functional
inequality is imported: drift under deterministic Wilson flow is
not a Lyapunov drift estimate for the actual conditional diffusion
generator. The proof below uses the already identified native law.
The scheduled reviews remain V95 for companions and V100 for the
broad literature sweep.

\subsection{The inherited rooted law, now with arbitrary coarse field}

Let \(n\ge5\), \(v=n^3\), fine side \(2n\), and fix any coarse field
\(g\). V33's rooted representative \(U_g^0(y)\) sets forest links to
\(I\), last-chain links to \(g\), and its \(14v\) physical chords to
\(y\). V33--V38 give the exact product-Haar disintegration, including
the full inverse-flow Jacobian. Set
\[
 \begin{gathered}
 \Psi_s=\Phi_{-s},\qquad
 E_{g,s}(y)=S(\Psi_sU_g^0(y)),\qquad S=24v-W,\\
 \mathfrak R_{g,s}(y)=|F|^2(\Psi_sU_g^0(y)),\qquad F=\nabla W,\\
 \mathcal U_{g,t}(y)=12\int_0^t E_{g,s}(y)\,ds .
 \end{gathered}
 \tag{V91.1}\label{c18v91:data}
\]
We suppress \(n\) on these functions, not in pressure normalizations.
At \(\beta=0\), Liouville's formula from V38 reads
\[
 J_{-t}(U_g^0(y))=\exp(288vt-\mathcal U_{g,t}(y)).
\]
Thus, after integrating the independent root-fixed gauge Haar
variables, the actual conditional chord law has density
\(e^{-\mathcal U_{g,t}}/\mathcal Z_{n,g}(t,0)\).
This argument uses neither V89's \(g=I\) coercive envelope nor an
induced-metric replacement. The latter remains relevant to kinetic
and inverse estimates.

For real source \(h\), define
\[
 \begin{aligned}
 \mathcal Z_{n,g}(t,h)&=\int
            e^{-\mathcal U_{g,t}+hE_{g,t}}\,dy,&
 p_{n,g}(t,h)&=v^{-1}\log\mathcal Z_{n,g}(t,h),\\
 d\mu_{n,g,t,h}&=\mathcal Z_{n,g}(t,h)^{-1}
            e^{-\mathcal U_{g,t}+hE_{g,t}}\,dy,&
 G_{n,g}(t)&=p_{n,g}(t,0).
 \end{aligned}
 \tag{V91.2}\label{c18v91:law}
\]
All integrals use the same normalized chord Haar measure. At \(h=0\)
this is the native conditional probability, not a surrogate law.

V86's all-configuration force bound and the chain rule give
\[
 0\le E_{g,s}\le48v,\qquad
 \partial_s E_{g,s}=\mathfrak R_{g,s},\qquad
 0\le\mathfrak R_{g,s}\le384v.
 \tag{V91.3}\label{c18v91:force}
\]
Indeed \(\nabla S=-F\) and \(\dot\Psi_s=-F\).
These statements hold at every finite time and every \(g\).

\subsection{A finite source shift with a positive quadratic remainder}

\begin{theorem}[Exact tilt--time comparison]
\label{c18v91:shift}
Let \(t\ge0\), \(h\in\mathbb R\), and
\(\tau=t-h/12\ge0\). On the same chord coordinates put
\[
 \mathfrak B_{g;t,h}
   =\frac{h^2}{12}\int_0^1(1-u)
              \mathfrak R_{g,t-uh/12}\,du .
 \tag{V91.4}\label{c18v91:B}
\]
Then, pointwise and without a small-field assumption,
\[
 -\mathcal U_{g,t}+hE_{g,t}
       =-\mathcal U_{g,\tau}+\mathfrak B_{g;t,h},
 \qquad 0\le\mathfrak B_{g;t,h}\le16vh^2 .
 \tag{V91.5}\label{c18v91:pointwise}
\]
Consequently
\[
 \begin{gathered}
 0\le p_{n,g}(t,h)-G_{n,g}(\tau)\le16h^2,\\
 \frac{d\mu_{n,g,t,h}}{d\mu_{n,g,\tau,0}}
      =\frac{e^{\mathfrak B_{g;t,h}}}
             {\mu_{n,g,\tau,0}(e^{\mathfrak B_{g;t,h}})},
 \qquad
 e^{-16vh^2}\le
 \frac{d\mu_{n,g,t,h}}{d\mu_{n,g,\tau,0}}\le e^{16vh^2}.
 \end{gathered}
 \tag{V91.6}\label{c18v91:normalized}
\]
Both directed relative entropies between these chord laws are at
most \(16vh^2\).
\end{theorem}
\begin{proof}
Write \(\delta=\tau-t=-h/12\). The fundamental theorem of calculus,
with oriented integrals if \(\delta<0\), gives
\[
 \begin{split}
 \mathcal U_{g,\tau}-\mathcal U_{g,t}+hE_{g,t}
   &=12\int_0^\delta(E_{g,t+s}-E_{g,t})\,ds\\
   &=12\delta^2\int_0^1(1-u)\mathfrak R_{g,t+u\delta}\,du
     =\mathfrak B_{g;t,h}.
 \end{split}
\]
This proves the identity for either sign of \(h\).
Since \(\int_0^1(1-u)\,du=1/2\), (V91.3) gives
\(0\le\mathfrak B\le h^2(384v)/24=16vh^2\).
Integrate the identity against Haar and divide by
\(\mathcal Z_{n,g}(\tau,0)\). The logarithm of the resulting
expectation lies in \([0,16vh^2]\), proving the pressure bound.
The normalized density formula follows from the same identity.
Its logarithm has absolute value at most \(16vh^2\);
integration under either law proves the relative-entropy bounds.
\end{proof}

This is a comparison on a common \emph{initial chord coordinate}
space. The physical points \(\Psi_tU_g^0(y)\) and
\(\Psi_\tau U_g^0(y)\) lie on different moving fibres. No equality
of physical states, physical times, or induced kinetic metrics is
being asserted. For \(h=o(v^{-1/2})\) the displayed density ratio
does tend uniformly to one. At \(h=a/\sqrt v\), with fixed
\(a\ne0\), it is only bounded between \(e^{-16a^2}\) and
\(e^{16a^2}\); asymptotic normality does not follow.

There is also an exact finite-regulator pressure equation:
\[
 (\partial_t+12\partial_h)p_{n,g}(t,h)=h\,q_{n,g}(t,h),
 \qquad
 q_{n,g}(t,h)=v^{-1}\mu_{n,g,t,h}(\mathfrak R_{g,t})
                 \in[0,384].
 \tag{V91.7}\label{c18v91:PDE}
\]
To prove it, differentiate (V91.2) on the compact domain:
the exponent has time derivative \(-12E_{g,t}+h\mathfrak R_{g,t}\)
and source derivative \(E_{g,t}\).
The pointwise proof above is stronger than simply integrating
this characteristic equation, since it also compares normalized laws.

\begin{lemma}[The remainder retains its actual force observable]
\label{c18v91:cubic}
In the same range,
\[
 \left|\mathfrak B_{g;t,h}-\frac{h^2}{24}\mathfrak R_{g,t}\right|
                  \le\frac{128}{9}v|h|^3.
 \tag{V91.8}\label{c18v91:cubicbound}
\]
\end{lemma}
\begin{proof}
V90's covariant Hessian block bounds imply
\(\|\nabla^2W\|_{2\to2}\le16\). Along the inverse flow,
\[
 \partial_s\mathfrak R_{g,s}
   =-2\nabla^2W(F,F)(\Psi_sU_g^0),\qquad
 |\partial_s\mathfrak R_{g,s}|\le32\mathfrak R_{g,s}\le12288v .
\]
Insert
\(|\mathfrak R_{g,t-uh/12}-\mathfrak R_{g,t}|
 \le12288v\,u|h|/12\)
in (V91.4), and use \(\int_0^1u(1-u)\,du=1/6\).
The coefficient is \(12288/(12^2\cdot6)=128/9\).
\end{proof}
This pointwise cubic error does not replace the random
\(\mathfrak R_{g,t}\) by its mean, nor does it bound the response
of that mean to a change of probability.

\subsection{The thermodynamic consequence: at most countably many cusp times}

Now specialize to \(g=I\). Write \(p_n=p_{n,I}\),
\(G_n=G_{n,I}\), and use V90 to define
\[
 p(t,h)=\lim_{n\to\infty}p_n(t,h),\qquad G(t)=p(t,0).
\]
At a fixed regulator let
\[
 e_n(t)=v^{-1}\mu_{n,I,t,0}(E_{I,t}),\qquad
 \chi_n(t)=v^{-1}\operatorname{Var}_{\mu_{n,I,t,0}}(E_{I,t}).
\]
V86, or direct differentiation of (V91.2), gives
\[
 G_n'=-12e_n,\qquad
 G_n''=144\chi_n-12q_{n,I}(t,0)\ge-4608,\qquad
 -576\le G_n'\le0 .
 \tag{V91.9}\label{c18v91:timecurvature}
\]
This is the inherited mean/variance identity, now in the
inverse-flow source notation.

\begin{theorem}[Native first-order response away from a countable set]
\label{c18v91:regular}
The function \(G\) is nonincreasing and \(576\)-Lipschitz,
and \(H(t)=G(t)+2304t^2\) is convex. For every \(t>0\),
\[
 \partial_h^-p(t,0)=-\frac1{12}G_+'(t),\qquad
 \partial_h^+p(t,0)=-\frac1{12}G_-'(t).
 \tag{V91.10}\label{c18v91:sides}
\]
In particular the set
\(\mathcal D=\{t>0:G_-'(t)<G_+'(t)\}\) is at most countable.
For every \(t\notin\mathcal D\), \(t>0\), both
\[
 e(t):=\partial_hp(t,0)=-G'(t)/12,\qquad
                    \lim_{n\to\infty}e_n(t)=e(t)
 \tag{V91.11}\label{c18v91:emean}
\]
exist. The set of cusp times is exactly the set where the limiting
source pressure is not differentiable at zero.
\end{theorem}
\begin{proof}
The derivative bounds in (V91.9) and pointwise convergence imply
local uniform convergence of \(G_n\) to \(G\) (use a finite mesh
and the common Lipschitz bound). Each
\(H_n=G_n+2304t^2\) is convex, so its pointwise limit is convex.
Thus \(G\) has finite left and right derivatives at interior
times, with \(G_-'\le G_+'\).

Pass to the limit in the first line of (V91.6):
\[
             0\le p(t,h)-G(t-h/12)\le16h^2
             \quad(t-h/12\ge0).
 \tag{V91.12}\label{c18v91:limitshift}
\]
For \(h>0\), divide by \(h\) after subtracting \(G(t)\), and use
the left derivative of \(G\). For \(h<0\) use its right derivative.
The error divided by \(|h|\) tends to zero, giving (V91.10).

The left and right derivatives of a convex function differ at
at most countably many points: the disjoint open intervals
\((H_-',H_+')\) at distinct such points each contain a distinct
rational number. Adding a smooth quadratic does not change
their jumps. At a differentiability point, convex secant bounds
squeeze \(H_n'(t)\) between the backward and forward difference
quotients of \(H_n\). First let \(n\to\infty\), then let the
secant length tend to zero. This proves \(G_n'(t)\to G'(t)\)
and (V91.11).
\end{proof}

For \(0<a<b<\infty\), the total possible source-slope jump obeys
\[
 \sum_{t\in\mathcal D\cap(a,b)}
  \bigl(\partial_h^+p(t,0)-\partial_h^-p(t,0)\bigr)
                         \le48+384(b-a).
 \tag{V91.13}\label{c18v91:jumps}
\]
Indeed the summands are \((H_+'-H_-')/12\), whose sum is bounded
by \([H_-'(b)-H_+'(a)]/12\). Use
\(-576\le G_\pm'\le0\).
Consequently there are only finitely many jumps larger than any
specified positive size in a compact interval. This is a bound
on possible limiting first-order jumps, not a bound on
finite-volume susceptibility peaks.

\subsection{A native concentration theorem, with its exact quantifiers}

\begin{corollary}[Action-density self-averaging at each regular chart time]
\label{c18v91:LLN}
Fix \(t>0\) with \(t\notin\mathcal D\), and \(\epsilon>0\).
There are \(c(t,\epsilon)>0\) and \(n_0(t,\epsilon)\) such that
for all \(n\ge n_0\),
\[
 \mu_{n,I,t,0}\{|E_{I,t}/v-e(t)|\ge\epsilon\}
                        \le2e^{-c(t,\epsilon)v}.
 \tag{V91.14}\label{c18v91:concentration}
\]
In particular
\[
              \frac{\operatorname{Var}(W)}{v^2}
              =\frac{V_{I,n}(t)}{v^2}\longrightarrow0 .
 \tag{V91.15}\label{c18v91:subquadratic}
\]
These probabilities and variances are for the actual conditional
law, transported through the exact chart.
\end{corollary}
\begin{proof}
Differentiability of \(p(t,\cdot)\) at zero allows fixed
\(h_+>0\), \(h_-<0\), with \(h_+<12t\), such that
\[
 |p(t,h_\pm)-p(t,0)-h_\pm e(t)|
                        \le |h_\pm|\epsilon/4.
\]
V90's convergence at these three source values bounds the
finite-volume error in each pressure difference by
\(|h_\pm|\epsilon/4\) for sufficiently large \(n\).
Apply exponential Markov to the upper tail with \(h_+\)
and to the lower tail with \(h_-\). Their exponents are
at most \(-v|h_\pm|\epsilon/2\).
This proves (V91.14) with
\(c=\epsilon\min\{h_+,|h_-|\}/2\).

Since \(0\le E_{I,t}/v,e(t)\le48\),
\[
 \operatorname{Var}(E_{I,t}/v)
 \le\mu_{n,I,t,0}[(E_{I,t}/v-e(t))^2]
 \le\epsilon^2+2\cdot48^2e^{-cv}.
\]
Let \(n\to\infty\), then \(\epsilon\downarrow0\).
The equality of the Wilson and deficit variances proves
(V91.15).
\end{proof}

The constants are not uniform as \(\epsilon\downarrow0\) or as
\(t\) approaches \(\mathcal D\). Nothing proves
\(4\notin\mathcal D\), and even proving this would give only
\(V_{I,n}(4)=o(v^2)\), not \(O(v)\). This theorem is not an
infinite-volume central limit theorem, a unique-state theorem,
or a statement about the auxiliary diffusion's return integral.

\subsection{The precise extra time-pressure estimate that would pay variance}

Return to any fixed \(g\) at finite \(n\). Define
\[
 \begin{split}
 \Lambda_{n,g,t}(h)&=
  v^{-1}\log\mu_{n,g,t,0}
       \exp\{h(E_{g,t}-\mu_{n,g,t,0}E_{g,t})\},\\
 T_{n,g,t}(\delta)&=
       G_{n,g}(t+\delta)-G_{n,g}(t)-\delta G_{n,g}'(t).
 \end{split}
\]
With \(\delta=-h/12\), (V91.6) gives the finite inequality
\[
       0\le\Lambda_{n,g,t}(h)-T_{n,g,t}(-h/12)\le16h^2.
 \tag{V91.16}\label{c18v91:centered}
\]
Here \(\Lambda\ge0\) by Jensen; \(T\) need not be nonnegative.
If, on a two-sided interval independent of \(n\), one establishes
\[
 T_{n,g,t}(\delta)\le K\delta^2/2
 \quad\hbox{uniformly in \(n\)},\qquad K\ge0,
 \tag{V91.17}\label{c18v91:gate}
\]
then
\[
 \Lambda_{n,g,t}(h)\le(K/288+16)h^2,\qquad
          V_{g,n}(t)/v\le K/144+32 .
 \tag{V91.18}\label{c18v91:paidif}
\]
The first inequality is (V91.16); differentiating it twice
at \(h=0\), at each fixed regulator, proves the second.
Conversely a bound \(\Lambda_{n,g,t}(h)\le Ch^2\) on such
an interval implies \(T_{n,g,t}(\delta)\le144C\delta^2\).
Thus the two \emph{uniform quadratic-remainder} targets are
equivalent up to explicit constants. No upper bound (V91.17)
has been proved at \(t=4\).

For example, if (V91.18)'s first bound holds for \(|h|\le h_0\)
and \(C=K/288+16\), then for \(0<r\le2Ch_0\) exponential
Markov at \(h=\pm r/(2C)\) gives
\[
 \mu_{n,g,t,0}\{|E_{g,t}-\mu E_{g,t}|\ge rv\}
                         \le2e^{-vr^2/(4C)}.
\]
Unlike (V91.14), this controls the small-tolerance quadratic
scale with a common constant. It remains conditional on the
unpaid finite-volume upper Taylor bound, not on smoothness of
the limiting pressure alone.

\subsection{The new shift identity still does not remove susceptibility spikes}

We extend, rather than rebrand, V90's two-point counterexample.
It now satisfies the exact scalar shift relation and has a
vanishing integrated susceptibility budget.

Fix \(t_*>0\), \(v=n^3\), and on two states put
\(E_\sigma=v+\sigma n^2\), \(\sigma=\pm1\). Use reference weights
\(\pi_n(\sigma)\) proportional to \(e^{12t_*E_\sigma}\), and
potential \(\widetilde{\mathcal U}_t=12tE_\sigma\).
The energy is constant along chart time, so its scalar
\(\widetilde{\mathfrak R}\) is zero. Direct summation gives
\[
 \widetilde p_n(t,h)
  =h-12t+\frac1{n^3}
   \left[\log\cosh\bigl((h-12(t-t_*))n^2\bigr)
                  -\log\cosh(12t_*n^2)\right].
 \tag{V91.19}\label{c18v91:toy}
\]
It obeys
\(\widetilde p_n(t,h)=\widetilde p_n(t-h/12,0)\)
exactly and converges locally uniformly to the affine analytic
function \(h-12t\), with \(O(1/n)\) error on compact sets.
Nevertheless
\[
 \widetilde\chi_n(t)
   =n\,\operatorname{sech}^2(12(t-t_*)n^2),\qquad
 \widetilde\chi_n(t_*)=n,\qquad
 \int_{\mathbb R}\widetilde\chi_n(t)\,dt=\frac1{6n}.
 \tag{V91.20}\label{c18v91:toyvariance}
\]
The formulas follow by two \(h\)-derivatives, followed by
integration of the derivative of \(\tanh\).
The convergence bound uses \(|\log\cosh x|\le|x|\).

This is an abstract two-state family with a volume-dependent
reference law. It is \emph{not} the native Haar geometry and
does not furnish a Yang--Mills counterexample. It proves the
limited logical point required here: the scalar shift identity,
smooth limiting pressure, surface-order pressure convergence,
and an integrated susceptibility bound still do not imply a
uniform pointwise susceptibility. One may set \(t_*=4\);
the finite-volume upper Taylor curvature there is
\(144n\), despite zero limiting curvature.

\subsection{Status, verification, and continuation}

V91 proves an all-coarse finite-volume tilt--time comparison at
\(\beta=0\), including its actual normalized law and a pointwise
force remainder. Using the \(g=I\) pressure limit, it proves
native action-density concentration outside a countable set of
chart times and \(V=o(v^2)\) at each such time. It does not
establish the endpoint \(O(v)\) bound, the integrated orthogonal
return, the probability-strain estimate, or the interacting
four-dimensional continuum mass gap. The goal remains active.

The next load-bearing task is a finite-volume upper
time-pressure Taylor bound at the prescribed endpoint, or a
native connected-correlation estimate strong enough to imply it.
Another proof of limiting-pressure differentiability is not
sufficient. The independent return margin of V87 remains even
after this static estimate. The interacting-vacuum factor adds
a time-dependent term to the logarithmic density and is not
covered by the present \(\beta=0\) cancellation.

The checker verifies the oriented-integral identities for both
source signs, the force-remainder constants, an exact
one-plaquette flow regression, and the complete two-state
pressure/variance formulas. V90's compact quaternion Hessian
checks are rerun as inherited regressions. These are analytic
and finite algebraic checks, not a simulation or numerical
certificate of native thermodynamic concentration.
Removing this append and restoring the title recovers V90
byte for byte. Only the active predecessor is replaced; all
other notes and archives are preserved. Build products stay
outside \path{latex_notes}. The next scheduled companion review
is V95; the next broad literature review and archive/restart
point are V100.
% END CYCLE 18 VERSION 91 APPEND
\end{document}
